\documentclass[a4paper,10pt]{article}
\usepackage[top=25mm,bottom=24mm,left=27mm,right=27mm,headheight=15pt,headsep=8mm,footskip=12mm]{geometry}
\usepackage{fontspec,amsmath,amsthm,unicode-math}
\usepackage[spanish,es-nodecimaldot,es-noquoting]{babel}
\usepackage{microtype,xcolor,fancyhdr,titlesec,booktabs,longtable,array,tabularx}
\usepackage{graphicx,tikz,enumitem,float,needspace}
\usepackage{tocloft}
\setlength{\cftbeforesecskip}{1pt}
\usetikzlibrary{arrows.meta,calc,positioning,decorations.pathreplacing,matrix}
\usepackage{xurl}
\usepackage[unicode,hidelinks]{hyperref}
\usepackage[nameinlink,noabbrev]{cleveref}
\crefname{theorem}{teorema}{teoremas}\Crefname{theorem}{Teorema}{Teoremas}
\crefname{proposition}{proposición}{proposiciones}\Crefname{proposition}{Proposición}{Proposiciones}
\crefname{lemma}{lema}{lemas}\Crefname{lemma}{Lema}{Lemas}
\crefname{corollary}{corolario}{corolarios}\Crefname{corollary}{Corolario}{Corolarios}
\crefname{definition}{definición}{definiciones}\Crefname{definition}{Definición}{Definiciones}
\crefname{section}{sección}{secciones}\Crefname{section}{Sección}{Secciones}
\crefname{subsection}{sección}{secciones}\Crefname{subsection}{Sección}{Secciones}
\crefname{equation}{ecuación}{ecuaciones}\Crefname{equation}{Ecuación}{Ecuaciones}
\crefname{figure}{figura}{figuras}\Crefname{figure}{Figura}{Figuras}
\crefname{table}{tabla}{tablas}\Crefname{table}{Tabla}{Tablas}
\setmainfont{STIXTwoText-Regular.otf}[BoldFont=STIXTwoText-Bold.otf,ItalicFont=STIXTwoText-Italic.otf,BoldItalicFont=STIXTwoText-BoldItalic.otf]
\setmathfont{STIXTwoMath-Regular.otf}
\hypersetup{pdftitle={Extrema y media razón holográfica},pdfauthor={Oumar Haidara Fall; Rubén Ramos Balsa},pdfsubject={Derivación de la constante holográfica de acoplamiento aritmético-geométrico y leyes de conservación y reversibilidad de la información},pdfkeywords={Holografía Modular Triádica; Mecánica Dimensional; extrema y media razón holográfica; acoplamiento aritmético-geométrico; refinamiento reversible; memoria operatoria; transporte bilateral; incidencia excepcional; observables invariantes; recuperación estable; conservación exterior; entropía de entrelazamiento}}
\newcommand{\Autores}{Oumar Haidara Fall \textperiodcentered\ Rubén Ramos Balsa}
\newcommand{\Titulo}{Extrema y media razón holográfica}
\newcommand{\RR}{\mathbb R}\newcommand{\ZZ}{\mathbb Z}\newcommand{\QQ}{\mathbb Q}\newcommand{\CC}{\mathbb C}\newcommand{\FF}{\mathbb F}
\newcommand{\Id}{\operatorname{id}}\newcommand{\tr}{\operatorname{tr}}\newcommand{\dd}{\mathrm d}
\newcommand{\Span}{\operatorname{span}}\newcommand{\Cyl}{\operatorname{Cyl}}
\providecommand{\llbracket}{\lBrack}\providecommand{\rrbracket}{\rBrack}
\newcommand{\HMT}{\mathrm{HMT}}\newcommand{\Tr}{\operatorname{Tr}}\newcommand{\Spec}{\operatorname{Spec}}
\newcommand{\Cstar}{C^{*}}\newcommand{\R}{\mathbb R}\newcommand{\Z}{\mathbb Z}\newcommand{\I}{\mathrm I}
\definecolor{ink}{HTML}{183D4A}\definecolor{accent}{HTML}{8C3D2E}
\newcommand{\HMTNarrativePaper}{\clearpage\pagecolor{white}}
\newcommand{\HMTMathematicalPaper}{\clearpage\pagecolor{white}}
\newcommand{\HMTNarrativeHeading}[1]{\par\Needspace{8\baselineskip}\noindent{\fontsize{16}{21}\selectfont\bfseries\color{ink}#1\par}\vspace{8mm}}
\newcommand{\HMTDivision}[4][]{%
  \HMTNarrativePaper\phantomsection
  \addcontentsline{toc}{part}{#3}%
  \noindent{\fontsize{22}{28}\selectfont #2\par}\vspace{5mm}
  \noindent{\fontsize{17}{22}\selectfont\bfseries\raggedright #3\par}\vspace{4mm}
  \noindent{\fontsize{11}{15}\selectfont #4\par}\vspace{7mm}}
\providecommand{\dashrightarrow}{\mathrel{\rightdasharrow}}
\newcommand{\square}{\mdlgwhtsquare}
\newcommand{\TPK}{\mathrm{TPK}}\newcommand{\APP}{\mathrm{APP}}\newcommand{\TRIT}{\{-1,0,+1\}}
\newtheorem{theorem}{Teorema}[section]
\newtheorem{lemma}[theorem]{Lema}
\newtheorem{proposition}[theorem]{Proposición}
\newtheorem{corollary}[theorem]{Corolario}
\newtheorem{criterion}[theorem]{Criterio}
\providecommand{\index}[1]{}
\theoremstyle{definition}
\newtheorem{definition}[theorem]{Definición}\newtheorem{axiom}[theorem]{Axioma}
\newtheorem{principle}[theorem]{Principio}
\theoremstyle{remark}\newtheorem{remark}[theorem]{Observación}
\numberwithin{equation}{section}
\setlength{\parindent}{1em}\setlength{\parskip}{3pt plus 1pt minus .5pt}
\setlength{\emergencystretch}{2em}
\renewcommand{\normalsize}{\fontsize{10.5}{13.4}\selectfont}\normalsize
\setlist{nosep,topsep=5pt}
\titleformat{\section}{\fontsize{13}{16}\selectfont\bfseries\raggedright\hyphenpenalty=10000}{\thesection}{.65em}{}
\titleformat{\subsection}{\fontsize{11.2}{14}\selectfont\bfseries}{\thesubsection}{.65em}{}
\titleformat{\subsubsection}{\normalsize\bfseries}{\thesubsubsection}{.65em}{}
\newif\ifHMTSectionPage\HMTSectionPagetrue
\AddToHook{cmd/section/before}{\ifHMTSectionPage\clearpage\else\Needspace{8\baselineskip}\fi}
\AddToHook{cmd/subsection/before}{\Needspace{7\baselineskip}}
% Keep the differential heading with its first theorem in this edition.
% Los saltos de subsección se rigen por el espacio disponible, no por números heredados.
\AddToHook{cmd/subsubsection/before}{\Needspace{6\baselineskip}}
\AddToHook{env/theorem/before}{\Needspace{7\baselineskip}}
\AddToHook{env/proposition/before}{\Needspace{7\baselineskip}}
\setcounter{tocdepth}{1}
\setlength{\cftbeforepartskip}{9pt}
\renewcommand{\cftpartfont}{\bfseries}
\newcounter{HMTBibliographyCount}
\let\HMTOriginalTheBibliography\thebibliography
\renewcommand{\thebibliography}[1]{%
  \HMTOriginalTheBibliography{99}%
  \setcounter{enumiv}{\value{HMTBibliographyCount}}}
\AddToHook{env/thebibliography/end}{\setcounter{HMTBibliographyCount}{\value{enumiv}}}
\newenvironment{HMTCompactSection}{%
  \begingroup
  \setlength{\parskip}{1.5pt plus 1pt minus .5pt}%
  \setlength{\abovedisplayskip}{8pt plus 2pt minus 4pt}%
  \setlength{\belowdisplayskip}{8pt plus 2pt minus 4pt}%
}{\par\endgroup}
\titlespacing*{\section}{0pt}{17pt}{9pt}\titlespacing*{\subsection}{0pt}{13pt}{6pt}
\pagestyle{fancy}\fancyhf{}
\fancyhead[L]{\fontsize{8}{10}\selectfont Extrema y media razón holográfica}
\fancyfoot[C]{\thepage}\renewcommand{\headrulewidth}{.2pt}
\clubpenalty=10000\widowpenalty=10000\displaywidowpenalty=10000
\raggedbottom
% Content-bound divisions remain valid when the causal order changes.
\AddToHook{file/v_funcional_barbero.tex/before}{\HMTDivision{\(\gamma\)}{Parámetro de Barbero--Immirzi}{Incidencia hexádico--octádica, acción, área e información.}\HMTNarrativeHeading{Acción, área e información}

El parámetro de Barbero--Immirzi se presenta sobre los antecedentes
de acción, orientación e incidencia ya construidos. Las coordenadas
angulares y los canales conjugados actúan junto con la inclusión
marcada entre configuraciones hexádicas y octádicas. El funcional
recibe esos datos y determina su evaluación. La relación entre
aritmética y geometría permanece visible en cada uno de sus argumentos.

La estructura de incidencia organiza qué elementos se distinguen y
cómo se prolongan. Una hexada es un soporte de seis posiciones del
diseño; su elevación utiliza una dodecada y un alineamiento declarados.
La octada resultante pertenece a la realización binaria. Los lectores
de área conservan la marca que hace significativa esa correspondencia.
Las figuras ayudan a seguir el transporte de los conjuntos y sus
intersecciones.

La acción y el reloj permiten después expresar la energía asociada
a una lectura. La constante de Boltzmann interviene en la transducción
térmica; el parámetro de Barbero y la unidad areal intervienen en el
incremento geométrico. La constante gravitatoria entra en la
realización que relaciona acción y área de Planck. Al componer estas
expresiones, una misma sección de acción puede cancelarse en el
cociente entre energía e incremento de área. La cancelación exhibe
una dependencia estructural concreta.

La ley de cambio de estado conserva también el término de entropía
relativa y la referencia respecto de la cual se realiza la comparación.
La información de una observación reducida depende de las correlaciones
que mantiene el estado conjunto. Esta distinción prolonga el argumento
del archivo: el contenido de la totalidad y el de un lector parcial
se relacionan mediante operaciones explícitas.

La teoría de conservación adquiere así una aplicación física
determinada. El estado, su orientación, el área que se le asigna y
la lectura térmica forman una composición con antecedentes
identificados. El desarrollo conserva esos antecedentes y sus
pruebas. La síntesis posterior reunirá esta composición con los
balances operatorios y con el ejemplo de entrelazamiento calculable.

\HMTMathematicalPaper}
\AddToHook{cmd/@startpbox/after}{\raggedright\arraybackslash}
% Hierarchical contents: divisions, sections, subsections.
\setcounter{tocdepth}{2}
\hypersetup{bookmarksdepth=3}
\setlength{\cftbeforepartskip}{11pt plus 2pt}
\setlength{\cftbeforesecskip}{4pt plus 1pt}
\setlength{\cftbeforesubsecskip}{1.4pt}
\renewcommand{\cftpartfont}{\fontsize{11.2}{14}\selectfont\bfseries}
\renewcommand{\cftpartpagefont}{\fontsize{10}{12}\selectfont\bfseries}
\renewcommand{\cftsecfont}{\fontsize{9.8}{12}\selectfont\normalfont}
\renewcommand{\cftsecpagefont}{\fontsize{9.8}{12}\selectfont\normalfont}
\renewcommand{\cftsubsecfont}{\fontsize{9}{11.2}\selectfont\normalfont}
\renewcommand{\cftsubsecpagefont}{\fontsize{9}{11.2}\selectfont\normalfont}
\setlength{\cftsecindent}{1em}
\setlength{\cftsecnumwidth}{2.8em}
\setlength{\cftsubsecindent}{3.8em}
\setlength{\cftsubsecnumwidth}{3.8em}
\renewcommand{\cftsecleader}{\cftdotfill{2.5}}
\renewcommand{\cftsubsecleader}{\cftdotfill{3.5}}
\cftsetpnumwidth{2.8em}
\cftsetrmarg{3.8em}
\renewcommand{\cftpartafterpnum}{\par\nobreak}

\begin{document}
\thispagestyle{empty}
\begin{center}
{\fontsize{10}{12}\selectfont Manuscrito de recopilación para transferencia académica\par}
\vspace{6mm}
{\fontsize{20}{24}\selectfont\bfseries\mbox{Extrema y media razón holográfica}\par}
\vspace{3mm}
{\fontsize{12}{15}\selectfont
Derivación de la constante holográfica de acoplamiento aritmético-geométrico.\par
Conservación y reversibilidad de la información\par}
\vspace{5mm}
{\fontsize{11.5}{14}\selectfont\Autores\par}
\vspace{2mm}
{\fontsize{8}{10}\selectfont Coautoría sin prelación de contribución\par
Integración y recopilación documental generadas por ChatGPT - Astra.\par}
\end{center}
\vspace{3mm}

\noindent\textbf{Resumen.}\par\smallskip
Se determina la constante holográfica de acoplamiento aritmético--geométrico mediante un registro ordenado \(K\), su evaluación racional y sus representaciones incidenciales, y se demuestra una ley de conservación que permite recuperar información a través de sus transformaciones. Fijadas la base mil y la carta de fase, la lectura \(\kappa=N_K/(1000^{12}-1)\) conserva exactamente las doce componentes del registro. Su órbita cíclica forma una base y proporciona identificación estable de operadores lineales: se utilizan las respuestas a la base completa o una respuesta vectorial cuando existe covarianza cíclica. Así, la estructura que determina el valor participa también en la reconstrucción de las operaciones que actúan sobre ella.

La construcción parte de un núcleo formal denominado Holografía Modular Triádica, en adelante HMT. Comprende \emph{All Possible Paths} (APP), estructura aritmética de caminos sobre un soporte toroidal con grafo de Cayley y evaluaciones aditiva y multiplicativa distinguidas; el TRIT, firma orientada de valores \(\{-1,0,+1\}\); y el \emph{Topological Phase Kernel} (TPK), que compone selección, transporte y memoria. La prolongación nonádica conserva el estado enriquecido del que proceden las coordenadas aritméticas y las incidencias de frontera. El registro \(K\) interviene en la clausura de \(\alpha\) junto con las coordenadas HMT de \(\pi,\varphi,e\); su realización incidencial transporta la misma información ordenada a los espacios que se especifican.

El resultado general de conservación es una identidad bilateral para dos isometrías con dominio y codominio comunes, válida en dimensión arbitraria. El refinamiento entero con acarreo tiene inversa exacta sobre etiquetas admisibles y un levantamiento unitario que se compone con esa identidad. En la especialización nonádica, los pesos \(72/73\) y \(1/73\) determinan la transformación de observables; los invariantes son exactamente los que conmutan con el transporte relativo. Una política explícita de archivo transfiere toda la norma a la memoria infinita y ofrece reconstrucción por un adjunto de norma uno, además de inversas finitas con cotas uniformes en profundidad. Dos observaciones complementarias recuperan con cota óptima el registro centrado, del que se reconstruye el proyector dimensional especificado. La conservación se extiende a cada grado exterior admisible, incluidas las contribuciones mixtas de lectura y memoria.

La extensión reversible relaciona deformación, holonomía y acción; las realizaciones electromagnéticas y gravitatorias conservan sus operadores y normalizaciones explícitos. Para el protocolo simpléctico declarado, con dos entradas gaussianas puras, centradas e idénticas y un número finito de modos, la componente antilineal de la memoria determina el espectro simpléctico reducido, manteniendo todas las ocupaciones de Fock. Un lazo elíptico--parabólico da pureza \(9/11\), espectro \((9/10)10^{-j}\) y entropía adimensional \((10/9)\log10-\log9\). El resultado conjunto caracteriza qué conserva una transformación, qué modifica una observación y qué datos permiten invertirla.

\medskip
\noindent\textbf{Palabras clave:} Holografía Modular Triádica; Mecánica Dimensional; extrema y media razón holográfica; acoplamiento aritmético--geométrico; refinamiento reversible; memoria operatoria; holonomía; incidencia excepcional; conservación exterior; transporte bilateral; observables invariantes; recuperación estable; acción hamiltoniana; covarianza gaussiana; entropía de entrelazamiento.

\clearpage
\noindent\textbf{Abstract.}\par\smallskip
The holographic arithmetic--geometric coupling constant is determined through an ordered register \(K\), its rational evaluation and its incidence representations, and a conservation law is proved that permits information recovery through these transformations. With base one thousand and the phase chart fixed, the reading \(\kappa=N_K/(1000^{12}-1)\) preserves all twelve register components exactly. Its cyclic orbit forms a basis and provides stable identification of linear operators: the data are the responses to the complete basis, or one vector response under cyclic covariance. The structure that determines the value thus also participates in reconstructing the operations acting on it.

The construction starts from a formal kernel called Triadic Modular Holography, hereafter HMT. It comprises \emph{All Possible Paths} (APP), an arithmetic path structure on a toroidal support with a Cayley graph and distinct additive and multiplicative evaluations; TRIT, an oriented signature taking values in \(\{-1,0,+1\}\); and the \emph{Topological Phase Kernel} (TPK), which composes selection, transport and memory. Nonadic prolongation preserves the enriched state from which arithmetic coordinates and boundary incidences arise. The register \(K\) enters the closure of \(\alpha\) together with the HMT coordinates of \(\pi,\varphi,e\); its incidence realization transports the same ordered information to the specified spaces.

The general conservation result is a bilateral identity for two isometries with a common domain and codomain, valid in arbitrary dimension. Integer refinement with carry has an exact inverse on admissible labels and a unitary lift that composes with this identity. In the nonadic specialization, the weights \(72/73\) and \(1/73\) determine the transformation of observables; the invariants are precisely those commuting with the relative transport. An explicit archive policy transfers the entire norm to infinite memory and yields reconstruction by an adjoint of norm one, together with finite inverses whose bounds are uniform in depth. Two complementary observations recover the centred register with an optimal bound, from which the specified dimensional projector is reconstructed. Conservation extends to every admissible exterior degree, including mixed readout--memory contributions.

The reversible extension relates deformation, holonomy, and action; the electromagnetic and gravitational realizations retain their explicit operators and normalizations. For the declared symplectic protocol with two identical centred pure Gaussian inputs and finitely many modes, the antilinear component of memory determines the reduced symplectic spectrum, retaining all Fock occupations. An elliptic--parabolic loop gives purity \(9/11\), spectrum \((9/10)10^{-j}\), and dimensionless entropy \((10/9)\log10-\log9\). Taken together, the results characterize what a transformation preserves, what an observation changes, and which data permit its inversion.

\medskip
\noindent\textbf{Keywords:} Triadic Modular Holography; Dimensional Mechanics; holographic extreme and mean ratio; arithmetic--geometric coupling; reversible refinement; operator memory; holonomy; exceptional incidence; exterior conservation; bilateral transport; invariant observables; stable recovery; Hamiltonian action; Gaussian covariance; entanglement entropy.

\clearpage\begingroup\setlength{\parskip}{0pt}\tableofcontents\endgroup
\HMTMathematicalPaper
\section*{Introducción: determinación estructural y recuperación de información}
\addcontentsline{toc}{section}{Introducción: determinación estructural y recuperación de información}
\label{libro:prologo}
La extrema y media razón holográfica plantea una cuestión precisa: cómo puede una estructura producir una evaluación numérica, transformarse y conservar la información necesaria para reconocer e invertir esa transformación. El problema comprende tanto la formación del valor como la persistencia de sus relaciones. Una igualdad entre cantidades permite comprobar un balance; una ley de recuperación determina además qué estado lo realiza, qué datos lo individualizan y mediante qué operación se restituye. El presente trabajo desarrolla esa segunda exigencia desde una construcción aritmética de caminos y la lleva hasta sus representaciones lineales, incidenciales y físicas.

El primer resultado que organiza el argumento es la determinación de la constante holográfica de acoplamiento aritmético--geométrico. Su representación ordenada es el registro dodecafásico \(K\); su evaluación posicional es el racional \(\kappa\). La lectura conserva exactamente el registro cuando se fijan la base, la longitud y el origen de fase. El interés del objeto reside también en sus operaciones: sus desplazamientos cíclicos forman una base, sus respuestas permiten identificar operadores y sus imágenes incidenciales transportan las relaciones de la configuración inicial. Estructura, ecuación y valor aparecen así como momentos de una construcción cuya información recuperable se determina matemáticamente.

El segundo resultado proporciona la ley general de esa conservación. Dos transportes isométricos con dominio y codominio comunes producen dos componentes cuyo balance ponderado restituye exactamente el producto interior inicial. La especialización nonádica conserva una norma ternaria de valor \(73\) y determina una transformación de observables con pesos \(72/73\) y \(1/73\). Una política de prolongación y archivo convierte esa identidad local en recuperación estable a profundidad arbitraria. La tesis del artículo adquiere con ello un contenido verificable: el valor evaluado, el estado conservado y la lectura observada se relacionan por mapas explícitos y fórmulas inversas.

\subsection*{El núcleo generador y sus dos realizaciones}
La construcción se establece en un núcleo formal denominado \emph{Holografía Modular Triádica}, en adelante HMT. Su primer objeto es \emph{All Possible Paths}, en adelante APP: una estructura aritmética de caminos sobre el soporte toroidal \((\mathbb Z/9\mathbb Z)^2\), cuyo grafo de Cayley describe los desplazamientos cardinales. Sobre ese mismo soporte se distinguen la evaluación aditiva y la multiplicativa. La conservación del residuo junto con el cociente mantiene el entero evaluado y permite seguir la procedencia de cada publicación. La suma acumula; el producto organiza la segunda hoja mediante su propia evaluación y sus reglas de transporte.

El segundo objeto, el \emph{TRIT}, aporta la firma orientada \(\{-1,0,+1\}\) y distingue los regímenes locales. La ecuación de Euler extendida expresa, en sus álgebras cuadráticas, los regímenes elíptico, parabólico e hiperbólico mediante una ley exponencial común. El tercer objeto es el \emph{Topological Phase Kernel}, en adelante TPK. Su acción compone selección, transporte y actualización de memoria. Ruta, dirección, hoja, residuo, cociente, acarreo y frontera se conservan en el estado enriquecido pertinente; el retorno de fase prolonga su historia. La conexión y holonomía nonádicas pertenecen a esta dinámica.

De ese estado procede la estructura discreta conjunta del continuo y sus dos realizaciones: la contracción de cilindros compatibles determina coordenadas aritméticas, mientras que la incidencia de frontera conserva relaciones combinatorias y geométricas. Las cinco construcciones correlativas del continuo permanecen en ese mismo desarrollo. Los espacios lineales, las formas y los observables utilizados después son realizaciones de objetos ya generados, con sus dominios y normalizaciones. La continuidad del argumento reside en que cada operación conserva las relaciones que necesita la siguiente.

\subsection*{Completación de la unidad y refinamiento reversible}
La completación de una célula de nueve posibilidades por coordenada añade una posición de frontera distinguida. En tres coordenadas, la distribución según el número de posiciones de frontera da
\[
10^3=729+243+27+1=729+271.
\]
El interior, los estratos intermedios y el ápice componen una misma unidad normalizada. En dimensión finita arbitraria, la expansión de \((9+1)^d\) expresa la misma estratificación; sumar las contribuciones de la coordenada omitida demuestra su compatibilidad con la proyección.

El bloque central \(\bigl(\begin{smallmatrix}7&2\\2&7\end{smallmatrix}\bigr)\) proporciona un antecedente aritmético concreto. Sus autovalores son \(9\) y \(5\), y sus concatenaciones satisfacen \(72+27=77+22=99\). La restitución de una unidad y el refinamiento siguiente forman la cadena \((72,27)\mapsto(73,27)\mapsto(729,271)\). La primera operación incorpora la unidad registrada; la segunda cambia la escala y redistribuye una unidad mediante una corrección de suma cero. Es esta secuencia explícita la que permite relacionar centro, completación y proporción conservando la procedencia de cada término.

El refinamiento con acarreo introduce una evolución de esa distribución. La transformación \((x,y)\mapsto(Bx-c,By+c)\) conserva la suma después de multiplicarla por \(B\). Sobre etiquetas enteras admisibles, con \(0\leq c<B\), la división euclídea de la segunda componente recupera primero \(c\) y después el par precedente. La profundidad fija cuántas veces debe aplicarse esa inversa. Cuando se permiten acarreos levantados, su cociente adicional forma parte del registro conservado. La reversibilidad depende, por tanto, de datos precisos y tiene una operación de recuperación explícita.

Esta distinción permite dar un sentido operativo a la conservación holográfica. La lectura normalizada determina una coordenada; el registro conserva las relaciones que permiten reconstruir el estado o continuar su refinamiento. El desarrollo posterior hace lineal esa correspondencia: la biyección entre etiquetas admisibles induce un operador unitario entre sus bases y permite transportar amplitudes, correlaciones y observables.

\subsection*{La función aritmética e incidencial de \(K\)}
El registro \(K\) se obtiene del estado terminal mediante sus visitas, trazas, recuentos, diferencias y carga. La evaluación
\[
\kappa=\frac{N_K}{1000^{12}-1}
\]
es reversible en la carta fijada: multiplicar por el denominador restituye el entero y su separación en doce bloques recupera el registro completo, incluidos los ceros iniciales. La historia enriquecida tiene, además, las coordenadas de profundidad y transporte que utiliza cada prolongación. La recuperación de cada objeto se formula sobre los datos que lo individualizan.

En la lectura aritmética, \(\kappa\) interviene en
\[
\alpha_A=\pi_{\rm HMT}+e_{\rm HMT}-\varphi_{\rm HMT}-4-\kappa.
\]
Las tres coordenadas regionales y el registro proceden del mismo núcleo generador y actúan según su orden constructivo. La carta analítica completa de \(\alpha\) representa esa misma precoordenada mediante una raíz simple y una familia compatible de intervalos. En la lectura incidencial, el registro transporta sus relaciones de fase a la configuración excepcional. La unidad del objeto se manifiesta en la composición efectiva de sus lectores, sus inversas y sus acciones.

La órbita cíclica de \(K\) añade un resultado de identificación. Sus doce vectores forman una base de la realización lineal de doce componentes. Las respuestas a esa base determinan cualquier operador lineal de la clase considerada; para un operador que conmuta con el desplazamiento, una respuesta vectorial completa determina las restantes. El menor autovalor de Gram, \(589609\), proporciona una cota exacta de estabilidad. Dos observaciones complementarias de memoria recuperan también la parte centrada del registro; la carga uniforme completa su reconstrucción. La constante funciona así como un objeto de acoplamiento con capacidad de evaluación, transporte e identificación cuantificada.

\subsection*{El balance bilateral y la transformación de observables}
El entero \(73\) aparece como norma de dos elementos conjugados y como índice de una inclusión reticular de rango dos. La identidad \(10\cdot73=729+1\) relaciona esa norma ternaria con la completación decimal. El teorema bilateral extiende el balance a dos isometrías \(B,C\) entre espacios con dominio y codominio comunes. Para \(a>0\), las componentes \(aB-C\) y \((a+1)B+C\), ponderadas por \(a+1\) y \(a\), conservan el producto interior tras la normalización indicada. La expansión y cancelación de los términos cruzados prueban la identidad en dimensión arbitraria. La especialización \(a=8\) mantiene los operadores nonádicos y su información aritmética.

La conservación del estado completo permite calcular el efecto de una observación particular. Si el transporte relativo \(R\) es unitario y \(G\) es el observable expresado en la carta común, la lectura diagonal transforma
\[
\mathcal T(G)=\frac{72}{73}G+\frac1{73}R^*GR.
\]
Un observable permanece invariante exactamente cuando conmuta con \(R\). El observable transportado sobre el espacio completo conserva toda la lectura inicial. El resultado precisa, en una misma construcción, la información que persiste y el cambio introducido al observarla a través de un canal determinado.

\subsection*{Recuperación a profundidad arbitraria}
En cada paso se puede prolongar una componente y archivar su complemento. El balance finito suma las normas cuadradas de las memorias y del estado terminal. Una política explícita impone la cota geométrica \((729/1241)^N\) sobre la norma cuadrada del operador terminal nonádico. En el límite, toda la norma queda contenida en el archivo infinito; este archivo es una isometría y su adjunto reconstruye el estado con norma operatoria uno.

Los primeros bloques archivados admiten asimismo una inversa exacta corregida por su operador de Gram, con una cota de estabilidad uniforme para profundidades positivas. El archivo infinito, la inversa finita y la aproximación obtenida al truncar el adjunto tienen fórmulas y estimaciones propias. En las dinámicas generales donde subsiste un terminal, su contribución se conserva junto con la memoria. La política demostrada identifica el caso en que ese terminal tiende a cero.

El levantamiento unitario del refinamiento se compone con el balance bilateral. El ejemplo decimal que lleva \((73,27,1)\) a \((729,271)\) mantiene explícita la procedencia del acarreo; su comparación con un acarreo vecino determina el cambio exacto de una lectura. La conservación evolutiva de la unidad queda expresada mediante una secuencia de transformaciones recuperables.

\begin{figure}[htbp]
\centering
\begin{tikzpicture}[>=Stealth,every node/.style={font=\small,align=center},
block/.style={draw=ink,rounded corners=2pt,inner sep=6pt,text width=35mm},
wide/.style={draw=ink,inner sep=6pt,text width=47mm}]
\node[block] (d) {Estado aritmético\\admisible\\\((x,y,c)\)};
\node[block,right=14mm of d] (e) {Estado refinado\\\((10x-c,10y+c)\)};
\node[block,right=14mm of e] (j) {Bases de etiquetas\\\(J\) unitario};
\draw[<->] (d)--node[above,font=\scriptsize]{acarreo} (e);
\draw[->] (e)--(j);
\node[wide,below=17mm of e] (w) {Transporte bilateral\\\(W^*W=I\)};
\draw[->] (j.south)|-(w.east);
\node[block,below left=17mm and 4mm of w] (a) {Lectura conjunta\\\(\frac{72G+R^*GR}{73}\)};
\node[block,below right=17mm and 4mm of w] (b) {Observable transportado\\\(WGW^*\)};
\draw[->] (w)--(a);\draw[->] (w)--(b);
\node[font=\small,below=8mm of a] {Invariante si \(GR=RG\)};
\node[font=\small,below=8mm of b] {Recuperación mediante \(W^*\)};
\end{tikzpicture}
\caption{Composición de operaciones demostradas. El refinamiento conserva la suma normalizada y permite recuperar la etiqueta. Su levantamiento actúa sobre amplitudes de estados y conserva la norma. La lectura de un observable y su transporte completo conservan funciones diferentes, especificadas por las dos fórmulas inferiores.}
\label{libro:fig:composicion}
\end{figure}


\subsection*{Formas, acción y observación cuántica}
El transporte del producto interior se prolonga a cada grado exterior admisible. Áreas y volúmenes conservan también las contribuciones mixtas de lectura y memoria. La incidencia excepcional transporta estas identidades sobre sus imágenes especificadas; las cargas variacionales se obtienen de una acción declarada y de sus simetrías. Las realizaciones electromagnéticas y gravitatorias mantienen las escalas y operadores que intervienen en sus pruebas.

La realización gaussiana muestra por qué conservar el estado y conocer una observación reducida son cuestiones complementarias. En el protocolo declarado de dos entradas gaussianas puras, centradas e idénticas y un número finito de modos, el transporte simpléctico conserva la estructura global, mientras su reducción adquiere un espectro determinado por la componente antilineal de la memoria. El lazo elíptico--parabólico estudiado produce pureza \(9/11\), espectro \((9/10)10^{-j}\) y entropía adimensional \((10/9)\log10-\log9\), con todas las ocupaciones del espacio de Fock conservadas en el cálculo.

El argumento concluye, por ello, con una ley de transformación más rica que una igualdad escalar: se determina el valor, se conserva la estructura que interviene en su lectura y se calculan las condiciones de su recuperación. Las fórmulas de inversión, las cotas de estabilidad y las realizaciones especificadas expresan conjuntamente el significado matemático de la extrema y media razón holográfica.

\subsection*{Resultados rectores y secuencia de demostración}
\label{paper:resultados-rectores}
La exposición inicial anticipa la función de la constante de acoplamiento;
su construcción se efectúa después de definir las hojas de APP, la
orientación TRIT, los operadores TPK y las prolongaciones regionales.
La definición~\ref{lib01:definicion} especifica el registro ordenado, su
lectura posicional y la carta en la que se representan. La
sección~\ref{subsec:k-terminal} construye la extracción terminal y distingue
la información visible de la información firmada conservada en el estado.

La función aritmética del registro se establece en la construcción de
\(\alpha\). La sección~\ref{subsec:alpha-publicaciones-completas} contiene
la representación analítica completa, su recurrencia en todos los grados,
la convergencia de sus truncamientos y la identificación de su raíz con la
precoordenada generada. Su dependencia de esa precoordenada forma parte de
la definición del lector. Esta construcción precede a la escala de acción,
a las coordenadas angulares y a las realizaciones de área y gravitación.

La función de reconstrucción se expresa mediante resultados distintos y
componibles. La proposición~\ref{lib01:marco-ciclico} demuestra que la
órbita cíclica del registro constituye una base; la
proposición~\ref{lib01:operadores} determina qué respuestas permiten
recuperar un operador. El teorema~\ref{lib04:teo-balance} establece la
conservación bilateral para dos isometrías. El
teorema~\ref{lib04:teo-archivo} construye el archivo de profundidad
ilimitada y su inversa continua. Finalmente, los
teoremas~\ref{lib05:teo-biyeccion} y~\ref{lib05:teo-lector} incorporan
el refinamiento aritmético y calculan la transformación de los observables.

Estos resultados hacen explícita la jerarquía del artículo: generación del
estado, extracción de la constante, realizaciones aritméticas e
incidenciales, conservación de formas y recuperación efectiva de la
información. Las fórmulas de reconstrucción especifican en cada caso el
registro que reciben; la composición conserva el significado de sus
dominios, sus escalas y sus orientaciones.

\section*{Arquitectura generadora y conservación de la información}
\addcontentsline{toc}{section}{Arquitectura generadora y conservación de la información}
\label{libnar:troncal}

La Holografía Modular Triádica, abreviada HMT, sitúa el origen de la
construcción en un sistema finito de posiciones, operaciones y transportes.
De su prolongación se obtienen historias compatibles; de las historias,
lecturas numéricas e incidenciales; de estas realizaciones, las estructuras
con las que opera la Mecánica Dimensional. La pregunta que organiza este
artículo concierne a esa continuidad constructiva: qué permanece cuando una
unidad cambia de escala, cuando una operación modifica su representación
y cuando una parte del estado se transfiere a la memoria.

El contenido de la unidad se precisa a lo largo del recorrido. En una
partición aritmética, se conserva la suma normalizada de sus componentes.
En un espacio con producto interno, se conservan normas y correlaciones.
En una representación geométrica, se transportan formas, direcciones e
incidencias. Cada conservación tiene un objeto y una ecuación propios.
Su articulación permite seguir una misma construcción desde el acarreo
de una cifra hasta la recuperación de un estado distribuido en un archivo
de profundidad arbitraria.

La exposición siguiente presenta ese recorrido en su orden de dependencia.
Las definiciones y demostraciones completas aparecen después, dentro del
cuerpo matemático. Aquí interesa que el lector pueda reconocer el papel
de cada operación antes de estudiar su formulación general: cómo se
produce una historia, qué determina una lectura y de qué manera puede
reconstruirse la información que una representación visible ha dejado
de mostrar.

\subsection*{De las posiciones finitas a las historias operatorias}
\label{libnar:app}

La construcción denominada APP, por \emph{All Possible Paths}, parte de
dos ejes de nueve marcas y de dos operaciones sobre sus pares de
posiciones: suma y producto. Las posiciones tienen una identidad que se
conserva al evaluarlas. El par \((i,j)\) admite la lectura aditiva
\(i+j\) y la multiplicativa \(ij\); una ruta registra además cómo se
ha llegado a ese par, qué operación ha actuado y qué orientación ha
seguido el transporte. La relación entre acumulación aditiva y
composición multiplicativa queda así situada en un soporte común.

El paso modular conserva tanto el resto como el cociente. Para una
evaluación positiva \(n\), la marca nonádica \(r\), comprendida entre
uno y nueve, y el cociente \(q\) satisfacen
\[
 n=r+9q.
\]
Esta igualdad permite recuperar la evaluación entera desde su lectura
modular enriquecida. Dos posiciones pueden compartir la misma marca y
conservar cocientes diferentes. Por ejemplo, \((5,5)\) y \((8,2)\)
producen la suma diez; sus productos son veinticinco y dieciséis. Ambos
productos tienen marca siete, con cocientes dos y uno, respectivamente.
La diferencia operatoria queda localizada en el registro que acompaña a
la marca. Las definiciones de la \cref{sec:nucleo} desarrollan esta
conservación conjuntamente con las hojas de operación y las rutas.

La unidad ternaria orientada, denominada TRIT, discrimina tres regímenes
mediante los signos \(-1,0,+1\). Su acción incluye la elección de fase,
la orientación de las prolongaciones y el acarreo. La división ternaria
balanceada escribe un entero como \(n=r+3c\), con
\(r\in\{-1,0,+1\}\); el resto identifica el estado local y el
cociente conserva lo transportado al nivel siguiente. Las realizaciones
geométricas posteriores distinguen los regímenes elíptico, parabólico e
hiperbólico. En la lectura temporal del corpus, el retorno con memoria,
el umbral y la apertura bidireccional poseen acciones específicas sobre
el estado.

El \emph{Topological Phase Kernel}, o núcleo topológico de fase, se
abrevia TPK. Compone la selección trítica, el transporte y la
actualización de memoria. Una transición modifica posiciones, hojas o
fase, y conserva los datos necesarios para relacionar el estado anterior
con el siguiente. Su resultado es un estado enriquecido: contiene
residuos, cocientes, acarreos, orientación, frontera y prefijos de ruta.
El número que se leerá posteriormente procede de esa historia
operatoria. Su genealogía conserva las operaciones que permiten
continuarlo y las relaciones que permiten compararlo con otras
historias.

Esta organización da un significado preciso al retorno. El calendario
de ciento ocho posiciones se distribuye en doce ventanas de nueve. Al
completar una vuelta nonádica retorna la fase correspondiente y avanza
el registro de ventana. El regreso a una posición visible y la
recuperación del estado enriquecido son, por tanto, dos sucesos
diferenciados. La \cref{sec:extension} formaliza el calendario, la
holonomía y el incremento de memoria que acompaña a sus retornos.

\subsection*{La contracción del valor y la permanencia de la incidencia}
\label{libnar:continuo}

Una historia se prolonga mediante prefijos compatibles. Cada prefijo
conserva las operaciones efectuadas y delimita sus continuaciones
admisibles. La construcción conjunta del continuo mantiene en ese mismo
sistema el transporte graduado, los balances de memoria, las incidencias
ternarias, las relaciones de frontera y las composiciones
determinantales. Estas operaciones se enlazan por sus acciones sobre
las mismas historias y por sus compatibilidades con el refinamiento
(\cref{iv:continuo-conjunto}). La correlación entre ellas es parte de
la construcción del objeto.

La lectura arquimediana asocia a los prefijos cilindros de precisión
creciente. Sus intervalos compatibles se contraen y determinan un valor
cuando sus diámetros tienden a cero. La convención semiabierta distingue
los representantes de frontera durante el refinamiento; la información
de acarreo registra las continuaciones que alcanzan un mismo extremo.
En esta lectura aparece la realización numérica del continuo generado.
La secuencia de intervalos expresa cuánto se ha determinado del valor;
el estado enriquecido expresa cómo se ha determinado.

La lectura incidencial conserva otra información: qué elementos forman
un bloque, cómo se orienta una frontera, qué transportes relacionan los
prefijos y qué acciones preservan esas relaciones. A través de los mapas
desarrollados en la \cref{exc:seccion}, esta información se realiza en
matrices de Paley y Hadamard, diseños de Witt y sistemas de Steiner,
códigos de Golay, construcciones reticulares y el retículo de Leech.
La prolongación al álgebra de operadores de vértice \(V^\natural\)
sitúa allí la acción del grupo Monstruo. Cada paso conserva su portador,
su aplicación y sus condiciones de compatibilidad.

La relación entre número y simetría se vuelve así concreta. El valor
arquimediano identifica historias bajo una equivalencia numérica; la
incidencia conserva relaciones entre las historias y sus fronteras.
Ambas lecturas proceden de un mismo estado y cumplen funciones
complementarias. Una magnitud escalar puede permanecer fija mientras
la representación incidencial experimenta una transformación. Seguir
las dos lecturas permite estudiar conjuntamente la estabilización del
valor y la evolución de la estructura que lo realiza.

La unidad también se conserva al distribuirse entre prolongaciones.
Las medidas de los cilindros hijos suman la medida del cilindro padre.
En los espacios de funciones asociados, esta igualdad se convierte en
una isometría de refinamiento: una función del nivel anterior conserva
su norma al representarse sobre las nuevas posiciones. El mismo
principio relaciona partición, medida y transporte. Las demostraciones
de la construcción conjunta especifican la composición y el paso al
límite, de modo que cada nueva profundidad permanece vinculada a las
anteriores.

La realización de una medida añade una operación concreta a esta relación
entre estado y lectura. El sistema y el dispositivo se acoplan, producen
un registro y establecen qué distinciones resultan accesibles. En una
realización isométrica del acoplamiento, la descripción conjunta conserva
las correlaciones; el resultado registrado determina después la
actualización correspondiente. Evaluar un observable, condicionar el
estado a un resultado y preparar nuevamente el registro tienen funciones
diferenciadas. Esta secuencia permite seguir tanto la información que
llega a la lectura como la que permanece en el estado y en el dispositivo.
La memoria interviene en esa continuidad operativa entre una observación
y las siguientes (secciones~\ref{nuclear:3} y~\ref{lib04:archivo}).

\subsection*{La unidad dimensional y el acarreo recuperable}
\label{libnar:unidad}

La extrema y media razón holográfica se estudia aquí a través de las
operaciones que conservan la unidad durante su extensión y su
refinamiento. La posición añadida tiene una identidad propia. Si
\(E_9\) contiene las nueve clases nonádicas, su completación se escribe
\(E_9\sqcup\{\ast\}\). La clase residual cero ya pertenece a
\(E_9\) y tiene el representante positivo nueve; \(\ast\) designa
el estado adicional de frontera. La décima posibilidad conserva así una
función distinta de la reducción modular. En varias coordenadas, contar
cuántas han alcanzado \(\ast\) determina los estratos de la completación
y permite seguir la contribución de cada uno a la unidad.

La completación cúbica ofrece una primera representación
exacta. Al ampliar de nueve a diez las posiciones de cada coordenada,
la unidad cúbica se descompone en interior, estratos de ampliación y
vértice final:
\[
 10^3=9^3+3\cdot9^2+3\cdot9+1
     =729+243+27+1=729+271.
\]
Dividir por mil convierte el recuento en una partición de masa unidad.
Los términos conservan su procedencia dimensional: una, dos o tres
coordenadas ocupan la posición añadida. La expansión
\((9+1)^d\) conserva esta distinción para cada dimensión entera
positiva. El aumento de dimensión y el cambio de escala quedan
relacionados por mapas de coordenadas y sumas de estratos
(\cref{sec:revision-emrh,lib03:completacion}).

El acarreo introduce una transformación más fina. Consideremos la
partición \(73+27=100\). La multiplicación por diez produce
\(730+270=1000\). Transferir una unidad de la primera componente a
la segunda da
\[
 (73,27,1)\longmapsto(10\cdot73-1,10\cdot27+1)
                  =(729,271).
\]
El tercer dato es el acarreo. La suma final sigue siendo mil y la
partición normalizada sigue siendo uno. A la vez, el residuo decimal
de la segunda componente conserva la transferencia: de \(271\) se
recupera el resto uno, después \((271-1)/10=27\), y finalmente
\((729+1)/10=73\). El resultado lleva consigo los datos de la
operación que lo ha producido.

El teorema~\ref{lib03:teo-residuo} demuestra esta inversión para la aplicación
\[
 L_c(x,y)=(Bx-c,By+c),\qquad x+y=M,
\]
con base \(B\), acarreo canónico \(0\leq c<B\) y las condiciones
de positividad declaradas. Tras varias etapas, el registro acumulado
\(C_n\) reúne los acarreos como una palabra en base \(B\), y las
componentes satisfacen
\[
 x_n=B^n x_0-C_n,\qquad y_n=B^n y_0+C_n.
\]
La suma conserva \(B^nM\). Fijadas base y profundidad, el estado
final recupera la palabra completa de acarreos y la partición inicial
(\cref{lib03:cor-palabra}). La profundidad forma parte del dato de
reconstrucción, pues determina también las posiciones iniciales que
contienen ceros.

Las lecturas normalizadas convergen con cotas explícitas para la cola.
Así se enlazan tres niveles de descripción: una operación entera
reversible, una sucesión de particiones de la unidad y un sistema de
cilindros que determina una lectura límite. La elección efectiva de las
continuaciones sigue las reglas del transporte considerado. El ejemplo
\((729,271)\) muestra su primera transferencia; las pruebas de
refinamiento describen también la memoria de todas las transferencias
ulteriores y las coincidencias de frontera entre expansiones
(\cref{lib03:teo-limite,lib03:teo-cilindros}).

\subsection*{Las regiones numéricas y el registro de acoplamiento}
\label{libnar:registro}

En el corpus, las coordenadas \(\pi_{\rm HMT}\),
\(\varphi_{\rm HMT}\) y \(e_{\rm HMT}\) proceden de lectores
regionales sobre la prolongación APP--TRIT--TPK. La región de cierre,
la propagación y la relación estable entre modos poseen operadores
propios. Por ejemplo, el calendario de modos aditivo y multiplicativos
produce la matriz de incidencia
\(F_{\rm av}=\bigl(\begin{smallmatrix}0&1\\1&1\end{smallmatrix}\bigr)\).
Su dirección positiva estable determina la coordenada áurea en la
realización correspondiente. Las pruebas de generación precisan, para
cada profundidad solicitada, cómo se obtiene un prefijo certificado y
cómo se conserva su compatibilidad con los siguientes
(\cref{sec:generacion}).

El mismo estado terminal conserva un registro orientado de sus doce
ventanas. La extracción de los canales firmados, la carga total y la
inversión de Hadamard producen el vector ordenado \(K\). El término
\emph{constante holográfica de acoplamiento aritmético-geométrico}
designa ese registro con sus lectores y con el origen de fase fijado.
Su función consiste en relacionar una codificación aritmética exacta
con las realizaciones de incidencia y geometría que actúan sobre el
mismo objeto (\cref{lib01:definicion}).

En la representación integral canónica se obtiene
\[
 \begin{aligned}
 K=\bigl(&234,543,140,729,659,824,\\
         &621,58,914,794,146,601\bigr),
 \qquad\sum_{j=1}^{12}K_j=6263.
 \end{aligned}
\]
La \cref{lib01:registro} muestra la extracción completa, las
congruencias de sus canales y la divisibilidad que permite invertir
la matriz de Hadamard. El orden de sus componentes conserva la
posición de ventana. Esa información interviene tanto en la lectura
de cifras como en los operadores cíclicos y en las incidencias.

La suma total cumple una función de reconstrucción. Las diferencias
entre posiciones separadas por tres y cuatro pasos determinan
conjuntamente las variaciones relativas del registro. Añadir una misma
cantidad a las doce componentes conserva esas diferencias; la suma fija
precisamente ese nivel común. En \(K\), la carga \(Q(K)=6263\)
restituye la componente uniforme. Diferencias, carga y compatibilidades
enteras permiten recuperar el registro completo
(\cref{prop:k-transversal}). La relación global entre las posiciones y
la organización relativa de sus canales quedan conservadas por lecturas
complementarias.

En base mil, con doce posiciones declaradas, el registro se codifica
por la fracción periódica
\[
 \kappa_{\rm ag}=
 \frac{234543140729659824621058914794146601}{10^{36}-1}.
\]
Cada componente ocupa tres cifras; por ello, la componente cincuenta
y ocho se escribe \(058\) en la expansión. Multiplicar la fracción
por su denominador y efectuar doce divisiones euclídeas recupera el
registro exacto. La periodicidad de esta codificación describe una
lectura fijada del registro finito; la actualización de la historia
enriquecida conserva, por su parte, su propia dinámica de memoria.

La conexión con la coordenada \(\alpha_A\) se expresa en el lector
de clausura por
\[
 \alpha_A=
 \pi_{\rm HMT}+e_{\rm HMT}-\varphi_{\rm HMT}-4-\kappa_{\rm ag}.
\]
Las tres coordenadas regionales y el registro proceden del estado
generado y actúan en el orden constructivo descrito en
\cref{lib01:clausura}. La fórmula hace visible su acoplamiento
aritmético: la extracción de \(K\) conserva una parte de la
estructura que interviene en esa clausura. Las otras realizaciones
de \(K\) conservan sus lectores respectivos. De este modo, una
identidad escalar, una matriz de incidencia y un transporte de fase
pueden estudiarse desde la misma procedencia, con sus acciones
claramente diferenciadas.

\subsection*{Un registro capaz de identificar una transformación}
\label{libnar:identificacion}

La función del registro puede examinarse mediante una operación
concreta. Sea \(S\) el desplazamiento cíclico de sus doce
componentes. Las doce columnas
\(K,SK,\ldots,S^{11}K\) forman una matriz \(O_K\) invertible.
El resultado significa que la órbita cíclica de este registro alcanza
todas las direcciones de su espacio de doce componentes. La prueba
calcula el espectro del Gram \(O_K^*O_K\) y da cotas positivas
exactas para sus valores propios (\cref{lib01:marco-ciclico}).

Supongamos que un operador \(H\) respeta el desplazamiento cíclico,
es decir, \(HS=SH\). Al observar la respuesta vectorial \(y=HK\),
sus desplazamientos proporcionan las respuestas a todas las columnas
de \(O_K\). Se recupera entonces el operador completo mediante
\[
 H=O_yO_K^{-1},\qquad
 O_y=(y,Sy,\ldots,S^{11}y).
\]
La observación contiene doce componentes, y la hipótesis de
conmutación identifica la clase de transformaciones recuperables.
Dentro de esa clase, una sola respuesta vectorial determina todos
los coeficientes del filtro circular. Por ejemplo, la combinación
\(2I+3S-S^5\) queda identificada desde su acción sobre \(K\),
con las cotas de error del teorema~\ref{lib01:operadores}.

Este hecho ofrece una aplicación operativa de la estructura del
registro. Su orden y su espectro permiten usarlo como estado de
prueba para identificar transformaciones compatibles. Si el estado
y sus respuestas se transportan mediante una isometría completa de
memoria, el Gram de la órbita se conserva y la identificación sigue
siendo recuperable por el adjunto. La estabilidad procede de valores
propios calculados y de operadores inversores explícitos. El mismo
registro que participa en la clausura aritmética adquiere así una
función precisa en el reconocimiento de transportes.

La normalización también tiene una información asociada. El factor
unitario de \(O_K\) conserva orientación y ángulos; su Gram
conserva las amplitudes necesarias para reconstruir la matriz
original. El par formado por ambos factores recupera \(O_K\)
(\cref{lib01:polar}). Esta separación permite elegir una
representación normalizada y mantener, al mismo tiempo, la escala
exacta del registro.

\subsection*{Dos transportes y una conservación bilateral}
\label{libnar:bilateral}

El siguiente paso relaciona el refinamiento con una ley de
conservación de correlaciones. Sean \(\mathcal B\) y
\(\mathcal C\) dos isometrías que transportan el mismo espacio de
estados a un espacio de salida común. Para \(a>0\), se forman los
dos operadores
\[
 Q_-=a\mathcal B-\mathcal C,\qquad
 Q_+=(a+1)\mathcal B+\mathcal C.
\]
Las dos combinaciones reúnen una contribución común y una
contribución relativa con signos diferentes. Al sumar sus normas
con los pesos correspondientes, los términos cruzados se cancelan:
\[
 (a+1)\|Q_-v\|^2+a\|Q_+v\|^2
   =(2a+1)(a^2+a+1)\|v\|^2.
\]
El teorema~\ref{lib04:teo-balance} demuestra esta igualdad para cualquier
dimensión de los espacios y para los transportes isométricos
declarados. La misma cancelación conserva productos internos y
admite su formulación para formas preservadas por los transportes.

En el caso \(a=8\), el factor \(a^2+a+1\) vale setenta y tres.
La factorización de la isometría completa distribuye primero la
norma entre dos componentes con pesos
\[
 \frac{72}{73}\qquad\text{y}\qquad\frac{1}{73}.
\]
Después actúa una mezcla ortogonal que produce los canales
normalizados de \(Q_-\) y \(Q_+\). Los pesos tienen así un
origen operatorio preciso: \(72=8\cdot9\) y
\(73=8\cdot9+1\). Su función se conserva al cambiar la dimensión
o los transportes, dentro de las hipótesis de la ley bilateral
(\cref{lib04:factorizacion}).

La recuperación se ve directamente en las combinaciones de salida.
Si \(x=Q_-v\) e \(y=Q_+v\), entonces
\[
 x+y=(2a+1)\mathcal Bv,\qquad
 ay-(a+1)x=(2a+1)\mathcal Cv.
\]
Los dos transportes del estado inicial permanecen determinados por
el par de canales. El adjunto de cualquiera de las isometrías
recupera después el estado correspondiente. Con la normalización
completa, el operador \(W_a\) satisface \(W_a^*W_a=I\):
conserva la norma y ofrece una inversa exacta sobre su imagen.

Una lectura particular puede variar mientras se conserva toda esa
información. Cuando \(\mathcal C=\mathcal B R\), con
\(\mathcal B\) y \(R\) unitarios, este último como transporte relativo, la
lectura diagonal de un observable \(G\) en ambos canales produce
\[
 \mathcal T(G)=\frac{72G+R^*GR}{73}
 \qquad(a=8).
\]
Los observables compatibles con \(R\) permanecen fijos; los demás
registran la diferencia de sus dos acciones. Esta ecuación distingue
la conservación del estado completo y la respuesta de un lector
determinado (\cref{lib04:teo-lector}).

\subsection*{El mismo acarreo, ahora como estado transportable}
\label{libnar:ejemplo}

La relación anterior se aplica materialmente al acarreo. A cada
etiqueta admisible \((x,y,c)\) se asocia un vector de una base
ortonormal. La transformación entera \(L_c\) induce entonces el
operador \(J\) que lleva esa base a las etiquetas de salida
\((Bx-c,By+c)\). Su inversa recupera el acarreo mediante el
residuo. La biyección de etiquetas convierte a \(J\) en una
transformación unitaria (\cref{lib05:prop-unitaria}). Los valores
enteros designan las etiquetas; las normas corresponden a las
amplitudes de los vectores que las representan.

Para base diez y suma inicial cien, consideremos las dos etiquetas
\(d_1=(73,27,1)\) y \(d_2=(73,27,2)\). Sus imágenes son
\(e_1=(729,271)\) y \(e_2=(728,272)\). Fijemos el transporte
\(R\) que intercambia esas dos etiquetas y deja fijas las demás.
La composición usa \(\mathcal B=J\) y \(\mathcal C=JR\).
Esta elección proporciona un ejemplo explícito dentro del teorema
bilateral, con regla de transporte especificada antes de evaluarlo.

Sobre el estado \(|d_1\rangle\), los canales para \(a=8\)
resultan
\[
 Q_-|d_1\rangle=8|e_1\rangle-|e_2\rangle,
 \qquad
 Q_+|d_1\rangle=9|e_1\rangle+|e_2\rangle.
\]
Sus normas cuadradas son sesenta y cinco y ochenta y dos. El balance
exacto es \(9\cdot65+8\cdot82=1241=17\cdot73\), mientras
que la suma de ambos canales recupera \(17|e_1\rangle\).
Aplicar la inversa de etiquetas recupera la partición \((73,27)\)
y el acarreo uno.

La lectura fraccionaria de la primera componente ofrece además una
predicción exacta para este transporte declarado. Al evaluarla
diagonalmente en ambos canales normalizados se obtiene
\[
 \frac{72\cdot729+728}{73\cdot1000}
   =\frac{729}{1000}-\frac{1}{73000}.
\]
El lector diagonal de acarreo da \((72\cdot1+2)/73=74/73\).
El observable transportado desde la entrada, en cambio, conserva
el valor propio uno del acarreo original. La diferencia entre ambos
resultados permite identificar qué operación de lectura se ha
realizado. La información del estado permanece recuperable y la
lectura diagonal cuantifica cómo interviene el transporte relativo
(\cref{lib05:ejemplo}).

El ejemplo reúne las dos conservaciones del artículo. La suma
aritmética se conserva en las etiquetas refinadas y la norma se
conserva en sus representaciones unitarias. La conexión entre ambas
procede de la biyección y de su levantamiento a bases ortonormales.
Esta composición permite calcular lectores, invertir el transporte
y controlar el error manteniendo identificados los objetos sobre
los que actúa cada ley.

\subsection*{Memoria completa, simetrías y profundidad arbitraria}
\label{libnar:memoria}

La conservación bilateral se prolonga etapa a etapa. En cada nivel,
un canal continúa el estado y el otro registra un complemento de
memoria. Si \(R_Nv\) es el estado que alcanza el nivel \(N\) y
\(m_0,\ldots,m_{N-1}\) son los complementos registrados, el
balance acumulado es
\[
 \|v\|^2=\|R_Nv\|^2+\sum_{n=0}^{N-1}\|m_n\|^2.
\]
La igualdad identifica exactamente dónde permanece la norma
inicial. La recuperación finita utiliza el estado continuado y
todos los complementos disponibles; sus correlaciones se conservan
por la misma isometría conjunta.

Para la política de archivo del teorema~\ref{lib04:teo-archivo}, la
norma del canal continuado posee una cota geométrica uniforme.
En el caso \(a=8\), su norma cuadrada está acotada en cada etapa
por \(729/1241\). Al aumentar la profundidad, la componente
continuada tiende a cero y el archivo infinito conserva toda la
norma. Si \(A_n\) y \(D_n\) designan los bloques de
continuación y memoria, respectivamente, el estado se reconstruye
por la serie convergente
\[
 v=\sum_{n\geq0}R_n^*D_n^*m_n.
\]
El operador inversor tiene norma uno. Una perturbación de norma
\(\varepsilon\) en el archivo completo produce un error de
reconstrucción de norma a lo sumo \(\varepsilon\). Las
truncaciones finitas tienen sus propias cotas y, cuando se emplea
sólo su memoria, la inversa finita especificada en el mismo
desarrollo.

La incidencia del registro \(K\), su órbita cíclica y sus
direcciones distinguidas pueden transportarse mediante estas
isometrías. Se conservan los productos internos que determinan
ángulos y matrices de Gram; los proyectores y lectores se
transportan con sus dominios; la acción del adjunto recupera la
representación inicial. Las potencias exteriores transportan la
orientación y conservan los determinantes de Gram, que miden los
cuadrados de los volúmenes correspondientes. La dimensión del
portador puede ampliarse y la información geométrica de la
imagen permanece especificada por la isometría.

La relación con la conservación de Noether se formula en un nivel
adicional bien determinado. Para la acción discreta auxiliar
construida en la \cref{lib01:carga-noether}, una familia de
generadores que entrelaza los transportes produce una carga
conservada de la forma \(p_n^*A_nq_n\). Aquí intervienen una
acción variacional, sus ecuaciones y la simetría que las preserva.
El balance de norma y la carga variacional se articulan mediante
los transportes comunes, cada uno con sus hipótesis y con la
función que desempeña en la demostración.

El recorrido establece una relación operativa entre unidad,
número, memoria y geometría. La unidad aritmética admite
refinamientos reversibles; el estado enriquecido conserva el
registro de las operaciones; las realizaciones numérica e
incidencial extraen propiedades correlacionadas; el acoplamiento
bilateral mantiene información y correlaciones durante su
redistribución. La extrema y media razón holográfica adquiere así
un desarrollo demostrativo en el que partición, extensión,
transporte y recuperación pueden seguirse mediante aplicaciones
explícitas. Las páginas siguientes presentan esas aplicaciones,
sus pruebas y sus ejemplos completos.

\begin{HMTCompactSection}
\section*{Cinco principios holográficos de construcción y conservación}
\addcontentsline{toc}{section}{Cinco principios holográficos de construcción y conservación}

Los \emph{Elementos} ofrecen un antecedente preciso para presentar una
arquitectura matemática mediante operaciones declaradas y consecuencias
demostradas. El libro I distingue definiciones, postulados y nociones comunes.
Sus tres primeros postulados permiten trazar un segmento entre puntos,
prolongar un segmento en línea recta y describir un círculo con centro y radio
dados; el cuarto establece la igualdad de los ángulos rectos. El quinto
formula una condición de encuentro: dos rectas cortadas por una transversal
se encuentran, al prolongarse, del lado donde los ángulos interiores suman
menos de dos rectos. Son enunciados sobre los objetos y las construcciones de
esa geometría.\footnote{Euclides, \emph{Elementos}, libro I, postulados 1--5.
Texto griego de Heiberg y traducción de Richard Fitzpatrick, p.~7:
\url{https://farside.ph.utexas.edu/books/Euclid/Elements.pdf}. También pueden
consultarse por separado en la edición de David E. Joyce:
\url{https://mathcs.clarku.edu/~djoyce/elements/bookI/bookI.html}.}

La extrema y media razón aparece en el libro VI como una relación entre un
todo y sus partes: el todo es al segmento mayor como éste al menor. En una
traducción algebraica moderna, para longitudes positivas \(u>v\), se escribe
\((u+v)/u=u/v\); por tanto, la razón \(r=u/v\) satisface
\(r^2=r+1\). Esta escritura explica la proporción transmitida sin convertir
el lenguaje algebraico moderno en el punto de partida de Euclides. En HMT,
la coordenada de autoescala se obtiene desde una incidencia producida por el
calendario operativo; su reconocimiento mediante la misma ecuación pertenece
a una etapa posterior.\footnote{Euclides, \emph{Elementos}, libro VI,
definición 3, en la edición de Joyce:
\url{https://mathcs.clarku.edu/~djoyce/elements/bookVI/defVI3.html}.
La edición de Fitzpatrick imprime este mismo enunciado como definición 2
del libro VI, p.~156. Se conservan aquí ambas numeraciones editoriales.}

La comparación que sigue se sitúa en la organización de la construcción.
HMT declara un soporte aritmético, dos hojas relacionadas, una orientación
trítica y operaciones que prolongan estados conservando sus registros. Un
cilindro representa la información fijada por un prefijo; una recta
euclidiana es otro objeto, con otra definición y otras operaciones. Los
cinco principios de esta sección expresan relaciones efectivas entre los
objetos holográficos. Su número ofrece una exposición reconocible; no
establece una correspondencia término a término con los cinco postulados,
ni una refutación del quinto, ni una sustitución de los axiomas geométricos.

Estos principios tienen estatuto \emph{organizador}: reúnen condiciones y
teoremas que el cuerpo del artículo construye y prueba. Cada uno presenta primero
su sentido, después sus dominios y una identidad que permite comprobarlo.
La secuencia común es
\[
 \mathrm{APP}\longrightarrow\mathrm{TRIT}\longrightarrow\mathrm{TPK}
 \longrightarrow\text{estado enriquecido}
 \longrightarrow\text{estructura discreta conjunta del continuo}.
\]
Las realizaciones numéricas, lineales y geométricas leen después esa
estructura. Esta posición de las realizaciones mantiene visibles el origen
de los coeficientes y la información que cada lectura conserva.

\subsection*{Principio I. Unidad relacional de la generación}
\hypertarget{prinh:unidad}{}

\emph{Cada lectura se construye mediante operaciones tipadas sobre un
mismo soporte, conservando las relaciones que permiten continuarla.} Una
casilla APP admite simultáneamente una lectura aditiva y otra multiplicativa.
El paso a una cifra visible registra una parte de ambas operaciones; el
cociente, la orientación y la hoja conservan las condiciones de su
prolongación. La unidad consiste en esta procedencia relacional, que permite
seguir una misma ejecución a través de sus distintas realizaciones.

El dominio inicial es \(D_9^2\), con \(D_9=\{1,\ldots,9\}\). Sobre él
actúan \(S(i,j)=i+j\) y \(P(i,j)=ij\). Para un entero \(n\), la carta
positiva de residuo y cociente viene dada por
\[
 \rho_9(n)=1+((n-1)\bmod9),\qquad
 q_9(n)=\frac{n-\rho_9(n)}9,\qquad
 n=\rho_9(n)+9q_9(n).
\]
La evaluación levantada de APP conserva
\(\bigl((\rho_9(S),q_9(S)),(\rho_9(P),q_9(P))\bigr)\).
La identidad reconstruye las dos lecturas enteras; por ejemplo, distingue
los valores 10 y 19, que tienen el mismo residuo positivo y cocientes
diferentes. El resultado y su prueba se desarrollan en
\eqref{eq:residuo-cociente} y \eqref{eq:app-levantada}. La representación
por residuos \(0,\ldots,8\), utilizada en otras cartas del artículo, conserva
su cociente correspondiente: cambiar de representantes exige cambiar ambos
componentes de la división.

TRIT añade la representación orientada de \(\mathbb F_3\) por
\(\{-1,0,+1\}\), con su división exacta \(n=r+3c\). El selector operativo
\(M_{\rm ph}\) toma los valores \(+1,-1,0\) en las clases de fase
\(\{1,4,7\}\), \(\{2,5,8\}\), \(\{3,6,9\}\), respectivamente.
El transporte cardinal tiene otro dominio y otra función: mueve los
cursores sobre \(X_9=(\mathbb Z/9\mathbb Z)^2\). En el estado enriquecido
\(\widehat X\), TPK compone selección, transporte y actualización de
memoria,
\[
 U_t=\operatorname{Mem}_t\circ\operatorname{Tra}_t
                    \circ\operatorname{Sel}_t.
\]
La definición \eqref{eq:actualizacion} y el desarrollo de
\cref{tpk-int:estado,tpk-int:seleccion} especifican fase, cursores, hojas,
direcciones, residuos, cocientes y registros. El signo de fase y la dirección
cardinal se componen conservando sus tipos distintos.

La autoescala ofrece un ejemplo de generación que puede seguirse sin partir
de su valor final. El calendario \((+1,-1,0)\) produce, con la regla de
cambio y conservación de hoja, el ciclo modal
\((\Sigma,\Pi,\Pi)\). Sus tres aristas determinan
\[
 F_{\rm av}=\begin{pmatrix}0&1\\1&1\end{pmatrix}.
\]
En su realización lineal posterior, el rayo propio positivo normalizado
como \((1,x)^{\mathsf T}\) satisface \(x=\lambda\) y
\(1+x=\lambda x\). Así se obtiene \(x^2-x-1=0\); la raíz positiva se
reconoce como \(\varphi\). El teorema que sigue a \eqref{eq:fav} prueba
esta identificación después de construir las entradas de la matriz.
La relación con la proporción euclidiana es exacta en esta lectura de
autoescala y conserva la diferencia entre sus dos procedimientos de origen.

\subsection*{Principio II. Prolongación orientada y memoria de composición}
\hypertarget{prinh:prolongacion}{}

\emph{Una prolongación admisible añade un acontecimiento al estado y
transporta la memoria anterior mediante la misma operación.} La historia
permite distinguir recorridos que terminan en una posición o una fase
iguales. Este principio proporciona una manera precisa de hablar de
continuación: se conserva el prefijo efectivamente recorrido y se añade
una emisión permitida por el selector y el régimen presentes.

Sea \(X_k\) el conjunto de historias admisibles de longitud \(k\),
construidas desde las semillas del núcleo. Si \(\varepsilon\) es admisible
tras \(h\), la prolongación \(h\varepsilon\) pertenece a \(X_{k+1}\)
y la truncación satisface \(\rho_k(h\varepsilon)=h\). Se retienen la
etiqueta de emisión y las hojas anteriores, incluso cuando dos emisiones
tengan la misma lectura final. En una carta de memoria sobre un módulo
\(M\), la actualización tiene la forma
\(U_\varepsilon(m)=C_\varepsilon m+\kappa_\varepsilon\), con
\(C_\varepsilon\in\operatorname{End}(M)\) y
\(\kappa_\varepsilon\in M\). Si dos tramos son componibles en esa
carta, entonces
\[
 C_{\beta\alpha}=C_\beta C_\alpha,\qquad
 \kappa_{\beta\alpha}=C_\beta\kappa_\alpha+\kappa_\beta,
 \qquad U_{\beta\alpha}=U_\beta U_\alpha.
\]
La prueba de \cref{iv:cont:concatenacion} consiste en sustituir la primera
actualización en la segunda. La memoria se suma después del transporte que
le corresponde; la asociatividad hace independiente el resultado de la
subdivisión del trayecto.

La diferencia entre fase y memoria aparece ya en una lectura elemental.
Para \((w,r)\in\mathbb Z\times C_9\), con
\(\bar r\in\{0,\ldots,8\}\), se tiene
\[
 \sigma_9(w,r)=
 \left(w+\left\lfloor\frac{\bar r+1}{9}\right\rfloor,
       r+1\pmod9\right),\qquad
 \sigma_9^9(w,r)=(w+1,r).
\]
La segunda igualdad cuenta exactamente un cruce de 8 a 0. Su primera
coordenada registra el avance que la segunda ya no distingue. En el
calendario de doce ventanas, la sucesión exacta
\(0\to C_{12}\to C_{108}\to C_9\to0\) conserva el acarreo entre
ventanas. No se escinde: \(C_{108}\) tiene exponente 108 y
\(C_{12}\times C_9\), exponente 36. La prueba y la ley del acarreo están
en \cref{prop:k-extension}.

El retorno de fase adquiere así un significado operativo: permite
comparar lecturas realizadas en una misma posición del calendario mientras
el estado conserva lo sucedido entre ellas. La holonomía \(\Gamma_9\)
se obtiene de la dinámica enriquecida; \(K_{\rm ph}\) es su publicación
reducida posterior. El registro dodecafásico entero \(K\), que utiliza
además orientación e incidencias, tiene también su construcción y su dominio
propios. La igualdad de una fase, de un residuo o de una coordenada publicada
no identifica por sí sola las historias. Ésta es la información de partida
que se conserva al pasar al refinamiento.

\subsection*{Principio III. Compatibilidad del refinamiento y realización}
\hypertarget{prinh:refinamiento}{}

\emph{Refinar consiste en añadir resolución compatible con el prefijo,
manteniendo las relaciones de truncación y las masas de sus prolongaciones.}
La lectura a una escala más fina precisa qué aconteció dentro de un
cilindro; al olvidar ese detalle se recupera la descripción anterior.
La realización de un valor surge cuando los cilindros de su lectura se
contraen de manera compatible. La historia que determina esos cilindros
permanece en el dominio de la construcción.

En el árbol de \cref{iv:cont:medida}, el conjunto inicial es finito y no
vacío, y cada historia tiene un número finito y positivo de hijos
admisibles. Una masa inicial positiva de suma uno y probabilidades locales
\(p(\varepsilon\mid h)>0\), de suma uno en cada padre, producen
\[
 \mu(h\varepsilon)=\mu(h)p(\varepsilon\mid h),\qquad
 \sum_{\rho_k(y)=h}\mu(y)=\mu(h).
\]
Por tanto, en la realización posterior
\(H_k=L^2(X_k,\mu_k)\), el levantamiento
\((J_kf)(y)=f(\rho_k(y))\) es isométrico. Su adjunto es la esperanza
condicional
\[
 (E_kg)(h)=\sum_{\rho_k(y)=h}
                  \frac{\mu(y)}{\mu(h)}g(y),
 \qquad E_kJ_k=I.
\]
Agrupar la norma cuadrada por padres demuestra la isometría y agrupar el
producto interno demuestra la fórmula del adjunto
(\cref{iv:cont:esperanza}). El levantamiento de una lectura y su retorno
al nivel anterior forman así un par de operaciones explícitas.

Las hipótesis de ramificación finita y prolongación no vacía proporcionan
historias infinitas compatibles. Sus cilindros reciben las masas anteriores,
que determinan una medida sobre la sigma-álgebra cilíndrica. La elección de
\(p\) pertenece al lector declarado: el caso de hijos equiprobables da
\(p=1/d(h)\), sin imponer esa medida a cualquier otra lectura. Esta
construcción conserva multiplicidades de historias, incluidas las que
comparten un valor realizado. La prolongabilidad se refiere al árbol
admisible construido; una restricción adicional de rutas necesita conservar
expresamente esa propiedad.

Para la lectura numérica regional, cada bloque
\(b\in\mathbb F_3^6\) tiene un índice entero
\(\nu(b)=\sum_{j=1}^6 b_j3^{6-j}\), comprendido entre 0 y 728.
El lector fija \(N_0=m\) y produce
\[
 N_{n+1}=729N_n+\nu(b_{n+1}),\qquad
 I_n=\left[\frac{N_n}{729^n},\frac{N_n+1}{729^n}\right].
\]
Se cumple \(I_{n+1}\subseteq I_n\) y
\(\operatorname{diam}(I_n)=729^{-n}\). Los extremos izquierdos forman
una sucesión racional de Cauchy; su clase define primero la realización
interna de la lectura. Identificada posteriormente con la completación
usual de los racionales, corresponde al único punto de los intervalos
cerrados anidados. El teorema posterior a \eqref{eq:cilindro} demuestra
esta afirmación y trata por separado las convenciones de frontera. Para
los caracteres regionales, el lector efectivo de cada prefijo decimal
utiliza las horquillas calculables y la separación demostrada respecto de
sus fronteras, como precisa \cref{gen:prefijos-regionales}. Esa profundidad
arbitraria conserva el cuantificador «para cada longitud finita»; la
existencia de un límite por compatibilidad y contracción es una afirmación
distinta y no aporta, por sí sola, un algoritmo para cualquier historia
abstracta. Las historias que representan el mismo punto conservan sus
registros distintos.

\subsection*{Principio IV. Doble lectura correlativa y transporte de marco}
\hypertarget{prinh:doble}{}

\emph{Las lecturas numérica e incidencial proceden del mismo estado
enriquecido y conservan sus compatibilidades al prolongar y cambiar de
marco.} La primera precisa una realización mediante intervalos cada vez
más estrechos. La segunda retiene relaciones entre posiciones, signos y
fronteras. El vínculo holográfico consiste en sus mapas de origen y
transporte comunes, que permiten comparar qué información publica cada una.

Sea \(X_n^{\rm cal}\) el dominio de prefijos admisibles que incluyen
su marco inicial y los transportes sucesivos de marco. La lectura aritmética
utiliza los bloques de ese prefijo; la incidencial utiliza las palabras
\(w_j\in\mathbb F_3^6\) y los marcos \(f_j\). Para una matriz
\(A\in M_6(\mathbb F_3)\), el lector lineal
\[
 \gamma_A:\mathbb F_3^6\longrightarrow\mathbb F_3^{12},
 \qquad \gamma_A(w)=(w,wA)
\]
actúa sobre palabras escritas como filas. Su especialización en la matriz
Paley--Witt construida en \cref{exc:paley-codigo} produce el código
ternario y sus incidencias. La matriz se obtiene en esa construcción
finita; un número real terminal no selecciona sus entradas. El orden
\(w_6\to w_{12}\to w_{18}\to w_{24}\to w_{30}\to R_{36}\to G_9\)
mantiene las transformaciones y memorias utilizadas antes de cada lectura.

La compatibilidad puede comprobarse sobre el propio operador. Para
\(f\in\operatorname{GL}_6(\mathbb F_3)\), sea
\(A^f=fAf^{-1}\). Entonces
\[
 \gamma_{A^f}(wf^{-1})
      =\gamma_A(w)(f^{-1}\oplus f^{-1}),
\]
pues ambos lados son \((wf^{-1},wAf^{-1})\). En particular, la lectura
del prefijo
\[
 Q_n(x)=\bigl(f_j,
       \gamma_{f_jA_Wf_j^{-1}}(w_j)\bigr)_{0\leq j<n}
\]
satisface \(r_{n,m}Q_n=Q_m p_{n,m}\), donde \(p_{n,m}\) trunca
historias y \(r_{n,m}\) trunca registros. La causalidad de la
actualización de los marcos garantiza que ambos procedimientos retienen
exactamente los mismos términos. El resultado
\cref{exc:limite-inverso} construye su paso al límite. La covarianza lineal
es válida para el cambio de marco escrito; conservar además un soporte
coordenado o una métrica exige transportar esa estructura, o restringir los
cambios de marco al grupo que la preserva.

La construcción correlativa del continuo extiende esta compatibilidad al
transporte por hojas, al balance orientado, a la incidencia ternaria, al
complejo rectangular de puertos y a la memoria celular con sus rellenos y
líneas determinantes. Estos aspectos se construyen sobre las mismas
historias, no sobre cinco dominios elegidos después. Las identidades de
concatenación, naturalidad y borde aseguran su ensamblaje en
\cref{iv:rectangular-natural,iv:cont:memoria-natural,iv:cont:ensamblaje,iv:cont:mapa-correlativo}.
Los cinco principios expositivos de esta apertura no son una nueva
partición de esas cinco construcciones consustanciales. La recuperación
de una palabra mediante la primera proyección de \(\gamma_A\), por su
parte, se refiere a esa palabra: la recuperación de la historia enriquecida
completa exige conservar los demás registros que la distinguen.

\subsection*{Principio V. Conservación y recuperación del registro completo}
\hypertarget{prinh:conservacion}{}

\emph{La información recuperable se conserva conjuntamente en el estado
que continúa, en los registros producidos y en sus correlaciones.} Una
lectura visible puede contraerse mientras el registro completo conserva una
representación isométrica del estado de partida. Este principio permite
formular una reversibilidad verificable: se declara qué se almacena, en qué
espacio reside y mediante qué operador se reconstruye.

El transporte de historias levantadas de \cref{iv:cont:transporte} actúa
entre los espacios \(H_k=\ell^2(\widehat X_k)\) mediante
\[
 \mathcal U_ke_{\widehat h}=
 \sum_{\varepsilon\ {
m admisible}}
 \sqrt{p(\varepsilon\mid\widehat h)}\,
 e_{\widehat h\varepsilon},\qquad
 \mathcal U_k^*\mathcal U_k=I.
\]
La segunda igualdad procede de la normalización local y de que padres
distintos tienen hijos con etiquetas distintas. Si \(P_k\) lee la hoja
\(\Sigma\) y \(Q_k=I-P_k\) la hoja \(\Pi\), las operaciones
\[
 \begin{aligned}
 A_k&=P_{k+1}\mathcal U_kP_k+Q_{k+1}\mathcal U_kQ_k,\\
 \eta_k&=Q_{k+1}\mathcal U_kP_k+P_{k+1}\mathcal U_kQ_k
 \end{aligned}
 \qquad\text{satisfacen}\qquad A_k^*A_k+\eta_k^*\eta_k=I.
\]
La identidad de Gram de \cref{iv:cont:gram} se demuestra al anular los
productos de proyectores complementarios. Distingue permanecer en una
hoja de cambiar de hoja; el valor trítico cero continúa siendo una
orientación y no se convierte en una tercera hoja.

Escribamos \(F_0=I\), \(F_{n+1}=A_nF_n\). La telescopía de la
identidad anterior da, para todo horizonte \(N\),
\[
 \|x\|^2=\|F_Nx\|^2+
                 \sum_{n=0}^{N-1}\|\eta_nF_nx\|^2.
\]
El primer término es el residuo terminal que todavía continúa; los otros
son las salidas registradas. El teorema \cref{iv:cont:terminal} demuestra
que \(F_N^*F_N\) converge fuertemente a un operador positivo
\(T_\infty\). Se obtiene, por tanto, el registro completo
\[
 \mathcal V:H_0\longrightarrow H_0\oplus
                     \bigoplus_{n\geq0}H_{n+1},\qquad
 \mathcal Vx=\bigl(T_\infty^{1/2}x,
                         (\eta_nF_nx)_{n\geq0}\bigr),
 \qquad\mathcal V^*\mathcal V=I.
\]
En efecto, la suma de las normas cuadradas es \(\|x\|^2\); la
polarización da la última identidad. Su imagen \(\mathcal M\) es
cerrada y la inversa sobre ella es \(\mathcal V^*\). El terminal se
conserva hasta demostrar su extinción; \(T_\infty\) no necesita ser
un proyector ni se supone la convergencia de \(F_Nx\).

La conservación alcanza también los operadores que leen esa información.
Para un operador acotado \(B\) sobre \(H_0\), el transporte
\(B^{\mathcal V}=\mathcal VB\mathcal V^*|_{\mathcal M}\) conserva
productos, adjuntos y productos internos:
\[
 (B_1B_2)^{\mathcal V}=B_1^{\mathcal V}B_2^{\mathcal V},\qquad
 \langle\mathcal Vx,B^{\mathcal V}\mathcal Vy\rangle
                       =\langle x,By\rangle.
\]
Insertar \(\mathcal V^*\mathcal V=I\) prueba ambas fórmulas. En el
lector nonádico de \cref{lib01:isometria}, las mismas identidades se
realizan con \(T_n=(8I+C_n)/9\) y
\(D_n=\sqrt8(C_n-I)/9\), para unitarias \(C_n\) del portador
declarado. La composición mantiene su estatuto de lector y no identifica
automáticamente cada \(C_n\) con la holonomía generadora \(\Gamma_9\).
La recuperación preserva así las correlaciones del registro y los
observables transportados, dentro de la imagen precisa donde actúa.

El registro \(K\) permite ver esta distinción en un objeto finito. Su
codificación \(\kappa(K)\), en base \(B=1000\) y longitud doce, es
\[
 \kappa(K)=\frac{N(K)}{B^{12}-1},\qquad
 N(K)=\sum_{j=1}^{12}K_jB^{12-j},\qquad
 K\in\{0,\ldots,999\}^{12}.
\]
Multiplicar por \(B^{12}-1\) y efectuar doce divisiones euclídeas
recupera el registro exactamente, incluidos sus bloques iniciales nulos.
La proposición \cref{lib01:inversion-kappa} prueba además que un error
absoluto menor que \(1/[2(B^{12}-1)]\) permite recuperar el mismo entero
mediante redondeo. La reversibilidad se refiere al registro en su carta y
con sus metadatos; no identifica a éste con toda la historia TPK. La
\emph{constante holográfica de acoplamiento aritmético-geométrico} designa
en este artículo el objeto \(K\) tipado con sus lectores y relaciones. Su
publicación procede del estado enriquecido y participa después en la
clausura de alfa: no se utiliza como entrada generadora de \(\pi\), ni
su denominación convierte cada lectura escalar en un invariante universal.

\subsection*{Una lectura conjunta de los cinco principios}

La secuencia ofrece un criterio positivo para recorrer el artículo. La unidad
relacional fija qué se opera; la prolongación conserva lo recorrido; el
refinamiento organiza la resolución; la doble lectura mantiene juntos valor
e incidencia; la recuperación identifica las condiciones exactas para
reconstruir el registro y sus observables. Las demostraciones se enlazan
porque actúan sobre objetos con procedencia común y respetan sus mapas de
transporte. Allí donde una lectura olvida información, su dominio y su
fibra de pérdida permanecen definidos.

La extrema y media razón histórica muestra cómo una relación entre el todo
y sus partes puede determinar una proporción. La construcción holográfica
desarrollada aquí estudia, además, relaciones conservadas cuando cambian la
resolución, el marco y la distribución de información entre continuidad y
memoria. La ecuación de autoescala establece un enlace matemático preciso
con aquella proporción; las identidades de concatenación, refinamiento e
isometría proporcionan el contenido adicional que este artículo desarrolla.
La comparación histórica conserva así su función orientadora, mientras
cada conclusión holográfica descansa en sus operaciones y pruebas propias.

\subsection*{Fuentes primarias de la comparación histórica}

\begin{enumerate}[label=\arabic*.]
 \item Euclides, \emph{Elementos}, libro I, postulados 1--5. Texto
 griego establecido por J. L. Heiberg, con traducción de Richard Fitzpatrick,
 edición corregida de 2008, p.~7; libro VI, definición 2 en esta edición,
 p.~156. \url{https://farside.ph.utexas.edu/books/Euclid/Elements.pdf}.
 \item Euclides, \emph{Elementos}, edición digital de David E. Joyce,
 Clark University: libro I, postulados 1--5,
 \url{https://mathcs.clarku.edu/~djoyce/elements/bookI/bookI.html};
 libro VI, definición 3,
 \url{https://mathcs.clarku.edu/~djoyce/elements/bookVI/defVI3.html}.
 La diferencia de numeración del libro VI se conserva para facilitar el
 cotejo entre ediciones.
\end{enumerate}

\end{HMTCompactSection}
\HMTDivision{\(\mathcal A\)}{Conservación evolutiva de la unidad dimensional por doble proyección holográfica}{Las dos hojas aritméticas, la orientación trítica y el transporte con memoria constituyen el punto de partida de las realizaciones del continuo.}
\HMTNarrativeHeading{Una unidad, dos realizaciones}

La doble proyección organiza este artículo desde su primera construcción. Un
mismo estado conserva los datos que permiten precisar un valor y los datos
que permiten determinar una incidencia. La primera realización reúne
cilindros compatibles hasta individualizar su lectura numérica; la segunda
transporta las relaciones de frontera. La tesis estudia su procedencia
común y las operaciones que mantienen su correspondencia durante el
refinamiento. Número y configuración geométrica reciben así una genealogía
explícita.

La unidad se constituye mediante posiciones, operaciones y relaciones de
transporte. APP dispone las dos direcciones nonádicas sobre un mismo
soporte y evalúa sus hojas aditiva y multiplicativa. Cada evaluación
conserva el residuo y su cociente. TRIT distingue la orientación y el
régimen local; el TPK compone esos datos con el desplazamiento, el acarreo,
la frontera y la prolongación. La estructura discreta conjunta del
continuo reside en esa construcción enriquecida. Sus realizaciones
posteriores reciben operaciones ya determinadas.

La conexión nonádica tiene aquí una función constitutiva. Una vuelta
compone las transiciones del estado completo. La fase observable retorna
y el registro incorpora lo recorrido. Esta diferencia permite reconocer
una regularidad mientras la historia aumenta de profundidad. La
continuidad del refinamiento conserva los prefijos que hacen compatibles
las dos realizaciones. Cada precisión nueva pertenece a la misma historia
de la que proceden las anteriores.

El sentido de conservación evolutiva aparece en este punto. Una cantidad
normalizada puede permanecer igual mientras cambian la distribución, la
resolución o la representación. Recuperar el estado exige identificar qué
datos acompañan ese cambio. La aritmética conserva cocientes; el transporte
conserva orientación y vueltas; la incidencia conserva relaciones; el
archivo operatorio conserva los complementos necesarios para invertir la
lectura. El artículo desarrolla los mapas que articulan esas funciones.

La extrema y media razón holográfica expresa esta cuestión de unidad,
proporción y transformación. Su formulación incluye la completación
dimensional y el refinamiento reversible, y alcanza después los balances
de estados y observables. La constante de acoplamiento ocupa una posición
determinada dentro de ese recorrido: el registro generado vincula la
lectura aritmética, la incidencia y la recuperación operatoria. Las
demostraciones siguientes construyen los antecedentes de cada una de esas
funciones.

\HMTMathematicalPaper
\section{Generación, conservación y realización}
\label{nuclear:introduccion}
El objeto de la Holografía Modular Triádica es un estado generado por posiciones y operaciones que conserva la información de su evolución. La evaluación numérica y la incidencia de frontera son realizaciones correlacionadas de ese antecedente. La contracción de cilindros compatibles determina valores; la prolongación de sus incidencias mantiene relaciones que distinguen los estados generadores. Esta prioridad de la construcción gobierna también la Mecánica Dimensional: acción, duración, extensión y masa aparecen mediante lectores cuyo dominio conserva la orientación, el acarreo y la historia de las operaciones.

El presente artículo estudia una cuestión estructural precisa: qué se conserva cuando una lectura reduce el estado, dónde reside el complemento y cómo reconstruir a partir de él tanto el valor como la dinámica. Una respuesta completa exige tres niveles. Primero se construyen las emisiones y su refinamiento. Después se demuestra una identidad de transporte sobre el estado realizado. Finalmente se determina qué observa un lector concreto y qué consecuencias físicas produce su composición con los lectores de acción y escala. El desarrollo sigue ese orden en cada aplicación.

\subsection{Tesis y enunciados rectores}
La conservación evolutiva de la unidad adquiere una expresión operatoria. Para dos transportes \(B,C\) entre espacios provistos de formas \(b,b'\), con
\[
B^{*}b'B=C^{*}b'C=b,
\]
la representación colectiva de las nueve posiciones produce
\begin{equation}
\mathcal U(B,C)=\frac19
\begin{pmatrix}8B+C&\sqrt8(C-B)\\ \sqrt8(C-B)&B+8C\end{pmatrix},
\qquad
\mathcal U(B,C)^*(b'\oplus b')\mathcal U(B,C)=b\oplus b.
\label{nuclear:eq:rectora}
\end{equation}
La primera columna contiene la lectura y la memoria generadas por un estado sin memoria entrante. La segunda columna describe la contribución de memoria que vuelve a intervenir. La igualdad conserva una forma positiva cuando \(b\) es un producto interior, y conserva área orientada cuando \(b\) es una forma simpléctica. Sus dos interpretaciones utilizan estructuras explícitamente diferentes.

Esta ley permite demostrar cuatro consecuencias conectadas. La primera es la reconstrucción estable del registro completo, incluido su terminal. La segunda recupera de la memoria dodecafásica la dirección que determina el proyector dimensional. La tercera prolonga la conservación a áreas, volúmenes y grados exteriores superiores. La cuarta identifica, en una realización cuántica declarada, la parte de la memoria que altera la pureza de una lectura parcial. Cada consecuencia se obtiene por una operación matemática especificada sobre la anterior.

\begin{figure}[htbp]
\centering
\begin{tikzpicture}[every node/.style={font=\small,align=center},b/.style={draw,rounded corners,inner sep=6pt,text width=39mm},>=Stealth]
\node[b] (g) {Emisión y estado\\APP--TRIT--TPK};
\node[b,right=12mm of g] (u) {Transporte completo\\y forma conservada};
\node[b,right=12mm of u] (l) {Lectura y memoria\\con reconstrucción};
\node[b,below=12mm of g] (i) {Incidencia y\\refinamiento};
\node[b,below=12mm of u] (a) {Acción y\\realización dimensional};
\node[b,below=12mm of l] (s) {Covarianza, espectro\\y entropía reducida};
\draw[->] (g)--(u);\draw[->] (u)--(l);\draw[->] (g)--(i);
\draw[->] (u)--(a);\draw[->] (l)--(s);\draw[->] (i)--(a);\draw[->] (a)--(s);
\end{tikzpicture}
\caption{Dependencia de las realizaciones. Las flechas representan operaciones demostradas en el cuerpo del artículo; la lectura escalar conserva un antecedente más rico que su propio valor.}
\end{figure}

\subsection{Memoria completa y observación parcial}
Para una sucesión de compresiones \(T_n\) con complementos \(D_n\), se escribe
\(P_0=I\), \(P_{n+1}=T_nP_n\). La igualdad local
\(T_n^*T_n+D_n^*D_n=I\) implica
\[
I=A_\infty+\sum_{n\ge0}P_n^*D_n^*D_nP_n,
\qquad
A_\infty=\underset{N\to\infty}{\operatorname{s-lim}}P_N^*P_N.
\]
La reconstrucción utiliza \(A_\infty^{1/2}v\) y todos los \(D_nP_nv\). La convergencia fuerte del Gram terminal basta para la conservación; se demuestra por monotonía positiva, con independencia de la convergencia del vector terminal. El registro contiene amplitudes y correlaciones, mientras sus intensidades proporcionan el balance de norma. Esta distinción determina la información recuperable.

En un registro finito concreto, la incidencia de doce posiciones en las hexadas de Witt posee Gram
\(36I+30\mathbf1\mathbf1^*\). Su normalización conserva exactamente la parte centrada del registro. Al incorporar la media se recuperan sus doce coordenadas. El mapa dimensional posterior utiliza esa información transportada y conserva la dependencia respecto del registro ordenado \(K\).

\subsection{Una consecuencia espectral calculable}
Sea \(J^2=-I\), \(J^{\mathsf T}=-J\), y sea \(H\) simpléctico en una realización canónica de dimensión \(2d\). Se define
\[
D=\frac{\sqrt8}{9}(H-I),\qquad
D_{\mathrm{anti}}=\frac{D+JDJ}{2},\qquad
W=\frac{8I+HH^{\mathsf T}}9.
\]
El artículo demuestra
\begin{equation}
\boxed{-(JW)^2=I+[D,J][D,J]^{\mathsf T}
=I+4D_{\mathrm{anti}}D_{\mathrm{anti}}^{\mathsf T}.}
\label{nuclear:eq:espectro}
\end{equation}
Para dos entradas gaussianas puras idénticas y el transporte \(\mathcal U(I,H)\), la matriz \(W\) es la covarianza reducida normalizada. Por tanto,
\[
\operatorname{Tr}\rho_{\mathrm{vis}}^2
=\det\bigl(I+4D_{\mathrm{anti}}D_{\mathrm{anti}}^{\mathsf T}\bigr)^{-1/4}.
\]
La fórmula relaciona un operador de memoria con el funcional de pureza del estado reducido. La pureza global permanece igual a uno. En particular, la mezcla depende de la componente antilineal respecto de la estructura compleja transportada, y no de una identificación indiscriminada entre memoria y entropía.

La elección de la representación gaussiana, del estado de entrada y del subsistema leído forma parte del enunciado. La genealogía fija los transportes que se componen; el protocolo fija la observación cuya distribución se calcula. Esta especificación permite una comparación reproducible entre realizaciones y separa el teorema de cualquier atribución experimental ulterior.

\subsection{Orden de exposición}
Los capítulos iniciales desarrollan íntegramente las primitivas comunes, la prolongación, las regiones numéricas, el registro y la incidencia excepcional. Siguen los lectores de acción y las representaciones dimensionales. La segunda parte reúne las pruebas de conservación, recuperación, holonomía y extensión reversible. La parte final compone esos resultados con acciones físicas y obtiene el espectro de la memoria reducida. La bibliografía sitúa las herramientas de representación después de la construcción que realizan.

La simetría finita se expresa mediante acciones y entrelazadores. La simetría variacional continua se expresa mediante una acción, una variación y su carga conservada. El teorema de Noether se aplica en este segundo dominio; el transporte de incidencia excepcional conserva su propio contenido algebraico. La relación entre ambos se establece por los observables y las formas transportadas, manteniendo explícitas sus diferencias de tipo.

\section{Núcleo formal de la Holografía Modular Triádica}
\label{sec:nucleo}

La Holografía Modular Triádica, abreviada HMT, comienza con una estructura finita de posiciones y operaciones. El orden de construcción importa: las posiciones reciben evaluaciones aritméticas, las transiciones adquieren una firma ternaria y la dinámica conserva la información necesaria para prolongarlas. Sólo después se forman las realizaciones numéricas. Una constante reconocida al final de este proceso no sustituye a la regla que produjo su región ni a la historia que determina su evaluación.

La denominación \emph{All Possible Paths}, APP, expresa el papel de los caminos sobre el soporte. Se adopta como denominación de la construcción, con una referencia conceptual a la formulación por caminos de la mecánica cuántica; no supone que su ley aritmética sea la integral de caminos de Feynman. El \emph{Topological Phase Kernel}, TPK,\footnote{Los autores dedican asimismo las iniciales TPK a Tan Pengkiang. Esta dedicación es independiente de la definición matemática.} designa la composición de selección, transporte, actualización de memoria y prolongación. Las siglas se mantienen únicamente después de haber definido los objetos que representan.

\subsection{APP: soporte posicional y evaluaciones aritméticas}

Consideremos las nueve marcas interiores de una recta graduada entre cero y diez,
\[
 D_9=\{1,2,3,4,5,6,7,8,9\}.
\]
Su disposición horizontal y vertical produce las ochenta y una posiciones de \(D_9^2\). Las marcas interiores no deben confundirse con los diez intervalos unitarios de la recta de referencia. Esta disposición fija un alfabeto y dos direcciones; todavía no utiliza una expansión real para decidir una trayectoria.

La identificación cíclica de las posiciones conduce al soporte
\[
 X_9=(\ZZ/9\ZZ)^2.
\]
La carta de representantes positivos conserva el símbolo nueve para la clase nula. Para todo entero positivo se definen
\begin{equation}
 \rho_9(n)=1+((n-1)\bmod9),\qquad
 q_9(n)=\frac{n-\rho_9(n)}9,
 \qquad n=\rho_9(n)+9q_9(n).
 \label{eq:residuo-cociente}
\end{equation}
La raíz digital positiva es, por tanto, una elección de representante de la clase residual. No conserva por sí sola el entero anterior a la reducción; la coordenada \(q_9\) proporciona precisamente la información que falta.

En cada posición \((i,j)\), la acumulación de \(i\) repetida \(j\) veces produce \(ij\). Las dos evaluaciones enteras y sus proyecciones son
\[
 S(i,j)=i+j,\qquad P(i,j)=ij,\qquad
 \Sigma(i,j)=\rho_9(i+j),\quad \Pi(i,j)=\rho_9(ij).
\]
La simetría \((i,j)\mapsto(j,i)\) intercambia las direcciones de acumulación y conserva ambas evaluaciones. Es una propiedad de un único soporte, no una identificación posterior entre dos matrices independientes.

\begin{definition}[Evaluación aritmética levantada]
La evaluación de APP que conserva las dos hojas es
\begin{equation}
 \mathcal A(i,j)=
 \bigl((R^+_{ij},q^+_{ij}),(R^\times_{ij},q^\times_{ij})\bigr)
 =\bigl((\rho_9(i+j),q_9(i+j)),(\rho_9(ij),q_9(ij))\bigr).
 \label{eq:app-levantada}
\end{equation}
El término \emph{hoja} indica aquí una evaluación sobre el mismo soporte: aditiva o multiplicativa. Las dos coordenadas de cociente pertenecen al estado aritmético y no se descartan al reducir los residuos.
\end{definition}

\begin{figure}[tbp]
\centering
\begin{tikzpicture}[x=.53cm,y=.53cm,font=\small]
\fill[black!5] (1.5,-.5) rectangle (2.5,8.5);
\fill[black!5] (4.5,-.5) rectangle (5.5,8.5);
\fill[black!5] (-.5,2.5) rectangle (8.5,3.5);
\fill[black!5] (-.5,5.5) rectangle (8.5,6.5);
\fill[black!12] (7.5,-.5) rectangle (8.5,8.5);
\fill[black!12] (-.5,-.5) rectangle (8.5,.5);
\foreach \i in {1,...,9}{
 \node at (\i-1,9.05) {\(\i\)};
 \node at (-1.05,9-\i) {\(\i\)};
 \foreach \j in {1,...,9}{
  \pgfmathtruncatemacro{\v}{mod(\i*\j-1,9)+1}
  \node at (\j-1,9-\i) {\(\v\)};
 }
}
\foreach \k in {0,...,9}{
 \draw[black!25] (\k-.5,-.5)--(\k-.5,8.5);
 \draw[black!25] (-.5,\k-.5)--(8.5,\k-.5);
}
\draw[line width=.65pt] (2.5,3.5) rectangle (4.5,5.5);
\node[anchor=west,align=left,text width=6.5cm] at (10,7.2) {Un soporte: \(81\) posiciones.\\[5pt]Dos evaluaciones:\\ \(S(i,j)=i+j\), \(P(i,j)=ij\).\\[5pt]Cada fila se genera por acumulación:\\ \(P(i,j+1)-P(i,j)=i\).};
\node[anchor=west,align=left,text width=6.5cm] at (10,1.8) {La carta representa \(0\pmod9\) por \(9\).\\[5pt]El bloque señalado es\\ \(B_c=\left(\begin{smallmatrix}7&2\\2&7\end{smallmatrix}\right)\),\\ con autovalores \(9\) y \(5\).};
\end{tikzpicture}
\caption{Tabla multiplicativa de APP en representantes positivos. Las filas y columnas no unitarias se distinguen por trama gris; el marco identifica el bloque diagonal adyacente de diagonales iguales. La tabla dibujada es un observable del toro aritmético; sus bordes de representación no constituyen una frontera topológica intrínseca del toro.}
\label{fig:app}
\end{figure}


\begin{proposition}[Reconstrucción de las evaluaciones]
Los dos pares de \eqref{eq:app-levantada} determinan exactamente \(i+j\) e \(ij\). Los residuos aislados no determinan esas evaluaciones enteras ni sus cocientes.
\end{proposition}
\begin{proof}
La primera afirmación es \eqref{eq:residuo-cociente} aplicada a cada hoja. Para la segunda, \(10=1+9\) y \(19=1+2\cdot9\) tienen el mismo residuo y distinto cociente. En el soporte bidimensional, las posiciones \((5,5)\) y \((8,2)\) dan la misma suma y el mismo residuo producto, pero productos \(25=7+18\) y \(16=7+9\). La proyección residual identifica datos que la evaluación levantada distingue.
\end{proof}

\subsubsection{Grafo de Cayley, topología toroidal y grado de enrollamiento}

Las traslaciones cardinales están dadas, en coordenadas de fila y columna, por
\[
 \nu_N=(-1,0),\quad\nu_E=(0,1),\quad
 \nu_S=(1,0),\quad\nu_O=(0,-1),\qquad T_d(x)=x+\nu_d.
\]
La primera coordenada aumenta hacia el sur y la segunda hacia el este. El grafo
\[
 \mathscr G_9=\operatorname{Cay}(X_9,\{\nu_N,\nu_E,\nu_S,\nu_O\})
\]
es la retícula cíclica bidimensional: tiene ochenta y un vértices y cuatro aristas orientadas salientes por vértice. Es un grafo de Cayley para la estructura \emph{aditiva} del soporte. La multiplicación de residuos tiene divisores de cero y no convierte todas las filas de APP en un grupo multiplicativo.

Una palabra cardinal \(w=d_1\cdots d_r\) y un origen \(x_0\) determinan
\[
 x_k=x_0+\sum_{j=1}^k\nu_{d_j}\pmod9.
\]
Existen \(81\cdot4^r\) caminos orientados de longitud \(r\), incluidos retornos y retrocesos. Cuando el camino se cierra, su desplazamiento entero previo a la reducción define
\begin{equation}
 \operatorname{wind}(w)=\frac19
 \bigl(\#S(w)-\#N(w),\#E(w)-\#O(w)\bigr)\in\ZZ^2.
\end{equation}
La inversión del camino cambia el signo de este vector. El retorno a la misma posición residual no obliga, por tanto, a que el enrollamiento sea nulo. Éste es el primer ejemplo de una diferencia que reaparecerá en el retorno de fase: una coordenada puede cerrarse mientras otra registra una evolución no trivial.

\subsubsection{Cociente de transporte e identidad de cociclo}

Sea \(s_9:\ZZ/9\ZZ\to D_9\) la sección de representantes positivos. El cociente de transporte asociado a la composición aditiva, denominado también \emph{acarreo} en aritmética posicional, es
\[
 c_9(a,b)=\frac{s_9(a)+s_9(b)-s_9(a+b)}9.
\]
El carácter composicional de esta memoria queda expresado por una identidad exacta.

\begin{lemma}[Identidad de cociclo]
Para \(a,b,c\in\ZZ/9\ZZ\),
\[
 c_9(a,b)+c_9(a+b,c)=c_9(b,c)+c_9(a,b+c).
\]
\end{lemma}
\begin{proof}
Ambas sumas son iguales a
\[
 \frac{s_9(a)+s_9(b)+s_9(c)-s_9(a+b+c)}9.
\]
La asociatividad de la suma de representantes levantados requiere precisamente este término adicional al componer en el cociente.
\end{proof}

Con representantes positivos, este cociclo no está normalizado en la identidad: \(c_9(0,a)=1\). Si \(s_0\) toma representantes en \(\{0,\ldots,8\}\), su acarreo normalizado \(c_0\) satisface
\[
 c_9(a,b)=c_0(a,b)+\mathbf1_{a=0}+\mathbf1_{b=0}-\mathbf1_{a+b=0}.
\]
La diferencia es el cambio de sección; no altera la identidad de cociclo. La descomposición posterior del calendario en fase y ventana utilizará los representantes de cero a ocho.

Los totales de ambas evaluaciones distinguen la contribución visible y la conservada en el cociente:
\[
 \sum S=810,\quad\sum P=2025,\quad
 \sum\Sigma=405,\quad\sum\Pi=459,
\]
\[
 \sum q^+=45,\qquad\sum q^\times=174.
\]
De aquí se obtiene
\begin{equation}
 2025-810=(459-405)+9(174-45)=54+9\cdot129.
 \label{eq:defecto-integrado}
\end{equation}
La diferencia suma--producto no se agota en el defecto visible \(54\). La identidad conserva también el defecto de cociente \(129\). Suprimirlo cambiaría la aritmética que se transporta, aunque la tabla residual permaneciera idéntica.

\subsubsection{Reducción ternaria y radical multiplicativo}

El anillo \(R=\ZZ/9\ZZ\) tiene grupo de unidades \(\{1,2,4,5,7,8\}\) e ideal máximo
\[
 \mathfrak m=(3)=\{0,3,6\},\qquad\mathfrak m^2=0.
\]
Toda fila aditiva es una permutación de los nueve representantes. Una fila multiplicativa lo es exactamente cuando su índice es una unidad. Las filas tres y seis contienen tres valores repetidos; la fila nueve contiene únicamente el representante de la clase nula. Por ello el plegamiento multiplicativo ya distingue, antes de cualquier geometría continua, las unidades del sector nilpotente.

La aplicación \(x\mapsto3x\) tiene núcleo e imagen iguales a \(\mathfrak m\). Induce un isomorfismo de espacios vectoriales sobre \(\FF_3\),
\[
 R/\mathfrak m\longrightarrow\mathfrak m,\qquad [x]\longmapsto3x,
\]
pero no un isomorfismo de anillos. En particular, la secuencia visible \(3,6,9\) admite dos lecturas diferentes: su reducción directa módulo tres es \(0,0,0\), mientras que la extracción del factor tres seguida de la reducción de su coordenada interna da \(1,2,0\). Esta segunda operación conserva una coordenada del ideal; no es la misma proyección que reducir directamente el número original.

\subsubsection{Realización latina de la evaluación aditiva}

La diferencia entre las dos evaluaciones puede expresarse también mediante
operadores. Sea \((e_r)_{r\in\ZZ/9\ZZ}\) la base ortonormal de \(\CC^9\).
La asignación \(v_{ab}=e_{a+b}\) produce una base ortonormal en cada fila
y columna. Los proyectores \(P_{ab}=v_{ab}v_{ab}^{*}\) satisfacen
\[
 \sum_bP_{ab}=I_9=\sum_aP_{ab},\qquad P_{ab}^2=P_{ab}.
\]
En efecto, cada traslación \(b\mapsto a+b\) permuta los nueve residuos.
Se obtiene una realización por proyectores de rango uno de un cuadrado
latino, en la terminología de los cuadrados latinos cuánticos
\cite{mustovicary,delascuevas2023}. Esta construcción parte de la hoja
aditiva ya definida y explicita el mapa hacia su representación operatoria.

La evaluación multiplicativa conserva otro contenido: para \(a=3\), los
vectores \(e_{3b}\) repiten los residuos \(0,3,6\), y
\[
 \sum_{b\in\ZZ/9\ZZ}|e_{3b}\rangle\langle e_{3b}|
 =3\bigl(|e_0\rangle\langle e_0|+|e_3\rangle\langle e_3|
                    +|e_6\rangle\langle e_6|\bigr).
\]
La multiplicidad queda registrada en el operador y caracteriza el radical
multiplicativo. Por ello el soporte toroidal, la propiedad latina y la
resolución operatoria de la identidad son propiedades que requieren
definiciones diferentes. Los cuadrados mágicos cuánticos estudiados por
De las Cuevas, Drescher y Netzer exigen entradas positivas y sumas de filas
y columnas iguales a la identidad; su teoría distingue además las
dilataciones con entradas conmutativas\cite{delascuevas2020}. Esta referencia
proporciona un lenguaje preciso para comparar realizaciones de APP.

\subsection{TRIT: orientación y regímenes cuadráticos}

\begin{definition}[Firma ternaria y levantamiento local]
La firma TRIT utiliza la identificación balanceada
\[
 \FF_3\longrightarrow\TRIT,\qquad 0\mapsto0,\quad1\mapsto+1,\quad2\mapsto-1.
\]
Para \(n\in\ZZ\), su levantamiento conserva \(n=r+3c\), donde \(r\in\TRIT\) y \(c\in\ZZ\). En una transición de amplitud local acotada, los tres estados \((0,-1),(0,0),(0,+1)\) pueden presentar el mismo residuo nulo y distintos acarreos.
\end{definition}

La firma no reemplaza a la trayectoria. Un cero residual no significa ausencia de información de orientación, fase o cociente. Tampoco introduce por sí solo una magnitud física. Su función es distinguir algebraicamente las clases de transporte que se realizarán a continuación.

\subsubsection{Tres regímenes cuadráticos y una exponencial común}

Una vez obtenida la firma ternaria, su realización cuadrática se define por
\begin{equation}
 \mathbb A_\tau=\RR[X]/(X^2+\tau),\qquad \tau\in\TRIT.
 \label{eq:algebras-trit}
\end{equation}
Si \(J_\tau\) es la clase de \(X\), entonces \(J_\tau^2=-\tau I\). La estructura distingue el tipo elíptico \(\tau=+1\), el parabólico \(\tau=0\) y el hiperbólico \(\tau=-1\). La realización real no genera la firma: expresa los tipos que la reducción orientada ya ha producido.

\begin{proposition}[Exponenciación de los tres tipos]
En cada álgebra \eqref{eq:algebras-trit},
\begin{equation}
 \exp(tJ_\tau)=C_\tau(t)I+S_\tau(t)J_\tau,
\end{equation}
donde
\[
 C_\tau(t)=\sum_{n\ge0}\frac{(-\tau)^nt^{2n}}{(2n)!},\qquad
 S_\tau(t)=\sum_{n\ge0}\frac{(-\tau)^nt^{2n+1}}{(2n+1)!}.
\]
Las tres realizaciones son
\[
 \begin{array}{c|c|c}
 \tau&J_\tau^2&\exp(tJ_\tau)\\\hline
 +1&-I&\cos t\,I+\sin t\,J_\tau\\
 0&0&I+tJ_\tau\\
 -1&I&\cosh t\,I+\sinh t\,J_\tau.
 \end{array}
\]
\end{proposition}
\begin{proof}
Separar los términos pares e impares de la serie exponencial y sustituir \(J_\tau^{2n}=(-\tau)^nI\) y \(J_\tau^{2n+1}=(-\tau)^nJ_\tau\) da las fórmulas. Para \(\tau=0\), todos los términos de grado al menos dos se anulan. En los otros dos casos se recuperan las series trigonométricas o hiperbólicas.
\end{proof}

La exponencial es común a los tres regímenes: no se asigna exclusivamente a la rama parabólica. Los invariantes
\[
 C_\tau(t)^2+\tau S_\tau(t)^2=1
\]
se deducen de \(\exp(tJ_\tau)\exp(-tJ_\tau)=I\). El mismo principio algebraico conserva la unidad a través de realizaciones que difieren por su tipo cuadrático.

En el régimen hiperbólico, \(P_\pm=(I\pm J_{-1})/2\) son proyectores complementarios y
\[
 \exp(tJ_{-1})=e^tP_++e^{-t}P_-.
\]
El crecimiento y la contracción aparecen conjuntamente en dos direcciones propias. En el elíptico, el parámetro describe una rotación; en el parabólico, una traslación nilpotente. La interpretación temporal adoptada por HMT asocia \(+1\) al retorno con memoria, \(0\) al umbral y \(-1\) a la apertura bidireccional. Esta interpretación exige un orden de parametrización; no añade un cuarto tipo algebraico.

% Ampliación acumulativa del TRIT; se inserta antes de la figura de regímenes.
% No sustituye ni elimina el desarrollo precedente.
\subsubsection{División balanceada y composición de los levantamientos}
\label{trit-ampl-div-balanceada}

La reducción ternaria permite expresar una evaluación entera de APP en una
coordenada visible y una coordenada de prolongación. Su función no consiste
únicamente en abreviar la escritura de un entero: proporciona una regla de
composición que determina cómo se modifica el cociente cuando se suman o se
multiplican evaluaciones. La elección balanceada hace explícita, además, la
compatibilidad de esa regla con la inversión de signo.

\begin{definition}[Coordenadas ternarias balanceadas]
\label{trit-ampl-def-coordenadas}
Para \(n\in\mathbb Z\), sean
\[
 q_3(n)=\left\lfloor\frac{n+1}{3}\right\rfloor,\qquad
 r_3(n)=n-3q_3(n).
\]
La aplicación
\[
 \mathcal B:\mathbb Z\longrightarrow\TRIT\times\mathbb Z,\qquad
 n\longmapsto(r_3(n),q_3(n))
\]
se denomina descomposición ternaria balanceada. Su inversa es
\(\mathcal L(r,c)=r+3c\).
\end{definition}

En efecto, \(r_3(n)\in\{-1,0,+1\}\). Si dos pares representan el mismo entero,
\(r+3c=s+3d\), entonces \(r-s\) es múltiplo de \(3\) y pertenece al intervalo
\([-2,2]\); necesariamente \(r=s\) y \(c=d\). La representación es, por tanto,
única. La división puede repetirse sobre \(q_3(n)\), de manera que cada nivel
conserva el dato necesario para reconstruir el anterior. El cociente no es una
perturbación del residuo: constituye su coordenada complementaria.

La firma obtenida de una evaluación \(F\), con \(F=S\) o \(F=P\) en APP,
es \(r_3\circ F\). Cuando la reducción actúa sobre el cursor de fase, la misma
sección balanceada produce las clases \(\{1,4,7\}\), \(\{2,5,8\}\) y
\(\{3,6,9\}\), con valores \(+1,-1,0\). El selector del TPK que se define
a continuación utiliza esas clases para determinar el modo operativo.
La lectura de una evaluación y la selección de un modo tienen así la misma
reducción ternaria, pero dominios y funciones diferentes.

\Needspace{16\baselineskip}
\begin{proposition}[Aritmética de los pares balanceados]
\label{trit-ampl-prop-aritmetica}
Sean \(n=r+3c\) y \(m=s+3d\), con \(r,s\in\TRIT\). Defínase
\[
 \chi(r,s)=\frac{r+s-r_3(r+s)}3.
\]
Las operaciones enteras se transportan a los pares mediante
\begin{align}
 (r,c)\boxplus(s,d)
 &=\bigl(r_3(r+s),c+d+\chi(r,s)\bigr),
 \label{trit-ampl-eq-suma}\\
 (r,c)\boxtimes(s,d)
 &=\bigl(rs,rd+sc+3cd\bigr).
 \label{trit-ampl-eq-producto}
\end{align}
La aplicación \(\mathcal L\) conserva ambas operaciones.
\end{proposition}
\begin{proof}
La igualdad \(r+s=r_3(r+s)+3\chi(r,s)\) da
\[
 n+m=r_3(r+s)+3\bigl(c+d+\chi(r,s)\bigr).
\]
Por otra parte,
\[
 nm=rs+3(rd+sc+3cd).
\]
Como el producto de dos representantes balanceados pertenece de nuevo a
\(\{-1,0,+1\}\), no necesita una reducción residual adicional. La unicidad
de la descomposición demuestra las dos fórmulas.
\end{proof}

Por ejemplo, \(2=(-1)+3\) y \(4=1+3\). Su suma y su producto se reconstruyen
como
\[
 (-1,1)\boxplus(1,1)=(0,2),\qquad
 (-1,1)\boxtimes(1,1)=(-1,3).
\]
Las evaluaciones resultantes son \(6\) y \(8\). La suma de cocientes basta
en este ejemplo aditivo; en el producto intervienen los términos cruzados
\(rd\), \(sc\) y \(3cd\). La discrepancia entre ambas operaciones reside
también en la coordenada de prolongación, no solamente en el residuo.

El término \(\chi\) es un cociclo de la sección balanceada. Si
\(r\oplus s=r_3(r+s)\), satisface
\[
 \chi(r,s)+\chi(r\oplus s,t)
 =\chi(s,t)+\chi(r,s\oplus t).
\]
Ambos miembros representan el cociente de la misma suma de tres enteros.
Esta identidad garantiza que reagrupar las operaciones no cambia la
evaluación reconstruida. Al mismo tiempo, la proyección sobre el primer
componente suprime esa contabilidad: \(1+1=-1+3\), mientras que la lectura
residual conserva únicamente \(-1\). La información eliminada está
determinada, en este caso, por \(\chi(1,1)=1\).

La relación con el módulo \(9\) conserva una segunda escala discreta.
La aplicación
\[
 \beta:\TRIT\times\mathbb F_3\longrightarrow\mathbb Z/9\mathbb Z,
 \qquad (r,\bar c)\longmapsto r+3c\pmod9
\]
está bien definida y es biyectiva. Cambiar el representante de \(\bar c\)
añade un múltiplo de \(9\); recíprocamente, una igualdad de imágenes fija
primero \(r\) por reducción módulo \(3\) y después \(\bar c\). La suma
transportada por esta biyección contiene el término
\(\chi(r,s)\) de \eqref{trit-ampl-eq-suma}, reducido ahora módulo \(3\).
Así, los dos niveles ternarios no se componen mediante una suma
independiente de coordenadas: la normalización del primer nivel modifica
el segundo.

En particular, las clases \(3\), \(6\) y \(0\) de
\(\mathbb Z/9\mathbb Z\) tienen primer componente \(0\), pero segundo
componente \(1\), \(2\) y \(0\), respectivamente. La distinción ya observada
entre reducción directa y extracción de la coordenada del ideal aparece
como una propiedad de \(\beta\). Conservar esa coordenada impide que el
sector residual nulo se convierta artificialmente en una sola posición.
El módulo \(1000\), utilizado después para organizar bloques decimales,
cumple otra función de representación posicional; su cociente se conserva
mediante la división correspondiente. Un cambio de base conserva el entero
cuando retiene residuo y cociente, mientras que la coincidencia de residuos
en bases distintas no identifica por sí sola sus estados de transporte.

\Needspace{11\baselineskip}
\subsubsection{Inversión, clase residual y coordenadas no visibles}
\label{trit-ampl-inversion}

La inversión aditiva se levanta sin ambigüedad:
\[
 \mathcal B(-n)=(-r_3(n),-q_3(n)).
\]
En particular, \(3\), \(0\) y \(-3\) tienen coordenadas
\((0,1)\), \((0,0)\) y \((0,-1)\). Comparten firma nula, pero representan
evaluaciones distintas. En una transición restringida a un cociente de
amplitud máxima \(1\), son las tres posibilidades de la fibra visible sobre
cero; en un levantamiento entero general, esa fibra contiene todos los pares
\((0,c)\), con \(c\in\mathbb Z\).

La coordenada de cociente recupera el entero evaluado, pero no recupera por
sí sola la trayectoria que lo produjo. Dos caminos de APP pueden compartir
evaluación, residuo y cociente, y diferir en orientación, sucesión de modos o
enrollamiento. Por ello la descomposición \(\mathcal B\) se integra en el
estado de transporte, en vez de sustituirlo. La memoria aritmética de una
división y la memoria de una trayectoria son coordenadas relacionadas, no
sinónimos.

Para una sucesión de firmas de longitud \(n\), la inversión orientada es
\[
 C_{\rm rev}(\tau_1,\ldots,\tau_n)
   =(-\tau_n,\ldots,-\tau_1).
\]
Aplicarla dos veces restituye la sucesión. La operación invierte el orden de
las posiciones y el signo de cada firma; conserva su longitud y las
posiciones nulas con el orden invertido. En la parametrización adoptada,
intercambia el retorno con memoria, asociado a \(+1\), y la apertura
bidireccional, asociada a \(-1\), manteniendo el umbral \(0\). El transporte
de los restantes datos de una trayectoria se especifica mediante la
involución correspondiente del TPK.

Esta inversión de firmas debe distinguirse de la conjugación dentro de un
álgebra cuadrática. La primera puede cambiar el tipo \(\tau\); la segunda,
que se construye a continuación, actúa dentro del tipo ya fijado. La
distinción impide atribuir a una inversión de signo una equivalencia
algebraica entre regímenes diferentes.

\subsubsection{Producto cuadrático, conjugación y norma algebraica}
\label{trit-ampl-norma}

Un elemento de \(\mathbb A_\tau\) tiene expresión única \(aI+bJ_\tau\).
La relación \(J_\tau^2=-\tau I\) determina su producto:
\begin{equation}
 (aI+bJ_\tau)(cI+dJ_\tau)
   =(ac-\tau bd)I+(ad+bc)J_\tau.
 \label{trit-ampl-producto-cuadratico}
\end{equation}
Se trata de una realización real del tipo ya seleccionado. Las coordenadas
\(a,b\) no son residuos de APP, y el producto
\eqref{trit-ampl-producto-cuadratico} no sustituye a
\eqref{trit-ampl-eq-producto}; expresa otra operación, en el codominio
cuadrático declarado.

\begin{definition}[Conjugación y norma algebraica]
\label{trit-ampl-def-norma}
La conjugación del tipo \(\tau\) y su norma algebraica son
\[
 \overline{aI+bJ_\tau}=aI-bJ_\tau,\qquad
 N_\tau(aI+bJ_\tau)=a^2+\tau b^2.
\]
\end{definition}

\begin{proposition}[Multiplicatividad e invertibilidad]
\label{trit-ampl-prop-norma}
Para \(z,w\in\mathbb A_\tau\),
\[
 z\bar z=N_\tau(z)I,\qquad
 N_\tau(zw)=N_\tau(z)N_\tau(w).
\]
Además, \(z\) es invertible si y sólo si \(N_\tau(z)\ne0\), y en ese caso
\[
 z^{-1}=\frac{\bar z}{N_\tau(z)}.
\]
\end{proposition}
\begin{proof}
El producto de \(aI+bJ_\tau\) y \(aI-bJ_\tau\) es
\((a^2+\tau b^2)I\). La conjugación es multiplicativa y el álgebra es
conmutativa, por lo que
\((zw)\overline{zw}=(z\bar z)(w\bar w)\). Si la norma no se anula, la
fórmula indicada proporciona la inversa. Recíprocamente, si \(zw=I\), la
multiplicatividad implica \(N_\tau(z)N_\tau(w)=1\).
\end{proof}

La forma cuadrática \(N_\tau\) es positiva definida únicamente en el tipo elíptico.
Para \(\tau=0\), vale \(a^2\), y el elemento \(J_0\) es no nulo aunque su
cuadrado y su norma sean cero. Para \(\tau=-1\), vale \(a^2-b^2\); los
elementos no nulos \(I+J_{-1}\) e \(I-J_{-1}\) tienen norma cero y su
producto se anula. La norma algebraica es aquí una forma cuadrática
multiplicativa: no presupone una norma positiva de espacio vectorial.

Los tres tipos pueden realizarse fielmente mediante
\[
 aI+bJ_\tau\longmapsto
 \begin{pmatrix}a&-\tau b\\b&a\end{pmatrix}.
\]
El determinante de esta matriz es \(N_\tau(aI+bJ_\tau)\). Quedan
representadas así, con una misma expresión matricial, la estructura
compleja, la estructura de números duales y la estructura real escindida.
Cada una conserva su ley y sus divisores de cero. La forma de firma
\((1,1)\) del último caso tiene interpretación lorentziana en dimensión
\(1+1\); la identificación de una métrica física requiere además un espacio,
un campo y una regla de transporte específicos.

\subsubsection{Conservación multiplicativa y composición de evoluciones}
\label{trit-ampl-evoluciones}

La exponencial ya obtenida pertenece al conjunto de elementos de norma uno:
\[
 N_\tau\bigl(\exp(tJ_\tau)\bigr)=1.
\]
Esta conservación se mantiene al componer parámetros. En efecto, todos los
factores son funciones del mismo generador, y
\[
 \exp(tJ_\tau)\exp(sJ_\tau)=\exp((t+s)J_\tau).
\]
Al comparar las dos componentes se deducen las leyes de adición
\begin{align}
 C_\tau(t+s)&=C_\tau(t)C_\tau(s)-\tau S_\tau(t)S_\tau(s),\\
 S_\tau(t+s)&=S_\tau(t)C_\tau(s)+C_\tau(t)S_\tau(s).
\end{align}
La conjugación cambia \(t\) por \(-t\) y coincide con la inversión dentro
de este subgrupo. No cambia \(\tau\): invierte una evolución del régimen
elegido.

En el tipo elíptico, la forma conservada es \(a^2+b^2\), y la trayectoria
exponencial recorre el círculo de norma uno. En el hiperbólico, las
coordenadas propias son \(u=a+b\), \(v=a-b\), con \(uv=N_{-1}(z)\).
La exponencial da
\[
 (u,v)=(e^t,e^{-t}),\qquad uv=1.
\]
El crecimiento de una coordenada se acompaña de la contracción de la otra.
La conservación expresa una relación entre ambas coordenadas; no afirma que
cada una permanezca constante.

En el tipo parabólico,
\[
 (I+tJ_0)(I+sJ_0)=I+(t+s)J_0,\qquad
 (I+tJ_0)^{-1}=I-tJ_0.
\]
La evolución no se detiene porque \(J_0^2=0\). La nilpotencia elimina los
términos de orden al menos dos, pero conserva el término lineal que
transporta el parámetro. Son distintos el símbolo visible \(0\), el
elemento nulo del álgebra y el generador no nulo \(J_0\). Esta distinción
es necesaria para describir el umbral sin identificarlo con ausencia de
estado o de dinámica.

\begin{proposition}[Composición de coordenadas afines]
\label{trit-ampl-prop-pendientes}
En la carta donde la componente escalar no se anula, represéntese una
dirección mediante \(I+uJ_\tau\). Si \(1-\tau uv\ne0\), la dirección de su
producto con \(I+vJ_\tau\) tiene coordenada
\[
 u\star_\tau v=\frac{u+v}{1-\tau uv}.
\]
\end{proposition}
\begin{proof}
Se tiene
\[
 (I+uJ_\tau)(I+vJ_\tau)
   =(1-\tau uv)I+(u+v)J_\tau.
\]
Dividir por la componente escalar da la expresión anunciada.
\end{proof}

La carta afín conserva la razón entre componentes, no el factor escalar
dividido. Cuando \(1-\tau uv=0\), debe examinarse el producto antes de
cambiar de carta: puede tener dirección definida con componente escalar
nula, o ser el elemento cero en presencia de divisores de cero. La ley de
composición sigue perteneciendo al álgebra completa; la singularidad de
una coordenada no equivale automáticamente a una singularidad del objeto.
Esta fórmula reaparecerá en la lectura regional elíptica. Su uso requiere
mantener el denominador y, cuando corresponda, el factor de normalización.

\subsubsection{Transporte del tipo y realizaciones posteriores}
\label{trit-ampl-puentes}

Las álgebras \(\mathbb A_{+1}\), \(\mathbb A_0\) y
\(\mathbb A_{-1}\) no son isomorfas como álgebras reales unitarias.
La primera es un cuerpo; la segunda contiene un nilpotente no nulo; la
tercera posee idempotentes no triviales y es isomorfa a
\(\mathbb R\times\mathbb R\). Por ello, el cambio de firma que organiza una
ruta no puede interpretarse como un isomorfismo indiferenciado entre las
tres realizaciones. El TPK conserva el índice de régimen, los dominios y el
operador efectivo que actúa en cada transición.

Dentro del capítulo electrónico, el par de proyecciones de las hojas de APP
produce un conmutador normalizado cuyo cuadrado es \(-I\). El mismo bloque
contiene una involución de cuadrado \(I\) y dos direcciones nilpotentes. La
realización conjunta se demostrará sobre ese par concreto de proyecciones;
la clasificación precedente explica por qué las tres relaciones cuadráticas
pueden aparecer en una misma álgebra matricial sin identificarse entre sí.

En la prolongación excepcional, la matriz de Paley reducida módulo \(3\)
satisface \(A_W^2=-I_6\). La relación dota al espacio ternario de una
estructura de dimensión \(3\) sobre \(\mathbb F_9\). Su vínculo con una raíz
operatorial real se expresa mediante la presentación integral
\[
 \mathbb Z[u]/(u^2+1).
\]
Los cambios de base a \(\mathbb F_3\) y a \(\mathbb R\) producen
representaciones de esa misma relación. Cada una conserva su característica,
su dimensión y su información de incidencia. El desarrollo de los
respectivos operadores establece el enlace, no la igualdad de sus
cardinales.

La conservación de norma uno describe, por consiguiente, las realizaciones
cuadráticas de las evoluciones. La conservación de cocientes, orientación y
trayectoria pertenece al estado aritmético transportado. Ambas propiedades
intervienen en la construcción, con funciones diferentes: una determina las
identidades de las representaciones; la otra permite componer, reconstruir y
prolongar las operaciones que las originan.


\begin{figure}[htbp]
\centering
\begin{tikzpicture}[scale=.88,>=Stealth,font=\small]
\foreach \c/\titulo in {0/{Elíptico},4.6/{Parabólico},9.2/{Hiperbólico}}{
\begin{scope}[xshift=\c cm]
\draw[->,gray] (-1.65,0)--(1.7,0) node[right] {\(u\)};
\draw[->,gray] (0,-1.35)--(0,1.5) node[above] {\(v\)};
\node at (0,2.05) {\titulo};
\end{scope}}
\draw[thick] (0,0) circle[radius=1];
\draw[->,thick] (1,0) arc[start angle=0,end angle=60,radius=1];
\node at (0,-1.85) {\(u^2+v^2=1\)};
\node at (0,-2.35) {\(J^2=-I\)};
\draw[thick] (3.6,-1.2)--(3.6,1.2);
\draw[->,thick] (5.6,-1.2)--(5.6,1.2);
\node at (4.6,-1.85) {\(u^2=1\)};
\node at (4.6,-2.35) {\(J^2=0\)};
\draw[thick,domain=-.9:.9,samples=50] plot ({9.2+cosh(\x)},{sinh(\x)});
\draw[thick,domain=-.9:.9,samples=50] plot ({9.2-cosh(\x)},{sinh(\x)});
\node at (9.2,-1.85) {\(u^2-v^2=1\)};
\node at (9.2,-2.35) {\(J^2=I\)};
\end{tikzpicture}
\caption{Realizaciones cuadráticas de TRIT. Para
\(J_\tau=\left(\begin{smallmatrix}0&-\tau\\1&0\end{smallmatrix}\right)\),
la exponencial conserva \(u^2+\tau v^2\). La degeneración parabólica
conserva \(u\) y transporta \(v\) linealmente. Los tres diagramas representan
órbitas de formas cuadráticas; su tipo geométrico corresponde a la relación
algebraica del generador.}
\label{fig:trit-regimenes}
\end{figure}


\newpage
\subsection{TPK: núcleo topológico de fase}

La dinámica utiliza el selector de fase
\[
 M_{\rm ph}(r)=
 \begin{cases}
 +1,&r\in\{1,4,7\},\\
 -1,&r\in\{2,5,8\},\\
 0,&r\in\{3,6,9\}.
 \end{cases}
\]
Su valor decide qué evaluación opera. El transporte cardinal desplaza el cursor correspondiente; el registro conserva el valor anterior, el nuevo cociente, la orientación y la transición. Por tanto, selección y transporte son funciones distintas: un signo ternario no es un desplazamiento del toro.

Denotaremos por \(\widehat X\) el espacio de estados aritméticos con memoria de trayectoria. Un elemento contiene las posiciones y direcciones de los dos cursores, el modo activo, la firma TRIT, la fase, los residuos y cocientes de ambas evaluaciones, los acarreos, el prefijo construido y los datos de frontera. Cuando se introduce un cilindro de refinamiento, éste forma parte del estado. La notación evita sustituir el objeto por una versión reducida que conserve únicamente el símbolo visible.

\begin{definition}[Actualización del TPK]
Sean \(\mathsf{Sel}_t\) la selección de modo mediante \(M_{\rm ph}\), \(\mathsf{Tra}_t\) el transporte cardinal levantado y \(\mathsf{Mem}_t\) la actualización de los registros. La transición es
\begin{equation}
 \mathcal U_t=\mathsf{Mem}_t\circ\mathsf{Tra}_t\circ\mathsf{Sel}_t.
 \label{eq:actualizacion}
\end{equation}
Cada aplicación conserva los campos que no modifica y actualiza los restantes mediante la división euclídea y la concatenación de la trayectoria.
\end{definition}

La notación \eqref{eq:actualizacion} expresa el orden de composición, no reemplaza las reglas efectivas. La siguiente sección especifica completamente su realización finita de dos cursores. Después, la prolongación no se obtendrá borrando su memoria y repitiendo la primera emisión, sino transportando los prefijos y sus fronteras compatibles. De este modo, APP proporciona las evaluaciones, TRIT distingue sus regímenes y TPK compone su evolución.

% Recuperación expositiva desde los propietarios del núcleo TPK.
% La procedencia y el control de conservación se documentan en technical/TPK_INTEGRACION.md.
% Este archivo amplía la sección TPK; no sustituye extension.tex ni generacion.tex.

\subsubsection{Estado dinámico, observables y dependencia de las operaciones}
\label{tpk-int:estado}

El núcleo topológico de fase organiza la evolución de las dos hojas de APP
con la orientación determinada por TRIT. Su definición comprende el estado
sobre el que actúa, las reglas de transición y las relaciones que hacen
compatibles las prolongaciones. La ecuación \eqref{eq:actualizacion} adquiere
contenido operativo al especificar qué coordenadas recibe y modifica cada
uno de sus factores. La posición pertenece al soporte toroidal; la fase
pertenece al calendario; la orientación determina el transporte; los
cocientes registran la evaluación aritmética anterior a su reducción.
Estas coordenadas se relacionan mediante la transición y conservan funciones
matemáticas diferentes.

La proyección observable de un estado es
\begin{equation}
 q=(x_+,d_+,x_\times,d_\times,\theta)
 \in\mathcal Q_{\rm obs}
 :=(X_9\times\mathcal D)^2\times\Theta_{108},
 \qquad \mathcal D=\{N,E,S,O\},\quad
 \Theta_{108}=\ZZ/108\ZZ.
 \label{tpk-int:qobs}
\end{equation}
Los subíndices distinguen los cursores aditivo y multiplicativo. En las
fibras de estados completos sobre \(q\) se conservan la trayectoria ordenada, las evaluaciones
enteras de ambas hojas, sus residuos y cocientes, y los registros de
orientación y frontera. En una prolongación de precisión se añaden el
prefijo y el cilindro compatible. Por tanto, dos estados con la misma
proyección \(q\) pueden conservar registros distintos.

La dependencia de las operaciones tiene tres niveles. En el nivel local,
el selector activa una hoja, el transporte modifica su posición y la
actualización calcula las nuevas coordenadas aritméticas. En el nivel de
trayectorias, estas transiciones se componen y conservan los registros de
los pasos anteriores. En el nivel de refinamiento, las relaciones de
extensión determinan qué continuación conserva el prefijo, la firma y la
frontera del antecedente. El calendario finito describe la fase de esas
operaciones; la holonomía describe su composición sobre la historia.

\subsubsection{Selección trítica y transporte cardinal levantado}
\label{tpk-int:seleccion}

Sea \(\bar\theta\in\{0,\ldots,107\}\) el representante del calendario.
La fase operativa y el sector dodecafásico son
\begin{equation}
 r(\theta)=1+(\bar\theta\bmod9),\qquad
 \beta(\theta)=\left[\left\lfloor\frac{\bar\theta}{9}\right\rfloor\right]_{12}.
 \label{tpk-int:fase}
\end{equation}
El selector ya definido puede escribirse \(\varepsilon_9=M_{\rm ph}\).
Su vector de valores, en la carta ordenada de \(D_9\), es
\[
 m_{\rm ph}=(1,-1,0,1,-1,0,1,-1,0).
\]
La función \(M_{\rm ph}:D_9\to\TRIT\) y este vector son dos
presentaciones del mismo selector, con sus dominios respectivos.

Definamos los indicadores de activación
\[
 a_+(\theta)=\mathbf1_{\{1,4,7\}}(r(\theta)),\qquad
 a_\times(\theta)=\mathbf1_{\{2,5,8\}}(r(\theta)).
\]
El transporte de los cursores es entonces
\begin{equation}
 x_+'=x_++a_+(\theta)\nu_{d_+},\qquad
 x_\times'=x_\times+a_\times(\theta)\nu_{d_\times}
 \quad\text{en }X_9.
 \label{tpk-int:cursores}
\end{equation}
En las fases \(3,6,9\) permanecen ambas posiciones; el evento neutral
forma parte de la cronología y de su registro. Las direcciones cardinales
determinan los vectores \(\nu_d\), mientras la firma de fase determina
cuál de esos vectores se aplica.

El levantamiento entero conserva el cruce de frontera. Si
\(s:X_9\to D_9^2\) elige los representantes positivos y
\(x'=x+a\nu_d\), con \(a\in\{0,1\}\), se define
\begin{equation}
 b_d(x,a)=\frac{s(x)+a\nu_d-s(x')}{9}\in\ZZ^2.
 \label{tpk-int:borde}
\end{equation}
La divisibilidad por \(9\) se sigue de la igualdad en \(X_9\).
Si \(z\in\ZZ^2\) registra los cruces anteriores, la actualización es
\(z'=z+b_d(x,a)\). De este modo,
\[
 s(x')+9z'=s(x)+9z+a\nu_d.
\]
El dato entero de transporte permanece reconstruible después de reducir
la posición sobre el toro.

\begin{proposition}[Conservación del desplazamiento levantado]
\label{tpk-int:prop-desplazamiento}
Para una trayectoria cardinal de \(n\) pasos, cuyos indicadores de
activación son \(a_t\), se cumple
\[
 s(x_n)+9z_n=s(x_0)+9z_0+
 \sum_{t=0}^{n-1}a_t\nu_{d_t}.
\]
La misma igualdad se conserva al concatenar dos trayectorias compatibles.
\end{proposition}
\begin{proof}
La identidad de un paso es \eqref{tpk-int:borde}. Al sumarla para los
\(n\) pasos, las coordenadas intermedias se cancelan. Dividir el intervalo
de suma en dos tramos produce exactamente la ley de concatenación. En una
trayectoria cerrada, la diferencia \(z_n-z_0\) es el grado de
enrollamiento definido previamente.
\end{proof}

El mismo principio rige las evaluaciones de APP. Si una hoja toma el
valor entero \(N_t\), su registro conserva simultáneamente
\(N_t=R_t+9q_t\). La diferencia entre dos pasos verifica
\[
 N_{t+1}-N_t=(R_{t+1}-R_t)+9(q_{t+1}-q_t).
\]
La memoria de cruce espacial \(z_t\) y el cociente aritmético \(q_t\)
tienen definiciones diferentes; su conservación conjunta mantiene la
procedencia geométrica y aritmética de cada transición.

\subsubsection{Cronología de dos cursores y retornos compuestos}
\label{tpk-int:cronologia}

La involución cardinal \(\jmath\) intercambia \(N\) con \(S\) y
\(E\) con \(O\). Tras aplicar \eqref{tpk-int:cursores}, ambas
direcciones se invierten cuando finaliza un bloque de \(27\) pasos.
Escribiendo
\[
 h(\theta)=\mathbf1_{\{0,27,54,81\}}
             ((\bar\theta+1)\bmod108),
\]
la regla completa sobre la proyección observable es
\begin{equation}
 F_{\rm obs}(q)=
 \bigl(x_+',\jmath^{h(\theta)}d_+,
       x_\times',\jmath^{h(\theta)}d_\times,
       \theta+1\bigr).
 \label{tpk-int:fobs}
\end{equation}
La lectura que forma la emisión utiliza las posiciones anteriores al
transporte, según la recurrencia explícita de la
sección~\ref{sec:extension}. La actualización observable y la emisión
se componen, por tanto, con un orden fijado.

\begin{proposition}[Reversibilidad y periodo del calendario observable]
\label{tpk-int:periodo}
La aplicación \(F_{\rm obs}\) es una permutación de
\(\mathcal Q_{\rm obs}\) y satisface \(F_{\rm obs}^{108}=I\).
En un bloque de \(54\) pasos se restauran posiciones y direcciones;
la coordenada de \(\Theta_{108}\) avanza \(54\).
\end{proposition}
\begin{proof}
Dado el estado final, se recupera primero la fase anterior restando uno
en \(\Theta_{108}\). Esa fase determina si se invirtieron las
direcciones y qué cursor se desplazó. Aplicar la involución correspondiente
y restar el vector cardinal recupera el estado inicial de manera única.

Cada intervalo de \(27\) pasos contiene nueve activaciones de cada
cursor. El bloque siguiente conserva las activaciones e invierte las
direcciones, de modo que los desplazamientos enteros se cancelan en
\(54\) pasos. La regla tiene periodo \(54\) en sus incrementos de
posición y orientación; la misma cancelación vale con cualquier origen
del calendario. Tras \(108\) pasos se restaura también \(\theta\).
\end{proof}

Se reserva
\[
 \mathcal K_{27}=F_{\rm obs}^{27}:
 \mathcal Q_{\rm obs}\longrightarrow\mathcal Q_{\rm obs}
\]
para el iterado de \(27\) transiciones. Su cuarto iterado es la
identidad sobre el estado observable definido en \eqref{tpk-int:qobs}.
La proyección de la historia al calendario permite efectuar esta
comprobación finita. El registro de la trayectoria, sus evaluaciones y
los nuevos prefijos se conservan en el estado dinámico sobre el que
actúa \(\mathcal U_t\); la igualdad observable no los reinicializa.

\subsubsection{Transporte afín de la memoria y composición de estados}
\label{tpk-int:memoria}

La proyección observable admite fibras de datos aritméticos y geométricos.
Sobre una configuración \(q\), escribimos un estado en la forma
\[
 x=(q,o,h,c,\sigma,C,b,\lambda).
\]
Aquí \(o\) conserva la orientación y la hoja; \(h\), los residuos
y cocientes; \(c\), la memoria aditiva de transporte; \(\sigma\),
la firma de frontera; \(C\), el cilindro de precisión; \(b\), su
etiqueta de frontera; y \(\lambda\), la trayectoria registrada. La
firma está determinada por una aplicación
\[
 \operatorname{Sig}_q:
 \mathsf H_q\times\mathsf C_q\times\mathsf{Hist}_q
 \longrightarrow\mathsf S_q,
 \qquad \sigma=\operatorname{Sig}_q(h,c,\lambda).
\]
Los símbolos \(\mathsf H_q,\mathsf C_q,\mathsf{Hist}_q\) y
\(\mathsf S_q\) denotan, respectivamente, los conjuntos de datos
residuo--cociente, el grupo abeliano de memoria de transporte, las
trayectorias registradas y las firmas de frontera.
Las realizaciones de esta firma se especifican con sus componentes en
las construcciones regional y terminal. Esta dependencia impide tratar
la firma como un parámetro libre una vez fijados los datos que la producen.

La carta elemental de residuo y cociente es explícita. Para
\(n\in\ZZ\), sean
\[
 r=1+((n-1)\bmod9),\qquad u=\frac{n-r}{9}.
\]
Entonces \(n=r+9u\), con \(r\in D_9\) y \(u\in\ZZ\).
La unicidad se sigue de que dos representantes de \(D_9\) congruentes
módulo \(9\) son iguales. Esta carta proporciona la reconstrucción
entera local. Las torres de precisión incorporan, además, las aplicaciones
de truncamiento y los cilindros desarrollados en
\ref{sec:extension} y \ref{sec:generacion}.

Sea \(\gamma:q\to q'\) una trayectoria observable. El transporte
de las fibras se expresa mediante aplicaciones
\[
 \mathsf H(\gamma):\mathsf H_q\to\mathsf H_{q'},\qquad
 \mathsf C(\gamma):\mathsf C_q\to\mathsf C_{q'},\qquad
 \mathsf{Hist}(\gamma):\mathsf{Hist}_q\to\mathsf{Hist}_{q'}.
\]
Las aplicaciones respetan identidades y composición, y
\(\mathsf C(\gamma)\) es un homomorfismo de grupos abelianos.
El incremento \(k_{\rm tr}(\gamma)\in\mathsf C_{q'}\) cumple
\begin{equation}
 \begin{split}
 k_{\rm tr}(\operatorname{id}_q)&=0,\\
 k_{\rm tr}(\gamma_2\circ\gamma_1)
 &=k_{\rm tr}(\gamma_2)+
   \mathsf C(\gamma_2)k_{\rm tr}(\gamma_1).
 \end{split}
 \label{tpk-int:cociclo}
\end{equation}
El símbolo \(k_{\rm tr}\) designa aquí el cociclo de transporte;
la coordenada racional \(\kappa\) del registro \(K\) conserva
su notación en el capítulo de la clausura dodecafásica.

Las coordenadas transportables se actualizan por
\begin{equation}
 \begin{aligned}
 h'&=\mathsf H(\gamma)h,&
 c'&=\mathsf C(\gamma)c+k_{\rm tr}(\gamma),\\
 \lambda'&=\gamma\circ\lambda,&
 \sigma'&=\operatorname{Sig}_{q'}(h',c',\lambda').
 \end{aligned}
 \label{tpk-int:accion-memoria}
\end{equation}
La segunda ecuación es una acción afín: el término lineal transporta
la memoria anterior y el cociclo añade el incremento producido por la
trayectoria. En la realización de cruce espacial de cada cursor por separado,
\(\mathsf C(\gamma)=I\) y \(k_{\rm tr}(\gamma)\) es la suma
de los vectores de \eqref{tpk-int:borde}. Esta instancia concreta
vincula la ley abstracta con las operaciones cardinales ya definidas.

\begin{proposition}[Compatibilidad de la actualización con la concatenación]
\label{tpk-int:composicion-memoria}
La actualización \eqref{tpk-int:accion-memoria} por
\(\gamma_1\), seguida de la actualización por \(\gamma_2\),
coincide con la actualización por \(\gamma_2\circ\gamma_1\).
\end{proposition}
\begin{proof}
Para la coordenada afín, dos actualizaciones dan
\[
 c''=\mathsf C(\gamma_2)\mathsf C(\gamma_1)c
       +\mathsf C(\gamma_2)k_{\rm tr}(\gamma_1)
       +k_{\rm tr}(\gamma_2).
\]
La functorialidad de \(\mathsf C\) y
\eqref{tpk-int:cociclo} convierten esta expresión en
\(\mathsf C(\gamma_2\circ\gamma_1)c+
k_{\rm tr}(\gamma_2\circ\gamma_1)\).
La afirmación para \(h\) se sigue de la composición de
\(\mathsf H\), y para \(\lambda\), de la asociatividad de las
trayectorias. La firma final coincide porque se evalúa sobre las mismas
tres coordenadas reconstruidas.
\end{proof}

El cilindro y su frontera completan esta ley mediante las relaciones
\(\operatorname{Ext}_\gamma\), definidas en
\ref{tpk-int:bidireccionalidad}. Sus elementos determinan cuáles de
las actualizaciones anteriores conservan una prolongación admisible.
Los estados, junto con las ternas \((\gamma,x,x')\) admitidas por
esas relaciones, forman una categoría sobre las trayectorias observables:
\[
 (\gamma_2,x',x'')\circ(\gamma_1,x,x')
     =(\gamma_2\circ\gamma_1,x,x'').
\]
La identidad de \(x\) es \((\operatorname{id}_q,x,x)\).
La composición conserva a la vez el transporte aritmético y la
prolongabilidad de frontera. Esta es la estructura en la que se efectúan
los retornos con memoria y los refinamientos sucesivos.

\subsubsection{Registro dodecafásico, fase y memoria acumulada}
\label{tpk-int:dodecafase}

Las \(108\) transiciones se distribuyen en las ventanas ordenadas
\[
 E_m=\{9m-8,\ldots,9m\},\qquad 1\leq m\leq12.
\]
El registro conserva el origen de fase, la orientación y los datos
producidos durante cada ventana. Su aplicación es
\[
 \mathcal R_{12}(x)=
 \bigl((E_m,\tau_m,\sigma_m,\kappa_m,\Lambda_m)\bigr)_{m=1}^{12},
\]
donde \(\tau_m\) es la subruta de la ventana, \(\sigma_m\) su
orientación, \(\kappa_m\) el incremento entero de frontera y
\(\Lambda_m\) el registro local de incidencias. Se adopta la carta
que se desarrollará en \ref{subsec:k-registro-integral}; sus índices
de ventana corresponden a \(m=\beta+1\).
El dominio comprende las historias compatibles
del TPK y el codominio, sus registros temporales orientados. La
composición de segmentos y sus incrementos se rige por
\eqref{tpk-int:accion-memoria}. El grupo \(C_{12}\) describe, en
cambio, la traslación del índice de ventana; \(\Theta_{108}\) conserva
el calendario de pasos. Esta distinción permite precisar qué se transforma
y qué se observa al completar un ciclo.

En las coordenadas \(\theta=9b+r\), con \(0\leq r<9\), el avance
de \(9\) pasos conserva \(r\) y aumenta \(b\) en una unidad.
La adición de fases introduce el cociente
\[
 c_0(r,r')=\left\lfloor\frac{r+r'}9\right\rfloor,
\]
cuya identidad de cociclo asegura la asociatividad de la composición.
En el calendario finito se conserva \([b]_{12}\); en un registro de
vueltas se conserva \(b\in\ZZ\), o \(b\in\mathbb N_0\) para
una historia de origen fijado. La sección~\ref{sec:extension} demuestra
la sucesión exacta no escindida \eqref{eq:extension108}. Su función es
describir la relación entre estas coordenadas, antes de emplearlas en la
construcción del estado firmado.

El registro firmado recibe los datos de fase y memoria mediante la
composición que se desarrolla en la sección~\ref{sec:k-registro}:
\begin{equation}
 R_{\rm sgn}=\Pi_H\circ\mathcal E_{108}^{90,120}
                         \circ\mathcal R_{12}.
 \label{tpk-int:registro-causal}
\end{equation}
Su dominio es el estado terminal con memoria; su salida pertenece a
\(\ZZ^{12}\). La transformación reversible entre ese estado firmado
y \(K\), así como la lectura racional de \(K\), se demuestran en
el capítulo correspondiente. Las operaciones de registro anteceden a
esas transformaciones. La reconstrucción de un registro desde sus
diferencias cíclicas verifica su recuperabilidad; la procedencia dinámica
debe conservarse mediante los eventos que lo produjeron.

Para la traza codificada \((d_t)\) conservada por el registro, el
lector de eventos cuenta los códigos \(\{2,4,6,8\}\) de cada
ventana con la orientación
\((+,+,+,-,-,-,+,+,+,-,-,-)\). Estos códigos y los símbolos
cardinales de \(\mathcal D\) son coordenadas distintas del
registro. El capítulo~\ref{sec:k-registro} desarrolla este lector,
la descomposición integral \(1+5+6\), las congruencias que
caracterizan su imagen y su inversa exacta. La prueba de inversión
establece qué datos conserva esa carta; la definición de
\(\mathcal R_{12}\) especifica la historia orientada sobre la que
actúa.

\subsubsection{Incidencia de fases y construcción interna de los canales}
\label{tpk-int:incidencia}

El selector determina tres subconjuntos de \(D_9\):
\[
 H_+=\{1,4,7\},\qquad H_-=\{2,5,8\},\qquad H_0=\{3,6,9\}.
\]
Sea \(H_{\rm act}=H_+\sqcup H_-\) el conjunto de fases activas y
sea \(O_{\rm ph}=D_9\setminus\{9\}\) la carta anterior al retorno
marcado. Denotemos por \(P_2(H_{\rm act})\) el conjunto de pares
no ordenados de fases activas. La incidencia interna está dada por
\begin{equation}
 \mathcal F^{\rm ph}_6=H_{\rm act}\times P_2(H_{\rm act}),\qquad
 \mathcal F^{\rm ph}_8=O_{\rm ph}\times P_2(H_{\rm act}),\qquad
 \Theta_{\rm ph}=D_9\times P_2(H_{\rm act}).
 \label{tpk-int:fases90120}
\end{equation}
Los índices describen la cantidad de posiciones de la primera componente.
La construcción utiliza las fases y el origen del calendario ya fijados;
las realizaciones excepcionales posteriores transportarán estas
incidencias a sus propios dominios.

\begin{proposition}[Cardinales e incidencia orientada de fase]
\label{tpk-int:cardinales}
Los conjuntos de \eqref{tpk-int:fases90120} tienen cardinales
\(90,120,135\), respectivamente. Si
\[
 \partial_{\rm ph}=\Theta_{\rm ph}\times\{-1,+1\},\qquad
 E_{\rm ph}=\partial_{\rm ph}\sqcup\{\infty\},\qquad
 L_{\rm ph}=H_0\times P_2(H_{\rm act}),
\]
entonces
\[
 |\partial_{\rm ph}|=270,\quad |E_{\rm ph}|=271,\quad
 |L_{\rm ph}|=45,\quad
 |\partial_{\rm ph}\sqcup L_{\rm ph}|=315.
\]
Trasladar simultáneamente el origen de fase y el retorno marcado conserva
todos estos cardinales y sus inclusiones.
\end{proposition}
\begin{proof}
Hay seis fases activas, por lo que
\(|P_2(H_{\rm act})|=\binom62=15\). Los productos por
\(6,8,9\) dan \(90,120,135\). La orientación duplica \(135\);
la marca de retorno añade un elemento; las tres fases neutrales dan
\(3\cdot15=45\). Una traslación del calendario transporta cada
factor por una biyección y lleva el punto marcado a su imagen. Los
productos y las uniones disjuntas conservan así la incidencia.
\end{proof}

El cardinal \(271\) tiene aquí una realización por incidencia de
fase. La completación cúbica de la sección~\ref{sec:extension} tiene
su propia construcción de la corona \(1000-729\). La correspondencia
entre ambas realizaciones debe conservar los mapas y las marcas que las
identifican, además del cardinal. Análogamente, el selector
\(\varepsilon_9\), la incidencia \((90,120)\) y el extractor
\(\mathcal E_{108}^{90,120}\) conservan los dominios y operaciones
que se acaban de distinguir.

\subsubsection{Inversión de trayectorias y prolongaciones conjugadas}
\label{tpk-int:bidireccionalidad}

La bidireccionalidad comienza en una operación sobre trayectorias, antes
de asignarles valores escalares. Para
\(\gamma=(x_0;d_1,\ldots,d_n)\), con extremo \(x_n\), definimos
\[
 \operatorname{rev}\gamma
   =(x_n;\jmath d_n,\ldots,\jmath d_1).
\]
El origen de la trayectoria inversa es el extremo de la original.
La composición satisface
\[
 \operatorname{rev}^2\gamma=\gamma,\qquad
 \operatorname{rev}(\gamma_2\circ\gamma_1)
   =\operatorname{rev}\gamma_1\circ\operatorname{rev}\gamma_2,
\]
y el desplazamiento entero cambia de signo. Estas identidades se obtienen
invirtiendo el orden de los pasos y sustituyendo cada vector cardinal por
su opuesto. En un lazo se sigue
\(\operatorname{wind}(\operatorname{rev}\gamma)
=-\operatorname{wind}(\gamma)\).

Sobre los registros, la inversión requiere transportar asimismo el orden
de sus componentes y la orientación de sus eventos. La inversión de
ruta y una conjugación de firma son operaciones distinguibles: su
composición se especifica cuando se aplica a un registro concreto. En
particular, la involución regional que intercambia las componentes de
clausura y propagación conserva el eje de autoescala; su fórmula se
presenta en la sección~\ref{mono-ampl:regiones}. Esta operación regional
actúa sobre otro dominio que la inversión de una trayectoria cardinal.

Las extensiones sobre una ruta \(r:q\to q'\) se expresan mediante
relaciones entre fibras de estados compatibles. Si \(\mathcal F_q\)
denota la fibra de estados completos sobre \(q\), se exige
\[
 \operatorname{Ext}_r\subseteq\mathcal F_q\times\mathcal F_{q'},
 \qquad \Delta_{\mathcal F_q}\subseteq
        \operatorname{Ext}_{\operatorname{id}_q},
\]
\[
 \operatorname{Ext}_{r_2}\circ\operatorname{Ext}_{r_1}
       \subseteq\operatorname{Ext}_{r_2\circ r_1}.
\]
Su composición conserva las evaluaciones de ambas hojas, la orientación,
los cocientes, el prefijo y la frontera. La inversión de la ruta no
implica por sí sola que toda relación de extensión sea una biyección.
La prolongación conjugada se define en el dominio donde el transporte de
esos datos satisface las compatibilidades correspondientes.

\subsubsection{Holonomía nonádica y continuidad del refinamiento}
\label{tpk-int:holonomia}

Una vuelta de fase compone las transiciones sobre el estado con memoria:
\begin{equation}
 \Gamma_{9,k}=\mathcal U_{k+8}\circ\cdots\circ\mathcal U_k.
 \label{tpk-int:gamma}
\end{equation}
Sobre una historia individualizada, esta composición funcional realiza
la relación de holonomía; en el dominio completo se conserva
\(\Gamma_{9,k}\subseteq\widehat X_k\times\widehat X_{k+9}\),
con las prolongaciones admisibles desarrolladas más adelante.
La base cíclica tiene nueve fases, mientras
la fibra conserva la trayectoria, los cocientes y la frontera. La
monodromía es la acción de una vuelta de esa base sobre las fibras; la
holonomía \(\Gamma_{9,k}\) expresa su realización por las transiciones
efectivas del TPK.

Si \(g\) observa la fase y \(M\) el número de vueltas conservadas,
\[
 g(\Gamma_{9,k}x)=g(x),\qquad M(\Gamma_{9,k}x)=M(x)+1.
\]
La primera identidad expresa clausura de fase y la segunda identifica
la nueva profundidad temporal. La prolongación coinductiva conserva
además sus cilindros y prefijos. Las restricciones entre profundidades
componen un sistema inverso; una sección compatible reúne sus
antecedentes sin reemplazarlos por la última observación. El desarrollo
de esta aritmética de historias y de su unidad ocupa
\ref{hist:seccion}--\ref{hist:doble-varianza}; la construcción de las
regiones y sus lectores se realiza en la sección~\ref{sec:generacion}.

La transferencia de fase sobre emisiones ya construidas tiene otro
dominio. Sean \(\mathcal A_+\) y \(\mathcal A_\times\) los
alfabetos de firmas aditivas y multiplicativas, de cardinales \(18\)
y \(26\). Su producto identifica el catálogo \(\mathcal U_6\).
Para \(f:\mathcal A_+\times\mathcal A_\times\to\RR\),
las medias de hoja son
\[
 (E_+f)(s,e)=\frac1{18}\sum_{s'\in\mathcal A_+}f(s',e),\qquad
 (E_\times f)(s,e)=\frac1{26}\sum_{e'\in\mathcal A_\times}f(s,e').
\]
La construcción de estas emisiones y sus cardinales se efectúa mediante
la recurrencia de la sección~\ref{sec:extension}. Una vez obtenidas,
el operador de transferencia se define sobre
\(\mathcal U_6\times C_9\) por
\[
 P_+=E_+,\quad P_-=E_\times,\quad P_0=I,\qquad
 \mathcal K_{\rm ph}((u,r),(v,r+1))=P_{M_{\rm ph}(r)}(u,v),
\]
con peso cero para las restantes transiciones de fase. El índice
\(r\) utiliza el representante positivo de la clase cíclica.
El orden de fases \((+1,-1,0)^3\) produce entonces la renovación
\[
 \mathcal K_{\rm ph}^{9}\big|_r=E_+E_\times.
\]
En efecto, ambas medias son idempotentes, conmutan y promedian factores
diferentes; su composición asigna peso \(1/(18\cdot26)=1/468\)
a cada emisión. Esta igualdad caracteriza un observable de transferencia
posterior. El transporte de las historias sigue dado por
\eqref{tpk-int:gamma}, con su memoria y su frontera.

La jerarquía resultante establece la función de los capítulos siguientes.
La emisión finita construye el dominio y las regiones; las relaciones de
extensión conservan sus prefijos; la conexión nonádica los prolonga; el
registro dodecafásico determina los datos que intervienen en \(K\) y
en la clausura de \(\alpha\). La realización numérica y la realización
incidencial reciben esa estructura anterior. Este orden permite exponer
cada consecuencia en su capítulo conservando, desde el núcleo formal,
los objetos que la hacen posible.


\subsection{Memoria acumulada y resolventes del transporte}
\label{subsec:memoria-resolvente}

El retorno del calendario no elimina los incrementos producidos durante
el recorrido. Para expresar esta diferencia operatoriamente, se construye
primero un registro acumulativo de eventos de la trayectoria y después su
serie generadora. La reducción de variables de memoria se realizará sobre
ese transporte, conservando también la fuente y el lector de salida.

\subsubsection{Eventos orientados y actualización del registro}

Sea \((d_t)\) la secuencia de códigos enteros de dirección conservada
por una trayectoria del TPK. El código \(d_t\) y la dirección cardinal
del cursor son coordenadas distintas; no se supone una identificación
entre sus alfabetos. Numeramos desde cero las ventanas
\[
 I_m=\{9m+1,\ldots,9m+9\},\qquad m\geq0.
\]
Durante el primer ciclo, \(I_m\) es la ventana \(E_{m+1}\) del registro
dodecafásico. Su orientación es
\(\sigma=(1,1,1,-1,-1,-1,1,1,1,-1,-1,-1)\).
Para una historia prolongada, \(\sigma_m\in\{-1,1\}\) es el signo
conservado en la ventana correspondiente. El evento firmado es
\begin{equation}
 a_m=\sigma_m\sum_{t\in I_m}
       \mathbf1_{\{2,4,6,8\}}(d_t),\qquad |a_m|\leq9.
 \label{mem:evento}
\end{equation}
Cada sumando se incorpora cuando ocurre su transición. Esta definición
cuenta los eventos de la traza; no escoge direcciones mediante un valor
que deba obtenerse después.

Sean \(\mathbf e_0,\ldots,\mathbf e_{11}\) la base de \(\ZZ^{12}\)
y \(S\mathbf e_j=\mathbf e_{j+1\bmod12}\). Escribimos
\(R=S^{-1}\). El registro acumulado en coordenadas fijas satisface
\[
z_{m+1}=z_m+a_m\mathbf e_{m\bmod12}.
\]
Un prefijo sin historia acumulada empieza en \(z_0=0\); en otro
origen, \(z_0\) conserva el registro previo.
El marco que acompaña al calendario es \(y_m=S^{-m}z_m\).

\begin{proposition}[Actualización y retorno con memoria]
\label{mem:prop-actualizacion}
El registro en el marco móvil cumple
\begin{equation}
 y_{m+1}=Ry_m+\eta_m,\qquad
 \eta_m=a_mR\mathbf e_0.
 \label{mem:actualizacion}
\end{equation}
Para todo \(n\geq1\),
\begin{equation}
 y_n=R^ny_0+\sum_{j=0}^{n-1}R^{n-1-j}\eta_j.
 \label{mem:iteracion}
\end{equation}
En particular,
\begin{equation}
 y_{12}=y_0+\sum_{j=0}^{11}a_j\mathbf e_j.
 \label{mem:retorno12}
\end{equation}
\end{proposition}
\begin{proof}
Multiplicar la actualización de \(z_m\) por \(S^{-(m+1)}\) da
\(y_{m+1}=S^{-1}y_m+a_mS^{-1}\mathbf e_0\), porque
\(S^{-m}\mathbf e_{m\bmod12}=\mathbf e_0\).
La fórmula \eqref{mem:iteracion} se obtiene por inducción.
Para \(n=12\), \(R^{12}=I\) y
\(R^{11-j}\eta_j=a_jR^{12-j}\mathbf e_0=a_j\mathbf e_j\),
lo que prueba \eqref{mem:retorno12}.
\end{proof}

Así, doce ventanas restauran el marco y conservan el vector de eventos.
El recuento no permite recuperar por sí solo el orden de todos los pasos
dentro de cada ventana; ese orden permanece en la trayectoria registrada.

Para precisar el balance, definimos
\[
 D_3=I-S^3,\quad D_4=I-S^4,\quad
 B=D_3^{\mathsf T}D_3+D_4^{\mathsf T}D_4,
\]
\[
 \mathcal M(y)=\tfrac12y^{\mathsf T}By,
 \qquad q(y)=\sum_{j=0}^{11}y_j.
\]
La primera es una forma cuadrática del registro; la segunda registra
la suma de sus coordenadas. No se identifican aquí con energía ni carga físicas.

\begin{proposition}[Balance producido por el evento]
\label{mem:prop-balance}
La actualización \eqref{mem:actualizacion} satisface
\begin{equation}
 \begin{aligned}
 \mathcal M(y_{m+1})-\mathcal M(y_m)
   &=a_m(By_m)_0+2a_m^2,\\
 q(y_{m+1})-q(y_m)&=a_m.
 \end{aligned}
 \label{mem:balance}
\end{equation}
\end{proposition}
\begin{proof}
Al desarrollar los productos,
\(B=4I-S^3-S^{-3}-S^4-S^{-4}\); por tanto
\(R^{\mathsf T}BR=B\) y \(B_{00}=4\).
Como \(y_{m+1}=R(y_m+a_m\mathbf e_0)\), la diferencia cuadrática
es \(a_m\mathbf e_0^{\mathsf T}By_m+\tfrac12a_m^2B_{00}\).
La suma de coordenadas es invariante bajo la permutación \(R\),
lo que da la segunda igualdad.
\end{proof}

\Needspace{20\baselineskip}
\subsubsection{Serie generadora y control uniforme de la memoria}

La serie conserva todos los incrementos, no únicamente su suma al final
de un ciclo. Con una variable formal \(u\), sean
\[
 Y(u)=\sum_{m\geq0}y_mu^m,\qquad
 H_\eta(u)=\sum_{m\geq0}\eta_mu^m.
\]

\begin{proposition}[Resolvente y cota de cola]
\label{mem:prop-serie}
Se cumple la identidad formal
\begin{equation}
 Y(u)=(I-uR)^{-1}\bigl(y_0+uH_\eta(u)\bigr).
 \label{mem:serie}
\end{equation}
Ambas series convergen para \(|u|<1\). Si \(\ell:\RR^{12}\to\RR\)
es lineal, \(L=\|\ell\|_{1\to\RR}\), \(M_0=\|y_0\|_1\) y
\(0<r<1\), entonces, para todo entero \(N\geq0\) y \(|u|\leq r\),
\begin{equation}
 \left|\sum_{m>N}\ell(y_m)u^m\right|
 \leq Lr^{N+1}\left(
 \frac{M_0+9(N+1)}{1-r}+\frac{9r}{(1-r)^2}\right).
 \label{mem:cota}
\end{equation}
La serie de derivadas converge uniformemente en cada disco
\(|u|\leq r<1\).
\end{proposition}
\begin{proof}
Multiplicar \eqref{mem:actualizacion} por \(u^{m+1}\) y sumar da
\(Y-y_0=uRY+uH_\eta\). El término constante de \(I-uR\) es
invertible, de modo que su inversa formal es \(\sum_{k\geq0}u^kR^k\).
Además, \(R\) preserva la norma \(\ell^1\) y
\(\|\eta_m\|_1\leq9\), de donde
\(\|y_m\|_1\leq M_0+9m\).
Para la cola se suman las mayorantes
\(L(M_0+9m)r^m\), usando
\[
 \sum_{m>N}r^m=\frac{r^{N+1}}{1-r},\qquad
 \sum_{m>N}mr^m=r^{N+1}
 \left(\frac{N+1}{1-r}+\frac r{(1-r)^2}\right).
\]
Por último, \(\sum_{m\geq1}Lm(M_0+9m)r^{m-1}<\infty\)
mayora la serie derivada y prueba su convergencia uniforme.
\end{proof}

La variable \(u\) cuenta ventanas de nueve transiciones. Su identificación
con la variable de otro lector requiere conservar ese reloj y el mapa
que relaciona ambas lecturas.

\subsubsection{Eliminación de memoria: operador, fuente y lector}

En un dominio donde la actualización está fijada como \(U:X\to X\),
el espacio vectorial libre \(V=\QQ^{(X)}\) sobre los estados admite el operador
\(T\delta_x=\delta_{U(x)}\). El reloj se incluye en \(x\) si la
transición depende de él. Esta extensión lineal conserva estados
distintos aunque compartan una fase observable.

Consideremos una descomposición \(V=V_0\oplus V_1\), donde \(V_1\)
es el sector auxiliar que se desea eliminar, y escribamos
\[
 T=\begin{pmatrix}A&B_1\\C&D\end{pmatrix},\qquad
 v=\binom{v_0}{v_1},\qquad \ell=(\ell_0,\ell_1).
\]
Aquí \(A:V_0\to V_0\), \(B_1:V_1\to V_0\),
\(C:V_0\to V_1\) y \(D:V_1\to V_1\).
Las identidades siguientes son formales; también valen donde los
resolventes convergen.

\begin{proposition}[Reducción de Schur con fuente y lector]
\label{mem:prop-schur}
Sean
\[
 D_u=(I-uD)^{-1},\qquad
 \mathscr F(u)=I-uA-u^2B_1D_uC.
\]
La componente retenida de \((I-uT)^{-1}v\) es
\begin{equation}
 w_0=\mathscr F(u)^{-1}\bigl(v_0+uB_1D_uv_1\bigr),
 \qquad w_1=D_u(v_1+uCw_0).
 \label{mem:schur-fuente}
\end{equation}
Su lectura completa satisface
\begin{equation}
 \begin{split}
 \ell(I-uT)^{-1}v
 ={}&(\ell_0+u\ell_1D_uC)
       \mathscr F(u)^{-1}(v_0+uB_1D_uv_1)\\
    &+\ell_1D_uv_1.
 \end{split}
 \label{mem:schur-lector}
\end{equation}
\end{proposition}
\begin{proof}
Las dos ecuaciones de \((I-uT)w=v\) son
\(w_0=v_0+uAw_0+uB_1w_1\) y
\(w_1=v_1+uCw_0+uDw_1\).
Resolver la segunda y sustituirla en la primera da
\eqref{mem:schur-fuente}; aplicar \(\ell_0\) y \(\ell_1\)
a ambas componentes da \eqref{mem:schur-lector}.
Todas las inversas formales existen porque sus términos constantes
son identidades.
\end{proof}

El término \(u^2B_1D_uC=\sum_{k\geq0}u^{k+2}B_1D^kC\)
describe una salida hacia la memoria, \(k\) pasos dentro de ella y
un regreso. En cambio, \(u\ell_1D_uC\) lee la memoria antes de ese
regreso. Por ello, reducir sólo el operador no conserva necesariamente
la observación: también deben transformarse fuente y lector.

La diferencia es visible en el ciclo ya construido. Si \(P\) retiene
\(\mathbf e_0\), entonces \(PRP=0\), pero
\begin{equation}
 \left\langle\mathbf e_0,(I-uR)^{-1}\mathbf e_0\right\rangle
 =\frac1{1-u^{12}},\qquad
 \left\langle\mathbf e_{11},(I-uR)^{-1}\mathbf e_0\right\rangle
 =\frac u{1-u^{12}}.
 \label{mem:ejemplo-ciclo}
\end{equation}
En efecto, \(R^n\mathbf e_0=\mathbf e_{-n\bmod12}\): las
dos series recogen, respectivamente, \(n\equiv0\) y
\(n\equiv1\pmod{12}\). El resolvente de la compresión es \(I\),
mientras la compresión del resolvente conserva los retornos omitidos.
Este ejemplo corresponde al registro de doce ventanas, no a la totalidad
del estado del TPK.

\subsubsection{Composición de segmentos y coherencia de la reducción}

Sean tres transportes afines composables
\(F_i(y)=R_i y+b_i\), con dominios y codominios compatibles.
La memoria producida por el recorrido completo es
\begin{equation}
 F_3F_2F_1(y)
 =R_3R_2R_1y+b_3+R_3b_2+R_3R_2b_1.
 \label{mem:cociclo-tres}
\end{equation}
La prueba es la sustitución sucesiva de las tres fórmulas afines.
Agrupar primero los dos primeros segmentos o los dos últimos produce
la misma expresión; los factores transportan cada incremento desde
el punto donde se originó. En \eqref{mem:actualizacion},
\(R_i=R\) y \(b_i\) es el evento de su ventana.

La eliminación de varias capas conserva una coherencia semejante.
Para hacerla explícita, descomponemos ahora
\(V=V_0\oplus V_1\oplus V_2\) y escribimos
\[
 T=\begin{pmatrix}
 A&B_1&B_2\\C_1&D_{11}&D_{12}\\C_2&D_{21}&D_{22}
 \end{pmatrix},\qquad Q_2=(I-uD_{22})^{-1}.
\]
Al eliminar la tercera componente quedan los cuatro bloques
\begin{equation}
 \begin{aligned}
 A'&=A+uB_2Q_2C_2,&
 B_1'&=B_1+uB_2Q_2D_{21},\\
 C_1'&=C_1+uD_{12}Q_2C_2,&
 D_{11}'&=D_{11}+uD_{12}Q_2D_{21}.
 \end{aligned}
 \label{mem:schur-intermedio}
\end{equation}

\begin{proposition}[Transitividad de la eliminación]
Si \(D_*=(D_{ij})_{i,j=1}^2\), el complemento obtenido en dos etapas
coincide con el obtenido al eliminar juntas ambas memorias:
\begin{equation}
 \begin{split}
 I-uA'-u^2B_1'(I-uD_{11}')^{-1}C_1'
 ={}&I-uA\\
 &-u^2(B_1\ B_2)(I-uD_*)^{-1}\binom{C_1}{C_2}.
 \end{split}
 \label{mem:schur-transitivo}
\end{equation}
\end{proposition}
\begin{proof}
Resolver la tercera ecuación de \((I-uT)w=v\) y sustituirla en las
dos primeras produce \eqref{mem:schur-intermedio}. Eliminar después la
segunda componente da el miembro izquierdo de
\eqref{mem:schur-transitivo}. Resolver a la vez las dos ecuaciones de
memoria da el derecho, por la proposición~\ref{mem:prop-schur}.
La solución formal es única porque \(I-uT\) tiene término constante
\(I\). En ambos procedimientos se transforman también la fuente y
el lector según \eqref{mem:schur-fuente}--\eqref{mem:schur-lector}.
\end{proof}

El balance \eqref{mem:balance} describe los incrementos del registro.
Las identidades anteriores permiten transportarlos al reutilizar la memoria
en la clausura dodecafásica; no determinan por sí solas coeficientes
analíticos adicionales.


\section{Desarrollo del núcleo formal: condiciones iniciales, geometría y refinamiento}
\label{sec:extension}

El soporte de ochenta y una posiciones no coincide con el dominio de condiciones iniciales de la dinámica. Una condición inicial debe indicar además una orientación; y las hojas aditiva y multiplicativa disponen de cursores independientes. Esta distinción permite reconstruir exhaustivamente la emisión sin seleccionar de antemano una región numérica.

\subsection{Los estados iniciales y la recurrencia de emisión}

Para la enumeración utilizamos índices \(I_9=\{0,\ldots,8\}\), relacionados con las etiquetas positivas por \(i\mapsto i+1\). Sea
\[
 \mathcal S=I_9^2\times\{N,E,S,O\},\qquad
 \Omega_{\rm seed}=\mathcal S_+\times\mathcal S_\times.
\]
Cada condición inicial \(\omega=(s_+,s_\times)\) determina dos posiciones y dos direcciones. El dominio comprende todas las parejas ordenadas, por lo que
\begin{equation}
 |\mathcal S|=324,\qquad |\Omega_{\rm seed}|=324^2=104\,976.
 \label{eq:semillas}
\end{equation}
El censo de firmas del emisor residual, sección~\ref{app:firmas-emision}, enumera estas condiciones mediante identificadores y presenta sus fibras; aquí se define la ley que las transforma. Las filas de una tabla de emisiones son fibras de este dominio de condiciones iniciales.

La realización finita utiliza seis ventanas de nueve transiciones. En cada instante \(t\), la firma de fase es \(M_{\rm ph}(1+((t-1)\bmod9))\). Si vale \(+1\), se lee el residuo positivo \(\Sigma(i+1,j+1)=\rho_9(i+j+2)\) del primer cursor y después se desplaza ese cursor; si vale \(-1\), se lee el residuo positivo \(\Pi(i+1,j+1)=\rho_9((i+1)(j+1))\) del segundo cursor y después se desplaza el segundo. El valor cero incrementa el contador neutral. Al finalizar las transiciones veintisiete y cincuenta y cuatro, se invierten las orientaciones cardinales. Las evaluaciones enteras y sus cocientes se conservan en el registro, pero el emisor utiliza las lecturas residuales aquí especificadas.

En el sector de unidades, el lector multiplicativo utiliza el logaritmo discreto de base dos en \((\ZZ/9\ZZ)^\times\):
\[
 \ell_9(1)=0,\quad\ell_9(2)=1,\quad\ell_9(4)=2,\quad
 \ell_9(8)=3,\quad\ell_9(7)=4,\quad\ell_9(5)=5.
\]
Para la emisión finita se extiende con valor cero al sector no unitario. Esta extensión es un observable y no una identificación de los elementos \(3,6,9\) con la unidad multiplicativa. La hoja y su pertenencia al sector radical permanecen en el registro de trayectoria.

Sean \(S_m\) la suma de los tres residuos positivos aditivos \(\Sigma\) de la ventana \(m\), \(E_m\) la suma de los tres exponentes discretos de los residuos multiplicativos \(\Pi\) y \(Z_m=3\) el número de eventos neutrales. Así, \(S_m\) no suma las evaluaciones enteras \(S=i+j\), cuyos cocientes permanecen en el registro. Con \([a]_{10}\in\{0,\ldots,9\}\), la regla de emisión decimal es
\begin{align}
 d_{m,1}&=[S_m]_{10},\notag\\
 d_{m,2}&=[S_m+E_m]_{10},\notag\\
 d_{m,3}&=[3S_m+5E_m+7Z_m]_{10},\notag\\
 U_m&=100d_{m,1}+10d_{m,2}+d_{m,3}.
 \label{eq:emisor-seis}
\end{align}
Los coeficientes enteros y el calendario de esta recurrencia son parte explícita de la ley emisora declarada. No se reciben valores de \(\pi,\varphi,e\) ni de \(\alpha\). La distinción entre una regla de emisión y una constante reconocida a partir de su salida evita que la genealogía quede escondida en una tabla de decimales.

Si \(s_m=[S_m]_{10}\) y \(e_m=[E_m]_{10}\), la misma regla se escribe
\begin{equation}
 U_m=100s_m+10[s_m+e_m]_{10}+[3s_m+5e_m+1]_{10}.
 \label{eq:acoplamiento}
\end{equation}
El uno final es el residuo de \(7Z_m=21\) módulo diez. La reducción ternaria orientada del catálogo adopta la convención
\begin{equation}
 \Psi(U_1,\ldots,U_6)=([-U_1]_3,\ldots,[-U_6]_3).
 \label{eq:psi-catalogo}
\end{equation}
El signo de esta carta es explícito: cambiarlo modifica las palabras y exige transportar las orientaciones posteriores. La identificación balanceada de cada componente de \(\FF_3\) con el TRIT se realiza después de \eqref{eq:psi-catalogo}.

\begin{proposition}[Factorización inyectiva de la emisión]
La palabra \(U=(U_1,\ldots,U_6)\) determina las dos firmas \(s=(s_1,\ldots,s_6)\) y \(e=(e_1,\ldots,e_6)\). La imagen de la emisión tiene \(18\cdot26=468\) elementos. La proyección ternaria tiene \(243\) valores distintos en el ambiente de \(729\) palabras.
\end{proposition}
\begin{proof}
Las centenas de \(U_m\) recuperan \(s_m\). Si \(b_m\) es su cifra de decenas, entonces \(e_m=[b_m-s_m]_{10}\). La cifra de unidades comprueba la tercera congruencia de \eqref{eq:acoplamiento}. Por ello el acoplamiento es inyectivo en el producto de firmas.

La enumeración de las \(324\) condiciones de un cursor mediante la recurrencia anterior produce dieciocho firmas aditivas, cada una con dieciocho preimágenes; y veintiséis multiplicativas, de las cuales veinticuatro tienen ocho preimágenes, una tiene veinticuatro y una tiene ciento ocho. La sección~\ref{app:firmas-emision} reúne las firmas, sus multiplicidades y la regla exacta de recuperación de sus condiciones iniciales. La inyectividad del acoplamiento da \(18\cdot26\) emisiones y multiplicidades
\[
 \underbrace{432}_{18\cdot24}\text{ fibras de }144,
 \qquad18\text{ fibras de }432,
 \qquad18\text{ fibras de }1944.
\]
La suma es \(432\cdot144+18\cdot432+18\cdot1944=104\,976\). Aplicar \eqref{eq:psi-catalogo} a las \(468\) emisiones produce exactamente \(243\) palabras. Esta última cifra es el cardinal de la imagen, no el del espacio \(\FF_3^6\).
\end{proof}

Esta factorización explica la diferencia entre exhaustividad finita y profundidad ilimitada. El dominio \eqref{eq:semillas} es completo para el soporte y el emisor declarados; no contiene anticipadamente cada prefijo de cada prolongación. La profundidad pertenece al sistema de extensiones compatibles que se construye sobre esas condiciones.

\subsection{Clausura del calendario y avance de memoria}

El calendario de ciento ocho transiciones consta de doce ventanas nonádicas. Su estructura de grupo es más precisa que un producto informal de dos períodos:
\begin{equation}
 0\longrightarrow C_{12}\xrightarrow{\iota}C_{108}
 \xrightarrow{p}C_9\longrightarrow0,
 \quad\iota([b])=[9b],\quad p([t])=[t]_9.
 \label{eq:extension108}
\end{equation}
Si \(t=9b+r\) con \(0\le r<9\), la suma de dos fases produce el acarreo
\[
 (b,r)+(b',r')=
 \left(b+b'+\left\lfloor\frac{r+r'}9\right\rfloor,\ [r+r']_9\right),
\]
donde la primera coordenada se reduce módulo doce. La composición conserva el término que enlaza las dos coordenadas.

\begin{proposition}[No escisión del calendario]
La sucesión \eqref{eq:extension108} es exacta y no admite una sección que sea homomorfismo de grupos.
\end{proposition}
\begin{proof}
El núcleo de la reducción módulo nueve es el subgrupo de múltiplos de nueve, imagen de \(\iota\). Si la extensión se escindiera, por conmutatividad se tendría \(C_{108}\cong C_{12}\times C_9\). El primer grupo tiene exponente \(108\), mientras el segundo tiene exponente \(\operatorname{mcm}(12,9)=36\), contradicción.
\end{proof}

Nueve transiciones restauran la fase nonádica y avanzan una ventana; ciento ocho restauran el calendario finito. Ninguno de estos hechos obliga a borrar la trayectoria o a restaurar el prefijo de una prolongación. La holonomía nonádica \(\Gamma_9\) designa el transporte de una vuelta sobre estados con memoria, de modo que
\begin{equation}
 g(\Gamma_9x)=g(x),\qquad M(\Gamma_9x)=M(x)+1.
 \label{eq:holonomia-memoria}
\end{equation}
La igualdad de fase y el incremento de memoria son compatibles y constituyen dos observaciones del mismo transporte.

% Recuperación de 02_trit_geometrias y de la holonomía con memoria.
\subsubsection{Parametrización temporal de la compatibilidad y orientación trítica}
\label{trit-temporal-compatibilidad}

La interpretación temporal del TRIT se establece sobre las transiciones del
estado completo. Sea \(x\xrightarrow{\rm adm}y\) la relación de admisibilidad determinada
por las operaciones del TPK en \(\widehat X\). Una trayectoria parametrizada
por un conjunto discretamente ordenado \(\mathcal T\) es una aplicación
\[
 \gamma:\mathcal T\longrightarrow\widehat X,\qquad
 \gamma(t)\xrightarrow{\rm adm}\gamma(t^+),
\]
donde \(t^+\) es el sucesor de \(t\). La relación de compatibilidad precede a
la elección del parámetro: el orden permite expresar sus transiciones como
una evolución. Cuando se utiliza el término sección, se declara además la
proyección \(p:E\to\mathcal T\) y la condición
\(p\circ\gamma=\operatorname{id}_{\mathcal T}\). Así quedan especificados
el espacio de estados, el orden temporal y la información que se transporta.

La firma \(\tau\) selecciona un régimen cuadrático y admite una lectura
temporal relativa a esa parametrización. El tipo \(+1\), con
\(J_{+1}^{2}=-I\), representa retorno y memoria; el tipo \(0\), con
\(J_0^{2}=0\), representa el umbral de transición; el tipo \(-1\), con
\(J_{-1}^{2}=I\), representa apertura en dos direcciones propias. En la
terminología temporal del modelo, estas funciones corresponden,
respectivamente, a pasado, presente y futuro. La correspondencia se refiere
a la posición operacional de una transición dentro de una historia:
antecedente conservado, actualización y prolongación admisible.

Esta interpretación adquiere contenido mediante tres operaciones distintas.
La restricción de una trayectoria recupera el prefijo ya construido. La
actualización \(\mathcal U_t\) compone selección, transporte y registro en
la transición observada. La prolongación considera las continuaciones que
conservan las condiciones de frontera de ese prefijo. Si
\(\mathcal P_n(x)\) denota el conjunto de continuaciones admisibles de
longitud \(n\) desde \(x\), la supresión del último paso define
\[
 r_n:\mathcal P_{n+1}(x)\longrightarrow\mathcal P_n(x).
\]
Una continuación ilimitada es una familia \((\gamma_n)_{n\geq1}\) que
satisface \(r_n(\gamma_{n+1})=\gamma_n\). El sistema inverso conserva de
este modo las restricciones que cada profundidad impone sobre las
anteriores. La existencia de una familia compatible se establece en el
dominio de prolongación correspondiente; el mero establecimiento de los
conjuntos \(\mathcal P_n(x)\) no sustituye esa demostración.

La bidireccionalidad combina, por tanto, la reconstrucción de antecedentes
con la compatibilidad de continuaciones. Sobre la sucesión de firmas, la
involución \(C_{\rm rev}\) invierte el orden y cambia el signo de cada
componente. Sobre la trayectoria actúan, además, las transformaciones de
posición, hoja y frontera especificadas por el transporte. La firma
invertida es una observación de la trayectoria conjugada, no un reemplazo
de sus restantes coordenadas. El umbral permanece distinguido: el valor
visible \(0\) conserva su función de transición incluso cuando su cociente
o su registro histórico son distintos de cero.

La relación con las geometrías exponenciales es igualmente precisa. En
cada régimen fijo, \(\exp(tJ_\tau)\) compone parámetros y conserva la
norma algebraica. El cierre elíptico pertenece a esa representación local.
El retorno de fase de \(\Gamma_9\), en cambio, actúa sobre la trayectoria
con memoria, conforme a \eqref{eq:holonomia-memoria}. Una trayectoria puede
recuperar su fase y conservar un nuevo antecedente: el dato histórico ha
cambiado aunque la observación de fase sea la misma. El umbral parabólico
retiene un desplazamiento lineal; la apertura hiperbólica separa dos
direcciones propias de expansión y contracción. Estas tres realizaciones
proporcionan leyes de composición, mientras que el TPK determina su
selección y su concatenación efectiva.

La suspensión aperiódica de la sección \ref{sec:generacion} hará explícita
esta distinción mediante un factor de fase finito, una rotación irracional
y un grado de memoria. El vínculo entre TRIT y temporalidad queda así
expresado por operaciones sobre trayectorias y por sus representaciones
cuadráticas. Su contenido excede una clasificación de signos: conserva el
régimen, el sentido del transporte y la relación entre antecedente,
transición y continuación.


La representación monodrómica de este retorno se desarrolla en la
sección~\ref{mono-ampl:conexion}. Allí se distingue la holonomía del
estado de su representación sobre fase y contador. La suspensión
simbólica de la sección~\ref{mono-ampl:cristal} define el modelo de
cristal temporal aperiódico, y la sección~\ref{mono-ampl:euler}
compone esa temporalidad con las tres realizaciones de Euler.

\subsection{Compresión observable y conservación ortogonal}

Una realización isométrica permite expresar de manera cuantitativa qué pierde una observación reducida. Sea \(U:\mathcal H\to\mathcal H\) una realización isométrica del transporte de una vuelta, \(U^*U=I\), y \(J:\mathcal H_{\rm vis}\to\mathcal H\) una inclusión isométrica. Definimos
\[
 T=J^*UJ,\qquad D=(I-JJ^*)UJ.
\]
Entonces
\begin{equation}
 UJ=JT+D,\qquad J^*D=0,\qquad I-T^*T=D^*D.
 \label{eq:conservacion-ortogonal}
\end{equation}
La última igualdad se deduce tomando el producto adjunto de la primera y usando la ortogonalidad. La contracción de la componente visible no constituye por sí sola destrucción de información del estado total.

\begin{proposition}[Conservación a horizonte finito]
Si \(I-T^*T=D^*D\), entonces, para todo \(n\ge1\),
\begin{equation}
 I=\sum_{r=0}^{n-1}T^{*r}D^*DT^r+T^{*n}T^n.
 \label{eq:telescopia}
\end{equation}
\end{proposition}
\begin{proof}
Multiplicar \(D^*D=I-T^*T\) por \(T^{*r}\) y \(T^r\) y sumar produce una suma telescópica. Se cancelan los términos interiores y queda \(I-T^{*n}T^n\).
\end{proof}

El último sumando conserva el estado terminal. Suprimirlo antes de justificar un límite alteraría la igualdad. Esta distinción es el contenido operativo de la conservación de la unidad en la compresión: la norma se distribuye entre la componente visible, la memoria y el terminal, sin ser reemplazada por un único residuo.

\subsection{Desarrollo local: conmutación, firma lorentziana y curvatura discreta}
\label{vii:sec:bloque-compresion}
\label{iv:extension-curvatura-mixta}

El bloque diagonal adyacente con diagonales iguales módulo nueve es único. En efecto,
\[
 k^2\equiv(k+1)^2\pmod9\quad\Longleftrightarrow\quad2k+1\equiv0\pmod9
\]
da \(k=4\). Su matriz visible es
\[
 B_c=\begin{pmatrix}7&2\\2&7\end{pmatrix}=7I+2R,
 \quad R=\begin{pmatrix}0&1\\1&0\end{pmatrix},\quad R^2=I.
\]
Con \(P_\pm=(I\pm R)/2\), se obtiene \(B_c=9P_++5P_-\). La unidad de la normalización estocástica y la conservación de una forma indefinida son dos realizaciones de este bloque:
\[
 \frac{B_c}{9}=P_++\frac59P_-,\qquad
 \frac{B_c}{\sqrt{45}}=\exp(\eta R),\quad\eta=\frac12\log\frac95.
\]
Para \(Q=\operatorname{diag}(1,-1)\), la segunda matriz satisface \(B_c^{\mathsf T}QB_c=45Q\). De aquí procede una realización local en \(SO^+(1,1)\); no se requiere introducir de antemano un espacio-tiempo de cuatro dimensiones para demostrarla.

El promediado uniforme sobre las clases \(C_1=\{1,4,7\}\), \(C_2=\{2,5,8\}\), \(C_0=\{3,6,9\}\) da
\[
 M_\Pi=\begin{pmatrix}4&5&6\\5&4&6\\6&6&9\end{pmatrix},\quad
 M_\Sigma=\begin{pmatrix}5&6&4\\6&4&5\\4&5&6\end{pmatrix}.
\]
El polinomio característico de \(M_\Pi\) es
\[
 (\lambda+1)((\lambda-9)^2-72),
\]
de donde su forma cuadrática tiene firma \((2,1)\). El salto espectral local \(9-5=4\), el coeficiente \(2\), la diferencia agregada \(51-45=6\) y el defecto integrado \(9\cdot6=54\) describen niveles distintos de la misma comparación aritmética. No son índices intercambiables de un mismo operador.

La curvatura discreta puede describirse directamente antes de esta realización espectral. En residuos, sean \(S(x,y)=x+y\), \(P(x,y)=xy\) y \(\Delta_x,\Delta_y\) las diferencias progresivas. Las variables de enlace son diferencias de potencial, \(A_x^T=\Delta_xT\), \(A_y^T=\Delta_yT\), y se fija
\[
 F(T_x,T_y)=\Delta_x\Delta_yT_y-\Delta_y\Delta_xT_x.
\]
Como \(\Delta_xS=\Delta_yS=1\), \(\Delta_xP=y\) y \(\Delta_yP=x\),
\begin{equation}
 F(S,S)=F(P,P)=0,\qquad F(S,P)=+1,\quad F(P,S)=-1\pmod9.
\end{equation}
La suma y el producto, tomados separadamente, inducen conexiones planas; la mezcla ordenada cambia el signo de la curvatura de plaqueta. Al sumar sobre un rectángulo de lados \(a,b\), las aristas interiores se cancelan y la suma orientada de borde vale \(\pm ab\) módulo nueve. Así aparece una relación exacta entre orden de composición y orientación de frontera.

\subsection{Completación cúbica y dos nociones de frontera}

La ampliación del soporte cúbico de lado nueve al de lado diez admite una estratificación posicional exacta:
\begin{equation}
 10^3=9^3+3\cdot9^2+3\cdot9+1=729+243+27+1.
 \label{eq:cubo}
\end{equation}
El primer término corresponde a las tres coordenadas interiores; el segundo, a una coordenada nueva; el tercero, a dos; y el último, al vértice de tres coordenadas nuevas. La corona de ampliación contiene \(271\) posiciones. En cambio, la frontera interna del cubo discreto de lado nueve contiene \(9^3-7^3=386\). La diferencia entre estos recuentos es una diferencia de dominio, no un conflicto de valores.

Dividir \eqref{eq:cubo} por \(1000\) da una partición normalizada de la unidad. La incidencia del cubo distingue seis caras, doce aristas y ocho vértices. El enlace ulterior con hexadas y octadas de Steiner requiere los mapas combinatorios correspondientes; los cardinales del cubo no los sustituyen. Esta completación aporta aquí la relación entre las bases \(729\) y \(1000\), mientras la prolongación excepcional se desarrollará después de la generación numérica, del registro dodecafásico y de las construcciones de \(\alpha\).

\subsubsection{Estratificación de la completación y restitución del ápice}
\label{sec:revision-emrh}

La completación admite una descripción conjuntista que hace explícita su
operación. Sea \(E_9\) el conjunto de las nueve clases residuales, representadas
positivamente por \(1,\ldots,9\), y sea \(\ast\) un elemento adicional.
El símbolo \(\ast\) es distinto de la clase nula, ya representada por \(9\).
En consecuencia, la ampliación de \(E_9\) a
\(\widehat E=E_9\sqcup\{\ast\}\) aumenta el cardinal de nueve a diez
sin modificar la identificación modular interna. La completación cúbica es
\(\widehat C=\widehat E^3\), y el cubo inicial es \(C=E_9^3\).

\begin{definition}[Estratos de codimensión combinatoria]
Para \(0\leq r\leq3\), se define
\[
 S_r=\{(x_1,x_2,x_3)\in\widehat E^3:
             \#\{j:x_j=\ast\}=r\}.
\]
El ápice de la completación es \(a_\ast=(\ast,\ast,\ast)\), y la
completación perforada es \(\widehat C^{\circ}=\widehat C\setminus\{a_\ast\}\).
\end{definition}

\begin{proposition}[Descomposición exacta de la completación]
Los estratos son disjuntos y satisfacen
\begin{align}
 \widehat C&=S_0\sqcup S_1\sqcup S_2\sqcup S_3,
 &|S_r|&=\binom3r9^{3-r},\label{eq:revision-estratos}\\
 |\widehat C^{\circ}|&=729+243+27=999,
 &|\widehat C|&=999+1=1000.\label{eq:revision-999}
\end{align}
La ampliación perforada contiene \(270\) elementos adicionales al cubo
inicial; la ampliación completa contiene \(271\).
\end{proposition}
\begin{proof}
Un elemento de \(S_r\) se obtiene eligiendo las \(r\) coordenadas que
contienen \(\ast\), de \(\binom3r\) maneras, y asignando a las restantes
una de las nueve clases de \(E_9\). La clasificación por el número de
coordenadas adicionales es exhaustiva y sus clases son disjuntas. El único
elemento de \(S_3\) es \(a_\ast\). Retirarlo resta una unidad al cardinal
total; restituirlo recupera el producto cartesiano completo.
\end{proof}

La expresión \emph{extrema y media razón holográfica}, abreviada EMRH en
el corpus, designa aquí esta articulación entre el cuerpo inicial, los
estratos de ampliación y la restitución unitaria. Su contenido inmediato
es la partición exacta de un producto finito y su normalización. El número
\(999\) representa un objeto determinado, la completación perforada;
\(1000\) representa su clausura por restitución del ápice. La unidad
añadida tiene, por tanto, una residencia de incidencia explícita.

\begin{figure}[htbp]
\centering
\begin{tikzpicture}[x=1cm,y=1cm,>=Stealth,font=\small]
\draw[->] (0,0)--(12.4,0);
\foreach \x/\a/\b in {0/729/{S_0},4.0/243/{S_1},7.5/27/{S_2},10.6/1/{S_3}}{
  \draw (\x,0.12)--(\x,-0.12);
  \node[above=5pt] at (\x,0) {\(\a\)};
  \node[below=5pt] at (\x,0) {\(\b\)};
}
\draw[decorate,decoration={brace,amplitude=5pt}] (0,1)--(7.7,1);
\node[above=8pt] at (3.85,1) {Completación perforada: \(999\)};
\draw[decorate,decoration={brace,mirror,amplitude=5pt}] (0,-1)--(10.8,-1);
\node[below=8pt] at (5.4,-1) {Completación íntegra: \(1000=999+1\)};
\node[align=center] at (10.6,1.5) {Ápice\\\((\ast,\ast,\ast)\)};
\end{tikzpicture}
\caption{Estratos de la completación cúbica. La disposición representa
la descomposición conjuntista, no una escala métrica. Los \(243\) elementos
de \(S_1\) son tres estratos coordenados de \(81\) elementos; los \(27\)
de \(S_2\) son tres estratos de \(9\). El ápice completa la partición.}
\label{fig:revision-emrh}
\end{figure}

\subsubsection{Conservación de la unidad y compatibilidad dimensional}

La normalización no se limita a dividir una identidad aislada. En dimensión
\(d\), el producto \(\widehat E^d\) lleva la medida uniforme
\(\mu_d(\{x\})=10^{-d}\). La proyección coordenada
\(p_d:\widehat E^{d+1}\to\widehat E^d\) suprime la última coordenada.

\begin{proposition}[Compatibilidad proyectiva de las medidas normalizadas]
Para todo \(d\geq1\),
\[
 (p_d)_*\mu_{d+1}=\mu_d,
 \qquad \mu_d(\widehat E^d)=1.
\]
La masa del estrato con \(r\) coordenadas adicionales es
\(\binom dr9^{d-r}/10^d\), y la suma de las masas de todos los estratos
es uno.
\end{proposition}
\begin{proof}
Cada punto de \(\widehat E^d\) tiene diez preimágenes. Su masa total es
\(10\,10^{-(d+1)}=10^{-d}\). La igualdad para cualquier subconjunto
se obtiene sumando sobre sus puntos. La fórmula estratificada procede
del mismo recuento que \eqref{eq:revision-estratos}; el binomio
\((9+1)^d\) da la normalización total.
\end{proof}

En dimensión tres, retirar el ápice deja masa \(999/1000\), y
restituirlo añade \(1/1000\). Renormalizar antes de la restitución
definiría otra medida: la distinción es operacionalmente relevante.
En particular, conservación de masa probabilística, conservación de los
registros de transporte y disipación termodinámica son afirmaciones con
objetos diferentes. La proposición demuestra la primera y proporciona el
soporte compatible para estudiar la segunda; la tercera requiere una
realización dinámica y energética especificada.

También permanece diferenciada la frontera geométrica interna.
Sobre el representante cartesiano \(\{1,\ldots,9\}^3\), retirar
\(\{2,\ldots,8\}^3\) deja \(386\) puntos. Esta frontera pertenece al
cubo inicial, mientras \(\widehat C\setminus C\) pertenece a su
ampliación. La completación cúbica, la expansión posicional de base mil y
la lectura de incidencia excepcional comparten datos del soporte;
sus mapas respectivos determinan qué información de esos datos utiliza
cada representación.


\subsection{Incidencia cúbica y realización lorentziana}

La completación cúbica aporta una relación entre interior y frontera.
Su geometría puede desarrollarse mediante la incidencia de las seis caras
con los ocho vértices. Etiquetemos las caras por \((i,\varepsilon)\),
\(i\in\{1,2,3\}\), \(\varepsilon\in\{+1,-1\}\), y los vértices por
\(\nu\in\{+1,-1\}^3\). La matriz de incidencia es
\[
 M_{(i,\varepsilon),\nu}=\mathbf1_{\{\nu_i=\varepsilon\}}.
\]
Sean \(u=(1,\ldots,1)^{\mathsf T}\in\RR^6\), \(R_F\) la oposición de
caras y \(P_u=uu^{\mathsf T}/6\), \(P_-=(I-R_F)/2\).

\begin{proposition}[Cociente de incidencia de dimensión cuatro]
El Gramiano de incidencia satisface
\[
 MM^{\mathsf T}=12P_u+4P_-.
\]
Sus autovalores \(12,4,0\) tienen multiplicidades \(1,3,2\).
En consecuencia,
\[
 Q_\square=\RR^6/\ker(MM^{\mathsf T})
 \simeq\RR u\oplus E_-,\qquad E_-=\operatorname{im}P_-.
\]
\end{proposition}
\begin{proof}
Una cara contiene cuatro vértices; dos caras opuestas tienen intersección
vacía, y dos caras distintas no opuestas comparten dos vértices.
La matriz \(12P_u+4P_-\) tiene precisamente esas entradas.
El espacio uniforme tiene dimensión uno; las diferencias entre caras
opuestas generan un espacio de dimensión tres, ortogonal al uniforme.
Su complemento tiene dimensión dos y pertenece al núcleo.
\end{proof}

La descomposición \(1+3\) procede de la incidencia. La elección temporal
se declara mediante la forma bilineal
\[
 B_\square([x],[y])=\langle x,(P_--P_u)y\rangle.
\]
Está bien definida en el cociente, porque ambos proyectores anulan el
núcleo. En la base
\[
 t=[u]/\sqrt6,\qquad
 d_i=[e_{i,+}-e_{i,-}]/\sqrt2,
\]
su matriz es \(\operatorname{diag}(-1,1,1,1)\). Quedan explicitados
el dominio, la descomposición y la convención de signo que producen
la realización de Minkowski. El Gramiano positivo determina el cociente;
la asignación temporal determina la forma indefinida sobre él.

El operador
\[
 W_\square=\sqrt6P_u+\sqrt2P_-,\qquad
 W_\square e_{i,\varepsilon}=t+\varepsilon d_i
\]
envía las seis etiquetas de cara a direcciones nulas, puesto que
\(B_\square(t+\varepsilon d_i,t+\varepsilon d_i)=-1+1=0\).
La inversión de una dirección y la oposición de caras disponen así
de una realización geométrica concreta.

\subsubsection{Soldadura, inversión y subdivisión}

La realización celular utiliza un transporte invertible \(U_m(a)\),
que conserva la forma lorentziana, entre las fibras de los extremos de
una arista dual orientada \(a\), una etiqueta
de cara \(\lambda_m(a)\) y una orientación \(\sigma_m(a)\).
Con \(\ell_m=\ell_0/9^m\), la aplicación de soldadura es
\[
 \beta_m(a)=\ell_m\sigma_m(a)
 W_{\square,s(a)}e_{\lambda_m(a)}.
\]
La inversión compatible satisface
\[
 \beta_m(\bar a)=-U_m(a)\beta_m(a).
\]
Identifiquemos las fibras del origen en dos niveles consecutivos mediante
la isometría de refinamiento declarada. Si los nueve subsegmentos
\(a_1,\ldots,a_9\), de nivel \(m+1\), conservan la etiqueta
y la orientación al ser transportados a esa fibra común, entonces
\begin{equation}
 \beta_m(a)=\sum_{j=1}^9
 U_{m+1}(a_1\cdots a_{j-1})^{-1}\beta_{m+1}(a_j).
 \label{eq:soldadura-refinamiento}
\end{equation}
Cada término transportado equivale a \(\ell_m/9\) por la misma
dirección nula, lo que demuestra la identidad. La hipótesis de
compatibilidad determina qué subdivisiones realizan este refinamiento.
La ecuación reúne dirección, escala y transporte, en lugar de limitar
la correspondencia geométrica a un recuento de dimensiones.

Fijadas además la orientación y la forma lorentziana, la dualidad de
Hodge en \(\Lambda^2Q_\square^*\) satisface \(\star^2=-I\).
En dimensión cuatro, el signo es \((-1)^{2(4-2)+1}=-1\).
Se obtiene una estructura compleja real sobre las dos-formas y,
tras complejificación, los subespacios propios de valores \(+i\) y
\(-i\). La realización electromagnética utilizará este dominio
después de construir la sección electrónica y declarar su acoplamiento.

\subsection{Compatibilidad entre las reducciones modular y posicional}

La base decimal por bloques interviene mediante una división euclídea
que conserva simultáneamente residuo y cociente:
\[
 n=1000q+r,\qquad 0\le r<1000.
\]
Puesto que \(1000\equiv1\pmod9\) y \(1000\equiv1\pmod3\),
\[
 n\equiv q+r\pmod9,\qquad n\equiv q+r\pmod3.
\]
Por tanto, la reducción de un bloque requiere la contribución del
cociente para recuperar la clase del entero. Los números \(0\) y
\(1000\) tienen el mismo resto por \(1000\), pero clases diferentes
por \(9\) y por \(3\). Ésta es una razón operatoria para conservar
la memoria durante el cambio de resolución.

El refinamiento ternario y la determinación de un prefijo decimal
pueden coordinarse mediante un único residuo entero. Si \(N/3^L\)
es un extremo de cilindro y \(D/1000^r\) un extremo decimal, definimos
\[
 A=1000^rN-D3^L.
\]
Al añadir seis símbolos ternarios, \(N'=729N+\nu\),
\(L'=L+6\), con \(0\le\nu<729\). Si el prefijo decimal permanece
fijo, se obtiene
\[
 A'=729A+\nu1000^r.
\]
Si se incorpora además un bloque decimal \(\delta\), entonces
\(D'=1000D+\delta\), \(r'=r+1\), y
\begin{equation}
 A'=729000A+\nu1000^{r+1}-729\delta3^L.
 \label{eq:residuo-dos-bases}
\end{equation}
Ambas identidades se demuestran sustituyendo las definiciones. La
comparación de extremos se reduce así a operaciones enteras que
retienen la contribución de cada prolongación. Esta compatibilidad
es la que permite refinar una región antes de evaluar su coordenada real.

% Incorporación al final de extension.tex; no contiene una sección principal.
% Procedencia: RESULTADO_RECUPERADO; explicitación sectorial de las hipótesis.
% Fuentes leídas: monografía Doble proyección holográfica (cantera editorial),
% manuscrito/local/01_tesis_recuperada.tex y
% manuscrito/generated/algebra_holonomica.tex, especialmente §§1--7 y 21.
% Definición rectora de número: manuscrito/flattened/main_autosuficiente.tex,
% secciones que comienzan en líneas 22137 y 23000 de la cantera de 775 páginas.
% No se adopta su presentación global como prueba de todas las flechas TPK.

\subsection{El número como sección compatible y el valor como lectura}
\label{hist:seccion}

La distinción entre estado y observación permite precisar qué objeto se
prolonga cuando aumenta la resolución. A profundidad \(n\), sea \(X_n\)
el conjunto de estados admisibles producidos por APP--TRIT--TPK. Sus datos
incluyen posición, hojas suma y producto, residuos y cocientes, régimen,
orientación, acarreo, fase, ruta, frontera, memoria y supervivencia. No son
coordenadas independientes: las evaluaciones aritméticas y los transportes
anteriores imponen sus relaciones. Una prolongación conserva esas relaciones
en cada antecedente, aunque añada nuevos pasos y nuevos datos de memoria.
Éste es el sentido de la aritmética de historias desarrollada en
(secciones~\ref{sec:extension} y~\ref{mono-ampl:conexion}).

Sean \(r_{m,n}:X_m\to X_n\), para \(m\geq n\), las restricciones de
profundidad, con \(r_{n,n}=\operatorname{id}\) y
\(r_{\ell,n}=r_{m,n}\circ r_{\ell,m}\). Cada restricción conserva el
prefijo y la información ya determinada en él.

\begin{definition}[Sección numérica compatible]
Un número HMT es una historia compatible del sistema de estados:
\begin{equation}
 x=(x_n)_{n\geq0}\in X_\infty:=\varprojlim_n X_n,
 \qquad r_{m,n}(x_m)=x_n.
 \label{hist:limite-inverso}
\end{equation}
Un lector es una operación adicional sobre un dominio especificado de estas
historias. Una cifra es una observación finita; un valor escalar es una
evaluación del lector. El par formado por la sección y su lector constituye
una presentación leída, no una redefinición de la identidad de la sección.
\end{definition}

En el dominio arquimediano \(X_\infty^{\mathrm{Arch}}\), el lector asocia
a cada prefijo un cilindro cerrado, conserva su encajamiento y hace tender
su diámetro a cero. Así determina
\begin{equation}
 \operatorname{ev}_{\mathbb R}:X_\infty^{\mathrm{Arch}}\longrightarrow
 \mathbb R,\qquad
 \{\operatorname{ev}_{\mathbb R}(x)\}=\bigcap_n I_n(x).
 \label{hist:evaluacion}
\end{equation}
La construcción concreta de estos cilindros y la decisión de sus fronteras
se desarrollan en la sección~\ref{sec:generacion}. La definición
\eqref{hist:limite-inverso} no afirma por sí sola que la imagen de
\eqref{hist:evaluacion} sea toda la recta: una afirmación de cobertura
requiere su propio teorema. Tampoco una igualdad de valores establece, sin
un resultado de inyectividad, la igualdad de sus historias. Esta separación
permite hablar de precisión como profundidad conservando la diferencia
entre el antecedente aritmético y su realización escalar.

\subsection{Composición de historias y multiplicador de memoria}
\label{hist:algebra}

Fijemos un sector de estados a profundidad \(n\) y una categoría
\(\mathcal C\) cuyas flechas sean sus historias admisibles. La composición
se escribe \(vu=v\circ u\): primero se recorre \(u\) y después \(v\).
Los pasos elementales proceden de la selección, el transporte y la
actualización del TPK. Una fibra trítica local admite la escritura
\(\mathbb R[J_x]/(J_x^2+\tau(x)1_x)\), donde
\(\tau(x)\in\{+1,0,-1\}\). El régimen y la orientación pertenecen al
estado que se transporta; una transición entre regímenes no equivale a
identificar sus álgebras cuadráticas.

La siguiente formulación precisa un sector algebraico de esa composición.
Para cada objeto \(x\), se da un álgebra real asociativa unitaria \(A_x\),
que puede ser no conmutativa. Para cada \(u:x\to y\), se da un
homomorfismo unital \(\rho_u:A_x\to A_y\). Suponemos una acción estricta:
\begin{equation}
 \rho_{\operatorname{id}_x}=\operatorname{id}_{A_x},\qquad
 \rho_v\rho_u=\rho_{vu}.
 \label{hist:accion-estricta}
\end{equation}
Para cada pareja componible \(u:x\to y\), \(v:y\to z\), fijamos un
multiplicador central invertible
\(\Omega(v,u)\in Z(A_z)^\times\). Se exige la normalización
\(\Omega(\operatorname{id}_y,u)=\Omega(u,\operatorname{id}_x)=1_y\)
y, para cada triple componible, la identidad
\begin{equation}
 \rho_w\bigl(\Omega(v,u)\bigr)\Omega(w,vu)
   =\Omega(w,v)\Omega(wv,u).
 \label{hist:cociclo}
\end{equation}
Son hipótesis explícitas sobre los transportes y sus coeficientes, no una
consecuencia de llamar holonómica a la dinámica. Al aplicar esta presentación
a un sector TPK deben identificarse sus \(\rho_u\) y \(\Omega(v,u)\).
Las transiciones expresadas mediante bimódulos requieren conservar esa
estructura y no quedan reemplazadas por \eqref{hist:accion-estricta}.

El espacio de sumas finitas
\[
 \mathcal H(\mathcal C,A,\rho,\Omega)
   =\bigoplus_{u\in\operatorname{Mor}(\mathcal C)} A_{t(u)}T_u
\]
recibe el producto bilineal
\begin{equation}
 (aT_v)(bT_u)=a\rho_v(b)\Omega(v,u)T_{vu},
 \label{hist:producto}
\end{equation}
si las flechas son componibles, y cero en otro caso. Aquí \(t(u)\)
designa el estado terminal. La suma reúne historias con coeficientes; el
producto concatena historias conservando el multiplicador del registro.

\begin{proposition}[Asociatividad sectorial con multiplicador central]
Bajo \eqref{hist:accion-estricta} y \eqref{hist:cociclo}, el producto
\eqref{hist:producto} es asociativo. Si el conjunto de objetos es finito,
su unidad es \(\sum_x1_xT_{\operatorname{id}_x}\).
\end{proposition}
\begin{proof}
Para \(u:x\to y\), \(v:y\to z\), \(w:z\to t\), sean
\(a\in A_y\), \(b\in A_z\), \(c\in A_t\). Los dos coeficientes
de \(T_{wvu}\) son
\begin{align*}
 ((cT_w)(bT_v))(aT_u):\quad&
 c\rho_w(b)\Omega(w,v)\rho_{wv}(a)\Omega(wv,u),\\
 (cT_w)((bT_v)(aT_u)):\quad&
 c\rho_w(b)\rho_w\rho_v(a)\rho_w(\Omega(v,u))\Omega(w,vu).
\end{align*}
La acción estricta iguala \(\rho_w\rho_v(a)\) con \(\rho_{wv}(a)\).
La centralidad de \(\Omega(w,v)\) permite desplazarlo a la derecha de
ese factor; \eqref{hist:cociclo} iguala entonces los productos restantes.
No se permutan los coeficientes \(a,b,c\). Si el triple no es componible,
ambas parentizaciones son cero por sus estados iniciales y terminales.
La bilinealidad concluye la prueba. Las identidades locales y la
normalización de \(\Omega\) verifican la unidad indicada.
\end{proof}

El acarreo del calendario proporciona un ejemplo efectivo. En la categoría
de un objeto y flechas \(C_9\), tomemos coeficientes
\(A=\mathbb R[z,z^{-1}]\), acción identidad y
\[
 \Omega(b,a)=z^{c_0(a,b)},\qquad
 c_0(a,b)=\left\lfloor\frac{a+b}{9}\right\rfloor,
 \quad 0\leq a,b<9.
\]
Las dos sumas de acarreos de un triple son
\(\lfloor(a+b+c)/9\rfloor\); por ello se cumple
\eqref{hist:cociclo}. La variable formal \(z\) registra la vuelta sin
asignarle un valor numérico. Imponer después \(z^{12}=1\) recupera el
registro dodecafásico del calendario \eqref{eq:extension108}; conservar
\(z\) sin esa relación distingue vueltas sucesivas. El ejemplo realiza
la memoria de este calendario, sin identificarla con la totalidad de la
memoria TPK.

\subsection{Doble varianza del refinamiento y conservación de la unidad}
\label{hist:doble-varianza}

Al restringir estados se prolongan observables en sentido contrario.
Consideremos las álgebras de funciones, con producto puntual,
\(D_n=\operatorname{Fun}(X_n,\mathbb R)\). La precomposición define
\begin{equation}
 r_{m,n}^*:D_n\longrightarrow D_m,
 \qquad r_{m,n}^*f=f\circ r_{m,n}.
 \label{hist:precomposicion}
\end{equation}
Son morfismos unitales; son además inyectivos cuando las restricciones
correspondientes son sobreyectivas. El límite directo
\(D_{\mathrm{cyl}}=\varinjlim_n(D_n,r_{n+1,n}^*)\) reúne las
observaciones de profundidad finita. No se identifica aquí este álgebra
con toda el álgebra de rutas: prolongar los operadores de transporte exige
también sus compatibilidades con el refinamiento.

\begin{lemma}[Unidad y naturalidad de la evaluación]
Si \(e_x\) es el indicador de \(x\in X_n\), se cumple
\begin{equation}
 r_{m,n}^*e_x=\sum_{y:\,r_{m,n}(y)=x}e_y,
 \qquad r_{m,n}^*1_n=1_m,
 \label{hist:unidad}
\end{equation}
y, para todo \(f\in D_n\), \(y\in X_m\),
\[
 \langle r_{m,n}^*f,y\rangle=\langle f,r_{m,n}y\rangle.
\]
Cada sección \(x\in X_\infty\) determina una evaluación bien definida
\(D_{\mathrm{cyl}}\to\mathbb R\), dada por \([f]\mapsto f(x_n)\)
para \(f\in D_n\).
\end{lemma}
\begin{proof}
Las fibras de \(r_{m,n}\) forman una partición de \(X_m\). Sus
indicadores son ortogonales y suman \(1_m\), lo que prueba
\eqref{hist:unidad}. La naturalidad es la definición de precomposición.
Finalmente, \((r_{m,n}^*f)(x_m)=f(x_n)\), de modo que dos representantes
compatibles dan la misma evaluación en el límite directo.
\end{proof}

La unidad global es la función constante uno sobre todas las posibilidades;
la sección individual selecciona una historia compatible. Conservar la
primera no significa inmovilizar la segunda. El refinamiento puede aumentar
la cantidad de extensiones y la profundidad de cada historia manteniendo
la unidad y todas las evaluaciones de sus antecedentes.

\subsection{Retorno de fase, representación de Weyl y dos lecturas}
\label{hist:puente-lecturas}

La ley \eqref{eq:holonomia-memoria} adquiere una realización operatoria
en la suspensión bilateral que se construirá en la
sección~\ref{sec:generacion}. Allí \(T\) es el avance de un paso,
\(V_9\) observa la fase nonádica, \(W_\delta\) la fase de rotación y
\(D_M\) el grado de memoria. Escribiendo
\(\widehat\Gamma_9=T^9\) para la vuelta en esta representación,
las relaciones de Weyl \eqref{eq:weyl-relojes} implican, sobre vectores
de soporte finito,
\begin{equation}
 \begin{aligned}
 V_9\widehat\Gamma_9&=\widehat\Gamma_9V_9,\qquad
 W_\delta\widehat\Gamma_9=e^{2\pi i9\delta}
     \widehat\Gamma_9W_\delta,\\
 [D_M,\widehat\Gamma_9]&=\widehat\Gamma_9.
 \end{aligned}
 \label{hist:tres-retornos}
\end{equation}
La primera igualdad sigue de \(\zeta_9^9=1\); la segunda, de nueve
aplicaciones de la covariancia de \(W_\delta\); la tercera expresa
el incremento unitario de memoria. La fase finita retorna, mientras la
rotación irracional y la memoria distinguen el estado siguiente. Los
caracteres y el parámetro \(\delta\) se realizan después desde la
construcción numérica allí especificada: no seleccionan retrospectivamente
las semillas. Esta representación describe observables y traslaciones de
la suspensión; no reemplaza todas las coordenadas ni todos los operadores
del TPK.

La doble lectura prolonga la misma distinción entre historia y observable.
En su dominio arquimediano se aplica \eqref{hist:evaluacion}; en el
dominio incidencial calibrado, los mapas \(Q_n\) de la
sección~\ref{exc:doble-lectura} conservan las palabras y los marcos de
frontera. La igualdad allí demostrada
\(r_{n,m}Q_n=Q_mp_{n,m}\) determina su prolongación \(Q_\infty\).
Sobre el dominio común de ambas lecturas pueden considerarse conjuntamente
\(\operatorname{ev}_{\mathbb R}(x)\) y \(Q_\infty(x)\): la
primera retiene el valor; la segunda, la incidencia transportada que su
codominio conserva. La sección antecede a ambas. Esta articulación permite
pasar de la aritmética de historias a las constantes y a la lectura
excepcional sin convertir sus realizaciones en nuevos datos iniciales.


\section{Construcción correlativa del continuo: historias, transporte, incidencia y terminal}
\label{iv:continuo-conjunto}
La aritmética publicada por una forma normal conserva una lectura precisa del estado que la produce. Para situar esa lectura dentro de HMT es necesario mantener también las historias, las hojas y las operaciones que actúan sobre ellas. Este capítulo reúne ese sustrato común. Sus distintas construcciones se obtienen del mismo árbol de ejecuciones APP–TRIT–TPK; los cortes de exposición permiten estudiar sus propiedades sin convertirlas en dominios independientes elegidos por una aplicación posterior.

El punto de partida es el núcleo formal expuesto anteriormente. Una emisión admisible contiene la operación de APP, la orientación del TRIT, el cambio de hoja, el residuo y el cociente, junto con la actualización efectiva de memoria del TPK. Cuando se publica solamente un residuo, la historia de ejecución sigue perteneciendo al objeto de partida. Dos ejecuciones con idéntica lectura final no se identifican por ese solo motivo.

\subsection{Historias admisibles y memoria por composición}
\label{iv:cont:historias}
Sea \(X_k\) el conjunto de historias admisibles de longitud \(k\). Sus elementos se construyen por emisiones sucesivas desde el dominio inicial del núcleo. Escribimos \(h\varepsilon\) para la prolongación de \(h\) por una emisión admisible \(\varepsilon\), y \(\rho_k(h\varepsilon)=h\) para la supresión de la última emisión. Se conserva la etiqueta de la emisión, incluso cuando diferentes etiquetas producen un mismo estado publicado.

En las dos hojas aritméticas, los datos locales satisfacen las divisiones exactas

\[
\delta_\diamond=r_\diamond+9q_\diamond,
\qquad 0\leq r_\diamond<9,
\qquad \diamond\in\{\Sigma,\Pi\}.
\]
La coordenada orientada del TRIT admite, análogamente, una escritura \(\Delta=o_r+3\kappa\), con el representante orientado declarado por el núcleo. La conservación del cociente permite componer esas publicaciones sin sustituir una emisión por su residuo aislado.

En una carta de memoria, la actualización de una emisión se escribe

\[
U_\varepsilon(m)=C_\varepsilon m+\kappa_\varepsilon.
\]
Para una trayectoria \(\alpha=(\varepsilon_1,\ldots,\varepsilon_k)\), la operación efectiva es la composición ordenada de esas actualizaciones. Su parte lineal y su traslación son

\[
C_\alpha=C_{\varepsilon_k}\cdots C_{\varepsilon_1},
\qquad
\kappa_\alpha=
\sum_{j=1}^{k}C_{\varepsilon_k}\cdots
C_{\varepsilon_{j+1}}\kappa_{\varepsilon_j}.
\]
El producto vacío de la última suma es la identidad.

\begin{proposition}[Compatibilidad afín de la concatenación]
\label{iv:cont:concatenacion}
Si \(\alpha\) precede a \(\beta\), entonces

\[
C_{\beta\alpha}=C_\beta C_\alpha,
\qquad
\kappa_{\beta\alpha}=C_\beta\kappa_\alpha+\kappa_\beta,
\qquad
U_{\beta\alpha}=U_\beta U_\alpha.
\]
\end{proposition}
\begin{proof}
Sustituir \(U_\alpha(m)\) en \(U_\beta\) produce \(C_\beta C_\alpha m+C_\beta\kappa_\alpha+\kappa_\beta\). La descomposición en parte lineal y traslación da las tres identidades. La asociatividad de la composición demuestra que el resultado no depende de una subdivisión auxiliar de la trayectoria.
\end{proof}

Esta identidad de cociclo es el contenido preciso de la memoria acumulada. La cantidad transportada por una segunda trayectoria depende de la actualización de carta que ésta realiza; no es, en general, la suma sin transporte de dos registros.

El retorno nonádico conserva una segunda distinción. En las coordenadas de fase y vuelta, con \(\bar r\in\{0,\ldots,8\}\), su publicación es

\[
\sigma_9(w,r)=
\left(w+\left\lfloor\frac{\bar r+1}{9}\right\rfloor,
r+1\pmod9\right),
\qquad
\sigma_9^9(w,r)=(w+1,r).
\]
Durante nueve pasos se produce exactamente un cruce de \(8\) a \(0\). La fase retorna y la coordenada de vuelta aumenta. La holonomía enriquecida retiene además las actualizaciones de memoria que acompañan ese recorrido.

\subsection{Refinamiento del árbol y medida de historias}
\label{iv:cont:medida}
Se fija el conjunto inicial finito no vacío de semillas admisibles, con una masa inicial estrictamente positiva y de suma \(1\). En este apartado cada historia tiene un número finito y positivo de prolongaciones. La supresión de emisiones define las aplicaciones de refinamiento. La admisibilidad del núcleo proporciona las prolongaciones; cuando una residencia utiliza una emisión neutra, ésta da explícitamente una sección de \(\rho_k\). La existencia de esa sección se refiere al árbol que incluye dicha emisión, no a cualquier restricción arbitraria de rutas.

Sea \(d(h)\geq1\) el número de emisiones admisibles en \(h\). La elección equiprobable entre sus hijos define

\[
\mu(h\varepsilon)=\frac{\mu(h)}{d(h)}.
\]
Más generalmente, una distribución local positiva \(p(\varepsilon\mid h)\), cuya suma es \(1\), define \(\mu(h\varepsilon)=\mu(h)p(\varepsilon\mid h)\). Esta segunda formulación permite conservar las ponderaciones del lector cuando los números de prolongaciones no son uniformes.

\begin{proposition}[Compatibilidad cilíndrica y esperanza condicional]
\label{iv:cont:esperanza}
En cada nivel se cumple

\[
\sum_{\rho_k(y)=x}\mu(y)=\mu(x).
\]
En los espacios \(L^2(X_k,\mu_k)\), las aplicaciones

\[
(J_kf)(y)=f(\rho_k(y)),
\qquad
(E_kg)(x)=\sum_{\rho_k(y)=x}\frac{\mu(y)}{\mu(x)}g(y)
\]
satisfacen \(E_k=J_k^*\), \(E_kJ_k=I\) y \(\|J_kf\|=\|f\|\).

\end{proposition}
\begin{proof}
La primera identidad es la normalización de las probabilidades locales. Aplicándola a \(|f(x)|^2\) se obtiene la isometría. Al agrupar por padres la suma que define \(\langle J_kf,g\rangle\), se obtiene la fórmula de \(E_k\); sobre una función constante en cada fibra, esa fórmula devuelve su valor.
\end{proof}

Las masas compatibles definen una medida de probabilidad sobre el espacio de historias infinitas: el cilindro de una historia finita recibe su masa \(\mu(h)\). La medida se extiende desde la álgebra de cilindros, y su unicidad procede de que éstos generan la sigma-álgebra. Si el árbol es finitamente ramificado y cada historia tiene una prolongación, el límite inverso de sus niveles por supresión de emisiones es no vacío. En la topología cilíndrica, cada cilindro no vacío tiene medida positiva.

Esta medida distingue historias que convergen a una misma publicación. La medida uniforme sobre valores terminales sería una lectura distinta, pues eliminaría las multiplicidades que la construcción acaba de conservar.

\subsection{Transporte por hojas y contabilidad exacta de las salidas}
\label{iv:cont:transporte}
Consideremos primero el espacio con base ortonormal formada por las historias del nivel \(k\). El transporte que conserva la etiqueta de cada prolongación es

\[
\mathcal U_ke_h=
\sum_{\varepsilon\ \mathrm{admisible\ en}\ h}
\sqrt{p(\varepsilon\mid h)}\,e_{h\varepsilon}.
\]
Los conjuntos de hijos de padres distintos son disjuntos. La normalización de \(p\) demuestra, por tanto, \(\mathcal U_k^*\mathcal U_k=I\). La representación mediante masas de historias se relaciona con ésta por el cambio de base diagonal que incorpora \(\sqrt{\mu(h)}\).

La etiqueta de hoja forma parte del registro completo de una historia. Se escribe
\[
\widehat h=((h_0,s_0),\ldots,(h_k,s_k)),\qquad
s_j\in\{\Sigma,\Pi\},
\]
donde \(s_j\) es la hoja activa en el paso \(j\). La prolongación añade la emisión y la hoja \(s_{k+1}\) que determina el selector operativo del TPK, conservando íntegro el prefijo, incluida la hoja inicial. En la base de esas historias levantadas,
\[
\mathcal U_ke_{\widehat h}
=\sum_{\varepsilon\ \mathrm{admisible\ en}\ \widehat h}
\sqrt{p(\varepsilon\mid\widehat h)}\,e_{\widehat h\varepsilon},
\qquad
P_ke_{\widehat h}=\mathbf1_{s_k=\Sigma}e_{\widehat h},\quad
Q_k=I-P_k.
\]
Así, \(P_k\) lee la hoja actual, mientras la distinción entre historias se mantiene por su prefijo completo. Dos historias iniciales que lleguen a la misma hoja activa siguen teniendo hijos distintos. La prueba de isometría por familias disjuntas de prolongaciones se aplica literalmente a esta base.
La orientación trítica es otra coordenada del registro: el valor \(0\) conserva la lectura de umbral en cada hoja y no añade una tercera hoja a esta descomposición. Separamos ahora la conservación de hoja y el intercambio de hoja:

\[
A_k=P_{k+1}\mathcal U_kP_k+Q_{k+1}\mathcal U_kQ_k,
\qquad
\eta_k=Q_{k+1}\mathcal U_kP_k+P_{k+1}\mathcal U_kQ_k.
\]
\begin{proposition}[Identidad de Gram por hojas]
\label{iv:cont:gram}
Para un transporte lineal \(U\), sea o no isométrico, la descomposición anterior satisface

\[
A^*A+\eta^*\eta=PU^*UP+QU^*UQ.
\]
En particular, para \(U=\mathcal U_k\),

\[
A_k^*A_k+\eta_k^*\eta_k=I.
\]
\end{proposition}
\begin{proof}
En \(A^*A\) y \(\eta^*\eta\), los términos mixtos que contienen \(P_{k+1}Q_{k+1}\) se anulan. Al sumar los términos restantes aparece \(P_kU^*(P_{k+1}+Q_{k+1})UP_k\), y su análogo con \(Q_k\). La primera identidad resulta de la complementariedad de los proyectores; la segunda utiliza además \(U^*U=I\).
\end{proof}

La primera identidad es la fórmula general. Sustituir su lado derecho por \(U^*U\) requiere que los bloques cruzados de ese Gram se anulen. La conservación isométrica del transporte de historias proporciona aquí esa condición; no debe atribuirse automáticamente a una compresión posterior.

Escribamos \(F_{0,0}=I\), \(F_{0,n}=A_{n-1}\cdots A_0\) y \(T_N=F_{0,N}^*F_{0,N}\).

\begin{theorem}[Descomposición finita con residuo terminal]
\label{iv:cont:terminal}
Para cada horizonte \(N\),

\[
I=\sum_{n=0}^{N-1}
F_{0,n}^*\eta_n^*\eta_nF_{0,n}+T_N.
\]
La sucesión \(T_N\) es positiva y decreciente, y posee un límite fuerte positivo \(T_\infty\). La suma de salidas converge fuertemente a \(I-T_\infty\).

\end{theorem}
\begin{proof}
La proposición~\ref{iv:cont:gram} da \(T_n-T_{n+1}=F_{0,n}^*\eta_n^*\eta_nF_{0,n}\). La suma telescópica prueba la igualdad finita y la monotonía. Para cada vector, las formas cuadráticas decrecen a un límite; la polarización define una forma acotada positiva y, por el teorema de representación de formas acotadas, un operador \(T_\infty\). La convergencia es fuerte porque, para un operador positivo \(S\leq I\), \(\|Sv\|^2\leq\langle v,Sv\rangle\); se aplica a \(S=T_N-T_\infty\).
\end{proof}

El residuo terminal se conserva hasta demostrar su extinción en el dominio considerado. Esta distinción impide que un retorno de fase, por sí solo, se convierta en una hipótesis de agotamiento de todas las historias.

\subsection{Balance orientado y cancelación en el refinamiento}
\label{iv:cont:balance}
Una historia lleva su orientación acumulada \(o(h)\) y los recuentos de sus visitas a las dos hojas. El balance local correspondiente es

\[
b(h)=o(h)\bigl(N_\Sigma(h)-N_\Pi(h)\bigr),
\qquad
b^+=\frac{|b|+b}{2},\quad b^-=\frac{|b|-b}{2}.
\]
Cuando un lector publica la media sobre las prolongaciones, se define \(b_c=Eb_f\). La orientación sigue estando en el argumento de \(b_f\); no se elimina antes de promediar.

\begin{proposition}[Defecto exacto de cancelación]
\label{iv:cont:cancelacion}
La cantidad

\[
\kappa_{\mathrm{can}}=
\frac{E|b_f|-|Eb_f|}{2}
\]
es no negativa y satisface

\[
E(b_f^+)=(b_c)^++\kappa_{\mathrm{can}},
\qquad
E(b_f^-)=(b_c)^-+\kappa_{\mathrm{can}}.
\]
\end{proposition}
\begin{proof}
La desigualdad triangular ponderada da \(|Eb_f|\leq E|b_f|\). Al escribir las partes positiva y negativa mediante \((|b|\pm b)/2\), se obtienen ambas igualdades.
\end{proof}

La información omitida al publicar solamente el balance es la cancelación común de las contribuciones opuestas. Esta cantidad se determina desde las mismas prolongaciones utilizadas por el transporte anterior. El factor \(1/9\) aparece sólo cuando la fibra contiene nueve hijos con pesos iguales; la esperanza condicional es la regla que sigue siendo válida para ramificación y ponderaciones generales.

\subsection{Incidencia ternaria de las hojas y complejo rectangular}
\label{iv:cont:incidencia-rectangular}
Las dos hojas determinan el módulo \(V=\mathbb F_3e_\Sigma\oplus\mathbb F_3e_\Pi\). Sus cuatro direcciones proyectivas son las rectas generadas por \(e_\Sigma,e_\Pi,e_\Sigma+e_\Pi,e_\Sigma-e_\Pi\). De ellas proceden seis pares no ordenados. Los ocho vectores no nulos son sus representantes orientados. Así, los recuentos \(4,6,8\) pertenecen a operaciones distintas sobre un mismo módulo de hojas.

Sobre las seis posiciones de pares, definimos

\[
(Ba)_{\{L,L'\}}=a_L+a_{L'},
\qquad
s(q)=\sum_{\{L,L'\}}q_{\{L,L'\}}.
\]
Cada dirección interviene en tres pares; por ello \(sB=0\) en \(\mathbb F_3\). El lector \(W(q)=\mathbf1_{s(q)=0}-1/3\) es invariante bajo \(q\mapsto q+Ba\). Sobre el ambiente uniforme \(\mathbb F_3^6\), sus momentos son

\[
\mathbb EW=0,\qquad
\mathbb EW^2=\frac29,\qquad
\mathbb EW^3=\frac2{27}.
\]
En efecto, \(s\) es sobreyectiva y cada fibra contiene \(3^5\) elementos. El lector vale \(2/3\) con probabilidad \(1/3\) y \(-1/3\) con probabilidad \(2/3\). Estas medias describen el ambiente uniforme; la distribución inducida por un subconjunto de historias debe calcularse con sus propias multiplicidades, ya conservadas en \(X_k\).

Para fijar los puertos sin perder multiplicidades, a profundidad \(k\) se conservan las lecturas etiquetadas \(\lambda_\Sigma(h)=(h,\Sigma)\) y \(\lambda_\Pi(h)=(h,\Pi)\). Sus conjuntos son
\[
I_k=\{\lambda_\Sigma(h):h\in X_k\},\qquad
J_k=\{\lambda_\Pi(h):h\in X_k\}.
\]
Son finitos y no vacíos porque lo es \(X_k\). Los valores, cocientes y residuos se publican sobre esos puertos; una igualdad de valores no identifica sus etiquetas. El producto \(I_k\times J_k\) es el dominio de la lectura rectangular, dentro del cual se registra el soporte de las incidencias admisibles.

Fijado este nivel, escribimos \(I=I_k\), \(J=J_k\) y \(A_N=\mathbb Z/3^N\mathbb Z\). Elegimos puertos de referencia \(i_0,j_0\). Esta elección sólo proporciona coordenadas para las fórmulas de reconstrucción. Las aplicaciones

\[
\kappa(u,v)_{ij}=u_i-v_j,
\qquad
(\delta w)_{ij}=w_{ij}-w_{ij_0}-w_{i_0j}+w_{i_0j_0}
\]
definen la siguiente secuencia, válida sobre el anillo finito \(A_N\):

\[
0\longrightarrow A_N\longrightarrow
A_N^I\oplus A_N^J\mathrel{\mathop{\longrightarrow}^{\kappa}}
A_N^{I\times J}\mathrel{\mathop{\longrightarrow}^{\delta}}
A_N^{(I\setminus\{i_0\})\times(J\setminus\{j_0\})}
\longrightarrow0.
\]
\begin{proposition}[Exactitud de la incidencia rectangular]
\label{iv:cont:rectangular-exacta}
La secuencia anterior es exacta, con primera flecha dada por la diagonal constante.

\end{proposition}
\begin{proof}
La igualdad \(u_i-v_j=0\) para cada par obliga a que todas las coordenadas de \(u\) y \(v\) sean una misma constante. La sustitución directa da \(\delta\kappa=0\). Si \(\delta w=0\), tómense \(u_i=w_{ij_0}\) y \(v_j=w_{i_0j_0}-w_{i_0j}\). Entonces \(u_i-v_j=w_{ij}\). Finalmente, cualquier matriz del último término se prolonga con fila y columna de referencia nulas; su imagen por \(\delta\) es la matriz inicial. La prueba sólo utiliza sumas y restas, por lo que no exige dividir por una unidad del anillo.
\end{proof}

Reducir de \(A_{N+1}\) a \(A_N\), manteniendo fijos \(k,I_k,J_k,i_0,j_0\), conmuta con las dos aplicaciones. Las fórmulas de reconstrucción son las mismas en cada profundidad modular. Ésta es la compatibilidad aritmética demostrada en este apartado. La prolongación de historias \(k\to k+1\) conserva las etiquetas de prefijo; sus aplicaciones entre conjuntos de puertos son un dato distinto del cambio de anillo \(N\to N+1\).

Esas aplicaciones se obtienen de la supresión de emisiones:
\[
\rho_k^\Sigma(h\varepsilon,\Sigma)=(h,\Sigma),\qquad
\rho_k^\Pi(h\varepsilon,\Pi)=(h,\Pi).
\]
Se fija una historia de referencia con prolongaciones compatibles, cuya existencia se demostró para el árbol sin hojas terminales. Sus dos etiquetas definen puertos de referencia compatibles en todos los niveles. El transporte de coordenadas es la precomposición:
\[
(L_ku)_{i'}=u_{\rho_k^\Sigma(i')},\quad
(L_kv)_{j'}=v_{\rho_k^\Pi(j')},\quad
(L_kw)_{i'j'}=w_{\rho_k^\Sigma(i'),\rho_k^\Pi(j')}.
\]
Para el último módulo de la secuencia, se prolonga primero la matriz por cero en la fila y la columna de referencia, se aplica esta misma fórmula y se restringe al nuevo complemento de los puertos de referencia. Esta regla define \(L_k\) también cuando un puerto nuevo tiene por padre el puerto marcado.

\begin{proposition}[Naturalidad del complejo rectangular bajo prolongación]
\label{iv:rectangular-natural}
Las aplicaciones anteriores verifican
\(L_k\kappa_k=\kappa_{k+1}L_k\) y
\(L_k\delta_k=\delta_{k+1}L_k\).
Además conmutan con la reducción de coeficientes
\(A_{N+1}\to A_N\) y se componen conforme a la supresión de prefijos.
\end{proposition}
\begin{proof}
La primera identidad se obtiene de
\(u_{\rho_k^\Sigma(i')}-v_{\rho_k^\Pi(j')}\).
En la segunda, cada uno de los cuatro términos que define \(\delta\)
se transporta por la misma precomposición. La compatibilidad de los
puertos marcados identifica sus filas y columnas; cuando un padre es
el puerto marcado, los cuatro términos se cancelan y coinciden con
la prolongación por cero. Todas las fórmulas son sumas y restas con
coeficientes enteros, de modo que conmutan con la reducción modular.
La composición de precomposiciones es la precomposición con el
prefijo compuesto, lo que demuestra la última afirmación.
\end{proof}

\subsection{Complejo de memoria y compatibilidad de las incidencias}
\label{iv:cont:memoria}
La misma coordenada acumulada de memoria \(c_y\in\mathbb Z\), en una residencia \(y\), define un complejo libre entero. Su reducción a \(A_N\) utiliza el mismo registro entero antes de reducirlo; por tanto, las profundidades sucesivas son compatibles. Sus generadores son \(f_y\) en grado \(2\), \(e_y,b_y,r_y\) en grado \(1\), y \(g_y\) en grado \(0\):

\[
\partial f_y=3c_y e_y+b_y,\qquad
\partial e_y=g_y,\qquad
\partial b_y=-3c_y g_y,\qquad
\partial r_y=0.
\]
La identidad \(\partial^2=0\) se verifica sobre \(f_y\) mediante \(3c_yg_y-3c_yg_y=0\), y es inmediata en los otros generadores.

Si una trayectoria \(\alpha:y\to y'\) actualiza la memoria, su transporte en el complejo envía los generadores homónimos a sus nuevos representantes, salvo

\[
\Psi_\alpha(b_y)=b_{y'}+3(c_{y'}-c_y)e_{y'}.
\]
\begin{proposition}[Naturalidad del complejo de memoria]
\label{iv:cont:memoria-natural}
Se cumplen \(\partial\Psi_\alpha=\Psi_\alpha\partial\) y \(\Psi_{\beta\alpha}=\Psi_\beta\Psi_\alpha\). Las aplicaciones se prolongan linealmente sobre el anillo de coeficientes elegido.

\end{proposition}
\begin{proof}
Sobre \(f_y\), la imagen de su borde es \(3c_ye_{y'}+b_{y'}+3(c_{y'}-c_y)e_{y'} =3c_{y'}e_{y'}+b_{y'}\). Sobre \(b_y\), el borde de su imagen es \(-3c_{y'}g_{y'}+3(c_{y'}-c_y)g_{y'}=-3c_yg_{y'}\). Los otros dos casos son inmediatos. En una concatenación, las diferencias de memoria se suman telescópicamente, lo que prueba la segunda identidad.
\end{proof}

La relación entre dos representantes de una misma residencia también puede escribirse como un borde explícito. Para una coordenada entera \(v\), reducida en el mismo anillo cuando corresponda, sean

\[
z^-=r_y,\qquad z^+=r_y+3c_yv e_y+v b_y,
\qquad h_*=-vf_y.
\]
Entonces \(\partial h_*=z^--z^+\). Bajo un transporte que actualiza también \(v\), la corrección \(K_\alpha=(v_{y'}-v_y)f_{y'}\) satisface \(K_{\beta\alpha}=K_\beta+\Psi_\beta K_\alpha\). Los símbolos \(K_\alpha\) de esta homotopía local no designan el sello dodecafásico \(K\): pertenecen a otro dominio y su función está especificada por la identidad anterior. Su función puede comprobarse también en la igualdad

\[
z^+_{y'}-\Psi_\alpha z^+_y=\partial K_\alpha.
\]
Ambos miembros valen \((v_{y'}-v_y)(3c_{y'}e_{y'}+b_{y'})\).

\begin{proposition}[Homología y memoria del complejo local]
\label{iv:cont:memoria-homologia}
Sobre \(\mathbb Z\), o después de reducir a \(A_N\), se tiene

\[
H_0(\Gamma)=H_2(\Gamma)=0,
\qquad H_1(\Gamma)\cong A_N[r_y]
\]
en la realización modular, y \(H_1(\Gamma)\cong\mathbb Z[r_y]\) en la realización entera.

\end{proposition}
\begin{proof}
Defínase la homotopía de grado \(1\) mediante \(h(g_y)=e_y\), \(h(b_y)=f_y\), y cero en los otros generadores. Sean \(p\) la proyección al sumando generado por \(r_y\) e \(\iota\) su inclusión. La comprobación sobre cada generador da

\[
\partial h+h\partial=I-\iota p.
\]
En particular, sobre \(b_y\) los términos \(3c_ye_y\) y \(-3c_ye_y\) se cancelan. El complejo se retrae por homotopía al módulo generado por \(r_y\), concentrado en grado \(1\), lo que demuestra la afirmación.
\end{proof}

El transporte \(\Psi_\alpha\) es una cizalla de determinante \(1\). Así, este complejo conserva explícitamente cómo cambia \(c_y\) a nivel de cadenas, pero su homología local no distingue los valores de esa coordenada. La memoria completa permanece en las etiquetas y en los transportes del estado fuente; no se recupera a partir de \(H_1\) aisladamente. Una torsión numérica dependiente de memoria requiere el ensamblaje y la lectura determinantal especificados, además de este complejo local.

\subsection{Ensamblaje con caminos terminales}
\label{iv:cont:ensamblaje}
Para un horizonte finito, las residencias y sus prolongaciones forman una categoría de caminos. A cada residencia \(\nu\) se asigna el módulo libre \(W(\nu)=A_N[\operatorname{Path}(\nu,\mathrm{Term})]\), que conserva cada continuación terminal como generador. Una trayectoria inicial actúa en \(W\) por precomposición. El complejo \(\Gamma(\nu)\) del apartado anterior actúa covariantemente mediante \(\Psi\).

El ensamblaje utiliza ambos transportes en el mismo módulo bigraduado:

\[
\mathcal B_{q,r}=
\bigoplus_{\nu_0\to\cdots\to\nu_q}
W(\nu_q)\otimes_{A_N}\Gamma_r(\nu_0).
\]
Se utiliza el complejo de caminos normalizado: se omiten las cadenas degeneradas que contienen identidades. El horizonte finito acota entonces el número de prolongaciones estrictas y el grado horizontal. El borde horizontal suprime una residencia intermedia componiendo las flechas adyacentes. En el primer extremo transporta \(\Gamma\); en el último precompone \(W\). Con los signos alternados se obtiene el borde de caminos \(d_{\mathrm{bar}}\). El diferencial total es

\[
D=d_{\mathrm{bar}}+(-1)^q\partial_\Gamma.
\]
Las identidades de caras cancelan por pares los términos de \(d_{\mathrm{bar}}^2\). La naturalidad de la proposición~\ref{iv:cont:memoria-natural} implica que \(d_{\mathrm{bar}}\) conmuta con \(\partial_\Gamma\); el factor \((-1)^q\) cancela los términos cruzados del cuadrado. Por tanto, \(D^2=0\). La compatibilidad se demuestra sobre generadores de caminos y no se introduce por una etiqueta de pertenencia a un complejo.

La evaluación terminal tampoco requiere elegir una historia particular. Para cada continuación \(h\), se aplica la composición efectiva de sus actualizaciones. El resultado conserva el par formado por \(h\) y su lectura, con su multiplicidad. Si \(k<\ell<N\), las continuaciones se descomponen como una unión disjunta de prefijos hasta \(\ell\) y sus continuaciones hasta \(N\). La proposición~\ref{iv:cont:concatenacion} demuestra que evaluar directamente o evaluar esas dos partes produce la misma lectura. Así se obtiene el cuadrado de naturalidad del terminal.

Los módulos de caminos, las incidencias, el transporte de memoria y sus diferenciales proceden de las mismas prolongaciones. Su ensamblaje no consiste en añadir a un estado cinco coordenadas que se conservarían por definición: sus relaciones se construyen mediante sumas, bordes y composiciones efectivas, y se verifican sobre sus generadores.

\subsection{Determinantes, límite y lectura aritmética}
\label{iv:cont:determinantes}
En cada residencia, el complejo local libre \(\Gamma\) admite su línea determinante graduada

\[
\operatorname{Det}\Gamma=
\bigotimes_r\left(\bigwedge^{\mathrm{rk}\Gamma_r}\Gamma_r\right)^{(-1)^r}.
\]
El ensamblaje total del apartado anterior tiene, por su parte, los módulos \(\operatorname{Tot}_d\mathcal B=\bigoplus_{q+r=d}\mathcal B_{q,r}\). A horizonte finito y sobre el mismo anillo, son libres de rango finito y sólo un número finito de ellos es no nulo. Su línea es
\[
\operatorname{Det}(\operatorname{Tot}\mathcal B)=
\bigotimes_{q,r}
\left(\bigwedge^{\operatorname{rk}\mathcal B_{q,r}}
\mathcal B_{q,r}\right)^{(-1)^{q+r}}.
\]
Esta fórmula resulta de la descomposición por suma directa en cada grado. Distingue la línea del complejo local y la del ensamblaje con todos los caminos; sus bases se ordenan mediante las etiquetas conservadas. La reducción modular conmuta con estas potencias exteriores porque los módulos y sus bases proceden de los mismos módulos libres enteros a horizonte fijado.
La línea está definida por los módulos y sus grados, tanto sobre \(\mathbb Z\) como después de reducir a \(A_N\); el exponente negativo designa el módulo dual. Una torsión numérica no nula exige, adicionalmente, especificar el complejo sobre un cuerpo, las bases y las identificaciones de homología pertinentes. Para calcular factores de Smith se usa la matriz entera producida antes de la reducción, no un levantamiento arbitrario de una matriz modular. Un bloque cuadrado entero de determinante no nulo es invertible sobre \(\mathbb Q\); su reducción es invertible sobre \(A_N\) solamente cuando ese determinante es una unidad, es decir, cuando no es divisible por \(3\). El refinamiento de memorias e incidencias permite formular y comprobar la estabilidad de los factores invariantes pertinentes; estas condiciones no se deducen de la existencia de la línea determinante.

El paso al límite conserva los proyectivos compatibles de residuos, las historias y sus multiplicidades. Los cuadrados probados en los apartados anteriores permanecen conmutativos porque sus aplicaciones están dadas por las mismas fórmulas en cada nivel. Para una coordenada arquimediana, el paso adicional consiste en publicar intervalos cilíndricos no vacíos, compatibles por inclusión y de diámetro tendente a cero. Sus clausuras compactas anidadas tienen una intersección de un solo punto: la existencia se obtiene por compacidad y la unicidad por la cota del diámetro. Ése es el valor de la lectura límite. Si los intervalos son semiabiertos, el punto puede pertenecer sólo a sus clausuras; la regla de representación de frontera determina entonces su escritura posicional. El valor no sustituye la historia que lo ha producido. La otra publicación conserva las incidencias del mismo estado. Los dos destinos comparten su origen sin recibir por ello la misma información.

En los capítulos siguientes, las regiones coinductivas, el registro dodecafásico y la incidencia excepcional se publican desde ese sustrato con sus respectivos mapas. La realización hilbertiana conserva después las etiquetas de fase y especifica su propia inclusión unital; la dualidad y las pantallas actúan sobre los codominios que allí se construyen. Las operaciones conjuntas aquí reunidas transmiten historias, hojas, orientación y memoria a esas publicaciones. La identificación de un funcional concreto con una forma analítica externa requiere, a su vez, escribir su aplicación sobre los generadores y verificar sus productos escalares: las identidades de conservación del presente capítulo no reemplazan esa composición.

\subsection{Procedencia, comprobación y alcance}
\label{iv:cont:mapa-correlativo}
% RESULTADO_RECUPERADO: 04b2_operaciones_intrinsecas_correlativas.tex,
% eq:pdfv2-cor-accion-sp/sol/gau/coh/det y eq:pdfv2-cor-memorias.
% Las marcas internas sp/sol/gau/coh/det corresponden, en ese orden,
% a las cinco filas siguientes; no son cinco dominios independientes.
Las cinco operaciones se reconocen por su acción sobre las mismas
historias y por las relaciones demostradas en los apartados anteriores.
El transporte afín \(U_\alpha\) de la memoria y el transporte
lineal \(\mathcal U_k\) de las historias tienen los dominios
respectivos de las secciones~\ref{iv:cont:historias}
y~\ref{iv:cont:transporte}. La graduación de hojas es
\(\sigma_k=P_k-Q_k\). Para una historia
\(h_j=(\varepsilon_1,\ldots,\varepsilon_j)\), la acumulación del
balance queda explícita mediante
\[
 M_0^\pm=0,\qquad
 M_k^\pm(h_k)=\sum_{j=1}^k b^\pm(h_j),\qquad
 M_{k+1}^\pm(h_k\varepsilon)
 =M_k^\pm(h_k)+b^\pm(h_k\varepsilon).
\]
Cada sumando es la lectura firmada del prefijo correspondiente,
definida en la sección~\ref{iv:cont:balance}. Separar el último
sumando demuestra la actualización; las identidades
\(M_k^+-M_k^-=\sum_{j=1}^kb(h_j)\) y
\(M_k^++M_k^-=\sum_{j=1}^k|b(h_j)|\) se deducen de las
definiciones de las partes positiva y negativa. La acumulación y la
cancelación condicional son, así, operaciones distintas sobre las
mismas contribuciones.

\begin{center}
\small
\begin{tabular}{@{}p{.21\linewidth}p{.48\linewidth}p{.23\linewidth}@{}}
\toprule
Operación conjunta & Objetos y acción efectiva & Desarrollo y prueba\\
\midrule
Transporte y graduación de hojas
& \(\mathcal U_k:\ell^2(\widehat X_k)\to
\ell^2(\widehat X_{k+1})\), con \(\widehat X_k\) formado por
las historias levantadas; \(\sigma_k=P_k-Q_k\), \(A_k\) y
\(\eta_k\) conservan o intercambian la hoja actual.
& Sección~\ref{iv:cont:transporte}; proposición~\ref{iv:cont:gram}
y teorema~\ref{iv:cont:terminal}.\\[5pt]
Balance, cancelación y memoria
& \(b:X_k\to\ZZ\), \(b^\pm:X_k\to\mathbb N_0\);
acumulación \(M_k^\pm\) sobre prefijos y cancelación
\(\kappa_{\mathrm{can}}=(E|b_f|-|Eb_f|)/2\) al publicar
la media condicional.
& Sección~\ref{iv:cont:balance};
proposición~\ref{iv:cont:cancelacion}.\\[5pt]
Incidencia ternaria de las hojas
& \(B:\FF_3^4\to\FF_3^6\),
\((Ba)_{\{L,L'\}}=a_L+a_{L'}\); se conserva el ambiente
completo de \(q\), con \(sB=0\) y \(W(q+Ba)=W(q)\).
& Primera parte de la
sección~\ref{iv:cont:incidencia-rectangular}.\\[5pt]
Puertos y complejo rectangular
& \(\kappa:A_N^{I_k}\oplus A_N^{J_k}\to
A_N^{I_k\times J_k}\), \(\kappa(u,v)_{ij}=u_i-v_j\);
la diferencia \(\delta\) y la precomposición \(L_k\)
conservan las etiquetas y sus multiplicidades.
& Proposiciones~\ref{iv:cont:rectangular-exacta}
y~\ref{iv:rectangular-natural}.\\[5pt]
Memoria celular, relleno y línea determinante
& \(\Psi_\alpha:\Gamma_y\to\Gamma_{y'}\),
\(K_\alpha\in\Gamma_{y',2}\),
\(\partial h_*=z^--z^+\); las líneas
\(\operatorname{Det}\Gamma\) y
\(\operatorname{Det}(\operatorname{Tot}\mathcal B)\)
se forman sobre los complejos basados correspondientes.
& Secciones~\ref{iv:cont:memoria}--\ref{iv:cont:determinantes};
proposiciones~\ref{iv:cont:memoria-natural}
y~\ref{iv:cont:memoria-homologia}.\\
\bottomrule
\end{tabular}
\end{center}

El cuadro identifica cinco acciones de una misma construcción: cada
prefijo conserva sus hojas, orientación, residuos, cocientes, memoria
y prolongaciones terminales. La precomposición de puertos lleva
coordenadas de un nivel al refinado; la reducción de coeficientes
lleva \(A_{N+1}\) a \(A_N\). Sus sentidos y dominios son
distintos y sus cuadrados conmutativos están explicitados en la
proposición~\ref{iv:rectangular-natural}. El ensamblaje de la
sección~\ref{iv:cont:ensamblaje} utiliza simultáneamente las
continuaciones terminales y el transporte del complejo de memoria;
sus fórmulas de composición reúnen las acciones del cuadro sobre el
mismo sistema de historias. La lectura de una línea determinante
como escalar conserva las condiciones de bases, homología y
no singularidad de la sección~\ref{iv:cont:determinantes}.

Este desarrollo reúne las construcciones de estado, árbol y medida común; las operaciones intrínsecas correlativas; y el ensamblaje terminal con naturalidad y límite demostrados en esta sección. Se incorporan sus identidades de balance, incidencia y memoria, con demostraciones dentro de este capítulo. Las cinco construcciones del continuo se mantienen como operaciones sobre un mismo árbol. La exposición amplía sus pruebas locales y explicita las condiciones de cada especialización: transporte isométrico, pesos uniformes, extinción del terminal y no singularidad de un bloque determinantal.

El comprobador adjunto contrasta instancias exactas de los cociclos, la cancelación ponderada, la secuencia rectangular y el complejo de memoria. Esos controles acompañan las demostraciones generales del texto y verifican los operadores aquí definidos en sus dominios respectivos. El paso posterior hacia las representaciones excepcionales y la realización hilbertiana conserva sus aplicaciones sobre generadores: la sección~\ref{iv:hilbert} especifica la fibra de fase, los operadores de Weyl y las inclusiones que determinan sus clausuras. La conservación documental y los controles finitos mantienen su función de trazabilidad y contraste de esas pruebas.

\HMTDivision{\(\pi\quad\varphi\quad e\)}{Generación coinductiva}{Las coordenadas regionales se construyen en prolongaciones compatibles de profundidad arbitraria. El estado conserva las relaciones que cada lectura utiliza.}
\HMTNarrativeHeading{La profundidad de un número}

La generación coinductiva organiza la determinación de un número mediante
prolongaciones compatibles. Una etapa produce un prefijo y conserva las
condiciones para continuarlo. La profundidad puede aumentar manteniendo
la concordancia de todas las lecturas anteriores. El valor aparece al
realizar esa compatibilidad; la historia conserva, además, el modo en que
la precisión fue obtenida.

Las regiones de las coordenadas circular, áurea y exponencial proceden del
mismo desarrollo APP--TRIT--TPK. Sus lectores cumplen funciones
diferenciadas y comparten la estructura que permite relacionarlos. La
extensión triádica de Euler explicita las realizaciones de los regímenes;
el calendario areal--volumétrico determina la autoescala; las emisiones
regionales y sus fronteras conservan las condiciones de evaluación.
La sucesión de estos pasos hace explícito el orden de construcción de las
coordenadas que intervienen después en alfa.

El calendario de vacancias permite comprender una particularidad del
refinamiento. La historia avanza en cada transición, mientras la capacidad
de una lectura puede permanecer fija. La palabra que registra esas
transiciones distingue la capacidad adquirida de la vacancia conservada.
Su balance cuenta ambas contribuciones. La regularidad del calendario y
la aperiodicidad de su itinerario describen aspectos relacionados de una
misma evolución.

La monodromía prolonga esta distinción. Volver a una fase significa
reconocer una posición del ciclo; recuperar una historia significa
restituir también sus vueltas, orientación y registros. El espejo regional
conserva el eje y el agregado de las dos coordenadas que intercambia,
mientras invierte su diferencia. Este dato permite reconstruir el
reparto. La figura nonádica sitúa las fases y sus operaciones; sus marcas
son posiciones del calendario, con una semántica diferente del valor real
de una constante.

La quinta frontera, las regiones posteriores y el retorno se explican,
por tanto, mediante sus acciones sobre el estado. Las figuras regionales
muestran simultáneamente coincidencia de una lectura y diferenciación de
sus antecedentes. Esa diferenciación transmite a la realización
excepcional la información incidencial que utiliza. La generación
numérica continúa así la doble proyección planteada al comienzo: cada
lectura obtiene su valor de una estructura cuya organización permanece
disponible para las otras realizaciones.

\HMTMathematicalPaper
\section{Generación coinductiva de \texorpdfstring{\(\pi,\varphi,e\)}{pi, phi, e} a profundidad arbitraria}
\label{sec:generacion}

La emisión de seis bloques es una etapa finita de la construcción, no una expansión decimal completa. Su función consiste en producir regiones con multiplicidad, orientación y reglas de transporte. La prolongación debe conservar esos datos y determinar una familia compatible de prefijos de longitud arbitraria. La palabra \emph{generación} designa esta operación; \emph{genealogía} designa la reconstrucción de sus dependencias. Ambas nociones intervienen, pero cumplen funciones distintas.

\subsection{Órbitas, regiones y caracteres de clausura, propagación y autoescala}

Sobre una palabra \(w=(w_1,\ldots,w_6)\), definimos
\[
 r(w)=(w_3,w_1,w_2,w_5,w_6,w_4),\qquad
 s(w)=(w_4,w_5,w_6,w_1,w_2,w_3).
\]
Se cumple \(r^3=s^2=I\) y \(srs=r^{-1}\); por tanto estas permutaciones realizan el grupo diedral \(D_3\). La acción conserva el catálogo y las multiplicidades y conmuta con \(\Psi\). El cociente de las \(243\) palabras visibles tiene \(43\) órbitas: treinta y ocho de cardinal seis y cinco de cardinal tres.

La clase de clausura se reconoce por un predicado de fibra: sus seis orientaciones tienen cinco emisiones por palabra, multiplicidad total \(1008\) y distribución \(432+4\cdot144\). Existe una sola órbita con esas propiedades en el catálogo. Para precisar las otras dos selecciones, escribimos \(f(w)=\#\Psi^{-1}(w)\), \(m(w)=\sum_{U\in\Psi^{-1}(w)}\operatorname{mult}(U)\) y \(\nu(w)=(n_0,n_1,n_2)\). La clase diagonal aditiva de una emisión es \(d(U)=[i+j]_9\) en los índices de cero a ocho; su firma aditiva conserva esta clase y la orientación \(ES\) o \(NO\). Sea \(\mathcal D(\mathcal O)\) el conjunto de esas clases para todas las emisiones sobre una órbita.

Entre las órbitas de tamaño seis, el predicado conjunto
\[
 f(w)=1,\qquad m(w)=144\quad(w\in\mathcal O),\qquad
 \mathcal D(\mathcal O)=\{0,3,6\}
\]
deja exactamente seis órbitas. En este sector, el censo \(\nu=(2,3,1)\), igual al de clausura, selecciona únicamente la propagación; el censo \(\nu=(2,2,2)\) selecciona únicamente la autoescala. El soporte diagonal no puede omitirse: sin él hay cuatro órbitas unipuntuales equilibradas de multiplicidad \(144\), no una. Finalmente, la existencia de una emisión con \(d=6\) y orientación \(ES\) selecciona un representante de cada órbita y conduce a
\begin{equation}
 w^{\rm cl}=010211,\qquad w^{\rm pr}=201101,\qquad w^{\rm au}=121200.
 \label{eq:palabras-regionales}
\end{equation}
Estos símbolos designan primero regiones del catálogo. La identificación escalar se demostrará mediante sus lectores; no se eligen porque coincidan con cifras de una constante conocida.

La palabra de clausura \(010211\) tiene las cinco preimágenes siguientes:
\begin{center}\small
\begin{tabular}{@{}lll r@{}}\toprule
Emisión de seis bloques&Firma multiplicativa&Papel&Multiplicidad\\\midrule
\(501,614,498,169,272,272\)&\(555555\)&transversal&432\\
\(810,923,870,169,272,272\)&\(339555\)&orientación positiva&144\\
\(810,983,810,169,272,272\)&\(393555\)&orientación positiva&144\\
\(870,923,810,169,272,272\)&\(933555\)&orientación positiva&144\\
\(870,923,810,418,674,521\)&\(933717\)&retorno orientado&144\\\bottomrule
\end{tabular}
\end{center}
La multiplicidad \(432\) y la firma constante \(555555\) caracterizan la región transversal. No identifican la posición central \((5,5)\) de APP. Confundir una firma de seis ventanas con una posición del soporte eliminaría la distinción entre región numérica y estructura electrónica.

La microfibra posee además una morfología bipartita precisa. Cada emisión es un producto de una clase de cursor aditivo y una clase multiplicativa. Las tres regiones positivas comparten la misma clase aditiva de dieciocho condiciones; juntas contienen tres clases multiplicativas de ocho condiciones cada una. La región transversal posee un producto \(18\times24\); las tres positivas reunidas, otro \(18\times24\); y el retorno, un producto \(18\times8\). La descomposición por emisiones tiene cinco partes, mientras la descomposición por componentes conexas del grafo bipartito tiene tres. Ambas conservan las mismas \(1008\) semillas sin reducir la forma a un cardinal.

\subsection{Transporte finito y frontera de prolongación}

Los endomorfismos \(L_0,L_1\) del espacio de palabras ternarias actúan antes de la realización arquimediana. Su determinación comienza por las transiciones del estado aritmético, no por la prescripción de dos matrices. Si \(U_t\) es el paso compuesto de selección, transporte y actualización, el bloque observado satisface
\[
 b_j^\chi=\Pi_6(x_j^\chi),\qquad
 x_{j+1}^\chi=U_{9j+9}\cdots U_{9j+1}(x_j^\chi),
 \qquad \chi\in\{\mathrm{cl},\mathrm{pr},\mathrm{au}\}.
\]
Se forman seis transiciones por cada régimen: dos consecutivas en cada una de las tres regiones. Los pares de matrices de orígenes y destinos obtenidos en el registro son
\[
\begin{array}{c|c|c|c}
X_0&Y_0&X_1&Y_1\\\hline
010211&012222&010211&002111\\
012222&010211&002111&110221\\
201101&121221&102011&012222\\
121221&102011&012222&102011\\
121200&112202&121020&010210\\
112202&121020&010210&010200
\end{array}
\]
Cada cadena de seis dígitos representa una fila en \(\FF_3^6\). Puesto que \(\det X_0=1\) y \(\det X_1=-1\) en \(\FF_3\), las ecuaciones \(X_aL_a=Y_a\) determinan unívocamente \(L_a=X_a^{-1}Y_a\). En esa carta de vectores fila resultan
\[
 L_0=\begin{pmatrix}
 2&2&2&1&2&1\\2&1&2&2&1&1\\1&0&0&1&0&2\\
 0&0&1&1&1&0\\2&2&2&0&2&1\\2&1&2&1&0&0
 \end{pmatrix},\quad
 L_1=\begin{pmatrix}
 2&1&2&0&2&2\\1&2&1&1&2&0\\1&2&1&1&1&2\\
 0&1&0&0&2&1\\2&0&2&1&0&1\\0&2&2&2&1&1
 \end{pmatrix}\quad(\bmod3).
\]
El calendario \((L_0,L_0,L_1,L_1)\) produce cuatro nuevos bloques de seis componentes. Sus concatenaciones forman
\[
 w_6\longrightarrow w_{12}\longrightarrow w_{18}
 \longrightarrow w_{24}\longrightarrow w_{30}.
\]
Los subíndices de \(w\) indican longitud acumulada. La frontera \(R_{36}\) conserva la estructura terminal y abre la relación de prolongación; no se renombra como un último bloque periódico que pudiera repetirse indefinidamente.

La identidad \(L=X^{-1}Y\) demuestra unicidad una vez fijadas las transiciones; no sustituye la producción de esas transiciones por \(U_t\). La procedencia de su registro y el examen de los dieciséis calendarios binarios se documentan en el suplemento técnico. El calendario indicado es el único de esos dieciséis que reproduce conjuntamente las tres sucesiones de cinco bloques.

\subsection{Cinco regiones de clausura y compensación angular}

La región de clausura contiene una componente transversal y cuatro acompañantes. El lector simétrico actúa en el espacio de sus cinco posiciones mediante
\[
 P_5=\frac15\mathbf1\mathbf1^{\mathsf T},\qquad
 \mathbf1=(1,1,1,1,1)^{\mathsf T}.
\]
La invariancia \(P_5\lambda=\lambda\) y la conservación \(\mathbf1^{\mathsf T}\lambda=1\) fuerzan \(\lambda=\mathbf1/5\). Esta normalización distribuye la unidad entre posiciones del lector; no sustituye las multiplicidades de semillas \(432,144,144,144,144\) por frecuencias uniformes. La carta elíptica del lector realiza cada coordenada mediante \(1+\lambda_jJ\); su pendiente es \(\lambda_j\), por lo que la pendiente primitiva resulta \(1/5\). Se hace explícita así la regla que pasa de una coordenada normalizada a una pendiente, en lugar de identificar ambas nociones sin mapa. El lector compone las cuatro posiciones acompañantes por la ley
\begin{equation}
 a\oplus b=\frac{a+b}{1-ab}.
 \label{eq:suma-pendientes}
\end{equation}
Esta ley se obtiene al multiplicar \((1+aJ)(1+bJ)\) en el régimen \(J^2=-I\) y dividir por su componente escalar. Por ello la composición de pendientes es una realización de la estructura cuadrática de TRIT, no una suma ordinaria de cuatro fracciones.

\begin{proposition}[Compensador de la clausura]
La composición de cuatro pendientes \(1/5\) es \(120/119\). La condición de retorno a pendiente uno determina un único compensador positivo de la forma \(1/q\), con \(q>1\), y da \(q=239\).
\end{proposition}
\begin{proof}
La primera duplicación da \((1/5)\oplus(1/5)=5/12\); la segunda, \((5/12)\oplus(5/12)=120/119\). La condición
\[
 \frac{120/119-1/q}{1+(120/119)/q}=1
\]
equivale a \((120-119)q=120+119\), por lo que \(q=239\).
\end{proof}

El carácter de clausura resultante tiene realización
\begin{equation}
 \Pi_{\rm cl}=4\left(4\arctan\frac15-\arctan\frac1{239}\right).
 \label{eq:lector-pi}
\end{equation}
La identidad angular de Machin reconoce \(\Pi_{\rm cl}=\pi\). El sentido de la construcción es preciso: la pendiente y el número de composiciones están declarados por el lector de la pentafibra; la ecuación de compensación fuerza \(239\). Una vez producido este carácter, su evaluación por series alternadas no selecciona de nuevo la órbita de semillas.

Para \(0<x<1\), las sumas alternadas
\[
 A_n(x)=\sum_{j=0}^n\frac{(-1)^jx^{2j+1}}{2j+1}
\]
encierran \(\arctan x\) entre dos truncamientos consecutivos, con diferencia \(x^{2n+3}/(2n+3)\). Aplicadas a las dos pendientes de \eqref{eq:lector-pi}, proporcionan extremos racionales de error explícito. El valor numérico se calcula así como consecuencia del carácter exacto, sin introducir un prefijo objetivo.

\subsection{Lector de propagación y normalización de la unidad}

La región de propagación se realiza mediante una familia formal \(P_t\in\mathcal A[[t]]\), donde \(\mathcal A\) es un álgebra conmutativa unitaria sobre \(\QQ\). Su composición satisface \(P_{t+s}=P_tP_s\), con unidad \(P_0=1\), y el transporte elemental orientado fija el generador escalar \(P'_0=+1\). La condición fija su signo, no sólo su magnitud. Escribiendo
\[
 P_t=\sum_{n\ge0}a_nt^n,
\]
la normalización determina \(a_0=a_1=1\). La comparación del coeficiente de \(s^nt\) en \(P_{s+t}=P_sP_t\) da \((n+1)a_{n+1}=a_na_1\); equivalentemente, \(P'_t=P_t\). Así,
\begin{equation}
 (n+1)a_{n+1}=a_n,\qquad a_n=\frac1{n!}.
 \label{eq:propagacion}
\end{equation}
El carácter en una unidad de parámetro es \(\Pi_{\rm pr}=P_1\), reconocido posteriormente como \(e\).

La normalización del generador importa. La ley semigrupal sin ese dato admitiría \(P_t=e^{ct}\) para distintos \(c\); por ello no se omite la condición que fija la unidad del parámetro en el lector. Esta condición es anterior a la comparación decimal y pertenece a la exposición de la regla operacional, no a la elección de un valor experimental.

\begin{proposition}[Cotas racionales de propagación]
Para \(n\ge1\), sea \(E_n=\sum_{j=0}^n1/j!\). Entonces
\[
 E_n<\Pi_{\rm pr}<E_n+\frac{n+2}{(n+1)(n+1)!}.
\]
\end{proposition}
\begin{proof}
El primer término omitido es \(1/(n+1)!\). Cada cociente posterior es a lo sumo \(1/(n+2)\). La suma geométrica mayorante de la cola es
\[
 \frac1{(n+1)!}\frac1{1-1/(n+2)}
 =\frac{n+2}{(n+1)(n+1)!}.
\]
La desigualdad estricta se sigue de la disminución de los cocientes sucesivos.
\end{proof}

\subsection{Incidencia modal y generación de la autoescala}

El selector trítico produce el ciclo \((+1,-1,0)\). En los modos \(\Sigma,\Pi\), la regla es \(+1:m\mapsto\Sigma\), \(-1:m\mapsto\Pi\), \(0:m\mapsto m\). Con cierre cíclico, el modo observado es \((\Sigma,\Pi,\Pi)\). Sus aristas son \(\Sigma\to\Pi\), \(\Pi\to\Pi\) y \(\Pi\to\Sigma\). Por tanto su incidencia, en ese orden de modos, es
\begin{equation}
 F_{\rm av}=\begin{pmatrix}0&1\\1&1\end{pmatrix}.
 \label{eq:fav}
\end{equation}
Las cuatro entradas se obtienen del calendario, no de una ecuación que haya recibido previamente el valor de \(\varphi\). La repetición del calendario da conteos \(3F_{\rm av},9F_{\rm av},36F_{\rm av}\) en longitudes nueve, veintisiete y ciento ocho. Estos conteos no son las potencias de la matriz.

\begin{theorem}[Carácter positivo de autoescala]
La incidencia \eqref{eq:fav} posee un único rayo propio estrictamente positivo. Al normalizarlo como \((1,x)^{\mathsf T}\), su coordenada \(x\) satisface \(x^2-x-1=0\), \(x>0\), y es \(\varphi=(1+\sqrt5)/2\).
\end{theorem}
\begin{proof}
La ecuación propia da \(x=\lambda\) y \(1+x=\lambda x\). Eliminar \(\lambda\) produce el polinomio. De sus dos raíces sólo una es positiva. Como \(F_{\rm av}^2\) tiene todas sus entradas positivas, el rayo positivo es único y simple. La expresión radical reconoce la coordenada obtenida; no participa en la producción de \(F_{\rm av}\).
\end{proof}

Con \(F_0=0,F_1=1,F_{n+1}=F_n+F_{n-1}\), la incidencia induce, para \(n\geq1\),
\[
 F_{\rm av}^n=\begin{pmatrix}F_{n-1}&F_n\\F_n&F_{n+1}\end{pmatrix}.
\]
Los cocientes consecutivos \(F_{n+1}/F_n\) alternan alrededor de \(\varphi\). La identidad de determinante
\[
 F_{n+1}^2-F_nF_{n+2}=(-1)^n
\]
da una anchura exacta \(1/(F_nF_{n+1})\) entre dos cocientes consecutivos. Su crecimiento hace arbitrariamente pequeña esta horquilla racional.

La medida sobre todas las historias compatibles no coincide con la frecuencia de modos del calendario de tres eventos. La frecuencia cronológica es \((1/3,2/3)\). La medida de máxima entropía del envolvente simbólico definido por \(F_{\rm av}\) se construye con sus vectores propios y tiene razón de pesos \(\varphi^{-2}\). Esta última interviene en la lectura areal--volumétrica; no se obtiene sustituyendo sin más una frecuencia de conteo por otra.

\subsubsection{Espacio de trayectorias y medida invariante entre hojas}
\label{sec:medida-retorno-hojas}

La topología toroidal de APP fija el soporte de las traslaciones cardinales
y sus clases de enrollamiento. La correspondencia entre hojas procede de
otra operación sobre ese mismo soporte: la selección temporal de la
evaluación aditiva o multiplicativa. Su tipado geométrico es
\[
 \eta_{\rm av}(\Sigma)=\mathscr H_{\rm ar},\qquad
 \eta_{\rm av}(\Pi)=\mathscr H_{\rm vol}.
\]
La primera hoja corresponde a la lectura aditiva de frontera; la segunda,
a la lectura multiplicativa de interior. Este mapa transporta las etiquetas
de modo y conserva la incidencia \eqref{eq:fav}. El cuadrado de la razón
de autoescala se determina mediante las trayectorias de esa incidencia.

\Needspace{8\baselineskip}
\begin{definition}[Envolvente simbólico de memoria uno]
Sea
\[
 X_{\rm av}=\{(s_k)_{k\in\mathbb Z}\in
 \{\mathscr H_{\rm ar},\mathscr H_{\rm vol}\}^{\mathbb Z}:
 (F_{\rm av})_{s_k,s_{k+1}}=1\text{ para todo }k\}.
\]
El desplazamiento de índices actúa sobre este espacio. Es el menor sistema
simbólico de memoria uno que contiene las sucesiones modales y conserva
sus pares consecutivos: admite las dos transiciones cruzadas y la
autorretención volumétrica. La fase, la orientación cardinal y los registros
de la trayectoria permanecen en el estado TPK del que se obtiene este
envolvente.
\end{definition}

\begin{proposition}[Ponderación de las trayectorias entre hojas]
En el orden areal--volumétrico, la medida invariante de entropía máxima
de \(X_{\rm av}\) tiene matriz de transición y distribución estacionaria
\begin{equation}
 P_{\rm av}=\begin{pmatrix}0&1\\\varphi^{-2}&\varphi^{-1}\end{pmatrix},
 \qquad
 (\mu_{\rm ar},\mu_{\rm vol})
 =\frac{(1,\varphi^2)}{1+\varphi^2}.
 \label{eq:medida-hojas-exacta}
\end{equation}
En particular, \(\mu_{\rm ar}/\mu_{\rm vol}=\varphi^{-2}\).
\end{proposition}
\begin{proof}
La incidencia ya construida tiene vector propio positivo
\(r=(1,\varphi)^{\mathsf T}\) y autovalor \(\varphi\). La
normalización
\[
 (P_{\rm av})_{ij}
 =\frac{(F_{\rm av})_{ij}r_j}{\varphi r_i}
\]
da la matriz indicada. Sus filas suman uno por
\(\varphi^{-2}+\varphi^{-1}=1\). La simetría de \(F_{\rm av}\)
implica que los pesos proporcionales a \(r_i^2\) satisfacen balance
detallado; por tanto, \(\mu P_{\rm av}=\mu\). La irreducibilidad y la
autorretención volumétrica determinan una única distribución estacionaria.

Para verificar la maximalidad, en toda transición permitida se cumple
\(\log(P_{\rm av})_{ij}=\log r_j-\log r_i-\log\varphi\).
Bajo cualquier medida invariante, los términos de extremo se cancelan
en promedio. Si \(\nu_i\) son sus pesos de vértice y
\(q_i(j)=\nu(s_1=j\mid s_0=i)\), para \(\nu_i>0\), entonces
\[
 h_\nu\leq H_\nu(s_1\mid s_0)
 =\log\varphi-\sum_{i:\nu_i>0}\nu_iD(q_i\Vert P_i)
 \leq\log\varphi,
\]
donde \(D(q\Vert p)=\sum_jq_j\log(q_j/p_j)\) es la entropía
relativa en los pares permitidos y se adopta \(0\log0=0\).
La medida markoviana de \eqref{eq:medida-hojas-exacta} alcanza la cota.
La igualdad en la primera desigualdad exige la propiedad de Markov;
la igualdad en la segunda exige \(q_i=P_i\). La estacionariedad y
la irreducibilidad fijan entonces \(\mu\), lo que establece la
unicidad. Esta es la medida de Parry del envolvente definido.
\end{proof}

La frecuencia \((1/3,2/3)\) cuenta los instantes del calendario primitivo;
\(\mu\) pondera las trayectorias del envolvente. Su razón cuadrática
expresa el producto de las amplitudes de incidencia de entrada y salida.
La memoria y las restricciones de fase del TPK conservan su dominio
propio; la entropía del envolvente corresponde al sistema simbólico
especificado en esta definición.

\subsubsection{Operador de clausura y límite de los retornos marcados}
\label{sec:operador-retorno-marcado}

La coordenada de clausura \(\pi\), generada por el lector regional,
interviene después de la construcción de la incidencia modal. Sean
\(e_{\rm ar}=(1,0)^{\mathsf T}\),
\(e_{\rm vol}=(0,1)^{\mathsf T}\) y
\(P_i=e_ie_i^{\mathsf T}\) los proyectores de hoja. Las condiciones
del lector son conservación de ambas hojas, normalización volumétrica
unitaria y marca areal igual a la coordenada de clausura. Estas condiciones
determinan el operador positivo
\begin{equation}
 D_\pi=\pi P_{\rm ar}+P_{\rm vol}
       =\begin{pmatrix}\pi&0\\0&1\end{pmatrix}.
 \label{eq:operador-marca-pi}
\end{equation}
La diagonalidad expresa la conservación de hojas y las dos entradas
expresan sus normalizaciones; de este modo queda explícita la unicidad
del operador para el lector fijado.

\Needspace{8\baselineskip}
\begin{proposition}[Razón de retornos de frontera]
Para \(n\geq1\), la evaluación marcada
\begin{equation}
 \Theta_n=
 \frac{\langle e_{\rm ar},D_\pi F_{\rm av}^{\,n}e_{\rm ar}\rangle}
 {\langle e_{\rm vol},F_{\rm av}^{\,n}e_{\rm vol}\rangle}
 =\pi\frac{F_{n-1}}{F_{n+1}}
 \label{eq:retorno-marcado-finito}
\end{equation}
satisface
\begin{equation}
 \lim_{n\to\infty}\Theta_n
 =\pi\frac{\mu_{\rm ar}}{\mu_{\rm vol}}
 =\frac{\pi}{\varphi^2}.
 \label{eq:retorno-marcado-limite}
\end{equation}
\end{proposition}
\begin{proof}
Las entradas diagonales de \(F_{\rm av}^{\,n}\) cuentan los caminos
de longitud \(n\) que regresan a su hoja inicial. La fórmula de potencias
ya demostrada proporciona \(F_{n-1}\) y \(F_{n+1}\). El operador
\(D_\pi\) multiplica por \(\pi\) el primer recuento. Finalmente,
\(F_{n-1}/F_{n+1}\) es el producto de dos cocientes consecutivos
que convergen a \(\varphi^{-1}\); su límite es
\(\varphi^{-2}\). La identidad con la razón de pesos se sigue de
\eqref{eq:medida-hojas-exacta}.
\end{proof}

El operador de clausura aparece una sola vez, en la evaluación terminal
del retorno. La transferencia \(F_{\rm av}\) permanece determinada
por el calendario. Así se compone la información de dos lectores del
mismo estado: la incidencia modal determina la ponderación cuadrática
y la región de clausura aporta su coordenada a la frontera.

La composición conserva también el refinamiento. Para cilindros positivos
\(I=[a,b]\) y \(J=[c,d]\), la imagen del lector
\(\mathcal R(x,y)=x/y^2\) es el intervalo
\[
 \mathcal Q(I,J)=\left[\frac{a}{d^2},\frac{b}{c^2}\right].
\]
La monotonía en cada argumento prueba que los cilindros encajados de
clausura y autoescala producen intervalos encajados de la razón; la
positividad de sus límites y la contracción de sus diámetros determinan
\(\pi/\varphi^2\). La expresión escalar retiene así una cadena
de operaciones compatible con la profundidad, que se aplica a
continuación en sus dos orientaciones.


Después de generados los caracteres, las dos orientaciones del lector
\[
 R(x,y)=\frac{x}{y^2}
\]
proporcionan
\begin{equation}
 \rho_A=\frac\pi{\varphi^2},\qquad
 \rho_E=\frac\varphi{\pi^2},\qquad
 \varphi^3\rho_A^2\rho_E=1.
 \label{eq:areal-electron}
\end{equation}
La segunda expresión interviene en el exponente de la construcción electrónica de la sección~\ref{el-centro:registros}. La primera se obtiene también como límite de la lectura marcada \(\pi F_{n-1}/F_{n+1}\). El intercambio de argumentos relaciona ambas, pero no convierte una reparametrización de la ecuación electrónica en una segunda derivación independiente de todos sus coeficientes.

% Articulación posterior a la generación de pi, phi y e.
% Contenido acumulativo; no reemplaza el desarrollo previo del TRIT.
\subsection{Realizaciones exponenciales de las coordenadas generadas}
\label{euler-trit-articulacion}

Las coordenadas regionales ya determinadas permiten reconocer las
realizaciones concretas de la familia cuadrática del TRIT. El orden es
sustantivo: la construcción de las regiones precede a la evaluación de una
exponencial en los parámetros \(\pi\) o \(2\log\varphi\). Estas expresiones
identifican la función de las coordenadas en operadores construidos desde
APP y TPK; no intervienen en la selección de los prefijos que las producen.

\subsubsection{Ciclo ternario, transporte nilpotente e incidencia de hojas}

La multiplicación \(a(r)=4r\) en \(\mathbb Z/9\mathbb Z\) tiene orden
\(3\). Sobre las unidades determina las órbitas \(\{1,4,7\}\) y
\(\{2,8,5\}\). En el módulo de aumento de cualquiera de ellas, formado por
las combinaciones enteras cuyos coeficientes suman cero, la base
\((e_r-e_{4r},e_{4r}-e_{7r})\) da
\[
 C=\begin{pmatrix}0&-1\\1&-1\end{pmatrix},
 \qquad C^2+C+I=0.
\]
La forma restringida tiene matriz de Gram
\(\bigl(\begin{smallmatrix}2&-1\\-1&2\end{smallmatrix}\bigr)\).
En las evaluaciones residuales módulo \(9\), la acción distingue las hojas:
\(S(ax,ay)=aS(x,y)\), mientras \(P(ax,ay)=a^{-1}P(x,y)\).
La realización cíclica conserva, por tanto, una orientación contragrediente
entre suma y producto. Los cocientes del levantamiento entero se conservan
en el registro aritmético correspondiente.

El transporte parabólico elemental se representa mediante
\[
 N=\begin{pmatrix}0&1\\0&0\end{pmatrix}.
\]
La incidencia areal--volumétrica ya construida en \eqref{eq:fav} proporciona
\[
 F_{\rm av}=\begin{pmatrix}0&1\\1&1\end{pmatrix},
 \qquad A=F_{\rm av}^2=\begin{pmatrix}1&1\\1&2\end{pmatrix}.
\]
Los tres operadores actúan en espacios reales declarados: el módulo de
aumento tensorizado con \(\mathbb R\), el plano del transporte nilpotente y
el plano de incidencia modal. Compartir dimensión matricial no identifica
estos espacios ni sus coordenadas.

\begin{proposition}[Generadores concretos de los tres tipos]
\label{euler-trit-generadores}
Los operadores
\[
 J_{\mathrm e}=\frac{2C+I}{\sqrt3},\qquad
 J_{\mathrm p}=N,\qquad
 J_{\mathrm h}=\frac{2A-3I}{\sqrt5}
\]
satisfacen, respectivamente,
\[
 J_{\mathrm e}^{\,2}=-I,\qquad
 J_{\mathrm p}^{\,2}=0,\qquad
 J_{\mathrm h}^{\,2}=I.
\]
\end{proposition}
\begin{proof}
La relación de \(C\) implica \((2C+I)^2=-3I\).
La multiplicación directa da \(N^2=0\).
Finalmente, \(A\) tiene traza \(3\) y determinante \(1\), por lo que
\(A^2-3A+I=0\) y \((2A-3I)^2=5I\).
\end{proof}

\begin{proposition}[Tres evaluaciones de la identidad exponencial]
\label{euler-trit-evaluaciones}
En las realizaciones anteriores,
\begin{equation}
 \boxed{
 \exp(\pi J_{\mathrm e})+I=0,\qquad
 \exp(J_{\mathrm p})=I+J_{\mathrm p},\qquad
 \exp(2\log\varphi\,J_{\mathrm h})=F_{\rm av}^{\,2}.}
 \label{euler-trit-tres-evaluaciones}
\end{equation}
\end{proposition}
\begin{proof}
La especialización elíptica de la exponencial trítica da
\(\cos\pi\,I+\sin\pi\,J_{\mathrm e}=-I\).
La nilpotencia elimina exactamente los términos de grado al menos \(2\)
en la segunda expresión. Para la tercera, \(\varphi^2=\varphi+1\) implica
\[
 \cosh(2\log\varphi)=\frac32,\qquad
 \sinh(2\log\varphi)=\frac{\sqrt5}{2}.
\]
Por tanto, la exponencial vale
\(\frac32I+\frac{\sqrt5}{2}J_{\mathrm h}=A\).
\end{proof}

La primera igualdad representa la media vuelta elíptica; la vuelta completa
satisface \(\exp(2\pi J_{\mathrm e})=I\). La segunda mantiene un término
lineal no nulo: el régimen parabólico expresa transporte, no ausencia de
evolución. La tercera realiza un avance hiperbólico unimodular. Sus
autovalores son \(\varphi^2\) y \(\varphi^{-2}\); una dirección se expande y
la otra se contrae conservando el determinante. Las tres evaluaciones
emplean parámetros distintos de una familia exponencial común.

En esa familia, \(e\) interviene mediante la evaluación escalar
\(\exp(I)=eI\) y la ley \(\exp(sI)\exp(tI)=\exp((s+t)I)\).
Su papel no se restringe al tipo parabólico. La coordenada \(e\), obtenida
previamente de su región, es reconocida aquí como la evaluación unitaria
de la exponencial; esta identidad no sustituye su genealogía coinductiva.

\subsubsection{Resolución polinómica de los regímenes}

Las tres realizaciones admiten una presentación compacta sobre la suma
directa de sus espacios. Defínanse
\[
 \mathbb J=J_{\mathrm e}\oplus J_{\mathrm p}\oplus J_{\mathrm h},
 \qquad Q=\mathbb J^2.
\]
Las relaciones anteriores implican
\[
 \mathbb J^6=\mathbb J^2,\qquad Q^3=Q.
\]
El uso de \(Q\) mantiene separado este cuadrado operatorial del registro
dodecafásico \(K\).

\begin{proposition}[Proyectores polinómicos de régimen]
\label{euler-trit-proyectores}
Las aplicaciones
\[
 p_{\mathrm e}=\frac{Q^2-Q}{2},\qquad
 p_{\mathrm p}=I-Q^2,\qquad
 p_{\mathrm h}=\frac{Q^2+Q}{2}
\]
son idempotentes mutuamente ortogonales, suman \(I\) y recuperan los tres
sumandos de la representación.
\end{proposition}
\begin{proof}
Los tres polinomios toman los valores \((1,0,0)\), \((0,1,0)\) y
\((0,0,1)\) en los puntos \(-1,0,+1\), respectivamente. Las diferencias
entre sus cuadrados y ellos mismos, así como sus productos cruzados, son
múltiplos de \(X^3-X\). Sustituir \(Q\), que anula ese polinomio, demuestra
las identidades. En los tres bloques, \(Q\) actúa como \(-I,0,I\).
\end{proof}

Se obtiene así
\begin{align}
 \exp(t\mathbb J)
 ={}&p_{\mathrm e}(\cos t\,I+\sin t\,\mathbb J)
   +p_{\mathrm p}(I+t\mathbb J)\notag\\
   &+p_{\mathrm h}(\cosh t\,I+\sinh t\,\mathbb J).
 \label{euler-trit-exponencial-total}
\end{align}
Cada sumando conserva su relación cuadrática y la conjugación invierte
su evolución. La representación total reúne los tres regímenes, pero no
define por sí sola un operador de transición entre ellos. Tal transición
requiere los mapas de transporte del TPK con sus dominios, orientación y
memoria. La articulación exponencial permite reconocer sus tipos sin
reemplazar esa dinámica por una identificación de las fibras.

% PROCEDENCIA FOCAL (no impresa)
% Resultado recuperado: tres generadores y tres evaluaciones exponenciales.
% Propietario completo:
% BASE/manuscrito/sections/hmt/09b_algebras_cuadraticas.tex:289-516.
% Antecedente causal de C y acción contragrediente:
% BASE/manuscrito/sections/hmt/03d_esqueleto_ciclotomico_app.tex:13-47,145-246.
% Los tres tipos y la orientación están en:
% BASE/manuscrito/sections/hmt/02_trit_geometrias.tex:83-196.
% Resolución polinómica recuperada:
% output/INVESTIGACION_HOLONOMIA_NONADICA_CRISTAL_APERIODICO_20260903/
% ALGEBRA_HOLONOMICA_UNIVERSAL_Y_EULER_TRITICO.md, secciones 8-11 y 22.
% Los cuatro propietarios se leyeron completos. BASE designa:
% /Users/ruben/Documents/HMT2/
% EL_CIERRE_HOLOGRAFICO_DEL_INFINITO_SUCESOR_102_CAPITULOS_2026-08-28_EN_TRABAJO/
% No se incorpora la identidad P9(alpha)=delta del documento especializado:
% esa identificación no pertenece a este bloque.
% Tampoco se identifica una dirección nilpotente con la vacancia electrónica,
% ni se presume una representación global del transporte por la suma directa.
% COMPROBACIONES: Python estándar, enteros y Fraction.
% C^2+C+I=0; (2C+I)^2=-3I; N^2=0; (2Fav^2-3I)^2=5I.
% 81 pares de residuos verifican las dos acciones contragredientes.
% Los 3 proyectores se verificaron en Q=-1,0,+1; prueba general en texto.
% 6 labels únicos y entornos equilibrados. Sin compilación.


\subsection{Cilindros compatibles y profundidad arbitraria}

La realización de un bloque ternario \(b=(b_1,\ldots,b_6)\) es el entero
\[
 \nu(b)=\sum_{j=1}^6b_j3^{6-j}\in\{0,\ldots,728\}.
\]
Una sucesión de bloques define prefijos enteros por \(N_{n+1}=729N_n+\nu(b_{n+1})\), comenzando en \(N_0=m\), donde \(m\) es la parte entera determinada por el lector. Las horquillas anteriores sitúan los caracteres de clausura, propagación y autoescala en \((3,4)\), \((2,3)\) y \((1,2)\), respectivamente; por tanto, sus valores iniciales son \(m=3,2,1\). Los bloques siguientes refinan la parte fraccionaria. El cilindro correspondiente es
\begin{equation}
 I_n=\left[\frac{N_n}{729^n},\frac{N_n+1}{729^n}\right],
 \qquad |I_n|=729^{-n}.
 \label{eq:cilindro}
\end{equation}
La convención semiabierta se utiliza para seleccionar un representante cuando un valor cae en una frontera. Para el paso al límite se utilizan las clausuras de estos intervalos y se identifican las dos expansiones fronterizas equivalentes. Esta precisión impide suponer que toda intersección de intervalos semiabiertos encajados sea automáticamente no vacía.

\begin{theorem}[Determinación de la realización escalar]
Una historia de bloques compatible determina un único límite de los extremos izquierdos de \eqref{eq:cilindro}. Si un carácter exacto dispone de horquillas racionales de anchura arbitrariamente pequeña y se separa de la frontera de cada cilindro decimal examinado, cada prefijo decimal de longitud finita se determina después de un cálculo finito. Los prefijos así obtenidos son compatibles entre sí.
\end{theorem}
\begin{proof}
La extensión de un bloque conserva el prefijo anterior por división euclídea. Los cilindros cerrados están encajados y su diámetro tiende a cero; sus extremos izquierdos son de Cauchy y tienen un único límite. Para una longitud decimal fijada, un valor separado de las fronteras de la partición posee distancia positiva a ellas. Una horquilla suficientemente estrecha queda contenida en una única celda y determina su prefijo. La unicidad de la celda y el encajamiento de las particiones prueban la compatibilidad por truncamiento. En los puntos con expansión terminal debe añadirse la decisión exacta de frontera y su convención de representación; no se deduce esa decisión de una anchura numérica solamente.
\end{proof}

\begin{corollary}[Obtención finita y compatibilidad de los prefijos regionales]
\label{gen:prefijos-regionales}
Para cada una de las realizaciones \(\pi,\varphi,e\) y para toda longitud
decimal finita, las horquillas construidas determinan el prefijo correspondiente
mediante un cálculo finito. Los prefijos obtenidos a distintas profundidades
son compatibles por truncamiento.
\end{corollary}
\begin{proof}
Los caracteres de clausura, propagación y autoescala poseen las horquillas
anteriores. Tras su reconocimiento, la irracionalidad de \(\pi,e,\varphi\)
garantiza su separación de las fronteras racionales de cada partición decimal
finita. Se aplica el teorema precedente a cada carácter.
\end{proof}
La profundidad arbitraria expresa el cuantificador del corolario: para todo
horizonte finito existe un cálculo que determina su prefijo. El objeto
matemático es la familia compatible completa; cada ejecución obtiene una
proyección finita de ella.

\subsection{Involución especular y sincronización de las regiones}

La conexión nonádica transporta conjuntamente los tres caracteres. En la carta de la novena fase, el eje de autoescala y la suma ternaria de clausura y propagación permanecen fijados bajo el intercambio especular. La inversión helicoidal cambia el sentido del transporte; el espejo de los dos canales cambia su orientación relativa. Son involuciones diferentes, aunque ambas actúen sobre el mismo sistema de trayectorias.

Las coincidencias de censo ilustran esa coordinación:
\begin{center}
\begin{tabular}{@{}c rrr l@{}}\toprule
Fase&\(N_\pi\)&\(N_e\)&\(N_\varphi\)&Igualdad de conteos\\\midrule
4&207&207&122&\(N_\pi=N_e\)\\
5&39&151&151&\(N_e=N_\varphi\)\\
6&110&110&69&\(N_\pi=N_e\)\\
7&81&29&81&\(N_\varphi=N_\pi\)\\
8&59&59&59&sincronización triple\\
9&43&43&43&sincronización triple\\\bottomrule
\end{tabular}
\end{center}
Estos enteros cuentan coordenadas compatibles de las regiones. No afirman una igualdad entre \(\pi\), \(e\) y \(\varphi\) como números reales. Si
\[
 \partial(a,b,c)=(a-b,b-c,c-a),
\]
su imagen pertenece al retículo \(A_2=\{(u,v,w)\in\ZZ^3:u+v+w=0\}\). En las fases ocho y nueve la componente relativa se anula, pero la componente diagonal cambia de \((59,59,59)\) a \((43,43,43)\). La sincronización relativa coexiste con un desplazamiento de \(-16\) por canal. El cierre que conducirá a \(\alpha\) debe conservar ambas informaciones: coincidencia relativa y evolución de la memoria.

\subsection{Suspensión helicoidal esturmiana}

La relación entre las bases nueve y diez de la completación define el parámetro adimensional
\begin{equation}
 \delta=\frac{\log(10/9)}{\log10}.
 \label{eq:delta-sturm}
\end{equation}
No se obtiene de una constante física objetivo. Es la escala logarítmica relativa de dos bases ya presentes en la construcción. Su empleo como lectura simbólica del refinamiento tiene una consecuencia exacta.

\begin{lemma}[Irracionalidad de la pendiente]
El número \(\delta\) es irracional.
\end{lemma}
\begin{proof}
Si \(\delta=p/q\), \(q>0\), entonces \((10/9)^q=10^p\) y \(9^q=10^{q-p}\). La factorización del primer miembro contiene el primo tres y la del segundo sólo los primos dos y cinco. La igualdad es imposible.
\end{proof}

La palabra mecánica inferior, con intercepto \(\theta\), es
\[
 z_n(\theta)=\lfloor(n+1)\delta+\theta\rfloor-\lfloor n\delta+\theta\rfloor\in\{0,1\}.
\]
Se obtiene al observar la rotación \(\theta\mapsto\theta+\delta\pmod1\) mediante el intervalo \([1-\delta,1)\). La sucesión de eventos tiene densidad \(\delta\), y la suma sobre cualquier ventana de longitud \(L\) vale \(\lfloor L\delta\rfloor\) o \(\lceil L\delta\rceil\). Como \(21<1/\delta<22\), las separaciones entre eventos son veintiuno o veintidós. La palabra que indica cuáles son los intervalos cortos es otra palabra mecánica de pendiente \(\sigma=22-1/\delta\); sus eventos están separados por seis o siete intervalos primarios. La segunda escala es inducida por la primera y no un segundo parámetro libre.

% Desarrollo reunido de DESARROLLO_MATEMATICO, apartados 38--40 y 50.
\subsubsection{Índice de capacidad, vacancias e inducción del régimen trítico}
\label{vac-capacidad-trit}

La sucesión de vacancias está vinculada al refinamiento mediante una
identidad de incrementos. Para mostrarla se fija el origen de fase
\(\theta=0\). Sean
\[
 \rho=\log_{10}9=1-\delta,\qquad
 \mathcal C_t=\lfloor(t+1)\rho\rfloor\quad(t\geq0).
\]
El índice \(\mathcal C_t\) expresa la capacidad decimal del refinamiento
en la normalización que compara un bloque ternario de \(6\) posiciones,
con \(729\) posibilidades, y un bloque decimal de \(3\) posiciones, con
\(1000\) posibilidades. En efecto,
\(\log_{1000}729=\log_{10}9=\rho\). Se reserva la letra \(K\) para
el registro dodecafásico que se construirá posteriormente.

\begin{proposition}[Complementariedad de capacidad y vacancia]
\label{prop:capacidad-vacancia}
Para \(t\geq1\),
\begin{equation}
 \mathcal C_t-\mathcal C_{t-1}=1-z_t(0).
 \label{eq:capacidad-vacancia}
\end{equation}
En consecuencia, si \(0\leq a<b\),
\[
 \mathcal C_b-\mathcal C_a+
 \sum_{t=a+1}^{b}z_t(0)=b-a.
\]
\end{proposition}
\begin{proof}
La irracionalidad de \(\delta\) da
\[
 \mathcal C_t
 =t+1-\lceil(t+1)\delta\rceil
 =t-\lfloor(t+1)\delta\rfloor.
\]
La resta de dos términos consecutivos produce
\eqref{eq:capacidad-vacancia}. La suma sobre el intervalo es telescópica.
\end{proof}

Cada transición aumenta en una unidad el contador de bloques. Si
\(z_t=0\), aumenta también la capacidad; si \(z_t=1\), la capacidad se
mantiene mientras se incorpora una transición a la historia. La vacancia
es aquí un evento de retención de capacidad, con una posición y un
antecedente determinados. El balance anterior conserva exactamente la
longitud entre capacidad adquirida y vacancias registradas. La fase
observada \(g_t^{\rm cap}=[\mathcal C_t]_9\) satisface
\[
 g_t^{\rm cap}-g_{t-1}^{\rm cap}\equiv1-z_t\pmod9.
\]
Esta fase de capacidad y la fase cronológica \([t]_9\) son dos
observaciones diferentes. Una retención local de capacidad y una vuelta
completa de la holonomía tienen mecanismos distintos, aunque ambas
permitan conservar la observación de fase con modificación del estado.

La relación con TRIT puede seguirse en las propias posiciones de vacancia.
Para \(j\geq1\), defínanse
\[
 p_j=\left\lfloor\frac j\delta\right\rfloor,\qquad
 v_j=\mathcal C_{p_j},\qquad h_j=p_{j+1}-p_j.
\]
La desigualdad \(p_j\delta<j<(p_j+1)\delta\) identifica \(p_j\)
como la posición del evento \(j\). Como \(0<\delta<1\), da también
\(\lfloor(p_j+1)\delta\rfloor=j\), y por tanto
\[
 v_j=p_j-j,\qquad
 v_{j+1}-v_j=h_j-1.
\]
Si \(s_j=22-h_j\), entonces \(s_j\in\{0,1\}\) y
\begin{equation}
 [v_{j+1}]_3=[v_j-s_j]_3.
 \label{eq:vacancia-regimen}
\end{equation}
Un intervalo de longitud \(22\) conserva la clase; uno de longitud
\(21\) la desplaza una unidad en sentido negativo. El selector canónico
se aplica a la fase positiva \(r_j=1+[v_j]_9\):
\[
 \tau_j=M_{\rm ph}(r_j).
\]
Queda determinado así un régimen trítico por el índice de capacidad
retenido en cada evento, además de la codificación binaria de la vacancia.

Las posiciones de los intervalos cortos forman la codificación mecánica
inducida de pendiente \(22-1/\delta\). Sus separaciones \(6\) y \(7\)
determinan las longitudes de las mesetas de la clase \([v_j]_3\), una
vez fijados el origen y la convención de extremos. Las primeras clases
son \(2\) durante \(6\) eventos, \(1\) durante \(7\) y \(0\) durante
\(7\); el primer tramo de \(6\) es la restricción inicial de una meseta.
Sus firmas recorren \(0,-1,+1\). Las ecuaciones
\eqref{eq:capacidad-vacancia} y \eqref{eq:vacancia-regimen} explicitan el
enlace entre refinamiento, vacancias y selección del régimen. La lectura
cuadrática del TRIT realiza después esos tipos mediante sus respectivos
generadores; el acontecimiento temporal y su representación geométrica
conservan así una dependencia identificable.

La inversión del índice monótono proporciona otra observación útil:
\[
 \mathcal T(k)=\min\{t:\mathcal C_t=k\}
   =\left\lceil\frac k\rho\right\rceil-1\qquad(k\geq1).
\]
Si \(\sigma=\rho^{-1}-1\), el número de bloques necesarios para ganar
\(m\) unidades de capacidad es
\[
 \mathcal T(k+m)-\mathcal T(k)
 =m+\lceil(k+m)\sigma\rceil-\lceil k\sigma\rceil
 \in\{m+\lfloor m\sigma\rfloor,m+\lceil m\sigma\rceil\}.
\]
En particular, \(9\) unidades de capacidad requieren \(9\) o \(10\)
bloques. El período de una fase finita permanece exacto, mientras el
tiempo de retorno medido por el otro índice admite dos valores. Esta
distinción es parte de la dinámica del refinamiento y no una corrección
de sus valores numéricos a posteriori.


\begin{theorem}[Complejidad de las lecturas de fase]
Para \(m\ge1\), la palabra mecánica tiene complejidad factorial \(p(m)=m+1\). Al conservar la fase \(n\bmod9\), la complejidad es \(p_9(m)=9(m+1)\); al conservar sólo su clase ternaria, es \(p_3(m)=3(m+1)\). Las tres entropías topológicas son nulas.
\end{theorem}
\begin{proof}
Los predecesores de los dos extremos del intervalo de codificación forman, para factores de longitud \(m\), \(m+1\) puntos distintos de la circunferencia: la coincidencia entre extremos desplazados elimina las duplicaciones y la irracionalidad impide coincidencias adicionales. Esos puntos determinan \(m+1\) intervalos de itinerario constante y distintos. Toda órbita irracional visita cada intervalo, de donde \(p(m)=m+1\).

Fijada una fase módulo nueve, las posiciones de esa fase recorren una rotación de paso \(9\delta\), también irracional. Por densidad aparecen todos los factores en cada fase, y la primera coordenada distingue sus nueve copias. El mismo argumento con paso \(3\delta\) da tres copias tras la reducción ternaria. Finalmente, \(m^{-1}\log(c(m+1))\to0\) para \(c=1,3,9\).
\end{proof}

El producto de fases \(C_9\times\mathbb T\), trasladado por \((1,\delta)\), tiene caracteres de frecuencia
\[
 \frac a9+k\delta\pmod1,\qquad a\in\ZZ/9\ZZ,\ k\in\ZZ.
\]
El módulo espectral no es \(\log10\): es el módulo aditivo generado por \(1/9\) y \(\delta\), módulo los enteros. El factor \(\log10\) pertenece a la normalización de \eqref{eq:delta-sturm}.

Una presentación operatoria hace explícita la memoria. En \(\ell^2(\ZZ)\), sea \(T|n\rangle=|n+1\rangle\), \(V_9|n\rangle=\zeta_9^n|n\rangle\), \(W_\delta|n\rangle=e^{2\pi i n\delta}|n\rangle\) y \(D_M|n\rangle=\lfloor n/9\rfloor|n\rangle\). Sobre vectores de soporte finito,
\begin{equation}
 V_9T=\zeta_9TV_9,\qquad
 W_\delta T=e^{2\pi i\delta}TW_\delta,\qquad
 [D_M,T^9]=T^9.
 \label{eq:weyl-relojes}
\end{equation}
Las dos primeras relaciones son covariancias de Weyl; la tercera expresa que una vuelta nonádica aumenta el grado de memoria. La rotación de base es abeliana, mientras el álgebra de observables y traslaciones no lo es. La inversión \(|n\rangle\mapsto|-n\rangle\) conjuga el avance con el retroceso. Esta realización proporciona una suspensión helicoidal aperiódica con retorno finito de fase, no una palabra periódica de longitud nueve.

\begin{figure}[htbp]
\centering
\begin{tikzpicture}[x=1cm,y=1cm,>=Latex,font=\small]
\begin{scope}[xshift=1.8cm,yshift=.3cm]
\draw[->,gray] (0,-.5)--(0,3.65) node[above] {\(n/9\)};
\foreach \k in {0,1,2,3}{
 \draw[gray!60,dashed] (-1.1,\k)--(1.1,\k);
 \node[left,font=\scriptsize] at (-1.1,\k){\(\k\)};
}
\draw[black,line width=.65pt,domain=0:27,samples=260,smooth,variable=\n]
 plot ({cos(40*\n)},{\n/9+.14*sin(40*\n)});
\foreach \n in {0,...,27}{
 \fill ({cos(40*\n)},{\n/9+.14*sin(40*\n)}) circle (.026);
}
\foreach \k in {0,1,2,3}{\fill (1,\k) circle (.045);}
\node[align=center,text width=3.3cm,font=\scriptsize] at (0,-1.02)
 {Misma fase en \(n=0,9,18,27\);\\grado de memoria distinto.};
\end{scope}
\begin{scope}[xshift=4.2cm,yshift=.5cm]
\node[anchor=west] at (0,3.1){\(z_n(0)=\lfloor(n+1)\delta\rfloor-\lfloor n\delta\rfloor\)};
\draw[->] (0,1.7)--(9,1.7) node[right]{\(n\)};
\foreach \n in {21,43,65,87,109,131,152,174,196,218}{
 \draw[line width=.6pt] ({\n*.039},1.7)--({\n*.039},2.25);
 \fill ({\n*.039},2.25) circle (.033);
 \node[below,font=\tiny] at ({\n*.039},1.7){\n};
}
\foreach \a/\b/\g in {21/43/22,43/65/22,65/87/22,87/109/22,109/131/22,131/152/21,152/174/22,174/196/22,196/218/22}{
 \node[font=\scriptsize] at ({(\a+\b)*.0195},2.55){\g};
}
\node[anchor=west,align=left,text width=8.3cm] at (0,.55)
 {La codificación irracional determina los intervalos entre eventos. El retorno de la fase nonádica no convierte esta secuencia en una palabra periódica.};
\end{scope}
\end{tikzpicture}
\caption{Representaciones complementarias del transporte. A la izquierda, proyección de la hélice \((\cos(2\pi n/9),\sin(2\pi n/9),n/9)\): el avance de nueve pasos recupera la fase e incrementa la altura. A la derecha, los diez primeros eventos de la palabra inferior para \(\delta=\log_{10}(10/9)\), con separaciones exactas \(21\) o \(22\). La representación helicoidal no introduce una métrica física adicional.}
\label{fig:holonomia-esturmiana}
\end{figure}


Para este recuento fijamos el intercepto \(\theta=0\) y los extremos acumulados \((N_0,\ldots,N_4)=(0,729,972,999,1000)\), correspondientes al orden cúbico \((729,243,27,1)\). Los conteos inferiores son \(\lfloor N_j\delta\rfloor-\lfloor N_{j-1}\delta\rfloor\), y dan \((33,11,1,0)\). Los superiores son \(\lceil N_j\delta\rceil-\lceil N_{j-1}\delta\rceil\), con la misma frontera inicial cero, y dan \((34,11,1,0)\). Estos valores dependen del intercepto y del orden de los estratos; no se afirman para una fase inicial arbitraria. Sus totales son cuarenta y cinco y cuarenta y seis. El incremento corresponde a una única unidad de frontera, situada en el primer estrato de esta partición ordenada. Las cifras treinta y cuatro y once aparecen así en un mismo recuento previo a toda carta de unidades; identificarlas con exponentes físicos exige, además, la derivación dimensional correspondiente.

La denominación \emph{cristal aperiódico} se emplea aquí para este orden simbólico, con complejidad, espectro y retorno especificados. La entropía topológica nula no equivale por sí sola a disipación física nula: una realización termodinámica necesita una dinámica energética y una operación de borrado definidas. Del mismo modo, esta lectura esturmiana no sustituye a todos los operadores del TPK ni demuestra por sí sola la periodicidad o aperiodicidad de la salida completa de \(\alpha\). Su función es hacer explícita una estructura de orden que el simple retorno del calendario no describe.

% Adición acumulativa. Insertar después de la suspensión esturmiana existente.
% Conservar íntegro el texto y la figura fig:holonomia-esturmiana.
\subsection{Conexión nonádica, holonomía y representación monodrómica}
\label{mono-ampl:conexion}

La conexión regional ya construida se estudia ahora mediante tres objetos:
el transporte del estado completo, la representación de sus vueltas y la
suspensión de sus itinerarios. Las operaciones de APP--TRIT--TPK y el
catálogo inicial conservan las definiciones anteriores. La distinción
necesaria concierne aquí a lo que cada representación retiene de una
misma historia (secciones~\ref{sec:extension} y~\ref{sec:generacion}).

En el índice de prolongación \(k\), una carta del estado conserva
\[
 \widehat x_k=(B_k,C_k,p_k,c_k,\Sigma_k,W_k,g_k,m_k).
\]
Aquí \(B_k\) reúne los tres bloques regionales, \(C_k\) es el cilindro,
\(p_k\) el prefijo, \(c_k\) el registro de acarreo, \(\Sigma_k\) la firma
de frontera y \(W_k\) la información incidencial. La fase es
\(g_k=1+[(k-1)]_9\); \(m_k\) cuenta las vueltas completas. La reducción
\(\beta_k=[m_k]_{12}\) conserva la posición dodecafásica, mientras el
entero \(m_k\) distingue sus sucesivas vueltas.

\begin{definition}[Transporte de una vuelta nonádica]
\label{mono-ampl:def-vuelta}
Sean \(\mathcal R_k\subseteq\widehat X_k\times\widehat X_{k+1}\)
las relaciones de extensión admisible del TPK. La realización de la
holonomía en el nivel \(k\) es la composición relacional
\[
 \Gamma_{9,k}=\mathcal R_{k+8}\circ\cdots\circ\mathcal R_k.
\]
Para \(s\geq1\), su prolongación a \(s\) vueltas es
\[
 \Gamma_{9s,k}=
 \Gamma_{9,k+9(s-1)}\circ\cdots\circ\Gamma_{9,k}.
\]
\end{definition}

La formulación relacional conserva las continuaciones admisibles. Sobre
una historia individualizada, cada flecha lleva al estado efectivamente
prolongado; el bloque visible aislado puede seguir admitiendo más de una
continuación. Esta diferencia es precisamente la función de la memoria y
de las condiciones de frontera en la determinación de la trayectoria.

\begin{proposition}[Iteración del retorno y memoria acumulada]
\label{mono-ampl:vueltas}
En toda composición admisible de \(s\) vueltas,
\[
 g_{k+9s}=g_k,\qquad m_{k+9s}=m_k+s,\qquad
 \beta_{k+9s}=\beta_k+s\pmod{12}.
\]
\end{proposition}
\begin{proof}
La regla de una vuelta conserva la fase e incrementa el contador en una
unidad. Su aplicación al estado ya transportado repite exactamente esas
dos operaciones. La inducción da las dos primeras igualdades; reducir la
segunda módulo \(12\) da la tercera.
\end{proof}

El contador proporciona así un testigo entero de la diferencia entre las
vueltas. Tras \(12\) vueltas la pareja de fases finitas se restaura, pero
el contador aumenta \(12\). La extensión no escindida
\eqref{eq:extension108} expresa la composición del calendario mediante su
cociclo; la trayectoria conserva además la elevación entera de ese dato.
La conservación significa aquí transporte de las contribuciones previas,
no inmovilidad de todas las coordenadas.

\subsubsection{Representación bilateral del índice de retorno}

El desplazamiento \(T\) de \eqref{eq:weyl-relojes} admite una
expresión en coordenadas de fase y vuelta. Numeramos las fases de \(0\) a
\(8\), mediante un desplazamiento del origen respecto de \(g_k\), y
escribimos \(k=9m+r\), con \(0\leq r<9\), en la recta entera
que recubre el ciclo de fases. Sobre
\(\mathcal H_{\rm ind}=\ell^2(\mathbb Z)\), la base
\(|k\rangle=|r,m\rangle\) expresa ese mismo desplazamiento unitario
\[
 T|r,m\rangle=
 \begin{cases}
 |r+1,m\rangle,&r<8,\\
 |0,m+1\rangle,&r=8.
 \end{cases}
\]
La inversa reduce \(r\) en una unidad cuando \(r>0\), y lleva
\(|0,m\rangle\) a \(|8,m-1\rangle\). El paso por el origen de fase
incorpora, por tanto, la unidad exacta de cociente.

\Needspace{8\baselineskip}
\begin{proposition}[Representación de la monodromía del ciclo]
\label{mono-ampl:representacion}
Sea \(S=T^9\). Entonces
\[
 S|r,m\rangle=|r,m+1\rangle,\qquad
 \mathbb Z\longrightarrow\mathcal U(\mathcal H_{\rm ind}),
 \quad s\longmapsto S^s
\]
es una representación unitaria fiel del grupo de vueltas del ciclo.
\end{proposition}
\begin{proof}
Nueve desplazamientos conservan el resto de la división por \(9\) e
incrementan el cociente. Las potencias satisfacen
\(S^{s+t}=S^sS^t\), y \(S^s|r,m\rangle=|r,m+s\rangle\) distingue
todo entero \(s\). La permutación de una base ortonormal es unitaria.
\end{proof}

Esta representación corresponde al recubrimiento del ciclo, cuyo grupo
de vueltas es \(\mathbb Z\), y representa la fase con su contador. El
estado del TPK contiene además las coordenadas de trayectoria ya descritas.
La extensión bilateral admite índices negativos como coordenadas del
recubrimiento; una trayectoria iniciada en una condición inicial utiliza
su semieje admisible. De este modo se dispone de una inversa operatoria
del índice sin convertir cada relación de extensión en una biyección del
estado completo.

Con \(D_M|k\rangle=\lfloor k/9\rfloor|k\rangle\), la identidad
\([D_M,S]=S\) expresa el incremento de memoria en el dominio denso de
vectores de soporte finito. Si \(\mathcal I|k\rangle=|-k\rangle\) y
\(P_0\) proyecta sobre los índices divisibles por \(9\), se obtiene
además
\begin{equation}
 \mathcal I T\mathcal I=T^{-1},\qquad
 \mathcal I D_M\mathcal I=-D_M-(I-P_0).
 \label{mono-ampl:inversion-contador}
\end{equation}
En efecto, para \(k=9m+r\),
\(\lfloor-k/9\rfloor=-m\) si \(r=0\) y vale \(-m-1\) si \(r>0\).
La inversión conserva así el término de frontera de la división euclídea.

\subsection{Reloj de bloques y reloj de capacidad}
\label{mono-ampl:relojes}

El índice de una aplicación de refinamiento y el número de posiciones
decimales disponibles tienen leyes relacionadas, pero diferentes. Para
expresarlas reservamos \(n\) al índice de bloques ternarios y definimos
\[
 \rho=\log_{1000}729=\log_{10}9,\qquad
 \delta=1-\rho,\qquad
 \mathcal C_n=\lfloor(n+1)\rho\rfloor,\quad n\geq0.
\]
La letra \(\mathcal C_n\) designa aquí la capacidad, no el registro
dodecafásico \(K\). En el origen de fase canónico, la palabra mecánica
anterior proporciona
\[
 Z_n=\sum_{j=0}^{n}z_j(0)=\lfloor(n+1)\delta\rfloor.
\]

\begin{proposition}[Balance exacto entre capacidad y vacancias]
\label{mono-ampl:balance-capacidad}
Para \(n\geq0\) y, en la segunda identidad, \(n\geq1\),
\[
 \mathcal C_n=n-Z_n,\qquad
 \mathcal C_n-\mathcal C_{n-1}=1-z_n(0).
\]
\end{proposition}
\begin{proof}
La irracionalidad de \(\delta\) implica
\(\lceil(n+1)\delta\rceil=\lfloor(n+1)\delta\rfloor+1\).
Entonces
\[
 \lfloor(n+1)(1-\delta)\rfloor
 =n+1-\lceil(n+1)\delta\rceil=n-Z_n.
\]
Restar dos identidades consecutivas concluye la prueba.
\end{proof}

Un evento \(z_n=1\) añade un bloque a la historia manteniendo la capacidad
\(\mathcal C_n\); un evento \(z_n=0\) incrementa ambos contadores. En
cualquier intervalo, los incrementos de capacidad y las vacancias suman
exactamente su longitud. Reduciendo módulo \(q\), para \(q=9\) o
\(q=108\), resulta
\[
 [\mathcal C_n]_q=[n]_q-[Z_n]_q.
\]
La fase uniforme de bloques y la fase de capacidad se relacionan mediante
la memoria integrada de vacancias. En particular, una retención de
capacidad y una vuelta completa de la conexión tienen mecanismos
distintos, aunque ambas puedan conservar una observación de fase.

La inversión monótona del reloj da un tiempo de primera llegada:
\[
 T_{\rm cap}(a)=\min\{n:\mathcal C_n=a\}
 =\left\lceil\frac a\rho\right\rceil-1,
 \qquad a\geq1.
\]
Con \(\sigma=\rho^{-1}-1\), el tiempo empleado en ganar \(b\) unidades
de capacidad es
\[
 T_{\rm cap}(a+b)-T_{\rm cap}(a)
 =b+\lceil(a+b)\sigma\rceil-\lceil a\sigma\rceil.
\]
La diferencia de techos toma uno de los dos enteros adyacentes a
\(b\sigma\). Así se obtienen \(9\) o \(10\) bloques para ganar
\(9\) unidades de capacidad, y \(113\) o \(114\) para ganar \(108\).
Son tiempos del reloj inverso, medidos desde una primera llegada; el
contador de vueltas se interpreta siempre en el índice al que pertenece.

\subsection{Suspensión simbólica y cristal temporal aperiódico}
\label{mono-ampl:cristal}

La codificación por rotación reúne periodicidad de fase, recurrencia local
y aperiodicidad global. Sea \(\mathcal X_\delta\subseteq\{0,1\}^{\mathbb Z}\)
la clausura, en la topología producto, de las traslaciones de la palabra
bilateral \(z_n(\theta)\). Denotemos por \(\sigma_{\rm sh}\) el
desplazamiento de sus sucesiones. Conservando una fase nonádica, el paso es
\[
 F(r,z)=([r+1]_9,\sigma_{\rm sh}z),\qquad
 (r,z)\in C_9\times\mathcal X_\delta.
\]

\begin{definition}[Cristal temporal aperiódico simbólico]
\label{mono-ampl:def-cristal}
La suspensión de techo unitario del sistema anterior es el espacio
\[
 \mathcal S_\delta=
 \bigl((C_9\times\mathcal X_\delta)\times\mathbb R\bigr)/\sim,
 \qquad (y,t+1)\sim(Fy,t),
\]
con flujo \(\Phi_s[y,t]=[y,t+s]\). Se denomina \emph{cristal temporal
aperiódico simbólico} a este modelo de orden temporal, junto con su
codificación de vacancias y la representación del contador de vueltas.
\end{definition}

El adjetivo temporal designa la acción del flujo y de su sección de retorno;
la aperiodicidad caracteriza sus itinerarios. La frecuencia de unos de cada
sucesión es \(\delta\): la cota uniforme de discrepancia menor que uno,
obtenida por telescopía en cualquier ventana, persiste al tomar la clausura.
Una sucesión periódica tendría frecuencia racional. Por tanto,
\(\mathcal X_\delta\) no contiene órbitas periódicas. La suspensión
tampoco tiene órbitas cerradas: un retorno del flujo exigiría un tiempo
entero y una periodicidad del itinerario.

Las palabras finitas sí reaparecen. Los itinerarios de longitud \(m\)
proceden de una partición de la circunferencia por \(m+1\) puntos; cada
intervalo es visitado recurrentemente por la rotación irracional. Incluso
al restringir los tiempos a una fase fija, el paso \(9\delta\) sigue
siendo irracional. Ésta es la razón de que cada palabra aparezca en las
\(9\) fases y de las complejidades ya establecidas
\[
 p(m)=m+1,\qquad p_9(m)=9(m+1),\qquad p_3(m)=3(m+1).
\]
La recurrencia de configuraciones locales coexiste con la ausencia de un
período global. Su entropía topológica es nula porque estas complejidades
crecen linealmente.

La realización por los observables \(V_9,W_\delta\) y el desplazamiento
\(T\) expresa el mismo orden mediante las covariancias
\eqref{eq:weyl-relojes}. Las frecuencias del reloj uniforme pertenecen a
\(\frac19\mathbb Z+\delta\mathbb Z\), módulo \(1\). El reloj de
capacidad conserva adicionalmente la relación
\([\mathcal C_n]_9=[n]_9-[Z_n]_9\); no se sustituye por la fase uniforme
al interpretar los retornos. El término \emph{cristal temporal} queda así
definido mediante un sistema dinámico matemático. Su realización física
requiere especificar una dinámica energética, sus observables temporales
y su estabilidad, datos distintos de la construcción simbólica presente.

\subsection{Simetría especular de las regiones y doble realización}
\label{mono-ampl:regiones}

El intercambio especular actúa en el bloque regional
\(B=(u^\pi,u^e,u^\varphi)\in(\mathbb F_3^6)^3\) por
\[
 \mathfrak cB=(u^e,u^\pi,u^\varphi).
\]
Su eje fijo y su agregado son, respectivamente, \(u^\varphi\) y
\(s=u^\pi+u^e\). Esta fijación se refiere a la involución en una fibra;
ambas coordenadas pueden evolucionar durante el transporte nonádico.
Con \(q=u^\pi+u^e+u^\varphi\),
\(a=u^\pi+u^e-u^\varphi\) y \(d=u^\pi-u^e\), se recupera
\[
 u^\varphi=2(q-a),\qquad s=q-u^\varphi,\qquad
 u^\pi=2(s+d),\qquad u^e=2(s-d).
\]
Las igualdades son componente a componente en \(\mathbb F_3\), donde
\(2^{-1}=2\). El espejo conserva \(q,a,s\) e invierte \(d\); esta
coordenada antisimétrica individualiza el reparto que el agregado conserva
sin distinguir. La inversión \(\mathcal I\) del índice temporal y la
involución \(\mathfrak c\) de canales actúan, por consiguiente, sobre
datos diferentes.

El retorno de fase posee un testigo explícito:
\[
 \begin{aligned}
 B_1&=(012222\mid121221\mid112202),\\
 B_{10}&=(101222\mid112221\mid112002).
 \end{aligned}
\]
Las dos posiciones tienen fase \(1\); el bloque, el eje y la memoria han
cambiado. Asimismo, la sincronización de los censos en las fases \(8\) y
\(9\) anula sus diferencias relativas, mientras la componente diagonal
pasa de \((59,59,59)\) a \((43,43,43)\). La igualdad de un observable
relativo no agota el estado que lo determina.

La doble realización conserva esa arquitectura. En la sección
\ref{exc:doble-lectura}, el mismo prefijo determina por una parte cilindros
encajados y por otra una sucesión de palabras codificadas con sus marcos.
El primer mapa pasa al límite por contracción de diámetros; el segundo
satisface \(r_{n,m}Q_n=Q_mp_{n,m}\) y define \(Q_\infty\).
La figura \ref{mono-ampl:fig-regional} anticipa esas dos aplicaciones.
El recorrido incidencial posterior construye códigos, retículos y la
realización de Moonshine en sus dominios propios; el grupo de automorfismos
actúa sobre esa realización, no sobre las cifras por una identificación
implícita.

% Figura acumulativa; destino figures/monodromia_regional_ampliacion.tex.
\begin{figure}[p]
\centering
\begin{tikzpicture}[x=1cm,y=1cm,>=Latex,font=\small,
  flow/.style={->,line width=.55pt},
  ann/.style={align=center,font=\scriptsize}]
\node[font=\bfseries] at (6.9,9.2)
  {Conexión nonádica y coordenadas regionales};
\begin{scope}[shift={(3.4,6.8)}]
 \foreach \k in {1,...,9}{
  \pgfmathsetmacro{\ang}{90-40*(\k-1)}
  \coordinate (p\k) at (\ang:1.52);
  \fill (p\k) circle (1.4pt);
  \node[fill=white,inner sep=1.1pt,font=\scriptsize]
    at (\ang:1.83) {\(\k\)};
 }
 \foreach \k in {1,...,8}{
  \pgfmathtruncatemacro{\next}{\k+1}
  \draw[flow,shorten >=3pt,shorten <=3pt] (p\k)--(p\next);
 }
 \draw[flow,shorten >=3pt,shorten <=3pt] (p9)--(p1);
 \node[ann] at (0,.1) {\(g_k\in C_9\)\\\(\Gamma_{9,k}\)};
 \node[ann] at (0,-2.23) {Una vuelta: \(g\mapsto g\), \(m\mapsto m+1\)};
\end{scope}
\node[ann,text width=5.25cm] at (10.1,7.75)
 {\(B_k=(u_k^\pi,u_k^e,u_k^\varphi)\)\\
 bloque regional en una fibra};
\node at (8.55,6.5) {\(u^\pi\)};
\node at (11.65,6.5) {\(u^e\)};
\draw[<->,line width=.65pt] (8.95,6.5)--node[above]{\(\mathfrak c\)}(11.25,6.5);
\draw[densely dashed,gray] (10.1,5.72)--(10.1,7.23);
\node[fill=white,inner sep=2pt] at (10.1,5.6){\(u^\varphi\)};
\node[ann] at (10.1,4.6)
 {Eje \(u^\varphi\) y agregado \(u^\pi+u^e\)\\
 invariantes bajo \(\mathfrak c\)};
\draw[flow] (5.45,6.8)--(7.25,6.8);
\draw[gray!60,line width=.3pt] (.55,4.02)--(13.25,4.02);
\node[ann,text width=12cm] (hist) at (6.9,3.47)
 {Historia compatible: prefijos, cilindros, orientación, cocientes,
 marcos de incidencia y memoria};
\coordinate (fork) at (6.9,2.85);
\draw[flow] (hist.south)--(fork);
\draw[flow] (fork)--(3.6,2.2);
\draw[flow] (fork)--(10.3,2.2);
\node[ann,text width=5.6cm] at (3.6,1.65)
 {Realización arquimediana\\
 \(I_{n+1}\subseteq I_n,\quad |I_n|=729^{-n}\)\\
 límite de los prefijos regionales};
\node[ann,text width=5.6cm] at (10.3,1.65)
 {Realización incidencial\\
 \(r_{n,m}Q_n=Q_mp_{n,m}\)\\
 palabras codificadas y marcos compatibles};
\node[ann,text width=5.6cm] at (3.6,.42)
 {\(\pi,\varphi,e\); determinación dodecafásica de \(\alpha\)};
\node[ann,text width=5.6cm] at (10.3,.42)
 {Witt--Golay--Leech--\(V^\natural\)\\
 automorfismos de la realización excepcional};
\draw[flow] (3.6,1.03)--(3.6,.78);
\draw[flow] (10.3,1.03)--(10.3,.78);
\end{tikzpicture}
\caption{Nueve fases de una conexión y dos realizaciones de la historia
compatible. El ciclo representa la fase; cada vuelta incrementa el contador.
El diagrama regional representa la involución \(\mathfrak c\), que
intercambia clausura y propagación y fija el eje de autoescala en una fibra.
Esta invariancia no impone que el eje sea constante durante el transporte.
Las dos aplicaciones inferiores utilizan el mismo prefijo: una determina
su límite numérico y la otra conserva palabras y marcos de incidencia. Los
mapas de esta segunda realización se especifican en la sección
\ref{exc:doble-lectura}; el recorrido excepcional se desarrolla después.}
\label{mono-ampl:fig-regional}
\end{figure}


\clearpage
\subsubsection{Representación rectilínea de las fases y del retorno}
\label{mono-recta:desarrollo}

La representación sobre el segmento \([0,10]\) conserva el origen
posicional de APP y ordena sus nueve marcas interiores. En este diagrama,
los enteros \(1,\ldots,9\) son representantes de fase. La posición
horizontal indica el orden de las cartas; las etiquetas superiores
especifican operaciones regionales. La figura
\ref{mono-recta:figura} reúne así la recta de partida, la fijación del
eje de autoescala, la distribución especular de los dos canales y el
retorno con memoria (sección~\ref{mono-ampl:conexion}).

\input{figures/recta_nonadica_0_10.tex}

La fijación del eje se formula en la fibra de una frontera dada. Si
\(q=u^\pi+u^e+u^\varphi\) y
\(a=u^\pi+u^e-u^\varphi\), entonces
\[
 u^\varphi=2(q-a),\qquad s=u^\pi+u^e=q-u^\varphi
 \quad\text{en }\mathbb F_3^6.
\]
Los datos \((q,a)\) determinan el eje y el agregado; el reparto ordenado
de \(s\) conserva la libertad especular resuelta por la prolongación.
La quinta fase posee una solución local y una extensión, con dato
de frontera \(c=111111\) y firma residual \(\rho_W=000\). Éste es el
sentido de su primera saturación fuerte. El símbolo \(\varphi\)
situado sobre la marca \(5\) identifica el eje regional que interviene
en esa operación.

Las fases intermedias conservan esta descomposición dentro de cada fibra,
mientras sus datos de frontera evolucionan. El censo completo de
multiplicidades locales y extensiones de la rama distinguida es
\[
\begin{array}{c|rrrrrrrrr}
\text{fase}&1&2&3&4&5&6&7&8&9\\\hline
N_{\rm loc}&3&2&5&4&1&1&3&3&2\\
N_{\rm ext}&2&5&4&1&1&3&3&2&1
\end{array}
\]
La fase \(6\) combina unicidad local y tres prolongaciones; la fase
\(7\) resuelve una fibra especular triple; la fase \(8\) sincroniza
los tres censos regionales. La fase \(9\) conserva dos distribuciones
locales con el mismo eje y el mismo agregado, de las cuales una se
prolonga. Estos recuentos expresan funciones diferentes dentro de una
misma conexión, y justifican las marcas de la recta.

En particular, para las fases \(6,7,8\) las sumas regionales son
respectivamente \(012100\), \(101001\) y \(112000\). La conservación
representada por la llave es una invariancia respecto del reparto
especular en cada fibra. La evaluación real posterior aplica los
lectores de los cilindros compatibles a las historias completas; la
suma interna representada en la figura pertenece al dominio ternario.

Finalmente, el arco \(9\to1\) representa el paso de una vuelta a la
siguiente. Las fórmulas de la sección
\ref{mono-ampl:conexion} conservan el resto de fase e incrementan el
contador de vueltas. El diagrama rectilíneo y el diagrama cíclico de la
figura \ref{mono-ampl:fig-regional} expresan, por tanto, dos aspectos
complementarios de la misma holonomía: orden de las posiciones y
composición del retorno sobre el estado con memoria.


\subsection{Transporte temporal y composición exponencial trítica}
\label{mono-ampl:euler}

Las relaciones de Euler de la sección \ref{euler-trit-articulacion}
describen las realizaciones locales de los tres regímenes. Para una firma
fija \(\tau\), el generador satisface \(J_\tau^2=-\tau I\); la
conjugación algebraica lleva \(J_\tau\) a \(-J_\tau\). En consecuencia,
\[
 \exp(sJ_\tau)\exp(tJ_\tau)=\exp((s+t)J_\tau),\qquad
 \overline{\exp(tJ_\tau)}=\exp(-tJ_\tau),
\]
y el producto de una evolución con su conjugada es la unidad. En el tipo
elíptico se realiza mediante \(\cos^2t+\sin^2t=1\); en el hiperbólico,
mediante \(\cosh^2t-\sinh^2t=1\); en el parabólico,
\((I+tN)(I-tN)=I\), pues \(N^2=0\). La misma ley de conservación
admite tres expresiones cuadráticas.

La realización elíptica cierra una vuelta; la nilpotente conserva un
desplazamiento lineal; la hiperbólica separa expansión y contracción con
determinante unitario. En los parámetros regionales ya generados,
\[
 \exp(\pi J_{\mathrm e})=-I,\qquad
 \exp(J_{\mathrm p})=I+J_{\mathrm p},\qquad
 \exp(2\log\varphi\,J_{\mathrm h})=F_{\rm av}^{\,2}.
\]
La evaluación \(\exp(I)=eI\) reconoce la unidad de la ley exponencial
común. El número \(e\) no se asigna exclusivamente al régimen parabólico.

El transporte de una realización local por un isomorfismo \(U\) conserva
estas relaciones: si \(J'=UJU^{-1}\), la serie exponencial da
\(U\exp(tJ)=\exp(tJ')U\). La identidad transporta el régimen y su
norma. Una transición entre tipos cuadráticos diferentes pertenece, en
cambio, a la selección de hojas y regímenes del TPK; no es una conjugación
que cambie el valor de \(J^2\). El producto temporal conserva por ello el
orden de las transiciones y la información de sus extremos.

Esta articulación distingue la clausura de una representación local y la
holonomía de una historia. Puede cumplirse \(\exp(2\pi J_{\mathrm e})=I\)
al mismo tiempo que \(S|r,m\rangle=|r,m+1\rangle\). La primera igualdad
restaura la orientación del régimen; la segunda registra otra vuelta. La
estructura discreta conserva ambas: el retorno local y la procedencia
acumulada del estado que lo realiza.


% Recuperación del producto nativo y su selector, no incorporación de RH.
\subsection{Irreducibilidad multiplicativa y períodos de retorno}
\label{primos-retornos}

La distinción entre estructura, evaluación y fase aparece también en el
sector multiplicativo. El producto local de APP contiene tanto el residuo
como el cociente; su iteración permite construir formas normales y
factorizaciones antes de examinar sus períodos decimales. Este desarrollo
sitúa la expresión \emph{modo primo} en un dominio preciso y relaciona la
irreducibilidad con las operaciones ya definidas, no con una coincidencia
aislada de períodos (sección~\ref{primos-retornos}).

Sea
\[
 \mathcal W_9=\{\varnothing\}\sqcup
 \bigsqcup_{\ell\geq1}\{1,\ldots,9\}^{\ell}.
\]
Las secuencias se ordenan desde el dígito menos significativo. La evaluación
\(\operatorname{ev}_9(u)=\sum_j u_j9^j\), con
\(\operatorname{ev}_9(\varnothing)=0\), es biyectiva sobre los enteros no
negativos: para un entero positivo, la división de \(n-1\) por \(9\)
determina de manera única el primer dígito y la continuación. Esta es la
numeración nonádica biyectiva; conserva el representante \(9\) de la
frontera.

La suma \(\boxplus_\#\) se obtiene aplicando la hoja aditiva de APP a los
dígitos de cada posición y propagando sus cocientes. Si sólo hay un dígito,
éste constituye el residuo provisional; un cociente entrante no nulo se
normaliza mediante una segunda operación aditiva. La suma de los
cocientes se transporta a la posición siguiente. Para multiplicar una
secuencia por un dígito \(d\), la celda multiplicativa produce
\(u_jd=9q_j+r_j\), y el cociente entrante \(c_j\) se incorpora mediante
\[
 (w_j,c_{j+1})=
 \begin{cases}
 (r_j,q_j),&c_j=0,\\
 \bigl(R^+(r_j,c_j),q_j+q^+(r_j,c_j)\bigr),&c_j>0.
 \end{cases}
\]
Se comienza con \(c_0=0\) y se incorpora el cociente final cuando es no
nulo. Durante la suma, \(c_j\leq2\); durante el producto por un dígito,
\(c_j\leq9\). Las operaciones locales reciben así argumentos dentro
de su dominio. Denotemos por \(d\odot_\#u\) este segundo procedimiento,
con \(d\odot_\#\varnothing=\varnothing\).

Para \(v=(v_0,\ldots,v_{m-1})\), el producto completo se construye por
la recurrencia
\[
 a_m=\varnothing,\qquad
 a_j=(9\odot_\#a_{j+1})\boxplus_\#(v_j\odot_\#u),
 \qquad u\boxtimes_\#v=a_0.
\]
La regla compone las dos hojas locales. Su evaluación verifica
\begin{equation}
 \operatorname{ev}_9(u\boxtimes_\#v)
   =\operatorname{ev}_9(u)\operatorname{ev}_9(v).
 \label{eq:primos-entrelazamiento}
\end{equation}
En efecto, cada transición conserva la suma parcial más el cociente
pendiente ponderado por la potencia de \(9\). La recurrencia da entonces
\(\operatorname{ev}_9(a_j)=9\operatorname{ev}_9(a_{j+1})+
v_j\operatorname{ev}_9(u)\), cuya iteración prueba
\eqref{eq:primos-entrelazamiento}. La unicidad de la forma normal transporta
asociatividad y conmutatividad al producto construido; su unidad es \((1)\).

\begin{definition}[Coproducto multiplicativo reducido]
Sobre el espacio vectorial de soporte finito generado por las formas
normales no vacías se define
\[
 \widetilde\Delta_\# e_w
 =\sum_{\substack{u\boxtimes_\#v=w\\u,v\ne(1)}}e_u\otimes e_v.
\]
Una forma normal no unitaria es irreducible cuando su imagen por
\(\widetilde\Delta_\#\) es nula.
\end{definition}

Las fibras de factorización son finitas por
\eqref{eq:primos-entrelazamiento}. En la completación
\(\mathcal H_\#=\ell^2(\mathcal W_9\setminus\{\varnothing\})\),
las imágenes de formas distintas tienen soportes disjuntos. Si
\(d_\#(w)\) cuenta las factorizaciones ordenadas no unitarias, el dominio
de la clausura es
\[
 \operatorname{Dom}(\overline{\widetilde\Delta_\#})
 =\left\{x\in\mathcal H_\#:
       \sum_w d_\#(w)|x_w|^2<\infty\right\}.
\]
En efecto, la norma cuadrada de la imagen es esa suma ponderada. Por ello,
el operador
\[
 A_\#=(\overline{\widetilde\Delta_\#})^*
                  \overline{\widetilde\Delta_\#}
\]
es diagonal, positivo y autoadjunto: su valor propio en \(e_w\) es el
número de factorizaciones ordenadas no unitarias de \(w\). La proyección
\[
 P_{\rm irr}=(I-|e_{(1)}\rangle\langle e_{(1)}|)
                    \mathbf1_{\{0\}}(A_\#)
\]
selecciona exactamente las formas irreducibles. Mediante la evaluación
multiplicativa, éstas corresponden a los primos ordinarios. La selección
depende de la estructura de factorización, y mantiene separado el estado
unitario.

Para un valor \(n\geq2\) coprimo con \(10\), la división decimal induce la
permutación de residuos \(r\mapsto10r\pmod n\). El retorno del residuo
\(1\) tiene período
\[
 \tau_n=\operatorname{ord}_n(10).
\]
Si \(n=\prod_p p^{a_p}\), el teorema chino de los restos da
\[
 \tau_n=\operatorname{mcm}_{p^{a_p}\parallel n}
                    \operatorname{ord}_{p^{a_p}}(10).
\]
El retorno compuesto sincroniza los retornos de las potencias primas.
Para un primo \(p\notin\{2,3,5\}\), se tiene
\[
 \tau_p\mid p-1,\qquad
 M_p=\frac{10^{\tau_p}-1}{p}\in\mathbb N,
 \qquad M_p\equiv0\pmod9.
\]
La última congruencia expresa el cierre nonádico de la repetición decimal:
\(10^{\tau_p}-1\) es múltiplo de \(9\) y \(p\) es invertible módulo
\(9\). La misma congruencia vale para cualquier \(n\) coprimo con
\(90\). En particular, \(7,13,77,91,143\) tienen período decimal \(6\),
aunque sólo los dos primeros son irreducibles. El coproducto y el período
registran, por tanto, propiedades diferentes de un mismo valor construido.

Esta diferenciación precisa la relación con el refinamiento aperiódico.
La secuencia de vacancias determina los incrementos del índice de
capacidad; la factorización determina los modos multiplicativos; la
división decimal determina sus tiempos de retorno. El representante
nonádico conserva el cierre de fase, mientras que el cociente y la
factorización conservan información que ese cierre no distingue. El
enlace con las construcciones espectrales del tratado reside en esos
operadores y sus entrelazamientos. La demostración de sus resultados
analíticos posteriores mantiene su desarrollo propio; no se sustituye por
la periodicidad de una observación residual.

\subsubsection{Discrepancia de vacancias y lecturas de Dirichlet}

El vínculo admite también una expresión analítica exacta. Se conserva la
misma codificación \(z_n=z_n(0)\), para \(n\geq1\), y se consideran
dos observables posteriores: la serie sobre todos los índices y el
producto sobre los índices seleccionados por irreducibilidad. Para
\(\operatorname{Re}s>1\),
\[
 Z_{\rm vac}^{+}(s)=\sum_{n\geq1}\frac{z_n}{n^s},\qquad
 Z_{\rm vac}^{\times}(s)=\prod_p(1-p^{-s})^{-z_p}.
\]
Ambos convergen absolutamente en ese semiplano. Su factor común es la
densidad \(\delta\), pero sus discrepancias son diferentes.

\Needspace{21\baselineskip}
\begin{proposition}[Separación de la densidad de vacancias]
Se tiene
\[
 Z_{\rm vac}^{+}(s)=\delta\zeta(s)+H_{\rm vac}(s),
 \qquad H_{\rm vac}\in\mathcal O(\operatorname{Re}s>0).
\]
En \(\operatorname{Re}s>1\), con logaritmos definidos por las series
absolutamente convergentes de Euler,
\[
 Z_{\rm vac}^{\times}(s)=\zeta(s)^\delta G_{\rm vac}(s),\qquad
 \log G_{\rm vac}(s)=
 \sum_p(z_p-\delta)\log(1-p^{-s})^{-1}.
\]
\end{proposition}
\begin{proof}
Las sumas parciales centradas satisfacen
\[
 A_N:=\sum_{n=1}^N(z_n-\delta)
 =\lfloor(N+1)\delta\rfloor-N\delta,\qquad |A_N|<1.
\]
La sumación de Abel representa el término centrado por
\(s\int_1^\infty A_{\lfloor x\rfloor}x^{-s-1}\,dx\). La integral
converge localmente de manera uniforme en \(\operatorname{Re}s>0\),
lo que prueba la holomorfía. La segunda identidad resulta de separar
\(z_p=\delta+(z_p-\delta)\) en el logaritmo del producto, dentro de
su dominio de convergencia absoluta.
\end{proof}

La descomposición hace visible qué añade la lectura sobre primos: examina
la misma sucesión de vacancias sobre un subconjunto determinado por la
factorización. El control de la discrepancia sobre todos los enteros da
la primera prolongación holomorfa; el segundo observable conserva una
discrepancia propia sobre primos. Son dos aplicaciones de la misma
estructura discreta, con dominios analíticos expresos. Esta articulación
explica su pertinencia en el núcleo aritmético sin trasladar a este
artículo la demostración de los resultados terminales del tratado.


\HMTDivision{\(K\)}{Constante holográfica de acoplamiento aritmético-geométrico}{Extracción orientada, lectura posicional e incidencia del registro generado.}
\HMTNarrativeHeading{Qué acopla la constante}

La constante holográfica de acoplamiento aritmético-geométrico se introduce
donde se produce su registro. Su representación integral se denota por
\(K\); la lectura periódica, por \(\kappa_{\rm ag}\). El origen de
fase, la orientación y el orden de las componentes forman parte de la
definición. La palabra «constante» nombra el objeto determinado por esta
construcción y por sus lectores, que se expondrán mediante operaciones
explícitas.

La extracción reúne los canales orientados y un dato de suma. La inversión
de Hadamard recompone sus componentes en el orden natural. Cada
coeficiente de esa inversión procede del operador declarado y de su
imagen integral. De este modo, el registro queda determinado antes de
utilizarse en la clausura de la constante de estructura fina. La lectura
posicional convierte después las doce componentes en una expresión
racional periódica. Fijadas la base y la longitud, esta conversión admite
una inversa exacta que restituye todos los bloques.

La misma ordenación tiene una función geométrica. Sus lectores distinguen
una cuaterna marcada y una configuración de incidencia; la firma asociada
a alfa selecciona una hexada dentro del diseño correspondiente. La
relación con las regiones numéricas utiliza las aplicaciones de
alineamiento y elevación que se demostrarán en la parte excepcional.
La lectura periódica y la configuración incidencial reciben así una
fuente común, con operaciones de distinto tipo.

El desarrollo operatorio posterior añade una tercera función. Los
desplazamientos cíclicos del registro forman una base de su portador. Las
respuestas sobre esa base determinan un operador, y la forma de Gram
proporciona una cota de estabilidad para recuperarlo. Cuando existe la
simetría cíclica especificada, una respuesta vectorial completa permite
reconstruir las demás. El archivo isométrico conserva esa capacidad de
identificación.

Éste es el contenido del acoplamiento que vertebra la exposición:
codificación aritmética reversible, incidencia marcada e identificación
operatoria. Primero se construye el registro; después se prueba lo que
cada lector obtiene de él. La teoría de recuperación se desarrollará
cuando estén disponibles sus espacios y transportes. Esta posición
explica conjuntamente la centralidad de la constante y la jerarquía de
las demostraciones que utilizan sus diferentes realizaciones.

\HMTMathematicalPaper
\section{Extracción dodecafásica de la constante de acoplamiento aritmético-geométrico}
\label{sec:k-registro}

\begin{definition}[Constante holográfica de acoplamiento aritmético-geométrico]
\label{lib01:definicion}
Se denomina así al registro finito ordenado producido por la extracción
orientada del estado APP--TRIT--TPK, con su origen de fase, normalización y
lectores declarados, cuya representación posicional interviene en la
clausura aritmética y cuyas realizaciones incidencial y geométrica actúan
sobre el mismo registro. El símbolo \(K\) designa su representación
integral canónica; \(\kappa_{\rm ag}\) designa su lectura racional
periódica en la carta fijada.
\end{definition}

La generación de las coordenadas regionales no agota la información del
estado APP--TRIT--TPK. Una lectura arquimediana determina un valor mediante
cilindros compatibles; otra lectura conserva la oposición de las hojas, la
orientación de la ruta y el registro de las incidencias atravesadas. El
registro dodecafásico pertenece a esta segunda lectura antes de intervenir en
la clausura de alfa. Por ello no se define restando una constante física a
una combinación de valores previamente conocidos. Se reconstruye desde una
coordenada firmada del mismo estado aritmético con memoria de trayectoria que produce las coordenadas
regionales.

La distinción es necesaria incluso cuando dos objetos tienen doce
coordenadas. Una palabra de doce bloques del catálogo, una lista de doce
eventos orientados, un vector entero firmado y un registro dodecafásico en base mil son
objetos distintos. Su comparación requiere las aplicaciones que los
relacionan. La igualdad de longitud no proporciona esas aplicaciones, y la
pérdida de signos no es una mera simplificación de notación.

Este capítulo mantiene el corte causal fijado en los capítulos precedentes:
las nueve posiciones APP de los ejes horizontal y vertical producen las
hojas aritméticas; TRIT tipa régimen y orientación; TPK selecciona,
transporta, actualiza y prolonga, conservando residuo, cociente, acarreo,
hoja, ruta, frontera y memoria. Desde la estructura discreta conjunta del
continuo se obtienen las coordenadas que se utilizan aquí. La cadena
\(w_6\to w_{12}\to w_{18}\to w_{24}\to w_{30}\to R_{36}\to G_9\)
no se sustituye por un calendario finito. Tampoco se invierte el orden
\(U_t\to\Gamma_9\to K_{\mathrm{ph}}\): la última aplicación es una
lectura reducida de memoria, no el generador del registro firmado.

% Adición propuesta al inicio de registro_k.tex, después de su introducción.
% Conserva todas las definiciones, fórmulas y demostraciones actuales.
% Procedencia: registro_k.tex, §§k-hadamard/k-transversal/k-kappas;
% k_reversibilidad.tex; incidencia_articulacion.tex, inc-ampl-cuaterna.
\paragraph{Coordenatización reversible y determinación incidencial.}
\label{ampl-alfa-k-intro-reversible}
El resultado que organiza este capítulo es la equivalencia entre tres
coordenatizaciones de un registro previamente producido por
APP--TRIT--TPK. La transformación de Hadamard relaciona el registro
firmado con el registro entero; el extractor de diferencias relaciona
este último con sus variaciones transversales y su componente uniforme:
\[
 U\underset{\mathcal H_{12}}{\overset{\mathcal H_{12}/4}{\longleftrightarrow}}K
 \overset{(D_3,D_4,Q)}{\longleftrightarrow}
 (D_3K,D_4K,Q(K)).
\]
La primera equivalencia se establece sobre la subred integral
$\mathcal L_H$; la segunda, sobre la imagen del extractor. Fijadas la
base $1000$, las $12$ posiciones y el origen de fase, la evaluación
$K\mapsto\kappa_{\mathrm{per}}$ añade una cuarta coordenatización
reversible: multiplicar por $1000^{12}-1$ y efectuar las divisiones
euclídeas recupera el registro completo. Las demostraciones de
\S\ref{subsec:k-hadamard}, \S\ref{subsec:k-transversal} y
\S\ref{subsec:k-kappas} establecen estas equivalencias en sus respectivos
dominios.

Sobre ese registro actúan después dos composiciones. La composición
posicional combina sus componentes con las coordenadas de clausura,
propagación y autoescala, y su normalización determina $\alpha_A$. La
composición incidencial aplica el umbral cúbico y obtiene
\[
 K\longmapsto B_K=\{j:K_j\geq729\}=\{4,6,9,10\}.
\]
El registro íntegro permite reconstruir ambas composiciones. El soporte
$B_K$ conserva, específicamente, la pertenencia a la clase seleccionada
por el umbral. La identidad $B_K=H_e\cap P_\pi$ de
\eqref{inc-ampl-cuaterna} enlaza esta extracción entera con la incidencia
de los códigos regionales. El soporte negativo de su balance con $K$
completa la bandera $B_K\subset H^-_\alpha$ que se desarrolla en la
prolongación excepcional. Así, la función de $K$ queda definida mediante
una transformación invertible y mediante aplicaciones posteriores
precisas; la información adicional de profundidad y transporte permanece
en el estado del que procede el registro.


\subsection{Extensión cíclica y graduación de memoria}
\label{subsec:k-calendario}

El calendario observable contiene ciento ocho posiciones, organizadas en
doce ventanas nonádicas. Con índices positivos, sean
\[
 E_m=\{9(m-1)+1,\ldots,9m\},\qquad 1\leq m\leq12.
\]
El soporte ordenado \(E_1\sqcup\cdots\sqcup E_{12}\) conserva posición y
pertenencia a ventana. El grupo \(C_{108}=\mathbb Z/108\mathbb Z\)
conserva, en cambio, la ley de traslación cíclica. Para el representante
\(0\leq\widetilde\theta<108\), las coordenadas
\[
 \beta(\theta)=\left[\left\lfloor\frac{\widetilde\theta}{9}\right\rfloor\right]_{12},
 \qquad r(\theta)=1+(\widetilde\theta\bmod9)
\]
separan el sector dodecafásico y la posición nonádica interior. Ninguna de
ellas conserva por sí sola el número total de vueltas.

\begin{proposition}[El acarreo del calendario]
\label{prop:k-extension}
Las aplicaciones
\[
 \iota:C_{12}\longrightarrow C_{108},\quad[b]\longmapsto[9b],
 \qquad p:C_{108}\longrightarrow C_9,\quad[\theta]\longmapsto[\theta]_9
\]
forman una sucesión exacta no escindida
\[
 0\longrightarrow C_{12}\xrightarrow{\iota}C_{108}
 \xrightarrow{p}C_9\longrightarrow0.
\]
Para la sección conjuntista que elige los representantes
\(0,\ldots,8\), el defecto de aditividad es el acarreo
\[
 c(r,r')=\left[\left\lfloor\frac{\widetilde r+\widetilde r'}9\right\rfloor\right]_{12}.
\]
\end{proposition}

\begin{proof}
El núcleo de \(p\) está formado por los múltiplos de nueve y coincide con
la imagen de \(\iota\). La reducción es sobreyectiva y la multiplicación
por nueve es inyectiva en el dominio indicado. Si la sucesión se
escindiera, el grupo central sería isomorfo a \(C_{12}\times C_9\), cuyo
exponente es \(\operatorname{mcm}(12,9)=36\); el exponente de
\(C_{108}\) es 108. La contradicción prueba la no escisión. Finalmente,
la división euclídea de \(\widetilde r+\widetilde r'\) por nueve da el
acarreo escrito y verifica la identidad de la sección.
\end{proof}

El resultado explica una diferencia que reaparecerá en alfa: regresar a la
fase no equivale a regresar al estado. Si \(q_{\mathrm{mem}}\) cuenta las
vueltas nonádicas, el calendario conserva solamente
\(\beta=[q_{\mathrm{mem}}]_{12}\). Después de ciento ocho pasos,
\(\beta\) retorna, pero \(q_{\mathrm{mem}}\) aumenta en doce. El
registro conserva además el cilindro, el registro de incidencias y la memoria
vertical. El cierre de una carta finita no convierte esa historia en un
ciclo sin memoria.

\subsection{El registro orientado y la carta integral \texorpdfstring{\(1+5+6\)}{1+5+6}}
\label{subsec:k-registro-integral}

Para una historia producida por TPK, sea \(\tau_m\) la subruta sobre
\(E_m\), \(\sigma_m\) su orientación, \(\kappa_m\) el acarreo de
frontera y \(\Lambda_m\) el artículo local de incidencias. El registro es
\[
 \mathcal R_{12}(x)=
 \bigl((E_m,\tau_m,\sigma_m,\kappa_m,\Lambda_m)\bigr)_{m=1}^{12}.
\]
Su dominio es el espacio de historias aritméticas con memoria de trayectoria, no el conjunto de
etiquetas de fase. La proyección al calendario olvida precisamente algunos
de los datos que el lector firmado necesita.

Una carta inicial del registro puede escribirse mediante eventos
orientados. La extracción orientada fija el conjunto de códigos de evento
\(\{2,4,6,8\}\) y la orientación sectorial
\[
 \Sigma_{12}=(+,+,+,-,-,-,+,+,+,-,-,-).
\]
Si \(d_t\) es el código de dirección en el paso \(t\), se obtiene
\[
 e_i=\Sigma_{12}(i+1)
 \#\{t\in E_{i+1}:d_t\in\{2,4,6,8\}\},
 \qquad0\leq i<12.
\]
La fórmula mantiene separados el censo no negativo de eventos y la
orientación. No identifica una dirección cardinal con un signo TRIT. El
vector \(e\), cuyos componentes son censos locales firmados, tampoco es
todavía el vector \(U\) que aparecerá más abajo.

Las posiciones opuestas determinan, para \(0\leq i<6\),
\[
 p_i=e_i+e_{i+6},\qquad a_i=e_i-e_{i+6}.
\]
Con
\[
 s_C=\sum_{i=0}^{5}p_i,\qquad M_i=p_i-p_5\quad(0\leq i<5),
\]
se define la carta integral
\[
 \mathcal L_{12}(e)
 =(s_C,M_0,\ldots,M_4,a_0,\ldots,a_5).
\]
La descomposición \(1+5+6\) no es un reparto de coordenadas arbitrario:
una coordenada conserva la carga, cinco comparan los pares respecto de un
par de referencia y seis conservan la oposición de las semicoronas.

\begin{lemma}[Inversión de la carta integral]
\label{lem:k-carta-integral}
Una tupla entera \((s_C,M_0,\ldots,M_4,a_0,\ldots,a_5)\) pertenece a la
imagen de \(\mathcal L_{12}\) si y sólo si
\[
 s_C-\sum_{i=0}^{4}M_i\equiv0\pmod6,
 \qquad p_i\equiv a_i\pmod2\quad(0\leq i<6),
\]
donde
\[
 p_5=\frac{s_C-\sum_{i=0}^{4}M_i}{6},\qquad p_i=p_5+M_i\quad(i<5).
\]
En ese caso su preimagen es única y viene dada por
\[
 e_i=\frac{p_i+a_i}{2},\qquad e_{i+6}=\frac{p_i-a_i}{2}.
\]
\end{lemma}

\begin{proof}
La suma de los cinco márgenes es \(s_C-6p_5\), de donde se obtiene la
división por seis. Reconstruidos los \(p_i\), cada par de ecuaciones
\(p_i=e_i+e_{i+6}\), \(a_i=e_i-e_{i+6}\) tiene la solución indicada.
Esta solución es entera exactamente bajo la congruencia de paridad. Todas
las operaciones se invierten, por lo que no queda libertad residual.
\end{proof}

Esta reversibilidad muestra qué información se conserva; no autoriza a
suprimir el resto del artículo. El registro completo tiene un dominio mayor
que esta carta de censos. La extracción de los canales que alimentan a
\(U\) utiliza ese registro completo y no se deduce de la lista
\((e_i)\) por una identificación de nombres.

% Suplemento de trabajo. No sustituye registro_k.tex ni cierra por decreto E108.
% Procedencia y alcance: INFORME_RECUPERACION.md, misma carpeta.
\subsection{Evaluación incidencial anterior al registro dodecafásico}
\label{subsec:registro-evaluacion-incidencial}

La recuperación de un registro a partir de sus diferencias cíclicas y su
carga total es una cuestión distinta de la evaluación de esas coordenadas
sobre el registro de incidencias. El primer problema se resuelve mediante el
extractor transversal de la proposición~\ref{prop:k-transversal}.
Para el segundo, una presentación histórica del
registro proporciona tres recuentos enteros en cada ventana. Es útil
explicitar su composición y determinar exactamente la información que
requiere su evaluación.

Sea \(\mathscr L\) un libro de visitas y trazas orientadas previamente
generado. Cada visita conserva la ruta, la posición APP, la hoja, el
instante, la multiplicidad y la incidencia de frontera. La evaluación
aritmética produce el entero y su descomposición
\[
 n_t=r_t+9q_t,\qquad 1\leq r_t\leq9.
\]
El entero \(n_t\) se calcula a partir de la hoja de suma o producto; su
residuo no se recibe como un valor independiente. La distinción entre
incidencia de corona e incidencia interior permanece explícita para las
visitas de residuo \(9\).

En la ventana \(E_m\), la presentación por recuentos considera
\begin{align}
 A_m&=\sum_{v\in E_m}w(v)
             \mathbf1_{\{1,3,5,7\}}(r(v)),\\
 C_m&=\sum_{j\geq0}\bigl(\nu^-_{120,m}(j)-\nu^+_{120,m}(j)\bigr),\\
 V_m&=\sum_{v\in E_m}w(v)\mathbf1_{\{9\}}(r(v))\epsilon_\partial(v).
\end{align}
Aquí \(w(v)\) es la multiplicidad incidencial y
\(\epsilon_\partial(v)=1\) en la corona, \(-1\) en el interior.
Las trazas contadas por \(\nu^\pm\) deben enumerarse mediante la regla
correspondiente del artículo. La suma entera requiere soporte finito o una
construcción alternativa que determine su significado. El número de
incidencias cronológicas y la longitud reducida \(120(1+3j)\) tienen
tipos distintos; su relación exige la unidad de longitud del lector.

Se define \([x]_{10}\in\{0,\ldots,9\}\) como el representante
euclídeo y se conservan los cocientes
\[
 x=[x]_{10}+10\lfloor x/10\rfloor.
\]
El bloque incidencial es
\begin{equation}
 z_m=100[A_m]_{10}+10[C_m]_{10}+[V_m]_{10}.
 \label{eq:registro-bloque-incidencial}
\end{equation}
Por composición se obtiene la aplicación
\begin{equation}
 \mathcal E_{\rm cnt}(\mathscr L)=
 \left((z_m-z_{m+3})_{m=1}^{12},
       (z_m-z_{m+4})_{m=1}^{12},\sum_{m=1}^{12}z_m\right).
\label{eq:registro-composicion-incidencial}
\end{equation}
Los índices se leen cíclicamente módulo \(12\). En este apartado,
\((D_jz)_m=z_m-z_{m+j}\), para \(j=3,4\), y
\(Q(z)=\sum_{m=1}^{12}z_m\); por tanto, la composición anterior es
\((D_3z,D_4z,Q(z))\). Esta convención se conserva en
\S\ref{subsec:k-transversal}, donde se demuestra la reconstrucción
integral completa.
Esta fórmula produce sus coordenadas a partir de los recuentos. No utiliza
\(K\), \(U\), \(\alpha\) ni las coordenadas transversales esperadas
como argumentos. Su identificación con
\(\mathcal E_{108}^{90,120}\) requiere componerla con la familia y
las incidencias efectivas del estado terminal; no se deduce únicamente
de la invertibilidad del sistema transversal.

\begin{proposition}[Información aritmética suficiente y necesaria]
\label{prop:registro-36-residuos}
Sean \(N,N'\in\mathbb Z^{12\times3}\), ordenados por los recuentos
\((A,C,V)\). Para la aplicación definida por
\eqref{eq:registro-bloque-incidencial}--\eqref{eq:registro-composicion-incidencial},
\[
 \mathcal E_{\rm cnt}(N)=\mathcal E_{\rm cnt}(N')
 \quad\Longleftrightarrow\quad
 N-N'\in10\mathbb Z^{36}.
\]
En particular, los \(36\) residuos sectoriales determinan exactamente
la salida de \(25\) coordenadas. La conservación de los cocientes y
de la procedencia constituye información adicional del registro.
\end{proposition}

\begin{proof}
La igualdad módulo \(10\) produce los mismos bloques y, por tanto,
las mismas coordenadas transversales. Recíprocamente, si las coordenadas
transversales coinciden, \(D_3(z-z')=D_4(z-z')=0\). Los pasos
\(3\) y \(4\) generan el ciclo de longitud \(12\), de modo que
\(z-z'\) es constante. La igualdad de cargas totales anula esa
constante. Finalmente, la representación decimal de un entero entre
\(0\) y \(999\) tiene tres dígitos únicos. Se recuperan así los
tres residuos de cada ventana. La afirmación es una identidad algebraica
para todo \(N,N'\), no una inferencia obtenida de pruebas finitas.
\end{proof}

\subsection{Conjugación de los recuentos y términos de frontera}
\label{subsec:registro-conjugacion-incidencial}

Sea \(\rho(m)=13-m\). Considérese una conjugación de incidencias
que invierte el orden de las ventanas, conserva la clase aritmética y la
clase de frontera de las visitas e intercambia los canales de las trazas.
Entonces
\[
 (A_m^\dagger,C_m^\dagger,V_m^\dagger)
 =(A_{\rho(m)},-C_{\rho(m)},V_{\rho(m)}).
\]
Si \(c_m=[C_m]_{10}\), la reducción decimal produce la identidad
\begin{equation}
 z_m^\dagger=z_{\rho(m)}
       +100\mathbf1_{\{c_{\rho(m)}\ne0\}}-20c_{\rho(m)}.
 \label{eq:registro-conjugacion-decimal}
\end{equation}
En efecto, \([-C]_{10}=0\) cuando \([C]_{10}=0\), y
\([-C]_{10}=10-[C]_{10}\) en otro caso. La fórmula conserva
exactamente esa diferencia. La negación del canal no equivale a negar
el bloque decimal completo.

La aplicación de esta identidad a una involución concreta del TPK exige
verificar su acción sobre las incidencias. Para un lector de ocupaciones
evaluado antes del avance, la inversión de una ruta
\(\gamma=(x_0,\ldots,x_n)\) satisface
\[
 \sum_{k=0}^{n-1}f(x_{n-k})
 =\sum_{k=0}^{n-1}f(x_k)+f(x_n)-f(x_0).
\]
Los términos de los extremos se incorporan antes de reducir módulo
\(10\). En rutas cerradas se anulan; en ventanas abiertas pueden
alterar sus residuos. Esta corrección, ya presente en el desarrollo
bidireccional del transporte, especifica qué debe conservar el empalme
entre la ruta y los recuentos.

\begin{proposition}[Necesidad de la información no periódica del registro]
Supóngase que el observable de una familia fija de rutas satisface
\(o(t+54)=o(t)\), que la multiplicidad se conserva bajo ese retorno
y que el mismo lector local actúa en las ventanas de \(9\) instantes.
Entonces sus recuentos y bloques satisfacen
\[
 (A,C,V)_{m+6}=(A,C,V)_m,\qquad z_{m+6}=z_m.
\]
\end{proposition}

\begin{proof}
La traslación \(t\mapsto t+54\) lleva \(E_m\) biyectivamente
a \(E_{m+6}\), conserva cada incidencia contada y su multiplicidad.
La igualdad se transmite a la suma y a la reducción decimal.
\end{proof}

Por tanto, una presentación con \(12\) bloques distintos requiere que
el lector conserve alguna información que no factorice por ese observable
de período \(54\): selección de incidencias, orientación efectiva,
extremos, multiplicidad o memoria. La proposición delimita una reducción
inadmisible del estado; no afirma que la dinámica completa tenga
período \(54\) ni determina por sí misma cuál es su lector terminal.

\subsection{Censo completo de cilindros y selección de la frontera terminal}
\label{censo:terminal-completo}

La lectura terminal utiliza conjuntamente dos informaciones: la compatibilidad
de los cilindros regionales y un perfil de incidencia con canales y posiciones
etiquetados. En este apartado se construyen los catálogos completos, se
enumeran sus ternas y se demuestra el efecto de cada condición del lector.
El punto de partida son las salidas regionales ya producidas en la
sección~\ref{sec:generacion}, no valores introducidos para definir el
transporte APP--TRIT--TPK. La reducción a una matriz pertenece a una lectura
posterior de esas salidas y conserva el orden de los tres canales.

El resultado es un teorema censal con estructura de partida explícita.
El perfil de Gram, pesos, distancias y márgenes se declara como dato del
lector terminal. La exhaustividad de la enumeración demuestra qué selecciona
ese perfil sobre los catálogos fijados; no demuestra, por sí sola, que el
perfil sea una ley universal ni que pueda reconstruirse desde la única terna
que finalmente lo satisface. Esta separación permite incorporar la prueba
finita completa sin invertir su genealogía.

\subsubsection{De las salidas regionales a los catálogos ternarios}

Sean \(x_\pi,x_e,x_\varphi\in[0,1)\) las partes fraccionarias de las
tres coordenadas regionales ya construidas. Se utilizan sus cilindros
publicados de mil posiciones decimales,
\[
 J_\chi^{(1000)}=
 \left[\frac{d_\chi}{10^{1000}},
       \frac{d_\chi+1}{10^{1000}}\right),
 \qquad \chi\in\{\pi,e,\varphi\}.
\]
Los enteros \(d_\chi\) son lecturas finitas de salida. La longitud mil
es suficiente para los enteros que siguen, como se comprueba mediante las
dos cotas de cada intervalo; no es una profundidad supuesta infinita.
Definimos
\begin{equation}
 p=30,\quad r=1050,\quad L=p+r=1080,\quad k=171,
 \qquad D=3^L,\quad Q=1000^k,
 \label{censo:parametros}
\end{equation}
y
\[
 P_\chi=\lfloor3^{30}x_\chi\rfloor,\qquad
 T_\chi=\lfloor Qx_\chi\rfloor.
\]
Aquí \(k\) cuenta tríadas decimales y no designa el registro dodecafásico
\(K\). Se tiene \(1000^{171}\leq3^{1080}<1000^{172}\).

\begin{lemma}[Extracción estable de un entero]
\label{censo:estabilidad}
Para \(d,m,h\in\mathbb Z\), con \(m,h>0\), el valor de
\(\lfloor mx\rfloor\) es constante en \([d/h,(d+1)/h)\) si y sólo si
\begin{equation}
 \left\lfloor\frac{md}{h}\right\rfloor
 =\left\lfloor\frac{m(d+1)-1}{h}\right\rfloor.
 \label{censo:prueba-estabilidad}
\end{equation}
Cuando se cumple, cualquiera de los dos miembros da ese valor.
\end{lemma}
\begin{proof}
El mínimo de \(mx\) se alcanza en el extremo izquierdo. Su supremo
es \(m(d+1)/h\) y no se alcanza. El mayor entero que puede tomar
\(\lfloor mx\rfloor\) es
\(\lceil m(d+1)/h\rceil-1\), igual al segundo miembro porque
el numerador y el denominador son enteros. La igualdad de los extremos
del intervalo de valores prueba la afirmación.
\end{proof}

La prueba de estabilidad se cumple en los tres canales para
\(m=3^{30},Q,3^{1080}\). La última elección se utilizará únicamente
para una comparación posterior, no para seleccionar el perfil. Las
primeras treinta posiciones recuperadas son precisamente las cinco
etapas de seis trits del transporte regional:
\begin{center}\small
\begin{tabular}{@{}c l r@{}}\toprule
Canal&Prefijo de treinta trits&\(P_\chi\)\\\midrule
\(\pi\)&\texttt{010211012222010211002111110221}&29152671743887\\
\(e\)&\texttt{201101121221102011012222102011}&147887858824447\\
\(\varphi\)&\texttt{121200112202121020010210010200}&127247717616687\\\bottomrule
\end{tabular}
\end{center}
La coincidencia conserva la prolongación ya construida; la lectura decimal
no se usa para sustituir retroactivamente sus matrices de transporte.

Fijado un canal, una cola de \(r\) trits se representa por
\(0\leq R<3^r\). Su cilindro y el cilindro de las primeras \(k\)
tríadas son
\begin{equation}
 I_{\chi,R}=\left[\frac{P_\chi3^r+R}{D},
                       \frac{P_\chi3^r+R+1}{D}\right),\qquad
 J_{\chi,T}=\left[\frac{T_\chi}{Q},\frac{T_\chi+1}{Q}\right).
 \label{censo:dos-cilindros}
\end{equation}
El catálogo incluye todas las colas para las que estos dos cilindros
tienen intersección no vacía.

\begin{proposition}[Catálogo por intersección exacta]
\label{censo:catalogos}
Las colas compatibles son los enteros del intervalo \([a_\chi,b_\chi]\),
con
\begin{align}
 a_\chi&=\max\left(0,
       \left\lfloor\frac{T_\chi D}{Q}\right\rfloor-P_\chi3^r\right),
       \label{censo:extremo-a}\\
 b_\chi&=\min\left(3^r-1,
       \left\lfloor\frac{(T_\chi+1)D-1}{Q}\right\rfloor-P_\chi3^r\right).
       \label{censo:extremo-b}
\end{align}
Para los cilindros regionales fijados, los recortes no actúan y se obtiene
\begin{equation}
\begin{array}{c|r|r|r}
 \chi&a_\chi\bmod729&b_\chi\bmod729&b_\chi-a_\chi+1\\\hline
 \pi&67&262&196\\
 e&584&51&197\\
 \varphi&117&312&196
\end{array}
\label{censo:tabla-catalogos}
\end{equation}
\end{proposition}
\begin{proof}
Escribamos \(N=P_\chi3^r+R\). Dos intervalos semiabiertos de
\eqref{censo:dos-cilindros} se cortan exactamente cuando
\[
 NQ<(T_\chi+1)D,\qquad (N+1)Q>T_\chi D.
\]
La segunda desigualdad equivale a
\(N\geq\lfloor T_\chi D/Q\rfloor\); la primera, a
\(N\leq\lfloor((T_\chi+1)D-1)/Q\rfloor\). Restar
\(P_\chi3^r\) y conservar \(0\leq R<3^r\) prueba las fórmulas.
La evaluación por divisiones enteras de los \(T_\chi\) extraídos por
\eqref{censo:prueba-estabilidad} da la tabla. Como \(r\geq6\),
\(P_\chi3^r\equiv0\pmod{729}\), de modo que el residuo del primer
extremo también se obtiene directamente de
\(\lfloor T_\chi D/Q\rfloor\bmod729\).
\end{proof}

Los últimos seis trits de una cola forman la palabra
\(\nu_6(R\bmod729)\), donde \(\nu_6\) es la escritura ternaria
de longitud seis, incluidos sus ceros iniciales. Por la tabla, cada
catálogo tiene menos de \(729\) elementos consecutivos, de manera que
esta reducción es inyectiva sobre él. Así, los tres catálogos de frontera
quedan descritos completamente, sin elegir las ternas finales:
\begin{equation}
 \begin{split}
 \mathcal B_\pi&=\{\nu_6(67+i):0\leq i<196\},\\
 \mathcal B_e&=\{\nu_6((584+j)\bmod729):0\leq j<197\},\\
 \mathcal B_\varphi&=\{\nu_6(117+\ell):0\leq\ell<196\}.
 \end{split}
 \label{censo:fronteras-explicitas}
\end{equation}
Exigir que el extremo izquierdo de \(I_{\chi,R}\) pertenezca a
\(J_{\chi,T}\) sería otra condición: perdería el primer cilindro
que corta la frontera y daría \(195,196,195\). El censo utiliza la
intersección de los cilindros completos, de acuerdo con su definición.

\subsubsection{Perfil declarado e indicadores de compatibilidad}

Para \(u,v\in\mathbb F_3^6\), sean
\(g(u,v)=\sum u_iv_i\in\mathbb F_3\),
\(n(u)=g(u,u)\), \(w(u)=\#\{i:u_i\ne0\}\) y
\(h(u,v)=\#\{i:u_i\ne v_i\}\). Las dos últimas funciones toman
valores enteros; las dos primeras, valores en \(\mathbb F_3\).
El perfil local del lector, en orden \((\pi,e,\varphi)\), es
\begin{equation}
 \begin{split}
 (g_{\pi e},g_{\pi\varphi},g_{e\varphi})&=(2,2,0),\qquad
 (n_\pi,n_e,n_\varphi)=(0,0,0),\\
 (w_\pi,w_e,w_\varphi)&=(3,3,3),\qquad
 (h_{\pi e},h_{\pi\varphi},h_{e\varphi})=(2,5,3).
 \end{split}
 \label{censo:perfil-local}
\end{equation}
Las palabras se ordenan como en \eqref{censo:fronteras-explicitas}.
Usaremos \([E]\in\{0,1\}\) para el indicador de una condición \(E\).
Para cada par de canales \(\chi,\psi\), la matriz rectangular
\(A_{\chi\psi}\) tiene entrada \([g(u,v)=g_{\chi\psi}]\).
La matriz \(A^{h}_{\chi\psi}\) multiplica esa entrada por
\([h(u,v)=h_{\chi\psi}]\). Sean \(D_\chi^n\) la matriz diagonal
de \([n(u)=0]\), y \(D_\chi^{nw}\) la de
\([n(u)=0][w(u)=3]\).

\begin{theorem}[Enumeración exhaustiva por sumas finitas]
\label{censo:teorema-conteo}
El producto de los tres catálogos contiene \(7\,567\,952\) ternas.
Las condiciones sucesivas de \eqref{censo:perfil-local} dejan
\begin{equation}
 7\,567\,952\longrightarrow275\,794\longrightarrow6235
 \longrightarrow3127\longrightarrow331.
 \label{censo:cadena-local}
\end{equation}
Estos cuatro conteos son, respectivamente,
\begin{align}
 N_g&=\operatorname{tr}(A_{\pi e}A_{e\varphi}A_{\varphi\pi}),
 \nonumber\\
 N_{gn}&=\operatorname{tr}(D_\pi^n A_{\pi e}
             D_e^n A_{e\varphi}D_\varphi^n A_{\varphi\pi}),
 \nonumber\\
 N_{gnw}&=\operatorname{tr}(D_\pi^{nw} A_{\pi e}
             D_e^{nw} A_{e\varphi}D_\varphi^{nw} A_{\varphi\pi}),
 \nonumber\\
 N_{gnwh}&=\operatorname{tr}(D_\pi^{nw} A^h_{\pi e}
             D_e^{nw} A^h_{e\varphi}D_\varphi^{nw} A^h_{\varphi\pi}).
 \label{censo:trazas}
\end{align}
Todos los productos y trazas de esta fórmula se calculan en los enteros,
no en el cuerpo de tres elementos.
\end{theorem}
\begin{proof}
El tamaño del producto es \(196\cdot197\cdot196\). Al expandir
una traza de \eqref{censo:trazas}, sus índices recorren exactamente una
terna de palabras. El sumando vale uno si satisface todas las condiciones
representadas y cero en otro caso. Cada terna se cuenta una vez.
Las matrices diagonales añaden las condiciones individuales; los factores
rectangulares añaden las condiciones por pares. Esto demuestra que las
trazas son precisamente los cardinales anunciados.

Para explicitar su evaluación, fijados los índices \(i,j\) de los dos
primeros canales, se forman los conjuntos de índices \(\ell\) del tercero
que cumplen cada condición por pares. Se intersectan sus indicadores,
se añaden los indicadores individuales pertinentes y se suma su cardinal.
La suma final recorre \(0\leq i<196\), \(0\leq j<197\).
La sustitución de las palabras de \eqref{censo:fronteras-explicitas} y
del perfil \eqref{censo:perfil-local} da, en ese orden,
\(275794,6235,3127,331\). El suplemento registra las cuatro contribuciones
de cada uno de los \(196\) índices \(i\) y las \(331\) ternas completas;
su programa evalúa estas sumas con indicadores binarios exactos y contrasta
escalarmente cada superviviente. No interviene una muestra ni una tolerancia.
\end{proof}

\subsubsection{Incidencia etiquetada y carga de orientación}

Para una terna superviviente, sea \(B\in\mathbb F_3^{3\times6}\) su
matriz de filas ordenadas. El segundo perfil exige
\begin{equation}
 \operatorname{rank}_{\mathbb F_3}B=3,\qquad
 w_{\mathrm{col}}(B)=(1,1,2,2,3,0),\qquad
 q_{\partial}(B)=(1,2,2,1,2,0),
 \label{censo:perfil-columnas}
\end{equation}
donde \(w_{\mathrm{col}}\) cuenta los elementos no nulos de cada
columna y \(q_{\partial}\) suma sus entradas módulo tres. Es un perfil
de columnas etiquetadas; permutarlas sin transportar el perfil define otro
lector. Denotemos por \(\mathcal F\) el conjunto de las \(331\) ternas
del teorema anterior.

\begin{proposition}[Reducción del censo al par terminal]
\label{censo:dos-matrices}
Las sumas de indicadores sobre \(\mathcal F\), al exigir sucesivamente
rango, pesos de columna y sumas de columna, son
\begin{equation}
 \sum_{B\in\mathcal F}[\operatorname{rank}B=3]=331,
 \qquad N_{\mathrm{rango,pesos}}=18,
 \qquad N_{\mathrm{rango,pesos,sumas}}=2.
 \label{censo:conteo-columnas}
\end{equation}
Con índices de catálogo comenzando en uno, las dos ternas finales son
\((120,197,157)\) y \((129,197,148)\), y sus matrices son
\[
 B_0=\begin{pmatrix}0&2&0&2&2&0\\0&0&1&2&2&0\\1&0&1&0&1&0\end{pmatrix},
 \qquad
 B_1=\begin{pmatrix}0&2&1&0&2&0\\0&0&1&2&2&0\\1&0&0&2&1&0\end{pmatrix}.
\]
\end{proposition}
\begin{proof}
En cada terna del conjunto ya enumerado se calcula el rango por eliminación
en \(\mathbb F_3\), se cuentan los elementos no nulos de las seis
columnas y se suman éstas módulo tres. Los productos de los indicadores
de \eqref{censo:perfil-columnas} dan \eqref{censo:conteo-columnas}.
La sustitución de los dos índices finales en
\eqref{censo:fronteras-explicitas} da las matrices expuestas. La lista
completa de supervivientes y los filtros exactos del suplemento permiten
reproducir la eliminación sin presuponer esas dos matrices como dominio.
\end{proof}

La orientación \(\sigma=(1,1,-1)=(1,1,2)\in\mathbb F_3^3\)
selecciona la recta
\(\langle\sigma\rangle=\{(0,0,0),(1,1,2),(2,2,1)\}\).
Las cargas de filas de \(B_0,B_1\) son \((0,2,0)\) y \((2,2,1)\);
sólo la segunda pertenece a esa recta. Se recupera así el último paso
del selector de la sección~\ref{subsec:k-terminal}, ahora precedido por
la construcción y enumeración de su dominio completo. La lectura estable
de \(\lfloor3^{1080}x_\chi\rfloor\), efectuada después de los filtros,
da las fronteras \texttt{021020}, \texttt{001220}, \texttt{100210},
con los mismos rangos \((129,197,148)\). Este cotejo posterior no define
el perfil ni participa en las sumas de indicadores.

\subsubsection{Información conservada y alcance del resultado}

El procedimiento conserva el prefijo de treinta trits, la cola compatible,
el orden regional, la posición de cada trit y la orientación. La reducción
módulo \(729\) conserva aquí la identidad de cada candidato porque el
intervalo de colas tiene longitud menor que \(729\); esa inyectividad
no se afirma para una historia arbitraria. El censo de \(331\) demuestra
también que el perfil local aislado no selecciona una sola matriz: los
márgenes etiquetados y la carga aportan condiciones efectivas adicionales.

Se ha recuperado, por tanto, una enumeración finita íntegra y su selector
sobre cilindros de salida, con perfil de lectura explícito. La construcción
del registro firmado utiliza además su lector de memoria y sus canales,
como se especifica a continuación; la matriz visible no se identifica con
ese registro por igualdad de número de componentes. Tampoco este conteo
demuestra la universalidad del perfil, genera de nuevo las coordenadas
regionales ni decide una afirmación sobre alfa o sobre el continuo completo.


\subsection{Dos lecturas del estado terminal}
\label{subsec:k-terminal}

Denotemos por \(x_{\mathrm{term}}\) el estado producido por la composición
de selección, transporte y actualización hasta el corte terminal. La
escritura operatoria es
\[
 x_{\mathrm{term}}=
 (\mathcal U_{108}\circ\cdots\circ\mathcal U_1)(x_{\mathrm{seed}}),
 \qquad \mathcal U_t=\mathsf{Upd}_t\circ\mathsf{Tra}_t\circ\mathsf{Sel}_t.
\]
Sus dos lecturas pertinentes son
\[
 x_{\mathrm{term}}
 \xrightarrow{(\pi_{\mathrm{term}},R_{\mathrm{sgn}})}
 (B_{\mathrm{term}},U).
\]
La matriz \(B_{\mathrm{term}}\) retiene la lectura ternaria visible; el
vector \(U\) conserva otra coordenada de la misma historia. No se inserta
entre ellos una flecha \(B_{\mathrm{term}}\to U\).

El selector terminal opera sobre los tres catálogos regionales
de tamaños 196, 197 y 196, construidos en la
sección~\ref{censo:terminal-completo}. Las firmas locales reducen sus
\(196\cdot197\cdot196=7\,567\,952\) ternas a 331, y el margen de
columnas \(q_{\partial}=(1,2,2,1,2,0)\) deja dos matrices:
\[
 B_0=\begin{pmatrix}0&2&0&2&2&0\\0&0&1&2&2&0\\1&0&1&0&1&0\end{pmatrix},
 \qquad
 B_1=\begin{pmatrix}0&2&1&0&2&0\\0&0&1&2&2&0\\1&0&0&2&1&0\end{pmatrix}.
\]
La orientación de filas es \(\sigma=(1,1,-1)=(1,1,2)\) en
\(\mathbb F_3^3\). Las cargas de filas de las candidatas son
\((0,2,0)\) y \((2,2,1)\). Sólo la segunda pertenece a
\(\langle\sigma\rangle=\{(0,0,0),(1,1,2),(2,2,1)\}\); por tanto
\(B_{\mathrm{term}}=B_1\). Esta comprobación prueba el último paso del
selector a partir del censo declarado. No convierte el margen de columnas
en una consecuencia de la carga de filas, ni reemplaza la enumeración
anterior por la inspección de esas dos matrices.

El lector firmado tiene la composición tipada
\begin{equation}
 R_{\mathrm{sgn}}
 =\Pi_H\circ\mathcal E_{108}^{90,120}\circ\mathcal R_{12}:
 \mathcal X_{\mathrm{TPK}}^{\mathrm{term,mem}}
 \longrightarrow\mathbb Z^{12}.
 \label{eq:k-lector-firmado}
\end{equation}
El registro de canales determina
\[
 \mathcal E_{108}^{90,120}(\mathcal R_{12}(x_{\mathrm{term}}))
 =(b^{90},b^{120},Q_{\mathrm{TPK}}),
\]
con las coordenadas
\begin{align*}
 b^{90}={}&(-495,-116,-684,108,601,-90,\\
            &\hspace{19mm}-173,-88,313,560,-397,461),\\
 b^{120}={}&(-425,-281,-481,671,-255,30,\\
             &\hspace{19mm}475,-543,680,251,6,-128),\\
 Q_{\mathrm{TPK}}={}&6263.
\end{align*}
Estas cantidades se utilizan como coordenadas del registro previo, no
como parámetros que se ajusten mediante alfa. La reconstrucción demostrada
en los apartados siguientes comienza exactamente en esta salida tipada.
El rastro material del extractor anterior y el alcance de su transcripción
autónoma se conservan en la nota técnica de procedencia; la aritmética
posterior no se presenta como sustituto de ese rastro.

\subsection{Del registro de canales al vector firmado}
\label{subsec:k-pi-h}

El lector armónico \(\Pi_H\) puede escribirse sin utilizar ni \(K\) ni
alfa. Para evitar confundir sus sumas orbitales con los bloques de alfa,
las denotaremos por \(s_1,s_2,s_3\). Sean
\[
 B_r=\sum_{j=0}^{3}b^{120}_{r+3j},\qquad
 d_{r,j}=b^{90}_{r+3j},\qquad1\leq r\leq3,
\]
y
\begin{equation}
 s_1=\frac{Q_{\mathrm{TPK}}+2B_1+B_2}{3},\qquad
 s_2=s_1-B_1,\qquad s_3=s_2-B_2.
 \label{eq:k-sumas-armonicas}
\end{equation}
La división por tres pertenece al lector de las tres órbitas. Su
integralidad se comprueba sobre las salidas compatibles del registro; no
se supone para un elemento arbitrario de \(\mathbb Z^{25}\).

Los tres bloques firmados son
\begin{equation}
 U^{(r)}=(s_r,d_{r,1}-d_{r,3},d_{r,0}+d_{r,2},d_{r,0}-d_{r,2}).
 \label{eq:k-bloques-firmados}
\end{equation}
La sustitución entera da
\[
 (B_1,B_2,B_3)=(972,-1073,101),\qquad
 (s_1,s_2,s_3)=(2378,1406,2479),
\]
y, por tanto,
\begin{align*}
 U^{(1)}&=(2378,-452,-668,-322),\\
 U^{(2)}&=(1406,998,-204,-28),\\
 U^{(3)}&=(2479,-551,-371,-997).
\end{align*}
Intercalando por columnas se obtiene
\begin{equation}
 U=(2378,1406,2479,-452,998,-551,
       -668,-204,-371,-322,-28,-997).
 \label{eq:k-u-firmado}
\end{equation}
La presencia de valores negativos y de componentes superiores a 999
muestra que \(U\) no es una palabra de dígitos en base mil. Tampoco es
\(U_6\Vert U_6\), concatenación catalogal de período seis. La diferencia
es de dominio y de información conservada, no solamente de valores.

El orden constructivo hasta aquí es
\[
 x_{\mathrm{term}}\longrightarrow\mathcal R_{12}
 \longrightarrow(b^{90},b^{120},Q_{\mathrm{TPK}})
 \longrightarrow U.
\]
La inversión que sigue produce el registro dodecafásico. Que después podamos recuperar
los canales desde el registro dodecafásico constituye un control de reversibilidad; no
invierte este orden causal.

\subsection{Hadamard por órbitas y reconstrucción entera}
\label{subsec:k-hadamard}

En el ciclo de doce posiciones, el paso tres determina tres órbitas:
\[
 (1,4,7,10),\qquad(2,5,8,11),\qquad(3,6,9,12).
\]
La separación de cuatro sectores de cada órbita es la que, una vez fijados
origen y orientación, admite la lectura angular de noventa grados. La
transformación no actúa sobre tres cuaternas contiguas del orden natural.

Sea
\[
 H_4=\begin{pmatrix}
 1&1&1&1\\1&1&-1&-1\\1&-1&1&-1\\1&-1&-1&1
 \end{pmatrix}.
\]
Si \(P_{\mathrm{orb}}\) reordena por las tres órbitas indicadas,
definimos
\[
 \mathcal H_{12}=P_{\mathrm{orb}}^{-1}
 (H_4\oplus H_4\oplus H_4)P_{\mathrm{orb}}.
\]

\begin{lemma}[Inversión de Hadamard e integralidad]
\label{lem:k-hadamard}
Se tiene
\[
 H_4^2=4I_4,\qquad H_4^{-1}=\frac14H_4,\qquad\det H_4=-16.
\]
Para \(u\in\mathbb Z^4\), la preimagen \(H_4^{-1}u\) es entera si y
sólo si \(H_4u\equiv0\pmod4\).
\end{lemma}

\begin{proof}
Las filas son ortogonales y tienen norma cuadrada cuatro; la matriz es
simétrica. Esto da \(H_4^2=4I_4\) y la inversa. Una eliminación entera,
o la expansión del determinante, da el signo negativo y el valor \(-16\).
La condición de integralidad es exactamente la divisibilidad por cuatro
de las cuatro coordenadas de \(H_4u\).
\end{proof}

\begin{definition}[Transformación integral dodecafásica]
La subred de registros firmados y su coordenatización integral son
\[
 \mathcal L_H=\{u\in\ZZ^{12}:\mathcal H_{12}u\equiv0\pmod4\},
 \qquad \mathscr K(u)=\tfrac14\mathcal H_{12}u.
\]
La aplicación \(\mathscr K:\mathcal L_H\to\ZZ^{12}\) es un isomorfismo
de grupos abelianos, con inversa \(k\mapsto\mathcal H_{12}k\).
El registro canónico \(K\) es la imagen del registro firmado generado
por las operaciones anteriores, con origen de fase y orientación fijados.
\end{definition}

\begin{proof}
La congruencia define exactamente la integralidad de \(\mathscr K(u)\).
La identidad \(\mathcal H_{12}^2=4I\) demuestra las dos composiciones
inversas. En particular, todo \(k\in\ZZ^{12}\) tiene preimagen
\(\mathcal H_{12}k\in\mathcal L_H\).
\end{proof}

\begin{proposition}[Registro integral canónico]
\label{prop:k-sello}
La inversión \(K^{(r)}=H_4U^{(r)}/4\) de los bloques de
\eqref{eq:k-bloques-firmados} produce
\begin{align*}
 K^{(1)}&=(234,729,621,794),\\
 K^{(2)}&=(543,659,58,146),\\
 K^{(3)}&=(140,824,914,601).
\end{align*}
En el orden natural, el registro dodecafásico es
\begin{equation}
 K=(234,543,140,729,659,824,621,058,914,794,146,601).
 \label{eq:k-vector}
\end{equation}
Es la única preimagen racional de \(U\), y pertenece a
\(\{0,\ldots,999\}^{12}\).
\end{proposition}

\begin{proof}
Los tres productos antes de dividir son
\begin{align*}
 H_4U^{(1)}&=(936,2916,2484,3176),\\
 H_4U^{(2)}&=(2172,2636,232,584),\\
 H_4U^{(3)}&=(560,3296,3656,2404).
\end{align*}
Todas las coordenadas son divisibles por cuatro. Dividir y deshacer la
permutación de órbitas da \eqref{eq:k-vector}. La unicidad procede de la
invertibilidad racional, no de una búsqueda entre candidatos próximos a
un valor físico.
\end{proof}

La integralidad también puede formularse en términos de redes. Los
divisores determinantes de \(H_4\) son \(1,2,4,16\): las entradas
tienen máximo común divisor uno; los menores de orden dos son pares y
alguno tiene módulo dos; los de orden tres son múltiplos de cuatro y
alguno tiene módulo cuatro. De ahí se deduce
\[
 \operatorname{SNF}(H_4)=\operatorname{diag}(1,2,2,4),
\]
y para la transformación global,
\[
 \operatorname{coker}\mathcal H_{12}
 \cong(\mathbb Z/2\mathbb Z)^6\oplus(\mathbb Z/4\mathbb Z)^3.
\]
El índice de la imagen es \(16^3=4096\). Así, no todo vector firmado
entero admite una preimagen entera; el vector obtenido satisface las
congruencias precisas que la inversión exige. El factor \(1/4\) está
forzado por la matriz, y no tiene el estatuto de un coeficiente libre.

\subsection{Diferencias cíclicas y componente uniforme}
\label{subsec:k-transversal}

Con índices cíclicos módulo doce, sean
\[
 (D_3z)_m=z_m-z_{m+3},\qquad(D_4z)_m=z_m-z_{m+4},
 \qquad Q(z)=\sum_{m=1}^{12}z_m.
\]
El primer operador compara sectores separados por tres posiciones; el
segundo, por cuatro. Los nombres de canal 90 y 120 se refieren aquí a
esas separaciones en el ciclo orientado, no a una identificación con un
operador físico de otro dominio.

\begin{proposition}[Completitud del sistema transversal]
\label{prop:k-transversal}
El operador
\[
 \mathcal D=(D_3,D_4,Q):\mathbb Z^{12}\longrightarrow\mathbb Z^{25}
\]
tiene rango doce y factores invariantes no nulos
\[
 \operatorname{diag}(1,\ldots,1,12),
\]
con once unidades. En particular, sus salidas determinan un único registro dodecafásico.
\end{proposition}

\begin{proof}
Si \(D_3z=D_4z=0\), el vector es constante a lo largo de ambos pasos.
Como \(3\) y \(4\) generan el ciclo módulo doce,
\(\ker D_3\cap\ker D_4=\langle\mathbf1\rangle\). La carga total
elimina esta última libertad porque \(Q(\mathbf1)=12\).

Para la afirmación integral, las filas de \((D_3,D_4)\) son incidencias
orientadas de un grafo conexo. Un árbol generador proporciona once
factores invariantes unitarios. Todo menor de orden doce debe contener la
fila \(Q\). Sumando las primeras once columnas a la última, esta última
se anula en las filas de incidencia y vale doce en la fila \(Q\). Todos
los menores máximos son divisibles por doce y el árbol produce uno de
módulo exactamente doce. El último factor es, por tanto, doce.
\end{proof}

La comprobación sobre \eqref{eq:k-vector} da
\[
 D_3K=b^{90},\qquad D_4K=b^{120},\qquad Q(K)=6263.
\]
Las fórmulas \eqref{eq:k-sumas-armonicas} y
\eqref{eq:k-bloques-firmados} reconstruyen a su vez \(U\) desde esas
diferencias. Se cierra así un triángulo de representaciones exactas del
mismo objeto: registro transversal, vector firmado y registro dodecafásico. La igualdad
de resultados no convierte sus codominios en uno solo; demuestra que
las aplicaciones declaradas recuperan la información pertinente.

% Formalización nueva integrada en la revisión III; originales preservados.
% Propietarios: U016, ecuaciones extractor-diferencias-k y lector-armonico-*;
% registro_k.tex, subsecciones k-pi-h, k-hadamard y k-transversal.
\subsection{Compatibilidad integral y determinación del registro transversal}
\label{img09:seccion}

El dominio de esta construcción es la publicación del registro de una
historia APP--TRIT--TPK. Se distinguen la evaluación de los eventos del
registro, la compatibilidad de sus canales y la inversión de sus
coordenadas. El resultado siguiente da un criterio integral completo para
el segundo nivel y una factorización explícita para conectar el primero con él.
La demostración utiliza los operadores del registro y establece sus
relaciones para todo elemento de los dominios enteros indicados.

\subsubsection{Caracterización por incrementos adyacentes}

Los índices pertenecen a $\mathbb Z/12\mathbb Z$. Sean
\[
 (Sz)_i=z_{i+1},\qquad D_3=I-S^3,\qquad D_4=I-S^4,
 \qquad Q(z)=\sum_{i=0}^{11}z_i.
\]
Se conserva el extractor transversal
\[
 \mathcal D:\mathbb Z^{12}\longrightarrow\mathbb Z^{25},
 \qquad k\longmapsto(D_3k,D_4k,Q(k)).
\]
Para una terna de coordenadas enteras $(b,c,q)$, escribimos
\begin{equation}
 w_i=c_i-b_{i+1},\qquad
 P_0=0,\quad P_i=\sum_{j=0}^{i-1}w_j\ (1\leq i\leq12),
 \qquad T(w)=\sum_{i=0}^{11}P_i.
 \label{img09:incremento}
\end{equation}

\begin{theorem}[Imagen integral del extractor transversal]
\label{img09:imagen}
Una terna $(b,c,q)\in\mathbb Z^{25}$ pertenece a
$\operatorname{im}\mathcal D$ si y sólo si satisface
\begin{align}
 b_i&=w_i+w_{i+1}+w_{i+2}&& (0\leq i<12),
 \label{img09:relaciones}\\
 \sum_{i=0}^{11}w_i&=0,
 \label{img09:cierre}\\
 q+T(w)&\equiv0\pmod {12}.
 \label{img09:congruencia}
\end{align}
Su única preimagen entera es
\begin{equation}
 k_i=\frac{q+T(w)}{12}-P_i,\qquad0\leq i<12.
 \label{img09:inversa}
\end{equation}
Las doce ecuaciones de \eqref{img09:relaciones} y la ecuación
\eqref{img09:cierre} son linealmente independientes sobre $\mathbb Q$.
\end{theorem}
\begin{proof}
Para una salida de $\mathcal D$ se tiene
\[
 c_i-b_{i+1}=(k_i-k_{i+4})-(k_{i+1}-k_{i+4})=k_i-k_{i+1}.
\]
La suma de tres incrementos da $b_i$; la suma cíclica de todos ellos
es cero. Además, $P_i=k_0-k_i$, por lo que
$q=12k_0-T(w)$. Se obtienen todas las condiciones y la fórmula inversa.

Recíprocamente, \eqref{img09:congruencia} hace entera la fórmula
\eqref{img09:inversa}; \eqref{img09:cierre} permite prolongarla
cíclicamente. Sus diferencias adyacentes son $w_i$. Por tanto,
$D_3k=b$ y
\[
 (D_4k)_i=w_i+w_{i+1}+w_{i+2}+w_{i+3}
 =w_i+b_{i+1}=c_i.
\]
La suma de la inversa da $Q(k)=q$. La misma fórmula prueba la unicidad.
Para la independencia, el cambio entero invertible
$(b,c)\mapsto(b,w=c-Sb)$ convierte las doce primeras relaciones en
$b=(I+S+S^2)w$. Cada una despeja una coordenada diferente de $b$.
La condición $\sum w_i=0$ es una relación adicional independiente.
\end{proof}

La imagen racional tiene dimensión doce. Su estructura integral añade
una congruencia de orden doce: ésta es la manifestación constructiva del
último factor de Smith de $\mathcal D$. Los once incrementos
$w_0,\ldots,w_{10}$ y la carga $q$ son coordenadas suficientes, con
$w_{11}=-\sum_{i=0}^{10}w_i$ y la congruencia
\eqref{img09:congruencia}. Las veinticinco coordenadas originales
conservan esos mismos doce grados de libertad mediante trece relaciones.

\subsubsection{Conmutatividad con la lectura armónica}

Sea $\mathcal H_{12}$ la transformación de Hadamard del registro,
con órbitas $(r,r+3,r+6,r+9)$ para $r=0,1,2$. Su bloque es
\[
 H_4=\begin{pmatrix}
 1&1&1&1\\1&1&-1&-1\\1&-1&1&-1\\1&-1&-1&1
 \end{pmatrix},\qquad H_4^2=4I_4.
\]
Para $(b,c,q)$ compatible, se definen
\[
 B_r=\sum_{j=0}^3c_{r+3j},\qquad
 A_0=\frac{q+2B_0+B_1}{3},\quad A_1=A_0-B_0,
 \quad A_2=A_1-B_1.
\]
El lector $\Pi_H$ de \eqref{eq:k-bloques-firmados} tiene los bloques
\[
 (\Pi_H(b,c,q))^{(r)}=
 (A_r,b_{r+3}-b_{r+9},b_r+b_{r+6},b_r-b_{r+6}).
\]

\begin{proposition}[Concordancia de las dos reconstrucciones]
\label{img09:triangulo}
Sobre $\operatorname{im}\mathcal D$ se cumple
\[
 \Pi_H\mathcal D=\mathcal H_{12},\qquad
 \mathcal D^{-1}=\tfrac14\mathcal H_{12}\Pi_H.
\]
La imagen de $\Pi_H$ restringida a ese dominio es exactamente
\[
 \mathcal L_H=
 \{u\in\mathbb Z^{12}:\mathcal H_{12}u\equiv0\pmod4\}.
\]
\end{proposition}
\begin{proof}
Para $k=\mathcal D^{-1}(b,c,q)$, sea
$a_r=\sum_{j=0}^3k_{r+3j}$. El desplazamiento de cuatro posiciones
lleva cada órbita a la siguiente, y por ello $B_r=a_r-a_{r+1}$.
Junto con $q=a_0+a_1+a_2$, estas identidades dan $A_r=a_r$.
En una órbita, escribamos $(x_0,x_1,x_2,x_3)$ para las coordenadas
de $k$ y $d_j=x_j-x_{j+1}$. Entonces
\[
 d_1-d_3=x_0+x_1-x_2-x_3,\quad
 d_0+d_2=x_0-x_1+x_2-x_3,\quad
 d_0-d_2=x_0-x_1-x_2+x_3.
\]
Son las tres filas restantes de $H_4k^{(r)}$. La inversión y la
caracterización integral se siguen de $\mathcal H_{12}^2=4I$.
\end{proof}

La condición $\Pi_H(b,c,q)\in\mathcal L_H$ sólo controla la
salida armónica. La pertenencia de las veinticinco coordenadas a
$\operatorname{im}\mathcal D$ exige además
\eqref{img09:relaciones}--\eqref{img09:congruencia}.
Por ejemplo, añadir a $c$ el vector $e_0-e_3$ conserva las tres sumas
orbitales y deja $\Pi_H$ inalterado. Partiendo de $(0,0,0)$ se
obtiene así una terna con $\Pi_H=0$ que incumple
\eqref{img09:relaciones}. Éste es un control negativo exacto de la
compatibilidad del canal, independiente de toda evaluación terminal.

\subsubsection{Determinación y transporte sobre las órbitas del registro}

El Teorema~\ref{img09:imagen} determina $k$ cuando se han producido
sus canales y su carga. Como condición sobre salidas aún no
individualizadas, la misma imagen es un grupo isomorfo a
$\mathbb Z^{12}$. Fijada únicamente una carga $q$, las soluciones
forman una clase afín de
\[
 L_0=\{v\in\mathbb Z^{12}:\sum v_i=0\}
 =\bigoplus_{i=0}^{10}\mathbb Z(e_i-e_{11}).
\]
Por ejemplo, $100\mathbf1+e_0$ y $100\mathbf1+e_1$ son dos
registros distintos, de carga $1201$, con componentes en
$\{0,\ldots,999\}$. Cada uno satisface todas las condiciones de
imagen y todas las congruencias de Hadamard mediante sus propios
canales. Estos testigos comparan las ecuaciones de compatibilidad;
no se les atribuye la condición de historias terminales HMT.

\begin{proposition}[Covariancia cíclica y estabilizador del registro]
\label{img09:covariancia}
Sean $\rho$ la traslación de ventanas en un dominio de registros
$\mathscr R$ y $S$ la traslación anterior en $\mathbb Z^{12}$.
Sobre la órbita de $R\in\mathscr R$, un lector covariante queda
determinado por un vector
\[
 k\in\operatorname{Fix}(S^p),
\]
donde $p\mid12$ es el período mínimo de $R$ bajo $\rho$.
La condición de carga $Q(k)=q$ admite soluciones enteras si y sólo si
$12/p$ divide a $q$; cuando existen, forman una clase afín de rango
$p-1$. En particular, sobre una órbita libre, covariancia y carga
dejan once grados de libertad.
\end{proposition}
\begin{proof}
La regla $F(\rho^jR)=S^jk$ está bien definida exactamente cuando
$S^pk=k$. Esos vectores repiten $12/p$ veces una palabra de
longitud $p$. Su carga es $\frac{12}{p}\sum_{i=0}^{p-1}k_i$.
La divisibilidad y el rango se siguen de esta fórmula. La
covariancia de los canales resulta de $D_aS=SD_a$ y $QS=Q$.
Sobre $\mathcal L_H$, la acción correspondiente es
$\widehat S_H=\mathcal H_{12}S\mathcal H_{12}/4$;
por la Proposición~\ref{img09:triangulo}, ambas reconstrucciones
entrelazan estas acciones.
\end{proof}

Si el dominio lleva un origen marcado, debe utilizarse la acción que
transporta también esa marca. La proposición se refiere a esa acción
cíclica declarada; no presupone la periodicidad de la historia
enriquecida. La naturalidad frente a truncamientos posteriores al corte
$108$ se obtiene, para cualquier lector $F_{108}$ ya definido, mediante
$F_N=F_{108}\circ\tau_{N,108}$. Esta naturalidad conserva la regla
de evaluación inicial; por sí sola no elige una de las reglas posibles
sobre los primeros ciento ocho eventos.

\subsubsection{Evaluación por eventos y composición de contribuciones}

La caracterización anterior proporciona una factorización de la extracción mediante dos
lecturas efectivas del registro:
\begin{equation}
 \omega:\mathscr R\longrightarrow
 \{w\in\mathbb Z^{12}:\sum w_i=0\},\qquad
 q_{\rm reg}:\mathscr R\longrightarrow\mathbb Z,
 \quad q_{\rm reg}(R)+T(\omega(R))\equiv0\pmod {12}.
 \label{img09:contrato-productor}
\end{equation}
Estas aplicaciones deben evaluarse sobre los eventos, la orientación,
la incidencia y la memoria que produjo APP--TRIT--TPK. Una vez dadas sus
acciones sobre esos datos, la composición
\[
 R\longmapsto(w,q)\longmapsto
 ((I+S+S^2)w,(I+S+S^2+S^3)w,q)
 \xrightarrow{\Pi_H}u\xrightarrow{\mathcal H_{12}/4}k
\]
está determinada y satisface todos los controles inversos. La fórmula
$w=c-Sb$ también permite conectar un productor que ya entregue ambos
canales: es una comprobación sobre su salida, anterior a toda comparación
con el testigo terminal.

\begin{proposition}[Acumulación de contribuciones y unicidad prospectiva]
\label{img09:acumulacion}
Supóngase que el estado inicial y cada evento $a_t$ del registro
producen, mediante reglas fijadas previamente, contribuciones
$(\delta w_t,\delta q_t)$ con
\[
 \sum_i\delta w_{t,i}=0,\qquad
 \delta q_t+T(\delta w_t)\equiv0\pmod {12}.
\]
Con condiciones iniciales compatibles $(w^{(0)},q^{(0)})$, la
actualización
\[
 w^{(t+1)}=w^{(t)}+\delta w_t,\qquad
 q^{(t+1)}=q^{(t)}+\delta q_t
\]
produce un único registro en cada etapa. Su incremento es
\[
 \delta k_{t,i}=
 \frac{\delta q_t+T(\delta w_t)}{12}-P_i(\delta w_t).
\]
La construcción respeta la concatenación de registros, y la lectura
de un prefijo permanece igual bajo prolongaciones posteriores.
\end{proposition}
\begin{proof}
Las condiciones de compatibilidad se conservan por suma. La fórmula
\eqref{img09:inversa} es aditiva en $(w,q)$ y produce el
incremento indicado. La inducción fija sucesivamente cada estado.
La suma sobre dos segmentos es la suma concatenada, y un prefijo recibe
exactamente los mismos sumandos antes y después de prolongar la historia.
\end{proof}

La proposición proporciona un criterio suficiente, no una exigencia de
linealidad para todo lector TPK. Si los eventos se transportan por acciones
no triviales, sus contribuciones pueden expresarse en la carta terminal
mediante el cociclo afín correspondiente. La ecuación de cociclo conserva
entonces la misma independencia de la partición de la trayectoria.
En ambos casos, las reglas de evento y el estado inicial constituyen los
datos constructivos que han de provenir del artículo: escogerlos para
reproducir un vector esperado sustituiría esa procedencia por otra tarea.

El resultado de este desarrollo es la condición exacta de determinación
\eqref{img09:contrato-productor}, su inversión demostrada y sus identidades
de transporte. Su alcance es el registro transversal; las restantes
construcciones de HMT conservan sus dominios y sus pruebas propias.


\subsection{Dos lecturas racionales: \texorpdfstring{\(\kappa_{36}\)}{kappa36} y \texorpdfstring{\(\kappa_{\mathrm{per}}\)}{kappa periódica}}
\label{subsec:k-kappas}

Una vez construido \(K\), la carta en base mil permite dos operaciones
distintas. La primera lee una sola vuelta; la segunda prolonga el registro dodecafásico
periódicamente. Sea \(B=1000\) y definamos su entero de concatenación
\begin{equation}
 N_K=\sum_{m=1}^{12}K_mB^{12-m}
 =234543140729659824621058914794146601.
 \label{eq:k-numerador}
\end{equation}
Los ceros iniciales de un bloque, como el de 058, son parte de su escritura
de tres cifras. No cambian su valor entero, pero sí permiten concatenar
sin ambigüedad los doce bloques.

\begin{definition}[Lecturas finita y periódica del registro dodecafásico]
\label{def:k-kappas}
La lectura de una vuelta es
\[
 \kappa_{36}=\sum_{m=1}^{12}K_mB^{-m}=\frac{N_K}{B^{12}}.
\]
La prolongación periódica es
\[
 K_j^{\mathrm{per}}=K_{1+((j-1)\bmod12)},\qquad
 \kappa_{\mathrm{per}}=\sum_{j\geq1}K_j^{\mathrm{per}}B^{-j}.
\]
\end{definition}

\begin{proposition}[Valor y período exactos del registro dodecafásico prolongado]
\label{prop:k-periodo}
Se cumple
\begin{equation}
 \kappa_{\mathrm{per}}=\frac{N_K}{B^{12}-1},\qquad
 \kappa_{\mathrm{per}}-\kappa_{36}
 =\frac{N_K}{B^{12}(B^{12}-1)}
 =B^{-12}\kappa_{\mathrm{per}}.
 \label{eq:k-diferencia-kappas}
\end{equation}
La palabra infinita del registro dodecafásico tiene período mínimo doce en base mil y
período mínimo treinta y seis en base diez.
\end{proposition}

\begin{proof}
La suma de la primera vuelta es \(N_KB^{-12}\); cada nueva vuelta
multiplica su contribución por \(B^{-12}\). La serie geométrica produce
\(N_K/(B^{12}-1)\), y la resta da la segunda identidad.

Los doce componentes de \(K\) son distintos. Un período propio del ciclo
identificaría dos de ellos; por tanto el período mínimo de bloques es
doce. La expansión decimal es canónica, pues el registro dodecafásico no tiene una cola
de bloques 999. Si su período mínimo decimal fuera \(d\), dividiría
36. Agrupar de tres en tres produciría un período de bloques
\(d/\gcd(d,3)\). Ningún divisor propio de 36 da doce mediante esa
operación. Luego el período decimal mínimo es 36.
\end{proof}

Ambas lecturas son racionales, pero no son el mismo racional. La primera
termina después de doce bloques; la segunda repite la vuelta. Además,
\(0<\kappa_{\mathrm{per}}-\kappa_{36}<B^{-12}\). Su cercanía no
permite sustituir una por otra en una identidad exacta. En particular,
el cierre finito de alfa utiliza \(\kappa_{36}\), mientras que la
suma de la sección prolongada utiliza \(\kappa_{\mathrm{per}}\).

\subsection{Reversibilidad de la constante racional y cambio de fase}

La constante racional \(\kappa_{\mathrm{per}}\) determina íntegramente
el registro \(K\) cuando se conservan la base, la longitud y el origen
de lectura. Su función numérica y su función estructural se relacionan
mediante una inversión exacta.

\begin{proposition}[Recuperación integral y transformación cíclica]
Sea \(I(K)=\sum_{j=1}^{12}K_j1000^{12-j}\). Entonces
\[
 I(K)=(1000^{12}-1)\kappa_{\mathrm{per}}.
\]
La expansión entera de \(I(K)\) en doce bloques de base \(1000\)
recupera todos los \(K_j\), incluidos los ceros iniciales de cada bloque.
Si \(\rho(K)=(K_2,\ldots,K_{12},K_1)\), se cumple
\begin{equation}
 \kappa(\rho K)=1000\kappa(K)-K_1.
 \label{eq:k-rotacion-racional}
\end{equation}
\end{proposition}
\begin{proof}
La primera identidad invierte la definición de la evaluación periódica.
La división euclídea iterada por \(1000\), con doce posiciones fijadas,
recupera el vector. Por otra parte,
\[
 I(\rho K)=1000I(K)-K_1(1000^{12}-1).
\]
Dividir por \(1000^{12}-1\) da \eqref{eq:k-rotacion-racional}.
\end{proof}

La recuperación compone, por tanto, \(\kappa\mapsto K\mapsto U\)
y las transformaciones \(D_3K,D_4K,Q\) ya construidas. La evaluación
racional conserva el registro firmado y las diferencias cíclicas que
dependen de él. La profundidad total y los operadores de transporte de
una historia permanecen en sus coordenadas propias: la inversión anterior
tiene por dominio el registro integral dodecafásico.

Escribamos \(\kappa_j=\kappa(\rho^jK)\), con índices módulo \(12\).
La misma identidad da
\[
 K_{j+1}=1000\kappa_j-\kappa_{j+1},\qquad
 \sum_{j=0}^{11}\kappa_j=\frac{\sum_{j=1}^{12}K_j}{999}
 =\frac{6263}{999}.
\]
La fase actúa de manera exacta sobre la constante: su valor es fijo
en la lectura canónica, mientras sus rotaciones siguen una ley afín.
La suma de la órbita es invariante bajo la elección del origen.

También la incidencia posterior dispone de una expresión en estas
coordenadas. El umbral cúbico determina
\[
 B_K=\{j+1:1000\kappa_j-\kappa_{j+1}\ge729\}
 =\{4,6,9,10\}.
\]
Se recupera la tétrada a partir de la órbita racional porque ésta
recupera primero las componentes enteras. El soporte firmado asociado
a \(\alpha\) utiliza además sus registros regionales y la normalización
de frontera. La relación entre coordenada e incidencia queda así
expresada mediante operaciones identificables.

Esta es una definición operacional de la función de \(K\): proporciona
una coordenatización integral reversible del registro orientado, cuya
evaluación periódica y cuyas funciones de incidencia participan en
construcciones posteriores. El valor, la transformación y el dominio
se conservan conjuntamente.


\subsection{Periodicidad y conúcleo del operador de transporte}
\label{subsec:k-clase-acarreo}

La periodicidad que acabamos de demostrar pertenece al registro dodecafásico. No implica
periodicidad de las coordenadas regionales, de los acarreos de la clausura
ni de alfa. Es útil precisar esta separación mediante un segundo
cociente, puramente integral.

Sea \(S\) el desplazamiento cíclico de doce coordenadas, con orientación
fijada, y sea \(b\geq2\) una base entera. El operador de acarreo
periódico es \(\partial_b=S-bI\). Si una identidad periódica se escribe
\[
 A=P+E-\Phi-K+\partial_bC,
\]
entonces las transformaciones
\[
 K\longmapsto K+\partial_bh,\qquad C\longmapsto C+h
\]
conservan su lado derecho. Esta invariancia compara representantes de una
clase de acarreo; no es una autorización para modificar el registro dodecafásico ya
seleccionado por el registro.

\begin{proposition}[Cociente del acarreo periódico]
\label{prop:k-cociente-acarreo}
Con la orientación anterior,
\[
 \operatorname{coker}(S-bI)\cong\mathbb Z/(b^{12}-1)\mathbb Z.
\]
\end{proposition}

\begin{proof}
El módulo cíclico se presenta como
\(\mathbb Z[t]/(t^{12}-1)\). Imponer \(S=bI\) equivale a imponer
\(t=b\), y deja la relación \(b^{12}-1=0\). La evaluación de los
coeficientes en \(b\), con el orden de concatenación elegido, representa
la clase resultante.
\end{proof}

El denominador \(B^{12}-1\) de \(\kappa_{\mathrm{per}}\) y este
cociente proceden de la misma longitud cíclica, pero sus objetos siguen
siendo distintos: uno es una evaluación racional y el otro un cociente
de un módulo de acarreos. La clausura finita de alfa será una cadena
orientada con frontera derecha. Su prolongación transportará una
frontera remota compatible, no un acarreo que se identifique por decreto
entre la primera y la decimotercera posición.

Queda así preparado el capítulo siguiente. El registro dodecafásico se ha fijado antes
de resolver la recurrencia de alfa; su reconstrucción por Hadamard es
entera y única; sus diferencias transversales recuperan el registro; y
sus dos evaluaciones racionales tienen dominios y usos separados. Ninguno
de estos hechos utiliza alfa como entrada, y ninguno decide todavía el
tipo aritmético de la salida de la clausura completa.

\HMTDivision{\(\alpha\)}{Constante de la estructura fina}{Las coordenadas regionales y el registro ordenado concurren en las dos lecturas de la constante de estructura fina.}
\HMTNarrativeHeading{La constante de estructura fina en la generación común}

La constante de estructura fina ocupa el punto donde las coordenadas
regionales y el registro ordenado concurren en una lectura determinada.
Las coordenadas circular, áurea y exponencial ya han sido construidas.
El registro firmado procede del mismo estado y conserva su orientación.
La clausura dodecafásica recibe esos resultados con los acarreos que
permiten componerlos.

La vía aritmética proporciona una caracterización precisa: su coordenada
de alfa es la diferencia entre el agregado regional declarado y la
lectura periódica de la constante de acoplamiento. Esa igualdad enlaza
objetos producidos previamente. Su explicación exige seguir los bloques,
las posiciones y los acarreos que convierten el registro en una expansión.
La demostración conservará tanto el valor como los datos que hacen
reproducible la operación.

La segunda vía emplea el resolvente de vacancias \(E_{21/22}\) y
sus prolongaciones graduadas. Su raíz positiva se estudia con su dominio
y sus condiciones de selección. La comparación de las dos vías se
formula mediante cilindros comunes a toda profundidad.
El origen compartido de las
constantes se expresa por esa composición de resultados, y cada
evaluación conserva su función particular dentro de ella.

La incidencia añade otra lectura del mismo cierre. La firma distingue
una hexada sobre la cuaterna determinada por el registro. Al prolongarse
la incidencia hacia el diseño binario, la configuración marcada
participa en la relación entre las regiones y sus octadas. Esta
conexión prepara el tránsito hacia las realizaciones excepcionales y,
posteriormente, hacia los lectores de área e información.

La importancia de alfa en el artículo procede de esta posición causal. La
coordenada obtenida interviene en la valoración de la acción y en los
lectores angulares; éstos transmiten a las construcciones siguientes
orientación y normalización. La exposición avanza desde el origen
regional hasta esos usos, manteniendo visible el registro que acompaña
al valor. Así, la relación entre alfa y la constante de acoplamiento
adquiere una definición operativa y un lugar inequívoco en la doble
proyección.
\HMTMathematicalPaper
\section{Derivación de la constante de estructura fina}
\label{sec:alpha}

La constante de estructura fina recibe en este desarrollo una determinación
aritmética que reúne las tres coordenadas regionales y el registro integral
dodecafásico. La operación combina sus bloques con orientación fijada y
transporta los cocientes de la normalización posicional. El parámetro
resultante expresa la compatibilidad de esa diferencia orientada con
la prolongación de los cilindros. Su valor y los datos de incidencia
utilizados en su construcción permanecerán relacionados durante el
desarrollo excepcional.

El orden interno es importante. Las coordenadas \(\pi\),
\(e\) y \(\varphi\), ya construidas en la sección precedente, son coordenadas
previas del mismo estado aritmético con memoria de trayectoria, y pueden actuar legítimamente como
entradas de un lector posterior. No son constantes convencionales
inyectadas para seleccionar alfa. La generación correlacionada de las cuatro coordenadas comparte origen y
compatibilidad coinductiva, sin que ello imponga evaluar sus cuatro
componentes en un solo paso.

Las dos vías expresan funciones matemáticas diferentes del mismo parámetro.
La primera realiza una clausura entera por división euclídea en base
\(1000\). La segunda estudia el equilibrio analítico de vacancias mediante
el resolvente \(E_{21/22}\) y sus prolongaciones graduadas. Una vía
determina bloques y cocientes; la otra determina una raíz mediante una
ecuación y un dominio de unicidad. Su identificación se formula después
en los cilindros comunes a toda profundidad. La exposición distingue
estos tres resultados: construcción posicional, determinación analítica
y compatibilidad de ambas.

\subsection{Primera vía: determinación posicional por clausura entera}
\subsubsection{Dominio y orientación de la clausura entera}
\label{subsec:alpha-dominio}

Sean \(P_j,E_j,\Phi_j\in\{0,\ldots,999\}\) los bloques canónicos de
las tres coordenadas regionales, y prolonguemos el registro dodecafásico por
\[
 K_j=K_{1+((j-1)\bmod12)}.
\]
La orientación del canal es \((1,1,-1)\). La contribución del registro dodecafásico se
sustrae según la clausura dodecafásica ya tipada. Para \(B=1000\), la
operación local distingue el balance firmado
\[
 Z_j:=P_j+E_j-\Phi_j-K_j
\]
de la suma que incorpora el cociente entrante. La normalización es
\begin{equation}
 s_j=Z_j+c_{j+1},\qquad
 A_j=s_j-B\left\lfloor\frac{s_j}{B}\right\rfloor,
 \qquad c_j=\left\lfloor\frac{s_j}{B}\right\rfloor.
 \label{eq:alpha-recurrencia}
\end{equation}
Así, \(A_j\in\{0,\ldots,B-1\}\), incluso cuando \(s_j\) es
negativo. El acarreo que entra desde la derecha pertenece al dominio de
la operación; no puede omitirse porque la tabla se imprima de izquierda
a derecha.

\begin{definition}[Clausura de profundidad finita]
\label{def:alpha-finita}
Para una profundidad \(N\) y una frontera entera \(t\), la aplicación
\(\mathcal C_N\) recibe
\((P_{1:N},E_{1:N},\Phi_{1:N},K_{1:N},c_{N+1}=t)\) y aplica
\eqref{eq:alpha-recurrencia} en el orden \(N,N-1,\ldots,1\). Su salida
es la pareja de palabras
\[
 (A_{1:N},c_{1:N+1}).
\]
\end{definition}

La división euclídea no posee una elección de signo residual: el resto se
fija en \(\{0,\ldots,999\}\) y el cociente queda determinado. Por
ejemplo, \(-431=569-1000\) y \(-1200=800-2\cdot1000\). Los
acarreos negativos de la clausura no son anomalías; son la información
entera necesaria para mantener la igualdad orientada.

\subsubsection{Ventana dodecafásica y condiciones de acarreo}
\label{subsec:alpha-ventana}

Las primeras doce tríadas obtenidas son
\begin{align*}
 P={}&(141,592,653,589,793,238,462,643,383,279,502,884),\\
 E={}&(718,281,828,459,045,235,360,287,471,352,662,497),\\
 \Phi={}&(618,033,988,749,894,848,204,586,834,365,638,117).
\end{align*}
En el canal de clausura, las cinco regiones que determinan la coordenada
\(P\) conservan sus antecedentes distintos. Compartir esta lectura no
las identifica como regiones o como historias.

Con frontera \(c_{13}=0\), la recurrencia produce la tabla siguiente.
Se muestra cada entrada, el acarreo entrante, el total antes de normalizar,
el bloque obtenido y el acarreo saliente. La tabla es un cálculo entero
reproducible, no una lista de bloques seleccionados por su parecido con
un objetivo.

\begin{table}[htbp]
\centering
\small
\setlength{\tabcolsep}{3.5pt}
\begin{tabular}{r|rrrr|rr|rr}
\hline
\(j\)&\(P_j\)&\(E_j\)&\(\Phi_j\)&\(K_j\)&\(c_{j+1}\)&\(s_j\)&\(A_j\)&\(c_j\)\\
\hline
1&141&718&618&234&0&7&007&0\\
2&592&281&033&543&0&297&297&0\\
3&653&828&988&140&\(-1\)&352&352&0\\
4&589&459&749&729&\(-1\)&\(-431\)&569&\(-1\)\\
5&793&045&894&659&\(-2\)&\(-717\)&283&\(-1\)\\
6&238&235&848&824&\(-1\)&\(-1200\)&800&\(-2\)\\
7&462&360&204&621&0&\(-3\)&997&\(-1\)\\
8&643&287&586&058&\(-1\)&285&285&0\\
9&383&471&834&914&\(-1\)&\(-895\)&105&\(-1\)\\
10&279&352&365&794&0&\(-528\)&472&\(-1\)\\
11&502&662&638&146&0&380&380&0\\
12&884&497&117&601&0&663&663&0\\
\hline
\end{tabular}
\caption{Clausura inicial con frontera derecha nula. La dependencia se propaga de la última fila a la primera.}
\label{tab:alpha-ventana}
\end{table}

La palabra de acarreos es
\[
 (c_1,\ldots,c_{13})=(0,0,0,-1,-1,-2,-1,0,-1,-1,0,0,0),
\]
y la lectura de la ventana es el racional
\begin{equation}
 \alpha_{12}^{(0)}
 =\frac{7297352569283800997285105472380663}{10^{36}}
 =0.007297352569283800997285105472380663.
 \label{eq:alpha-ventana-cero}
\end{equation}
El superíndice recuerda la frontera elegida. Esta ventana es exacta en su
dominio finito; no se la identifica sin más con el truncamiento de la
sección infinita. La diferencia entre ambas lecturas se hará visible al
transportar el acarreo de la cola remota.

\subsubsection{Telescopado y reversibilidad en una fibra de frontera}
\label{subsec:alpha-telescopia}

La forma local de la igualdad es
\begin{equation}
 P_j+E_j-\Phi_j-K_j-A_j=Bc_j-c_{j+1}.
 \label{eq:alpha-identidad-local}
\end{equation}
El término de acarreo es un enlace entre posiciones contiguas. Si se lo
elimina coordenada a coordenada, se obtiene una igualdad falsa. Sólo una
evaluación que conserve sus extremos permite telescoparlo.

\begin{proposition}[Identidad entera de profundidad arbitraria]
\label{prop:alpha-telescopia}
Para \(\iota_N(z)=\sum_{j=1}^{N}z_jB^{N-j}\), toda ejecución de
\(\mathcal C_N\) satisface
\begin{equation}
 \iota_N(P)+\iota_N(E)-\iota_N(\Phi)-\iota_N(K)-\iota_N(A)
 =B^Nc_1-c_{N+1}.
 \label{eq:alpha-telescopia}
\end{equation}
\end{proposition}

\begin{proof}
Multiplicar \eqref{eq:alpha-identidad-local} por \(B^{N-j}\) y sumar
produce en el lado derecho
\[
 \sum_{j=1}^{N}(Bc_j-c_{j+1})B^{N-j}=B^Nc_1-c_{N+1}.
\]
Cada acarreo interior aparece dos veces con signo opuesto; sólo sobreviven
los dos extremos.
\end{proof}

Para la tabla~\ref{tab:alpha-ventana}, ambos extremos son cero. Si
\(p_{12},e_{12},f_{12}\) son las tres sumas fraccionarias truncadas,
\[
 \alpha_{12}^{(0)}=p_{12}+e_{12}-f_{12}-\kappa_{36}.
\]
La constante del registro dodecafásico que aparece es \(\kappa_{36}\), no
\(\kappa_{\mathrm{per}}\). El dominio de esta identidad contiene
truncamientos y una frontera finita; su forma no autoriza a borrar ambas
condiciones al pasar a la sección completa.

La inversión local es también exacta:
\[
 K_j=P_j+E_j-\Phi_j+c_{j+1}-A_j-Bc_j.
\]
Dados \(P,E,\Phi\), los bloques \(A\) y los acarreos completos,
recupera el registro dodecafásico. La concatenación recupera además
\[
 \iota_N(K)=\iota_N(P)+\iota_N(E)-\iota_N(\Phi)-\iota_N(A)
             -B^Nc_1+c_{N+1}.
\]
Esta reversibilidad es una propiedad de la fibra con datos y fronteras
fijados. No debe convertirse en una nueva receta causal que empiece por
alfa y fabrique después el registro. La dirección de construcción sigue
siendo registro firmado, registro dodecafásico, clausura y salida.

\subsubsection{Prolongación y normalización canónica de la sección completa}
\label{subsec:alpha-completa}

La prolongación conserva el registro dodecafásico periódico y recibe las secuencias
regionales completas. Sean
\[
 p=\sum_{j\geq1}P_jB^{-j},\qquad
 e_0=\sum_{j\geq1}E_jB^{-j},\qquad
 f=\sum_{j\geq1}\Phi_jB^{-j}.
\]
Cada serie es una realización del carácter interno correspondiente.
Las partes enteras ya obtenidas son 3, 2 y 1, de modo que
\(p=\pi-3\), \(e_0=e-2\) y
\(f=\varphi-1\).

\begin{proposition}[Existencia y unicidad de la normalización completa]
\label{prop:alpha-normalizacion}
La serie
\begin{equation}
 y=p+e_0-f-\kappa_{\mathrm{per}}
   =\sum_{j\geq1}(P_j+E_j-\Phi_j-K_j)B^{-j}
 \label{eq:alpha-serie}
\end{equation}
converge absolutamente. En el dominio de los bloques obtenidos,
\(0<y<10^{-2}\). Existe una única expansión canónica \((A_j)\) de
\(y\), sin cola finalmente constante 999, y una única sucesión entera
acotada \((c_j)\) que satisface \eqref{eq:alpha-identidad-local} con
\(c_1=0\). Estos acarreos pertenecen al conjunto invariante
\(\mathcal C=\{-2,-1,0,1,2\}\).
\end{proposition}

\begin{proof}
Los coeficientes de la serie están acotados por \(2(B-1)\) en módulo,
de modo que domina una serie geométrica. Si
\(Z_j=P_j+E_j-\Phi_j-K_j\) y
\(R_j(Z)=\sum_{r\geq0}Z_{j+r}B^{-r-1}\), cada cola es la diferencia
de dos sumas de dos colas canónicas y pertenece a \((-2,2)\).
El primer coeficiente es
\(Z_1=141+718-618-234=7\); por tanto
\(y=(7+R_2(Z))/1000\in(0.005,0.009)\).

Sea \((A_j)\) la expansión canónica de \(y\). Definamos
\[
 c_j=R_j(Z)-R_j(A).
\]
Al multiplicar la igualdad de las dos series por \(B^{j-1}\) y separar
sus primeros \(j-1\) términos se obtiene
\[
 c_j=\sum_{r=1}^{j-1}(A_r-Z_r)B^{j-1-r}\in\mathbb Z,
\]
y \(c_1=0\). Además, \(R_j(A)\in[0,1)\), luego
\(-3<c_j<2\), lo que incluye todos los acarreos en \(\mathcal C\).
Las identidades de desplazamiento de las colas dan
\(Bc_j=Z_j-A_j+c_{j+1}\), que es la recurrencia requerida.

Recíprocamente, cualquier solución acotada con salida canónica y
\(c_1=0\) telescopa en el límite a la misma serie \(y\). La expansión
canónica es única y, una vez fijada, la fórmula de las colas determina
todos los acarreos. Finalmente, si se aplica una ejecución finita con
\(c_{j+1}\in\mathcal C\), se tiene \(-2000\leq s_j\leq2000\),
por lo que su cociente euclídeo vuelve a pertenecer a \(\mathcal C\).
\end{proof}

La proposición explicita la normalización del lector posterior: no
reemplaza la generación de las tres secuencias por una serie convencional
introducida como dato. Una vez producidas esas secuencias por el estado
coinductivo, su combinación y la frontera de la normalización quedan
determinadas sin consultar alfa.

El procedimiento finito conservado en el corpus propaga las cinco
fronteras de \(\mathcal C\) y conserva solamente el prefijo común a
todas ellas. Si dos profundidades han coalescido antes de una frontera,
la recurrencia determinista de derecha a izquierda da los mismos bloques
a su izquierda. Ésa es la compatibilidad con el truncamiento. No se
escoge una frontera porque dé cifras favorables, ni se deduce la
compatibilidad a toda profundidad de haber observado una corrida larga.

La comparación inicial muestra por qué esta precaución importa. En la
prolongación estabilizada se obtiene \(c_{13}^{\infty}=-1\), mientras
que la ventana aislada imponía \(c_{13}=0\). Así,
\[
 884+497-117-601-1=662
\]
en el último bloque de la primera vuelta. Las primeras once tríadas no
cambian, pero la prolongación comienza
\[
 (007,297,352,569,283,800,997,285,105,472,380,662,999,\ldots).
\]
Una aparición del bloque 999 no viola la convención canónica: lo
excluido es una cola que sea finalmente constante 999. La ventana con
frontera cero termina en 663; el truncamiento de la sección completa
termina en 662. Ambas afirmaciones describen operaciones diferentes y
exactas. No deben fusionarse en una sola notación de proyección.

\subsubsection{Definición operativa y ecuación aditiva completa}
\label{subsec:alpha-definicion}

\begin{definition}[Coordenada alfa de la clausura]
\label{def:alpha-operativa}
La vía entera determina la coordenada
\begin{equation}
 \alpha_A=\sum_{j\geq1}A_jB^{-j}
 =p+e_0-f-\kappa_{\mathrm{per}},
 \label{eq:alpha-via-a}
\end{equation}
donde \((A_j,c_j)\) es la normalización canónica compatible de
\eqref{eq:alpha-recurrencia}. En la realización escalar de las
coordenadas HMT,
\begin{equation}
 \boxed{\alpha_A=
 \pi+e-\varphi
 -4-\kappa_{\mathrm{per}}.}
 \label{eq:alpha-identidad-completa}
\end{equation}
\end{definition}

El entero cuatro no se calibra: es
\(3+2-1\), combinación de las partes enteras de las tres coordenadas.
La base mil procede de la carta cúbica y el período doce del registro dodecafásico. Los
signos son los del canal orientado y del cierre, y la frontera remota es
la que impone la sección compatible. Éstos son los orígenes efectivos de
los coeficientes de la ecuación.

En este nivel puede describirse alfa como la coordenada escalar del
defecto de clausura entre las tres realizaciones regionales y el retorno
racional dodecafásico, calculado con acarreo compatible. El término
\emph{defecto} designa una diferencia algebraica determinada, no una
pérdida arbitraria ni una imperfección del sistema. Esta descripción no
suprime la genealogía: sin las regiones, el registro firmado, la
orientación y la ley de acarreo, la igualdad escalar no sería una
exposición suficiente de la construcción.

También deben mantenerse separados tres objetos que la notación de alfa
puede acompañar en lectores posteriores: \(A\) es una palabra de
bloques; el soporte negativo de preacarreo se obtiene del signo de los
totales antes de normalizar; y un vector radial excepcional pertenece a
un codominio de incidencia. Ninguno de ellos es intercambiable con el
valor real \(\alpha_A\).

\paragraph{Balance firmado y suma con cociente entrante.}
\label{ampl-alfa-z-s-aclaracion}
La división posicional consta de dos etapas consecutivas. Con
\(K_j^{\mathrm{per}}=K_{1+((j-1)\bmod12)}\), que designa la
prolongación periódica denotada también por \(K_j\) en
\eqref{eq:alpha-recurrencia}, el balance anterior al transporte del
cociente es
\[
 Z_j:=P_j+E_j-\Phi_j-K_j^{\mathrm{per}}.
\]
La variable $s_j$ de \eqref{eq:alpha-recurrencia} designa la suma que ya
incluye la contribución entrante, por lo que
\[
 s_j=Z_j+c_{j+1},\qquad
 A_j=Z_j+c_{j+1}-1000c_j.
\]
La configuración incidencial
$H^-_\alpha$ utiliza los signos del balance $Z_j$ de la ventana inicial;
la determinación de las cifras utiliza además los cocientes de enlace
$c_{j+1}$ y $c_j$. Esta distinción conserva conjuntamente el soporte
firmado y la normalización de la expansión.

% Borrador focal de adición, NO demostración de la igualdad de dos vías.
% Puede incorporarse su proposición positiva sin alterar fórmulas actuales.
% La nota PROCEDENCIA_Y_ALCANCE.md delimita la compatibilidad no reconstruida.
\paragraph{Dependencia exacta respecto de la frontera de truncamiento.}
\label{ampl-alfa-frontera-exacta}
Las coordenadas regionales y el registro dodecafásico proceden de la
generación común APP--TRIT--TPK. Una vez producidas, su composición
posicional posee una dependencia explícita respecto de la profundidad
de lectura. Sean $B=1000$,
\[
 Z_j=P_j+E_j-\Phi_j-K_j^{\mathrm{per}},\qquad
 y_N=\sum_{j=1}^{N}Z_jB^{-j},\qquad
 R_{N+1}(Z)=\sum_{r\geq0}Z_{N+1+r}B^{-r-1}.
\]
La serie completa de \eqref{eq:alpha-serie} satisface exactamente
\[
 \alpha_A=y_N+B^{-N}R_{N+1}(Z),\qquad
 -2<R_{N+1}(Z)<2.
\]
Esta identidad reúne la contribución de la ventana y la de su
prolongación, ambas expresadas en las mismas coordenadas.

\begin{proposition}[Transporte de la frontera y cilindro canónico]
\label{ampl-alfa-prop-frontera}
Sea $A_j^{(N,t)}$ la salida de la división posicional finita con
$c_{N+1}=t\in\mathcal C=\{-2,-1,0,1,2\}$, y sea
$a_N^{(t)}=\sum_{j=1}^{N}A_j^{(N,t)}B^{-j}$. Para los registros del
dominio considerado, se cumple
\begin{equation}
 \alpha_A-a_N^{(t)}
 =B^{-N}\bigl(R_{N+1}(Z)-t\bigr),
 \qquad
 |\alpha_A-a_N^{(t)}|<4B^{-N}.
 \label{ampl-alfa-eq-frontera}
\end{equation}
El cociente de la prolongación completa es
$c_{N+1}^{\infty}=\lfloor R_{N+1}(Z)\rfloor$. Con esta frontera,
\[
 0\leq\alpha_A-a_N^{(c_{N+1}^{\infty})}<B^{-N},
\]
de modo que la ventana obtenida es el prefijo canónico de orden $N$.
\end{proposition}
\begin{proof}
El telescopado de \eqref{eq:alpha-identidad-local} da
$a_N^{(t)}=y_N+tB^{-N}-c_1$. Como $Z_1=7$, cada estado entrante
pertenece a $\mathcal C$ y el primer dividendo queda entre $5$ y $9$,
se tiene $c_1=0$. Sustituir esta igualdad en la descomposición de la
serie prueba \eqref{ampl-alfa-eq-frontera}. La cota resulta de
$R_{N+1}(Z)\in(-2,2)$ y $t\in\mathcal C$.

Para la normalización canónica completa,
$c_{N+1}^{\infty}=R_{N+1}(Z)-R_{N+1}(A)$ y
$R_{N+1}(A)\in[0,1)$. La integralidad del cociente implica
$c_{N+1}^{\infty}=\lfloor R_{N+1}(Z)\rfloor$. La fórmula de error
se convierte entonces en
$\alpha_A-a_N^{(c_{N+1}^{\infty})}=B^{-N}R_{N+1}(A)$, que prueba
la pertenencia al cilindro canónico.
\end{proof}

La proposición identifica el significado de la frontera derecha: es la
parte entera de la cola firmada expresada en la escala de la posición
siguiente. El cálculo sobre las cinco fronteras conserva esta información
como una familia finita de posibilidades hasta que los bloques a la
izquierda se estabilizan. La contracción cuantitativa procede de la
profundidad y del transporte de la cola; su objeto es la normalización de
las coordenadas ya generadas.

\paragraph{Dominio graduado del resolvente analítico.}
\label{ampl-alfa-dominio-resolvente}
La composición analítica exhibida en
\S\ref{subsec:alpha-resolvente} actúa sobre
$\operatorname{span}\{e_7,e_8,e_9\}$ y satisface $Ne_9=0$. Su
evaluación produce exactamente $T_{7:9}$; el grado $9$ es la última
componente de ese dominio. La identidad $N^3=0$ explica por qué el
resolvente se expresa mediante tres sumandos. Para describir una
prolongación en grados superiores hay que especificar, además, el
transporte entre dominios graduados sucesivos y su acción sobre la
componente inicial de cada dominio. La nilpotencia local y la
compatibilidad global desempeñan así funciones matemáticas distintas.


\subsubsection{Determinación constructiva y función del registro dodecafásico}
\label{alpha-ampl-dependencias}

La definición anterior permite precisar qué significa determinar alfa desde
la estructura. El dominio de la operación contiene cuatro registros de
procedencia conocida: las tres secuencias regionales y el registro integral
de fase. La normalización incorpora, además, una condición de frontera
compatible con la prolongación. Una vez fijados esos datos, la división
euclídea determina cada resto y su cociente. La proposición
\ref{prop:alpha-normalizacion} extiende esa determinación al sistema
completo. La independencia respecto de un valor físico objetivo es, por
tanto, una propiedad del orden de construcción: el valor que se compara al
final no interviene en la selección de ninguno de los argumentos de esta
aplicación.

Esta independencia causal debe distinguirse de la independencia algebraica
entre números. La primera se establece identificando el dominio, las
operaciones y la procedencia de los coeficientes; la segunda sería una
afirmación sobre relaciones polinómicas. El resultado utilizado aquí es el
primero. La construcción determina una coordenada mediante registros
producidos por APP--TRIT--TPK y permite comprobar posteriormente las
relaciones que satisface. Ese orden hace posible utilizar una coordenada ya
generada en una operación ulterior sin convertirla en un parámetro externo.

El registro $K$ tiene una función más precisa que la de una corrección
numérica añadida a una fórmula. Sus doce posiciones proceden de la
transformación reversible del registro firmado. La ecuación
$U=\mathcal H_{12}K$ conserva la posibilidad de volver a ese registro,
mientras las diferencias transversales y la suma total permiten otra
reconstrucción. En la evaluación periódica, el orden de las posiciones
determina el peso posicional de cada componente. El mismo conjunto de
componentes, dispuesto en otro orden, produce en general otra coordenada
racional. Por ello el origen de fase y la orientación forman parte de la
definición del número que interviene en \eqref{eq:alpha-identidad-completa}.

La normalización de alfa compone esa información con las regiones. Su
identidad local distingue dos operaciones sucesivas. Primero se forma el
balance orientado de bloques; después se representa ese balance mediante
un resto canónico y un cociente que enlaza posiciones contiguas. La primera
operación conserva el signo del balance y permite la extracción posterior
de un soporte. La segunda permite formar los prefijos de la coordenada
real. Ambas lecturas permanecen relacionadas porque el estado conserva los
datos anteriores y posteriores a la normalización. El soporte firmado y la
expansión posicional son, así, realizaciones diferentes de una composición
aritmética explícita.

La identidad telescópica expresa el contenido global de esta composición.
Los cocientes interiores se cancelan por pares al evaluar un prefijo,
mientras sus extremos permanecen en la igualdad. Esa cancelación es una
propiedad de la suma ponderada; no elimina los cocientes del estado ni
permite suprimirlos antes de efectuar la suma. La frontera derecha de una
ventana recibe la contribución de la continuación. Por esa razón, la
evaluación de una ventana aislada y la restricción de una sección
prolongada son operaciones distintas. El cambio del bloque final $663$ a
$662$, demostrado anteriormente, ofrece una realización explícita de esa
diferencia sin alterar los bloques ya estabilizados a su izquierda.

En este sentido, la bidireccionalidad tiene una formulación aritmética
concreta. La generación proporciona prolongaciones compatibles; la
normalización transmite la información entera de su frontera desde las
posiciones de menor peso hacia las de mayor peso. La ecuación conserva la
contribución de ambas direcciones mediante sus índices y sus condiciones
de contorno. La reversibilidad en la fibra fijada se comprueba reconstruyendo
$K$ a partir de los demás registros y de los cocientes completos. Esa
inversión certifica la consistencia de la operación original; la dirección
genealógica continúa siendo la que produce primero el registro y después
la coordenada de alfa.

La segunda determinación tiene un contenido complementario. En ella, el
objeto inmediato es un balance de vacancias y el parámetro aparece como
su raíz en un dominio aislante. La vía posicional proporciona una
normalización canónica de registros; la vía analítica estudia la anulación
de un operador. La exposición de ambas requiere conservar precisamente
esa diferencia de funciones. Sus relaciones con el estado común han de
expresarse mediante las aplicaciones de refinamiento y la comparación de
sus realizaciones completas, tal como se desarrolla en la subsección
correspondiente. El acuerdo de una aproximación finita constituye un
control localizado; la identificación de las construcciones completas
tiene el dominio y el alcance de su propio enunciado.

Finalmente, esta organización explica la posición de alfa en el artículo.
Su ecuación aparece después de la generación regional y de la construcción
de $K$, porque utiliza ambos resultados. La sección electrónica utiliza
después esa coordenada como un factor de su composición. La prolongación
excepcional recibe los registros y sus soportes, donde se conserva la
información de incidencia. Estas dos continuaciones retoman una misma
procedencia y desarrollan funciones matemáticas distintas: evaluación de
una sección central y construcción de una configuración de frontera.
El desarrollo de cada una permite seguir la estructura a través de la
ecuación, en lugar de limitar la comparación a su valor escalar.


\subsection{Segunda vía: determinación analítica mediante vacancias}
\subsubsection{Invariantes del transporte y resolvente graduado}
\label{subsec:alpha-via-b}

La vía analítica recibe invariantes del estado orientado, no la salida
\(\alpha_A\). Su truncamiento de orden seis es
\begin{equation}
 P_6(x)=2\pi x-\frac74x^2
 +\frac{x^3}{2\pi}
 +\frac{x^4}{20}-\frac{2x^5}{21}-\frac{x^6}{46},
 \label{eq:alpha-p6}
\end{equation}
y la condición de vacancia utiliza
\[
 \delta=\log_{10}\frac{10}{9}.
\]
La razón \(10/9\) compara la carta decimal con las nueve posiciones
APP de la construcción nonádica. Su logaritmo pertenece a esta lectura
posterior de compactificación; no es un valor objetivo de alfa.

La genealogía de los coeficientes debe conservarse completa, pero con su
fuerza lógica exacta. Los desarrollos recientes reúnen una firma marcada
del estado. La matriz de sincronización y el censo asociado son
\[
 D_{\mathrm{syn}}=\begin{pmatrix}85&112\\41&52\end{pmatrix},
 \qquad N_9=43.
\]
De su determinante se extrae
\[
 a=-\frac{\det D_{\mathrm{syn}}}{N_9}
   =\frac{172}{43}=4.
\]
El reloj inducido sobre los huecos cortos tiene espectro secundario
\(\{6,7\}\); su extremo superior da \(b=7\). La cardinalidad
trítica \(m=3\) determina
\[
 t_*=mb=21,\qquad k_*=mb-1=20.
\]
La lectura de vacancia en la carta mil da
\[
 v_-=\lfloor1000\delta\rfloor=45,\qquad
 v_+=\lceil1000\delta\rceil=46.
\]
Estos enteros son antecedentes tipados de la evaluación del truncamiento. No se
obtienen de su raíz ni de una aproximación metrológica a ella.

La rama inferior coincide además con el determinante del bloque APP
\[
 B_c=\begin{pmatrix}7&2\\2&7\end{pmatrix},\qquad\det B_c=45.
\]
El rango transversal del registro del capítulo anterior da
\(r_{\mathrm{dod}}=11\), y la semicorona de fase proporciona
\(f_{\mathrm{ph}}=8\binom62=120\). En particular,
\(45r_{\mathrm{dod}}=495\).

Con \(\omega=2\pi\), la firma
\[
 \mathcal I_9=(\omega,m,a,b,k_*,t_*,v_-,v_+,f_{\mathrm{ph}},r_{\mathrm{dod}})
\]
alimenta el lector graduado
\begin{align}
 \mathfrak E_9(\mathcal I_9)(x)
 ={}&\omega x-\frac ba x^2+\omega^{-1}x^3
 +\frac{x^4}{k_*}-\frac{2x^5}{t_*}-\frac{x^6}{v_+}\notag\\
 &-\frac{x^7}{f_{\mathrm{ph}}}+\frac{x^8}{v_-}
 +\frac{2x^9}{v_-r_{\mathrm{dod}}}.
 \label{eq:alpha-firma-coeficientes}
\end{align}
La sustitución recupera los coeficientes de \(P_6\) y de su
prolongación de orden nueve. El par \(45/46\) atraviesa la separación
entre ambos: la rama superior ocupa el grado seis y la inferior reaparece
en los grados ocho y nueve.

La factorización de \eqref{eq:alpha-firma-coeficientes} está formulada
para un estado que conserva origen, orden de canales y orientación. La
naturalidad respecto de morfismos que preservan esos datos no equivale a
la unicidad de todo lector después de olvidarlos. El dominio de esta
factorización está determinado por la firma
\eqref{eq:alpha-firma-coeficientes}; la prolongación en todos los grados
se construye en la sección~\ref{subsec:alpha-publicaciones-completas},
con su clase de recurrencia y su dependencia de la precoordenada del estado
expresamente definidas.

\subsubsection{Resolvente nilpotente en los grados 7, 8 y 9}
\label{subsec:alpha-resolvente}

La primera prolongación admite una construcción operatoria finita
especialmente explícita. Sobre el espacio con base \(e_7,e_8,e_9\),
sea
\[
 Ne_7=-\frac83e_8,\qquad Ne_8=\frac2{11}e_9,\qquad Ne_9=0,
 \qquad v_7=-\frac1{120}e_7.
\]
El símbolo \(N\) designa aquí un operador nilpotente, no una profundidad
de truncamiento. Su dominio y graduación están fijados antes de evaluar
el polinomio en un número.

\begin{lemma}[Prolongación nilpotente exacta]
\label{lem:alpha-resolvente}
Se cumple \(N^3=0\) y
\[
 (I-N)^{-1}v_7
 =-\frac{e_7}{120}+\frac{e_8}{45}+\frac{2e_9}{495}.
\]
La evaluación \(e_n\mapsto x^n\) produce
\begin{equation}
 T_{7:9}(x)=-\frac{x^7}{120}+\frac{x^8}{45}+\frac{2x^9}{495}.
 \label{eq:alpha-t79}
\end{equation}
\end{lemma}

\begin{proof}
El operador aumenta el grado y se anula en el último vector de la base,
por lo que su cubo es cero. La inversa es el resolvente finito
\(I+N+N^2\). Además,
\[
 Nv_7=\frac1{45}e_8,\qquad N^2v_7=\frac2{495}e_9.
\]
Sumar los tres términos y evaluar da la fórmula.
\end{proof}

Aquí los denominadores 120, 45 y 495 no son ajustes decimales: proceden
de la semicorona, del determinante APP y de su producto por el rango
dodecafásico. Las transferencias \(-8/3\) y \(2/11\), junto con el
signo y el grado de la semilla, pertenecen al lector tipado. Conservar
sólo la simetría del ciclo y sustituir este operador por otro operador
simétrico no conserva necesariamente los coeficientes.

\subsubsection{Raíz de multiplicidad uno del truncamiento de orden nueve}
\label{subsec:alpha-jet}

Definamos
\[
 P_9=P_6+T_{7:9},\qquad E_{21/22}^{(9)}(x)=P_9(x)-\delta,
 \qquad I_\alpha=(0,10^{-2}).
\]
Este operador finito tiene un resultado propio de existencia y unicidad.

\begin{proposition}[Raíz de multiplicidad uno del truncamiento]
\label{prop:alpha-raiz-jet}
El operador \(E_{21/22}^{(9)}\) tiene una única raíz
\(\xi_9\) en \(I_\alpha\), y esa raíz es simple. Se obtiene el aislamiento racional
\begin{equation}
 \frac{729735256928380099}{10^{20}}
 <\xi_9<\frac{729735256928380100}{10^{20}}.
 \label{eq:alpha-intervalo-jet}
\end{equation}
\end{proposition}

\begin{proof}
Las cotas de las coordenadas internas
\(3.14<\pi<3.15\) y
\(0.045<\delta<0.046\) son suficientes para la existencia. En cero,
\(E_{21/22}^{(9)}(0)=-\delta<0\). En \(10^{-2}\), la contribución
lineal menos los términos negativos es mayor que \(0.0626\), luego el
valor del operador es positivo.

Para \(0\leq x\leq10^{-2}\), al descartar los términos positivos de
la derivada se obtiene
\[
 (E_{21/22}^{(9)})'(x)
 \geq2\pi-\frac72\cdot10^{-2}
 -\frac{10}{21}\cdot10^{-8}-\frac6{46}\cdot10^{-10}
 -\frac7{120}\cdot10^{-12}>6.24.
\]
La continuidad da existencia; la cota derivativa da unicidad y
simplicidad. El intervalo más estrecho
\eqref{eq:alpha-intervalo-jet} resulta de una evaluación racional dirigida.
Sean \(x_-\) y \(x_+\) sus extremos. Para calcular
\(\delta=\log(10/9)/\log 10\), se emplea, para \(z>1\),
\[
 L_N(z)=2\sum_{j=0}^{N-1}\frac{y^{2j+1}}{2j+1},\qquad
 y=\frac{z-1}{z+1},\qquad
 0<\log z-L_N(z)<\frac{2y^{2N+1}}{(2N+1)(1-y^2)}.
\]
El resto se acota mediante la serie geométrica de los términos positivos.
Se aplica esta fórmula a \(z=10/9\) y \(z=3\), con \(N=520\),
usando \(\log 10=\log(10/9)+2\log3\). La evaluación de
\(\pi\) utiliza el cilindro regional generado de mil cifras, incluyendo
su anchura \(10^{-1000}\). La aritmética de intervalos en
\eqref{eq:alpha-firma-coeficientes} da
\[
 -\frac{4559}{10^{23}}<E_{21/22}^{(9)}(x_-)< -\frac{4558}{10^{23}},
 \qquad
 \frac{1698}{10^{23}}<E_{21/22}^{(9)}(x_+)<\frac{1699}{10^{23}}.
\]
Estos signos y la monotonía probada establecen el aislamiento. Los extremos
son resultados de la evaluación del operador; sus coeficientes proceden de
la firma previamente construida.
\end{proof}

El resultado concierne a \(\xi_9\). No autoriza a escribir
\(\xi_9=\alpha_A\) como una igualdad exacta de números completos.
La raíz de un truncamiento puede compartir un prefijo con la raíz del
operador prolongado y diferir en una posición posterior. La simplicidad
permite controlar el movimiento de la raíz bajo perturbaciones; no
anula esas perturbaciones.

\subsection{Compatibilidad de las dos determinaciones a toda profundidad}
\label{subsec:alpha-limite-dos-vias}

\begin{remark}[Truncamiento, prolongación y compatibilidad]
Las operaciones de los grados \(6\) a \(9\) determinan un jet finito.
La normalización posicional determina la coordenada completa mediante
su sucesión compatible de acarreos. La proposición siguiente expresa
el criterio general de identificación por cilindros. La
sección~\ref{subsec:alpha-publicaciones-completas} construye después
una carta analítica correlacionada de todos los grados, con recurrencia,
convergencia uniforme, unicidad de la raíz y compatibilidad proyectiva.
La construcción recibe la coordenada compartida del estado y la clase
de prolongaciones declarada; conserva así la diferencia entre una
representación analítica de esa coordenada y una selección independiente.
\end{remark}

La vía analítica completa se organiza en una familia de lectores
graduados compatibles
\[
 E_{21/22}^{(6)}\longleftarrow E_{21/22}^{(9)}
 \longleftarrow E_{21/22}^{(12)}\longleftarrow\cdots.
\]
El símbolo
\[
 E_{21/22}=\varprojlim_m E_{21/22}^{(6+3m)}
\]
designa ese lector completo, con sus transportes y truncamientos, no una
serie formal elegida libremente para hacerla pasar por \(\alpha_A\).
La segunda vía determina su raíz positiva única en \(I_\alpha\),
denotada por \(\alpha_B\), con la unicidad y la compatibilidad del
lector completo demostradas en la sección~\ref{subsec:alpha-publicaciones-completas}.

Las demostraciones de la sección~\ref{subsec:alpha-publicaciones-completas}
establecen que ambos sistemas determinan, a cada profundidad, el mismo
cilindro superviviente \(C_N^\alpha\), y que sus diámetros tienden a
cero. Conviene exponer por separado la consecuencia matemática de esta
compatibilidad.

\begin{proposition}[Identificación de los límites compatibles]
\label{prop:alpha-dos-limites}
Si los valores de las vías entera y analítica pertenecen a los
mismos cilindros \(C_N^\alpha\) para todo \(N\), y
\(\operatorname{diam}C_N^\alpha\to0\), entonces
\[
 \alpha_A=\alpha_B=\alpha.
\]
\end{proposition}

\begin{proof}
Para todo \(N\), la distancia entre ambos valores está acotada
por \(\operatorname{diam}C_N^\alpha\). Hacer tender \(N\) a infinito
da distancia cero. La existencia de las secciones compatibles asegura
que no se trata de la intersección vacía de un sistema de aproximaciones.
\end{proof}

Esta prueba identifica la consecuencia de los cuadrados de compatibilidad.
Su construcción a profundidad arbitraria aparece en la
sección~\ref{subsec:alpha-publicaciones-completas}. Las evaluaciones
finitas contrastan realizaciones de esa construcción; la recurrencia
y sus cotas uniformes establecen el paso al límite.

La diferencia puede expresarse algebraicamente. Sea
\(F=\mathbb Q(\pi,\delta)\) y sea
\(j_9:F[[x]]\to F[x]/(x^{10})\) el truncamiento. Entonces
\[
 \ker j_9=x^{10}F[[x]].
\]
Todos los operadores
\[
 E_h(x)=E_{21/22}^{(9)}(x)+x^{10}h(x)
\]
tienen el mismo truncamiento de orden nueve. Sin embargo, tomando
\(h=0\) y \(h=1/729\),
\[
 E_{1/729}(\xi_9)=\frac{\xi_9^{10}}{729}>0,
\]
de modo que no comparten esa raíz. La observación no permite elegir una
cola ajena al transporte HMT; precisamente demuestra por qué la cola
generada y su compatibilidad deben conservarse. Tampoco basta una serie
formal arbitraria para hablar de raíces analíticas sin su dominio de
convergencia o sin la realización del sistema de cilindros.
\subsection{Compatibilidad entre las publicaciones dodecafásica y analítica de alfa}
\label{subsec:alpha-publicaciones-completas}

La compatibilidad se establece a toda profundidad. La publicación
dodecafásica reversible y la representación analítica de la
coordenada de alfa
\(E_{21/22}\) son dos publicaciones internas con codominios distintos,
producidas a partir del mismo estado enriquecido APP--TRIT--TPK. El jet
\(E_{21/22}^{(9)}\) es un testigo finito exacto, pero no determina por sí solo
la cola. La completación que se construye después no se deduce del jet
desnudo: recibe la coordenada compartida del estado y publica su imagen
analítica mediante una recurrencia explícita. Por ello no constituye una
selección causal independiente; constituye una representación analítica
tipada del mismo punto generado.

Esta separación conserva el orden causal
\[
 \mathrm{APP}\longrightarrow\mathrm{TRIT}\longrightarrow\mathrm{TPK}
 \longrightarrow\mathcal X_{\alpha}^{\mathrm{enr}}
 \longrightarrow (P,E,\Phi,U_{12}^{\mathrm{sgn}},K)
 \longrightarrow\alpha_{\mathrm{HMT}}.
\]
Las coordenadas \(P,E,\Phi\) y el registro \(K\) son salidas anteriores del
mismo estado enriquecido. Su uso en la clausura de alfa es, por tanto, una
composición interna posterior y no una inyección de constantes
convencionales.

Para \(r\in6\mathbb N\), \(r\ge36\), sea
\(\widehat{\mathcal X}^{\mathrm{enr},\alpha}_r\) el sector de estados
truncados que prolongan la semilla interna de alfa fijada en \(R_{36}\). Se
escribe
\[
 x_r^\alpha=(\widetilde x_r,\chi_\alpha,
              \varepsilon_r^\alpha,N_r^\alpha),
\]
donde \(\chi_\alpha\) es el carácter interno del sector, no el valor escalar
publicado. Cada estado conserva residuo, cociente, acarreo, hoja, orientación,
fase, ruta, cilindro, frontera, firma, registro, memoria, supervivencia,
procedencia y profundidad. Si
\[
 d_r^\alpha=\operatorname{Dig}_{729}(\varepsilon_r^\alpha),\qquad
 \varepsilon_{r+6}^\alpha=729\varepsilon_r^\alpha-d_r^\alpha,
 \qquad N_{r+6}^\alpha=729N_r^\alpha+d_r^\alpha,
\]
la transición efectiva es
\[
 \mathcal U_r^\alpha=
 \operatorname{Upd}_r^\alpha\circ
 \operatorname{Tra}_r^\alpha\circ
 \operatorname{Sel}_r^\alpha,
 \qquad
 \operatorname{Ext}^{\alpha,\to}_{r+6,r}
 =\operatorname{graph}(\mathcal U_r^\alpha).
\]
La funcionalidad de estas extensiones, con la semilla \(R_{36}\) fijada,
produce una única sección compatible: la fase retorna mientras memoria y
profundidad avanzan. El dominio global de ambas determinaciones es
\begin{equation}
 \mathcal X_\alpha^{\mathrm{enr}}:=\widehat\Omega_\alpha
 :=\varprojlim_{r\ge36}
 \bigl(\widehat{\mathcal X}^{\mathrm{enr},\alpha}_r,
       \operatorname{Ext}^{\alpha,\to}_{r+6,r}\bigr).
 \label{eq:alpha-dominio-enriquecido-completo}
\end{equation}

En esta sección, \(\kappa_{\mathrm{per}}\) es la lectura periódica \(\kappa_{\mathrm{ag}}\) del registro ya construido; ambas notaciones designan la misma fracción en la carta de base mil fijada.

\subsubsection{La clausura dodecafásica como sección completa}

Para \(N\geq1\), pongamos
\[
 \mathcal W_N:=\{0,\ldots,999\}^{N},
 \qquad
 \mathsf T_{\alpha,N}(\mathsf s):=(A_1,\ldots,A_N),
 \qquad \mathsf s\in\mathcal X_{\alpha}^{\mathrm{enr}},
\]
donde los bloques \(A_j\) se obtienen mediante
\eqref{eq:alpha-recurrencia} con la frontera compatible
\(c_{N+1}^{\infty}=\lfloor R_{N+1}(Z)\rfloor\). Si \(N'\geq N\), sea
\(\rho_{N',N}:\mathcal W_{N'}\to\mathcal W_N\) la restricción a los
primeros \(N\) bloques.

\begin{theorem}[Salida dodecafásica completa y compatibilidad proyectiva]
\label{thm:alpha-salida-dodecafasica-completa}
La familia \((\mathsf T_{\alpha,N})_{N\geq1}\) es compatible:
\begin{equation}
 \rho_{N',N}\circ\mathsf T_{\alpha,N'}
 =\mathsf T_{\alpha,N}
 \qquad(N'\geq N).
 \label{eq:alpha-via-a-proyectiva-exacta}
\end{equation}
Sus cilindros canónicos
\begin{equation}
 C_N^A(\mathsf s):=
 \left[
  \sum_{j=1}^{N}A_j1000^{-j},
  \sum_{j=1}^{N}A_j1000^{-j}+1000^{-N}
 \right)\cap I_\alpha
 \label{eq:alpha-cilindros-via-a-exactos}
\end{equation}
son anidados, tienen diámetro a lo sumo \(1000^{-N}\) y satisfacen
\begin{equation}
 \bigcap_{N\geq1}C_N^A(\mathsf s)
 =\{\alpha_A(\mathsf s)\},
 \qquad
 \alpha_A
 =\pi_{\mathrm{HMT}}+e_{\mathrm{HMT}}-\varphi_{\mathrm{HMT}}
  -4-\kappa_{\mathrm{per}}.
 \label{eq:alpha-salida-a-completa}
\end{equation}
En consecuencia, la primera vía publica una única sección infinita, que
denotamos \(\alpha_{\mathrm{HMT}}:=\alpha_A\).
\end{theorem}

\begin{proof}
La proposición~\ref{prop:alpha-normalizacion} proporciona la expansión
canónica única y la sucesión acotada única de acarreos. La identidad local
\eqref{eq:alpha-identidad-local} y su telescopado
\eqref{eq:alpha-telescopia} conservan la frontera remota; la proposición
\ref{ampl-alfa-prop-frontera} prueba que restringir una ejecución prolongada
recupera exactamente el prefijo anterior. Esto da
\eqref{eq:alpha-via-a-proyectiva-exacta}. La pertenencia de la suma completa
a cada cilindro se sigue de la identidad de colas, y
\(\operatorname{diam}C_N^A\leq1000^{-N}\to0\); la intersección contiene un
solo punto. Finalmente, \eqref{eq:alpha-identidad-completa} identifica ese
punto con la expresión de \eqref{eq:alpha-salida-a-completa}. Todos sus
coeficientes y signos se fijan antes de la publicación de alfa.
\end{proof}

La ventana \(\alpha_{12}^{(0)}\), cerrada con \(c_{13}=0\), sigue siendo un
cálculo finito exacto; no debe confundirse con la restricción de la sección
completa, cuya frontera compatible satisface \(c_{13}^{\infty}=-1\). En
particular,
\begin{equation}
 \rho_{36,12}(\alpha_{A,36})=\alpha_{A,12}
 \label{eq:alpha-rho-36-12-proyectiva}
\end{equation}
es una igualdad de palabras compatibles, no una igualdad entre palabras de
longitudes distintas ni una identificación de la ventana aislada con el
límite.


% Representación analítica correlacionada de alfa: acción, convergencia e
% intervalos.
\subsubsection{Representación analítica de la coordenada de alfa}
\label{subsec:alpha-lector-completo-rev08}

La prolongación analítica se especifica mediante todos sus coeficientes.
Sea \(\mathsf s\in\mathcal X_\alpha^{\mathrm{enr}}\). Antes
de la normalización decimal, el mismo estado ya publica la precoordenada
\begin{equation}
 a(\mathsf s):=p(\mathsf s)+e_0(\mathsf s)-f(\mathsf s)
                 -\kappa_{\mathrm{per}}(\mathsf s)\in I_\alpha,
 \label{eq:alpha-precoordenada-comun}
\end{equation}
con \(p=\pi_{\mathrm{HMT}}-3\),
\(e_0=e_{\mathrm{HMT}}-2\) y
\(f=\varphi_{\mathrm{HMT}}-1\). Equivalentemente,
\(a=\pi_{\mathrm{HMT}}+e_{\mathrm{HMT}}-
\varphi_{\mathrm{HMT}}-4-\kappa_{\mathrm{per}}\). No se introduce aquí
\(\alpha\) convencional: \(p,e_0,f\) son los tres caracteres regionales ya
generados y \(\kappa_{\mathrm{per}}\) es la lectura racional del registro
terminal. La clausura entera aplica a
\eqref{eq:alpha-precoordenada-comun} la normalización de
\eqref{eq:alpha-recurrencia}; la carta que sigue publica sobre la misma
precoordenada una representación analítica compatible. Su grafo de
dependencia recibe \((p,e_0,f,\kappa_{\mathrm{per}})\) directamente: no
contiene una arista \(\alpha_A\to\lambda\), aunque la igualdad demostrada más
abajo permita identificar después ambos nombres de publicación.

Escribamos \(q=1/729\), sea \(F_{\mathsf s}\) el cuerpo ordenado generado por
las coordenadas de \(\mathsf s\), y pongamos
\(E_{9,\mathsf s}(X)=E_{21/22}^{(9)}(\mathsf s;X)\). Definimos el coeficiente
de enlace
\begin{equation}
 \lambda(\mathsf s):=
 -\frac{E_{9,\mathsf s}(a(\mathsf s))
         \bigl(1-q\,a(\mathsf s)^3\bigr)}{a(\mathsf s)^{10}}.
 \label{eq:alpha-lambda-estado}
\end{equation}
Todos sus argumentos pertenecen al estado antes del reconocimiento
metrológico. En particular, \(\lambda\) no se obtiene ajustando una cadena
decimal objetivo. La ecuación~\eqref{eq:alpha-lambda-estado} debe leerse con
su tipo exacto: una vez que la primera publicación del estado ha producido
\(a(\mathsf s)\), fija la única prolongación de la clase geométrica que se
define a continuación. No constituye una segunda selección independiente de
\(a(\mathsf s)\), sino una representación analítica correlacionada del mismo
estado enriquecido.

En efecto, para \(\mu\in F_{\mathsf s}\) pongamos
\[
 E_{\mathsf s,\mu}(X)
 :=E_{9,\mathsf s}(X)+\mu\frac{X^{10}}{1-qX^3}.
\]
Como \(0<a(\mathsf s)<10^{-2}\) y \(qa(\mathsf s)^3<1\), se tiene
\begin{equation}
 E_{\mathsf s,\mu}(a(\mathsf s))=0
 \quad\Longleftrightarrow\quad
 \mu=\lambda(\mathsf s).
 \label{eq:alpha-unicidad-lambda-clase-geometrica}
\end{equation}
Así, la semilla de coeficientes no contiene libertad una vez fijados el estado
y la clase de prolongaciones \(q\)-geométricas. Esta unicidad es relativa a la
clase declarada; no afirma que el jet de orden nueve seleccione por sí solo una
cola entre todas las series analíticas posibles.

La acción local sobre bloques de tres coeficientes es
\begin{equation}
 \mathcal C_\alpha:F_{\mathsf s}^{3}\longrightarrow F_{\mathsf s}^{3},
 \qquad \mathcal C_\alpha(u,v,w)=q(u,v,w),
 \qquad c^{(0)}=(\lambda(\mathsf s),0,0).
 \label{eq:alpha-accion-coeficientes}
\end{equation}
Por tanto, para todo \(n\geq0\),
\begin{equation}
 b_{10+3n}=q^n\lambda(\mathsf s),\qquad
 b_{11+3n}=b_{12+3n}=0.
 \label{eq:alpha-recurrencia-coeficientes}
\end{equation}
Ésta es la regla que calcula cada coeficiente posterior al grado nueve. No
queda un parámetro libre en cada prolongación.

Para \(M\geq0\), sea
\begin{equation}
 E_{\mathsf s,M}(X):=E_{9,\mathsf s}(X)
 +\lambda(\mathsf s)\sum_{n=0}^{M}q^nX^{10+3n}.
 \label{eq:alpha-truncamientos-completos}
\end{equation}
La función completa es
\begin{equation}
 E_{\mathsf s,\infty}(X)
 :=E_{9,\mathsf s}(X)
 +\lambda(\mathsf s)\frac{X^{10}}{1-X^3/729}.
 \label{eq:alpha-funcion-completa-explicita}
\end{equation}

Para formular el sistema proyectivo sin tipos implícitos, fijamos
\(d_M=12+3M\) y definimos la fibración de jets
\begin{equation}
 \mathfrak J_M:=
 \coprod_{\mathsf s\in\mathcal X_\alpha^{\mathrm{enr}}}
 \{\mathsf s\}\times F_{\mathsf s}[X]/(X^{d_M+1}).
 \label{eq:alpha-espacio-jets-rev08}
\end{equation}
Si \(M'\ge M\), la restricción es
\begin{equation}
 \tau_{M',M}(\mathsf s,[P]_{d_{M'}})
 :=(\mathsf s,[P]_{d_M}).
 \label{eq:alpha-truncamiento-jets-rev08}
\end{equation}

\begin{proposition}[Naturalidad de la recurrencia analítica]
\label{prop:alpha-naturalidad-recurrencia-explicita}
Sea
\[
 \Gamma_M(\mathsf s)
 =\bigl(\mathsf s,[E_{\mathsf s,M}]_{d_M}\bigr).
\]
Si \(M'\geq M\), el truncamiento \(\tau_{M',M}\) satisface
\begin{equation}
 \tau_{M',M}\circ\Gamma_{M'}=\Gamma_M,
 \qquad
 \Gamma_{E21}(\mathsf s)=\varprojlim_M\Gamma_M(\mathsf s).
 \label{eq:alpha-naturalidad-explicita}
\end{equation}
La proyección sobre los tres grados nuevos es precisamente
\(\mathcal C_\alpha^{M+1}(c^{(0)})\).
\end{proposition}

\begin{proof}
La ecuación~\eqref{eq:alpha-recurrencia-coeficientes} fija el bloque de
grados \(10+3n,11+3n,12+3n\). Al pasar de \(M\) a \(M+1\) sólo se añade
el bloque siguiente; truncarlo recupera literalmente el polinomio anterior.
La compatibilidad para \(M'\geq M\) se obtiene por composición. Así, el
límite inverso no es una sección escogida del jet desnudo, sino la órbita
única de \(c^{(0)}\) bajo \(\mathcal C_\alpha\).
\end{proof}

\subsubsection{Convergencia uniforme, raíz simple e intervalos racionales}
\label{subsec:alpha-cotas-completas-rev08}

Los cilindros regionales y el registro periódico proporcionan cotas mediante
aritmética de intervalos racionales. Para cada coordenada
\[
 c\in\{\pi_{\mathrm{HMT}},e_{\mathrm{HMT}},\varphi_{\mathrm{HMT}}\},
\]
sea \(P_c\) su prefijo nonádico certificado de mil cifras. El dato utilizado
no es el racional \(P_c\) como si fuese el valor completo, sino el cilindro
semiabierto y su cierre exterior
\[
 \mathcal C_c^{(1000)}=[P_c,P_c+\varepsilon),
 \qquad
 \overline{\mathcal C}_c^{(1000)}=[P_c,P_c+\varepsilon],
 \qquad \varepsilon=10^{-1000}.
\]
La aritmética dirigida opera conservadoramente con esos cierres. Si
\(P_\pi,P_e,P_\varphi\) denotan los tres prefijos, se obtiene
\begin{equation}
\begin{aligned}
 a_-&:=P_\pi+P_e-(P_\varphi+\varepsilon)-4-
       \kappa_{\mathrm{per}},\\
 a_+&:=(P_\pi+\varepsilon)+(P_e+\varepsilon)-P_\varphi-4-
       \kappa_{\mathrm{per}},\\
 A_{1000}&:=[a_-,a_+],
 \qquad a(\mathsf s)\in[a_-,a_+]=A_{1000},
 \qquad \operatorname{diam}A_{1000}=3\varepsilon.
\end{aligned}
\label{eq:alpha-propagacion-colas-rev08}
\end{equation}
La extensión intervalar natural de
\eqref{eq:alpha-lambda-estado} transporta después \(A_{1000}\), el cilindro
de \(\pi_{\mathrm{HMT}}\) y la cota racional dirigida de \(\delta\) a un
intervalo cerrado \(\Lambda_{1000}\ni\lambda(\mathsf s)\); no se congela
ninguna de las tres colas. En particular,
\begin{equation}
\begin{aligned}
 \frac{7297352569283800997285105472380662}{10^{36}}
 &<a(\mathsf s)<
 \frac{7297352569283800997285105472380663}{10^{36}},\\
 -8.711\mathbin{\times}10^{-6}
 &<\lambda(\mathsf s)<-8.709\mathbin{\times}10^{-6}.
\end{aligned}
 \label{eq:alpha-cotas-precoordenada-lambda}
\end{equation}
Estas desigualdades usan prefijos certificados con su cota racional de
cola; no reciben una tabla de alfa. El verificador del paquete las reproduce
desde las tres salidas nonádicas y desde la carta terminal.

Sea
\[
 r:=\frac{10^{-6}}{729}<1.372\mathbin{\times}10^{-9}.
\]
Para \(0\leq x\leq10^{-2}\), la cola posterior a \(M\) satisface
\begin{equation}
 \left|E_{\mathsf s,\infty}(x)-E_{\mathsf s,M}(x)\right|
 \leq
 8.711\mathbin{\times}10^{-26}\,\frac{r^{M+1}}{1-r}.
 \label{eq:alpha-cota-cola-uniforme}
\end{equation}
Además,
\begin{equation}
 \left|\frac{\mathrm d}{\mathrm dx}
  \left(\lambda\frac{x^{10}}{1-x^3/729}\right)\right|
 \leq 8.711\mathbin{\times}10^{-6}
 \left(\frac{10\,10^{-18}}{1-r}
 +\frac{3\,10^{-24}/729}{(1-r)^2}\right)
 <8.712\mathbin{\times}10^{-23}.
 \label{eq:alpha-cota-derivada-cola}
\end{equation}
La cola de las derivadas posterior al bloque \(M\) admite, además, la
estimación dependiente de la profundidad
\begin{equation}
 \sup_{0\le x\le10^{-2}}
 \left|E'_{\mathsf s,\infty}(x)-E'_{\mathsf s,M}(x)\right|
 \le 8.711\mathbin{\times}10^{-24}\,r^{M+1}
 \left(\frac{13+3M}{1-r}+\frac{3r}{(1-r)^2}\right)
 \xrightarrow[M\to\infty]{}0.
 \label{eq:alpha-cota-resto-derivadas-rev08}
\end{equation}

\begin{theorem}[Función límite y raíz simple única de la carta correlacionada]
\label{thm:alpha-raiz-completa-explicita}
La sucesión \(E_{\mathsf s,M}\) converge uniformemente, junto con sus
derivadas, a \(E_{\mathsf s,\infty}\) en
\(\overline I_\alpha=[0,10^{-2}]\). Se cumple
\begin{equation}
 E_{\mathsf s,\infty}(a(\mathsf s))=0,
 \qquad
 \inf_{x\in\overline I_\alpha}
 E'_{\mathsf s,\infty}(x)>6.239.
 \label{eq:alpha-raiz-y-derivada-completas}
\end{equation}
Por consiguiente, \(a(\mathsf s)\) es la raíz simple única de la función
completa en \(I_\alpha\).
\end{theorem}

\begin{proof}
La razón de dos términos consecutivos de la cola está acotada por \(r<1\),
lo que da \eqref{eq:alpha-cota-cola-uniforme}; la derivación término a
término y la suma de las dos series geométricas dan
\eqref{eq:alpha-cota-derivada-cola} y
\eqref{eq:alpha-cota-resto-derivadas-rev08}. La prueba del jet establece
\(E'_{9,\mathsf s}>6.24\) en \(\overline I_\alpha\); por tanto, la
perturbación completa deja la derivada por encima de \(6.239\).
Finalmente, al sustituir \eqref{eq:alpha-lambda-estado} en
\eqref{eq:alpha-funcion-completa-explicita}, los dos sumandos evaluados en
\(a(\mathsf s)\) se cancelan exactamente. La función es estrictamente
creciente; su única raíz es, pues, simple y coincide con esa precoordenada.
\end{proof}

Para hacer efectiva la identificación a toda profundidad, las evaluaciones
dirigidas de los extremos certifican primero
\[
 E_{\mathsf s,\infty}(0)<0<
 E_{\mathsf s,\infty}(10^{-2}),
 \qquad
 E'_{\mathsf s,\infty}\geq\mu:=6239/1000>c:=6
 \quad\hbox{en }[0,10^{-2}].
\]
Por continuidad existe una raíz en el intervalo y, por la segunda desigualdad,
es única. Sólo después de fijar esta existencia, la monotonía y la invariante
de que la horquilla contiene la raíz se aplica el rescate por el teorema del
valor medio. Partimos de \(D_0=[0,10^{-2}]\). Si
\(D_j=[r_j,s_j]\), sea \(m_j=(r_j+s_j)/2\), y sea
\([F_j^-,F_j^+]\) una inclusión racional dirigida de
\(E_{\mathsf s,\infty}(m_j)\), obtenida propagando los cilindros de
\(\pi_{\mathrm{HMT}},e_{\mathrm{HMT}},\varphi_{\mathrm{HMT}}\),
\(\delta\) y \(\lambda\), con
\begin{equation}
 0\leq F_j^+-F_j^-\leq \frac{c}{4}\,(s_j-r_j).
 \label{eq:alpha-anchura-evaluacion-rescate-rev08}
\end{equation}
Cuando \(F_j^+<0\) se conserva la mitad derecha;
cuando \(F_j^->0\), la mitad izquierda. Si
\(F_j^-\leq0\leq F_j^+\), no se intenta decidir si el punto medio es
exactamente la raíz. Sea \(w_j=s_j-r_j\). La cota uniforme
\(E'_{\mathsf s,\infty}\geq c=6\) y el teorema del valor medio implican que
cualquier raíz contenida en \(D_j\) pertenece a
\begin{equation}
 D_{j+1}=D_j\cap
 \left[m_j-\frac{w_j}{4},
             m_j+\frac{w_j}{4}\right].
 \label{eq:alpha-rescate-derivada-rev08}
\end{equation}
En efecto, sea \(\xi_j\) la raíz y sea
\(y_j=E_{\mathsf s,\infty}(m_j)\in[F_j^-,F_j^+]\). Como el recinto contiene
cero y satisface \eqref{eq:alpha-anchura-evaluacion-rescate-rev08}, se tiene
\(|y_j|\leq cw_j/4\). El teorema del valor medio proporciona
\(|m_j-\xi_j|\leq |y_j|/c\leq w_j/4\), lo que prueba
\eqref{eq:alpha-rescate-derivada-rev08}. El argumento incluye el caso
\(y_j=0\): la raíz está entonces en el punto medio, pero el algoritmo no
necesita reconocer esa igualdad para conservarla.

Si el recinto disponible no satisface
\eqref{eq:alpha-anchura-evaluacion-rescate-rev08}, se refinan primero los
cilindros fuente y se repite la evaluación; no se declara un signo ni se
acepta una contracción insuficiente. En cualquiera de las tres ramas el nuevo
intervalo conserva la raíz y tiene diámetro a lo sumo \(w_j/2\). Para cada
profundidad objetivo finita, el refinamiento dirigido permite satisfacer la
cota de anchura sin comparar de forma no certificada un número real con cero.
La monotonía demuestra inductivamente
\begin{equation}
 a(\mathsf s)\in D_{j+1}\subseteq D_j,\qquad
 \operatorname{diam}D_{j+1}\leq
 \tfrac12\operatorname{diam}D_j,
 \qquad \operatorname{diam}D_j\longrightarrow0,
 \label{eq:alpha-biseccion-racional}
\end{equation}
donde la desigualdad es una cota de contracción y no una igualdad impuesta.
El certificado ejecutable aplica
esta regla a veinticuatro niveles en base mil y comprueba que el cierre
exterior completo \(A_{1000}\), no sólo su extremo inferior, queda dentro de
cada cilindro publicado. Sus pruebas negativas rechazan un recinto demasiado
ancho, un recinto desordenado, un punto que no sea el medio, un dominio sin
cambio de signo y una horquilla degenerada. La prolongación coinductiva permite
reemplazar mil por cualquier profundidad requerida.

Concretamente, para \(S_N=1000^N\) el verificador calcula
\begin{equation}
 j_N^-:=\lfloor S_Na_-\rfloor,
 \qquad j_N^+:=\lfloor S_Na_+\rfloor
 \label{eq:alpha-indices-cilindro-dirigidos-rev08}
\end{equation}
y sólo publica el nivel si \(j_N^-=j_N^+\). Esta igualdad, que incluye el
extremo superior del cierre exterior, prueba
\(A_{1000}\subset[j_N^-/S_N,(j_N^-+1)/S_N)\) y evita atribuir a la celda
izquierda un extremo situado exactamente sobre su frontera derecha.

Sea \(C_N^A(\mathsf s)\) el cilindro entero de
\eqref{eq:alpha-cilindros-via-a-exactos}. Elegimos recursivamente
\(m(N)>m(N-1)\) de modo que
\(\operatorname{diam}D_{m(N)}<1000^{-N}/4\), y definimos
\begin{equation}
I_{B,N}:=D_{m(N)}\cap C_N^A(\mathsf s).
 \label{eq:alpha-intervalos-b-completos}
\end{equation}
Definimos
\(\ell_N:=\inf I_{B,N}\) y \(u_N:=\sup I_{B,N}\). Ambos son racionales;
el intervalo es no vacío porque los dos factores contienen
\(a(\mathsf s)\). Puesto que el cilindro es semiabierto, \(u_N\) puede no
pertenecer a \(I_{B,N}\). La evaluación en ese extremo se entiende mediante
la extensión continua de \(E_{\mathsf s,\infty}\) al cierre exterior. Así,
\begin{equation}
 I_{B,N+1}\subseteq I_{B,N}\subseteq C_N^A(\mathsf s),
 \qquad \operatorname{diam}I_{B,N}\longrightarrow0,
 \qquad E_{\mathsf s,\infty}(\ell_N)\leq0\leq
 E_{\mathsf s,\infty}(u_N).
 \label{eq:alpha-inclusion-signos-completa}
\end{equation}
La cota~\eqref{eq:alpha-cota-cola-uniforme} permite escoger un
truncamiento que conserve el signo en cada extremo donde la función completa
no se anula. Si un extremo coincide con la raíz, se utiliza la identidad de
la función completa y la pertenencia ya demostrada; no se atribuye ese mismo
signo a sus truncamientos. En efecto, el resto exacto es
\[
 E_{\mathsf s,\infty}(x)-E_{\mathsf s,M}(x)
 =\lambda(\mathsf s)x^{10}
   \frac{(qx^3)^{M+1}}{1-qx^3}.
\]
En el dominio positivo considerado, donde \(\lambda(\mathsf s)<0\), se
obtiene en la raíz
\[
 E_{\mathsf s,M}(a(\mathsf s))
 =-\lambda(\mathsf s)a(\mathsf s)^{10}
   \frac{(qa(\mathsf s)^3)^{M+1}}{1-qa(\mathsf s)^3}>0
 \qquad(M<\infty).
\]
Por ello, la certificación de los signos débiles en
\eqref{eq:alpha-inclusion-signos-completa} se refiere a
\(E_{\mathsf s,\infty}\). Su evaluación utiliza la función completa o el
truncamiento acompañado del resto con signo. Esto cubre tanto el ínfimo
perteneciente como el supremo eventualmente excluido del cilindro.

\begin{proposition}[Cuadrados de compatibilidad y cilindros comunes]
\label{prop:alpha-cuadrados-compatibilidad}
Las restricciones de palabras y de publicaciones analíticas conmutan:
\[
\begin{array}{ccc}
 \mathcal X_\alpha^{\mathrm{enr}}&
 \xrightarrow{\ \mathsf T_{\alpha,N'}\ }&\mathcal W_{N'}\\[1mm]
 \Big\Vert&&\Big\downarrow\rho_{N',N}\\[1mm]
 \mathcal X_\alpha^{\mathrm{enr}}&
 \xrightarrow{\ \mathsf T_{\alpha,N}\ }&\mathcal W_N
\end{array}
\qquad
\begin{array}{ccc}
 \mathcal X_\alpha^{\mathrm{enr}}&
 \xrightarrow{\ \Gamma_{M'}\ }&\mathfrak J_{M'}\\[1mm]
 \Big\Vert&&\Big\downarrow\tau_{M',M}\\[1mm]
 \mathcal X_\alpha^{\mathrm{enr}}&
 \xrightarrow{\ \Gamma_M\ }&\mathfrak J_M .
\end{array}
\]
La familia \(I_{B,N}\) proporciona la conexión efectiva entre ambos
cuadrados y satisface materialmente
\(I_{B,N}\subseteq C_N^A\) para todo \(N\).
\end{proposition}

\begin{proof}
El primer cuadrado es la compatibilidad de la normalización entera; el segundo
es la proposición~\ref{prop:alpha-naturalidad-recurrencia-explicita}. La
inclusión se sigue de la definición~\eqref{eq:alpha-intervalos-b-completos},
y la no vacuidad y los signos se siguen del
teorema~\ref{thm:alpha-raiz-completa-explicita}. Ningún paso identifica la
raíz de un truncamiento con la de una cola arbitraria.
\end{proof}

Definimos ahora, sin dejar implícito el lector de prefijos,
\begin{equation}
 \alpha_B(\mathsf s):=
 \operatorname*{arg\,zero}_{x\in I_\alpha}E_{\mathsf s,\infty}(x),
 \qquad
 \operatorname{pref}_N(x):=
 \frac{\lfloor1000^Nx\rfloor}{1000^N},
 \qquad
 \alpha_{B,N}:=\operatorname{pref}_N(\alpha_B).
 \label{eq:alpha-definicion-b-y-prefijos-rev08}
\end{equation}

\begin{theorem}[Igualdad proyectiva de las dos publicaciones correlacionadas]
\label{thm:alpha-dos-publicaciones-completas}
Para todo \(N\), la vía entera y la vía analítica publican el mismo prefijo:
\begin{equation}
 \boxed{\alpha_{A,N}=\operatorname{pref}_N(\alpha_A)
 =\operatorname{pref}_N(\alpha_B)=\alpha_{B,N}.}
 \label{eq:alpha-dos-vias-cada-nivel}
\end{equation}
En el límite,
\begin{equation}
 \boxed{\alpha_A
 =\operatorname*{arg\,zero}_{x\in I_\alpha}
       E_{\mathsf s,\infty}(x)
 =\alpha_B=\alpha_{\mathrm{HMT}}.}
 \label{eq:alpha-dos-vias-limite-coinductivo}
\end{equation}
\end{theorem}

\begin{proof}
La normalización entera publica exactamente \(a(\mathsf s)\). El
teorema~\ref{thm:alpha-raiz-completa-explicita} demuestra que la carta
analítica representa ese mismo punto como raíz simple única. Los
cilindros semiabiertos \(C_N^A\) fijan un prefijo canónico y
\(I_{B,N}\subseteq C_N^A\) contiene \(\alpha_B=\alpha_A\). Por tanto, sus
prefijos coinciden para cada \(N\); el anidamiento y los diámetros tendentes a
cero dan la igualdad de los límites. El reconocimiento
metrológico ocurre después.
\end{proof}

La función analítica constituye una publicación correlacionada del mismo
estado, no una segunda derivación causalmente independiente de la clausura
entera. Esta precisión no reduce la igualdad demostrada: identifica su
mecanismo exacto, impide atribuir al jet una cola que no determina y evita la
promoción de un valor convencional a generador.



\HMTDivision{\(h\quad\hbar\)}{Constantes de Planck y orientación}{La estructura determinantal fija la escala de acción; los coeficientes de retorno determinan las coordenadas angulares.}
\HMTNarrativeHeading{De la estructura determinantal a la orientación}

Las constantes de Planck se presentan después de los objetos que
determinan sus secciones de acción. El registro dodecafásico contiene el
bloque local de Hadamard; su cociente integral tiene un orden calculable.
El lector determinantal compone ese dato con las coordenadas ya generadas
y fija una valoración decimal. La escala y el coeficiente de una sección
cumplen funciones diferentes, que la demostración mantiene explícitas.

El retorno distingue las secciones anterior y posterior a la vuelta.
Sus coeficientes conservan la procedencia del carácter y de la
continuación regional. La expresión de la acción en unidades físicas
se introduce mediante la carta dimensional correspondiente. La
construcción transmite así una sección, su coeficiente, su escala y la
regla de lectura. El uso posterior de la constante conserva esas
relaciones.

Ángulo y contraángulo se desarrollan a continuación como coordenadas de
orientación. La primera lectura desplaza los bloques de alfa en la base
de completación; la segunda compone el coeficiente adimensional de
retorno con alfa. La carta circular expresa el resultado en grados o
radianes y conserva el número de vueltas cuando se utiliza el
levantamiento. Esta secuencia explica por qué las coordenadas angulares
son intermediarias importantes entre la acción y los lectores
geométricos posteriores.

Los canales conjugados reciben esa orientación y evalúan, junto con
la incidencia marcada, el funcional de Barbero--Immirzi.
La realización electrónica utiliza su trayectoria y sus lectores de
retorno, y recibe el cociente de las secciones de acción. La narración conserva
el orden de sus antecedentes y distingue el coeficiente adimensional
de la magnitud de acción que contribuye a representar. Una misma
constante puede actuar en varias etapas sin perder el mapa por el que
fue obtenida.

El recorrido prepara una composición física concreta. Acción y reloj
determinan una energía; la transducción térmica relaciona esa energía
con la información; el lector de área incorpora la incidencia marcada.
La realización gravitatoria completa después el cociente entre energía
y área. El artículo conserva los cálculos de esas composiciones para
mostrar cómo los resultados aritméticos y geométricos participan en
una lectura física común.
\HMTMathematicalPaper
% RESULTADO_RECUPERADO. II REV10, 05_accion.tex:1--195,
% revision_planck.tex:13--93 y definicion_planck_rev06.tex:4--34.
% Selección: antecedentes de acción de la elipse; no incorpora aplicaciones
% electrónicas ni el contraste metrológico, que no determinan el radio.
\section{Estructura determinantal y secciones orientadas de acción}
\label{iv:action:section}

Las publicaciones de \(\pi,\varphi,e,\alpha\) proceden del estado
APP--TRIT--TPK de las secciones
\ref{sec:nucleo}--\ref{sec:alpha}. El registro dodecafásico de
\ref{sec:k-registro} conserva su orientación, sus residuos y
cocientes y la memoria del retorno
\eqref{eq:holonomia-memoria}. Sobre ese registro se definen dos operaciones
sucesivas: una norma local determina la valoración de la sección de
acción, y un carácter con continuación regional determina sus
coordenadas anterior y posterior al retorno. Ambas operaciones
transmiten a la construcción angular sus reglas y normalizaciones.

\subsection{Cociente local y norma determinantal}
\label{iv:action:determinantal}

El bloque local de paso tres del registro dodecafásico, identificado
en \ref{subsec:k-hadamard}, es
\begin{equation}
 H_4=\begin{pmatrix}
 1&1&1&1\\
 1&1&-1&-1\\
 1&-1&1&-1\\
 1&-1&-1&1
 \end{pmatrix},
 \qquad G_H=\mathbb Z^4/H_4\mathbb Z^4.
 \label{iv:action:h4}
\end{equation}
El cociente corresponde a este bloque local. Los tres bloques del
registro completo conservan sus posiciones; la operación de norma
que sigue actúa sobre las hojas de \(G_H\).

\begin{lemma}[Estructura del cociente local]
\label{iv:action:smith}
La forma normal de Smith de \(H_4\) es
\(\operatorname{diag}(1,2,2,4)\). Por consiguiente,
\[
 G_H\simeq\mathbb Z/2\mathbb Z\oplus\mathbb Z/2\mathbb Z
       \oplus\mathbb Z/4\mathbb Z,\qquad |G_H|=16.
\]
\end{lemma}
\begin{proof}
El máximo común divisor de las entradas es uno. Los menores de
orden dos son pares y uno tiene valor \(-2\), por lo que su máximo
común divisor es dos. Las identidades
\(H_4H_4^{\mathsf T}=4I\) y \(\det H_4=-16\) muestran que los
cofactores valen \(\pm4\); el tercer divisor determinantal es
cuatro. El cuarto es \(16\). Los cocientes sucesivos de
\(1,2,4,16\) son \(1,2,2,4\); su producto determina el orden
del cociente.
\end{proof}

\begin{definition}[Lector determinantal de acción]
\label{iv:action:norma}
La sección escalar \(\alpha\), constante sobre las hojas de \(G_H\),
y una aplicación de la autoescala \(\varphi\) determinan
\begin{equation}
 \Sigma_H=\varphi\,\mathrm N_{G_H}(\alpha),\qquad
 \mathrm N_{G_H}(\alpha)=\prod_{g\in G_H}\alpha=\alpha^{16}.
 \label{iv:action:sigma}
\end{equation}
Su valoración decimal es
\[
 \operatorname{ord}_{10}(\Sigma_H)
   =\lfloor\log_{10}\Sigma_H\rfloor.
\]
\end{definition}
La potencia dieciséis procede del número de hojas una vez fijada
la norma multiplicativa. La norma y la aplicación única de
\(\varphi\) forman parte del lector declarado: el orden del grupo,
considerado aisladamente, admite otros funcionales.

\begin{theorem}[Valoración decimal de la sección de acción]
\label{iv:action:decada}
El lector anterior, aplicado a las coordenadas HMT ya generadas,
satisface
\begin{equation}
 10^{-34}<\Sigma_H<10^{-33},
 \qquad \operatorname{ord}_{10}(\Sigma_H)=-34.
 \label{iv:action:orden-decimal}
\end{equation}
\end{theorem}
\begin{proof}
Se utilizan las inclusiones racionales de las publicaciones
anteriores,
\begin{equation}
 \frac{1618}{1000}<\varphi<\frac{1619}{1000},
 \qquad
 \frac{7297}{10^6}<\alpha<\frac{7298}{10^6}.
 \label{iv:action:intervalos}
\end{equation}
La función \((x,a)\mapsto xa^{16}\) es creciente en ambas
variables positivas. Las desigualdades enteras
\[
 1618\cdot7297^{16}>10^{65},
 \qquad 1619\cdot7298^{16}<10^{66}
\]
dan
\[
 10^{-34}
 <\frac{1618}{1000}\left(\frac{7297}{10^6}\right)^{16}
 <\Sigma_H
 <\frac{1619}{1000}\left(\frac{7298}{10^6}\right)^{16}
 <10^{-33}.
\]
La definición de la valoración concluye la prueba. Las inclusiones
son controles posteriores de las coordenadas ya producidas; la
operación determinantal precede a cualquier carta de unidades.
\end{proof}

\subsection{Carácter de acción y continuación regional}
\label{iv:action:caracter}

La carta central conserva la matriz
\[
 B_c=\begin{pmatrix}7&2\\2&7\end{pmatrix},
 \qquad\lambda_+=9,\quad\lambda_-=5.
\]
El calendario tiene longitud \(12\), la órbita del paso tres tiene
longitud cuatro y el cociente censal utilizado es
\(104976/1944=54\), con los censos definidos en la
sección~\ref{sec:extension}. Estos invariantes determinan los coeficientes
asignados a los grados \(2,3,4,5\) del carácter:
\begin{equation}
 \frac9{16}=\frac{\lambda_+}{4^2},\qquad
 \frac59=\frac{\lambda_-}{\lambda_+},\qquad
 \frac7{48}=\frac{\operatorname{tr}(B_c)/2}{4\cdot12},
 \qquad
 \frac1{54}=\frac{1944}{104976}.
 \label{iv:action:coeficientes}
\end{equation}
La asignación a esos grados y la alternancia de signos pertenecen
a la definición del carácter analítico. El contenido operativo
reúne el soporte, la orientación y esa regla de asignación.

\begin{definition}[Carácter local y lector regional de retorno]
\label{iv:action:lectores}
Sobre las publicaciones \(\alpha,\varphi,\pi\), el carácter de
quinto orden, la transferencia regional y su continuación son
\begin{align}
 H_5={}&\frac12\sqrt{\frac{1000\alpha}{\varphi}}-\alpha
       +\frac9{16}\alpha^2-\frac59\alpha^3
       +\frac7{48}\alpha^4-\frac1{54}\alpha^5,
       \label{iv:action:h5}\\
 D_A={}&\exp\!\left(-\frac{100\pi\alpha}{9}\right),
       \label{iv:action:da}\\
 \mathcal C_\pi(\alpha)={}&169\alpha^6+
       \frac{D_A\alpha^7}{1-D_A\alpha},
       \label{iv:action:cpi}\\
 \eta_{\mathrm{ret}}={}&H_5-\frac{90}{\pi}\mathcal C_\pi(\alpha).
       \label{iv:action:eta-ret}
\end{align}
El entero \(169\) es el selector regional de las publicaciones
construidas en \ref{sec:generacion}: se conserva la emisión regional
\(169\,|\,272\,|\,272\), distinguida de
\(418\,|\,674\,|\,521\) en la construcción excepcional.
La primera prolongación estricta de su continuación
tiene peso \(D_A\) y cada concatenación multiplica el peso por
\(D_A\). La normalización inicial es unitaria.
\end{definition}

La continuación se formula antes de evaluar en \(z=\alpha\).
Con \(D_A\) ya producido por su lector, escribimos
\(\mathscr C(z)=169z^6+\mathscr R(z)\), donde
\(\mathscr R(z)\in z^7\mathbb R[[z]]\).
La concatenación aumenta el grado en una unidad; el peso inicial
y la regla de transferencia dan
\begin{equation}
 \mathscr R(z)=D_Az^7+D_Az\mathscr R(z).
 \label{iv:action:renovacion}
\end{equation}

\begin{proposition}[Unicidad de la continuación normalizada]
\label{iv:action:unicidad-renovacion}
La ecuación \eqref{iv:action:renovacion} determina una única serie:
\begin{equation}
 \mathscr C(z)=169z^6+\frac{D_Az^7}{1-D_Az}
             =169z^6+\sum_{n=7}^{\infty}D_A^{n-6}z^n.
 \label{iv:action:serie-renovacion}
\end{equation}
Su realización es analítica para \(|D_Az|<1\).
\end{proposition}
\begin{proof}
El término independiente de \(1-D_Az\) es uno, de modo que ese
elemento es invertible en el anillo de series formales.
Resolver \eqref{iv:action:renovacion} da la expresión racional.
Equivalentemente, el coeficiente de grado siete es \(D_A\) y cada
coeficiente posterior es \(D_A\) veces el anterior. La serie
geométrica converge absolutamente en el disco indicado.
\end{proof}

La normalización del primer peso interviene en la unicidad.
El impulso de grado seis y la razón de concatenación, sin ese
peso inicial, también son compatibles con
\[
 169z^6+\lambda\frac{D_Az^7}{1-D_Az}.
\]
La regla unitaria fija \(\lambda=1\) antes de la evaluación.
Para \(0<\alpha<1\), se tiene \(0<D_A<1\) y \(D_A\alpha<1\);
por tanto,
\begin{equation}
 \frac{D_A\alpha^7}{1-D_A\alpha}
   =\sum_{n=0}^{\infty}D_A^{n+1}\alpha^{n+7}.
 \label{iv:action:cola}
\end{equation}
Después del término de índice \(N\), el resto exacto es
\(D_A^{N+2}\alpha^{N+8}/(1-D_A\alpha)\). La expresión racional
conserva así la continuación completa a cualquier precisión.

\subsection{Secciones positivas y acción por vuelta}
\label{iv:action:secciones}

Sea \(\mathcal L_S\) la recta orientada de acción y
\(\mathcal U_S\) una base positiva. Sus secciones anterior y
posterior al retorno son
\begin{equation}
 \hbar_{\mathrm{pre}}=H_5\,10^{-34}\mathcal U_S,\qquad
 \hbar_{\mathrm{ret}}=\eta_{\mathrm{ret}}\,10^{-34}\mathcal U_S.
 \label{iv:action:hbar}
\end{equation}
El retorno conserva la compensación regional
\(90\mathcal C_\pi(\alpha)/\pi\) y la razón
\begin{equation}
 \delta_{\mathrm{ret}}=H_5-\eta_{\mathrm{ret}},\qquad
 R_{\mathrm{act}}=\frac{\hbar_{\mathrm{pre}}}
                        {\hbar_{\mathrm{ret}}}
                =\frac{H_5}{\eta_{\mathrm{ret}}}.
 \label{iv:action:razon}
\end{equation}

\begin{proposition}[Positividad y covariancia de la razón de acción]
\label{iv:action:positividad}
En los intervalos \eqref{iv:action:intervalos},
\(0<\eta_{\mathrm{ret}}<H_5\), y \(R_{\mathrm{act}}>1\).
La razón permanece fija al cambiar la base positiva común
\(\mathcal U_S\).
\end{proposition}
\begin{proof}
Los términos de \(\mathcal C_\pi\) son positivos. Para acotar
\(H_5\), bastan \(\alpha>0.007\), \(\alpha<0.008\) y
\(\varphi<1.619\). El término de raíz de \eqref{iv:action:h5}
es mayor que uno y la suma de sus términos negativos es menor
que \(0.008001\); luego \(H_5>0.99\).
Con \(D_A<1\) y \(\pi>3\),
\[
 0<\frac{90}{\pi}\mathcal C_\pi
 <30\left(169(0.008)^6+
             \frac{(0.008)^7}{1-0.008}\right)<10^{-6}.
\]
Así, \(\eta_{\mathrm{ret}}>0.989999>0\) y
\(\eta_{\mathrm{ret}}<H_5\). El cambio de base divide ambas
coordenadas de acción por el mismo factor positivo; éste se
cancela en su cociente.
\end{proof}

La determinación estructural de la acción angular es el par
\((\hbar_{\mathrm{pre}},\hbar_{\mathrm{ret}})\) con su
identificación por retorno. Conserva la valoración de la norma
local, el carácter de acción y la continuación regional.
En cada sección \(s\in\{\mathrm{pre},\mathrm{ret}\}\), la acción
de una vuelta completa es
\begin{equation}
 h_s=2\pi\hbar_s.
 \label{iv:action:vuelta}
\end{equation}
La carta posterior \(\mathcal U_S\mapsto\mathrm{J\,s}\) expresa
esas secciones en unidades de acción. El radio normalizado que
se construirá conserva su independencia respecto de esta elección
común de unidades.

% RESULTADO_RECUPERADO. II REV10,
% 05b_coordenadas_angulares.tex:1--125.
% Se conserva la publicación angular, sus cartas y pruebas; el funcional
% ulterior de Barbero no es una dependencia de la forma elíptica.
\section{Publicación angular y orientación de los canales}
\label{iv:angular:section}

La sección \(\hbar_{\mathrm{ret}}\) pertenece a la recta de
acción \(\mathcal L_S\); \(\eta_{\mathrm{ret}}\) es su
coeficiente adimensional en la base declarada en
\eqref{iv:action:hbar}. La publicación angular actúa después
sobre \((\alpha,\eta_{\mathrm{ret}})\). Escribiendo
\([t]_{360}\) para la clase de \(t\) en
\(\mathbb R/(360\mathbb Z)\), se define
\begin{equation}
 \mathcal P_{\mathrm{ang}}(\alpha,\eta_{\mathrm{ret}})
 =\bigl([1000\alpha]_{360},
        [2(\eta_{\mathrm{ret}}+\alpha)]_{360}\bigr).
 \label{iv:angular:publicacion}
\end{equation}
La completación nonádica determina el factor \(1000\); la
clausura bidireccional determina el factor dos. La carta circular
de \(360\) unidades conserva las particiones
\(12\cdot30=4\cdot90=3\cdot120\). En la rama positiva, los
representantes levantados son
\begin{equation}
 A_{\deg}=1000\alpha,\qquad
 C^*_{\deg}=2(\eta_{\mathrm{ret}}+\alpha).
 \label{iv:angular:representantes}
\end{equation}
El subíndice distingue estas coordenadas angulares de los bloques
del registro dodecafásico.

\begin{proposition}[Desplazamiento de la publicación en base mil]
\label{iv:angular:desplazamiento}
Sea \(\alpha=\sum_{n\geq1}b_n1000^{-n}\), con
\(b_n\in\{0,\ldots,999\}\), la expansión compatible ya
generada. Su coordenada levantada satisface
\[
 A_{\deg}=b_1+\sum_{n\geq1}b_{n+1}1000^{-n}.
\]
La publicación conserva los bloques a partir del segundo y
sitúa el bloque inicial en la parte entera.
\end{proposition}
\begin{proof}
La serie converge absolutamente. Multiplicar cada término por
\(1000\) y separar el término inicial proporciona la identidad.
La operación actúa sobre la expansión construida, con sus
bloques y memoria; la publicación circular posterior toma
su clase módulo \(360\).
\end{proof}

\subsection{Cartas circulares y covariancia de la acción}
\label{iv:angular:cartas}

\begin{definition}[Cartas de la coordenada angular]
\label{iv:angular:unidades}
La carta en grados expresa una vuelta completa mediante \(360\)
unidades; la carta en radianes, mediante \(2\pi\). El cambio
de coordenadas y su inversa son
\begin{equation}
 \theta_A=\frac{\pi}{180}A_{\deg},\qquad
 \theta_C=\frac{\pi}{180}C^*_{\deg},\qquad
 A_{\deg}=\frac{180}{\pi}\theta_A,\qquad
 C^*_{\deg}=\frac{180}{\pi}\theta_C.
 \label{iv:angular:cambio-carta}
\end{equation}
Un recorrido con número de vueltas conservado utiliza las
coordenadas levantadas. El cociente circular publica su
clase módulo la vuelta correspondiente.
\end{definition}

El registro retiene orientación, vueltas y memoria del transporte.
La fase final y el levantamiento son publicaciones diferentes:
el segundo distingue recorridos que comparten la primera.
Cada expresión de acción declara cuál de esas coordenadas utiliza.

\begin{proposition}[Covariancia del lector de acción]
\label{iv:angular:accion}
Para una sección \(s\in\{\mathrm{pre},\mathrm{ret}\}\) y una
coordenada angular levantada, se cumple
\begin{equation}
 \mathcal S_s(\theta)
 =\hbar_s\theta_{\mathrm{rad}}
 =h_s\frac{\theta_{\deg}}{360}.
 \label{iv:angular:accion-cartas}
\end{equation}
Una vuelta positiva publica \(h_s\).
\end{proposition}
\begin{proof}
Sustituir
\(\theta_{\mathrm{rad}}=\pi\theta_{\deg}/180\) y
\(h_s=2\pi\hbar_s\) da
\[
 \hbar_s\frac{\pi\theta_{\deg}}{180}
 =\frac{h_s\theta_{\deg}}{360}.
\]
En una vuelta \(\theta_{\deg}=360\), con lo que resulta \(h_s\).
\end{proof}

La periodicidad de una representación mantiene su condición
propia: si \(U(2\pi)=-I\), la restitución operatoria ocurre
en \(4\pi\). La normalización de la acción por vuelta de
\eqref{iv:angular:accion-cartas} y el retorno del estado
representado conservan así sus dominios respectivos.

\subsection{Canales positivos y reversión de la coordenada impar}
\label{iv:angular:canales}

Los canales se publican, en cualquiera de las dos cartas, mediante
\begin{equation}
 q_\pm=\exp[-(\theta_A\pm\theta_C)]
 =\exp\!\left[-\frac{\pi}{180}
                   (A_{\deg}\pm C^*_{\deg})\right].
 \label{iv:angular:q}
\end{equation}

\begin{proposition}[Recuperación de las coordenadas angulares]
\label{iv:angular:inversion}
En la carta real levantada, los canales positivos determinan
unívocamente
\begin{equation}
 A_{\deg}=-\frac{90}{\pi}\log(q_+q_-),\qquad
 C^*_{\deg}=\frac{90}{\pi}\log\frac{q_-}{q_+}.
 \label{iv:angular:recuperacion}
\end{equation}
El intercambio \((q_+,q_-)\mapsto(q_-,q_+)\) conserva
\(A_{\deg}\) y revierte \(C^*_{\deg}\).
\end{proposition}
\begin{proof}
La suma y la diferencia de los logaritmos de
\eqref{iv:angular:q} son \(-\pi A_{\deg}/90\) y
\(-\pi C^*_{\deg}/90\), respectivamente. El logaritmo real
es unívoco en los canales positivos. Su producto es simétrico
y el cociente se invierte al intercambiarlos.
\end{proof}

La involución antipodal \(r\mapsto r+6\pmod {12}\) del
registro forma seis parejas. La distribución uniforme de la
coordenada impar sobre esas parejas asigna
\(C^*_{\deg}/6\) a cada una; la suma recupera
\(C^*_{\deg}\). La conversión de unidades se aplica una
sola vez a cada argumento. En particular,
\(C^*_{\deg}/A_{\deg}=\theta_C/\theta_A\) cuando
\(A_{\deg}\ne0\), mientras el estado anterior conserva
los ángulos levantados y su memoria.

\HMTDivision{\(e^-\)}{Estructura algebraica del electrón}{Persistencia central, orientación y secciones de acción en la realización electrónica.}
\HMTNarrativeHeading{Persistencia y trayectoria electrónica}

La estructura electrónica empieza en la posición central persistente
de APP. La inversión celular distingue esa posición de las parejas
conjugadas que la rodean. Su evaluación conserva las hojas, los
residuos y los cocientes; la trayectoria prolonga esos datos mediante
las operaciones del TPK. La exposición de la masa recibe después las
secciones de acción y los lectores que utiliza.

La vacancia central permite ver una diferencia decisiva entre lectura
residual e historia aritmética. Dos posiciones pueden compartir la
suma y el residuo multiplicativo, mientras sus productos enteros
conservan cocientes distintos. El cambio de cociente registra la
transferencia; la orientación distingue sus posiciones conjugadas.
El ejemplo hace visible qué información necesita una reconstrucción
que restituya la operación realizada.

La cadena electrónica añade una cronología explícita. La recurrencia
celular, las inversiones y las ventanas de lectura conservan los
registros de la sección persistente. El calendario registra los pasos
en que permanece la capacidad; el defecto central registra la diferencia
de cociente entre evaluaciones con igual residuo. La trayectoria
electrónica conserva este segundo dato junto con su orientación y su
cronología. La vacancia de capacidad y el defecto de cociente cuentan
operaciones distintas, con sus respectivos registros.

Esta continuidad vuelve pertinente la distinción entre retorno y
recuperación. Una fase puede repetirse mientras el registro de la
sección incorpora el trayecto. La evaluación posterior utiliza esa
información transportada y conserva identificada la sección central.
La masa, la orientación y las expresiones de acción se presentan con
sus antecedentes, de modo que cada paso pueda leerse como una
composición determinada.

El electrón ofrece así un caso desarrollado de la tesis del artículo.
La persistencia corresponde a una estructura; la trayectoria describe
su evolución; los registros distinguen los cambios; el lector
dimensional obtiene una magnitud. Su interés para la conservación
reside en esa sucesión completa. La unidad que se mantiene durante
el transporte conserva relaciones cuyo contenido se hace accesible
en lecturas diferentes. El estudio operatorio posterior abstraerá
las condiciones de recuperación y mostrará sus balances exactos.
\HMTMathematicalPaper
\begin{HMTCompactSection}
\section{Estructura algebraica del electrón y escala de acción}
\label{sec:electron}

La construcción electrónica recibe el estado aritmético con memoria de
trayectoria producido por APP--TRIT--TPK y sus coordenadas numéricas. En este punto del desarrollo,
\(\pi,\varphi,e\) y \(\alpha\) designan las coordenadas ya obtenidas de
la misma genealogía, con el orden constructivo expuesto anteriormente. No
son parámetros tomados de una tabla física. La sección electrónica es una
realización posterior de esa estructura discreta conjunta del continuo:
conserva la celda de origen, las dos hojas aritméticas, la orientación, el
acarreo y la memoria de retorno antes de producir un escalar.

Conviene separar tres objetos desde el comienzo. El primero es la sección
central y su historia de transporte; el segundo, la coordenada
adimensional que se obtiene de ella; el tercero, su realización en una
recta de energía o de masa. La igualdad entre dos coordenadas numéricas no
identifica estos tres niveles. En particular, el número de Euler \(e\),
la coordenada electrónica \(\varepsilon_e\) y una energía de reposo
\(E_e\) llevarán notaciones distintas.

El objetivo de esta sección es establecer la composición
\begin{equation}
 \varepsilon_e^{\beta}
 =\frac{\sqrt3}{4}
  \left[\exp\!\left(\frac{\varphi}{\pi^2}\right)-22\alpha^3\right]
  (1+15\Delta_4)
  R_{\mathrm{act}}^{\,80-54\alpha+6\Delta_4},
 \label{eq:el-formula-completa}
\end{equation}
explicando por separado el origen del prefactor, los dos enteros de la
etapa basal, la normalización de \(\Delta_4\), la razón de acción y el
triple de retorno. La carta dimensional se abre solamente después de
esta composición. El resultado no consiste en hacer que un número
adimensional adquiera unidades por su mera semejanza con una magnitud
tabulada.

La composición reúne factores con funciones distintas. El cuadro
siguiente identifica su significado y remite a la construcción
correspondiente; las demostraciones se desarrollan en el orden de sus
dependencias.

\begin{table}[htbp]
\centering
\small
\setlength{\tabcolsep}{4pt}
\renewcommand{\arraystretch}{1.18}
\begin{tabularx}{\linewidth}{@{}p{0.20\linewidth}Xp{0.24\linewidth}@{}}
\toprule
Factor & Significado en la composición & Construcción y demostración\\
\midrule
\(\sqrt3/4\)
 & Norma del conmutador de los proyectores en el plano principal.
 & Ecuación~\eqref{eq:el-prefactor}.\\
\(\exp(\varphi/\pi^2)\)
 & Carácter exponencial de la orientación electrónica del cociente regional.
 & Ecuación~\eqref{eq:el-bisagra}; apartados~\ref{sec:medida-retorno-hojas}
   y~\ref{sec:operador-retorno-marcado}.\\
\(-22\alpha^3\)
 & Contribución del complemento regional en el retorno central.
 & Ecuaciones~\eqref{eq:el-selector-enteros} y~\eqref{eq:el-censo}.\\
\(1+15\Delta_4\)
 & Normalización torsional compuesta con la multiplicidad regional
   y trítica.
 & Ecuaciones~\eqref{eq:el-selector-enteros}, \eqref{eq:el-delta}
   y~\eqref{eq:el-basal}.\\
\(R_{\mathrm{act}}^{\kappa_e}\)
 & Cociente de las secciones de acción, elevado al carácter determinado
   por los dos lectores del retorno electrónico.
 & Ecuaciones~\eqref{eq:el-ratio}, \eqref{eq:el-triple},
   \eqref{eq:el-kappa} y~\eqref{eq:el-beta}.\\
\bottomrule
\end{tabularx}
\caption{Factores de la composición electrónica y procedencia de sus reglas.}
\label{tab:el-factores}
\end{table}

\paragraph{Identidad de la sección y composición de sus observables.}
\label{el-ampl-articulacion}
La identidad estructural del electrón se introduce mediante la sección
central y sus operaciones de transporte. La inversión de la carta APP
selecciona la celda $(5,5)$ por su propiedad de punto fijo. Las dos hojas
aritméticas permiten estudiar las deformaciones que conservan la suma y
el residuo del producto. El resultado distingue las dos orientaciones
$(8,2)$ y $(2,8)$ y conserva la diferencia de una unidad de cociente.
Éste es el primer nivel de información: una celda, una involución, una
fibra conservada y un defecto entero mínimo. Ninguna evaluación de masa
es necesaria para enunciarlo o demostrarlo.

El paso a las proyecciones aritméticas introduce un segundo nivel. Las
fibras de suma y producto determinan esperanzas condicionales sobre el
mismo espacio finito. Su conmutador mide la dependencia respecto del
orden de esas dos proyecciones. La norma $\sqrt3/4$ resulta del espectro
de su emparejamiento; el bloque que alcanza esa norma conserva además
las relaciones cuadráticas demostradas a continuación. Por tanto, el
prefactor escalar y el álgebra local tienen una procedencia común, pero
contienen información diferente. La norma retiene la magnitud del
conmutador, mientras las matrices conservan su orientación y su ley de
composición.

La transferencia areal--volumétrica incorpora las coordenadas regionales
en un tercer nivel. El calendario de modos determina la matriz de
incidencia y su acción cuadrática. La marca de clausura produce la razón
$\pi/\varphi^2$; el intercambio de los dos argumentos proporciona la
orientación $\varphi/\pi^2$ utilizada en el carácter exponencial
electrónico. Las dos expresiones quedan relacionadas por una involución
explícita. Esa relación explica por qué una coordenada regional puede
aparecer en la ecuación electrónica con una potencia o en una posición
diferente de la que tenía en la primera lectura: la operación aplicada
es parte del dato matemático que se conserva.

El término cúbico de alfa pertenece al complemento regional de un retorno
marcado. La cuenta de sus posiciones produce $27-5=22$, mientras las
tres hojas de cada una de las cinco regiones producen $3\cdot5=15$.
Las demostraciones conservan el dominio de esos recuentos y la
orientación que determina sus signos. La fórmula basal compone,
por consiguiente, una norma de incompatibilidad, un carácter de
transferencia, un déficit de posiciones y una memoria torsional.
La escritura en una sola línea expresa la composición de esas
operaciones; su explicación exige reconstruirlas en ese orden.

La transición de la coordenada basal a la coordenada completa conserva
la misma sección central. El factor de retorno utiliza los registros
de la historia, y los dos lectores declarados recuperan el triple
$(80,54,6)$ sobre esa trayectoria. La multiplicación por
$R_{\mathrm{act}}^{80-54\alpha+6\Delta_4}$ modifica la evaluación,
pero mantiene identificados su origen y sus antecedentes. El valor
basal constituye una etapa intermedia de la composición; el valor
completo incluye también esa contribución del retorno. La comparación
física debe utilizar la etapa especificada por la fórmula que se evalúa.

La escala de acción interviene después con otra función. Una razón de
dos secciones de acción es adimensional y cancela su escala común.
La realización absoluta utiliza, en cambio, una base de acción, un
lector temporal y el mapa hacia energía o masa. Esta separación permite
conservar simultáneamente la construcción de la escala y la identidad
de la sección electrónica. El cuanto cronológico proporciona una
referencia de escala; la sección central determina el factor estructural
que se expresa respecto de ella. Las ecuaciones dimensionales de este
capítulo mantienen explícita esa composición.

La pertinencia del electrón en el argumento general procede de esta
convergencia de operaciones. La misma aritmética que produjo las
regiones determina ahora una celda central, sus deformaciones y sus
proyecciones. La constante de estructura fina, construida antes, actúa
en esa composición como coordenada de clausura. La información de
retorno determina la evaluación completa, y la realización dimensional
permite el contraste posterior. Explicar el electrón desde la estructura
consiste en conservar esta sucesión de objetos, mapas y valores en cada
paso de la demostración.


\subsection{Centro, orientación y vacancia sobre una fibra conservada}

El soporte celular considerado es \(\{1,\ldots,9\}^2\), con ambas
evaluaciones APP, \(r+c\) y \(rc\). La inversión central es
\begin{equation}
 \rho(r,c)=(10-r,10-c).
 \label{eq:el-inversion-central}
\end{equation}
El centro no se elige por una masa: se caracteriza por una propiedad del
soporte anterior a su lectura física.

\begin{proposition}[Unicidad del centro]
La inversión \(\rho\) tiene un único punto fijo, \((5,5)\).
\end{proposition}
\begin{proof}
Las ecuaciones \(r=10-r\) y \(c=10-c\) equivalen a
\(2r=2c=10\). Su única solución en la carta nonádica es \((5,5)\).
\end{proof}

La propiedad anterior identifica una celda, no una historia profunda
única. Una ruta añade la elección inicial de orientación y sus
prolongaciones compatibles. En la carta utilizada aquí la sección
orientada central comienza en \((5,5)\) con las dos orientaciones
positivas; se denotará por \(\Gamma_e\). El retorno de su proyección
celular no borra las orientaciones transportadas ni el registro de
acarreo de la sección con su memoria de trayectoria.

La vacancia se considera ahora, después del cierre de \(\alpha\), como
una operación local de esta sección. Su aritmética puede aislarse en un
lema finito; ese aislamiento no la convierte en un sustituto del
generador anterior de constantes.

\begin{definition}[Fibra central y defecto de cociente]
La fibra aditiva central es
\[
 \mathcal F_{10}
 =\{(5+t,5-t):t\in\mathbb Z,\ |t|\leq4\}.
\]
En la carta \(n=\rho_9(n)+9q_9(n)\), con
\(\rho_9(n)\in\{1,\ldots,9\}\), se llama defecto de cociente respecto
al centro a
\[
 v(t)=q_9(25)-q_9(25-t^2),
\]
cuando ambos productos poseen el mismo residuo nonádico.
\end{definition}

\begin{lemma}[Vacancia central mínima]
Los únicos puntos de \(\mathcal F_{10}\) que conservan el residuo
multiplicativo del centro son
\[
 (5,5),\qquad (8,2),\qquad (2,8).
\]
En los dos puntos no centrales el defecto de cociente es exactamente uno.
\end{lemma}
\begin{proof}
La suma es constantemente diez, mientras
\[
 (5+t)(5-t)=25-t^2.
\]
La conservación del residuo equivale a \(t^2\equiv0\pmod9\), es decir,
a \(3\mid t\). En el intervalo entero \([-4,4]\) las soluciones son
\(t=-3,0,3\). Para las dos no nulas,
\[
 25=7+9\cdot2,\qquad16=7+9\cdot1.
\]
Se conserva el residuo siete y se pierde una unidad de cociente. No hay
otro desplazamiento admisible no central en ese dominio, de modo que el
defecto positivo es mínimo. El intercambio \(t\mapsto-t\) permuta las
dos salidas y conserva el defecto.
\end{proof}

La vacancia no es aquí un cero numérico: es una diferencia entre dos
estados que tienen la misma lectura residual y distinto cociente. Si se
conservara únicamente el residuo, esa diferencia desaparecería. La
orientación distingue además los dos puntos de salida. Ni el defecto de
cociente ni esa orientación se identifican automáticamente con una carga
eléctrica o con una masa; su realización física requiere los lectores
correspondientes.

\subsubsection{Coordenadas centrales, orientación y registros regionales}
\label{el-centro:registros}

La estructura central relaciona posición, evaluación y transporte. En la
celda \((5,5)\), las dos hojas de APP dan
\begin{equation}
 \mathcal A(5,5)=\bigl((1,1),(7,2)\bigr),\qquad
 5+5=1+9\cdot1,\qquad5\cdot5=7+9\cdot2.
 \label{el-centro:app}
\end{equation}
Los dos unos de la hoja aditiva son, en su orden, residuo y cociente.
Su coincidencia numérica conserva dos funciones aritméticas distintas.
Sobre la fibra de suma \(10\), las dos vacancias orientadas verifican
\begin{equation}
 \mathcal A(8,2)=\mathcal A(2,8)=\bigl((1,1),(7,1)\bigr).
 \label{el-centro:vacancias}
\end{equation}
Se conservan la hoja aditiva completa y el residuo multiplicativo;
la coordenada de cociente del producto disminuye exactamente en
\(1\). El transporte retiene, además, cuál de las dos posiciones
conjugadas se ha alcanzado.

\begin{figure}[H]
\centering
\begin{tikzpicture}[x=.67cm,y=.67cm,font=\footnotesize,>=Latex]
 \draw[step=1,black!12,line width=.3pt] (1,1) grid (9,9);
 \draw[->,black!60] (.7,1)--(9.5,1) node[right] {$i$};
 \draw[->,black!60] (1,.7)--(1,9.5) node[above] {$j$};
 \foreach \k in {1,...,9}{
  \node[below=2pt] at (\k,1) {$\k$};
  \node[left=2pt] at (1,\k) {$\k$};
 }
 \draw[black!65] (1,9)--(9,1);
 \draw[<->,line width=.9pt] (2,8)--(8,2);
 \foreach \x/\y in {2/8,5/5,8/2}{
  \filldraw[fill=white,line width=.8pt] (\x,\y) circle (3pt);
 }
 \fill (5,5) circle (1.5pt);
 \node[above right=3pt,fill=white,inner sep=1pt] at (2,8) {$(2,8)$};
 \node[above right=3pt,fill=white,inner sep=1pt] at (5,5) {$(5,5)$};
 \node[above right=3pt,fill=white,inner sep=1pt] at (8,2) {$(8,2)$};
 \node[anchor=west,align=left] at (10.5,7.6)
   {$i+j=10$\\$R^+=1,\quad q^+=1$};
 \node[anchor=west,align=left] at (10.5,5.4)
   {centro: $ij=25$\\$R^\times=7,\quad q^\times=2$};
 \node[anchor=west,align=left] at (10.5,3.1)
   {vacancias: $ij=16$\\$R^\times=7,\quad q^\times=1$};
\end{tikzpicture}
\caption{Centro electrónico y vacancias conjugadas en la carta APP.
La diagonal es la fibra aditiva \(\mathcal F_{10}\); el centro
es el punto fijo de \((i,j)\mapsto(10-i,10-j)\). Los dos
desplazamientos de módulo \(3\) conservan la suma y el residuo
producto y reducen en una unidad el cociente multiplicativo. Las
flechas representan las deformaciones sobre la fibra.}
\label{el-centro:figura}
\end{figure}


La simetría central caracteriza también su persistencia bajo cambios de
carta. Si una transformación admisible \(f\) del atlas entrelaza la
inversión, \(f\rho=\rho f\), entonces
\[
 \rho f(5,5)=f\rho(5,5)=f(5,5).
\]
La unicidad del punto fijo obliga a \(f(5,5)=(5,5)\). Ésta es la
propiedad interna que hace estable la localización central: toda
transformación de ese tipo conserva la celda antes de evaluar su
carácter electrónico (sección~\ref{el-centro:registros}).

El bloque matricial central pertenece a la evaluación sobre las dos
posiciones adyacentes \(4,5\). Su descomposición, establecida en el
desarrollo local del núcleo, es
\[
 B_c=\begin{pmatrix}7&2\\2&7\end{pmatrix}=9P_++5P_-.
\]
Aquí \(9\) y \(5\) son autovalores de dos modos, mientras que
\((5,5)\) es una posición del atlas. Las concatenaciones de sus
entradas dan \(72+27=77+22=99\) y
\(77-22=55=54+1\). La primera igualdad conserva la suma total bajo
dos disposiciones; la segunda expresa el defecto entre las lecturas
diagonal y antidiagonal. Estas operaciones preceden a la composición
electrónica y conservan su procedencia matricial.

La orientación de partida es otro dato. En la representación del
retorno electrónico se fija \((h_0,v_0)=(1,1)\), es decir, avance
positivo en ambas coordenadas. La inversión simultánea cada \(27\)
pasos, leída en ventanas de \(9\), determina
\[
 \tau_e=(+1,+1,+1,-1,-1,-1,+1,+1,+1,-1,-1,-1).
\]
Los tramos \((1,1,1)\) pertenecen a esta sucesión de orientaciones:
cada tramo comprende tres ventanas del mismo sentido. La recurrencia
de la sección de lectores demuestra este registro en
\eqref{eq:el-tau}, y sus dos evaluaciones recuperan el triple
\((80,54,6)\) que interviene en
\eqref{eq:el-formula-completa}. La coordenada de posición, los
cocientes APP y la sucesión de signos quedan así conectados mediante
operaciones explícitas.

La firma \(555555\) pertenece a una lectura regional de otro
dominio. Para una emisión decimal \(U=(U_1,\ldots,U_6)\), escribimos
\(U_j=100d_{1j}+10d_{2j}+d_{3j}\), con tres cifras, incluidos los
ceros iniciales. El lector de firma multiplicativa es
\[
 \mathcal E_\times(U)
   =\bigl([d_{2j}-d_{1j}]_{10}\bigr)_{j=1}^{6}.
\]
En la región transversal
\(U_\pi^\perp=(501,614,498,169,272,272)\), las diferencias son
\((-5,-5,5,5,5,5)\). Por consiguiente,
\begin{equation}
 \mathcal E_\times(U_\pi^\perp)=(5,5,5,5,5,5),\qquad
 \Psi(U_\pi^\perp)=010211.
 \label{el-centro:firma-transversal}
\end{equation}
La primera aplicación conserva una firma multiplicativa del catálogo;
la segunda conserva la palabra ternaria común a las cinco regiones.
La figura \ref{pi-estrella:figura} muestra la posición transversal de
esta emisión y su multiplicidad \(432\).

\paragraph{Memoria necesaria para transportar la firma constante.}
La firma \(555555\) proporciona, además, un ejemplo verificable de
la información de trayectoria que exige TPK. Sea
\(\mathcal X_\times=D_9^2\times\{N,E,S,O\}\), con \(324\)
posiciones orientadas del canal producto. El avance \(P\) desplaza
una casilla en la dirección almacenada y la media vuelta \(H\)
invierte esa dirección. En esta representación, el lector
\(E_{\rm sig}\) suma por ventanas los exponentes discretos
\[
 \ell_9(1,2,4,8,7,5)=(0,1,2,3,4,5),\qquad
 \ell_9(3)=\ell_9(6)=\ell_9(9)=0.
\]
Durante cada una de las seis ventanas de nueve pasos, se evalúa
\(\ell_9(\rho_9(ij))\) antes de cada avance producto, seleccionado
por TRIT en los tiempos \(2,5,8\pmod9\). La suma de la ventana se
reduce módulo \(10\). Tras los pasos \(27\) y \(54\) se invierte
la dirección. Esta regla determina las seis componentes de
\(E_{\rm sig}\) desde la posición y orientación iniciales.

La enumeración de este dominio produce \(26\) firmas. La fibra de
\(555555\) contiene \(24\) posiciones orientadas; después de un
avance, sus imágenes se distribuyen en tres firmas, con ocho
representantes cada una:
\[
 555177,\qquad555717,\qquad555771.
\]
Tres testigos, en coordenadas de \(1\) a \(9\), permiten comprobar
la distinción directamente:
\[
\begin{array}{c|c|c}
x&E_{\rm sig}(x)&E_{\rm sig}(Px)\\\hline
(3,5,N)&555555&555177\\
(3,4,N)&555555&555717\\
(3,4,S)&555555&555771
\end{array}
\]
Por tanto, el valor común de la firma no determina su siguiente valor:
la clase de transporte dentro de esa fibra forma parte del dato
operativo. La partición estable mínima se calcula partiendo de la
equivalencia de firmas \(\sim_0\) y refinando mediante
\[
 x\sim_{n+1}y\quad\Longleftrightarrow\quad
 x\sim_n y,\quad Px\sim_n Py,\quad Hx\sim_n Hy.
\]
El procedimiento termina al ser finito el dominio. Toda congruencia
estable que refine \(\sim_0\) refina cada \(\sim_n\), por inducción;
su punto fijo es, por ello, la congruencia estable mínima requerida.
El censo da \(26\to28\to28\), con histograma
\(27\cdot8+108=324\). La única firma que se descompone es
\(555555\), en tres clases de ocho elementos. El certificado adjunto
reproduce la enumeración, las tres clases y la estabilidad por ambos
operadores (sección~\ref{el-centro:registros}).

La multiplicidad \(24\) pertenece aquí a la fibra del canal producto;
la multiplicidad \(432\) pertenece a la emisión regional completa
\(U_\pi^\perp\). La primera prueba por qué la firma necesita memoria
de transporte; la segunda sitúa esa firma en el catálogo conjunto de
las dos hojas. Ambas lecturas participan de la misma arquitectura,
con dominios y operaciones especificados.

Esta distinción de dominios permite reunir las apariciones de la unidad
y del centro con sus respectivos significados:
\begin{center}
\small
\begin{tabularx}{\textwidth}{@{}lX@{}}
\toprule
Registro & Objeto y función \\
\midrule
\((5,5)\) & Celda fija de la inversión central del atlas APP.\\
\((R^+,q^+)=(1,1)\) & Residuo y cociente de la evaluación aditiva central.\\
\((h_0,v_0)=(1,1)\) & Orientación inicial de los dos desplazamientos.\\
\((+1,+1,+1)\) & Tres ventanas consecutivas de \(\tau_e\) con orientación positiva.\\
\(c=111111\) & Dato de la quinta frontera nonádica, sección \ref{mono-recta:desarrollo}.\\
\(\mathcal E_\times=555555\) & Firma de la región transversal de \(\pi\), con seis posiciones.\\
\bottomrule
\end{tabularx}
\end{center}

La coordinación electrónica se obtiene componiendo estos niveles en su
orden de dependencia. APP determina el centro y las evaluaciones;
TRIT selecciona regímenes y orientaciones; TPK prolonga la trayectoria
y conserva sus registros. Las regiones previamente generadas aportan
los caracteres \(\pi,\varphi,e\); el registro dodecafásico aporta
\(\alpha\); los lectores electrónicos determinan los factores de
la ecuación \eqref{eq:el-formula-completa}. Esta composición conserva
la identidad de la sección a lo largo de sus distintas
representaciones, hasta la evaluación dimensional y el contraste
metrológico posterior.


\subsection{Norma del conmutador de las proyecciones aritméticas}

Se obtiene el prefactor electrónico sin utilizar \(\alpha\), la torsión
o una unidad dimensional. El dominio es la tercera realización de APP,
\(\mathcal A_3=\{1,\ldots,9\}^3\), provista de medida uniforme.
Definimos
\[
 \Sigma(x)=\rho_9(x_1+x_2+x_3),\qquad
 \Pi(x)=\rho_9(x_1x_2x_3).
\]
Las esperanzas condicionales \(P_\Sigma\) y \(P_\Pi\) son las
proyecciones ortogonales sobre las funciones constantes en cada fibra
aditiva y multiplicativa, respectivamente. Así, para una función \(f\),
\[
 (P_\Sigma f)(x)
 =\frac{1}{|\Sigma^{-1}(\Sigma(x))|}
   \sum_{y:\,\Sigma(y)=\Sigma(x)}f(y),
\]
y análogamente para \(P_\Pi\). Las dos proyecciones actúan sobre el
mismo espacio; no representan dos soportes aritméticos independientes.

\begin{lemma}[Dos proyecciones y ángulos principales]
Si \(s\in[0,1]\) es un valor singular del emparejamiento entre bases
ortonormales de los rangos de dos proyecciones \(P,Q\), su bloque no
trivial contribuye a \(\|[P,Q]\|\) con
\(s\sqrt{1-s^2}\).
\end{lemma}
\begin{proof}
En el plano principal correspondiente se pueden escribir
\[
 P=\begin{pmatrix}1&0\\0&0\end{pmatrix},\qquad
 Q=\begin{pmatrix}
 s^2&s\sqrt{1-s^2}\\
 s\sqrt{1-s^2}&1-s^2
 \end{pmatrix}.
\]
Su conmutador es
\[
 [P,Q]=s\sqrt{1-s^2}
 \begin{pmatrix}0&1\\-1&0\end{pmatrix}.
\]
La matriz final es ortogonal, por lo que su norma es uno. Los bloques
comunes u ortogonales tienen conmutador nulo. La descomposición ortogonal
en planos principales da la afirmación y reduce la norma total al máximo
de estas contribuciones.
\end{proof}

\begin{proposition}[Norma exacta de incompatibilidad]
Sobre \(\ell^2(\mathcal A_3)\),
\begin{equation}
 \|[P_\Sigma,P_\Pi]\|=\frac{\sqrt3}{4}.
 \label{eq:el-prefactor}
\end{equation}
\end{proposition}
\begin{proof}
Sea \(N_{ab}=|\Sigma^{-1}(a)\cap\Pi^{-1}(b)|\). La enumeración entera
de los tres símbolos proporciona
\[
 N=9\begin{pmatrix}
 0&1&1&0&1&1&0&1&4\\
 1&0&1&1&0&1&1&0&4\\
 1&0&2&0&1&2&0&0&3\\
 0&1&1&0&1&1&0&1&4\\
 1&0&1&1&0&1&1&0&4\\
 0&0&2&1&0&2&0&1&3\\
 0&1&1&0&1&1&0&1&4\\
 1&0&1&1&0&1&1&0&4\\
 0&1&2&0&0&2&1&0&3
 \end{pmatrix}.
\]
Cada fila suma \(81\); las sumas de columna son
\[
 (m_1,\ldots,m_9)=(36,36,108,36,36,108,36,36,297).
\]
Las indicatrices normalizadas de las fibras tienen matriz de
emparejamiento \(C_{ab}=N_{ab}/\sqrt{81m_b}\). Por consiguiente,
\[
 M=CC^{\mathsf T},\qquad
 M_{ac}=\sum_{b=1}^9\frac{N_{ab}N_{cb}}{81m_b}
\]
es una matriz racional cuyo espectro consta de los cuadrados de los
valores singulares requeridos.

La multiplicación exacta muestra que \((1,\ldots,1)\),
\((1,-1,0)^3\) y \((1,1,-2)^3\) son autovectores de autovalores
\(1\), \(1/4\) y \(7/132\), respectivamente. El subespacio de
vectores soportados en \(\{3,6,9\}\) cuya suma es cero tiene dimensión
dos y autovalor \(1/18\). Los vectores de suma cero soportados en
\(\{1,4,7\}\) o en \(\{2,5,8\}\) forman un núcleo de dimensión
cuatro. Estos subespacios son independientes y sus dimensiones suman
nueve. Queda, pues, probado el espectro completo
\[
 \operatorname{spec}(M)=
 \left\{1,\frac14,\frac7{132},\frac1{18},\frac1{18},0,0,0,0\right\}.
\]
El máximo de \(\lambda(1-\lambda)\) en este conjunto es \(3/16\),
alcanzado en \(\lambda=1/4\). El lema anterior da
\(\sqrt{3/16}=\sqrt3/4\).
\end{proof}

El modo \((1,-1,0)^3\) es precisamente la distribución de valores del
selector trítico de fase en la carta nonádica. La norma escalar retiene
una medida de incompatibilidad, pero no conserva por sí sola la
orientación de ese modo ni reconstruye los dos proyectores. Esa memoria
permanece en el estado de ruta que acompaña la lectura escalar.

\begin{proposition}[Álgebra del bloque electrónico local]
En el plano principal que alcanza \eqref{eq:el-prefactor}, sean
\[
 P=\begin{pmatrix}1&0\\0&0\end{pmatrix},\qquad
 Q=\begin{pmatrix}1/4&\sqrt3/4\\\sqrt3/4&3/4\end{pmatrix},\qquad
 S=2P-I,\qquad J=\frac4{\sqrt3}[P,Q].
\]
Entonces \(S^2=I\), \(J^2=-I\), \(SJ=-JS\), y el álgebra real
generada por \(S,J\) es \(M_2(\mathbb R)\), realización de
\(\mathrm{Cl}_{1,1}\). Además,
\[
 n_\pm=\frac{S\pm J}{2},\qquad
 n_\pm^2=0,\qquad n_+n_-+n_-n_+=I.
\]
\end{proposition}
\begin{proof}
Se tiene \(S=\operatorname{diag}(1,-1)\) y
\(J=\bigl(\begin{smallmatrix}0&1\\-1&0\end{smallmatrix}\bigr)\).
La multiplicación verifica las tres relaciones. Las matrices
\(I,S,J,SJ\) son linealmente independientes y forman una base de
\(M_2(\mathbb R)\), lo que prueba la afirmación sobre el álgebra.
Finalmente, \((S\pm J)^2=S^2+J^2\pm(SJ+JS)=0\), y la suma de los
productos cruzados es \((S^2-J^2)/2=I\).
\end{proof}

En este bloque conviven una dirección elíptica, una hiperbólica y dos
direcciones nulas parabólicas. El resultado describe una representación
matricial local de la aritmética de las dos hojas. No identifica toda APP
con \(M_2(\mathbb R)\), ni con un cuadrado latino cuántico; tales
afirmaciones exigirían otros dominios y otras relaciones. Tampoco la
exponencial pertenece exclusivamente al régimen parabólico: puede
aplicarse a cualquiera de los generadores de este bloque.

% Adición focal propuesta; no sustituye el álgebra electrónica existente.
% Inserción: después de la proposición «Álgebra del bloque electrónico local»
% y su párrafo explicativo, antes de «Transferencia entre hojas...».
% Procedencia y alcance: LOCALIZADORES_Y_ALCANCE.md, en esta carpeta.
\subsection{Representación espinorial del bloque electrónico y helicidad}
\label{sec:el-representacion-espinorial}

La estructura electrónica dispone de una representación que conserva más
información que su evaluación escalar. El plano principal de los dos
proyectores de APP tridimensional, seleccionado por el modo trítico
\((1,-1,0)^3\), ya ha proporcionado los operadores \(S\) y \(J\).
El primero es una involución simétrica; el segundo es el conmutador
normalizado y conserva la orientación de las hojas. Su composición permite
construir explícitamente una representación de rotaciones y separar sus
dos componentes de helicidad. Esta construcción actúa sobre el bloque
electrónico; el registro TPK de la ruta y su memoria permanecen como datos
de procedencia del bloque.

Sea \(E_e\) el plano real principal, dotado del producto escalar heredado
de las esperanzas condicionales, y sea
\(\mathscr S_e=E_e\otimes_{\mathbb R}\mathbb C\) su complejificación.
Los operadores ya obtenidos satisfacen
\[
 S^*=S,\qquad J^*=-J,\qquad
 S^2=I,\qquad J^2=-I,\qquad SJ=-JS.
\]
La complejificación extiende el cuerpo de representación. Conserva los
operadores reales y permite distinguir la unidad escalar \(i\) del
endomorfismo \(J\), aunque ambos tengan cuadrado \(-1\) en sus
respectivos dominios.

\begin{proposition}[Terna hermítica inducida por las dos hojas]
\label{prop:el-terna-hermitica}
Los operadores
\begin{equation}
 \sigma_1=SJ,\qquad \sigma_2=-iJ,\qquad \sigma_3=S
 \label{eq:el-terna-hermitica}
\end{equation}
son hermíticos y de traza nula. Satisfacen
\begin{equation}
 \sigma_a\sigma_b
 =\delta_{ab}I+i\sum_{c=1}^3\epsilon_{abc}\sigma_c,
 \qquad
 \frac12\operatorname{tr}(\sigma_a\sigma_b)=\delta_{ab},
 \label{eq:el-algebra-terna}
\end{equation}
donde \(\epsilon_{123}=1\). Por tanto constituyen una base ortonormal
del espacio real \(V_e\) de endomorfismos hermíticos de traza nula de
\(\mathscr S_e\), para el producto
\(\langle X,Y\rangle=\frac12\operatorname{tr}(XY)\).
\end{proposition}
\begin{proof}
Las relaciones anteriores dan \((SJ)^*=(-J)S=SJ\) y
\((-iJ)^*=-iJ\). Los cuadrados de las tres matrices son \(I\).
Además,
\[
 (SJ)(-iJ)=iS,\qquad (-iJ)S=iSJ,\qquad S(SJ)=J=i(-iJ).
\]
Al invertir el orden cambia el signo. Esto prueba la primera identidad
de \eqref{eq:el-algebra-terna}. En la base ortonormal del plano principal,
las matrices son
\[
 \sigma_1=\begin{pmatrix}0&1\\1&0\end{pmatrix},\quad
 \sigma_2=\begin{pmatrix}0&-i\\i&0\end{pmatrix},\quad
 \sigma_3=\begin{pmatrix}1&0\\0&-1\end{pmatrix}.
\]
Su traza es nula, y tomar trazas en la identidad de producto prueba la
ortonormalidad. Todo endomorfismo hermítico de traza nula tiene la forma
\(\bigl(\begin{smallmatrix}z&x-iy\\x+iy&-z\end{smallmatrix}\bigr)\),
con \(x,y,z\) reales, por lo que pertenece a su envolvente real.
\end{proof}

Esta terna recibe, en la representación matricial usual, el nombre de
matrices de Pauli. Su empleo aquí es posterior a la construcción de
\(P,Q,S,J\): se obtiene por los productos
\eqref{eq:el-terna-hermitica}. El espacio de direcciones de esta
representación es la esfera unidad de \(V_e\), cuya forma positiva
acaba de construirse mediante la traza.

\begin{proposition}[Generadores de rotación y descomposición direccional]
\label{prop:el-helicidad}
Defínanse \(L_a=\sigma_a/2\). Entonces
\begin{equation}
 [L_a,L_b]=i\sum_c\epsilon_{abc}L_c,
 \qquad \sum_{a=1}^3L_a^2=\frac34 I.
 \label{eq:el-generadores-spin}
\end{equation}
Para \(n=(n_1,n_2,n_3)\) con \(\sum_an_a^2=1\), sean
\[
 H_e(n)=\sum_an_a\sigma_a,\qquad
 h_e(n)=\frac12H_e(n),\qquad
 \Pi_\pm(n)=\frac12\bigl(I\pm H_e(n)\bigr).
\]
La fórmula lineal \(H_e(v)=\sum_av_a\sigma_a\) se utiliza también
para \(v\) de norma arbitraria.
Se tiene
\begin{equation}
 \mathscr S_e=\operatorname{im}\Pi_+(n)
             \mathbin{\oplus}\operatorname{im}\Pi_-(n),
 \qquad
 h_e(n)\Pi_\pm(n)=\pm\frac12\Pi_\pm(n).
 \label{eq:el-descomposicion-helicidad}
\end{equation}
Los dos sumandos son líneas complejas ortogonales. El transporte de
rotación alrededor de \(n\) es
\begin{equation}
 U_n(\theta)=\exp\!\left(-\frac{i\theta}{2}H_e(n)\right)
 =\cos\frac\theta2\,I-i\sin\frac\theta2\,H_e(n),
 \qquad
 U_n(2\pi)=-I,\quad U_n(4\pi)=I.
 \label{eq:el-retorno-spin}
\end{equation}
\end{proposition}
\begin{proof}
La identidad de producto de la terna da los conmutadores de
\eqref{eq:el-generadores-spin}; sus cuadrados suman \(3I/4\).
Los términos antisimétricos se cancelan al multiplicar
\(H_e(n)\) por sí mismo, de modo que \(H_e(n)^2=I\).
Se deduce
\[
 \Pi_\pm(n)^2=\Pi_\pm(n),\qquad
 \Pi_+(n)\Pi_-(n)=0,\qquad \Pi_+(n)+\Pi_-(n)=I.
\]
Son hermíticos y tienen traza \(1\); por ello su rango es \(1\).
Multiplicar por \(H_e(n)/2\) prueba
\eqref{eq:el-descomposicion-helicidad}. La serie de la exponencial se
separa en potencias pares e impares gracias a \(H_e(n)^2=I\), lo que
produce \eqref{eq:el-retorno-spin}. También prueba directamente la
unitariedad y la ley \(U_n(\theta+\psi)=U_n(\theta)U_n(\psi)\).

La normalización angular se verifica sobre \(V_e\). Para todo
\(v\in\mathbb R^3\), la misma identidad de producto da
\[
 U_n(\theta)H_e(v)U_n(\theta)^*
 =H_e\!\left(v\cos\theta+(n\mathbin{\times}v)\sin\theta
       +n(n\mathbin{\cdot}v)(1-\cos\theta)\right).
\]
La acción inducida es una rotación de ángulo \(\theta\) y eje \(n\).
El factor \(1/2\) en el generador espinorial es, por tanto, el que
corresponde a esa parametrización angular. Una rotación completa restaura
la dirección en \(V_e\) y cambia el signo del vector de
\(\mathscr S_e\); dos rotaciones completas restauran ambos.
\end{proof}

La representación es irreducible: un subespacio complejo invariante por
las tres matrices sería invariante por toda \(M_2(\mathbb C)\), el
álgebra que generan. Sus únicos subespacios invariantes son \(0\) y
\(\mathscr S_e\). Las relaciones
\eqref{eq:el-generadores-spin} identifican así la representación de
espín \(1/2\). Si \(n\) es la dirección de propagación en una
realización cinemática, \(h_e(n)\) es su helicidad adimensional, con
valores \(\pm1/2\). La dirección cinemática se declara al efectuar esa
realización; los dos proyectores están definidos para cada dirección
antes de seleccionar un estado de propagación particular.

\paragraph{Conservación de la estructura electrónica y tipos de inversión.}
La inversión de dirección tiene una ley exacta:
\[
 H_e(-n)=-H_e(n),\qquad \Pi_\pm(-n)=\Pi_\mp(n).
\]
Esta operación intercambia las etiquetas de helicidad relativas a una
dirección. La conjugación de la carga actúa, en cambio, sobre la firma
orientada de la ruta material; la inversión celular actúa sobre las
coordenadas del atlas. Sus dominios y acciones permanecen distintos.
La unicidad de la celda \((5,5)\) es compatible con las dos líneas de
\eqref{eq:el-descomposicion-helicidad}: una celda del atlas no es un
vector de la representación espinorial que porta.

La evaluación electrónica \(\varepsilon_e^\beta\) conserva su
fórmula y sus lectores. Sobre \(\mathscr S_e\), su realización escalar
es \(\varepsilon_e^\beta I\), que conmuta con los generadores y con
\(\Pi_\pm(n)\). La descomposición de helicidad añade la información
rotacional sin modificar la evaluación ni convertir el espín en un
coeficiente multiplicativo de la masa.

La quiralidad corresponde a la graduación de la representación relativista
que se utilice: una involución \(\Gamma_{\rm ch}\) determina sus
proyectores \((I\mp\Gamma_{\rm ch})/2\). La helicidad de
\eqref{eq:el-descomposicion-helicidad} depende además de \(n\).
Por su parte, la holonomía de una banda real de Möbius describe el cambio
de signo de una línea al recorrer un lazo de su base. La ley de esa línea,
el retorno espinorial \eqref{eq:el-retorno-spin} y la holonomía nonádica
de fase con avance de memoria son leyes de transporte sobre objetos
especificados. Su comparación conserva esos objetos: la forma helicoidal
de una trayectoria no sustituye la representación que acaba de
construirse, ni fija por sí misma la quiralidad o la carga.


\subsection{Transferencia entre hojas y coeficientes de la etapa basal}

El calendario trítico primitivo visita los modos
\((\Sigma,\Pi,\Pi)\). Al contar los pares consecutivos, incluido el
cierre del bloque, se obtienen una transición \(\Sigma\to\Pi\), una
\(\Pi\to\Pi\) y una \(\Pi\to\Sigma\). Su incidencia primitiva es
\[
 F_{\mathrm{av}}=\begin{pmatrix}0&1\\1&1\end{pmatrix}.
\]
Esta realización matricial se obtiene después del calendario. No se
emplea para reemplazar el lector coinductivo que ya generó \(\varphi\).
Su polinomio característico es \(X^2-X-1\), y el reconocimiento de las
dos raíces como \(\varphi\) y \(-\varphi^{-1}\) describe la acción de
esa salida en la transferencia modal. La acción cuadrática tiene autovalores
\(\varphi^2,-1,\varphi^{-2}\); el último es el modo estable positivo.
La marca regional \(\pi\), aplicada una sola vez, da
\(\pi/\varphi^2\).

El espacio de trayectorias, su medida y el operador de clausura están
construidos en los apartados~\ref{sec:medida-retorno-hojas}
y~\ref{sec:operador-retorno-marcado}. La fórmula
\eqref{eq:retorno-marcado-finito} conserva su procedencia como cociente
de retornos; la aplicación electrónica compone después esta lectura
con la involución siguiente. La firma regional \(555555\) acompaña
a la región de clausura, cuyo transporte conserva el estado con memoria;
el centro \((5,5)\) fija la celda de origen de la sección electrónica
orientada. Ambas informaciones convergen en la composición y mantienen
sus dominios respectivos.

\begin{lemma}[Intercambio de las dos orientaciones]
Sean \(\mathcal R(x,y)=x/y^2\) y
\(\mathscr S_{\pi\varphi}(x,y)=(y,x)\), para \(x,y>0\).
Entonces
\begin{equation}
 \begin{aligned}
 \mathscr S_{\pi\varphi}^{\,2}&=I,\qquad
 \mathcal R(\pi,\varphi)=\frac{\pi}{\varphi^2},\\
 (\mathcal R\circ\mathscr S_{\pi\varphi})(\pi,\varphi)
 &=\frac{\varphi}{\pi^2}
 =\frac1{\varphi^3(\pi/\varphi^2)^2}.
 \end{aligned}
 \label{eq:el-bisagra}
\end{equation}
\end{lemma}
\begin{proof}
El intercambio aplicado dos veces es la identidad. Además,
\(\varphi^3(\pi/\varphi^2)^2=\pi^2/\varphi\); tomar su inversa prueba
la última igualdad.
\end{proof}

El lector electrónico utiliza la orientación \(\varphi/\pi^2\),
seguida del carácter exponencial. La identidad anterior explica el
intercambio, pero no sustituye la regla que elige ese carácter ni produce
por sí sola los enteros de la corrección. Estos poseen otra procedencia.

El retorno central compatible con la fase tiene longitud mínima
veintisiete. Con el registro de fase y hoja, su espacio finito es
\[
 \mathcal K_e=\mathbb Z/9\mathbb Z\times\mathbb F_3.
\]
Las cinco regiones de \(\pi\) son identidades genealógicas distintas,
no cinco copias intercambiables de una expansión visible. En el calibre
regional fijado por la construcción, su inyección en el retorno tiene
los lugares
\begin{equation}
 (4,0),\quad(5,0),\quad(6,0),\quad(6,1),\quad(6,2).
 \label{eq:el-cinco-lugares}
\end{equation}
La orientación transversal, el retorno mutado y las tres regiones
positivas permanecen separados por ese registro. Los cambios admisibles
de calibre transportan las etiquetas conjuntamente; no vuelven a
seleccionar los lugares a partir de un valor electrónico.

\begin{proposition}[Déficit y memoria retenida]
Fijados el retorno central y la inyección regional anterior, el déficit
orientado de frontera y la memoria torsional son
\begin{equation}
 n_e=-|\mathcal K_e\setminus\iota(\mathscr R_\pi)|=-22,
 \qquad
 m_e^{\mathrm{tor}}=|\mathscr R_\pi\times\mathbb F_3|=15.
 \label{eq:el-selector-enteros}
\end{equation}
\end{proposition}
\begin{proof}
Los cinco lugares de \eqref{eq:el-cinco-lugares} son distintos. Por
tanto, \(|\mathcal K_e|=9\cdot3=27\) y el complemento de la imagen
tiene cardinal \(27-5=22\). La convención de orientación registra como
negativo el déficit retirado de la frontera. La memoria retenida posee
tres hojas por región, de modo que su cardinal positivo es
\(5\cdot3=15\).
\end{proof}

La misma pareja posee el control entero de la carta APP de la familia:
\begin{equation}
 396=18+72+108+198,
 \qquad \frac{396}{18}=22,
 \qquad \frac{270}{18}=15=3\cdot5=\binom62.
 \label{eq:el-censo}
\end{equation}
El peso \(18\) pertenece al núcleo central \(2\times2\), y \(270\)
al canal global declarado. Estas igualdades controlan la normalización
del recuento; el argumento de retorno fija además los signos. El
cociente de los anillos alternos,
\(396/(18+108)=22/7\), es racional y no se identifica con \(\pi\).
Tampoco \(\binom62\) significa \(6/2\): cuenta los pares no ordenados
de un conjunto de seis elementos.

La evaluación del déficit en la tercera realización simétrica emplea
\(\alpha^3\). Esta potencia es el producto de tres copias del escalar
de cierre ya obtenido; no afirma una regla universal que haga
corresponder cualquier dimensión con la misma potencia de \(\alpha\).
El dominio y el lector fijados son esenciales para la composición.

\subsection{La torsión global y la coordenada basal}

\begin{definition}[Normalización torsional electrónica]
Se distinguen el defecto reducido \(d_4\) y el lector completo
\(\Delta_4\):
\begin{equation}
 d_4=\frac{\pi-e\log\pi}{270},\qquad
 \Delta_4=d_4\left(1+\frac{\pi}{729}\right).
 \label{eq:el-delta}
\end{equation}
El denominador \(270\) corresponde al canal del censo, mientras el
factor \(1+\pi/729\) retiene el acoplamiento con la celda cúbica de
\(729\) lugares. La lectura electrónica utiliza \(\Delta_4\) tanto en
su etapa basal como en el exponente de retorno.
\end{definition}

Esta definición actúa sobre \(\pi\) y \(e\) previamente generados. El
logaritmo es una operación posterior en su realización real, no una
entrada retrospectiva de valores convencionales. La torsión es global:
no se determina independientemente para cada sección mediante una masa
observada.

\begin{lemma}[Las dos normalizaciones no son intercambiables]
Se cumple
\begin{equation}
 \Delta_4-d_4=\frac{\pi}{729}d_4.
 \label{eq:el-delta-diferencia}
\end{equation}
En particular, \(d_4>0\) y \(\Delta_4>d_4\) para las coordenadas
consideradas.
\end{lemma}
\begin{proof}
La identidad es distributiva. Para la positividad, la función
\(f(x)=x-e\log x\), sobre \(x>0\), tiene derivada \(1-e/x\) y
segunda derivada \(e/x^2>0\). Su mínimo único está en \(x=e\), donde
vale cero. Como \(\pi\ne e\), se tiene \(f(\pi)>0\). Ambos factores
de \eqref{eq:el-delta} son positivos y \(1+\pi/729>1\).
\end{proof}

\begin{remark}[Conservación de la normalización]
Suprimir el segundo factor de \eqref{eq:el-delta} cambia el lector; no
es una simplificación algebraica. Si se introduce \(d_4\) como
coordenada auxiliar, el mismo objeto debe escribirse
\(\Delta_4=(1+\pi/729)d_4\) en todas sus apariciones. Mantener los
enteros \(15\) y \(6\) y sustituir sólo \(\Delta_4\) por \(d_4\)
produce otra fórmula.
\end{remark}

\begin{definition}[Coordenada electrónica basal]
La composición de incompatibilidad, transferencia modal, déficit y memoria retenida
es
\begin{equation}
 \varepsilon_e^{(0)}
 =\frac{\sqrt3}{4}
  \left[\exp\!\left(\frac{\varphi}{\pi^2}\right)-22\alpha^3\right]
  (1+15\Delta_4).
 \label{eq:el-basal}
\end{equation}
\end{definition}

La ecuación reúne tres operaciones de naturaleza distinta. La norma
\(\sqrt3/4\) cuantifica la incompatibilidad de las proyecciones de
las hojas; su origen es un ángulo principal de las fibras de APP.
La exponencial actúa sobre la orientación \(\varphi/\pi^2\) del
lector de transferencia. La resta cúbica registra las \(22\) posiciones
complementarias del retorno marcado, y el factor torsional conserva las
\(15\) componentes correspondientes a tres hojas por región. La
composición conserva estos dominios aunque su evaluación final sea escalar.

La presencia de \(\pi\), \(\varphi\) y \(e\) tiene así significados
operativos. El número de Euler interviene por la exponencial y por
la normalización logarítmica de \(\Delta_4\); \(\varphi\) y \(\pi\)
intervienen en la transferencia orientada. La constante \(\alpha\)
actúa después de su propia construcción mediante el término cúbico
de clausura. El retorno cronológico aún debe operar sobre esta coordenada
basal. La separación de etapas permite explicar qué cambia al modificar
una orientación, una normalización o un registro, conservando la estructura
central que identifica la sección electrónica.

Los paréntesis de esta ecuación son parte de su contenido. El término
\(-22\alpha^3\) se resta \emph{fuera} de la exponencial; la memoria
\(1+15\Delta_4\) multiplica todo el corchete. Introducir el déficit en
el exponente o aplicarle una torsión distinta alteraría la composición.

\begin{proposition}[Positividad basal]
Para \(0<\alpha<10^{-2}\), \(\varphi>0\) y la normalización
\eqref{eq:el-delta}, se tiene \(\varepsilon_e^{(0)}>0\).
\end{proposition}
\begin{proof}
La exponencial es mayor que uno, y
\(22\alpha^3<22\cdot10^{-6}<1\). Por ello el corchete es positivo.
También \(1+15\Delta_4>1\) y \(\sqrt3/4>0\).
\end{proof}

La evaluación posterior de \eqref{eq:el-basal} comienza por
\[
 \varepsilon_e^{(0)}=0.5109989224880660725\ldots.
\]
Esta cifra es todavía una coordenada sin dimensión y no constituye la
última etapa: falta aplicar la memoria de retorno de la propia ruta.

\subsection{Escala determinantal de acción y valoración decimal}

La escala absoluta y el corrector relativo se originan en objetos
distintos. El primero utiliza el bloque local del registro dodecafásico
de paso tres,
\begin{equation}
 H_4=\begin{pmatrix}
 1&1&1&1\\1&1&-1&-1\\1&-1&1&-1\\1&-1&-1&1
 \end{pmatrix},\qquad
 G_H=\mathbb Z^4/H_4\mathbb Z^4.
 \label{eq:el-h4}
\end{equation}
La operación utiliza un cociente local. Los tres bloques del registro
completo no se sustituyen por una tensorización que multiplicara de
nuevo el número de hojas de este lector.

\begin{lemma}[Número de hojas del cociente local]
La forma normal de Smith de \(H_4\) es
\(\operatorname{diag}(1,2,2,4)\). En consecuencia,
\[
 G_H\simeq\mathbb Z/2\mathbb Z\oplus\mathbb Z/2\mathbb Z
 \oplus\mathbb Z/4\mathbb Z,\qquad |G_H|=16.
\]
\end{lemma}
\begin{proof}
El máximo común divisor de las entradas es uno. Los menores de orden dos
son pares y hay uno de valor \(-2\), de modo que su máximo común divisor
es dos. Como \(H_4H_4^{\mathsf T}=4I\) y \(\det H_4=-16\), los
cofactores son \(\pm4\); el tercer divisor determinantal es cuatro.
El cuarto es \(16\). Los cocientes sucesivos de los divisores
\(1,2,4,16\) son \(1,2,2,4\). Su producto da el orden del cociente.
\end{proof}

\begin{definition}[Lector determinantal de acción]
La norma multiplicativa de la sección escalar \(\alpha\), constante
sobre las hojas de \(G_H\), y una aplicación de la autoescala
\(\varphi\) producen
\begin{equation}
 \Sigma_H=\varphi\,\mathrm N_{G_H}(\alpha),\qquad
 \mathrm N_{G_H}(\alpha)=\prod_{g\in G_H}\alpha=\alpha^{16}.
 \label{eq:el-lector-determinantal}
\end{equation}
La carta decimal asociada determina
\(\operatorname{ord}_{10}(\Sigma_H)=\lfloor\log_{10}\Sigma_H\rfloor\).
\end{definition}

La potencia dieciséis procede así del número de hojas del cociente, una
vez fijado este lector local. El factor \(\varphi\) actúa una vez sobre
la base de la recta de acción. El cálculo de Smith, por sí solo, no
seleccionaría cualquier funcional posible sobre el cociente: la norma
multiplicativa y la autoescala forman parte de la operación declarada.

\begin{theorem}[Década de la sección de acción]
El lector \eqref{eq:el-lector-determinantal}, aplicado a las coordenadas
HMT ya generadas, satisface
\begin{equation}
 10^{-34}<\Sigma_H<10^{-33},\qquad
 \operatorname{ord}_{10}(\Sigma_H)=-34.
 \label{eq:el-decada}
\end{equation}
\end{theorem}
\begin{proof}
Bastan los intervalos racionales, mucho más gruesos que los cilindros
disponibles,
\[
 \frac{1618}{1000}<\varphi<\frac{1619}{1000},\qquad
 \frac{7297}{10^6}<\alpha<\frac{7298}{10^6}.
\]
La función \((x,a)\mapsto xa^{16}\) es creciente en ambos argumentos
positivos. Las dos desigualdades enteras
\[
 1618\cdot7297^{16}>10^{65},\qquad
 1619\cdot7298^{16}<10^{66}
\]
implican
\[
 10^{-34}
 <\frac{1618}{1000}\left(\frac{7297}{10^6}\right)^{16}
 <\Sigma_H
 <\frac{1619}{1000}\left(\frac{7298}{10^6}\right)^{16}
 <10^{-33}.
\]
La definición del orden decimal concluye la prueba. No interviene una
unidad SI ni una masa de referencia en estas desigualdades.
\end{proof}

La sección normalizada comienza por
\(\Sigma_H=1.0462438381\ldots\times10^{-34}\). Su mantisa no se
identifica con la coordenada de acción anterior o posterior al retorno.
El lector determinantal fija el orden decimal; las dos coordenadas de
retorno se obtienen mediante los lectores siguientes.


\subsection{Secciones de acción y cociente adimensional de retorno}

Sea \(\mathcal L_S\) una recta orientada de acción y \(\mathcal U_S\)
su base positiva. La composición electrónica utiliza dos secciones de
esta recta y su cociente adimensional. Su construcción requiere la
década determinantal ya obtenida, el carácter de acción de quinto orden
y la contribución regional del retorno.

\subsubsection{Carácter de acción y procedencia de sus coeficientes}
\label{sec:revision-planck}

La constante de Planck reducida, \(\hbar\), es el objeto físico que
motiva la realización de la sección de acción angular. Su posición en
la dependencia electrónica debe especificarse en dos niveles. La norma
determinantal \(\Sigma_H\) determina la década de la escala. El carácter
\(H_5\) y el retorno regional determinan después las coordenadas de las
secciones de acción. Esta separación permite exponer el resultado
necesario para el electrón sin anticipar los desarrollos completos de
otros sectores de la mecánica dimensional.

La procedencia del carácter de quinto orden puede explicitarse mediante
invariantes ya fijados en el soporte. En la carta central se considera
\[
 B_c=\begin{pmatrix}7&2\\2&7\end{pmatrix},
 \qquad \lambda_+=9,\quad\lambda_-=5.
\]
El calendario tiene longitud \(12\), la órbita del paso \(3\) tiene
longitud \(4\), y el cociente censal pertinente es
\(104976/1944=54\). Los coeficientes del carácter satisfacen
\begin{equation}
 \frac9{16}=\frac{\lambda_+}{4^2},\quad
 \frac59=\frac{\lambda_-}{\lambda_+},\quad
 \frac7{48}=\frac{\operatorname{tr}(B_c)/2}{4\cdot12},\quad
 \frac1{54}=\frac{1944}{104976}.
 \label{eq:revision-h5-coeficientes}
\end{equation}
La asignación de estos invariantes a los grados \(2,3,4,5\), con la
alternancia de signos de \eqref{eq:el-h5}, pertenece a la definición del
carácter analítico utilizado. Explicitarla es esencial: el conjunto de
cardinales, considerado sin el carácter y sin el orden de sus grados,
admitiría otras expresiones. El contenido constructivo reside en la
composición del soporte, la orientación y la regla analítica declarada.

Sobre esta base, el carácter de quinto orden y el término regional de
retorno tienen las expresiones
\begin{align}
 H_5={}&\frac12\sqrt{\frac{1000\alpha}{\varphi}}-\alpha
   +\frac9{16}\alpha^2-\frac59\alpha^3
   +\frac7{48}\alpha^4-\frac1{54}\alpha^5,
   \label{eq:el-h5}\\
 D_A={}&\exp\!\left(-\frac{100\pi\alpha}{9}\right),
   \label{eq:el-da}\\
 \mathcal C_\pi(\alpha)={}&169\alpha^6+
       \frac{D_A\alpha^7}{1-D_A\alpha},
   \label{eq:el-cpi}\\
 \eta_{\mathrm{ret}}={}&H_5-\frac{90}{\pi}\mathcal C_\pi(\alpha).
   \label{eq:el-eta}
\end{align}
Estas son lecturas posteriores del cierre dodecafásico y de su
incidencia regional. El entero \(169\) se conserva como el selector
regional ya construido, y los coeficientes de \(H_5\) como los de su
carácter local; no se reemplazan por una interpolación de la constante de
Planck. Las fórmulas explicitan exactamente qué lecturas se necesitan
para la etapa electrónica. No es necesario introducir aquí otros
funcionales de la teoría de masas.

Para \(0<\alpha<1\), se tiene \(0<D_A<1\), de modo que
\(D_A\alpha<1\). El término racional de \eqref{eq:el-cpi} conserva la
suma completa de la cola geométrica
\begin{equation}
 \frac{D_A\alpha^7}{1-D_A\alpha}
 =\sum_{n=0}^{\infty}D_A^{n+1}\alpha^{n+7}.
 \label{eq:el-cola-regional}
\end{equation}
No es una truncación que pueda modificarse silenciosamente al cambiar la
precisión. Tras el término de índice \(N\), la cola restante es
\(D_A^{N+2}\alpha^{N+8}/(1-D_A\alpha)\), lo que permite controlar la
evaluación sin acudir a una magnitud objetivo.

\begin{definition}[Secciones de acción y lector de retorno]
Las secciones anterior y posterior al retorno son
\begin{equation}
 \hbar_{\mathrm{pre}}=H_5\,10^{-34}\mathcal U_S,
 \qquad
 \hbar_{\mathrm{ret}}=\eta_{\mathrm{ret}}\,10^{-34}\mathcal U_S.
 \label{eq:el-secciones-accion}
\end{equation}
Siempre que \(\eta_{\mathrm{ret}}>0\), se definen
\begin{equation}
 \delta_{\mathrm{ret}}=H_5-\eta_{\mathrm{ret}},\qquad
 R_{\mathrm{act}}=\frac{\hbar_{\mathrm{pre}}}{\hbar_{\mathrm{ret}}}
 =\frac{H_5}{\eta_{\mathrm{ret}}}.
 \label{eq:el-ratio}
\end{equation}
\end{definition}

\begin{proposition}[Positividad e invariancia de la razón]
En los intervalos usados en la prueba de \eqref{eq:el-decada},
\(0<\eta_{\mathrm{ret}}<H_5\), y por tanto
\(R_{\mathrm{act}}>1\). La razón no cambia al sustituir la base
\(\mathcal U_S\) por cualquier otra base positiva común.
\end{proposition}
\begin{proof}
Los términos de \(\mathcal C_\pi\) son positivos. Para una cota
elemental basta usar \(\alpha<0.008\), \(\varphi<1.619\),
\(\alpha>0.007\), \(D_A<1\) y \(\pi>3\). La raíz en
\eqref{eq:el-h5} es mayor que uno, mientras la suma de los términos
negativos es menor que \(0.008001\); en particular, \(H_5>0.99\).
Asimismo,
\[
 0<\frac{90}{\pi}\mathcal C_\pi
 <30\left(169(0.008)^6+
           \frac{(0.008)^7}{1-0.008}\right)<10^{-6}.
\]
Así \(\eta_{\mathrm{ret}}>0.989999>0\) y
\(\eta_{\mathrm{ret}}<H_5\). Finalmente, el cambio de base divide
ambas coordenadas de acción por el mismo factor positivo, que cancela en
su cociente.
\end{proof}

La evaluación posterior proporciona
\begin{align*}
 H_5&=1.0545718176461544696\ldots,\\
 \eta_{\mathrm{ret}}&=1.0545718169150390611\ldots,\\
 R_{\mathrm{act}}&=1.0000000006932817630\ldots.
\end{align*}
La diferencia y el cociente son dos lecturas de las mismas secciones,
una aditiva y otra multiplicativa. No representan dos constantes de
Planck independientes. El paso de acción angular a acción por vuelta
completa se expresa, en cada sección, mediante
\(h=2\pi\hbar\). La posterior carta
\(\mathcal U_S\mapsto\mathrm{J\,s}\) asigna una unidad, pero no
retroactúa sobre el orden decimal ya demostrado.

El retorno regional puede escribirse también en una carta angular.
Esta reparametrización no es una dependencia adicional indispensable:
\eqref{eq:el-h5}--\eqref{eq:el-eta} suministran directamente
\(H_5\), \(\eta_{\mathrm{ret}}\) y \(R_{\mathrm{act}}\). Si se
introducen ángulo y contraángulo, ambos deben transportarse en la misma
carta; no se mezclan grados y radianes dentro de una diferencia.

\subsubsection{Continuación regional y ecuación de renovación}

El selector \(169\) determina el impulso de grado seis. La prolongación
posterior se organiza por la transferencia \(D_A\) y una normalización
unitaria de la línea determinantal. Estas reglas admiten una formulación
en series formales antes de sustituir \(z=\alpha\), con \(D_A\) ya
construido. Escribamos
\[
 \mathscr C(z)=169z^6+R(z),\qquad R(z)\in z^7\mathbb R[[z]].
\]
La concatenación de una prolongación aumenta el grado en una unidad y
multiplica su peso por \(D_A\). La primera prolongación estricta tiene
peso \(D_A\). Por tanto, su ecuación de renovación es
\begin{equation}
 R(z)=D_Az^7+D_AzR(z).
 \label{eq:revision-renovacion}
\end{equation}

\begin{proposition}[Unicidad formal de la continuación normalizada]
La ecuación \eqref{eq:revision-renovacion} determina una única serie y
\[
 \mathscr C(z)=169z^6+\frac{D_Az^7}{1-D_Az}
             =169z^6+\sum_{n=7}^{\infty}D_A^{n-6}z^n.
\]
La realización es analítica en \(|D_Az|<1\).
\end{proposition}
\begin{proof}
El elemento \(1-D_Az\) tiene término independiente uno y es invertible
en el anillo de series formales. Resolver la ecuación produce la serie
exhibida. Equivalentemente, el coeficiente de grado siete es \(D_A\)
y cada coeficiente posterior es \(D_A\) veces el anterior. La serie
geométrica converge absolutamente en el disco indicado.
\end{proof}

La normalización de la primera prolongación es un dato operativo
indispensable. Si se conservaran únicamente el impulso de grado seis y
la razón de concatenación, las series
\[
 169z^6+\lambda\frac{D_Az^7}{1-D_Az}
\]
seguirían compartiendo ambos datos para cualquier \(\lambda\).
La regla unitaria fija \(\lambda=1\) antes de efectuar la evaluación.
Así queda localizada la contribución de cada regla a la unicidad del
lector, y la expresión racional conserva la continuación completa.

\subsubsection{Constante de Planck reducida y retorno de la sección de acción}

En la carta dimensional de \eqref{eq:el-secciones-accion}, la sección
anterior al retorno es la construcción del modelo destinada al contraste
con la constante de Planck reducida:
\[
 \hbar_{\mathrm{pre}}=H_5\,10^{-34}\mathcal U_S.
\]
La sección posterior conserva, además, la compensación regional
\(90\mathscr C(\alpha)/\pi\). Ambas pertenecen a la misma recta de
acción y su cociente \(R_{\mathrm{act}}\) mide el efecto multiplicativo
del retorno. El paso a \(h\) mediante \(2\pi\) utiliza la relación
entre la parametrización angular y una vuelta completa.

Esta construcción interviene efectivamente en
\(R_{\mathrm{act}}^{80-54\alpha+6\Delta_4}\). El electrón conserva
su coordenada central y su trayectoria; la escala cronológica no
reemplaza esa estructura por un cuanto indiferenciado. Tampoco se
multiplica de nuevo por \(10^{-34}\) una masa ya expresada en MeV:
la década pertenece a la sección de acción, cancela en su cociente
adimensional y reaparece únicamente donde actúa la realización
dimensional correspondiente.

El contraste numérico de \(H_5\), de la sección retornada y de
\(h/(2\pi)\) se presenta en la sección~\ref{sec:revision-planck-contraste}.
Así se distingue la denominación del objeto físico, la ecuación que el
modelo le asigna y la precisión con la que dicha ecuación lo reproduce.


\subsection{Dos lectores del retorno electrónico}

La memoria electrónica se representa en una palabra trítica de longitud
\(108\) y en una palabra firmada de doce fases. Para que el cálculo del
exponente sea reproducible, explicitamos la proyección cronológica
utilizada. Sean inicialmente
\((r_0,c_0)=(5,5)\) y \((h_0,v_0)=(1,1)\). En cada paso el selector
repite \((+1,-1,0)\). Con orientación fija dentro de cada bloque de
veintisiete pasos, la regla visible es
\[
 \begin{array}{c|c|c}
 \text{modo}&\text{entero leído}&\text{transporte celular}\\\hline
 +1&r+c&c\mapsto1+(c-1+h)\bmod9\\
 -1&rc&r\mapsto1+(r-1+v)\bmod9\\
 0&9&\text{sin desplazamiento}
 \end{array}
\]
El símbolo emitido es \(\rho_9(\text{entero leído})\bmod3\).
Después de los pasos \(27,54,81\) se invierten simultáneamente
\(h\) y \(v\). Esta regla calcula la proyección observable; la
actualización del estado conserva adicionalmente residuo, cociente,
acarreo, hoja, frontera y memoria. No se identifica el registro de
\(108\) símbolos con el estado completo.

En cada bloque completo de tres pasos se ha avanzado una vez en cada
coordenada; hacen falta nueve de estos bloques para volver al centro en
la misma fase. La longitud mínima de este retorno registrado es, por
tanto, \(27\). Si se olvida la fase, la celda ya se alcanza tras el
paso \(26\), antes del umbral neutro que completa el bloque; ése no es
el retorno tipado que cuenta \(\mathcal K_e\). La lectura explícita
de los dos primeros bloques de retorno da
\begin{equation}
 a=100000220,\qquad b=120200000,\qquad
 \gamma_e=(a^3b^3)^2.
 \label{eq:el-palabra}
\end{equation}
Aquí las potencias denotan concatenación de palabras. En particular,
\(\gamma_{e,t+54}=\gamma_{e,t}\). En el registro de nueve pasos por
fase, las orientaciones centrales dan
\begin{equation}
 \tau_e=(1,1,1,-1,-1,-1,1,1,1,-1,-1,-1).
 \label{eq:el-tau}
\end{equation}
La periodicidad de estas representaciones no equivale al reinicio de la
historia que las produce.

\begin{definition}[Lector de discrepancias cíclicas]
Para una palabra \(\gamma\in\mathbb F_3^{108}\), se define
\[
 D_k(\gamma)=\#\{t\in\mathbb Z/108\mathbb Z:
                   \gamma_t\ne\gamma_{t+k}\}.
\]
En el dominio de palabras de período divisor de \(54\), el lector
entero es
\begin{equation}
 K_D(\gamma)=\left(D_{27}+D_{36},D_1,\frac{D_4-D_3}{2}\right).
 \label{eq:el-kd}
\end{equation}
\end{definition}

\begin{lemma}[Integridad y evaluación central]
Los \(D_k\) del dominio anterior son pares. Para
\eqref{eq:el-palabra},
\[
 \begin{array}{c|rrrrr}
 k&1&3&4&27&36\\\hline
 D_k&54&56&68&48&32
 \end{array}
\]
y \(K_D(\gamma_e)=(80,54,6)\).
\end{lemma}
\begin{proof}
La aplicación \(t\mapsto t+54\) empareja sin puntos fijos los índices
de discrepancia y conserva su condición; cada cardinal es par. Para la
palabra explícita \(a^3b^3\) de longitud \(54\), el recuento de
discrepancias con cada desplazamiento es, en el mismo orden,
\(27,28,34,24,16\). Duplicar la palabra duplica esos cinco cardinales.
Finalmente,
\[
 (D_{27}+D_{36},D_1,(D_4-D_3)/2)
 =(48+32,54,(68-56)/2)=(80,54,6).
\]
\end{proof}

Los desplazamientos \(27\) y \(36\) son las elevaciones
cronológicas de los canales de \(90\) y \(120\) grados en una carta
donde nueve pasos corresponden a treinta grados. La lectura angular
explica esa representación, pero los enteros se computan sobre la
palabra, sin introducir una masa ni un valor real buscado.

\begin{definition}[Lector de ventanas dodecafásicas]
Para \(\tau\in\{-1,0,1\}^{12}\), con índices cíclicos, sean
\begin{align*}
 C_k(\tau)&=\sum_{r=0}^{11}\tau_r\tau_{r+k},\\
 A_\ell(\tau)&=\sum_{r=0}^{11}
           \left(\sum_{j=0}^{\ell-1}\tau_{r+j}\right)^2,\\
 P_6(\tau)&=\sum_{r=0}^{5}\tau_r\tau_{r+6}.
\end{align*}
Se define
\begin{equation}
 \Omega(\tau)=4A_3(\tau)-3A_4(\tau),\qquad
 K_\Omega(\tau)=(\Omega(\tau),54,P_6(\tau)).
 \label{eq:el-komega}
\end{equation}
El \(54\) de este lector registra la semivuelta cronológica declarada.
No se confunde con el salto local de valor cuatro de APP.
\end{definition}

\begin{lemma}[Contraste primitivo y forma de correlación]
Entre los contrastes lineales enteros \(uA_3+vA_4\) que se anulan
siempre que \(A_3=3a\) y \(A_4=4a\), la pareja primitiva, con
orientación \(90-120\), es \((u,v)=(4,-3)\). Además,
\begin{equation}
 \Omega=-2C_1-4C_2-6C_3.
 \label{eq:el-omega-corr}
\end{equation}
\end{lemma}
\begin{proof}
La anulación para todo \(a\) exige \(3u+4v=0\). Por coprimalidad,
las soluciones son \((u,v)=n(4,-3)\), y la orientación fija el signo
del generador primitivo. Al desarrollar los cuadrados,
\[
 A_\ell=\ell C_0+2\sum_{k=1}^{\ell-1}(\ell-k)C_k.
\]
Por tanto,
\(A_3=3C_0+4C_1+2C_2\) y
\(A_4=4C_0+6C_1+4C_2+2C_3\). Su combinación \(4A_3-3A_4\)
cancela \(C_0\) y produce \eqref{eq:el-omega-corr}.
\end{proof}

La unicidad demostrada pertenece a esa clase de contrastes lineales
enteros; no excluye por sí sola todo funcional no lineal de las ventanas.
La especificación de la clase funcional evita convertir una identidad
de dos canales en una afirmación de selección universal.

\begin{theorem}[Triple de retorno de la sección electrónica]
Los dos lectores aplicados a los registros de \(\Gamma_e\)
coinciden y dan
\begin{equation}
 K_D(\gamma_e)=K_\Omega(\tau_e)=(80,54,6).
 \label{eq:el-triple}
\end{equation}
Su evaluación en las coordenadas ya generadas es
\begin{equation}
 \kappa_e=80-54\alpha+6\Delta_4.
 \label{eq:el-kappa}
\end{equation}
\end{theorem}
\begin{proof}
El primer lector se calculó arriba. Para \(\tau_e\), el producto
antipodal es siempre uno, y así \(P_6=6\). El recuento cíclico da
\(C_1=4\), \(C_2=-4\), \(C_3=-12\); luego
\[
 \Omega=-2\cdot4-4(-4)-6(-12)=80.
\]
Equivalentemente, \(A_3=44\), \(A_4=32\) y
\(4\cdot44-3\cdot32=80\). Por tanto ambos triples son
\((80,54,6)\). Evaluar el funcional
\((A,B,C)\mapsto A-\alpha B+\Delta_4C\) prueba la última fórmula.
\end{proof}

Esta coincidencia corresponde a la sección central. No identifica los
dos lectores en todos los dominios de rutas: uno compara símbolos
tríticos y el otro compara ventanas de un registro firmado. Tampoco el
\(80\) se obtiene del período de una palabra asignada al número de
Euler. Su origen en esta composición es el triple de la ruta
electrónica.

\subsection{Composición electrónica completa y normalización}

\begin{definition}[Retorno multiplicativo]
La acción de la memoria de retorno sobre la coordenada basal es
\begin{equation}
 \varepsilon_e^\beta
 =\varepsilon_e^{(0)}\exp\!\left(\kappa_e\log R_{\mathrm{act}}\right).
 \label{eq:el-beta}
\end{equation}
\end{definition}

\begin{theorem}[Factorización electrónica positiva]
Con los lectores, la orientación y la normalización declarados,
\eqref{eq:el-beta} coincide con \eqref{eq:el-formula-completa}. Todos
sus factores son adimensionales y la coordenada resultante es positiva.
\end{theorem}
\begin{proof}
La norma de incompatibilidad es \eqref{eq:el-prefactor}; la transferencia modal,
el déficit y la memoria retenida componen \eqref{eq:el-basal}. El
triple \eqref{eq:el-triple} produce \eqref{eq:el-kappa}, y la razón de
las dos secciones de acción es \eqref{eq:el-ratio}. Sustituyendo esas
tres identidades en \eqref{eq:el-beta} se obtiene exactamente
\eqref{eq:el-formula-completa}. La etapa basal es positiva y la
exponencial de un número real también. El factor común de dimensión de
acción ya canceló al formar \(R_{\mathrm{act}}\).
\end{proof}

La evaluación de esta composición comienza por
\begin{equation}
 \varepsilon_e^\beta
 =0.5109989506900008086082571648\ldots.
 \label{eq:el-beta-decimal}
\end{equation}
Las cifras son una evaluación posterior de la fórmula, no datos usados
para seleccionar \(22\), \(15\) o \((80,54,6)\). La representación
numérica del basal mediante un prefijo finito introduce únicamente su
error de truncación; no modifica la fórmula exacta que lo define.

La incidencia de una normalización torsional distinta puede localizarse
sin compararla con una medida. Escribamos
\(B=\frac{\sqrt3}{4}[\exp(\varphi/\pi^2)-22\alpha^3]\),
\(R=R_{\mathrm{act}}\) y
\[
 \mathcal E(d)=B(1+15d)R^{80-54\alpha+6d}.
\]
Para los dos defectos de \eqref{eq:el-delta},
\begin{equation}
 \frac{\mathcal E(\Delta_4)}{\mathcal E(d_4)}
 =\frac{1+15\Delta_4}{1+15d_4}
   R^{6(\Delta_4-d_4)}>1.
 \label{eq:el-normalizacion-beta}
\end{equation}
La primera desigualdad torsional y \(R>1\) prueban que no son la misma
coordenada. Si sólo se cambia el defecto del exponente, manteniendo el
basal completo, la razón de las dos salidas es
\(R^{6(\Delta_4-d_4)}\), también distinta de uno. Una diferencia muy
pequeña sigue siendo una diferencia de fórmula; no puede atribuirse a
una simple conversión de unidades.

\subsection{Realización dimensional: energía, masa y escala cronológica}

La coordenada electrónica actúa finalmente sobre una recta dimensional.
Sea \(\mathcal L_E\) una recta de energía con base positiva
\(\mathcal U_E\). La realización es
\begin{equation}
 E_e^{(0)}=\varepsilon_e^{(0)}\mathcal U_E,\qquad
 E_e^\beta=\varepsilon_e^\beta\mathcal U_E,
 \qquad m_e^\beta=\frac{E_e^\beta}{c_{\mathrm{int}}^2}.
 \label{eq:el-dimensional}
\end{equation}
La última expresión pertenece a la recta de masa cuando
\(c_{\mathrm{int}}\) lleva la dimensión de velocidad en la realización
declarada. En la carta comparativa \(\mathcal U_E=1\,\mathrm{MeV}\),
\eqref{eq:el-beta-decimal} expresa una energía en MeV; la masa
correspondiente se expresa en \(\mathrm{MeV}/c^2\). Es frecuente
abreviar ambas magnitudes tomando \(c=1\), pero esa abreviatura no
debe ocultar el mapa dimensional.

La escala de acción también admite una realización acción--reloj. Si
\(t_0\) es una sección temporal positiva, el reloj de \(108\) pasos
define
\begin{equation}
 \omega_0=\frac{2\pi}{108t_0},\qquad
 E_{\mathrm{clk}}=\hbar_{\mathrm{ret}}\omega_0,\qquad
 m_{\mathrm{clk}}=E_{\mathrm{clk}}c_{\mathrm{int}}^{-2}.
 \label{eq:el-cuanto-reloj}
\end{equation}
La acción por frecuencia tiene dimensión de energía. Este cuanto
cronológico es una sección dimensional común; no es, por definición,
la sección electrónica. La palabra de ruta y sus factores no se
obtienen de la uniformidad del reloj.

\begin{proposition}[Cancelación de la escala común]
El factor \(10^{-34}\mathcal U_S\) de
\eqref{eq:el-secciones-accion} no forma parte de \(\kappa_e\). En
una misma recta de energía,
\begin{equation}
 \frac{E_e^\beta}{E_e^{(0)}}=R_{\mathrm{act}}^{\kappa_e}.
 \label{eq:el-razon-energia}
\end{equation}
El resultado es independiente de la elección de la base dimensional
común.
\end{proposition}
\begin{proof}
Al formar \(\hbar_{\mathrm{pre}}/\hbar_{\mathrm{ret}}\), el factor
\(10^{-34}\mathcal U_S\) cancela exactamente. Al dividir las dos
expresiones de energía de \eqref{eq:el-dimensional}, cancela además
\(\mathcal U_E\), y \eqref{eq:el-beta} da la razón indicada.
\end{proof}

Por tanto, una coordenada ya expresada en MeV no se multiplica otra vez
por \(10^{-34}\). Si se elige \(E_{\mathrm{clk}}\) como base de
\(\mathcal L_E\), deben transportarse las coordenadas por el factor
de cambio de base; conservar sin cambio la cifra de la carta de un MeV
sería identificar dos bases sin demostrarlo. Éste es el punto preciso
en que la escala de acción entra en la realización absoluta y deja de
intervenir en el corrector relativo.

Existe asimismo una realización lineal que hace visible la separación
entre reloj y genealogía. En
\(\ell^2(\mathbb Z/108\mathbb Z;\mathbb R)\), el desplazamiento
cíclico \(S\) tiene espacio fijo
\[
 \ker(S-I)=\mathbb R n_0,\qquad
 n_0=\frac1{\sqrt{108}}\sum_{t=0}^{107}\delta_t.
\]
Si \(g_{\Gamma_e}\) es el símbolo de la ruta en la linealización libre
de las genealogías y \(\mathcal U_M\) una base de la recta de masa,
queda definido el mapa de rango uno
\begin{equation}
 \mathfrak C_e:
 \mathbb R n_0\otimes\mathbb R g_{\Gamma_e}\longrightarrow\mathcal L_M,
 \qquad
 \mathfrak C_e(n_0\otimes g_{\Gamma_e})
 =\varepsilon_e^\beta\mathcal U_M.
 \label{eq:el-rango-uno}
\end{equation}
El mapa reúne la factorización anterior una vez declaradas sus bases.
El vector uniforme \(n_0\) no selecciona \(\varepsilon_e^\beta\):
esa coordenada procede de APP, del retorno orientado y de sus lectores.
Así se distingue una realización algebraica del modo nulo de una
identificación del electrón con el cuanto del reloj.

La conclusión interna es la factorización positiva de una sección
central, con sus operaciones, memoria y normalizaciones explícitas. La
asignación a un electrón físico añade la representación de carga, el
espín, el acoplamiento y la carta de unidades correspondiente; no se
deduce exclusivamente de una coincidencia decimal. El contraste
metrológico se presenta después de la prolongación excepcional y no altera
las entradas de ninguna de las construcciones aquí realizadas.

\subsubsection{Dependencia estructural de la coordenada electrónica}

La ecuación electrónica expresa una composición jerárquica. Su primer
factor, \(\sqrt3/4\), pertenece a la estructura matricial central y fija
una normalización geométrica. El exponente \(\varphi/\pi^2\) utiliza las
coordenadas regionales ya generadas en la orientación indicada del lector
areal--volumétrico. El término \(22\alpha^3\) incorpora el acoplamiento
dodecafásico en la corrección cúbica. El factor \(1+15\Delta_4\) conserva
la normalización torsional global. Finalmente, el cociente de las secciones
de acción actúa con el exponente determinado por los lectores de retorno.
Cada operación tiene así un antecedente distinto dentro de la misma
construcción.

La articulación resulta más precisa al considerar sus dependencias.
La norma central puede calcularse antes de evaluar las coordenadas
regionales. La exponencial y la corrección cúbica necesitan esas
coordenadas y el cierre de \(\alpha\). La memoria multiplicativa necesita
además las secciones de acción y la trayectoria de \(108\) posiciones.
Por ello, conocer la expresión basal no equivale a haber evaluado la
coordenada completa. La sucesión de operaciones no agrega una masa
externa al punto \((5,5)\); desarrolla los lectores atribuidos a su
sección persistente.

La distinción entre escala y estructura proporciona también una
interpretación de la comparación física. La masa es la realización
dimensional de la coordenada completa, mientras el centro, la involución
y la vacancia describen las relaciones que la producen. Un cambio común
de unidad modifica las coordenadas dimensionales, pero deja inalteradas
las razones adimensionales y las identidades de transporte demostradas.
Éste es el sentido en que el electrón conserva su carácter estructural
al expresarse respecto de una escala cronológica.


\newpage
\subsection{Transporte lorentziano y representación electromagnética}

La sección electrónica proporciona una estructura y una coordenada
antes de su aplicación a una interacción. La realización lorentziana
de la incidencia cúbica determina, con la convención temporal declarada,
el espacio de formas en el que puede expresarse esa interacción.
Para formularla se fija una realización local orientada \((\mathcal M,g)\)
de dimensión cuatro, con fibras tangentes identificadas con el modelo
\(Q_\square\). Se especifican además una conexión \(\nabla\) con
\(\nabla g=0\), una uno-forma potencial \(A\in\Omega^1(\mathcal M)\),
su curvatura \(F=dA\) y una representación de carga de la sección.
Estos datos son los argumentos de la representación física que sigue.

La dualidad de Hodge satisface \(\star^2=-I\) sobre las dos-formas.
La separación eléctrica y magnética corresponde a una elección temporal
en la descomposición \(1+3\); \(F\) y \(\star F\) conservan su
relación dentro del mismo espacio de dos-formas. La neutralidad de una
representación de carga significa \(q=0\), mientras la sección puede
conservar orientación, fase y transporte interno. La ausencia de
acoplamiento eléctrico en esa representación tiene un contenido distinto
de la ausencia de información de fase.

En esta realización dimensional, sean \(m_e>0\), \(q\) la carga,
\(\gamma(s)\) una trayectoria parametrizada y \(u=\dot\gamma\)
su cuadrivelocidad. La fuerza de Lorentz se expresa por
\begin{equation}
 m_e\nabla_u u^\mu=qF^\mu{}_{\nu}u^\nu.
 \label{eq:lorentz-electron}
\end{equation}
La antisimetría de \(F\) implica
\(F_{\mu\nu}u^\mu u^\nu=0\). Por consiguiente,
\[
 \frac{d}{ds}g(u,u)=2g(\nabla_u u,u)=0.
\]
Se conserva la normalización de la trayectoria siempre que \(\nabla\)
sea métrica y se cumpla \eqref{eq:lorentz-electron}. La masa utilizada
es la realización dimensional de la sección electrónica precedente;
el parámetro de acoplamiento se identifica en la convención electromagnética
adoptada. La deducción concreta del potencial y de una configuración
de campo constituye un problema adicional con condiciones de frontera
propias. El resultado aquí expuesto determina el dominio, los datos
de la representación y la conservación que su ecuación implica.

La inversión de rutas del TPK y la orientación de una línea electrónica
ofrecen entonces una comparación precisa: cada una tiene su involución,
su transporte y su representación. El antecedente de Wheeler y Feynman
sitúa históricamente la pregunta; la equivalencia de dos realizaciones
se decide mediante sus aplicaciones y condiciones de compatibilidad.
Esta disposición permite mantener unido el origen aritmético del electrón
con la formulación de su interacción, haciendo explícita la función
de cada etapa.


\end{HMTCompactSection}
\HMTDivision{\(V^{\natural}\)}{Incidencia y realización dimensional}{Las relaciones excepcionales, las formas y los transportes dimensionales conservan sus mapas y sus condiciones de compatibilidad.}
\HMTNarrativeHeading{La forma conservada por la segunda proyección}

La realización excepcional desarrolla la incidencia que acompaña a
la generación numérica. Cada prefijo transmite relaciones de frontera
compatibles con los truncamientos. La codificación de Paley conserva
la palabra y su imagen; el diseño de Witt organiza sus soportes; la
elevación hacia Golay introduce la incidencia binaria con sus marcas.
El recorrido posterior construye las realizaciones reticulares y
operatoriales correspondientes.

Las estrellas de incidencia permiten observar este enlace. El
registro de acoplamiento determina una cuaterna. Las hexadas que la
contienen distinguen cuatro pares complementarios y una firma
selecciona la hexada marcada. En la figura, las uniones representan
operaciones entre conjuntos. La elevación a octadas conserva esa
configuración y distingue una quinta octada transversal. Las cinco
regiones numéricas mantienen así una correspondencia con relaciones
de incidencia especificadas.

La igualdad de una lectura ternaria puede reunir regiones cuya
orientación, multiplicidad y posición siguen siendo distintas. Esos
datos intervienen en la correspondencia excepcional. La segunda
proyección conserva precisamente relaciones que la lectura de valor
trata de otro modo. La compatibilidad entre ambas se expresa sobre
el estado enriquecido común y sus mapas de refinamiento.

La teoría M y la dualidad T se presentan mediante las realizaciones
dimensionales y sus condiciones de compatibilidad. La acción del
Monstruo tiene su portador en el espacio construido para ella. Los
lectores numéricos, las configuraciones finitas, los retículos y las
representaciones reciben dominios diferenciados dentro de una misma
genealogía. Esta precisión permite seguir cómo una estructura
incidencial se transmite hasta sus consecuencias posteriores.

La parte siguiente conserva las construcciones, tablas y pruebas
que hacen efectivo el recorrido. La narración devuelve cada
realización a la cuestión rectora: qué relaciones sobreviven a la
proyección y cómo se recuperan en el objeto al que llegan. La
conservación de una forma de Gram, la covarianza de una codificación
y la invariancia de una acción tendrán pruebas propias. Su
articulación aporta el contenido geométrico de la conservación
evolutiva de la unidad.

\HMTMathematicalPaper
\section{Prolongación excepcional: de la incidencia dodecafásica a Moonshine}
\label{exc:seccion}

La expresión numérica no agota el estado que la ha generado. En la
construcción precedente, APP conserva las dos hojas y la división con
residuo y cociente; el TRIT determina el régimen local y su orientación;
el TPK compone selección, transporte, acarreo, memoria y retorno. La
lectura arquimediana contrae cilindros compatibles. La lectura que ahora
estudiamos conserva, en cambio, la incidencia de las posiciones y su
transporte entre marcos. Ambas parten del mismo estado aritmético con
memoria de trayectoria. No se
reconstruye una geometría excepcional a partir de los decimales finales
de $\pi$, $\varphi$, $e$ o $\alpha$.

Este punto permite precisar el sentido de la prolongación
$\alpha/K\longrightarrow\text{incidencia excepcional}$. La flecha no
afirma que un número real aislado determine un código, un retículo o un
grupo. Transporta el registro dodecafásico y los datos firmados que han
intervenido en la determinación de $\alpha$, junto con las palabras y los
marcos procedentes de APP--TRIT--TPK. En particular, el registro ordenado
$K$ puede recuperarse desde su constante racional $\kappa$ con las
convenciones de base, longitud y origen ya fijadas. La proyección a un
soporte conserva, en cambio, sólo las incidencias seleccionadas.
Por su parte, $K_{\rm ph}$
designa otra lectura, de memoria reducida, posterior a la holonomía
$\Gamma_9$.

La exposición separa tres operaciones matemáticas: construir los objetos
finitos y sus mapas; prolongarlos a un retículo y a un álgebra de operadores
de vértice; reconocer después los objetos obtenidos mediante los teoremas
clásicos pertinentes. Los nombres Paley, Witt, Mathieu, Golay, Leech y
Monstruo designan esos reconocimientos y las realizaciones explícitas
que los hacen verificables. No son selectores retrospectivos de las
semillas ni de los coeficientes de la generación.

El argumento comienza con la compatibilidad de prefijos y marcos,
construye el código ternario y determina la bandera mediante \(K\) y el
registro firmado de \(\alpha\). La incidencia marcada selecciona después
el origen reticular y el vecino de Leech. El último paso realiza el álgebra
de vértices y sus automorfismos. En cada etapa se precisará qué estructura
se conserva: combinaciones lineales, soportes, clases discriminantes, forma
bilineal o producto graduado. Las explicaciones de la función de \(K\),
de la normalización y del origen de norma \(54\) acompañan sus respectivas
construcciones en las secciones \ref{exc:bandera} y \ref{exc:leech}.




\subsection{Compatibilidad de las realizaciones arquimediana e incidencial}
\label{exc:doble-lectura}

Sea $X_n^{\rm cal}$ el espacio de prefijos admisibles de longitud $n$ del
estado aritmético con memoria de trayectoria, con truncaciones
$p_{n,m}:X_n^{\rm cal}\to X_m^{\rm cal}$ para $m\leq n$. El superíndice
recuerda que se conserva el marco inicial declarado y su transporte, no
que se ajuste un parámetro contra un valor experimental. Escribimos
$(w_j,f_j)$ para la palabra visible de seis trits y el marco en el paso
$j$. Los marcos satisfacen una ley causal
\[
 f_{j+1}=g_j(x)f_j,
\]
donde $g_j(x)$ depende del prefijo ya construido. En consecuencia, una
prolongación no altera los marcos anteriores.

La cadena
\[
 w_6\longrightarrow w_{12}\longrightarrow w_{18}
 \longrightarrow w_{24}\longrightarrow w_{30}
 \longrightarrow R_{36}\longrightarrow G_9
\]
se conserva en esta lectura. El retorno de la fase no reinicia ni el
marco ni la memoria del estado. Asimismo, se mantienen distintos el
censo de $468$ emisiones decimales, las $243$ palabras visibles de seis
trits y el ambiente $\mathbb F_3^6$, de $729$ palabras. El ambiente se
utilizará para construir un código lineal; no sustituye al dominio
admisible de emisiones.

La lectura arquimediana toma, para un bloque $b$ de seis trits, su
dirección entera $\nu(b)\in\{0,\ldots,728\}$. Un prefijo determina
\[
 q_n=m+\sum_{j=0}^{n-1}\frac{\nu(b_j)}{729^{j+1}},
 \qquad I_n=[q_n,q_n+729^{-n}].
\]
La compatibilidad de prefijos da cilindros encajados cuyo diámetro tiende
a cero. Su clase de Cauchy pertenece primero a la completación interna
$\mathbb R_{\rm HMT}$; su identificación posterior con la recta real
usual no interviene en la elección de los bloques.

Para la lectura incidencial introducimos, sobre el mismo prefijo, el
mapa de gráfico
\begin{equation}
 \gamma_A:\mathbb F_3^6\longrightarrow\mathbb F_3^{12},
 \qquad w\longmapsto(w,wA).
 \label{exc:gamma}
\end{equation}
Las palabras se escriben como vectores fila. Si $A^f=fAf^{-1}$, el
cambio de marco satisface la identidad concreta
\begin{equation}
 \gamma_{A^f}(wf^{-1})
   =\gamma_A(w)(f^{-1}\oplus f^{-1}).
 \label{exc:covariancia}
\end{equation}
El codominio calibrado conserva $f$ junto a la palabra codificada.
Definimos entonces
\begin{equation}
 Q_n(x)=\bigl((f_j,\gamma_{f_jA_Wf_j^{-1}}(w_j))\bigr)_{0\leq j<n}.
 \label{exc:qn}
\end{equation}

\begin{proposition}[Compatibilidad y prolongación de la lectura incidencial]
\label{exc:limite-inverso}
Si $r_{n,m}$ trunca las sucesiones codificadas, se cumple
\[
 r_{n,m}Q_n=Q_m p_{n,m}.
\]
Por tanto, los $Q_n$ determinan un único mapa $Q_\infty$ entre los
respectivos límites inversos. En cada marco, el mapa conserva exactamente
la palabra visible, pues $\operatorname{pr}_1\gamma_A=\operatorname{id}$.
\end{proposition}

\begin{proof}
Dos prefijos compatibles tienen las mismas palabras y los mismos
operadores de transporte hasta el índice $m-1$. Tienen, por inducción,
los mismos marcos en esos índices. Las primeras $m$ componentes de
\eqref{exc:qn} coinciden, lo que prueba el cuadrado. La propiedad
universal del límite inverso da $Q_\infty$. La afirmación de recuperación
es la proyección sobre las seis primeras coordenadas de cada gráfico.
\end{proof}

Esta inyectividad local tiene un dominio preciso. No demuestra que
$Q_\infty$ recupere toda la memoria profunda del estado TPK si sólo se
ha retenido la sucesión de palabras visibles y marcos. Mucho menos
demuestra inyectividad después de pasar de palabras a soportes.
Además, un cambio general de base en $\operatorname{GL}_6(\mathbb F_3)$
no conserva por sí mismo el peso de Hamming ordinario. En una carta
general se transportan también la incidencia y la métrica de la carta
inicial. Los cálculos de soportes que siguen se realizan en la carta
Paley declarada; los relabelados monomiales conservan directamente su
interpretación coordenada.

\subsection{La realización finita y el código ternario}
\label{exc:paley-codigo}

El operador elíptico del TRIT admite la siguiente realización finita en
la carta de seis posiciones. Éstas coordinatizan las seis unidades
$\{1,2,4,5,7,8\}$ módulo nueve, conservadas por el transporte modular.
Sobre $\mathbb P^1(\mathbb F_5)$, una vez
declarado el orden $\infty,0,1,2,3,4$, sea $C$ la matriz con diagonal
nula, $C_{\infty,a}=C_{a,\infty}=1$ y
$C_{a,b}=\chi(b-a)$ en las cinco posiciones finitas. Aquí
$\chi(1)=\chi(4)=1$, $\chi(2)=\chi(3)=-1$ y $\chi(0)=0$.
Es una realización en coordenadas de la estructura cuadrática ya
transportada; la identificación de las seis posiciones con la recta
proyectiva es un dato de carta, no una nueva entrada numérica en el
generador.

Las identidades de la suma cuadrática dan $C^2=5I_6$. Al reducir módulo
tres obtenemos
\begin{equation}
 A_W=
 \begin{pmatrix}
 0&1&1&1&1&1\\
 1&0&1&2&2&1\\
 1&1&0&1&2&2\\
 1&2&1&0&1&2\\
 1&2&2&1&0&1\\
 1&1&2&2&1&0
 \end{pmatrix},
 \qquad A_W^2=2I_6=-I_6.
 \label{exc:aw}
\end{equation}
La igualdad cuadrática realiza el régimen elíptico; no identifica el
cuerpo $\mathbb F_3$ con una realización arquimediana de la unidad
imaginaria. Son realizaciones distintas de una relación operatoria
común, con sus correspondientes dominios.

\Needspace{8\baselineskip}
\begin{definition}[Código de la carta dodecafásica]
\label{exc:codigo}
El código ternario asociado a \eqref{exc:aw} es
\[
 C_W=\{(w,wA_W):w\in\mathbb F_3^6\}
       \subset\mathbb F_3^{12}.
\]
Su conjunto de hexadas es
\[
 \mathcal H=
 \{\operatorname{supp}(c):c\in C_W,
                           \operatorname{wt}(c)=6\}.
\]
\end{definition}

\begin{proposition}[Autodualidad, distancia y diseño de Witt]
\label{exc:codigo-witt}
El código $C_W$ es autodual de parámetros $[12,6,6]_3$. Su enumerador
de pesos es
\begin{equation}
 \sum_{c\in C_W}z^{\operatorname{wt}(c)}
   =1+264z^6+440z^9+24z^{12}.
 \label{exc:enumerador}
\end{equation}
El conjunto $\mathcal H$ tiene $132$ elementos y constituye un sistema
$S(5,6,12)$: cada conjunto de cinco posiciones está contenido en una
única hexada.
\end{proposition}

\begin{proof}
La matriz generadora $[I_6\mid A_W]$ tiene rango seis y, por simetría
de $A_W$, satisface
\[
 [I_6\mid A_W][I_6\mid A_W]^{\mathsf T}
   =I_6+A_W^2=0.
\]
El código es autoortogonal de dimensión mitad de la longitud y, por
tanto, autodual. La evaluación finita de las $3^6$ palabras da
\eqref{exc:enumerador}; es una enumeración sobre la matriz explícita,
sin valores reales ni objetivos externos. En particular, la distancia
mínima es seis.

Si dos palabras de peso seis tienen el mismo soporte, entre sus seis
cocientes de coordenadas no nulas uno de los signos $1,-1$ aparece al
menos tres veces. Restando el múltiplo correspondiente se obtendría
una palabra no nula de peso a lo sumo tres, imposible. Luego un soporte
corresponde exactamente a la pareja $\{c,-c\}$ y hay $264/2=132$
soportes. Si dos soportes distintos compartieran cinco posiciones,
el mismo argumento cancelaría al menos tres coordenadas de una unión
de tamaño siete, produciendo peso a lo sumo cuatro. Por ello una
quintupla pertenece a lo sumo a una hexada. Finalmente,
\[
 132\binom65=792=\binom{12}5,
\]
de modo que pertenece a una y sólo una.
\end{proof}

El reconocimiento posterior es el código de Golay ternario extendido y
el pequeño diseño de Witt \cite{conwaysloane}. La prueba anterior explica qué objeto se
reconoce. No reemplaza la genealogía APP--TRIT--TPK por una búsqueda
entre códigos conocidos.

\subsection{Bandera dodecafásica asociada a \texorpdfstring{$K$ y $\alpha$}{K y alfa}}
\label{exc:bandera}

\paragraph{Articulación de las coordenadas regionales y la incidencia de clausura.}
\label{inc-ampl-articulacion}
La función de esta prolongación se comprende a partir de una distinción
entre tres operaciones: conservar un registro, evaluar una coordenada y
extraer una relación de incidencia. Las tres actúan sobre información
producida por APP--TRIT--TPK. Su articulación hace explícita la tesis
estructural del artículo: una realización numérica posee una procedencia
operatoria, y determinados datos de esa procedencia admiten, además, una
realización combinatoria y reticular. Las ecuaciones permiten seguir ambos
recorridos con sus dominios y con la información que permanece en cada uno.

El punto de partida es el estado terminal con sus prolongaciones
compatibles. En él se conservan los residuos y cocientes de las dos hojas
APP, el régimen y la orientación TRIT, y el transporte TPK con su memoria.
Las regiones de clausura, propagación y autoescala se determinan en ese
dominio antes de que sus lectores establezcan las coordenadas
$\pi,e,\varphi$. Al mismo tiempo, el registro orientado de las doce
ventanas determina el vector firmado $U$. La aparición posterior de $K$
procede de la transformación integral demostrada en
\S\ref{subsec:k-hadamard}:
\begin{equation}
 U=\mathcal H_{12}K,\qquad
 K=\frac14\mathcal H_{12}U,\qquad
 \mathcal H_{12}^{2}=4I_{12}.
 \label{inc-ampl-registro}
\end{equation}
La congruencia que define $\mathcal L_H$ garantiza la integralidad de
la inversa. De este modo, $K$ constituye una coordenatización reversible
del registro firmado en la subred pertinente. Sus posiciones conservan el
orden de fase; sus diferencias $D_3K,D_4K$ conservan las relaciones
transversales; la suma $Q(K)$ determina la componente uniforme. Esa
estructura explica por qué el registro interviene tanto en la composición
numérica como en la selección de una configuración de frontera.

La evaluación periódica
\[
 K\longmapsto\kappa_{\mathrm{per}}
 =\frac{\sum_{m=1}^{12}K_m1000^{12-m}}{1000^{12}-1}
\]
mantiene recuperables los doce bloques cuando se fijan base, longitud y
origen de fase. La extracción de un soporte realiza una operación
diferente: conserva la pertenencia de cada posición a una clase definida
por un umbral y agrupa los registros que tienen el mismo soporte. La
distinción entre estas dos aplicaciones resulta decisiva. La precisión
escalar de $\kappa_{\mathrm{per}}$ y la incidencia de $B_K$ son dos
contenidos matemáticos específicos, relacionados mediante el registro
completo del que proceden.



En el marco heredado del registro dodecafásico, los datos son
\begin{align}
 K={}&(234,543,140,729,659,824,621,\mathtt{058},914,794,146,601),
 \label{exc:k}\\
 w_\pi={}&010211, &w_e={}&201101,\nonumber\\
 w_\varphi={}&121200, &w_{\rm APP}={}&022020.\nonumber
\end{align}
Al aplicar $\gamma_{A_W}$ se obtienen las siguientes incidencias:
\begin{align*}
 P_\pi&=\{2,4,5,6,7,8,9,10,11\},\\
 H_e&=\{1,3,4,6,9,10\},\\
 H_\varphi&=\{1,2,3,4,7,9\},\\
 H_{\rm APP}&=\{2,3,5,10,11,12\}.
\end{align*}
La primera es el soporte de una palabra de peso nueve; las otras tres
son hexadas. El registro aporta, por su lectura de bloque alto,
\begin{equation}
 B_K=\{i:K_i\geq729\}=\{4,6,9,10\}=H_e\cap P_\pi.
 \label{exc:bk}
\end{equation}
El entero $729$ señala aquí una frontera del formato de bloques
declarado. No se identifica por el mero cardinal con el ambiente
ternario usado en \eqref{exc:codigo}. En este registro concreto, todos
los umbrales enteros entre $660$ y $729$ producen el mismo soporte:
el soporte no recupera el umbral ni los doce bloques.

La vía aritmética de $\alpha$ conserva, antes del acarreo, el vector
firmado
\begin{equation}
 Z_0=(7,297,353,-430,-715,-1199,-3,286,-894,-528,380,663).
 \label{exc:precarry}
\end{equation}
Su soporte negativo es
\begin{equation}
 H^-_\alpha=\{4,5,6,7,9,10\}.
 \label{exc:halpha}
\end{equation}
No se extraen estos signos de la expansión decimal del escalar
$\alpha$; se conservan en el estado que precede a su expresión numérica.

\paragraph{La normalización de alfa y su configuración firmada.}
\label{inc-ampl-clausura}
La coordenada de $\alpha$ se determina en la vía posicional mediante la
composición de los bloques regionales y del registro dodecafásico. Antes
de la división posicional, la combinación orientada tiene componentes
$Z_j=P_j+E_j-\Phi_j-K_j^{\mathrm{per}}$ y el balance con cociente entrante
es $s_j=Z_j+c_{j+1}$. La normalización satisface
\begin{equation}
 A_j=s_j-1000c_j=Z_j+c_{j+1}-1000c_j,\qquad 0\leq A_j<1000.
 \label{inc-ampl-normalizacion}
\end{equation}
Esta identidad conserva simultáneamente el bloque normalizado y la
contribución del cociente transportado. El límite de los prefijos
compatibles determina la coordenada escalar de la clausura. La
configuración firmada conserva, por su parte, qué posiciones de la
combinación anterior a la normalización presentan defecto negativo.
En la ventana inicial, esa configuración es el vector $Z_0$ de
\eqref{exc:precarry}, y su soporte negativo es $H^-_\alpha$.

Así quedan definidos dos contenidos de una misma operación. La
normalización posicional determina una coordenada real mediante sus
restos y cocientes compatibles; el soporte negativo determina una
configuración de posiciones en el ciclo de doce componentes. El signo se
refiere al balance de registros anterior a la normalización. Su lectura
incidencial permanece disponible porque el estado conserva ese balance
junto con el resultado normalizado. Esta articulación corresponde a la
vía posicional demostrada; la comparación con los operadores analíticos
de vacancias conserva el alcance preciso establecido en el capítulo de
$\alpha$.


\begin{proposition}[Coincidencia de los dos selectores de la bandera]
\label{exc:selector-bandera}
Dentro de $\mathcal H$, el soporte \eqref{exc:halpha} es la única hexada
$H$ que satisface
\[
 B_K\subset H\subset P_\pi,
 \qquad
 (|H\cap H_e|,|H\cap H_\varphi|,|H\cap H_{\rm APP}|)=(4,3,2).
\]
Así, la lectura firmada y la lectura incidencial seleccionan la misma
bandera $B_K\subset H^-_\alpha$.
\end{proposition}

\begin{proof}
Las cuatro hexadas que contienen $B_K$ se obtienen añadiendo,
respectivamente, los pares
\[
 \{1,3\},\qquad\{2,11\},\qquad\{5,7\},\qquad\{8,12\}.
\]
Sólo los pares segundo y tercero están contenidos en $P_\pi$. Para el
segundo, la intersección con $H_\varphi$ tiene tamaño tres y con
$H_{\rm APP}$ tamaño tres; no cumple la última condición. Para el
tercero, las tres intersecciones tienen tamaños $4,3,2$. Este último
produce precisamente el soporte negativo de \eqref{exc:precarry}.
\end{proof}

Conviene conservar tanto la bandera como su marco marcado:
\[
 \mathcal I_*=(B_K\subset H^-_\alpha),\qquad
 \widetilde{\mathcal I}_*
   =(\mathcal I_*;P_\pi,H_e,H_\varphi,H_{\rm APP}).
\]
Las operaciones de unión, intersección y diferencia son equivariantes
bajo un relabelado simultáneo del diseño. La clase marcada, y no los
números de las posiciones tomados aisladamente, es el objeto transportado.

\begin{lemma}[Origen seleccionado por la incidencia marcada]
\label{exc:origen}
Para el marco anterior,
\[
 o_*=(H_e\cap H_\varphi)
           \setminus(H^-_\alpha\cup H_{\rm APP})=\{1\}.
\]
La regla es equivariante bajo todo automorfismo aplicado simultáneamente
a los datos de $\widetilde{\mathcal I}_*$.
\end{lemma}

\begin{proof}
La primera intersección es $\{1,3,4,9\}$. La unión que se sustrae
contiene $3,4,9$ y no contiene $1$. La equivariancia se sigue de que una
biyección conmuta con las operaciones de conjuntos.
\end{proof}

Este origen seleccionará una dirección del vecino reticular cuando se
declaren la métrica $A_2$ y la cámara positiva. La incidencia por sí sola
no contiene esa métrica. El paso reticular conservará explícitamente
ambos datos.

\paragraph{De la codificación ternaria a la bandera marcada.}
\label{inc-ampl-bandera}
Las regiones aportan también los códigos $w_\pi,w_e,w_\varphi$ y
$w_{\rm APP}$ en el marco heredado. La aplicación
\[
 \gamma_{A_W}:w\longmapsto(w,wA_W)
\]
conserva la palabra de seis componentes como primera coordenada y añade
su transformación por $A_W$. La relación $A_W^2=-I_6$ fija la
compatibilidad cuadrática de esta realización. El paso siguiente,
$c\mapsto\operatorname{supp}(c)$, conserva las posiciones no nulas.
En las palabras de peso seis, el soporte identifica exactamente la pareja
$\{c,-c\}$, como prueba la distancia mínima del código. De esta forma,
el cociente por signo produce las hexadas, mientras el código completo
sigue disponible para las operaciones que requieren sus símbolos.

En esta carta, las dos construcciones de la cuaterna coinciden:
\begin{equation}
 \{i:K_i\geq729\}
 =H_e\cap P_\pi=B_K=\{4,6,9,10\}.
 \label{inc-ampl-cuaterna}
\end{equation}
El lado izquierdo procede de una comparación entera de los componentes
de $K$; el derecho, de la intersección de los soportes codificados de
propagación y clausura. El umbral pertenece a la carta de bloques
declarada y permanece explicitado en la definición del lector. El
registro íntegro conserva los valores sobre los que actúa, incluso cuando
distintos umbrales producen el mismo subconjunto. La igualdad expresa,
por tanto, una compatibilidad entre aplicaciones especificadas.

El soporte $H^-_\alpha$ completa esa cuaterna hasta una hexada. La
condición de pertenencia al soporte $P_\pi$, junto con los tamaños de
intersección $(4,3,2)$ respecto de $H_e,H_\varphi,H_{\rm APP}$,
selecciona una única hexada. Las pruebas siguientes reúnen las dos
determinaciones: la obtenida mediante el signo de $Z_0$ y la obtenida
mediante incidencia. El resultado se organiza como
\begin{equation}
 \widetilde{\mathcal I}_*
 =\bigl(B_K\subset H^-_\alpha;
         P_\pi,H_e,H_\varphi,H_{\rm APP}\bigr).
 \label{inc-ampl-marco}
\end{equation}
La inclusión es la bandera; los cuatro soportes adicionales constituyen
su marco marcado. Las permutaciones simultáneas transportan el marco y
conservan sus inclusiones e intersecciones. La información relevante reside
en esta configuración transportada y en su clase de isomorfía.

Las cinco regiones de $\pi$ recuperan aquí su diferenciación interna.
Comparten $w_\pi$ como código ternario visible, pero conservan emisiones,
firmas y multiplicidades distintas. En la realización binaria sobre una
dodecada marcada, la región transversal corresponde a la octada cuya
intersección con la dodecada es la cuaterna; las otras cuatro regiones
corresponden a las elevaciones de las cuatro hexadas de su estrella. La
región negativa queda distinguida por $H^-_\alpha$ y el orden de carta
individualiza las tres regiones positivas. Esta correspondencia conserva
los papeles regionales y la distribución $432+4\cdot144=1008$, además
de realizar la partición incidencial $1+4$.


\subsection{Entrelazamiento isométrico de las representaciones de incidencia}
\label{exc:entrelazador}

La incidencia define un mapa lineal cuya acción puede comprobarse
posición por posición:
\begin{equation}
 I:\mathbb R^{12}\longrightarrow\mathbb R^{\mathcal H},
 \qquad (Iv)(H)=\sum_{p\in H}v_p.
 \label{exc:incidencia}
\end{equation}
Sea $V_{11}=\mathbf1^\perp\subset\mathbb R^{12}$. Las dos realizaciones
euclídeas son aquí lenguajes posteriores para el mismo diseño finito.

\begin{proposition}[Isometría de la pantalla dodecafásica]
\label{exc:isometria}
Se cumple
\[
 I^{\mathsf T}I=36I_{12}+30J_{12},
\]
donde $J_{12}$ es la matriz cuyas entradas son uno. En consecuencia,
\[
 U_W=\frac16 I\big|_{V_{11}}:
 V_{11}\xrightarrow{\ \sim\ }E_{\rm inc}:=I(V_{11})
\]
es una isometría equivariante para las acciones por permutación de los
automorfismos del diseño.
\end{proposition}

\begin{proof}
Cada punto pertenece a
$r=\binom{11}{4}/\binom54=66$ hexadas, y cada pareja a
$\lambda_2=\binom{10}{3}/\binom43=30$. Son, respectivamente, las
entradas diagonales y no diagonales de $I^{\mathsf T}I$. Sobre
$\mathbf1^\perp$, el término $J_{12}$ se anula y queda $36I$.
Por último, si $g$ es un automorfismo, $p\in H$ equivale a
$g(p)\in g(H)$; por ello $I\rho_{12}(g)=\rho_{\mathcal H}(g)I$.
\end{proof}

La información de incidencia también tiene una lectura espectral
concreta. Definamos
\[
 G_{H,H'}=\binom{|H\cap H'|}{4}.
\]

\Needspace{8\baselineskip}
\begin{lemma}[Acción del núcleo de intersección]
\label{exc:espectro}
Si $J_{132,12}$ es la matriz rectangular de unos,
\[
 GI=30I+15J_{132,12}.
\]
Por tanto, $G$ actúa por el escalar $30$ en $E_{\rm inc}$.
\end{lemma}

\begin{proof}
La entrada $(H,p)$ de $GI$ cuenta parejas $(T,H')$ con
$T\subset H\cap H'$, $|T|=4$ y $p\in H'$. Si $p\in H$, hay diez
cuaternas $T$ que contienen $p$, cada una contenida en cuatro hexadas,
y cinco que no lo contienen, cada una con una única extensión de la
quintupla $T\cup\{p\}$. El total es $10\cdot4+5=45$.
Si $p\notin H$, las quince cuaternas de $H$ dan una extensión cada
una: el total es $15$. De ahí la identidad. El término rectangular
se anula en $V_{11}$.
\end{proof}

Aquí sí existe un entrelazador material: tiene dominio, codominio,
entradas y normalización determinados. Por ejemplo, cualquier
proyector ortogonal $P$ ya construido sobre $V_{11}$ se transporta a
$Q=U_WPU_W^*$ sobre $E_{\rm inc}$, con
$U_WP=QU_W$. Esta afirmación no identifica por su sola traza a todos
los operadores que actúan en esos espacios.

\subsection{La dirección seleccionada por el registro dodecafásico}
\label{exc:k-direccion}

El registro \(K\) tiene una función geométrica que complementa su
evaluación racional y su lectura de soportes. Una vez construida la
representación de las doce posiciones, sus componentes determinan una
dirección en el espacio de variaciones de media nula. La reducción
\(12\to11\) elimina el modo uniforme; la reducción siguiente
\(11\to10\) elimina la dirección que selecciona el propio registro.
Las dos operaciones tienen, por tanto, antecedentes diferentes.

Para declarar completamente la carta utilizada, consideremos el grupo
\(A_5\) de permutaciones pares de cinco objetos y el subgrupo
\(C_5=\langle(1\,2\,3\,4\,5)\rangle\). Las doce clases izquierdas
\(r_pC_5\) se ordenan mediante estos representantes, escritos como
listas de imágenes:
\[
\begin{array}{c|c@{\qquad}c|c}
p&r_p&p&r_p\\ \hline
1&(4,3,5,2,1)&7&(5,4,3,2,1)\\
2&(4,2,1,3,5)&8&(1,3,4,2,5)\\
3&(1,2,4,5,3)&9&(4,2,3,5,1)\\
4&(2,5,3,4,1)&10&(3,2,5,4,1)\\
5&(4,3,1,5,2)&11&(4,5,2,3,1)\\
6&(5,1,2,3,4)&12&(5,3,2,4,1).
\end{array}
\]
La acción por multiplicación izquierda define una representación
ortogonal \(\rho:A_5\to O(\RR^{12})\):
\(\rho(a)e_p=e_q\) cuando \(ar_pC_5=r_qC_5\).
Esta carta de representación actúa sobre el registro ya obtenido;
los caracteres del grupo no se utilizan para elegir sus doce valores.

Sean \(\chi_3\) los valores del carácter tridimensional
\[
\begin{array}{c|ccccc}
\text{clase}&1&2^2&3&5_a&5_b\\ \hline
\chi_3&3&-1&0&(1+\sqrt5)/2&(1-\sqrt5)/2.
\end{array}
\]
La clase \(5_a\) contiene el generador de \(C_5\) mostrado y
\(5_b\) contiene su cuadrado. Esta convención distingue los dos
sectores tridimensionales conjugados de la representación. Definimos
\begin{equation}
 P_3=\frac3{60}\sum_{a\in A_5}\chi_3(a^{-1})\rho(a),
 \qquad P_{11}=I_{12}-\frac1{12}J_{12}.
 \label{eq:k-proyectores-incidencia}
\end{equation}
Ambas matrices quedan determinadas por sumas finitas en
\(\QQ(\sqrt5)\). La aparición de este cuerpo cuadrático pertenece
a la realización de la representación, posterior a la generación
regional de \(\varphi\).

\begin{lemma}[Proyección del registro sobre el sector tridimensional]
\label{lem:k-sector-no-nulo}
Se tienen
\[
 P_3=P_3^{\mathsf T}=P_3^2,\quad
 P_3\mathbf1=0,\quad \operatorname{tr}P_3=3,
 \qquad P_3P_{11}=P_{11}P_3=P_3.
\]
Para el registro \eqref{exc:k}, el vector
\begin{equation}
 u_K=P_3P_{11}K
 \label{eq:k-vector-direccional}
\end{equation}
es no nulo. Su norma exacta es
\begin{equation}
 \lVert u_K\rVert^2
 =\frac{6638585+2275584\sqrt5}{20}>0.
 \label{eq:k-norma-direccion}
\end{equation}
\end{lemma}
\begin{proof}
La suma de carácter en \eqref{eq:k-proyectores-incidencia} es el
proyector isótipo. El carácter real y la sustitución \(a\mapsto a^{-1}\)
dan la simetría. La multiplicidad del sector en la acción sobre
\(A_5/C_5\) es la dimensión de sus vectores fijos por \(C_5\),
\[
 \frac15\left(3+2\frac{1+\sqrt5}{2}
                 +2\frac{1-\sqrt5}{2}\right)=1.
\]
Por ello el rango es tres. Su ortogonalidad al carácter trivial
da \(P_3\mathbf1=0\) y las dos composiciones con \(P_{11}\).
También pueden verificarse estas identidades multiplicando las
matrices de la suma finita.

Finalmente, \(\sum K_p=6263\) y la sustitución de las doce
componentes en esa misma suma da
\[
 \lVert u_K\rVert^2=K^{\mathsf T}P_3K
 =\frac1{20}\sum_{a\in A_5}
                \chi_3(a^{-1})K^{\mathsf T}\rho(a)K
 =\frac{6638585+2275584\sqrt5}{20}.
\]
Esta es una evaluación exacta de sesenta términos especificados,
reproducida por el certificado adjunto con aritmética racional
cuadrática. Los dos coeficientes positivos prueban la no nulidad.
\end{proof}

\begin{definition}[Dirección dodecafásica y complemento transversal]
La dirección seleccionada por \(K\) en la carta declarada es
\(\ell_K=\RR u_K\). Su complemento transversal dentro de
\(V_{11}=\mathbf1^\perp\) es
\[
 V_{10}(K)=V_{11}\cap u_K^\perp.
\]
\end{definition}

\begin{proposition}[Reducción ortogonal determinada por \(K\)]
\label{prop:k-reduccion-diez}
El operador
\begin{equation}
 P_{10}(K)=P_{11}
       -\frac{u_Ku_K^{\mathsf T}}{\lVert u_K\rVert^2}
 \label{eq:k-proyector-diez}
\end{equation}
es el proyector ortogonal de rango diez sobre \(V_{10}(K)\). Se cumple
\[
 V_{11}=\ell_K\mathbin{\oplus^{\perp}}V_{10}(K),
 \qquad P_{10}P_{11}=P_{10},\quad P_{10}u_K=0.
\]
\end{proposition}
\begin{proof}
Sea \(E_K=u_Ku_K^{\mathsf T}/\lVert u_K\rVert^2\). La
no nulidad demostrada garantiza que está definido. Es simétrico,
\(E_K^2=E_K\), tiene rango uno y, como \(u_K\in V_{11}\),
\(P_{11}E_K=E_KP_{11}=E_K\). De aquí
\[
 (P_{11}-E_K)^2=P_{11}-2E_K+E_K=P_{11}-E_K.
\]
Su imagen es \(V_{11}\cap u_K^\perp\) y su traza es
\(11-1=10\). La descomposición y las identidades restantes se
obtienen por proyección sobre los dos sumandos ortogonales.
\end{proof}

Este resultado precisa una función de \(K\) que su denominación como
registro no agotaba: además de representar reversiblemente datos
orientados y de seleccionar la tétrada \(B_K\), determina una
polarización lineal. La dirección depende del registro completo a
través de \(P_3K\), no únicamente de los cuatro puntos de su soporte.
Se ha mantenido así el orden registro, representación, dirección y
reducción dimensional (proposición~\ref{prop:k-reduccion-diez}).

La transformación tiene una ley de covariancia precisa. Si se cambia
simultáneamente la carta por una matriz ortogonal de permutación \(S\),
\(K'=SK\) y \(P'_3=SP_3S^{-1}\), entonces
\[
 u_{K'}'=Su_K,\qquad P'_{10}(K')=SP_{10}(K)S^{-1}.
\]
Por otra parte, sustituir \(u_K\) por \(a u_K\), con \(a\ne0\),
no cambia \(E_K\). La dirección retiene una polarización y no
permite recuperar por sí sola ni la amplitud de \(u_K\) ni el
registro íntegro. La reversibilidad de \(\kappa\leftrightarrow K\)
permanece en su dominio propio.

La realización geométrica aquí establecida es una reducción de espacios
con producto interno. Su empleo posterior como dirección compacta
requiere transportar esta recta a la realización física correspondiente.
El apartado \ref{exc:pantallas-dualidad} expone el mapa algebraico
que mantiene coordinadas esta reducción y la pantalla reticular.


\subsection{La estrella de Witt y la generación de \texorpdfstring{$M_{12}$}{M12}}
\label{exc:estrella}

\begin{definition}[Estrella sobre una cuaterna]
\label{exc:def-estrella}
Para $B\in\binom{[12]}4$, la estrella de Witt es la familia de las
cuatro hexadas que contienen $B$. Sus pétalos son los pares $H\setminus B$.
La permutación $\sigma_B$ fija $B$ e intercambia los dos puntos de cada
pétalo.
\end{definition}

La existencia de cuatro hexadas se deduce de
\[
 \lambda_4=\frac{\binom81}{\binom21}=4.
\]
Dos pétalos no pueden compartir un punto, pues entonces dos hexadas
contendrían la misma quintupla. Son, por tanto, cuatro pares disjuntos
que particionan $[12]\setminus B$. Esta descripción no es una estrella
de cinco puntas ni una metáfora geométrica: es una incidencia
cuatro-a-uno del diseño.

\begin{figure}[htbp]
\centering
\begin{tikzpicture}[x=1cm,y=1cm, every node/.style={font=\small}]
 \node[draw,rounded corners,minimum width=3.8cm,minimum height=.7cm]
  (b) at (0,0) {$B_K=\{4,6,9,10\}$};
 \node[draw,minimum width=2.7cm,minimum height=.65cm]
  (p1) at (-4.6,-1.7) {$\{1,3\}$};
 \node[draw,minimum width=2.7cm,minimum height=.65cm]
  (p2) at (-1.55,-1.7) {$\{2,11\}$};
 \node[draw,very thick,minimum width=2.7cm,minimum height=.65cm]
  (p3) at (1.55,-1.7) {$\{5,7\}$};
 \node[draw,minimum width=2.7cm,minimum height=.65cm]
  (p4) at (4.6,-1.7) {$\{8,12\}$};
 \draw (b) -- (p1); \draw (b) -- (p2);
 \draw[very thick] (b) -- (p3); \draw (b) -- (p4);
 \node at (-4.6,-2.35) {$H_e$};
 \node at (-1.55,-2.35) {$B_K\cup\{2,11\}$};
 \node at (1.55,-2.35) {$H^-_\alpha$};
 \node at (4.6,-2.35) {$B_K\cup\{8,12\}$};
 \node at (0,-3.05)
  {$\sigma_{B_K}=(1\ 3)(2\ 11)(5\ 7)(8\ 12)$};
\end{tikzpicture}
\caption{La estrella concreta de la bandera dodecafásica. Cada línea
indica la unión de la cuaterna central con el par de su extremo; la
línea gruesa señala la hexada seleccionada también por los signos
anteriores al acarreo de $\alpha$. Los cuatro pares agotan las ocho
posiciones exteriores.}
\label{exc:fig-estrella}
\end{figure}

\begin{proposition}[La familia de estrellas]
\label{exc:m12}
Las $\binom{12}4=495$ permutaciones $\sigma_B$ preservan $\mathcal H$.
El grupo que generan tiene orden $95040$ y es el grupo de Mathieu
$M_{12}$ en su acción sobre el pequeño diseño de Witt.
\end{proposition}

\begin{proof}
La afirmación de preservación es una comprobación finita sobre las
$132$ hexadas de \eqref{exc:codigo}: para cada cuaterna se construyen
los cuatro pétalos y se verifica
$\{\sigma_B(H):H\in\mathcal H\}=\mathcal H$.
La clausura por composición de estas permutaciones produce $95040$
elementos. En la enumeración lexicográfica basta añadir sucesivamente
seis estrellas no contenidas aún en el grupo; los órdenes intermedios
son $2,6,18,360,7920,95040$, y las $495$ estrellas pertenecen al grupo
final. El procedimiento y los seis generadores están explicitados en
la nota técnica de procedencia. Finalmente se aplica el reconocimiento
del grupo de automorfismos de $S(5,6,12)$ como $M_{12}$, cuyo orden es
$12\cdot11\cdot10\cdot9\cdot8=95040$.
\end{proof}

Una estrella individual genera un grupo cíclico de orden dos. No
genera por sí sola $M_{12}$, y menos aún el Monstruo. Tampoco la
existencia de $495$ generadores de un grupo de permutaciones prueba que
cada uno sea el transporte de un lazo específico del TPK. Esa afirmación
requiere el mapa de lazos correspondiente; el resultado aquí establecido
es la acción incidencial finita.

\begin{remark}[La igualdad de trazas no identifica operadores]
Sobre $V_{11}$, una estrella tiene espectro $+1$ con multiplicidad siete
y $-1$ con multiplicidad cuatro. Su traza normalizada es $3/11$.
Un proyector de rango tres en el mismo espacio tiene también traza
normalizada $3/11$, pero espectro $1^3,0^8$. No son conjugados: sus
polinomios mínimos y sus espectros son diferentes. El entrelazador
\eqref{exc:incidencia} transporta operadores definidos; no convierte
esta coincidencia escalar en una equivalencia de operadores.
\end{remark}

\subsection{El residuo binario y las cinco regiones de \texorpdfstring{$\pi$}{pi}}
\label{exc:cinco-regiones}

El siguiente reconocimiento utiliza el código de Golay binario
extendido $G_{24}$, de parámetros $[24,12,8]_2$. Sus pesos son
$0,8,12,16,24$, y los soportes de peso ocho forman el diseño
$S(5,8,24)$. Fijamos una dodecada $D$, soporte de una palabra de peso
doce. Este dato y el isomorfismo de incidencia que se declarará a
continuación pertenecen a la carta de realización; no se usan para
generar ni seleccionar los decimales de $\pi$.

\begin{proposition}[Residuo de una dodecada y elevación de hexadas]
\label{exc:residuo-binario}
El conjunto
\[
 \mathcal H(D)=
 \{O\cap D: O\text{ es una octada},\ |O\cap D|=6\}
\]
es un sistema $S(5,6,12)$ sobre $D$. Cada una de sus hexadas tiene
una única elevación a una octada. Para cualquier cuaterna $B\subset D$,
las cinco octadas que la contienen se distribuyen en cuatro elevaciones
residuales y una única octada $O_B^\perp$ con $O_B^\perp\cap D=B$.
\end{proposition}

\begin{proof}
Una quintupla $A\subset D$ pertenece a una única octada $O$. Si
$j=|O\cap D|$, la palabra binaria suma de $O$ y $D$ tiene peso
$20-2j$. Entre los pesos permitidos, y con $0\leq j\leq8$, sólo
son posibles $j=2,4,6$. Como $A\subset O\cap D$, se tiene $j\geq5$
y por tanto $j=6$. Esto da existencia y unicidad de la hexada
residual, así como unicidad de su elevación al escoger cinco de sus
puntos. Por otra parte,
\[
 \lambda_4(S(5,8,24))=\frac{\binom{20}1}{\binom41}=5,
 \qquad
 \lambda_4(S(5,6,12))=4.
\]
Las cuatro hexadas residuales sobre $B$ dan cuatro octadas distintas.
La quinta no puede cortar $D$ en seis puntos, pues sería otra hexada
residual; como contiene $B$, su intersección tiene tamaño cuatro y
coincide con $B$.
\end{proof}

Una carta $\ell:[12]\to D$ que preserve las hexadas identifica
$\mathcal H$ con $\mathcal H(D)$. Su existencia es el reconocimiento
del diseño ya construido como el pequeño diseño de Witt; una elección
concreta de $\ell$ fija el relabelado. La elevación
\[
 \iota_D:\mathcal H\longrightarrow
       \{\text{octadas de }G_{24}\},
 \qquad H\longmapsto O\text{ tal que }O\cap D=\ell(H),
\]
es inyectiva porque la intersección con $D$ recupera $\ell(H)$.
Esta inyectividad tiene por dominio las hexadas, no el registro $K$.

La relación con las cinco regiones de $\pi$ conserva datos más ricos
que el número cinco. Las emisiones y sus firmas son:
\begin{center}
\small
\begin{tabular}{@{}crll@{}}
\toprule
Región & Multiplicidad & Emisión $U_6$ & Firma \\
\midrule
$\perp$ & $432$ & $(501,614,498,169,272,272)$ & $555555$ \\
$++_1$ & $144$ & $(810,923,870,169,272,272)$ & $339555$ \\
$++_2$ & $144$ & $(810,983,810,169,272,272)$ & $393555$ \\
$++_3$ & $144$ & $(870,923,810,169,272,272)$ & $933555$ \\
$--$ & $144$ & $(870,923,810,418,674,521)$ & $933717$ \\
\bottomrule
\end{tabular}
\end{center}
En cada bloque decimal de tres cifras, la firma usa la diferencia
modular de las dos primeras cifras. Las cinco emisiones tienen la
misma palabra trítica visible $w_\pi=010211$, pero no la misma
genealogía regional. Las tres regiones positivas conservan la cola
$169|272|272$; la región negativa la cambia por $418|674|521$;
la región transversal tiene firma constante y multiplicidad mayor.

La extensión marcada de la carta regional envía la región transversal
a $O_{\ell(B_K)}^\perp$, la negativa a $\iota_D(H^-_\alpha)$ y las
tres positivas a las otras tres elevaciones de la estrella. Así se
realiza la partición $4+1$ de las cinco octadas sobre $\ell(B_K)$.
La correspondencia individual entre los tres nombres positivos y
los tres pétalos restantes exige un orden de carta; no la determina
la incidencia sin ese dato. Por ello la afirmación precisa es una
correspondencia de regiones marcada, compatible con sus roles y con
la bandera, no una biyección canónica deducida del cardinal cinco.

La identidad $432+4\cdot144=1008$ conserva las multiplicidades
regionales. También $1008=36\cdot28$ y una octada tiene
$\binom82=28$ parejas. Estas igualdades son controles aritméticos de
catálogo, no sustitutos del mapa $\iota_D$ ni de la carta regional.

\subsubsection{Representación pentádica de la microfibra y de su incidencia}
\label{pi-estrella:desarrollo}

La disposición de las cinco regiones permite visualizar simultáneamente
la unidad del cociente ternario y la diferenciación que conserva el
catálogo. La aplicación relevante es la restricción
\[
 \Psi\big|_{\mathscr R_\pi^{\rm micro}}:
 \mathscr R_\pi^{\rm micro}\longrightarrow\{010211\},
 \qquad \Psi(U)=(-U)\bmod3.
\]
Sus cinco antecedentes son los vectores \(U_6\) de la tabla anterior.
La figura \ref{pi-estrella:figura} recupera su organización radial,
conservando el orden de carta, las firmas y las multiplicidades del
catálogo (ecuación~\eqref{eq:palabras-regionales}).

% Orden de carta y datos recuperados de fig_cinco_regiones_pi.tex.
\begin{figure}[H]
\centering
\begin{tikzpicture}[x=1cm,y=1cm,font=\small,>=Latex]
 \coordinate (a) at (90:2.2);
 \coordinate (b) at (18:2.2);
 \coordinate (c) at (-54:2.2);
 \coordinate (d) at (-126:2.2);
 \coordinate (e) at (162:2.2);
 \draw[black!35,line width=.45pt]
    (a)--(c)--(e)--(b)--(d)--cycle;
 \node[draw,fill=white,circle,minimum size=1.6cm,align=center,
    inner sep=3pt] (z) at (0,0) {$w_\pi$\\$010211$};
 \foreach \v in {a,b,c,d,e}{
    \fill (\v) circle (2pt);
    \draw[->,line width=.6pt] (\v)--(z);
 }
 \node[above=3mm,align=center] at (a)
   {$U_\pi^\perp$\\$555555$\\multiplicidad $432$};
 \node[right=3mm,align=left] at (b)
   {$U_{\pi,1}^{++}$\\$339555$\\multiplicidad $144$};
 \node[below right=2mm and 1mm,align=left] at (c)
   {$U_{\pi,2}^{++}$\\$393555$\\multiplicidad $144$};
 \node[below left=2mm and 1mm,align=right] at (d)
   {$U_{\pi,3}^{++}$\\$933555$\\multiplicidad $144$};
 \node[left=3mm,align=right] at (e)
   {$U_\pi^{--}$\\$933717$\\multiplicidad $144$};
 \node[align=center,font=\footnotesize] at (0,-3.55)
   {$1_\perp+3_{++}+1_{--}$\qquad $432+4\cdot144=1008$};
\end{tikzpicture}
\caption{Las cinco regiones de \(\pi\) en disposición pentádica.
Cada posición conserva su firma multiplicativa y su multiplicidad;
las flechas radiales representan la proyección común
\(\Psi(U)=010211\). El contorno estrellado organiza las cinco
posiciones de la carta. Su realización en las cinco octadas de Witt
se efectúa mediante el alineamiento regional y la descomposición
\(4+1\) del texto.}
\label{pi-estrella:figura}
\end{figure}


La proyección al centro identifica las cinco imágenes ternarias; el
registro de las posiciones exteriores retiene la información regional
que hace posible la realización incidencial. La distribución
\(1_\perp+3_{++}+1_{--}\) especifica una región transversal, tres
regiones positivas y una negativa. Sus pesos normalizados son
\(3/7\) y \(1/7\) para cada una de las cuatro restantes: la
multiplicidad transversal es triple, aun cuando las cinco imágenes bajo
\(\Psi\) coincidan.

\Needspace{11\baselineskip}
El enlace con la estrella de Witt se formula mediante la bandera
\(B_K\subset H^-_\alpha\), el alineamiento \(\ell\) y la
elevación \(\iota_D\) ya construidos. Las cinco posiciones regionales
se realizan en las cinco octadas sobre \(\ell(B_K)\):
\[
\begin{aligned}
 U_\pi^\perp&\longmapsto O^\perp_{\ell(B_K)},\\
 U_\pi^{--}&\longmapsto\iota_D(H^-_\alpha),\\
 \{U_{\pi,1}^{++},U_{\pi,2}^{++},U_{\pi,3}^{++}\}
 &\longmapsto
 \{\iota_D(B_K\cup P):P\in\mathcal P_+\},\\
 \mathcal P_+&=\bigl\{\{1,3\},\{2,11\},\{8,12\}\bigr\}.
\end{aligned}
\]
La asignación de las tres posiciones positivas a los tres pares se
efectúa en una carta ordenada. La fila transversal y la fila negativa
quedan distinguidas por sus predicados regionales y por la bandera.

Así se articulan las dos figuras: la figura
\ref{exc:fig-estrella} muestra las cuatro hexadas del pequeño diseño
de Witt sobre \(B_K\); la figura \ref{pi-estrella:figura} muestra
las cinco regiones que, bajo el alineamiento, ocupan las cuatro octadas
residuales y la octada transversal del diseño binario. La ley
\(4+1\), demostrada en la proposición
\ref{exc:residuo-binario}, especifica el cambio de incidencia. En esta
composición, \(K\) determina la cuaterna marcada y la firma de
\(\alpha\) distingue una de sus hexadas. Las regiones conservan,
además de su valor ternario común, la estructura necesaria para ese
transporte hacia la realización excepcional.


\subsection{Del código ternario al retículo de Niemeier}
\label{exc:niemeier}

El paso a dimensión veinticuatro no necesita recorrer primero el
código binario. Se realiza directamente desde $C_W$ por pegado de
discriminantes. Así quedan distinguidas dos prolongaciones compatibles:
una conserva el residuo $W_{12}\subset W_{24}$; la otra construye el
retículo sobre el que se realizará el álgebra de operadores de vértice.

Sea $A_2=\mathbb Z a_1\oplus\mathbb Z a_2$ con matriz de Gram
\[
 \begin{pmatrix}2&-1\\-1&2\end{pmatrix},
 \qquad \rho=a_1+a_2,\qquad \rho^2=2.
\]
Representamos las tres clases de $A_2^*/A_2$ mediante
\[
 \varpi_0=0,\qquad
 \varpi_1=\frac{2a_1+a_2}{3},\qquad
 \varpi_2=\frac{a_1+2a_2}{3}.
\]
La aplicación $j\mapsto\varpi_j+A_2$ identifica aditivamente
$\mathbb F_3$ con ese discriminante. Para $c\in C_W$, escribimos
$\phi(c)=(\varpi_{c_1},\ldots,\varpi_{c_{12}})$ y definimos
\begin{equation}
 N(C_W)=\bigcup_{c\in C_W}\bigl(A_2^{12}+\phi(c)\bigr).
 \label{exc:niemeier-def}
\end{equation}

\begin{theorem}[Pegado ternario]
\label{exc:niemeier-teorema}
El objeto \eqref{exc:niemeier-def} es un retículo positivo, par y
unimodular de rango $24$. Su sistema de raíces es exactamente
$A_2^{12}$; por tanto, se reconoce como el retículo de Niemeier con
ese sistema de raíces.
\end{theorem}

\begin{proof}
La linealidad del código hace que la unión de clases sea estable por
suma. La autoortogonalidad anula el emparejamiento discriminante:
éste es un múltiplo de $c\cdot d/3$ módulo $\mathbb Z$. Así, el
retículo es integral. La norma de una clase de pegado es
$2\operatorname{wt}(c)/3$ módulo $2\mathbb Z$. Los pesos de
\eqref{exc:enumerador} son múltiplos de tres, de modo que es par.

El índice de $A_2^{12}$ en $N(C_W)$ es $|C_W|=3^6$. En consecuencia,
\[
 \det N(C_W)=\frac{3^{12}}{(3^6)^2}=1.
\]
La positividad procede de la suma ortogonal de las doce métricas $A_2$.
En una clase no nula de $A_2^*/A_2$ la norma mínima es $2/3$;
toda palabra no nula tiene al menos seis posiciones no nulas.
Por tanto, un vector de una clase de pegado no trivial tiene norma
al menos $6\cdot2/3=4$. No aparecen raíces de norma dos fuera de
$A_2^{12}$. Dentro de éste, las raíces son precisamente las de sus
doce componentes.
\end{proof}

\subsection{Selección incidencial del vecino unimodular sin raíces}
\label{exc:leech}

\paragraph{Del origen incidencial a la modificación reticular.}
\label{inc-ampl-reticulo}
La bandera marcada determina una posición mediante una operación
conjuntista equivariante:
\[
 o_*=(H_e\cap H_\varphi)
       \setminus(H^-_\alpha\cup H_{\rm APP})=\{1\}.
\]
Su función consiste en distinguir una componente entre las doce que
intervienen en la realización reticular. El código ternario completo
actúa como código de pegado en el módulo discriminante
$(A_2^*/A_2)^{12}$. Su preimagen define $N(C_W)$, y la
autodualidad, los pesos y la distancia del código se traducen,
respectivamente, en propiedades de integralidad y determinante, paridad
y cotas de norma de las clases de pegado. Aquí la información simbólica
del código cumple una función que el soporte aislado ya no desempeña.

La métrica $A_2$ y la cámara radial positiva se fijan en el desarrollo
reticular. Dentro de esa cámara, la congruencia necesaria para un vecino
par de índice tres tiene doce realizaciones de norma mínima $54$,
permutadas por el cambio de componente distinguida. El origen $o_*$
selecciona
\[
 v_\alpha=4\rho_1+\rho_2+\cdots+\rho_{12},\qquad
 v_\alpha^2=2(4^2+11)=54.
\]
El coeficiente $4$ pertenece a esa solución de la congruencia en la
cámara declarada. El subíndice $\alpha$ registra la procedencia
incidencial de la selección: su contenido es el marco que intervino en
la clausura y que ahora individualiza una dirección reticular.

El emparejamiento con $v_\alpha$ determina el subretículo
$N_{\alpha,0}$ de índice tres, y la adición de la clase
$v_\alpha/3$ produce el vecino $N_\alpha$ de
\eqref{exc:vecino}. Esta modificación conserva el rango y recupera
unimodularidad y paridad; la selección por congruencia elimina las raíces
anteriores, y las cotas locales excluyen raíces en las nuevas clases.
La identificación con Leech aparece después de esas verificaciones. La
realización binaria de las cinco regiones y esta construcción reticular
son dos prolongaciones del marco común: la segunda recibe directamente
el código ternario como dato de pegado.



La etiqueta $o_*=1$ de la bandera se realiza ahora como una de las
doce componentes $A_2$. Esta operación declara la métrica anterior y
la cámara radial positiva
\[
 v(a)=\sum_{i=1}^{12}a_i\rho_i,
 \qquad a_i=1+3b_i,\quad b_i\in\mathbb Z_{\geq0}.
\]
El requisito de paridad del vecino de índice tres es
$v(a)^2\equiv0\pmod{18}$. Puesto que
\[
 v(a)^2=2\sum_i(1+3b_i)^2,
\]
esa congruencia equivale a $\sum_i b_i\equiv1\pmod3$. El menor
valor de la norma en la cámara se obtiene con un único $b_i=1$ y
todos los demás nulos; es $54$. El origen marcado selecciona
\begin{equation}
 v_\alpha=4\rho_1+\rho_2+\cdots+\rho_{12},
 \qquad v_\alpha^2=54.
 \label{exc:vector}
\end{equation}
La palabra ``menor'' se refiere a esta cámara positiva. No afirma
minimalidad en toda la clase residual: el vector
$(a_1-2a_2,\rho,\ldots,\rho)$ representa la misma clase módulo tres
y tiene norma $14+11\cdot2=36$. La restricción de cámara es parte
de los datos de la construcción, no una conclusión omitible.

Sea $N=N(C_W)$. Definimos
\begin{equation}
 N_{\alpha,0}=\{x\in N:\langle x,v_\alpha\rangle\equiv0\pmod3\},
 \qquad
 N_\alpha=N_{\alpha,0}+\mathbb Z\frac{v_\alpha}{3}.
 \label{exc:vecino}
\end{equation}

\begin{theorem}[Vecino de Leech determinado por la bandera marcada]
\label{exc:teorema-leech}
Con la métrica y la cámara declaradas, $N_\alpha$ es un retículo
positivo, par y unimodular de rango $24$, sin vectores de norma dos.
Por el teorema de clasificación correspondiente, es isométrico al
retículo de Leech $\Lambda_{24}$.
\end{theorem}

\begin{proof}
El carácter $x\mapsto\langle x,v_\alpha\rangle\bmod3$ no es nulo:
una raíz simple de una componente se empareja con $a_i\rho_i$ en
$a_i\equiv1\pmod3$. Por ello $[N:N_{\alpha,0}]=3$ y
$\det N_{\alpha,0}=9$. Además,
$v_\alpha/3\in N_{\alpha,0}^*$ y $v_\alpha\in N_{\alpha,0}$.
Si $v_\alpha/3$ perteneciera a $N$, la integridad de $N$ haría que
todos los emparejamientos con $v_\alpha$ fueran divisibles por tres,
contradiciendo el carácter no nulo. La clase de $v_\alpha/3$ tiene
por tanto orden tres. Se deduce
\[
 [N_\alpha:N_{\alpha,0}]=3,
 \qquad \det N_\alpha=9/3^2=1.
\]
La extensión es integral, y para $x\in N_{\alpha,0}$,
\[
 \left(x+m\frac{v_\alpha}{3}\right)^2
   =x^2+2m\frac{\langle x,v_\alpha\rangle}{3}+6m^2
   \in2\mathbb Z.
\]
Luego es par.

Queda excluir las raíces, que es el paso decisivo. Las raíces de $N$
pertenecen a $A_2^{12}$. Sus emparejamientos con $v_\alpha$ son
congruentes con $\pm1$ o $\pm2$ módulo tres; ninguna pertenece a
$N_{\alpha,0}$. Para las otras dos clases del vecino, cada componente
pertenece a alguna de las clases locales
\[
 A_2+\varpi_j\pm\rho/3,\qquad j=0,1,2.
\]
Como $4\rho/3\equiv\rho/3\pmod{A_2}$, esta descripción vale
también en la componente distinguida. Escribiendo un vector local
como $(u a_1+v a_2)/3$, su norma es
$2(u^2-uv+v^2)/9$. Las seis clases anteriores tienen residuos no
nulos y mínimo $2/9$: los representantes de mínima forma cuadrática
son $(\pm1,0),(0,\pm1),(1,1),(-1,-1)$ en los residuos apropiados.
Un vector de una clase nueva tiene, por tanto, norma al menos
$12\cdot2/9=8/3>2$. No hay raíces en ninguna de las tres clases.
La clasificación de los retículos pares unimodulares positivos de
rango veinticuatro identifica el único caso sin raíces con
$\Lambda_{24}$.
\end{proof}

El argumento no reconoce Leech por una coincidencia de dimensión.
Construye el pegado, prueba integridad, paridad y determinante, y
elimina las raíces mediante una cota local. Sólo entonces aplica el
teorema de clasificación. La dependencia de $\alpha$ se conserva en
la bandera que selecciona $o_*$ y en el vector \eqref{exc:vector};
no se convierte el valor escalar de la constante en una longitud del
retículo.

\subsection{El álgebra de operadores de vértice y el Monstruo}
\label{exc:voa}

Sea $\Lambda=N_\alpha$, ya reconocido como Leech. El paso siguiente
utiliza su retículo integral, no una expansión decimal. La construcción
de álgebra de operadores de vértice de retículo se escribe
\[
 V_\Lambda=M(1)\otimes\mathbb C_\varepsilon[\Lambda].
\]
Aquí $M(1)$ es el espacio de osciladores de Heisenberg en rango
veinticuatro; $\mathbb C_\varepsilon[\Lambda]$ es el álgebra de grupo
torcida por el cociclo de la extensión central del retículo. El
cociclo y el levantamiento de las isometrías forman parte de la
construcción: una suma vectorial desnuda no determina los productos
de vértice.

La negación $\lambda\mapsto-\lambda$ admite el levantamiento estándar
$\theta$. Con su módulo torcido y la elección de paridad compatible,
la construcción de Frenkel--Lepowsky--Meurman produce
\begin{equation}
 V^\natural_{\rm HMT}
   =V_\Lambda^+\oplus(V_\Lambda^T)^+.
 \label{exc:flm}
\end{equation}
El signo $+$ denota la parte par de la acción levantada. La expresión
incluye el producto de vértice del orbifold; no se limita a una suma
de espacios graduados.

\begin{theorem}[Prolongación hasta el módulo Moonshine]
\label{exc:moonshine}
El álgebra \eqref{exc:flm}, construida sobre el vecino
\eqref{exc:vecino}, es una realización del álgebra Moonshine
$V^\natural$. Tiene carga central $24$, componente de peso uno nula
y grupo de automorfismos isomorfo al Monstruo $\mathbb M$. Su carácter
graduado es
\begin{equation}
 \operatorname{Tr}_{V^\natural_{\rm HMT}}q^{L_0-1}
   =J(\tau)=j(\tau)-744
   =q^{-1}+196884q+21493760q^2+\cdots,
 \qquad q=e^{2\pi i\tau}.
 \label{exc:j}
\end{equation}
\end{theorem}

\begin{proof}
El teorema \ref{exc:teorema-leech} verifica las hipótesis reticulares:
positividad, paridad, unimodularidad, rango veinticuatro y ausencia de
raíces. Se aplica entonces la construcción y el teorema de
Frenkel--Lepowsky--Meurman para el retículo de Leech. Éstos identifican
el orbifold con $V^\natural$, establecen su carácter y su grupo de
automorfismos. La composición de la bandera, el vecino y el functor
de construcción de VOA da la realización aquí indicada. No se deduce
el grupo a partir del coeficiente $196884$.
\end{proof}

El teorema utilizado es un reconocimiento clásico posterior con
hipótesis verificadas, no una nueva demostración en este artículo de
la construcción FLM \cite{flm}.
Además, la notación $q=e^{2\pi i\tau}$ es la variable convencional
de Fourier empleada para reconocer el carácter. No introduce $\pi$
ni $e$ como entradas del generador coinductivo que ya los ha determinado.

Puede verse una parte del mecanismo de traza antes del orbifold.
Sólo el vector reticular cero queda fijo por la negación; cada
oscilador contribuye con signo alterno. De ahí
\[
 \operatorname{Tr}_{V_\Lambda}\theta q^{L_0-1}
   =t_2(\tau)
   :=q^{-1}\prod_{n\geq1}(1+q^n)^{-24}.
\]
Esta traza en el álgebra de retículo no debe confundirse con el
carácter sin inserción de $V^\natural$. Tampoco identifica por sí
sola una clase de conjugación del Monstruo: el espacio, la inserción
y el levantamiento deben especificarse antes de comparar funciones.

\paragraph{Graduación, unidad y automorfismos.}
\label{inc-ampl-automorfismos}
La construcción de Frenkel--Lepowsky--Meurman se aplica al retículo
$N_\alpha$ una vez verificadas sus propiedades. Incorpora el espacio
de osciladores, la extensión central del retículo, el producto de vértice
y el sector torcido por el levantamiento de la negación. El resultado es
el álgebra $V^\natural_{\rm HMT}$ de \eqref{exc:flm}, cuyo grupo
de automorfismos es isomorfo al Monstruo. Su dominio de acción queda así
fijado por una estructura algebraica con graduación y producto. El teorema
de Moonshine de Borcherds determina posteriormente las trazas graduadas
con inserción de esos automorfismos \cite{flm,borcherds}.

La unidad adquiere en este nivel una realización precisa: la componente
de peso cero es la línea del vacío y aporta el coeficiente $1$ de
$q^{-1}$ en el carácter graduado. En el nivel cilíndrico, las
aplicaciones de refinamiento conservan la función constante mediante
precomposición. Cada afirmación de conservación tiene, por tanto, su
objeto y su aplicación. El desarrollo de las pruebas hace visible esa
continuidad de construcción, desde las operaciones aritméticas y la
memoria hasta la incidencia, la métrica y los productos de vértice.

Esta articulación sitúa a $K$ y a la clausura de $\alpha$ en el argumento
principal. $K$ conserva reversiblemente una coordenada integral de la
historia; la normalización de $\alpha$ compone los registros regionales
con esa coordenada; la configuración firmada determina una bandera; y
la bandera marcada selecciona la modificación reticular sobre la que se
realiza Moonshine. Las demostraciones siguientes desarrollan estas
aplicaciones, distinguiendo las equivalencias reversibles, los cocientes
de información y las realizaciones que añaden métrica o estructura
algebraica. Así se expresa una relación demostrable entre estructura,
ecuación y valor, conservando la atribución de cada construcción y el
alcance de sus teoremas.


\subsection{Construcción estatal de las doce fibras ciclotómicas}
\label{nuclear:fibras-app-witt}
Sobre las marcas de APP, la multiplicación \(a:r\mapsto4r\) en
\(\ZZ/9\ZZ\) satisface \(a^3=1\). Las hojas verifican
\[
 S(ax,ay)=aS(x,y),\qquad P(ax,ay)=a^{-1}P(x,y),
\]
pues \(16\equiv7\equiv4^{-1}\pmod9\). Las órbitas unitarias
ordenadas son \(\mathcal O_0=(1,4,7)\), \(\mathcal O_1=(2,8,5)\).
En cada una, el módulo de aumento
\(A(\mathcal O)=\ker(\ZZ[\mathcal O]\xrightarrow{\sum}\ZZ)\)
hereda el producto que hace ortonormales los \(e_r\). En la base
\(b_1=e_r-e_{4r}\), \(b_2=e_{4r}-e_{7r}\), sus matrices son
\[
 G_{A_2}=\begin{pmatrix}2&-1\\-1&2\end{pmatrix},\qquad
 C=\begin{pmatrix}0&-1\\1&-1\end{pmatrix},\qquad
 C^{\mathsf T}G_{A_2}C=G_{A_2}.
\]
La forma reticular procede, por tanto, de las órbitas de las marcas.
Para los pares \((1,2),(2,3),(3,1)\) y las hojas
\(\mu\in\{\Sigma,\Pi\}\), se obtiene
\[
 W_{\APP}=\bigoplus_{(i,j)}\bigoplus_\mu\bigoplus_{\varepsilon=0}^1
 A(\mathcal O_\varepsilon),\qquad
 \ell^\Sigma_{ij}(x)=x_i+x_j,\quad
 \ell^\Pi_{ij}(x)=x_ix_j\pmod9.
\]
Cada sumando conserva par, hoja y órbita. La aplicación estatal es
\[
 \Psi:(\ZZ/9\ZZ)^3\longrightarrow W_{\APP},\qquad
 \Psi(x)_{ij,\mu,\varepsilon}=\beta_\varepsilon(\ell^\mu_{ij}(x)),
 \quad
 \beta_\varepsilon(r)=
 \begin{cases}e_r-e_{4r},&r\in\mathcal O_\varepsilon,\\0,&r\notin\mathcal O_\varepsilon.
 \end{cases}
\]
\begin{proposition}[Equivarianza y rango del módulo estatal]
\label{nuclear:psi-rango}
La acción de \(C\) en las seis fibras aditivas y de \(C^{-1}\) en
las seis multiplicativas cumple \(\Psi(ax)=\rho(a)\Psi(x)\).
Las imágenes generan \(W_{\APP}\otimes\QQ\), de dimensión veinticuatro.
\end{proposition}
\begin{proof}
Las identidades de las hojas y \(\beta(4r)=C\beta(r)\) prueban la
equivarianza. Para el rango, se ordenan las coordenadas por los tres
pares indicados, por \(\Sigma,\Pi\), por \(\mathcal O_0,\mathcal O_1\)
y por \(b_1,b_2\). En cada órbita, \(\beta\) toma sucesivamente
\((1,0),(0,1),(-1,-1)\). La matriz de las imágenes de estos estados,
leídos por filas, tiene determinante \(9\):
\[
\begin{array}{rrrr}
(0,0,1)&(0,0,2)&(0,0,4)&(0,0,5)\\
(0,1,0)&(0,1,1)&(0,1,2)&(0,1,3)\\
(0,1,4)&(0,1,5)&(0,1,6)&(0,2,0)\\
(0,2,2)&(0,2,3)&(0,4,0)&(0,5,0)\\
(1,0,1)&(1,0,2)&(1,0,4)&(1,0,5)\\
(1,1,0)&(1,2,0)&(1,4,0)&(1,5,0).
\end{array}
\]
La eliminación sucesiva, reduciendo cada fila por las anteriores y
normalizando su primer coeficiente no anulado, da columnas pivote
numeradas desde cero
\[
\begin{split}
&(8,10,9,11,0,12,14,16,13,15,17,2,\\
&\hspace{2em}18,19,1,3,20,22,21,23,4,6,5,7)
\end{split}
\]
y pivotes
\[
(1,1,1,-1,1,1,1,1,1,-1,3,1,1,3,1,-1,1,1,1,-1,1,1,1,-1).
\]
Su producto, multiplicado por el signo de la permutación de columnas,
es \(9\). El menor invertible acredita el rango total.
\end{proof}

\subsection{Alineamiento dodecafásico con las posiciones de Witt}
\label{nuclear:alineamiento-witt}
El índice del par es \(i\in\ZZ/3\ZZ\), el de la hoja
\(\mu\in\{0,1\}\) y el de la órbita \(\varepsilon\in\{0,1\}\).
Fijado el origen dodecafásico, \(\lambda(i,\mu,\varepsilon)
=4i+2\mu+\varepsilon\pmod{12}\) es biyectiva: los cuatro valores
de cada \(i\) forman un bloque y los tres bloques son disjuntos.
Sea \(R=\left(\begin{smallmatrix}0&1\\1&0\end{smallmatrix}\right)\).
La multiplicación prueba \(RCR=C^{-1}\) y
\(R^{\mathsf T}G_{A_2}R=G_{A_2}\). Para las trivializaciones
\(\iota_b\) y los marcos \(P_b\) transportados conjuntamente con
la incidencia de Witt, ponemos
\[
g_b=\iota_b^{-1}P_bCP_b^{-1}\iota_b,\qquad
T_{i\mu\varepsilon}=\iota_{\lambda(i,\mu,\varepsilon)}^{-1}
P_{\lambda(i,\mu,\varepsilon)}R^\mu.
\]
Entonces
\[
g_{\lambda(i,\mu,\varepsilon)}T_{i\mu\varepsilon}
=T_{i\mu\varepsilon}C^{(-1)^\mu}.
\]
La suma directa es un isomorfismo \(C_3\)-equivariante entre las doce
fibras estatales y las doce posiciones reticulares. Conserva par,
hoja, órbita y forma bajo el transporte simultáneo del origen,
orden y orientación. Estos datos permanecen en el pegado reticular
y el paso al vecino que se construyen a continuación.

% Incorporación aditiva para el Artículo I; no modifica la fuente anterior.
% RESULTADO_RECUPERADO: integral, capítulo 22, propietario extenso
% colaboracion/partes_i_ii/source/base_residencias/
% capitulo_22_c20_lineas_2613_3270.tex, líneas 176--576.
% Equivalencia 2/3: public/residencias/capitulo_22_voa_leech_moonshine.tex,
% líneas 55--127. Orden dos y Fricke: hmt/15d_voa_moonshine_orden_dos_rev11.tex.
% Prerrequisitos residentes en I: código C_W, matriz A_W, N(A_2^12),
% origen o_*, vector v_alpha de norma54, N_alpha reconocido como Leech,
% VOA reticular y presentación FLM de las secciones precedentes.
% Bibliografía nueva: chenlamshimakura2016, abelamyamada2017,
% kmconwaynorton1979 (bibitems listos para incorporar al final de este archivo).

\subsection{Orientación ternaria del vecino de Leech}
\label{km:orden-tres-reticular}

La presentación de orden dos construida anteriormente utiliza la negación
del retículo de Leech. El mismo retículo, obtenido de la incidencia marcada
por \(K\) y por el soporte dodecafásico de \(\alpha\), admite una segunda
presentación compatible con el carácter ternario de la construcción.
Para hacer explícita esta continuidad se conserva el pegado, se construye
su isometría de orden tres y se comprueba que ésta preserva el vecino ya
seleccionado. El paso no requiere alterar \(K\), su lectura racional ni
la definición previa de \(\alpha\). Se desarrolla aquí la prolongación de orden tres, incluyendo las operaciones
reticulares y las condiciones precisas de los teoremas clásicos utilizados.

Trabajamos en la base de raíces simples de cada factor \(A_2\), con
\begin{equation}
 B_2=\begin{pmatrix}2&-1\\-1&2\end{pmatrix},\qquad
 \mathsf C=\begin{pmatrix}0&-1\\1&-1\end{pmatrix},\qquad
 \mathsf R=\begin{pmatrix}0&1\\1&0\end{pmatrix},\qquad
 \rho=\begin{pmatrix}1\\1\end{pmatrix}.
 \label{km:coxeter-reflexion}
\end{equation}
La multiplicación da
\[
 \mathsf C^3=I_2,\quad
 \mathsf C^2+\mathsf C+I_2=0,\quad
 \mathsf C^{\mathsf T}B_2\mathsf C=B_2,\quad
 \mathsf R\mathsf C\mathsf R=\mathsf C^{-1},\quad
 \mathsf R^{\mathsf T}B_2\mathsf R=B_2.
\]
Estas identidades fijan tanto la métrica como la orientación de la acción.
En el discriminante \(A_2^*/A_2\cong\mathbb F_3\) tomamos como
generador la clase de \(\varpi=(2,1)^{\mathsf T}/3\). Se cumple
\(\mathsf C\varpi-\varpi=(-1,0)^{\mathsf T}\), por lo que
\(\mathsf C\) actúa trivialmente sobre el discriminante; la reflexión
\(\mathsf R\) cambia su signo.

La matriz de Paley--Witt ya construida produce un elemento de soporte
completo mediante la siguiente operación:
\begin{equation}
 \begin{gathered}
 q_*=(1,1,1,1,1,1),\qquad q_*A_W=(2,1,1,1,1,1),\\
 c_*=(q_*,q_*A_W)
 =(1,1,1,1,1,1,2,1,1,1,1,1)\in C_W.
 \end{gathered}
 \label{km:palabra-total}
\end{equation}
Cada coordenada de \(c_*\) es no nula. Para las doce posiciones
definimos los marcos y la isometría
\begin{equation}
 P_b(c_*)=
 \begin{cases}\mathsf R,&(c_*)_b=1,\\I_2,&(c_*)_b=2,\end{cases}
 \qquad
 g_c=\bigoplus_{b=1}^{12}P_b(c_*)\mathsf C P_b(c_*)^{-1}.
 \label{km:isometria-ternaria}
\end{equation}
El símbolo \(P_b(c_*)\) designa aquí un marco invertible de un factor
\(A_2\), no el proyector regional \(P_3\).

\begin{proposition}[Isometría de orden tres sobre el vecino marcado]
\label{km:estabilidad-vecino}
El operador \(g_c\) es una isometría de orden tres sin vectores fijos
no nulos. Preserva el retículo de Niemeier \(N=N(C_W)\) y el vecino
\(N_\alpha\) construido en \eqref{exc:vecino}.
\end{proposition}
\begin{proof}
Cada bloque de \(g_c\) es \(\mathsf C\) o \(\mathsf C^{-1}\);
su polinomio mínimo es \(X^2+X+1\). Esto prueba el orden y la
ausencia de vectores fijos. Ambos bloques actúan trivialmente sobre el
discriminante, de modo que conservan cada clase de pegado y, por tanto,
el retículo \(N\).

Escribamos \(v=v_\alpha=4\rho_{o_*}+\sum_{b\ne o_*}\rho_b\),
con los marcos y el origen incidencial ya fijados. Sus coeficientes
radiales son \(a_{o_*}=4\) y \(a_b=1\) para \(b\ne o_*\).
Todos satisfacen \(a_b\equiv1\pmod3\). En un factor,
\[
 (\mathsf C-I_2)\rho=(-2,-1)^{\mathsf T},\qquad
 (\mathsf C^{-1}-I_2)\rho=(-1,-2)^{\mathsf T}.
\]
Después de dividir por tres, sus clases discriminantes son,
respectivamente, \(2\) y \(1\). En consecuencia, los vectores
\begin{equation}
 y=\frac{g_cv-v}{3},\qquad z=\frac{g_c^{-1}v-v}{3}
 \label{km:desplazamientos}
\end{equation}
tienen palabras discriminantes \(c_*\) y \(-c_*\), ambas en
\(C_W\). Por ello \(y,z\in N\). El producto local es
\[
 \left\langle
 \frac{(\mathsf C^{\pm1}-I_2)a_b\rho}{3},a_b\rho
 \right\rangle=-a_b^2,
\]
y su suma da
\(\langle y,v\rangle=\langle z,v\rangle=-(16+11)=-27\).
Así, \(y,z\in N_0=\{x\in N:\langle x,v\rangle\equiv0\pmod3\}\).
Para \(x\in N_0\),
\[
 \langle g_cx,v\rangle
 =\langle x,g_c^{-1}v\rangle
 =\langle x,v\rangle+3\langle x,z\rangle\equiv0\pmod3.
\]
El argumento aplicado a \(g_c^{-1}\) prueba \(g_cN_0=N_0\).
Finalmente, \(g_c(v/3)=v/3+y\) conserva la clase que genera
\(N_\alpha=N_0+\mathbb Z(v/3)\). Se obtiene
\(g_cN_\alpha=N_\alpha\).
\end{proof}

La prueba conserva más que la dimensión del retículo: utiliza el código
de pegado, su orientación, el origen de la bandera y la congruencia del
vector radial. Estos datos explican por qué la segunda presentación
permanece vinculada al registro dodecafásico del que partió la construcción.

\subsection{Traza ternaria, peso torcido y extensión de orden tres}
\label{km:orbifold-tres}

Sea \(\Lambda=N_\alpha\) y sea \(\widehat g_c\) el levantamiento
estándar de \(g_c\) al álgebra reticular
\(V_\Lambda=M(1)\otimes\mathbb C_\varepsilon[\Lambda]\), con la
extensión central y el cociclo ya declarados. El orden impar permite
elegir este levantamiento de orden tres. Cada factor complejo \(A_2\)
aporta los autovalores \(\zeta_3\) y \(\zeta_3^2\), donde
\(\zeta_3^2+\zeta_3+1=0\).

\begin{proposition}[Traza graduada y tipo del levantamiento]
\label{km:traza-tipo}
Para \(j=1,2\),
\begin{equation}
 \operatorname{Tr}_{V_\Lambda}
       \bigl(\widehat g_c^{\,j}q^{L_0-1}\bigr)
 =t_3(\tau)
 :=q^{-1}\prod_{n\ge1}(1+q^n+q^{2n})^{-12}.
 \label{km:tres-traza}
\end{equation}
El peso mínimo de los módulos torcidos no triviales es \(4/3\), y
el levantamiento es de tipo cero.
\end{proposition}
\begin{proof}
En la parte reticular de la traza sólo contribuye el vector cero,
porque \(g_c\) y \(g_c^2\) carecen de vectores fijos no nulos.
En cada grado oscilador, la traza de potencias simétricas es el inverso
del determinante \(I-g_c^jq^n\). Cada factor \(A_2\) contribuye
\(1+q^n+q^{2n}\), lo que demuestra \eqref{km:tres-traza}.

La fórmula del peso del vacío torcido para un automorfismo reticular de
orden tres, aplicada a los dos espacios propios de dimensión doce, da
\[
 \rho_{\widehat g_c}
 =\frac{1}{4\cdot3^2}
   \bigl(1\cdot2\cdot12+2\cdot1\cdot12\bigr)=\frac43.
\]
El tipo es la clase de \(3^2\rho_{\widehat g_c}=12\) módulo
tres, que es cero. La fórmula de peso es un resultado de la teoría de
módulos torcidos; las multiplicidades y la isometría que la alimentan
se han construido aquí.
\end{proof}

\begin{theorem}[Segunda presentación del álgebra Moonshine]
\label{km:moonshine-tres}
La extensión cíclica de tipo cero
\begin{equation}
 V_{(3)}=\operatorname{Orb}_{\mathbb Z/3}
                 (V_\Lambda,\widehat g_c)
 \label{km:extension-orden3}
\end{equation}
es isomorfa al álgebra de operadores de vértice \(V^\natural\).
\end{theorem}
\begin{proof}
El teorema \ref{exc:teorema-leech} proporciona un retículo positivo,
par, unimodular, de rango veinticuatro y sin raíces. La
Proposición~\ref{km:estabilidad-vecino} construye sobre ese mismo
retículo una isometría de orden tres sin vectores fijos; la
Proposición~\ref{km:traza-tipo} verifica el levantamiento y su tipo.
Se aplica entonces el teorema de Chen--Lam--Shimakura sobre el
orbifold de orden tres del álgebra reticular de Leech
\cite[Teorema~1.1]{chenlamshimakura2016}. La identificación también queda comprendida
en el tratamiento uniforme de Abe--Lam--Yamada
\cite[Teorema~4.4 y secciones~2--4]{abelamyamada2017}. Esta aplicación es posterior a la construcción
de los datos reticulares; la igualdad de caracteres no sustituye las
hipótesis del teorema utilizado.
\end{proof}

Para precisar la simetría dual, escribamos
\(V_{(3)}=W_0\oplus W_1\oplus W_2\), donde \(W_i\) es el sector
integral seleccionado en el módulo \(\widehat g_c^{\,i}\)-torcido.
El producto satisface
\(Y(W_i,z)W_j\subset W_{i+j\bmod3}((z))\). Por ello
\begin{equation}
 g^\vee|_{W_i}=\zeta_3^i I
 \label{km:dual-ciclico}
\end{equation}
define un automorfismo de orden tres. En efecto, el factor sobre un
producto es \(\zeta_3^{i+j}=\zeta_3^i\zeta_3^j\), y el vacío y el
vector conforme pertenecen a \(W_0\). La identificación clásica de
este automorfismo en el Monstruo es la clase \(3B\), con serie
\(T_{3B}=t_3+12\) \cite{kmconwaynorton1979,borcherds}.
Esta normalización puede comprobarse en la construcción. El carácter
reticular es \(\operatorname{ch}V_\Lambda=J+24\): su polo inicial
es \(q^{-1}\), su espacio de peso uno tiene dimensión 24 y el
carácter modular de peso cero queda determinado por esos datos.
Si \(F\) es el carácter del subespacio fijo y \(H\) el de cada
sector torcido integral, la proyección sobre invariantes y la conjugación
de los dos sectores dan
\[
 F=\frac{J+24+2t_3}{3},\qquad F+2H=J,
 \qquad \operatorname{Tr}g^\vee q^{L_0-1}=F-H=t_3+12.
\]
Su dominio es el álgebra extendida, mientras
\(\widehat g_c\) actúa sobre el álgebra reticular anterior.

\subsection{Equivalencia de las presentaciones de órdenes dos y tres}
\label{km:equivalencia-dos-tres}

Denotemos por
\[
 V_{(2)}=(V_\Lambda)^+\oplus(V_\Lambda^T)^+
\]
la presentación FLM de \eqref{exc:flm}. Ambas presentaciones conservan
el mismo vecino de Leech como antecedente. Sus extensiones emplean
automorfismos diferentes y productos de vértice determinados mediante
los respectivos datos de levantamiento.

\begin{theorem}[Equivalencia graduada y transporte de automorfismos]
\label{km:equivalencia-voa}
Fijados isomorfismos de álgebras de operadores de vértice
\[
 \phi_2:V_{(2)}\xrightarrow{\ \cong\ }V^\natural,
 \qquad
 \phi_3:V_{(3)}\xrightarrow{\ \cong\ }V^\natural,
\]
la composición
\begin{equation}
 \Phi_{23}=\phi_2^{-1}\phi_3:
 V_{(3)}\xrightarrow{\ \cong\ }V_{(2)}
 \label{km:Phi23}
\end{equation}
conserva el vacío, el vector conforme, la gradación y el producto de
vértice. La conjugación por \(\Phi_{23}\) identifica los grupos de
automorfismos y conserva sus trazas graduadas.
\end{theorem}
\begin{proof}
La composición de los dos isomorfismos preserva las estructuras
declaradas. En particular,
\[
 \Phi_{23}Y_{(3)}(a,z)b
 =Y_{(2)}(\Phi_{23}a,z)\Phi_{23}b,
 \qquad
 \Phi_{23}L_0^{(3)}=L_0^{(2)}\Phi_{23}.
\]
Si \(h\) es un automorfismo de \(V_{(3)}\), entonces
\(h' =\Phi_{23}h\Phi_{23}^{-1}\) es un automorfismo de
\(V_{(2)}\). La igualdad de trazas en cada componente homogénea
finito-dimensional da
\[
 \operatorname{Tr}_{V_{(3)}}h q^{L_0-1}
 =\operatorname{Tr}_{V_{(2)}}h' q^{L_0-1}.
\]
\end{proof}

La elección de \(\phi_2\) y \(\phi_3\) forma parte del enunciado;
\(\Phi_{23}\) no se presenta como una identificación canónica
independiente de esas elecciones. Tampoco transforma un automorfismo
de orden tres en la involución de orden dos: la conjugación conserva
el orden. Su contenido es que dos extensiones distintas realizan la
misma estructura graduada, con transporte explícito de sus productos,
representaciones y trazas.

\subsection{Funciones eta, involución de Fricke y normalización del vacío}
\label{km:fricke}

Para \(\operatorname{Im}\tau>0\), las trazas de los dos automorfismos
reticulares admiten las expresiones
\begin{equation}
 \eta(\tau)=q^{1/24}\prod_{n\ge1}(1-q^n),\qquad
 t_2=\left(\frac{\eta(\tau)}{\eta(2\tau)}\right)^{24},\qquad
 t_3=\left(\frac{\eta(\tau)}{\eta(3\tau)}\right)^{12}.
 \label{km:eta-trazas}
\end{equation}
Se deducen directamente de las factorizaciones
\(1-q^{2n}=(1-q^n)(1+q^n)\) y
\(1-q^{3n}=(1-q^n)(1+q^n+q^{2n})\). El orden en \(q\) del
cociente de nivel tres, antes de elevarlo a la potencia doce, es
\[
 \frac1{24}-\frac3{24}=-\frac1{12};
 \qquad 12\left(-\frac1{12}\right)=-1.
\]
Aquí \(-1/12\) designa la diferencia exacta de los exponentes
iniciales de dos funciones eta. La operación está especificada y no
requiere interpretar una suma divergente como una suma convergente.
La coincidencia con otros defectos normalizados del corpus conserva
las realizaciones respectivas; no identifica sus operadores por
igualdad del escalar.

\begin{proposition}[Involución de Fricke sobre la traza de orden dos]
\label{km:fricke-prop}
Sea \(w_2\tau=-1/(2\tau)\). Entonces
\begin{equation}
 t_2(w_2\tau)=\frac{2^{12}}{t_2(\tau)},\qquad
 T_{2A}=t_2+24+\frac{2^{12}}{t_2}
 \label{km:fricke-t2}
\end{equation}
y \(T_{2A}(w_2\tau)=T_{2A}(\tau)\).
\end{proposition}
\begin{proof}
La ley de transformación
\(\eta(-1/\tau)=(-i\tau)^{1/2}\eta(\tau)\) da
\[
 \frac{\eta(-1/(2\tau))}{\eta(-1/\tau)}
 =\sqrt2\,\frac{\eta(2\tau)}{\eta(\tau)}.
\]
La potencia veinticuatro produce la primera identidad. La sustitución
intercambia los dos términos variables de \(T_{2A}\) y deja fijo
el término constante. El reconocimiento de esta expresión como serie
de McKay--Thompson de la clase \(2A\) corresponde a la
simetrización de nivel dos de Conway--Norton
\cite[\S\S4 y 6, tabla~3]{kmconwaynorton1979} y al teorema
de Moonshine \cite{borcherds}.
\end{proof}

Las uniformizaciones modulares de niveles dos y tres, en las
normalizaciones anteriores, son
\begin{align}
 j&=\frac{(t_2+256)^3}{t_2^2},
 \label{km:j-nivel2}\\
 J=j-744
 &=t_3+12+\frac{10\cdot3^9}{t_3}
       +\frac{4\cdot3^{14}}{t_3^2}
       +\frac{3^{18}}{t_3^3}.
 \label{km:j-nivel3}
\end{align}
Se utilizan como identidades modulares clásicas sobre las trazas ya
construidas, no como reglas para seleccionar el código o el vecino.
La segunda también se escribe
\(j=(t_3+27)(t_3+243)^3/t_3^3\).
El desarrollo directo del producto de \eqref{km:tres-traza} comienza por
\[
 t_3=q^{-1}-12+54q-76q^2-243q^3+1188q^4+O(q^5).
\]
En particular, \([q]J=54+10\cdot3^9=196884\); ésta es una
consecuencia de las trazas y la uniformización. La construcción de
\(V^\natural\) conserva su prueba reticular y de orbifold,
independientemente de esta comprobación del primer coeficiente.

Las tres operaciones que comparecen en esta prolongación tienen
dominios distintos. La simetría dual \(g^\vee\) actúa sobre sectores
del álgebra de vértices; Fricke actúa sobre el parámetro modular y
su función; la dualidad de radio intercambia momento y enrollamiento
en el sector reticular compacto. Sus fórmulas de reciprocidad pueden
coordinarse mediante los mapas correspondientes, pero el nombre
«dualidad» no sustituye esos mapas. Esta distinción permite continuar
la construcción hacia las pantallas dimensionales manteniendo la
incidencia, la gradación y el origen de cada operación.

% BIBITEMS NUEVOS: copiar al entorno thebibliography del artículo.
% Verificación primaria: arXiv 1606.05961v2 (teorema 1.1);
% arXiv 1705.09022v4 (secciones 2--4);
% Conway--Norton, pp. 311 y 314--315, tabla 3.
% \bibitem{chenlamshimakura2016}
% H.-Y. Chen, C. H. Lam y H. Shimakura,
% ``$\mathbb Z_3$-orbifold construction of the Moonshine vertex operator
% algebra and some maximal $3$-local subgroups of the Monster'',
% prepublicación, arXiv:1606.05961 (2016), versión 2 (2017).
% \href{https://arxiv.org/abs/1606.05961}{arXiv:1606.05961}.
% \bibitem{abelamyamada2017}
% T. Abe, C. H. Lam y H. Yamada,
% ``A remark on $\mathbb Z_p$-orbifold constructions of the Moonshine
% vertex operator algebra'', prepublicación, arXiv:1705.09022 (2017).
% \href{https://arxiv.org/abs/1705.09022}{arXiv:1705.09022}.
% \bibitem{kmconwaynorton1979}
% J. H. Conway y S. P. Norton, ``Monstrous Moonshine'',
% \emph{Bulletin of the London Mathematical Society} \textbf{11},
% 308--339 (1979).
% \href{https://doi.org/10.1112/blms/11.3.308}{doi:10.1112/blms/11.3.308}.


% REV02: pruebas reunidas en 02_dualidad_t.tex y 03_pantallas_coordinadas.tex.
\subsection{Incidencia excepcional y coordinación de las pantallas}
\label{exc:pantallas-dualidad}

La dirección seleccionada por \(K\) y el retículo de Leech proceden
de representaciones distintas de la incidencia construida. La primera
determina la reducción ortogonal \(12\to11\to10\); la segunda
interviene en la reducción isotrópica \(26\to24\). La coordinación
conserva el espacio previo, sus relaciones y la información transmitida
por cada proyección.

Las secciones \ref{iv:dualidad} y \ref{iv:pantallas} desarrollan
conjuntamente esa coordinación. La primera reúne la involución de
coordenadas enteras, la conjugación unitaria con sus dominios y el
cambio de sección residual con memoria del cociente. La segunda
construye el morfismo \(\mathcal R_M^{\mathrm{alg}}\) de
\eqref{iv:eq:rm-alg} y demuestra el isomorfismo de las presentaciones
cocientadas en el teorema \ref{iv:thm:pantallas-isomorfas}.
El grafo conserva la coordenada en \(V_{11}\) junto con su proyección
en \(V_{10}\), mientras el cociente isotrópico recupera el retículo
de Leech. Así se preservan los respectivos núcleos y codominios.

La elipse de acción se construye antes de aplicar la dualidad y fija
su radio adimensional junto con el radio recíproco
(sección \ref{iv:ellipse:section}). El sector compacto enlaza entonces
la incidencia excepcional con una forma cuadrática determinada. Los
campos, la acción y los operadores de la realización física posterior
se presentan en su dominio propio, conservando este antecedente
algebraico y sus demostraciones internas.


\subsection{Trazas graduadas de Moonshine y representación del Monstruo}
\label{exc:alcance}

Para $g\in\operatorname{Aut}(V^\natural_{\rm HMT})$, la traza
con inserción es
\[
 T_g(\tau)=
 \operatorname{Tr}_{V^\natural_{\rm HMT}}gq^{L_0-1}.
\]
El teorema de Moonshine monstruoso identifica estas series como
Hauptmoduln de los grupos de género cero correspondientes
\cite[teorema~1.1]{borcherds}.
El carácter $J$, el grupo $\mathbb M$ y la familia $g\mapsto T_g$
son realizaciones del álgebra ya construida. No forman una sucesión
causal en la que primero se adivina $J$ y después se deduce el Monstruo.

La unidad tiene aquí una localización algebraica precisa. El subespacio
de peso cero es la línea del vacío, \(V^\natural_0=\CC\mathbf1\), y
el coeficiente de \(q^{-1}\) de \(J\) es por ello \(1\). En peso uno
la dimensión es cero, de donde procede la ausencia del término constante
en \(J=j-744\). Un morfismo de álgebras de vértices conserva el vector
vacío. Esta propiedad proporciona el sentido técnico de conservación
de la unidad en este nivel; la unidad de los observables cilíndricos
se conserva por la precomposición demostrada anteriormente. La relación
entre ambos niveles debe seguir los mapas de construcción con sus unidades.

En peso dos aparece la igualdad
\[
 196884=1+196883=196560+324=54(5\cdot729+1).
\]
La primera separación corresponde a la línea conforme trivial y al
componente irreducible de dimensión $196883$ de la representación
del Monstruo. Las otras expresiones tienen lecturas reticulares o
aritméticas. Su igualdad no identifica automáticamente sus sumandos
con regiones, palabras o rutas del TPK. Para hacerlo habría que
presentar los mapas entre esos espacios y probar que respetan las
acciones relevantes.

En particular, una isometría $N_\alpha\simeq\Lambda_{24}$ induce
una realización de \eqref{exc:flm}, y una elección de isomorfismo
con $V^\natural$ transporta su grupo de automorfismos. No se ha
construido con ello una acción del Monstruo sobre los dígitos de
$\pi$, sobre los bloques de $K$ o sobre todas las historias TPK.
Tal afirmación adicional requeriría una representación
$\rho_{\rm TPK}$ y un mapa $F$ con un cuadrado explícito
\[
 F\rho_{\rm TPK}(g)=\rho_{V^\natural}(g)F,
\]
incluyendo dominio, codominio e información conservada. La ausencia
de ese paso en la afirmación presente no retira el entrelazador
incidencial $U_W$, la bandera ni el vecino ya construidos; sólo
delimita una promoción distinta, de acción sobre el estado fuente.

La cadena establecida en esta sección puede resumirse, manteniendo
visibles sus datos de transporte, como
\begin{align*}
 &\text{APP--TRIT--TPK y estado aritmético con memoria de trayectoria}
   \longrightarrow (w_j,f_j;K,Z_0)\\
 &\longrightarrow (C_W,\mathcal H;B_K\subset H^-_\alpha,
                   P_\pi,H_e,H_\varphi,H_{\rm APP})\\
 &\longrightarrow \bigl(N(C_W),o_*,\text{métrica y cámara}\bigr)
   \longrightarrow N_\alpha
   \longrightarrow V^\natural_{\rm HMT}.
\end{align*}
En el nivel de incidencia, las estrellas generan $M_{12}$ y la
carta residual realiza la estructura regional $4+1$ en $W_{24}$.
En el nivel reticular y de VOA, la construcción da Leech y el módulo
Moonshine. Son prolongaciones de una misma fuente marcada, con
objetos y mapas diferentes. La continuidad causal no obliga a
identificar sus codominios ni a borrar las elecciones explícitas
de marco, cámara y levantamiento.

El resultado de la sección es, por tanto, una derivación reunida:
las pruebas finitas y reticulares hacen efectivo el tránsito desde
la incidencia dodecafásica hasta las hipótesis de los teoremas de
reconocimiento. La conclusión comprende las estructuras y los
transportes especificados en las proposiciones precedentes, con sus
dominios y elecciones explícitos. La comparación de constantes y la
validación experimental corresponden a afirmaciones distintas de esta
construcción algebraica.

\subsection{Orientaciones de soporte completo y equivariancia reticular}
\label{paper:24-orientaciones}

El enumerador de pesos de \eqref{exc:enumerador} contiene
\(24y^{12}\). Por tanto, el conjunto
\[
 C_W^\times:=
 \{u\in C_W:u_b\in\{1,2\}\text{ para todo }b\}
\]
posee 24 elementos. Para cada \(u\in C_W^\times\), definimos
\[
 P_b(u)=
 \begin{cases}
  \mathsf R,&u_b=1,\\
  I,&u_b=2,
 \end{cases}
 \qquad
 g_u:=\bigoplus_{b=0}^{11}P_b(u)\mathsf C P_b(u)^{-1}.
\]

\begin{proposition}[Robustez de las orientaciones de soporte total]
\label{paper:robustez-orientaciones}
Para toda \(u\in C_W^\times\),
\[
 \frac{g_uv-v}{3},
 \quad
 \frac{g_u^{-1}v-v}{3}
 \in N_0.
\]
En particular, cada \(g_u\) preserva el vecino construido con la
orientación transportada.
\end{proposition}

\begin{proof}
En coordenadas de raíces simples,
\[
 \mathsf C\rho-\rho=(-2,-1),
 \qquad
 \mathsf C^{-1}\rho-\rho=(-1,-2).
\]
Las clases discriminantes de los dos desplazamientos son \(u\) y \(-u\).
Para el coeficiente radial \(a_b\in\{4,1\}\),
\[
 \left\langle
 \frac{(\mathsf C^{\pm1}-I)a_b\rho}{3},a_b\rho
 \right\rangle=-a_b^2.
\]
La suma sobre los doce factores vuelve a dar \(-27\). La prueba del
\cref{km:estabilidad-vecino} se aplica componente a
componente.
\end{proof}

Para tipar la recalibración, sea
\(h\in\operatorname{Aut}(C_W)\) monomial: una permutación \(\sigma\) de
las doce coordenadas seguida de multiplicadores
\(\epsilon_b\in\{1,2\}\). Definimos su elevación reticular
\begin{equation}
 \widetilde h
 :=P_\sigma\circ
 \bigoplus_{b=0}^{11}\mathsf R^{\nu_b},
 \qquad
 \nu_b=
 \begin{cases}
 0,&\epsilon_b=1,\\
 1,&\epsilon_b=2.
 \end{cases}
 \label{paper:elevacion-monomial}
\end{equation}
Su acción en el discriminante es exactamente \(h\). Si el cambio transporta
simultáneamente código, bandera, origen incidencial, vector \(v\), vecino y
marcos, entonces
\begin{equation}
 g_{hu}=\widetilde h\,g_u\,\widetilde h^{-1}.
 \label{paper:naturalidad-isometrias}
\end{equation}
La matriz ternaria \(h\) no actúa sin elevación sobre \(A_2^{12}\) ni sobre
Leech.


\section{Holonomía finita, configuración de Schläfli y simetría \texorpdfstring{\(E_6\)}{E6}}
\label{p12-83-c20-s6:chap:holonomia-schlafli-e6-nuevo}
\index{holonomía finita}
\index{Schläfli}
\index{E6@\(E_6\)}

La rama de Witt no agota las estructuras excepcionales accesibles desde
APP. La rama paralela de este capítulo parte de las dos polarizaciones
mixtas suma--producto y producto--suma, junto con una orientación de
plaqueta y una sección central declaradas. La curvatura calculada determina
la forma alternante; ésta define la extensión central. No se infiere una
extensión a partir de la sola cardinalidad 27.

\subsection{Curvaturas mixtas de APP}

Sobre \(T_9=(\mathbb Z/9\mathbb Z)^2\), con
\(S(x,y)=x+y\) y \(P(x,y)=xy\), sean
\[
 \Delta_xf=f(x+1,y)-f(x,y),\qquad
 \Delta_yf=f(x,y+1)-f(x,y).
\]
Las conexiones mixtas son
\[
 A^{SP}=(\Delta_xS,\Delta_yP)=(1,x),\qquad
 A^{PS}=(\Delta_xP,\Delta_yS)=(y,1).
\]
Su curvatura discreta es
\[
 F_A=\Delta_xA_y-\Delta_yA_x.
\]

\begin{proposition}[Curvaturas orientadas opuestas]
\label{p12-83-c20-s6:prop:curvaturas-app-e6-nuevo}
\[
 F_{A^{SP}}=1,\qquad F_{A^{PS}}=-1\pmod9.
\]
La circulación sobre un rectángulo orientado de lados \(r,s\) es,
respectivamente, \(rs\) y \(-rs\).
\end{proposition}

\begin{proof}
\(\Delta_xx-\Delta_y1=1\) y
\(\Delta_x1-\Delta_yy=-1\). Un rectángulo contiene \(rs\) plaquetas; al
sumar, las aristas interiores se cancelan y queda la circulación del borde.
\end{proof}

Para desplazamientos \(u=(a,b)\), \(v=(c,d)\), la clase alternante es
\begin{equation}
 \sigma(u,v)=ad-bc\pmod9.
 \label{p12-83-c20-s6:eq:forma-alternante-e6-nueva}
\end{equation}
La orientación fija el signo de \(\sigma\). La sección central empleada
abajo pertenece al calibre de esta realización; cambiarla por un coborde
produce una extensión isomorfa, no un nuevo invariante.

\subsection{Extensión central y reducción ternaria}

Definimos el grupo de Heisenberg nonádico
\begin{equation}
 \mathcal H_9=\{(a,b,z):a,b,z\in\mathbb Z/9\mathbb Z\},
 \quad
 (a,b,z)(c,d,w)=(a+c,b+d,z+w+bc).
 \label{p12-83-c20-s6:eq:heisenberg-9-nuevo}
\end{equation}
El conmutador de dos desplazamientos horizontales es
\[
 [(a,b,0),(c,d,0)]=(0,0,bc-ad),
\]
que registra \(-\sigma(u,v)\). Reduciendo cada coordenada módulo tres se
obtiene
\begin{equation}
 H_3=\{(a,b,z):a,b,z\in\mathbb F_3\},\qquad |H_3|=27,
 \label{p12-83-c20-s6:eq:heisenberg-3-nuevo}
\end{equation}
con la misma ley. El núcleo de
\(\mathcal H_9\to H_3\) es
\((3\mathbb Z/9\mathbb Z)^3\), de cardinal 27.

\subsection{Representación de Weyl--Heisenberg y fase central}

La extensión nonádica admite una realización unitaria explícita. Sea
\[
 \zeta_9:=\exp(2\pi i/9),
 \qquad
 \mathscr H_9^{\rm Sch}:=\mathbb C^9,
\]
con base \(\{|n\rangle:n\in\mathbb Z/9\mathbb Z\}\). Definimos
\begin{equation}
 X|n\rangle:=|n+1\rangle,
 \qquad
 Z|n\rangle:=\zeta_9^n|n\rangle.
 \label{p12-83-c20-s6:eq:weyl-desplazamiento-fase-u029}
\end{equation}
Se tiene
\[
 ZX=\zeta_9XZ
\]
y, por multiplicación,
\begin{equation}
 (X^aZ^b)(X^cZ^d)
 =\zeta_9^{bc}X^{a+c}Z^{b+d}.
 \label{p12-83-c20-s6:eq:cociclo-weyl-nonadico-u029}
\end{equation}

\begin{theorem}[Representación fiel de Weyl--Heisenberg]
\label{p12-83-c20-s6:thm:representacion-weyl-heisenberg-u029}
La aplicación
\begin{equation}
 \rho_{\rm WH}:\mathcal H_9\longrightarrow
 U(\mathscr H_9^{\rm Sch}),
 \qquad
 \rho_{\rm WH}(a,b,z):=\zeta_9^zX^aZ^b,
 \label{p12-83-c20-s6:eq:representacion-weyl-heisenberg-u029}
\end{equation}
es una representación unitaria fiel. Para
\(u=(a,b)\) y \(v=(c,d)\), su conmutador satisface
\begin{equation}
 [\rho_{\rm WH}(u,0),\rho_{\rm WH}(v,0)]
 =\zeta_9^{\,bc-ad}I
 =\zeta_9^{-\sigma(u,v)}I.
 \label{p12-83-c20-s6:eq:conmutador-weyl-holonomia-u029}
\end{equation}
\end{theorem}

\begin{proof}
La identidad \eqref{p12-83-c20-s6:eq:cociclo-weyl-nonadico-u029} reproduce el término
\(bc\) de \eqref{p12-83-c20-s6:eq:heisenberg-9-nuevo}; de ahí la propiedad de
homomorfismo. Si \(\zeta_9^zX^aZ^b=I\), la acción sobre la base obliga
sucesivamente a \(a=0\), \(b=0\) y \(z=0\). Esto prueba fidelidad. Comparar
los productos en los dos órdenes da
\(\zeta_9^{bc-ad}\), que es la fase alternante de
\eqref{p12-83-c20-s6:eq:forma-alternante-e6-nueva}.
\end{proof}

La representación distingue dos observables conjugados. Sean
\[
 Q_{\rm H}:=\operatorname{diag}(0,1,\ldots,8),
 \qquad
 F_{jk}:=\frac13\zeta_9^{jk},
 \qquad
 P_{\rm H}:=F^\dagger Q_{\rm H}F.
\]
Las bases propias de \(Q_{\rm H}\) y \(P_{\rm H}\) son mutuamente
insesgadas, pues \(|F_{jk}|^2=1/9\). Para todo vector unitario
\(\psi\in\mathscr H_9^{\rm Sch}\),
\begin{align}
 \operatorname{Var}_\psi(Q_{\rm H})
 \operatorname{Var}_\psi(P_{\rm H})
 &\ge
 \frac14\left|
 \langle\psi,[Q_{\rm H},P_{\rm H}]\psi\rangle
 \right|^2,
 \label{p12-83-c20-s6:eq:incertidumbre-weyl-u029}\\
 H_{Q_{\rm H}}(\psi)+H_{P_{\rm H}}(\psi)
 &\ge\log9.
 \label{p12-83-c20-s6:eq:incertidumbre-entropica-weyl-u029}
\end{align}
Aquí \(H_{Q_{\rm H}}\) y \(H_{P_{\rm H}}\) son las entropías de Shannon
de las distribuciones de Born en las respectivas bases propias,
con logaritmo natural y convención \(0\log0=0\).
Estos operadores son adimensionales. Una carta física de unidades es una
realización posterior y no interviene en la extensión central.

\subsection{Acción de 10 desplazamientos sobre el grafo de Cayley}

Para un funcional lineal
\(\ell(a,b)=pa+qb\), ponemos
\[
 D_\ell(a,b)=\left(a,b,\frac{ab}{2}+\ell(a,b)\right),
 \qquad Z=(0,0,1),
\]
donde \(2^{-1}=2\) en \(\mathbb F_3\). El conjunto inverso-cerrado
\begin{equation}
 S_\ell=
 \{D_\ell(v):v\in\mathbb F_3^2\setminus\{0\}\}
 \cup\{Z,Z^{-1}\}
 \label{p12-83-c20-s6:eq:conjunto-conexion-e6}
\end{equation}
tiene diez elementos y no contiene la identidad.

\begin{theorem}[Independencia del calibre lineal]
\label{p12-83-c20-s6:thm:independencia-calibre-e6}
Los nueve funcionales \(\ell\) producen grafos de Cayley isomorfos. Los
nueve conjuntos \(S_\ell\) forman una sola órbita bajo
\(\operatorname{Aut}(H_3)\).
\end{theorem}

\begin{proof}
La aplicación
\[
 \phi_\ell(a,b,z)=(a,b,z+\ell(a,b))
\]
es un automorfismo central y satisface
\(\phi_\ell(S_0)=S_\ell\). Para escribir la familia completa en las
coordenadas polarizadas de \cref{p12-83-c20-s6:eq:heisenberg-9-nuevo}, sea
\(c(v,w)=b c\) para \(v=(a,b)\), \(w=(c,d)\). Dado
\(M\in\operatorname{GL}(2,3)\), el defecto
\[
 B_M(v,w)=c(Mv,Mw)-\det(M)c(v,w)
\]
es simétrico, porque las partes alternantes de los dos términos coinciden.
Definimos la corrección cuadrática
\begin{equation}
 q_M(v)=2B_M(v,v),
 \qquad
 q_M(v+w)-q_M(v)-q_M(w)=B_M(v,w),
 \label{p12-83-c20-s6:eq:correccion-cuadratica-automorfismos-h3}
\end{equation}
donde \(2=2^{-1}\) en \(\mathbb F_3\). La familia completa de
automorfismos se escribe entonces
\begin{equation}
 (v,t)\longmapsto
 \bigl(Mv,\det(M)t+q_M(v)+\ell(v)\bigr),
 \qquad M\in\operatorname{GL}(2,3),\
 \ell\in(\mathbb F_3^2)^*.
 \label{p12-83-c20-s6:eq:automorfismos-h3-completos}
\end{equation}
La identidad polar de \cref{p12-83-c20-s6:eq:correccion-cuadratica-automorfismos-h3}
prueba directamente que el producto central se conserva. La familia tiene
\(48\cdot9=432\) elementos. La enumeración exhaustiva de los
\(\binom{13}{5}=1287\) subconjuntos inverso-cerrados de cardinal diez
encuentra exactamente esos nueve con los parámetros que se derivan a
continuación.
\end{proof}

Escribamos \(\underline S=\sum_{s\in S}s\) en
\(\mathbb Z[H_3]\). El conteo de productos ordenados da
\begin{equation}
 \underline S^{\,2}
 =10e+\underline S+5(\underline{H_3}-e-\underline S)
 =5\underline{H_3}+5e-4\underline S.
 \label{p12-83-c20-s6:eq:identidad-grupo-schlafli-nueva}
\end{equation}
En efecto, la identidad recibe diez factorizaciones; un elemento de
\(S\) recibe una; cada uno de los otros dieciséis recibe cinco.

\begin{theorem}[Grafo de las veintisiete rectas]
\label{p12-83-c20-s6:thm:grafo-27-rectas-nuevo}
La matriz de adyacencia \(A_\ell\) de
\(\Gamma_\ell=\operatorname{Cay}(H_3,S_\ell)\) satisface
\begin{equation}
 A_\ell^2=5I-4A_\ell+5J.
 \label{p12-83-c20-s6:eq:identidad-schlafli-nueva}
\end{equation}
Por tanto,
\[
 \Gamma_\ell=\operatorname{srg}(27,10,1,5),
 \qquad
 \operatorname{Spec}(A_\ell)=\{10^1,1^{20},(-5)^6\}.
\]
Es isomorfo al grafo de intersección de las veintisiete rectas de una
superficie cúbica lisa.
\end{theorem}

\begin{proof}
La representación regular de
\cref{p12-83-c20-s6:eq:identidad-grupo-schlafli-nueva} produce
\cref{p12-83-c20-s6:eq:identidad-schlafli-nueva}. Sus entradas dicen que todo vértice
tiene diez vecinos, que un par adyacente tiene un vecino común y un par no
adyacente tiene cinco. Sobre la recta constante, \(A_\ell\) actúa por 10;
en su complemento, \(J=0\) y
\((A_\ell-I)(A_\ell+5I)=0\). Traza cero fija las multiplicidades 20 y 6.

En el modelo de una cúbica obtenida al soplar seis puntos de
\(\mathbb P^2\), las 27 clases son
\[
 E_i,\qquad L_{ij}=H-E_i-E_j,\qquad
 Q_i=2H-\sum_{j\ne i}E_j.
\]
Con \(H^2=1\), \(E_i^2=-1\), \(H\cdot E_i=0\) y
\(E_i\cdot E_j=0\) para \(i\ne j\), la matriz de adyacencia definida
por intersección entre clases distintas tiene los mismos parámetros.
La matriz completa de intersección tiene diagonal \(-1\).
La identificación con la configuración clásica usa el
teorema de unicidad de esta realización de las 27 rectas
\cite{Manivel2005}, con el dominio y la acción allí especificados.
\end{proof}

La afirmación de isomorfía puede hacerse completamente efectiva mediante
una biyección calculada. El etiquetado no es absoluto: otra sección lineal
o una automorfía de la superficie produce otra tabla de la misma clase.
\begin{longtable}{@{}c@{\,}c@{\,}c@{\qquad}l@{}}
\caption{Bijección explícita entre \(H_3\) y las veintisiete clases de
rectas.}
\label{p12-83-c20-s6:tab:h3-27-rectas-u029}\\
\toprule
\(a\)&\(b\)&\(z\)&clase de recta\\
\midrule
\endfirsthead
\multicolumn{4}{@{}l}{\small\itshape Cuadro \thetable\ (continuación)}\\
\toprule
\(a\)&\(b\)&\(z\)&clase de recta\\
\midrule
\endhead
0&0&0&\(E_1\)\\
0&0&1&\(L_{12}\)\\
0&0&2&\(Q_2\)\\
0&1&0&\(L_{13}\)\\
0&1&1&\(Q_1\)\\
0&1&2&\(E_3\)\\
0&2&0&\(Q_3\)\\
0&2&1&\(E_2\)\\
0&2&2&\(L_{23}\)\\
1&0&0&\(L_{14}\)\\
1&0&1&\(L_{35}\)\\
1&0&2&\(L_{26}\)\\
1&1&0&\(L_{36}\)\\
1&1&1&\(L_{24}\)\\
1&1&2&\(L_{15}\)\\
1&2&0&\(L_{25}\)\\
1&2&1&\(L_{16}\)\\
1&2&2&\(L_{34}\)\\
2&0&0&\(Q_4\)\\
2&0&1&\(L_{46}\)\\
2&0&2&\(E_6\)\\
2&1&0&\(E_5\)\\
2&1&1&\(Q_6\)\\
2&1&2&\(L_{56}\)\\
2&2&0&\(L_{45}\)\\
2&2&1&\(E_4\)\\
2&2&2&\(Q_5\)\\
\bottomrule
\end{longtable}

Como control local, los diez vecinos de \(e=(0,0,0)\) son los elementos de
\(S_0\). La tabla los transforma en las diez clases que cortan a \(E_1\):
las cinco \(L_{1j}\) y las cinco \(Q_j\), con \(2\le j\le6\). La
verificación completa compara las 729 entradas de ambas matrices de
adyacencia.

El complemento \(B_\ell=J-I-A_\ell\) satisface
\begin{equation}
 B_\ell^2=8I+2B_\ell+8J,
 \qquad
 \operatorname{srg}(27,16,10,8),
 \qquad
 \operatorname{Spec}(B_\ell)=\{16^1,4^6,(-2)^{20}\}.
 \label{p12-83-c20-s6:eq:complemento-schlafli}
\end{equation}
Éste es el grafo de Schläfli en la convención de adyacencia complementaria;
\(A_\ell\) es el grafo de intersección de parámetros
\((27,10,1,5)\).

Cada arista pertenece a un triángulo único, pues \(\lambda=1\). El número
de triángulos es
\[
 \frac{27\cdot10}{6}=45.
\]
El grupo completo de automorfismos tiene orden
\[
 27\cdot1920=51840
\]
y, por su acción sobre la configuración de rectas, se identifica con
\(W(E_6)\), no sólo por igualdad de órdenes \cite{Manivel2005}.

\subsection{El retículo de raíces que realiza \texorpdfstring{\(E_6\)}{E6}}

Para hacer explícito el objeto sobre el que actúa ese grupo, sea
\begin{equation}
 \operatorname{Pic}(X)
 =\mathbb ZH\oplus\bigoplus_{i=1}^6\mathbb ZE_i,
 \qquad
 H^2=1,\qquad E_i^2=-1,\qquad H\cdot E_i=E_i\cdot E_j=0\ (i\ne j),
 \label{p12-83-c20-s6:eq:picard-cubica-e6}
\end{equation}
el retículo de Picard de la explosión de seis puntos generales de
\(\mathbb P^2\). Su clase canónica es
\[
 K_X=-3H+E_1+\cdots+E_6.
\]
El complemento ortogonal
\begin{equation}
 L_{E_6}=K_X^\perp\subset\operatorname{Pic}(X),
 \qquad (x,y)_{E_6}:=-x\cdot y,
 \label{p12-83-c20-s6:eq:reticulo-raices-e6}
\end{equation}
es positivo definido de rango seis. Una base de raíces simples es
\begin{equation}
 \begin{aligned}
 \alpha_1&=E_1-E_2,& \alpha_2&=E_2-E_3,&
 \alpha_3&=E_3-E_4,\\
 \alpha_4&=E_4-E_5,& \alpha_5&=E_5-E_6,&
 \alpha_6&=H-E_1-E_2-E_3.
 \end{aligned}
 \label{p12-83-c20-s6:eq:raices-simples-e6}
\end{equation}
Cada \(\alpha_i\) es ortogonal a \(K_X\), tiene norma dos y sus productos
son \(-1\) exactamente para los pares adyacentes del diagrama
\(E_6\): la cadena \(1-2-3-4-5\), con el nodo \(6\) unido al nodo \(3\).
Por tanto, su matriz de Gram es la matriz de Cartan de \(E_6\).

\begin{proposition}[Las setenta y dos raíces y sus reflexiones]
\label{p12-83-c20-s6:prop:setenta-dos-raices-e6}
Las raíces de \(L_{E_6}\) son
\begin{equation}
 \begin{aligned}
 &E_i-E_j &&(i\ne j),\\
 &\pm(H-E_i-E_j-E_k) &&(1\le i<j<k\le6),\\
 &\pm(2H-E_1-\cdots-E_6),
 \end{aligned}
 \label{p12-83-c20-s6:eq:enumeracion-raices-e6}
\end{equation}
en total \(30+40+2=72\). Para cada raíz \(\alpha\),
\begin{equation}
 s_\alpha(x)=x+(x\cdot\alpha)\alpha
 \label{p12-83-c20-s6:eq:reflexion-raiz-e6}
\end{equation}
preserva la intersección y \(K_X\). Las seis reflexiones simples generan
\(W(E_6)\) y permutan fielmente las 27 clases
\(E_i,L_{ij},Q_i\) de \cref{p12-83-c20-s6:thm:grafo-27-rectas-nuevo}.
\end{proposition}

\begin{proof}
Sustituir cada clase de \cref{p12-83-c20-s6:eq:enumeracion-raices-e6} en
\cref{p12-83-c20-s6:eq:picard-cubica-e6} da \(K_X\cdot\alpha=0\) y
\(\alpha^2=-2\). La cuenta de clases es la indicada. La fórmula de reflexión
es la reflexión ortogonal para la forma \(-(\cdot)\); preserva \(K_X\) porque
\(K_X\cdot\alpha=0\). El sistema simple de
\cref{p12-83-c20-s6:eq:raices-simples-e6} genera todas las raíces y su grupo de Coxeter es
\(W(E_6)\). Su acción sobre las 27 clases de curvas excepcionales conserva la
intersección, de modo que coincide con la acción sobre el grafo ya
construido \cite{Manivel2005}.
\end{proof}

\subsection{Estratificación de las 729 palabras}

Escribimos una palabra de seis trits como
\[
 w=(\ell_1,h_1,\ell_2,h_2,\ell_3,h_3),
\]
y definimos dos elementos de \(H_3\),
\[
 L(w)=(\ell_1,\ell_2,\ell_3),\qquad
 H(w)=(h_1,h_2,h_3),
\]
usando la ley central en sus productos. Según
\(L(w)^{-1}H(w)\), las 729 palabras se dividen en
\begin{align*}
 \mathcal D_\ell&=\{w:L(w)^{-1}H(w)=e\},\\
 \mathcal I_\ell&=\{w:L(w)^{-1}H(w)\in S_\ell\},\\
 \mathcal S_\ell&=\{w:L(w)^{-1}H(w)\notin\{e\}\cup S_\ell\}.
\end{align*}

\begin{theorem}[Partición Hensel--Schläfli]
\label{p12-83-c20-s6:thm:particion-hensel-schlafli-nueva}
\begin{equation}
 729=27+270+432,
 \label{p12-83-c20-s6:eq:particion-729-e6-nueva}
\end{equation}
correspondientes a diagonal, incidencia y alabeamiento.
\end{theorem}

\begin{proof}
Para cada uno de los 27 valores de \(L(w)\), hay un único valor de
\(H(w)\) con diferencia identidad, diez con diferencia en \(S_\ell\) y
dieciséis restantes. Se multiplica \(27\) por \(1,10,16\).
\end{proof}

El cardinal \(432\) del estrato alabeado coincide numéricamente con la
multiplicidad de la región transversal de \(\pi\) en
\cref{pi-estrella:figura}. No se identifica por ello ambos objetos: aquí
se cuentan palabras clasificadas por \(L(w)^{-1}H(w)\), mientras allí se
cuentan preimágenes de un registro \(U_6\). Sin un entrelazador explícito y
equivariante entre esos dominios, \(432=432\) es sólo un control cardinal.

\subsection{Necesidad de la fase central}

La enumeración de todos los subconjuntos inverso-cerrados de cardinal diez
en los cinco grupos de orden 27 da
\[
\begin{array}{c|c|c}
\text{grupo}&\text{abeliano}&
\#\operatorname{srg}(27,10,1,5)\\ \hline
C_{27}&\text{sí}&0\\
C_9\times C_3&\text{sí}&0\\
C_3^3&\text{sí}&0\\
C_9\rtimes C_3&\text{no}&9\\
H_3&\text{no}&9.
\end{array}
\]
Así, conservar 27 etiquetas pero abelianizar la composición destruye la
configuración. No bastan la cardinalidad ni el grado.

Si \(\mathcal T\) es el conjunto de 45 triángulos, el cúbico
\[
 \mathcal C_\Gamma(x_1,\ldots,x_{27})
 =\sum_{\{i,j,k\}\in\mathcal T}x_ix_jx_k
\]
recupera la adyacencia, porque cada arista está en un triángulo único. Su
estabilizador permutacional es exactamente
\[
 \operatorname{Aut}(\Gamma_\ell)\cong W(E_6).
\]

Para el diagrama siguiente se escriben
\(\Gamma_{27}=\Gamma_\ell\) y \(\mathcal T_{45}=\mathcal T\).
\begin{figure}[htbp]
\centering
\[
 \begin{gathered}
 \text{curvaturas APP }\pm1
 \longrightarrow\sigma
 \longrightarrow\mathcal H_9
 \longrightarrow H_3,\\[3pt]
 H_3
 \longrightarrow\Gamma_{27}
 \longrightarrow\mathcal T_{45}
 \longrightarrow W(E_6).
 \end{gathered}
\]
\caption{Cadena causal de la rama de holonomía mixta. Es paralela a
\(C_W\to W_{12}\to M_{12}\) y no un eslabón de ella.}
\label{p12-83-c20-s6:fig:cadena-holonomia-e6-u029}
\end{figure}

\begin{criterion}[Criterios de refutación de la rama \(E_6\)]
\label{p12-83-c20-s6:crit:refutacion-rama-e6-u029}
La rama queda refutada si las curvaturas no son opuestas, si el cociclo
central se borra, si la corrección cuadrática
\eqref{p12-83-c20-s6:eq:correccion-cuadratica-automorfismos-h3} no conserva el producto,
si \eqref{p12-83-c20-s6:eq:identidad-grupo-schlafli-nueva} no produce los coeficientes
\((10,1,5)\), si un grupo abeliano produce la misma configuración, si la
partición no suma 729, si la tabla
\ref{p12-83-c20-s6:tab:h3-27-rectas-u029} no entrelaza las matrices de adyacencia o si
\(W(E_6)\) se identifica sólo por su orden y no por su acción sobre las
veintisiete clases y las 72 raíces.
\end{criterion}

% Procedencia: integral 2249, cap. 23, pp. 601--608; propietario base_c21_sin_encabezado.tex.
% Recuperación y desarrollo de pruebas. Requiere el núcleo común APP--TRIT--TPK.
\section{Realización hilbertiana de la fase nonádica y conservación de la unidad}
\label{iv:hilbert}

La construcción de Holografía Modular Triádica proporciona primero los
estados, sus transiciones y los mapas de restricción. La realización
hilbertiana añade a sus coordenadas discretas una estructura lineal y un
producto interno. Este orden determina qué representa cada operador:
las etiquetas de la base proceden de la construcción nonádica; las
operaciones lineales actúan después sobre esas etiquetas. La memoria
genealógica permanece en el estado del TPK y acompaña a las proyecciones
que se utilicen para representarlo.

\subsection{Fibra de fase, base etiquetada y producto interno}

Para una profundidad entera \(d\geq1\), sea
\begin{equation}
 \mathcal C_{9,d}=(\ZZ/9\ZZ)^d,
 \qquad \mathcal H_d=\ell^2(\mathcal C_{9,d}),
 \qquad \mathcal A_d=B(\mathcal H_d).
 \label{iv:eq:hilbert-finito}
\end{equation}
El producto interno es
\(\langle f,g\rangle=\sum_x\overline{f(x)}g(x)\). Las funciones
\(\delta_x\) forman una base ortonormal etiquetada por las coordenadas
de fase. Así, \(\dim\mathcal H_d=9^d\) y
\(\mathcal A_d\cong M_{9^d}(\CC)\). El álgebra matricial representa
esta fibra finita; las rutas, los cocientes y la memoria que una
proyección de fase omita siguen perteneciendo al estado de partida.

La factorización de coordenadas induce el unitario
\begin{equation}
 \mathcal H_{d+1}\longrightarrow\mathcal H_d\otimes\CC^9,
 \qquad \delta_{(x,r)}\longmapsto\delta_x\otimes\delta_r.
 \label{iv:eq:factorizacion-hilbert}
\end{equation}
Definimos la traza normalizada
\(\tau_d(a)=9^{-d}\operatorname{Tr}(a)\). La cuenta finita y la
normalización se distinguen: \(\operatorname{Tr}(I)=9^d\), mientras
que \(\tau_d(I)=1\). Asimismo, el grupo \((\ZZ/9\ZZ)^3\) y el
espacio vectorial \(\FF_3^6\) tienen ambos \(729\) elementos, pero
conservan leyes algebraicas diferentes. La realización de una fibra de
fase no sustituye la codificación ternaria de las emisiones APP.

\subsection{Traslación y fase: representación finita de Weyl}

En la realización compleja posterior a las construcciones de
\(\pi\) y de la estructura \(\mathrm i^2=-1\), fijamos
\(\omega=\exp(2\pi\mathrm i/9)\). Sobre \(\mathcal H_1\), sean
\begin{equation}
 (Xf)(r)=f(r-1),\qquad (Zf)(r)=\omega^r f(r).
 \label{iv:eq:weyl-operadores}
\end{equation}
La primera operación desplaza la coordenada de fase; la segunda la
evalúa mediante un carácter. Ambas son unitarias y satisfacen
\begin{equation}
 X^9=Z^9=I,\qquad ZX=\omega XZ.
 \label{iv:eq:weyl-relacion}
\end{equation}
En efecto, \((ZXf)(r)=\omega^r f(r-1)\), mientras que
\((XZf)(r)=\omega^{r-1}f(r-1)\). El carácter aparece aquí como
realización de la fase ya construida y no selecciona los registros
numéricos del TPK.

% RESULTADO_RECUPERADO: 14_numero_heisenberg_y_d108.tex:741--788.
% Prueba algebraica desarrollada con la convención X^a Z^b ya fijada aquí.
\subsubsection{Área orientada y extensión central de las traslaciones}
\label{iv:weyl-area-extension}

El soporte de traslaciones es \(X_9=(\ZZ/9\ZZ)^2\).
La curvatura mixta de APP construida en la
sección~\ref{iv:extension-curvatura-mixta} asigna \(+1\) a la
plaqueta orientada de la pareja suma--producto y \(-1\) a la
pareja de orden contrario. Su extensión bilineal a dos desplazamientos
\(u=(a,b)\) y \(v=(c,d)\) es la forma alternante
\begin{equation}
 \sigma_{\mathrm{ar}}(u,v)=ad-bc\quad\text{en }\ZZ/9\ZZ.
 \label{iv:eq:weyl-area-orientada}
\end{equation}
En efecto, la alternancia anula la evaluación sobre dos desplazamientos
del mismo eje. Los valores sobre la base orientada son
\(\sigma_{\mathrm{ar}}((1,0),(0,1))=1\) y
\(\sigma_{\mathrm{ar}}((0,1),(1,0))=-1\); distribuir ambos
argumentos da \eqref{iv:eq:weyl-area-orientada}. La misma fórmula es
la suma de plaquetas de cualquier cadena celular orientada que realice
ese paralelogramo: las aristas interiores se cancelan. Modificar un
representante entero por un múltiplo de nueve conserva su publicación
en \(\ZZ/9\ZZ\).

El orden \(X^aZ^b\) de los operadores ya definidos fija una sección
de la extensión central. Denotemos su factor de composición por
\begin{equation}
 c_{\mathrm W}(u,v)=bc\quad\text{en }\ZZ/9\ZZ.
 \label{iv:eq:weyl-cociclo-polarizado}
\end{equation}
Este cociclo conserva el orden de las dos operaciones, mientras
\(\sigma_{\mathrm{ar}}\) conserva la orientación del paralelogramo.

\begin{proposition}[Extensión central y signo del conmutador]
\label{iv:prop:heisenberg-orientado}
El conjunto \(\mathfrak H_9=(\ZZ/9\ZZ)^3\), con producto
\begin{equation}
 (a,b,t)(c,d,s)=(a+c,b+d,t+s+bc),
 \label{iv:eq:heisenberg-ley}
\end{equation}
es un grupo de orden \(729\). La aplicación
\[
 \varrho(a,b,t)=\omega^tX^aZ^b
\]
es una representación unitaria fiel sobre \(\mathcal H_1\).
Con el convenio \([g,h]=ghg^{-1}h^{-1}\), se tiene
\begin{equation}
 [\varrho(a,b,0),\varrho(c,d,0)]
 =\omega^{bc-ad}I
 =\omega^{-\sigma_{\mathrm{ar}}((a,b),(c,d))}I.
 \label{iv:eq:weyl-conmutador-orientado}
\end{equation}
El centro y el grupo de traslaciones forman la sucesión exacta
\[
 0\longrightarrow\ZZ/9\ZZ
 \xrightarrow{\ t\mapsto(0,0,t)\ }\mathfrak H_9
 \xrightarrow{\ (a,b,t)\mapsto(a,b)\ }X_9
 \longrightarrow0.
\]
\end{proposition}
\begin{proof}
Para \(w=(e,f)\), la identidad de cociclo se escribe
\[
 c_{\mathrm W}(u,v)+c_{\mathrm W}(u+v,w)
 =bc+(b+d)e
 =de+b(c+e)
 =c_{\mathrm W}(v,w)+c_{\mathrm W}(u,v+w).
\]
Prueba la asociatividad de \eqref{iv:eq:heisenberg-ley}; la identidad
es \((0,0,0)\) y la inversa de \((a,b,t)\) es
\((-a,-b,-t+ab)\). Los tres residuos son libres, de donde el orden
\(9^3=729\). De \(ZX=\omega XZ\) resulta
\[
 (X^aZ^b)(X^cZ^d)=\omega^{bc}X^{a+c}Z^{b+d},
\]
que demuestra la propiedad de representación. Si
\(\omega^tX^aZ^b=I\), su acción sobre \(\delta_r\) obliga primero
a \(a=0\). La evaluación en \(r=0\) da \(t=0\), y la
evaluación en \(r=1\) da \(b=0\), pues \(\omega\) tiene orden
nueve. Por ello la representación es fiel. La comparación de los
productos en los dos órdenes da el factor \(bc-ad\).
Un elemento central debe satisfacer \(bc-ad=0\) para todo \((c,d)\).
Tomando sucesivamente \((c,d)=(1,0),(0,1)\), se obtiene
\(b=a=0\). Esto identifica el centro y prueba la exactitud.
\end{proof}

La fase de Wilson para la orientación
\eqref{iv:eq:weyl-area-orientada} es
\(\omega^{\sigma_{\mathrm{ar}}(u,v)}\); el conmutador de la
sección ordenada \(X^aZ^b\) publica su inversa. En particular,
\(c_{\mathrm W}(u,v)-c_{\mathrm W}(v,u)
=-\sigma_{\mathrm{ar}}(u,v)\).
Esta identidad relaciona dos convenciones determinadas, sin identificar
el cociclo polarizado con la forma alternante. La extensión pertenece
al grupo de traslaciones del soporte APP. Su representación actúa sobre
la fibra de fase; el registro de la trayectoria y los acarreos conservan
su residencia en el estado TPK que produjo esa fibra. El orden \(729\)
del grupo y el cardinal \(729\) del ambiente ternario conservan los
tipos distintos establecidos en \eqref{iv:eq:hilbert-finito}.
El centro \(\ZZ/9\ZZ\) de esta extensión registra la fase
del conmutador. La memoria \(C_{12}\) de la sucesión
\(0\to C_{12}\to C_{108}\to C_9\to0\), construida en la
proposición~\ref{prop:k-extension}, pertenece al calendario
dodecafásico. Cada extensión conserva su ley de composición y
su proyección propias.

\begin{proposition}[Álgebra generada por fase y traslación]
\label{iv:prop:weyl-completo}
Los \(81\) operadores \(X^aZ^b\), con \(0\leq a,b<9\), forman
una base de \(\mathcal A_1=M_9(\CC)\). En particular, la
representación de \eqref{iv:eq:weyl-operadores} es irreducible.
\end{proposition}
\begin{proof}
Si \(a\neq0\pmod9\), la matriz \(X^aZ^b\) carece de entradas
diagonales. Si \(a=0\), su traza es
\(\sum_{r=0}^8\omega^{br}\), igual a \(9\) para \(b=0\) y a
\(0\) en los demás casos. Por \eqref{iv:eq:weyl-relacion}, el
producto \((X^aZ^b)^*X^{a'}Z^{b'}\) es un factor de módulo uno
por \(X^{a'-a}Z^{b'-b}\). Por tanto,
\[
 \tau_1\bigl((X^aZ^b)^*X^{a'}Z^{b'}\bigr)
 =\delta_{a,a'}\delta_{b,b'}.
\]
Son \(81\) matrices linealmente independientes en un espacio de
dimensión \(81\). Generan toda el álgebra matricial, cuyo conmutante
está formado por los escalares; de ello resulta la irreducibilidad.
\end{proof}

\begin{theorem}[Unicidad con carácter central primitivo fijado]
\label{iv:thm:weyl-unicidad}
Toda representación unitaria irreducible del grupo presentado por
\(X^9=Z^9=c^9=I\), \(c\) central y \(ZX=cXZ\), en la que
\(c\) actúa como \(\omega I\), es unitariamente equivalente a
\eqref{iv:eq:weyl-operadores}.
\end{theorem}
\begin{proof}
Sea \(v\neq0\) un autovector de \(Z\), de autovalor \(\lambda\).
Como \(Z^9=I\), \(\lambda\) es una potencia de \(\omega\).
Los vectores \(X^kv\), \(0\leq k<9\), tienen autovalores
\(\omega^k\lambda\), distintos porque el carácter es primitivo.
Son ortogonales, y su espacio generado es invariante por \(X\),
\(Z\) y el centro. La irreducibilidad obliga a que sea todo el
espacio. Reordenando la base para comenzar en el autovalor \(1\),
\(Z\) es diagonal con entradas \(1,\omega,\ldots,\omega^8\),
y \(X\) es la traslación cíclica. La normalización de la base
produce la equivalencia unitaria afirmada.
\end{proof}

La primitividad del carácter y la irreducibilidad forman parte del
enunciado. A profundidad \(d\), los operadores que actúan en factores
distintos conmutan, y los productos tensoriales de las bases anteriores
generan \(\mathcal A_d\). La estructura de fase queda así
representada a cada nivel sin identificar el retorno de fase con el
retorno del estado genealógico completo.

\subsection{Dos elevaciones y dos normalizaciones diferentes}

La factorización \eqref{iv:eq:factorizacion-hilbert} permite conservar
toda la fibra nueva o fijar una de sus etiquetas. Para
\(p_j=|\delta_j\rangle\langle\delta_j|\), definimos
\begin{equation}
 \iota_d(a)=a\otimes I_9,
 \qquad \jmath_{d,j}(a)=a\otimes p_j.
 \label{iv:eq:elevaciones}
\end{equation}
\begin{proposition}[Conservación de unidad y traza]
\label{iv:prop:unidad-traza}
La inclusión \(\iota_d\) es unital y conserva la traza normalizada.
La inclusión \(\jmath_{d,j}\) es una inclusión en una esquina y
satisface
\[
 \tau_{d+1}(\iota_d(a))=\tau_d(a),\quad \iota_d(I)=I,
 \qquad
 \tau_{d+1}(\jmath_{d,j}(a))=\frac{\tau_d(a)}9,
 \quad \jmath_{d,j}(I)=I\otimes p_j.
\]
\end{proposition}
\begin{proof}
La traza de un producto tensorial es el producto de trazas.
Como \(\operatorname{Tr}(I_9)=9\) y
\(\operatorname{Tr}(p_j)=1\), dividir por \(9^{d+1}\)
produce las dos fórmulas. Las imágenes de la identidad se leen
directamente en \eqref{iv:eq:elevaciones}.
\end{proof}

La inclusión unital conserva el conjunto de las \(9\) fases y su
normalización global. La esquina conserva una continuación marcada y
su proyector de soporte. Ambas operaciones son legítimas, pero
transmiten datos distintos a la siguiente profundidad.

\subsection{Cierre inductivo y realización tracial}

\begin{theorem}[Álgebra uniformemente hiperfinita nonádica]
\label{iv:thm:uhf}
El cierre normado de la unión mediante las inclusiones \(\iota_d\)
es
\[
 \mathcal A_{9^\infty}
 =\varinjlim(\mathcal A_d,\iota_d)
 \cong\bigotimes_{n=1}^{\infty}M_9(\CC).
\]
Es un álgebra UHF de tipo \(9^\infty=3^\infty\), con una única
traza normalizada \(\tau\), cuya restricción a \(\mathcal A_d\)
es \(\tau_d\).
\end{theorem}
\begin{proof}
Las inclusiones son inyectivas y unitales. La clausura de la unión
creciente de álgebras matriciales es precisamente la construcción UHF.
La compatibilidad de la proposición \ref{iv:prop:unidad-traza}
define \(\tau\) sobre la unión; positividad y norma uno permiten
extenderla al cierre. Toda traza normalizada del límite se restringe
a la única traza normalizada de cada álgebra matricial. Como la unión
es densa, esas restricciones determinan la traza de manera única.
\end{proof}

Sea \(\mathcal H_\tau\) la completación de
\(\mathcal A_{9^\infty}\) para
\(\langle a,b\rangle_\tau=\tau(a^*b)\), y sea \(\pi_\tau\)
su representación por multiplicación izquierda. El vector definido
por la identidad se denota \(\Omega_\tau\).

\begin{theorem}[Factor tracial del refinamiento completo]
\label{iv:thm:factor-tracial}
La clausura débil
\(\mathcal R_9=\pi_\tau(\mathcal A_{9^\infty})''\)
es un factor hiperfinito de tipo \(II_1\). Las inclusiones
nonádicas conservan en él la unidad y la traza normalizada.
\end{theorem}
\begin{proof}
La traza producto se extiende a una traza normal fiel y finita
en la representación tracial. Para ver que el centro es trivial,
sea \(E_d\) la expectativa tracial sobre \(\pi_\tau(\mathcal A_d)\).
En un tensor de longitud \(n>d\), toma la traza de los últimos
\(n-d\) factores. Estas aplicaciones se prolongan por continuidad
a las proyecciones ortogonales correspondientes en \(L^2\).
Si \(z\) es central, \(E_d(z)\) conmuta con
\(\pi_\tau(\mathcal A_d)\), por la bimodularidad de \(E_d\).
Es por ello escalar e igual a \(\tau(z)I\). La densidad de la
unión matricial implica \(E_d(z)\to z\) en \(L^2\), de donde
\(z=\tau(z)I\). El álgebra es, pues, un factor finito.
Contiene unitariamente matrices de tamaños \(9^d\) sin cota,
por lo que no es un factor matricial finito. Es de tipo \(II_1\).
Es hiperfinita porque esa misma unión de matrices es débilmente
densa. La conservación de unidad y traza procede de cada inclusión
finita y se mantiene al tomar las clausuras.
\end{proof}

\subsection{Rangos enteros y dimensión tracial continua}

En \(\mathcal A_d\), una proyección \(P\) tiene dimensión
normalizada \(\tau_d(P)=\operatorname{rank}(P)/9^d\).
Los rangos siguen siendo enteros a toda profundidad finita. Sus
normalizaciones compatibles permiten el siguiente paso al límite.

\begin{proposition}[Realización de cada dimensión tracial]
\label{iv:prop:dimension-tracial}
Para cada \(t\in[0,1]\) existe una proyección
\(P_t\in\mathcal R_9\) con \(\tau(P_t)=t\).
\end{proposition}
\begin{proof}
Sea \(m_d=\lfloor9^dt\rfloor\). Se tiene
\(0\leq m_{d+1}-9m_d\leq8\). Elegida una proyección diagonal
de rango \(m_d\), su inclusión tiene rango \(9m_d\); se
añaden \(m_{d+1}-9m_d\) posiciones diagonales en su complemento.
Se obtiene así una sucesión creciente de proyecciones dentro de
\(\mathcal R_9\). Su límite fuerte es una proyección \(P_t\),
y la normalidad de la traza da
\[
 \tau(P_t)=\lim_{d\to\infty}\frac{\lfloor9^dt\rfloor}{9^d}=t.
\]
\end{proof}

Aquí \(t\) designa la coordenada de la dimensión que se representa
en el cierre, no una entrada del generador APP. La afirmación
describe la amplitud de la realización tracial resultante. Para
compararla con un observable determinado sigue siendo necesario
identificar la proyección y la información genealógica que conserva.

Por contraste, las isometrías
\(V_{d,j}\xi=\xi\otimes\delta_j\) realizan
\(\jmath_{d,j}(a)=V_{d,j}aV_{d,j}^*\). En el límite de una
ruta marcada, los espacios finitos son densos. Las esquinas
aproximan todos los operadores de rango finito; su cierre normado
es \(\mathcal K(\mathcal H_\infty)\), y su bicommutante es
\(B(\mathcal H_\infty)\), de tipo \(I_\infty\). Esta
diferencia conserva el contenido operatorio de las dos elecciones
de \eqref{iv:eq:elevaciones}: la selección de una continuación y
la prolongación unital de toda la fibra tienen límites diferentes.

La realización hilbertiana, sus operadores de fase y el cierre
tracial quedan construidos en esta sección, con sus elecciones de
inclusión explícitas. Constituyen el
soporte operatorio que se utilizará después para el intercambio de
coordenadas y para el Hamiltoniano compacto.

\section{Núcleos genealógicos de caminos: composición finita, realización unitaria y límite cilíndrico}
\label{sec:integral-genealogica-caminos-rev2}
\index{Feynman!integral genealógica de caminos}
\index{camino!categoría enriquecida}
\index{integral!cilíndrica}

La composición matricial que se demuestra en esta sección establece exactamente que cada entrada de
\(U^R\) es la suma de los pesos ordenados de las rutas de longitud \(R\).
El paso siguiente conserva la información que una ruta visible no publica y
organiza los refinamientos finitos en un espacio compacto de genealogías. De
este modo se distinguen tres resultados: la identidad combinatoria finita, la
integral cilíndrica sobre el límite de caminos y la realización de esa
integral como propagador de un sistema físico concreto.

\subsection{Categoría de rutas y peso functorial}

Sea \(\mathsf P_{\rm TPK}\) la categoría libre generada por las transiciones
admisibles del TPK. Sus objetos son estados enriquecidos de frontera y sus
morfismos son secuencias composables de transiciones. La composición es la
concatenación y el camino vacío es la identidad. Sea \(\mathsf P_X\)
la categoría de rutas visibles obtenida al aplicar el lector a los
objetos y transiciones; sus composiciones son las concatenaciones visibles.
La aplicación
\[
 q_{\rm vis}:\mathsf P_{\rm TPK}\longrightarrow\mathsf P_X
\]
publica la ruta observable y conserva en su fibra las diferencias de hoja,
orientación, acarreo, bloque, frontera y memoria.

\begin{definition}[Peso genealógico]
\label{def:peso-genealogico-caminos-rev2}
Sea \(\mathcal B\) un álgebra real normada, posiblemente no conmutativa, y
sea \(\mathbf B_{\mathcal B}\) la categoría de un solo objeto \(\ast\) cuyo
monoide de endomorfismos es
\(\operatorname{End}(\ast)=(\mathcal B,\cdot,1)\). Un peso genealógico es un funtor
\[
 W:\mathsf P_{\rm TPK}\longrightarrow\mathbf B_{\mathcal B}
\]
tal que
\[
 W(\operatorname{id}_x)=1,
 \qquad
 W(\gamma_2\circ\gamma_1)=W(\gamma_2)W(\gamma_1).
\]
El orden del producto coincide con el orden causal de la ruta.
\end{definition}

Para estados enriquecidos \(\widetilde i,\widetilde f\), el kernel a
profundidad \(R\) es
\begin{equation}
 \mathcal K_R(\widetilde f,\widetilde i)
 =\sum_{\substack{
   \gamma\in\mathsf P_{\rm TPK}(\widetilde i,\widetilde f)\\
   |\gamma|=R}}
 W(\gamma).
 \label{eq:kernel-genealogico-finito-rev2}
\end{equation}
Dos rutas con la misma palabra visible permanecen como sumandos diferentes
cuando pertenecen a fibras genealógicas distintas.

\begin{proposition}[Composición por corte enriquecido]
\label{prop:composicion-kernel-enriquecido-rev2}
Si toda ruta de longitud \(R+S\) posee un corte único después de \(R\)
transiciones, entonces
\[
 \mathcal K_{R+S}(\widetilde f,\widetilde i)
 =\sum_{\widetilde x}
 \mathcal K_S(\widetilde f,\widetilde x)
 \mathcal K_R(\widetilde x,\widetilde i).
\]
La fórmula conserva el orden causal de los factores cuando \(\mathcal B\) no
es conmutativa.
\end{proposition}

\begin{proof}
Se particiona el conjunto de rutas completas por el estado enriquecido
\(\widetilde x\) alcanzado en el corte. La factorización functorial del peso
convierte la suma sobre cada clase en el producto de los kernels parciales.
La unicidad del corte garantiza que cada ruta aparece exactamente una vez.
\end{proof}

\subsection{Construcción de la fase interna mediante el carácter real de acción sobre rutas genealógicas}

Sea \((E,J_E)\) un espacio euclídeo real provisto de un endomorfismo
ortogonal \(J_E\) tal que \(J_E^2=-I\). Si
\(\mathcal A_*:\mathsf P_{\rm TPK}\to\mathbb R\) es aditiva por
concatenación, la fase
\begin{equation}
 \Phi_J(\gamma)
 :=\exp[-J_E\mathcal A_*(\gamma)]
 =\cos\mathcal A_*(\gamma)\,I
  -\sin\mathcal A_*(\gamma)\,J_E
 \label{eq:fase-interna-caminos-rev2}
\end{equation}
es multiplicativa. Fijemos una sección positiva no nula
\(\hbar_{\rm int}\) de la línea de acción \(\mathcal L_S\)
construida anteriormente; puede utilizarse la sección
\(\hbar_{\rm ret}\) del perfil declarado.
Para una acción dimensional
\(S_{\rm int}(\gamma)=\hbar_{\rm int}\mathcal A_*(\gamma)\), esta
expresión realiza internamente la fase
\(\exp[-J_E\,S_{\rm int}(\gamma)/\hbar_{\rm int}]\). El símbolo escalar
\(\mathrm i\) aparece únicamente al elegir una carta compleja para el par
\((E,J_E)\).

\begin{proposition}[Functor de fase]
\label{prop:fase-interna-functor-rev2}
La aplicación \(\Phi_J\) satisface
\[
 \Phi_J(\operatorname{id}_x)=I,
 \qquad
 \Phi_J(\gamma_2\circ\gamma_1)
 =\Phi_J(\gamma_2)\Phi_J(\gamma_1).
\]
\end{proposition}

\begin{proof}
La acción del camino vacío es cero. Para rutas composables, la aditividad de
\(\mathcal A_*\) y la conmutación de los múltiplos de un mismo endomorfismo
\(J_E\) permiten separar la exponencial del valor sumado.
\end{proof}

\subsection{Realización unitaria sobre el registro cronológico
\texorpdfstring{\(27\times27\)}{27 x 27}}
\label{sec:realizacion-unitaria-registro-27-rev3}
\index{registro cronológico!realización unitaria}
\index{Fourier!bloques del registro cronológico}
\index{Hamiltoniano!efectivo}

Antes de pasar al límite cilíndrico, un núcleo finito admite una realización
unitaria explícita. Sea
\[
 \mathbb T_{27}^{\,2}:=(\mathbb Z/27\mathbb Z)^2,
 \qquad
 \mathcal D:=\{\mathrm N,\mathrm S,\mathrm E,\mathrm W\},
\]
y considérese el espacio real
\[
 \mathcal H_{27}^{\mathbb R,\mathrm{reg}}
 :=\ell^2(\mathbb T_{27}^{\,2};\mathbb R)
   \otimes\mathbb R^{\mathcal D}\otimes E.
\]
El endomorfismo \(J_E\) se extiende como la identidad sobre los factores de
posición y dirección, y se conserva para esa extensión la misma notación.

En la base ordenada
\((\mathrm N,\mathrm S,\mathrm E,\mathrm W)\), el mezclador direccional es
\begin{equation}
 C_4:=\frac12
 \begin{pmatrix}
  1&1&1&1\\
  1&1&-1&-1\\
  1&-1&1&-1\\
  1&-1&-1&1
 \end{pmatrix},
 \qquad C_4^{\mathsf T}C_4=I_4.
 \label{eq:mezclador-walsh-registro-27-rev3}
\end{equation}
Aquí \(C_4\) designa exclusivamente el mezclador de Walsh--Hadamard sobre la
fibra direccional; cuando actúa en
\(\mathcal H_{27}^{\mathbb R,\mathrm{reg}}\) se entiende
\(C_4\otimes I_E\).

Fíjese una carta traslacional en la que la acción de una arista elemental
\(\gamma_d\) dependa sólo de su dirección:
\[
 s_d:=S_{\rm int}(\gamma_d)\in\mathcal L_S,
 \qquad
 a_d:=\frac{s_d}{\hbar_{\rm int}}\in\mathbb R.
\]
El operador real diagonal de fase es
\[
 D_J:=\operatorname{diag}
 \left(
  e^{-J_Ea_{\mathrm N}},
  e^{-J_Ea_{\mathrm S}},
  e^{-J_Ea_{\mathrm E}},
  e^{-J_Ea_{\mathrm W}}
 \right).
\]
El desplazamiento condicionado transporta cada componente en su dirección:
\begin{align*}
 (S\psi)(x,y,\mathrm N)&=\psi(x,y-1,\mathrm N),&
 (S\psi)(x,y,\mathrm S)&=\psi(x,y+1,\mathrm S),\\
 (S\psi)(x,y,\mathrm E)&=\psi(x-1,y,\mathrm E),&
 (S\psi)(x,y,\mathrm W)&=\psi(x+1,y,\mathrm W),
\end{align*}
con todas las coordenadas tomadas módulo \(27\). El paso elemental y su
composición de doce pasos son
\begin{equation}
 U_{\rm micro}:=S D_J C_4,
 \qquad
 U_{12}^{\mathrm{reg}}:=U_{\rm micro}^{12}.
 \label{eq:paso-unitario-registro-27-rev3}
\end{equation}
El superíndice \(\mathrm{reg}\) conserva el tipo: este operador actúa sobre
el toro direccional \(27\times27\) y no se identifica con los demás objetos
dodecafásicos de la construcción ni con el endomorfismo temporal de orden \(108\).

\begin{proposition}[Unitariedad y suma exacta de caminos en el registro]
\label{prop:unitariedad-suma-registro-27-rev3}
El operador \(U_{\rm micro}\) es ortogonal sobre la realificación y unitario
en la carta compleja determinada por \(J_E\). En particular,
\(U_{12}^{\mathrm{reg}}\) es unitario. Si
\(\Xi:=\mathbb T_{27}^{\,2}\times\mathcal D\), entonces, para todo
\(n\geq0\) y \(\xi_i,\xi_f\in\Xi\),
\[
 \bigl(U_{\rm micro}^{n}\bigr)_{\xi_f,\xi_i}
 =
 \sum_{\substack{\gamma:\xi_i\to\xi_f\\|\gamma|=n}}
 W_{U_{\rm micro}}(\gamma),
\]
donde, para \(\gamma=(\xi_0,\ldots,\xi_n)\),
\[
 W_{U_{\rm micro}}(\gamma)
 :=
 (U_{\rm micro})_{\xi_n,\xi_{n-1}}
 \cdots
 (U_{\rm micro})_{\xi_1,\xi_0}.
\]
La suma recorre exactamente los caminos permitidos por los bloques no nulos
del operador.
\end{proposition}

\begin{proof}
Como \(J_E\) es ortogonal y \(J_E^2=-I\), se tiene
\(J_E^{\mathsf T}=-J_E\). Por tanto, cada bloque
\(e^{-J_Ea_d}\) es una rotación ortogonal y \(D_J\) es ortogonal.
El operador \(S\) permuta una base ortonormal y
\eqref{eq:mezclador-walsh-registro-27-rev3} demuestra que \(C_4\) es
ortogonal. En consecuencia,
\[
 U_{\rm micro}^{\mathsf T}U_{\rm micro}
 =C_4^{\mathsf T}D_J^{\mathsf T}S^{\mathsf T}
   S D_JC_4
 =I.
\]
Los tres factores conmutan con la extensión de \(J_E\); así,
\(U_{\rm micro}\) es complejo lineal en la carta
\(J_E\leftrightarrow\mathrm i\), y su ortogonalidad real equivale a
unitariedad. Toda potencia, en particular
\(U_{12}^{\mathrm{reg}}\), conserva esa propiedad.

Para \(n\geq1\), la multiplicación por bloques da
\[
 \bigl(U_{\rm micro}^{n}\bigr)_{\xi_f,\xi_i}
 =
 \sum_{\xi_1,\ldots,\xi_{n-1}}
 (U_{\rm micro})_{\xi_f,\xi_{n-1}}
 \cdots
 (U_{\rm micro})_{\xi_1,\xi_i}.
\]
Cada elección de estados intermedios determina un camino de longitud \(n\),
y cada camino determina un único producto ordenado. Para \(n=0\), el único
término es el camino vacío y su peso es la identidad.
\end{proof}

La invariancia por traslaciones descompone el problema en bloques de cuatro
componentes direccionales, tensorizados con \(E\). Para
\[
 k_x=\frac{2\pi a}{27},
 \qquad
 k_y=\frac{2\pi b}{27},
 \qquad 0\leq a,b<27,
\]
adoptamos la convención de Fourier realificada
\[
 \chi_{a,b}(x,y)
 :=\exp[-J_E(k_xx+k_yy)].
\]
La sustitución directa en el desplazamiento condicionado produce
\[
 S(k_x,k_y)
 =\operatorname{diag}
 \bigl(
  e^{J_Ek_y},e^{-J_Ek_y},
  e^{J_Ek_x},e^{-J_Ek_x}
 \bigr),
\]
y, como \(D_J\) y \(C_4\) no dependen de la posición,
\begin{equation}
 \begin{aligned}
 U_{\rm micro}(k_x,k_y)
   &=S(k_x,k_y)D_JC_4,\\
 U_{12}^{\mathrm{reg}}(k_x,k_y)
   &=U_{\rm micro}(k_x,k_y)^{12}.
 \end{aligned}
 \label{eq:bloques-fourier-registro-27-rev3}
\end{equation}
Por consiguiente,
\[
 \operatorname{Spec}U_{12}^{\mathrm{reg}}
 =
 \bigcup_{a,b=0}^{26}
 \operatorname{Spec}
 U_{12}^{\mathrm{reg}}
 \left(\frac{2\pi a}{27},\frac{2\pi b}{27}\right).
\]
La potencia matricial enumera caminos y la descomposición de Fourier publica
las cuasienergías de la misma dinámica finita.

Fijada una duración elemental \(\delta t>0\), sea
\(\Delta t_{12}:=12\,\delta t\). En cada bloque unitario elíjase una rama
espectral \(\mathfrak b\) del logaritmo. El operador
\(\operatorname{Log}_{\mathfrak b}U_{12}^{\mathrm{reg}}(k_x,k_y)\) es
anti-autoadjunto y conmuta con \(J_E\). Se define
\begin{equation}
 \begin{aligned}
 H_{\rm eff}(k_x,k_y)
 &:=
 \frac{\hbar_{\rm int}}{\Delta t_{12}}\,
 J_E\,
 \operatorname{Log}_{\mathfrak b}
 U_{12}^{\mathrm{reg}}(k_x,k_y),\\
 U_{12}^{\mathrm{reg}}(k_x,k_y)
 &=
 \exp\!\left[
  -\frac{J_E\Delta t_{12}}{\hbar_{\rm int}}\,
  H_{\rm eff}(k_x,k_y)
 \right].
 \end{aligned}
 \label{eq:hamiltoniano-efectivo-registro-27-rev3}
\end{equation}
La primera línea es autoadjunta en la carta compleja. La segunda se sigue de
\(-J_E^2=I\) y demuestra que la exponencial recupera exactamente el
propagador. La rama \(\mathfrak b\) fija la zona de cuasienergía.

Esta realización construye un propagador finito, su suma exacta de caminos y
su generador efectivo. Su identificación con un Hamiltoniano físico continuo
requiere todavía verificar las hipótesis de representación, acción y límite
que se enuncian a continuación.

\subsection{Límite cilíndrico sobre el espacio de genealogías}

Fijados los extremos, sean \(\mathcal P_n\) conjuntos finitos de caminos a
resolución \(n\), con aplicaciones de restricción sobreyectivas. El espacio
\[
 \Omega^{\rm path}_{fi}:=\varprojlim_n\mathcal P_n
\]
es compacto y totalmente disconexo. Una familia compatible de medidas
probabilísticas \(\mu_n\) sobre los cilindros determina una medida de Borel
\(\mu\) en \(\Omega^{\rm path}_{fi}\).

Sea \(\Pi_n\) la partición en cilindros de nivel \(n\). Para
\(C\in\Pi_n\), elíjanse un representante \(\gamma_C\), una acción aproximada
\(S_n(C)\) y el peso \(\mu_n(C)\). Para
\(F:\Omega^{\rm path}_{fi}\to\operatorname{End}(E)\), definimos
\begin{equation}
 \mathcal I_n(F)
 :=\sum_{C\in\Pi_n}\mu_n(C)F(\gamma_C)
 \exp\!\left[-J_E\frac{S_n(C)}{\hbar_{\rm int}}\right].
 \label{eq:integral-cilindrica-aproximante-rev2}
\end{equation}

\begin{theorem}[Integral genealógica de caminos]
\label{thm:integral-genealogica-caminos-rev2}
Supóngase que:
\begin{enumerate}
 \item las medidas cilíndricas \(\mu_n\) son proyectivamente compatibles;
 \item \(F\) es continua y acotada;
 \item existe una acción continua
 \(S:\Omega^{\rm path}_{fi}\to\mathcal L_S\) para la cual
 \[
  \max_{C\in\Pi_n}\sup_{\gamma\in C}
  \left\|\frac{S_n(C)-S(\gamma)}{\hbar_{\rm int}}\right\|
  \longrightarrow0;
 \]
 \item la malla de \(\Pi_n\) tiende a cero.
\end{enumerate}
Entonces \(\mathcal I_n(F)\) converge en norma de operador, el límite es
independiente de los representantes y define la integral de Bochner
\begin{equation}
 \int_{\Omega^{\rm path}_{fi}}
 F(\gamma)
 \exp\!\left[-J_E\frac{S(\gamma)}{\hbar_{\rm int}}\right]
 \mathrm d\mu(\gamma).
 \label{eq:integral-genealogica-caminos-rev2}
\end{equation}
\end{theorem}

\begin{proof}
La compatibilidad de las medidas finitas construye \(\mu\) sobre el álgebra
de cilindros y, por extensión, sobre la sigma-álgebra de Borel. La función
\[
 G(\gamma)=F(\gamma)
 \exp[-J_E\,S(\gamma)/\hbar_{\rm int}]
\]
es continua sobre un compacto y, por tanto, uniformemente continua y
acotada. La ortogonalidad de \(J_E\) proporciona la estimación
\[
 \|e^{-J_E a}-e^{-J_E b}\|\leq |a-b|.
\]
La aproximación uniforme de la acción y la contracción de la malla hacen
converger las sumas cilíndricas a la integral de Bochner de \(G\)
\cite{Bochner1955}.
\end{proof}

El teorema construye una integral rigurosa sobre el espacio compacto de
genealogías. Su realización como integral oscilatoria de Feynman para un
Hamiltoniano concreto se obtiene cuando los kernels finitos, la acción y la
medida de cilindros satisfacen estas hipótesis de convergencia. La expansión
matricial finita, la integral genealógica y el propagador físico quedan así
relacionados por morfismos explícitos, sin identificarse por abreviación.

\subsection{Valor publicado y memoria de la fibra}

Sea
\(\operatorname{Val}_{\rm path}:\Omega^{\rm path}_{fi}\to\mathcal Y_{\rm val}\subset\mathbb R\)
un lector de caminos Borel-medible, donde \(\mathcal Y_{\rm val}\) lleva su estructura
boreliana estándar, y sea
\(\nu=(\operatorname{Val}_{\rm path})_*\mu\). Para toda función que
factoriza por el valor, \(F=\widehat F\circ\operatorname{Val}_{\rm path}\),
se cumple
\[
 \int_{\Omega^{\rm path}_{fi}}F\,\mathrm d\mu
 =\int_{\mathcal Y_{\rm val}}\widehat F(x)\,\mathrm d\nu(x).
\]
Como \(\Omega^{\rm path}_{fi}\) es compacto metrizable y \(\mu\) es una
probabilidad de Borel, existe una familia de probabilidades condicionales
regulares \((\mu_x)_{x\in\mathcal Y_{\rm val}}\), definida para \(\nu\)-casi todo \(x\). Para
un integrando genealógico general, la desintegración de \(\mu\) adopta la forma
\begin{equation}
 \int_{\Omega^{\rm path}_{fi}}G(\gamma)\,\mathrm d\mu(\gamma)
 =\int_{\mathcal Y_{\rm val}}\left[
   \int_{\operatorname{Val}_{\rm path}^{-1}(x)}
   G(\gamma)\,\mathrm d\mu_x(\gamma)
  \right]\mathrm d\nu(x).
 \label{eq:desintegracion-genealogica-caminos-rev2}
\end{equation}
La integral interior conserva las historias distintas que publican la misma
coordenada. Este es el enlace preciso entre la suma sobre todos los caminos,
la evaluación real y la tesis según la cual el continuo es una proyección de
un dominio genealógico más rico.

\paragraph{Dominio de la integral genealógica de caminos: compacidad cilíndrica, consistencia proyectiva, convergencia de núcleos y desintegración sobre fibras}
La expansión finita, la realización unitaria sobre
\(\mathbb T_{27}^{\,2}\) y su reducción de Fourier son exactas en sus dominios
finitos. La construcción de \(H_{\rm eff}\) es
exacta una vez declaradas la estructura compleja, la acción elemental
traslacional, la duración y la rama espectral; no selecciona por sí sola un
Hamiltoniano físico continuo. La categoría enriquecida de rutas y la fase
interna preceden al ensamblaje cilíndrico. El
\cref{thm:integral-genealogica-caminos-rev2} es exacto bajo las cuatro
premisas declaradas. La realización de un modelo cuántico continuo concreto
requiere verificar esas premisas para su acción y su Hamiltoniano.

% RESULTADO_RECUPERADO. II REV10, 07_elipse.tex:1--61,
% 07b_geometria_elipse.tex:73--115 y 12_dualidad_t.tex:27--107.
% Explicitaciones algebraicas: cota H5<2 para verificar la carta y
% sustitución Rstar^4=(499 alpha-eta_ret)/(501 alpha+eta_ret).
\section{Elipse de acción y determinación del radio normalizado}
\label{iv:ellipse:section}

El lector angular y la sección de acción conservan el mismo
antecedente APP--TRIT--TPK. La forma que se construye a
continuación utiliza las publicaciones
\eqref{iv:angular:representantes} y una sección positiva
\(\hbar=\hbar_s\), con
\(s\in\{\mathrm{pre},\mathrm{ret}\}\), de
\eqref{iv:action:hbar}. La razón angular determinará la forma;
la sección de acción conservará su escala común.

\subsection{Dominio y semiejes etiquetados}
\label{iv:ellipse:dominio}

Para abreviar, sean \(A=A_{\deg}\) y \(C^*=C^*_{\deg}\).
En la carta \(A\ne0\), \(\hbar>0\) y \(|C^*/A|<1\),
definimos
\begin{equation}
 \xi=\frac{C^*}{A},\qquad
 \eta_{\mathrm{el}}=\operatorname{artanh}\xi,\qquad
 a_S=\sqrt{2\hbar}\,e^{\eta_{\mathrm{el}}/2},\qquad
 b_S=\sqrt{2\hbar}\,e^{-\eta_{\mathrm{el}}/2}.
 \label{iv:ellipse:semiejes}
\end{equation}
Los semiejes y las coordenadas equilibradas \(Q,P\) tienen
unidad de raíz de acción. Los ejes permanecen etiquetados:
el signo de \(\eta_{\mathrm{el}}\) determina cuál es el
mayor. La curva y su producto de semiejes son
\begin{equation}
 \gamma(\theta)=(a_S\cos\theta,b_S\sin\theta),\qquad
 \frac{Q^2}{a_S^2}+\frac{P^2}{b_S^2}=1,\qquad
 a_Sb_S=2\hbar.
 \label{iv:ellipse:curva}
\end{equation}

\begin{proposition}[Pertenencia de las publicaciones a la carta elíptica]
\label{iv:ellipse:pertenencia}
La rama positiva de \eqref{iv:angular:representantes}, con los
intervalos \eqref{iv:action:intervalos}, satisface
\(A>0\), \(0<C^*<A\) y \(\hbar_s>0\).
La reversión de la coordenada impar conserva \(|C^*/A|<1\).
\end{proposition}
\begin{proof}
La proposición \ref{iv:action:positividad} da
\(0<\eta_{\mathrm{ret}}<H_5\) y la positividad de ambas
secciones de acción. Para una cota superior de \(H_5\),
\eqref{iv:action:intervalos} implica
\(\alpha<1/125\) y \(\varphi>8/5\). Al conservar sólo
los términos positivos de \eqref{iv:action:h5}, se obtiene
\[
 H_5<
 \frac{\sqrt5}{2}+
 \frac{9}{16\cdot125^2}+
 \frac{7}{48\cdot125^4}.
\]
Como \(5<81/16\), se tiene \(\sqrt5/2<9/8\). Además,
\[
 \frac{9}{16\cdot125^2}+
 \frac{7}{48\cdot125^4}
 =\frac{421882}{11718750000}<\frac1{1000}.
\]
Por tanto, \(H_5<9/8+1/1000<2\). Usando ahora
\(\alpha>7/1000\),
\[
 0<C^*=2(\eta_{\mathrm{ret}}+\alpha)
 <2\left(2+\frac1{125}\right)=\frac{502}{125}
 <7<1000\alpha=A.
\]
El intercambio de canales de
\ref{iv:angular:inversion} cambia \(C^*\) por \(-C^*\)
y conserva \(A\); por ello conserva la desigualdad absoluta.
\end{proof}

\subsection{Fase paramétrica y área orientada}
\label{iv:ellipse:area}

\begin{proposition}[Ley areal y transporte de las coordenadas angulares]
\label{iv:ellipse:ley-areal}
La fase paramétrica \(\theta\) de \(\gamma\) y su ángulo
polar \(\phi\), levantados continuamente desde cero, satisfacen
\begin{equation}
 \phi(\theta)=\arg(a_S\cos\theta+i b_S\sin\theta),\qquad
 \frac{d\phi}{d\theta}
 =\frac{a_Sb_S}{a_S^2\cos^2\theta+b_S^2\sin^2\theta}>0.
 \label{iv:ellipse:fases}
\end{equation}
En una carta donde ambas tangentes estén definidas,
\(\tan\phi=(b_S/a_S)\tan\theta\). El área orientada barrida
desde la fase cero es
\begin{equation}
 \mathcal A(\theta)=
 \frac12\int_0^\theta(QP'-PQ')\,dt=\hbar\theta,
 \qquad
 \mathcal A(1)=\hbar,\quad\mathcal A(2\pi)=h_s.
 \label{iv:ellipse:area-fase}
\end{equation}
\end{proposition}
\begin{proof}
La derivada del argumento de un vector no nulo es
\((QP'-PQ')/(Q^2+P^2)\). Para la curva
\eqref{iv:ellipse:curva}, el numerador vale
\(a_Sb_S=2\hbar\); esto demuestra la derivada y su
positividad. El cociente \(P/Q\) da la relación entre tangentes;
el levantamiento continuo conserva los cuadrantes y las vueltas.
La integral vale \(a_Sb_S\theta/2=\hbar\theta\).
Finalmente, \(h_s=2\pi\hbar_s\) da la identidad de la
vuelta completa.
\end{proof}

Un radián de fase paramétrica barre \(\hbar\). El ángulo
polar conserva su propia derivada en
\eqref{iv:ellipse:fases}. En coordenadas polares, la ley
equivalente es
\[
 d\mathcal A=\frac12\rho_{\partial}(\phi)^2\,d\phi,\qquad
 \rho_{\partial}(\phi)=
 \left(\frac{\cos^2\phi}{a_S^2}
       +\frac{\sin^2\phi}{b_S^2}\right)^{-1/2}.
\]
La fórmula radial se obtiene sustituyendo
\((Q,P)=\rho(\cos\phi,\sin\phi)\) en
\eqref{iv:ellipse:curva}. Con la orientación de
\(\gamma\), la integral de la uno-forma canónica conserva
el signo
\begin{equation}
 \oint_\gamma P\,dQ=-\pi a_Sb_S=-h_s.
 \label{iv:ellipse:accion-orientada}
\end{equation}
En efecto, \(P\,dQ=-a_Sb_S\sin^2\theta\,d\theta\) y
la integral de \(\sin^2\theta\) en una vuelta es \(\pi\).

\subsection{Reciprocidad de forma y orientación simpléctica}
\label{iv:ellipse:reciprocidad}

Definimos las matrices
\begin{equation}
 D(\eta)=\begin{pmatrix}e^\eta&0\\0&e^{-\eta}\end{pmatrix},
 \qquad
 S_2=\begin{pmatrix}0&1\\1&0\end{pmatrix},
 \qquad
 J_2=\begin{pmatrix}0&-1\\1&0\end{pmatrix}.
 \label{iv:ellipse:transportes}
\end{equation}

\begin{proposition}[Forma cuadrática y acción bajo reciprocidad]
\label{iv:ellipse:accion-reciproca}
Ambos transportes satisfacen
\(T^{\mathsf T}D(-\eta)T=D(\eta)\).
Para \(\omega=dQ\wedge dP\), el intercambio \(S_2\) es
antisimpléctico y el giro \(J_2\) es simpléctico:
\[
 S_2^*\omega=-\omega,\qquad J_2^*\omega=\omega.
\]
Si \(\mathcal I(\gamma)=\oint_\gamma P\,dQ\) sobre una
curva cerrada, entonces
\begin{equation}
 \mathcal I(S_2\gamma)=-\mathcal I(\gamma),\qquad
 \mathcal I(J_2\gamma)=\mathcal I(\gamma).
 \label{iv:ellipse:signos-accion}
\end{equation}
\end{proposition}
\begin{proof}
La multiplicación matricial intercambia las dos entradas
diagonales de \(D\), lo que demuestra ambas congruencias.
Los determinantes de \(S_2,J_2\) son \(-1,1\);
de ahí se obtiene la transformación de la forma de área.
Para \(\lambda=P\,dQ\), el cálculo directo da
\[
 S_2^*\lambda=d(PQ)-\lambda,\qquad
 J_2^*\lambda=\lambda-d(PQ).
\]
La integral de la diferencial exacta sobre la curva cerrada
es cero, y quedan probadas las dos identidades de acción.
\end{proof}

La orientación distingue los dos transportes aunque publiquen
la misma forma cuadrática, incluso en el caso circular.
El marcado de los ejes y el signo de la acción pertenecen,
por tanto, a los datos que se conservan al pasar a la
representación reticular.

\Needspace{12\baselineskip}
\subsection{Radio normalizado y matriz de la lectura reticular}
\label{iv:ellipse:radio}

\begin{proposition}[Determinación y recuperabilidad del radio de forma]
\label{iv:ellipse:radio-prop}
La elipse etiquetada determina un único radio positivo
\begin{equation}
 R_\star=\sqrt{\frac{b_S}{a_S}}
 =e^{-\eta_{\mathrm{el}}/2}
 =\left(\frac{A-C^*}{A+C^*}\right)^{1/4},
 \qquad \tau_{\mathrm{el}}=iR_\star^2.
 \label{iv:ellipse:rstar}
\end{equation}
La razón angular se recupera mediante
\begin{equation}
 \frac{C^*}{A}=\frac{1-R_\star^4}{1+R_\star^4},
 \qquad
 R_\star^4=
 \frac{499\alpha-\eta_{\mathrm{ret}}}
      {501\alpha+\eta_{\mathrm{ret}}}.
 \label{iv:ellipse:rstar-compuesto}
\end{equation}
El intercambio de semiejes transforma
\((\eta_{\mathrm{el}},R_\star,\tau_{\mathrm{el}})\) en
\((-\eta_{\mathrm{el}},R_\star^{-1},-1/\tau_{\mathrm{el}})\)
y conserva \(a_Sb_S=2\hbar\).
\end{proposition}
\begin{proof}
La división de los semiejes da \(b_S/a_S=e^{-\eta_{\mathrm{el}}}\).
La aplicación \(R\mapsto R^2\) es biyectiva sobre los reales
positivos, de modo que la raíz positiva es única.
Para \(\xi=C^*/A\in(-1,1)\),
\[
 2\operatorname{artanh}\xi
 =\log\frac{1+\xi}{1-\xi};
\]
por tanto, \(R_\star^4=(1-\xi)/(1+\xi)\).
Despejar \(\xi\) demuestra la primera identidad de
\eqref{iv:ellipse:rstar-compuesto}. Sustituir
\(A=1000\alpha\) y
\(C^*=2(\eta_{\mathrm{ret}}+\alpha)\) da
\[
 \frac{A-C^*}{A+C^*}
 =\frac{998\alpha-2\eta_{\mathrm{ret}}}
        {1002\alpha+2\eta_{\mathrm{ret}}}
 =\frac{499\alpha-\eta_{\mathrm{ret}}}
        {501\alpha+\eta_{\mathrm{ret}}}.
\]
El intercambio invierte la razón de semiejes y cambia el
signo de \(\eta_{\mathrm{el}}\). La identidad
\(-1/(iR_\star^2)=iR_\star^{-2}\) da el transporte
modular. El producto de los semiejes permanece fijo.
\end{proof}

En la rama positiva, \(0<R_\star<1\); la carta de forma
revertida publica \(R_\star^{-1}>1\). La razón angular
se recupera a partir del radio, mientras los ángulos absolutos
y su historia permanecen en el estado enriquecido.
La forma circular corresponde a \(C^*/A=0\), con
\(R_\star=1\). Cambiar la sección positiva de acción,
manteniendo el par angular, modifica la escala común de
los semiejes y conserva su razón.

\begin{proposition}[Forma reticular determinada por la elipse]
\label{iv:ellipse:matriz-prop}
La matriz de semiejes cuadrados normalizados es
\begin{equation}
 D(\eta_{\mathrm{el}})
 =\frac1{2\hbar}
    \begin{pmatrix}a_S^2&0\\0&b_S^2\end{pmatrix}
 =\begin{pmatrix}R_\star^{-2}&0\\0&R_\star^2\end{pmatrix}.
 \label{iv:ellipse:matriz}
\end{equation}
Su evaluación sobre un par entero produce
\begin{equation}
 E_\star(m,w)=
 \begin{pmatrix}m&w\end{pmatrix}
 D(\eta_{\mathrm{el}})
 \begin{pmatrix}m\\w\end{pmatrix}
 =\frac{m^2}{R_\star^2}+w^2R_\star^2.
 \label{iv:ellipse:energia}
\end{equation}
La ecuación de la elipse multiplicada por \(2\hbar\)
utiliza, en cambio, la matriz inversa \(D(-\eta_{\mathrm{el}})\).
\end{proposition}
\begin{proof}
Los cuadrados de los semiejes de
\eqref{iv:ellipse:semiejes}, divididos por \(2\hbar\),
son \(e^{\eta_{\mathrm{el}}}\) y
\(e^{-\eta_{\mathrm{el}}}\).
Por \eqref{iv:ellipse:rstar}, éstos son
\(R_\star^{-2}\) y \(R_\star^2\).
La multiplicación por el vector entero demuestra
\eqref{iv:ellipse:energia}. Finalmente,
\(2\hbar/a_S^2=e^{-\eta_{\mathrm{el}}}\) y
\(2\hbar/b_S^2=e^{\eta_{\mathrm{el}}}\), lo que
identifica la matriz de la ecuación elíptica.
\end{proof}

El radio \(R_\star\) es una publicación adimensional de
la forma etiquetada. El radio polar
\(\rho_{\partial}\) pertenece a la recta de raíz de acción.
La representación por un radio espacial requiere su carta
dimensional posterior; ésta conserva la forma ya determinada
y la independencia de su construcción respecto de una
longitud prescrita.

\subsection{Realización inductiva--capacitiva de la forma de acción}
\label{nuclear:realizacion-lc}

La forma de acción ya construida admite una realización eléctrica en
una carta provista de un reloj \(\omega>0\) y una impedancia de
referencia \(Z_*>0\). Estas dos magnitudes conservan sus unidades y
acompañan al parámetro de forma \(\eta\). Se definen
\begin{equation}
 L_\eta=\frac{Z_*}{\omega}e^{-\eta},\qquad
 C_\eta=\frac{e^\eta}{Z_*\omega},\qquad
 Z_\eta=\sqrt{L_\eta/C_\eta}.
 \label{hmtII:eq:ii-elipse-lc}
\end{equation}
Aquí \(C_\eta\) denota capacitancia, distinta del contraángulo \(C^*\).

\begin{proposition}[Conservación del producto dinámico]
\label{nuclear:lc-producto}
La realización \eqref{hmtII:eq:ii-elipse-lc} satisface
\begin{equation}
 L_\eta C_\eta=\omega^{-2},\qquad
 \frac{Z_\eta}{Z_*}=e^{-\eta}=\frac{b_S}{a_S},
 \label{hmtII:eq:ii-elipse-impedancia}
\end{equation}
donde \(a_S=\sqrt{2\hbar}\,e^{\eta/2}\) y
\(b_S=\sqrt{2\hbar}\,e^{-\eta/2}\) son los semiejes de
\cref{iv:ellipse:semiejes}.
Con \(Q=\sqrt{Z_*}q\) y \(P=\Phi/\sqrt{Z_*}\), la energía toma la forma
\[
 H_{LC}=\frac{q^2}{2C_\eta}+\frac{\Phi^2}{2L_\eta}
       =\frac\omega2(e^{-\eta}Q^2+e^\eta P^2).
\]
Sobre \((Q,P)=(a_S\cos\theta,b_S\sin\theta)\), las contribuciones
eléctrica y magnética son \(\hbar\omega\cos^2\theta\) y
\(\hbar\omega\sin^2\theta\); la energía total es \(\hbar\omega\).
\end{proposition}
\begin{proof}
La multiplicación de \(L_\eta\) y \(C_\eta\) cancela los factores
\(Z_*\) y \(e^{\pm\eta}\), mientras su cociente es
\(Z_*^2e^{-2\eta}\). La raíz positiva prueba la identidad de impedancia.
Sustituir \(q=Q/\sqrt{Z_*}\) y \(\Phi=\sqrt{Z_*}P\) da la forma
de \(H_{LC}\). Los cuadrados de los semiejes cancelan los factores
de forma restantes; la suma utiliza \(\cos^2\theta+\sin^2\theta=1\).
\end{proof}

La convención hamiltoniana
\(\dot Q=\partial_PH\), \(\dot P=-\partial_QH\) asigna
\(\dot\theta=-\omega\) a esta parametrización. La forma, el reloj
y la carta dimensional permanecen diferenciados. La respuesta
constitutiva del vacío y la impedancia \(Z_\eta\) quedan relacionadas
por la inversión de canales de \cref{iii:prop:forma-LC}; cada una
conserva su definición y su normalización.

% Procedencia: integral 2249, cap. 24, pp. 609--618; capitulo_24_dualidad_t.tex.
% Las fórmulas de dominio operatorio se explicitan sobre la realización ya construida.
\section{Dualidad T: coordenadas enteras, cambio de sección y equivalencia operatoria}
\label{iv:dualidad}

La descomposición aritmética de APP conserva el representante residual
y su cociente. La evolución TPK transporta además su orientación y
la memoria de los cruces de frontera. El intercambio de las dos
coordenadas exige conservar esta distinción y precisar el dominio
en que la operación vuelve a estar definida. Sobre ese dominio se
construye primero una involución aritmética; su forma cuadrática
admite después una representación por operadores autoadjuntos.

\subsection{Dominio estable del intercambio}

En la sección residual positiva,
\begin{equation}
 r_9^+(n)=1+((n-1)\bmod9),\qquad
 Q_9^+(n)=\frac{n-r_9^+(n)}9,
 \qquad n=r_9^+(n)+9Q_9^+(n).
 \label{iv:eq:division-positiva}
\end{equation}
Un intercambio puede situar un cociente arbitrario en la primera
coordenada. Por ello su dominio estable es
\(\mathscr D_0=\ZZ^2\times\RR_{>0}\), no solamente la sección
\(1\leq r\leq9\). La prolongación del dominio permite definir
una operación total; la admisibilidad de una pareja como lectura de
una ruta TPK conserva por separado sus reglas de selección.

\begin{definition}[Forma compacta normalizada e intercambio]
\label{iv:def:dualidad}
Para \(d=(m,w,R_0)\in\mathscr D_0\), sean
\begin{equation}
 E_{R_0}(m,w)=\frac{m^2}{R_0^2}+w^2R_0^2,
 \qquad
 \mathsf T_0(m,w,R_0)=(w,m,R_0^{-1}).
 \label{iv:eq:energia-t}
\end{equation}
Los símbolos \(m,w\) nombran las dos coordenadas enteras
transportadas; \(R_0\) es la escala positiva adimensional de
la forma. Su interpretación mediante momento, enrollamiento y
radio físico se establece por un mapa de realización posterior.
\end{definition}

\begin{theorem}[Involución y conservación de la forma cuadrática]
\label{iv:thm:dualidad-escalar}
La aplicación \(\mathsf T_0\) es una involución de
\(\mathscr D_0\), y
\[
 E_{R_0^{-1}}(w,m)=E_{R_0}(m,w).
\]
\end{theorem}
\begin{proof}
Dos intercambios restituyen \((m,w)\), y dos inversiones restituyen
\(R_0\). Además,
\[
 E_{R_0^{-1}}(w,m)=w^2R_0^2+\frac{m^2}{R_0^2}
 =E_{R_0}(m,w).
\]
\end{proof}

La equivalencia vale para cada radio y su recíproco. En \(R_0=1\),
el radio queda fijo y la operación actúa sobre una misma forma.
Una pareja concreta es fija por la involución completa precisamente
cuando, además, \(m=w\).

\subsection{Hamiltoniano compacto y dominio autoadjunto}

Sea \(\mathcal H_{\mathrm{lat}}=\ell^2(\ZZ^2)\), con base
ortonormal \(\{|m,w\rangle\}\). Definimos
\begin{align}
 H(R_0)|m,w\rangle&=E_{R_0}(m,w)|m,w\rangle,\label{iv:eq:hamiltoniano}\\
 D(H(R_0))&=\left\{\psi:\sum_{m,w\in\ZZ}
 E_{R_0}(m,w)^2|\psi_{m,w}|^2<\infty\right\}.\label{iv:eq:dominio-h}
\end{align}
Los vectores de soporte finito pertenecen al dominio y forman un
núcleo del operador diagonal. La forma cuadrática asociada tiene
dominio
\[
 D(H(R_0)^{1/2})=\left\{\psi:\sum_{m,w}
 E_{R_0}(m,w)|\psi_{m,w}|^2<\infty\right\}.
\]

\begin{proposition}[Autoadjunción y resolvente compacto]
\label{iv:prop:h-autoadjunto}
El operador \(H(R_0)\) es autoadjunto y no negativo. Sus
autovalores son \(E_{R_0}(m,w)\), contados con multiplicidad,
y \((H(R_0)+I)^{-1}\) es compacto.
\end{proposition}
\begin{proof}
Un operador de multiplicación por una sucesión real es simétrico
en \eqref{iv:eq:dominio-h}. Si \(\eta\) pertenece al dominio
del adjunto, probar la identidad del adjunto contra cada vector
\(|m,w\rangle\) obliga a que su imagen tenga coordenadas
\(E_{R_0}(m,w)\eta_{m,w}\). La condición de pertenecer a
\(\ell^2\) es exactamente \eqref{iv:eq:dominio-h}; por tanto
el adjunto tiene el mismo dominio y la misma acción. La no
negatividad se obtiene sumando los pesos no negativos.

Con \(c_{R_0}=\min(R_0^{-2},R_0^2)>0\), se tiene
\(E_{R_0}(m,w)\geq c_{R_0}(m^2+w^2)\). Fuera de cualquier
conjunto finito creciente, los coeficientes
\((1+E_{R_0}(m,w))^{-1}\) tienden uniformemente a cero.
Truncar el resolvente en dichos conjuntos produce operadores
de rango finito que convergen en norma, lo cual prueba su
compacidad. La base diagonal determina el espectro anunciado.
\end{proof}

\begin{theorem}[Conjugación unitaria, incluidos los dominios]
\label{iv:thm:t-unitaria}
\label{prop:dualidad-radios-k}
El intercambio
\(U_T|m,w\rangle=|w,m\rangle\) es unitario,
\(U_T^2=I\), y satisface
\begin{equation}
 U_TD(H(R_0))=D(H(R_0^{-1})),
 \qquad U_TH(R_0)U_T^{-1}=H(R_0^{-1}).
 \label{iv:eq:conjugacion-h}
\end{equation}
También transporta los dominios de las formas cuadráticas.
\end{theorem}
\begin{proof}
El intercambio es una permutación involutiva de la base
ortonormal. Para una sucesión \(\psi\), la identidad del
teorema \ref{iv:thm:dualidad-escalar} da
\[
 \sum_{m,w}E_{R_0^{-1}}(m,w)^2|(U_T\psi)_{m,w}|^2
 =\sum_{m,w}E_{R_0}(m,w)^2|\psi_{m,w}|^2.
\]
Queda establecida la igualdad de los dominios. En cada coordenada,
\(U_TH(R_0)U_T^{-1}\) multiplica por
\(E_{R_0}(w,m)=E_{R_0^{-1}}(m,w)\). La misma suma, con
la primera potencia del peso, demuestra el resultado para las
formas cuadráticas.
\end{proof}

La igualdad escalar y la igualdad operatoria son dos realizaciones
de la misma involución. La segunda conserva los dominios que
permiten interpretar la primera como equivalencia espectral,
no sólo como una igualdad sobre vectores de soporte finito.

\subsection{Especialización al radio de la elipse de acción}

El radio \(R_\star\) determinado en la sección
\ref{iv:ellipse:section} fija la realización de esta familia que
corresponde a las publicaciones angulares del estado. El símbolo
\(R_0\) permite expresar los resultados de transporte para toda la
familia positiva; la composición con la elipse evalúa esa familia en
\(R_0=R_\star\) y conserva también su imagen recíproca.

\begin{corollary}[Familia elíptica estable por dualidad]
\label{iv:cor:radio-especializado}
El conjunto
\[
 \mathscr D_{\mathrm{el}}
 =\ZZ^2\times\{R_\star,R_\star^{-1}\}
\]
es estable por \(\mathsf T_0\). Sus Hamiltonianos son autoadjuntos,
no negativos y de resolvente compacto. El intercambio \(U_T\)
transporta el dominio de \(H(R_\star)\) sobre el de
\(H(R_\star^{-1})\), incluidos los dominios de sus formas.
\end{corollary}
\begin{proof}
La positividad de \(R_\star\) y su determinación única se establecen
en la sección \ref{iv:ellipse:section}. La inversión permuta
\(R_\star\) y \(R_\star^{-1}\), mientras el intercambio preserva
\(\ZZ^2\). Se aplican, en estas dos fibras, la proposición
\ref{iv:prop:h-autoadjunto} y el teorema \ref{iv:thm:t-unitaria}.
Si \(R_\star=1\), las dos fibras coinciden; en cualquier otro caso
el operador entrelaza dos fibras distintas de la misma familia.
\end{proof}

La escala común de acción determina el área de la elipse; su cociente
de semiejes determina \(R_\star\). Esta separación permite conservar
el área y la escala común de acción al intercambiar los semiejes y obtener la reciprocidad
del radio. El intercambio de índices sobre \(\ell^2(\ZZ^2)\)
y los transportes del plano de acción actúan en espacios diferentes;
sus propiedades respectivas están demostradas en sus dominios propios.

\subsection{Cambio de representante y memoria del cociente}

La sección cero usa \(\bar r\in\{0,\ldots,8\}\). El cambio
que conserva el entero de \eqref{iv:eq:division-positiva} es
\begin{equation}
 C(r,Q)=(\bar r,Q+\mathbf1_{\{r=9\}}).
 \label{iv:eq:cambio-seccion}
\end{equation}
La modificación del cociente forma parte de la operación.
Para \(n=9\), las dos parejas son \((9,0)\) y \((0,1)\).
Sus energías no corregidas son \(81/R_0^2\) y \(R_0^2\).
Representar el mismo entero no implica, por tanto, igualdad de
esas dos inserciones en una forma que distingue las coordenadas.

La relación de equivalencia se conserva definiendo
\[
 L_9=\ZZ(9,-1),\qquad S(x_1,x_2)=(x_2,x_1),
 \qquad q_{R_0}(x)=\frac{x_1^2}{R_0^2}+x_2^2R_0^2.
\]
Para \(x\in\ZZ^2\), sea
\begin{equation}
 \mathcal E_{R_0}^{L_9}([x])
 =\min_{k\in\ZZ}q_{R_0}(x+k(9,-1)).
 \label{iv:eq:energia-cociente}
\end{equation}
Esta construcción opera sobre clases. La memoria que distingue
representantes individuales permanece en el estado TPK anterior
al cociente; la función de \eqref{iv:eq:energia-cociente}
determina la evaluación invariante de la clase elegida.

\begin{theorem}[Dualidad covariante de las secciones]
\label{iv:thm:dualidad-cociente}
La energía de \eqref{iv:eq:energia-cociente} está bien definida.
La transformación
\[
 \mathsf T_{\mathrm{cov}}([x]_{L_9},R_0)
 =([Sx]_{SL_9},R_0^{-1})
\]
es una biyección entre
\((\ZZ^2/L_9)\times\RR_{>0}\) y
\((\ZZ^2/SL_9)\times\RR_{>0}\), con inversa dada por el
segundo intercambio, y
\[
 \mathcal E_{R_0^{-1}}^{SL_9}([Sx])
 =\mathcal E_{R_0}^{L_9}([x]).
\]
\end{theorem}
\begin{proof}
Reemplazar \(x\) por \(x+k_0(9,-1)\) reindexa los candidatos
al mínimo. Éste existe: la función de \(k\) es cuadrática y
tiene coeficiente principal \(81/R_0^2+R_0^2>0\). El cambio
\((9,Q)\mapsto(0,Q+1)\) tiene diferencia en \(L_9\), por
lo que representa una misma clase. El intercambio transporta
la relación de equivalencia a \(SL_9\), y
\[
 \min_{\ell\in L_9}q_{R_0^{-1}}(Sx+S\ell)
 =\min_{\ell\in L_9}q_{R_0}(x+\ell).
\]
Finalmente, \(S^2=I\) recupera el cociente y el radio iniciales.
\end{proof}

La evaluación del mínimo puede efectuarse exactamente con dos
candidatos. Si \(x=(m,w)\), el mínimo real de la cuadrática
ocurre en
\[
 k_* =\frac{wR_0^4-9m}{81+R_0^4}.
\]
El mínimo entero se alcanza en \(\lfloor k_*\rfloor\) o en
\(\lceil k_*\rceil\). Esta fórmula no altera la equivalencia;
explicita su evaluación finita y permite comprobar el efecto
del cambio de sección sin descartar el cociente.

\subsection{Normalización de radio en la realización de cuerda}

Sea \(\alpha'>0\) una escala de la realización de cuerda.
Este símbolo designa la escala de cuerda y se mantiene distinto
de la constante de estructura fina \(\alpha\). Definimos
\begin{equation}
 \mathsf T_{\alpha'}(m,w,R)
 =\left(w,m,\frac{\alpha'}R\right),\qquad
 \mathcal N_{\alpha'}(m,w,R)
 =\left(m,w,\frac R{\sqrt{\alpha'}}\right).
 \label{iv:eq:normalizacion-cuerda}
\end{equation}

\begin{theorem}[Conjugación exacta por la escala]
\label{iv:thm:normalizacion-cuerda}
La normalización satisface
\[
 \mathcal N_{\alpha'}\mathsf T_{\alpha'}
 =\mathsf T_0\mathcal N_{\alpha'}.
\]
Si
\[
 E_{\mathrm{str}}(m,w,R)
 =\frac{m^2}{R^2}+\frac{w^2R^2}{\alpha'^2},
\]
entonces
\(E_{\mathrm{str}}=\alpha'^{-1}E\circ\mathcal N_{\alpha'}\).
Las combinaciones
\[
 p_L=\frac mR+\frac{wR}{\alpha'},\qquad
 p_R=\frac mR-\frac{wR}{\alpha'}
\]
se transforman mediante \(p_L\mapsto p_L\) y \(p_R\mapsto-p_R\).
\end{theorem}
\begin{proof}
Las dos composiciones de la primera igualdad dan
\((w,m,\sqrt{\alpha'}/R)\). Por otra parte,
\[
 E_{R/\sqrt{\alpha'}}(m,w)
 =\frac{\alpha'm^2}{R^2}+\frac{w^2R^2}{\alpha'}
 =\alpha'E_{\mathrm{str}}(m,w,R).
\]
Sustituir \((w,m,\alpha'/R)\) en las expresiones de
\(p_L,p_R\) da, respectivamente, las dos transformaciones
afirmadas.
\end{proof}

La conjugación fija qué significa expresar el mismo resultado en
una carta dimensional. En la familia elíptica, los radios son
\[
 R=\sqrt{\alpha'}R_\star,\qquad
 R^\vee=\sqrt{\alpha'}R_\star^{-1}.
\]
Su producto es \(\alpha'\). La determinación interna del cociente
adimensional de semiejes y la escala de longitud al cuadrado de la
realización cumplen funciones distintas: la primera fija la forma;
la segunda expresa sus radios en una unidad dimensional.

La normalización anterior permite expresar el mismo resultado en
radio normalizado o físico. La asignación de una tensión y una
unidad a \(\alpha'\) pertenece al mapa de realización y tiene
su propio origen; no es un dato que intervenga en la producción
de los registros APP--TRIT--TPK.

\subsection{Relación con la transformación modular}

La inversión del radio tiene una representación precisa en el
eje imaginario del semiplano superior. Sean
\(\iota(R_0)=\mathrm iR_0^2\) y
\(\mathsf S(\tau)=-1/\tau\). Entonces
\begin{equation}
 \mathsf S\circ\iota(R_0)=\iota(R_0^{-1}).
 \label{iv:eq:semiconjugacion-modular}
\end{equation}
La igualdad resulta de
\(-1/(\mathrm iR_0^2)=\mathrm iR_0^{-2}\). Para la función
modular \(J=j-744\), cuya invariancia se conserva en el
desarrollo de Moonshine, se deduce
\(J(\mathrm iR_0^2)=J(\mathrm iR_0^{-2})\).

El mapa \(\iota\) expresa una semiconjugación de parámetros.
El intercambio \(U_T\) actúa en las coordenadas de
\(\ell^2(\ZZ^2)\); \(\mathsf S\) actúa en el semiplano
superior; el automorfismo dual de un orbifold actúa en sus
sectores graduados. La escritura de sus dominios permite
componer relaciones específicas entre ellos conservando la
información de cada representación. Los resultados aquí
reunidos establecen el intercambio de coordenadas con sus
dominios operatorios y demostraciones explícitos.

% Procedencia: integral 2249, cap. 25, pp. 619--640, especialmente §25.13.
% Requiere la incorporación material de exc:k-direccion, exc:leech y el núcleo común.
\section{Pantallas dimensionales y reducción coordinada del estado}
\label{iv:pantallas}

Los rangos que aparecen en la prolongación excepcional y en la
compactificación pertenecen a objetos diferentes. Para relacionarlos
se conserva su presentación: generadores, relaciones, producto
interno y dirección seleccionada. La dimensión es entonces una
consecuencia del módulo de canales y de sus relaciones. Su
transporte se formula mediante mapas que actúan sobre ese módulo,
además de actuar sobre el estado que lo determina.

\subsection{Canales, relaciones y dimensión efectiva}

Sea \(\mathcal S_n\) el conjunto de estados supervivientes de
profundidad \(n\), con restricciones compatibles
\(r_{m,n}:\mathcal S_m\to\mathcal S_n\), \(m\geq n\).
El dato de frontera \(\beta_n(x)\) conserva la información
necesaria para decidir la compatibilidad de una prolongación:
hoja, orientación, cociente, memoria y clase de incidencia según
el lector considerado. Las restricciones satisfacen
\[
 r_{\ell,n}=r_{m,n}r_{\ell,m},\qquad
 \beta_n r_{m,n}=b_{m,n}\beta_m.
\]
Estas igualdades son las condiciones de naturalidad de la
presentación y mantienen los mismos datos al restringir por
etapas o directamente.

\begin{definition}[Pantalla de canales]
\label{iv:def:pantalla}
Fijado un cuerpo \(\mathbb K_n\), sea \(M_n(x)\) el espacio
generado por los canales de prolongación y \(\mathcal R_n(x)\)
el subespacio de relaciones. La pantalla es
\[
 \Sigma_n(x)=\bigl(\mathbb K_n,M_n(x),\beta_n(x),
                    \mathcal R_n(x)\bigr),
\]
y su dimensión efectiva es
\begin{equation}
 d_{\Sigma_n}(x)=\dim_{\mathbb K_n}
                    \bigl(M_n(x)/\mathcal R_n(x)\bigr).
 \label{iv:eq:dimension-pantalla}
\end{equation}
Si \(M_n(x)\cong\mathbb K_n^{g_n}\) y las columnas de
\(R_n(x)\) generan las relaciones, entonces
\(d_{\Sigma_n}(x)=g_n-\operatorname{rank}R_n(x)\).
\end{definition}

\begin{proposition}[Invariancia de presentación y transporte]
\label{iv:prop:pantalla-transporte}
La dimensión \eqref{iv:eq:dimension-pantalla} se conserva
bajo cambios de bases invertibles. Para dos pantallas sobre un mismo cuerpo
\(\mathbb K_n=\mathbb K_m=\mathbb K\), si la aplicación lineal
\(F_\gamma:M_n(x)\to M_m(T_\gamma x)\) envía
\(\mathcal R_n(x)\) en \(\mathcal R_m(T_\gamma x)\), induce
de manera única una aplicación lineal entre los cocientes de pantalla.
\end{proposition}
\begin{proof}
Un cambio de generadores y de relaciones transforma \(R\) en
\(URV\), con \(U,V\) invertibles, y conserva su rango. Para
el descenso, definimos
\(\overline F_\gamma([v])=[F_\gamma(v)]\). Si dos
representantes difieren en una relación, la linealidad implica que sus imágenes difieren
en una relación de la pantalla terminal; la aplicación está
bien definida. La proyección al cociente es sobreyectiva y
garantiza la unicidad.
\end{proof}

El cambio dimensional durante el transporte queda determinado
por
\[
 \Delta d(\gamma;x)=d_{\Sigma_m}(T_\gamma x)
                         -d_{\Sigma_n}(x).
\]
Esta expresión conserva el estado y su presentación de canales.
Un rango aislado no permite reconstruir ni las relaciones ni
el mapa que lo produjo.

\subsection{Perforación ternaria y pantalla de síndromes}

El código \(C_W=[12,6,6]_3\), construido a partir de la
incidencia excepcional, produce otra pantalla mediante
perforación de una coordenada. Denotemos por
\(p:C_W\to\FF_3^{11}\) ese mapa y por
\(G_{11}=p(C_W)\) su imagen.

\begin{proposition}[Código perforado y partición perfecta]
\label{iv:prop:codigo-once}
El mapa \(p\) es inyectivo. El código \(G_{11}\) tiene
parámetros \([11,6,5]_3\), y sus bolas de Hamming de radio
\(2\) particionan \(\FF_3^{11}\).
\end{proposition}
\begin{proof}
Un vector del núcleo de \(p\) tendría soporte contenido en
la única coordenada retirada. La distancia mínima \(6\)
de \(C_W\) excluye todo vector no nulo de esa forma.
Así, la dimensión sigue siendo \(6\). Perforar disminuye
el peso a lo sumo en una unidad, luego la distancia mínima
de la imagen es al menos \(5\). Las bolas de radio \(2\)
son por ello disjuntas. Cada una contiene
\[
 1+2\binom{11}{1}+2^2\binom{11}{2}
 =1+22+220=243
\]
vectores. Como \(|G_{11}|=3^6\) y
\(3^6\cdot243=3^{11}\), las bolas cubren el ambiente.
Para comprobar que la distancia es exactamente \(5\),
tómese un vector de peso \(3\). Su bola decodificadora
no puede tener centro cero y posee un centro de peso a lo
sumo \(3+2=5\); dicho centro debe tener peso al menos
\(5\). Se obtiene así una palabra de peso \(5\).
\end{proof}

El entero \(243=3^{11-6}\) es aquí el número de clases
del cociente de síndromes \(\FF_3^{11}/G_{11}\). Su
función se fija por la perforación y la distancia del
código; la igualdad de cardinal con otra fibra no sustituye
el mapa entre ambas. Esta pantalla combinatoria y la
siguiente pantalla real conservan sus operaciones propias.

\subsection{La dirección dodecafásica y la reducción de rango}

La construcción de \ref{exc:k-direccion} parte del registro
\(K\) ya obtenido por APP--TRIT--TPK y de la carta de
representación declarada sobre \(\RR^{12}\). En esa carta,
\begin{equation}
 V_{11}=\mathbf1^\perp,\qquad
 P_{11}=I_{12}-\frac1{12}J_{12},\qquad
 u_K=P_3P_{11}K.
 \label{iv:eq:p11-uk}
\end{equation}
Aquí \(J_{12}\) es la matriz de unos y \(P_3\) es el
proyector isótipo especificado por la acción de \(A_5\).
La suma finita que define \(P_3\), sus representantes y la
generación de \(K\) se conservan en la sección citada de
este mismo artículo. En particular,
\[
 \|u_K\|^2=\frac{6638585+2275584\sqrt5}{20}>0.
\]
La norma positiva define \(\ell_K=\RR u_K\), y el
proyector
\begin{equation}
 P_{10}=P_{11}-\frac{u_Ku_K^{\mathsf T}}{\|u_K\|^2}
 \label{iv:eq:p10}
\end{equation}
tiene imagen \(V_{10}=V_{11}\cap u_K^\perp\).

\begin{proposition}[Dos reducciones ortogonales]
\label{iv:prop:12-11-10}
Se cumplen
\[
 \RR^{12}=\RR\mathbf1\mathbin{\oplus^\perp}V_{11},\qquad
 V_{11}=\ell_K\mathbin{\oplus^\perp}V_{10},\qquad
 \dim V_{10}=10.
\]
Además, \(P_{10}P_{11}=P_{10}\), \(P_{10}u_K=0\), y
\(P_{10}\) es un proyector ortogonal.
\end{proposition}
\begin{proof}
La proposición \ref{prop:k-reduccion-diez}, demostrada antes de
construir las pantallas, establece la idempotencia, simetría, rango y
las dos identidades operatorias para este mismo \(u_K\) y este mismo
\(P_{10}\). La primera descomposición se obtiene de
\(P_{11}=I-J_{12}/12\); la segunda es la descomposición ortogonal
de aquella proposición. Así se conserva la demostración única de la
reducción y se explicitan aquí sus dos etapas.
\end{proof}

La primera reducción elimina exclusivamente el modo uniforme.
La segunda elimina la dirección determinada por \(K\) en
la representación fijada. Si se cambia simultáneamente la carta
por una matriz ortogonal \(S\), el vector se transforma por
\(u_K\mapsto Su_K\) y el proyector por
\(P_{10}\mapsto SP_{10}S^{-1}\). Esta covariancia conserva
la operación geométrica, aunque cambien sus coordenadas.

\subsection{Pantalla reticular lorentziana}

Sea \(\Lambda=N_\alpha\) el retículo de Leech construido
en \ref{exc:leech}. Se considera el plano hiperbólico integral
\(U=\ZZ e_+\oplus\ZZ e_-\), con
\[
 \langle e_+,e_+\rangle=\langle e_-,e_-\rangle=0,
 \qquad\langle e_+,e_-\rangle=1.
\]
El retículo \(L_{26}=\Lambda\oplus U\) es par y
unimodular, de firma \((25,1)\). Cada vector se escribe
\(y=\lambda+ae_++be_-\). Como
\(\langle y,e_+\rangle=b\),
\begin{equation}
 e_+^\perp=\Lambda\oplus\ZZ e_+,
 \qquad e_+^\perp/\ZZ e_+\cong\Lambda.
 \label{iv:eq:26-24}
\end{equation}
El mapa del cociente es
\(q_{24}(\lambda+ae_+)=\lambda\); es sobreyectivo y tiene
núcleo \(\ZZ e_+\). La disminución de rango en dos procede
de una restricción de ortogonalidad seguida por un cociente
isotrópico. Es, por tanto, una operación de tipo diferente a
\eqref{iv:eq:p10}.

\subsection{Morfismo conjunto de pantallas}

Con el dominio \(\mathscr D_0\) de \ref{iv:def:dualidad},
sean
\[
 \mathscr G_K=\{(v_{11},P_{10}v_{11}):v_{11}\in V_{11}\},
\]
\[
 \mathscr S_{\mathrm{HMT}}=\RR^{12}\times\mathscr D_0\times e_+^\perp,
 \qquad
 \mathscr S_M^{\mathrm{alg}}=\mathscr G_K\times\mathscr D_0\times\Lambda.
\]
\begin{definition}[Mapa algebraico de pantallas coordinadas]
\label{iv:def:rm-alg}
Se define
\begin{equation}
 \mathcal R_M^{\mathrm{alg}}(v,d,y)
 =\bigl((P_{11}v,P_{10}v),d,q_{24}(y)\bigr).
 \label{iv:eq:rm-alg}
\end{equation}
\end{definition}

\begin{theorem}[Isomorfismo cocientado y entrelazamiento de dualidad]
\label{iv:thm:pantallas-isomorfas}
\label{thm:k-pantallas-coordinadas}
El mapa \eqref{iv:eq:rm-alg} induce un isomorfismo
\[
 \overline{\mathcal R}_M^{\mathrm{alg}}:
 (\RR^{12}/\RR\mathbf1)\times\mathscr D_0
 \times(e_+^\perp/\ZZ e_+)
 \longrightarrow\mathscr S_M^{\mathrm{alg}}.
\]
Entrelaza las transformaciones que actúan por \(\mathsf T_0\)
en el factor \(\mathscr D_0\) y por la identidad en los
otros dos factores. La forma \(E_{R_0}(m,w)\) se conserva
bajo esa dualidad.
\end{theorem}
\begin{proof}
La aplicación \([v]\mapsto P_{11}v\) es un isomorfismo
de \(\RR^{12}/\RR\mathbf1\) sobre \(V_{11}\).
Como \(P_{10}P_{11}=P_{10}\), el mapa
\([v]\mapsto(P_{11}v,P_{10}v)\) identifica ese cociente
con \(\mathscr G_K\); su inversa toma la primera componente
y su clase. Por \eqref{iv:eq:26-24}, \(q_{24}\) induce el
isomorfismo reticular del tercer factor. El segundo se
transporta mediante la identidad. El producto de los tres
mapas prueba el isomorfismo.

Los proyectores y \(q_{24}\) no modifican \(d\), mientras
la dualidad actúa únicamente en \(d\). Por ello las dos
composiciones son idénticas. La conservación de la forma
es el teorema \ref{iv:thm:dualidad-escalar}.
\end{proof}

El grafo \(\mathscr G_K\) conserva la coordenada anterior
\(v_{11}\) junto con su proyección. El teorema demuestra
la equivalencia de las presentaciones cocientadas; la reducción
aislada \(V_{11}\to V_{10}\) conserva su núcleo
\(\ell_K\). Esta precisión establece qué información
retiene el isomorfismo y qué información elimina cada
proyección que lo compone.

\begin{corollary}[Pantallas sobre la familia elíptica]
El mismo morfismo se restringe a un isomorfismo
\[
 (\RR^{12}/\RR\mathbf1)\times\mathscr D_{\mathrm{el}}
 \times(e_+^\perp/\ZZ e_+)
 \ \cong\ \mathscr G_K\times\mathscr D_{\mathrm{el}}\times\Lambda,
\]
y conserva el entrelazamiento de dualidad.
\end{corollary}
\begin{proof}
El morfismo y su inversa conservan exactamente el segundo factor.
El corolario \ref{iv:cor:radio-especializado} prueba que la dualidad
preserva \(\mathscr D_{\mathrm{el}}\). Por tanto se restringen
tanto la biyección como el cuadrado de entrelazamiento ya demostrado.
\end{proof}

\subsection{Coordinación con la realización excepcional}

Sobre el dominio \(\mathcal X_{\mathrm{HMT}}\), con fase, orientación,
cociente, frontera, incidencia, gradación y memoria, las lecturas
\(\varepsilon_{\mathrm{exc}}\) y \(\varepsilon_{\mathrm{scr}}\) extraen
los datos excepcionales y de pantalla. El corredor construye el álgebra
graduada \(\mathfrak V(x)\cong V^\natural\), con su producto de vértices
y sus automorfismos; \(\mathfrak S_M(x)\) reúne las pantallas, cocientes
y operadores anteriores. La asignación conjunta
\(x\mapsto(\mathfrak V(x),\mathfrak S_M(x))\) tiene, por tanto, un
álgebra como primer componente y una familia de espacios y mapas como segundo.
La acción del Monstruo pertenece al primero; la reducción dirigida por \(K\)
y la dualidad compacta, al segundo. Las dos lecturas fijan la procedencia
común y cada construcción conserva su tipo.
El isomorfismo de esta sección establece la correspondencia de
pantallas con sus factores y su prueba
explícitos.

% Procedencia: integral 2249, §25.10--13 y cap. 89; propietarios 02_pantalla_accion,
% 37_cuerdas_branas_teoria_m, terna_espectral_teoria_m, capitulo_25_pantallas_teoria_m.
\section{Acción sobre trayectorias y realización dinámica undecadimensional}
\label{iv:accion-realizacion}

La extensión de la estructura a una dinámica dimensional requiere
transportar conjuntamente el estado, los canales y los datos de
frontera. Esta sección comienza por un funcional definido sobre
las trayectorias del TPK y por su ley exacta de composición.
La superficie de mundo y el codominio undecadimensional se
introducen después, con los mapas que deben conservar esa
composición y los operadores ya construidos.

\subsection{Descenso de una evolución y suficiencia del dato observado}

Sea \(q:\mathcal X_{\mathrm{HMT}}\to Y\) una lectura
sobreyectiva de los estados con memoria, y sea
\(U:\mathcal X_{\mathrm{HMT}}\to\mathcal X_{\mathrm{HMT}}\)
una evolución completa. Una evolución sobre los valores
observados existe precisamente cuando el dato conservado por
\(q\) basta para determinar el siguiente valor observado.

\begin{theorem}[Criterio exacto de descenso por fibras]
\label{iv:thm:descenso-fibras}
Existe una única aplicación \(\overline U:Y\to Y\)
que satisfaga \(qU=\overline Uq\) si y sólo si
\begin{equation}
 q(x)=q(y)\quad\Longrightarrow\quad q(Ux)=q(Uy).
 \label{iv:eq:descenso-fibras}
\end{equation}
\end{theorem}
\begin{proof}
Si existe \(\overline U\), aplicar la misma función a
\(q(x)=q(y)\) prueba la condición. Recíprocamente, se
define \(\overline U(q(x))=q(Ux)\). La condición
\eqref{iv:eq:descenso-fibras} garantiza independencia del
representante. La sobreyectividad de \(q\) define la
aplicación en todo \(Y\) y prueba su unicidad.
\end{proof}

El teorema permite decidir qué memoria debe conservar una
lectura. Cuando dos representantes tienen imágenes observadas
diferentes, se puede ampliar el dato retenido, restringir
el dominio o utilizar una correspondencia multivaluada.
Son elecciones matemáticas distintas. En particular, el
retorno de una coordenada de fase no autoriza a identificar
las evoluciones de dos historias que conserven memorias
diferentes.

\subsection{Cociclo de cambios direccionales y aritméticos}

Una trayectoria admisible del TPK lleva, en cada arista
\(e_j\), una dirección \(d_j\) y un tipo aritmético
\(t_j\). Para contabilizar los cambios en una concatenación,
se conserva también el dato precedente. La arista aumentada es
\(\widetilde e_j=(e_j,d_{j-1},t_{j-1})\); la pareja
\((d_0,t_0)\) forma parte de la frontera inicial. Definimos
\[
 w_{\mathrm{giro}}(\widetilde e_j)
       =\mathbf1_{\{d_j\ne d_{j-1}\}},\qquad
 w_{\mathrm{tipo}}(\widetilde e_j)
       =\mathbf1_{\{t_j\ne t_{j-1}\}}.
\]
Para \(\widetilde\gamma=(\widetilde e_1,\ldots,
\widetilde e_R)\), sean
\begin{equation}
 \mathbf I(\widetilde\gamma)=
 \left(\sum_{j=1}^R w_{\mathrm{giro}}(\widetilde e_j),
       \sum_{j=1}^R w_{\mathrm{tipo}}(\widetilde e_j)\right)
 =\bigl(N_{\mathrm{giro}},N_{\mathrm{tipo}}\bigr).
 \label{iv:eq:cociclo-accion}
\end{equation}
Dos trayectorias se componen cuando coinciden sus extremos
y el dato precedente de la segunda coincide con el dato
terminal de la primera.

\begin{proposition}[Aditividad con frontera conservada]
\label{iv:prop:accion-aditiva}
Para toda pareja composable,
\begin{equation}
 \mathbf I(\widetilde\gamma\widetilde\eta)
 =\mathbf I(\widetilde\gamma)+\mathbf I(\widetilde\eta).
 \label{iv:eq:accion-aditiva}
\end{equation}
Por tanto, \(\mathbf I\) es un cociclo aditivo con valores
en \(\mathbb N^2\) sobre la categoría de trayectorias
aumentadas. Para cada \(\lambda>0\), el funcional
\[
 I_\lambda(\widetilde\gamma)
 =N_{\mathrm{giro}}(\widetilde\gamma)
      +\lambda N_{\mathrm{tipo}}(\widetilde\gamma)
\]
es también aditivo.
\end{proposition}
\begin{proof}
La concatenación no cambia el dato precedente de ninguna
arista: el de la primera arista de \(\widetilde\eta\)
coincide, por la condición de composición, con el terminal
de \(\widetilde\gamma\). La suma sobre todas las aristas
se divide exactamente en las dos sumas de
\eqref{iv:eq:accion-aditiva}. Aplicar el funcional lineal
\((a,b)\mapsto a+\lambda b\) prueba la última afirmación.
\end{proof}

La información de frontera es parte efectiva de la prueba.
Si dos segmentos de una arista se evalúan por separado
ignorando el enlace entre ellos, ambos pueden tener cuenta
de giro nula aunque sus direcciones sean diferentes. El
camino concatenado contiene entonces un giro. Conservar
\((d_{j-1},t_{j-1})\) registra ese término en la primera
arista del segundo segmento y restituye la aditividad.

Para extremos y datos de frontera fijados, el principio de
interacción mínima del corpus selecciona minimizadores de
\(I_\lambda\) entre las rutas admisibles. El principio
es una regla de selección, mientras que el siguiente
resultado establece una propiedad de existencia finita.

\begin{proposition}[Existencia finita de minimizadores]
Si el conjunto \(\Gamma(a,b)\) de trayectorias admisibles
con la frontera fijada es finito y no vacío, existe
\(\gamma_*\in\Gamma(a,b)\) tal que
\(I_\lambda(\gamma_*)\leq I_\lambda(\gamma)\)
para todo \(\gamma\in\Gamma(a,b)\).
\end{proposition}
\begin{proof}
La imagen de \(\Gamma(a,b)\) por \(I_\lambda\) es
un conjunto finito no vacío de números reales, que posee
un elemento mínimo y al menos una preimagen.
\end{proof}

La proposición no selecciona el valor de \(\lambda\),
ni confunde existencia con unicidad. El valor y la regla
de selección utilizados en una realización determinada
deben formar parte de sus datos declarados. El funcional
aditivo recuperado aquí cuantifica cambios del propio
camino; su transporte a una acción dimensional constituye
una operación ulterior.

\subsection{Persistencia extendida y superficie de mundo}

La realización de una ruta como objeto extendido distribuye
su genealogía sobre una coordenada interna. Para una cuerda,
la historia tiene dos coordenadas locales \(\sigma^a\):
una espacial interna y otra temporal. Una realización se
describe por un mapa \(X:\Sigma\to\mathcal M\), una
métrica \(h_{ab}\) en la superficie de mundo y campos
de fondo en el espacio objetivo. Las condiciones sobre
\(\partial\Sigma\) forman parte del dominio variacional.
Para una brana de dimensión espacial \(p\), la historia
es un volumen de mundo de dimensión \(p+1\).

El codominio variacional de cuerda que utiliza el corpus
es la acción de Polyakov
\begin{equation}
 S_{\mathrm P}[X,h]
 =-\frac{T_{\mathrm s}}2\int_\Sigma d^2\sigma\,
 \sqrt{-h}\,h^{ab}g_{MN}(X)
             \partial_aX^M\partial_bX^N,
 \label{iv:eq:polyakov}
\end{equation}
donde \(T_{\mathrm s}\) es la tensión. La notación
\(T_{\mathrm s}\) la distingue de la dualidad
\(\mathsf T_0\), del transporte cardinal \(T_d\)
y de la evolución completa \(U\).

La composición de trayectorias aumentadas tiene ya la
propiedad demostrada en \ref{iv:prop:accion-aditiva}.
Una realización continua de esa composición debe
especificar el mapa de caminos a configuraciones,
su transporte de las condiciones de frontera y la
normalización que compara \(I_\lambda\) con la acción
continua. La integral \eqref{iv:eq:polyakov} precisa
el codominio y sus campos; sus parámetros no eligen
retroactivamente las emisiones ni los registros del TPK.

\subsection{Realización undecadimensional y operadores graduados}

La reducción dirigida por \(K\) proporciona
\(V_{11}=\ell_K\oplus V_{10}\), y el teorema
\ref{iv:thm:pantallas-isomorfas} coordina esa reducción
con la pantalla reticular y la dualidad. La realización
física se tipa mediante
\[
 \Xi_M:\mathscr S_M^{\mathrm{alg}}
                  \longrightarrow\mathcal X_{11}^{\mathrm{phys}}.
\]
El codominio lleva métrica \(g_{MN}\), tres-forma
\(C_{MNP}\), gravitino \(\psi_M\), acción,
restricciones de calibre y transformaciones de
supersimetría. Una realización operatoria incluye
\[
 \mathfrak S_M=(\mathcal H_{\mathrm{phys}},H,Q),\qquad
 \mathcal H_{\mathrm{phys}}=\mathcal H_+\oplus\mathcal H_-.
\]
La graduación se expresa mediante una involución
\(\Gamma\), igual a \(+I\) en \(\mathcal H_+\)
y a \(-I\) en \(\mathcal H_-\). La supercarga
es impar, \(\Gamma Q=-Q\Gamma\), y el
Hamiltoniano es no negativo. La relación
\(Q^2=H\) se entiende como igualdad de operadores,
incluyendo sus dominios.

Sea \(\mathcal H_{\mathrm{HMT}}\) una realización
hilbertiana de los estados genealógicos pertinentes,
con operadores \(\widehat H,\widehat Q\). Una
isometría \(W:\mathcal H_{\mathrm{HMT}}\to
\mathcal H_{\mathrm{phys}}\) debe transportar
dominios, graduación y observables. Las condiciones
operatorias se escriben
\begin{equation}
 W\widehat H=HW,\qquad
 W\widehat Q=QW,\qquad
 WU_T^{\mathrm{HMT}}=U_T^{\mathrm{phys}}W.
 \label{iv:eq:entrelazamiento-fisico}
\end{equation}

\begin{theorem}[Transporte compatible de la supercarga]
\label{iv:thm:supercarga-transporte}
Supóngase \(\widehat Q^{\,2}=\widehat H\), que
\(W D(\widehat Q)\subset D(Q)\), y que los
entrelazamientos de \eqref{iv:eq:entrelazamiento-fisico}
se satisfacen en sus dominios. Entonces, para
\(\psi\in D(\widehat H)=D(\widehat Q^{\,2})\),
\[
 W\psi\in D(Q^2),\qquad Q^2W\psi=HW\psi.
\]
Si además el transporte físico de los radios respeta
\eqref{iv:eq:conjugacion-h}, el sector compacto
transportado conserva su espectro bajo radios recíprocos.
\end{theorem}
\begin{proof}
Como \(\psi\in D(\widehat Q^2)\), tanto
\(\psi\) como \(\widehat Q\psi\) pertenecen
a \(D(\widehat Q)\). Por tanto, \(W\psi\in D(Q)\)
y \(QW\psi=W\widehat Q\psi\in D(Q)\). Se puede
aplicar \(Q\) de nuevo, y
\[
 Q^2W\psi=W\widehat Q^{\,2}\psi
          =W\widehat H\psi=HW\psi.
\]
La igualdad espectral del sector compacto se transporta
por la isometría y el entrelazamiento unitario de T;
sobre la imagen cerrada de \(W\), éste es una
equivalencia unitaria de las realizaciones correspondientes.
\end{proof}

El teorema conserva explícitamente los operadores y el
dominio sobre el que se prueba la identidad. La
graduación organiza dos sectores; la construcción de
campos fermiónicos y sus relaciones de ocupación posee
su propia estructura. Las distribuciones térmicas y el
catálogo de especies no se utilizan en la demostración
de las pantallas ni en la dualidad compacta de este
artículo.

\subsection{Compactificación compatible y conclusión del desarrollo}

Una compactificación del estado se describe mediante
una sección prolongable \(x_\infty\), un mapa
\(\kappa_{\mathrm{comp}}\) hacia el espacio de módulos
elegido y un transporte \(\mathcal U\) de los datos
de métrica, orientación, frontera y memoria. El
subíndice distingue este mapa de la lectura racional
\(\kappa\) del registro dodecafásico. Su compatibilidad
exige cuatro propiedades concretas:
\begin{enumerate}
\item las restricciones finitas de \(x_\infty\)
      pertenecen a las prolongaciones admisibles del TPK;
\item el mapa de compactificación entrelaza las dualidades,
      \(\kappa_{\mathrm{comp}}\mathsf T
        =\mathsf T_M\kappa_{\mathrm{comp}}\);
\item los observables y el espectro descienden a los
      cocientes elegidos conforme al criterio
      \ref{iv:thm:descenso-fibras};
\item la realización transporta la dirección
      \(\ell_K\), las condiciones de frontera y los
      dominios operatorios, y conserva las ecuaciones
      de acción y supersimetría declaradas.
\end{enumerate}

Las pruebas precedentes proporcionan la realización
hilbertiana de fase, el cierre tracial, las pantallas
con sus mapas de cociente, la involución de radio,
su equivalencia unitaria y el cociclo aditivo de
trayectorias. La realización undecadimensional
especifica los campos y las compatibilidades que
utiliza al recibir esa estructura. Su identificación
física completa conserva, además, las condiciones
de tensión, osciladores, amplitudes, anomalías y
límite de baja energía que enumera el corpus.
Estas condiciones no alteran las identidades
algebraicas demostradas: delimitan la afirmación
posterior que se obtiene al transportarlas.

La prolongación excepcional y la dimensional
permanecen vinculadas por su origen común. El
Monstruo actúa en la construcción graduada
\(V^\natural\); el Hamiltoniano compacto y las
pantallas reciben los datos de otros mapas del
mismo estado. Esta formulación conserva el
contenido positivo de la doble realización,
con sus operaciones, sus pruebas y la información
transmitida en cada etapa (teorema~\ref{iv:thm:supercarga-transporte}).

\HMTDivision{\(\mathsf U\)}{Conservación y recuperación de la información}{Refinamiento, recuperación, holonomía y simetrías de una evolución con registro completo.}
\HMTNarrativeHeading{Lo que conserva una transformación}

La teoría operatoria da una forma calculable a la conservación del
estado y de sus registros. Una lectura parcial obtiene una componente;
el complemento recoge la contribución que permite restituir el
balance completo. La suma de las normas cuadradas de la componente leída
y de su complemento reproduce la norma cuadrada del estado inicial.
Las fórmulas de recuperación especifican qué
datos y qué operadores hacen posible invertir la lectura.

Al iterar el transporte, cada etapa incorpora un complemento y
prolonga una componente terminal. El balance finito conserva
explícitamente ese terminal. El paso a profundidad arbitraria
estudia la convergencia y determina cuándo el archivo reúne toda la
norma. Este orden permite interpretar la memoria como una operación
de conservación y recuperación, con condiciones verificables.

La extensión exterior muestra cómo se transportan áreas y volúmenes
de la realización lineal. Las contribuciones mixtas participan en
los productos exteriores y en sus balances. El cambio de
representación conserva las relaciones que esos objetos miden.
La recuperación mantiene el enlace con la orientación y las formas
que utiliza cada prueba.

La monodromía y los lazos entre regímenes geométricos añaden una
dinámica determinada. El artículo compone transformaciones elípticas,
parabólicas e hiperbólicas, conserva su memoria y calcula su
respuesta. Las simetrías de la acción de transporte producen las
cargas variacionales de ese dominio. La invariancia de un observable
se formula, por su parte, mediante su relación con el transporte.
Estas dos formas de conservación se enlazan conservando sus
definiciones particulares.

Las realizaciones físicas recuperan después las secciones de acción,
los lectores de área y las escalas previamente generadas. El ejemplo
gaussiano permite calcular la pureza y el espectro de una lectura
reducida desde una preparación conjunta especificada. Sus
correlaciones conservan la información del estado completo. Este
resultado muestra una consecuencia operativa del argumento: el
transporte determina tanto lo que una observación recibe como la
información necesaria para comprender su relación con la totalidad.
\HMTMathematicalPaper
\section{Refinamiento de historias y conservación de la unidad}
\label{nuclear:1}
\subsection{Generación de emisiones y dominio de realización}

La construcción parte de las emisiones de APP y de sus prolongaciones enriquecidas. Una emisión conserva las dos hojas y sus divisiones

\[
\delta^\diamond=\operatorname{res}^\diamond+9\operatorname{quo}^\diamond,
\qquad \diamond\in\{\Sigma,\Pi\}.
\]

TRIT produce orientación y acarreo mediante

\[
\Delta(\epsilon)=\operatorname{or}(\epsilon)+3\kappa(\epsilon),
\qquad \operatorname{or}(\epsilon)\in\{-1,0,+1\}.
\]

TPK transporta la pareja completa de hojas y actualiza la memoria por

\[
U_{\epsilon_2\epsilon_1}=U_{\epsilon_2}U_{\epsilon_1},\qquad
\mathsf U_\epsilon(m)=\mathsf C_\epsilon m+\kappa(\epsilon),\qquad
\kappa(\epsilon_2\epsilon_1)
=\kappa(\epsilon_2)+\mathsf C_{\epsilon_2}\kappa(\epsilon_1).
\]

El cociclo hace asociativa la composición afín. La frontera de la ruta telescopa, pero su registro ordenado es una palabra ordenada, no la suma de sus entradas. El odómetro incrementa la memoria vertical al retornar la fase; el retorno no identifica los estados completos. La transferencia reducida del catálogo no sustituye la acción de las emisiones sobre el árbol; la construcción de \cref{iv:cont:historias,iv:cont:concatenacion} conserva esta distinción.

El estado enriquecido conserva las células intermedias, ambas hojas, residuos, cocientes, orientación, acarreo, reloj, memoria, frontera, ruta y registro ordenado. Una emisión admisible actualiza conjuntamente esos campos. El truncamiento borra exactamente su última contribución. Una prolongación neutra da una sección del truncamiento, pero registra un paso: no es la ausencia de historia. De ahí proceden la finitud, la no vacuidad y la sobreyectividad de todos los niveles, y después la supervivencia total.

El corte de este desarrollo es la realización analítica de ese sistema y su memoria. No intervienen valores numéricos objetivo ni constantes convencionales como entradas. Los coeficientes probabilísticos siguientes provienen únicamente del número de emisiones salientes efectivas.

\subsection{Refinamiento sobre la multisección completa}

Fijemos un prefijo enriquecido real \(x\) de profundidad \(k\). Escribimos

\[
H_n=\operatorname{Term}_{k,k+n}(x),\qquad
\rho_n:H_{n+1}\longrightarrow H_n,
\qquad \Omega_x=\varprojlim_n(H_n,\rho_n).
\]

Cada elemento de \(H_n\) es una prolongación completa hasta ese horizonte, con sus testigos de emisión. Para \(h\in H_n\), los hijos son \(h\epsilon\), para todas las emisiones \(\epsilon\in\operatorname{Out}(h)\). Se conservan multiplicidades de operaciones, incluso cuando coinciden datos terminales reducidos. Por \cref{iv:cont:medida}:

\[
p_h(\epsilon)=\frac1{d(h)}>0,\qquad
d(h)=\#\operatorname{Out}(h),\qquad
\mu_{n+1}(h\epsilon)=\mu_n(h)p_h(\epsilon).
\]

Por tanto,

\begin{equation}
\sum_{\epsilon\in\operatorname{Out}(h)}\mu_{n+1}(h\epsilon)=\mu_n(h),
\qquad (\rho_n)_*\mu_{n+1}=\mu_n.
\label{nuc01:eq:1}
\end{equation}

La proposición~\ref{iv:cont:medida} construye la única probabilidad \(\mu_x\) en \(\Omega_x\) con esas masas de cilindro; todo cilindro no vacío tiene masa positiva. La medida no selecciona una historia ni impone un filtro posterior a la existencia de las cinco construcciones.

\begin{proposition}[El refinamiento de observables es isométrico y unital]
\label{nuc01:pullback}

Sean \(\mathscr L_n=L^2(H_n,\mu_n)\) y

\[
R_n:\mathscr L_n\longrightarrow\mathscr L_{n+1},\qquad
(R_nf)(h\epsilon)=f(h).
\]

Entonces \(R_n^*R_n=I\), \(R_n1=1\), y su adjunto es

\begin{equation}
(E_ng)(h)=\sum_{\epsilon\in\operatorname{Out}(h)}p_h(\epsilon)g(h\epsilon).
\label{nuc01:eq:2}
\end{equation}

En particular, \(E_nR_n=I\), \(E_n1=1\), y \(R_nE_n\) es la proyección ortogonal sobre los observables que dependen sólo del prefijo.
\end{proposition}

\begin{proof}
Aplicando \eqref{nuc01:eq:1},

\[
\|R_nf\|_{n+1}^2
=\sum_h\mu_n(h)|f(h)|^2\sum_\epsilon p_h(\epsilon)
=\|f\|_n^2.
\]

El mismo agrupamiento en el producto interno da \eqref{nuc01:eq:2}. Las afirmaciones sobre la unidad son puntuales; la afirmación sobre la proyección se sigue de \(E_n=R_n^*\) y \(E_nR_n=I\). Además, \(R_n(fg)=R_nf\,R_ng\) y \(R_n\bar f=\overline{R_nf}\): en los observables finitos es también un homomorfismo unital de álgebras.
\end{proof}

Al iterar, \(R_{m,n}=R_{m-1}\cdots R_n\) coincide con el pullback por el truncamiento compuesto. Los pullbacks cilíndricos \(R_{\infty,n}\) realizan el límite inductivo de estos Hilbert en \(L^2(\Omega_x,\mu_x)\). Su unión es densa: los cilindros forman un álgebra que genera la sigma-álgebra, y las funciones simples de ese álgebra son densas en \(L^2\). No se requiere elegir una trayectoria para construir el límite.

\subsubsection{La misma isometría en la base de historias}

La identificación unitaria \(W_nf(h)=\sqrt{\mu_n(h)}f(h)\) lleva \(\mathscr L_n\) a \(\ell^2(H_n)\). Conjugando \(R_n\), se obtiene

\begin{equation}
\widehat R_n\delta_h
=\sum_{\epsilon\in\operatorname{Out}(h)}\sqrt{p_h(\epsilon)}\,\delta_{h\epsilon}.
\label{nuc01:eq:3}
\end{equation}

Los vectores de la derecha tienen soportes disjuntos para padres distintos, porque \(\rho_n(h\epsilon)=h\). Cada columna tiene norma uno por normalización. Esto demuestra directamente la isometría sobre las historias enriquecidas ya generadas, sin imponer que una proyección numérica conserve toda la memoria. El vector unidad se transforma en

\[
\Omega_n=\sum_{h\in H_n}\sqrt{\mu_n(h)}\,\delta_h,
\qquad \widehat R_n\Omega_n=\Omega_{n+1}.
\]

No se añade una copia artificial de \(h\) a un resultado para recuperarlo: \(h\epsilon\) es precisamente la prolongación que el dominio APP–TRIT–TPK ya produjo.

\subsection{Prolongación correlativa de las cinco operaciones}

La construcción de \cref{iv:cont:mapa-correlativo} define un único registro \(\mathfrak e_\epsilon\) por emisión. Sus operaciones espectral, solenoidal, de calibre, cohomológica y determinantal dependen de esa misma emisión, con iguales fuente, destino, orientación, acarreo, frontera, ruta y registro ordenado. La acción celular satisface

\[
\partial\operatorname{Cell}(\epsilon)=\operatorname{Cell}(\epsilon)\partial,
\qquad
\operatorname{Cell}(\epsilon_2\epsilon_1)
=\operatorname{Cell}(\epsilon_2)\operatorname{Cell}(\epsilon_1).
\]

El complejo y su línea determinante son los del mismo objeto; el relleno construido verifica \(\partial h_*=z_--z_+\). La totalidad del constructor no requiere que una trivialización especializada sea acíclica o invertible. Así se conserva la existencia conjunta antes de cualquier lector especializado.

La naturalidad de \cref{iv:rectangular-natural,iv:cont:memoria-natural} se expresa conjuntamente como

\begin{equation}
u_{(\ell,N'),(k,N)}\mathfrak D_{\ell,N'}(\widehat x_\ell)
=\mathfrak D_{k,N}(\tau_{\ell,k}\widehat x_\ell).
\label{nuc01:eq:4}
\end{equation}

Su prueba no identifica una arista fina aislada con una arista gruesa ficticia: restringe la historia de emisiones y conserva las acciones de sus entradas. El transporte compone sobre las mismas probabilidades de cilindro; los balances utilizan la esperanza condicional y su corrector; incidencia y puertos usan sus mapas celulares; homología y determinante utilizan el mismo mapa de complejos. No son cinco límites de historias escogidas por separado.

Hay dos direcciones que deben distinguirse. Los mapas de estado y de coeficientes van de fino a grueso. Las imágenes inversas \(R_n\) de la proposición~\ref{nuc01:pullback} van de los observables gruesos a los finos. La esperanza \(E_n\) vuelve en la dirección de reducción. Esta distinción permite realizar \eqref{nuc01:eq:4} analíticamente sin invertir sus dominios.

La construcción terminal de \cref{iv:cont:terminal,iv:cont:ensamblaje,iv:cont:determinantes} prolonga los terminales completos y conserva simultáneamente la telescopía polar, Gram y pérdidas, incidencia con medida, cono de retorno y datos Smith–Bockstein. La promoción de una de estas acciones a un operador acotado requiere la norma y la cota correspondiente (véase el paso al límite demostrado al final de esta sección). Una familia finita compatible no se convierte sólo por su nombre en un operador acotado del límite.

\subsection{Fibras de memoria: criterio de Gram y conservación de la unidad}

Para una realización no escalar, asignemos a cada prefijo \(h\) una fibra de Hilbert \(\mathscr M_h\), y a su emisión un transporte lineal tipado
\(I_\epsilon:\mathscr M_h\to\mathscr M_{h\epsilon}\). No se identifica este operador lineal con la actualización afín \(\mathsf U_\epsilon\) sin declarar la representación que la realiza.

Sobre \(\mathscr H_n=\bigoplus_{h\in H_n}\mathscr M_h\), se define

\[
(V_n\xi)_{h\epsilon}=\sqrt{p_h(\epsilon)}I_\epsilon\xi_h.
\]

La disjunción de hijos demuestra la identidad exacta de \cref{iv:cont:gram}:

\begin{equation}
\left.V_n^*V_n\right|_{\mathscr M_h}
=\sum_{\epsilon\in\operatorname{Out}(h)}p_h(\epsilon)I_\epsilon^*I_\epsilon.
\label{nuc01:eq:5}
\end{equation}

El criterio necesario y suficiente para \(V_n\) isométrico es que la suma de \eqref{nuc01:eq:5} sea \(I\) en cada padre. Que cada \(I_\epsilon\) sea isométrico es suficiente. También es necesario si previamente todos los \(I_\epsilon\) son contracciones: entonces cada \(I-I_\epsilon^*I_\epsilon\) es positivo, y una suma con pesos estrictamente positivos sólo puede anularse si se anula cada sumando. El criterio de Gram caracteriza la isometría de la aplicación total; la isometría de cada transporte constituye una condición suficiente, que resulta también necesaria en la clase de contracciones.

En las historias escalares de \eqref{nuc01:eq:3}, \(I_\epsilon=1\), por lo que la condición queda demostrada. En fibras analíticas adicionales se conserva \eqref{nuc01:eq:5} y se identifica el supuesto concreto que permite pasar a la isometría.

Si las fibras poseen vectores unitarios distinguidos \(u_h\) y \(I_\epsilon u_h=u_{h\epsilon}\), entonces

\begin{equation}
V_n\Bigl(\bigoplus_h\sqrt{\mu_n(h)}u_h\Bigr)
=\bigoplus_{h\epsilon}\sqrt{\mu_{n+1}(h\epsilon)}u_{h\epsilon}.
\label{nuc01:eq:6}
\end{equation}

La isometría por sí sola no afirma \eqref{nuc01:eq:6}: conservar la norma y conservar un vector unidad distinguido son propiedades distintas.

\subsection{Simetrías y lectores sobre el dominio común}

Una simetría reversible coherente del árbol transporta todas las emisiones y sus multiplicidades. La invariancia de las multiplicidades en \cref{iv:cont:medida} prueba que conserva los grados y las probabilidades de los cilindros. Puede llevar \(\Omega_x\) a \(\Omega_{gx}\); sólo si fija el prefijo actúa dentro de una misma fibra de prolongaciones.

Por ello

\[
(K_gf)(g\omega)=f(\omega)
\]

define una unitaria \(L^2(\Omega_x,\mu_x)\to L^2(\Omega_{gx},\mu_{gx})\). La naturalidad del truncamiento implica \(K_{g,n+1}R_n=R'_nK_{g,n}\), y por adjunción la igualdad correspondiente de esperanzas. En fibras de memoria se exige además el cuadrado

\[
S_{g,h\epsilon}I_\epsilon=I_{g\epsilon}S_{g,h}
\]

con \(S_{g,h}\) unitarios. Es el entrelazador que convierte una simetría del árbol en simetría de esa representación de memoria; no se atribuye por nominalidad a todo grupo excepcional posterior.

Para la doble lectura se toma el calibre inicial y su transporte prefijal declarado en \cref{exc:doble-lectura,exc:limite-inverso}. Se trabaja sobre una multisección común \(\Omega_x\) en la que los lectores estén definidos; para el arquimediano esto incluye los intervalos compactos encajados y la contracción de diámetros declarados en \eqref{eq:cilindro}. Es una premisa sobre la realización del lector, no un nuevo criterio de existencia del constructor conjunto. Las identidades finitas son

\begin{equation}
r_{n,m}Q_n=Q_m\rho_{n,m},\qquad
a_{n,m}R^{\rm arq}_n=R^{\rm arq}_m\rho_{n,m}.
\label{nuc01:eq:7}
\end{equation}

Determinan la lectura conjunta al límite. Si \(L=(\operatorname{ev}_{\rm arq},Q_\infty)\) y \(\nu=L_*\mu_x\), entonces

\[
C_L:L^2(Y,\nu)\longrightarrow L^2(\Omega_x,\mu_x),
\qquad C_Lf=f\circ L,
\]

es una isometría, pues la igualdad de normas es la definición de medida imagen. Conserva la unidad y sus rangos son los observables medibles respecto de la información publicada por \(L\). Para lectores finitos compatibles, \eqref{nuc01:eq:7} implica también la consistencia de las medidas imagen.

Esta isometría no equivale a que el lector distinga todos los estados. La construcción de \cref{exc:limite-inverso} distingue que igual corriente visible con distinta memoria oculta puede tener la misma incidencia. El adjunto de \(C_L\) es la esperanza respecto del lector; \(C_LC_L^*\) conserva exactamente esa parte observable. Si una realización ampliada \((L,M)\) es inyectiva, Borel y con inversa medible sobre su imagen —por ejemplo, una inyección continua desde el espacio compacto de historias a un espacio Hausdorff apropiado—, su pullback sí identifica todo el \(L^2\). La memoria requerida es la ya conservada en el estado, no una nueva selección retrospectiva.

\subsection{Teorema de paso al límite para la carga observable/memoria}

Esta sección tipa el enlace con una identidad de conservación de memoria. El refinamiento canónico de observables proporciona los mapas isométricos necesarios. Una realización adicional en fibras de memoria utiliza el criterio de Gram de \eqref{nuc01:eq:5}.

Sean dos sistemas inductivos de Hilbert \((\mathscr H_n,R_n)\) y \((\mathscr K_n,S_n)\), con inclusiones isométricas. Para cada \(n\), sean

\[
J_n:\mathscr H_n\to\mathscr K_n\text{ isométrico},\quad
P_n=J_nJ_n^*,\quad U_n\text{ unitario en }\mathscr K_n,
\]

y \(\widetilde A_n\) autoadjunto acotado, con \(\sup_n\|\widetilde A_n\|<\infty\). Definimos

\[
A_n=J_n^*\widetilde A_nJ_n,\quad
T_n=J_n^*U_nJ_n,\quad
\eta_n=(I-P_n)U_nJ_n.
\]

Supongamos la conservación y separación a cada nivel,

\begin{equation}
U_n^*\widetilde A_nU_n=\widetilde A_n,
\qquad [\widetilde A_n,P_n]=0,
\label{nuc01:eq:8}
\end{equation}

y los siguientes cuadrados de refinamiento:

\begin{equation}
J_{n+1}R_n=S_nJ_n,\qquad
P_{n+1}S_n=S_nP_n,
\label{nuc01:eq:9}
\end{equation}

\begin{equation}
U_{n+1}S_n=S_nU_n,\qquad
\widetilde A_{n+1}S_n=S_n\widetilde A_n.
\label{nuc01:eq:10}
\end{equation}

\begin{theorem}
Estas familias inducen operadores acotados en los límites inductivos, con \(J_\infty\) isométrico y \(U_\infty\) unitario, y conservan

\begin{equation}
A_\infty
=T_\infty^*A_\infty T_\infty
 +\eta_\infty^*\widetilde A_\infty\eta_\infty.
\label{nuc01:eq:11}
\end{equation}

La misma identidad vale a cada profundidad finita. Si \(\widetilde A_n\geq0\), la carga retenida es positiva. Sin positividad se conserva una forma autoadjunta, no necesariamente una energía positiva.
\end{theorem}

\begin{proof}
Descomponemos \(U_nJ_n=J_nT_n+\eta_n\). Los dos términos pertenecen a \(P_n\mathscr K_n\) y \((I-P_n)\mathscr K_n\). La conmutación de \eqref{nuc01:eq:8} anula los términos cruzados de la forma \(\widetilde A_n\). Aplicando la otra identidad de \eqref{nuc01:eq:8},

\[
A_n=J_n^*U_n^*\widetilde A_nU_nJ_n
=T_n^*A_nT_n+\eta_n^*\widetilde A_n\eta_n.
\]

La segunda igualdad de \eqref{nuc01:eq:9}, junto con la primera, implica
\(J_{n+1}^*S_n=R_nJ_n^*\): basta escribir
\(P_{n+1}S_n=S_nJ_nJ_n^*=J_{n+1}R_nJ_n^*\) y multiplicar por \(J_{n+1}^*\).
Por \eqref{nuc01:eq:9}–\eqref{nuc01:eq:10},

\[
T_{n+1}R_n=R_nT_n,\qquad
\eta_{n+1}R_n=S_n\eta_n,\qquad
A_{n+1}R_n=R_nA_n.
\]

Así los operadores se definen sobre las uniones algebraicas de los sistemas. Las isometrías y la cota uniforme permiten extenderlos por continuidad a las completaciones. La compatibilidad de \(U_n^*\) también se deduce de la de \(U_n\), porque ambos niveles son unitarios; el límite es, por ello, unitario y no sólo isométrico. Las proyecciones \(P_n\) inducen una proyección cuyo rango es la clausura de la unión de los rangos de \(J_n\); es \(J_\infty J_\infty^*\). Todas las identidades persisten primero en los vectores de nivel finito y después por densidad. Esto prueba \eqref{nuc01:eq:11}.
\end{proof}

La compatibilidad de \(J_n\) sola no sustituye la segunda igualdad de \eqref{nuc01:eq:9}. Un refinamiento que conserve la inclusión observable pero mezcle su complemento con ella puede alterar el defecto de memoria.

\subsubsection{Realización nonádica explícita compatible}

Supongamos \(C_n\) unitario en \(\mathscr H_n\), natural respecto de \(R_n\), y \(A_n\) autoadjunto acotado, uniformemente acotado, natural y con \([A_n,C_n]=0\). Sobre \(\mathscr K_n=\mathscr H_n^{\oplus9}\), definimos

\[
J_nv=\tfrac13(v,\ldots,v),\quad
S_n=R_n^{\oplus9},\quad
U_n(v_1,\ldots,v_9)=(C_nv_9,v_1,\ldots,v_8),
\quad \widetilde A_n=I_9\otimes A_n.
\]

El factor \(1/3\) normaliza las nueve componentes. Las fórmulas verifican directamente \eqref{nuc01:eq:8}–\eqref{nuc01:eq:10}, incluida la proyección. La compresión es \(T_n=(8I+C_n)/9\), y

\begin{equation}
\eta_n^*\widetilde A_n\eta_n
=A_n-T_n^*A_nT_n
=\frac8{81}(I-C_n)^*A_n(I-C_n).
\label{nuc01:eq:12}
\end{equation}

La última igualdad usa \(C_n^*A_nC_n=A_n\). Por el teorema, \eqref{nuc01:eq:12} vale también al límite. El refinamiento nonádico o el refinamiento completo de historias transporta la carga y su memoria cuando \(C_n\) y \(A_n\) satisfacen los cuadrados indicados.

Para \(A=I\), la telescopía finita es

\[
\|v\|^2=\|T^Nv\|^2+\sum_{j=0}^{N-1}\|\eta T^jv\|^2.
\]

Por tanto \(v\mapsto(\eta v,\eta Tv,\ldots,\eta T^{N-1}v,T^Nv)\) es una isometría. La fórmula de carga correspondiente sustituye las normas por las formas de \(A\) y \(\widetilde A\). No se elimina el término terminal sin probar su convergencia a cero; la compatibilidad de profundidad no prueba por sí sola esa estabilidad temporal.

Si \(D_{N,n}\) denota esa isometría a profundidad \(n\) y horizonte temporal \(N\), los cuadrados anteriores prueban además

\[
(S_n^{\oplus N}\oplus R_n)D_{N,n}
=D_{N,n+1}R_n.
\]

En efecto, cada componente cumple \(\eta_{n+1}T_{n+1}^jR_n=S_n\eta_nT_n^j\), y la terminal cumple \(T_{n+1}^NR_n=R_nT_n^N\). La acumulación de memoria a horizonte finito y el refinamiento de resolución conmutan. El mismo cálculo prueba equivariancia para una simetría representada por unitarios que entrelacen \(J,U,P,\widetilde A\); no extiende la acción a simetrías para las que esos cuadrados no se hayan construido.

\begin{HMTCompactSection}
\section{Conservación de observables y registro completo}
\label{nuclear:2}
\subsection{Balance de un observable conservado}

Sean \(\mathcal H,\mathcal K\) las realizaciones de un espacio de lectura y de su extensión enriquecida. Sean \(J:\mathcal H\to\mathcal K\) una isometría, \(P=JJ^*\), \(U:\mathcal K\to\mathcal K\) un transporte acotado y \(\widetilde A\) un observable acotado autoadjunto. Supóngase

\[
U^*\widetilde A U=\widetilde A,\qquad [\widetilde A,P]=0.
\]

Definimos

\[
A=J^*\widetilde A J,\qquad T=J^*UJ,
\qquad \eta=(I-P)UJ.
\]

\begin{proposition}
\label{nuc02:balance}
La carga se descompone exactamente entre lectura y complemento:

\begin{equation}
\boxed{A=T^*AT+\eta^*\widetilde A\eta.}\label{nuc02:eq:1}
\end{equation}
\end{proposition}

\begin{proof}
La identidad \(UJ=JT+\eta\) es una descomposición ortogonal. Al expandir \(J^*U^*\widetilde A UJ\), los términos cruzados son cero porque \(\widetilde A\) reduce los subespacios \(P\mathcal K\) y \((I-P)\mathcal K\). El término izquierdo es \(A\); los dos términos diagonales son los de \eqref{nuc02:eq:1}. Esto prueba también la identidad de las formas sesquilineales para dos estados distintos. Si \(\widetilde A\geq0\), ambos sumandos son positivos. Para \(\widetilde A=I\), la hipótesis de conservación es \(U^*U=I\).
\end{proof}

Si la lectura mezcla sectores del observable, el balance conserva adicionalmente

\[
T^*J^*\widetilde A\eta+\eta^*\widetilde A JT.
\]

Estos términos representan la interferencia entre las dos residencias de la carga. La condición de reducción escrita arriba garantiza su anulación para toda la descomposición. En realizaciones particulares pueden anularse también por otras razones, por ejemplo cuando \(\eta=0\).

\subsection{Especialización a la conexión nonádica}

Sea \(C\) la realización unitaria del transporte de memoria sobre \(\mathcal H\). En \(\mathcal H^9\), el transporte de las nueve fases y su inclusión uniforme son

\[
S_C(v_0,\ldots,v_8)=(Cv_8,v_0,\ldots,v_7),
\qquad Jv=\tfrac13(v,\ldots,v).
\]

El factor \(1/3\) resulta de la normalización de nueve copias. La compresión y su complemento son

\[
T=\frac{8I+C}{9},\qquad
\eta v=\frac1{27}\bigl(8(C-I)v,-(C-I)v,\ldots,-(C-I)v\bigr).
\]

Si \(A=A^*\) es acotado y \([A,C]=0\), el observable \(\widetilde A=I_9\otimes A\) satisface las hipótesis de la proposición~\ref{nuc02:balance}. Por tanto,

\begin{equation}
\boxed{A=T^*AT+\frac8{81}(I-C)^*A(I-C).}\label{nuc02:eq:2}
\end{equation}

Los coeficientes proceden del reparto entre las nueve componentes: \(8^2+8=72\), dividido por \(27^2\), da \(8/81\). La identidad es el balance sesquilineal ya documentado en el artículo espectral, expresado aquí para un observable general.

La iteración telescópica proporciona, para cada \(N\geq1\),

\begin{equation}
A=T^{*N}AT^N+
\sum_{j=0}^{N-1}T^{*j}\eta^*\widetilde A\eta T^j.\label{nuc02:eq:3}
\end{equation}

Para \(A=I\), la aplicación

\[
V_Nv=(T^Nv,\eta v,\eta Tv,\ldots,\eta T^{N-1}v)
\]

es isométrica. El paso \(N\to N+1\) descompone exclusivamente el terminal \(T^Nv\) en su nueva lectura y complemento; mantiene todos los registros previos. Así, la conservación se verifica a cualquier profundidad.

Distinción de operaciones. Se cumple \(S_C^9=I_9\otimes C\). La iteración \(T^N\) describe compresiones sucesivas con almacenamiento de sus complementos. El transporte enriquecido sin compresiones intermedias es \(S_C^N\). Ambas operaciones tienen dominios definidos y funciones diferentes; la prueba conserva esa distinción.

\subsection{Sector terminal y registro a profundidad ilimitada}

Sea \(P_{\rm fix}\) la proyección ortogonal sobre \(\ker(I-C)\).

\begin{theorem}
Para todo transporte unitario \(C\), la compresión nonádica satisface

\begin{equation}
\underset{N\to\infty}{\operatorname{s-lim}}T^N=P_{\rm fix},\qquad
\boxed{I=P_{\rm fix}+\sum_{j=0}^{\infty}T^{*j}\eta^*\eta T^j.}\label{nuc02:eq:4}
\end{equation}

La serie converge en la topología fuerte. Para cada observable acotado autoadjunto \(A\) que conmuta con \(C\), se obtiene además

\begin{equation}
\boxed{A=P_{\rm fix}AP_{\rm fix}
+\sum_{j=0}^{\infty}T^{*j}\eta^*(I_9\otimes A)\eta T^j.}\label{nuc02:eq:5}
\end{equation}
\end{theorem}

\begin{proof}
En la representación espectral del operador unitario \(C\), la compresión corresponde a \(t(z)=(8+z)/9\). Sobre el círculo unitario,

\[
|t(z)|^2=1-\frac8{81}|1-z|^2.
\]

Por tanto \(t(z)^N\to0\) para \(z\ne1\), y \(t(1)^N=1\). La convergencia dominada respecto de cada medida espectral vectorial prueba \(T^N\to P_{\rm fix}\) fuertemente. El mismo argumento vale para \(T^{*N}\). Como \(A\) es acotado, \(T^{*N}AT^N\to P_{\rm fix}AP_{\rm fix}\) fuertemente. El límite de \eqref{nuc02:eq:3} da \eqref{nuc02:eq:5}, y su especialización \(A=I\) da \eqref{nuc02:eq:4}.
\end{proof}

En particular,

\[
V_\infty v=(P_{\rm fix}v,\eta v,\eta Tv,\eta T^2v,\ldots)
\]

es una isometría. Conserva el producto interior completo, además de la norma.

\begin{corollary}[Memoria bilateral]
Para el desplazamiento \(Ce_n=e_{n+1}\) en \(\ell^2(\mathbb Z)\), \(P_{\rm fix}=0\): una sucesión fija sería constante y su única realización cuadrado-sumable es cero. En consecuencia, toda la norma inicial se distribuye en el registro de complementos. Esto precisa el comportamiento ilimitado de la realización bilateral de memoria presente en el corpus.
\end{corollary}

\begin{corollary}[Calendario finito]
En un ciclo de \(d\) posiciones, la sucesión uniforme determina un subespacio fijo unidimensional. Ese sector permanece como terminal y el complemento acumula el resto. El calendario finito y el registro bilateral poseen, por tanto, terminales diferentes.
\end{corollary}

En el caso bilateral, la convergencia fuerte puede coexistir con \(\|T^N\|=1\) para todo \(N\). La demostración no supone una tasa uniforme de contracción. También distingue esta compresión del modo \((5/9)C\), cuya contracción geométrica es uniforme y está estudiada separadamente en las fuentes.

\subsection{Conservación de simetría y transporte de carta}

Sea \(R\) unitario y definamos \(C'=RCR^*\), \(A'=RAR^*\). La fórmula explícita de la compresión da

\begin{equation}
T(C')R=RT(C),\qquad
\eta(C')R=R^{\oplus9}\eta(C).\label{nuc02:eq:6}
\end{equation}

\begin{proof}
La primera igualdad se obtiene distribuyendo \((8I+RCR^*)R/9\). La segunda resulta de \((C'-I)R=R(C-I)\) en cada una de las nueve componentes de \(\eta\).
\end{proof}

La iteración de \eqref{nuc02:eq:6} entrelaza los registros finitos. Además, \(P_{\rm fix}(C')=RP_{\rm fix}(C)R^*\), por transporte del núcleo de \(I-C\), de modo que también queda entrelazado el registro ilimitado. Así, el balance de observables es covariante cuando se transportan conjuntamente la carta, el operador y el observable. Las simetrías que conmutan con \(C\) actúan además en una misma carta fija. Esta distinción permite incorporar las transformaciones excepcionales con sus marcos, banderas y orientaciones efectivamente transportados.

La composición de \cref{lem:k-hadamard,exc:entrelazador,exc:isometria} utiliza la matriz de Hadamard del registro, \(H_{12}^2=4I\), y la incidencia \(\mathcal I\) de las doce posiciones en las 132 hexadas:

\[
J_H=H_{12}/2,\qquad
\mathcal I^*\mathcal I=36I+30\mathbf1\mathbf1^*,\qquad
U_W=\mathcal I/6\big|_{\mathbf1^\perp}.
\]

Ambos operadores normalizados son isométricos en los dominios indicados. La composición \(B=U_WJ_H\) sobre \(W=J_H(\mathbf1^\perp)\) conserva productos interiores y entrelaza la representación transportada de \(M_{12}\) con la acción en las hexadas. El registro entero \(K\) conserva, además, la inversión integral \(H_{12}/4\) sobre su red de congruencias; esa inversión y la normalización isométrica tienen funciones distintas. La demostración, la prolongación APP–Witt–Leech y sus dominios están desarrollados aquí.

En esta composición, la conservación de incidencia tiene un contenido operatorio concreto. Su extensión por otros transportes usa \eqref{nuc02:eq:6} y los entrelazadores efectivos de las fuentes, respetando sus grupos y dominios respectivos.

\subsubsection{Transporte íntegro de las doce coordenadas de K}

La separación del subespacio centrado puede completarse manteniendo también su media. Sean \(P_0=I-\mathbf1\mathbf1^*/12\), \(\mu(k)=\mathbf1^*k/12\) y

\[
F(k)=\left(\mu(k),\frac16\mathcal I P_0k\right)\in\mathbb R\oplus E_{\rm inc}.
\]

Equipamos el codominio con \(\|(\mu,y)\|^2=12|\mu|^2+\|y\|^2\). El Gram de incidencia demuestra

\begin{equation}
\|F(k)\|^2=\|k\|^2,\qquad
\boxed{F^{-1}(\mu,y)=\mu\mathbf1+\frac16P_0\mathcal I^*y.}\label{nuc02:eq:6a}
\end{equation}

En efecto, \(\mathcal I^*\mathcal I P_0=36P_0\), de modo que la fórmula inversa recupera \(P_0k\), y el primer término recupera su media. La acción de \(M_{12}\) es trivial sobre la media y permuta las hexadas; por ello \(F\) es equivariante. Componer \(FJ_H\) transporta las coordenadas firmadas normalizadas. Así se conservan las doce coordenadas del registro, no sólo su parte centrada. El dominio de esta reconstrucción es el registro dodecafásico; la memoria completa de la historia mantiene sus otras coordenadas y su prolongación.

\subsubsection{Especialización al operador excepcional de orden tres}

La construcción excepcional produce el operador ortogonal \(g_c\) sobre la realización del vecino de Leech, mediante las doce fibras APP–Witt de tipo \(A_2\), y demuestra \(g_c^3=I\) y \(\ker(I-g_c)=0\). Su construcción, entrelazamiento y preservación reticular están localizados en \cref{nuclear:fibras-app-witt,nuc04:c3}.

Construimos ahora, como composición operatoria de esa salida con la carta de nueve fases,

\[
T_3=\frac{8I+g_c}{9},\qquad
\eta_3=(I-JJ^*)S_{g_c}J.
\]

\begin{corollary}
Esta composición satisface

\begin{equation}
\boxed{T_3^*T_3=\frac{19}{27}I,\qquad
\eta_3^*\eta_3=\frac8{27}I,\qquad
\|T_3^Nv\|^2=\left(\frac{19}{27}\right)^N\|v\|^2.}\label{nuc02:eq:6b}
\end{equation}
\end{corollary}

\begin{proof}
El vector \((I+g_c+g_c^2)v\) es fijo por \(g_c\); la ausencia de vectores fijos implica \(I+g_c+g_c^2=0\). Como \(g_c^*=g_c^2\), tenemos \(g_c+g_c^*=-I\). Sustituyendo,

\[
T_3^*T_3=\frac{65I+8(g_c+g_c^*)}{81}
=\frac{57}{81}I=\frac{19}{27}I.
\]

El balance \eqref{nuc02:eq:2} con \(A=I\) da el segundo coeficiente. La aplicación reiterada de la igualdad de normas da la potencia para todo \(N\).
\end{proof}

El registro ilimitado es isométrico y su terminal tiende a cero con la tasa explícita de \eqref{nuc02:eq:6b}. Aquí la estructura excepcional determina los dos coeficientes complementarios; ambos proceden de la relación polinómica de \(g_c\) y de las nueve fases. Esta es una composición matemática desarrollada en esta investigación. Su definición no identifica \(g_c\) con \(\Gamma_9\), ni afirma que \(S_{g_c}\) sea el operador temporal canónico del TPK: conserva la distinción entre prolongación excepcional y holonomía temporal mientras las compone explícitamente.

\subsection{Refinamiento de historias y compatibilidad de los balances}

Para las historias enriquecidas de profundidad \(n\), sea \(p_h(\varepsilon)\) la probabilidad de una extensión superviviente, normalizada por \(\sum_\varepsilon p_h(\varepsilon)=1\). El refinamiento canónico sobre las bases de historias es

\[
R_n\delta_h=\sum_{\varepsilon\in\operatorname{Out}(h)}
\sqrt{p_h(\varepsilon)}\,\delta_{h\varepsilon}.
\]

Las extensiones de historias diferentes tienen prefijos diferentes y soportes disjuntos. Por ello \(R_n^*R_n=I\). Esta es la realización isométrica de la consistencia proyectiva de la medida: los coeficientes proceden de los pesos de supervivencia ya construidos.

Para transportar \eqref{nuc02:eq:1} entre profundidades, se conservan conjuntamente las inclusiones, las proyecciones, los transportes y los observables. Con isometrías de refinamiento \(R_n,\widetilde R_n\), los cuadrados necesarios son

\[
J_{n+1}R_n=\widetilde R_nJ_n,\quad
P_{n+1}\widetilde R_n=\widetilde R_nP_n,\quad
U_{n+1}\widetilde R_n=\widetilde R_nU_n,\quad
\widetilde A_{n+1}\widetilde R_n=\widetilde R_n\widetilde A_n.
\]

Estos cuadrados implican \(T_{n+1}R_n=R_nT_n\) y \(\eta_{n+1}R_n=\widetilde R_n\eta_n\). Con cotas uniformes, las identidades pasan al límite inductivo de las realizaciones. Las historias, en cambio, forman el sistema proyectivo de prefijos; las direcciones contrarias de ambos tipos de mapas quedan diferenciadas. Se incluyen las pruebas de medida, Gram y paso al límite.

\subsection{Especialización variacional: acción y orientación}

La acción elíptica de \cref{iv:ellipse:section,nuclear:7} proporciona una especialización física del principio, con los parámetros construidos en sus etapas HMT. En sus coordenadas canónicas,

\[
H_\eta(Q,P)=\frac{\omega}{2}
\left(e^{-\eta}Q^2+e^\eta P^2\right),\qquad
I_\eta=H_\eta/\omega.
\]

La transformación canónica \(Q=e^{\eta/2}q\), \(P=e^{-\eta/2}p\) convierte \(I_\eta\) en \((q^2+p^2)/2\). Las rotaciones de \((q,p)\), transportadas a \((Q,P)\), constituyen una simetría continua. Para el lagrangiano \(L=P\dot Q-\omega I_\eta\), su carga de Noether es \(I_\eta\). En el contorno de la fuente \(I_\eta=\hbar\), la acción orientada es \(\oint P\,dQ=2\pi\hbar=h\), con la orientación de la órbita hamiltoniana. La prueba distingue la transformación simpléctica y la antisimpléctica: pueden conservar la misma energía y actuar con signos diferentes sobre la acción orientada. La demostración variacional completa se presenta en \cref{nuclear:7}.

\subsubsection{Actualización de la forma y término de conexión}

La matriz de forma \(M_\eta=\operatorname{diag}(e^{\eta/2},e^{-\eta/2})\) permite desarrollar un transporte entre elipses con parámetros diferentes. Fijamos \(J_2=\begin{pmatrix}0&-1\\1&0\end{pmatrix}\) y \(R(\theta)=\exp(-\theta J_2)\), la rotación hamiltoniana en el plano \((q,p)\). El transporte

\[
z'=M_{\eta'}R(\delta\theta)M_\eta^{-1}z,
\qquad z=(Q,P)^\mathsf T,
\]

es simpléctico y satisface exactamente \(I_{\eta'}(z')=I_\eta(z)\). La prueba consiste en llevar ambos estados a sus coordenadas normalizadas y usar la conservación de \(q^2+p^2\) por la rotación. Conserva la acción mientras cambia la forma.

Para una realización diferenciable \(\eta(t)\), con avance angular \(\dot\theta=\omega\), la diferenciación de ese transporte produce

\[
\dot Q=\omega e^\eta P+\tfrac12\dot\eta Q,\qquad
\dot P=-\omega e^{-\eta}Q-\tfrac12\dot\eta P.
\]

Su generador hamiltoniano, obtenido de estas dos ecuaciones, es

\begin{equation}
\boxed{H_{\rm trans}=\omega I_\eta+\frac{\dot\eta}{2}QP.}\label{nuc02:eq:7}
\end{equation}

El segundo sumando es el término de conexión de la forma variable. Con la convención \(\{f,g\}=\partial_Qf\partial_Pg-\partial_Pf\partial_Qg\),

\[
\partial_tI_\eta=\frac{\dot\eta}{2}(-e^{-\eta}Q^2+e^\eta P^2),
\quad
\left\{I_\eta,\frac{\dot\eta}{2}QP\right\}
=\frac{\dot\eta}{2}(e^{-\eta}Q^2-e^\eta P^2).
\]

Ambas contribuciones se cancelan y \(dI_\eta/dt=0\). Para el transporte prescrito, \eqref{nuc02:eq:7} es único salvo una función del tiempo. La realización diferenciable es una especialización declarada del transporte entre formas; el resultado discreto anterior vale directamente entre dos parámetros generados. La actualización TPK específica conserva su propia regla para seleccionar esos parámetros y sus tiempos.

Así se obtiene una conexión precisa: un observable conservado por la dinámica completa se distribuye entre lectura y memoria mediante \eqref{nuc02:eq:1}. Cuando ese observable es el generador de una simetría continua de la acción construida, la cantidad distribuida es una carga de Noether. La norma, la carga de una simetría y el funcional energético quedan relacionados por sus operadores específicos.

Una especialización acotada efectiva se obtiene del calendario de 108 estados. Su Hamiltoniano realizado es \(H_{\rm clk}=\hbar\omega_0N_{108}\), con \(N_{108}=\sum_{k=0}^{107}kP_k\) y \(P_k\) sus proyectores de Fourier. La acción del reloj tiene carga \(\hbar\langle\psi,N_{108}\psi\rangle=E_{\rm clk}/\omega_0\). Los transportes y las inclusiones que satisfacen las hipótesis de \eqref{nuc02:eq:1} distribuyen esa carga entre lectura y complemento. Aquí todos los operadores son finitos y acotados. La acción elíptica de \eqref{nuc02:eq:7} se ha probado en su realización clásica simpléctica; una cuantización con observables no acotados conserva, además, sus condiciones de dominio propias.

\end{HMTCompactSection}
\section{Recuperación del estado y de su dinámica}
\label{nuclear:3}
\subsection{Representación recuperable del transporte}

Sea \(C\) la realización unitaria del transporte de memoria sobre \(\mathcal H\). La carta de nueve fases, su inclusión uniforme y su compresión son

\[
S_C(v_0,\ldots,v_8)=(Cv_8,v_0,\ldots,v_7),\qquad
Jv=(v,\ldots,v)/3,\qquad T=(8I+C)/9.
\]

Sea \(\eta=(I-JJ^*)S_CJ\). Su acción es

\[
\eta v=\frac1{27}\bigl(8(C-I)v,-(C-I)v,\ldots,-(C-I)v\bigr).
\]

La normalización de nueve copias determina \(1/3\); el reparto de la costura determina \(8/81\) en

\begin{equation}
I-T^*T=\eta^*\eta=\frac8{81}(I-C)^*(I-C).
\label{nuc03:eq:1}
\end{equation}

Escribimos \(P=\operatorname{proj}\ker(I-C)\). El teorema precedente prueba la convergencia fuerte \(T^N\to P\). Por tanto,

\begin{equation}
\mathcal Vv=(Pv,\eta v,\eta Tv,\eta T^2v,\ldots)
\quad\text{en}\quad
\mathcal R=P\mathcal H\oplus\ell^2(\mathbb N_0;\mathcal H^9)
\label{nuc03:eq:2}
\end{equation}

es una isometría. Denotamos su imagen cerrada por \(\mathcal M\).

En el espacio de registros definimos el desplazamiento de bloques

\begin{equation}
\mathcal B(a,m_0,m_1,m_2,\ldots)=(a,m_1,m_2,\ldots).
\label{nuc03:eq:3}
\end{equation}

El primer componente conserva el sector fijo; los demás componentes son los registros vectoriales completos, incluidas sus correlaciones.

\begin{theorem}
\label{nuc03:transporte}
El transporte original tiene la representación exacta

\begin{equation}
\boxed{\mathcal V C=(9\mathcal B-8I)\mathcal V.}
\label{nuc03:eq:4}
\end{equation}

La restricción \(\mathcal C=(9\mathcal B-8I)|_{\mathcal M}\) es unitaria sobre \(\mathcal M\), y

\begin{equation}
C=\mathcal V^*\mathcal C\mathcal V.
\label{nuc03:eq:5}
\end{equation}
\end{theorem}

\begin{proof}
Como \(PT=P\), la definición~\eqref{nuc03:eq:2} da componente a componente

\[
\mathcal VT v=(Pv,\eta Tv,\eta T^2v,\ldots)=\mathcal B\mathcal Vv.
\]

La identidad algebraica \(C=9T-8I\) implica \eqref{nuc03:eq:4}. La isometría \(\mathcal V:\mathcal H\to\mathcal M\) es sobreyectiva por definición de su imagen. La igualdad \eqref{nuc03:eq:4} identifica la restricción con \(\mathcal V C\mathcal V^{-1}\), que es unitaria porque \(C\) lo es. Esto demuestra además que \(\mathcal C\mathcal M=\mathcal M\) y \eqref{nuc03:eq:5}.
\end{proof}

La restricción al subespacio de registros compatibles es esencial: \(9\mathcal B-8I\) sobre todo \(\mathcal R\) no es unitario. El teorema caracteriza el transporte sobre registros que proceden efectivamente de una misma trayectoria lineal de la carta documentada.

\begin{corollary}
Para \(v,w\in\mathcal H\) y \(k\in\mathbb Z\),

\begin{equation}
\langle C^kv,w\rangle
=\langle\mathcal C^k\mathcal Vv,\mathcal Vw\rangle.
\label{nuc03:eq:6}
\end{equation}

Así se conservan todos los momentos cruzados de la dinámica realizada. Las potencias negativas pertenecen a la inversa de \(\mathcal C\) en \(\mathcal M\), no a una inversión del desplazamiento unilateral sobre el espacio ambiente.
\end{corollary}

\begin{corollary}
Para todo observable acotado \(A\), su realización en la imagen es \(\widehat A=\mathcal V A\mathcal V^*|_{\mathcal M}\). Se preservan productos, adjuntos y conmutadores. En particular,

\begin{equation}
[A,C]=0\quad\Longleftrightarrow\quad[\widehat A,\mathcal C]=0.
\label{nuc03:eq:7}
\end{equation}
\end{corollary}

La prueba consiste en usar \(\mathcal V^*\mathcal V=I\) y que \(\mathcal V\mathcal V^*\) es la identidad de \(\mathcal M\). La conservación de simetría queda expresada mediante el transporte íntegro de sus operadores, además de mediante una igualdad de normas.

\subsection{Reconstrucción explícita y profundidad del registro}

Sea \(m_j=\eta T^jv\). La identidad telescópica da para cada \(N\)

\begin{equation}
\boxed{v=T^{*N}T^Nv+\sum_{j=0}^{N-1}T^{*j}\eta^*m_j.}
\label{nuc03:eq:8}
\end{equation}

Al pasar al límite fuerte,

\begin{equation}
\boxed{v=Pv+\sum_{j=0}^{\infty}T^{*j}\eta^*m_j.}
\label{nuc03:eq:9}
\end{equation}

La reconstrucción parcial con terminal fijo,
\(v_N=Pv+\sum_{j<N}T^{*j}\eta^*m_j\), tiene error exacto

\begin{equation}
v-v_N=T^{*N}T^N(I-P)v.
\label{nuc03:eq:10}
\end{equation}

La serie converge para cada estado. Si el registro completo presenta un error \(\delta r\) en norma de \(\mathcal R\), su reconstrucción mediante \(\mathcal V^*\) cambia en una norma a lo sumo \(\|\delta r\|\), porque \(\|\mathcal V^*\|=1\). Ésta es una estabilidad uniforme de la lectura completa.

\subsubsection{Diferencia entre profundidad finita e ilimitada}

Para el desplazamiento bilateral \(Ce_k=e_{k+1}\) en \(\ell^2(\mathbb Z)\), se tiene \(P=0\). Considérese

\[
v_L=L^{-1/2}\sum_{k=1}^{L}e_k.
\]

Entonces \(\|(I-C)v_L\|^2=2/L\). La conmutación de \(T\) con \(I-C\), la contracción de \(T\) y \eqref{nuc03:eq:1} implican

\begin{equation}
\sum_{j=0}^{N-1}\|\eta T^jv_L\|^2
\leq\frac{16N}{81L}.
\label{nuc03:eq:11}
\end{equation}

Para cualquier profundidad fija \(N\), el lado derecho tiende a cero cuando \(L\to\infty\), aunque \(\|v_L\|=1\). El registro de memoria de profundidad fija, separado del terminal, carece de una cota inferior uniforme. En cambio, el registro ilimitado tiene norma exactamente uno para cada \(v_L\).

La diferencia corresponde a dos cuantificadores distintos: cada estado queda conservado al completar su registro; una profundidad prefijada no recupera uniformemente todos los estados sin su terminal. La conservación a profundidad arbitraria posee aquí una expresión funcional exacta.

\subsubsection{Información vectorial e intensidad}

El teorema conserva \(m_j\), no únicamente \(\|m_j\|^2\). Por ejemplo, en la realización bilateral, \(v=e_0\) y \(w=e_1\) tienen idénticas intensidades en cada profundidad porque el desplazamiento conmuta con \(T\) y transporta \(\eta\) unitariamente. Sus registros vectoriales son diferentes y \eqref{nuc03:eq:9} recupera los dos estados distintos. Las intensidades sirven para el balance; los registros con sus correlaciones sirven además para la reconstrucción.

\subsection{Naturalidad de la recuperación bajo refinamiento}

Sean \(\iota_n:\mathcal H_n\to\mathcal H_{n+1}\) las isometrías de refinamiento obtenidas de las emisiones comunes, y \(C_n\) realizaciones unitarias con

\begin{equation}
C_{n+1}\iota_n=\iota_nC_n.
\label{nuc03:eq:12}
\end{equation}

La relación vale también para los adjuntos: multiplicar por \(C_{n+1}^*\) y \(C_n^*\) da \(C_{n+1}^*\iota_n=\iota_nC_n^*\). Por tanto las imágenes son reductoras. Las fórmulas polinómicas de \(T_n\), \(\eta_n\) y el límite fuerte de \(T_n^N\) dan

\begin{equation}
T_{n+1}\iota_n=\iota_nT_n,\quad
\eta_{n+1}\iota_n=\iota_n^{\oplus9}\eta_n,\quad
P_{n+1}\iota_n=\iota_nP_n.
\label{nuc03:eq:13}
\end{equation}

Sea \(\mathfrak I_n\) el refinamiento que aplica \(\iota_n\) al terminal y \(\iota_n^{\oplus9}\) a cada registro. Entonces

\begin{equation}
\boxed{\mathcal V_{n+1}\iota_n=\mathfrak I_n\mathcal V_n,\qquad
\mathcal B_{n+1}\mathfrak I_n=\mathfrak I_n\mathcal B_n.}
\label{nuc03:eq:14}
\end{equation}

\begin{proof}
Cada componente de la primera igualdad es una identidad de \eqref{nuc03:eq:13}, iterada a la potencia correspondiente de \(T_n\). La segunda igualdad afirma que refinar cada bloque conmuta con desplazar el índice de bloques, lo cual se comprueba directamente. Ambas aplicaciones son acotadas, de modo que las igualdades pasan de los registros de soporte finito a su completación.
\end{proof}

Por \eqref{nuc03:eq:4} y \eqref{nuc03:eq:14}, la reconstrucción de la dinámica también conmuta con el refinamiento. Se mantienen distintos el índice de profundidad genealógica \(n\) y el de compresiones sucesivas \(j\). La igualdad expresa su compatibilidad cuando actúan mediante los cuadrados documentados, en lugar de identificarlos nominalmente.

\subsection{Dos regímenes exactos de distribución de memoria}

Para un estado unitario \(v\) con \(Pv=0\), definimos

\begin{equation}
a_N=\|T^Nv\|^2,\qquad b_N=\|\eta T^Nv\|^2=a_N-a_{N+1}.
\label{nuc03:eq:15}
\end{equation}

Se tiene \(a_0=1\), \(a_N\to0\), \(b_N\geq0\) y \(\sum b_N=1\). Los pesos \(b_N\) describen la distribución entre profundidades del registro de memoria. Este índice cuenta compresiones; su identificación con una duración física requeriría aplicar el lector temporal correspondiente.

\subsubsection{Transporte excepcional de orden tres}

El operador ortogonal fijo-libre \(g_c\) satisface \(g_c^2+g_c+I=0\). Para la composición explícita \(C=g_c\), el desarrollo precedente obtuvo

\begin{equation}
T^*T=\frac{19}{27}I,
\qquad a_N=\left(\frac{19}{27}\right)^N,
\qquad b_N=\frac8{27}\left(\frac{19}{27}\right)^N.
\label{nuc03:eq:16}
\end{equation}

La profundidad media, contando el primer registro como profundidad uno, es

\begin{equation}
\sum_{N\geq0}(N+1)b_N=\frac{27}{8}.
\label{nuc03:eq:17}
\end{equation}

Es un resultado de esa composición de la carta nonádica con el transporte excepcional ya construido; conserva el dominio y no identifica por su orden a \(g_c\) con la holonomía temporal completa.

\subsubsection{Registro bilateral de vueltas}

Para \(C e_k=e_{k+1}\) y \(v=e_0\), el binomio determina todos los términos:

\begin{equation}
T^Ne_0=9^{-N}\sum_{k=0}^{N}\binom{N}{k}8^{N-k}e_k,
\quad
a_N=81^{-N}\sum_{k=0}^{N}\binom{N}{k}^{\!2}64^{N-k}.
\label{nuc03:eq:18}
\end{equation}

Esta fórmula es exacta a cualquier profundidad. La representación circular posterior proporciona

\begin{equation}
a_N=\frac1{2\pi}\int_{-\pi}^{\pi}
\left(\frac{65+16\cos\theta}{81}\right)^N\,d\theta.
\label{nuc03:eq:19}
\end{equation}

\begin{theorem}
En esta realización bilateral,

\begin{equation}
\boxed{a_N\sim\frac9{4\sqrt{2\pi}}N^{-1/2},\qquad
b_N\sim\frac9{8\sqrt{2\pi}}N^{-3/2}.}
\label{nuc03:eq:20}
\end{equation}
\end{theorem}

\begin{proof}
Escribimos \(r(\theta)=(65+16\cos\theta)/81\) y \(q=8/81\). En torno a cero,

\[
r(\theta)=1-q\theta^2+O(\theta^4),\qquad
\log r(\theta)=-q\theta^2+O(\theta^4).
\]

Fijado un entorno suficientemente pequeño, existen constantes positivas \(c_1,c_2\) que acotan \(-\log r(\theta)\) entre \(c_1\theta^2\) y \(c_2\theta^2\). Fuera de ese entorno, \(r\leq\rho<1\). El cambio \(x=\sqrt N\theta\), la cota gaussiana y la convergencia dominada dan

\[
\lim_{N\to\infty}\sqrt N\,a_N
=\frac1{2\pi}\int_{\mathbb R}e^{-qx^2}\,dx
=\frac1{2\sqrt{\pi q}}=\frac9{4\sqrt{2\pi}}.
\]

Para el segundo límite se aplica el mismo procedimiento directamente a

\[
b_N=\frac1{2\pi}\int_{-\pi}^{\pi}r(\theta)^N(1-r(\theta))\,d\theta.
\]

En la variable \(x\), \(N(1-r(x/\sqrt N))\to qx^2\), dominado localmente por una constante por \(x^2\); fuera del entorno la contribución es exponencialmente pequeña incluso tras multiplicar por \(N^{3/2}\). Por tanto

\[
\lim_{N\to\infty}N^{3/2}b_N
=\frac q{2\pi}\int_{\mathbb R}x^2e^{-qx^2}\,dx
=\frac1{4\sqrt{\pi q}}=\frac9{8\sqrt{2\pi}}.
\]

Esta prueba del segundo límite usa su integral propia, sin diferenciar una equivalencia asintótica no controlada.
\end{proof}

Por Tonelli y la identidad de colas de una sucesión positiva,

\begin{equation}
\sum_{N\geq0}(N+1)b_N=\sum_{N\geq0}a_N=+\infty.
\label{nuc03:eq:21}
\end{equation}

Toda la norma se conserva en memoria, pero su profundidad media diverge en este estado bilateral. La ley de reparto depende de la representación efectiva del transporte: el sector excepcional de orden tres posee tasa geométrica, mientras el registro bilateral aquí considerado posee tasa algebraica.

\subsection{Composición con la identidad de Euler y la incidencia}

La identidad de Euler extendida de la sección~\ref{euler-trit-generadores} realiza los tres regímenes mediante el Coxeter de \(A_2\), un transporte nilpotente y \(F_{\rm av}^2\), con sus antecedentes declarados. Sus generadores normalizados satisfacen \(J_\tau^2=-\tau I\) y tienen traza cero. En sus cartas bidimensionales,

\[
C_\tau(t)=\exp(tJ_\tau),\qquad
\Omega=\begin{pmatrix}0&1\\-1&0\end{pmatrix},\qquad
C_\tau(t)^\mathsf T\Omega C_\tau(t)=\Omega.
\]

La traza cero da \(\det C_\tau(t)=1\), y la identidad \(M^\mathsf T\Omega M=(\det M)\Omega\) en dimensión dos prueba la última igualdad. Para

\[
T_\tau=\frac{8I+C_\tau}{9},\qquad
D_\tau=\frac{\sqrt8}{9}(C_\tau-I),
\]

se obtiene una ley de conservación orientada común:

\begin{equation}
\boxed{\Omega=T_\tau^\mathsf T\Omega T_\tau+D_\tau^\mathsf T\Omega D_\tau.}
\label{nuc03:eq:22}
\end{equation}

\begin{proof}
Al expandir los dos términos, las contribuciones cruzadas de \(C_\tau\) tienen coeficientes opuestos. Quedan \(72\Omega/81+9C_\tau^\mathsf T\Omega C_\tau/81=\Omega\).
\end{proof}

Esta forma alternante corresponde al área orientada del plano canónico. Su conservación permite comparar los tres regímenes respetando sus diferencias. Para los representantes concretos \(C\) de Coxeter, \(I+N\) y \(F_{\rm av}^2\), respectivamente,

\begin{equation}
(\det T,\det D)=
\left(\frac{19}{27},\frac8{27}\right),\qquad
(1,0),\qquad
\left(\frac{89}{81},-\frac8{81}\right).
\label{nuc03:eq:23}
\end{equation}

La suma es uno en los tres casos. En el parabólico, \(D\) es de rango uno y puede conservar información aunque su área sea cero. En el hiperbólico, el complemento tiene orientación opuesta; el balance es de área firmada. El teorema~\ref{nuc03:transporte} utiliza una realización unitaria y positiva: se conservan expresamente las hipótesis de cada resultado.

El desarrollo de acción y forma variable prueba la versión entre cartas diferentes, su integración sobre curvas cerradas y su compatibilidad con refinamiento. El desarrollo de incidencia compone la reconstrucción con \(F(k)=(\mu(k),\mathcal I P_0k/6)\), manteniendo las doce coordenadas de \(K\); transporta también los operadores y sus marcos locales. Así, una misma identidad de conservación admite realizaciones tipadas de norma, incidencia y acción orientada.

\section{Incidencia excepcional y transporte de formas}
\label{nuclear:4}
\subsection{Generación del registro y conservación de incidencia}

\subsubsection{Registro dodecafásico de la historia enriquecida}

La sección~\ref{subsec:k-registro-integral} define

\[
\mathcal R_{12}(x)=((E_m,\tau_m,\sigma_m,\kappa_m,\Lambda_m))_{m=1}^{12}.
\]

Conserva ventana, subruta, orientación, acarreo y libro de incidencias. La sección~\ref{subsec:k-terminal} fija

\[
x_{\rm term}=(\mathcal U_{108}\cdots\mathcal U_1)x_{\rm seed},
\quad \mathcal U_t=\mathsf{Upd}_t\mathsf{Tra}_t\mathsf{Sel}_t,
\]

\[
R_{\rm sgn}=\Pi_H\mathcal E_{108}^{90,120}\mathcal R_{12}:
\mathcal X_{\rm TPK}^{\rm term,mem}\to\mathbb Z^{12}.
\]

La reconstrucción algebraica siguiente comienza en esa salida tipada. No sustituye la evaluación de los eventos anteriores ni obtiene sus coordenadas imponiendo \(K\) o alfa como objetivo.

\subsubsection{Normalización integral y normalizaciones isométricas}

Para las órbitas de paso tres \((1,4,7,10)\), \((2,5,8,11)\), \((3,6,9,12)\), el lema~\ref{lem:k-hadamard} fija la matriz simétrica \(H=H_{12}\) y prueba \(H^2=4I\).

\[
\mathcal L_H=\{u\in\mathbb Z^{12}:Hu\equiv0\pmod4\},
\qquad \mathscr K(u)=\tfrac14Hu,
\qquad \mathscr K^{-1}(k)=Hk.
\]

Es un isomorfismo integral entre la subred de registros firmados y \(\mathbb Z^{12}\). El factor \(1/4\) realiza la inversión, no una isometría euclídea.

Como consecuencia inmediata de la misma identidad, \(J_H=H/2\) satisface \(J_H^*J_H=I\): éste es el transporte isométrico real. Si \(u=Hk\), se tiene \(u/2=J_Hk\) y \(\|u\|^2=4\|k\|^2\). Confundir \(H/4\) con \(H/2\) cambiaría el balance por un factor cuatro.

En \cref{exc:entrelazador,exc:isometria}, el diseño construido determina

\[
I:\mathbb R^{12}\to\mathbb R^{132},\qquad
(Iv)(H')=\sum_{p\in H'}v_p,\qquad I^*I=36I_{12}+30J_{12}.
\]

Los coeficientes proceden de las incidencias: cada punto pertenece a 66 hexadas y cada pareja a 30. Sobre \(V=\mathbf1^\perp\) desaparece \(J_{12}\), por lo que \(U_W=I/6:V\to E_{\rm inc}\) es isometría. El factor \(1/6\) es el normalizador de esta pantalla, no el inversor de \(K\).

\subsubsection{Refinamiento de prefijos y particiones de eventos}

La sección~\ref{exc:doble-lectura} construye marcos causales \(f_{j+1}=g_j(x)f_j\) y codificaciones \(Q_n\) del prefijo. Prueba

\[
r_{n,m}Q_n=Q_mp_{n,m},\qquad
\gamma_{A^f}(wf^{-1})=\gamma_A(w)(f^{-1}\oplus f^{-1}).
\]

La prueba usa que una prolongación conserva las palabras, transportes y marcos anteriores. El límite inverso produce \(Q_\infty\). La recuperación de la palabra visible se debe a la primera componente de \(\gamma\); el texto distingue expresamente ese alcance de la recuperación de toda memoria profunda y del paso a soportes.

La composición aditiva y afín del registro, descrita en \cref{subsec:k-registro-integral,iv:cont:concatenacion}, prueba una segunda conservación: con contribuciones prospectivas compatibles \((\delta w_t,\delta q_t)\), la inversión del registro es aditiva. Concatenar segmentos o subdividir la misma lista de eventos da el mismo registro final. La versión con transporte no trivial usa el cociclo afín. Conservar el total al cambiar la partición no equivale a afirmar que el registro deja de crecer cuando se añaden nuevos eventos.

\subsection{Transporte K–Witt de normas, operadores y balances}\label{nuc04:kwitt}

Se reúne aquí una formalización de los resultados recuperados, con la prueba algebraica que sigue. Este ensamblaje no atribuye prioridad nueva a sus identidades de partida.

Sean \(P_0=I-J_{12}/12\), \(V=P_0\mathbb R^{12}\), \(W_H=J_HV\) y

\[
B:=U_WJ_H|_{W_H}:W_H\longrightarrow E_{\rm inc}.
\]

Como \(J_H^2=I\), éste es un isomorfismo isométrico entre espacios de dimensión once. Para un registro \(k\), su coordenada firmada normalizada es \(J_Hk\). Si \(\widehat P_0=J_HP_0J_H\),

\[
B\widehat P_0(J_Hk)=U_WP_0k,
\]

\[
\|\widehat P_0(J_Hk)\|^2
=\|U_WP_0k\|^2
=\sum_{p=1}^{12}k_p^2-\frac{(\sum_pk_p)^2}{12}.
\]

No se ha introducido ningún valor de constante para obtener esta igualdad: procede de \(H^2\) y del Gramiano de incidencia. La media eliminada debe conservarse aparte si se quiere recuperar todo \(k\).

Para \(g\in\operatorname{Aut}(\mathcal H)\), sea \(\rho_{12}(g)\) su permutación de puntos y \(\rho_{\rm hex}(g)\) la de hexadas. La acción sobre las coordenadas firmadas es

\[
\widehat\rho_H(g)=J_H\rho_{12}(g)J_H
=\tfrac14H\rho_{12}(g)H.
\]

Es ortogonal, conserva \(W_H\) y también conserva la subred \(\mathcal L_H\): si \(u=Hk\), su imagen es \(H\rho_{12}(g)k\). El entrelazamiento material de \cref{exc:entrelazador} da

\[
B\widehat\rho_H(g)=\rho_{\rm hex}(g)B.
\]

\begin{proof}
sustituir las definiciones, cancelar \(J_H^2\) y aplicar \(U_W\rho_{12}(g)=\rho_{\rm hex}(g)U_W\). Se preservan así el registro integral, el producto interno centrado y la acción incidencial, cada uno en su dominio.
\end{proof}

Para un operador \(A\) ya construido sobre \(W_H\) se obtiene \(A_{\rm inc}=BAB^*\). Por sustitución,

\[
BA=A_{\rm inc}B,\quad
\langle z,Az\rangle=\langle Bz,A_{\rm inc}Bz\rangle.
\]

Por tanto, si un transporte \(C\) sobre \(W_H\) conserva \(A\), \(C^*AC=A\), su transporte \(C_{\rm inc}=BCB^*\) conserva \(A_{\rm inc}\). La condición de que un \(C\) particular sea una publicación de la dinámica TPK permanece en la composición causal que lo produce. La isometría no transforma por sí sola una permutación de incidencia en el ciclo nonádico.

No hace falta imponer conmutación para comparar cartas transportadas. Si \(R\) es unitario, \(C'=RCR^*\) y \(T(C)=(8I+C)/9\), entonces \(T(C')R=RT(C)\) por sustitución. Junto con \(A'=RAR^*\), el balance en una carta se conjuga al de la otra. La conmutación \([R,C]=0\) sólo corresponde al problema adicional de mantener fija esa misma costura \(C\). Esta distinción coincide con la naturalidad reticular de la sección~\ref{nuc04:c3}.

Una acción finita concreta disponible es la estrella seleccionada por la bandera:

\[
\sigma_{B_K}=(1\ 3)(2\ 11)(5\ 7)(8\ 12).
\]

La sección~\ref{exc:estrella} prueba su construcción y la distingue del grupo generado por las 495 estrellas. La estrella individual da \(C_2\); la familia da \(M_{12}\). Transportar la marca y el origen da covariancia entre marcos. Si se pretende una simetría de un marco fijo, el dominio adecuado es su estabilizador, no automáticamente \(M_{12}\) entero. Si además debe conmutar con \(C\), se utiliza el subgrupo que satisfaga la ecuación concreta \([\widehat\rho_H(g),C]=0\).

\subsection{Acción de orden tres sobre las fibras reticulares}\label{nuc04:c3}

La construcción de \cref{nuclear:fibras-app-witt,nuclear:alineamiento-witt} produce doce fibras \(A_2\) etiquetadas por \((i,\mu,\varepsilon)\in\mathbb Z/3\mathbb Z\times\{0,1\}\times\{0,1\}\). El alineamiento \(\lambda(i,\mu,\varepsilon)=4i+2\mu+\varepsilon\pmod{12}\) es biyectivo. Si \(C_{A_2}\) es el Coxeter y \(R\) la reflexión con \(RC_{A_2}R=C_{A_2}^{-1}\), las trivializaciones y marcos transportados fijan

\[
g_b=\iota_b^{-1}P_bC_{A_2}P_b^{-1}\iota_b,
\quad T_{i\mu\varepsilon}=\iota_{\lambda(i,\mu,\varepsilon)}^{-1}P_{\lambda(i,\mu,\varepsilon)}R^\mu,
\]

\[
g_\lambda T_{i\mu\varepsilon}
=T_{i\mu\varepsilon}C_{A_2}^{(-1)^\mu}.
\]

La suma directa \(T_\lambda:W_{\rm APP}\otimes\mathbb C\to\bigoplus_{b=1}^{12}(A_{2,b}\otimes\mathbb C)\) es un isomorfismo \(C_3\)-equivariante. La prueba consiste en la conjugación y en la biyección de las doce etiquetas: conserva par, hoja y órbita. No es una coincidencia de dimensiones.

Las construcciones de \cref{km:palabra-total,km:estabilidad-vecino} utilizan \(q^\star=(1,1,1,1,1,1)\) y su palabra \(c^\star=(q^\star,q^\star A_W)\) de soporte total para orientar los doce bloques. El \(g_c\) resultante tiene orden tres, no tiene vectores fijos no nulos y preserva el pegado de Niemeier. Para el vector incidencial \(v\) de norma cuadrática \(54\),

\[
N_0=\ker(\langle\cdot,v\rangle\bmod3),\qquad
N_v=N_0+\mathbb Z\frac v3,
\]

\[
y_*=(g_cv-v)/3\in N_0,\qquad
z_*=(g_c^{-1}v-v)/3\in N_0.
\]

La prueba identifica sus clases discriminantes como \(c^\star\) y \(-c^\star\), y sus productos con \(v\) como \(-27\). De ahí \(g_cN_0=N_0\) y \(g_c(v/3)=v/3+y_*\), luego \(g_cN_v=N_v\). La estructura conservada incluye el pegado, el núcleo de congruencia y la clase que genera el vecino de Leech.

La prolongación a las 24 orientaciones de soporte total y la elevación reticular de un automorfismo monomial \(h\) del código se desarrollan en la sección~\ref{paper:24-orientaciones}. Con todas las marcas transportadas,

\[
g_{hu}=\widetilde h\,g_u\,\widetilde h^{-1}.
\]

Este \(C_3\) es un grupo permisible realmente entrelazado con las fibras APP y la estructura excepcional. No se lo identifica aquí con \(C_9\) ni con su cociente por una mera relación de órdenes. La elevación al retículo tampoco se reemplaza por la matriz ternaria sin su acción sobre las fibras \(A_2\).

\subsection{Refinamiento de la energía y de su diferencial}
\label{paper:refinamiento-energia}

El operador de energía de \eqref{vii:eq:energia-memoria-discreta} tiene la forma \(D_{9,m}^*W_mD_{9,m}\), con \(U_a\) transportando orientación y acarreo. Los dominios de \cref{vii:subsec:funcional-memoria} son:

\[
\mathcal H_m=\bigoplus_vH_v,\quad
\mathcal K_m=\bigoplus_{a\in\mathcal E_m^+}H_{t(a)},\quad
(D_mf)_a=f_{t(a)}-U_af_{s(a)},\quad
E_m=\langle D_mf,\mathsf W_mD_mf\rangle.
\]

Las aplicaciones fijas \(J_m\) y \(K_m\) satisfacen, para la familia \(C^1\) considerada,

\[
D_{m+1}(t)J_m=K_mD_m(t),\qquad
K_m^*\mathsf W_{m+1}(t)K_m=\mathsf W_m(t).
\]

El \cref{vii:thm:refinamiento-respuesta} prueba \(E_{m+1}(J_mf)=E_m(f)\) por sustitución, y la igualdad de diferenciales se obtiene diferenciando ambas identidades. Para un refinamiento dependiente del parámetro se conserva también el término \(D_{m+1}\dot J_mf\).

Con \(v_a=U_af_s\) y \(q=\mathsf W_mD_mf\), la respuesta a \(\dot U_a=A_aU_a\) es

\[
\mathfrak j_m(A)=-2\operatorname{Re}\sum_a\langle q_a,A_av_a\rangle,
\quad \mathcal J_a=v_aq_a^*-q_av_a^*.
\]

El producto de transportes \(U_\gamma=U_N\cdots U_1\) cumple

\[
\dot U_\gamma U_\gamma^{-1}
=\sum_{k=1}^N B_kA_kB_k^{-1},\qquad B_k=U_N\cdots U_{k+1}.
\]

Esta fórmula retiene el orden de cada incidencia y transporta contragredientemente su covector. El corolario~\ref{vii:cor:refinamiento-corriente-conexion} compone ese diferencial con las variaciones de conexión y prueba, para los refinamientos compatibles de la acción,

\[
s_m=P_m^{\mathsf T}s_{m+1},\qquad
M_m\sigma_m=P_m^{\mathsf T}M_{m+1}\sigma_{m+1}.
\]

Aquí \(P_m\) lleva variaciones de conexión y \(M_m\) es el emparejamiento celular; no son el proyector de media nula ni la matriz de Hadamard.

\subsubsection{Ensamblaje con la incidencia}

En una realización de fibras \(W_H\) cuyas acciones de marco sean \(\widehat\rho_H(g_a)\), la conjugación por \(B\) lleva cada transporte, campo y peso a \(E_{\rm inc}\). Por el entrelazamiento demostrado en la sección~\ref{nuc04:kwitt}, los operadores de diferencias satisfacen \(D_{\rm inc}B_{\rm vertices}=B_{\rm edges}D_H\). Transportando también \(J_m\), \(K_m\) y \(\mathsf W_m\), los dos cuadrados de refinamiento anteriores siguen siendo los mismos cuadrados conjugados. Por sustitución conservan la energía; para \(B\) fijo, conservan además el diferencial y los covectores.

La ley de balance y sus covectores se transportan a la representación excepcional conservando su compatibilidad con el refinamiento. Cada aplicación utiliza los transportes de la dinámica especificada y los subespacios que preservan. El dominio de las variaciones está definido en \cref{vii:subsec:funcional-memoria}.

\section{Reconstrucción incidencial y covariancia}
\label{nuclear:5}
\subsection{La isometría completa y su dominio}\label{nuc05:isometria}

Sea \(H=\mathbb R^{12}\), con producto euclídeo; \(P_0=I-\mathbf1\mathbf1^*/12\), \(\mu(k)=\mathbf1^*k/12\), y \(\mathcal I:H\to\mathbb R^{132}\) la matriz de incidencia de \cref{exc:entrelazador,exc:isometria}. Escribimos \(E_{\rm inc}=\mathcal I(\mathbf1^\perp)\) y

\[
Y=\mathbb R\oplus E_{\rm inc},\qquad
\langle(\mu,y),(\nu,z)\rangle_Y=12\mu\nu+\langle y,z\rangle.
\]

La isometría completa construida en la sección~\ref{nuclear:2} es

\[
F:H\longrightarrow Y,\qquad Fk=(\mu(k),\mathcal I P_0k/6),
\]

\[
F^{-1}(\mu,y)=\mu\mathbf1+P_0\mathcal I^*y/6.
\]

El Gram \(\mathcal I^*\mathcal I=36I+30\mathbf1\mathbf1^*\) prueba ambas composiciones inversas y \(F^*F=I_H\), \(FF^*=I_Y\), con el producto ponderado indicado. Por ello \(F\) es ortogonal entre \(H\) e \(Y\). Puede complejificarse sin cambiar las pruebas.

El codominio es \(Y\), de dimensión doce, no todo \(\mathbb R\oplus\mathbb R^{132}\). Extender un operador desde \(Y\) a su complemento en ese ambiente requiere una elección adicional; no forma parte de la dinámica transportada de manera única. En coordenadas firmadas normalizadas se usa \(FJ_H\), con \(J_H=H_{12}/2\). La inversión integral sigue siendo \(H_{12}/4\) en su subred, no un normalizador métrico alternativo.

\subsection{Teorema de levantamiento covariante en la imagen}\label{nuc05:levantamiento}

\begin{theorem}
Para un operador acotado \(D\) sobre \(H\), existe un único operador sobre \(Y\) que satisface \(D_YF=FD\). Es

\[
D_Y=FDF^{-1}.
\]

La correspondencia conserva sumas, productos, adjuntos, norma, positividad e identidad. En particular, lleva unitarios a unitarios y observables autoadjuntos a observables autoadjuntos. Para un unitario \(C\) y un observable \(A\),

\[
C^*AC=A\quad\Longleftrightarrow\quad C_Y^*A_YC_Y=A_Y.
\]

Si \(R\) es un cambio unitario de carta, \(R_Y=FRF^{-1}\), entonces

\[
F(RDR^{-1})F^{-1}=R_YD_YR_Y^{-1}.
\]
\end{theorem}

\begin{proof}
La primera fórmula se obtiene multiplicando el entrelazamiento por \(F^{-1}\). La unicidad se sigue de la sobreyectividad hacia \(Y\). Las identidades de productos y adjuntos resultan de \(F^{-1}=F^*\), y la igualdad de normas de la isometría sobreyectiva. La última igualdad es una composición asociativa de esas aplicaciones. No interviene ninguna hipótesis de conmutación entre \(R\) y \(D\).
\end{proof}

Para \(g\in\operatorname{Aut}(\mathcal H)\), la proposición~\ref{exc:isometria} proporciona

\[
F\rho_{12}(g)=\bigl(1\oplus\rho_{\rm hex}(g)\bigr)F.
\]

Así, el cambio de carta posee una acción combinatoria concreta sobre las hexadas. En cambio, para un unitario general \(D\), \(D_Y\) es un operador de la representación incidencial, no necesariamente una permutación del diseño. Preservar el producto interno de esa representación no convierte a todos sus unitarios en automorfismos combinatorios ni en elementos del Monstruo.

\subsubsection{Covariancia bajo un cambio de carta no conmutante}

Sea \(C\) el avance \(Ce_p=e_{p+1}\) cíclico en las doce posiciones, es decir, el sentido inverso de la convención \(S\) de \cref{sec:k-registro}. Es una operación de la carta de ventanas, no la holonomía temporal completa. Sea \(R\) la estrella ya construida

\[
R=(1\ 3)(2\ 11)(5\ 7)(8\ 12).
\]

Entonces \(RCe_1=e_{11}\), mientras \(CRe_1=e_4\). No conmutan. Sin embargo, para \(C'=RCR^{-1}\), se cumplen exactamente \(T(C')R=RT(C)\) y \(\eta(C')R=R^{\oplus9}\eta(C)\), y lo mismo en \(Y\) mediante \(F\). La covariancia de carta sí funciona en un ejemplo donde la conmutación en carta fija falla.

\subsection{Recuperación explícita desde el registro incidencial de complementos}\label{nuc05:recuperacion}

Sean \(C\) unitario, \(T=(8I+C)/9\), \(\eta=(I-JJ^*)S_CJ:H\to H^9\), y \(P\) la proyección sobre \(\operatorname{Fix}(C)\), como en la sección~\ref{nuclear:2}. Escribimos \(F_9=F^{\oplus9}\). La versión incidencial del registro finito es

\[
\mathcal W_Nk=\left(FT^Nk,F_9\eta k,F_9\eta Tk,\ldots,F_9\eta T^{N-1}k\right).
\]

\begin{theorem}
Esta aplicación es isométrica y posee la inversa izquierda explícita

\[
\mathcal W_N^*(a,b_0,\ldots,b_{N-1})
=T^{*N}F^{-1}a+
\sum_{j=0}^{N-1}T^{*j}\eta^*F_9^{-1}b_j.
\]

En particular, aplicar esta fórmula a \(\mathcal W_Nk\) recupera exactamente \(k\). El registro infinito

\[
\mathcal W_\infty k=\left(FPk,(F_9\eta T^jk)_{j\ge0}\right)
\]

también es isométrico, y su inversa izquierda es

\[
k=PF^{-1}a+\sum_{j\ge0}T^{*j}\eta^*F_9^{-1}b_j
\quad\text{para }(a,(b_j))=\mathcal W_\infty k.
\]

La serie converge en norma. Para datos arbitrarios cuadrado-sumables la misma fórmula define el adjunto acotado, aunque no todo dato del codominio pertenece a la imagen de \(\mathcal W_\infty\).
\end{theorem}

\begin{proof}
Aplicar \(F\) y \(F_9\) a cada componente de \(V_N\) o \(V_\infty\) conserva el producto interno. Tomar el adjunto de esa composición da las fórmulas. La identidad telescópica de la sección~\ref{nuclear:2} prueba \(\mathcal W_N^*\mathcal W_N=I\), y su límite fuerte prueba \(\mathcal W_\infty^*\mathcal W_\infty=I\). La convergencia de la serie procede de aplicar un operador acotado a las truncaciones de una sucesión en suma directa hilbertiana.
\end{proof}

Esta reconstrucción no conserva \(k\) como una primera coordenada intacta: el primer componente es el terminal comprimido o su sector fijo; el resto procede de los complementos sucesivos. Cuando \(P=0\), los complementos incidenciales por sí solos recuperan \(k\). La contribución recuperada a partir de los primeros \(N\) bloques es

\[
k_N=\sum_{j<N}T^{*j}\eta^*F_9^{-1}b_j
=\bigl(I-T^{*N}T^N\bigr)k.
\]

El residuo es \(T^{*N}T^Nk\), que tiende a cero si \(P=0\). Con un sector fijo no nulo debe conservarse \(FPk\); la memoria de complementos sola no recupera ese sector. Las tasas de reconstrucción y la dinámica sobre el índice de bloques están desarrolladas en la sección~\ref{nuclear:3}.

Si \(A\) es autoadjunto y \([A,C]=0\), la forma del observable se conserva también:

\[
\langle k,Ak\rangle
=\langle FPk,A_YFPk\rangle_Y
+\sum_{j\ge0}\langle F_9\eta T^jk,(I_9\otimes A_Y)F_9\eta T^jk\rangle.
\]

La serie escalar converge absolutamente porque \(A\) es acotado y la suma de las normas cuadradas de los componentes es finita. La hipótesis expresa conservación de \(A\) por \(C\); no exige conmutación de \(C\) con los cambios de carta excepcionales.

\subsection{Levantamiento natural del refinamiento con marcos locales}\label{nuc05:refinamiento}

La construcción de \cref{iv:cont:historias,iv:cont:medida} aporta conjuntos de estados cilíndricos enriquecidos \(Z_n\), truncamientos y medidas proyectivas de soporte completo. Para un descendiente \(z'\) de \(z\),

\[
w_n(z'|z)=\sqrt{\mu_{n+1}(z')/\mu_n(z)},\qquad
\sum_{\tau z'=z}w_n(z'|z)^2=1.
\]

La etiqueta \(z'\) conserva registro ordenado, orientación, acarreo, frontera, memoria, ruta y terminal. El refinamiento de un estado no identifica descendientes distintos.

Sean \(q(z)\) el registro entero de memoria, \(C\) un unitario sobre \(H\), y \(R_z\) cambios ortogonales de marco de la fibra. Para marcos incidenciales concretos puede usarse \(R_z=\rho_{12}(g_z)\), con \(g_z\in\operatorname{Aut}(\mathcal H)\), transportando sus marcas. Definimos

\[
G_z=R_zC^{q(z)},\qquad
U_{z',z}=G_{z'}G_z^*
=R_{z'}C^{q(z')-q(z)}R_z^*.
\]

Esta construcción prolonga la equivalencia de calibre explícita de \cref{exc:doble-lectura}. No requiere \([R_z,C]=0\), y deja visible el orden de las operaciones.

Sobre \(\mathscr H_n=\ell^2(Z_n)\otimes H\),

\[
\mathscr U_n(e_z\otimes v)
=\sum_{\tau z'=z}w_n(z'|z)e_{z'}\otimes U_{z',z}v.
\]

Definimos \(\mathscr F_n=I_{\ell^2(Z_n)}\otimes F\), con codominio \(\mathscr Y_n=\ell^2(Z_n)\otimes Y\).

\begin{theorem}
El refinamiento es isométrico y satisface

\[
\mathscr U^Y_n\mathscr F_n=\mathscr F_{n+1}\mathscr U_n,
\quad
\mathscr U^Y_n=\mathscr F_{n+1}\mathscr U_n\mathscr F_n^{-1}.
\]

La fórmula de cada término es la anterior con \(U^Y_{z',z}=FU_{z',z}F^{-1}\). Los sistemas de varias profundidades conservan composición y admiten una isometría sobreyectiva entre sus límites inductivos. Se preservan exactamente todos los coeficientes de fibra con sus etiquetas \(z\).
\end{theorem}

\begin{proof}
Los descendientes de padres distintos son ortogonales; cada \(U_{z',z}\) es unitario; y los pesos tienen suma cuadrática uno. Esto prueba la isometría. En una trayectoria, los factores intermedios \(G_z^*G_z\) se cancelan y los cocientes de medidas telescopan. La composición no depende de la partición de la trayectoria. El cuadrado con \(F\) se verifica sobre \(e_z\otimes v\), término a término. Las aplicaciones y sus inversas son isometrías compatibles, por lo que se extienden inversamente a las completaciones inductivas.
\end{proof}

\subsubsection{Observables y dinámica compatibles}

Si \(A_0\) es un observable acotado de la fibra, sea

\[
(\mathscr A_n\psi)_z=G_zA_0G_z^*\psi_z.
\]

Entonces \(\mathscr A_{n+1}\mathscr U_n=\mathscr U_n\mathscr A_n\), de donde \(\mathscr U_n^*\mathscr A_{n+1}\mathscr U_n=\mathscr A_n\). Se obtiene sustituyendo \(U_{z',z}=G_{z'}G_z^*\). Esta es conservación de un campo de observables transportado por refinamiento; no afirma que \(A_0\) conmute con \(C\) en una carta fija.

La dinámica unitaria por nivel

\[
(\mathscr C_n\psi)_z=G_zCG_z^*\psi_z
\]

satisface \(\mathscr C_{n+1}\mathscr U_n=\mathscr U_n\mathscr C_n\). También entrelazan sus adjuntos: como ambas \(\mathscr C\) son unitarias, la identidad se multiplica por sus inversas. Por tanto entrelazan \(T\), sus complementos nonádicos y sus proyecciones fijas. El registro finito y el infinito de la sección~\ref{nuc05:recuperacion} son naturales con respecto a este refinamiento. Al aplicar \(\mathscr F_n\), esos mismos cuadrados valen en las fibras incidenciales. Para que \(\mathscr A_n\) sea conservado además por esa dinámica en un nivel fijo, la condición precisa es \([A_0,C]=0\); la covariancia entre marcos sigue sin requerir \([R_z,C]=0\).

\subsubsection{Especialización de memoria bilateral con marcos variables}

En la órbita reversible de \cref{nuclear:3} puede representarse una publicación de coordenadas por \(\psi=(\psi_m)_{m\in\mathbb Z}\in\ell^2(\mathbb Z;H)\). Su avance es \((C_0\psi)_m=\psi_{m-1}\). \(F\) actúa punto a punto. Para marcos \(R_m=\rho_{12}(g_m)\), el cambio unitario \((\mathscr R\psi)_m=R_m\psi_m\) produce

\[
(C'\psi)_m=R_mR_{m-1}^*\psi_{m-1}.
\]

El operador incidencial tiene la misma fórmula con \(1\oplus\rho_{\rm hex}(g_mg_{m-1}^{-1})\). No se imponen marcos constantes ni conmutación con el desplazamiento. Como \(C'\) es conjugado a \(C_0\), no tiene vectores fijos no nulos en este espacio cuadrado-sumable. La sección~\ref{nuc05:recuperacion} recupera la publicación lineal completa desde sus complementos incidenciales. Se trata de la órbita y de las coordenadas representadas: no de una identificación de todas las historias TPK con \(\mathbb Z\).

\subsection{Extensión directa a las doce fibras \texorpdfstring{\(A_2\)}{A2}}\label{nuc05:a2}

La construcción excepcional conserva doce fibras \(A_2\) y el operador \(g_c\) de orden tres desarrollado en la sección~\ref{nuc04:c3}. Sea \(V_{A_2}=A_2\otimes_{\mathbb Z}\mathbb R\) su realización euclídea bidimensional. La misma isometría posicional puede tensorizarse sobre \(\mathbb R\) con ese espacio:

\[
F_{A_2}=F\otimes I_{V_{A_2}}:
\mathbb R^{12}\otimes_{\mathbb R}V_{A_2}
\longrightarrow Y\otimes_{\mathbb R}V_{A_2}.
\]

Su primera componente es la media vectorial y su segunda son las sumas incidenciales centradas, ahora con valores en \(V_{A_2}\). El Gram de incidencia, tensorizado con la identidad, prueba de nuevo la isometría completa y la inversa. Por ello \(g_c^Y=F_{A_2}g_cF_{A_2}^{-1}\) es ortogonal, de orden tres y fijo-libre. La composición nonádica de la sección~\ref{nuclear:2} puede publicarse y recuperarse en esta pantalla con valores en \(V_{A_2}\), mediante las secciones~\ref{nuc05:levantamiento} y~\ref{nuc05:recuperacion}. No se confunden los doce escalares de \(K\) con las veinticuatro coordenadas de las fibras \(A_2\), ni se convierte la acción en una permutación de hexadas cuando la elevación también actúa dentro de las fibras.

\subsection{Alcance exacto respecto de las historias}

Sea \(k:X\to H\) el lector del registro sobre un dominio de historias, y sea \(\ell=F\circ k\). Entonces

\[
\ell(x)=\ell(y)\quad\Longleftrightarrow\quad k(x)=k(y).
\]

La pantalla completa añade cero pérdida respecto de \(K\), pero tampoco cambia las fibras de ese lector. Un mapa \(d:X\to X\) desciende a una dinámica sobre los valores de \(K\) si y sólo si

\[
k(x)=k(y)\Longrightarrow k(d x)=k(d y).
\]

La misma condición es necesaria y suficiente para descender sobre la imagen incidencial. La prueba consiste en definir \(D(k(x))=k(dx)\); la condición equivale a que la definición sea independiente del representante. Si ese \(D\) es lineal y unitario en la realización declarada, se aplica la sección~\ref{nuc05:levantamiento}. Si es no lineal, \(F\) sigue conjugándolo como aplicación, pero no lo convierte en un unitario.

El criterio distingue el descenso a \(K\) de la dinámica enriquecida. La sección~\ref{subsec:k-terminal} establece que el lector utiliza subruta, orientación, acarreo e incidencias; la sección~\ref{exc:limite-inverso} distingue recuperación de palabra visible, soporte y memoria profunda. En \(\ell^2(Z_n)\otimes H\), la sección~\ref{nuc05:refinamiento} conserva las etiquetas enriquecidas y transporta la representación de la fibra. La equivalencia actúa sobre el dominio de historias previamente generado.

Para un corte ya construido, la naturalidad del registro en \cref{prop:k-extension} fija \(K_N=K_{108}\circ\tau_{N,108}\): \(F\) conserva exactamente esa naturalidad. Una sucesión de nuevos eventos puede actualizar \(K\); la composición aditiva o afín de \cref{sec:k-registro} conserva la independencia de la partición, no la constancia del valor al añadir eventos.

\section{Holonomía afín y simetrías de la memoria}
\label{nuclear:6}
\subsection{Composición efectiva de la memoria}

Para una emisión admisible \(\epsilon\), la construcción de \cref{iv:cont:concatenacion} define

\begin{equation}
\mathsf U_\epsilon(m)=C_\epsilon m+\kappa_\epsilon,
\qquad
\kappa_{\epsilon_2\epsilon_1}
=\kappa_{\epsilon_2}+C_{\epsilon_2}\kappa_{\epsilon_1}.
\label{nuc06:eq:1}
\end{equation}

El dominio es el módulo de memoria de la fuente; el codominio, el de la emisión siguiente. El producto conserva el orden. En una ruta \(\gamma\) de longitud \(r\) se obtiene

\begin{equation}
C_\gamma=C_r\cdots C_1,\qquad
\kappa_\gamma=\sum_{j=1}^r C_r\cdots C_{j+1}\kappa_j,
\qquad
m_\gamma=C_\gamma m_0+\kappa_\gamma.
\label{nuc06:eq:2}
\end{equation}

El producto vacío de la suma es la identidad. La demostración es la inducción de \eqref{nuc06:eq:1}: añadir la emisión siguiente multiplica el registro anterior por su transporte lineal y suma su incremento. La frontera posicional puede cancelarse mientras permanece la suma transportada de \eqref{nuc06:eq:2}.

\begin{proposition}[Clase de memoria de un retorno]

Sea \(M\) un módulo abeliano y sea una ruta cerrada en la base cuyo transporte de memoria sea \(H(m)=Cm+\kappa\), con \(C\) un automorfismo de \(M\). Su clase de retorno es

\begin{equation}
\mathfrak m(H)=[\kappa]\in M/(I-C)M.
\label{nuc06:eq:3}
\end{equation}

Un cambio de coordenadas de memoria \(m'=Rm+b\), con \(R\) invertible, transforma el retorno en

\begin{equation}
C'=RCR^{-1},\qquad
\kappa'=R\kappa+(I-C')b.
\label{nuc06:eq:4}
\end{equation}

La aplicación inducida por \(R\) lleva \([\kappa]\) a \([\kappa']\). El transporte \(H\) tiene un punto fijo exactamente cuando su clase \eqref{nuc06:eq:3} es cero.
\end{proposition}

\begin{proof}
Sustituir \(m=R^{-1}(m'-b)\) en \(H\) da \eqref{nuc06:eq:4}. Como \(R(I-C)=(I-C')R\), \(R\) induce un isomorfismo entre los cocientes. El término \((I-C')b\) representa cero en el cociente destino. Finalmente, la ecuación de punto fijo \(Cm+\kappa=m\) equivale a \((I-C)m=\kappa\). Esto prueba la última afirmación.
\end{proof}

La fórmula \eqref{nuc06:eq:3} permite decidir si un cambio de origen elimina el incremento o sólo lo redistribuye entre coordenadas. En particular, cuando \(C=I\), la clase es el propio vector \(\kappa\). Una traslación no nula permanece no nula bajo cualquier cambio afín invertible.

\begin{proposition}[Simetría del retorno frente a cambio de representación]

Una transformación afín \(A(m)=Rm+b\) conmuta con \(H\) exactamente cuando

\begin{equation}
RC=CR,\qquad (I-C)b=(I-R)\kappa.
\label{nuc06:eq:5}
\end{equation}
\end{proposition}

\begin{proof}
Las partes lineales de \(AH\) y \(HA\) son \(RC\) y \(CR\); sus partes constantes son \(R\kappa+b\) y \(Cb+\kappa\). Igualarlas da \eqref{nuc06:eq:5}. Si sólo se transporta \(H\) mediante \eqref{nuc06:eq:4}, no se requiere conmutación: se obtiene la misma ley en otra representación. Para una simetría del retorno fijado, en cambio, se requieren ambas igualdades de \eqref{nuc06:eq:5}.
\end{proof}

Así se distinguen con precisión la covarianza de la ley y el estabilizador de una memoria concreta. Esta diferencia es pertinente para estudiar variaciones de simetría del estado sin modificar las leyes de composición.

\subsection{Aplicación al registro efectivo de la vuelta observable}

La sección~\ref{vii:sec:retorno-memoria-gravedad} define dos contadores ordenados: \(g\) cuenta inversiones de orientación y \(s\) cambios efectivos de hoja. Su cómputo canónico por bloque de 27 pasos es

\[
c_{27}=(1,18).
\]

La fase observable completa retorna después de cuatro de esos bloques. En el grupoide de rutas, el módulo de registro es \(M=\mathbb Z^2\), y

\begin{equation}
c_{108}=(4,72)=4v,\qquad v=(1,18),\qquad
H(g,s)=(g+4,s+72).
\label{nuc06:eq:6}
\end{equation}

La identidad \eqref{nuc06:eq:6} es un resultado del corpus que se utiliza aquí; el comprobador verifica sus consecuencias, no vuelve a contar las emisiones originales. El símbolo \(c_{108}\) designa este incremento de dos contadores. Es distinto del sello dodecafásico \(K\) y del cociclo central de la representación de Heisenberg.

\begin{theorem}[Invariantes completos del registro bajo vueltas completas]

Las órbitas de \(H\) y su inversa sobre \(\mathbb Z^2\) quedan clasificadas exactamente por

\begin{equation}
\boxed{\quad \mathcal I(g,s)=s-18g,\qquad \nu(g,s)=[g]_4.\quad}
\label{nuc06:eq:7}
\end{equation}

En consecuencia,

\begin{equation}
\mathbb Z^2/\mathbb Z(4,72)\cong\mathbb Z\oplus\mathbb Z/4\mathbb Z.
\label{nuc06:eq:8}
\end{equation}
\end{theorem}

\begin{proof}
Añadir \((4,72)\) conserva ambas expresiones de \eqref{nuc06:eq:7}. Si dos registros tienen iguales invariantes, \(g'-g=4N\) para un entero \(N\) y \(s'-s=18(g'-g)=72N\); luego difieren exactamente en \(N(4,72)\). La aplicación de \eqref{nuc06:eq:7} es sobreyectiva: dados \(j\) entero y \(r\) módulo cuatro, elegir un representante \(g\) de \(r\) y poner \(s=j+18g\). Esto demuestra \eqref{nuc06:eq:8}.
\end{proof}

La descripción cubre todas las órbitas de este lector de memoria. Las órbitas del estado enriquecido completo conservan, además, sus restantes coordenadas. Para una prolongación dirigida hacia adelante, un registro posterior se obtiene cuando el entero N anterior es no negativo; la equivalencia bilateral de \eqref{nuc06:eq:8} pertenece al grupoide.

La magnitud \(\mathcal I\) es la única carga lineal entera primitiva, salvo signo, invariante bajo esas traslaciones. En efecto, para \(L(g,s)=ag+bs\), la invariancia exige \(4a+72b=0\), por lo que \((a,b)=b(-18,1)\). Una vuelta acumula memoria longitudinal y conserva una relación transversal determinada por el incremento.

El factor cuatro tiene aquí un origen preciso: es \(\gcd(4,72)\) y, en el cómputo anterior, el número de bloques de 27 que completa el retorno observable. La clase \(\nu\) conserva la información que se pierde al dividir el incremento por cuatro y guardar únicamente la dirección primitiva \(v\).

\begin{corollary}[Registro conjunto de bloque y fase]

Escríbase la fase de esos cuatro bloques como \(a\in\mathbb Z/4\mathbb Z\). La actualización

\[
(a,g,s)\longmapsto(a+1,g+1,s+18)
\]

conserva \((s-18g,[g-a]_4)\). Ambos datos clasifican sus órbitas bilaterales: la diferencia \(N=g'-g\) satisface a la vez \(a'-a=[N]_4\) y \(s'-s=18N\). Esto reúne el avance de memoria y el retorno de la fase sin reemplazar ninguno por el otro.
\end{corollary}

\subsection{Generador parabólico determinado por la memoria}

La pareja ordenada de contadores dota al módulo \(M\) de la forma alternante unimodular

\[
\Omega_M((g,s),(g',s'))=gs'-sg'.
\]

Su dominio es el módulo de registro. No se identifica esta forma con un área física dimensional, ni con la forma del toro APP, que tiene otro dominio. La elección de orden \((g,s)\) fija su orientación.

\begin{theorem}[Estabilizador unimodular del incremento]

Defínase, directamente desde el vector primitivo \(v\) de \eqref{nuc06:eq:6},

\begin{equation}
N_v m=v\,\Omega_M(v,m)=v\mathcal I(m),\qquad
N_v=\begin{pmatrix}-18&1\\-324&18\end{pmatrix}.
\label{nuc06:eq:9}
\end{equation}

Entonces \(N_v^2=0\), \(N_v\ne0\) y

\begin{equation}
\boxed{\quad
\operatorname{Stab}_{SL_2(\mathbb Z)}(c_{108})
=\{I+nN_v:n\in\mathbb Z\}.
\quad}
\label{nuc06:eq:10}
\end{equation}

Cada transformación de \eqref{nuc06:eq:10} conserva \(\Omega_M\), el incremento \(c_{108}\) y la carga \(\mathcal I\). Para \(n\) no nulo es una matriz unipotente no identidad, de traza dos, perteneciente al tipo parabólico.
\end{theorem}

\begin{proof}
\(\Omega_M(v,v)=0\) implica \(N_v^2m=v\Omega_M(v,v)\Omega_M(v,m)=0\). Su entrada superior derecha es uno, de modo que es no nulo. La matriz

\[
B=\begin{pmatrix}1&0\\18&1\end{pmatrix}
\]

tiene determinante uno, lleva el primer vector de base a \(v\) y satisface

\begin{equation}
N_v=BNB^{-1},\qquad N=\begin{pmatrix}0&1\\0&0\end{pmatrix}.
\label{nuc06:eq:11}
\end{equation}

Un automorfismo entero fija \(4v\) exactamente cuando fija \(v\), porque \(M\) carece de torsión. Una matriz de determinante uno que fija el primer vector de base tiene necesariamente la forma \(\begin{pmatrix}1&n\\0&1\end{pmatrix}\). Conjugarla por \(B\) prueba \eqref{nuc06:eq:10}. El resto se deduce de esa expresión o directamente de \eqref{nuc06:eq:9}.
\end{proof}

La fórmula \eqref{nuc06:eq:9} hace que la construcción sea independiente del segundo vector elegido para completar \(v\) a una base orientada unimodular: cualquier otro complemento difiere en un múltiplo entero de \(v\) y produce el mismo \(N_v\). Bajo una transformación orientada \(R\in SL_2(\mathbb Z)\), se tiene \(N_{Rv}=RN_vR^{-1}\).

La acción puede escribirse con especial claridad:

\begin{equation}
(I+nN_v)m=m+n\mathcal I(m)v.
\label{nuc06:eq:12}
\end{equation}

La carga transversal determina cuánto desplaza esa simetría a lo largo de la dirección de memoria. La suma de vueltas produce el desplazamiento \(4Nv\); las simetrías estabilizadoras preservan esa ley de avance.

Sobre el cociente \eqref{nuc06:eq:8}, la transformación \(I+nN_v\) actúa como \((\mathcal I,[g]_4)\mapsto(\mathcal I,[g+n\mathcal I]_4)\). Por tanto, conserva la familia de órbitas y puede permutar sus clases residuales; \(\nu\) es invariante del retorno \(H\), no de cada transformación estabilizadora. Esta diferencia se deduce directamente de \eqref{nuc06:eq:12}.

La conjugación \eqref{nuc06:eq:11} proporciona un mapa algebraico explícito desde el modelo parabólico del TRIT, definido en \cref{euler-trit-generadores}, al módulo de registro extendido a escalares reales. Se obtiene además \(\exp(tN_v)=I+tN_v\) y \(B\exp(tN)=\exp(tN_v)B\). Este entrelazamiento concierne a operadores: el movimiento particular seleccionado por el TPK utiliza también sus emisiones y calendario. El grupo aritmético completo de \eqref{nuc06:eq:10} fija este lector; el grupo de simetrías admisibles del estado total se obtiene intersectándolo con las acciones que preservan los demás datos y relaciones.

En el sistema trasladado \eqref{nuc06:eq:6}, los cambios afines \(m\mapsto Rm+b\) con \(R\) en \eqref{nuc06:eq:10} y \(b\) arbitrario conmutan con \(H\). Los que preservan además el registro inicial fijado deben satisfacer \(Rm_0+b=m_0\). Distinguir ambas condiciones separa simetría de la ley y simetría de una trayectoria señalada.

\begin{theorem}[Realización variacional de la carga y del estabilizador]
\label{nuc06:variacional}

La extensión escalar del registro, \(M_{\mathbb R}=M\otimes_{\mathbb Z}\mathbb R\), conserva las coordenadas generadas \((g,s)\) y la forma \(\Omega_M=dg\wedge ds\). Es una realización posterior del módulo entero. Fijando la convención \(\iota_{X_H}\Omega_M=dH\), se tiene

\begin{equation}
X_{\mathcal I}=v,\qquad
X_{\frac12\mathcal I^2}=\mathcal I v=N_vm.
\label{nuc06:eq:12a}
\end{equation}

Así, la traslación y la familia parabólica construidas admiten generadores hamiltonianos explícitos. El flujo de \(\mathcal I\) es \(m\mapsto m+uv\); en \(u\) igual a cuatro reproduce \eqref{nuc06:eq:6}. El flujo de \(\mathcal I^2/2\) es \(m\mapsto(I+uN_v)m\). El parámetro \(u\) corresponde a esta realización continua; el paso discreto de retorno y las emisiones intermedias mantienen su calendario TPK propio.
\end{theorem}

\begin{proof}
Para un vector \(X=(X_g,X_s)\), \(\iota_X(dg\wedge ds)=X_g\,ds-X_s\,dg\). Como \(d\mathcal I=ds-18\,dg\), el primer campo es \((1,18)=v\). La regla de la cadena da el segundo. \(\mathcal I\) permanece constante a lo largo de ambos campos, porque \(d\mathcal I(v)=0\). Por tanto el segundo flujo integra \(dm/du=\mathcal I(m_0)v\), equivalente a \(\exp(uN_v)=I+uN_v\). La no degeneración de \(\Omega_M\) implica que el Hamiltoniano que genera \(N_v\) es \(\mathcal I^2/2\) más una constante aditiva: su diferencial queda fijado por \(\iota_{N_vm}\Omega_M\).
\end{proof}

La acción matemática

\begin{equation}
\mathscr A[m]=\int_{u_0}^{u_1}
\left[s\frac{dg}{du}-\frac12(s-18g)^2\right]du
\label{nuc06:eq:12b}
\end{equation}

produce por variación las ecuaciones \(dg/du=\mathcal I\), \(ds/du=18\mathcal I\). La traslación \(g\mapsto g+\varepsilon\), \(s\mapsto s+18\varepsilon\) conserva el Hamiltoniano y cambia el lagrangiano en la derivada total \(d(18\varepsilon g)/du\), para \(\varepsilon\) constante. Su carga de Noether es

\begin{equation}
\underbrace{s\,\delta g}_{s}
-\underbrace{18g}_{\text{término de borde}}
=s-18g=\mathcal I.
\label{nuc06:eq:12c}
\end{equation}

La variación se interpreta por unidad de \(\varepsilon\). Las ecuaciones comprueban directamente \(d\mathcal I/du=18\mathcal I-18\mathcal I=0\). La propia traslación discreta \eqref{nuc06:eq:6} se realiza también mediante el Hamiltoniano \(H_{\rm ret}=4\mathcal I\), cuyo flujo de parámetro uno es \(H\); ambos Hamiltonianos conmutan en el sentido de Poisson. El flujo estabilizador \(I+uN_v\) y el retorno \(m+4v\) son acciones distintas que conmutan: en la recta \(\mathcal I=0\) el primero fija cada punto, mientras el segundo sigue acumulando memoria.

Más generalmente, todo Hamiltoniano \(C^1\) sobre este plano que sea invariante bajo las traslaciones \(m\mapsto m+uv\) tiene la forma \(F(\mathcal I)\). Para probarlo, las coordenadas globales \((g,\mathcal I)\) escriben esa simetría como traslación de \(g\); la invariancia hace independiente de \(g\) a \(H\). Su campo es \(F'(\mathcal I)v\). Esta clasificación explicita tanto la conservación común como la elección adicional de una dinámica concreta.

La acción \eqref{nuc06:eq:12b} es adimensional y pertenece al plano de registro realizado. Su demostración proporciona una relación variacional exacta entre memoria, carga y estabilizador. La conversión en acción física y en un tiempo dimensional se realiza mediante los lectores de acción y reloj correspondientes; no se asignan esas unidades al contador u.

\subsection{Curvatura de APP y conservación de la fase central}

La sección~\ref{iv:extension-curvatura-mixta} trabaja sobre \(X_9=(\mathbb Z/9\mathbb Z)^2\). Las hojas son \(S(x,y)=x+y\) y \(P(x,y)=xy\). Las conexiones mixtas utilizan diferencias de esos potenciales, no los potenciales como enlaces:

\begin{equation}
F(S,P)=\Delta_x\Delta_yP-\Delta_y\Delta_xS=1,
\qquad F(P,S)=-1\pmod9.
\label{nuc06:eq:13}
\end{equation}

La suma orientada de enlaces en el contorno de un rectángulo es \(ab\) o \(-ab\) según el orden. El área alternante es \(\sigma(u,v)=ad-bc\) para \(u=(a,b)\), \(v=(c,d)\). La sección~\ref{iv:weyl-area-extension} representa esos transportes mediante \(D(u)=X^aZ^b\), \(ZX=\omega XZ\), con \(\omega\) la raíz novena en la representación posterior. Entonces

\begin{equation}
D(u)D(v)=\omega^{bc}D(u+v),\qquad
D(u)D(v)D(u)^{-1}D(v)^{-1}=\omega^{-\sigma(u,v)}I.
\label{nuc06:eq:14}
\end{equation}

La posición resultante del conmutador retorna. La fase central retiene el área orientada con el signo fijado en \eqref{nuc06:eq:14}. Es una retención módulo nueve: un área múltiplo de nueve tiene fase central trivial, aunque un levantamiento con cociente y registro pueda distinguir sus recorridos. La fase no reemplaza ese levantamiento.

\begin{proposition}[Covarianza exacta del transporte central]

Como dos es invertible módulo nueve, defínase

\begin{equation}
W(a,b)=\omega^{ab/2}D(a,b),\qquad
W(u)W(v)=\omega^{-\sigma(u,v)/2}W(u+v).
\label{nuc06:eq:15}
\end{equation}

Para \(R\in SL_2(\mathbb Z/9\mathbb Z)\), la aplicación

\begin{equation}
\omega^zW(u)\longmapsto\omega^zW(Ru)
\label{nuc06:eq:16}
\end{equation}

es un automorfismo de la extensión central. Conserva sus productos y sus fases de conmutación.
\end{proposition}

\begin{proof}
La fase al multiplicar los dos \(W\) es
\(\frac{ab}{2}+\frac{cd}{2}+bc-\frac{(a+c)(b+d)}2=\frac{bc-ad}{2}\). Esto prueba \eqref{nuc06:eq:15}. La identidad \(\sigma(Ru,Rv)=\det(R)\sigma(u,v)=\sigma(u,v)\) muestra que \eqref{nuc06:eq:16} conserva la ley; \(R\) inversible proporciona el automorfismo inverso. Si \(\det R=-1\), el automorfismo correspondiente cambia también \(z\mapsto-z\); el signo de área y el centro se invierten conjuntamente.
\end{proof}

El registro íntegro y su fase central tienen tipos diferentes. La covarianza permite conservar la ley de área mientras cambia la presentación; la clase afín de \eqref{nuc06:eq:3} permite conservar el contenido de memoria mientras cambia su origen. Ambas propiedades actúan sobre lectores tipados de rutas, y sus mapas deben mantenerse separados al componerlos con la lectura excepcional del estado completo.

\begin{proposition}[Persistencia bajo marcos locales y subdivisiones]

Sea un lazo \(\gamma\) basado en \(x\) y sean sus enlaces invertibles \(U_\epsilon\). El cambio local \(U'_\epsilon=G_tU_\epsilon G_s^{-1}\) transforma su holonomía en \(H'_\gamma=G_xH_\gamma G_x^{-1}\): los marcos de los vértices interiores se cancelan sucesivamente. Si \(H_\gamma=\omega^kI\), la fase central permanece exactamente igual. Para enlaces refinados cuyo producto reproduce el enlace anterior, la subdivisión conserva asimismo \(H_\gamma\), por asociatividad. En las plaquetas APP la cancelación de aristas interiores materializa esta propiedad mediante Stokes finito.
\end{proposition}

Un rectángulo elemental tiene Wilson aditivo igual a uno y fase central no trivial. Por tanto, ese transporte mixto distingue una holonomía real de la identidad que producirían enlaces de la forma \(G_tG_s^{-1}\). La conservación bajo subdivisión concierne a lazos y enlaces compatibles: añadir otra área cambia el transporte.

El levantamiento entero de la holonomía distingue el área orientada de su residuo nonádico: para un rectángulo de lados cuatro y siete, \(28=1+9\cdot3\), mientras que \(-28=8+9(-4)\) corresponde a su orientación inversa. La fase residual coincide para áreas uno y veintiocho, mientras sus cocientes distinguen las acumulaciones. Este ejemplo exhibe la diferencia entre fase finita y registro íntegro; su cociente de área y el contador temporal \(q_{\rm mem}\) conservan dominios distintos.

\subsection{Composición entre los tres regímenes y continuidad del balance}

Las realizaciones de \cref{euler-trit-generadores,euler-trit-evaluaciones} construyen generadores \(J_k^2=-\tau_kI\) sobre sus fibras propias, de dimensión real dos. Se toman enlaces efectivamente declarados \(B_k:V_k\to V_{k+1}\) que preservan las formas alternantes, \(B_k^*\Omega_{k+1}B_k=\Omega_k\). El signo \(*\) en esta sección significa transposición/pullback bilineal. Los operadores \(E_k=\exp(t_kJ_k)\) preservan \(\Omega_k\): en dimensión dos, la traza nula del generador implica \(J_k^*\Omega_k+\Omega_kJ_k=0\). La identidad simpléctica es la hipótesis que se conservaría al extender el lema a dimensiones mayores.

Para la compresión nonádica y su complemento registrado,

\[
T_k=B_k(8I+E_k)/9,\qquad
D_k=B_k\sqrt8(E_k-I)/9,
\]

la expansión algebraica da

\begin{equation}
\Omega_k=T_k^*\Omega_{k+1}T_k+D_k^*\Omega_{k+1}D_k.
\label{nuc06:eq:17}
\end{equation}

En efecto, los términos cruzados se cancelan y el numerador es \(72\Omega_k+9E_k^*\Omega_kE_k=81\Omega_k\). Para \(P_0=I\), \(P_{k+1}=T_kP_k\), se obtiene, por sustitución sucesiva y en el orden de la ruta,

\begin{equation}
\boxed{\quad
\Omega_0=P_N^*\Omega_NP_N+
\sum_{k=0}^{N-1}(D_kP_k)^*\Omega_{k+1}(D_kP_k).
\quad}
\label{nuc06:eq:18}
\end{equation}

El terminal y cada complemento permanecen presentes. La identidad es bilineal orientada; los resultados de positividad se aplican en los perfiles positivos demostrados en las secciones anteriores. La forma conservada no obliga a que un enlace entre regímenes distintos entrelace sus generadores: si \(BJ_\tau=J_{\tau'}B\), elevar al cuadrado con \(B\) invertible impondría \(\tau=\tau'\). La conservación de una estructura común permite mantener diferenciados los tres regímenes.

El transporte íntegro \(U_N=(B_{N-1}E_{N-1})\cdots(B_0E_0)\) lleva \(V_0\) a \(V_N\). Una monodromía matricial requiere un cierre tipado \(K:V_N\to V_0\); sólo entonces \(KU_N\) es un endomorfismo cuya traza puede calcularse. Cambios de base compatibles conservan esa traza por conjugación. Retorno de fase, cierre de fibra y retorno del vector siguen siendo condiciones diferentes.

Si los enlaces tienen exclusivamente la forma \(B_{yx}=G_yG_x^{-1}\), su producto a lo largo de un camino es \(G_{\rm final}G_{\rm inicial}^{-1}\). Al cerrar con el mismo marco da identidad. La holonomía de \eqref{nuc06:eq:14} y el incremento de \eqref{nuc06:eq:6} muestran mecanismos concretos que esa mera sustitución de marcos no puede representar por sí sola.

\subsection{Compatibilidad con refinamiento y alcance arbitrario}

La acción de memoria es compatible con un homomorfismo de módulos de restricción \(p:M'\to M\) cuando

\begin{equation}
pC'=Cp,\qquad p\kappa'=\kappa.
\label{nuc06:eq:19}
\end{equation}

Entonces \(pH'=Hp\) y \(p\) induce un mapa de cocientes

\begin{equation}
M'/(I-C')M'\longrightarrow M/(I-C)M,
\qquad [\kappa']\longmapsto[\kappa].
\label{nuc06:eq:20}
\end{equation}

\begin{proof}
La primera igualdad de \eqref{nuc06:eq:19} lleva \((I-C')M'\) a \((I-C)M\); la segunda fija la imagen del incremento. La compatibilidad de \(H\) resulta por sustitución. Concatenaciones y refinamientos componen los mismos cuadrados gracias a la ley \eqref{nuc06:eq:1}.
\end{proof}

Una clase de retorno no nula que aparece en un nivel conservado permanece no nula en cualquier levantamiento compatible: una clase cero se proyectaría a cero. Esta afirmación se aplica al sistema de memorias y restricciones que satisface \eqref{nuc06:eq:19}; el constructor conjunto y los mapas de su sistema inverso se conservan en \cref{iv:continuo-conjunto}.

En el registro de \eqref{nuc06:eq:6}, para cada entero \(N\) se tiene \(H^N(m)=m+N(4,72)\). La prueba es algebraica para \(N\) arbitrario, no una inferencia desde una cantidad finita de iteraciones. En la prolongación hacia adelante, cada prefijo conserva su registro finito; la historia completa se describe mediante la familia compatible de esos registros, sin reemplazar sus contadores por una cantidad entera «infinita».

Si una dinámica completa tiene el lector \(q_{\rm mem}\) con \(q_{\rm mem}(\Gamma_9x)=q_{\rm mem}(x)+1\), tampoco puede poseer un estado periódico de período \(N\) positivo en ese dominio: aplicar el lector produciría \(q_{\rm mem}(x)=q_{\rm mem}(x)+N\). Este corolario utiliza el grado entero completo; sus cocientes residuales pueden tener períodos finitos, como corresponde al calendario observable del corpus.

\subsection{Consecuencia fundacional y alcance de la composición}

El enlace específico obtenido es

\begin{equation}
\text{incremento de retorno }4v
\ \longrightarrow\ \mathcal I(m)=\Omega_M(v,m)
\ \longrightarrow\ N_vm=v\mathcal I(m)
\ \longrightarrow\ I+nN_v.
\label{nuc06:eq:21}
\end{equation}

El registro acumulado determina una carga conservada y una familia parabólica que estabiliza su dirección. La evolución y la conservación pertenecen a la misma acción: cambian \(g\) y \(s\), permanece su relación \(s-18g\), y el incremento conserva un estabilizador explícito. Esto precisa una parte del vínculo entre memoria y simetría que motivó la investigación.

La igualdad \eqref{nuc06:eq:18} conecta esa investigación con los balances anteriores: una evolución puede redistribuir la forma entre lectura y complementos conservando el total. Las identidades de APP \eqref{nuc06:eq:13}–\eqref{nuc06:eq:16} muestran, en otro lector concreto del mismo origen, cómo el orden del transporte produce un dato central aun cuando retorna la posición.

La carga \eqref{nuc06:eq:7} es un invariante algebraico del transporte y, mediante la acción explícita \eqref{nuc06:eq:12b}, una carga de Noether de la realización del registro. El teorema~\ref{nuc06:variacional} y su desarrollo variacional construyen esa relación, su simetría y su término de borde. Los desarrollos anteriores de acción y corriente físicas conservan sus dominios y unidades; transportarlos a esta realización exige mantener sus mapas dimensionales.

La composición con los lectores dimensionales debe transportar el par \((v,\Omega_M)\) y la carga \(\mathcal I\), preservando sus unidades y la corriente variacional. Toda relación posterior con \(\pi,\phi,e,\alpha\) y \(K\) ha de proceder de esa composición; las fórmulas presentes fijan la información que debe conservarse.

\section{Conservación variacional y acción orientada}
\label{nuclear:7}
\subsection{Simetría de la acción elíptica y carga de Noether}

Se fija una sección positiva de acción \(\hbar\), el par angular \(A,C^*\) en una misma unidad, \(A\ne0\), \(|C^*/A|<1\), y

\[
\eta=\operatorname{artanh}(C^*/A).
\]

El plano realizado \(\mathcal P_\eta\) tiene coordenadas \((Q,P)\), ambas de dimensión raíz de acción, con

\[
\lambda=P\,dQ,\qquad \Omega=dQ\wedge dP=-d\lambda.
\]

La construcción de \cref{iv:ellipse:section} da

\begin{equation}
\mathcal I_\eta(Q,P)=\frac12(e^{-\eta}Q^2+e^\eta P^2),
\qquad H_\eta=\omega\mathcal I_\eta,\qquad \omega>0.
\label{nuc07:eq:1}
\end{equation}

La elipse HMT es \(\mathcal I_\eta=\hbar\). Su área es \(2\pi\hbar\); \(\mathcal I_\eta\) tiene dimensión de acción y \(H_\eta\) de energía. Definimos las matrices ya realizadas

\begin{equation}
M_\eta=\operatorname{diag}(e^{\eta/2},e^{-\eta/2}),\quad
J_2=\begin{pmatrix}0&-1\\1&0\end{pmatrix},\quad
R_\eta(\varepsilon)=M_\eta e^{-\varepsilon J_2}M_\eta^{-1}.
\label{nuc07:eq:2}
\end{equation}

\begin{proposition}
El grupo \(R_\eta\) es una simetría canónica de la acción

\begin{equation}
\mathscr S_\eta[z]=\int_{t_a}^{t_b}
\left(P\dot Q-\omega\mathcal I_\eta(Q,P)\right)dt
\label{nuc07:eq:3}
\end{equation}

y su carga de Noether es \(\mathcal I_\eta\).
\end{proposition}

\begin{proof}
El determinante de \(M_\eta\) es uno y \(M_\eta^{\mathsf T}J_2M_\eta=J_2\). Además,
\(\mathcal I_\eta(M_\eta y)=\|y\|^2/2\), invariante bajo \(e^{-\varepsilon J_2}\).
El generador

\[
X_\eta=e^\eta P\,\partial_Q-e^{-\eta}Q\,\partial_P
\]

satisface \(\iota_{X_\eta}\Omega=d\mathcal I_\eta\). Por la fórmula de Cartan,

\[
\mathcal L_{X_\eta}\lambda=dF_\eta,\qquad
F_\eta=\lambda(X_\eta)-\mathcal I_\eta=e^\eta P^2-\mathcal I_\eta.
\]

Como \(X_\eta H_\eta=0\), la variación de la densidad de acción es \(dF_\eta/dt\). Su carga es

\[
\frac{\partial L}{\partial\dot Q}X_\eta Q-F_\eta
=Pe^\eta P-F_\eta=\mathcal I_\eta.
\]

Las ecuaciones \(\dot Q=\omega e^\eta P\), \(\dot P=-\omega e^{-\eta}Q\) verifican directamente \(d\mathcal I_\eta/dt=0\).
\end{proof}

La parametrización positiva de la elipse es \(Q=a_S\cos\theta\), \(P=b_S\sin\theta\). Estas ecuaciones le asignan \(\dot\theta=-\omega\). Por tanto, en una vuelta recorrida según el flujo dinámico, \(\oint P\,dQ=2\pi\mathcal I_\eta\); con \(\theta\) creciente, la integral es \(-2\pi\mathcal I_\eta\). Es la distinción de orientación establecida en \cref{iv:ellipse:area}. En el contorno radional, el módulo es \(h=2\pi\hbar\).

\subsection{Cambio de forma: transporte y conexión determinados}

Para dos formas ya publicadas \(\eta,\eta'\) y un avance de fase \(\vartheta\), se define

\begin{equation}
\mathcal T_{\eta',\eta}^{\vartheta}:\mathcal P_\eta\to\mathcal P_{\eta'},
\qquad
\mathcal T_{\eta',\eta}^{\vartheta}
=M_{\eta'}e^{-\vartheta J_2}M_\eta^{-1}.
\label{nuc07:eq:4}
\end{equation}

\begin{proposition}
Este transporte es simpléctico y cumple

\begin{equation}
\mathcal I_{\eta'}(\mathcal T_{\eta',\eta}^{\vartheta}z)
=\mathcal I_\eta(z),\qquad
\mathcal T_{\eta'',\eta'}^{\vartheta_2}
\mathcal T_{\eta',\eta}^{\vartheta_1}
=\mathcal T_{\eta'',\eta}^{\vartheta_1+\vartheta_2}.
\label{nuc07:eq:5}
\end{equation}
\end{proposition}

\begin{proof}
Todos los factores son simplécticos. La primera igualdad se reduce a la invariancia de \(\|y\|^2/2\) bajo una rotación. En la segunda se cancelan las matrices contiguas \(M_{\eta'}^{-1}M_{\eta'}\) y se suman las fases.
\end{proof}

La fórmula recibe las formas y los avances del lector; no selecciona sus valores. La lista completa de transiciones puede retenerse junto con el transporte compuesto: dos listas que den el mismo mapa compuesto no se identifican como historias.

\begin{proposition}
Una realización diferenciable \(\eta(t),\vartheta(t)\in C^1\), con \(\dot\vartheta=\omega(t)\), del transporte \eqref{nuc07:eq:4} tiene generador

\begin{equation}
\boxed{
H_{\rm tr}(Q,P,t)=
\omega(t)\mathcal I_{\eta(t)}(Q,P)
+\frac{\dot\eta(t)}2 QP.}
\label{nuc07:eq:6}
\end{equation}

Este generador conserva \(\mathcal I_{\eta(t)}(z(t))\).
\end{proposition}

\begin{proof}
Al derivar
\(z(t)=M_{\eta(t)}e^{-\vartheta(t)J_2}y_0\) se obtiene

\[
\dot Q=\omega e^\eta P+\frac{\dot\eta}2Q,\qquad
\dot P=-\omega e^{-\eta}Q-\frac{\dot\eta}2P.
\]

Son las ecuaciones de Hamilton de \eqref{nuc07:eq:6}. El término de conexión realiza exactamente \(\dot M_\eta M_\eta^{-1}z\), y su Hamiltoniano queda determinado salvo una función exclusiva de \(t\). En efecto,

\begin{equation}
\begin{aligned}
\frac{d\mathcal I_\eta}{dt}
&=\frac{\dot\eta}{2}(-e^{-\eta}Q^2+e^\eta P^2)
+\left\{\mathcal I_\eta,\frac{\dot\eta}{2}QP\right\}\\
&=\frac{\dot\eta}{2}(-e^{-\eta}Q^2+e^\eta P^2)
+\frac{\dot\eta}{2}(e^{-\eta}Q^2-e^\eta P^2)=0.
\end{aligned}
\label{nuc07:eq:7}
\end{equation}

Equivalentemente, \(Q=e^{\eta/2}q\), \(P=e^{-\eta/2}p\) dan

\begin{equation}
P\dot Q-H_{\rm tr}
=p\dot q-\frac{\omega(t)}2(q^2+p^2).
\label{nuc07:eq:8}
\end{equation}

La simetría de rotación y su carga se transportan exactamente.
\end{proof}

Aquí se ha derivado el término de conexión del transporte, no de una fuerza externa. El Hamiltoniano LC de forma fija conserva \(\mathcal I_\eta\). Durante una deformación, omitir el segundo término de \eqref{nuc07:eq:6} deja sin cancelar el primer término de \eqref{nuc07:eq:7}. La conservación variacional exige transportar también la forma canónica.

La proposición determina el generador de la interpolación declarada; no asigna una interpolación física única a toda historia TPK. La identidad discreta \eqref{nuc07:eq:4}–\eqref{nuc07:eq:5} funciona sin interpolación.

\subsection{Orientación y frontera de la acción}

Las transformaciones de \cref{iv:ellipse:reciprocidad} distinguen \(S(Q,P)=(P,Q)\) y \(J_2(Q,P)=(-P,Q)\). Ambos intercambian la forma \(\eta\leftrightarrow-\eta\), pero

\begin{equation}
S^*\Omega=-\Omega,\quad J_2^*\Omega=\Omega,\qquad
S^*\lambda=-\lambda+d(PQ),\quad
J_2^*\lambda=\lambda-d(PQ).
\label{nuc07:eq:9}
\end{equation}

Asimismo,

\begin{equation}
S R_\eta(\varepsilon)S^{-1}=R_{-\eta}(-\varepsilon),\qquad
J_2R_\eta(\varepsilon)J_2^{-1}=R_{-\eta}(\varepsilon).
\label{nuc07:eq:10}
\end{equation}

Se comprueban usando \(SM_\eta=M_{-\eta}S\), \(J_2M_\eta=M_{-\eta}J_2\), \(SJ_2S^{-1}=-J_2\). Para una curva abierta orientada \(z:[a,b]\to\mathcal P_\eta\),

\begin{equation}
\int_{J_2z}\lambda=\int_z\lambda-[PQ]_a^b,\qquad
\int_{Sz}\lambda=-\int_z\lambda+[PQ]_a^b.
\label{nuc07:eq:11}
\end{equation}

Los términos de frontera se anulan en un ciclo y se cancelan al concatenar extremos compatibles. La energía cuadrática por sí sola no determina estos signos ni esos términos.

La monodromía nonádica registra retorno de fase y \(q_{\rm mem}\mapsto q_{\rm mem}+1\). La constancia de una carga no obliga al contador a permanecer fijo. En la realización elíptica, una vuelta conserva \(\mathcal I_\eta\) mientras incrementa la integral acumulada orientada en \(2\pi\mathcal I_\eta\). Carga, acción integrada, fase y registro tienen tipos distintos. No se identifica aquí una vuelta angular con cualquier ciclo TPK sin especificar su lector.

\subsection{Carga variacional del reloj y balance entre lectura y memoria}

\subsubsection{Una carga de acción explícita del reloj}

El calendario se realiza sobre \(\mathcal H_{\rm clk}=\mathbb C^{108}\), con proyectores Fourier \(P_k=|\widetilde k\rangle\langle\widetilde k|\), y determina

\begin{equation}
N_{108}=\sum_{k=0}^{107}kP_k,\quad
H_{\rm clk}=\hbar\omega_0N_{108},\quad
\omega_0=\frac{2\pi}{108t_0}.
\label{nuc07:eq:12}
\end{equation}

Su propagador durante \(t_0\) es el desplazamiento de una celda. La acción del reloj es

\begin{equation}
\mathscr S[\psi]=\int
\left[\frac{i\hbar}{2}
(\langle\psi,\dot\psi\rangle-\langle\dot\psi,\psi\rangle)
-\langle\psi,H_{\rm clk}\psi\rangle\right]dt.
\label{nuc07:eq:13}
\end{equation}

El producto es antilineal en la primera variable. La ecuación es \(i\hbar\dot\psi=H_{\rm clk}\psi\).

\begin{proposition}
Para \(B=B^*\) con \([B,H_{\rm clk}]=0\), la simetría \(\psi\mapsto e^{-i\varepsilon B}\psi\) tiene carga

\begin{equation}
\mathcal Q_B(\psi)=\hbar\langle\psi,B\psi\rangle.
\label{nuc07:eq:14}
\end{equation}
\end{proposition}

\begin{proof}
Al localizar el parámetro, \(\delta\psi=-i\varepsilon(t)B\psi\), se obtiene
\(\delta L=\hbar\dot\varepsilon\langle\psi,B\psi\rangle\).
La integración por partes en las ecuaciones variacionales da \(d\mathcal Q_B/dt=0\). Directamente,

\[
\frac{d}{dt}\langle\psi,B\psi\rangle
=\frac{i}{\hbar}\langle\psi,[H_{\rm clk},B]\psi\rangle=0.
\]

Con \(B=I\) se conserva norma; con \(B=N_{108}\) la carga es la acción \(E_{\rm clk}/\omega_0\).
\end{proof}

La forma simpléctica del espacio real subyacente es
\(\Omega_{\rm clk}(u,v)=2\hbar\operatorname{Im}\langle u,v\rangle\);
su uno-forma es
\(\Theta_\psi(v)=-\hbar\operatorname{Im}\langle\psi,v\rangle\),
con \(\Omega_{\rm clk}=-d\Theta\).

\subsubsection{Balance exacto de carga}

Sean \(\mathcal H,\mathcal K\) espacios hermíticos finitos, \(J:\mathcal H\to\mathcal K\) isometría, \(U\) unitario sobre \(\mathcal K\), y

\begin{equation}
\Pi=JJ^*,\quad T=J^*UJ,\quad
\eta_{\rm mem}=(I-\Pi)UJ.
\label{nuc07:eq:15}
\end{equation}

El operador \(\eta_{\rm mem}\) es distinto del parámetro escalar \(\eta\).

\begin{proposition}
Si

\begin{equation}
\widetilde B=\widetilde B^*,\quad U^*\widetilde BU=\widetilde B,
\quad [\widetilde B,\Pi]=0,\quad B=J^*\widetilde BJ,
\label{nuc07:eq:16}
\end{equation}

entonces

\begin{equation}
B=T^*BT+\eta_{\rm mem}^*\widetilde B\eta_{\rm mem},
\qquad
\mathcal Q_B(\psi)=
\mathcal Q_B(T\psi)+\mathcal Q_{\widetilde B}(\eta_{\rm mem}\psi).
\label{nuc07:eq:17}
\end{equation}
\end{proposition}

\begin{proof}
Se expande \(UJ\psi=JT\psi+\eta_{\rm mem}\psi\) en
\(J^*U^*\widetilde BUJ=J^*\widetilde BJ\).
La conmutación con \(\Pi\) anula los términos cruzados. Evaluar la identidad operatoria y multiplicar por \(\hbar\) da el balance de acción.
\end{proof}

Por iteración,

\begin{equation}
\mathcal Q_B(\psi)=\mathcal Q_B(T^n\psi)
+\sum_{j=0}^{n-1}\mathcal Q_{\widetilde B}
(\eta_{\rm mem}T^j\psi).
\label{nuc07:eq:18}
\end{equation}

Esta suma registra complementos sucesivos de la dinámica comprimida. La dinámica completa tiene otra fórmula exacta, con \(z_n=U^nJ\psi\):

\begin{equation}
\mathcal Q_B(\psi)=
\mathcal Q_B(J^*z_n)+
\mathcal Q_{\widetilde B}((I-\Pi)z_n).
\label{nuc07:eq:19}
\end{equation}

No se identifica \(T^n\) con \(J^*U^nJ\) sin otra propiedad: el estado completo puede reabsorber memoria. Si \(B,\widetilde B\ge0\), todos los sumandos de \eqref{nuc07:eq:18} son no negativos; para cargas firmadas permanece la igualdad, sin esa interpretación de positividad.

Si \(U\) es propagador de una acción como \eqref{nuc07:eq:13}, la conmutación de \(\widetilde B\) con su Hamiltoniano aporta la conservación de Noether. La condición de reducción por \(\Pi\) especifica las lecturas que separan carga sin términos de interferencia.

Sin imponer \eqref{nuc07:eq:16}, la identidad general conserva los términos adicionales:

\begin{equation}
\begin{aligned}
&\mathcal Q_B(T\psi)+\mathcal Q_{\widetilde B}(\eta_{\rm mem}\psi)
+2\hbar\operatorname{Re}
\langle JT\psi,\widetilde B\eta_{\rm mem}\psi\rangle\\
&=\mathcal Q_B(\psi)+
\hbar\langle\psi,J^*(U^*\widetilde BU-\widetilde B)J\psi\rangle.
\end{aligned}
\label{nuc07:eq:20}
\end{equation}

La carga física utiliza su lector, por ejemplo \(N_{108}\); una identidad de normas corresponde a \(B=I\). Las secciones \(\hbar_{\rm pre}\) y \(\hbar_{\rm ret}\) conservan el lector de retorno de \cref{iv:action:secciones,iv:action:razon}; la diferencia entre ellas y el complemento de memoria poseen sus respectivos dominios de definición.

\subsection{Balance covariante del tensor total}

La eliminación variacional de la contorsión y la identidad de Noether de la acción covariante, desarrolladas en la proposición~\ref{vii:prop:fuente-metrica-balance}, proporcionan

\begin{equation}
\mathcal T_{\mu\nu}
=T^{\rm phys,LC}_{\mu\nu}+U^{\rm phys,spin}_{\mu\nu},
\qquad
\mathring\nabla^\mu\mathcal T_{\mu\nu}=0.
\label{nuc07:eq:21}
\end{equation}

Su dominio es la carta HMT–MD realizada, con los acoplamientos mantenidos constantes durante la variación. Para un campo de Killing \(\xi\),

\begin{equation}
j^\mu_\xi=\mathcal T^{\mu\nu}\xi_\nu,\qquad
\mathring\nabla_\mu j^\mu_\xi=0.
\label{nuc07:eq:22}
\end{equation}

La prueba usa \eqref{nuc07:eq:21}, la simetría de \(\mathcal T\) y la antisimetría de \(\mathring\nabla_\mu\xi_\nu\). Al integrar, la diferencia entre cargas de las hipersuperficies es el flujo por la frontera lateral.

La conservación pertenece al tensor total; cada componente puede intercambiar carga. El teorema~\ref{vii:thm:intercambio-traza-residual} determina explícitamente \(\mathcal J^{\rm tr}_\nu=\partial_\nu\operatorname{tr}_g(T^{\rm res})/4\) y los balances opuestos de traza y parte de traza nula. Esto se recupera como instancia covariante del criterio de balance conjunto, no como nueva ecuación gravitatoria.

\section{Conservación de formas de acción y del área orientada en los regímenes tríticos}
\label{nuclear:8}
\subsection{Balance entre dos formas de acción}

Las salidas angulares \(A,C^*\), en una misma unidad, determinan la forma \(\eta=\operatorname{artanh}(C^*/A)\) sobre el dominio \(A\ne0\), \(|C^*/A|<1\). Recibidas dos formas sucesivas \(\eta_n,\eta_{n+1}\), definimos

\[
M_n=\operatorname{diag}(e^{\eta_n/2},e^{-\eta_n/2}),\qquad
G_n=M_n^{-\mathsf T}M_n^{-1},\qquad
I_n(z)=\tfrac12z^\mathsf TG_nz.
\]

Aquí \(z=(Q,P)^\mathsf T\), con ambas coordenadas de dimensión raíz de acción. La escala \(\hbar\), cuando fija el contorno \(I_n=\hbar\), es la salida de acción previamente construida; no interviene como objetivo de ajuste.

Sea \(R_n\) la rotación canónica del avance angular \(\theta_n\). El transporte sin avance angular y el transporte con dicho avance son, respectivamente,

\begin{equation}
B_n=M_{n+1}M_n^{-1},\qquad
C_n=M_{n+1}R_nM_n^{-1}.
\label{nuc08:eq:1}
\end{equation}

Ambos van de la carta \(n\) a la carta \(n+1\). Definimos la lectura comprimida y su registro complementario normalizado:

\begin{equation}
T_n=\frac{8B_n+C_n}{9},\qquad
D_n=\frac{\sqrt8}{9}(C_n-B_n).
\label{nuc08:eq:2}
\end{equation}

El operador \(D_n\) es una realización isométricamente normalizada del único patrón de diferencias entre las nueve fases. El registro de nueve componentes es \((8z,-z,\ldots,-z)/27\), con \(z=(C_n-B_n)v\); la suma de sus productos interiores da el mismo factor \(8/81\). Se mantiene el registro de nueve fases cuando las etiquetas individuales intervienen en otros lectores.

\begin{theorem}
Se conserva la acción entre lectura y memoria:

\begin{equation}
\boxed{G_n=T_n^\mathsf TG_{n+1}T_n+D_n^\mathsf TG_{n+1}D_n.}
\label{nuc08:eq:3}
\end{equation}
\end{theorem}

\begin{proof}
La normalización de las cartas da
\(B_n^\mathsf TG_{n+1}B_n=G_n=C_n^\mathsf TG_{n+1}C_n\).
Al expandir \eqref{nuc08:eq:2}, las contribuciones cruzadas tienen coeficientes \(+8\) y \(-8\). Los coeficientes diagonales restantes son \(72\) para \(B_n\) y \(9\) para \(C_n\), divididos por \(81\). La suma es \(G_n\).
\end{proof}

Con \(z_{n+1}=T_nz_n\), \(m_n=D_nz_n\), la forma escalar es

\begin{equation}
\boxed{I_n(z_n)=I_{n+1}(z_{n+1})+I_{n+1}(m_n).}
\label{nuc08:eq:4}
\end{equation}

La suma telescópica, reteniendo cada memoria en la carta donde fue emitida, da

\begin{equation}
I_0(z_0)=I_N(z_N)+\sum_{n=0}^{N-1}I_{n+1}(m_n).
\label{nuc08:eq:5}
\end{equation}

El transporte \(B_n\) es material en esta composición: sustituirlo por la identidad cuando las formas difieren altera los dominios y puede romper \eqref{nuc08:eq:3}.

\subsection{Ley exacta de acumulación con fase variable}

En coordenadas normalizadas \(y_n=M_n^{-1}z_n\), la lectura es
\(y_{n+1}=(8I+R_n)y_n/9\). Para una rotación plana,

\begin{equation}
\left(\frac{8I+R_n}{9}\right)^\mathsf T
\left(\frac{8I+R_n}{9}\right)=r_n I,\qquad
r_n=1-\frac{32}{81}\sin^2(\theta_n/2).
\label{nuc08:eq:6}
\end{equation}

Por tanto,

\begin{equation}
I_N(z_N)=I_0(z_0)\prod_{n=0}^{N-1}r_n.
\label{nuc08:eq:7}
\end{equation}

\begin{theorem}
Para \(I_0(z_0)>0\), el terminal tiende a cero exactamente cuando

\begin{equation}
\sum_{n\ge0}\sin^2(\theta_n/2)=+\infty.
\label{nuc08:eq:8}
\end{equation}
\end{theorem}

\begin{proof}
Sea \(c=32/81\) y \(u_n=\sin^2(\theta_n/2)\in[0,1]\). Se cumple
\(cu_n\leq-\log(1-cu_n)\leq cu_n/(1-c)\), por integración de \(1/(1-s)\) entre cero y \(cu_n\). La suma de estos logaritmos diverge exactamente cuando diverge \(\sum u_n\). Aplicando el logaritmo al producto de \eqref{nuc08:eq:7}, se obtiene la equivalencia.
\end{proof}

Se conserva \eqref{nuc08:eq:5} en ambos casos. Una sucesión de avances angulares puede dejar un terminal positivo o transferir toda la acción al registro, según la suma \eqref{nuc08:eq:8}. El cambio de forma por sí solo conserva la acción; la dispersión angular determina su reparto en esta realización. Las fases \(\theta_n\) deben proceder del lector concreto de la trayectoria: la prueba no selecciona retrospectivamente una sucesión para imponer un resultado.

\subsection{Euler extendido y área orientada común}

Los generadores de los tres regímenes cuadráticos son

\[
J_{+1}=(2C_{A_2}+I)/\sqrt3,\qquad
J_0=N,\qquad
J_{-1}=(2F_{\rm av}^2-3I)/\sqrt5.
\]

En sus cartas bidimensionales tienen traza cero y satisfacen
\(J_\tau^2=-\tau I\). El Coxeter procede del antecedente APP, el nilpotente del transporte de umbral y \(F_{\rm av}\) del descenso de modos TPK. Las funciones exponenciales constituyen la realización analítica posterior:

\[
E_\tau(t)=\exp(tJ_\tau)=c_\tau(t)I+s_\tau(t)J_\tau,
\]

donde \((c_\tau,s_\tau)\) son \((\cos,\sin)\), \((1,t)\) y \((\cosh,\sinh)\), respectivamente. La traza cero da determinante uno. Con

\[
\Omega=\begin{pmatrix}0&1\\-1&0\end{pmatrix},
\]

se obtiene \(E_\tau^\mathsf T\Omega E_\tau=\Omega\). En cartas variables se toma
\(C_n=M_{n+1}E_{\tau_n}(t_n)M_n^{-1}\) y el mismo \(B_n\) de \eqref{nuc08:eq:1}.

\begin{theorem}
Los tres regímenes satisfacen

\begin{equation}
\boxed{\Omega=T_n^\mathsf T\Omega T_n+D_n^\mathsf T\Omega D_n,}
\label{nuc08:eq:9}
\end{equation}

con \(T_n,D_n\) de \eqref{nuc08:eq:2}.
\end{theorem}

\begin{proof}
En dimensión dos, \(L^\mathsf T\Omega L=(\det L)\Omega\). Los determinantes de \(M_n\), \(B_n\) y \(C_n\) son uno. La expansión usada en \eqref{nuc08:eq:3}, aplicada a la forma alternante, demuestra \eqref{nuc08:eq:9}.
\end{proof}

En una carta fija, los dos determinantes son

\begin{equation}
\det T_\tau=\frac{65+16c_\tau(t)}{81},\qquad
\det D_\tau=\frac{16(1-c_\tau(t))}{81}.
\label{nuc08:eq:10}
\end{equation}

Esto resulta de \(\det(aI+C)=a^2+a\operatorname{tr}C+\det C\) y \(\operatorname{tr}E_\tau=2c_\tau\). Su suma es uno. En los representantes \(C_{A_2}\), \(I+N\), \(F_{\rm av}^2\), los pares son \((19/27,8/27)\), \((1,0)\), \((89/81,-8/81)\).

El régimen parabólico puede producir una memoria no nula de rango uno con área cero. El hiperbólico conserva área orientada mediante un complemento de signo opuesto. El elíptico permite además el balance positivo de acción entre las formas de las dos fibras. La forma alternante común y las formas métricas de cada régimen mantienen sus tipos respectivos.

\subsubsection{Acción integrada y orientación}

Sea \(\gamma\) una curva cerrada a trozos \(C^1\) en el plano canónico. Para cualquier matriz real \(L\) de orden dos,

\[
\oint_{L\gamma}P\,dQ=(\det L)\oint_\gamma P\,dQ.
\]

En efecto, al escribir \(Q'=aQ+bP\), \(P'=cQ+dP\), las integrales de \(Q\,dQ\) y \(P\,dP\) son cero, y \(\oint Q\,dP=-\oint P\,dQ\). Queda el coeficiente \(ad-bc\). Aplicando \eqref{nuc08:eq:10},

\begin{equation}
\boxed{\oint_\gamma P\,dQ
=\oint_{T_n\gamma}P\,dQ+\oint_{D_n\gamma}P\,dQ.}
\label{nuc08:eq:11}
\end{equation}

Los signos son los de las curvas transportadas. La identidad se conserva para los tres regímenes y para cambios de forma con determinante uno.

\subsection{El factor común del lector exponencial}

El lector radional de \cref{iv:ellipse:section,euler-trit-evaluaciones} contiene también \(e^{-x}\exp(-yJ_\tau)\). Se conserva ese factor cuando actúa en el lector. Sea

\[
C_n=e^{-x_n}M_{n+1}E_{\tau_n}(t_n)M_n^{-1}.
\]

Ahora \(C_n^\mathsf T\Omega C_n=e^{-2x_n}\Omega\). La misma expansión da

\begin{equation}
T_n^\mathsf T\Omega T_n+D_n^\mathsf T\Omega D_n
=\frac{8+e^{-2x_n}}9\Omega.
\label{nuc08:eq:12}
\end{equation}

Para \(x_n\geq0\), definimos un registro complementario de escala

\[
D_{{\rm esc},n}=\frac{\sqrt{1-e^{-2x_n}}}{3}B_n.
\]

Entonces la identidad completa es

\begin{equation}
\Omega=T_n^\mathsf T\Omega T_n+D_n^\mathsf T\Omega D_n
+D_{{\rm esc},n}^\mathsf T\Omega D_{{\rm esc},n}.
\label{nuc08:eq:13}
\end{equation}

Es una composición definida aquí para representar explícitamente el defecto de escala del lector. Su coeficiente procede de \eqref{nuc08:eq:12}. No se atribuye a este registro por su nombre una identidad con otro sector físico del corpus; se conserva el mapa para que esa comparación pueda hacerse con sus operadores y dominios respectivos.

\subsection{Refinamiento de acción sobre las emisiones completas}

Para cada prefijo enriquecido \(h\), sea \(M_h\) su carta y \(G_h=M_h^{-\mathsf T}M_h^{-1}\). Para cada emisión superviviente \(\epsilon\), sea \(R_\epsilon\) una rotación realizada y

\[
U_\epsilon=M_{h\epsilon}R_\epsilon M_h^{-1},\qquad
(\mathscr U z)_{h\epsilon}=\sqrt{p_h(\epsilon)}U_\epsilon z_h.
\]

Los pesos proceden de la medida de emisiones de la construcción común; \(\sum_\epsilon p_h(\epsilon)=1\), con multiplicidades preservadas. De \(U_\epsilon^\mathsf TG_{h\epsilon}U_\epsilon=G_h\) se sigue

\begin{equation}
\sum_{h,\epsilon}I_{h\epsilon}((\mathscr U z)_{h\epsilon})
=\sum_h I_h(z_h).
\label{nuc08:eq:14}
\end{equation}

La prueba suma primero los pesos de los descendientes de un padre. En una trayectoria, los factores \(M_h^{-1}M_h\) intermedios se cancelan, y las rotaciones se componen en el orden registrado. Esto conserva las etiquetas y la memoria sin mezclar hijos de padres distintos. La forma alternante en suma directa satisface la misma identidad de transporte para los tres regímenes simplécticos, sustituyendo \(R_\epsilon\) por \(E_{\tau_\epsilon}(t_\epsilon)\).

Las completaciones métricas se toman con \(\sum_h z_h^\mathsf TG_hz_h\), de modo que estos mapas son isométricos. Identificar además esa norma con una suma directa euclídea fija requiere cotas uniformes de \(M_h\) y \(M_h^{-1}\); esa identificación adicional no se utiliza en \eqref{nuc08:eq:14}.

\section{Reconstrucción explícita desde los retornos completos}
\label{ret:seccion}

\subsection{Dominio de estados y extensión bilateral}

Se considera una carta del dominio superviviente generado por
APP--TRIT--TPK. Un estado de esta carta conserva las hojas operatorias,
la orientación trítica, el residuo, el cociente, el acarreo, la ruta,
la frontera, el registro de memoria y el origen del calendario.
La actualización actúa sobre este estado completo.

Sean \(X_0\) un espacio compacto de Hausdorff no vacío y
\(F:X_0\to X_0\) una actualización continua y sobreyectiva.
Estas son las hipótesis del resultado en la carta considerada.
La extensión bilateral de la actualización es
\begin{equation}
 \mathscr X_F=
 \bigl\{(x_k)_{k\in\ZZ}\in X_0^{\ZZ}:
                  F(x_k)=x_{k+1}\text{ para todo }k\in\ZZ\bigr\},
 \label{ret:extension}
\end{equation}
con la topología inducida por el producto.

\begin{lemma}[Existencia y transporte bilateral]
\label{ret:existencia}
El espacio \(\mathscr X_F\) es compacto, de Hausdorff y no vacío.
Cada proyección coordenada \(\mathscr X_F\to X_0\) es sobreyectiva.
El desplazamiento
\[
 \sigma((x_k)_{k\in\ZZ})=(x_{k+1})_{k\in\ZZ}
\]
es un homeomorfismo de \(\mathscr X_F\).
\end{lemma}
\begin{proof}
Cada condición \(F(x_k)=x_{k+1}\) define un cerrado del producto
compacto de Hausdorff \(X_0^{\ZZ}\).
Fijados un índice \(k_0\) y un estado \(z\in X_0\), cualquier familia
finita de estas condiciones, junto con \(x_{k_0}=z\), admite solución:
se completan hacia delante mediante \(F\) y hacia atrás mediante
preimágenes, que existen por sobreyectividad. Las condiciones poseen,
por tanto, la propiedad de intersección finita. La compacidad produce
una historia bilateral con \(x_{k_0}=z\).

Los desplazamientos \(\sigma\) y
\(\sigma^{-1}((x_k))=(x_{k-1})\) conservan las relaciones de
\eqref{ret:extension}; ambos son continuos por su acción coordenada.
\end{proof}

\subsection{Muestreo trítico y reconstrucción de todos los pasos}

Fijado el origen del calendario, los instantes tipados de retorno son
\[
 \mathcal T_+
 =\{k\in\ZZ:1+(k\bmod 9)\in\{1,4,7\}\}=3\ZZ.
\]
En cada instante se conserva el estado completo \(x_k\).
Definimos el espacio de registros de retorno por
\begin{equation}
 \mathscr Y_+
 =\bigl\{(y_j)_{j\in\ZZ}\in X_0^{\ZZ}:
                 F^3(y_j)=y_{j+1}\text{ para todo }j\in\ZZ\bigr\}.
 \label{ret:registros}
\end{equation}

\begin{theorem}[Recuperación desde los retornos completos]
\label{ret:homeomorfismo}
La aplicación
\[
 r=\operatorname{Ret}_+:\mathscr X_F\longrightarrow\mathscr Y_+,
 \qquad r(x)_j=x_{3j},
\]
es un homeomorfismo. Su inversa está dada explícitamente por
\begin{equation}
 \boxed{
  \bigl(r^{-1}(y)\bigr)_{3j+s}=F^s(y_j),
  \qquad j\in\ZZ,\quad s\in\{0,1,2\}.}
 \label{ret:inversa}
\end{equation}
\end{theorem}
\begin{proof}
Para \(x\in\mathscr X_F\),
\(F^3(x_{3j})=x_{3j+3}\), de modo que \(r(x)\in\mathscr Y_+\).
Recíprocamente, cada entero \(k\) posee una única expresión
\(k=3j+s\), con \(j\in\ZZ\) y \(0\le s<3\), también cuando
\(k<0\). La fórmula \eqref{ret:inversa} define, por tanto, todas las
coordenadas de una sucesión \(x\).

Si \(s=0\) o \(s=1\), se tiene
\(F(x_{3j+s})=F^{s+1}(y_j)=x_{3j+s+1}\).
En la frontera entre dos bloques,
\[
 F(x_{3j+2})=F^3(y_j)=y_{j+1}=x_{3j+3}.
\]
Así, \(x\in\mathscr X_F\).
La lectura de sus coordenadas de índice \(3j\) devuelve \(y_j\).
En sentido inverso, las relaciones de \eqref{ret:extension} dan
\(x_{3j+s}=F^s(x_{3j})\), por lo que la reconstrucción de \(r(x)\)
devuelve \(x\) íntegramente.

La aplicación \(r\) es continua por ser una selección de coordenadas.
Cada coordenada de su inversa es la composición de una proyección
coordenada con una de las aplicaciones continuas \(I,F,F^2\).
La inversa es, por consiguiente, continua.
\end{proof}

Una coordenada discreta adicional, por ejemplo el índice \(m\in\ZZ\)
de una carta de realización, se conserva mediante
\[
 \widetilde r:\ZZ\times\mathscr X_F
 \longrightarrow\ZZ\times\mathscr Y_+,
 \qquad \widetilde r(m,x)=(m,r(x)).
\]
Su inversa es \((m,y)\mapsto(m,r^{-1}(y))\).
El mismo resultado vale para un origen de calendario fijo \(a\):
se sustituyen \(x_{3j}\) y \(x_{3j+s}\) por
\(x_{a+3j}\) y \(x_{a+3j+s}\).

\subsection{Conjugación del transporte y conservación de las lecturas}

Sea \(\sigma_+(y)_j=y_{j+1}\) el desplazamiento de los registros.
La reconstrucción anterior da
\begin{equation}
 r\sigma^3=\sigma_+r,\qquad
 (r\sigma r^{-1}(y))_j=F(y_j).
 \label{ret:conjugacion}
\end{equation}
La aplicación de la segunda igualdad es reversible en
\(\mathscr Y_+\); su inversa es
\[
 y\longmapsto\bigl(F^2(y_{j-1})\bigr)_{j\in\ZZ}.
\]
En efecto, las dos composiciones se reducen a
\(F^3(y_{j-1})=y_j\).
Así, la reversibilidad de la historia bilateral se conserva aun
cuando la actualización \(F\) sobre un estado aislado sea no inyectiva.
Si la vuelta nonádica se realiza como \(\Gamma_9=\sigma^9\), entonces
\begin{equation}
 r\Gamma_9=\sigma_+^3r.
 \label{ret:vuelta-nonadica}
\end{equation}

\begin{corollary}[Lecturas conjuntas y transporte recuperado]
\label{ret:lecturas}
Sean \(a:\mathscr X_F\to\mathcal R\) y
\(q:\mathscr X_F\to\mathcal E\) dos lectores continuos del mismo
estado, y sea \(d:\mathscr X_F\to\mathscr X_F\) un homeomorfismo del
dominio calibrado. Entonces
\[
 a_+=a\circ r^{-1},\qquad
 q_+=q\circ r^{-1},\qquad
 d_+=r\circ d\circ r^{-1}
\]
son continuos, \(d_+\) es un homeomorfismo y
\begin{equation}
 (a,q)=(a_+,q_+)\circ r,\qquad rd=d_+r.
 \label{ret:factorizacion}
\end{equation}
Para todo entero \(N>0\),
\[
 d^Nx=x\quad\Longleftrightarrow\quad d_+^Nr(x)=r(x).
\]
\end{corollary}
\begin{proof}
Se componen las aplicaciones con el homeomorfismo del
\cref{ret:homeomorfismo}. La inversa de \(d_+\) es
\(rd^{-1}r^{-1}\). La identidad
\(d_+^Nr=rd^N\), junto con la inyectividad de \(r\), prueba la
última equivalencia.
\end{proof}

La aplicación a la doble lectura toma
\(\mathcal R=R_{\rm HMT}\) y el codominio incidencial construido en
\cref{exc:doble-lectura,exc:limite-inverso}.
El resultado establece una presentación recuperable de las mismas
historias, de sus lectores y de su transporte. Su objeto es esta
equivalencia dinámica de presentaciones.


\section{Retorno observable y persistencia del registro}
\label{nuclear:9}
\subsection{Recuperación conjunta desde los retornos completos}

Sea \(X\) el dominio de historias completas calibradas de
\cref{ret:extension}, incluido el origen del calendario y, cuando se
utilice, el índice discreto de la carta. Los estados conservan las
transiciones y el registro de memoria. Escribimos
\(r=\operatorname{Ret}_+:X\to Y_+\) para el muestreo en los retornos
completos, con \(Y_+=r(X)\).

El \cref{ret:homeomorfismo} proporciona su inversa explícita:
\(x_{3j+s}=F^s(y_j)\), para \(s=0,1,2\). La reconstrucción utiliza
el estado completo conservado en cada retorno y la actualización
efectiva; \eqref{ret:conjugacion} conserva también el transporte.

Sean \(a:X\to R_{\rm HMT}\) la evaluación numérica por cilindros y
\(q:X\to E_\infty\) la incidencia con su calibre, según
\cref{sec:generacion,exc:doble-lectura,exc:limite-inverso}.
Sea \(d:X\to X\) un transporte reversible y continuo, con el calibre
transportado. Sobre \(Y_+\) se definen



\[
a_+=a\circ r^{-1},\qquad q_+=q\circ r^{-1},\qquad d_+=r\circ d\circ r^{-1}.
\]



Entonces



\begin{equation}
\boxed{(a,q)=(a_+,q_+)\circ r,\qquad rd=d_+r.}
\label{nuc09:lecturas}
\end{equation}



\begin{proof}
Es el \cref{ret:lecturas}: se cancelan las composiciones con \(r\)
y \(r^{-1}\), y la inversa de \(d_+\) es \(rd^{-1}r^{-1}\).
Ambos lectores se transportan desde el mismo estado completo.
\end{proof}

En particular, para \(N\in\ZZ_{>0}\),
\(d^Nx=x\) si y sólo si \(d_+^Nr(x)=r(x)\).
La conjugación conserva la periodicidad o aperiodicidad del estado
completo. La dinámica de un lector parcial se determina mediante
su propio factor de transporte.

\subsection{Representación posterior explícita del transporte}

Sobre el conjunto de historias ya generado se considera la realización
posterior \(H=\ell^2(X)\), con base ortonormal
\((\delta_x)_{x\in X}\) y medida de conteo. Cada vector posee soporte
a lo sumo numerable. El transporte reversible determina



\[
C\delta_x=\delta_{d x},\qquad C^{-1}\delta_x=\delta_{d^{-1}x}.
\]



La biyección \(d\) prueba que \(C\) es unitario.
Esta representación conserva las etiquetas completas de las historias.
El espacio \(X\) y su transporte preceden a esta realización lineal.

La carta nonádica de \cref{nuclear:2} determina



\[
T=\frac{8I+C}{9},\qquad
\eta v=\frac1{27}(8(C-I)v,-(C-I)v,\ldots,-(C-I)v),
\]





\[
I-T^*T=\eta^*\eta=\frac8{81}(I-C)^*(I-C).
\]



Sea \(P\) la proyección ortogonal sobre \(\ker(I-C)\).
La prueba espectral de \cref{nuclear:2,nuclear:3} da
\(T^N\to P\) fuertemente. La telescopía proporciona la isometría



\[
\mathcal Vv=(Pv,\eta v,\eta Tv,\eta T^2v,\ldots),\qquad
\mathcal M=\operatorname{Im}\mathcal V.
\]



El desplazamiento de registros \(B\) conserva el primer componente
y elimina el primer bloque de memoria. Sobre la imagen compatible,



\begin{equation}
\boxed{\widehat C=\mathcal VC\mathcal V^{-1}
=(9B-8I)|_{\mathcal M}.}
\label{nuc09:transporte-recuperado}
\end{equation}



La inversa de la representación está dada por



\begin{equation}
v=Pv+\sum_{j\ge0}T^{*j}\eta^*(\eta T^jv).
\label{nuc09:inversa-registro}
\end{equation}



\begin{proof}
La identidad de defecto telescopa hasta \(N\); su límite conserva
exactamente la norma. Las identidades
\(\mathcal VT=B\mathcal V\) y \(C=9T-8I\) dan
\eqref{nuc09:transporte-recuperado}.
La isometría \(\mathcal V\) es sobreyectiva sobre \(\mathcal M\);
tomar su adjunto da \eqref{nuc09:inversa-registro}, con convergencia
en norma de la serie. La unitariedad del transporte recuperado
se establece en \(\mathcal M\).
\end{proof}

Para \(A\in\mathcal B(H)\), sea
\(\widehat A=\mathcal VA\mathcal V^{-1}\) en \(\mathcal M\).
Esta correspondencia preserva norma, productos, adjuntos y conmutadores,
pues \(\mathcal V^{-1}=\mathcal V^*|_{\mathcal M}\).
Representa conjuntamente los observables numéricos, los de incidencia
y sus relaciones con \(C\).

\subsection{Convergencia numérica y correlaciones con incidencia}

Los cilindros de \eqref{eq:cilindro} proporcionan, en la carta de
índice entero \(m(x)\), seis trits por bloque. Escribimos



\[
\xi_n(x)=\sum_{k=1}^{6n}d_k(x)3^{-k},\qquad
\xi(x)=a(x)-m(x),\qquad 0\le\xi(x)-\xi_n(x)\le3^{-6n}.
\]



Aquí \(d_k\in\{0,1,2\}\) es la lectura ternaria posterior del trit;
su orientación balanceada conserva su tipado propio. La carta
retiene \(m\) por separado, incluida la convención de extremos,
y reconstruye \(a=m+\xi\). En extremos de doble representación,
\(\xi\) se interpreta en esa carta.

Sea \(M_f\) la multiplicación por \(f\) en \(\ell^2(X)\).
La normalización \(0\le\xi_n\le\xi\le1\) hace que
\(M_\xi\) y \(M_{\xi_n}\) sean contracciones.
Para un conjunto \(B\) del codominio incidencial finito de \(q_j\),
definimos
\(E_{j,B}=M_{\mathbf1_{q_j^{-1}(B)}}\).
Es el proyector sobre las historias que satisfacen esa incidencia.

\begin{theorem}[Conservación simultánea de precisión y transporte]
\label{nuc09:precision}
Para todo \(t\in\ZZ\),



\begin{equation}
\boxed{
\left\|\widehat C^{-t}\widehat M_\xi\widehat C^t
-\widehat C^{-t}\widehat M_{\xi_n}\widehat C^t\right\|
=\|M_\xi-M_{\xi_n}\|\le729^{-n}.}
\label{nuc09:error-observable}
\end{equation}



Las relaciones entre los proyectores incidenciales, sus productos
y su transporte se conservan exactamente.
Para un producto ordenado de contracciones que contenga \(r\)
observables numéricos \(M_\xi\), evaluados a tiempos enteros
arbitrarios, y cualquier número finito de factores incidenciales
o unitarios exactos, reemplazar esos \(r\) factores por
\(M_{\xi_n}\) cambia el producto en norma a lo sumo



\begin{equation}
\boxed{r\,729^{-n}.}
\label{nuc09:error-producto}
\end{equation}



La misma cota vale para cualquier elemento de matriz entre vectores
unitarios. Los factores pueden ser no conmutativos; se conserva su orden.
\end{theorem}

\begin{proof}
La norma de un multiplicador sobre \(\ell^2(X)\) es
\(\sup_{x\in X}|f(x)|\); la cola ternaria prueba la primera cota.
Conjugar por \(C^t\) y por \(\mathcal V\) conserva la norma.
Para productos \(A_1\cdots A_s\) y \(B_1\cdots B_s\), se usa



\[
\prod_{i=1}^s A_i-\prod_{i=1}^s B_i
=\sum_{j=1}^s B_1\cdots B_{j-1}(A_j-B_j)A_{j+1}\cdots A_s.
\]



Los factores exactos dan diferencia cero; cada uno de los \(r\)
restantes aporta norma a lo sumo \(729^{-n}\).
La submultiplicatividad y la cota uno de los demás factores prueban
\eqref{nuc09:error-producto}.
Cauchy--Schwarz da la afirmación sobre elementos de matriz.
\end{proof}

Esta uniformidad conserva la precisión al transportar el observable
aproximado completo. El operador \(C\) y la incidencia permanecen
exactos en \eqref{nuc09:error-observable}.
El índice \(n\) mide profundidad de cilindro; \(t\) mide iteraciones
del transporte \(d\); \(j\), en \eqref{nuc09:inversa-registro},
mide iteraciones de compresión.

\subsection{Covariancia del operador de memoria bajo retorno}

Se restringen \(X\) y su representación atómica al sector reversible
graduado \(X_{\rm rev,grad}\) del reloj de
\eqref{eq:weyl-relojes}. Se conserva en cada historia el grado entero
de memoria \(q_{\rm mem}\). Para \(d=\Gamma_9\), su actualización es
\(q_{\rm mem}(dx)=q_{\rm mem}(x)+1\).
El operador de memoria \(D\) es la multiplicación por
\(q_{\rm mem}\), con dominio



\[
\operatorname{Dom}D=\{v:\sum_{x\in X}|q_{\rm mem}(x)|^2|v_x|^2<\infty\}.
\]



El cambio de índice \(x\mapsto dx\) y las desigualdades
\((q\pm1)^2\le2q^2+2\) prueban que \(C\) y \(C^{-1}\)
conservan ese dominio. En él,



\begin{equation}
C^*DC=D+I,\qquad
\widehat C^*\widehat D\widehat C=\widehat D+I_{\mathcal M},
\qquad\operatorname{Dom}\widehat D=\mathcal V\operatorname{Dom}D.
\label{nuc09:memoria}
\end{equation}



\begin{proof}
Sobre \(\delta_x\), el primer miembro vale
\(q_{\rm mem}(dx)\delta_x=(q_{\rm mem}(x)+1)\delta_x\).
La identidad se extiende al dominio indicado por truncación del
soporte en la norma del grafo de \(D\).
La segunda igualdad se obtiene por conjugación unitaria hacia
\(\mathcal M\).
\end{proof}

Por inducción,
\(q_{\rm mem}(d^Nx)=q_{\rm mem}(x)+N\).
Cada entero \(N>0\) cambia el estado completo.
Toda órbita de \(d\) es infinita; una función cuadrado-sumable
constante sobre cada una de esas órbitas es cero.
En esta realización bilateral \(P=0\).
La fórmula \eqref{nuc09:inversa-registro} reconstruye entonces
el estado lineal completo desde sus complementos de memoria.
Esta conclusión utiliza la graduación efectiva del reloj.

El retorno de fase coexiste con el incremento de memoria de
\eqref{nuc09:memoria}.
La aperiodicidad esturmiana y su complejidad pertenecen al factor
simbólico de \cref{sec:generacion}; el presente resultado conserva
la dinámica graduada del estado completo.

\subsection{Incidencia, refinamiento y tres regímenes}

El registro \(K\) posee además la realización completa de
\cref{nuclear:5},
\(\mathcal Fk=(\mu(k),IP_0k/6)\).
Su inversa conserva las doce coordenadas, incluida la media.
Aplicar \(\mathcal F\) a las fibras correspondientes y conservar
las etiquetas de historia transporta esta representación.
La aplicación \(\mathcal F\circ K\) conserva exactamente las fibras
del lector \(K\); la recuperación de historias completas procede
de \(r\) o de sus etiquetas efectivas.

Los sistemas de refinamiento isométrico de
\cref{nuclear:1,nuclear:3,nuclear:5} conservan
\(\mathcal V\) y \(C\) mediante sus cuadrados de entrelazamiento.
La prueba utiliza las realizaciones por nivel allí construidas.
La representación atómica \(\ell^2(X)\) y las realizaciones
\(L^2\) de medida proyectiva conservan sus correspondientes
dominios y medidas.

El enlace con pliegue y despliegue se desarrolla en
\cref{euler-trit-generadores,euler-trit-evaluaciones}:
Coxeter, transporte nilpotente y acción de \(F_{\rm av}^2\).
Para cada realización bidimensional \(E\) que conserva la forma
alternante \(\Omega\), la carta nonádica satisface



\begin{equation}
\Omega=T_E^*\Omega T_E+D_E^*\Omega D_E,
\qquad T_E=(8I+E)/9,\quad D_E=\sqrt8(E-I)/9.
\label{nuc09:balance-alternante}
\end{equation}



La expansión algebraica cancela los términos cruzados y utiliza
\(E^*\Omega E=\Omega\). En la realización real,
\(*\) significa transpuesta.
La ley orientada reúne los tres regímenes mediante las formas
y transportes de \cref{nuclear:7,nuclear:8}.
Su realización como acción utiliza además las normalizaciones
y conexiones allí especificadas.

La relación con la dinámica de conexión se formula mediante
las construcciones de \cref{nuclear:6,nuclear:7,nuclear:8}.
Las identidades
\eqref{nuc09:transporte-recuperado}--\eqref{nuc09:balance-alternante}
conservan su representación, la precisión de los observables
y las formas transportadas.

\begin{HMTCompactSection}
\section{Acoplamiento reversible y transporte no estacionario}
\label{nuclear:10}
\subsection{Extensión reversible del balance entre dos formas}

Sean \(B,C:V\to V'\) isomorfismos que conservan una forma no degenerada \(b\) entre dos realizaciones:



\[
B^*b'B=C^*b'C=b.
\]



La forma puede ser una métrica positiva o una forma alternante. En espacios reales, \(*\) denota transposición respecto de las coordenadas elegidas; la igualdad expresa el transporte de la forma entre las realizaciones especificadas. En espacios de Hilbert, los operadores y sus inversos son acotados.

Con \(s=\sqrt8\) y \(Q_V=\frac13\left(\begin{smallmatrix}sI&-I\\I&sI\end{smallmatrix}\right)\), definimos



\begin{equation}
\boxed{
\mathcal U(B,C)=Q_{V'}^*\operatorname{diag}(B,C)Q_V
=\frac19\begin{pmatrix}8B+C&\sqrt8(C-B)\\\sqrt8(C-B)&B+8C\end{pmatrix}.}
\label{nuc10:eq:1}
\end{equation}



Escribimos \(T=(8B+C)/9\), \(D=\sqrt8(C-B)/9\) y \(R=(B+8C)/9\). La actualización completa es



\begin{equation}
\binom{x'}{m'}=\begin{pmatrix}T&D\\D&R\end{pmatrix}\binom{x}{m}.
\label{nuc10:eq:2}
\end{equation}



\begin{theorem}
La aplicación \eqref{nuc10:eq:1} es reversible y conserva la forma \(b\oplus b\):



\begin{equation}
\mathcal U(B,C)^*(b'\oplus b')\mathcal U(B,C)=b\oplus b,
\quad\mathcal U(B,C)^{-1}=\mathcal U(B^{-1},C^{-1}).
\label{nuc10:eq:3}
\end{equation}
\end{theorem}



\begin{proof}
Se cumple \(Q^*Q=I\) por \(8+1=9\). Sus bloques escalares hacen que \(Q\) preserve \(b\oplus b\) para cualquier forma \(b\). La aplicación diagonal conserva las formas por hipótesis. La composición prueba \eqref{nuc10:eq:3}, y la inversión de los tres factores da la inversa declarada.
\end{proof}

Para memoria entrante \(m=0\), se recupera el balance \(b=T^*b'T+D^*b'D\). La segunda columna determina la respuesta a memoria entrante arbitraria. La conservación de normas positivas y la conservación de área orientada son las especializaciones respectivas de \eqref{nuc10:eq:3}.

\subsubsection{Realización colectiva en las nueve copias}

En \(V^9\), sea \(W_E=\operatorname{diag}(I,\ldots,I,E)\), con \(E\) en la novena posición. Sean \(A(a,b)=(a/\sqrt8,\ldots,a/\sqrt8,b)\) y \(L=AQ\). Entonces \(L(x,0)=Jx\) y



\begin{equation}
W_E L=L\mathcal U(I,E).
\label{nuc10:eq:4}
\end{equation}



\begin{proof}
\(W_EA=A\operatorname{diag}(I,E)\); multiplicar por \(Q\) y usar \(AQ=L\) da \eqref{nuc10:eq:4}. La ortogonalidad de la media de las ocho primeras posiciones con sus siete diferencias muestra además que el complemento queda fijo bajo \(W_E\).
\end{proof}

El paso cíclico original es \(S_E=P_{\rm cyc}W_E\). Por ello \(S_EL=(P_{\rm cyc}L)U(I,E)\), con cartas de entrada y salida distintas. La reducción colectiva conserva el patrón correspondiente: \(S_E^9=I_9\otimes E\), mientras \(W_E^9=\operatorname{diag}(I_8,E^9)\). Cuando un lector utiliza las nueve etiquetas individuales, éstas se conservan en su representación original.

\subsection{Composición exacta y reentrada de memoria}

Para \(B_1,C_1:V_0\to V_1\) y \(B_2,C_2:V_1\to V_2\),



\begin{equation}
\boxed{\mathcal U(B_2,C_2)\mathcal U(B_1,C_1)
=\mathcal U(B_2B_1,C_2C_1).}
\label{nuc10:eq:5}
\end{equation}



\begin{proof}
En \eqref{nuc10:eq:1} se cancelan \(Q_{V_1}Q_{V_1}^*\), y los dos bloques diagonales se multiplican en el orden indicado. La inducción extiende la fórmula a cualquier número finito de transportes, conservando su orden.
\end{proof}

El bloque de lectura contiene la identidad



\begin{equation}
\boxed{T_2T_1+D_2D_1=\frac{8B_2B_1+C_2C_1}{9}.}
\label{nuc10:eq:6}
\end{equation}



Desde \((x,0)\), la primera etapa produce \((T_1x,D_1x)\). La segunda lectura es \(T_2T_1x+D_2D_1x\). El segundo término es exactamente la contribución de la memoria reintroducida. Al almacenar \(D_1x\) fuera del siguiente acoplamiento y continuar sólo con \(T_1x\) se obtiene el procedimiento de registros separados, cuya lectura es \(T_2T_1x\). La igualdad \eqref{nuc10:eq:6} precisa la diferencia entre ambos procedimientos.

La realización \(U\) conserva el estado del acoplamiento y los productos ordenados \(B_N\cdots B_1\), \(C_N\cdots C_1\). La conservación de todos los prefijos incluye además el registro de emisiones de HMT, que mantiene cada factor junto con su orden.

\subsubsection{Ecuación covariante de incremento}

Para el registro separado \(x_{n+1}=T_nx_n\), \(m_n=D_nx_n\), se tiene



\begin{equation}
x_{n+1}-B_nx_n=m_n/\sqrt8.
\label{nuc10:eq:7}
\end{equation}



Con \(G_0=I\) y \(G_{n+1}=B_nG_n\), la telescopía da



\begin{equation}
x_0=G_N^{-1}x_N-\frac1{\sqrt8}\sum_{n<N}G_{n+1}^{-1}m_n.
\label{nuc10:eq:8}
\end{equation}



La memoria es el incremento respecto del transporte entre formas \(B_n\), con la normalización declarada. La igualdad \eqref{nuc10:eq:8} es una reconstrucción finita algebraica; su paso a una serie infinita requiere convergencia. El teorema siguiente construye la inversa estable mediante el límite de operadores positivos.

\subsection{Memoria de una sucesión no estacionaria}

En una realización de Hilbert \(H\), sean \(C_n\) unitarios, con su orden de composición declarado. Escribimos



\[
T_n=(8I+C_n)/9,\quad D_n=\sqrt8(C_n-I)/9,
\quad P_0=I,\quad P_{n+1}=T_nP_n,\quad A_N=P_N^*P_N.
\]



El índice \(N\) cuenta compresiones con almacenamiento separado, a diferencia del transporte completo con reentrada. Para métricas variables con \(B_n\) isomorfismos isométricos, la conjugación por \(G_n\) lleva los transportes a un Hilbert fijo y produce este mismo caso, con \(\widetilde C_n=G_{n+1}^{-1}C_nG_n\). La afirmación utiliza las métricas de las fibras y las isometrías declaradas entre ellas.

\begin{theorem}[Terminal positivo y recuperación no estacionaria]
\label{nuclear:terminal-no-estacionario}
Existe un operador positivo \(0\leq A_\infty\leq I\), límite fuerte de \(A_N\), y



\begin{equation}
\boxed{I=A_\infty+\sum_{n\ge0}P_n^*D_n^*D_nP_n.}
\label{nuc10:eq:9}
\end{equation}



La aplicación



\begin{equation}
\boxed{\mathcal Vv=(A_\infty^{1/2}v,(D_nP_nv)_{n\ge0})}
\label{nuc10:eq:10}
\end{equation}



es isométrica y tiene inversa izquierda estable



\begin{equation}
\mathcal V^*(a,(m_n))=A_\infty^{1/2}a+
\sum_{n\ge0}P_n^*D_n^*m_n,\qquad\|\mathcal V^*\|=1.
\label{nuc10:eq:11}
\end{equation}



\end{theorem}
\begin{proof}
La identidad local \(T_n^*T_n+D_n^*D_n=I\) da
\(A_N-A_{N+1}=P_N^*D_N^*D_NP_N\geq0\). Las formas cuadráticas de
la sucesión positiva decreciente convergen a la de un operador positivo
\(A_\infty\). Para \(B_N=A_N-A_\infty\), la desigualdad
\(0\leq B_N^2\leq B_N\) convierte esa convergencia en fuerte.
La telescopía finita y el límite prueban \eqref{nuc10:eq:9}. Evaluar sobre \(v\)
prueba \eqref{nuc10:eq:10}. Su adjunto es \eqref{nuc10:eq:11}; la serie converge en norma porque
las truncaciones de un registro cuadrado-sumable convergen y el adjunto es acotado.
\end{proof}

La reconstrucción con el terminal límite y \(N\) registros tiene error
exacto \((A_N-A_\infty)v\). El vector \(A_\infty^{1/2}v\) es la
realización positiva canónica del terminal de norma. Su existencia procede
del límite de Gram, con independencia de la convergencia vectorial de \(P_Nv\).

\subsection{Tres propiedades diferentes: agotamiento, distinción y estabilidad}

Sea \(Mv=(D_nP_nv)_n\). De \eqref{nuc10:eq:9}, \(M^*M=I-A_\infty\). Por tanto:

\begin{enumerate}
\item Toda la norma de \(v\) pasa al registro exactamente cuando \(A_\infty v=0\).
\item El registro distingue cada estado exactamente cuando \(\ker(I-A_\infty)=\{0\}\).
\item La recuperación desde memoria sola tiene inversa acotada exactamente cuando
\(I-A_\infty\geq\varepsilon I\) para algún \(\varepsilon>0\).
En ese caso una inversa izquierda es \((I-A_\infty)^{-1}M^*\).
\end{enumerate}

Además,



\begin{equation}
\boxed{\ker M=\bigcap_{n\ge0}\operatorname{Fix}(C_n).}
\label{nuc10:eq:12}
\end{equation}



\begin{proof}
Si \(Mv=0\), entonces \(D_0v=0\), de donde \(C_0v=v\)
y \(P_1v=v\). Inductivamente, \(D_nP_nv=0\) da \(C_nv=v\) y
\(P_{n+1}v=v\). El recíproco es inmediato. Las tres equivalencias se
deducen del Gram \(M^*M\) y de la caracterización de los operadores
acotados inferiormente.
\end{proof}

\subsubsection{Ejemplo exacto que separa las tres propiedades}

En el ejemplo algebraico \(C_0=-I\) y \(C_n=I\) para \(n\geq1\),



\begin{equation}
A_\infty=\frac{49}{81}I,\qquad M^*M=\frac{32}{81}I,
\qquad m_0=-\frac{2\sqrt8}{9}v,
\qquad v=-\frac9{2\sqrt8}m_0.
\label{nuc10:eq:13}
\end{equation}



La memoria contiene \(32/81\) de la norma al cuadrado y permite recuperar el vector completo. El terminal retiene \(49/81\) y es un operador positivo distinto de un proyector. La recuperabilidad depende del núcleo y del condicionamiento del lector. Lectura y memoria constituyen componentes correlacionadas de una isometría.

Este ejemplo separa las tres consecuencias lógicas del teorema. La
aplicación a una trayectoria física conserva además la selección y la
procedencia efectiva de sus \(C_n\) mediante el TPK.

\subsection{Criterio cuantitativo para transferencia completa}

Para \(P_{j,n}=T_{j-1}\cdots T_n\) y \(P_{n,n}=I\), una ventana de longitud \(L\) tiene Gram



\[
G_{n,L}=\sum_{j=n}^{n+L-1}P_{j,n}^*D_j^*D_jP_{j,n}
=I-P_{n+L,n}^*P_{n+L,n}.
\]



Si \(G_{n,L}\geq\kappa I\), con \(\kappa>0\), en cada bloque consecutivo de longitud \(L\), entonces



\begin{equation}
\|P_N\|^2\le(1-\kappa)^{\lfloor N/L\rfloor},\qquad A_\infty=0.
\label{nuc10:eq:14}
\end{equation}



Cada bloque contrae la norma al cuadrado por \(1-\kappa\); la composición y la contractividad de los pasos restantes prueban \eqref{nuc10:eq:14}.

Un criterio suficiente anterior a los productos es



\begin{equation}
\sum_{j=n}^{n+L-1}D_j^*D_j\ge\alpha I
\quad\Longrightarrow\quad
G_{n,L}\ge\frac{\alpha}{[1+2(L-1)/9]^2}I.
\label{nuc10:eq:15}
\end{equation}



\begin{proof}
Para \(y_j=P_{j,n}v\), se tiene
\(y_{j+1}-y_j=D_jy_j/\sqrt8\) y \(\|D_j\|\leq2\sqrt8/9\). Por ello
\[
\|D_jv\|\leq\|D_jy_j\|+\frac29\sum_{i=n}^{j-1}\|D_iy_i\|.
\]
La matriz triangular de sumas parciales tiene norma a lo sumo \(L-1\).
Aplicar la norma \(\ell^2\) de la ventana y la hipótesis de Gram prueba \eqref{nuc10:eq:15}.

\end{proof}

La ausencia de un vector fijo común da distinción, mientras \eqref{nuc10:eq:15} añade una separación uniforme cuantitativa. Ninguna de las dos condiciones se presume por mencionar simetrías: se aplica a los transportes efectivos de la secuencia.

\subsection{Realización incidencial íntegra y cambios de régimen}

La sección~\ref{nuc05:isometria} construye, para el registro de doce
coordenadas, \(Fk=(\mu(k),\mathcal IP_0k/6)\), con inversa y
producto interno ponderado en \(Y=\RR\oplus\mathcal I(\mathbf1^\perp)\).
La transformación \(F\oplus F\) lleva el acoplamiento completo a



\begin{equation}
\mathcal U_Y=(F'\oplus F')\mathcal U(B,C)(F^{-1}\oplus F^{-1})
=\mathcal U(F'BF^{-1},F'CF^{-1}).
\label{nuc10:eq:16}
\end{equation}



Esta identidad conserva la segunda columna, la reentrada y la composición.
Para formas variables se transporta también la forma mediante \(F\).
En la realización real de las doce fibras \(A_2\), se usa
\(F\otimes\operatorname{id}_{V_{A_2}}\), con
\(V_{A_2}=A_2\otimes_{\ZZ}\RR\). Su dominio es el espacio real de
fibras; la preservación del retículo entero corresponde al transporte
reticular ya demostrado. La aplicación \(F\) recupera \(K\) y
conserva las restantes etiquetas genealógicas en el estado completo.

La versión positiva de \eqref{nuc10:eq:9}–\eqref{nuc10:eq:15} utiliza transportes unitarios en las
métricas declaradas. La versión de forma alternante de \eqref{nuc10:eq:1}–\eqref{nuc10:eq:8} admite
los tres regímenes de Euler. Con
\(E=F_{\rm av}^2=\left(\begin{smallmatrix}1&1\\1&2\end{smallmatrix}\right)\),
se obtiene \(\det T=89/81\), \(\det D=-8/81\). Las sumas
orientadas finitas telescopan y los terminales de área crecen como
\((89/81)^N\). El límite positivo de Hilbert utiliza las hipótesis
unitarias específicas del teorema~\ref{nuclear:terminal-no-estacionario}.

El balance reúne formas y dominios, evolución reversible con reentrada
y recuperación ilimitada en las realizaciones positivas. Las
aplicaciones físicas siguientes conservan los lectores, conexiones y
escalas de esa composición operatoria.

\end{HMTCompactSection}
\section{Reconstrucción dimensional y conservación exterior}
\label{nuclear:11}
\subsection{Dos memorias recuperan el registro centrado}

\subsubsection{Operadores de lectura del registro}

En la carta ordenada del registro \(K\) de
\cref{sec:k-registro,img09:imagen}, sea \(V=\RR^{12}\) su
realización lineal posterior y sea \(S\) la permutación
\((Sx)_i=x_{i+1}\), con índices módulo \(12\).
Los lectores \(I-S^3\) e \(I-S^4\) son los del extractor transversal
de \cref{img09:imagen}.
La compresión nonádica de \cref{nuclear:10} proporciona,
para \(j=3,4\),



\begin{equation}
T_j=\frac{8I+S^j}{9},\qquad D_j=\frac{\sqrt8}{9}(S^j-I),
\qquad T_j^*T_j+D_j^*D_j=I.
\label{nuc11:eq:1}
\end{equation}



La partición \(8+1\) y la normalización \(9\) proceden de las nueve
posiciones y de su contraste colectivo.
Se considera la siguiente composición explícita de lectores,
aplicada al registro generado:



\begin{equation}
m_0=D_3K,\qquad m_1=D_4T_3K,
\qquad M K=(m_0,m_1).
\label{nuc11:eq:2}
\end{equation}



La segunda memoria actúa sobre la lectura producida por la primera
etapa. Los complementos se archivan por separado.
La evolución completa con reentrada corresponde al acoplamiento
reversible de \cref{nuclear:10}.

\subsubsection{Inversa exacta}

Puesto que \((S^3)^4=I\),



\begin{equation}
T_3^{-1}=\frac{512I-64S^3+8S^6-S^9}{455}.
\label{nuc11:eq:3}
\end{equation}



En efecto,
\((8I+S^3)(512I-64S^3+8S^6-S^9)=4095I\)
y \(4095=9\cdot455\).
Por tanto, las memorias proporcionan



\begin{equation}
b=-\frac9{\sqrt8}m_0=(I-S^3)K,
\quad c=-\frac9{\sqrt8}T_3^{-1}m_1=(I-S^4)K.
\label{nuc11:eq:4}
\end{equation}



La conmutación de \(T_3\) con \(S^4\) justifica la segunda igualdad.
Definamos



\[
w_i=c_i-b_{i+1},\quad B_0=0,\quad
B_i=\sum_{j<i}w_j.
\]



La identidad \(w_i=K_i-K_{i+1}\) prueba que \(B_i=K_0-K_i\).
Si \(q=\sum_{i=0}^{11}K_i\), se recupera exactamente



\begin{equation}
K_i=\frac{q+\sum_{j=0}^{11}B_j}{12}-B_i.
\label{nuc11:eq:5}
\end{equation}



Es la inversa \eqref{img09:inversa}, compuesta ahora con la
compresión y la memoria. Con \(q=0\), la fórmula devuelve
\(P_{11}K\), donde



\[
P_{11}=I-\frac1{12}\mathbf1\mathbf1^*.
\]



Las dos memorias recuperan las once componentes de media nula.
La carga uniforme se recupera mediante \(q\).
El núcleo de \(M\) es precisamente la recta uniforme:
anular ambas diferencias implica fijación por \(S^3\) y \(S^4\);
como \(S=S^4(S^3)^{-1}\), implica fijación por \(S\).

\subsubsection{Cota óptima de estabilidad}

El Gram del procedimiento es



\begin{equation}
G=M^*M=D_3^*D_3+T_3^*D_4^*D_4T_3
=I-(T_4T_3)^*(T_4T_3).
\label{nuc11:eq:6}
\end{equation}



En el modo \(k\) de \(S\), el valor propio de \(T_j^*T_j\) es
\[
 \frac{65+16\cos(2\pi jk/12)}{81}.
\]
Su sustitución en el Gram anterior da, con multiplicidades,



\begin{equation}
\operatorname{Spec}G=\left\{
0^{[1]},\ (16/81)^{[2]},\ (8/27)^{[2]},\ (32/81)^{[1]},
(952/2187)^{[4]},\ (1256/2187)^{[2]}\right\}.
\label{nuc11:eq:7}
\end{equation}



El modo nulo es la media. Sea \(L\) la inversa de mínimos cuadrados
de \(M\) sobre \(\mathbf1^\perp\), extendida al espacio de
registros por



\[
L=(G|_{\mathbf1^\perp})^{-1}M^*.
\]



Entonces



\begin{equation}
LM=P_{11},\qquad \|L\|=\frac94.
\label{nuc11:eq:8}
\end{equation}



La cota es óptima porque los modos \(k=3,9\) alcanzan \(16/81\).
Para datos perturbados y reconstrucción con \(L\),



\begin{equation}
\|\delta K\|^2\le \frac{|\delta q|^2}{12}
+\frac{81}{16}\bigl(\|\delta m_0\|^2+\|\delta m_1\|^2\bigr).
\label{nuc11:eq:9}
\end{equation}



Las componentes uniforme y centrada son ortogonales. La fórmula \eqref{nuc11:eq:5} constituye una inversa exacta sobre datos compatibles; \eqref{nuc11:eq:8} especifica la extensión estable usada para datos con ruido, que no tienen por qué satisfacer las compatibilidades de \eqref{nuc11:eq:5}.

\subsubsection{Ejemplo con el registro publicado}

El registro de \eqref{eq:k-vector} es



\begin{equation}
K=(234,543,140,729,659,824,621,58,914,794,146,601).
\label{nuc11:eq:10}
\end{equation}



Es una salida del generador anterior, utilizada como caso de
aplicación. Los coeficientes y los lectores se fijan antes de
evaluar este registro. Para almacenar las memorias con aritmética
racional, escribimos \(z_i=m_i/\sqrt8\).
Los dos registros efectivos son



\[
9z_0=(495,116,684,-108,-601,90,173,88,-313,-560,397,-461),
\]





\begin{equation}
81z_1=(2729,2503,3818,-5843,2583,-920,-4051,4338,-5312,-1583,233,1505).
\label{nuc11:eq:11}
\end{equation}



La sustitución de estos registros en las fórmulas de reconstrucción,
con \(q=6263\), reproduce las doce componentes de \(K\).
Con \(q=0\) reproduce \(K-(6263/12)\mathbf1\).
El certificado suplementario reconstruye primero el registro desde
las memorias y después lo compara con \eqref{eq:k-vector}.

Añadir una tercera memoria \(D_3T_4T_3K\) eleva el menor valor
propio positivo del Gram a \(8/27\).
La amplificación óptima desciende de \(9/4\) a \(\sqrt{27/8}\).
Una tercera memoria \(D_4T_4T_3K\) mantiene \(16/81\).
La elección entre repetir el paso \(3\) o el paso \(4\) tiene
un efecto de estabilidad calculable antes de evaluar el registro.

\subsection{Reconstrucción del proyector dimensional desde dos memorias}

\subsubsection{Composición con la incidencia ya construida}

La carta excepcional de \cref{exc:k-direccion} define la
representación ortogonal de \(A_5\) sobre las doce clases de
\(A_5/C_5\), con representantes y carácter explícitos.
Su proyector isótipo tridimensional es
\eqref{eq:k-proyectores-incidencia}:



\begin{equation}
P_3=\frac3{60}\sum_{a\in A_5}\chi_3(a^{-1})\rho(a),
\quad P_3P_{11}=P_3,\quad P_3\mathbf1=0.
\label{nuc11:eq:12}
\end{equation}



El registro generado precede a esta representación.
El \cref{lem:k-sector-no-nulo} y
\eqref{eq:k-norma-direccion} y \eqref{eq:k-proyector-diez} establecen



\[
u_K=P_3P_{11}K,\qquad
\|u_K\|^2=\frac{6638585+2275584\sqrt5}{20}>0,
\]





\begin{equation}
P_{10}(K)=P_{11}-\frac{u_Ku_K^*}{\|u_K\|^2}.
\label{nuc11:eq:13}
\end{equation}



La composición nueva de \eqref{nuc11:eq:8}, \eqref{nuc11:eq:12} y \eqref{nuc11:eq:13} es



\begin{equation}
\boxed{
u_K=P_3L(m_0,m_1),\qquad
P_{10}(K)=P_{11}-
\frac{P_3L(m)\,[P_3L(m)]^*}{\|P_3L(m)\|^2}.
}
\label{nuc11:eq:14}
\end{equation}



\begin{proof}
Las identidades \(LM=P_{11}\) y \(P_3P_{11}=P_3\) dan
\(P_3LMK=P_3K=u_K\).
La no nulidad está evaluada en \eqref{eq:k-norma-direccion}.
El cociente de rango uno coincide, por tanto, con el de
\eqref{eq:k-proyector-diez}.
La composición recupera el proyector ortogonal íntegro de rango
diez a partir de las dos memorias, con independencia de la carga
\(q\) y de la lectura terminal \(T_4T_3K\).
\end{proof}

La reducción \(12\to11\) elimina la dirección uniforme y la
reducción \(11\to10\) utiliza el vector seleccionado por \(K\).
La composición recupera la matriz completa del proyector en el
portador lineal de incidencia.
Sus realizaciones dimensionales mantienen los mapas de
\cref{prop:k-reduccion-diez}.

\subsubsection{Error del proyector}

Sea \(\varepsilon\) la norma conjunta del error de las dos
memorias y sea \(\widehat u=P_3L(m+\delta m)\).
Entonces \(\|\widehat u-u_K\|\le9\varepsilon/4\).
Si \(9\varepsilon/4<\|u_K\|\), el vector \(\widehat u\) es no
nulo y



\begin{equation}
\boxed{
\|\widehat P_{10}-P_{10}\|
\le\frac{9\varepsilon}{4\|u_K\|}.
}
\label{nuc11:eq:15}
\end{equation}



\begin{proof}
La distancia en norma operatoria entre dos proyectores ortogonales
sobre rectas es el seno del ángulo entre ellas.
La distancia de \(u_K\) a la recta de \(\widehat u\) es a lo sumo
\(\|u_K-\widehat u\|\). Dividir por \(\|u_K\|\) prueba la cota.
Aquí \(\|u_K\|\approx765{,}733\), y el coeficiente es
aproximadamente \(0{,}00293836\).
\end{proof}
Para registros próximos a \(\ker P_3\), la estabilidad de la
dirección depende de \(\|u_K\|\).
La cota \(9/4\) para el registro centrado permanece válida.

\subsubsection{Un descriptor escalar de K, con invariancias explícitas}

Se distinguen el registro \(K\), sus lecturas escalares y sus
invariantes. Para \(K\notin\RR\mathbf1\), definimos



\begin{equation}
\eta_K=\frac{\|P_3P_{11}K\|^2}{\|P_{11}K\|^2}.
\label{nuc11:eq:16}
\end{equation}



Para el registro \eqref{eq:k-vector},
\(\sum_iK_i^2=4251257\) y
\(\|P_{11}K\|^2=11789915/12\). Así,



\begin{equation}
\eta_K=
\frac{3(6638585+2275584\sqrt5)}{5\cdot11789915}
=\frac{13089003+13653504\varphi_{\rm HMT}}{58949575}
\approx0,5967954228259091.
\label{nuc11:eq:17}
\end{equation}



La segunda escritura utiliza el reconocimiento cuadrático posterior
de \(\varphi_{\rm HMT}\).
Se cumple \(0\le\eta_K\le1\).
El cociente es invariante bajo
\(K\mapsto aK+b\mathbf1\), con \(a\ne0\), y bajo cambios
ortogonales simultáneos de \(K,P_{11},P_3\).
En la carta fija es invariante bajo \(\rho(A_5)\), pues ambos
proyectores conmutan con esa representación.
Es recuperable desde las memorias porque ambas normas que lo
definen lo son.
Mide la fracción cuadrática del registro centrado situada en el
sector isótipo elegido.
El registro completo conserva además la información que este
descriptor escalar identifica.
La procedencia de esta formalización adicional se documenta en
el suplemento; las invariancias afines aquí probadas pertenecen
a la realización lineal.

\subsection{Aplicación sobre un vector efectivo de la construcción excepcional}

La construcción de
\cref{km:orden-tres-reticular,km:estabilidad-vecino}
utiliza doce fibras \(A_2\), su pegado y el vecino reticular.
La palabra trítica es
\(c^\star=(1,1,1,1,1,1,2,1,1,1,1,1)\).
Sus bloques actúan por \(C^{-1}\) cuando la componente es \(1\)
y por \(C\) cuando es \(2\), en las cartas especificadas por
\cref{km:palabra-total,km:isometria-ternaria}.
El operador \(g\) resultante cumple
\(g^3=I\), \(g^2+g+I=0\), y es ortogonal para la suma de las
métricas \(A_2\).
El descenso al pegado y al vecino se demuestra en
\cref{km:estabilidad-vecino}.
Los cálculos siguientes actúan sobre su portador real de
dimensión \(24\).

En cada fibra,
\[
 C=\begin{pmatrix}0&-1\\1&-1\end{pmatrix},
 \qquad
 G_{A_2}=\begin{pmatrix}2&-1\\-1&2\end{pmatrix}.
\]
Sobre la suma de las doce fibras se define



\[
T=\frac{8I+g}{9},\qquad D=\frac{\sqrt8}{9}(g-I).
\]



Con el adjunto relativo a esa métrica,
\(g^*=g^{-1}=-g-I\), y se obtiene



\begin{equation}
T^*T=\frac{19}{27}I,
\qquad D^*D=\frac8{27}I,
\qquad D^{-1}=-\frac3{\sqrt8}(g+2I).
\label{nuc11:eq:18}
\end{equation}



La inversa procede de \((g-I)(g+2I)=-3I\).
La construcción del vecino fija el vector de
\eqref{exc:vector},
\(v=4\rho_1+\rho_2+\cdots+\rho_{12}\), con una raíz
\(\rho_b\) de norma cuadrática \(2\) en cada fibra.
Por tanto,



\begin{equation}
\|v\|^2=54,\qquad \|Tv\|^2=38,\qquad \|Dv\|^2=16,
\qquad v=-\frac3{\sqrt8}(g+2I)Dv.
\label{nuc11:eq:19}
\end{equation}



Un solo complemento recupera este vector completo.
La lectura conserva \(38\) unidades de norma cuadrática y el
registro contiene \(16\); su suma es \(54\).
El resultado utiliza el vector de \eqref{exc:vector} y sus fibras
reticulares.

\subsection{Ley de conservación en cada grado exterior}

\subsubsection{Enunciado y prueba}

Sean \(P_0=I\), \(P_n=T_{n-1}\cdots T_0\) y



\[
\mathcal J_Nv=(P_Nv,D_0P_0v,\ldots,D_{N-1}P_{N-1}v).
\]



La ley de \cref{nuclear:8,nuclear:10} conserva la forma \(G_0\)
mediante la suma ortogonal de las formas correspondientes:
\(\mathcal J_N^*G_{\rm sal}\mathcal J_N=G_0\).
Para cada grado exterior finito admisible \(p\), se deduce



\begin{equation}
\boxed{
(\Lambda^p\mathcal J_N)^*(\Lambda^pG_{\rm sal})
(\Lambda^p\mathcal J_N)=\Lambda^pG_0.
}
\label{nuc11:eq:20}
\end{equation}



La forma inducida se define por



\[
\langle x_1\wedge\cdots\wedge x_p,
y_1\wedge\cdots\wedge y_p\rangle_{\Lambda^pG}
=\det(\langle x_i,y_j\rangle_G)_{i,j=1}^p.
\]



\begin{proof}
Sustituir el balance de grado uno en cada entrada de ese determinante. La igualdad se cumple sobre vectores descomponibles y se extiende bilinealmente porque éstos generan la potencia exterior.
\end{proof}

La descomposición relevante es



\[
\Lambda^p\Bigl(\bigoplus_aH_a\Bigr)
\simeq\bigoplus_{\sum n_a=p}\bigotimes_a\Lambda^{n_a}H_a.
\]



En el caso positivo, si
\(\mathcal J_{p,\mathbf n}\) designa las componentes de ese mapa,



\begin{equation}
\boxed{\|\xi\|^2=\sum_{\sum n_a=p}
\|\mathcal J_{p,\mathbf n}\xi\|^2.}
\label{nuc11:eq:21}
\end{equation}



Se conservan las componentes puramente visibles, las puramente
registradas y las mixtas.
Para \(p=2\) y \(p=3\), se conservan, respectivamente, las normas
cuadráticas de área y volumen.
Para una forma alternante o indefinida, la identidad exterior
conserva su interpretación bilineal.

\subsubsection{Distribución exacta entre los sectores exteriores}

En la realización Coxeter anterior,
\(a=19/27\) y \(b=8/27\).
Normalizar \(T\) y \(D\) da dos isometrías, de modo que la
componente con \(k\) direcciones en la memoria tiene norma cuadrática



\begin{equation}
\boxed{
\|\mathcal J_{p,k}\xi\|^2=
\binom pk a^{p-k}b^k\|\xi\|^2.
}
\label{nuc11:eq:22}
\end{equation}



Se prueba expandiendo exteriormente
\(x\mapsto(\sqrt a\,x,\sqrt b\,x)\).
Las componentes con posiciones distintas de memoria son
ortogonales; las isometrías normalizadas \(T\) y \(D\)
conservan sus normas.

Para un bivector de norma uno,

\begin{center}
\begin{tabular}{@{}lr@{}}
\toprule
Sector & Norma cuadrática\\
\midrule
Dos direcciones en la lectura & \(361/729\)\\
Una dirección en lectura y otra en memoria & \(304/729\)\\
Dos direcciones en memoria & \(64/729\)\\
\bottomrule
\end{tabular}
\end{center}

La suma es uno. El sector mixto contiene exactamente \(304/729\),
aproximadamente \(41{,}7\%\), de la identidad de área al cuadrado.
En grado tres, las cuatro componentes son
\((6859,8664,3648,512)/19683\); las componentes mixtas suman
\(12312/19683\).

La ley se aplica a \(0\le p\le24\) en el portador de las doce
fibras \(A_2\), y functorialmente a otros portadores generados.
El grado exterior, la resolución de refinamiento, la escala de
una métrica y la dimensión física poseen tipados distintos.
El cambio \(G\mapsto s^2G\), con \(s>0\), multiplica la norma
cuadrática de grado \(p\) por \(s^{2p}\); ambas partes de la
identidad exterior se transforman de igual manera.
La realización de las dimensiones \(10\) y \(11\) conserva los
mapas dimensionales previamente especificados.

\subsubsection{Naturalidad y profundidad arbitraria}

El refinamiento de historias generado conserva la medida mediante



\[
R\delta_h=\sum_{\epsilon\in\operatorname{Out}(h)}
\sqrt{p_h(\epsilon)}\,\delta_{h\epsilon}.
\]



Las probabilidades proceden de las multiplicidades y de la
normalización de las emisiones APP--TRIT--TPK de
\cref{iv:cont:medida,nuclear:1}.
El registro identifica el estado precedente de cada emisión.
Si el cuadrado de transporte es
\(\mathcal J'_NR=(\bigoplus_aR_a)\mathcal J_N\), entonces



\begin{equation}
(\Lambda^pJ'_N)(\Lambda^pR)
=\Lambda^p(\bigoplus R_a)(\Lambda^pJ_N).
\label{nuc11:eq:23}
\end{equation}



Se aplica
\(\Lambda^p(AB)=\Lambda^p(A)\Lambda^p(B)\).
Los cuadrados de las cinco operaciones correlativas y su límite
común son los de
\cref{iv:rectangular-natural,iv:cont:memoria-natural,iv:cont:mapa-correlativo}; todas actúan en la misma construcción
del continuo.

En el caso positivo no estacionario,
\cref{nuclear:10} demuestra
\(P_N^*P_N\downarrow A_\infty\) fuertemente y



\[
\mathcal Jv=(A_\infty^{1/2}v,(D_nP_nv)_{n\ge0}),
\qquad \mathcal J^*\mathcal J=I.
\]



Por densidad, \(\Lambda^p\mathcal J\) es isométrico para cada
grado finito \(p\).
La suma exterior se extiende a multiíndices de soporte finito
sobre la suma de Hilbert numerable; Parseval justifica el paso
al límite.
El terminal \(A_\infty^{1/2}\) y las memorias conservan
conjuntamente la forma inicial a profundidad arbitraria.

\subsubsection{Precisión de reconstrucción de áreas y volúmenes}

Para el lector \(M\), ordenemos los once valores propios
positivos de \(G\), con multiplicidad, como
\(\lambda_1\le\cdots\le\lambda_{11}\).
En los grados \(1\le p\le11\) de \(\mathbf1^\perp\),
la descomposición singular induce



\begin{equation}
\|(\Lambda^pM)^\dagger\|
=(\lambda_1\cdots\lambda_p)^{-1/2}.
\label{nuc11:eq:24}
\end{equation}



Los modos exteriores utilizan \(p\) direcciones distintas.
El espectro calculado da las constantes óptimas
\(9/4\), \(81/16\), \(243\sqrt6/64\) y \(2187/128\)
para \(p=1,2,3,4\).
En grado doce, \(\Lambda^{12}M=0\); la recuperación del volumen
completo incorpora también la carga uniforme.
El producto tensorial de \(p\) lectores independientes sobre
media nula tiene cota \((9/4)^p\).
Tensorizar el lector por una identidad conserva la cota \(9/4\).
El registro completo \(\mathcal J\) tiene inversa de norma uno
sobre su imagen en cada grado.

\subsection{Memoria del transporte y holonomía relativa}

El acoplamiento reversible de \cref{nuclear:10} es



\[
U(B,C)=\frac19
\begin{pmatrix}8B+C&\sqrt8(C-B)\\
\sqrt8(C-B)&B+8C\end{pmatrix}
=Q^*\operatorname{diag}(B,C)Q,
\]



con \(Q\) constante y ortogonal, fijada por la partición \(8+1\).
Por tanto,
\(U(B_2,C_2)U(B_1,C_1)=U(B_2B_1,C_2C_1)\),
conservando el orden de factores.
Sobre un lazo \(\gamma\), sean \(H_B,H_C\) los productos
ordenados de ambos transportes.
El bloque de memoria de la evolución completa es



\[
D_\gamma=\frac{\sqrt8}{9}(H_C-H_B).
\]



Se reconstruye exactamente



\begin{equation}
\boxed{H_B^{-1}H_C=I+\frac9{\sqrt8}H_B^{-1}D_\gamma.}
\label{nuc11:eq:25}
\end{equation}



La igualdad es operatorial: registrar la respuesta del bloque
sobre una base recupera la holonomía relativa completa.
Su evaluación sobre un vector recupera la acción correspondiente.
Bajo cambio de carta en el punto base, ambos miembros se conjugan.
Si \(H_B=I\), el complemento determina \(H_C-I\).
Así se relacionan memoria y defecto de transporte sobre un
circuito; la realización diferencial utiliza adicionalmente
su conexión y el límite de circuitos.

\subsubsection{Un lazo elíptico–parabólico produce autoescala hiperbólica}

La construcción de \cref{euler-trit-generadores} contiene el
Coxeter \(C\) y el nilpotente
\[
 N_0=\begin{pmatrix}0&1\\0&0\end{pmatrix}
\]
en sus portadores respectivos.
Se consideran sus realizaciones en un plano orientado común,
mediante las bases indicadas y la forma
\[
 \Omega=\begin{pmatrix}0&1\\-1&0\end{pmatrix}.
\]
El Coxeter \(C\) y el transporte \(I+N_0\) tienen determinante
uno y conservan \(\Omega\).
Esta carta permite componerlos conservando la distinción entre
los generadores de sus regímenes y sus métricas positivas propias.

Sea
\[
 P=I+N_0=\begin{pmatrix}1&1\\0&1\end{pmatrix}.
\]
Como \(N_0^2=0\), se cumple \(P^n=I+nN_0\) para todo
\(n\in\ZZ\).
El lazo ordenado de cuatro transportes da



\begin{equation}
H_n=C^{-1}P^{-n}CP^n
=\begin{pmatrix}1+n&n^2\\n&n^2-n+1\end{pmatrix}.
\label{nuc11:eq:28}
\end{equation}



Para \(n\ne0\), su diferencia normalizada es



\[
F_n=\frac{H_n-I}{n}
=\begin{pmatrix}1&n\\1&n-1\end{pmatrix}.
\]



Multiplicando esas matrices se obtiene, sin introducir valores espectrales,



\begin{equation}
\boxed{F_n^2=H_n,\qquad F_n^2-nF_n-I=0.}
\label{nuc11:eq:29}
\end{equation}



Con \(H_B=I\), la memoria es
\(D_{\gamma,n}=(\sqrt8/9)nF_n\).
En consecuencia,



\begin{equation}
\boxed{H_n=\left(\frac9{n\sqrt8}D_{\gamma,n}\right)^2.}
\label{nuc11:eq:30}
\end{equation}



La respuesta matricial de memoria reconstruye la autoescala
hiperbólica. Para \(n>0\), el reconocimiento espectral posterior
identifica



\begin{equation}
\rho_n=\frac{n+\sqrt{n^2+4}}2,\quad
\operatorname{Spec}F_n=\{\rho_n,-\rho_n^{-1}\},\quad
\operatorname{Spec}H_n=\{\rho_n^2,\rho_n^{-2}\}.
\label{nuc11:eq:31}
\end{equation}



El caso \(n=1\) produce el polinomio áureo.
La relación con el operador
\(F_{\rm av}=\begin{pmatrix}0&1\\1&1\end{pmatrix}\)
generado por el calendario en \eqref{eq:fav} conserva el cambio
de carta:



\begin{equation}
F_1=PF_{\rm av}P^{-1},\qquad
H_1=PF_{\rm av}^{\,2}P^{-1}.
\label{nuc11:eq:32}
\end{equation}



El transporte \(P\), de determinante uno, entrelaza ambas
matrices y conserva \(\Omega\).
El caso \(n=2\) produce las lecturas espectrales
\(1+\sqrt2\) y \(3+2\sqrt2\).
La composición relaciona estas realizaciones conservando sus
genealogías respectivas; la selección de una palabra de cuatro
factores permanece especificada como parte del lector.

En el mismo ejemplo \(n=1\), se obtiene el balance orientado



\begin{equation}
D_\gamma^{\mathsf T}\Omega D_\gamma=-\frac8{81}\Omega,
\qquad
T_\gamma^{\mathsf T}\Omega T_\gamma=\frac{89}{81}\Omega.
\label{nuc11:eq:33}
\end{equation}



La suma de ambas contribuciones es \(\Omega\).
El término de memoria tiene orientación opuesta:
la contribución visible \(89/81\) se compensa con \(-8/81\).
Este balance alternante y el reparto positivo
\(19/27+8/27=1\) conservan sus formas respectivas.

La expresión de \(H_n\) vale para todo \(n\in\ZZ\).
La definición de \(F_n\) y la reconstrucción mediante su cuadrado
utilizan \(n\ne0\); la parametrización espectral anterior utiliza
\(n>0\).
En general,
\[
 D_{\gamma,n}^{\mathsf T}\Omega D_{\gamma,n}
 =-\frac{8n^2}{81}\Omega,\qquad
 T_{\gamma,n}^{\mathsf T}\Omega T_{\gamma,n}
 =\left(1+\frac{8n^2}{81}\right)\Omega.
\]
Las identidades \(N_0^2=0\) y \(C^2+C+I=0\) demuestran la
familia completa; el programa suplementario comprueba sus
especializaciones exactas.
La procedencia de esta formalización adicional se documenta
en el suplemento.

\subsection{Simetría, acción y carga conservada}

Sean \((E_n)_{n\in\ZZ}\) espacios vectoriales reales de dimensión
finita y \(U_n:E_n\to E_{n+1}\) los transportes invertibles de
la evolución completa.
La elevación cotangente actúa sobre
\(q_n\in E_n\) y \(p_n\in E_n^*\).
En bases duales, definimos el lagrangiano discreto auxiliar
\[
 L_n(q_n,q_{n+1},p_{n+1})
 =p_{n+1}^{\mathsf T}(q_{n+1}-U_nq_n).
\]
La acción se utiliza en sentido local:
para cada variación de soporte finito, se suma \(L_n\) sobre un
intervalo finito que contenga todos los términos afectados.
La primera variación es independiente de cualquier ampliación
de ese intervalo. En concreto,
\begin{equation}
 \delta\mathscr S
 =\sum_{n\in\ZZ}
 \left[
  \delta p_n^{\mathsf T}(q_n-U_{n-1}q_{n-1})
  +\delta q_n^{\mathsf T}(p_n-U_n^{\mathsf T}p_{n+1})
 \right],
 \label{nuc11:accion-local}
\end{equation}
donde la suma es finita para cada variación admitida.
La estacionariedad respecto de todas estas variaciones equivale a
\[
 q_{n+1}=U_nq_n,\qquad
 p_n=U_n^{\mathsf T}p_{n+1}
 \quad(n\in\ZZ).
\]
La segunda ecuación es
\(p_{n+1}=U_n^{-\mathsf T}p_n\), es decir, la elevación
cotangente del transporte.

\begin{proposition}[Carga variacional del transporte entrelazado]
\label{nuc11:noether}
Sea \(A_n\in\operatorname{End}(E_n)\) una familia que satisface
\[
 A_{n+1}U_n=U_nA_n.
\]
Sobre las soluciones estacionarias de
\eqref{nuc11:accion-local}, la carga
\begin{equation}
 \boxed{J_n=p_n^{\mathsf T}A_nq_n=J_{n+1}}
 \label{nuc11:carga}
\end{equation}
es constante.
\end{proposition}
\begin{proof}
Las ecuaciones de evolución y el entrelazamiento dan
\[
 J_{n+1}
 =p_{n+1}^{\mathsf T}A_{n+1}U_nq_n
 =p_{n+1}^{\mathsf T}U_nA_nq_n
 =p_n^{\mathsf T}A_nq_n.
\]
La conservación posee además una derivación variacional explícita.
Para parámetro constante \(\varepsilon\), las transformaciones
\[
 q_n\mapsto e^{\varepsilon A_n}q_n,\qquad
 p_n\mapsto e^{-\varepsilon A_n^{\mathsf T}}p_n
\]
conservan cada \(L_n\), porque
\(e^{\varepsilon A_{n+1}}U_n=U_ne^{\varepsilon A_n}\).
Si el parámetro se sustituye por una sucesión
\((\varepsilon_n)\) de soporte finito, la primera variación es
\[
 \delta L_n
 =(\varepsilon_{n+1}-\varepsilon_n)
   p_{n+1}^{\mathsf T}U_nA_nq_n.
\]
Sobre una solución, el último factor es \(J_n\).
La suma finita por partes da
\[
 \delta\mathscr S
 =\sum_{n\in\ZZ}\varepsilon_n(J_{n-1}-J_n).
\]
La estacionariedad para toda sucesión de soporte finito produce
\eqref{nuc11:carga}.
\end{proof}

Para los lectores cíclicos anteriores se toma
\[
 A_0=S-S^{-1},\qquad
 \mathcal A=\operatorname{diag}(A_0,A_0)
\]
en el portador completo \(V\oplus V\).
La matriz \(\mathcal A\) conmuta con \(U(I,S^3)\) y
\(U(I,S^4)\), pues sus bloques son polinomios en \(S\).
En consecuencia, \(p_n^{\mathsf T}\mathcal A q_n\) conserva
una correlación orientada entre posiciones vecinas.

La construcción especifica un lagrangiano, su simetría y la carga
conservada correspondiente, en el sentido variacional desarrollado
en \cref{nuclear:7}.
Su dominio es la realización cotangente de los lectores.
La acción física y sus unidades conservan los operadores y
normalizaciones de las construcciones correspondientes.
El lagrangiano auxiliar se obtiene después del transporte \(U_n\).

\section{Realización canónica de la covarianza}
\label{nuclear:covarianza-antecedente}
% Propuesta de inserción después de 07b_geometria_elipse.tex.
% RESULTADO_RECUPERADO: c52_06_01_incertidumbre.tex, líneas 10--91.
% Dependencias: c52_04_01_ley_areal.tex, líneas 137--198;
% c52_06_02_intercambio_constitutivo.tex, líneas 5--29.
% Explicitación local: dominio operatorio, prueba de covarianza,
% transporte recíproco y distinción dimensional posición--momento.
% Este fragmento no introduce una constante ni un generador adicional.

\subsection{Covarianza canónica y reciprocidad de la elipse de acción}
\label{sec:ii-area-covarianza}

La construcción precedente determina dos invariantes complementarios:
el producto de los semiejes fija el área de acción y su cociente conserva
la anisotropía del par angular. Precisamos ahora la realización canónica
en la que esos semiejes son dos desviaciones estándar. La estructura de
partida es la elipse ya obtenida de las publicaciones angulares y de
acción de APP--TRIT--TPK; la realización operatoria se establece después
de esa construcción.

En este apartado escribimos \(\eta=\eta_{\rm el}
=\operatorname{artanh}(C^*/A)\), con \(|C^*/A|<1\) y
\(\hbar>0\). El cociente angular utiliza una misma unidad para sus dos
componentes. En las coordenadas equilibradas \((Q,P)\), ambas de
dimensión raíz de acción, definimos
\begin{equation}
 V_\eta:=\frac14
 \begin{pmatrix}a_S^2&0\\0&b_S^2\end{pmatrix}
 =\frac{\hbar}{2}
 \begin{pmatrix}e^\eta&0\\0&e^{-\eta}\end{pmatrix}.
 \label{eq:ii-covarianza-geometrica}
\end{equation}
Esta matriz positiva tiene entradas de dimensión acción. Su interpretación
como covarianza cuántica queda especificada por la realización siguiente.

\begin{proposition}[Realización gaussiana de la covarianza de acción]
\label{prop:ii-covarianza-gaussiana}
Considérese la representación canónica en \(L^2(\mathbb R,dQ)\), con
\(\widehat Qf(Q)=Qf(Q)\) y
\(\widehat Pf(Q)=-i\hbar f'(Q)\), inicialmente sobre el espacio de
Schwartz \(\mathcal S(\mathbb R)\). El estado normalizado
\begin{equation}
 \psi_\eta(Q)=\bigl(\pi\hbar e^\eta\bigr)^{-1/4}
 \exp\!\left(-\frac{Q^2}{2\hbar e^\eta}\right)
 \label{eq:ii-gaussiana-accion}
\end{equation}
tiene medias nulas y matriz de covarianza \(V_\eta\). En particular,
\begin{equation}
 (\Delta Q)^2=\frac{\hbar e^\eta}{2},\qquad
 (\Delta P)^2=\frac{\hbar e^{-\eta}}{2},\qquad
 \operatorname{Cov}(Q,P)=0,\qquad
 \det V_\eta=\frac{\hbar^2}{4}.
 \label{eq:ii-varianzas-accion}
\end{equation}
Esta realización satura la desigualdad de Robertson--Schrödinger.
\end{proposition}

\begin{proof}
Los operadores anteriores son esencialmente autoadjuntos sobre
\(\mathcal S(\mathbb R)\), que constituye un núcleo común invariante;
sobre este espacio se cumple
\([\widehat Q,\widehat P]=i\hbar I\). Para
\(\widehat Y=(\widehat Q,\widehat P)\), la covarianza simétrica es
\[
 V_{jk}=\frac12\left\langle
 (\widehat Y_j-\langle\widehat Y_j\rangle)
 (\widehat Y_k-\langle\widehat Y_k\rangle)
 +(\widehat Y_k-\langle\widehat Y_k\rangle)
 (\widehat Y_j-\langle\widehat Y_j\rangle)
 \right\rangle.
\]
Con \(s=\hbar e^\eta>0\), sea
\(I_s=\int_{\mathbb R}e^{-Q^2/s}\,dQ\). El producto de dos
integrales iguales, evaluado en coordenadas polares, da
\(I_s^2=2\pi\int_0^\infty r e^{-r^2/s}\,dr=\pi s\).
Además, la integral de
\(\frac{d}{dQ}(Qe^{-Q^2/s})\) es cero, por el decaimiento en ambos
extremos; de ahí
\(\int_{\mathbb R}Q^2e^{-Q^2/s}\,dQ=sI_s/2\).
Estas identidades y la paridad prueban la normalización, la media nula y
\(\langle\widehat Q^2\rangle=s/2\). Además,
\(\widehat P\psi_\eta=i e^{-\eta}Q\psi_\eta\).
Por tanto, \(\langle\widehat P\rangle=0\),
\(\|\widehat P\psi_\eta\|^2=\hbar e^{-\eta}/2\) y
\(\operatorname{Re}\langle\widehat Q\psi_\eta,
\widehat P\psi_\eta\rangle=0\).

Para cualquier estado normalizado de ese dominio, Cauchy--Schwarz aplicada
a \((\widehat Q-\langle\widehat Q\rangle)\psi\) y
\((\widehat P-\langle\widehat P\rangle)\psi\) da
\[
 (\Delta Q)^2(\Delta P)^2
 \geq \operatorname{Cov}(Q,P)^2
       +\frac14\left|\langle[\widehat Q,\widehat P]\rangle\right|^2
 =\operatorname{Cov}(Q,P)^2+\frac{\hbar^2}{4}.
\]
Las varianzas calculadas alcanzan la igualdad. Equivalentemente, el
operador
\[
 a_\eta=\frac{e^{-\eta/2}\widehat Q
                  +i e^{\eta/2}\widehat P}{\sqrt{2\hbar}}
\]
satisface \([a_\eta,a_\eta^\dagger]=I\) y
\(a_\eta\psi_\eta=0\), por sustitución directa.
\end{proof}

\begin{proposition}[Área, dispersión y transporte recíproco]
\label{prop:ii-area-covarianza-reciprocidad}
La elipse de acción es exactamente el subnivel de covarianza
\begin{equation}
 \mathcal E_\eta
 =\left\{y=(Q,P)^{\mathsf T}:y^{\mathsf T}V_\eta^{-1}y\leq4\right\}
 =\left\{\frac{Q^2}{a_S^2}+\frac{P^2}{b_S^2}\leq1\right\}.
 \label{eq:ii-elipse-covarianza}
\end{equation}
Su área es \(h=2\pi\hbar\), sus semiejes son
\(a_S=2\Delta Q\), \(b_S=2\Delta P\) y
\begin{equation}
 \frac{\Delta P}{\Delta Q}=e^{-\eta}
 =\frac{b_S}{a_S}=R_{\rm el}^2.
 \label{eq:ii-covarianza-forma}
\end{equation}
La rotación canónica \(J_2\) de la sección~\ref{iv:ellipse:section} transporta
\(V_\eta\) a \(V_{-\eta}\) y conserva la forma simpléctica.
\end{proposition}

\begin{proof}
La inversión de la matriz diagonal
\eqref{eq:ii-covarianza-geometrica} da
\eqref{eq:ii-elipse-covarianza}. Sus semiejes son las raíces positivas de
cuatro veces sus varianzas, y
\[
 \operatorname{Area}(\mathcal E_\eta)
 =4\pi\sqrt{\det V_\eta}=2\pi\hbar=h,
\]
en concordancia con la ley areal~\eqref{iv:ellipse:area-fase}.
La razón de las raíces positivas prueba
\eqref{eq:ii-covarianza-forma}.
Si \(M_\eta=\operatorname{diag}(e^{\eta/2},e^{-\eta/2})\), entonces
\[
 V_\eta=M_\eta V_0M_\eta^{\mathsf T},\qquad
 M_\eta^{\mathsf T}J_2M_\eta=J_2,\qquad
 J_2V_\eta J_2^{\mathsf T}=V_{-\eta}.
\]
Estas identidades se obtienen por multiplicación. En términos operatorios,
\((\widehat Q',\widehat P')=(-\widehat P,\widehat Q)\) conserva
\([\widehat Q',\widehat P']=i\hbar I\). La permutación
\(S_2:(Q,P)\mapsto(P,Q)\) también intercambia las varianzas, pero invierte
el signo del conmutador y de la forma simpléctica. Así, el transporte de
covarianza respeta la distinción entre reciprocidad canónica e inversión
de orientación ya establecida en la sección~\ref{iv:ellipse:section}.
\end{proof}

\paragraph{Posición, momento y transducción dimensional.}
Sea \(\lambda_{xp}>0\) un factor de la carta dimensional, de dimensión
momento por unidad de longitud. La transformación
\[
 Q=\sqrt{\lambda_{xp}}\,x,\qquad
 P=\frac{p_x}{\sqrt{\lambda_{xp}}}
\]
realiza el par equilibrado en una posición \(x\) y su momento conjugado
\(p_x\); conserva \(dQ\wedge dP=dx\wedge dp_x\) y
\([\widehat x,\widehat p_x]=i\hbar I\). Por transformación de la
covarianza,
\begin{equation}
 V_\eta^{(x,p_x)}=\frac{\hbar}{2}
 \begin{pmatrix}e^\eta/\lambda_{xp}&0\\
                 0&\lambda_{xp}e^{-\eta}\end{pmatrix},\qquad
 \Delta x\,\Delta p_x=\frac{\hbar}{2},\qquad
 \frac{\Delta p_x}{\Delta x}=\lambda_{xp}e^{-\eta}.
 \label{eq:ii-covarianza-dimensiones}
\end{equation}
El cociente de dispersión adquiere aquí la dimensión momento por longitud.
En cambio, \eqref{eq:ii-covarianza-forma} pertenece a la carta equilibrada.
Para la realización inductivo--capacitiva ya definida,
\eqref{hmtII:eq:ii-elipse-impedancia} aporta la identidad adimensional
\(Z_\eta/Z_*=e^{-\eta}=\Delta P/\Delta Q\).
Cada lectura utiliza su transducción declarada y conserva el mismo
determinante de acción.

\paragraph{Energía del contorno y energía media del estado.}
La realización inductivo--capacitiva tiene el Hamiltoniano
\[
 H_\eta(Q,P)=\frac{\omega}{2}
 \left(e^{-\eta}Q^2+e^\eta P^2\right),\qquad \omega>0.
\]
Sobre \(\partial\mathcal E_\eta\), su valor es
\(H_\eta=\hbar\omega\), por las ecuaciones de los semiejes. En la
realización cuántica de la proposición~\ref{prop:ii-covarianza-gaussiana},
\begin{equation}
 \langle\widehat H_\eta\rangle_{\psi_\eta}
 =\frac{\omega}{2}
 \left(e^{-\eta}\frac{\hbar e^\eta}{2}
       +e^\eta\frac{\hbar e^{-\eta}}{2}\right)
 =\frac{\hbar\omega}{2}.
 \label{eq:ii-energia-contorno-covarianza}
\end{equation}
El contorno de dos desviaciones estándar tiene, por tanto, el doble de la
energía media de este estado gaussiano. La ley areal fija la normalización
de acción; el par angular fija su distribución entre las dos coordenadas;
la representación canónica y el estado explícito realizan esa distribución
como covarianza. La saturación de incertidumbre queda así acreditada en
la realización especificada por sus operadores y su estado.

\section[Retorno de hoja, incidencia areal--volumétrica y orientación constitutiva]{Retorno de hoja, incidencia areal--volumétrica\\y orientación constitutiva}
\label{iii:sec:hojas}
\label{nuclear:constitutiva-antecedente}

La estructura de hojas interviene antes de la evaluación de las magnitudes
electromagnéticas. Su función consiste en distinguir el cierre de una proyección,
la restitución de una orientación y la evolución del registro transportado.
Estas operaciones actúan sobre objetos relacionados, pero de distinto tipo: una
trayectoria proyectada puede haberse cerrado mientras su levantamiento permanece
en la hoja opuesta; la orientación puede haberse restituido mientras la memoria
conserva las dos vueltas efectuadas. La construcción constitutiva debe recibir
esta información con la estructura que la determina.

\subsection{Retorno proyectado y restitución de la cubierta orientada}

Sea $\ell$ un lazo de retorno de una extensión superviviente y sea
$\varepsilon:\langle\ell\rangle\to\{+1,-1\}$ el carácter de su
levantamiento, con $\varepsilon(\ell)=-1$. El fibrado de intervalos asociado
tiene la presentación
\begin{equation}
 \mathcal M=([0,1]\times[-1,1])/{(0,u)\sim(1,-u)}.
 \label{iii:eq:mobius}
\end{equation}
El dato decisivo es el signo del transporte sobre el lazo especificado. La
identificación de los extremos realiza geométricamente ese carácter.

\begin{proposition}[Dos órdenes de clausura]
Una vuelta del lazo cierra la base e invierte la orientación de la fibra. Dos
vueltas restituyen la orientación. En la parametrización angular del lector de
borde, estos retornos corresponden, respectivamente, al cierre areal $2\pi$ y
al cierre volumétrico $4\pi$ de la cubierta orientada.
\end{proposition}
\begin{proof}
La relación de equivalencia de \eqref{iii:eq:mobius} identifica la fibra terminal
con la inicial mediante $u\mapsto-u$. La composición de dos transportes es
$u\mapsto-(-u)=u$. Parametrizar una vuelta de la base por $2\pi$ asigna $4\pi$
a su cuadrado. La afirmación se refiere a este levantamiento y a esta lectura.
\end{proof}

La memoria del TPK conserva, además, la trayectoria recorrida y las actualizaciones
del registro. El cuadrado del carácter de signo es la identidad, pero no impone
que se anulen esas actualizaciones. La representación espinorial, cuando se
utilice, deberá transportar este dato mediante su aplicación propia.

\subsection{Descenso de la incidencia desde el calendario trítico}

En el bloque primitivo $(+1,-1,0)$, el TPK actualiza el modo aritmético mediante
$+1:m\mapsto\Sigma$, $-1:m\mapsto\Pi$ y $0:m\mapsto m$. Su órbita cíclica
es $(\Sigma,\Pi,\Pi)$. El lector de hojas asigna $\Sigma$ a la hoja areal y
$\Pi$ a la volumétrica. En el orden de bases $(e_{\rm ar},e_{\rm vol})$, la
incidencia de aristas queda representada por
\begin{equation}
 F_{\rm av}=\begin{pmatrix}0&1\\1&1\end{pmatrix}.
 \label{iii:eq:fav}
\end{equation}
Los tres pares consecutivos son $\Sigma\to\Pi$, $\Pi\to\Pi$ y
$\Pi\to\Sigma$, cada uno con multiplicidad uno. Los calendarios de longitudes
$9$, $27$ y $108$ contienen $3$, $9$ y $36$ copias del bloque primitivo. Sus
matrices de incidencia son, por tanto, esos múltiplos de $F_{\rm av}$.

Esta construcción conserva la diferencia entre frecuencia cronológica y medida
del espacio de continuaciones. La primera asigna pesos $(1/3,2/3)$ al ciclo
de modos. La segunda se define sobre la envolvente simbólica mínima de memoria
uno que admite exactamente los pares anteriores. Su matriz de adyacencia es
$F_{\rm av}$; los recuentos de trayectorias pertenecen a esa envolvente.

\Needspace{10\baselineskip}
\begin{proposition}[Pesos de retorno de las dos hojas]
Sean $F_0=0$, $F_1=1$ y $F_{n+1}=F_n+F_{n-1}$. Para $n\geq1$,
\begin{equation}
 F_{\rm av}^{n}=
 \begin{pmatrix}F_{n-1}&F_n\\F_n&F_{n+1}\end{pmatrix}.
 \label{iii:eq:fav-powers}
\end{equation}
La medida de Parry de esta envolvente asigna a las hojas pesos proporcionales
a $(1,\varphi^2)$, y la razón de recuentos diagonales converge a $\varphi^{-2}$.
\end{proposition}
\begin{proof}
La fórmula de potencias se prueba por inducción utilizando la recurrencia.
El cuadrado de $F_{\rm av}$ es estrictamente positivo. Su vector positivo de
Perron es $(1,\varphi)$, con autovalor $\varphi$, pues
$\varphi^2=\varphi+1$. Al ser simétrica la matriz, los pesos de Parry son
proporcionales a los cuadrados de las componentes de ese vector. La otra raíz
del polinomio característico es $-\varphi^{-1}$; su módulo es estrictamente
menor que $\varphi$, lo que da el límite de los cocientes diagonales.
\end{proof}

La identificación espectral de esta incidencia es posterior a su construcción
desde el calendario. La coordenada regional $\varphi$ conserva su genealogía
coinductiva anterior: el reconocimiento del autovalor no selecciona el calendario.

\subsection{Marcado de la frontera y razón areal--volumétrica}

Sea $\Pi_{\rm HMT}(X_\infty)$ el valor del lector regional de clausura ya
construido. Con los proyectores diagonales de hoja, se define
\begin{equation}
 D_\pi=\Pi_{\rm HMT}(X_\infty)P_{\rm ar}+P_{\rm vol},\qquad
 \Theta_n=
 \frac{\langle e_{\rm ar},D_\pi F_{\rm av}^{n}e_{\rm ar}\rangle}
 {\langle e_{\rm vol},F_{\rm av}^{n}e_{\rm vol}\rangle}.
 \label{iii:eq:marked-return}
\end{equation}
La normalización volumétrica es unitaria; el cierre regional marca una sola vez
el retorno areal. Aplicando \eqref{iii:eq:fav-powers},
\begin{equation}
 \Theta_n=\Pi_{\rm HMT}(X_\infty)\frac{F_{n-1}}{F_{n+1}}
 \longrightarrow\frac{\pi}{\varphi^2}.
 \label{iii:eq:av-limit}
\end{equation}
La razón de retorno queda determinada por la incidencia, la medida y el marcado.
El cociente $2\pi/4\pi$ describe, en cambio, los dos órdenes de clausura de la
cubierta. Las dos relaciones participan de la lectura de hojas sin sustituirse.

\subsection{Transporte de la orientación a los canales constitutivos}

En una representación de dos canales, sea $\mathcal R$ una involución
autoadjunta y $\mathcal S$ un operador unitario de intercambio que satisface
$\mathcal S\mathcal R\mathcal S^{-1}=-\mathcal R$. Las coordenadas angulares
ya construidas se escriben $x=A_{\rm rad}$ e $y=C^*_{\rm rad}$, con $x>|y|$.
Entonces
\begin{equation}
 B=xI+y\mathcal R,\qquad T=e^{-B},\qquad
 \mathcal ST(y)\mathcal S^{-1}=T(-y).
 \label{iii:eq:channel-conjugation}
\end{equation}
En efecto, la conjugación transforma $B(y)$ en $B(-y)$ y conmuta con su serie
exponencial convergente. Los autovalores $q_\pm=e^{-x\mp y}$ se intercambian.
Para fijar el orden de las respuestas se utiliza la cámara $x>y>0$, donde
$0<q_+<q_-<1$. La aplicación $(x,y)\mapsto(x,-y)$ la transporta a la
cámara opuesta $x>-y>0$, donde $0<q_-<q_+<1$. Ambas pertenecen al dominio
contractivo $x>|y|$. En la frontera $y=0$ los dos canales coinciden.
La conjugación intercambia las cámaras, no conserva el orden estricto de
sus etiquetas; sobre los proyectores actúa mediante
$\mathcal SP_\pm\mathcal S^{-1}=P_\mp$, con
$P_\pm=(I\pm\mathcal R)/2$.
Esta aplicación concreta transmite la orientación al operador constitutivo que
se desarrolla a continuación.

% Ampliación propia de III; no modifica el núcleo común copiado de I.
% Residencia propuesta: antes de Retorno de hoja, incidencia areal--volumétrica.
% Procedencia y controles: VACANCIAS_PREFACTOR_PROCEDENCIA.md, en esta carpeta.
\section{Polarización discreta y modo trítico de las proyecciones aritméticas}
\label{iii:sec:vacancias-prefactor}

El refinamiento y la comparación entre las hojas de APP proporcionan dos
observaciones complementarias. La primera sigue la discrepancia acumulada
entre el número de transiciones y la capacidad adquirida; la segunda conserva
el modo del selector trítico en el emparejamiento de las fibras aditivas y
multiplicativas. La polarización discreta y el prefactor electrónico se
obtienen de estas operaciones respectivas. Su articulación con el vacío
requiere mantener tanto la información temporal como los espacios sobre los
que actúan los operadores.

\subsection{Discrepancia del refinamiento y convención de extremos}

Los cardinales del bloque ternario y de su representación decimal fijan
\[
 \rho=\log_{1000}729=\log_{10}9,\qquad
 \delta=1-\rho=\log_{10}(10/9).
\]
Se tiene $0<\delta<1$. Además, $\delta$ es irracional: una igualdad
$\delta=p/q$, con $q>0$, implicaría $10^{q-p}=9^q$, incompatible con
la descomposición en factores primos. Este desajuste ya interviene en la
codificación del refinamiento. Aquí se explicita su lectura centrada,
conservando el origen de fase y la convención de extremos.

\begin{definition}[Codificaciones inferior y superior]
Para $\theta\in\mathbb R$ y $n\in\mathbb Z$, se definen
\begin{align}
 z_n^{\downarrow}(\theta)
 &=\lfloor(n+1)\delta+\theta\rfloor-\lfloor n\delta+\theta\rfloor,
 \label{iii:vp:lower}\\
 z_n^{\uparrow}(\theta)
 &=\lceil(n+1)\delta+\theta\rceil-\lceil n\delta+\theta\rceil.
 \label{iii:vp:upper}
\end{align}
Ambas toman valores en $\{0,1\}$. La flecha distingue exclusivamente la
convención de extremos de los intervalos; la orientación trítica conserva
su definición anterior.
\end{definition}

La codificación inferior coincide con la utilizada en el desarrollo
del refinamiento. La superior permite recuperar la formulación de la
polarización con función techo. Las dos expresiones se relacionan por
un término de frontera exactamente determinado.

\begin{proposition}[Cambio de convención y término de frontera]
Sea $b_n(\theta)=\mathbf1_{\mathbb Z}(n\delta+\theta)$. Entonces
\begin{equation}
 z_n^{\uparrow}(\theta)-z_n^{\downarrow}(\theta)
 =b_n(\theta)-b_{n+1}(\theta).
 \label{iii:vp:endpoint}
\end{equation}
En particular, para $N\geq0$,
\[
 \sum_{n=0}^{N-1}(z_n^{\uparrow}-z_n^{\downarrow})
 =b_0-b_N.
\]
\end{proposition}
\begin{proof}
Para todo real $x$, $\lceil x\rceil=\lfloor x\rfloor+1-
\mathbf1_{\mathbb Z}(x)$. Aplicar esta identidad a los dos extremos
de cada incremento da \eqref{iii:vp:endpoint}; su suma es telescópica.
La irracionalidad de $\delta$ implica que, fijada $\theta$, como máximo
un índice entero satisface $n\delta+\theta\in\mathbb Z$.
\end{proof}

Con $\theta=0$, el primer símbolo superior es $1$ y el inferior es $0$;
para $n\geq1$ coinciden. El término inicial pertenece al registro y se
conserva cuando se cambia la convención. Las densidades y las cotas de
discrepancia se transportan con ese término, sin modificar la actualización
del TPK ni reiniciar la historia en cada retorno de fase.

\subsection{Polarización acotada y balance temporal exacto}

\begin{definition}[Polarización discreta centrada]
La polarización de la codificación superior, respecto de la densidad
$\delta$, es
\begin{equation}
 P_N^{\uparrow}(\theta)
 =\sum_{n=0}^{N-1}\bigl(z_n^{\uparrow}(\theta)-\delta\bigr),
 \qquad P_0^{\uparrow}=0.
 \label{iii:vp:polarization}
\end{equation}
Su incremento temporal es
\begin{equation}
 J_N^{\uparrow}(\theta)
 =P_{N+1}^{\uparrow}(\theta)-P_N^{\uparrow}(\theta)
 =z_N^{\uparrow}(\theta)-\delta.
 \label{iii:vp:current}
\end{equation}
\end{definition}

El centrado separa el crecimiento uniforme de densidad $\delta$ de la
información de fase. Así, $P_N^{\uparrow}$ mide un desequilibrio acumulado,
no el número total de vacancias. El incremento $J_N^{\uparrow}$ conserva
el pulso local y el fondo respecto del cual se evalúa.

\begin{theorem}[Cota uniforme y continuidad discreta temporal]
Para $N\geq0$ y toda fase $\theta$,
\begin{equation}
 P_N^{\uparrow}(\theta)
 =\lceil N\delta+\theta\rceil-\lceil\theta\rceil-N\delta,
 \qquad |P_N^{\uparrow}(\theta)|<1.
 \label{iii:vp:bound}
\end{equation}
Para cualesquiera enteros $0\leq a<b$,
\begin{equation}
 \sum_{n=a}^{b-1}J_n^{\uparrow}(\theta)
 =P_b^{\uparrow}(\theta)-P_a^{\uparrow}(\theta),
 \qquad
 \left|\sum_{n=a}^{b-1}J_n^{\uparrow}(\theta)\right|<1.
 \label{iii:vp:continuity}
\end{equation}
La densidad de la sucesión de vacancias es $\delta$.
\end{theorem}
\begin{proof}
La suma de los incrementos de techo de \eqref{iii:vp:upper} produce la
primera identidad. Escribiendo $r(x)=\lceil x\rceil-x\in[0,1)$,
se obtiene $P_N^{\uparrow}=r(N\delta+\theta)-r(\theta)$; su valor
absoluto es menor que uno. Sobre el intervalo $[a,b)$ se obtiene
análogamente $r(b\delta+\theta)-r(a\delta+\theta)$, lo que prueba
la segunda cota, más precisa que la suma de dos cotas individuales.
Finalmente,
$N^{-1}\sum_{n=0}^{N-1}z_n^{\uparrow}
=\delta+N^{-1}P_N^{\uparrow}\to\delta$.
\end{proof}

La codificación inferior admite las mismas demostraciones sustituyendo
$r(x)$ por $\lfloor x\rfloor-x\in(-1,0]$. En particular,
$P_N^{\uparrow}-P_N^{\downarrow}=b_0-b_N$. La discrepancia acotada
utilizada anteriormente en las lecturas de Dirichlet y la polarización
presente conservan, por tanto, una relación explícita de origen y extremos.
La aperiodicidad tampoco introduce una deriva secular de este desequilibrio:
la historia puede prolongarse indefinidamente mientras la discrepancia
centrada permanece uniformemente acotada.

\subsection{Realización en las rectas de carga y tiempo}

Para realizar dimensionalmente el balance se especifican una recta de carga
$\mathcal L_Q$ con sección no nula $u_Q$ y una sucesión creciente de
instantes $t_N$ en una recta temporal orientada. Con
$\Delta t_N=t_{N+1}-t_N>0$, la carga centrada y la corriente media de
cada intervalo quedan definidas por
\begin{equation}
 Q_N=P_N^{\uparrow}u_Q,\qquad
 I_N=\frac{Q_{N+1}-Q_N}{\Delta t_N}
     =\frac{u_Q}{\Delta t_N}\bigl(z_N^{\uparrow}-\delta\bigr).
 \label{iii:vp:dimensional}
\end{equation}
Así, $Q_N\in\mathcal L_Q$ e
$I_N\in\mathcal L_Q\otimes\mathcal L_T^{-1}$. La sección $u_Q$ y
los intervalos temporales son datos de esta realización; la construcción
aritmética precedente ya ha determinado $P_N^{\uparrow}$ y
$J_N^{\uparrow}$.

\begin{proposition}[Conservación del balance bajo realización dimensional]
Para $a<b$,
\begin{equation}
 Q_b-Q_a=\sum_{n=a}^{b-1}\Delta t_n I_n,
 \qquad
 \left|\frac{Q_b-Q_a}{u_Q}\right|<1.
 \label{iii:vp:dimensional-balance}
\end{equation}
Si todos los intervalos tienen longitud $t_0$, las dos corrientes posibles
son $-\delta u_Q/t_0$ y $(1-\delta)u_Q/t_0$.
\end{proposition}
\begin{proof}
Multiplicar cada identidad de \eqref{iii:vp:dimensional} por
$\Delta t_n$ y sumar telescopa las cargas. La cota es
\eqref{iii:vp:continuity}; los dos valores locales se obtienen de
$z_n^{\uparrow}\in\{0,1\}$.
\end{proof}

La continuidad demostrada es temporal y está referida a la carga centrada.
Una realización espacial incorpora además la localización de las cargas,
la incidencia entre celdas y la identificación de los flujos de frontera.
En una descripción constitutiva por canales, el acoplamiento se especifica
mediante una aplicación de excitación
$\iota:\mathcal L_Q\to\mathcal H_{\rm ch}\otimes\mathcal L_Q$
y una aplicación de observación
$\ell:\mathcal H_{\rm ch}\otimes\mathcal L_Q\to\mathcal L_Q$,
donde $\mathcal H_{\rm ch}$ es una realización real del espacio de canales.
Si ambas son
lineales e independientes del índice temporal, la respuesta de un operador
$\mathcal V$ fijo se transporta por
\[
 Q_N^{\rm resp}=\ell(\mathcal V\otimes I_{\mathcal L_Q})\iota(Q_N),
 \qquad
 I_N^{\rm resp}=\frac{Q_{N+1}^{\rm resp}-Q_N^{\rm resp}}{\Delta t_N}.
\]
La linealidad conserva el balance telescópico. Esta composición hace
explícitos los datos de una identificación constitutiva: la secuencia
escalar determina el pulso, mientras $\iota$, $\mathcal V$ y $\ell$
determinan su excitación, transporte y observación por canales. Las
identidades \eqref{iii:vp:polarization}--\eqref{iii:vp:dimensional-balance}
establecen la primera construcción con independencia de esa elección de
realización. El operador del vacío se construye en la
sección~\ref{iii:sec:constitutiva} sobre sus canales y proyectores propios.

\subsection{Emparejamiento de las fibras aditivas y multiplicativas}

La segunda observación procede de la realización tridimensional de APP,
$\mathcal A_3=D_9^3$, con medida uniforme y
$D_9=\{1,\ldots,9\}$. Se conservan los observables
\[
 \Sigma(x)=\rho_9(x_1+x_2+x_3),\qquad
 \Pi(x)=\rho_9(x_1x_2x_3),
\]
donde $\rho_9(n)=1+[(n-1)\bmod9]$ para $n>0$. En
$\ell^2(\mathcal A_3)$, sean $P_{\Sigma,3}$ y $P_{\Pi,3}$ las
esperanzas condicionales sobre sus fibras. Por ejemplo,
\[
 (P_{\Sigma,3}f)(x)=
 \frac1{|\Sigma^{-1}(\Sigma(x))|}
 \sum_{y:\,\Sigma(y)=\Sigma(x)}f(y).
\]
Son proyecciones ortogonales definidas sobre el mismo soporte. La matriz
de incidencias $N_{ab}=|\Sigma^{-1}(a)\cap\Pi^{-1}(b)|$ resulta de
enumerar los $729$ triples, sumando y multiplicando sus coordenadas antes
de aplicar $\rho_9$:
\begin{equation}
 N=9\begin{pmatrix}
 0&1&1&0&1&1&0&1&4\\
 1&0&1&1&0&1&1&0&4\\
 1&0&2&0&1&2&0&0&3\\
 0&1&1&0&1&1&0&1&4\\
 1&0&1&1&0&1&1&0&4\\
 0&0&2&1&0&2&0&1&3\\
 0&1&1&0&1&1&0&1&4\\
 1&0&1&1&0&1&1&0&4\\
 0&1&2&0&0&2&1&0&3
 \end{pmatrix}.
 \label{iii:vp:joint-matrix}
\end{equation}
Las filas suman $81$ y las sumas de columna son
\[
 (m_1,\ldots,m_9)=(36,36,108,36,36,108,36,36,297).
\]
Por ello, el emparejamiento de las indicatrices normalizadas de las fibras
es $C_{ab}=N_{ab}/\sqrt{81m_b}$. La normalización recoge las
multiplicidades producidas por APP, no una elección de valores singulares.

\begin{proposition}[Modo singular del selector trítico]
En la carta canónica, el vector del selector es
$m_{\rm ph}=(1,-1,0)^3$. Se cumple
\begin{equation}
 Cm_{\rm ph}=C^{\mathsf T}m_{\rm ph}=-\tfrac12m_{\rm ph}.
 \label{iii:vp:trit-mode}
\end{equation}
Así, el valor singular correspondiente es $s=1/2$.
\end{proposition}
\begin{proof}
En las seis columnas donde $m_{\rm ph}$ es no nulo, $m_b=36$ y
$C_{ab}=N_{ab}/54$. Multiplicar las filas de
\eqref{iii:vp:joint-matrix} por $(1,-1,0)^3$ produce
$-m_{\rm ph}/2$. Para el producto transpuesto, las sumas firmadas
de las columnas $3,6,9$ son cero; las otras seis dan los mismos
valores $-m_{{\rm ph},b}/2$. Esto prueba ambas identidades y, por
composición, $CC^{\mathsf T}m_{\rm ph}=m_{\rm ph}/4$.
\end{proof}

\subsection{Norma del conmutador y prefactor electrónico}

\begin{theorem}[Norma exacta de incompatibilidad aritmética]
Las dos proyecciones satisfacen
\begin{equation}
 \bigl\|[P_{\Sigma,3},P_{\Pi,3}]\bigr\|=\frac{\sqrt3}{4}.
 \label{iii:vp:prefactor}
\end{equation}
El máximo se alcanza en el plano principal correspondiente al modo
trítico de \eqref{iii:vp:trit-mode}.
\end{theorem}
\begin{proof}
La matriz racional $M=CC^{\mathsf T}$ tiene entradas
$M_{ac}=\sum_bN_{ab}N_{cb}/(81m_b)$. Su espectro completo se
comprueba mediante los siguientes subespacios independientes: los
vectores $(1,\ldots,1)$, $(1,-1,0)^3$ y $(1,1,-2)^3$ tienen
autovalores $1$, $1/4$ y $7/132$, respectivamente; el subespacio de
vectores soportados en $\{3,6,9\}$ cuya suma es cero tiene dimensión
$2$ y autovalor $1/18$; los vectores de suma cero soportados en
$\{1,4,7\}$ o en $\{2,5,8\}$ forman un núcleo de dimensión $4$.
La sustitución en la matriz verifica estas acciones y las dimensiones
suman $9$. En consecuencia,
\[
 \operatorname{spec}M=
 \{1,1/4,7/132,1/18,1/18,0,0,0,0\}.
\]
Para dos proyecciones ortogonales, un valor singular $s\in(0,1)$ del
emparejamiento determina un plano principal en el que sus matrices son
\[
 P=\begin{pmatrix}1&0\\0&0\end{pmatrix},\qquad
 Q=\begin{pmatrix}s^2&s\sqrt{1-s^2}\\
 s\sqrt{1-s^2}&1-s^2\end{pmatrix}.
\]
Esta forma se obtiene tomando un vector unitario $u$ de la primera
imagen con $PQu=s^2u$ y completándolo con
$(Qu-s^2u)/(s\sqrt{1-s^2})$. Su conmutador es
$s\sqrt{1-s^2}\bigl(\begin{smallmatrix}0&1\\-1&0\end{smallmatrix}\bigr)$.
Los sectores de valor singular $0$ o $1$ contribuyen con cero.
Por tanto, la norma total es el máximo de
$\sqrt{\lambda(1-\lambda)}$ sobre el espectro de $M$. Dicho máximo
es $\sqrt{3/16}$, alcanzado en $\lambda=1/4$.
\end{proof}

La norma es el prefactor que interviene en la composición electrónica.
Su obtención conserva el soporte tridimensional, las fibras y el modo
del selector; la evaluación escalar no sustituye esa estructura. En
particular, conocer $\sqrt3/4$ permite conocer la magnitud máxima de
la incompatibilidad, mientras los proyectores conservan las direcciones
en que actúa.

\clearpage
\subsection{Defecto cuadrático y respuesta de los sectores 90/120}

El modo trítico determina simultáneamente
\begin{equation}
 1-s^2=\frac34=\frac{90}{120},\qquad s=\frac12.
 \label{iii:vp:defect-ratio}
\end{equation}
La primera expresión es el defecto cuadrático del plano principal de APP.
La última conserva las longitudes de los dos sectores del calendario.
La respuesta constitutiva utiliza esos sectores sobre otro objeto: un
canal contractivo $q\in(0,1)$. Su evaluación es
\[
 r(q)=\frac{1-q^{90}}{1-q^{120}}.
\]
La relación entre el cociente de sectores y esta respuesta puede
precisarse sin identificar sus operadores.

\begin{proposition}[Límite del cociente de respuesta]
Para $0<q<1$,
\begin{equation}
 \frac34<r(q)<1,\qquad
 \lim_{q\uparrow1}r(q)=\frac34=1-s^2.
 \label{iii:vp:response-limit}
\end{equation}
\end{proposition}
\begin{proof}
Con $u=q^{30}\in(0,1)$, la factorización de $1-u^3$ y
$1-u^4$ da
$r(q)=(1+u+u^2)/(1+u+u^2+u^3)$. El límite en $u=1$ es
$3/4$. Además,
$4(1+u+u^2)-3(1+u+u^2+u^3)
=(1-u)(3u^2+2u+1)>0$, y el denominador supera al numerador
en $u^3>0$. Se obtienen las dos desigualdades.
\end{proof}

La coincidencia exacta concierne al defecto del modo singular y al límite
de la respuesta de los sectores. En el interior contractivo, la respuesta
depende de $q$ y es estrictamente mayor que $3/4$. El desarrollo
constitutivo conserva sus dos canales $q_\pm$ y sus proyectores
$P_\pm$; los proyectores $P_{\Sigma,3}$ y $P_{\Pi,3}$ conservan,
por su parte, las fibras aritméticas que determinaron el prefactor.
Una equivalencia operatoria entre ambas realizaciones exige una aplicación
entre sus espacios que entrelace las acciones correspondientes.
Las ecuaciones \eqref{iii:vp:trit-mode}, \eqref{iii:vp:prefactor} y
\eqref{iii:vp:response-limit} precisan la información compartida y
mantienen visibles los datos de cada realización.

\section{Operador constitutivo y reconstrucción de los canales del vacío}
\label{iii:sec:constitutiva}

Las magnitudes constitutivas se obtienen mediante una operación sobre el
transporte angular. La información de partida comprende los canales, sus
proyectores y la orientación que los distingue. Esta organización permite
seguir por separado la construcción del operador, la evaluación de sus
funciones espectrales y su realización dimensional posterior.

\subsection{Respuesta positiva sobre dos hojas}

En la cámara orientada $x>y>0$ de la sección~\ref{iii:sec:hojas}, sean
$P_\pm=(I\pm\mathcal R)/2$ y
$T=q_+P_++q_-P_-$, con $0<q_+<q_-<1$. Las incidencias construidas por el
calendario fijan los sectores $90$ y $120$. La respuesta constitutiva es
\begin{equation}
 \mathcal V=(I-T^{90})(I-T^{120})^{-1}.
 \label{iii:eq:vacuum-op}
\end{equation}
Esta definición recibe el transporte procedente del TPK y conserva las
asignaciones de hoja. La regla de respuesta y su dominio son parte explícita
de la construcción.

El cociente procede de comparar las dos incidencias del calendario sobre
un mismo transporte, no de ajustar por separado cuatro magnitudes. Si
$S=T^{30}$ y $\Sigma_j(S)=I+S+\cdots+S^{j-1}$, la factorización de
$I-S^j$ da
\begin{equation}
 \mathcal V\,\Sigma_4(S)=\Sigma_3(S),\qquad
 \mathcal V=\Sigma_3(S)\Sigma_4(S)^{-1}.
 \label{iii:eq:completion-response}
\end{equation}
Los factores $\Sigma_3$ y $\Sigma_4$ expresan respectivamente los
sectores $90=3\cdot30$ y $120=4\cdot30$. Como el espectro de $S$ está
contenido en $(0,1)$, $\Sigma_4(S)$ es invertible: la ecuación de
completación determina una única respuesta para el transporte fijado.
Esta unicidad concierne a la regla constitutiva declarada y a sus
incidencias. La identificación electromagnética añade la asignación
orientada de las hojas y las rectas dimensionales que se precisan más
abajo.

\begin{proposition}[Descomposición y positividad]
La respuesta está bien definida y
\begin{equation}
 \mathcal V=r_+P_++r_-P_-,\qquad
 r_\pm=\frac{1-q_\pm^{90}}{1-q_\pm^{120}}.
 \label{iii:eq:rchannels}
\end{equation}
Sus autovalores satisfacen $3/4<r_-<r_+<1$ en esta cámara.
\end{proposition}
\begin{proof}
Cada coeficiente $1-q_\pm^{120}$ es positivo. La inversa de $I-T^{120}$
es la suma de los proyectores divididos por esos coeficientes, lo que prueba
\eqref{iii:eq:rchannels}. Con $s=q^{30}$ se obtiene
\begin{equation}
 f(s)=\frac{1+s+s^2}{1+s+s^2+s^3},\qquad
 f'(s)=-\frac{s^2(3+2s+s^2)}{(1+s+s^2+s^3)^2}<0.
 \label{iii:eq:monotone}
\end{equation}
Los límites en $0$ y $1$ son $1$ y $3/4$, respectivamente. La monotonía
y el orden de los canales concluyen la prueba.
\end{proof}

\subsection{Cuatro coordenadas de una respuesta común}

Se definen las evaluaciones normalizadas
\begin{equation}
 \widehat\varepsilon=r_+^2,\qquad
 \widehat\mu=r_-^2,\qquad
 \widehat Z=\frac{r_-}{r_+},\qquad
 \widehat c=\frac1{r_+r_-}.
 \label{iii:eq:four}
\end{equation}
La primera pareja conserva las respuestas cuadráticas asignadas a cada hoja.
El producto recoge el modo común; el cociente distingue su orientación relativa.
Se verifican exactamente
\begin{equation}
 \widehat\mu\widehat\varepsilon=\widehat c^{-2},\qquad
 \frac{\widehat\mu}{\widehat\varepsilon}=\widehat Z^2,
 \qquad
 r_+=\frac1{\sqrt{\widehat Z\widehat c}},\quad
 r_-=\sqrt{\frac{\widehat Z}{\widehat c}}.
 \label{iii:eq:four-identities}
\end{equation}
Las raíces positivas fijan la inversión. Bajo la conjugación de hojas,
leída respecto de las marcas $P_\pm$ fijas, se
intercambian $\widehat\varepsilon$ y $\widehat\mu$, se invierte
$\widehat Z$ y se conserva $\widehat c$. La respuesta se transporta entonces
a la cámara $x>-y>0$: las fórmulas siguen siendo válidas, con los órdenes
$q_-<q_+$ y $r_+<r_-$ invertidos. Por consiguiente, el modo común
y la orientación son datos recuperables mediante observables complementarios.

\begin{proposition}[Determinación de las evaluaciones cuadráticas]
Sea $W=r_+P_++r_-P_-$ una respuesta positiva sobre las dos hojas.
Fijadas las lecturas de producto y cociente
\[
 \widehat c^{-1}=\det W=r_+r_-,
 \qquad \widehat Z=r_-/r_+,
\]
existe una única pareja positiva $(\widehat\varepsilon,\widehat\mu)$
que satisface
$\widehat\mu\widehat\varepsilon=\widehat c^{-2}$ y
$\widehat\mu/\widehat\varepsilon=\widehat Z^2$.
Es precisamente $(r_+^2,r_-^2)$.
\end{proposition}
\begin{proof}
La positividad permite tomar las raíces sin ambigüedad:
\[
 \widehat\mu=\frac{\widehat Z}{\widehat c}
 =\frac{r_-}{r_+}r_+r_-=r_-^2,\qquad
 \widehat\varepsilon=\frac1{\widehat Z\widehat c}
 =\frac{r_+}{r_-}r_+r_-=r_+^2.
\]
Las expresiones verifican las dos relaciones y son las únicas raíces
positivas de ellas.
\end{proof}

Así, los cuadrados no constituyen dos elecciones independientes añadidas
a los canales: quedan determinados por las lecturas multiplicativa y
orientada, una vez fijadas las relaciones constitutivas. Esta proposición
establece la unicidad dentro de ese sistema de reglas. La interpretación
de las reglas como respuesta electromagnética conserva su contenido
adicional de representación física.

La lectura determinantal del transporte y la respuesta constitutiva
también deben conservar su orden. Para $T=e^{-xI-y\mathcal R}$ se tiene
$\det T=e^{-2x}$, mientras que
\[
 \det\mathcal V=f(q_+^{30})f(q_-^{30})=\widehat c^{-1}.
\]
Ambas proceden del mismo operador de dos hojas, pero evalúan funciones
diferentes de su espectro. En general, aplicar primero el determinante
y después una función racional no equivale a tomar el determinante de
esa función del operador. Mantener los dos autovalores hasta completar
la respuesta conserva la información de orientación que un determinante
aislado ya no distingue.

\subsection{Inversión cúbica y recuperación angular}

\begin{theorem}[Inversión constitutiva en el dominio contractivo]
Para cada $r\in(3/4,1)$ existe una única $s\in(0,1)$ tal que $f(s)=r$.
Esta $s$ es la raíz única en dicho intervalo de
\begin{equation}
 r s^3+(r-1)(1+s+s^2)=0.
 \label{iii:eq:cubic}
\end{equation}
La respuesta etiquetada reconstruye $s_\pm$, $q_\pm=s_\pm^{1/30}$ y
\begin{equation}
 x=-\frac12\log(q_+q_-),\qquad
 y=\frac12\log\frac{q_-}{q_+}.
 \label{iii:eq:recover-angular}
\end{equation}
\end{theorem}
\begin{proof}
La derivada y los límites de \eqref{iii:eq:monotone} prueban que $f$ es una
biyección decreciente entre los intervalos indicados. Multiplicar $f(s)=r$
por su denominador produce \eqref{iii:eq:cubic}. La raíz positiva de orden
$30$ recupera el canal. Finalmente, suma y diferencia de
$\log q_\pm=-x\mp y$ dan \eqref{iii:eq:recover-angular}.
\end{proof}

La reconstrucción conserva las etiquetas de hoja. Con el operador completo
$\mathcal V$ se conservan también los proyectores y se aplica la función
inversa por cálculo espectral. Recuperar dos coordenadas angulares es una
operación definida sobre esta representación; la traza histórica completa
del TPK conserva datos adicionales.

El intervalo $r\in(3/4,1)$ y la función inversa de este teorema pertenecen
a la normalización interna de \eqref{iii:eq:four}. Si se dispone de
coordenadas expresadas en otra carta de unidades, se recupera primero
esta normalización mediante la transformación de la
ecuación~\eqref{iii:eq:dimensions}; después se aplica la inversión cúbica.

\subsection{Coordenadas logarítmicas y refinamiento}

Sean $\mathcal E=\log(r_+r_-)$ y $\mathcal B=\log(r_-/r_+)$. Entonces
\begin{equation}
 \log\mathcal V=\tfrac12\mathcal E I-\tfrac12\mathcal B\mathcal R.
 \label{iii:eq:logvac}
\end{equation}
Esta descomposición exhibe la componente invariante y la componente orientada.
En la amplificación $\mathcal V_m=\mathcal V\otimes I_{9^m}$, las
expectativas parciales normalizadas satisfacen
$E_m(\mathcal V_{m+1})=\mathcal V_m$. Para todo polinomio $p$,
$p(\mathcal V_m)=p(\mathcal V)\otimes I_{9^m}$; de ahí
$E_m(p(\mathcal V_{m+1}))=p(\mathcal V_m)$. La continuidad uniforme
extiende la identidad a las funciones continuas del espectro. Se conservan
así los proyectores, las potencias y el logaritmo en la torre indicada.

\subsection{Tipos dimensionales de las cuatro evaluaciones}

Con las rectas de acción, carga y velocidad, provistas de secciones
$u_S,u_Q,u_C$, la realización dimensional se escribe
\begin{align}
 Z_{\rm vac}&=\widehat Z u_Su_Q^{-2},&
 c_{\rm vac}&=\widehat c u_C,\\
 \mu_{\rm vac}&=\widehat\mu u_Su_C^{-1}u_Q^{-2},&
 \varepsilon_{\rm vac}&=\widehat\varepsilon u_S^{-1}u_C^{-1}u_Q^2.
 \label{iii:eq:dimensions}
\end{align}
Multiplicando y dividiendo estas expresiones se preservan
$\mu_{\rm vac}\varepsilon_{\rm vac}=c_{\rm vac}^{-2}$ y
$\mu_{\rm vac}/\varepsilon_{\rm vac}=Z_{\rm vac}^2$.
La carta SI evalúa después las secciones. La comparación física requiere
exponer su construcción y normalización junto a la evaluación; las cuatro
coordenadas normalizadas de \eqref{iii:eq:four} conservan su tipo adimensional.

% INSERCIÓN FOCAL AUTORIZADA: no introduce un capítulo ni el concepto de radión.
% Incluir dentro de la sección constitutiva, después de la inversión cúbica
% y de la declaración de las cartas dimensionales. Etiquetas fuera de grupos
% de prefijado automático. Conserva íntegramente el bloque geométrico anterior.
\subsection{Reconstrucción de la forma de acción desde la respuesta constitutiva}
\label{iii:subsec:puente-forma-LC}

La respuesta constitutiva y la realización inductiva--capacitiva reciben
el mismo par angular construido desde APP--TRIT--TPK. La primera evalúa
funciones racionales de los dos canales; la segunda utiliza la razón de
sus exponentes para determinar la forma de la elipse de acción. La
inversión cúbica permite componer ambas lecturas sin identificar sus
cocientes numéricos.

Sean $x=A_{\rm rad}$ e $y=C^*_{\rm rad}$, con $x>|y|$, las coordenadas
levantadas ya construidas, y consérvense sus etiquetas de hoja. Escribimos
\begin{equation}
 q_\pm=e^{-(x\pm y)},\qquad s_\pm=q_\pm^{30},\qquad
 r_\pm=f(s_\pm),\qquad
 f(s)=\frac{1+s+s^2}{1+s+s^2+s^3}.
 \label{iii:eq:forma-canales}
\end{equation}
El dominio contractivo da $s_\pm\in(0,1)$ y
$r_\pm\in(3/4,1)$. En la cámara $x>y>0$ se cumplen
$s_+<s_-$ y $r_+>r_-$; el intercambio de hojas transporta esta cámara
a la orientación opuesta.

La imagen de las dos coordenadas constitutivas normalizadas es el dominio
\begin{equation}
 \mathcal U=\left\{(z,v)\in(0,\infty)^2:
   (zv)^{-1/2}\in(3/4,1),\quad
   (z/v)^{1/2}\in(3/4,1)\right\}.
 \label{iii:eq:forma-dominio}
\end{equation}
Aquí $(z,v)=(\widehat Z,\widehat c)$. Las restricciones expresan la
pertenencia de los canales reconstruidos al dominio del teorema de
inversión; no constituyen una elección ulterior de sus valores.

\Needspace{23\baselineskip}
\begin{proposition}[Composición constitutiva de la forma y de la impedancia LC]
\label{iii:prop:forma-LC}
Para $(\widehat Z,\widehat c)\in\mathcal U$, sean $s_\pm\in(0,1)$ las
raíces únicas de
\begin{equation}
 r_\pm s_\pm^3+(r_\pm-1)(1+s_\pm+s_\pm^2)=0,
 \qquad
 r_+=\frac1{\sqrt{\widehat Z\widehat c}},\quad
 r_-=\sqrt{\frac{\widehat Z}{\widehat c}}.
 \label{iii:eq:forma-inversion}
\end{equation}
La razón angular, la anisotropía de la elipse de acción y la impedancia
LC relativa se reconstruyen exactamente mediante
\begin{align}
 \frac{y}{x}
 &=\frac{\log s_+-\log s_-}{\log s_++\log s_-},
 \label{iii:eq:forma-razon}\\
 \eta_{\rm el}
 &=\frac12\log\frac{\log s_+}{\log s_-},
 \label{iii:eq:forma-eta}\\
 \frac{Z_{\rm LC}}{Z_*}
 &=e^{-\eta_{\rm el}}
  =\sqrt{\frac{\log s_-}{\log s_+}}.
 \label{iii:eq:forma-LC}
\end{align}
En la última identidad $Z_{\rm LC}=Z_{\eta_{\rm el}}$ es la impedancia
de la carta LC de \eqref{hmtII:eq:ii-elipse-lc}.
\end{proposition}
\begin{proof}
La inversión cúbica de \eqref{iii:eq:cubic} determina $s_\pm$ de manera
unívoca, porque $f$ es una biyección decreciente de $(0,1)$ sobre
$(3/4,1)$. Sus raíces positivas de orden $30$ recuperan $q_\pm$.
Las identidades
\[
 \log s_+=-30(x+y),\qquad \log s_-=-30(x-y)
\]
prueban \eqref{iii:eq:forma-razon} por diferencia y suma. Ambos logaritmos
son negativos; su cociente es positivo y su suma no es nula. Por ello,
\[
 \frac12\log\frac{\log s_+}{\log s_-}
   =\frac12\log\frac{x+y}{x-y}
   =\operatorname{artanh}(y/x)=\eta_{\rm el}.
\]
La relación $Z_{\eta_{\rm el}}/Z_*=e^{-\eta_{\rm el}}$, demostrada en
\eqref{hmtII:eq:ii-elipse-impedancia}, da \eqref{iii:eq:forma-LC}.
\end{proof}

La composición conserva la información necesaria para cada reconstrucción.
Con $\hbar$ se recuperan
$a_S=\sqrt{2\hbar}\,e^{\eta_{\rm el}/2}$ y
$b_S=\sqrt{2\hbar}\,e^{-\eta_{\rm el}/2}$; con la carta
$(\omega,Z_*)$ se recuperan la inductancia y la capacitancia ya definidas.
La acción, el reloj y la referencia dimensional acompañan la forma como
datos tipados. La inversión de dos coordenadas adimensionales no
reconstruye sus unidades ni toda la memoria del estado.

La normalización interviene antes de resolver las ecuaciones cúbicas.
Si se reciben $Z_{\rm vac}$ y $c_{\rm vac}$ en una carta dimensional,
las relaciones \eqref{iii:eq:dimensions} determinan primero
\[
 \widehat Z=\frac{Z_{\rm vac}}{u_Su_Q^{-2}},\qquad
 \widehat c=\frac{c_{\rm vac}}{u_C}.
\]
Esas son las coordenadas a las que se aplica
\eqref{iii:eq:forma-inversion}. La referencia $Z_*$ pertenece a la
realización LC y se conserva expresamente en \eqref{iii:eq:forma-LC}.

Los dos cocientes están relacionados mediante los canales, pero evalúan
funciones diferentes:
\begin{equation}
 \widehat Z=\frac{f(s_-)}{f(s_+)},\qquad
 \frac{Z_{\rm LC}}{Z_*}
       =\sqrt{\frac{\log s_-}{\log s_+}}.
 \label{iii:eq:forma-dos-cocientes}
\end{equation}
La reconstrucción anterior proporciona el mapa entre ambas lecturas.
En particular, intercambiar las hojas invierte ambos cocientes, cambia
$\eta_{\rm el}$ por $-\eta_{\rm el}$ y conserva $\widehat c$.
En el límite simétrico $y=0$, los canales coinciden y ambos cocientes
valen $1$. Fuera de ese caso, la igualdad de origen no impone igualdad
de las dos funciones. La forma de acción es recuperable desde la
respuesta constitutiva conservando orientación y normalización.

% RESULTADO_RECUPERADO: funtor dimensional integral, eq:cZ-internos,
% eq:longitud-ruta; composición reunida en Artículo V REV02, sección 69.
% Incorporación autónoma a III: mismo transporte, mismas hojas y misma carta.
\subsection{Velocidad constitutiva y realización cinemática del reloj}
\label{iii:sec:enlace-velocidad-constitutiva}

La respuesta constitutiva y el lector cinemático realizan una misma
sección de velocidad. Su identificación conserva el orden de construcción:
las hojas de APP, su orientación TRIT y el transporte angular publicado
por el TPK anteceden a las lecturas de velocidad y a la elección de unidades.
En la representación de dos hojas de la
sección~\ref{iii:sec:hojas}, con $x>|y|$ y proyectores de rango uno
$P_\pm$, el transporte es
\[
 T=e^{-xI-y\mathcal R}=q_+P_++q_-P_-,
 \qquad q_\pm=e^{-x\mp y},\qquad 0<q_\pm<1.
\]
Sus bloques temporales de longitudes $90$ y $120$ producen la respuesta
$\mathcal V$ de \eqref{iii:eq:vacuum-op}. Las coordenadas $x,y$ y los
canales son salidas del desarrollo angular previo; esta composición
aplica sus lectores a ese transporte ya construido.

El lector logarítmico simétrico del transporte temporal es
\begin{equation}
 \mathcal E=
 \log\!\bigl[(1-q_+^{90})(1-q_-^{90})\bigr]
 -\log\!\bigl[(1-q_+^{120})(1-q_-^{120})\bigr].
 \label{iii:eq:velocidad-lector-simetrico}
\end{equation}
En la recta dimensional de velocidad, con sección positiva $u_C$, el
lector cinemático publica $c_{\mathrm{int}}=\exp(-\mathcal E)u_C$.
La respuesta constitutiva publica
$c_{\mathrm{vac}}=\widehat c u_C$ mediante \eqref{iii:eq:four}.

\begin{proposition}[Identidad constitutiva y cinemática]
\label{iii:prop:identidad-velocidad}
En la misma carta dimensional se verifica
\begin{equation}
 c_{\mathrm{int}}=c_{\mathrm{vac}}=\widehat c\,u_C,
 \qquad \widehat c=(r_+r_-)^{-1}=\exp(-\mathcal E).
 \label{iii:eq:identidad-velocidad}
\end{equation}
La identificación conserva el intercambio de hojas.
\end{proposition}
\begin{proof}
Todos los factores de \eqref{iii:eq:velocidad-lector-simetrico} son
positivos. La ley del logaritmo y \eqref{iii:eq:rchannels} dan
\[
 \mathcal E=\log(r_+r_-)=\log\det\mathcal V,
 \qquad e^{-\mathcal E}=(\det\mathcal V)^{-1}.
\]
El determinante corresponde a la realización de dos hojas de rango uno.
Las dos definiciones publican, por tanto, la misma sección. El intercambio
de hojas permuta $r_+$ y $r_-$ y conserva su producto. El determinante
$\det T=e^{-2x}$ corresponde al transporte elemental, mientras que
$\det\mathcal V=r_+r_-$ corresponde a su respuesta temporal;
\eqref{iii:eq:identidad-velocidad} utiliza el segundo.
\end{proof}

\subsubsection{Longitud de trayectoria y duración elemental}

El marco dimensional positivo $(u_S,u_C,u_T,u_Q)$ induce la base de
longitud $u_L=u_Cu_T$. Para una trayectoria enriquecida $\gamma$ de
$n=|\gamma|$ transiciones y canal constante, los lectores aditivos son
\begin{equation}
 T_{\mathrm{int}}(\gamma)=nu_T,
 \qquad L_{\mathrm{int}}(\gamma)=n\widehat c\,u_L.
 \label{iii:eq:velocidad-longitud-reloj}
\end{equation}
La concatenación suma el número de transiciones y sus realizaciones.
Las rutas enriquecidas conservan, además, su orientación y su memoria.

\begin{corollary}[Normalización de la longitud elemental]
Para $n>0$, $L_{\mathrm{int}}(\gamma)/T_{\mathrm{int}}(\gamma)
=c_{\mathrm{vac}}$. En la normalización $t_0=u_T$,
\begin{equation}
 \ell_0=c_{\mathrm{int}}t_0=\widehat c\,u_L.
 \label{iii:eq:longitud-elemental-constitutiva}
\end{equation}
En una carta temporal $t_0=\tau_0u_T$, con $\tau_0>0$, se obtiene
$\ell_0=\tau_0\widehat c\,u_L$.
\end{corollary}
\begin{proof}
Dividir las dos expresiones de \eqref{iii:eq:velocidad-longitud-reloj}
cancela $n$ y utiliza $u_L/u_T=u_C$. Multiplicar la sección de velocidad
por $t_0$ prueba ambas fórmulas elementales.
\end{proof}

Cuando el canal depende de la arista, el lector conserva la forma aditiva
$L_{\mathrm{int}}(\gamma)=\sum_{e\in\gamma}\widehat c(e)u_L$.
Su cociente por $nu_T$ es la media de las velocidades de esas aristas.
La expresión de canal constante es la especialización homogénea de
esta lectura.

\subsubsection{Cambio de carta y normalización adaptada}

\begin{proposition}[Covariancia de la identificación]
Sean $u_C'=d u_C$ y $u_T'=\eta u_T$, con $d,\eta>0$.
Las coordenadas de las mismas secciones satisfacen
\begin{equation}
 \widehat c'=d^{-1}\widehat c,
 \qquad \tau_0'=\eta^{-1}\tau_0,
 \qquad u_L'=d\eta u_L.
 \label{iii:eq:velocidad-cambio-carta}
\end{equation}
Si otra publicación escribe $c'=\widehat c' u_C'$ con $u_C'=d u_C$,
la igualdad de secciones equivale exactamente a $d\widehat c'=\widehat c$.
\end{proposition}
\begin{proof}
La igualdad de cada sección en las dos bases da
$\widehat c'u_C'=\widehat c u_C$ y $\tau_0'u_T'=\tau_0u_T$.
En particular,
\[
 \tau_0'\widehat c'u_L'
 =\frac{\tau_0}{\eta}\frac{\widehat c}{d}\,d\eta u_L
 =\tau_0\widehat c u_L=\ell_0.
\]
La última equivalencia procede de la misma sustitución para la velocidad.
\end{proof}

La elección adaptada $d=\widehat c$ da $u_C'=c_{\mathrm{vac}}$ y
$\widehat c'=1$: expresa la sección construida en una base adaptada.
El coeficiente de la carta interna sigue siendo $(r_+r_-)^{-1}$, y esa
es la normalización utilizada por la inversión cúbica. Manteniendo las
bases de acción y carga, \eqref{iii:eq:dimensions} produce
\[
 \widehat\mu'=d\widehat\mu,\qquad
 \widehat\varepsilon'=d\widehat\varepsilon,\qquad
 \widehat\mu'\widehat\varepsilon'=(\widehat c')^{-2}.
\]
En la carta adaptada se tiene $\widehat\mu'=r_-/r_+$ y
$\widehat\varepsilon'=r_+/r_-$, por sustitución de
$\widehat\mu=r_-^2$ y $\widehat\varepsilon=r_+^2$.

\subsubsection{Conservación del lector bajo amplificación}

Para el refinamiento $\mathcal V_m=\mathcal V\otimes I_{9^m}$, la
traza normalizada $\tau_m=\operatorname{Tr}/(2\cdot9^m)$ satisface
\begin{equation}
 \exp[-2\tau_m(\log\mathcal V_m)]
 =\exp[-\log r_+-\log r_-]=\widehat c.
 \label{iii:eq:velocidad-traza-normalizada}
\end{equation}
En efecto, cada autovalor aparece $9^m$ veces; esa multiplicidad se
cancela con el denominador de la traza normalizada. El determinante
ordinario ampliado es $(r_+r_-)^{9^m}$. La normalización de dos hojas,
equivalentemente expresada por \eqref{iii:eq:velocidad-traza-normalizada},
conserva el lector constitutivo durante la amplificación. Esta
realización conserva el transporte enriquecido que la origina.

\section{Sección de carga y relaciones de cuantización eléctrica}
\label{iii:sec:carga}

La respuesta constitutiva establece una relación entre las dos hojas y una
escala relativa de propagación. La sección de carga se obtiene al componer
esta respuesta con la constante de estructura fina y la acción, ambas
construidas en los capítulos antecedentes. El orden tiene una consecuencia
precisa: la carga no determina retrospectivamente el registro dodecafásico.
Es el registro, junto con la acción y la respuesta del vacío, el que
determina la sección de carga en la recta correspondiente.

\subsection{Determinación positiva de la sección de carga}

Consideremos las secciones positivas $\hbar_a$, con
$a\in\{\mathrm{pre},\mathrm{ret}\}$, y $h_a=2\pi\hbar_a$.
La impedancia $Z_{\rm vac}$ ya está construida mediante
\eqref{iii:eq:dimensions}. El cociente $h_a/Z_{\rm vac}$ pertenece a
la recta cuadrática de carga. Esta compatibilidad dimensional permite
formular la siguiente definición sin elegir todavía una carta SI.

\begin{definition}[Sección positiva de carga]
La sección de carga correspondiente a la orientación de acción $a$ es
\begin{equation}
 e_a=\left(\frac{2\alpha h_a}{Z_{\rm vac}}\right)^{1/2}.
 \label{iii:eq:charge}
\end{equation}
La raíz se toma en la semirrecta positiva de carga. Su sección conjugada
es $-e_a$; el signo de esta última no se identifica con el índice de
retorno de la acción.
\end{definition}

\begin{proposition}[Existencia, unicidad y compatibilidad del acoplamiento]
La ecuación $Z_{\rm vac}e_a^2=2\alpha h_a$ determina una única sección
positiva. En la realización constitutiva se cumple también
\begin{equation}
 \alpha=\frac{Z_{\rm vac}e_a^2}{2h_a}
 =\frac{e_a^2}{4\pi\varepsilon_{\rm vac}\hbar_a c_{\rm vac}}.
 \label{iii:eq:coupling}
\end{equation}
\end{proposition}
\begin{proof}
La positividad de $\alpha$, $h_a$ y $Z_{\rm vac}$ garantiza una raíz
positiva única. Las identidades constitutivas dan
$\varepsilon_{\rm vac}c_{\rm vac}=Z_{\rm vac}^{-1}$.
Sustituir esta igualdad y $h_a=2\pi\hbar_a$ en la primera expresión
de \eqref{iii:eq:coupling} produce la segunda.
\end{proof}

La última expresión es la forma habitual del acoplamiento electromagnético
en una carta física. Aquí aparece como identidad posterior de las secciones
construidas. La prueba no vuelve a calcular $\alpha$ a partir de una carga
medida: establece la coherencia algebraica del resultado anterior con la
sección que acaba de definirse. La identificación física exige, además,
mantener la carta de unidades y el régimen del acoplamiento empleados en la
comparación. Una identidad entre magnitudes y una comparación experimental
tienen funciones diferentes en ese reconocimiento.

\subsection{Resistencia, conductancia, flujo y frecuencia}

Las siguientes cuatro magnitudes utilizan la misma sección $e_a$. Su
dependencia puede resolverse por completo antes de evaluar ningún decimal.

\begin{theorem}[Composición de los cuantos eléctricos]
Las secciones
\begin{align}
 R_K&=\frac{h_a}{e_a^2}=\frac{Z_{\rm vac}}{2\alpha},
 &G_0&=\frac{2e_a^2}{h_a}=\frac{4\alpha}{Z_{\rm vac}},
 \label{iii:eq:rk-g0}\\
 \Phi_{0,a}&=\frac{h_a}{2e_a}
 =\sqrt{\frac{h_aZ_{\rm vac}}{8\alpha}},
 &K_{J,a}&=\frac{2e_a}{h_a}=\Phi_{0,a}^{-1}
 \label{iii:eq:flux-josephson}
\end{align}
satisfacen exactamente
\begin{equation}
 R_KG_0=2,\qquad G_0Z_{\rm vac}=4\alpha,\qquad
 K_{J,a}\Phi_{0,a}=1,\qquad
 K_{J,a}^{2}R_K=\frac4{h_a}.
 \label{iii:eq:quanta-identities}
\end{equation}
\end{theorem}
\begin{proof}
La sustitución de $e_a^2=2\alpha h_a/Z_{\rm vac}$ en las dos primeras
definiciones cancela $h_a$ y da \eqref{iii:eq:rk-g0}. Al elevar al
cuadrado la sección positiva $h_a/(2e_a)$ se obtiene
$h_aZ_{\rm vac}/(8\alpha)$; la positividad fija la raíz en
\eqref{iii:eq:flux-josephson}. Su inversa es $2e_a/h_a$.
Los tres primeros productos de \eqref{iii:eq:quanta-identities} se
reducen, respectivamente, a $2$, $4\alpha$ y $1$. Finalmente,
$(4e_a^2/h_a^2)(h_a/e_a^2)=4/h_a$.
\end{proof}

La realización física reconoce en estas secciones la resistencia de von
Klitzing, el cuanto de conductancia, el cuanto de flujo y la constante de
Josephson. El factor $2$ de la conductancia pertenece a la normalización
de la magnitud aquí definida; su identificación con una multiplicidad de
canales de un sistema material exige especificar esa representación.
Las cuatro ecuaciones no constituyen cuatro generadores independientes:
comparten $\alpha$, la acción y la respuesta del vacío. Precisamente esa
dependencia produce las identidades cruzadas y permite comprobar una
evaluación sin alterar las operaciones que la originaron.

\subsection{Retorno de la acción y transporte de las secciones}

El cociente $R_{\rm act}=h_{\rm pre}/h_{\rm ret}$, obtenido antes del
par angular, conserva la diferencia entre las dos secciones de acción.
No cambia las asignaciones de hoja de la respuesta constitutiva. Son dos
operaciones distintas: retorno de la acción e intercambio de los canales.

\begin{proposition}[Covariancia bajo el retorno]
Con $\alpha$ y $Z_{\rm vac}$ comunes, se tiene
\begin{align}
 \frac{e_{\rm pre}}{e_{\rm ret}}&=R_{\rm act}^{1/2},
 &R_K^{\rm pre}&=R_K^{\rm ret},
 &G_0^{\rm pre}&=G_0^{\rm ret},\label{iii:eq:charge-return}\\
 \frac{\Phi_{0,\rm pre}}{\Phi_{0,\rm ret}}&=R_{\rm act}^{1/2},
 &\frac{K_{J,\rm pre}}{K_{J,\rm ret}}&=R_{\rm act}^{-1/2}.
 \label{iii:eq:flux-return}
\end{align}
\end{proposition}
\begin{proof}
La primera identidad sigue de dividir las dos versiones de
\eqref{iii:eq:charge} y tomar la raíz positiva. En
\eqref{iii:eq:rk-g0} la acción ya se ha cancelado. Las otras dos
identidades siguen de las potencias $h_a^{1/2}$ y $h_a^{-1/2}$ en
\eqref{iii:eq:flux-josephson}.
\end{proof}

Resistencia y conductancia conservan el cociente constitutivo, mientras
que flujo y frecuencia retienen también la orientación de acción. Esta
separación da una lectura estructural de sus dependencias: unas magnitudes
son invariantes del cociente de hojas y otras miden, además, la sección de
acción utilizada. El contraste metrológico debe declarar qué sección evalúa;
la proximidad numérica de las dos acciones no autoriza identificarlas.

% RESULTADO_RECUPERADO desde II REV09. Véase MAPA_PROCEDENCIA.json.
% Selección de líneas y cambios mecánicos explícitos; ninguna prueba nueva.
\section{Incidencia hexádico--octádica y funcional de Barbero--Immirzi}
\label{v:ant:sec:barbero}
Sea $H\hookrightarrow O$ la inclusión marcada de una hexada en una octada,
$|H|=6$, $|O|=8$, construida mediante el residuo binario de la
sección~\ref{exc:residuo-binario}. Se conservan la hexada, la dodecada
y la carta regional que determinan esa inclusión.
Los objetos de incidencia
\[
\mathcal F_6=H\times\binom H2,\qquad
\mathcal F_8=O\times\binom H2
\]
tienen cardinales $90$ y $120$. Su diferencia normalizada da
\[
\frac{|\mathcal F_6|-|\mathcal F_8|}{24\binom62}=-\frac1{12}.
\]
\begin{figure}[htbp]
\centering
\begin{tikzpicture}[>=Latex,every node/.style={font=\small}]
\draw[rounded corners,draw=ink] (-1.8,-1.35) rectangle (1.8,1.35);
\foreach \a in {0,60,...,300}{
\fill[ink] (\a:0.83) circle(2pt);}
\node at (0,-1.7) {$H:\ 6$ puntos};
\draw[->,thick,ink] (2.1,0)--(3.6,0);
\node[above] at (2.85,1.5) {inclusión marcada};
\draw[rounded corners,draw=ink] (4,-1.35) rectangle (7.6,1.35);
\foreach \a in {0,60,...,300}{
\fill[ink] ($(5.8,0)+(\a:0.83)$) circle(2pt);}
\fill[accent] (4.5,0) circle(2pt);
\fill[accent] (7.1,0) circle(2pt);
\node at (5.8,-1.7) {$O:\ 8$ puntos};
\node at (0,-2.35) {$6\binom62=90$};
\node at (5.8,-2.35) {$8\binom62=120$};
\end{tikzpicture}
\caption{Incidencia marcada de seis a ocho puntos. El dibujo representa
los conjuntos y la inclusión; no pretende convertir las caras o vértices
de un cubo en bloques de Steiner sin su mapa de incidencia.}
\end{figure}

\subsection{Los pesos preceden al coeficiente de polo}
En el módulo libre de sucesos orientados se fija la cadena
fundamental de la vuelta
\[
c_{12}^{+}=
\sum_{j\in\Z/12\Z}[j,\mathcal F_6]
-[\partial_{\rm ret},\mathcal F_8].
\]
La multiplicidad doce corresponde a los sectores y la multiplicidad uno
a la costura orientada. La involución cambia la orientación de la cadena
completa. Sea
\[
S_m(q)=\frac{q^m}{1-q^{3m}},\qquad 0<q<1.
\]
El lector lineal asigna $S_{90}$ a cada sector y $S_{120}$ a la costura;
en consecuencia,
\[
\mathscr R(c_{12}^{+})=12S_{90}-S_{120}.
\]

\begin{definition}[Evaluación orientada de la incidencia hexádico--octádica]
\label{paper:barbero-definicion}
El parámetro \(\gamma_{\HMT}\) es la evaluación conjugada de la cadena
\(c_{12}^{+}\), completada por su componente angular. En el dominio
\(0<q_\pm<1\), se define mediante
\[
\gamma_{\HMT}=\frac{\pi AC^*}{180}
+[\mathscr R(c_{12}^{+})](q_-)-[\mathscr R(c_{12}^{+})](q_+).
\]
La incidencia marcada determina los espacios de sucesos; la cadena
orientada fija las multiplicidades, y los canales de la
sección~\ref{iv:angular:inversion} determinan la evaluación.
\end{definition}

\begin{theorem}[Coeficiente de polo y evaluación conjugada]
La imagen de la cadena fundamental tiene coeficiente $1/24$
respecto de $(1-q)^{-1}$ y residuo $-1/24$ respecto de $(q-1)^{-1}$.
Su evaluación en los canales previamente generados determina
\begin{equation}
\gamma_{\HMT}=
\frac{\pi AC^*}{180}
+12\bigl[S_{90}(q_-)-S_{90}(q_+)\bigr]
-\bigl[S_{120}(q_-)-S_{120}(q_+)\bigr].
\label{v:ant:eq:gamma}
\end{equation}
\end{theorem}
\begin{proof}
La factorización $1-q^{3m}=(1-q)\sum_{j=0}^{3m-1}q^j$ implica
$\lim_{q\uparrow1}(1-q)S_m(q)=1/(3m)$. Por linealidad,
\[
\frac{12}{270}-\frac1{360}=\frac1{24}.
\]
La coordenada $q-1$ invierte el signo del residuo.
La evaluación conjugada y la componente angular de la definición~\ref{paper:barbero-definicion}
producen \eqref{v:ant:eq:gamma}.
\end{proof}

El certificado diofántico $4a-3b=45$ admite los pares
$(12+3t,1+4t)$, $t\ge0$; certifica la minimalidad de $(12,1)$,
pero no reemplaza el origen de las multiplicidades en la cadena.
La evaluación de \eqref{v:ant:eq:gamma} es
\[
\gamma_{\HMT}
=0.27400767269445927474725526077398041940\ldots.
\]
Esta cifra es una salida del funcional orientado; no es un valor externo
usado para seleccionar $A$, $C^*$ o sus coeficientes.

\subsection{Paridad y sensibilidad estructural}
Al invertir $C^*$ se intercambian $q_+$ y $q_-$ y cambia de signo
la componente angular. Por ello
\[
\gamma_{\HMT}(A,-C^*)=-\gamma_{\HMT}(A,C^*).
\]
El funcional conserva orientación, mientras que otros lectores,
como $q_+q_-$, son pares. Esta diferencia precede a su interpretación física.

% RESULTADO_RECUPERADO desde II REV09. Véase MAPA_PROCEDENCIA.json.
% Selección de líneas y cambios mecánicos explícitos; ninguna prueba nueva.
\section{Acción, gravedad e información térmica}
\label{v:ant:sec:gravity}
\subsection{Período dodecafásico y realización circular de la acción}
El calendario ya construido determina $T_{108}=108t_0$.
La velocidad constitutiva, la duración y la longitud de trayectoria se
obtienen mediante las identidades de la
proposición~\ref{iii:prop:identidad-velocidad} y de las
ecuaciones~\eqref{iii:eq:velocidad-longitud-reloj} y
\eqref{iii:eq:longitud-elemental-constitutiva}.
En su realización circular,
\[
\omega_{108}=\frac{2\pi}{108t_0},\qquad
R_{108}=\frac{c_{\rm int}}{\omega_{108}}
=\frac{54}{\pi}\ell_0,\qquad \ell_0=c_{\rm int}t_0.
\]
La unidad $\ell_0$ y las cartas de tiempo y velocidad se mantienen
explícitas; la carta SI no selecciona el ciclo.
Para cada sección orientada de acción $a\in\{\mathrm{pre},\mathrm{ret}\}$,
\[
E_{108,a}=\hbar_a\omega_{108},\qquad
m_{108,a}=\frac{E_{108,a}}{c_{\rm int}^2}.
\]
La longitud reducida y la completa son respectivamente
\[
\frac{\hbar_a}{m_{108,a}c_{\rm int}}=R_{108},
\qquad
\frac{h_a}{m_{108,a}c_{\rm int}}=2\pi R_{108}.
\]
La misma acción y la misma duración intervienen en ambas lecturas.

\begin{proposition}[Selector de coincidencia de radios]
En la familia dimensional
\[
G_a(\vartheta)=\vartheta\frac{c_{\rm int}^5t_0^2}{\hbar_a},
\qquad \vartheta>0,
\]
la condición de la realización circular
$G_a(\vartheta)m_{108,a}/c_{\rm int}^2=R_{108}$
tiene la única solución
\[
\vartheta^\circ_{108}=(54/\pi)^2.
\]
\end{proposition}
\begin{proof}
El primer radio es
$r_g(\vartheta)=\vartheta(\pi/54)\ell_0$.
Compararlo con $(54/\pi)\ell_0$, con $\ell_0>0$, reduce la
igualdad a una ecuación lineal estrictamente creciente.
\end{proof}
La condición de coincidencia es parte explícita de esta realización;
el teorema resuelve su coeficiente sin usar un valor metrológico de $G$.
El radio aquí empleado es $Gm/c^2$, por lo que no se traslada
inadvertidamente el factor dos de la convención $2Gm/c^2$.

\begin{proposition}[Acción y gravedad recíprocas, área conservada]
En las dos cartas anteriores,
\[
\frac{G_{\rm pre}}{G_{\rm ret}}=R_{\rm act}^{-1},\qquad
\hbar_aG_a=\vartheta^\circ_{108}c_{\rm int}^5t_0^2,
\qquad
\ell_P^2=\frac{\hbar_aG_a}{c_{\rm int}^3}
=\vartheta^\circ_{108}\ell_0^2.
\]
\end{proposition}
\begin{proof}
La definición de $G_a$ contiene $\hbar_a^{-1}$ y el resto de factores
es común a las dos cartas. Tomar el cociente y el producto da
las tres igualdades.
\end{proof}
Esta reciprocidad explica por qué la información de orientación puede
variar sin que cambie el área. Se dispone así de una relación precisa
entre las secciones recíprocas de acción, la escala de Planck y el operador areal
de la sección~\ref{v:ant:sec:area}.

% RESULTADO_RECUPERADO desde II REV09. Véase MAPA_PROCEDENCIA.json.
% Selección de líneas y cambios mecánicos explícitos; ninguna prueba nueva.
\section{Operador de área y reconstrucción de la información sectorial}
\label{v:ant:sec:area}
El reagrupamiento de seis trits construye una biyección entre
$(\Z/9\Z)^3$ y $(\Z/27\Z)^2$, ambos de cardinal $729$.
La biyección no se presenta como isomorfismo aditivo. Sobre un bloque
de borde $9\times9$, la fuente descompone 36 enlaces en dos
conjuntos de 18, etiquetados por los canales $90$ y $120$.

\subsection{Reagrupamiento trítico e incidencia del borde}
\label{v:ant:sec:ii-area-incidencia-explicita}

La correspondencia conserva una palabra ordenada, además de su cardinal.
Escribamos cada coordenada \(x_j\in\mathbb Z/9\mathbb Z\), para
\(j=1,2,3\), mediante sus representantes posicionales únicos
\(x_j=r_{2j-1}+3r_{2j}\), con \(r_k\in\{0,1,2\}\). El mapa es
\begin{equation}
 \beta_{729}(x_1,x_2,x_3)
 =\bigl(r_1+3r_2+9r_3,\ r_4+3r_5+9r_6\bigr)
 \in(\mathbb Z/27\mathbb Z)^2.
 \label{v:ant:eq:ii-area-beta729}
\end{equation}

\begin{proposition}[Correspondencia posicional entre interior y borde]
\label{v:ant:prop:ii-area-beta729}
La aplicación \(\beta_{729}\) es biyectiva y conserva el orden de los
seis trits. Su inversa extrae las dos ternas de dígitos y las reagrupa
en tres parejas.
\end{proposition}
\begin{proof}
La expansión posicional de los representantes \(0,\ldots,8\) y
\(0,\ldots,26\) es única. Extraer los seis dígitos, reagruparlos y
extraerlos de nuevo devuelve la misma palabra. La operación inversa
restituye así cada una de las tres coordenadas iniciales.
\end{proof}

El acarreo conserva su posición en la palabra enriquecida. Las operaciones
aditivas de los dos empaquetados tienen fronteras diferentes: por ejemplo,
\(\beta_{729}(0,2,0)=(18,0)\) y
\(\beta_{729}(0,1,0)=(9,0)\), cuya suma en el toro de orden
veintisiete es \((0,0)\), mientras
\(\beta_{729}(0,3,0)=(0,1)\). Esta distinción permite reutilizar
la correspondencia posicional manteniendo tipadas las operaciones.

En el toro \((\mathbb Z/27\mathbb Z)^2\), el bloque
\(B_\partial=\{0,\ldots,8\}^2\) tiene los conjuntos de enlaces
orientados
\begin{align}
 \partial B_{90}
 &=\{((-1,j),(0,j)),\ ((8,j),(9,j)):0\le j\le8\},\\
 \partial B_{120}
 &=\{((i,-1),(i,0)),\ ((i,8),(i,9)):0\le i\le8\}.
 \label{v:ant:eq:ii-area-enlaces}
\end{align}
Todas las coordenadas se leen módulo veintisiete. Cada familia contiene
dieciocho enlaces y ambas son disjuntas. Sus etiquetas conservan los
canales de incidencia; el cardinal de enlaces es \(18+18\).

\begin{theorem}[Capacidad mínima del corte cúbico]
\label{v:ant:thm:ii-area-corte-minimo}
En el grafo cúbico de vértices \(\{0,\ldots,N-1\}^3\), con
\(N\ge2\), enlaces entre vecinos y capacidad unitaria por enlace,
el corte mínimo entre dos caras opuestas tiene capacidad \(N^2\).
Para \(N=9\), un estado producto máximamente mezclado sobre los
trits de una capa mínima tiene información
\begin{equation}
 \operatorname{cap}_{\min}=81,\qquad s_{\rm corte}=81\log3.
 \label{v:ant:eq:ii-area-corte-informacion}
\end{equation}
\end{theorem}
\begin{proof}
Las \(N^2\) rectas paralelas a la normal de las caras forman caminos
disjuntos por enlaces. Todo corte separador intercepta cada camino, de
modo que contiene al menos \(N^2\) enlaces. Los enlaces que atraviesan
una capa transversal constituyen un corte de ese cardinal. Cada trit
uniforme aporta \(\log3\) nats; la entropía es aditiva para el estado
producto de los ochenta y un trits de la capa.
\end{proof}

La capacidad del corte cúbico y el estado de los treinta y seis enlaces
de \eqref{v:ant:eq:ii-area-enlaces} pertenecen a observables definidos sobre
cortes diferentes. La genealogía conserva ambos dominios y sus medidas.

\subsection{Transporte de los modos colectivos de incidencia}
\label{v:ant:sec:ii-area-transporte-colectivo}

Los espacios \(\mathcal F_6\) y \(\mathcal F_8\) de la
sección~\ref{v:ant:sec:barbero} proporcionan los vectores uniformes
\[
 u_{90}=\frac1{\sqrt{90}}\sum_{f\in\mathcal F_6}(\delta_f,0),
 \qquad
 u_{120}=\frac1{\sqrt{120}}\sum_{f\in\mathcal F_8}(0,\delta_f)
\]
en \(\ell^2(\mathcal F_6)\oplus\ell^2(\mathcal F_8)\).
En \(\ell^2(\partial B_{90})\oplus\ell^2(\partial B_{120})\)
se construyen
\[
 v_{90}=\frac1{\sqrt{18}}\sum_{e\in\partial B_{90}}(\delta_e,0),
 \qquad
 v_{120}=\frac1{\sqrt{18}}\sum_{e\in\partial B_{120}}(0,\delta_e).
\]
Las dos parejas son ortonormales. Denotemos por
\(\mathscr H_{\rm inc}^{\rm coll}\) y
\(\mathscr H_\partial^{\rm coll}\) sus respectivos planos generados.

\begin{theorem}[Entrelazamiento colectivo de las etiquetas sectoriales]
\label{v:ant:thm:ii-area-entrelazamiento-colectivo}
Existe una única isometría lineal sobreyectiva entre esos planos tal que
\begin{equation}
 \mathcal J_{6\to8}u_m=v_m\quad(m=90,120).
 \label{v:ant:eq:ii-area-j-colectivo}
\end{equation}
Para los operadores
\[
 \begin{aligned}
 D_{\rm inc}&=90|u_{90}\rangle\langle u_{90}|
             +120|u_{120}\rangle\langle u_{120}|,\\
 D_\partial&=90|v_{90}\rangle\langle v_{90}|
             +120|v_{120}\rangle\langle v_{120}|,
 \end{aligned}
\]
se cumple
\begin{equation}
 \mathcal J_{6\to8}D_{\rm inc}
 =D_\partial\mathcal J_{6\to8}.
 \label{v:ant:eq:ii-area-j-entrelaza}
\end{equation}
Se transportan también los proyectores sectoriales y su combinación
orientada de coeficientes \(12:-1\).
\end{theorem}
\begin{proof}
La imagen de una base fija la aplicación lineal única. Las normas y los
productos internos se conservan porque ambas bases son ortonormales.
La identidad \eqref{v:ant:eq:ii-area-j-entrelaza} vale sobre cada vector
\(u_m\), pues sus dos miembros son \(m v_m\). Los proyectores se
expresan como \((120I-D)/30\) y \((D-90I)/30\); su cálculo funcional
lineal y cualquier combinación de ambos conservan el entrelazamiento.
\end{proof}

Este transporte actúa sobre los dos modos uniformes y sus etiquetas
espectrales. Las incidencias individuales y las fluctuaciones dentro de
cada sector permanecen en sus espacios completos. El enlace conserva
exactamente el dato colectivo que utilizará la ponderación areal.

\subsection{Carácter angular y operador de área}
\label{v:ant:sec:ii-area-caracter}

Sea $N_e=\operatorname{diag}(0,1,2)$ en cada enlace y
$N_{90},N_{120}$ las sumas por sector. Fijada la escala
\[
\theta_{\rm clk}=\frac{C^*}{6}\frac{\pi}{180},\qquad
a_0=\frac{1+2\cos(10^\circ)}3\theta_{\rm clk},
\]
el factor angular es el carácter real normalizado del triplete
\begin{equation}
 \chi_{10}=\frac13\Tr\operatorname{diag}
       (1,e^{i\pi/18},e^{-i\pi/18})
 =\frac{1+2\cos(10^\circ)}3,
 \qquad a_0=\chi_{10}\theta_{\rm clk}.
 \label{v:ant:eq:ii-area-caracter-triplete}
\end{equation}
La suma de las fases conjugadas prueba la igualdad. En la rama positiva
considerada, \(\theta_{\rm clk}>0\), \(\chi_{10}>0\) y
\(\gamma_{\HMT}>0\); por tanto, \(a_0\) y la escala areal son
positivos. El triplete angular y su normalización preceden a la
distribución entrópica sobre los enlaces. Con estas cantidades,
el operador areal normalizado es
\[
\widehat{\mathcal A}/\ell_P^2
=\gamma_{\HMT}a_0(N_{90}+N_{120}).
\]
Su espectro tiene valores $n\gamma_{\HMT}a_0$, $0\le n\le72$,
con multiplicidades dadas por $(1+z+z^2)^{36}$.
La prueba consiste en sumar los autovalores $0,1,2$ de los
36 factores tensoriales.

\begin{theorem}[Recuperación de los sectores de borde]
Sean
\[
N=N_{90}+N_{120},\qquad D=90N_{90}+120N_{120}.
\]
Entonces
\begin{equation}
N_{90}=4N-\frac D{30},\qquad
N_{120}=\frac D{30}-3N.
\label{v:ant:eq:recover}
\end{equation}
Por tanto, el total y el lector ponderado recuperan conjuntamente
las dos ocupaciones.
\end{theorem}
\begin{proof}
La matriz $\left(\begin{smallmatrix}1&1\\90&120\end{smallmatrix}\right)$
tiene determinante $30$. Su inversión da \eqref{v:ant:eq:recover}.
En el dominio de ocupaciones enteras, la imagen satisface
$D\in30\Z$ y $90N\le D\le120N$.
\end{proof}

El Hamiltoniano modular de la fuente es
$K_A=-\log\rho_A=cI+\Delta_{\rm mod}D$.
Para $\Delta_{\rm mod}\ne0$ y normalización conocida,
$D=(K_A-cI)/\Delta_{\rm mod}$ publica el segundo observable.
Este operador describe la estructura espectral del estado reducido;
no se identifica por notación con el generador de evolución temporal.
El área conserva la ocupación total y la pareja añade la procedencia.
Por ejemplo, $(N_{90},N_{120})=(1,0)$ y $(0,1)$ comparten $N=1$,
pero tienen $D=90$ y $D=120$.

% RESULTADO_RECUPERADO desde II REV09. Véase MAPA_PROCEDENCIA.json.
% Selección de líneas y cambios mecánicos explícitos; ninguna prueba nueva.
\subsection{Normalización del estado y significado del Hamiltoniano modular}
El segundo observable se explicita mediante el estado que lo define.
Se mantiene la descomposición tensorial anterior: dieciocho qutrits
por sector, cada uno con operador número $\operatorname{diag}(0,1,2)$;
$N_m$ es la suma de esos operadores y los sectores actúan en factores
distintos. Sea
\[
 z_m=1+e^{-m\Delta_{\rm mod}}+e^{-2m\Delta_{\rm mod}},
 \qquad m\in\{90,120\}.
\]
La familia de estados de la fuente tiene la realización producto
\[
 \rho_A=Z^{-1}e^{-\Delta_{\rm mod}(90N_{90}+120N_{120})},
 \qquad Z=z_{90}^{18}z_{120}^{18}.
\]
La conmutatividad de los operadores número y la factorización tensorial
prueban la expresión de $Z$. El estado es de rango completo para
$\Delta_{\rm mod}$ real y finito, y su logaritmo espectral da
\[
 K_A=-\log\rho_A=(\log Z)I+\Delta_{\rm mod}D.
\]
Para $\Delta_{\rm mod}\ne0$, $K_A$ pesa de forma distinta las excitaciones
de los dos sectores. La lectura conjunta con $N$ recupera las dos ocupaciones
por la ecuación~\eqref{v:ant:eq:recover}.

La necesidad de la lectura conjunta se ve en dos ejemplos complementarios.
Los estados de ocupación $(1,0)$ y $(0,1)$ tienen igual total $N=1$
y diferente peso modular. A la inversa, $(4,0)$ y $(0,3)$ tienen
el mismo $D=360$ y diferentes totales. La pareja $(N,D)$ separa ambos
casos. Esta recuperación se refiere a ocupaciones sectoriales, no a
cada vector microscópico dentro de una degeneración de esas ocupaciones.

Cuando $\Delta_{\rm mod}=0$, el estado es proporcional a la identidad
y $K_A$ es escalar: esa lectura deja de distinguir los sectores.
Para $\Delta_{\rm mod}\ne0$, la normalización conocida permite recuperar
$D$ y componerlo con el área. Así queda definido el papel del Hamiltoniano
modular: caracteriza el estado reducido y su distribución de información.
Su inclusión en el artículo está motivada por esa función matemática,
mientras el Hamiltoniano de evolución pertenece al desarrollo dinámico
de la acción.


\subsection{Estado reducido, descomposición de Schmidt y entropía de entrelazamiento}
Si $\rho_A=\sum_\nu p_\nu|\nu\rangle\langle\nu|$,
la purificación
\[
|\Psi\rangle=\sum_\nu\sqrt{p_\nu}\,
|\nu\rangle_A|\nu\rangle_B
\]
satisface $\Tr_B|\Psi\rangle\langle\Psi|=\rho_A$.
Así, $-\Tr(\rho_A\log\rho_A)$ es la entropía de entrelazamiento
de esta bipartición pura concreta. El estado ampliado declara
qué sistema porta esa interpretación.

Para una involución de dos hojas con un autovector de cada signo,
\[
\rho_y=\frac{e^{-yR}}{2\cosh y},\qquad R^2=I,
\]
la diagonalización proporciona
\[
S(\rho_y)=\log(2\cosh y)-y\tanh y,\qquad
D(\rho_y\Vert\rho_{-y})=2y\tanh y.
\]
La entropía es par en $y$; la distancia compara explícitamente
las dos orientaciones. No se identifica esta entropía de dos niveles
con la capacidad $81\log3$ del corte triádico.

% RESULTADO_RECUPERADO desde II REV09. Véase MAPA_PROCEDENCIA.json.
% Selección de líneas y cambios mecánicos explícitos; ninguna prueba nueva.
% Fragmento aditivo. No sustituye 10b_boltzmann.tex.
% Procedencia y distinción entre recuperación y formalización:
% 10c_boltzmann_procedencia.md. Etiquetas locales prefijadas ii-kb-aut.

\section{Memoria genealógica, reducción informacional y transducción térmica}
\label{v:ant:sec:ii-kb-aut-desarrollo}

La construcción térmica utiliza dos publicaciones del mismo desarrollo
APP--TRIT--TPK. La primera conserva los estados, las rutas admisibles y
la información del registro; la segunda proporciona la acción orientada
y la duración del retorno. Su composición determina una energía por
decisión trítica. Las coordenadas de temperatura y energía se introducen
después mediante aplicaciones dimensionales explícitas. Esta secuencia
permite identificar qué información se conserva, qué información publica
un lector y qué magnitud convierte esa información en entropía dimensional.

\subsection{Lector, memoria y recuperación del estado}

Fijemos un nivel finito del desarrollo y denotemos por \(\Omega\) el
conjunto de estados enriquecidos admisibles. El registro conserva hoja,
orientación, residuo, cociente, acarreo, ruta, fase y memoria en los
dominios donde interviene cada coordenada. Sean
\[
 q:\Omega\longrightarrow Y,\qquad
 M:\Omega\longrightarrow\mathcal M
\]
la publicación visible y la memoria retenida, respectivamente. Para una
distribución \(p\) sobre \(\Omega\), escribimos
\(\mathsf H(p)=-\sum_xp_x\log p_x\), con \(0\log0=0\).
El logaritmo se aplica a probabilidades adimensionales.

\begin{definition}[Recuperación genealógica de una publicación]
El par \(\widehat q=(q,M)\) es un lector recuperable sobre el soporte
\(\Omega_p\) cuando existe
\[
 d:\widehat q(\Omega_p)\longrightarrow\Omega_p,
 \qquad d\circ\widehat q=\operatorname{id}_{\Omega_p}.
\]
La información de fibra de \(q\) es
\(\mathsf H_{\mathrm{fib}}(q;p)=\mathsf H(X\mid q(X))\), donde
\(X\) tiene ley \(p\).
\end{definition}

\begin{proposition}[Identidad de recuperación y coste de la reducción]
Para un lector recuperable,
\begin{align}
 \mathsf H(X\mid q(X),M(X))&=0,\\
 \mathsf H(X)&=\mathsf H(q(X))+\mathsf H(X\mid q(X)),\\
 \mathsf H(X\mid q(X))&=\mathsf H(M(X)\mid q(X)).
 \label{v:ant:eq:ii-kb-aut-fibra}
\end{align}
\end{proposition}
\begin{proof}
La aplicación \(d\) recupera \(X\) del par \((q(X),M(X))\),
por lo que la primera entropía condicional es nula. Como \(q\) es
función de \(X\), la regla de la cadena da la segunda igualdad.
Para la tercera, se desarrolla de dos maneras
\(\mathsf H(X,M(X)\mid q(X))\). Una proporciona
\(\mathsf H(X\mid q(X))\), pues \(M\) es función de \(X\);
la otra proporciona
\(\mathsf H(M(X)\mid q(X))+\mathsf H(X\mid q(X),M(X))\).
\end{proof}

La división euclídea de las hojas de APP muestra el mecanismo
aritmético elemental. Con \(i,j\in\{1,\ldots,9\}\), se conserva
\[
 i+j=9q^+_{ij}+R^+_{ij},\qquad
 ij=9q^\times_{ij}+R^\times_{ij},\qquad
 1\leq R^\sigma_{ij}\leq9\quad(\sigma\in\{+,\times\}).
\]
El par residuo--cociente recupera el valor entero de la operación.
Al sumar las diferencias sobre las ochenta y una posiciones se obtiene
\[
 2025-810=54+9\cdot129.
\]
La memoria del cociente restaura la contribución que no publica la
reducción residual. La recuperabilidad del estado completo requiere,
además, conservar las coordenadas de posición y de ruta pertinentes;
el par residuo--cociente identifica aquí el resultado de la operación.

\subsection{Actualización reversible, decisión trítica y borrado}

Sea \(U:\Omega\to\Omega\) la actualización TPK en un dominio donde
el transporte enriquecido posee inversa. La ley transportada es
\(p'_y=p_{U^{-1}y}\). La biyección conserva el multiconjunto de
probabilidades y, por tanto,
\begin{equation}
 \mathsf H(U(X))=\mathsf H(X),\qquad
 \mathsf H(X\mid U(X))=0.
 \label{v:ant:eq:ii-kb-aut-reversible}
\end{equation}
Si la salida observada es \(q(U(X))\), la información no publicada
se cuantifica por \(\mathsf H(X\mid q(U(X)))\). El estado completo,
su publicación y el reinicio de un registro son así tres mapas distintos.

Para los tres estados orientados del TRIT, sea
\(p=(p_{-1},p_0,p_{+1})\) una distribución normalizada. Su información es
\begin{equation}
 I_{\mathrm{TRIT}}(p)=-\sum_{a=-1}^{1}p_a\log p_a,
 \qquad 0\leq I_{\mathrm{TRIT}}(p)\leq\log3.
 \label{v:ant:eq:ii-kb-aut-trit}
\end{equation}
El extremo superior corresponde a \(p_a=1/3\). En efecto,
\[
 D(p\Vert(1/3,1/3,1/3))
 =\sum_ap_a\log(3p_a)=\log3-I_{\mathrm{TRIT}}(p)\geq0;
\]
la desigualdad se obtiene de la convexidad de \(t\log t\), y la
igualdad caracteriza la distribución uniforme. El extremo inferior
corresponde a una distribución concentrada en un estado.

La realización ideal de Landauer utiliza una memoria con estados lógicos
energéticamente degenerados, un entorno isotermo de temperatura \(\Theta\)
y un balance que conserva la salida del mapa lógico y contabiliza los
registros auxiliares. En ese modelo, la disminución de entropía lógica es
\(k_B[\mathsf H(X)-\mathsf H(F(X))]
=k_B\mathsf H(X\mid F(X))\) para un mapa determinista \(F\).
El balance entrópico exige transferir al entorno al menos esa cantidad;
su expresión energética es \(Q\geq k_B\Theta\mathsf H(X\mid F(X))\).
La distinción entre evolución lógica reversible y borrado físico
corresponde a la realización de Landauer--Bennett \cite{bennett2003}.
La actualización invertible
\eqref{v:ant:eq:ii-kb-aut-reversible} tiene contribución lógica nula; el
reinicio de una terna uniforme requiere transferir \(\log3\) nats
al entorno y satisface \(Q\geq k_B\Theta\log3\).
La energía de control y la disipación de una realización física
concreta pertenecen a su balance termodinámico adicional.

\subsection{Acción, frecuencia y energía de decisión}

Sean \(\mathcal L_{\rm act}\) y \(\mathcal L_T\) las rectas
dimensionales de acción y duración; la recta energética es
\(\mathcal L_E=\mathcal L_{\rm act}\otimes\mathcal L_T^*\).
El desarrollo anterior aporta las secciones positivas
\(\hbar_a\in\mathcal L_{\rm act}\), con
\(a\in\{\mathrm{pre},\mathrm{ret}\}\), y la duración elemental
\(t_0\in\mathcal L_T\). El calendario completo proporciona
\[
 T_{108}=108t_0,\qquad
 \omega_{108}=\frac{2\pi}{T_{108}},\qquad
 h_a=2\pi\hbar_a.
\]
La contracción dimensional entre acción y frecuencia define
\begin{equation}
 \epsilon_a:=\hbar_a\omega_{108}
 =\frac{h_a}{108t_0}
 =\frac{\hbar_a}{t_P},\qquad
 t_P=\frac{54}{\pi}t_0.
 \label{v:ant:eq:ii-kb-aut-energia}
\end{equation}
El último miembro emplea el mismo período circular reducido que la
sección gravitatoria. Las tres expresiones corresponden a una única
sección de energía, con una orientación \(a\) mantenida en todos
los factores. La información normalizada de una decisión uniforme
produce la sección de energía por nat
\begin{equation}
 \varepsilon_a:=\frac{\epsilon_a}{\log3}\in\mathcal L_E^+.
 \label{v:ant:eq:ii-kb-aut-energia-nat}
\end{equation}
La acción previamente construida determina la escala de energía al
fijar la duración; el cardinal de alternativas del TRIT determina su
normalización informacional. El valor de \(\log3\) entra por el
conteo y la medida declarados en \eqref{v:ant:eq:ii-kb-aut-trit}.

\subsection{Definición dimensional de Boltzmann y selección de coordenadas}

Sea \(\mathcal L_\Theta\) la recta orientada de temperatura.
El espacio dimensional de entropía es
\(\mathcal L_{\rm ent}=\mathcal L_E\otimes\mathcal L_\Theta^*\).
Una aplicación lineal positiva
\(k_B:\mathcal L_\Theta\to\mathcal L_E\) es, equivalentemente,
un elemento positivo de \(\mathcal L_{\rm ent}\). Por eso el mismo
objeto puede actuar sobre una temperatura o multiplicar una información
adimensional.

\begin{definition}[Normalización térmica de la decisión trítica]
Un par térmico compatible con la sección orientada \(a\) consta de
un transductor positivo \(k_B\) y una sección positiva
\(\Theta_{{\rm clk},a}\in\mathcal L_\Theta\) tales que
\begin{equation}
 k_B(\Theta_{{\rm clk},a})=\varepsilon_a,
 \qquad
 k_B\Theta_{{\rm clk},a}\log3=\epsilon_a.
 \label{v:ant:eq:ii-kb-aut-par-termico}
\end{equation}
La constante de Boltzmann desempeña la función de transductor
energía--temperatura de la información trítica en esta realización.
\end{definition}

\begin{theorem}[Unicidad relativa a la sección térmica y covariancia dimensional]
Fijada una sección positiva \(\Theta_{{\rm clk},a}\), existe una
única aplicación lineal positiva que satisface
\eqref{v:ant:eq:ii-kb-aut-par-termico}. Para bases positivas \(u_E,u_\Theta\),
si
\[
 \epsilon_a=\epsilon_a^{\rm num}u_E,\quad
 \Theta_{{\rm clk},a}=\theta_a u_\Theta,\quad
 k_B(u_\Theta)=b\,u_E,
\]
su coordenada es
\begin{equation}
 b=\frac{\epsilon_a^{\rm num}}{\theta_a\log3}.
 \label{v:ant:eq:ii-kb-aut-coordenada}
\end{equation}
Bajo \(u'_E=s_Eu_E\), \(u'_\Theta=s_\Theta u_\Theta\), con
\(s_E,s_\Theta>0\), se cumple
\[
 b'=\frac{s_\Theta}{s_E}b,\qquad
 \theta'_a=\frac{\theta_a}{s_\Theta},\qquad
 (\epsilon_a^{\rm num})'=\frac{\epsilon_a^{\rm num}}{s_E}.
\]
\end{theorem}
\begin{proof}
Una sección no nula genera una recta. Prescribir su imagen determina
la aplicación lineal única; la positividad de ambas secciones prueba
su positividad. Sustituir las coordenadas en
\eqref{v:ant:eq:ii-kb-aut-par-termico} da
\(b\theta_a\log3=\epsilon_a^{\rm num}\). Las tres leyes de
transformación se obtienen expresando las mismas secciones y la misma
aplicación en las bases nuevas.
\end{proof}

Una misma constante de Boltzmann debe actuar en las dos orientaciones.
Puesto que \(t_P\) es común, la compatibilidad simultánea equivale a
\begin{equation}
 \frac{\Theta_{{\rm clk},{\rm pre}}}{\Theta_{{\rm clk},{\rm ret}}}
 =\frac{\epsilon_{\rm pre}}{\epsilon_{\rm ret}}
 =\frac{\hbar_{\rm pre}}{\hbar_{\rm ret}}.
 \label{v:ant:eq:ii-kb-aut-orientaciones}
\end{equation}
En efecto, un isomorfismo de rectas conserva cocientes de secciones.
Recíprocamente, si la igualdad anterior se cumple, el isomorfismo
determinado en una orientación satisface también la normalización de
la otra. Así, el retorno modifica la sección térmica del ciclo en la
misma proporción que la energía, mientras el transductor permanece
común a ambas lecturas.

La ecuación anterior explicita la elección requerida para publicar un
valor individual. Si la sección térmica todavía no ha sido seleccionada
por un lector independiente, para cada \(\lambda>0\) el par
\[
 (k_B,\Theta_{{\rm clk},a})
 \longmapsto(\lambda k_B,\lambda^{-1}\Theta_{{\rm clk},a})
\]
conserva \(\epsilon_a\). Por consiguiente, la presentación de una
coordenada individual internamente seleccionada exige dar el lector de
\(\Theta_{{\rm clk},a}\), las bases y su transporte dimensional.
Las ecuaciones de acción e información permanecen determinadas antes
de esa especificación. La identificación posterior de las bases con
joule y kelvin pertenece a la publicación metrológica del resultado.

\subsection{Entropía dimensional e información del estado reducido}

Para una distribución de rutas admitidas \(p\), o para un estado
reducido de matriz de densidad \(\rho\), sean
\[
 s(p)=-\sum_xp_x\log p_x,\qquad
 s(\rho)=-\operatorname{Tr}(\rho\log\rho).
\]
Sus publicaciones dimensionales son
\begin{equation}
 S(p)=k_Bs(p),\qquad S(\rho)=k_Bs(\rho)
 \quad\text{en }\mathcal L_{\rm ent}.
 \label{v:ant:eq:ii-kb-aut-entropia}
\end{equation}
La evaluación en \(\Theta\) produce una energía
\(\Theta S=k_B\Theta\,s\). Para la terna uniforme a la temperatura
del ciclo, \eqref{v:ant:eq:ii-kb-aut-par-termico} da exactamente
\[
 S_{\rm TRIT}=k_B\log3,\qquad
 \Theta_{{\rm clk},a}S_{\rm TRIT}=\epsilon_a.
\]
Si \(\rho_A\) es la reducción de la purificación bipartita
construida en la sección areal, \(s(\rho_A)\) mide el entrelazamiento
de esa bipartición. El transporte de unidades conserva sus
autovalores, su orientación y las ocupaciones sectoriales. La escala
areal \(\gamma a_0\ell_P^2\) convierte ocupaciones en área; el
transductor \(k_B\) convierte información en entropía dimensional.
Ambas lecturas utilizan el mismo registro mediante operadores distintos.

\subsection{Normalización espectral de las historias térmicas}

Sea \(\mathcal G\) un grafo finito de extensiones admisibles ya
construido, con matriz de adyacencia \(A\). Cada arista \(e\) conserva
sus coordenadas de ruta y recibe un incremento adimensional de acción
\(a_*(e)>0\) del lector correspondiente. Elegida la sección de energía
\(\epsilon_*\), pongamos \(t=\beta\epsilon_*\), donde
\(\beta=(k_B\Theta)^{-1}\). El operador térmico es
\begin{equation}
 (L_tf)(y)=\sum_{e:x\to y}\exp[-t a_*(e)]f(x).
 \label{v:ant:eq:ii-kb-aut-operador}
\end{equation}
Todos sus exponentes son adimensionales. La condición
\(\rho(L_t)=1\) normaliza el crecimiento ponderado de las historias.

\begin{proposition}[Existencia y unicidad de la presión normalizada]
Supongamos que \(A\) es irreducible, \(\rho(A)>1\), y que
\(0<m\leq a_*(e)\leq M\) para todas las aristas. Existe un único
\(t_*>0\) tal que \(\rho(L_{t_*})=1\), y
\begin{equation}
 \frac{\log\rho(A)}{M}\leq t_*\leq
 \frac{\log\rho(A)}{m}.
 \label{v:ant:eq:ii-kb-aut-presion}
\end{equation}
\end{proposition}
\begin{proof}
La comparación entrada a entrada proporciona
\(e^{-tM}A\leq L_t\leq e^{-tm}A\) para \(t\geq0\).
La monotonía del radio espectral de matrices no negativas implica
las cotas correspondientes para \(\rho(L_t)\). El radio espectral
es continuo, vale \(\rho(A)>1\) en cero y tiende a cero en infinito.
Para \(s>0\), además,
\(L_{t+s}\leq e^{-sm}L_t\), de modo que
\(\rho(L_{t+s})\leq e^{-sm}\rho(L_t)<\rho(L_t)\).
La continuidad y la disminución estricta prueban existencia y
unicidad. Sustituir el valor uno en las cotas iniciales demuestra
\eqref{v:ant:eq:ii-kb-aut-presion}.
\end{proof}

La realización local de tres alternativas equiponderadas, representada
por un vértice con tres lazos y \(a_*(e)=1\), satisface
\(\rho(L_t)=3e^{-t}\). Su normalización da \(t_*=\log3\).
Con \(\epsilon_*=\epsilon_a\), resulta
\(\beta_{\rm clk}=\log3/\epsilon_a\), exactamente la relación
\eqref{v:ant:eq:ii-kb-aut-par-termico}. Este cálculo identifica la
confluencia local trítica. Para un grafo de historias con restricciones
adicionales, la confluencia con el mismo reloj exige comprobar
\[
 \rho\!\left(L_{(\log3/\epsilon_a)\epsilon_*}\right)=1
\]
con sus incrementos efectivos. La normalización local y la presión
del sistema completo conservan así sus respectivos dominios.

La relación obtenida reúne estructura, valor y significado. El registro
determina la información recuperable; las probabilidades sobre las
continuaciones determinan su lectura entrópica; la acción y el período
determinan \(\epsilon_a\); el transductor energía--temperatura
publica la entropía dimensional. El área y el entrelazamiento se
relacionan a través del estado reducido y de sus operadores de borde,
manteniendo visible qué dato produce cada magnitud.

% FORMALIZACION_REUNIDA: representación pentacomponente de II,
% con antecedente Paley y demostración directa de los momentos del marco.
% El retorno transporta la carta; no se identifica un ciclo de orden doce con A5.
\section{Incidencia orientada y representación tensorial gravitatoria}
\label{vii:sec:pentacomponente}
La publicación excepcional del estado conserva la incidencia que acompaña
a su lectura aritmética. La matriz ternaria \(A_W\) de
\eqref{exc:aw} proporciona aquí un antecedente material preciso. Su
levantamiento balanceado \(C=\operatorname{bal}(A_W)\), con representantes
\(0,1,-1\), satisface
\[
 C=C^{\mathsf T},\qquad C^2=5I_6,\qquad \operatorname{Tr}C=0.
\]
Las seis posiciones tienen el orden \((\infty,0,1,2,3,4)\).
Esta incidencia produce primero un espacio tridimensional; su lectura
cuadrática produce después el espacio de cinco componentes sobre el que
actúa el acoplamiento gravitatorio.

\subsection{Recubrimiento firmado y transporte del calendario}
Definimos
\[
 E_C=\frac12(I_6+C/\sqrt5),\quad
 \mathcal H_C=\operatorname{im}E_C,\quad
 v_{u,\sigma}=\sigma\sqrt2E_Ce_u,\quad \sigma\in\{1,-1\}.
\]
\begin{proposition}[Doce posiciones orientadas]
\label{vii:prop:doce-orientadas}
Los vectores \(v_{u,\sigma}\) son doce vectores unitarios distintos
en un espacio de dimensión tres. Su incidencia está determinada por
\(\langle v_{u,\sigma},v_{v,\tau}\rangle=1/\sqrt5\).
La inversión de la orientación actúa por antipodía.
\end{proposition}
\begin{proof}
La identidad cuadrática de \(C\) implica \(E_C^2=E_C=E_C^*\);
su traza es tres. Puesto que la diagonal de \(C\) es nula,
\[
 \langle v_{u,\sigma},v_{v,\tau}\rangle
 =\begin{cases}\sigma\tau,&u=v,\\
 \sigma\tau C_{uv}/\sqrt5,&u\ne v.
 \end{cases}
\]
Esto prueba normalización, distinción de los doce vectores e incidencia.
Cambiar \(\sigma\) por \(-\sigma\) cambia el signo del vector.
\end{proof}

Una expresión euclidiana explícita se obtiene poniendo
\(\nu=\sqrt{1+\varphi^2}\) y
\[
 B=\frac1\nu
 \begin{pmatrix}
 0&-1&-\varphi&0&\varphi&1\\
 -1&-\varphi&0&1&0&-\varphi\\
 -\varphi&0&-1&-\varphi&-1&0
 \end{pmatrix}.
\]
La relación \(\varphi^2=\varphi+1\), para la coordenada ya generada,
da \(B^*B=2E_C\).
En consecuencia \(B/\sqrt2|_{\mathcal H_C}\) es una isometría y lleva
\(v_{u,\sigma}\) a \(\sigma Be_u\). La imagen es
\[
 \mathcal V=\frac1\nu\bigl(
 \{0\}\times\{\pm1\}\times\{\pm\varphi\}\ \cup\
 \{\pm1\}\times\{\pm\varphi\}\times\{0\}\ \cup\
 \{\pm\varphi\}\times\{0\}\times\{\pm1\}\bigr).
\]
Así quedan relacionados la coordenada interna, la incidencia y sus
doce representantes geométricos.

La coordenada del calendario se escribe \((b,r)\in
\mathbb Z/12\mathbb Z\times\{0,\ldots,8\}\). Una ruta de longitud \(n\)
induce
\[
 k(r,n)=\left\lfloor\frac{r+n}{9}\right\rfloor,\qquad
 (b,r)\longmapsto
 \bigl(b+k(r,n)\bmod12,\ (r+n)\bmod9\bigr).
\]
El alineamiento explícito de sus posiciones \(p=b+1\) con la carta firmada es
\[
\begin{array}{c|rrrrrrrrrrrr}
p&1&2&3&4&5&6&7&8&9&10&11&12\\ \hline
u&\infty&0&0&\infty&1&3&3&1&2&4&4&2\\
\sigma&+&-&+&-&-&+&-&+&-&+&-&+
\end{array}
\]
Esta tabla fija una trivialización del portador. Cada cambio de posición
transporta conjuntamente los vectores y su matriz de incidencia.

\begin{proposition}[Covariancia del marco bajo retorno]
Si \(S e_p=e_{p+1\bmod12}\), la representación matricial de un
operador \(P\) del portador se transporta como \(P_k=S^kPS^{-k}\).
Este transporte conserva productos escalares, incidencias y composición
de rutas. Tras doce retornos vuelve la carta, conservándose la memoria
acumulada en el estado fuente.
\end{proposition}
\begin{proof}
\(S\) es ortogonal y \(S^{12}=I\); de aquí se siguen las propiedades
métricas y la periodicidad de la carta. Para la composición, con
\(r'=(r+n)\bmod9\), la división euclidiana da
\(k(r,n+m)=k(r,n)+k(r',m)\). Los exponentes de \(S\) se suman
exactamente con el acarreo. La operación actúa sobre la coordenada
del calendario; las actualizaciones de memoria de la ruta siguen
presentes en su dominio enriquecido.
\end{proof}

\subsection{Lectura cuadrática y espacio de cinco componentes}
Sea \(\mathbb R^{\mathcal V}\) el espacio de las doce coordenadas,
\(\mathsf A e_v=e_{-v}\) su involución antipodal y
\(\mathsf J=\mathbf1\mathbf1^{\mathsf T}\). Se definen
\[
 P_5=\frac12(I+\mathsf A)-\frac1{12}\mathsf J,\qquad
 Q_v=vv^{\mathsf T}-\frac13I_3,\qquad
 \Gamma_0x=\sum_{v\in\mathcal V}x_vQ_v.
\]
\begin{theorem}[Proyector pentacomponente e isometría tensorial]
\label{vii:thm:portador-pentacomponente}
\(P_5\) es el proyector ortogonal de rango cinco sobre el espacio de
coordenadas pares bajo antipodía y de suma nula. Además,
\[
 \Gamma_0^*\Gamma_0=\frac85P_5,\qquad
 \Gamma_0\Gamma_0^*=\frac85I_{\operatorname{Sym}^2_0(\mathbb R^3)}.
\]
Por tanto
\[
 \Gamma_5=\sqrt{\frac58}\,\Gamma_0|_{\operatorname{im}P_5}
\]
es una isometría sobre el espacio de tensores simétricos sin traza.
\end{theorem}
\begin{proof}
El espacio par tiene una coordenada por cada par antipodal, por lo que
su dimensión es seis. El vector constante pertenece a él; restar su
proyector de rango uno da \(P_5\), de rango cinco.

Como
\(\langle Q_v,Q_w\rangle_F=(v^{\mathsf T}w)^2-1/3\), las entradas
de \(\Gamma_0^*\Gamma_0\) valen \(2/3\) para \(w=\pm v\) y
\(-2/15\) en los otros diez lugares de cada fila. Éstas son
exactamente las entradas de \((8/5)P_5\).

También puede comprobarse la identidad en el espacio tensorial sin
recurrir a una clasificación de representaciones. La lista de vértices da
\[
 \sum_vv_i^4=\frac{12}{5},\qquad
 \sum_vv_i^2v_j^2=\frac45\quad(i\ne j),
\]
y los momentos con alguna potencia impar se anulan por los cambios de
signo. En efecto,
\((1+\varphi^2)^2=5\varphi^2\) y \(1+\varphi^4=3\varphi^2\).
Para un tensor simétrico \(Q\), estas sumas dan
\[
 \sum_v(v^{\mathsf T}Qv)\,vv^{\mathsf T}
 =\frac85Q+\frac45(\operatorname{tr}Q)I.
\]
Con \(\operatorname{tr}Q=0\), y usando
\(\sum_vv^{\mathsf T}Qv=4\operatorname{tr}Q=0\), resulta
\(\Gamma_0\Gamma_0^*Q=(8/5)Q\). Las dos identidades prueban
la isometría normalizada y su sobreyectividad.
\end{proof}

La sección gravitatoria de \(\ref{v:ant:sec:gravity}\) determina el
operador \(G_aP_5\). Su espectro tiene \(G_a\) con multiplicidad
cinco y cero con multiplicidad siete. La orientación se conserva en
el portador y en su transporte; \(G_a\) fija el acoplamiento común.

\subsection{Proyección transversal y helicidades}
Para \(n\in S^2\), pongamos \(\Pi=I-nn^{\mathsf T}\) y
\[
 \Lambda_nQ=\Pi Q\Pi-\frac12\operatorname{tr}(\Pi Q\Pi)\Pi.
\]
\begin{proposition}[Reducción transversal de rango dos]
\(\Lambda_n\) es el proyector ortogonal sobre
\(\{Q\in\operatorname{Sym}^2_0(\mathbb R^3):Qn=0\}\), de dimensión dos.
\end{proposition}
\begin{proof}
\(\Pi n=0\) da transversalidad y \(\operatorname{tr}\Pi=2\) da
traza nula. Un tensor transversal sin traza queda fijo.
La fórmula es autoadjunta para el producto de Frobenius.
En un marco \((e_1,e_2,n)\), la imagen tiene base ortonormal
\[
 E_+=\frac{e_1e_1^{\mathsf T}-e_2e_2^{\mathsf T}}{\sqrt2},\qquad
 E_\times=\frac{e_1e_2^{\mathsf T}+e_2e_1^{\mathsf T}}{\sqrt2}.
\]
\end{proof}
Los otros tres vectores de una base ortonormal del portador son
\[
 E_1=\frac{e_1n^{\mathsf T}+ne_1^{\mathsf T}}{\sqrt2},\quad
 E_2=\frac{e_2n^{\mathsf T}+ne_2^{\mathsf T}}{\sqrt2},\quad
 E_0=\frac{2nn^{\mathsf T}-e_1e_1^{\mathsf T}-e_2e_2^{\mathsf T}}{\sqrt6}.
\]
Al rotar el marco transversal un ángulo \(\theta\), las parejas
\((E_+,E_\times)\) y \((E_1,E_2)\) giran \(2\theta\) y
\(\theta\), respectivamente, mientras \(E_0\) queda fijo. La
complejificación tiene así los pesos \(2,-2,1,-1,0\).
\(\Lambda_n\) conserva los dos primeros y anula los restantes.
Queda establecida la cadena
\[
 \mathbb R^{12}\xrightarrow{P_5}\operatorname{im}P_5
 \xrightarrow{\Gamma_5}\operatorname{Sym}^2_0(\mathbb R^3)
 \xrightarrow{\Lambda_n}\mathcal H_{\rm TT}(n),
 \qquad 12\longrightarrow5\longrightarrow5\longrightarrow2.
\]
La representación pentacomponente y su reducción son resultados
algebraicos anteriores a la elección de ecuaciones de campo.
En la realización gravitatoria sin masa, las restricciones de calibre
seleccionan la imagen transversal. Esta distinción conserva tanto las
cinco coordenadas estructurales como las dos helicidades radiativas.


% Artículo VII, recuperación aditiva, revisión 04, 10-09-2026.
% RESULTADO_RECUPERADO: S15, 11c_gravedad_holonomia_rev6.tex:16--126.
% El cociclo canónico c27=(1,18) se conserva como antecedente del cómputo
% de rutas; este fragmento prueba su composición y representación.
% La memoria traslacional finita y la amplitud cosmológica s0 tienen
% tipos distintos. El cociclo no evalúa por sí solo una densidad local.

\section{Retorno de fase y representación de la memoria acumulada}
\label{vii:sec:retorno-memoria-gravedad}

La cronología de las rutas enriquecidas distingue el retorno de posición,
el retorno de orientación y el retorno de fase. La publicación observable
identifica ciertos estados de llegada, mientras el registro acumulativo
conserva la secuencia recorrida. La representación geométrica debe
transportar ambos datos conjuntamente.

\subsection{Cociclo de inversión y cambio de hoja}
\label{vii:subsec:cociclo-memoria}

Sea \(F_{\rm obs}\) la publicación cronológica del transporte TPK.
Sus bloques canónicos tienen la jerarquía
\[
\begin{array}{c|l}
27&\text{retorno posicional con inversión de orientación},\\
54&\text{retorno posicional y orientacional},\\
108&\text{retorno completo de la fase observable}.
\end{array}
\]
El registro cuenta las inversiones de orientación \(g(a)\) y los
cambios efectivos de hoja \(s(a)\). Para una ruta,
\begin{equation}
 c(\gamma)=\sum_{a\subset\gamma}(g(a),s(a)),\qquad
 c(\gamma_2\gamma_1)=c(\gamma_2)+c(\gamma_1).
 \label{vii:eq:cociclo-ruta}
\end{equation}
El cómputo canónico del bloque de longitud \(27\) publica
\(c_{27}=(1,18)\). La composición conserva esos incrementos en
los cuatro bloques que completan el período observable.
En rutas orientadas hacia adelante el registro es acumulativo;
su extensión al grupoide de rutas utiliza \(\mathbb Z^2\) y
\(c(\bar\gamma)=-c(\gamma)\).

\begin{proposition}[Memoria de la vuelta observable completa]
\label{vii:prop:memoria-vuelta-completa}
Escribiendo \(\mathcal K_{27}=F_{\rm obs}^{27}\), la elevación
al registro, para \(x\in\mathcal O_{108}\) y \(z\in\mathbb Z^2\), satisface
\[
 \widetilde{\mathcal K}_{27}(x,z)
   =(\mathcal K_{27}x,z+(1,18)),\qquad
 \widetilde{\mathcal K}_{27}^{\,4}(x,z)
   =(x,z+(4,72)).
\]
Tras \(N\) vueltas completas, el incremento es \(N(4,72)\).
\end{proposition}

\begin{proof}
La proyección observable de \(\mathcal K_{27}\) tiene orden cuatro.
El cómputo canónico anterior fija \(c_{27}=(1,18)\) de manera constante
sobre esa órbita observable. La aditividad de
\eqref{vii:eq:cociclo-ruta} suma los cuatro incrementos y da
\(c_{108}=4(1,18)=(4,72)\).
La repetición aplica la misma ley de concatenación y produce
\(Nc_{108}\). El retorno de la primera coordenada conserva el
incremento de la segunda.
\end{proof}

\subsection{Representación afín del cociclo}
\label{vii:subsec:representacion-afin-cociclo}

La realización discreta utiliza un elemento \(R_4\) de orden cuatro en
\(\operatorname{Spin}^+(1,3)\), un vector temporal unitario \(u\)
fijado por su acción vectorial y una unidad de longitud \(\ell_0>0\).
La escala \(\ell_0\) conserva su procedencia en la sección dimensional.
Definimos
\begin{equation}
 \rho_{\rm C}^{\rm disc}(\gamma)=
 \bigl(R_4^{c_g(\gamma)},\,\ell_0c_s(\gamma)u\bigr)
 \in\operatorname{Spin}^+(1,3)\ltimes\mathbb R^{1,3}.
 \label{vii:eq:representacion-discreta-memoria}
\end{equation}
La potencia actúa sobre el vector mediante la representación vectorial
asociada del grupo de espín.

\begin{theorem}[Composición y amplitud traslacional del retorno]
\label{vii:thm:traslacion-memoria}
La aplicación \(\rho_{\rm C}^{\rm disc}\) conserva la identidad,
la concatenación y la inversión. En los retornos completos,
\begin{equation}
 \rho_{\rm C}^{\rm disc}(\gamma_{108})=(I,72\ell_0u),\qquad
 \rho_{\rm C}^{\rm disc}(\gamma_{108}^{\,N})=(I,72N\ell_0u).
 \label{vii:eq:traslacion-retorno-108}
\end{equation}
La inversión cambia el signo de la traslación.
\end{theorem}

\begin{proof}
El producto semidirecto es
\((R,a)(S,b)=(RS,a+Rb)\). Como \(R_4u=u\),
\[
\begin{split}
 \rho_{\rm C}^{\rm disc}(\gamma_2)
 \rho_{\rm C}^{\rm disc}(\gamma_1)
 &=
 \bigl(R_4^{c_g(\gamma_2)+c_g(\gamma_1)},
 \ell_0(c_s(\gamma_2)+c_s(\gamma_1))u\bigr)\\
 &=\rho_{\rm C}^{\rm disc}(\gamma_2\gamma_1).
\end{split}
\]
El cociclo de la identidad es cero y el de la inversa es el opuesto.
Para la vuelta completa se sustituyen \(c_{108}=(4,72)\) y
\(R_4^4=I\). La concatenación de \(N\) vueltas mantiene la parte
rotacional identidad y suma sus traslaciones.
\end{proof}

La amplitud \(72\ell_0\) es una longitud asociada a una ruta finita.
Una realización mediante cotetrada y conexión la transporta a su
holonomía afín. La contribución local a las ecuaciones gravitatorias
se obtiene después mediante la curvatura de esa conexión, la
variación de la acción material y la normalización de su densidad.
La representación conserva así el contenido acumulativo exacto que
debe intervenir en esa realización.

% Artículo VII. Incorporación de resultados preexistentes, 10-09-2026.
% Propietarios íntegros preservados en fuentes_conservadas:
% 17_pantalla_cubica_soldadura.tex (P), 18_clausura_energetica.tex (E).
% P: incidencia, cociente, soldadura celular y refinamiento.
% E: respuesta, unicidad, extensión mínima y expectativa tracial.
% Dependencias anteriores de VII: núcleo APP--TRIT--TPK; realización
% cúbica de la ruta, cociclo de transporte y compresión q_vis=5/9.
% Este fragmento no constituye por sí solo el artículo autónomo.
% La identidad de pantalla se distingue de la contracción material de espín.

\section{Soldadura de frontera y respuesta energética}
\label{vii:sec:pantalla-respuesta}

La realización cúbica del transporte conserva seis incidencias orientadas.
Su cociente geométrico, la forma lorentziana y la respuesta energética
se obtienen de la misma matriz de incidencias. El transporte procede de
la ruta enriquecida APP--TRIT--TPK: la cara, el signo, el acarreo y la
memoria acompañan a cada arista. En esta sección se desarrolla la
operación lineal que actúa sobre esa publicación de frontera.

\subsection{Incidencia cúbica y cociente de rango cuatro}
\label{vii:sec:cociente-pantalla}

Sea \(H_F=\mathbb R^F\), con
\(F=\{(i,\epsilon):1\leq i\leq3,\ \epsilon\in\{-1,1\}\}\),
su producto escalar euclidiano y su base \(e_{i,\epsilon}\).
La oposición de caras es la involución ortogonal
\(Re_{i,\epsilon}=e_{i,-\epsilon}\). Definimos
\[
 u=\sum_{i,\epsilon}e_{i,\epsilon},\qquad
 P_u=\frac{uu^{\mathsf T}}6,\qquad P_-=\frac{I-R}{2},\qquad
 P_0=\frac{I+R}{2}-P_u,\qquad P_\square=P_u+P_-.
\]
Los vértices son \(V=\{-1,1\}^3\). La matriz de incidencia
\(M:\mathbb R^V\longrightarrow H_F\) queda determinada por
\[
 M_{(i,\epsilon),\nu}=\mathbf1_{\{\nu_i=\epsilon\}}.
\]

\begin{proposition}[Descomposición de la incidencia]
\label{vii:prop:incidencia}
Se tiene
\[
 MM^{\mathsf T}=12P_u+4P_-,
 \qquad
 \ker M^{\mathsf T}=P_0H_F.
\]
Por consiguiente, el cociente
\(Q=H_F/P_0H_F\), identificado ortogonalmente con
\(P_\square H_F\), tiene dimensión cuatro.
\end{proposition}

\begin{proof}
Una cara contiene cuatro vértices. Dos caras opuestas carecen de vértices
comunes y dos caras de ejes diferentes tienen dos vértices comunes.
Escribiendo \(\mathbf J=uu^{\mathsf T}\), resulta
\[
 MM^{\mathsf T}=4I+2(\mathbf J-I-R)
 =2(I-R)+2\mathbf J=4P_-+12P_u.
\]
Los tres proyectores \(P_u,P_-,P_0\) son ortogonales y suman \(I\);
sus rangos son, respectivamente, uno, tres y dos. La identidad
\(\|M^{\mathsf T}x\|^2=\langle x,MM^{\mathsf T}x\rangle\)
identifica el núcleo anunciado. El rango complementario es cuatro.
\end{proof}

Sobre \(Q\), la forma
\[
 B=P_--P_u
\]
tiene signatura \((3,1)\). La orientación temporal se fija por \(u\).
Con
\[
 t=\frac{u}{\sqrt6},\qquad
 d_i=\frac{e_{i,+}-e_{i,-}}{\sqrt2},\qquad
 W=\sqrt6P_u+\sqrt2P_-,
\]
se cumple \(We_{i,\epsilon}=t+\epsilon d_i\). Esos seis vectores
son nulos para \(B\), mientras que
\[
 3t+\nu_1d_1+\nu_2d_2+\nu_3d_3
\]
tiene norma cuadrática \(-6\). La misma incidencia satisface
\(MM^{\mathsf T}=2W^*W\) sobre \(Q\). Las rotaciones propias del cubo
conmutan con los proyectores, preservan ambas formas y fijan \(t\).
Estas identidades determinan la métrica de la publicación cúbica.

\subsection{Soldadura celular y transporte de memoria}
\label{vii:sec:soldadura-celular}

Para cada arista orientada \(a\), la publicación de la ruta proporciona
la cara \(\lambda(\underline a)\), el signo
\(\sigma(a)\), las fibras de origen y destino y el transporte
\(U(a):Q_{s(a)}\longrightarrow Q_{t(a)}\).
La inversión verifica \(U(\bar a)=U(a)^{-1}\).
Las semiaristas de frontera conservan su inversión formal y su incidencia
en el borde. Con \(\ell_m=\ell_0\,9^{-m}>0\), la soldadura es
\begin{equation}
 \beta_m(a)=\ell_m\sigma_m(a)
              n_{\lambda(\underline a)}^{s(a)},\qquad
 n_{i,\epsilon}=t+\epsilon d_i.
 \label{vii:eq:soldadura-arista}
\end{equation}
La cara y el signo se transportan junto con la memoria de la ruta.
La convención de inversión exige
\begin{equation}
 \beta_m(\bar a)=-U_m(a)\beta_m(a).
 \label{vii:eq:inversion-soldadura}
\end{equation}
Aplicando dos veces esta relación se recupera \(\beta_m(a)\);
el registro de acarreo conserva, por separado, la composición ordenada.

\begin{proposition}[Naturalidad de la soldadura por refinamiento]
\label{vii:prop:refinamiento-soldadura}
Sea \(a=a_1\cdots a_9\) un refinamiento compatible de una arista.
Supóngase que las nueve incidencias hijas, transportadas a la fibra del
padre, conservan la cara y el signo, y que
\(\ell_{m+1}=\ell_m/9\). Sea
\(\mathcal U_j:Q_{m,s(a)}\longrightarrow Q_{m+1,s(a_j)}\)
el transporte acumulado que incluye la identificación por refinamiento
y el trayecto hasta el origen del hijo \(j\). Entonces
\[
 \beta_m(a)=
 \sum_{j=1}^{9}\mathcal U_j^{-1}\beta_{m+1}(a_j).
\]
La igualdad es compatible con la inversión y con la concatenación
del registro de memoria.
\end{proposition}

\begin{proof}
Cada sumando, expresado en la fibra del origen del padre, es
\((\ell_m/9)\sigma_m(a)n_{\lambda(\underline a)}\).
La suma de nueve términos reproduce la soldadura del padre.
La inversión invierte el orden de los transportes y sustituye cada uno
por su inverso; la relación \eqref{vii:eq:inversion-soldadura} proporciona
el signo global. La memoria se concatena en ese mismo orden, por lo
que el refinamiento conserva tanto la publicación vectorial como su
registro enriquecido.
\end{proof}

\subsection{Respuesta positiva, unicidad y extensión mínima}
\label{vii:sec:respuesta-positiva}

La soldadura lineal de pantalla es
\(\bar\beta_m=\ell_mW|_Q\). Es invertible, con
\[
 \bar\beta_m^{-1}
 =\ell_m^{-1}\left(\frac1{\sqrt6}P_u+\frac1{\sqrt2}P_-\right).
\]
El autovalor \(5/9\) del sector antisimétrico del bloque normalizado
\(B_c/9\), construido en la sección~\ref{vii:sec:bloque-compresion},
determina el coeficiente de
compresión \(q_{\mathrm{vis}}=5/9\). Reservamos esta notación para
distinguirlo del cociente euclídeo \(q_9(n)\) del núcleo aritmético.
Su realización sobre la pantalla utiliza una isometría
\(V_9:Q\longrightarrow\mathcal K\). La contribución visible
\(C_9=q_{\mathrm{vis}}V_9\) y su complemento
\(D_9=\sqrt{1-q_{\mathrm{vis}}^2}\,V_9\) satisfacen
\[
 C_9^*C_9+D_9^*D_9=I_Q,\qquad D_9^*D_9=\frac{56}{81}I_Q.
\]
El coeficiente \(56\) es \(9^2-5^2\); pertenece al defecto cuadrático
de ese lector de compresión.

\begin{theorem}[Respuesta energética determinada por la soldadura]
\label{vii:thm:respuesta-unica}
Existe un único operador \(\Lambda_m:Q\longrightarrow Q\) tal que
\(\bar\beta_m^*\Lambda_m\bar\beta_m=D_9^*D_9\). Es positivo definido y
\begin{equation}
 \Lambda_m=\frac{28}{243\ell_m^2}(P_u+3P_-)
 =\frac{112}{81\ell_m^2}
       (MM^{\mathsf T}|_Q)^{-1}.
 \label{vii:eq:lambda-pantalla}
\end{equation}
Sobre \(H_F\), la identidad correspondiente tiene como segundo miembro
\((56/81)P_\square\).
\end{theorem}

\begin{proof}
La congruencia por la soldadura invertible determina necesariamente
\[
 \Lambda_m=\bar\beta_m^{-*}\frac{56}{81}I_Q\bar\beta_m^{-1}.
\]
Sustituyendo la inversa, el coeficiente de \(P_u\) es
\(56/(486\ell_m^2)=28/(243\ell_m^2)\); el de \(P_-\) es
\(28/(81\ell_m^2)\). Ambos son positivos. Por la
proposición \ref{vii:prop:incidencia},
\((MM^{\mathsf T}|_Q)^{-1}=P_u/12+P_-/4\),
lo que da la segunda expresión. La extensión a seis caras anula
\(P_0H_F\); por ello la identidad se extiende mediante \(P_\square\).
\end{proof}

\begin{proposition}[Caracterización por mínima extensión]
\label{vii:prop:minima-extension}
Para todo \(b\in Q\),
\[
 \langle b,\Lambda_m b\rangle
 =\frac{112}{81\ell_m^2}
   \min_{Mx=b}\|x\|^2.
\]
La extensión minimizante es única:
\(x_{\min}=M^{\mathsf T}(MM^{\mathsf T}|_Q)^{-1}b\).
\end{proposition}

\begin{proof}
El vector anunciado satisface \(Mx_{\min}=b\) y pertenece a
\(\operatorname{im}M^{\mathsf T}=(\ker M)^\perp\).
Toda otra solución es \(x_{\min}+z\), con \(z\in\ker M\); por
ortogonalidad, su norma al cuadrado es \(\|x_{\min}\|^2+\|z\|^2\).
Además,
\[
 \|x_{\min}\|^2
 =\langle b,(MM^{\mathsf T}|_Q)^{-1}b\rangle.
\]
La fórmula de \(\Lambda_m\) concluye la prueba.
\end{proof}

El operador \(MM^{\mathsf T}|_Q\) realiza la respuesta
Dirichlet--a--Neumann de esta formulación de incidencia; su inversa
es la respuesta Neumann--a--Dirichlet. La energía anterior corresponde
a esta última, con la normalización de la soldadura y del defecto de
compresión. El complemento de Schur de
\(\left(\begin{smallmatrix}I&M^{\mathsf T}\\M&0\end{smallmatrix}\right)\)
es \(-MM^{\mathsf T}\), lo que ofrece la misma eliminación variacional.

\subsection{Conservación de la densidad por refinamiento}
\label{vii:sec:densidad-refinamiento}

Consideremos el sector de transportes cúbicos, ortogonales para el producto
escalar energético y compatibles con la forma \(B\). Fijamos un
espacio de etiquetas común \(Q^{\mathrm{lab}}\), identificado
isométricamente con la pantalla del padre. Las soldaduras tienen
los tipos
\(\bar\beta_m:Q^{\mathrm{lab}}\longrightarrow Q_m\) y
\(\bar\beta_{m+1,j}:Q^{\mathrm{lab}}\longrightarrow Q_{m+1,j}\);
los operadores \(\Lambda\) actúan en sus respectivas fibras de llegada.
Si
\(U_j:Q_j\longrightarrow Q\) transporta la fibra hija a la del padre,
entonces
\[
 \bar\beta_{m+1,j}=\frac19U_j^{-1}\bar\beta_m,\qquad
 \Lambda_{m+1,j}=81U_j^{-1}\Lambda_mU_j.
\]
El factor \(81\) procede de \(\ell_{m+1}^{-2}=81\ell_m^{-2}\).

\begin{theorem}[Invariancia de la densidad normalizada]
\label{vii:thm:densidad}
Cada hijo conserva la respuesta cuadrática del padre:
\[
 \bar\beta_{m+1,j}^*\Lambda_{m+1,j}\bar\beta_{m+1,j}
 =\frac{56}{81}I_{Q^{\mathrm{lab}}}.
\]
Para \(r\) refinamientos, la respuesta conjunta, expresada en las fibras
del padre, es
\[
 E_{m+r}=\frac{56}{81}I_{Q^{\mathrm{lab}}}\otimes I_{9^r}.
\]
La expectativa con traza normalizada
\(\tau_{9^r}(A)=9^{-r}\operatorname{Tr}(A)\) satisface
\[
 (\operatorname{id}\otimes\tau_{9^r})(E_{m+r})
 =\frac{56}{81}I_{Q^{\mathrm{lab}}}
\]
a toda profundidad finita.
\end{theorem}

\begin{proof}
La congruencia de un hijo contiene los factores \(1/9,81,1/9\);
su producto es uno. La ortogonalidad cancela los transportes
intermedios. La suma directa de nueve respuestas iguales identifica
el primer refinamiento con \((56/81)I_Q\otimes I_9\).
La iteración produce el tensor con \(I_{9^r}\).
Finalmente \(\tau_{9^r}(I_{9^r})=1\). Las expectativas se componen
por la identidad
\(\tau_{9^{r+s}}=\tau_{9^r}\otimes\tau_{9^s}\).
\end{proof}

Las inclusiones unitales \(A\mapsto A\otimes I_9\) son compatibles con
esa familia de trazas. En el límite inductivo de las álgebras matriciales,
la respuesta es el elemento central
\((56/81)I_Q\otimes\mathbf1\). La suma extensiva sin normalizar produce,
en cambio, \(9^r(56/81)I_Q\). Esta distinción especifica qué magnitud
se conserva: la densidad cuadrática de pantalla.

\subsection{Realización variacional de la respuesta}
\label{vii:sec:realizacion-respuesta}

La realización material usa una isometría
\(J_{\mathrm{scr}}:Q\longrightarrow\mathcal H_\partial\)
compatible con la acción de espín y con el refinamiento.
La cotetrada de frontera es \(e=J_{\mathrm{scr}}\bar\beta_m\).
Sea \(\Lambda_{\mathrm{CE}}\) la respuesta de frontera obtenida
por segunda variación de la acción material y geométrica, con sus
condiciones de contorno y la eliminación de las variables interiores.
El defecto de compatibilidad es
\begin{equation}
 \mathcal R_{\mathrm{grav}}
 =\bar\beta_m^*
   (J_{\mathrm{scr}}^*\Lambda_{\mathrm{CE}}J_{\mathrm{scr}}-\Lambda_m)
   \bar\beta_m
 =e^*\Lambda_{\mathrm{CE}}e-D_9^*D_9.
 \label{vii:eq:defecto-respuesta}
\end{equation}
Por la invertibilidad de \(\bar\beta_m\),
\[
 \mathcal R_{\mathrm{grav}}=0
 \quad\Longleftrightarrow\quad
 J_{\mathrm{scr}}^*\Lambda_{\mathrm{CE}}J_{\mathrm{scr}}=\Lambda_m.
\]
La isometría \(J_{\mathrm{scr}}\) transporta los estados de frontera;
la corriente variacional de espín es una forma de grado tres con valores
en bivectores. Cada una interviene con su dominio propio en la
realización de la respuesta.

La construcción determina una forma energética positiva, su extensión
interior minimizante y la densidad conservada por refinamiento.
La amplitud de una contribución material concreta se obtiene al evaluar
la corriente de esa realización en su forma cuadrática variacional.
Así quedan diferenciados el operador que fija la respuesta, la corriente
sobre la que actúa y el observable gravitatorio que publica la evaluación.

% RESULTADO_RECUPERADO: S15, 11c_gravedad_holonomia_rev6.tex:128--212;
% S01, 02_dinamica_tres_hojas.tex:155--219; pantalla y respuesta de30.
% Pruebas de curvatura, covariancia y Bianchi reunidas para autonomía.
% La realización suave conserva sus condiciones de holonomía explícitas.

\section{Realización geométrica del transporte y de la soldadura}
\label{vii:sec:realizacion-cartan}

La representación discreta del retorno y la soldadura de pantalla
determinan dos aspectos del transporte: la acumulación de memoria en
las rutas y la comparación de sus fibras. Su realización geométrica
conserva ambas operaciones. En esta sección se fijan el espacio de
llegada, la normalización dimensional de la conexión y las identidades
que utilizará su variación.

\subsection{Pantalla, cotétrada y conexión}

El cociente de pantalla \(Q_\square\), su forma lorentziana y la
soldadura \(\beta_m\) proceden de
\ref{vii:sec:pantalla-respuesta}. En cada celda, una realización de
pantalla \(J_{{\rm scr},m}\) transporta la forma y la soldadura
a la fibra geométrica correspondiente. Su imagen celular es
\[
 e_m=J_{{\rm scr},m}\beta_m.
\]
La escala \(\ell_m\) permanece dentro de la soldadura. Los
cambios de marco actúan sobre \(J_{{\rm scr},m}\), la conexión
y el transporte de las celdas conjuntamente.

Una realización suave de esos datos sobre una variedad orientada
\(\mathcal U\) de dimensión cuatro consiste en una cotétrada
no degenerada \(e=(e^I)\), una conexión de Lorentz
\(\omega=(\omega^I{}_J)\) y una realización de rutas
\(\mathfrak R_{\rm C}\) que satisfacen
\begin{equation}
 \operatorname{Hol}_{\mathcal A_{\ell_0}}
       (\mathfrak R_{\rm C}(\gamma))
 =\rho_{\rm C}(\operatorname{Hol}_{\rm TPK}(\gamma)).
 \label{vii:eq:compatibilidad-holonomia-geometrica}
\end{equation}
En esta igualdad matricial, la representación se normaliza mediante
\[
 N_{\ell_0}(R,a)=(\Lambda(R),a/\ell_0),\qquad
 \rho_{\rm C}=N_{\ell_0}\circ\rho_{\rm C}^{\rm disc},
\]
donde \(\Lambda:\operatorname{Spin}^+(3,1)\to SO^+(3,1)\)
es la representación vectorial. La aplicación \(N_{\ell_0}\)
preserva el producto semidirecto: dividir la traslación
\(a+\Lambda(R)b\) por \(\ell_0\) equivale a componer las
dos traslaciones normalizadas. Por ello la vuelta completa publica
\(72u\) en la conexión normalizada y \(72\ell_0u\) en la
carta dimensional. El levantamiento de espín conserva, por separado,
\(R_4^{c_g(\gamma)}\); la representación vectorial tiene el
núcleo central \(\{1,-1\}\), de modo que se retiene ese
levantamiento al utilizar campos espinoriales.
Identidad, concatenación, inversión y subdivisión se mantienen en la
realización correspondiente. El emparejamiento de
pantalla y la condición de holonomía especifican los datos de la
realización utilizada en las ecuaciones de campo.

Con \(\eta=\operatorname{diag}(-1,1,1,1)\), se define
\(g=\eta_{IJ}e^I\otimes e^J\). Para una escala elemental
positiva \(\ell_0\), constante en la carta, la conexión afín es
\begin{equation}
 \mathcal A_{\ell_0}
 =\begin{pmatrix}\omega&\ell_0^{-1}e\\0&0\end{pmatrix}.
 \label{vii:eq:conexion-afin-realizada}
\end{equation}
La cotétrada tiene dimensión de longitud; el factor
\(\ell_0^{-1}\) homogeneiza su componente traslacional.

\begin{proposition}[Curvatura y torsión de la conexión realizada]
\label{vii:prop:curvatura-afin}
La curvatura de \eqref{vii:eq:conexion-afin-realizada} es
\begin{equation}
 \mathcal F_{\ell_0}
 =\mathrm d\mathcal A_{\ell_0}
   +\mathcal A_{\ell_0}\wedge\mathcal A_{\ell_0}
 =\begin{pmatrix}F_\omega&\ell_0^{-1}T\\0&0\end{pmatrix},
 \quad F_\omega=\mathrm d\omega+\omega\wedge\omega,
 \quad T=\mathrm de+\omega\wedge e.
 \label{vii:eq:curvatura-afin-realizada}
\end{equation}
\end{proposition}
\begin{proof}
El producto de matrices de formas produce en el bloque diagonal
\(\omega\wedge\omega\) y en el traslacional
\(\ell_0^{-1}\omega\wedge e\). La derivada exterior aporta
\(\mathrm d\omega\) y \(\ell_0^{-1}\mathrm de\), pues
\(\ell_0\) es constante. La suma demuestra la igualdad.
\end{proof}

La curvatura describe el defecto de transporte de marcos y la torsión
el de la cotétrada. Ambas pertenecen a la misma conexión realizada;
sus componentes conservan dominios tensoriales distintos.

\begin{proposition}[Covariancia e identidades de Bianchi]
\label{vii:prop:bianchi-cartan}
Bajo un cambio de marco \(g_L:\mathcal U\to SO^+(3,1)\),
\[
 e'=g_Le,\qquad
 \omega'=g_L\omega g_L^{-1}-\mathrm dg_L\,g_L^{-1},
\]
se cumple \(T'=g_LT\) y \(F_{\omega'}=g_LF_\omega g_L^{-1}\).
Además,
\begin{equation}
 D_\omega T=F_\omega\wedge e,\qquad D_\omega F_\omega=0.
 \label{vii:eq:bianchi-realizacion}
\end{equation}
\end{proposition}
\begin{proof}
En \(\mathrm d(g_Le)+\omega'\wedge g_Le\), los términos
que contienen \(\mathrm dg_L\wedge e\) se cancelan.
La misma sustitución en la curvatura da su transformación por
conjugación. Para la primera identidad de Bianchi, se expande
\[
 D_\omega(\mathrm de+\omega\wedge e)
 =(\mathrm d\omega+\omega\wedge\omega)\wedge e.
\]
En \(\mathrm dF_\omega+\omega\wedge F_\omega-F_\omega\wedge\omega\),
los términos que contienen \(\mathrm d\omega\) y los productos
cúbicos se cancelan por pares. Esto prueba la segunda identidad.
\end{proof}

\subsection{Las tres realizaciones cuadráticas de curvatura}

La hoja \(\tau\in\{+1,0,-1\}\) del TRIT se representa en
el álgebra de Cartan mediante generadores \(J_{ab}=-J_{ba}\)
y \(P_a^{(\tau)}\), con relaciones
\begin{align*}
 [J_{ab},J_{cd}]&=
 \eta_{bc}J_{ad}-\eta_{ac}J_{bd}
 -\eta_{bd}J_{ac}+\eta_{ad}J_{bc},\\
 [J_{ab},P_c^{(\tau)}]&=
 \eta_{bc}P_a^{(\tau)}-\eta_{ac}P_b^{(\tau)},\\
 [P_a^{(\tau)},P_b^{(\tau)}]&=-\tau J_{ab}.
\end{align*}
Para \(\ell>0\) constante, se escribe
\[
 \mathcal A_\tau^{\rm Cart}
 =\tfrac12\omega^{ab}J_{ab}+\ell^{-1}e^aP_a^{(\tau)}.
\]

\begin{proposition}[Curvatura de la realización tritificada]
\label{vii:prop:curvatura-tritificada}
Con las relaciones anteriores,
\begin{equation}
 \mathcal F_\tau^{\rm Cart}
 =\tfrac12(F_\omega^{ab}-\tau\ell^{-2}e^a\wedge e^b)J_{ab}
     +\ell^{-1}T^aP_a^{(\tau)}.
 \label{vii:eq:curvatura-tritificada}
\end{equation}
\end{proposition}
\begin{proof}
Se expande \(\mathrm d\mathcal A+\tfrac12[\mathcal A,\mathcal A]\).
El corchete del sector de Lorentz produce \(F_\omega\), el
mixto produce \(D_\omega e\), y el traslacional aporta
\(-\tfrac12\tau\ell^{-2}e^a\wedge e^bJ_{ab}\).
\end{proof}

En la hoja central se recupera la conexión afín; las hojas extremas
conservan los dos signos del término de curvatura constante. La
representación mantiene así el régimen trítico de la ruta antes de
aplicar las ecuaciones variacionales. La igualdad de holonomía
\eqref{vii:eq:compatibilidad-holonomia-geometrica} corresponde a la
hoja central. En cada hoja extrema se utiliza la representación
\(\rho_{{\rm C},\tau}\) en su grupo de Cartan y su condición
de holonomía respectiva. La corriente y el tensor métrico utilizan
los campos y la representación de la misma realización.

% Artículo VII. Desarrollo aditivo, revisión 04, 10-09-2026.
% ARQUITECTURA_AUTORAL_PREEXISTENTE: transporte enriquecido, energía
% D_9^* W D_9 y compatibilidad por refinamiento, propietario S01,
% 02_dinamica_tres_hojas.tex, líneas 221--250.
% FORMALIZACION_NUEVA / CERTIFICADO_NUEVO: diferencial explícito de ese
% funcional, representante antihermítico, composición e inversión ponderada.
% Las variaciones pertenecen al espacio de realizaciones unitarias del
% transporte. No se presume que toda variación continua sea una ruta TPK.
% Su realización como corriente material de espín se trata por separado.
% \mathsf W_m es el peso energético; W de la sección anterior es soldadura.

\section{Respuesta variacional de la energía de memoria}
\label{vii:sec:variacion-memoria}

El transporte de una ruta identifica las fibras de sus extremos y conserva
la orientación y la memoria que intervienen en esa identificación. La energía
asociada compara el campo final con el campo inicial transportado. Su
diferencial cuantifica cómo responde esa comparación a una variación del
transporte. Esta sección obtiene la respuesta directamente del funcional
discreto y desarrolla su composición, covariancia e inversión.

\subsection{Funcional discreto y espacio de variaciones}
\label{vii:subsec:funcional-memoria}

En un nivel finito \(m\), sean \(\mathcal V_m\) los vértices de la
publicación de rutas y \(\mathcal E_m^+\) una orientación de cada arista.
Cada vértice lleva una fibra hermítica finito-dimensional \(H_v\).
Definimos
\[
 \mathcal H_m=\bigoplus_{v\in\mathcal V_m}H_v,\qquad
 \mathcal K_m=\bigoplus_{a\in\mathcal E_m^+}H_{t(a)}.
\]
El producto hermítico es antilineal en la primera variable.
Para \(U_a:H_{s(a)}\to H_{t(a)}\) unitario, el operador de diferencias es
\begin{equation}
 (D_mf)_a=f_{t(a)}-U_af_{s(a)},\qquad
 E_m(f,U)=\langle D_mf,\mathsf W_mD_mf\rangle,\qquad
 \mathsf W_m=\mathsf W_m^*>0.
 \label{vii:eq:energia-memoria-discreta}
\end{equation}
Es la energía del término \(D_m^*\mathsf W_mD_m\) del generador discreto.
El peso \(\mathsf W_m\) puede acoplar distintas aristas. El operador \(W\)
de la soldadura de la sección~\ref{vii:sec:pantalla-respuesta} y este peso
energético tienen dominios y funciones diferentes.

El transporte \(U\) es una publicación del estado enriquecido. El campo
\(f\) y el peso \(\mathsf W_m\) permanecen explícitos como argumentos
del funcional; cada aplicación material conserva las reglas que los
determinan. Se consideran familias diferenciables de las realizaciones
unitarias del transporte,
\[
 \dot U_a=A_aU_a,\qquad A_a^*=-A_a.
\]
Las variaciones físicamente utilizadas forman el subespacio que conserve
las condiciones de la realización correspondiente. El diferencial se
calcula primero en el espacio tangente unitario y después se restringe
a ese subespacio.

\begin{proposition}[Covector de respuesta al transporte]
\label{vii:prop:covector-memoria}
Escribamos
\[
 v_a=U_af_{s(a)},\qquad r=D_mf,\qquad q=\mathsf W_mr.
\]
A campo y peso fijos, el diferencial de la energía es
\begin{equation}
 \mathfrak j_m(A)=-2\operatorname{Re}
       \sum_{a\in\mathcal E_m^+}\langle q_a,A_av_a\rangle .
 \label{vii:eq:covector-memoria}
\end{equation}
Su representante para el producto real de Hilbert--Schmidt es
\begin{equation}
 \mathcal J_a=v_aq_a^*-q_av_a^*,\qquad
 \mathcal J_a^*=-\mathcal J_a,\qquad
 \mathfrak j_m(A)=
 \sum_a\operatorname{Re}\operatorname{Tr}(\mathcal J_a^*A_a).
 \label{vii:eq:representante-corriente-discreta}
\end{equation}
Para variaciones simultáneas de campo, transporte y peso se tiene
\begin{equation}
 \dot E_m=
 2\operatorname{Re}\sum_a
 \left\langle q_a,\dot f_{t(a)}-U_a\dot f_{s(a)}-A_av_a\right\rangle
 +\langle r,\dot{\mathsf W}_mr\rangle .
 \label{vii:eq:variacion-completa-memoria}
\end{equation}
\end{proposition}

\begin{proof}
La actualización de primer orden, escrita en el álgebra de números
duales con \(\epsilon^2=0\), da
\[
 r+\epsilon\dot r,\qquad
 \dot r_a=\dot f_{t(a)}-U_a\dot f_{s(a)}-A_av_a.
\]
El coeficiente de \(\epsilon\) en la energía es
\[
 \langle\dot r,\mathsf W_mr\rangle+
 \langle r,\mathsf W_m\dot r\rangle+
 \langle r,\dot{\mathsf W}_mr\rangle.
\]
La autoadjunción del peso produce
\eqref{vii:eq:variacion-completa-memoria}. Al fijar campo y peso queda
\eqref{vii:eq:covector-memoria}. Además,
\[
 \operatorname{Tr}(\mathcal J_a^*A_a)
 =v_a^*A_aq_a-q_a^*A_av_a.
\]
La antihermiticidad de \(A_a\) implica
\(v_a^*A_aq_a=-\overline{q_a^*A_av_a}\), de donde se obtiene
el representante anunciado.
\end{proof}

Así queda evaluada la respuesta: en cada arista se transporta el campo,
se calcula su diferencia con el campo final y se aplica el peso energético.
El producto antisimetrizado de esos dos vectores determina
\(\mathcal J_a\in\mathfrak u(H_{t(a)})\). Este operador representa,
mediante el producto real de Hilbert--Schmidt, el covector de respuesta
sobre las variaciones unitarias de la fibra de destino.

\subsection{Composición y transporte de la respuesta}
\label{vii:subsec:composicion-respuesta}

\begin{proposition}[Regla de composición sobre una ruta]
\label{vii:prop:composicion-respuesta}
Para una ruta formada por \(N\) transportes, sea
\(U_\gamma=U_N\cdots U_1\) y
\(B_k=U_N\cdots U_{k+1}\), con \(B_N=I\).
Si \(\dot U_k=A_kU_k\), entonces
\begin{equation}
 \dot U_\gamma U_\gamma^{-1}
 =\sum_{k=1}^N B_kA_kB_k^{-1}.
 \label{vii:eq:variacion-transporte-compuesto}
\end{equation}
Por tanto, un representante \(\mathcal J_\gamma\) en la fibra final
induce en el factor \(k\) la contribución
\(B_k^*\mathcal J_\gamma B_k\). Cuando un factor pertenece a varias
rutas del funcional, sus contribuciones se suman con las incidencias
que éste utiliza.
\end{proposition}

\begin{proof}
La regla del producto da
\[
 \dot U_\gamma=\sum_k
 U_N\cdots U_{k+1}A_kU_k\cdots U_1.
\]
Multiplicar por \(U_\gamma^{-1}\) elimina los factores a la derecha
de \(A_k\) y produce \eqref{vii:eq:variacion-transporte-compuesto}.
Por unitariedad y ciclicidad de la traza,
\[
 \operatorname{Re}\operatorname{Tr}
   (\mathcal J_\gamma^*B_kA_kB_k^{-1})
 =\operatorname{Re}\operatorname{Tr}
   ((B_k^*\mathcal J_\gamma B_k)^*A_k).
\]
La linealidad del diferencial demuestra la suma de contribuciones.
\end{proof}

Esta fórmula conserva la posición de cada operación dentro de la ruta.
El transporte posterior a una incidencia determina cómo se representa
su contribución en la fibra final; la respuesta local se recupera
transportándola a la fibra de esa incidencia.

\begin{proposition}[Covariancia bajo cambios unitarios de marco]
\label{vii:prop:covariancia-corriente-memoria}
Sean \(g_v\) cambios unitarios de marco, independientes del parámetro
de variación, y \(G_{\mathcal E}=\bigoplus_a g_{t(a)}\).
Bajo las transformaciones
\[
 f'_v=g_vf_v,\quad U'_a=g_{t(a)}U_ag_{s(a)}^{-1},\quad
 \mathsf W'_m=G_{\mathcal E}\mathsf W_mG_{\mathcal E}^{-1},\quad
 A'_a=g_{t(a)}A_ag_{t(a)}^{-1},
\]
se tiene
\begin{equation}
 \mathcal J'_a=g_{t(a)}\mathcal J_ag_{t(a)}^{-1},
 \qquad \mathfrak j'_m(A')=\mathfrak j_m(A).
 \label{vii:eq:covariancia-respuesta}
\end{equation}
\end{proposition}

\begin{proof}
Se obtiene \(r'=G_{\mathcal E}r\), \(q'=G_{\mathcal E}q\) y
\(v'_a=g_{t(a)}v_a\). La sustitución en
\eqref{vii:eq:representante-corriente-discreta} demuestra la primera
identidad. La invariancia del producto hermítico, o la ciclicidad de
la traza, demuestra la segunda.
\end{proof}

Cuando los cambios de marco dependan del parámetro, sus derivadas
actúan también sobre el campo y el peso. La expresión covariante
correspondiente es entonces el diferencial completo
\eqref{vii:eq:variacion-completa-memoria}.

\subsection{Inversión orientada y transporte del peso}
\label{vii:subsec:inversion-respuesta}

La inversión puede verse explícitamente en el caso de una energía
local por arista, \(\mathsf W_m=\bigoplus_a W_a\).
Para \(a:s\to t\), su inversa lleva
\[
 U_{\bar a}=U_a^{-1},\qquad
 r_{\bar a}=-U_a^{-1}r_a,\qquad
 W_{\bar a}=U_a^{-1}W_aU_a.
\]
Estas reglas conservan exactamente la energía de la arista.

\begin{proposition}[Compatibilidad diferencial con la inversión]
\label{vii:prop:inversion-respuesta}
Si \(W_a\) y los campos de ambos extremos permanecen fijos, entonces
\[
 A_{\bar a}=-U_a^{-1}A_aU_a,\qquad
 \dot W_{\bar a}=U_a^{-1}[W_a,A_a]U_a.
\]
La variación completa de la energía de la arista invertida coincide con
la original. En particular, si \([W_a,A_a]=0\), el covector parcial
de transporte satisface
\begin{equation}
 \mathfrak j_{\bar a}(U_a^{-1}A_aU_a)
 =-\mathfrak j_a(A_a).
 \label{vii:eq:imparidad-transporte}
\end{equation}
\end{proposition}

\begin{proof}
La derivada de \(U_a^{-1}\) es \(-U_a^{-1}A_a\), de donde resultan
las dos fórmulas. Escribiendo \(v_a=U_af_s\) y \(f_t=r_a+v_a\),
la parte debida exclusivamente al transporte invertido vale
\[
 \mathfrak j_{\bar a}(A_{\bar a})
 =-2\operatorname{Re}\langle r_a,W_aA_af_t\rangle
 =\mathfrak j_a(A_a)
   -2\operatorname{Re}\langle r_a,W_aA_ar_a\rangle.
\]
La variación del peso aporta
\[
 \langle r_{\bar a},\dot W_{\bar a}r_{\bar a}\rangle
 =\langle r_a,[W_a,A_a]r_a\rangle
 =2\operatorname{Re}\langle r_a,W_aA_ar_a\rangle.
\]
Ambos términos se compensan. Si el conmutador es nulo, esa contribución
desaparece; la linealidad en \(A_a\) da
\eqref{vii:eq:imparidad-transporte}.
\end{proof}

La energía se suma sobre una orientación por arista. Una suma sobre ambas
orientaciones lleva el factor \(1/2\). Esta convención conserva el mismo
funcional y, por tanto, la misma respuesta.

\subsection{Compatibilidad con el refinamiento}
\label{vii:subsec:refinamiento-respuesta}

\begin{theorem}[Transporte del diferencial por refinamiento]
\label{vii:thm:refinamiento-respuesta}
Sean \(J_m:\mathcal H_m\to\mathcal H_{m+1}\) y
\(K_m:\mathcal K_m\to\mathcal K_{m+1}\) aplicaciones fijas.
Para una familia \(C^1\), supongamos las identidades de refinamiento
\begin{equation}
 D_{m+1}(t)J_m=K_mD_m(t),\qquad
 K_m^*\mathsf W_{m+1}(t)K_m=\mathsf W_m(t).
 \label{vii:eq:familia-refinamiento-memoria}
\end{equation}
Entonces \(E_{m+1}(J_mf)=E_m(f)\).
Los covectores de transporte coinciden sobre las variaciones enlazadas
por la primera identidad. Cuando varían los pesos, sus contribuciones
al diferencial también coinciden.
\end{theorem}

\begin{proof}
Por sustitución,
\[
 \langle D_{m+1}J_mf,\mathsf W_{m+1}D_{m+1}J_mf\rangle
 =\langle D_mf,K_m^*\mathsf W_{m+1}K_mD_mf\rangle=E_m(f).
\]
Diferenciando la primera identidad de
\eqref{vii:eq:familia-refinamiento-memoria}, se obtiene
\(\dot D_{m+1}J_m=K_m\dot D_m\). Por consiguiente,
\[
 2\operatorname{Re}
 \langle D_{m+1}J_mf,\mathsf W_{m+1}\dot D_{m+1}J_mf\rangle
 =2\operatorname{Re}\langle D_mf,\mathsf W_m\dot D_mf\rangle.
\]
La derivada de la segunda identidad es
\(K_m^*\dot{\mathsf W}_{m+1}K_m=\dot{\mathsf W}_m\), que demuestra
la coincidencia de los términos de variación del peso.
\end{proof}

El enunciado utiliza la compatibilidad de la familia de transportes.
Si \(J_m\) depende del parámetro, la regla de la cadena conserva
también el término
\[
 2\operatorname{Re}
 \langle D_{m+1}J_mf,\mathsf W_{m+1}D_{m+1}\dot J_mf\rangle.
\]
La naturalidad obtenida transporta conjuntamente el funcional, sus
incidencias y su diferencial.

\subsection{Realización material de la respuesta}
\label{vii:subsec:realizacion-respuesta-discreta}

El resultado de esta sección es el covector explícito
\(\mathfrak j_m\), con su representante \(\mathcal J_a\).
La realización material compone ese covector con la variación del
transporte inducida por la conexión y con la normalización de la acción.
Si \(\mathcal L_m\) denota ese mapa diferencial, el covector sobre
variaciones de conexión es \(\mathcal L_m^*\mathfrak j_m\).
La regla de composición
\eqref{vii:eq:variacion-transporte-compuesto} evalúa su parte
correspondiente a productos finitos de transportes.

En la sección~\ref{vii:sec:corriente-respuesta-cartan}, la corriente
\(\sigma_{IJ}\) es una tres-forma definida por la variación de la acción
material. La identificación entre ambas realizaciones conserva el campo,
la representación, el peso, la normalización dimensional y el paso al
límite utilizados por esa acción. Estos datos permiten evaluar la
contracción que interviene en la densidad torsional. El desarrollo
presente aporta el diferencial discreto que se compone en ese cálculo.

% Artículo VII, incorporación aditiva, revisión 05.
% Arquitectura autoral: energía de memoria APP--TRIT--TPK, S01.
% Antecedente recuperado: diferencial y refinamiento de 30b.
% Formalización reunida: nota COMPOSICION_VARIACIONAL_TRIT_CONEXION.md.
% Se conserva íntegra 30b; su representante antihermítico es la
% restricción unitaria del diferencial que aquí se desarrolla.
% Esta sección evalúa la corriente de una acción y una realización
% declaradas. No selecciona retrospectivamente su campo ni su peso.

\section{Realización de la corriente en coordenadas de conexión}
\label{vii:sec:composicion-corriente-conexion}

La energía de memoria asigna a cada ruta una respuesta que conserva sus
incidencias y el orden de los transportes. Para insertarla en una acción
geométrica se necesitan sus componentes respecto de las variaciones de
conexión. El paso se obtiene diferenciando el transporte que ya realiza
la ruta: la corriente de conexión es la composición de ese diferencial
con el covector energético de la sección anterior.

Los regímenes elíptico, parabólico e hiperbólico del TRIT requieren
conservar el diferencial completo antes de restringirlo a una
representación particular. Las rotaciones unitarias constituyen un
dominio de esa construcción. La ampliación siguiente conserva también
las direcciones de dilatación y transporte que admite la realización
geométrica utilizada.

\subsection{Diferencial sobre transportes invertibles}

Se mantienen las fibras, los campos y el peso de
\eqref{vii:eq:energia-memoria-discreta}. Consideremos ahora transportes
invertibles y familias diferenciables con
\(\dot U_a=A_aU_a\), donde \(A_a\) pertenece al espacio tangente
de la representación elegida. Escribimos
\[
 v_a=U_af_{s(a)},\qquad r_a=f_{t(a)}-v_a,\qquad
 q=\mathsf W_mr,\qquad E_m=\langle r,\mathsf W_mr\rangle.
\]

\begin{proposition}[Representante completo del diferencial]
\label{vii:prop:diferencial-invertible}
A campo y peso fijos, se cumple
\begin{equation}
 \mathrm d_UE_m(A)=-2\operatorname{Re}\sum_a
       \langle q_a,A_av_a\rangle
 =\sum_a\operatorname{Re}\operatorname{Tr}(\mathcal G_a^*A_a),
 \qquad \mathcal G_a=-2q_av_a^*.
 \label{vii:eq:gradiente-transporte-completo}
\end{equation}
Las proyecciones antihermítica y hermítica de ese representante son
\begin{equation}
 \frac{\mathcal G_a-\mathcal G_a^*}{2}
   =v_aq_a^*-q_av_a^*=\mathcal J_a,
 \qquad
 \frac{\mathcal G_a+\mathcal G_a^*}{2}
   =-(q_av_a^*+v_aq_a^*).
 \label{vii:eq:proyecciones-gradiente-memoria}
\end{equation}
La restricción a generadores antihermíticos recupera exactamente
\eqref{vii:eq:representante-corriente-discreta}.
\end{proposition}

\begin{proof}
La identidad \(\dot r_a=-A_av_a\) y la autoadjunción del peso dan
\(\dot E_m=2\operatorname{Re}\langle q,\dot r\rangle\).
Además,
\[
 \operatorname{Tr}(\mathcal G_a^*A_a)
 =-2\operatorname{Tr}(v_aq_a^*A_a)=-2q_a^*A_av_a.
\]
Las dos proyecciones se obtienen tomando el adjunto. Para el producto
real de Hilbert--Schmidt, los subespacios hermítico y antihermítico son
ortogonales; por ello la parte hermítica tiene contracción nula con
un generador antihermítico.
\end{proof}

La fórmula permite pesos que acoplan aristas: su acción se calcula en
\(q=\mathsf W_mr\) antes de separar las contribuciones de cada
incidencia. Si varían campo o peso, se conserva la expresión completa
\eqref{vii:eq:variacion-completa-memoria}, válida también para los
transportes invertibles aquí considerados.

\subsection{Representaciones cuadráticas del TRIT}

En una realización matricial bidimensional de
\(\mathbb A_\tau\), pueden tomarse
\[
 J_{+}=\begin{pmatrix}0&1\\-1&0\end{pmatrix},\qquad
 J_{-}=\begin{pmatrix}0&1\\1&0\end{pmatrix},\qquad
 J_{0}=\begin{pmatrix}0&1\\0&0\end{pmatrix}.
\]
Por multiplicación directa,
\[
 J_+^2=-I,\qquad J_-^2=I,\qquad J_0^2=0.
\]
Sus exponenciales realizan, respectivamente, los regímenes elíptico,
hiperbólico y parabólico. El primero es ortogonal para el producto
positivo fijado. El segundo preserva la forma
\(B=\operatorname{diag}(-1,1)\), pues
\(J_-^{\mathsf T}B+BJ_-=0\). El tercero representa el álgebra
de números duales. Estas representaciones fijan ejemplos de los tipos
cuadráticos; su uso en una pantalla de otra dimensión conserva el mapa
de representación correspondiente.

Un cálculo distingue las componentes del diferencial. En una arista,
tómense \(U=I\), \(f_s=(1,1)^{\mathsf T}\),
\(f_t=2f_s\) y \(\mathsf W=I\). Entonces \(q=v=f_s\),
\(\mathcal J=0\), mientras
\begin{equation}
 \mathrm dE(J_-)=-4,\qquad \mathrm dE(J_0)=-2.
 \label{vii:eq:control-regimenes-corriente}
\end{equation}
El representante completo conserva ambas respuestas. La distinción
identifica el espacio de variaciones antes de efectuar la contracción
material.

\subsection{Composición de rutas e inversión orientada}

\begin{proposition}[Transporte contragrediente del diferencial]
\label{vii:prop:pullback-invertible}
Para \(U_\gamma=U_N\cdots U_1\), sean
\(B_k=U_N\cdots U_{k+1}\) y \(B_N=I\). Se cumple
\[
 \dot U_\gamma U_\gamma^{-1}
   =\sum_kB_kA_kB_k^{-1}.
\]
El representante \(\mathcal G_\gamma\) en la fibra final induce
en la incidencia \(k\)
\begin{equation}
 \mathcal G_k=B_k^*\mathcal G_\gamma(B_k^{-1})^*.
 \label{vii:eq:transporte-contragrediente}
\end{equation}
\end{proposition}

\begin{proof}
La regla del producto, seguida de multiplicación por
\(U_\gamma^{-1}\), demuestra la primera igualdad. La ciclicidad
de la traza da
\[
 \operatorname{Re}\operatorname{Tr}
  (\mathcal G_\gamma^*B_kA_kB_k^{-1})
 =\operatorname{Re}\operatorname{Tr}
  ((B_k^*\mathcal G_\gamma(B_k^{-1})^*)^*A_k).
\]
Para \(B_k\) unitario se recupera la conjugación de
\ref{vii:prop:composicion-respuesta}.
\end{proof}

\begin{proposition}[Inversión por congruencia energética]
\label{vii:prop:inversion-congruencia}
Para una energía local por arista, la inversión se representa por
\begin{equation}
 U_{\bar a}=U_a^{-1},\qquad
 r_{\bar a}=-U_a^{-1}r_a,\qquad
 W_{\bar a}=U_a^*W_aU_a.
 \label{vii:eq:inversion-peso-general}
\end{equation}
Conserva la energía. Si \(W_a\) permanece fijo,
\[
 A_{\bar a}=-U_a^{-1}A_aU_a,\qquad
 \dot W_{\bar a}=U_a^*(A_a^*W_a+W_aA_a)U_a.
\]
La variación completa de la energía invertida coincide con la original.
\end{proposition}

\begin{proof}
La sustitución en \(r_{\bar a}^*W_{\bar a}r_{\bar a}\)
produce \(r_a^*W_ar_a\). La derivada de la inversa es
\(\dot U_a^{-1}=-U_a^{-1}A_a\); la del peso se obtiene por
la regla del producto. La igualdad de energías vale para toda la
familia, de modo que sus derivadas coinciden, incluyendo el término
\(\langle r_{\bar a},\dot W_{\bar a}r_{\bar a}\rangle\).
\end{proof}

La congruencia conserva positividad para todo transporte invertible.
En la representación unitaria se reduce a la regla de inversión de
\ref{vii:prop:inversion-respuesta}. La suma utiliza una orientación
por arista, o la mitad de la suma sobre ambas orientaciones.

\subsection{Evaluación de las componentes de corriente}

Sea \(\theta=(\theta_b)\) una variación real de las coordenadas
de conexión de la realización utilizada. El diferencial efectivo del
transporte define los operadores
\begin{equation}
 A_a(\theta)=\dot U_aU_a^{-1}
     =\sum_bB_{ab}\theta_b.
 \label{vii:eq:diferencial-conexion-transporte}
\end{equation}
Para una composición finita, sus contribuciones se obtienen por
\eqref{vii:eq:transporte-contragrediente}. Los operadores
\(B_{ab}\) son las derivadas de la representación de conexión que
realiza la ruta.

\begin{theorem}[Corriente de conexión de la acción de memoria]
\label{vii:thm:corriente-conexion-evaluada}
Para un término material \(S_{\rm mem}=\mathfrak a_m E_m\),
con \(\mathfrak a_m\) su normalización en la recta de acción,
a campo, peso y normalización fijos se tiene
\begin{equation}
 \delta S_{\rm mem}=-\frac12\sum_b\theta_b s_b,\qquad
 s_b=4\mathfrak a_m\operatorname{Re}
        \sum_a\langle q_a,B_{ab}v_a\rangle.
 \label{vii:eq:componentes-corriente-conexion}
\end{equation}
Si el emparejamiento celular entre las bases de conexión y corriente
es una matriz invertible \(M\), definida por
\(\delta S_{\rm mem}=-\tfrac12\theta^{\mathsf T}M\sigma\),
las coordenadas de la corriente son
\begin{equation}
 \sigma=M^{-1}s.
 \label{vii:eq:coordenadas-corriente-material}
\end{equation}
\end{theorem}

\begin{proof}
Se sustituye \eqref{vii:eq:diferencial-conexion-transporte} en
\eqref{vii:eq:gradiente-transporte-completo}, se multiplica por
\(\mathfrak a_m\) y se agrupan los coeficientes de
\(\theta_b\). La convención variacional \(-1/2\) fija el
factor \(4\). Comparar ambas expresiones para toda variación
\(\theta\) da \(M\sigma=s\); la invertibilidad del
emparejamiento demuestra existencia y unicidad de esas coordenadas.
\end{proof}

La fórmula realiza explícitamente la composición mencionada en
\ref{vii:subsec:realizacion-respuesta-discreta}. Para evaluar cada
componente se transporta el campo, se calcula el residuo, se aplica el
peso y se contrae con la variación de conexión de la misma realización.
El emparejamiento \(M\) conserva orientación y volumen de las celdas.
Un emparejamiento diagonal da la correspondiente división por sus
volúmenes; otras bases conservan la matriz completa.

Si la conexión interviene también en \(f\), \(\mathsf W_m\)
o \(\mathfrak a_m\), se aplica
\eqref{vii:eq:variacion-completa-memoria} y se añade
\(E_m\delta\mathfrak a_m\). Las contribuciones de todos los
términos materiales se suman antes de componer la respuesta de Cartan.

\begin{corollary}[Compatibilidad de la corriente con el refinamiento]
\label{vii:cor:refinamiento-corriente-conexion}
Para las familias compatibles de
\ref{vii:thm:refinamiento-respuesta}, sea \(P_m\) el mapa fijo
que lleva variaciones de conexión del nivel \(m\) al nivel
\(m+1\). Si las acciones de memoria coinciden sobre esos campos
y transportes refinados, sus covectores satisfacen
\[
 s_m=P_m^{\mathsf T}s_{m+1},\qquad
 M_m\sigma_m=P_m^{\mathsf T}M_{m+1}\sigma_{m+1}.
\]
\end{corollary}

\begin{proof}
Diferenciar la igualdad de acciones en la dirección
\(P_m\theta\) da
\(-\tfrac12\theta^{\mathsf T}s_m
 =-\tfrac12\theta^{\mathsf T}P_m^{\mathsf T}s_{m+1}\).
La arbitrariedad de \(\theta\) y las representaciones celulares
de ambos covectores prueban las dos identidades. La composición de
estos mapas itera la igualdad a todo número finito de refinamientos.
\end{proof}

\subsection{Inserción en la ecuación de Cartan}

La sección~\ref{vii:sec:corriente-respuesta-cartan} normaliza la
tres-forma material mediante
\(\delta_\omega S_m=-\tfrac12\int\delta\omega^{IJ}\wedge
\sigma_{IJ}\). El teorema anterior proporciona sus componentes
celulares para la acción de memoria y el emparejamiento declarados.
La respuesta de Cartan utiliza esa corriente junto a la cotétrada,
el parámetro de Barbero--Immirzi y el acoplamiento gravitatorio de la
misma realización.

La dependencia material respecto de la contorsión determina cómo se
resuelve esa ecuación. Para materia afín, la corriente es independiente
de la contorsión y se aplica la inversión lineal y la eliminación
cuadrática demostradas en la sección siguiente. Para una acción de
holonomía, el diferencial se evalúa en la conexión actual y se conserva
esa dependencia al resolver la ecuación. Por ejemplo, una arista con
\(f_s=f_t=1\), peso uno y \(U(t)=e^{it}\) tiene
\[
 E(t)=2-2\cos t,\qquad E'(t)=2\sin t.
\]
La corriente depende aquí de \(t\). La fórmula completa mantiene
esa dependencia en vez de sustituirla por su valor en una conexión
particular.

Finalmente, la densidad gravitatoria se obtiene de la variación
métrica de la acción efectiva, con su campo y su normalización.
La representación de la corriente, la contorsión que produce y la
densidad resultante son las sucesivas operaciones de esa composición.
La ley de dilución posterior transporta la amplitud así evaluada en
la sección material correspondiente.

% Artículo VII. Incorporación aditiva, 10-09-2026.
% RESULTADO_RECUPERADO: fuentes_conservadas/18_clausura_energetica.tex,
% líneas 357--485: Gramiano, peso, extensión mínima y eliminación de Schur;
% líneas 149--180 y 235--295: normalización y unicidad de la respuesta.
% Residencia anterior en VII: 30_pantalla_y_respuesta.tex,
% vii:prop:incidencia, vii:eq:lambda-pantalla, vii:prop:minima-extension.
% FORMALIZACION_NUEVA / CERTIFICADO_NUEVO: diferencial del mínimo con
% incidencia, frontera y normalización variables; incidencia transportada
% por bloques y composición con coordenadas de conexión.
% La familia de rutas que realiza cada incidencia se mantiene explícita.
% No se atribuyen al propietario histórico ni esa familia diferenciable
% completa ni la evaluación de una amplitud material de espín.
% Dependencias: 30, 30c y la convención variacional de 31. No modifica
% esos fragmentos ni constituye por sí solo una realización material.

\section{Diferencial de la extensión mínima de frontera}
\label{vii:sec:corriente-extension-minima}

La incidencia cúbica determina simultáneamente el campo de extensión
mínima y el peso de su respuesta. Para un dato de frontera, ambos se
calculan mediante el inverso del Gramiano de incidencia. Se desarrolla
ahora el diferencial de esa misma operación, conservando las variaciones
de la frontera, del transporte y de la escala de normalización. De este
modo, el campo que interviene en la respuesta queda seleccionado por
el problema variacional de pantalla.

\subsection{Campo minimizante y peso de respuesta}
\label{vii:subsec:extension-minima-parametrizada}

Sean \(H\) y \(Q\) espacios hermíticos de dimensión finita, con
productos fijos y antilineales en la primera variable. El caso real se
obtiene sustituyendo los adjuntos por transpuestas. En un abierto de
parámetros reales \(\Omega\), consideremos familias de clase \(C^1\)
\[
 \mathsf C(\theta):H\longrightarrow Q,\qquad
 b(\theta)\in Q,\qquad \chi(\theta)>0,
 \qquad \operatorname{Ran}\mathsf C(\theta)=Q.
\]
La sobreyectividad equivale a rango fila completo en bases ortonormales.
Definimos el Gramiano, el potencial conjugado y la extensión
\begin{equation}
 \mathsf L=\mathsf C\mathsf C^*,\qquad
 y=\mathsf L^{-1}b,\qquad f_{\min}=\mathsf C^*y,
 \qquad \Lambda=\chi\mathsf L^{-1}.
 \label{vii:eq:gramiano-extension-minima}
\end{equation}
La incidencia \(\mathsf C\) y el emparejamiento celular \(M\) de
\eqref{vii:eq:coordenadas-corriente-material} son operadores diferentes.

\begin{proposition}[Selección variacional de la extensión]
\label{vii:prop:minimo-parametrizado}
La extensión \(f_{\min}\) es la única minimizante de la norma entre
los campos \(f\in H\) que satisfacen \(\mathsf C f=b\). Su energía es
\begin{equation}
 E(\theta)=\chi\min_{\mathsf C f=b}\|f\|^2
 =\chi\langle b,\mathsf L^{-1}b\rangle
 =\langle b,\Lambda b\rangle.
 \label{vii:eq:energia-extension-minima}
\end{equation}
El campo \(f_{\min}\), el potencial \(y\) y la energía \(E\)
dependen de manera \(C^1\) de los parámetros.
\end{proposition}

\begin{proof}
La sobreyectividad de \(\mathsf C\) hace inyectivo a \(\mathsf C^*\),
por lo que \(\langle z,\mathsf Lz\rangle=\|\mathsf C^*z\|^2>0\)
para \(z\ne0\). Así, \(\mathsf L\) es invertible y
\(\mathsf C f_{\min}=\mathsf L y=b\). Toda solución es
\(f_{\min}+z\), con \(z\in\ker\mathsf C\). Como
\(f_{\min}\in\operatorname{Ran}\mathsf C^*=(\ker\mathsf C)^\perp\),
\[
 \|f_{\min}+z\|^2=\|f_{\min}\|^2+\|z\|^2,
 \qquad
 \|f_{\min}\|^2=\langle y,\mathsf L y\rangle
 =\langle b,\mathsf L^{-1}b\rangle.
\]
Estas identidades prueban el mínimo, su unicidad y su valor. La inversión
es una operación diferenciable en el abierto de matrices invertibles;
las expresiones de \eqref{vii:eq:gramiano-extension-minima} proporcionan
la regularidad afirmada.
\end{proof}

La realización cúbica de \ref{vii:prop:incidencia} corresponde a
\(H=\mathbb R^{\{-1,1\}^3}\), al cociente de pantalla
\(Q=P_\square\mathbb R^6\) y a la incidencia \(\mathsf C=M:H\to Q\).
Su Gramiano y su normalización son
\begin{equation}
 \mathsf L_\square=12P_u+4P_-,\qquad
 \chi_m=\frac{112}{81\ell_m^2},\qquad
 \Lambda_m=\frac{28}{243\ell_m^2}(P_u+3P_-).
 \label{vii:eq:normalizacion-minimo-cubico}
\end{equation}
El factor \(112/81\) procede de la soldadura y del defecto cuadrático
de compresión de \eqref{vii:eq:lambda-pantalla}. Si \(b\) tiene dimensión
de longitud, también la tienen \(y\) y \(f_{\min}\), la incidencia es
adimensional y \(E\) es adimensional. La conversión a acción utiliza
la normalización en la recta de acción que se especifica más adelante.

\Needspace{15\baselineskip}
\subsection{Diferencial completo del mínimo}
\label{vii:subsec:diferencial-minimo}

\begin{theorem}[Respuesta a la variación de incidencia y frontera]
\label{vii:thm:diferencial-minimo}
Para toda dirección tangente real \(\delta\theta\), se tiene
\begin{equation}
 \delta E
 =\delta\chi\,\langle b,y\rangle
   +2\chi\operatorname{Re}
       \langle y,\delta b-(\delta\mathsf C)f_{\min}\rangle.
 \label{vii:eq:diferencial-minimo-completo}
\end{equation}
Si \(\chi=112/(81\ell_m^2)\), con \(\ell_m>0\) variable,
\begin{equation}
 \delta E=-2E\frac{\delta\ell_m}{\ell_m}
 +2\chi_m\operatorname{Re}
       \langle y,\delta b-(\delta\mathsf C)f_{\min}\rangle.
 \label{vii:eq:diferencial-minimo-escala}
\end{equation}
\end{theorem}

\begin{proof}
Diferenciar \(\mathsf L\mathsf L^{-1}=I_Q\) da
\[
 \delta\mathsf L^{-1}
 =-\mathsf L^{-1}(\delta\mathsf L)\mathsf L^{-1},\qquad
 \delta\mathsf L=(\delta\mathsf C)\mathsf C^*
                      +\mathsf C(\delta\mathsf C)^*.
\]
Aplicando la regla del producto a \eqref{vii:eq:energia-extension-minima}
y usando la autoadjunción de \(\mathsf L^{-1}\), resulta
\[
 \delta E=\delta\chi\,\langle b,y\rangle
 +2\chi\operatorname{Re}\langle y,\delta b\rangle
 -\chi\langle y,(\delta\mathsf L)y\rangle.
\]
Los dos términos del último producto son conjugados entre sí y
\(\mathsf C^*y=f_{\min}\); por tanto,
\[
 \langle y,(\delta\mathsf L)y\rangle
 =2\operatorname{Re}\langle y,(\delta\mathsf C)f_{\min}\rangle.
\]
Se obtiene \eqref{vii:eq:diferencial-minimo-completo}.
Finalmente, \(\delta\chi_m=-2\chi_m\delta\ell_m/\ell_m\)
produce \eqref{vii:eq:diferencial-minimo-escala}.
\end{proof}

La fórmula incorpora la variación del campo minimizante mediante la
ecuación que lo determina. También se comprueba directamente:
\(\mathsf C\delta f_{\min}=\delta b-(\delta\mathsf C)f_{\min}\)
y \(f_{\min}=\mathsf C^*y\) implican
\[
 \delta\|f_{\min}\|^2
 =2\operatorname{Re}\langle f_{\min},\delta f_{\min}\rangle
 =2\operatorname{Re}
     \langle y,\delta b-(\delta\mathsf C)f_{\min}\rangle.
\]
Por ejemplo, si \(\mathsf C\) es fijo y \(b=\ell_m b_0\), con
\(b_0\) fijo, los dos términos de
\eqref{vii:eq:diferencial-minimo-escala} se cancelan y la energía
permanece constante. Este control conserva conjuntamente la dimensión
de frontera y su normalización.

\subsection{Cociente fijo y cambios de marco}
\label{vii:subsec:marcos-minimo}

Las derivadas anteriores están tomadas en espacios \(H,Q\) y productos
hermíticos fijos. En la presentación de seis caras, la inversa siempre
actúa sobre \(Q=P_\square H_F\). Para una familia de cocientes
\(Q(\theta)\), se fija primero una trivialización isométrica
\(V_Q(\theta):Q\to Q(\theta)\), y análogamente para el dominio.
Las derivadas de esos mapas se incluyen al diferenciar las expresiones
transportadas a los espacios fijos. Las fórmulas se aplican en cada
abierto donde se conserve el rango. Los cambios del producto energético
requieren diferenciar además ese producto; una métrica lorentziana
y el producto positivo usado para minimizar conservan aquí sus tipos
distintos.

\begin{proposition}[Covariancia del mínimo]
\label{vii:prop:covariancia-minimo}
Sean \(g_H,g_Q\) unitarios. Bajo
\[
 \mathsf C'=g_Q\mathsf C g_H^*,\qquad b'=g_Qb,\qquad \chi'=\chi,
\]
se cumple \(y'=g_Qy\), \(f'_{\min}=g_Hf_{\min}\) y \(E'=E\).
Cuando los cambios de marco dependen del parámetro, el diferencial
completo conserva esa igualdad. En particular, una variación pura de
marco tiene \(\delta E=0\).
\end{proposition}

\begin{proof}
El Gramiano transformado es \(\mathsf L'=g_Q\mathsf Lg_Q^*\),
de donde se obtienen las transformaciones de \(y,f_{\min}\) y
la igualdad de energías. Para una variación pura de marco, en un punto
con marcos identidad, sean \(A_Q^*=-A_Q\) y \(A_H^*=-A_H\)
sus generadores. Entonces
\[
 \delta\mathsf C=A_Q\mathsf C-\mathsf C A_H,\qquad
 \delta b=A_Qb,\qquad
 \delta b-(\delta\mathsf C)f_{\min}=\mathsf C A_Hf_{\min}.
\]
Su contribución a \eqref{vii:eq:diferencial-minimo-completo} es
\(2\chi\operatorname{Re}\langle f_{\min},A_Hf_{\min}\rangle=0\).
La prueba incluye, así, los términos debidos a los marcos móviles.
\end{proof}

\subsection{Incidencia transportada y componentes de conexión}
\label{vii:subsec:incidencia-conexion-minima}

La dependencia de conexión se especifica antes de efectuar la
contracción. Una realización explícita de la incidencia por bloques
se obtiene como sigue. Sean \(\mathcal K\) una fibra hermítica común,
\(F=\{(i,\epsilon):1\leq i\leq3,\ \epsilon=\pm1\}\) y
\(V=\{-1,1\}^3\). Para cada incidencia \(\nu_i=\epsilon\),
la realización utilizada declara una ruta \(\gamma_{i,\epsilon,\nu}\)
desde la fibra del vértice hasta la de la cara, con transporte
\(U_{i,\epsilon,\nu}\). Las trivializaciones identifican esas fibras
con \(\mathcal K\). Definimos
\begin{equation}
 \begin{split}
 (\mathsf B_Uf)_{i,\epsilon}
   &=\sum_{\nu_i=\epsilon}U_{i,\epsilon,\nu}f_\nu,\\
 \mathsf C_U&=(q_\square\otimes I_{\mathcal K})\mathsf B_U:
   \bigoplus_{\nu\in V}\mathcal K\longrightarrow Q\otimes\mathcal K,
 \qquad q_\square=P_\square:H_F\to Q.
 \end{split}
 \label{vii:eq:incidencia-transportada-minima}
\end{equation}
Esta es una continuación variacional explícita de la incidencia cúbica
en la representación de rutas declarada. Para transportes identidad,
\(\mathsf C_U=M\otimes I_{\mathcal K}\) y se recupera
\eqref{vii:eq:normalizacion-minimo-cubico}, tensorizada con la fibra.
El Gramiano sigue siendo positivo en un entorno de esa realización,
pues la positividad definida es una condición abierta. En cada dominio
de rango completo, el peso y el campo se fijan por
\eqref{vii:eq:gramiano-extension-minima}. Las rutas que realizan las
incidencias, sus orientaciones y sus memorias son datos de esa
realización del transporte APP--TRIT--TPK y acompañan a la construcción.

Si \(\omega=(\omega_j)\) son coordenadas reales de conexión,
la derivada de \(\mathsf C_U\) se calcula sustituyendo en cada bloque
\begin{equation}
 \partial_j U_\gamma
 =\sum_{k=1}^{N}
    U_N\cdots U_{k+1}(\partial_jU_k)U_{k-1}\cdots U_1,
 \qquad U_\gamma=U_N\cdots U_1.
 \label{vii:eq:derivada-incidencia-rutas}
\end{equation}
Los productos vacíos son identidades. La fórmula procede de la regla
del producto y conserva el lugar de cada operación en la ruta.
Así, \(\mathsf C_j:=\partial_j\mathsf C_U\) se obtiene de una
suma finita de matrices conocidas en la realización elegida, seguida
de la proyección \(q_\square\otimes I\). Si se utiliza un cociente
móvil, su derivada se incorpora conforme a
\ref{vii:subsec:marcos-minimo}. La matriz cúbica entera fija corresponde
al caso \(\mathsf C_j=0\); la respuesta a conexión procede de la
dependencia de los transportes efectivamente declarada en
\eqref{vii:eq:incidencia-transportada-minima}.

\begin{theorem}[Componentes variacionales de la extensión mínima]
\label{vii:thm:componentes-minimo-conexion}
Para la acción de pantalla
\(S_{\min}=\mathfrak a_m E\), donde \(\mathfrak a_m\) es su
normalización en la recta de acción, escribamos
\[
 b_j=\partial_jb,\qquad \chi_j=\partial_j\chi,
 \qquad a_j=\partial_j\mathfrak a_m.
\]
Con la convención \(\delta S_{\min}=-\tfrac12\sum_j s_j\delta\omega_j\),
las componentes son
\begin{equation}
 \begin{split}
 s_j={}&4\mathfrak a_m\chi\operatorname{Re}
             \langle y,\mathsf C_jf_{\min}-b_j\rangle\\
      &-2\mathfrak a_m\chi_j\langle b,y\rangle-2a_jE.
 \end{split}
 \label{vii:eq:componentes-minimo-conexion}
\end{equation}
Para la normalización cúbica, el segundo término equivale a
\(4\mathfrak a_m E\,\partial_j\ell_m/\ell_m\).
Si el emparejamiento celular de esta misma acción tiene matriz
invertible \(\mathsf M_{\rm cel}\), sus coordenadas de corriente
son \(\sigma=\mathsf M_{\rm cel}^{-1}s\).
\end{theorem}

\begin{proof}
La regla del producto da
\(\partial_jS_{\min}=a_jE+\mathfrak a_m\partial_jE\).
Sustituir \eqref{vii:eq:diferencial-minimo-completo} y multiplicar
por \(-2\) prueba \eqref{vii:eq:componentes-minimo-conexion}.
La derivada de \(\chi_m\) produce la expresión de escala.
Finalmente, comparar
\(\delta S_{\min}=-\tfrac12\delta\omega^{\mathsf T}s\)
con \(-\tfrac12\delta\omega^{\mathsf T}\mathsf M_{\rm cel}\sigma\)
para toda variación da \(\mathsf M_{\rm cel}\sigma=s\),
y su inversa determina las coordenadas.
\end{proof}

La construcción recupera un campo y un peso seleccionados conjuntamente
por la incidencia, y evalúa su respuesta a las variaciones de la misma
realización. El dato de frontera \(b\), la familia de rutas, la escala
\(\ell_m\), la normalización \(\mathfrak a_m\) y el emparejamiento
celular conservan sus reglas de publicación. La interpretación como
corriente material utiliza esos objetos y la acción completa. La
respuesta de Cartan y su densidad métrica posterior se componen con
ella en el dominio de materia correspondiente, conservando cualquier
dependencia de la corriente respecto de la conexión.

\begin{HMTCompactSection}
% Artículo VII. Fragmento de cuerpo: corriente y respuesta Cartan--Holst.
% RESULTADO_RECUPERADO: acción, corriente variacional, inversa de Holst,
% torsión algebraica, respuesta cuadrática y normalización métrica.
% Fuentes leídas íntegramente:
% [II10] output/ARTICULO_II_REV10_EDICION_INTEGRADA_20260910/manuscrito/
%        sections/10_gravedad.tex, líneas 109--155.
% [G11c] integral activo 2249, fuente/manuscrito/sections/md/
%        11c_gravedad_holonomia_rev6.tex, líneas 214--334.
% [B028] integral activo, fuente/manuscrito/integracion_83/parte_iv/body/
%        B028_ch68_realizacion_cartan_holst_barbero_7df2afd9f675_11_gravedad_cartan_holst.tex,
%        líneas 76--218; copia idéntica conservada en II10/fuentes/propietarios_II/G03/.
%        SHA-256 e1e1f3325399b3ef09c331ed8268e3883f7bcc561cd7a560cf80d1be7eb1e91c.
% CERTIFICADO_NUEVO de resultados recuperados: inversión explícita del mapa
% torsión -> D(e wedge e), prueba de unicidad de contorsión y eliminación
% homogénea del término cuadrático. Las dos inversas se comprobaron sobre
% las 24 bases correspondientes mediante aritmética racional exacta.
% Precisión de dominio: la respuesta lineal en corriente y su eliminación
% cuadrática se enuncian para materia afín en la contorsión. Se conserva
% la dependencia del coeficiente respecto de la acción material concreta.
%
% DEPENDENCIAS DE INTEGRACIÓN (no constituyen remisiones bibliográficas):
% 1. vii:sec:realizacion-cartan: cotetrada, conexión y realización del
%    transporte enriquecido; incluir su prueba en el cuerpo precedente.
% 2. vii:sec:acoplamientos: genealogía de gamma_HMT, G_int, c_int,
%    homogeneidad dimensional y elección de la carta, con pruebas.
% 3. Para evaluar una especie material: acción S_m concreta, representación
%    de sus campos y cálculo de su corriente sigma. Este fragmento recibe
%    esa acción declarada y demuestra su respuesta variacional.
% 4. La identificación de una corriente de pantalla con sigma, su
%    contracción física y su normalización material tienen residencia propia.
%    Aquí no se introduce ni se selecciona una amplitud s_0.
% 5. El balance métrico utiliza la covariancia de la acción material y
%    sus ecuaciones de movimiento; no supone conservación separada del
%    tensor de espín y del tensor material de Levi--Civita.

\section{Corriente de espín y respuesta algebraica de Cartan--Holst}
\label{vii:sec:corriente-respuesta-cartan}

La realización del transporte enriquecido de la
sección~\ref{vii:sec:realizacion-cartan} proporciona la cotetrada y la
conexión; las secciones dimensionales y el parámetro de incidencia de la
secciones~\ref{v:ant:sec:barbero} y~\ref{v:ant:sec:gravity} fijan sus acoplamientos. Su composición
determina un problema variacional de primer orden. La corriente material
ocupa en él una posición precisa: es la derivada de la acción respecto de
la conexión y determina la respuesta torsional con la misma normalización.

\subsection{Dominio geométrico y normalización de la corriente}
\label{vii:subsec:dominio-corriente-cartan}

Trabajamos en una variedad lisa orientada \(M\) de dimensión cuatro,
con fibrado lorentziano interno \(E\), métrica
\(\eta=\operatorname{diag}(-1,1,1,1)\) y cotetrada no degenerada
\(e\in\Omega^1(M;E)\). En un marco local escribimos
\(e^I\), con \(I=0,1,2,3\). La conexión métrica
\(\omega^{IJ}=-\omega^{JI}\) actúa mediante \(D_\omega\).
Cuando los campos \(\Psi\) sean espinoriales, se conserva una estructura
de espín y la representación material correspondiente. Definimos
\begin{equation}
 \Sigma^{IJ}=e^I\wedge e^J,\qquad
 F^{IJ}=\mathrm d\omega^{IJ}+\omega^I{}_{K}\wedge\omega^{KJ},
 \qquad T^I=D_\omega e^I.
 \label{vii:eq:campos-cartan}
\end{equation}
El operador \(\star\) actúa sobre los índices bivectoriales internos,
\begin{equation}
 (\star X)^{IJ}=\tfrac12\epsilon^{IJ}{}_{KL}X^{KL},
 \qquad \star^2=-I,\qquad
 \mathsf P_\gamma=\star+\gamma^{-1}I.
 \label{vii:eq:operador-holst}
\end{equation}
Se distingue así esta dualidad de la dualidad exterior de formas sobre
\(M\). Los índices internos se elevan y contraen con \(\eta\).

Fijamos una carta con \(G_{\rm int}>0\), \(c_{\rm int}>0\),
\(\gamma=\gamma_{\rm HMT}\in\mathbb R\setminus\{0\}\) y
\(\Lambda_{\rm int}\) constantes sobre ella. Abreviamos
\(\kappa=8\pi G_{\rm int}/c_{\rm int}^4\) y \(c=c_{\rm int}\).
Estas secciones permanecen fijas durante la variación. La acción es
\begin{equation}
 \begin{split}
 S[e,\omega,\Psi]={}&
 \frac1{2\kappa c}\int_M\Sigma_{IJ}\wedge
                         (\mathsf P_\gamma F)^{IJ}
 -\frac{\Lambda_{\rm int}}{\kappa c}
                         \int_M\operatorname{vol}_e
 +S_{\rm m}[e,\omega,\Psi].
 \end{split}
 \label{vii:eq:accion-cartan-holst}
\end{equation}
La convención de volumen y orientación es la que da
\(\Sigma_{IJ}\wedge(\star F)^{IJ}=R\operatorname{vol}_e\)
en el sector de Levi--Civita. El tipo dimensional previo hace de cada
sumando una sección de acción.

Para la variación de la conexión se utilizan variaciones de soporte
compacto en el interior, o se fija
\(\delta\omega|_{\partial M}=0\). Puede emplearse también una
completación de borde cuyo problema variacional anule el mismo término
fronterizo. A cotetrada y campos materiales fijos, la corriente de espín
es la forma de grado tres, con valores en bivectores internos duales,
determinada por
\begin{equation}
 \delta_\omega S_{\rm m}
 =-\tfrac12\int_M\delta\omega^{IJ}\wedge\sigma_{IJ}.
 \label{vii:eq:corriente-variacional}
\end{equation}
El emparejamiento con una variación arbitraria de soporte compacto fija
\(\sigma\) de manera única. Esta definición conserva la unidad de
acción de la corriente y el factor de normalización que interviene en
su respuesta geométrica.

\subsection{Variación de la conexión e inversión de Holst}
\label{vii:subsec:variacion-holst}

\begin{theorem}[Ecuación variacional de Cartan--Holst]
\label{vii:thm:ecuacion-cartan-holst}
En el dominio anterior, la estacionariedad respecto de la conexión
equivale a
\begin{equation}
 D_\omega(\mathsf P_\gamma\Sigma)^{IJ}
       =\kappa c\,\sigma^{IJ}.
 \label{vii:eq:ecuacion-cartan-holst}
\end{equation}
\end{theorem}
\begin{proof}
La variación de la curvatura es
\(\delta F=D_\omega\delta\omega\). La dualidad interna es paralela
y autoadjunta para la contracción bivectorial; como \(\gamma\) es
constante, \(D_\omega\mathsf P_\gamma=\mathsf P_\gamma D_\omega\).
Poniendo \(B_{IJ}=(\mathsf P_\gamma\Sigma)_{IJ}\), se obtiene
\[
 B_{IJ}\wedge D_\omega\delta\omega^{IJ}
 =\mathrm d(B_{IJ}\wedge\delta\omega^{IJ})
   -D_\omega B_{IJ}\wedge\delta\omega^{IJ}.
\]
La última forma cambia de signo al permutar sus factores de grados tres
y uno. Las condiciones de borde eliminan la integral del término
exacto. El término cosmológico es independiente de \(\omega\).
Por tanto,
\[
 \delta_\omega S=
 \frac1{2\kappa c}\int_M\delta\omega^{IJ}\wedge
       \bigl[D_\omega B_{IJ}-\kappa c\,\sigma_{IJ}\bigr].
\]
La arbitrariedad de la variación prueba la ecuación, incluido su signo
y su coeficiente.
\end{proof}

\begin{lemma}[Inversa real del operador de Holst]
\label{vii:lem:inversa-holst}
Para \(\gamma\in\mathbb R\setminus\{0\}\),
\begin{equation}
 \mathsf P_\gamma^{-1}
 =\frac{\gamma^2}{1+\gamma^2}
                      (-\star+\gamma^{-1}I).
 \label{vii:eq:inversa-holst}
\end{equation}
\end{lemma}
\begin{proof}
Los dos factores son polinomios en \(\star\) y conmutan. Su producto
antes de normalizar es
\[
 (\star+\gamma^{-1}I)(-\star+\gamma^{-1}I)
 =-\star^2+\gamma^{-2}I=(1+\gamma^{-2})I.
\]
El denominador \(1+\gamma^2\) es positivo en el dominio real fijado.
La multiplicación por \(\gamma^2/(1+\gamma^2)\) da la identidad.
\end{proof}

La corriente determina así la forma bivectorial de grado tres
\begin{equation}
 Y^{IJ}:=\kappa c\,
   \frac{\gamma^2}{1+\gamma^2}
   (-\star+\gamma^{-1}I)\sigma^{IJ},
 \qquad D_\omega\Sigma^{IJ}=Y^{IJ}.
 \label{vii:eq:corriente-dualizada}
\end{equation}

\subsection{Reconstrucción puntual de la torsión y de la contorsión}
\label{vii:subsec:reconstruccion-torsion}

La regla de Leibniz da
\begin{equation}
 D_\omega\Sigma^{IJ}
 =T^I\wedge e^J-e^I\wedge T^J
 =:\mathsf L_e(T)^{IJ}.
 \label{vii:eq:mapa-torsion}
\end{equation}
En cada punto, \(\mathsf L_e\) actúa de
\(\Lambda^2T^*M\otimes E\) en
\(\Lambda^3T^*M\otimes\Lambda^2E\). Ambos espacios tienen dimensión
veinticuatro. Sea \(E_I\) el marco dual de la cotetrada, de modo que
\(\iota_{E_I}e^J=\delta_I^J\).

\begin{theorem}[Inversión algebraica de la ecuación torsional]
\label{vii:thm:inversion-torsion}
Para una cotetrada no degenerada, \(\mathsf L_e\) es un isomorfismo.
Su inversa sobre la forma \(Y\) se escribe
\begin{equation}
 B^I=\sum_J\iota_{E_J}Y^{IJ},\qquad
 b=\tfrac14\sum_I\iota_{E_I}B^I,\qquad
 T^I=B^I-e^I\wedge b.
 \label{vii:eq:torsion-explicita}
\end{equation}
La ecuación~\eqref{vii:eq:ecuacion-cartan-holst} determina por tanto
la torsión puntualmente a partir de \(e\) y de la corriente que figura
en ella.
\end{theorem}
\begin{proof}
Sea primero \(Y=\mathsf L_e(T)\) y pongamos
\(t=\sum_J\iota_{E_J}T^J\). La contracción de
\eqref{vii:eq:mapa-torsion}, utilizando que \(T^I\) tiene grado dos
y que hay cuatro elementos en la cotetrada, da
\[
 \sum_J\iota_{E_J}(T^I\wedge e^J)=2T^I,
 \qquad
 \sum_J\iota_{E_J}(e^I\wedge T^J)=T^I-e^I\wedge t.
\]
De aquí \(B^I=T^I+e^I\wedge t\). Una segunda contracción produce
\[
 \sum_I\iota_{E_I}B^I=t+3t=4t.
\]
Se recupera \(t=b\) y después \(T^I=B^I-e^I\wedge b\).
En particular, el núcleo de \(\mathsf L_e\) es nulo. La igualdad
de dimensiones prueba su sobreyectividad y hace válida la fórmula
para toda forma \(Y\) del codominio.
\end{proof}

La conexión se descompone de manera única como
\begin{equation}
 \omega^{IJ}=\mathring\omega^{IJ}(e)+K^{IJ},\qquad
 K^{IJ}=-K^{JI},\qquad T^I=K^I{}_{J}\wedge e^J,
 \label{vii:eq:contorsion}
\end{equation}
donde \(\mathring\omega(e)\) es la conexión de Levi--Civita y
\(K\) la contorsión. Para hacer explícita la inversa, escribimos
\(K_{IJ}=K_{IJK}e^K\) y
\(T_I=\tfrac12T_{IJK}e^J\wedge e^K\). Entonces
\begin{equation}
 T_{IJK}=K_{IKJ}-K_{IJK},\qquad
 K_{IJK}=\tfrac12(T_{KIJ}-T_{IJK}-T_{JKI}).
 \label{vii:eq:contorsion-componentes}
\end{equation}
La segunda igualdad se obtiene combinando cíclicamente la primera y
usando \(K_{IJK}=-K_{JIK}\); su sustitución recupera las componentes
de \(T\). La torsión y la contorsión son, por consiguiente, dos
presentaciones equivalentes de la respuesta algebraica con cotetrada fija.

\begin{corollary}[Sector sin corriente y respuesta lineal mínima]
\label{vii:cor:desacoplamiento-torsional}
Si \(\sigma=0\), se obtiene \(T=0\) y
\(\omega=\mathring\omega(e)\). Para un acoplamiento material cuya
corriente \(\sigma[e,\Psi]\) sea independiente de \(K\),
las ecuaciones~\eqref{vii:eq:corriente-dualizada},
\eqref{vii:eq:torsion-explicita} y
\eqref{vii:eq:contorsion-componentes} proporcionan una respuesta única
\(K_*(e,\Psi)\), lineal en \(\sigma\).
\end{corollary}
\begin{proof}
Las tres aplicaciones de reconstrucción son lineales con \(e\),
\(\gamma\) y \(\kappa c\) fijos. La corriente nula produce
\(Y=0\), después \(T=0\) y finalmente \(K=0\).
\end{proof}

En la acción mínima, la ecuación de conexión contiene esta respuesta
puntual. Su evolución se obtiene al evolucionar \(e\) y \(\Psi\):
\(K_*=K_*(e,\Psi)\) cambia con ellos. Si una acción material conserva
dependencia algebraica adicional en \(K\), la ecuación de Cartan
mantiene su forma variacional y debe resolverse con esa corriente;
la respuesta lineal anterior corresponde al dominio mínimo especificado.

\subsection{Eliminación de la contorsión y contribución cuadrática}
\label{vii:subsec:respuesta-cuadratica}

Consideremos ahora el acoplamiento mínimo afín en \(K\), expresado por
\begin{equation}
 S_{\rm m}[e,\mathring\omega+K,\Psi]
 =S_{\rm m}[e,\mathring\omega,\Psi]
   -\tfrac12\int_M K^{IJ}\wedge\sigma_{IJ}[e,\Psi].
 \label{vii:eq:materia-afin-contorsion}
\end{equation}
Esta condición identifica la acción material a la que se aplica la
eliminación cuadrática; la corriente sigue siendo la definida por
\eqref{vii:eq:corriente-variacional}.

La expansión de la curvatura es
\begin{equation}
 F(\mathring\omega+K)
 =F(\mathring\omega)+D_{\mathring\omega}K+K\wedge K,
 \qquad (K\wedge K)^{IJ}=K^I{}_{L}\wedge K^{LJ}.
 \label{vii:eq:curvatura-contorsion}
\end{equation}
Como \(D_{\mathring\omega}\Sigma=0\), el término lineal es la
derivada exterior de \(\Sigma_{IJ}\wedge(\mathsf P_\gamma K)^{IJ}\).
La contribución de borde correspondiente se mantiene explícita:
\begin{equation}
 S_\partial[e,K]=\frac1{2\kappa c}
     \int_{\partial M}\Sigma_{IJ}\wedge(\mathsf P_\gamma K)^{IJ}.
 \label{vii:eq:borde-contorsion}
\end{equation}
En un dominio interior, o con la completación de borde compatible,
la densidad de acción que depende de la contorsión es
\begin{equation}
 \mathcal Q_{e,\gamma}(K)-\tfrac12K^{IJ}\wedge\sigma_{IJ},
 \qquad
 \mathcal Q_{e,\gamma}(K)
  =\frac1{2\kappa c}\Sigma_{IJ}\wedge
                         \mathsf P_\gamma(K\wedge K)^{IJ}.
 \label{vii:eq:forma-cuadratica-contorsion}
\end{equation}
La densidad \(\mathcal Q_{e,\gamma}\) es una forma diferencial de
grado cuatro, homogénea de grado dos en \(K\).

\begin{theorem}[Respuesta efectiva cuadrática en la corriente]
\label{vii:thm:respuesta-cuadratica}
Bajo \eqref{vii:eq:materia-afin-contorsion}, la sustitución de la
solución algebraica \(K_*\) produce la densidad efectiva de espín
\begin{equation}
 \mathcal L_{\rm spin}^{\rm eff}
 =-\tfrac14 K_*^{IJ}\wedge\sigma_{IJ}
 =-\mathcal Q_{e,\gamma}(K_*).
 \label{vii:eq:accion-efectiva-spin}
\end{equation}
Con cotetrada y acoplamientos fijos, es homogénea de grado dos en
la corriente. La acción efectiva conserva además
\(S_\partial[e,K_*]\), o su completación de borde declarada.
\end{theorem}
\begin{proof}
Por el corolario~\ref{vii:cor:desacoplamiento-torsional},
\(K_*\) es lineal en \(\sigma\). La ecuación de conexión es
precisamente la estacionariedad de
\eqref{vii:eq:forma-cuadratica-contorsion}. Como ésta es algebraica,
se aplica puntualmente la variación radial \(\delta K=K_*\);
también puede multiplicarse por una función de soporte compacto.
La identidad de Euler para una forma cuadrática da
\[
 2\mathcal Q_{e,\gamma}(K_*)
             -\tfrac12K_*^{IJ}\wedge\sigma_{IJ}=0.
\]
La sustitución en la densidad de acción prueba
\eqref{vii:eq:accion-efectiva-spin}. Si
\(\sigma\mapsto\lambda\sigma\), entonces
\(K_*\mapsto\lambda K_*\) y la densidad efectiva se multiplica
por \(\lambda^2\). La expansión
\eqref{vii:eq:curvatura-contorsion} identifica separadamente el
término de borde que acompaña a esa eliminación.
\end{proof}

La corriente concreta, su contracción lorentziana, la dualidad y las
convenciones de la acción material determinan el signo y el coeficiente
de esta contribución. La identidad cuadrática proporciona la regla de
composición que permite evaluarlos una vez especificados esos datos.

\subsection{Fuente métrica efectiva y balance covariante}
\label{vii:subsec:fuente-metrica-efectiva}

La normalización de la fuente métrica se fija en la misma recta de
acción. Tras eliminar \(K\), definimos
\begin{align}
 \delta_g S_{\rm m}[e,\mathring\omega,\Psi]
 &=-\tfrac12\int_M\operatorname{vol}_e\,
                   \mathcal T^{S,\rm LC}_{\mu\nu}\,\delta g^{\mu\nu},
 \\
 \delta_g S_{\rm spin}^{\rm eff}
 &=-\tfrac12\int_M\operatorname{vol}_e\,
                   \mathcal U^{S,\rm spin}_{\mu\nu}\,\delta g^{\mu\nu}.
 \label{vii:eq:fuentes-normalizadas-accion}
\end{align}
Aquí se consideran variaciones interiores, con los términos de borde
tratados conforme al problema variacional fijado. Para el acoplamiento
afín, la segunda fuente es la variación métrica de
\eqref{vii:eq:accion-efectiva-spin}. La variación incluye la dependencia
de la corriente material respecto de \(e\) y de \(\Psi\).

\begin{proposition}[Ecuación métrica y conservación de la fuente total]
\label{vii:prop:fuente-metrica-balance}
Con las convenciones anteriores, las ecuaciones métricas son
\begin{equation}
 G_{\mu\nu}[g]+\Lambda_{\rm int}g_{\mu\nu}
 =\kappa c\bigl(\mathcal T^{S,\rm LC}_{\mu\nu}
                    +\mathcal U^{S,\rm spin}_{\mu\nu}\bigr)
 =\kappa\bigl(T^{\rm phys,LC}_{\mu\nu}
                    +U^{\rm phys,spin}_{\mu\nu}\bigr),
 \label{vii:eq:ecuacion-metrica-efectiva}
\end{equation}
donde \(T^{\rm phys,LC}=c\mathcal T^{S,\rm LC}\) y
\(U^{\rm phys,spin}=c\mathcal U^{S,\rm spin}\) cuando la
cotetrada toma valores en la recta de longitud. Toda solución satisface
\begin{equation}
 \mathring\nabla^\mu
   (T^{\rm phys,LC}_{\mu\nu}+U^{\rm phys,spin}_{\mu\nu})=0.
 \label{vii:eq:balance-metrico-total}
\end{equation}
\end{proposition}
\begin{proof}
La eliminación de \(K\) conserva las ecuaciones de las variables
restantes: en la regla de la cadena, el término que multiplica
\(\delta K_*\) se anula por la ecuación de conexión. En el sector
de Levi--Civita, la primera identidad de Bianchi anula el término
\(\Sigma_{IJ}\wedge F^{IJ}\), y la parte gravitatoria se escribe
\((2\kappa c)^{-1}\int_M(R-2\Lambda_{\rm int})\operatorname{vol}_e\).
Su variación interior es
\[
 \frac1{2\kappa c}\int_M\operatorname{vol}_e\,
        (G_{\mu\nu}+\Lambda_{\rm int}g_{\mu\nu})\delta g^{\mu\nu}.
\]
La suma con \eqref{vii:eq:fuentes-normalizadas-accion} da
\eqref{vii:eq:ecuacion-metrica-efectiva}. La identidad de Bianchi
contraída para \(\mathring\nabla\), la compatibilidad métrica y
la constancia de \(\kappa\) y \(\Lambda_{\rm int}\) en la carta
producen \eqref{vii:eq:balance-metrico-total}.
\end{proof}

Antes de eliminar la conexión, las identidades geométricas conservan
la forma
\begin{equation}
 D_\omega T^I=F^I{}_{J}\wedge e^J,\qquad D_\omega F^{IJ}=0.
 \label{vii:eq:bianchi-cartan}
\end{equation}
La primera es \(D_\omega^2e=F\wedge e\); la segunda se obtiene
expandiendo \(D_\omega(\mathrm d\omega+\omega\wedge\omega)\)
y cancelando los términos. En una acción material covariante, el
balance métrico total coincide con la identidad de Noether sobre las
ecuaciones materiales. Geometría y materia conservan así una fuente
conjunta; el reparto entre su contribución de Levi--Civita y su
contribución de espín queda determinado por la acción especializada.

\end{HMTCompactSection}
% FORMALIZACION_NUEVA / CERTIFICADO_NUEVO, 10-09-2026.
% Composición de las operaciones efectivas ya reunidas en 30e y 31.
% No modifica la selección de rutas, la frontera, los acoplamientos ni
% la normalización de la acción; no evalúa una amplitud cosmológica s_0.
% Propietarios y alcance: desarrollo/COMPOSICION_METRICA_REV06.md.
% Control independiente: pruebas/verificar_eliminacion_torsional_no_lineal.py.

\section{Eliminación torsional y respuesta métrica de la extensión mínima}
\label{vii:sec:eliminacion-torsional-no-lineal}

La extensión mínima proporciona un campo, un peso de respuesta y una
corriente mediante una sola operación variacional. Su incidencia
transportada depende de la conexión a través de los productos de ruta de
\eqref{vii:eq:derivada-incidencia-rutas}. Esa dependencia se conserva al
componer la corriente con la ecuación de Cartan--Holst. La eliminación
de la conexión se realiza, por tanto, sobre la acción completa de esa
extensión. El caso afín de
\ref{vii:thm:respuesta-cuadratica} queda contenido como una
especialización exacta de la construcción siguiente.

\subsection{Composición variacional en un nivel celular}
\label{vii:subsec:composicion-torsional-celular}

Fijemos un nivel finito, sus rutas, la representación de transporte y
el emparejamiento celular de
\ref{vii:thm:componentes-minimo-conexion}. Sean \(p\) coordenadas
reales de la cotetrada y de los campos materiales en una trivialización
local declarada, y \(k\in\mathbb R^d\) las coordenadas de la
contorsión \(K\). La familia \(p\mapsto e(p)\) conserva la
no degeneración de la cotetrada. La conexión que actúa en cada ruta es
\begin{equation}
 \omega(p,k)=\mathring\omega(e(p))+K(k),\qquad
 F(p,k)=S_{\min}[e(p),\omega(p,k),\Psi(p)]
       =\mathfrak a(p,k)E(p,k).
 \label{vii:eq:accion-minima-en-contorsion}
\end{equation}
Los espacios y las bases de esta expresión se transportan a espacios
fijos antes de derivar. Se trabaja en un abierto donde la incidencia
\(\mathsf C(p,k)\) es sobreyectiva y las familias son \(C^2\).
Esta regularidad adicional permite examinar el Hessiano de la acción.
Los acoplamientos \(G_{\rm int},c_{\rm int},\gamma_{\rm HMT}\)
conservan la carta y la convención de variación de
\ref{vii:subsec:dominio-corriente-cartan}.

Sea \(Q(p,k)\) la integral, o la realización celular declarada, de
la forma cuadrática
\eqref{vii:eq:forma-cuadratica-contorsion}. En bases reales,
\begin{equation}
 Q(p,k)=\tfrac12 k^{\mathsf T}\mathsf A(p)k,
 \qquad \mathsf A(p)^{\mathsf T}=\mathsf A(p).
 \label{vii:eq:cuadratica-celular-contorsion}
\end{equation}
El operador \(\mathsf A\) procede de la acción gravitatoria y de
su emparejamiento, mientras que el Gramiano
\(\mathsf C\mathsf C^*\) procede del problema de extensión.
La positividad del segundo garantiza la existencia del campo
minimizante; la forma gravitatoria conserva su propia signatura.
Los términos de borde se mantienen según
\eqref{vii:eq:borde-contorsion}; las fórmulas interiores siguientes
se aplican con variaciones admisibles que anulan su contribución, o
con la completación de borde correspondiente.

Definamos el covector de respuesta mediante
\begin{equation}
 d_kF(p,k)[h]= -\tfrac12 h^{\mathsf T}s(p,k),
 \qquad s(p,k)=\mathsf M_{\rm cel}(p,k)\sigma(p,k).
 \label{vii:eq:covector-torsional-minimo}
\end{equation}
Aquí \(s\) y \(\sigma\) conservan los tipos distintos de
\eqref{vii:eq:componentes-minimo-conexion}: el primero contiene los
coeficientes del diferencial y el segundo las coordenadas de la
corriente en el emparejamiento material. Si se cambia de base en
\(k\), también se transporta ese emparejamiento. Al diferenciar
\(s=\mathsf M_{\rm cel}\sigma\), se incluyen las derivadas de
ambos factores cuando el emparejamiento depende del parámetro.

\begin{proposition}[Ecuación compuesta de la extensión y la contorsión]
\label{vii:prop:ecuacion-torsional-no-lineal}
La estacionariedad interior de la acción respecto de la contorsión es
\begin{equation}
 \mathcal R(p,k):=\mathsf A(p)k-\tfrac12s(p,k)=0.
 \label{vii:eq:ecuacion-torsional-no-lineal}
\end{equation}
Cada componente de \(s\) se calcula a partir de la misma extensión:
\begin{equation}
 \begin{split}
 s_j={}&4\mathfrak a\chi\operatorname{Re}
       \langle y,(\partial_{k_j}\mathsf C)f_{\min}
                         -\partial_{k_j}b\rangle\\
 &-2\mathfrak a(\partial_{k_j}\chi)\langle b,y\rangle
       -2(\partial_{k_j}\mathfrak a)E.
 \end{split}
 \label{vii:eq:seleccion-variacional-corriente-k}
\end{equation}
Cuando el emparejamiento celular realiza la ecuación de conexión de
\ref{vii:thm:ecuacion-cartan-holst}, esta ecuación equivale a componer
\(\sigma(p,k)\) con la inversa de Holst, la inversa de
\(\mathsf L_e\) y la reconstrucción de contorsión:
\begin{equation}
 K=\mathsf C_e^{\rm tor}\,\mathsf L_e^{-1}
           \bigl(\kappa c\,\mathsf P_\gamma^{-1}
             \sigma[e,\mathring\omega(e)+K,\Psi]\bigr).
 \label{vii:eq:composicion-implicita-cartan}
\end{equation}
El símbolo \(\mathsf C_e^{\rm tor}\) designa exclusivamente la
aplicación \(T\mapsto K\) de
\eqref{vii:eq:contorsion-componentes}, con cotetrada fija.
\end{proposition}

\begin{proof}
El diferencial de \(Q+F\), en la dirección arbitraria \(h\), es
\[
 d_k(Q+F)[h]
   =h^{\mathsf T}\bigl(\mathsf A k-\tfrac12s\bigr).
\]
Su anulación para toda dirección prueba
\eqref{vii:eq:ecuacion-torsional-no-lineal}.
El teorema~\ref{vii:thm:componentes-minimo-conexion}, aplicado a
\(\partial_{k_j}\), da
\eqref{vii:eq:seleccion-variacional-corriente-k}. Con \(p\) fijo,
\(\delta\omega=\delta K\). La ecuación de Cartan se transforma,
sucesivamente, mediante \eqref{vii:eq:inversa-holst},
\eqref{vii:eq:torsion-explicita} y
\eqref{vii:eq:contorsion-componentes}; cada una de estas operaciones
es invertible en su dominio. Esto prueba la equivalencia con
\eqref{vii:eq:composicion-implicita-cartan} cuando ambas expresiones
usan la misma realización material y celular.
\end{proof}

La composición conserva la conexión en los dos miembros. Las rutas
que recorren varias celdas pueden acoplar las componentes de la
corriente entre ellas. La reconstrucción puntual de \(T\) a partir
de una corriente dada y la resolución conjunta de
\eqref{vii:eq:ecuacion-torsional-no-lineal} son, en ese caso,
operaciones sucesivas. Esta formulación finita mantiene explícita
la representación de rutas que interviene en el funcional.

\subsection{Existencia local de una rama estacionaria}
\label{vii:subsec:rama-torsional-local}

\begin{theorem}[Continuación local de la solución torsional]
\label{vii:thm:rama-torsional-local}
Supongamos que \((p_0,k_0)\) pertenece al abierto anterior,
\(\mathcal R(p_0,k_0)=0\), y que el Hessiano total
\begin{equation}
 \mathsf H_0
 =\mathsf A(p_0)+\partial_k^2F(p_0,k_0)
 =\mathsf A(p_0)-\tfrac12\partial_ks(p_0,k_0)
 \label{vii:eq:hessiano-total-torsional}
\end{equation}
es invertible. Existen entornos de \(p_0\) y \(k_0\) en los que
una única rama \(C^1\), \(k=k_*(p)\), satisface
\eqref{vii:eq:ecuacion-torsional-no-lineal} y pasa por \(k_0\).
En esa rama,
\begin{equation}
 \partial_{p_i}k_*
 =-\bigl[\mathsf A+\partial_k^2F\bigr]^{-1}
     \left[(\partial_{p_i}\mathsf A)k_*
                      +\partial_{p_i}\partial_kF\right].
 \label{vii:eq:derivada-rama-torsional}
\end{equation}
\end{theorem}

\begin{proof}
La inversión del Gramiano es \(C^2\) en el abierto de rango
completo. Por ello \(F\) es \(C^2\) y \(\mathcal R\) es
\(C^1\). Su derivada respecto de \(k\), en \((p_0,k_0)\), es
\(\mathsf H_0\). El teorema de la función implícita proporciona
la rama, su regularidad y su unicidad en un entorno del punto.
También puede verse la existencia mediante la aplicación
\[
 \mathcal T_p(k)=k-\mathsf H_0^{-1}\mathcal R(p,k).
\]
Su derivada respecto de \(k\) se anula en \((p_0,k_0)\).
En una bola suficientemente pequeña tiene norma a lo sumo
\(q<1\). Reduciendo el entorno de \(p_0\), se consigue
\(\|\mathcal T_p(k_0)-k_0\|<(1-q)r\), donde \(r\) es el
radio de la bola. Así, \(\mathcal T_p\) lleva la bola cerrada
en sí misma y es contractiva. Su punto fijo es la solución local.
Diferenciar \(\mathcal R(p,k_*(p))=0\) y despejar
\(\partial_{p_i}k_*\) da
\eqref{vii:eq:derivada-rama-torsional}.
\end{proof}

La hipótesis se verifica sobre la acción compuesta: el rango del
Gramiano asegura el mínimo de pantalla y el Hessiano total controla
la rama torsional. La unicidad afirmada es local alrededor de la
solución especificada; otras ramas pueden coexistir en el mismo
dominio de rango completo.

\subsection{Acción efectiva y contribución no afín}
\label{vii:subsec:accion-efectiva-no-afin}

\begin{theorem}[Identidad exacta de eliminación]
\label{vii:thm:eliminacion-torsional-no-lineal}
Sea \(k_*(p)\) una rama estacionaria. Supongamos, además, que el
segmento \(t\mapsto (p,tk_*)\), \(0\leq t\leq1\), permanece
en el abierto donde está definido \(F\). La corrección de acción
respecto de la realización de Levi--Civita es
\begin{equation}
 \begin{split}
 \Delta S(p)
  &:=Q(p,k_*)+F(p,k_*)-F(p,0)\\
  &=\tfrac14 k_*^{\mathsf T}s(p,k_*)
     -\tfrac12\int_0^1
           k_*^{\mathsf T}s(p,tk_*)\,dt.
 \end{split}
 \label{vii:eq:accion-efectiva-no-lineal}
\end{equation}
Equivalentemente,
\begin{equation}
 \Delta S=-\tfrac14 k_*^{\mathsf T}s(p,k_*)
  +\tfrac12\int_0^1 k_*^{\mathsf T}
       \bigl[s(p,k_*)-s(p,tk_*)\bigr]\,dt.
 \label{vii:eq:correccion-no-afin-integrada}
\end{equation}
Estas identidades se refieren a la acción integrada o celular.
En una realización con densidad local, se aplican a la densidad
con su emparejamiento de formas correspondiente. La acción efectiva
conserva separadamente el término de borde declarado.
\end{theorem}

\begin{proof}
Por la regla de la cadena y
\eqref{vii:eq:covector-torsional-minimo},
\[
 \frac{d}{dt}F(p,tk_*)
      =-\tfrac12 k_*^{\mathsf T}s(p,tk_*).
\]
Integrar entre \(0\) y \(1\) da la diferencia de acciones
materiales. La ecuación de estacionariedad, evaluada en la
dirección \(k_*\), y la homogeneidad cuadrática de \(Q\) dan
\[
 2Q(p,k_*)=\tfrac12 k_*^{\mathsf T}s(p,k_*).
\]
La suma de ambas identidades demuestra
\eqref{vii:eq:accion-efectiva-no-lineal}; sumar y restar
\(\tfrac12k_*^{\mathsf T}s(p,k_*)\) proporciona
\eqref{vii:eq:correccion-no-afin-integrada}.
\end{proof}

\begin{corollary}[Especialización afín]
\label{vii:cor:recuperacion-eliminacion-afin}
Si \(s(p,k)=s(p)\), entonces
\[
 \Delta S=-\tfrac14 k_*^{\mathsf T}s(p)=-Q(p,k_*).
\]
Es la realización integrada de
\eqref{vii:eq:accion-efectiva-spin}.
\end{corollary}
\begin{proof}
La diferencia dentro de la integral de
\eqref{vii:eq:correccion-no-afin-integrada} se anula. La segunda
igualdad procede de la identidad radial de estacionariedad.
\end{proof}

\subsection{Diferencial métrico de la acción eliminada}
\label{vii:subsec:diferencial-metrico-no-lineal}

\begin{theorem}[Respuesta de las variables geométricas]
\label{vii:thm:envolvente-metrica}
Sobre la rama del teorema~\ref{vii:thm:rama-torsional-local},
\begin{equation}
 \partial_{p_i}\Delta S
 =\tfrac12 k_*^{\mathsf T}(\partial_{p_i}\mathsf A)k_*
  +(\partial_{p_i}F)(p,k_*)-(\partial_{p_i}F)(p,0),
 \label{vii:eq:envolvente-metrica-no-lineal}
\end{equation}
donde las derivadas parciales mantienen fijas las coordenadas \(k\).
En productos hermíticos fijos, cada derivada material se calcula por
\begin{equation}
 \begin{split}
 \partial_{p_i}F={}&(\partial_{p_i}\mathfrak a)E
       +\mathfrak a(\partial_{p_i}\chi)\langle b,y\rangle\\
   &+2\mathfrak a\chi\operatorname{Re}
       \langle y,\partial_{p_i}b
                    -(\partial_{p_i}\mathsf C)f_{\min}\rangle.
 \end{split}
 \label{vii:eq:diferencial-geometrico-minimo}
\end{equation}
Si \(p_i\) varía la cotetrada, \(\partial_{p_i}\mathsf C\)
incluye la variación de
\(\mathring\omega(e(p))+K(k)\), de las rutas realizadas y de
las trivializaciones utilizadas, según sus dependencias declaradas.
\end{theorem}

\begin{proof}
La derivada total de \(Q(p,k_*)+F(p,k_*)\) contiene, además
de sus derivadas parciales, el término
\[
 \bigl[\partial_kQ(p,k_*)+\partial_kF(p,k_*)\bigr]
                      \partial_{p_i}k_*.
\]
El corchete es cero por la ecuación de conexión. Restar la derivada
de \(F(p,0)\) prueba
\eqref{vii:eq:envolvente-metrica-no-lineal}.
La regla del producto y
\eqref{vii:eq:diferencial-minimo-completo} prueban
\eqref{vii:eq:diferencial-geometrico-minimo}. La composición
\eqref{vii:eq:accion-minima-en-contorsion} obliga a diferenciar
la conexión de Levi--Civita al variar \(e\) con \(k\) fijo.
\end{proof}

Cuando el producto positivo de extensión también varía, su
contribución se conserva de forma explícita. En coordenadas, sea
\(\mathsf h(p,k)>0\) su matriz hermítica y considérese el mínimo
de \(f^*\mathsf h f\) bajo \(\mathsf C f=b\). Entonces
\begin{equation}
 \mathsf L=\mathsf C\mathsf h^{-1}\mathsf C^*,\quad
 y=\mathsf L^{-1}b,\quad
 f_{\min}=\mathsf h^{-1}\mathsf C^*y,\quad
 E=\chi b^*y.
 \label{vii:eq:minimo-producto-variable}
\end{equation}
Aquí \(y\) representa el covector de la restricción en bases
fijas. La matriz positiva \(\mathsf h\) conserva su papel
energético, distinto del de la métrica lorentziana.

\begin{lemma}[Variación del producto de extensión]
\label{vii:lem:producto-variable-extension}
Con las convenciones anteriores,
\begin{equation}
 \delta E=(\delta\chi)b^*y
   +2\chi\operatorname{Re}\bigl[y^*(\delta b-\delta\mathsf C f_{\min})\bigr]
   +\chi f_{\min}^*(\delta\mathsf h)f_{\min}.
 \label{vii:eq:diferencial-producto-variable}
\end{equation}
\end{lemma}
\begin{proof}
La fórmula \(\delta\mathsf h^{-1}
=-\mathsf h^{-1}(\delta\mathsf h)\mathsf h^{-1}\) da
\[
 \delta\mathsf L=(\delta\mathsf C)\mathsf h^{-1}\mathsf C^*
 +\mathsf C\mathsf h^{-1}(\delta\mathsf C)^*
 -\mathsf C\mathsf h^{-1}(\delta\mathsf h)
                              \mathsf h^{-1}\mathsf C^*.
\]
Al diferenciar \(E=\chi b^*\mathsf L^{-1}b\), el término
\(-\chi y^*(\delta\mathsf L)y\) produce los dos términos
con \(\delta\mathsf C\) y el término positivo con
\(\delta\mathsf h\). Esto prueba la identidad completa.
\end{proof}

En una realización material covariante cuya variación admita el
emparejamiento métrico de
\eqref{vii:eq:fuentes-normalizadas-accion}, la corrección determina
\begin{equation}
 \delta_g\Delta S
 =-\tfrac12\int_M\operatorname{vol}_e\,
          \mathcal U^{S,\min}_{\mu\nu}\,\delta g^{\mu\nu},
 \qquad U^{\rm phys,\min}_{\mu\nu}
                      =c\mathcal U^{S,\min}_{\mu\nu}.
 \label{vii:eq:fuente-metrica-minimo-no-lineal}
\end{equation}
Las ecuaciones \eqref{vii:eq:envolvente-metrica-no-lineal} y
\eqref{vii:eq:diferencial-geometrico-minimo}, completadas por
\eqref{vii:eq:diferencial-producto-variable} cuando corresponde,
evalúan sus coeficientes en las variaciones métricas declaradas.
El mínimo de pantalla, la solución torsional y la derivada métrica
son tres operaciones consecutivas sobre la misma acción.
La acción reducida conserva los posibles acoplamientos entre
celdas presentes en las rutas originales.

La especialización homogénea de
\ref{vii:sec:dilucion-rebote} utiliza además la realización del
tensor obtenido: su simetría espacial, el signo de la contracción
con el observador y la ley de escala. Cuando esa evaluación da
\(\rho_{\rm corr}(a)=-s_0a^{-6}\), la amplitud se lee en el
estado de referencia como
\(s_0=-a_{\rm ref}^{6}\rho_{\rm corr}(a_{\rm ref})\).
Esta lectura conserva el estado material y su normalización; el
teorema de eliminación proporciona la acción que debe contraerse.

\subsection{Modelo exacto de comprobación no afín}
\label{vii:subsec:control-no-afin}

El siguiente modelo finito comprueba los coeficientes y la
dependencia no afín de las identidades. Sus parámetros pertenecen
al ejemplo algebraico y mantienen separado el problema de
realización material de las rutas.

\begin{proposition}[Eliminación exacta de una extensión de rango uno]
\label{vii:prop:ejemplo-no-afin}
Sean \(H=\mathbb R^2\), \(Q=\mathbb R\),
\(\mathsf C(k)=(1\ \ k)\), \(b=1\), \(\chi=1\),
\(\mathfrak a=\lambda>0\), y \(Q_{\rm ex}(k)=k^2/2\).
Entonces
\begin{equation}
 f_{\min}=\frac{(1,k)^{\mathsf T}}{1+k^2},\quad
 F(\lambda,k)=\frac{\lambda}{1+k^2},\quad
 s(\lambda,k)=\frac{4\lambda k}{(1+k^2)^2},\quad
 \mathcal R(\lambda,k)
    =k\left(1-\frac{2\lambda}{(1+k^2)^2}\right).
 \label{vii:eq:ejemplo-minimo-no-afin}
\end{equation}
Para \(\lambda=25/2\), las soluciones estacionarias son
\(0,2,-2\). En \(k_*=2\),
\begin{equation}
 \mathsf H_*=\frac{16}{5},\quad s_*=4,\quad
 Q_{\rm ex}(k_*)=2,\quad F(\lambda,k_*)=\frac52,
 \quad \Delta S=-8.
 \label{vii:eq:valores-ejemplo-no-afin}
\end{equation}
La expresión afín \(-k_*s_*/4=-2\) difiere de la corrección
exacta. En la rama positiva,
\begin{equation}
 \partial_\lambda\Delta S=-\frac45,
 \qquad \partial_\lambda k_* =\frac1{20}
 \quad\text{en }\lambda=25/2.
 \label{vii:eq:derivadas-ejemplo-no-afin}
\end{equation}
\end{proposition}

\begin{proof}
El Gramiano es \(1+k^2>0\), por lo que el mínimo tiene las
expresiones indicadas. Diferenciar \(F\) da
\(s=-2\partial_kF\) y después \(\mathcal R=k-s/2\).
Para \(\lambda=25/2\), la ecuación es
\(k[(1+k^2)^2-25]=0\); como \(1+k^2>0\), sus raíces
son precisamente \(0,\pm2\). El Hessiano total es
\[
 1-\frac{2\lambda(1-3k^2)}{(1+k^2)^3},
\]
que vale \(16/5\) en ambas raíces no nulas. La corrección es
\(2+5/2-25/2=-8\). En la identidad radial,
\[
 \int_0^1 k_*s(\lambda,tk_*)\,dt
    =\int_0^1\frac{200t}{(1+4t^2)^2}\,dt=20,
\]
pues una primitiva es \(-25/(1+4t^2)\).
Así, \(\Delta S=8/4-20/2=-8\), incluido el término no afín.
Para \(\lambda>1/2\), la rama positiva satisface
\(k_*^2=\sqrt{2\lambda}-1\) y
\(\Delta S=\sqrt{2\lambda}-1/2-\lambda\).
Sus derivadas dan \eqref{vii:eq:derivadas-ejemplo-no-afin}.
La fórmula de envolvente proporciona el mismo resultado:
\((\partial_\lambda F)(\lambda,k_*)
 -(\partial_\lambda F)(\lambda,0)=1/5-1=-4/5\).
\end{proof}

% Artículo VII. Incorporación aditiva autorizada, 10-09-2026.
% RESULTADO_RECUPERADO: representación real de Clifford y CAR:
% TEORIA_HOLOGRAFICA_INTEGRAL_MASA_HMT_MD_REV2_2026-08-20/fuente/modulos/
% 18_cuantica_estadistica_y_termodinamica.md, líneas 131--181 y 218--255;
% acción de primer orden: 14_topologia_exoticos_y_gravedad.md, 459--517;
% distinción corriente media/segundo momento y dilución:
% HMT_OBRAS/OBRA_II/chapters/11_cartan_bianchi_cosmologia.tex, 201--230.
% PDF REV2: b56bb38d10b9a4f4d0de47b7da8010fb31a6e8e8684a3bfd4237e37d7a3dddfa.
% FORMALIZACION_NUEVA: acción espinorial afín explícita, contracción
% normalizada mediante 31, observable cuártico CAR y funcional de amplitud.
% CERTIFICADO_NUEVO: cálculo racional de la forma cuadrática en las cuatro
% componentes de una trivector y evaluación de la realización CAR finita.
% Convención explicitada: las matrices g tienen firma (-+++); A_D=-i g^0
% y la acción cinética simétrica lleva hbar/2. La combinación equivalente
% i hbar Gamma^I usa Gamma^I=-i g^I y firma opuesta. No se mezclan ambas.
% Dependencias: cotetrada/conexión de 30d; constantes internas heredadas;
% corriente e inversas de 31. Estado material, celda y normal orden son
% datos de la realización declarada, no constantes universales extraídas.
% La acción afín de este fragmento y la acción de extensión mínima de 30e
% conservan sus problemas variacionales propios. No se identifican sus
% corrientes. No se cambia 50: su dominio de amplitud positiva se concreta.

\section{Corriente espinorial y amplitud torsional del estado material}
\label{vii:sec:corriente-espinorial-densidad}

El desarrollo precedente y la realización espinorial que sigue articulan
dos problemas variacionales con antecedentes geométricos comunes. La
acción de extensión mínima \(S_{\min}=\mathfrak a_m E\), construida en
la sección~\ref{vii:sec:corriente-extension-minima}, incorpora la
dependencia de la incidencia transportada y de su mínimo respecto de la
conexión; su eliminación torsional conserva esa dependencia completa en
la sección~\ref{vii:sec:eliminacion-torsional-no-lineal}. La acción
espinorial \(S_D\) de \eqref{vii:eq:accion-dirac-material} acopla la
misma conexión a los campos de su representación y es afín en la
contorsión al fijar la cotetrada y los campos. Cada corriente se obtiene
variando su propia acción. Una preparación minimizante puede determinar
el estado material inicial de la realización espinorial; la sustitución
de campos dependientes de la conexión conserva, además, los términos
del diferencial compuesto. Esta distinción sitúa el cálculo de amplitud
siguiente en su acción y en su dominio variacional específicos.

La representación espinorial del transporte y la respuesta algebraica
de Cartan--Holst permiten calcular la intensidad torsional como un
funcional del estado material. Su evaluación conserva tres datos:
la representación de los campos, el segundo momento de su corriente y
el volumen respecto del cual se define la densidad. El coeficiente de
acoplamiento resulta de la acción y de su inversión; el estado inicial
determina la amplitud material.

\subsection{Representación espinorial y acción material afín}
\label{vii:subsec:accion-espinorial-afin}

Conservamos la cotetrada, la orientación y los acoplamientos de la
sección~\ref{vii:sec:corriente-respuesta-cartan}. La realización
espinorial utiliza el módulo complejo de dimensión cuatro obtenido al
complejificar las matrices reales
\begin{equation}
 \begin{gathered}
 J=\begin{pmatrix}0&-1\\1&0\end{pmatrix},\qquad
 R=\begin{pmatrix}0&1\\1&0\end{pmatrix},\qquad
 Z=\begin{pmatrix}1&0\\0&-1\end{pmatrix},\\
 g^0=J\otimes I_2,\qquad g^1=R\otimes I_2,\qquad
 g^2=Z\otimes R,\qquad g^3=Z\otimes Z.
 \end{gathered}
 \label{vii:eq:clifford-material-explicito}
\end{equation}
La multiplicación da
\(\{g^I,g^J\}=2\eta^{IJ}I_4\), con
\(\eta=\operatorname{diag}(-1,1,1,1)\). Las matrices se distinguen
del parámetro escalar \(\gamma\) del operador de Holst. Definimos
\begin{equation}
 g_I=\eta_{IJ}g^J,\quad
 \mathbb B_{IJ}=\tfrac14[g_I,g_J],\quad
 A_D=-i g^0,\quad \bar\psi=\psi^\dagger A_D,\quad
 \Omega=\tfrac12\omega^{IJ}\mathbb B_{IJ}.
 \label{vii:eq:adjunto-conexion-espinorial}
\end{equation}
El producto de Clifford verifica
\[
 [\mathbb B_{IJ},g^K]=\delta_J^K g_I-\delta_I^K g_J.
\]
Así, la misma conexión métrica actúa en los campos mediante
\(D\psi=\mathrm d\psi+\Omega\psi\) y
\(D\bar\psi=\mathrm d\bar\psi-\bar\psi\Omega\).
El transporte sobre una ruta se obtiene integrando esta conexión en
su representación espinorial; conserva la conexión y la ruta de la
realización geométrica precedente.

Sea \(\eta_I=\iota_{E_I}\operatorname{vol}_e\), de modo que
\(e^J\wedge\eta_I=\delta_I^J\operatorname{vol}_e\). Para un
campo de tipo dimensional \([\psi]=L^{-3/2}\), tomamos la acción
material mínima
\begin{equation}
 S_D=\int_M\left\{
 \frac{\hbar}{2}
 \bigl[\bar\psi g^I D\psi-(D\bar\psi)g^I\psi\bigr]\wedge\eta_I
 -mc\,\bar\psi\psi\operatorname{vol}_e\right\}.
 \label{vii:eq:accion-dirac-material}
\end{equation}
Aquí \(\hbar=\hbar_{\rm int}\), \(c=c_{\rm int}\) y la masa
\(m\) es la lectura de la especie material en la carta elegida.
Estos valores son salidas anteriores a esta realización. La matriz
\(A_D\) es hermítica y \(A_Dg^I\) es antihermítica. Por ello el
término cinético simétrico tiene la convención real correspondiente a
\eqref{vii:eq:clifford-material-explicito}. En esta firma, la ecuación
sin torsión tiene término principal \(\hbar g^I D_I\); la notación
equivalente \(i\hbar\Gamma^I D_I\), con \(\Gamma^I=-i g^I\),
emplea la relación de Clifford de firma opuesta. Esta precisión fija
conjuntamente adjunto, firma y factor imaginario de la acción.

Los campos materiales y la cotetrada permanecen fijos en la variación
de la conexión. Una preparación procedente de la extensión mínima de
la sección~\ref{vii:sec:corriente-extension-minima} puede determinar
el dato inicial del campo o del estado. Cada acción conserva su
variación propia: si se sustituye en una acción un campo minimizante
dependiente de la conexión, se aplica además su diferencial compuesto,
o el teorema de la envolvente bajo sus hipótesis de estacionariedad.
La acción~\eqref{vii:eq:accion-dirac-material} realiza aquí un
acoplamiento afín en la contorsión.

\begin{proposition}[Corriente espinorial con normalización variacional]
\label{vii:prop:corriente-dirac-explicita}
Sea
\(B^{IJK}=\bar\psi g^{[I}g^Jg^{K]}\psi\), donde los corchetes
designan la antisimetrización normalizada. La corriente definida en
\eqref{vii:eq:corriente-variacional} por la acción
\eqref{vii:eq:accion-dirac-material} es
\begin{equation}
 \sigma_{JK}
 =-\frac{\hbar}{2}\,
   \bar\psi\{g^I,\mathbb B_{JK}\}\psi\,\eta_I
 =-\frac{\hbar}{2}\,B^I{}_{JK}\eta_I.
 \label{vii:eq:corriente-dirac-trivector}
\end{equation}
Es independiente de la contorsión cuando \(e\) y \(\psi\) están
fijos. La inversión de la sección~\ref{vii:subsec:reconstruccion-torsion}
se aplica por tanto a esta corriente.
\end{proposition}
\begin{proof}
La variación de \(\Omega\) es
\(\delta\Omega=\frac12\delta\omega^{JK}\mathbb B_{JK}\).
Los dos sumandos cinéticos dan, respectivamente, el producto con
\(g^I\) a la izquierda y a la derecha. En consecuencia,
\[
 \delta_\omega S_D
 =\frac{\hbar}{4}\int_M\delta\omega^{JK}\wedge\eta_I\,
                 \bar\psi\{g^I,\mathbb B_{JK}\}\psi.
\]
La comparación con \eqref{vii:eq:corriente-variacional} fija el
factor y el signo de \eqref{vii:eq:corriente-dirac-trivector}.
Al expandir el anticonmutador y usar la relación de Clifford se
cancelan las contracciones de grado uno, y queda
\(\{g^I,\mathbb B_{JK}\}=g^{[I}g_Jg_{K]}\).
La acción es lineal en \(\Omega\), lo que prueba la última
afirmación.
\end{proof}

\subsection{Contracción efectiva y coeficiente del segundo momento}
\label{vii:subsec:contraccion-espinorial}

La forma cuadrática lorentziana del trivector se escribe
\begin{equation}
 q(B)=\frac1{3!}B^{IJK}B_{IJK}
 =-(B^{012})^2-(B^{013})^2-(B^{023})^2+(B^{123})^2.
 \label{vii:eq:cuadratica-trivector}
\end{equation}
Esta contracción conserva la firma del módulo espinorial. Su segundo
momento depende del estado y de la realización de los productos de
campos.

\begin{theorem}[Interacción efectiva de la corriente espinorial]
\label{vii:thm:coeficiente-torsional-espinorial}
Con la acción material~\eqref{vii:eq:accion-dirac-material}, la
eliminación algebraica de la contorsión da
\begin{equation}
 \mathcal L_{\rm spin}^{\rm eff}
 =C_\gamma q(B)\operatorname{vol}_e,\qquad
 C_\gamma=\frac{3\kappa c\hbar^2}{16}
                      \frac{\gamma^2}{1+\gamma^2}.
 \label{vii:eq:coeficiente-cuadratico-espinorial}
\end{equation}
El coeficiente se obtiene de la corriente y de la inversa de Holst
con las normalizaciones de la acción precedente.
\end{theorem}
\begin{proof}
Por \eqref{vii:eq:accion-efectiva-spin}, basta evaluar
\(-\frac14K_*^{IJ}\wedge\sigma_{IJ}\), con
\[
 Y^{IJ}=\kappa c\frac{\gamma^2}{1+\gamma^2}
           (-\star+\gamma^{-1}I)\sigma^{IJ},\qquad
 K_*=\mathsf C_e(Y).
\]
Aquí \(\mathsf C_e\) es la composición explícita de
\eqref{vii:eq:torsion-explicita} con
\eqref{vii:eq:contorsion-componentes}. La linealidad de esas dos
inversas reduce el cálculo a las partes \(-\star\) e \(I\).
Para dar la contracción completa, ordenamos las componentes del
trivector como \((012,013,023,123)\), extraemos los factores
\(\kappa c\hbar^2\gamma^2/(1+\gamma^2)\), y sustituimos
\(\sigma_{JK}=-\frac12B^I{}_{JK}\eta_I\). Las dos matrices
de la forma cuadrática restante son
\begin{equation}
 \mathcal B_{-\star}
 =\frac3{16}\operatorname{diag}(-1,-1,-1,1),\qquad
 \mathcal B_I=0_{4\times4}.
 \label{vii:eq:matrices-contraccion-espinorial}
\end{equation}
Estas matrices se obtienen usando únicamente
\(e^I\wedge\eta_J=\delta_J^I\operatorname{vol}_e\): para
cada componente \(B^{abc}\), se forma \(\sigma\), se calcula
\(Y\), después
\(T^I=\sum_J\iota_{E_J}Y^{IJ}
 -\frac14e^I\wedge\sum_{L,J}\iota_{E_L}\iota_{E_J}Y^{LJ}\),
y finalmente
\(K_{IJK}=(T_{KIJ}-T_{IJK}-T_{JKI})/2\).
La contracción da las cuatro entradas diagonales mostradas; las
seis polarizaciones entre componentes distintas son cero para ambos
operadores. Es el cálculo de todos los coeficientes de las dos
formas cuadráticas. La parte \(\gamma^{-1}I\) se anula y la parte
\(-\star\) produce exactamente \eqref{vii:eq:cuadratica-trivector},
con el coeficiente de \eqref{vii:eq:coeficiente-cuadratico-espinorial}.
\end{proof}

\subsection{Realización fermiónica finita y segundo momento del estado}
\label{vii:subsec:segundo-momento-car}

Fijamos una celda comóvil periódica, un marco ortonormal adaptado a su
observador y el modo espacial homogéneo de las cuatro componentes
espinoriales. Esta realización finita tiene espacio de estados
\(\mathcal F=\bigoplus_{n=0}^4\Lambda^n\mathbb C^4\).
Sean \(a_r,a_r^\dagger\), \(r=1,\ldots,4\), sus operadores
canónicos, con
\(\{a_r,a_s^\dagger\}=\delta_{rs}\) y
\(\{a_r,a_s\}=0\). El vacío de ocupación fija el orden
normal. Para una matriz \(M\), escribimos
\[
 \mathrm d\Gamma(M)=\sum_{r,s}a_r^\dagger M_{rs}a_s,
 \qquad \widehat N=\mathrm d\Gamma(I_4).
\]
La letra \(\Gamma\) de esta notación designa la segunda
cuantización; la holonomía nonádica conserva su notación propia.

Las cuatro matrices de los bilineales se calculan directamente:
\begin{equation}
 \begin{array}{c|cccc}
 abc&012&013&023&123\\ \hline
 M_{abc}=A_Dg^ag^bg^c
       &iJ\otimes R&iJ\otimes Z&iI_2\otimes J&iZ\otimes J
 \end{array}
 \label{vii:eq:matrices-bilineales-espin}
\end{equation}
Todas son hermíticas, tienen traza cero y satisfacen
\(M_{abc}^2=I_4\). El observable cuártico correspondiente a
\eqref{vii:eq:cuadratica-trivector}, en este orden normal declarado,
es
\begin{equation}
 \begin{split}
 \widehat{\mathscr Q}_{\rm sp}
 &=\sum_{a<b<c}\eta_{aa}\eta_{bb}\eta_{cc}
       :\!\mathrm d\Gamma(M_{abc})^2\!:\\
 &=\sum_{a<b<c}\eta_{aa}\eta_{bb}\eta_{cc}
       \mathrm d\Gamma(M_{abc})^2+2\widehat N.
 \end{split}
 \label{vii:eq:observable-cuadratico-espin}
\end{equation}
El orden normal es parte de la realización finita especificada;
su definición hace explícita la sustracción de las contracciones de
una sola ocupación.

\begin{proposition}[Evaluación por las relaciones canónicas]
\label{vii:prop:operador-espin-car}
El operador \(\widehat{\mathscr Q}_{\rm sp}\) es autoadjunto,
conserva cada sector de ocupación y tiene restricciones
\begin{equation}
 \widehat{\mathscr Q}_{\rm sp}|_{\Lambda^0\oplus\Lambda^1}=0,
 \qquad
 \widehat{\mathscr Q}_{\rm sp}|_{\Lambda^3}=4I_4,
 \qquad
 \widehat{\mathscr Q}_{\rm sp}|_{\Lambda^4}=8I_1.
 \label{vii:eq:sectores-espin-car}
\end{equation}
En la base de \(\Lambda^2\mathbb C^4\) cuyos subconjuntos de
ocupación son \((12,13,23,14,24,34)\), su matriz es
\begin{equation}
 \begin{pmatrix}
 0&0&0&0&0&-4\\
 0&2&0&0&6&0\\
 0&0&2&-6&0&0\\
 0&0&-6&2&0&0\\
 0&6&0&0&2&0\\
 -4&0&0&0&0&0
 \end{pmatrix}.
 \label{vii:eq:sector-dos-espin-car}
\end{equation}
\end{proposition}
\begin{proof}
La relación \(a_s a_t^\dagger=\delta_{st}-a_t^\dagger a_s\)
aplicada a \(\mathrm d\Gamma(M)^2\) separa su contracción de
una ocupación y da
\[
 :\!\mathrm d\Gamma(M)^2\!:
 =\mathrm d\Gamma(M)^2-\mathrm d\Gamma(M^2).
\]
Como los tres primeros signos de
\eqref{vii:eq:matrices-bilineales-espin} son negativos y el último
positivo, la suma de las contracciones es \(-2\widehat N\).
Esto prueba \eqref{vii:eq:observable-cuadratico-espin}.
La autoadjunción y la conservación de la ocupación se siguen de
esas mismas expresiones.

En \(\Lambda^0\) todos los sumandos se anulan. En \(\Lambda^1\),
\(\mathrm d\Gamma(M)=M\), y los cuadrados suman \(-2I_4\),
que se cancela con \(2\widehat N\). En \(\Lambda^4\),
\(\mathrm d\Gamma(M)=\operatorname{tr}M=0\), quedando ocho.
La identificación por complemento entre \(\Lambda^3\mathbb C^4\)
y el dual de \(\mathbb C^4\) lleva
\(\mathrm d\Gamma(M)\) a
\((\operatorname{tr}M)I-M^{\mathsf T}=-M^{\mathsf T}\).
Sus cuadrados vuelven a sumar \(-2I_4\), mientras
\(2\widehat N=6I_4\); por tanto la restricción es \(4I_4\).
Finalmente, la aplicación de las cuatro matrices explícitas a los
seis productos exteriores de la base indicada produce
\eqref{vii:eq:sector-dos-espin-car}. Así queda evaluado el operador
en las dieciséis dimensiones de \(\mathcal F\).
\end{proof}

Para un operador de estado \(\varrho\geq0\),
\(\operatorname{Tr}\varrho=1\), definimos
\(q_\varrho=\operatorname{Tr}(\varrho\widehat{\mathscr Q}_{\rm sp})\).
Su dependencia de los segundos momentos es explícita:
\begin{equation}
 q_\varrho=\sum_{a<b<c}\eta_{aa}\eta_{bb}\eta_{cc}
 \left\{
  \langle\widehat B_{abc}\rangle_\varrho^2
  +\operatorname{Var}_\varrho(\widehat B_{abc})
  -\langle\widehat N\rangle_\varrho
 \right\},\qquad
 \widehat B_{abc}=\mathrm d\Gamma(M_{abc}).
 \label{vii:eq:covarianza-espin-material}
\end{equation}
La covarianza y la sustracción normal se conservan incluso cuando
las corrientes medias se anulan. En particular, si \(P_3\) es la
proyección sobre \(\Lambda^3\mathbb C^4\),
\begin{equation}
 \varrho_3=\tfrac14P_3,\qquad
 \langle\widehat B_{abc}\rangle_{\varrho_3}=0,\qquad
 q_{\varrho_3}=4>0.
 \label{vii:eq:testigo-positivo-espin}
\end{equation}
Es un estado positivo de traza uno y constituye un testigo explícito
del dominio positivo del funcional. El sector de dos ocupaciones
contiene también valores negativos; por ejemplo, el vector
\((|12\rangle+|34\rangle)/\sqrt2\) tiene valor \(-4\).
El signo de la intensidad procede así de la contracción y del
estado, y queda separado de la selección de condiciones cosmológicas.

\subsection{Densidad, presión y conservación de la amplitud comóvil}
\label{vii:subsec:amplitud-espinorial-comovil}

En la realización espacialmente homogénea, fijamos
\[
 e^0=N(t)c\,\mathrm dt,\qquad e^j=a(t)\,\mathrm dx^j,
 \qquad \mathcal V(t)=a(t)^3V_c,
\]
con lapse \(N(t)>0\), factor de escala adimensional \(a(t)>0\)
y volumen comóvil \(V_c>0\). El marco y el modo periódico de la
subsección~\ref{vii:subsec:segundo-momento-car} permanecen declarados.
La normalización del campo homogéneo es
\(\widehat\psi=\mathcal V^{-1/2}(a_1,a_2,a_3,a_4)^{\mathsf T}\).
Cada bilineal lleva un factor \(\mathcal V^{-1}\), y su producto
normal ordenado lleva \(\mathcal V^{-2}\). La forma temporal
cinética simétrica conserva esta normalización canónica: los términos
procedentes de derivar \(\mathcal V^{-1/2}\) se cancelan entre
sus dos sumandos.

\begin{theorem}[Amplitud torsional como funcional del estado]
\label{vii:thm:amplitud-material-explicita}
Manteniendo fijos el estado canónico y la realización del producto
durante las variaciones métricas, el término efectivo de espín tiene
acción reducida
\begin{equation}
 S_{\rm spin}^{\rm hom}
 =\int\!\mathrm dt\,N(t)
       \frac{cC_\gamma}{a(t)^3V_c}\,q_\varrho.
 \label{vii:eq:accion-homogenea-espin}
\end{equation}
Su densidad de energía y su presión isotrópica son
\begin{equation}
 \rho_{\rm spin}=p_{\rm spin}
 =-\frac{cC_\gamma}{V_c^2}\,q_\varrho a^{-6}.
 \label{vii:eq:densidad-presion-espin}
\end{equation}
En cualquier evolución que conserve \(q_\varrho\), la amplitud
constante correspondiente es
\begin{equation}
 \boxed{\displaystyle
 s_0[\varrho,V_c]
 =\frac{3\kappa c^2\hbar^2}{16}
    \frac{\gamma^2}{1+\gamma^2}
    \frac{\operatorname{Tr}
       (\varrho\widehat{\mathscr Q}_{\rm sp})}{V_c^2}.}
 \label{vii:eq:amplitud-s0-espinorial}
\end{equation}
\end{theorem}
\begin{proof}
La integración espacial de
\eqref{vii:eq:coeficiente-cuadratico-espinorial} multiplica el
factor \(\mathcal V^{-2}\) por
\(Nc\mathcal V\,\mathrm dt\), lo que da
\eqref{vii:eq:accion-homogenea-espin}. En variables canónicas,
la variación del lapse define la energía mediante
\(\delta_N S_{\rm m}=-\int E\,\delta N\,\mathrm dt\).
Por tanto,
\[
 E_{\rm spin}=-\frac{cC_\gamma q_\varrho}{a^3V_c},
 \qquad \rho_{\rm spin}=\frac{E_{\rm spin}}{a^3V_c}.
\]
La variación espacial isotrópica satisface
\(\delta_a S_{\rm m}=\int3Na^2V_c p\,\delta a\,\mathrm dt\).
Aplicada a \eqref{vii:eq:accion-homogenea-espin}, da
\[
 3Na^2V_c p_{\rm spin}
 =-\frac{3NcC_\gamma q_\varrho}{a^4V_c}.
\]
Esto prueba simultáneamente el signo, la presión y la densidad de
\eqref{vii:eq:densidad-presion-espin}. Al conservarse
\(q_\varrho\), el coeficiente de \(-a^{-6}\) es constante y
se obtiene \eqref{vii:eq:amplitud-s0-espinorial}. La dimensionalidad
es consistente: \([cC_\gamma]=\mathrm{energía}\,L^3\),
de modo que \(s_0\) tiene dimensión de densidad de energía.
\end{proof}

\begin{corollary}[Dominio conservado de amplitud positiva]
\label{vii:cor:dominio-positivo-conservado}
Sea una evolución unitaria de \(\mathcal F\) que preserve
\(\Lambda^3\mathbb C^4\). Para todo estado inicial de traza uno
soportado en ese sector, se cumple \(q_{\varrho(t)}=4\), y
\begin{equation}
 s_0=\frac{3\kappa c^2\hbar^2}{4V_c^2}
                   \frac{\gamma^2}{1+\gamma^2}>0.
 \label{vii:eq:amplitud-positiva-sector-tres}
\end{equation}
\end{corollary}
\begin{proof}
La invariancia del sector conserva el soporte del estado; en él,
\eqref{vii:eq:sectores-espin-car} da
\(\widehat{\mathscr Q}_{\rm sp}=4I\). La conservación de la
traza fija entonces su esperanza en cuatro, independientemente de
la evolución interna de sus coherencias. La positividad de los
acoplamientos, \(V_c>0\) y \(\gamma\ne0\) prueba la fórmula
y su signo.
\end{proof}

Para estados generales, la conservación de
\(\langle\widehat N\rangle\) determina la dilución de la densidad
de ocupación, \(n(t)=\langle\widehat N\rangle/(a^3V_c)\).
La conservación del momento cuártico se comprueba por separado.
Si \(i\hbar\dot\varrho=[H(t),\varrho]\), entonces
\begin{equation}
 \frac{\mathrm d q_\varrho}{\mathrm dt}
 =\frac{i}{\hbar}\operatorname{Tr}
       \bigl(\varrho[H(t),\widehat{\mathscr Q}_{\rm sp}]\bigr).
 \label{vii:eq:conservacion-momento-cuartico}
\end{equation}
La conmutación, la invariancia de un sector donde el observable es
escalar, o una conservación demostrada para el estado concreto
proporcionan las condiciones pertinentes. Si el momento evoluciona,
\eqref{vii:eq:densidad-presion-espin} conserva su evaluación
instantánea; su balance incluye
\[
 \dot\rho_{\rm spin}+6H\rho_{\rm spin}
 =-\frac{cC_\gamma}{V_c^2}\,a^{-6}\dot q_\varrho,
 \qquad H=\dot a/a
\]
en tiempo propio. Este término registra el intercambio con los
restantes grados de libertad materiales.

La composición queda determinada: representación espinorial del
transporte, acción material, corriente variacional, inversión de
Cartan--Holst, segundo momento y densidad normalizada. El dominio
\(s_0>0\) de la sección~\ref{vii:sec:dilucion-rebote} dispone así
de un funcional evaluable y de estados testigo explícitos. Su uso en
la dinámica cosmológica conserva además las condiciones materiales,
geométricas y de conservación enunciadas en aquella sección.

% Artículo VII: pruebas preexistentes incorporadas al cuerpo.
% Propietario íntegro: fuentes_conservadas/17_rebote_einstein_cartan.tex.
% Incluye la ley de escala y las pruebas de rebote y persistencia.
% La amplitud s_0 es la salida de la realización material antecedente;
% este fragmento no certifica por sí solo la evaluación de esa realización.

\section{Dilución de memoria y persistencia del rebote}
\label{vii:sec:dilucion-rebote}

La evolución de la memoria y la expansión espacial son dos lecturas de
una misma sucesión de estados. El índice de refinamiento puede eliminarse
entre ambas; esa eliminación determina el exponente de dilución que la
realización cosmológica transmite a su contribución torsional.
La amplitud inicial depende de la evaluación de la corriente material
concreta en su respuesta variacional. La ley de escala y el teorema
de rebote que siguen conservan esa dependencia.

El teorema~\ref{vii:thm:amplitud-material-explicita} determina esa
amplitud como funcional del estado material y del volumen comóvil.
El corolario~\ref{vii:cor:dominio-positivo-conservado} proporciona
un dominio explícito de estados con amplitud positiva y conservada.
En ese dominio, la evaluación de \(s_0\) precede a la integración
cosmológica que se desarrolla a continuación. La densidad ordinaria
\(\rho_0>0\) y su ley barotrópica se especifican adicionalmente en la
sección~\ref{vii:sec:umbral-rebote}: el testigo de positividad de la
amplitud no evalúa por sí mismo esa densidad material.

\subsection{Ley de dilución \texorpdfstring{\(a^{-6}\)}{a elevado a -6}}
\label{vii:sec:dilucion}

\begin{proposition}[Eliminación del índice de memoria]
\label{vii:prop:dilucion}
Sean \(a_{\mathrm{ref}}>0\), \(\varphi>1\) y los lectores de la ruta
\[
 \frac{a_n}{a_{\mathrm{ref}}}=\varphi^{n/3},\qquad
 M_n=\varphi^{-2n}.
\]
Entonces
\[
 M_n=\left(\frac{a_n}{a_{\mathrm{ref}}}\right)^{-6}
\]
para todo índice de la sucesión. Si la realización de la densidad
cuadrática es \(s_n^2=s_*^2M_n\), con \(s_*^2>0\), se obtiene
\[
 s_n^2=s_*^2\left(\frac{a_n}{a_{\mathrm{ref}}}\right)^{-6}.
\]
\end{proposition}

\begin{proof}
Elevando la primera igualdad a la potencia \(-6\) se obtiene
\(\varphi^{-2n}\), que es exactamente el segundo lector. Multiplicar
por la amplitud inicial reproduce la segunda afirmación. El exponente
queda determinado por el cociente entre los exponentes de ambos
lectores; la amplitud es la evaluación del estado de referencia.
\end{proof}

En la realización homogénea continua se conserva esta ley constitutiva:
la densidad torsional se escribe \(s_0a^{-6}\), donde la normalización
del factor \(a\), el coeficiente material y la amplitud de referencia
se han absorbido en \(s_0\). La ley de escala y la evaluación de
\(s_0\) son operaciones sucesivas de la construcción.

\subsection{Umbral homogéneo y transversalidad}
\label{vii:sec:umbral-rebote}

Consideremos la carta homogénea, espacialmente plana, en unidades \(c=1\),
con \(G>0\) constante, densidades positivas \(\rho_0,s_0\) y parámetro
barotrópico constante \(w<1\). La realización especificada satisface
\(\Lambda_{\mathrm{eff}}=0\) y las ecuaciones
\begin{align}
 \rho_{\mathrm m}(a)&=\rho_0a^{-3(1+w)},&
 \rho_{\mathrm s}(a)&=s_0a^{-6},\\
 H^2&=\frac{8\pi G}{3}
        \bigl(\rho_{\mathrm m}-\rho_{\mathrm s}\bigr)
        =:F_0(a),&
 \dot H&=-4\pi G
        \bigl((1+w)\rho_{\mathrm m}-2\rho_{\mathrm s}\bigr).
 \label{vii:eq:dinamica-homogenea}
\end{align}
Los acoplamientos pertenecen a la sección dimensional ya construida.
La expresión de \(\rho_{\mathrm s}\) recoge el signo de su contribución
en esta realización gravitatoria.

\begin{theorem}[Existencia y unicidad del umbral transversal]
\label{vii:thm:rebote}
Existe un único \(a_{\mathrm b}>0\) con \(F_0(a_{\mathrm b})=0\):
\[
 a_{\mathrm b}^{\,3(1-w)}=\frac{s_0}{\rho_0}.
\]
La región permitida por \(H^2\geq0\) es \(a\geq a_{\mathrm b}\).
En el umbral, con
\(\rho_{\mathrm b}=\rho_{\mathrm m}(a_{\mathrm b})
                  =\rho_{\mathrm s}(a_{\mathrm b})>0\),
\[
 \dot H_{\mathrm b}=4\pi G(1-w)\rho_{\mathrm b}>0.
\]
La solución que atraviesa \(H=0\) presenta un mínimo estricto de \(a\),
con expansión local
\[
 a(t)=a_{\mathrm b}\left[
 1+2\pi G(1-w)\rho_{\mathrm b}(t-t_{\mathrm b})^2
 +O((t-t_{\mathrm b})^3)\right].
\]
\end{theorem}

\begin{proof}
La factorización
\[
 F_0(a)=\frac{8\pi G}{3}a^{-6}
             \bigl(\rho_0a^{3(1-w)}-s_0\bigr)
\]
reduce sus ceros al de una función estrictamente creciente de
\((0,\infty)\) sobre \((-s_0,\infty)\). Esto demuestra existencia,
unicidad y el signo de \(F_0\) a ambos lados del umbral.
La segunda ecuación de \eqref{vii:eq:dinamica-homogenea}, evaluada
en la igualdad de densidades, da la fórmula de \(\dot H_{\mathrm b}\).
Como \(\dot a=aH\), en el umbral
\(\ddot a_{\mathrm b}=a_{\mathrm b}\dot H_{\mathrm b}>0\).
Las ecuaciones son suaves para \(a>0\), por lo que su desarrollo de
Taylor da la expansión indicada. El cambio de signo de \(H\) es
transversal: pasa de contracción a expansión.
\end{proof}

Para \(w=1\), ambos términos tienen el mismo exponente y la igualdad de
amplitudes determina un caso degenerado. Para \(w>1\), el signo de
\(\dot H\) en una coincidencia de densidades es negativo.
El dominio \(w<1\) del teorema expresa exactamente el orden entre los
dos exponentes de dilución.

\subsection{Persistencia local bajo perturbaciones controladas}
\label{vii:sec:persistencia}

La persistencia del umbral utiliza dos estimaciones independientes:
el control del signo en los extremos y el de la derivada en el
intervalo. Sea \(I=[a_-,a_+]\subset(0,\infty)\), con
\(a_-<a_{\mathrm b}<a_+\), elegido de forma que
\[
 F_0(a_-)<0<F_0(a_+),\qquad
 m_I:=\inf_{a\in I}F_0'(a)>0.
\]
Tal intervalo existe porque
\[
 F_0'(a_{\mathrm b})
   =\frac{8\pi G}{a_{\mathrm b}}(1-w)\rho_{\mathrm b}>0.
\]
Definimos
\(\delta_I=\min\{-F_0(a_-),F_0(a_+)\}>0\).

\begin{theorem}[Persistencia local del rebote transversal]
\label{vii:thm:persistencia}
Sea \(R\in C^1(I)\) una perturbación que satisface
\[
 \|R\|_{\infty,I}<\delta_I,\qquad
 \|R'\|_{\infty,I}<m_I.
\]
Entonces \(F_R=F_0+R\) tiene un único cero \(a_R\in(a_-,a_+)\) y
\(F_R'(a_R)>0\). La realización \(H^2=F_R(a)\) admite el rebote local
transversal, con
\[
 \dot H_R=\frac{a_R}{2}F_R'(a_R)>0
\]
en el umbral. Si \(R(a,\eta)\) es conjuntamente \(C^1\), los ceros
persistentes dependen localmente de manera \(C^1\) del parámetro \(\eta\).
\end{theorem}

\begin{proof}
La primera cota conserva
\(F_R(a_-)<0<F_R(a_+)\); el teorema del valor intermedio proporciona
un cero interior. La segunda da
\(F_R'(a)\geq m_I-\|R'\|_{\infty,I}>0\) en \(I\), por lo que el cero
es único y simple.

En \(a>a_R\) cercano al cero,
\(F_R(a)=F_R'(a_R)(a-a_R)+o(a-a_R)\).
La integral
\[
 t-t_{\mathrm b}=\int_{a_R}^{a}
          \frac{d\xi}{\xi\sqrt{F_R(\xi)}}
\]
es finita y estrictamente creciente. Su inversa define la rama
expansiva; el reflejo temporal define la rama contractiva.
Ambas se unen con \(\dot a(t_{\mathrm b})=0\) y
\[
 a(t)-a_R
 =\frac{a_R^2F_R'(a_R)}4(t-t_{\mathrm b})^2
   +o((t-t_{\mathrm b})^2).
\]
Fuera del umbral, la derivación de \(H^2=F_R(a)\) da
\(\dot H=aF_R'(a)/2\); su límite continuo da la expresión positiva
en \(a_R\). La regularidad \(C^1\) de \(F_R\) garantiza la continuidad
de la aceleración en esta unión. Finalmente,
\(\partial_aF_R(a_R)>0\) permite aplicar el teorema de la función
implícita cuando la perturbación depende conjuntamente de \((a,\eta)\).
\end{proof}

La primera cota controla la existencia; la segunda asegura unicidad
local y transversalidad. La estabilidad demostrada corresponde a
ese umbral en el intervalo seleccionado. La amplitud torsional
interviene en su localización, mientras que la relación entre
exponentes determina la orientación temporal del cruce.

% FORMALIZACION_NUEVA de la realización homogénea preexistente.
% Prioridad de la fórmula integrada: procedencia histórica pendiente de cotejo.
% Conserva la ecuación y el dominio del capítulo de rebote; no selecciona s_0.
\section{Solución homogénea explícita y completitud causal}
\label{vii:sec:solucion-homogenea}
La determinación de la amplitud permite desarrollar una realización
completa del rebote, además del argumento local. Se conserva la carta
espacialmente plana de \eqref{vii:eq:dinamica-homogenea}, en unidades
\(c=1\), con \(\rho_0,s_0,G>0\), y se especializa el sector material
a polvo, \(w=0\), con \(\Lambda_{\mathrm{eff}}=0\).
El coeficiente \(s_0\) es el obtenido del estado
material en la realización correspondiente; la integración siguiente
emplea ese coeficiente y conserva su dominio.

\begin{theorem}[Integración exacta del rebote de polvo]
\label{vii:thm:rebote-polvo}
Pongamos
\[
 a_{\rm b}^3=\frac{s_0}{\rho_0},\qquad
 \tau^2=\frac{s_0}{6\pi G\rho_0^2},\qquad u=t-t_{\rm b}.
\]
La solución del sistema homogéneo con
\(a(t_{\rm b})=a_{\rm b}\) y \(H(t_{\rm b})=0\) es, para todo
\(t\in\mathbb R\),
\begin{equation}
 a(t)=a_{\rm b}\left(1+\frac{u^2}{\tau^2}\right)^{1/3},
 \qquad H(t)=\frac{2u}{3(\tau^2+u^2)},\qquad
 \dot H(t)=\frac{2(\tau^2-u^2)}{3(\tau^2+u^2)^2}.
 \label{vii:eq:solucion-polvo}
\end{equation}
Tiene un único mínimo, en \(t_{\rm b}\); es contractiva antes y
expansiva después. La evolución es suave y satisface
\(a(t)\geq a_{\rm b}>0\).
\end{theorem}
\begin{proof}
Sea \(y=a^3\). La ecuación de Friedmann da
\[
 \dot y^{\,2}=24\pi G(\rho_0y-s_0).
\]
La función \(y=s_0/\rho_0+6\pi G\rho_0u^2\) verifica esa
identidad y las condiciones iniciales. Sus derivadas dan las
tres expresiones de \eqref{vii:eq:solucion-polvo}. Para comprobar
también la ecuación de evolución, escribamos
\[
 \rho_{\rm b}=\frac{\rho_0^2}{s_0},\qquad
 x=\frac{u}{\tau}.
\]
Entonces
\[
 \rho_{\rm m}=\frac{\rho_{\rm b}}{1+x^2},\qquad
 \rho_{\rm s}=\frac{\rho_{\rm b}}{(1+x^2)^2},\qquad
 4\pi G\rho_{\rm b}=\frac{2}{3\tau^2}.
\]
Sustituir en \(\dot H=-4\pi G(\rho_{\rm m}-2\rho_{\rm s})\)
da exactamente la tercera fórmula. Se satisface, por tanto,
el sistema completo \((\dot a,\dot H)\), cuyo campo es suave
para \(a>0\). Su unicidad local, prolongada sobre la solución
explícita, fija la evolución por esas condiciones iniciales.
El denominador es positivo para todo \(t\), y las afirmaciones
sobre el mínimo y el signo siguen de \(H\).
\end{proof}

\begin{figure}[htbp]
\centering
\begin{tikzpicture}[x=1.45cm,y=1.15cm,>=Latex]
\draw[->] (-3.25,0)--(3.3,0) node[right] {\(\displaystyle u/\tau\)};
\draw[->] (0,0)--(0,2.8) node[above] {\(\displaystyle a/a_{\rm b}\)};
\draw[densely dashed,gray] (-3.1,1)--(3.1,1);
\draw[very thick,ink,domain=-3:3,samples=121,smooth]
 plot (\x,{(1+\x*\x)^(1/3)});
\fill[accent] (0,1) circle (1.6pt);
\node[anchor=north east] at (-.08,.98) {\(1\)};
\node[ink] at (-2,2.65) {Contracción};
\node[ink] at (2,2.65) {Expansión};
\foreach \t in {-3,-2,-1,1,2,3}
{\draw (\t,.045)--(\t,-.045) node[below] {\(\t\)};}
\end{tikzpicture}
\caption{Perfil adimensional de la solución demostrada:
\(a/a_{\rm b}=(1+(u/\tau)^2)^{1/3}\).
Los parámetros \(a_{\rm b}\) y \(\tau\) se calculan desde
\(\rho_0,s_0,G\); el dibujo representa toda la familia reescalada.}
\label{vii:fig:rebote-polvo}
\end{figure}

\subsection{Curvatura finita y extensión de las geodésicas}
La métrica realizada en esta carta es
\(g=-dt^2+a(t)^2\,d\mathbf x^2\) sobre
\(\mathbb R\times\mathbb R^3\). Con la convención en que
el escalar de curvatura plano es \(R=6(\dot H+2H^2)\), la solución da
\[
 R(t)=\frac{4\tau^2+\frac43u^2}{(\tau^2+u^2)^2}.
\]
En particular, \(R(t_{\rm b})=4/\tau^2\) es finito.
En un marco ortonormal, las componentes de Riemann dependen de
\(H^2\) y de \(\dot H+H^2\), que son funciones acotadas en
\(\mathbb R\). Lo mismo ocurre con las contracciones polinómicas
de esas componentes; por ejemplo,
\[
 R_{\mu\nu\rho\sigma}R^{\mu\nu\rho\sigma}
 =12\bigl[(\dot H+H^2)^2+H^4\bigr].
\]

\begin{proposition}[Completitud de las geodésicas causales]
\label{vii:prop:completitud-causal-polvo}
La realización métrica precedente es geodésicamente completa
para las geodésicas temporales y nulas de su conexión de
Levi--Civita, en ambas direcciones.
\end{proposition}
\begin{proof}
Las traslaciones espaciales proporcionan las constantes
\(p_i=a^2\,dx^i/d\lambda\). La normalización causal da
\[
 \left|\frac{dt}{d\lambda}\right|
 =\sqrt{\epsilon+\frac{|\mathbf p|^2}{a(t)^2}},
 \qquad \epsilon=1\ \text{en el caso temporal},\quad
 \epsilon=0\ \text{en el caso nulo}.
\]
Una geodésica nula no constante tiene \(|\mathbf p|>0\).
Como \(a\geq a_{\rm b}\), en el caso temporal
\[
 \left|\frac{d\lambda}{dt}\right|
 \geq\left(1+\frac{|\mathbf p|^2}{a_{\rm b}^2}\right)^{-1/2}>0,
\]
y en el nulo
\(\left|d\lambda/dt\right|=a/|\mathbf p|
 \geq a_{\rm b}/|\mathbf p|>0\).
Así, alcanzar \(t=\pm\infty\) requiere longitud afín infinita.
En un intervalo temporal compacto, \(a\) y sus derivadas son
suaves y \(a\) está separado de cero. Las ecuaciones explícitas
anteriores y \(dx^i/d\lambda=p_i/a^2\) permiten continuar la
geodésica a través de sus extremos finitos. Quedan cubiertas
las dos posibilidades de prolongación.
\end{proof}

El resultado de completitud pertenece a esta realización homogénea de
polvo, con sus condiciones de espacio, materia y amplitud. Su
persistencia local bajo perturbaciones del umbral conserva las cotas
del teorema~\ref{vii:thm:persistencia}; esa propiedad local y la
completitud global de la solución explícita son conclusiones distintas.
La escala mínima procede de la relación entre densidad material y
respuesta torsional. La memoria interviene así en una evolución
geométrica verificable, con una realización regular de la contracción
y la expansión.

% Artículo VII. Recuperación por dependencia, 10-09-2026.
% RESULTADO_RECUPERADO: integral activo 2249, fuente/colaboracion/
% partes_iv_v/tex/residencias/86_cosmologia_cerrada.tex (S04),
% líneas 2--44 (cartas seleccionadas), 174--248 (tensor y aceleración).
% Fuente S04 leída íntegramente, SHA-256:
% feb2582c3316c8b30829cfb0476bfa2f4e95889dabfba16fe43f2d835a226673.
% Se recuperan definición, unicidad, conservación, descomposición,
% intercambio y aceleración, con sus condiciones explícitas.
% FORMALIZACION_REUNIDA: relación con la fuente efectiva de 31;
% se distingue conservación total de conservación material separada.
% CERTIFICADO_NUEVO: expansión de las pruebas de traza, balance y signos;
% cálculo local de Ricci para las tres cartas ya seleccionadas por S04.
% No se modifica el rebote de 50 ni se determina aquí una amplitud s_0.
%
% DEPENDENCIAS DE INTEGRACIÓN:
% - vii:sec:realizacion-cartan: la métrica procede de la cotetrada realizada.
% - vii:sec:acoplamientos: G_int y el reloj t_0 del mismo estado HMT;
%   c_int=1 es una elección de unidades posterior a su construcción.
% - vii:eq:ecuacion-metrica-efectiva y vii:eq:balance-metrico-total, de 31:
%   se usan únicamente para identificar el residual sobre una solución.
% - T_b es el tensor publicado por el transductor material declarado;
%   su conservación separada es una condición expresamente enunciada.
% - El selector TRIT de la carta y t_0(x)>0 son antecedentes de la
%   especialización de Sitter, no datos cosmológicos observacionales.
% El presente fragmento no certifica la clausura del conjunto de VII.

\section{Tensor residual, descomposición de traza y balance covariante}
\label{vii:sec:tensor-residual-balance}

La realización geométrica del estado enriquecido proporciona una métrica,
mientras el transductor material publica la fuente a la que se refiere la
lectura. El tensor residual conserva la diferencia entre ambas publicaciones
en un mismo espacio tensorial. Su descomposición permite distinguir el
sector proporcional a la métrica y el sector de traza nula, y determina
la corriente de intercambio que corresponde a ese reparto.

\subsection{Datos de la carta y construcción del tensor residual}
\label{vii:subsec:datos-tensor-residual}

Fijemos un estado HMT \(x\) y una carta conexa \(\mathcal U\) de su
realización geométrica. Sobre \(\mathcal U\), la cotetrada no degenerada
de la sección~\ref{vii:sec:realizacion-cartan} determina una métrica lisa
\(g=g_x\) de firma \((-+++)\). Denotamos por
\(\mathring\nabla\) su conexión de Levi--Civita y escribimos
\[
 \mathsf G_{\mu\nu}[g]
 :=\operatorname{Ric}_{\mu\nu}[g]-\tfrac12R[g]g_{\mu\nu}.
\]
Así se distingue el tensor de Einstein \(\mathsf G[g]\) del acoplamiento
gravitatorio \(G_{\rm int}\). La derivada utilizada en los balances
de esta sección es \(\mathring\nabla\), coherentemente con la ecuación
métrica efectiva de la sección~\ref{vii:subsec:fuente-metrica-efectiva}.

Trabajamos en la carta dimensional \(c_{\rm int}=1\). Suponemos
\(G_{\rm int}(x)>0\) y \(\Lambda_{\rm ref}(x)\) constantes sobre
\(\mathcal U\), y ponemos \(\kappa_x=8\pi G_{\rm int}(x)\).
El símbolo \(x\) identifica el estado que suministra esos lectores;
la diferenciación covariante actúa sobre los puntos de \(\mathcal U\).
La genealogía de los acoplamientos permanece en la
sección~\ref{v:ant:sec:gravity}.

Sea \(T_b\) el tensor simétrico liso publicado por el transductor material
de la carta; en la lectura bariónica, es su fuente bariónica. Definimos
\begin{equation}
 \mathcal E^{\rm ref}_{\mu\nu}
 :=\kappa_x^{-1}\bigl(\mathsf G_{\mu\nu}[g]
                            +\Lambda_{\rm ref}g_{\mu\nu}\bigr),
 \qquad
 T^{\rm res}_{\mu\nu}:=\mathcal E^{\rm ref}_{\mu\nu}-T_{b,\mu\nu}.
 \label{vii:eq:residual-definicion}
\end{equation}
El tensor \(T^{\rm res}\) es el tensor residual efectivo, también
denotado \(T_{\rm osc}\) al interpretar el sector geométrico que queda
por publicar respecto de \(T_b\). La construcción conserva como
argumentos efectivos \(g_x,G_{\rm int},\Lambda_{\rm ref}\) y \(T_b\).

\begin{theorem}[Unicidad del residual y balance con la fuente material]
\label{vii:thm:residual-unicidad-balance}
Para los datos anteriores existe un único tensor simétrico
\(T^{\rm res}\) que satisface
\(\mathcal E^{\rm ref}=T_b+T^{\rm res}\). Su divergencia es
\begin{equation}
 \mathring\nabla^\mu T^{\rm res}_{\mu\nu}
 =-\mathring\nabla^\mu T_{b,\mu\nu}.
 \label{vii:eq:balance-residual-material}
\end{equation}
En particular, si la fuente material se conserva por separado,
\(\mathring\nabla^\mu T_{b,\mu\nu}=0\), también se conserva
\(T^{\rm res}\).
\end{theorem}
\begin{proof}
La diferencia de dos tensores que satisfagan la ecuación anunciada es
cero; la expresión~\eqref{vii:eq:residual-definicion} proporciona su
existencia. La identidad de Bianchi contraída y la compatibilidad métrica
de Levi--Civita dan
\[
 \mathring\nabla^\mu\mathsf G_{\mu\nu}=0,
 \qquad \mathring\nabla g=0.
\]
Como \(\kappa_x\) y \(\Lambda_{\rm ref}\) son constantes sobre la
carta, \(\mathring\nabla^\mu\mathcal E^{\rm ref}_{\mu\nu}=0\).
Derivar la diferencia que define \(T^{\rm res}\) prueba
\eqref{vii:eq:balance-residual-material}, incluida la afirmación bajo
conservación material separada.
\end{proof}

La constancia espacial y temporal de los acoplamientos en la carta fija
el dominio de este balance. Una familia de estados puede publicar
acoplamientos diferentes en cartas diferentes; cada identidad anterior
se aplica a la carta y al estado que la especifican.

\subsection{Descomposición canónica y corriente de intercambio}
\label{vii:subsec:traza-corriente-residual}

Definimos la función de traza y los dos tensores
\begin{equation}
 \theta:=g^{\mu\nu}T^{\rm res}_{\mu\nu},\qquad
 T^{\Lambda}_{\mu\nu}:=\tfrac14\theta g_{\mu\nu},\qquad
 T^{0}_{\mu\nu}:=T^{\rm res}_{\mu\nu}-T^{\Lambda}_{\mu\nu}.
 \label{vii:eq:descomposicion-traza-residual}
\end{equation}
El superíndice \(\Lambda\) identifica la componente proporcional a la
métrica; \(\Lambda_{\rm ref}\) es el parámetro de referencia de
\eqref{vii:eq:residual-definicion}. Sus funciones son distintas.

\begin{proposition}[Proyectores de traza y unicidad de la descomposición]
\label{vii:prop:proyectores-traza-residual}
Sobre los tensores simétricos covariantes de orden dos, las aplicaciones
\[
 P_{\rm tr}(S)=\tfrac14\operatorname{tr}_g(S)g,
 \qquad P_0(S)=S-P_{\rm tr}(S)
\]
satisfacen
\begin{equation}
 P_{\rm tr}^2=P_{\rm tr},\quad P_0^2=P_0,\quad
 P_{\rm tr}P_0=P_0P_{\rm tr}=0,\quad P_{\rm tr}+P_0=I.
 \label{vii:eq:proyectores-traza-residual}
\end{equation}
En particular, \(T^{\rm res}=T^\Lambda+T^0\),
\(\operatorname{tr}_gT^0=0\), y esta descomposición es única.
\end{proposition}
\begin{proof}
En dimensión cuatro, \(g^{\mu\nu}g_{\mu\nu}=4\). Por tanto,
\(\operatorname{tr}_g(P_{\rm tr}S)=\operatorname{tr}_gS\), lo que
prueba la idempotencia de \(P_{\rm tr}\), la anulación de la traza
de \(P_0S\) y, por expansión, las otras identidades.
Si \(S=fg+Q\) con \(\operatorname{tr}_gQ=0\), tomar la traza
da \(f=\operatorname{tr}_gS/4\). Queda determinado también
\(Q=S-fg\), y con ello la unicidad.
\end{proof}

\begin{theorem}[Intercambio covariante entre los dos sectores]
\label{vii:thm:intercambio-traza-residual}
En el dominio de \eqref{vii:eq:residual-definicion}, sea
\(b_\nu:=\mathring\nabla^\mu T_{b,\mu\nu}\). La corriente
de intercambio de traza es la uno-forma
\begin{equation}
 \mathcal J^{\rm tr}_\nu
 :=\tfrac14\partial_\nu\theta
 =\mathring\nabla^\mu T^\Lambda_{\mu\nu}.
 \label{vii:eq:corriente-intercambio-traza}
\end{equation}
Los balances completos son
\begin{equation}
 \mathring\nabla^\mu T^0_{\mu\nu}
       =-b_\nu-\mathcal J^{\rm tr}_\nu,
 \qquad
 \mathring\nabla^\mu(T^\Lambda+T^0)_{\mu\nu}=-b_\nu.
 \label{vii:eq:balances-sectores-residuales}
\end{equation}
Si \(b=0\), ambos sectores intercambian corrientes exactamente opuestas.
En ese caso se conservan por separado si y sólo si \(\theta\) es
constante en la carta conexa.
\end{theorem}
\begin{proof}
La compatibilidad métrica permite calcular
\[
 \mathring\nabla^\mu\bigl(\tfrac14\theta g_{\mu\nu}\bigr)
 =\tfrac14g^{\mu\alpha}(\partial_\alpha\theta)g_{\mu\nu}
 =\tfrac14\partial_\nu\theta.
\]
Restar esta igualdad a \eqref{vii:eq:balance-residual-material} produce
\eqref{vii:eq:balances-sectores-residuales}. Con \(b=0\), la
conservación separada equivale a \(\mathrm d\theta=0\); una función
lisa con diferencial nulo es constante en cada componente conexa.
\end{proof}

La uno-forma \(\mathcal J^{\rm tr}\) mide el intercambio del reparto
tensorial. Conserva un tipo distinto del covector de transporte de la
sección~\ref{vii:sec:variacion-memoria} y de la corriente de espín
\(\sigma_{IJ}\in\Omega^3(\mathcal U;\Lambda^2E^*)\) de
\eqref{vii:eq:corriente-variacional}.

\subsection{Enlace con la fuente efectiva de Cartan--Holst}
\label{vii:subsec:residual-fuente-cartan}

Sobre una solución de \eqref{vii:eq:ecuacion-metrica-efectiva}, en las
mismas unidades \(c=1\), el residual tiene la expresión
\begin{equation}
 T^{\rm res}
 =T^{\rm phys,LC}+U^{\rm phys,spin}-T_b
   +\frac{\Lambda_{\rm ref}-\Lambda_{\rm int}}{\kappa_x}\,g.
 \label{vii:eq:residual-composicion-cartan}
\end{equation}
En efecto, sustituir
\(\mathsf G+\Lambda_{\rm int}g
  =\kappa_x(T^{\rm phys,LC}+U^{\rm phys,spin})\)
en \eqref{vii:eq:residual-definicion} da la igualdad término a término.
Cuando el transductor material publica precisamente
\(T_b=T^{\rm phys,LC}\) y \(\Lambda_{\rm ref}=\Lambda_{\rm int}\),
se obtiene \(T^{\rm res}=U^{\rm phys,spin}\).
La condición de conservación separada de \(T_b\) del
teorema~\ref{vii:thm:residual-unicidad-balance} conserva su función
propia: \eqref{vii:eq:balance-metrico-total} establece el balance de la
fuente total, y el reparto entre sus componentes depende de la acción
material utilizada.

\subsection{Sector de traza y aceleración de las cartas seleccionadas}
\label{vii:subsec:residual-aceleracion-ds}

El lector TRIT del estado \(x\) publica \(s=\tau(x)\in\{+1,0,-1\}\),
y su reloj positivo publica
\begin{equation}
 H_*=\frac{2\pi}{108t_0(x)}>0.
 \label{vii:eq:reloj-carta-ds}
\end{equation}
Sea \(h_s\) una métrica tridimensional de curvatura seccional constante
\(s\). Las tres cartas de la realización seleccionada se escriben
\begin{equation}
 g_s=-\mathrm dt^2+a_s(t)^2h_s,\qquad
 a_+(t)=H_*^{-1}\cosh(H_*t),\quad
 a_0(t)=\exp(H_*t),\quad
 a_-(t)=H_*^{-1}\sinh(H_*t).
 \label{vii:eq:cartas-ds-residual}
\end{equation}
Las cartas cerrada y plana se consideran para \(t\in\mathbb R\);
la carta abierta, para \(t>0\). El signo \(s\) y el reloj \(H_*\)
mantienen sus antecedentes en el estado que selecciona la realización.

\begin{lemma}[Curvatura de las tres cartas del mismo reloj]
\label{vii:lem:curvatura-cartas-ds}
Para las funciones de \eqref{vii:eq:cartas-ds-residual},
\begin{equation}
 \frac{\ddot a_s}{a_s}=H_*^2,\qquad
 \left(\frac{\dot a_s}{a_s}\right)^2+\frac{s}{a_s^2}=H_*^2,
 \qquad
 \operatorname{Ric}[g_s]=3H_*^2g_s.
 \label{vii:eq:curvatura-cartas-ds}
\end{equation}
En consecuencia, \(R[g_s]=12H_*^2\) y
\(\mathsf G[g_s]=-3H_*^2g_s\).
\end{lemma}
\begin{proof}
Las dos primeras igualdades se obtienen diferenciando las tres funciones:
en las cartas cerrada y abierta se utiliza
\(\cosh^2z-\sinh^2z=1\), y en la plana
\(\dot a_0=H_*a_0\). Para verificar la identificación tensorial,
los coeficientes de conexión que involucran la coordenada temporal son
\[
 \Gamma^0{}_{ij}=a_s\dot a_s(h_s)_{ij},\qquad
 \Gamma^i{}_{0j}=(\dot a_s/a_s)\delta^i_j,
\]
y los puramente espaciales coinciden con los de \(h_s\).
La contracción de la curvatura da
\[
 \operatorname{Ric}_{00}=-3\ddot a_s/a_s,\qquad
 \operatorname{Ric}_{0i}=0,\qquad
 \operatorname{Ric}_{ij}
  =(a_s\ddot a_s+2\dot a_s^2+2s)(h_s)_{ij}.
\]
Sustituir las dos primeras identidades de
\eqref{vii:eq:curvatura-cartas-ds} produce
\(\operatorname{Ric}_{00}=-3H_*^2\) y
\(\operatorname{Ric}_{ij}=3H_*^2a_s^2(h_s)_{ij}\).
Tomar la traza y formar el tensor de Einstein prueba las otras
afirmaciones, con la convención de curvatura de
\eqref{vii:eq:ecuacion-metrica-efectiva}.
\end{proof}

\begin{theorem}[Densidad, presión y signo de la aceleración seleccionada]
\label{vii:thm:residual-ds-aceleracion}
En las cartas anteriores, para \(\Lambda_{\rm ref}=0\) y \(T_b=0\),
el tensor residual es
\begin{equation}
 T^{\rm res}=-\rho_*g_s,\qquad
 \rho_*:=\frac{3H_*^2}{8\pi G_{\rm int}}>0,
 \qquad T^\Lambda=T^{\rm res},\quad T^0=0.
 \label{vii:eq:residual-ds-puro}
\end{equation}
Para un observador comóvil unitario \(u=\partial_t\), la densidad
y la presión isotrópica son \(\rho_*\) y \(p_*=-\rho_*\).
La ecuación de aceleración se evalúa como
\begin{equation}
 \frac{\ddot a_s}{a_s}
 =-\frac{4\pi G_{\rm int}}3(\rho_*+3p_*)=H_*^2>0.
 \label{vii:eq:aceleracion-residual-ds}
\end{equation}
\end{theorem}
\begin{proof}
El lema anterior y \eqref{vii:eq:residual-definicion} dan
\(T^{\rm res}=-(3H_*^2/\kappa_x)g_s\).
Como \(g_s(u,u)=-1\),
\(T^{\rm res}_{\mu\nu}u^\mu u^\nu=\rho_*\).
El proyector espacial \(q^{\mu\nu}=g_s^{\mu\nu}+u^\mu u^\nu\)
satisface \(q^{\mu\nu}(g_s)_{\mu\nu}=3\), y por ello
\[
 p_*:=\tfrac13q^{\mu\nu}T^{\rm res}_{\mu\nu}=-\rho_*.
\]
La traza vale \(-4\rho_*\), de modo que
\eqref{vii:eq:descomposicion-traza-residual} produce los dos sectores
de \eqref{vii:eq:residual-ds-puro}. Finalmente,
\(\rho_*+3p_*=-2\rho_*\), y la sustitución en el segundo miembro
de \eqref{vii:eq:aceleracion-residual-ds} da
\(\kappa_x\rho_*/3=H_*^2\), igual al valor geométrico ya calculado.
La positividad utiliza exactamente \(G_{\rm int}>0\) y \(H_*>0\).
\end{proof}

En esta realización, el sector de traza publica una densidad constante
y una presión opuesta, y el reloj interno fija la aceleración positiva.
En una carta general, \(T^\Lambda\) describe la componente isotrópica
puntual, mientras \(T^0\) conserva el residual de traza nula y ambos
obedecen el intercambio de \eqref{vii:eq:balances-sectores-residuales}.
La identificación de estos sectores con fuentes oscuras observacionales
requiere que una misma realización reproduzca conjuntamente expansión,
lentes, rotación, crecimiento de estructura y fondo cósmico. La
construcción presente fija el tensor, sus balances y la aceleración de
las cartas seleccionadas que intervienen en esa comparación.

% Recuperacion de 09c_entropia_sectorial.tex, 11_confluencia.tex y
% 85_agujero_negro_doble_hoja.tex; originales en fuentes_conservadas/
% horizontes_informacion. Procedencia: desarrollo/HORIZONTES_INFORMACION_FUENTES.md.
% Las referencias a antecedentes se acuerdan con el editor de VII.
% No incorpora una corriente material ni evalua la amplitud torsional.

\section{Estados bipartitos puros, entropía reducida y leyes de área del horizonte}
\label{vii:sec:horizontes-informacion}

El estado enriquecido conserva las rutas y sus registros durante la
evolución; una lectura de subsistema determina qué parte de esa información
queda accesible. La construcción de esta sección utiliza los canales
angulares de la sección~\ref{iv:angular:section}, la incidencia y el
operador de área de la sección~\ref{v:ant:sec:area}, y las unidades
de acción, área y temperatura de las
secciones~\ref{v:ant:sec:gravity}
y~\ref{v:ant:sec:ii-kb-aut-desarrollo}. Son publicaciones anteriores de la
misma construcción APP--TRIT--TPK. Cada una transmite aquí su dominio,
orientación, normalización y condiciones de realización.

Se distinguen tres espacios: el registro binario de hoja, los factores
tríticos de la incidencia sectorial y la realización geométrica de los
exteriores de un horizonte. La purificación, el transporte de observables
y la conversión dimensional precisan las relaciones entre ellos. Esta
distinción conserva la estructura que acompaña a cada valor entrópico.

\subsection{Registro de hoja y purificación bipartita}
\label{vii:subsec:hojas-purificacion}

Escribamos \(a=A_{\rm rad}\) y \(y=C^*_{\rm rad}\) para las
coordenadas angulares ya construidas. Los canales y sus probabilidades son
\begin{equation}
 \mathfrak q_\pm=e^{-a\mp y},\qquad
 p_\pm=\frac{\mathfrak q_\pm}{\mathfrak q_++\mathfrak q_-}
       =\frac{e^{\mp y}}{2\cosh y}.
 \label{vii:eq:probabilidades-hoja}
\end{equation}
Los signos indexan estos dos canales. En
\(\mathcal H_{\rm h}=\operatorname{span}\{|+\rangle,|-\rangle\}\)
definimos
\[
 \rho_{\rm h}(y)=p_+|+\rangle\langle+|+p_-|-\rangle\langle-|.
\]
El espacio auxiliar
\(\mathcal H_{\rm ar}=\operatorname{span}
\{|\mathrm{adv}\rangle,|\mathrm{ret}\rangle\}\)
porta dos etiquetas ortonormales. Fijada una fase \(\chi_0\), el estado
conjunto es
\begin{equation}
 |\Psi_{\rm ar}\rangle
 =\sqrt{p_+}|+\rangle\otimes|\mathrm{adv}\rangle
  +e^{i\chi_0}\sqrt{p_-}|-\rangle\otimes|\mathrm{ret}\rangle.
 \label{vii:eq:purificacion-hoja}
\end{equation}

\begin{proposition}[Información accesible en el registro de hoja]
\label{vii:prop:purificacion-hoja}
El vector~\eqref{vii:eq:purificacion-hoja} está normalizado y su traza
parcial sobre \(\mathcal H_{\rm ar}\) es \(\rho_{\rm h}(y)\).
Su entropía de entrelazamiento, expresada en nats, es
\begin{equation}
 s_{\rm h}(y)=\log(2\cosh y)-y\tanh y.
 \label{vii:eq:entropia-hoja}
\end{equation}
El intercambio de las etiquetas \(+\) y \(-\) conjuga
\(\rho_{\rm h}(y)\) con \(\rho_{\rm h}(-y)\), y conserva
\(s_{\rm h}\).
\end{proposition}
\begin{proof}
La norma al cuadrado es \(p_++p_-=1\). La ortogonalidad de las
etiquetas auxiliares elimina los términos cruzados de la traza parcial.
Los coeficientes de Schmidt al cuadrado son precisamente \(p_\pm\).
Por~\eqref{vii:eq:probabilidades-hoja},
\(p_+-p_-=-\tanh y\); sustituir
\(\log p_\pm=\mp y-\log(2\cosh y)\) en
\(-\sum_\pm p_\pm\log p_\pm\) demuestra la fórmula.
El intercambio de signos permuta los dos autovalores y conserva la entropía.
\end{proof}

El factor común \(e^{-a}\) se cancela al normalizar el registro binario.
La fase relativa \(\chi_0\) permanece en el estado conjunto aunque esta
reducción diagonal no la publique. La reversibilidad lógica se refiere
al transporte del registro completo, con sus etiquetas y correlaciones;
el retorno de fase conserva el avance de memoria descrito en la
sección~\ref{vii:sec:retorno-memoria-gravedad}. Una implementación térmica
tiene además su propia ley de energía y disipación. En particular,
\(s_{\rm h}\) mide este subsistema binario y su bipartición concreta.

\subsection{Purificación de los sectores de incidencia}
\label{vii:subsec:purificacion-sectorial}

La incidencia de borde construida anteriormente distingue dos conjuntos
de dieciocho enlaces, con etiquetas sectoriales \(m=90,120\). Cada enlace
porta \(\mathbb C^3\), con base \(|0\rangle,|1\rangle,|2\rangle\)
y operador de ocupación \(\operatorname{diag}(0,1,2)\). Sean
\(\mathsf N_m\) sus sumas sectoriales, y
\begin{equation}
 \mathsf N=\mathsf N_{90}+\mathsf N_{120},\qquad
 \mathsf D_\partial=90\mathsf N_{90}+120\mathsf N_{120}.
 \label{vii:eq:ocupaciones-sectoriales}
\end{equation}
En la rama positiva, \(\gamma_{\rm BI}\) denota el parámetro de
Barbero--Immirzi ya obtenido, y \(a_0>0\) su factor areal angular.
Esta notación distingue el parámetro de la conexión y de sus holonomías.
El área se publica como
\begin{equation}
 \widehat{\mathcal A}_m=\ell_P^2\gamma_{\rm BI}a_0\mathsf N_m,
 \qquad
 \widehat{\mathcal A}=\widehat{\mathcal A}_{90}
                         +\widehat{\mathcal A}_{120}.
 \label{vii:eq:area-sectores}
\end{equation}

Para el parámetro finito \(\Delta=\Delta_{\rm mod}\in\mathbb R\) del
estado reducido, ponemos
\begin{align}
 x_m&=e^{-m\Delta},& z_m&=1+x_m+x_m^2,&
 Z(\Delta)&=z_{90}^{18}z_{120}^{18},\\
 \rho_\partial(\Delta)&=Z(\Delta)^{-1}e^{-\Delta\mathsf D_\partial}.
 \label{vii:eq:estado-sectorial}
\end{align}
La traza en esta sección es la traza matricial ordinaria y los estados
tienen traza uno. El parámetro \(\Delta\) conserva la selección y la
normalización del estado antecedente.

\begin{proposition}[Purificación tensorial de la frontera]
\label{vii:prop:purificacion-sectorial}
Para un peso real finito \(w\), sea
\begin{equation}
 |\operatorname{TFD}(w)\rangle
 =\frac{|0\rangle|0\rangle+e^{-w/2}|1\rangle|1\rangle
                   +e^{-w}|2\rangle|2\rangle}
        {\sqrt{1+e^{-w}+e^{-2w}}}.
 \label{vii:eq:tfd-local}
\end{equation}
El producto
\begin{equation}
 |\Psi_\Delta\rangle
 =|\operatorname{TFD}(90\Delta)\rangle^{\otimes18}
  \otimes|\operatorname{TFD}(120\Delta)\rangle^{\otimes18}
 \label{vii:eq:tfd-sectorial}
\end{equation}
es una purificación normalizada de \(\rho_\partial(\Delta)\).
\end{proposition}
\begin{proof}
En cada enlace, la suma de los cuadrados de los tres coeficientes es uno.
La traza parcial elimina los términos con índices distintos y devuelve
\(\operatorname{diag}(1,e^{-w},e^{-2w})/(1+e^{-w}+e^{-2w})\).
La traza parcial conmuta con el producto tensorial. Multiplicar las
treinta y seis reducciones da el exponencial y la función de partición
de~\eqref{vii:eq:estado-sectorial}.
\end{proof}

\subsection{Ecuación informacional y fluctuaciones areales}
\label{vii:subsec:ecuacion-informacional}

Las medias y varianzas de ocupación por enlace son
\begin{equation}
 f_m=\frac{x_m+2x_m^2}{z_m},\qquad
 v_m=\frac{x_m+4x_m^2}{z_m}-f_m^2.
 \label{vii:eq:momentos-sectoriales}
\end{equation}
Escribimos \(s(\rho)=-\operatorname{Tr}(\rho\log\rho)\), con
\(0\log0=0\). En la purificación anterior, \(s(\rho_\partial)\)
es su entropía de entrelazamiento.

\begin{proposition}[Área media, entropía y respuesta sectorial]
\label{vii:prop:respuesta-sectorial}
Para \(\Delta\) real y finito se cumple
\begin{align}
 \langle\mathsf N\rangle_\Delta&=18(f_{90}+f_{120}),\\
 \langle\mathsf D_\partial\rangle_\Delta&=18(90f_{90}+120f_{120}),\\
 \langle\widehat{\mathcal A}\rangle_\Delta
 &=18\ell_P^2\gamma_{\rm BI}a_0(f_{90}+f_{120}),\\
 s(\Delta)&=\log Z(\Delta)+\Delta\langle\mathsf D_\partial\rangle_\Delta,
 \label{vii:eq:entropia-sectorial}\\
 \frac{\mathrm d\langle\mathsf D_\partial\rangle_\Delta}{\mathrm d\Delta}
 &=-\operatorname{Var}_\Delta(\mathsf D_\partial)
  =-18(90^2v_{90}+120^2v_{120}),\\
 \frac{\mathrm ds}{\mathrm d\Delta}
 &=-\Delta\operatorname{Var}_\Delta(\mathsf D_\partial),\\
 \frac{\mathrm d\langle\widehat{\mathcal A}\rangle_\Delta}{\mathrm d\Delta}
 &=-18\ell_P^2\gamma_{\rm BI}a_0(90v_{90}+120v_{120}).
 \label{vii:eq:respuestas-area-entropia}
\end{align}
\end{proposition}
\begin{proof}
Los pesos locales son \((1,x_m,x_m^2)/z_m\), cuyos dos primeros
momentos dan \(f_m\) y \(v_m\). La aditividad de los operadores
ocupación y la factorización del estado dan las tres primeras fórmulas.
Tomar la traza de
\(-\rho_\partial\log\rho_\partial
=\rho_\partial(\log Z+\Delta\mathsf D_\partial)\)
da la entropía. La diferenciación de las sumas finitas produce
\(\partial_\Delta\log Z=-\langle\mathsf D_\partial\rangle_\Delta\)
y \(\partial_\Delta^2\log Z
=\operatorname{Var}_\Delta(\mathsf D_\partial)\).
Los factores independientes suman sus varianzas. En la derivada de
\eqref{vii:eq:entropia-sectorial} se cancelan los dos términos de
primer momento. Finalmente, \(\partial_\Delta f_m=-mv_m\)
da la respuesta del área.
\end{proof}

La pérdida de uniformidad se cuantifica en dimensión \(d_\partial=3^{36}\):
\begin{equation}
 I_\partial(\Delta)
 =36\log3-s(\Delta)
 =\mathcal D\bigl(\rho_\partial(\Delta)\Vert I/d_\partial\bigr)\ge0.
 \label{vii:eq:deficit-informacional}
\end{equation}
Aquí \(\mathcal D\) denota la entropía relativa, definida y justificada
abajo. Sustituir \(\log(I/d_\partial)=-36\log3\,I\) prueba la
identidad. La convexidad de \(t\log t\) da la positividad; la igualdad
corresponde a la distribución uniforme, que en esta familia ocurre
exactamente en \(\Delta=0\). El déficit mide información adimensional;
la escala \(\gamma_{\rm BI}a_0\ell_P^2\) mide área por ocupación.

\begin{theorem}[Inversión de ocupaciones y área complementaria]
\label{vii:thm:inversion-sectorial}
Sea \(\mathcal R_\partial\) la involución que intercambia
\(|0\rangle\) y \(|2\rangle\) y fija \(|1\rangle\) en cada
factor. Entonces
\begin{align}
 \mathcal R_\partial\mathsf N\mathcal R_\partial^*&=72I-\mathsf N,\\
 \mathcal R_\partial\mathsf D_\partial\mathcal R_\partial^*
 &=7560I-\mathsf D_\partial,\\
 \rho_\partial(-\Delta)&=\mathcal R_\partial\rho_\partial(\Delta)
                          \mathcal R_\partial^*,\\
 s(-\Delta)&=s(\Delta),\\
 \langle\widehat{\mathcal A}\rangle_{-\Delta}
 +\langle\widehat{\mathcal A}\rangle_\Delta
 &=72\gamma_{\rm BI}a_0\ell_P^2,\\
 \mathcal D\bigl(\rho_\partial(\Delta)\Vert\rho_\partial(-\Delta)\bigr)
 &=\Delta\bigl(7560-2\langle\mathsf D_\partial\rangle_\Delta\bigr).
 \label{vii:eq:inversion-sectorial}
\end{align}
\end{theorem}
\begin{proof}
Cada ocupación se transforma por \(n\mapsto2-n\). En cada sector,
\(\mathsf N_m\mapsto36I-\mathsf N_m\), por lo que el término
constante del lector ponderado es \(36(90+120)=7560\).
Además, \(z_m(-\Delta)=e^{2m\Delta}z_m(\Delta)\) y
\(Z(-\Delta)=e^{7560\Delta}Z(\Delta)\). La sustitución en el
estado normalizado prueba la conjugación. Su espectro conserva la entropía;
la transformación de \(\mathsf N\) da el área complementaria. Restar
los logaritmos de los dos estados y tomar la esperanza en
\(\rho_\partial(\Delta)\) da la última igualdad.
\end{proof}

En \(\Delta=0\), \(s=36\log3\). Cuando \(\Delta\to+\infty\),
la ocupación nula concentra el estado y la entropía tiende a cero;
la inversión lleva ese límite a la ocupación máxima. Para \(\Delta>0\),
la entropía decrece estrictamente porque el lector ponderado tiene más
de un autovalor y el estado tiene rango completo. La configuración
completa, el área y la entropía conservan así funciones diferenciadas.

\subsection{Ley modular y variaciones finitas de área}
\label{vii:subsec:ley-modular-finita}

Fijemos \(\rho_0=\rho_\partial(\Delta_0)\), con
\(\Delta_0\) real finito, y mantengamos
\(\Delta_0,\gamma_{\rm BI},a_0,\ell_P\) constantes. Las áreas
normalizadas \(\mathcal A_m^{\rm n}=\gamma_{\rm BI}a_0\mathsf N_m\)
expresan el operador modular como
\begin{equation}
 \mathsf K_0=-\log\rho_0
 =c_0I+\sum_{m=90,120}\beta_m\mathcal A_m^{\rm n},
 \quad \beta_m=\frac{m\Delta_0}{\gamma_{\rm BI}a_0},
 \quad c_0=\log Z(\Delta_0).
 \label{vii:eq:operador-modular-areal}
\end{equation}
Se tiene \(\beta_{120}=\tfrac43\beta_{90}\), también cuando ambos
son cero. El cociente es \(4/3\) cuando \(\Delta_0\ne0\).
\(\mathsf K_0\) es el logaritmo de este estado de referencia; su tipo
lo distingue del transporte temporal y de la conexión gravitatoria.

\begin{theorem}[Primera ley modular y corrección finita]
\label{vii:thm:ley-modular-finita}
Para cualquier estado normalizado \(\sigma\) del mismo espacio, sea
\(\Delta_\sigma\langle B\rangle
=\operatorname{Tr}[(\sigma-\rho_0)B]\). Entonces
\begin{equation}
 s(\sigma)-s(\rho_0)
 =\sum_{m=90,120}\beta_m\Delta_\sigma\langle\mathcal A_m^{\rm n}\rangle
  -\mathcal D(\sigma\Vert\rho_0),
 \label{vii:eq:ley-modular-finita}
\end{equation}
donde
\(\mathcal D(\sigma\Vert\rho_0)
=\operatorname{Tr}[\sigma(\log\sigma-\log\rho_0)]\)
es finita y no negativa. Si \(\sigma(\varepsilon)\) es diferenciable,
normalizada y \(\sigma(0)=\rho_0\), se obtiene
\begin{equation}
 \left.\frac{\mathrm ds(\sigma(\varepsilon))}{\mathrm d\varepsilon}\right|_0
 =\sum_{m=90,120}\beta_m
   \left.\frac{\mathrm d\langle\mathcal A_m^{\rm n}\rangle_{
                 \sigma(\varepsilon)}}{\mathrm d\varepsilon}\right|_0.
 \label{vii:eq:primera-ley-modular}
\end{equation}
\end{theorem}
\begin{proof}
Por definición,
\(\mathcal D(\sigma\Vert\rho_0)=-s(\sigma)
+\operatorname{Tr}(\sigma\mathsf K_0)\), mientras
\(s(\rho_0)=\operatorname{Tr}(\rho_0\mathsf K_0)\).
La resta y~\eqref{vii:eq:operador-modular-areal} prueban la identidad
finita, pues la normalización cancela el término \(c_0I\).

Para la positividad, diagonalicemos \(\sigma\) con autovalores
\(p_i\) y vectores \(u_i\), y \(\rho_0\) con autovalores
\(q_j>0\) y vectores \(v_j\). Si
\(r_i=\langle u_i,\rho_0u_i\rangle>0\), la concavidad del
logaritmo da
\[
 \sum_j|\langle u_i,v_j\rangle|^2\log q_j\le\log r_i.
\]
En consecuencia,
\(\mathcal D(\sigma\Vert\rho_0)
\ge\sum_i p_i\log(p_i/r_i)\ge0\).
La última desigualdad es la convexidad de \(t\log t\) con pesos
\(r_i\), cuya suma es uno. La fidelidad de \(\rho_0\) mantiene
finitos sus logaritmos; los autovalores nulos de \(\sigma\) se
interpretan mediante \(0\log0=0\).

Para la fórmula infinitesimal, en el entorno del estado fiel \(\rho_0\),
la derivada de la traza de \(-\sigma\log\sigma\) es
\(-\operatorname{Tr}[\dot\sigma(\log\rho_0+I)]\).
Esta identidad se obtiene de la derivada espectral de la traza.
La condición \(\operatorname{Tr}\dot\sigma=0\) deja
\(\operatorname{Tr}(\dot\sigma\mathsf K_0)\), y sustituir
\eqref{vii:eq:operador-modular-areal} concluye la prueba.
\end{proof}

Para las áreas físicas, los coeficientes son \(\beta_m/\ell_P^2\).
La aplicación del transductor común \(k_B\) da \(S=k_Bs\) y
conserva íntegro el término \(k_B\mathcal D\). La derivada en
\eqref{vii:eq:respuestas-area-entropia} varía \(\Delta\); la ley
\eqref{vii:eq:primera-ley-modular} mantiene fijo \(\Delta_0\) y
varía el estado alrededor de la referencia. Son operaciones distintas.

\subsection{Acción, energía de decisión y unidad areal}
\label{vii:subsec:conversion-informacion-area}

Fijamos una misma sección orientada de acción y sus secciones gravitatoria
y térmica, en la rama positiva. Los antecedentes proporcionan
\begin{equation}
 \begin{aligned}
 \ell_P^2&=\frac{\hbar G}{c^3},& t_P&=\frac{\ell_P}{c},\\
 T_{108}&=2\pi t_P,& h&=2\pi\hbar,\\
 E_P&=\frac{\hbar}{t_P}=k_B\Theta_{\rm clk}\log3.&&
 \end{aligned}
 \label{vii:eq:unidad-accion-informacion}
\end{equation}
La igualdad térmica conserva la pareja y la normalización previamente
construidas; al cambiar su base dimensional se transportan conjuntamente
\(k_B\) y \(\Theta_{\rm clk}\). El incremento elemental de área es
\(\delta\mathcal A=\gamma_{\rm BI}a_0\ell_P^2\).

\begin{proposition}[Coeficiente de conversión entre decisión y área]
\label{vii:prop:tension-informacion-area}
El cociente de la energía de decisión y del incremento areal satisface
\begin{equation}
 \begin{aligned}
 \mathcal T_{I/A}:=\frac{E_P}{\delta\mathcal A}
 &=\frac{k_B\Theta_{\rm clk}\log3}{\gamma_{\rm BI}a_0\ell_P^2}\\
 &=\frac{c^3}{\gamma_{\rm BI}a_0G t_P}
  =\frac{c^4}{\gamma_{\rm BI}a_0G\ell_P}.
 \end{aligned}
 \label{vii:eq:tension-informacion-area}
\end{equation}
Su dimensión es energía por área.
\end{proposition}
\begin{proof}
La sustitución de \(E_P=\hbar/t_P\) y
\(\ell_P^2=\hbar G/c^3\) cancela la sección de acción común.
La identidad \(\ell_P=ct_P\) da la última expresión. La positividad
de los factores asegura que todos los cocientes están definidos.
\end{proof}

La capacidad \(81\) del corte cúbico y la bipartición de los treinta y
seis enlaces sectoriales corresponden a dominios diferentes. En el
grafo cúbico \(\{0,\ldots,N-1\}^3\), las \(N^2\) rectas
paralelas entre dos caras opuestas son disjuntas por enlaces: todo corte
intercepta las \(N^2\), y una capa transversal alcanza esa cota.
Para \(N=9\), el estado producto uniforme de los trits de esa capa
tiene \(s_{\rm corte}=81\log3\). La asignación de una decisión del
mismo ciclo a cada enlace publica
\begin{equation}
 S_{\rm corte}^{\rm fis}=81k_B\log3,\qquad
 E_{\rm corte}^{(1)}=81E_P
                  =\Theta_{\rm clk}S_{\rm corte}^{\rm fis}.
 \label{vii:eq:corte-energia-informacion}
\end{equation}
La prueba combina el corte mínimo, la aditividad de la entropía del
estado producto y~\eqref{vii:eq:unidad-accion-informacion}. En el estado
sectorial, la ecuación finita pertinente sigue siendo
\eqref{vii:eq:ley-modular-finita}, con su entropía relativa.

\subsection{Normalización cosmológica del horizonte}
\label{vii:subsec:balance-horizonte}

Para un radio areal \(R>0\), la realización cosmológica considerada
utiliza la tríada
\begin{equation}
 E_\Lambda(R)=\frac{c^4R}{2G},\qquad
 S_{\rm H}(R)=\frac{k_B\pi R^2}{\ell_P^2},\qquad
 T_{\rm cos}(R)=\frac{\hbar c}{2\pi k_BR}.
 \label{vii:eq:horizonte-cosmologico}
\end{equation}
El factor \(2\pi\) pertenece a esta normalización de temperatura.
El área \(\mathcal A_{\rm H}=4\pi R^2\) expresa la misma entropía
como \(S_{\rm H}/k_B=\mathcal A_{\rm H}/(4\ell_P^2)\).

\begin{theorem}[Celda gravitatoria de media acción]
\label{vii:thm:celda-media-accion}
En~\eqref{vii:eq:horizonte-cosmologico},
\begin{align}
 T_{\rm cos}(R)S_{\rm H}(R)&=E_\Lambda(R),\\
 T_{\rm cos}(\ell_P)S_{\rm H}(\ell_P)
 &=\frac{E_P}{2}=\frac{k_B\Theta_{\rm clk}\log3}{2},\\
 E_\Lambda(\ell_P)t_P&=\frac{\hbar}{2},&
 E_\Lambda(\ell_P)T_{108}&=\frac h2.
 \label{vii:eq:media-accion-horizonte}
\end{align}
Manteniendo fijas las constantes de la carta y variando \(R\),
\begin{equation}
 T_{\rm cos}\,\mathrm dS_{\rm H}=2\,\mathrm dE_\Lambda,
 \qquad S_{\rm H}\,\mathrm dT_{\rm cos}=-\mathrm dE_\Lambda.
 \label{vii:eq:diferencial-horizonte}
\end{equation}
\end{theorem}
\begin{proof}
Multiplicar temperatura y entropía da
\(\hbar cR/(2\ell_P^2)=c^4R/(2G)\).
En \(R=\ell_P\), el resultado es
\(\hbar/(2t_P)=E_P/2\). Las dos duraciones de
\eqref{vii:eq:unidad-accion-informacion} dan las relaciones de acción.
Para los diferenciales,
\(\mathrm dE_\Lambda=c^4\mathrm dR/(2G)\),
\(\mathrm dS_{\rm H}=2\pi k_BR\mathrm dR/\ell_P^2\) y
\(\mathrm dT_{\rm cos}=-T_{\rm cos}\mathrm dR/R\).
La sustitución prueba ambas identidades.
\end{proof}

El producto \(E_\Lambda=T_{\rm cos}S_{\rm H}\) expresa una
identidad de estado. Sus variaciones radiales conservan los dos términos
de~\eqref{vii:eq:diferencial-horizonte}; una interpretación termodinámica
completa especifica además el generador temporal, el trabajo y la
realización material correspondientes. En esta carta, acción y
gravitación fijan la unidad \(\ell_P^2\), mientras Boltzmann convierte
información en entropía dimensional. El transporte de secciones con
\(G\hbar\) fijo conserva esa unidad cuando la sección \(c\) se
mantiene fija, y conserva \(S_{\rm H}/k_B\) a radio fijo.

En la celda \(R=\ell_P\), \(S_{\rm H}/k_B=\pi\), de modo que
la multiplicidad efectiva es \(\Omega_{\rm eff}=e^\pi\).
Para la involución hiperbólica trítica \(H_{\rm trit}\), con espectro
\(\{1,-1\}\), la diagonalización da
\(\operatorname{spec}(e^{\pi H_{\rm trit}})=\{e^\pi,e^{-\pi}\}\).
Esta relación espectral especifica el valor expansivo; una identificación
microscópica de ambos espacios conserva adicionalmente su mapa de
realización y el estado.

\subsection{Carta de Kruskal e intercambio de los exteriores}
\label{vii:subsec:kruskal-dos-hojas}

Fijemos un estado publicado \(x\), su radio \(R=R(x)>0\) y las
secciones positivas \(c,G\). En la carta que sigue, \(x,R,c,G\)
son parámetros fijos; \(t,r\) son las coordenadas espacio-temporales.
La realización exterior de las dos hojas se especifica por
\begin{equation}
 f(r)=1-\frac Rr,\qquad
 \mathrm ds^2=-f(r)c^2\mathrm dt^2+f(r)^{-1}\mathrm dr^2
                 +r^2\mathrm d\Omega_2^2,\qquad r>R.
 \label{vii:eq:metrica-exterior}
\end{equation}
Para el lector de hoja \(\sigma\in\{+1,-1\}\), definimos
\begin{align}
 r_*&=r+R\log\left(\frac rR-1\right),\\
 \mathcal U&=-\sigma\exp\left(-\frac{ct-r_*}{2R}\right),&
 \mathcal V&=\sigma\exp\left(\frac{ct+r_*}{2R}\right).
 \label{vii:eq:carta-kruskal}
\end{align}
Los símbolos \(\mathcal U,\mathcal V\) son coordenadas y se
distinguen de los operadores de transporte de rutas.

\begin{theorem}[Realización geométrica e involución de hoja]
\label{vii:thm:kruskal-hojas}
La carta~\eqref{vii:eq:carta-kruskal} realiza las dos hojas como los
dos exteriores \(\mathcal U\mathcal V<0\) de la extensión de
Kruskal. Se cumple
\begin{equation}
 -\mathcal U\mathcal V=\left(\frac rR-1\right)e^{r/R},\qquad
 \mathrm ds^2=-\frac{4R^3}{r}e^{-r/R}\,
                    \mathrm d\mathcal U\,\mathrm d\mathcal V
                    +r^2\mathrm d\Omega_2^2.
 \label{vii:eq:metrica-kruskal}
\end{equation}
La involución
\begin{equation}
 \mathcal C_{\rm H}:(\sigma,\mathcal U,\mathcal V)
 \longmapsto(-\sigma,-\mathcal U,-\mathcal V)
 \label{vii:eq:involucion-horizonte}
\end{equation}
es isométrica, fija el radio areal, intercambia los exteriores y preserva
el conjunto de horizontes \(\mathcal U\mathcal V=0\).
La métrica es regular en \(r=R\) y satisface las ecuaciones de vacío
en el dominio regular \(r>0\).
\end{theorem}
\begin{proof}
Multiplicar las coordenadas da la primera identidad. Como
\(\mathrm dr_*/\mathrm dr=f^{-1}\),
\[
 \mathrm d\mathcal U=\frac{\mathcal U}{2R}
            (-c\,\mathrm dt+f^{-1}\mathrm dr),\qquad
 \mathrm d\mathcal V=\frac{\mathcal V}{2R}
            ( c\,\mathrm dt+f^{-1}\mathrm dr).
\]
Su producto reproduce la parte radial y temporal de la métrica.
La función \(z\mapsto(z-1)e^z\) tiene derivada \(ze^z>0\)
para \(z>0\), lo que determina \(r\) a partir del producto.
En el exterior se recuperan
\(\sigma=\operatorname{sign}\mathcal V\) y
\(t=(R/c)\log(-\mathcal V/\mathcal U)\).
En \(r=R\), el coeficiente radial de la métrica extendida es
\(-4R^2/e\), finito y distinto de cero. La inversión conserva el
producto, el radio y \(\mathrm d\mathcal U\mathrm d\mathcal V\),
e intercambia las dos elecciones de signo.

Para comprobar el vacío en la carta exterior, las componentes mixtas
independientes del tensor de Einstein se reducen a
\[
 \mathsf G^t{}_t=\mathsf G^r{}_r
 =\frac{rf'+f-1}{r^2},\qquad
 \mathsf G^\theta{}_\theta=\mathsf G^\phi{}_\phi
 =\frac{f'}r+\frac{f''}2.
\]
Con \(f'=R/r^2\) y \(f''=-2R/r^3\), todas se anulan.
La expresión regular~\eqref{vii:eq:metrica-kruskal} prolonga el
resultado a través de los horizontes en el dominio \(r>0\).
\end{proof}

Esta realización conserva
\(E_{\rm H}=c^4R/(2G)\) y
\(S_{\rm H}=k_B\pi R^2/\ell_P^2\), que dependen del radio areal.
La temperatura \(T_{\rm cos}\) de
\eqref{vii:eq:horizonte-cosmologico} pertenece a la realización
cosmológica allí declarada. La carta exterior
\eqref{vii:eq:metrica-exterior} conserva su propia normalización temporal.
El mapa geométrico establece el intercambio de exteriores; la comparación
con formación por colapso, espectros de radiación o modos cuasinormales
utiliza además sus condiciones físicas respectivas.

\subsection{Transporte de estados, observables y refinamientos}
\label{vii:subsec:transporte-informacion-horizonte}

Denotemos por \(F_x\) la carta de los dos exteriores recién construida.
En cada carta angular regular, su jacobiano temporal y radial es
\begin{equation}
 \left|\frac{\partial(\mathcal U,\mathcal V)}{\partial(t,r)}\right|
 =\frac{cr}{2R^3}e^{r/R}>0.
 \label{vii:eq:jacobiano-kruskal}
\end{equation}
La identidad resulta de los dos diferenciales de la prueba anterior.
Los restantes registros del estado se transportan como coordenadas de
fibra; la carta conserva su identidad y su memoria.

\begin{proposition}[Equivalencia unitaria de las lecturas de hoja y exterior]
\label{vii:prop:transporte-unitario-horizonte}
Sea \(\mu\) una medida de hoja y \(\nu=(F_x)_*\mu\) su medida
transportada. El cambio de variables define una unitaria
\begin{equation}
 W_{\rm H}:L^2(\mu)\longrightarrow L^2(\nu),\qquad
 (W_{\rm H}\psi)(z)=\psi(F_x^{-1}z).
 \label{vii:eq:unitaria-horizonte}
\end{equation}
Con densidades respecto de medidas coordenadas fijadas, la misma
aplicación incorpora el factor de jacobiano correspondiente.
La conjugación
\begin{equation}
 \mathfrak I_{\rm H}(B)=W_{\rm H}BW_{\rm H}^*
 \label{vii:eq:mapa-observables-horizonte}
\end{equation}
es un \(*\)-isomorfismo unital completamente positivo. Transporta el
estado, sus esperanzas y su cálculo funcional modular. Para la medida
simétrica de las dos hojas, entrelaza también sus involuciones.
\end{proposition}
\begin{proof}
La definición de medida imagen da
\(\int|\psi\circ F_x^{-1}|^2\,\mathrm d\nu
=\int|\psi|^2\,\mathrm d\mu\).
La carta es invertible en los exteriores, de modo que la isometría es
sobreyectiva. La conjugación unitaria preserva producto, adjunto e
identidad; en cada ampliación matricial conjuga por
\(I_n\otimes W_{\rm H}\), lo que prueba la positividad completa.
Para un estado \(\rho\), la ciclicidad de la traza da
\(\operatorname{Tr}(W_{\rm H}\rho W_{\rm H}^*
\,\mathfrak I_{\rm H}(B))=\operatorname{Tr}(\rho B)\).
El cálculo funcional implica
\(g(W_{\rm H}\rho W_{\rm H}^*)=W_{\rm H}g(\rho)W_{\rm H}^*\)
en su dominio, en particular para el logaritmo y el flujo modular de
un estado fiel en su soporte. La igualdad
\(F_x\circ\mathcal C_{\rm hoja}=\mathcal C_{\rm H}\circ F_x\)
prueba el entrelazamiento de involuciones; la medida simétrica las realiza
unitariamente.
\end{proof}

Si los refinamientos conservan \((\sigma,t,r,R)\) y utilizan la misma
isometría sobre los registros de fibra, las inclusiones
\(j_m^{\rm hoja}\) y \(j_m^{\rm ext}\) satisfacen
\begin{equation}
 W_{{\rm H},m+1}j_m^{\rm hoja}
 =j_m^{\rm ext}W_{{\rm H},m}.
 \label{vii:eq:naturalidad-horizonte}
\end{equation}
En efecto, el cambio de variables actúa sobre las coordenadas geométricas
conservadas y la inclusión sobre las mismas coordenadas de fibra; evaluar
ambas composiciones da la misma función y la misma densidad transportada.
La igualdad extiende el mapa al sistema refinado sin reiniciar los registros.

Para una bipartición \(\mathcal H_A\otimes\mathcal H_B\), un
transporte factorizado \(U_A\otimes U_B\) cumple
\[
 \rho'_A=U_A\rho_AU_A^*,\qquad s(\rho'_A)=s(\rho_A).
\]
La primera identidad sigue de la traza parcial y la segunda del espectro.
Para una unitaria general, la formulación equivalente transporta también
el álgebra del subsistema:
\(\mathcal B(\mathcal H_A)\otimes I\mapsto
U(\mathcal B(\mathcal H_A)\otimes I)U^*\).
Así se conserva la misma pregunta operacional acerca de la información
accesible. Las igualdades de entropía se aplican a estados con las entropías
finitas pertinentes.

La relación entre información y gravitación queda expresada mediante
operaciones distintas y composables: el estado y la bipartición fijan la
información accesible; la incidencia y Barbero--Immirzi fijan la resolución
areal; acción y gravitación construyen la unidad \(\ell_P^2\);
Boltzmann realiza dimensionalmente la entropía; la carta de horizonte
transporta la geometría y sus observables. Las leyes finitas conservan la
entropía relativa, y las realizaciones geométricas conservan sus medidas,
su normalización temporal y sus dominios.

% Recuperación y articulación de los propietarios documentados en
% PROCEDENCIA_OBSERVADOR_RESPUESTA.md. Originales conservados íntegros.
% Antecedentes de integración: núcleo APP--TRIT--TPK, realización del
% transporte de borde, lectores areal/volumétrico y reloj energético.
% Las condiciones de representación se enuncian antes de cada uso.

\section{Observador, memoria y respuesta relacional}
\label{vii:sec:observador-respuesta}

La operación de lectura actúa sobre el transporte que publica el estado
enriquecido APP--TRIT--TPK. Su dominio conserva ruta, hoja, orientación,
acarreo y memoria; el resultado accesible al observador es una aplicación
de ese dominio. Esta distinción permite describir conjuntamente tres
operaciones: la elección de un representante orientado del transporte,
la inscripción de su registro en un aparato y la respuesta local inducida
por la frontera global. Cada operación transmite una información distinta.
La construcción siguiente especifica sus mapas y las condiciones que
permiten componerlos.

\subsection{Transporte de borde y lectura intrínseca}
\label{vii:sec:observador-ciclo}

La publicación finita del transporte proporciona un ciclo orientado
\(C_n\), con \(n\geq1\), vértices \(0,\ldots,n-1\) y aristas
\(e_j:j\longrightarrow j+1\pmod n\). Fijamos la realización reticular
\(L\) del registro, su espacio vectorial
\(W=L\otimes_{\mathbb Z}\mathbb R\), un producto interno positivo y
el representante primitivo no nulo \(v\in L\) de la clase residual
distinguida, en la cámara declarada. Las hojas suma y producto de APP
aportan las polarizaciones del transporte; la orientación del TRIT y la
actualización del TPK preceden a esta publicación lineal. El ciclo,
la realización reticular y la clase residual son los datos que esta
sección recibe del transporte construido anteriormente.

Una \(0\)-cocadena es una función sobre los vértices con valores en
\(W\); una \(1\)-cocadena asigna un vector a cada arista orientada.
El coborde actúa por
\[
 (d\Phi)(e_j)=\Phi(j+1)-\Phi(j).
\]
Denotamos por \(\mathbf1_E\) la cocadena escalar constante en las
aristas. Para un transporte \(\omega\in C^1(C_n;W)\), su publicación
acumulada \(\Gamma_\omega\) se obtiene por sumas parciales e
interpolación afín; su defecto de cierre es
\begin{equation}
 B(\omega)=\Gamma_\omega(1)-\Gamma_\omega(0)
 =\sum_{j=0}^{n-1}\omega(e_j).
 \label{vii:eq:observador-defecto}
\end{equation}

\begin{theorem}[Descomposición exacta y modo armónico del ciclo]
\label{vii:thm:observador-descomposicion}
Todo transporte tiene una única descomposición
\[
 \omega=d\Phi+u\mathbf1_E,\qquad \Phi(0)=0,
 \qquad u=\frac1n\sum_j\omega(e_j).
\]
En particular, \(B(d\Phi)=0\) y \(B(\omega)=nu\).
\end{theorem}

\begin{proof}
Sea \(y_j=\omega(e_j)-u\). La suma de los \(y_j\) es cero.
Definimos \(\Phi(k)=\sum_{j=0}^{k-1}y_j\); el valor para \(k=n\)
coincide con \(\Phi(0)=0\), de modo que la función está bien definida
en el ciclo, incluida su arista final. Se cumple \(d\Phi(e_j)=y_j\).
Si \(d\Phi+u\mathbf1_E=d\Psi+w\mathbf1_E\), la suma de aristas da
\(nu=nw\). Entonces \(d(\Phi-\Psi)=0\); la conectividad del ciclo
y la normalización en el vértice inicial implican \(\Phi=\Psi\).
La identidad \(\sum_j(\Phi(j+1)-\Phi(j))=0\) prueba directamente
la cancelación telescópica.
\end{proof}

\begin{definition}[Transporte admisible y lectura normalizada]
\label{vii:def:observador-admisible}
El transporte es admisible cuando su modo armónico pertenece a
\(\mathbb Rv\). Sobre el espacio \(\mathcal T_{\rm adm}\) de esos
transportes definimos la lectura global y su densidad por fase:
\begin{equation}
 \mathfrak a_v(\omega)
 =\frac{\langle B(\omega),v\rangle}{\langle v,v\rangle},
 \qquad
 \mu_v(\omega)=\frac{\mathfrak a_v(\omega)}n.
 \label{vii:eq:observador-lectura}
\end{equation}
Se tiene \(B(\omega)=\mathfrak a_v(\omega)v\).
\end{definition}

\begin{theorem}[Determinación de la lectura intrínseca]
\label{vii:thm:observador-unicidad}
El cociente
\(\mathcal T_{\rm adm}/dC^0(C_n;W)\) es unidimensional, generado
por la clase de \(v\mathbf1_E\). La aplicación \(\mu_v\) es el
único funcional lineal sobre \(\mathcal T_{\rm adm}\) que satisface
\[
 F(\omega+d\Phi)=F(\omega),\qquad F(v\mathbf1_E)=1.
\]
\end{theorem}

\begin{proof}
La descomposición anterior escribe cada transporte admisible como
\(d\Phi+\lambda v\mathbf1_E\). Su clase módulo cocadenas exactas
queda determinada por \(\lambda\). La clase de \(v\mathbf1_E\)
es no nula porque su suma es \(nv\neq0\). Por linealidad e
invariancia, todo funcional de las características indicadas vale
\(F(\omega)=\lambda F(v\mathbf1_E)=\lambda\). La fórmula
\eqref{vii:eq:observador-lectura} da precisamente \(\mu_v=\lambda\).
La lectura global es \(\mathfrak a_v=n\mu_v\).
\end{proof}

\begin{definition}[Observador de borde]
\label{vii:def:observador-borde}
Fijada la polarización del transporte publicado, un observador de borde
es un par \(\mathcal O=(\varepsilon,\Psi)\), donde
\(\varepsilon\in\{+1,-1\}\) y \(\Psi\in C^0(C_n;W)\).
Su acción de lectura es
\begin{equation}
 \omega^{\mathcal O}=\varepsilon\omega+d\Psi.
 \label{vii:eq:observador-accion}
\end{equation}
El signo fija la orientación y la cocadena fija la trivialización.
Dos elecciones de polarización se comparan mediante su transporte
publicado; la relación \eqref{vii:eq:observador-accion} describe
exactamente los representantes que difieren por orientación y coborde.
\end{definition}

\begin{proposition}[Covarianza de la lectura]
\label{vii:prop:observador-covarianza}
La acción conserva la admisibilidad y verifica
\[
 B(\omega^{\mathcal O})=\varepsilon B(\omega),\quad
 \mathfrak a_v(\omega^{\mathcal O})=\varepsilon\mathfrak a_v(\omega),
 \quad \mu_v(\omega^{\mathcal O})=\varepsilon\mu_v(\omega).
\]
Las trivializaciones con igual orientación publican el mismo cierre.
\end{proposition}

\begin{proof}
La suma del coborde \(d\Psi\) es cero. Por tanto, el modo armónico
se transforma en \(\varepsilon u\), que sigue perteneciendo a
\(\mathbb Rv\), y el defecto se transforma en \(\varepsilon B\).
Las dos igualdades escalares se obtienen por la proyección normalizada.
\end{proof}

La unidad de esta lectura es el representante residual \(v\), con su
normalización interna. La libertad de trivialización actúa sobre el
representante de la cocadena; la información de cierre reside en su
clase. La descripción permite comparar observadores manteniendo
separados el transporte genealógico, su publicación lineal y el escalar
que evalúa la clase residual.

\subsection{Identidad de borde e interior celular}
\label{vii:sec:observador-stokes}

Sea \(K\) un complejo celular finito orientado de dimensión dos,
provisto de una \(2\)-cadena fundamental \(c_K\) tal que
\(\partial c_K=[C_n]\). Esta condición expresa que las incidencias
interiores se cancelan con sus multiplicidades y orientaciones.
Sea \(\Omega\in C^1(K;W)\) una extensión de \(\omega\) y
\(F_\Omega=d\Omega\) su curvatura discreta aditiva.

\begin{theorem}[Identidad holográfica celular de la lectura]
\label{vii:thm:observador-stokes}
Con los datos anteriores,
\begin{equation}
 B(\omega)=\langle\Omega,\partial c_K\rangle
          =\langle d\Omega,c_K\rangle,
 \qquad
 \mathfrak a_v(\omega)
 =\frac{\langle\langle d\Omega,c_K\rangle,v\rangle}
        {\langle v,v\rangle}.
 \label{vii:eq:observador-stokes}
\end{equation}
En particular, \(F_\Omega=0\) implica \(\mathfrak a_v(\omega)=0\).
Si el cierre es no nulo, toda extensión del transporte en ese complejo
tiene curvatura discreta no nula.
\end{theorem}

\begin{proof}
El emparejamiento entre cadenas y cocadenas satisface
\(\langle d\Omega,c_K\rangle=\langle\Omega,\partial c_K\rangle\).
Al expandirlo, cada arista interior aparece con la suma de los números
de incidencia de las celdas contiguas; la condición sobre
\(\partial c_K\) deja exactamente el ciclo orientado de borde.
La restricción de \(\Omega\) a ese ciclo es \(\omega\), cuya suma
es \(B(\omega)\). La proyección sobre \(v\) da la segunda fórmula
y la implicación de planitud.
\end{proof}

La igualdad compara dos evaluaciones de la misma cocadena extendida:
el defecto acumulado del borde y el coborde integrado sobre el interior.
Su curvatura toma valores en \(W\) y usa el coborde celular aditivo.
La realización por una conexión geométrica conserva, adicionalmente,
su propio espacio de coeficientes y su ley de transporte.

\subsection{Inscripción isométrica del registro en el aparato}
\label{vii:sec:observador-aparato}

Sea \(\mathfrak M_n\) un conjunto finito no vacío de registros
publicados a nivel \(n\), con resultado, ruta, acarreo, hoja y memoria.
Su espacio de registros es
\(\mathcal H_n=\ell^2(\mathfrak M_n)\), con base \(|m\rangle\).
La actualización visible \(F_n\) tiene un levantamiento de registro
\[
 \widehat F_n(m)=(F_n(m),m).
\]
Tomamos como registros del nivel siguiente un conjunto que contiene esas
parejas. La segunda coordenada conserva el antecedente y hace inyectivo
el levantamiento. Esta construcción especifica la memoria retenida por
la inscripción, también cuando la actualización visible reúne valores.

En un espacio finito de aparato \(\mathcal K_{A,n}\), fijamos una
familia ortonormal \((|a_m\rangle)_{m\in\mathfrak M_n}\) y definimos
\[
 \iota_n:\mathcal H_n\longrightarrow\mathcal K_{A,n},
 \quad \iota_n|m\rangle=|a_m\rangle,
 \qquad
 V_n:\mathcal H_n\longrightarrow\mathcal H_{n+1},
 \quad V_n|m\rangle=|\widehat F_n(m)\rangle.
\]
Sea \(\mathcal P_n=\iota_n\mathcal H_n\) el subespacio de puntero.

\begin{theorem}[Compatibilidad entre actualización e inscripción]
\label{vii:thm:observador-entrelazador}
Los mapas \(\iota_n\) y \(V_n\) son isometrías. Existe una única
isometría \(W_n^0:\mathcal P_n\longrightarrow\mathcal K_{A,n+1}\)
que verifica
\begin{equation}
 W_n^0\iota_n=\iota_{n+1}V_n.
 \label{vii:eq:observador-cuadrado}
\end{equation}
Si \(\dim\mathcal K_{A,n+1}\geq\dim\mathcal K_{A,n}\),
se prolonga a una isometría de todo \(\mathcal K_{A,n}\).
Mediante complementos auxiliares que igualen las dimensiones, se obtiene
una prolongación unitaria entre los espacios ampliados.
\end{theorem}

\begin{proof}
Las imágenes de las bases por \(\iota_n\) son ortonormales; las
imágenes por \(V_n\) también lo son por la inyectividad de
\(\widehat F_n\). Para \(z=\iota_n\psi\) definimos
\(W_n^0z=\iota_{n+1}V_n\psi\). La representación de \(z\) es única,
y los tres mapas conservan su norma. La fórmula determina el mapa en
una base de \(\mathcal P_n\) y prueba el cuadrado.
Si la dimensión de llegada es suficiente, el complemento ortogonal de
\(W_n^0\mathcal P_n\) contiene una familia ortonormal del tamaño de
\(\mathcal P_n^\perp\); enviando una base de este último a aquella
se construye la prolongación. Finalmente se amplían ambos espacios a
una misma dimensión y se completan las dos familias prescritas a bases
ortonormales. La aplicación entre esas bases es unitaria.
\end{proof}

Cuando el registro incluye rótulos de un instrumento finito, esta misma
inscripción realiza su ambiente: si los operadores
\(K_{y,a}:\mathcal H_S\to\mathcal H'_S\) satisfacen
\(\sum_{y,a}K_{y,a}^*K_{y,a}=I\), la aplicación
\[
 V_{\rm inst}\psi=\sum_{y,a}K_{y,a}\psi\otimes|y,a\rangle
\]
es isométrica, pues su norma al cuadrado es
\(\sum_{y,a}\langle\psi,K_{y,a}^*K_{y,a}\psi\rangle=\|\psi\|^2\).
La inscripción \(\iota\) transporta cada rótulo al puntero que le
corresponde. El modelo fija así el registro de la premedición; un
dispositivo material requiere además su interacción y su calibración.

\subsection{Aplicación de salida y condicionamiento del registro}
\label{vii:sec:salida-condicionamiento}

La inscripción anterior se compone con el lector del aparato sobre el
estado genealógico completo. El acoplamiento conserva la operación que
produce la correlación; la aplicación de salida publica el resultado; el
condicionamiento restringe las prolongaciones compatibles con ese registro.

Sean \((\Omega_S,\Sigma_S)\) y \((\Omega_M,\Sigma_M)\) los espacios
medibles de estados genealógicos del sistema y del aparato, obtenidos de
sus cilindros y registros compatibles. Un contexto de medición se escribe
\(\mathfrak M_a=(U_a,q_a,\mathcal C_a,m_0)\): \(m_0\) es la
preparación del aparato, \(\mathcal C_a\) fija las compatibilidades y el
horizonte de prolongación, y
\[
 U_a:\Omega_S\times\Omega_M\longrightarrow\Omega_S\times\Omega_M,
 \qquad q_a:\Omega_M\longrightarrow Y_a
\]
son, respectivamente, el acoplamiento medible y el lector medible hacia
el espacio de resultados \((Y_a,\Sigma_{Y_a})\). En un régimen reversible,
\(U_a\) es biyectivo y bimedible; en la realización de Hilbert se exige
asimismo la unitariedad correspondiente. La salida individual es la composición
\begin{equation}
 \operatorname{Out}_a(x)
 =q_a\bigl(\operatorname{pr}_M U_a(x,m_0)\bigr).
 \label{vii:eq:out-genealogico}
\end{equation}
Queda determinada por el estado completo, la preparación del aparato y
el contexto. La descripción estadística de una preparación constituye
otra lectura del conjunto de estados compatibles.

\begin{proposition}[Medida de resultados y restricción a la fibra]
\label{vii:prop:medida-salida}
Sea \(\Omega_a\in\Sigma_S\) el conjunto de estados compatibles con la
preparación y el contexto, provisto de una probabilidad \(\mu_a\) sobre
\((\Omega_a,\Sigma_S|_{\Omega_a})\).
La aplicación \(\operatorname{Out}_a|_{\Omega_a}\) es medible y produce
la probabilidad de resultados
\begin{equation}
 \nu_a(B)=\mu_a\bigl(\operatorname{Out}_a^{-1}(B)\cap\Omega_a\bigr),
 \qquad B\in\Sigma_{Y_a}.
 \label{vii:eq:medida-imagen-salida}
\end{equation}
Para un resultado \(y\) con singleton medible, la fibra
\(\Omega_{a,y}=\{x\in\Omega_a:\operatorname{Out}_a(x)=y\}\)
es medible. Cuando \(\mu_a(\Omega_{a,y})>0\),
\begin{equation}
 \mu_{a,y}(E)=
 \frac{\mu_a(E\cap\Omega_{a,y})}{\mu_a(\Omega_{a,y})}
 \label{vii:eq:condicionamiento-salida}
\end{equation}
para \(E\in\Sigma_S|_{\Omega_a}\) define una probabilidad concentrada
en esa fibra.
\end{proposition}
\begin{proof}
La inclusión \(x\mapsto(x,m_0)\), el acoplamiento, la proyección sobre
el aparato y su lector son medibles; su composición también lo es.
Las preimágenes conservan uniones disjuntas y el conjunto total, por lo
que \eqref{vii:eq:medida-imagen-salida} es numerablemente aditiva y tiene
masa uno. La fibra es la preimagen del singleton \(\{y\}\) dentro de
\(\Omega_a\). La restricción de \(\mu_a\) a esa fibra es una medida
positiva y numerablemente aditiva; dividir por su masa positiva la normaliza.
\end{proof}

Para un prefijo \(s\) de profundidad \(n\), las prolongaciones posteriores
al registro forman
\[
 \operatorname{Fut}_{a,y}(s)
 =\{x\in\Omega_{a,y}:x|_n=s\}.
\]
La restricción conserva el prefijo ya recorrido y el registro del
acoplamiento; modifica la familia de prolongaciones admitidas. La salida
\eqref{vii:eq:out-genealogico} actúa sobre estados completos; la ley de
resultados \eqref{vii:eq:medida-imagen-salida} se obtiene transportando
la medida de preparación por esa aplicación. La representación por una
matriz de densidad expresa el nivel
estadístico de la realización cuántica y conserva sus condiciones de
correspondencia con las fibras genealógicas.


\subsection{Álgebra accesible y complemento del puntero}
\label{vii:sec:observador-expectativa}

Definimos \(P_m=|a_m\rangle\langle a_m|\),
\(P=\sum_mP_m\) y \(P_\perp=I-P\). El álgebra del subespacio de
puntero es \(P\mathcal B(\mathcal K_{A,n})P\), cuya unidad es \(P\).
En ella, el lector diagonal es
\[
 \mathcal D_P(A)=\sum_mP_mAP_m.
\]

\begin{proposition}[Expectativa del registro y extensión unital]
\label{vii:prop:observador-pinzamiento}
La aplicación \(\mathcal D_P\) es una expectativa condicional
completamente positiva, preservadora de traza e idempotente sobre el
álgebra del puntero, con \(\mathcal D_P(P)=P\). Su imagen es
\(\bigoplus_m\mathbb CP_m\).
Sobre el álgebra del aparato completo, la aplicación
\begin{equation}
 \mathcal E_P(A)=\sum_mP_mAP_m+P_\perp AP_\perp
 \label{vii:eq:observador-complemento}
\end{equation}
es una expectativa condicional unital, completamente positiva y
preservadora de traza. Su imagen es
\[
 \left(\bigoplus_m\mathbb CP_m\right)
 \oplus\mathcal B(\mathcal P_n^\perp).
\]
Ambas aplicaciones son proyecciones ortogonales para el producto de
Hilbert--Schmidt. Sus espectros están contenidos en \(\{0,1\}\).
Cuando existen componentes eliminadas, la separación entre el
subespacio fijo y el núcleo es uno.
\end{proposition}

\begin{proof}
Los proyectores \(P_m\) son mutuamente ortogonales. En el aparato
completo, la familia que añade \(P_\perp\) suma \(I\). Cada sumando
\(A\mapsto P_mAP_m\) es completamente positivo: tras cualquier
ampliación matricial actúa por congruencia con \(I\otimes P_m\).
Lo mismo ocurre para \(P_\perp\). La suma conserva esa propiedad.
La ciclicidad de la traza y la suma de los proyectores prueban la
preservación de traza en cada dominio. La ortogonalidad anula los
términos cruzados al iterar el mapa, por lo que es idempotente.
Sus elementos fijos son precisamente los bloques anunciados y la
aplicación es bimodular respecto de esa subálgebra.
Además,
\(\operatorname{Tr}(A^*P_mBP_m)
 =\operatorname{Tr}((P_mAP_m)^*B)\),
de donde resulta la autoadjunción para Hilbert--Schmidt. Una proyección
autoadjunta tiene únicamente los valores espectrales cero y uno cuando
ambos subespacios son no nulos. Si sólo queda el subespacio fijo,
la aplicación es la identidad.
\end{proof}

La suma \(\sum_mP_mAP_m\), aplicada sin el complemento a todo el
aparato, envía \(I\) a \(P\). Su unidad natural pertenece al álgebra
del puntero. La fórmula \eqref{vii:eq:observador-complemento} conserva
también el bloque auxiliar; una resolución ortonormal de ese bloque
proporciona, cuando se desea, la diagonalización del aparato entero.

Si \(U\) es una transformación unitaria del aparato y
\(P'_m=UP_mU^*\), entonces
\[
 \mathcal E_{P'}(UAU^*)=U\mathcal E_P(A)U^*.
\]
La igualdad se obtiene sustituyendo cada proyector en la suma.
Asimismo, para \(A\in\mathcal B(\mathcal P_n)\), el cuadrado
\eqref{vii:eq:observador-cuadrado} implica
\[
 \mathcal D_{P_{n+1}}(W_n^0A(W_n^0)^*)
 =W_n^0\mathcal D_{P_n}(A)(W_n^0)^*.
\]
En efecto, cada unidad matricial \(|a_m\rangle\langle a_{m'}|\)
se transporta a la unidad correspondiente a los registros distintos
\(\widehat F_n(m),\widehat F_n(m')\); el lector diagonal conserva
exactamente los términos con \(m=m'\).

\subsection{Factorización de la respuesta por la frontera efectiva}
\label{vii:sec:observador-factorizacion}

Sea \(\mathcal X_\infty^{\rm enr}\) el dominio de estados compatibles
enriquecidos y sea
\(b:\mathcal X_\infty^{\rm enr}\longrightarrow B_\infty\) su traza
global de frontera. Denotamos por \(B_{\rm ef}=\operatorname{im}b\)
la frontera efectivamente publicada. Para un espacio de respuestas
\(\mathcal V_v\), consideramos
\(I_v:\mathcal X_\infty^{\rm enr}\longrightarrow\mathcal V_v\).

\begin{theorem}[Criterio de respuesta relacional]
\label{vii:thm:observador-factorizacion}
Existe una única aplicación \(\overline I_v:B_{\rm ef}\to\mathcal V_v\)
tal que \(I_v=\overline I_v\circ b\) si y sólo si
\begin{equation}
 b(x)=b(x')\quad\Longrightarrow\quad I_v(x)=I_v(x').
 \label{vii:eq:observador-fibra}
\end{equation}
La unicidad se extiende a todo \(B_\infty\) cuando \(b\) es
sobreyectiva.
\end{theorem}

\begin{proof}
Una composición por \(b\) es constante en cada fibra. Recíprocamente,
para \(\beta\in B_{\rm ef}\), se define
\(\overline I_v(\beta)=I_v(x)\) con \(b(x)=\beta\).
La condición \eqref{vii:eq:observador-fibra} garantiza independencia
del representante; todo valor de \(\overline I_v\) queda determinado
por algún estado de esa fibra. Si \(b\) es sobreyectiva,
\(B_{\rm ef}=B_\infty\). En otro caso, la ecuación de composición
determina únicamente los valores sobre \(B_{\rm ef}\).
\end{proof}

Este criterio es conjuntista. Para una factorización continua, se
conservan además las topologías y la condición de aplicación cociente
de \(b:\mathcal X_\infty^{\rm enr}\to B_{\rm ef}\); con ella, la
continuidad de \(I_v\) equivale a la de su factor. Esta precisión
separa la determinación de una lectura de sus propiedades de regularidad.

Sea \(\ell_v:\mathcal X_\infty^{\rm enr}\to\mathcal L_v\) la
restricción accesible a una vecindad local. El carácter global efectivo
se verifica mediante dos estados que satisfagan
\begin{equation}
 \ell_v(x)=\ell_v(x'),\qquad I_v(x)\neq I_v(x').
 \label{vii:eq:observador-testigo-global}
\end{equation}
Si \(I_v\) factoriza por \(b\), estos dos estados tienen fronteras
distintas. Además, la respuesta no puede escribirse como función de
\(\ell_v\), pues esa función asignaría el mismo valor a ambos estados.
La diferencia de fronteras adquiere así un efecto operatorio comprobable.

\subsection{Pesos de ambivalencia y respuesta autoinercial}
\label{vii:sec:observador-ambivalencia}

Para la fórmula general consideramos lectores positivos
\(A_\partial,V_\partial\) y un lector energético \(Q\)
definidos sobre el dominio efectivo de frontera declarado.
A continuación se explicita una realización sobre
\(B_{\rm av}\subseteq B_{\rm ef}\). La memoria permanece como coordenada
\(\mathsf M_\partial(\beta)\) de la frontera. La bisagra de
ambivalencia proporciona las dos orientaciones
\[
 r_{\rm av}=\frac{\pi_{\rm HMT}}{\varphi_{\rm HMT}^2},\qquad
 r_{\rm av}^{J}=\frac{\varphi_{\rm HMT}}{\pi_{\rm HMT}^2}.
\]
El superíndice \(J\) expresa intercambio del par generador.
\paragraph{Realización normalizada de los lectores de frontera.}
La incidencia modal APP--TRIT--TPK, desarrollada en la construcción
regional, proporciona el envolvente de caminos de matriz
\[
 F_{\rm av}=\begin{pmatrix}0&1\\1&1\end{pmatrix}.
\]
El vector positivo \(v=(1,\varphi_{\rm HMT})^{\mathsf T}\)
satisface \(F_{\rm av}v=\varphi_{\rm HMT}v\). Su transición
normalizada y sus pesos estacionarios son
\begin{equation}
 P_{ij}=\frac{(F_{\rm av})_{ij}v_j}{\varphi_{\rm HMT}v_i},
 \qquad
 P=\begin{pmatrix}0&1\\
 \varphi_{\rm HMT}^{-2}&\varphi_{\rm HMT}^{-1}\end{pmatrix},
 \qquad
 (\mu_{\rm ar},\mu_{\rm vol})
 =\frac{(1,\varphi_{\rm HMT}^{2})}{1+\varphi_{\rm HMT}^{2}}.
 \label{vii:eq:lectores-parry}
\end{equation}
Las filas de \(P\) suman uno por
\(\varphi^{-2}+\varphi^{-1}=1\); la ecuación estacionaria
\(\mu_{\rm ar}=\mu_{\rm vol}\varphi^{-2}\), junto con la
normalización, determina los pesos escritos. Esta medida de Parry
pondera los caminos del envolvente; la frecuencia cronológica
\((1/3,2/3)\) cuenta los instantes del calendario modal y conserva
su función distinta.

La realización de amplitud común de memoria se define sobre el sector
\(\mathcal X_{\rm av}\subset\mathcal X_\infty^{\rm enr}\) en el que
los lectores de las dos hojas tienen la factorización
\begin{equation}
 \mathcal L_{\rm vol}(x)=\mu_{\rm vol}\mathcal M(x),\qquad
 \mathcal L_{\rm ar}(x)=\pi_{\rm HMT}\mu_{\rm ar}\mathcal M(x),
 \qquad \mathcal M(x)>0.
 \label{vii:eq:lectores-amplitud-comun}
\end{equation}
Ambos valores pertenecen a la misma recta positiva de amplitud de
memoria. La publicación areal lleva la marca regional
\(\pi_{\rm HMT}\) del retorno. La lectura adimensional se obtiene
dividiendo ambas secciones por \(\mathcal L_{\rm vol}(x)\).
De este modo las sumas de los pesos se efectúan entre magnitudes del
mismo tipo; una eventual realización como áreas y volúmenes físicos
se compone después con sus respectivos transductores dimensionales.

\begin{proposition}[Descenso de los lectores normalizados y del reloj energético]
\label{vii:prop:lectores-frontera-explicitos}
Sobre \(B_{\rm av}=b(\mathcal X_{\rm av})\), las reglas
\begin{equation}
 A_\partial(b(x))
 =\frac{\mathcal L_{\rm ar}(x)}{\mathcal L_{\rm vol}(x)}
 =r_{\rm av},\qquad V_\partial(b(x))=1
 \label{vii:eq:lectores-normalizados}
\end{equation}
definen dos lectores positivos independientes del representante.
Fijada una orientación \(a\) de la sección de acción, la energía de
decisión de \eqref{v:ant:eq:ii-kb-aut-energia} determina sobre estados
\begin{equation}
 q_a(x)=\frac{h_a(x)}{108t_0(x)}
       =\frac{\hbar_a(x)}{t_P(x)}\in\mathcal L_E^+,
 \qquad t_P=\frac{54}{\pi_{\rm HMT}}t_0.
 \label{vii:eq:lector-q-accion}
\end{equation}
Escribimos
\[
 b_{\rm av}:=b|_{\mathcal X_{\rm av}}:
 \mathcal X_{\rm av}\longrightarrow B_{\rm av}.
\]
Existe un único lector \(Q:B_{\rm av}\to\mathcal L_E^+\) tal que
\(q_a|_{\mathcal X_{\rm av}}=Q\circ b_{\rm av}\)
si y sólo si, para cualesquiera \(x,x'\in\mathcal X_{\rm av}\),
\begin{equation}
 b(x)=b(x')\ \Longrightarrow\
 \frac{h_a(x)}{t_0(x)}=\frac{h_a(x')}{t_0(x')}.
 \label{vii:eq:descenso-reloj}
\end{equation}
En particular, una carta de acción y duración fijas satisface esta
condición y publica \(Q=h_a/(108t_0)\).
\end{proposition}
\begin{proof}
El cociente de \eqref{vii:eq:lectores-amplitud-comun} cancela la
amplitud común y vale
\(\pi_{\rm HMT}\mu_{\rm ar}/\mu_{\rm vol}
=\pi_{\rm HMT}/\varphi_{\rm HMT}^2\).
El resultado es el mismo para todos los representantes de una fibra
de \(b_{\rm av}\) dentro de \(\mathcal X_{\rm av}\) y es positivo.
Dividir ambos lectores por el mismo factor positivo
conserva los cocientes ponderados en los que intervienen.
Para el reloj, el criterio de factorización del
teorema~\ref{vii:thm:observador-factorizacion}, aplicado a \(q_a\),
equivale exactamente a \eqref{vii:eq:descenso-reloj}. Las secciones
fijas lo satisfacen; también lo satisfacen secciones variables cuyo
cociente descienda por \(b_{\rm av}\). La unicidad resulta de la definición
de \(B_{\rm av}\) como imagen efectiva.
\end{proof}

El valor del lector \(Q\) pertenece a la recta de energía
\(\mathcal L_E=\mathcal L_{\rm act}\otimes\mathcal L_T^*\).
Con \(E_1,E_2\) en esa misma recta,
\(E_1E_2/Q^2\) es adimensional e independiente de la base energética
positiva. El símbolo \(Q\) de esta sección designa exclusivamente
esa energía de reloj; la cantidad de calor de la realización de
Landauer pertenece al balance térmico correspondiente.
La factorización de amplitud común especifica una realización
concreta de los lectores generales. Cada otro dominio conserva sus
aplicaciones de frontera y su propia comprobación del descenso.
La memoria completa sigue en la historia enriquecida y en su
publicación de frontera, aunque se cancele en el cociente normalizado.


Definimos
\[
 w_A(\beta)=
 \frac{A_\partial(\beta)r_{\rm av}}
 {A_\partial(\beta)r_{\rm av}+V_\partial(\beta)r_{\rm av}^{J}},
 \qquad w_V(\beta)=1-w_A(\beta).
\]
En la realización \eqref{vii:eq:lectores-normalizados}, los pesos
admiten la evaluación exacta
\begin{equation}
 w_A=\frac{r_{\rm av}^{\,2}}{r_{\rm av}^{\,2}+r_{\rm av}^{J}}
     =\frac{\pi_{\rm HMT}^{4}}
            {\pi_{\rm HMT}^{4}+\varphi_{\rm HMT}^{5}},\qquad
 w_V=\frac{\varphi_{\rm HMT}^{5}}
            {\pi_{\rm HMT}^{4}+\varphi_{\rm HMT}^{5}}.
 \label{vii:eq:pesos-amplitud-comun}
\end{equation}
La sustitución conserva los dos factores de la definición: el lector
areal normalizado y su ponderación posterior por \(r_{\rm av}\).
La generación y el marcado regional del retorno se realizan una vez.
Multiplicar numerador
y denominador por \(\varphi_{\rm HMT}^4\pi_{\rm HMT}^2\)
da la última expresión. Estos pesos permanecen constantes en el
sector de amplitud común. Con proyectores y energías fijos, una
variación de la respuesta puede proceder del lector \(Q\);
el testigo de dependencia global requiere además los estados con
igual restricción local y distinta respuesta de
\eqref{vii:eq:observador-testigo-global}.


Sean \(P_A,P_V\) proyecciones ortogonales complementarias de un
espacio de respuesta \(\mathcal H_v\), y sean \(E_1,E_2>0\)
dos energías internas fijadas para la comparación. La respuesta es
\begin{equation}
 \overline I_v(\beta)=\frac{E_1E_2}{Q(\beta)^2}
       \bigl(w_A(\beta)P_A+w_V(\beta)P_V\bigr),
 \qquad I_v=\overline I_v\circ b.
 \label{vii:eq:observador-respuesta-ponderada}
\end{equation}
También se admiten energías que sean lectores explícitos de
\(\beta\); en ese caso se conservan como tales en la fórmula.

\begin{proposition}[Positividad y covarianza de la respuesta]
\label{vii:prop:observador-respuesta-positiva}
La respuesta \eqref{vii:eq:observador-respuesta-ponderada} es
autoadjunta, positiva definida y constante en las fibras de \(b\).
Sus valores propios en los sectores no nulos son
\[
 \lambda_A=E_1E_2Q^{-2}w_A,\qquad
 \lambda_V=E_1E_2Q^{-2}w_V.
\]
El intercambio completo
\((A_\partial,r_{\rm av},P_A)
 \leftrightarrow(V_\partial,r_{\rm av}^{J},P_V)\)
permuta los dos sumandos y conserva el operador. Un transporte
unitario del espacio de respuesta conjuga ese operador y conserva
su espectro con multiplicidades.
\end{proposition}

\begin{proof}
La positividad de los lectores da \(0<w_A,w_V<1\). Para
\(z=z_A+z_V\), con \(z_A=P_Az\), \(z_V=P_Vz\),
\[
 \langle z,\overline I_vz\rangle
 =\frac{E_1E_2}{Q^2}
   (w_A\|z_A\|^2+w_V\|z_V\|^2)>0
\]
si \(z\neq0\). La descomposición ortogonal prueba las expresiones
espectrales. La dependencia exclusiva de \(\beta\) da constancia
en las fibras y la factorización. El intercambio de las ternas
intercambia numeradores, pesos y proyectores a la vez, por lo que
conserva la suma. Para un transporte unitario \(U\), la sustitución
\(P_A\mapsto UP_AU^*\), \(P_V\mapsto UP_VU^*\) produce
\(U\overline I_vU^*\), con el mismo espectro.
\end{proof}

Un intercambio implementado dentro de un mismo espacio por una unitaria
que envíe \(P_A\) a \(P_V\) requiere que ambos sectores tengan
igual dimensión. El intercambio de las dos ternas como rótulos de la
fórmula es válido con las dimensiones propias de cada sector.

La condición \eqref{vii:eq:observador-testigo-global} exige comparar
las respuestas. Por ejemplo, multiplicar simultáneamente
\(A_\partial,V_\partial\) por un mismo factor positivo conserva
ambos pesos; si también se conservan \(E_1,E_2,Q\), la respuesta
permanece idéntica. En cambio, con iguales pesos y energías, sustituir
\(Q\) por \(qQ\), \(q>0\), multiplica la respuesta por \(q^{-2}\).
Si dos estados con igual restricción local realizan ese cambio con
\(q\neq1\), proporcionan el testigo global. Estas son condiciones
operatorias de separación; la existencia de los estados se comprueba
en el dominio de frontera que se esté realizando.

\subsection{Conexión del estado conjunto y aceleración publicada}
\label{vii:sec:observador-autoinercia}

En una realización diferenciable del estado conjunto
\(\Gamma=(x,\mathcal O,m)\), sean \(N\) su espacio de realización,
\(\nabla^{\rm tot}\) su conexión y \(\pi_S:N\to N_S\) una
proyección diferenciable hacia el subsistema, dotado de conexión
\(\nabla^S\). La conexión y la proyección conservan la procedencia
de los transportes y lectores que realizan. El marco observado se
obtiene mediante
\[
 \mathfrak F_{\mathcal O}(\Gamma)
 =\operatorname{Res}_{\mathcal O}(\Gamma,\nabla^{\rm tot}).
\]
Para una trayectoria dos veces diferenciable \(\Gamma(t)\),
escribimos \(\gamma_S=\pi_S\circ\Gamma\).

\begin{definition}[Autoinercia y defecto de proyección]
\label{vii:def:observador-autoinercia}
Una trayectoria es autoinercial respecto de la conexión declarada
cuando \(\nabla^{\rm tot}_{\dot\Gamma}\dot\Gamma=0\).
El defecto de su proyección es
\begin{equation}
 \mathcal K_{\pi_S}(\dot\Gamma,\dot\Gamma)
 =\nabla^S_{\dot\gamma_S}\dot\gamma_S
  -d\pi_S\bigl(\nabla^{\rm tot}_{\dot\Gamma}\dot\Gamma\bigr).
 \label{vii:eq:observador-defecto-proyeccion}
\end{equation}
\end{definition}

\begin{proposition}[Aceleración del subsistema]
\label{vii:prop:observador-aceleracion}
Para una trayectoria autoinercial,
\[
 a_S=\nabla^S_{\dot\gamma_S}\dot\gamma_S
     =\mathcal K_{\pi_S}(\dot\Gamma,\dot\Gamma).
\]
En coordenadas \(x^B\) de \(N\) y \(y^a\) de \(N_S\), el
defecto es cuadrático en la velocidad, con componentes
\begin{equation}
 (\mathcal K_{\pi_S})^a{}_{BC}
 =\partial_B\partial_C\pi_S^a
  +(\Gamma^S)^a{}_{bc}\,
        \partial_B\pi_S^b\partial_C\pi_S^c
  -\partial_D\pi_S^a(\Gamma^{\rm tot})^D{}_{BC}.
 \label{vii:eq:observador-defecto-coordenadas}
\end{equation}
\end{proposition}

\begin{proof}
La regla de la cadena da
\(\ddot\gamma_S^a=\partial_B\pi_S^a\ddot\Gamma^B
 +\partial_B\partial_C\pi_S^a\dot\Gamma^B\dot\Gamma^C\).
Sumando la contribución de \(\nabla^S\) y restando la imagen de
la aceleración total, los términos que contienen \(\ddot\Gamma\)
se cancelan. Los restantes son
\eqref{vii:eq:observador-defecto-coordenadas} contraídos con
\(\dot\Gamma^B\dot\Gamma^C\). Al anularse la aceleración total
resulta la igualdad anunciada.
\end{proof}

La respuesta geométrica factoriza por \(b\) cuando la conexión,
la proyección y el dato cinemático utilizado en la evaluación se
determinan por \(\beta=b(x)\), o cuando esos datos cinemáticos se
mantienen fijados en la comparación. Con esas condiciones,
\eqref{vii:eq:observador-defecto-coordenadas} da el mismo resultado
para cualesquiera representantes de una fibra; el
teorema~\ref{vii:thm:observador-factorizacion} construye su factor
único sobre \(B_{\rm ef}\). Si varía el dato cinemático, éste
forma parte del dominio de la respuesta que se compara.

La formulación autoinercial describe una trayectoria geodésica del
estado conjunto y una aceleración efectiva de su subsistema debida
al defecto de proyección. La realización gravitatoria identifica
después esa conexión, el marco y la respuesta con sus observables
físicos. La covarianza del observador, la inscripción de memoria y
la factorización por la frontera conservan así funciones complementarias:
la primera determina la clase de lectura, la segunda retiene la
genealogía y la tercera especifica la dependencia relacional de la
respuesta.

% Recuperación de 13_accion_interaccion_absorbedor y83_operador_absorbedor.
% Desarrollo del canal cronológico desde el transporte de rutas.
% Procedencia y condiciones en PROCEDENCIA_PROPAGACION_BIDIRECCIONAL.md.
% Los originales y la sección80 permanecen íntegros. Sin contenido de radión.

\section{Propagación bidireccional y condición de absorción en la frontera}
\label{vii:sec:propagacion-bidireccional}

El transporte APP--TRIT--TPK conserva la orientación de la ruta y la memoria
que distingue su ida de su retorno. Sobre ese antecedente se construyen tres
operaciones sucesivas: la acción de una historia enriquecida, la propagación
de una fuente a través de los transportes realizados y la imposición de una
condición común en la frontera. La inscripción del registro en un espacio de
aparato, establecida en la sección~\ref{vii:sec:observador-aparato}, permite
expresar esas operaciones mediante aplicaciones lineales sin confundir la
publicación observable con el estado que conserva sus antecedentes.

La dirección temporal, la conjugación de hojas y la reflexión espacial
tienen dominios distintos. Sus relaciones se fijan primero sobre el registro;
después se especifican los entrelazamientos que utiliza una realización
dinámica. Este orden determina qué información permanece en los propagadores
y qué condición introduce la elección de una frontera absorbente.

\subsection{Acción de rutas e involuciones del registro}
\label{vii:subsec:accion-bidireccional}

Una ruta finita contiene sus posiciones, direcciones cardinales
\((d_1,\ldots,d_N)\), hojas \((\mu_1,\ldots,\mu_N)\), con
\(\mu_t\in\{\Sigma,\Pi\}\), y los incrementos de su registro
enriquecido. Para \(\lambda\geq0\), definimos
\begin{equation}
 \mathcal A_\lambda(\gamma)
 =\sum_{t=2}^N
 \bigl(1-\delta_{d_t,d_{t-1}}
       +\lambda(1-\delta_{\mu_t,\mu_{t-1}})\bigr).
 \label{vii:eq:accion-ruta-bidireccional}
\end{equation}
El primer sumando cuenta cambios de dirección; el segundo cuenta cambios
de hoja con el peso declarado. Para \(N\leq1\), la suma es cero.
La acción general del sector conserva además sus términos de frontera y
memoria:
\begin{equation}
 \mathcal S[\gamma]
 =\sum_{a\subset\gamma}\mathcal L_a
  +\mathcal B_{\rm in}(\partial_-\gamma)
  +\mathcal B_{\rm out}(\partial_+\gamma)
  +\mathcal M(\gamma).
 \label{vii:eq:accion-enriquecida-frontera}
\end{equation}
Los términos de esta expresión especifican una clase variacional con
extremos tipados y memoria conservada. La dinámica del sector determina
\(\mathcal L_a\), las condiciones admisibles y el término \(\mathcal M\).

Sobre los registros definimos \(\mathcal C_{\rm r}\) como el intercambio
global de hojas; \(\mathcal P_{\rm r}\) como una reflexión de la carta
espacial y de sus direcciones; y \(\mathcal T_{\rm r}\) como la inversión
del orden temporal, la sustitución de cada dirección por su opuesta y el
intercambio de extremos. La reflexión cardinal es una permutación involutiva.
La inversión del registro firmado es
\(C_M(a_1,\ldots,a_N)=(-a_N,\ldots,-a_1)\).
Las hojas, los cocientes y la memoria se transportan con la ruta completa;
la lectura escalar de una palabra constituye una operación posterior.

\begin{proposition}[Invariancia de la acción local de ruta]
\label{vii:prop:invariancia-ruta-bidireccional}
En el dominio de registros anterior,
\[
 \mathcal A_\lambda(\gamma)
 =\mathcal A_\lambda(\mathcal C_{\rm r}\gamma)
 =\mathcal A_\lambda(\mathcal P_{\rm r}\gamma)
 =\mathcal A_\lambda(\mathcal T_{\rm r}\gamma)
 =\mathcal A_\lambda(\mathcal C_{\rm r}\mathcal P_{\rm r}
                              \mathcal T_{\rm r}\gamma).
\]
\end{proposition}
\begin{proof}
Una permutación global del alfabeto de direcciones conserva la igualdad
entre direcciones consecutivas. El intercambio de las dos hojas conserva
la igualdad entre hojas consecutivas. La inversión del orden establece una
biyección entre pares adyacentes, y la sustitución por direcciones opuestas
también conserva sus igualdades. Cada operación conserva los dos recuentos
de \eqref{vii:eq:accion-ruta-bidireccional}; su composición los conserva.
\end{proof}

La proposición determina la simetría del funcional local. Para que una
transformación actúe además sobre un problema dinámico con fronteras fijas,
debe transportar su clase de prolongaciones admisibles y los otros tres
términos de \eqref{vii:eq:accion-enriquecida-frontera}. En particular, dos
rutas con igual publicación final pueden conservar distintas clases
variacionales por su memoria.

Una interacción entre dos secciones realizadas en una frontera \(B\)
utiliza un mapa de acoplamiento
\[
 \mathcal I_B:E_B^{(1)}\otimes E_B^{(2)}\longrightarrow E_B^{(12)}
\]
y un término \(\mathcal S_{\rm int}\) de la misma acción. Su covariancia
se expresa por
\[
 \rho_{12}(g)\mathcal I_B
 =\mathcal I_B\bigl(\rho_1(g)\otimes\rho_2(g)\bigr).
\]
Al componer las historias, el registro de la frontera compartida se cuenta
una vez; las memorias de ambos lados se transportan a la sección compuesta.
El acoplamiento y esta condición de covariancia son datos de la interacción
concreta, distintos de las involuciones del registro.

\subsection{Reversión de la evolución realizada}
\label{vii:subsec:reversion-evolucion}

La realización del registro proporciona un espacio de Hilbert real
\(\mathcal H_{\mathbb R}\), una estructura compleja ortogonal
\(J_E\), con \(J_E^2=-I\), y un Hamiltoniano autoadjunto \(H\).
Su dominio \(\mathcal D(H)\) es invariante por \(J_E\); se supone
\(HJ_E=J_EH\) en ese dominio. Sea \(\mathcal C\) una involución
ortogonal que conserva \(\mathcal D(H)\) y satisface
\begin{equation}
 \mathcal C J_E\mathcal C^{-1}=-J_E,\qquad
 \mathcal C H\mathcal C^{-1}=H.
 \label{vii:eq:dominio-reversion-evolucion}
\end{equation}
La constante interna de acción \(\hbar_{\rm int}>0\) conserva su
generación anterior; aquí normaliza el grupo de evolución.

\begin{theorem}[Conjugación temporal del grupo]
\label{vii:thm:reversion-evolucion}
El generador \(A=-J_EH/\hbar_{\rm int}\), con dominio \(\mathcal D(H)\),
es antiadjunto. Su grupo ortogonal, unitario respecto de \(J_E\), cumple
\begin{equation}
 U(t)=\exp(-J_EHt/\hbar_{\rm int}),\qquad
 \mathcal C U(t)\mathcal C^{-1}=U(-t).
 \label{vii:eq:reversion-grupo-bidireccional}
\end{equation}
\end{theorem}
\begin{proof}
La ortogonalidad y \(J_E^2=-I\) implican \(J_E^*=-J_E\).
La invariancia del dominio y la conmutación con \(H\) permiten tratar
\(H\) como operador autoadjunto complejo en el espacio definido por
\(J_E\). Por cálculo funcional, \(-J_EH/\hbar_{\rm int}\) es
antiadjunto y genera el grupo anunciado. En dimensión finita, las mismas
afirmaciones siguen de la serie exponencial y de \(A^*=-A\).
Las hipótesis sobre \(\mathcal C\) dan
\(\mathcal C A\mathcal C^{-1}=-A\). Conjugar el grupo, o su cálculo
funcional, produce \(\exp(-tA)=U(-t)\).
\end{proof}

El mapa que realiza la involución firmada \(C_M\) como \(\mathcal C\)
conserva la orientación temporal y el dominio anterior. La conjugación de
carga de una representación de campos conserva, además, sus propios
caracteres y operadores. La igualdad
\eqref{vii:eq:reversion-grupo-bidireccional} fija la reversión temporal
del grupo bajo las condiciones expresadas.

\subsection{Propagadores del canal cronológico de una ruta}
\label{vii:subsec:green-cronologico}

La inscripción isométrica del registro admite las prolongaciones unitarias
del teorema~\ref{vii:thm:observador-entrelazador}. Fijemos una realización
bilateral de un canal cronológico con fibras \(\mathcal H_n\),
\(n\in\mathbb Z\), y transportes unitarios
\(U_n:\mathcal H_n\to\mathcal H_{n+1}\).
La elección incluye los espacios auxiliares cuando se requieren; conserva
la ruta y sus registros anteriores. Denotamos por \(U(n,k)\) el transporte
de \(k\) a \(n\), con
\[
 U(k,k)=I,\qquad U(n,k)U(k,l)=U(n,l),\qquad
 U(n,k)^*=U(k,n).
\]
Para \(n>k\), es el producto \(U_{n-1}\cdots U_k\); para \(n<k\),
es el inverso del transporte de \(n\) a \(k\).

El espacio de fuentes \(\mathcal J_c\) contiene las secciones de soporte
finito. El espacio de respuestas \(\mathcal F\) contiene todas las
secciones de esas fibras. La acción cuadrática de la ruta, en unidades
internas de esta realización, utiliza
\[
 (D\psi)_n=\psi_{n+1}-U_n\psi_n,\qquad
 E[\psi]=\sum_n\|(D\psi)_n\|^2
\]
para secciones de soporte finito. El operador cinético con la convención
de signo \(L=-D^*D\) es
\begin{equation}
 (L\psi)_n=U_n^*\psi_{n+1}-2\psi_n+U_{n-1}\psi_{n-1}.
 \label{vii:eq:diferencia-covariante-bidireccional}
\end{equation}
En efecto, \((D^*z)_n=z_{n-1}-U_n^*z_n\), y la sustitución produce
esa fórmula. La variación de \(E\) es
\(-2\operatorname{Re}\langle\delta\psi,L\psi\rangle\).

\begin{theorem}[Green exactos de la diferencia covariante]
\label{vii:thm:green-cronologico}
Escribiendo \(x_+=\max(x,0)\), las aplicaciones
\begin{align}
 (G_{\rm ret}j)_n&=\sum_k(n-k)_+U(n,k)j_k,\label{vii:eq:green-ret-cronologico}\\
 (G_{\rm adv}j)_n&=\sum_k(k-n)_+U(n,k)j_k\label{vii:eq:green-adv-cronologico}
\end{align}
de \(\mathcal J_c\) en \(\mathcal F\) satisfacen
\[
 LG_{\rm ret}j=LG_{\rm adv}j=j.
\]
La respuesta retardada se anula antes y en la primera posición de la
fuente; la avanzada se anula después y en la última. Son las únicas
soluciones que se anulan suficientemente lejos en la dirección respectiva.
Para \(\psi\) de soporte finito también se cumple
\(G_{\rm ret}L\psi=G_{\rm adv}L\psi=\psi\).
Respecto del emparejamiento con fuentes de soporte finito,
\(G_{\rm adv}=G_{\rm ret}^{\dagger}\).
\end{theorem}
\begin{proof}
Cada suma tiene tantos términos como posiciones de la fuente.
La composición de los transportes reduce la aplicación de \(L\) a
la identidad escalar
\[
 (n+1-k)_+-2(n-k)_++(n-1-k)_+=\delta_{nk}.
\]
La función \((k-n)_+\) satisface la misma identidad. Esto prueba
las dos ecuaciones y los soportes. Para la unicidad, una solución
homogénea que se anula en dos posiciones consecutivas queda determinada
como cero en la siguiente por \eqref{vii:eq:diferencia-covariante-bidireccional};
la recurrencia, usando transportes invertibles, se prolonga en ambas
direcciones. La diferencia de dos soluciones con el mismo soporte causal
se anula suficientemente lejos y, por tanto, es cero.
Aplicar este argumento a \(\psi\) y a \(G_{\rm ret}L\psi\), o al
caso avanzado, da las identidades por la derecha sobre soporte finito.
Finalmente, para dos fuentes compactas, la suma doble del emparejamiento
es finita y
\[
 \bigl((k-n)_+U(k,n)\bigr)^*=(k-n)_+U(n,k).
\]
Intercambiar los índices demuestra la adjunción indicada.
\end{proof}

Las respuestas pertenecen a \(\mathcal F\); en general crecen
linealmente fuera de la fuente y no pertenecen a \(\ell^2\).
El dominio de soporte finito y su emparejamiento especifican la adjunción
de Green utilizada aquí. El resultado construye la propagación del canal
cronológico desde su acción; una realización de campos espaciales conserva
además su operador, geometría y dominio de propagación propios.

Sea \(C_0\) una involución isométrica, lineal o antilineal, de
\(\mathcal H_0\). El transporte desde la fibra de referencia construye
\[
 C_n=U(-n,0)C_0U(0,n):\mathcal H_n\longrightarrow\mathcal H_{-n},
 \qquad (\mathcal C\psi)_n=C_{-n}\psi_{-n}.
\]
Se tiene \(C_{-n}C_n=I\) y
\(C_{n+1}U_n=U_{-n-1}^*C_n\), por composición.
En particular, \(\mathcal C^2=I\), y el cambio \((n,k)\mapsto(-n,-k)\)
en las sumas anteriores prueba
\begin{equation}
 \mathcal C L=L\mathcal C,\qquad
 \mathcal C G_{\rm ret}\mathcal C^{-1}=G_{\rm adv}.
 \label{vii:eq:intercambio-green-cronologico}
\end{equation}
Cuando \(C_0\) es la realización declarada de la involución del registro,
estas fórmulas transportan esa involución a todo el canal, sin identificar
el retorno de fase con un reinicio de memoria.

\subsection{Intercambio causal y descomposición de la respuesta}
\label{vii:subsec:green-realizacion-campos}

En una realización de campos, sea
\(L:\mathcal D(L)\subset\mathcal H_\Phi\to\mathcal H_J\)
un operador con Green retardado y avanzado únicos en los dominios de
fuentes admisibles. Denotamos por \(J^+(S)\) y \(J^-(S)\) los
dominios causales futuro y pasado de un soporte \(S\). Las condiciones son
\[
 \begin{gathered}
 LG_{\rm ret}=LG_{\rm adv}=I,\\
 \operatorname{supp}(G_{\rm ret}j)\subset J^+(\operatorname{supp}j),\\
 \operatorname{supp}(G_{\rm adv}j)\subset J^-(\operatorname{supp}j).
 \end{gathered}
\]
Sean \(C_\Phi,C_J\) involuciones isométricas, ambas lineales o ambas
antilineales, que transportan esos dominios,
invierten sus orientaciones causales y satisfacen
\(LC_\Phi=C_JL\). Se conserva un emparejamiento de Green para el cual
\(G_{\rm adv}=G_{\rm ret}^{\dagger}\).

\begin{proposition}[Conjugación de los Green y paridades]
\label{vii:prop:green-paridades}
Con las condiciones anteriores,
\[
 C_\Phi G_{\rm ret}C_J^{-1}=G_{\rm adv}.
\]
Sobre el espacio real de operadores de respuesta, definamos
\[
 \Theta(K)=C_\Phi K C_J^{-1},\qquad
 \mathsf P_\pm(K)=\tfrac12(K\pm\Theta(K)).
\]
Son proyecciones complementarias para la involución \(\Theta\), y
\begin{equation}
 G_{\rm sym}=\tfrac12(G_{\rm ret}+G_{\rm adv}),\qquad
 G_{\rm rad}=\tfrac12(G_{\rm ret}-G_{\rm adv})
 \label{vii:eq:descomposicion-green}
\end{equation}
satisfacen \(\Theta(G_{\rm sym})=G_{\rm sym}\),
\(\Theta(G_{\rm rad})=-G_{\rm rad}\),
\(G_{\rm sym}^{\dagger}=G_{\rm sym}\) y
\(LG_{\rm rad}=0\). El soporte de la respuesta simétrica está contenido
en la unión de los dos dominios causales.
\end{proposition}
\begin{proof}
El operador conjugado resuelve la ecuación de fuente, pues
\[
 LC_\Phi G_{\rm ret}C_J^{-1}=C_JLG_{\rm ret}C_J^{-1}=I.
\]
La inversión de la orientación lleva su soporte al dominio avanzado;
su unicidad lo identifica con \(G_{\rm adv}\). Las involuciones dan
\(\Theta^2=I\); expandir \((I\pm\Theta)^2\) y
\((I+\Theta)(I-\Theta)\) demuestra las afirmaciones sobre proyecciones.
La adjunción de Green demuestra la simetría y la diferencia de las dos
ecuaciones de fuente da \(LG_{\rm rad}=0\). La suma de dos respuestas
tiene soporte contenido en la unión de sus soportes.
\end{proof}

Si las involuciones son antilineales, \(\Theta\) es lineal sobre
escalares reales; ése es el espacio operatorio usado en la proposición.
La conjugación del grupo temporal, la condición de soporte y la
adjunción de Green son tres propiedades distintas, transmitidas por los
mapas de realización correspondientes.

La reducción causal de la actualización bilateral, desarrollada en la
sección sobre conservación y recuperación, identifica otro origen preciso
del retardo: la eliminación de una componente de memoria que sigue
participando en la evolución conjunta. Su núcleo discreto conserva el
orden de las interacciones y la contribución de la memoria inicial.
Los operadores de Green aquí construidos resuelven, en cambio, una
ecuación de propagación con condiciones de soporte. Ambas construcciones
pueden componerse cuando se fija el mapa de realización que identifica
sus dominios y operadores; sus respectivas condiciones de causalidad
permanecen explícitas.

\subsection{Proyección de corrientes absorbidas y variación de la acción}
\label{vii:subsec:proyeccion-absorcion}

Sea \(\gamma_\partial:\mathcal H_\Phi\to\mathcal H_\partial\)
la lectura de frontera, distinta de la ruta \(\gamma\). Definimos
\[
 \mathcal A_\partial=\gamma_\partial G_{\rm rad}.
\]
La ventana finita fija el espacio de fuentes y la lectura de frontera;
los Green conservan su dominio de respuestas. En un nivel finito con
productos internos fijados, su inversa de
Moore--Penrose determina
\begin{equation}
 P_{\rm abs}=I-\mathcal A_\partial^+\mathcal A_\partial.
 \label{vii:eq:proyector-absorcion}
\end{equation}

\begin{theorem}[Condición de absorción y respuesta de borde]
\label{vii:thm:proyeccion-absorcion}
\(P_{\rm abs}\) es la proyección ortogonal sobre
\(\ker\mathcal A_\partial\). Para cualquier fuente \(j\), la fuente
\(j_{\rm abs}=P_{\rm abs}j\) cumple
\begin{equation}
 \gamma_\partial G_{\rm ret}j_{\rm abs}
 =\gamma_\partial G_{\rm adv}j_{\rm abs}.
 \label{vii:eq:igualdad-borde-absorbido}
\end{equation}
Es la fuente absorbida más próxima a \(j\) para el producto interno
declarado. Existen fuentes absorbidas no nulas exactamente cuando
\(\operatorname{rank}\mathcal A_\partial<\dim\mathcal H_J\).
\end{theorem}
\begin{proof}
La descomposición ortogonal
\(\mathcal H_J=\ker\mathcal A_\partial\oplus
 \operatorname{im}\mathcal A_\partial^*\) muestra que
\(\mathcal A_\partial^+\mathcal A_\partial\) es la identidad
en el segundo sumando y cero en el primero. Su complemento es el
proyector anunciado. Por definición,
\(2\mathcal A_\partial j_{\rm abs}
 =\gamma_\partial(G_{\rm ret}-G_{\rm adv})j_{\rm abs}=0\).
Para \(k\in\ker\mathcal A_\partial\), la descomposición ortogonal da
\[
 \|j-k\|^2=\|(I-P_{\rm abs})j\|^2+\|P_{\rm abs}j-k\|^2.
\]
El mínimo único corresponde a \(k=P_{\rm abs}j\).
La equivalencia final es el teorema de rango y nulidad.
\end{proof}

En el emparejamiento de fuente y campo, la acción simétrica es
\(\mathscr S[j]=\tfrac12\langle j,G_{\rm sym}j\rangle\).
La adjunción declarada hace real esa cantidad. Para variaciones reales,
\[
 \mathrm d\mathscr S[j](h)
 =\operatorname{Re}\langle h,G_{\rm sym}j\rangle.
\]
La acción restringida a fuentes absorbidas se escribe sobre el espacio
ambiente como
\begin{equation}
 \mathscr S_{\rm abs}[j]
 =\tfrac12\langle P_{\rm abs}j,G_{\rm sym}P_{\rm abs}j\rangle,
 \quad
 \mathrm d\mathscr S_{\rm abs}[j](h)
 =\operatorname{Re}\langle P_{\rm abs}h,G_{\rm sym}P_{\rm abs}j\rangle.
 \label{vii:eq:variacion-restringida-absorcion}
\end{equation}
Cuando el emparejamiento identifica campo y fuente por Riesz, el gradiente
ambiente es \(P_{\rm abs}G_{\rm sym}P_{\rm abs}j\).
La respuesta de campo sigue siendo \(G_{\rm sym}j_{\rm abs}\); el
proyector del gradiente expresa la restricción de las variaciones.
Las fórmulas se demuestran expandiendo el término cuadrático de
\(j+th\) y utilizando la autoadjunción. Un término material
\(\mathscr S_{\rm matter}\) añade su propio diferencial.

\subsection{Condición absorbente explícita en una ventana cronológica}
\label{vii:subsec:absorcion-momentos}

El canal de la subsección~\ref{vii:subsec:green-cronologico} permite
evaluar esa condición sin introducir una respuesta objetivo. Sea una
fuente apoyada en \(0,\ldots,N\), con \(N\geq1\). Transportemos
sus componentes a la fibra de referencia:
\[
 \widehat j_k=U(0,k)j_k,\qquad
 M_0(j)=\sum_{k=0}^N\widehat j_k,\qquad
 M_1(j)=\sum_{k=0}^Nk\widehat j_k.
\]
La identidad \((n-k)_+-(k-n)_+=n-k\) da
\begin{equation}
 (G_{\rm ret}-G_{\rm adv})j\big|_n
 =U(n,0)\bigl(nM_0(j)-M_1(j)\bigr).
 \label{vii:eq:diferencia-green-momentos}
\end{equation}

\begin{proposition}[Absorción y dos momentos transportados]
\label{vii:prop:absorcion-momentos}
Si la traza evalúa el campo en dos posiciones distintas \(a,b\),
la igualdad de trazas avanzada y retardada equivale a
\(M_0(j)=M_1(j)=0\). En una fibra de dimensión \(d\), el espacio
de tales fuentes tiene dimensión \((N-1)d\). Para \(N\geq2\),
contiene las fuentes no nulas con componentes transportadas
\((z,-2z,z,0,\ldots,0)\), \(z\neq0\).
\end{proposition}
\begin{proof}
La invertibilidad de \(U(a,0)\) y \(U(b,0)\) reduce la condición
de borde a \(aM_0-M_1=0\) y \(bM_0-M_1=0\).
Restar las ecuaciones da \((a-b)M_0=0\), y después \(M_1=0\).
El mapa de los dos momentos es sobreyectivo: para valores prescritos
\((z_0,z_1)\), basta tomar
\(\widehat j_0=z_0-z_1\), \(\widehat j_1=z_1\) y las demás
componentes cero. Su rango es \(2d\); la nulidad es
\((N+1)d-2d\). La fuente indicada tiene ambos momentos nulos.
\end{proof}

En el marco transportado, escribamos
\[
 B_N=\begin{pmatrix}1&1&\cdots&1\\0&1&\cdots&N\end{pmatrix}.
\]
El proyector absorbente queda evaluado por
\begin{equation}
 P_N=\bigl[I-B_N^{\mathsf T}(B_NB_N^{\mathsf T})^{-1}B_N\bigr]
       \otimes I_{\mathcal H_0}.
 \label{vii:eq:proyector-momentos-cronologicos}
\end{equation}
Las dos filas de \(B_N\) son independientes para \(N\geq1\), por lo
que la inversa existe. El transporte unitario entre marcos lleva este
proyector a las fibras originales. Para \(N=2\),
\[
 P_2=\frac16
 \begin{pmatrix}1&-2&1\\-2&4&-2\\1&-2&1\end{pmatrix}
 \otimes I_{\mathcal H_0}.
\]
Esta construcción produce una clase absorbida no nula y su operación de
proyección en el canal cronológico. Su amplitud y su identificación con una
corriente material conservan las reglas del sector físico correspondiente.

\subsection{Refinamiento y continuidad de la condición de borde}
\label{vii:subsec:refinamiento-absorcion}

Sean \(R_m^\Phi,R_m^J,R_m^\partial\) los mapas de refinamiento de
campos, fuentes y fronteras. Supongamos que entrelazan los dos Green y la
traza. Por composición se obtiene
\begin{equation}
 R_m^\Phi G_{{\rm sym},m}=G_{{\rm sym},m+1}R_m^J,\qquad
 R_m^\partial\mathcal A_{\partial,m}
 =\mathcal A_{\partial,m+1}R_m^J.
 \label{vii:eq:refinamiento-lectura-absorcion}
\end{equation}
La segunda igualdad lleva una fuente absorbida a otra fuente absorbida.
La naturalidad de los proyectores conserva adicionalmente la estructura
ortogonal de las realizaciones.

\begin{proposition}[Naturalidad de la proyección absorbente]
\label{vii:prop:naturalidad-proyeccion-absorcion}
Para una isometría de fuentes \(R=R_m^J\), se tiene
\(P_{{\rm abs},m+1}R=RP_{{\rm abs},m}\) exactamente cuando
\[
 R(\ker\mathcal A_{\partial,m})
       \subset\ker\mathcal A_{\partial,m+1},\qquad
 R((\ker\mathcal A_{\partial,m})^\perp)
       \subset(\ker\mathcal A_{\partial,m+1})^\perp.
\]
Con esas condiciones en todos los niveles, los proyectores definen una
proyección ortogonal en la completación del límite inductivo isométrico
de los espacios de fuentes.
\end{proposition}
\begin{proof}
Se descompone cada fuente en sus componentes en el núcleo y en su
complemento ortogonal. Las dos inclusiones implican que aplicar la
proyección antes o después del refinamiento conserva respectivamente la
primera componente y anula la segunda. Recíprocamente, la igualdad de
operadores aplicada a cada componente da las dos inclusiones.
En el límite inductivo, la igualdad permite definir un operador sobre la
unión de los niveles sin depender del representante. Todos los proyectores
tienen norma a lo sumo uno, de modo que el operador se extiende a la
completación. Sus identidades de autoadjunción e idempotencia, válidas en
la unión densa, pasan por continuidad a la completación.
\end{proof}

Para prolongar también una respuesta de Green a una completación se
conservan sus dominios y las estimaciones de continuidad correspondientes.
La propagación cronológica de soporte finito y la factorización de la
respuesta por la imagen de borde, demostrada en
\ref{vii:thm:observador-factorizacion}, mantienen así sus dominios
respectivos durante el refinamiento.

La realización física de emisión y absorción compone estos operadores con
las corrientes materiales, la acción y las condiciones causales del sector.
El entrelazamiento debe conservar la involución de ruta, el emparejamiento
de fuente y campo, la acción transportada y las condiciones de frontera a
cada nivel. La construcción precedente aporta la reversión del grupo, los
Green del canal cronológico, la descomposición temporal y una proyección
absorbente explícita; las realizaciones espaciales emplean sus mapas y
condiciones de causalidad propios.

\section{Realización de la acción y comparación metrológica}
\subsection{Contraste de la constante de Planck reducida}
\label{sec:revision-planck-contraste}

El Sistema Internacional vigente desde el 20 de mayo de 2019 fija
exactamente \(h=6.62607015\times10^{-34}\,\mathrm{J\,s}\)
\cite{bipm2018}. Por tanto, su constante reducida es
\[
 \hbar_{\mathrm{SI}}=\frac{6.62607015}{2\pi}\,10^{-34}\,
 \mathrm{J\,s}
 =1.0545718176461563912\ldots\times10^{-34}\,\mathrm{J\,s}.
\]
Los puntos suspensivos corresponden a la expansión de una expresión
exacta, no a una incertidumbre experimental de \(h\) en el SI actual.
Este valor de referencia se utiliza aquí después de la evaluación de
los lectores de HMT; su función es comparativa.

Las dos coordenadas del retorno del artículo son
\begin{align*}
 H_5&=1.0545718176461544696\ldots,\\
 \eta_{\mathrm{ret}}&=1.0545718169150390611\ldots.
\end{align*}
En la identificación de \(\mathcal U_S\) con \(\mathrm{J\,s}\), la
diferencia relativa de la sección anterior respecto de
\(\hbar_{\mathrm{SI}}\) es aproximadamente \(-1.82\times10^{-15}\).
La de la sección posterior es aproximadamente \(-6.93\times10^{-10}\).
La segunda diferencia contiene el efecto del retorno que utiliza el
corrector electrónico; no equivale al error de evaluación de la
primera sección.

La precisión de estas comparaciones permite reconocer la proximidad
cuantitativa de la construcción de \(H_5\) y, simultáneamente, distinguir
proximidad de identidad exacta. En el presente manuscrito las expresiones
exhibidas no proporcionan una igualdad exacta con \(h/(2\pi)\).
El resultado demostrado sobre \(\Sigma_H\) es su década \(-34\);
el resultado analítico sobre \(H_5\) es la evaluación del carácter
declarado. La identificación física se examina en esta carta y con las
diferencias explícitas que acabamos de dar. El valor SI no interviene
en la definición de la norma determinantal ni en la recurrencia regional.


\section{Covarianza, forma normal y espectro gaussiano}
\label{app:gaussianas-espectro}

La sección positiva de acción \(\hbar_a\) y la realización canónica
anterior determinan la normalización empleada aquí. Se demuestra la
correspondencia entre covarianza, espectro, pureza y entropía para un
número finito de modos; el espacio de Fock conserva todas sus ocupaciones.
La reducción simpléctica y las fórmulas gaussianas son resultados
establecidos que se reproducen para hacer autónoma su aplicación al
transporte de memoria.

\subsection{Coordenadas canónicas y función característica}
En \(L^2(\RR^d)\), sean \(\widehat q_j=\widehat Q_j/\sqrt{\hbar_a}\),
\(\widehat p_j=\widehat P_j/\sqrt{\hbar_a}\) y
\(\widehat z=(\widehat q_1,\widehat p_1,\ldots,\widehat q_d,\widehat p_d)\).
Sobre Schwartz,
\[
 [\widehat z_j,\widehat z_k]=iJ_{jk}I,\qquad
 J=\bigoplus_{j=1}^d\begin{pmatrix}0&1\\-1&0\end{pmatrix}.
\]
Esta convención fija el signo de la forma canónica. Los operadores de
Weyl \(\mathcal W(\xi)=\exp(i\xi^{\mathsf T}\widehat z)\) se definen
por las clausuras autoadjuntas de las combinaciones lineales indicadas.
Para un estado centrado \(\rho\) de segundos momentos finitos, la
covarianza normalizada y su función característica son
\[
 W_{jk}=\Tr\rho(\widehat z_j\widehat z_k+\widehat z_k\widehat z_j),
 \qquad \chi_\rho(\xi)=\Tr(\rho\mathcal W(\xi)),\qquad
 V=\frac{\hbar_a}{2}W.
\]
El estado es gaussiano centrado cuando
\begin{equation}
 \chi_\rho(\xi)=\exp(-\xi^{\mathsf T}W\xi/4).
 \label{app:eq:caracteristica-gaussiana}
\end{equation}
La esperanza de \((c^{\mathsf T}\widehat z)^\dagger
(c^{\mathsf T}\widehat z)\) es \(c^\dagger(W+iJ)c/2\); así
\(W+iJ\geq0\). Además \(W>0\): si \(v\) es real y
\(v^{\mathsf T}Wv=0\), la positividad implica \((W+iJ)v=0\),
por tanto \(Jv=0\) y \(v=0\).

\begin{lemma}[Unicidad de la función característica]
\label{app:lem:unicidad-weyl}
Dos operadores de clase traza con la misma función característica son iguales.
\end{lemma}
\begin{proof}
Reagrupando posiciones y momentos como \(\xi=(u,v)\),
\[
 (\mathcal W(u,v)f)(x)=e^{iu\cdot(x+v/2)}f(x+v).
\]
Para un operador de rango finito con vectores de Schwartz y núcleo \(K_A\),
\[
 \Tr(A\mathcal W(u,v))=\int K_A(x+v/2,x-v/2)e^{iu\cdot x}\,dx.
\]
Plancherel y el cambio de variables de jacobiano absoluto uno dan
\[
 \|A\|_{\rm HS}^2=(2\pi)^{-d}\int_{\RR^{2d}}
 |\Tr(A\mathcal W(\xi))|^2\,d\xi.
\]
Los operadores anteriores son densos en clase traza. La convergencia
en esa norma implica convergencia en Hilbert--Schmidt y convergencia
uniforme de las funciones características, porque su diferencia está
acotada por \(\|A-B\|_1\). La igualdad se extiende en \(L^2\).
Aplicarla a la diferencia de los operadores prueba la afirmación.
\end{proof}

\subsection{Reducción simpléctica de la forma positiva}
\begin{theorem}[Forma normal de Williamson]
\label{app:thm:williamson}
Para una matriz real simétrica \(W>0\) de orden \(2d\), existen
\(S\in\operatorname{Sp}(2d,\RR)\) y \(\nu_j>0\) tales que
\[
 W=SDS^{\mathsf T},\qquad D=\bigoplus_j\nu_jI_2.
\]
Los autovalores de \(JW\) son \(\pm i\nu_j\). La condición
\(W+iJ\geq0\) equivale a \(\nu_j\geq1\) para todo \(j\).
\end{theorem}
\begin{proof}
La matriz \(B=W^{1/2}JW^{1/2}\) es antisimétrica e invertible.
Si \(-B^2u=\nu^2u\) y \(\|u\|=1\), el vector
\(v=-Bu/\nu\) es unitario y ortogonal a \(u\); cumple
\(Bu=-\nu v\), \(Bv=\nu u\). El complemento de esta pareja es
invariante. La iteración produce una matriz ortogonal \(O\) con
\(O^{\mathsf T}BO=JD\). Definiendo \(S=W^{1/2}OD^{-1/2}\),
\[
 S^{\mathsf T}JS=J,\qquad SDS^{\mathsf T}=W.
\]
La semejanza \(W^{1/2}(JW)W^{-1/2}=B\) determina los valores
\(\nu_j\), con multiplicidad. La congruencia por \(S^{-1}\)
transforma \(W+iJ\) en \(D+iJ\); cada bloque tiene autovalores
\(\nu_j+1\) y \(\nu_j-1\).
\end{proof}

\begin{lemma}[Implementación unitaria en un número finito de modos]
\label{app:lem:implementacion-simplectica}
Toda matriz simpléctica \(S\) admite un unitario \(U_S\) que satisface
\(U_S^\dagger\mathcal W(\xi)U_S=\mathcal W(S^{\mathsf T}\xi)\).
La conjugación \(\rho\mapsto U_S\rho U_S^\dagger\) transporta
la covarianza por \(W\mapsto SWS^{\mathsf T}\).
\end{lemma}
\begin{proof}
En la descomposición polar \(S=OP\), la identidad
\(J^{-1}(S^{\mathsf T}S)J=(S^{\mathsf T}S)^{-1}\), aplicada por
cálculo espectral a la raíz positiva, prueba \(PJP=J\).
Entonces \(O\) es ortogonal simpléctica. Los autoespacios de \(P\)
se emparejan mediante \(J\) con autovalores \(\lambda,\lambda^{-1}\);
el de autovalor uno es \(J\)-invariante. Por ello una matriz
ortogonal simpléctica \(K\) reduce \(P\) a
\(\bigoplus_j\operatorname{diag}(e^{r_j},e^{-r_j})\).

La dilatación se implementa por
\[
 (\mathcal D_rf)(x)=e^{-\frac12\sum_jr_j}
 f(e^{-r_1}x_1,\ldots,e^{-r_d}x_d).
\]
Un cambio de variables prueba su unitariedad y la conjugación de
\((\widehat q_j,\widehat p_j)\) a
\((e^{r_j}\widehat q_j,e^{-r_j}\widehat p_j)\).
Una transformación ortogonal simpléctica conmuta con \(J\) y
corresponde a una matriz \(u\in U(d)\) sobre
\(a_j=(\widehat q_j+i\widehat p_j)/\sqrt2\). Su segunda
cuantización \(\Gamma(u)=\bigoplus_{n\geq0}u^{\otimes_s n}\)
es unitaria en \(\mathcal F_s(\CC^d)\), y la identidad sobre
productos simétricos finitos da la transformación de los \(a_j\).

La base de Hermite identifica este espacio con \(L^2(\RR^d)\).
Para justificar su completitud, si \(f\) es ortogonal a todos los
polinomios multiplicados por \(e^{-|x|^2/2}\), la función entera
\(F(z)=\int f(x)e^{-|x|^2/2}e^{z\cdot x}\,dx\) tiene nulas
todas sus derivadas en cero. Así \(F=0\); restringiendo a
\(z=i\xi\), la unicidad de Fourier da \(f=0\). Las
normalizaciones y la ortogonalidad resultan de creación y aniquilación.
La composición de las implementaciones anteriores produce \(U_S\).
La igualdad de Weyl implica
\(\chi_{U_S\rho U_S^\dagger}(\xi)=\chi_\rho(S^{\mathsf T}\xi)\),
que demuestra la ley de covarianza.
\end{proof}

\subsection{Cálculo del espectro de la densidad isotrópica}
En un modo, los estados coherentes normalizados son
\[
 |z\rangle=e^{-|z|^2/2}\sum_{n\geq0}\frac{z^n}{\sqrt{n!}}|n\rangle.
\]
Su función de posición,
\(\pi^{-1/4}\exp[-(x-q_0)^2/2+ip_0(x-q_0/2)]\), donde
\(q_0=\sqrt2\Re z\), \(p_0=\sqrt2\Im z\), da por integración
\[
 \langle z|\mathcal W(u,v)|z\rangle
 =e^{-(u^2+v^2)/4}e^{i(uq_0+vp_0)}.
\]
Para \(s>0\), la integral
\[
 \rho^{(s)}=\frac1{\pi s}\int_{\CC} e^{-|z|^2/s}
 |z\rangle\langle z|\,d^2z
\]
converge en clase traza, es positiva y tiene traza uno. La integración
angular anula sus elementos fuera de la diagonal de ocupación; la
integración radial da
\[
 \langle n|\rho^{(s)}|n\rangle=\frac{s^n}{(1+s)^{n+1}},\qquad
 \chi_{\rho^{(s)}}(u,v)=e^{-(2s+1)(u^2+v^2)/4}.
\]
Con \(\nu=2s+1\), el único estado gaussiano de covarianza
\(\nu I_2\) es, por el lema~\ref{app:lem:unicidad-weyl},
\begin{equation}
 \rho_\nu=(1-t)\sum_{n\geq0}t^n|n\rangle\langle n|,
 \qquad t=\frac{\nu-1}{\nu+1}.
 \label{app:eq:estado-geometrico}
\end{equation}
El caso \(\nu=1\) es el proyector sobre el vacío.

\begin{theorem}[Espectro, pureza y entropía]
\label{app:thm:espectro-gaussiano}
Un estado gaussiano centrado con covarianza \(W\) y valores
simplécticos \(\nu_1,\ldots,\nu_d\) es unitariamente equivalente
a \(\bigotimes_j\rho_{\nu_j}\). Sus autovalores son
\[
 \lambda_{\symbf{n}}=\prod_j(1-t_j)t_j^{n_j},\qquad
 t_j=\frac{\nu_j-1}{\nu_j+1},\quad \symbf{n}\in\mathbb N^d.
\]
Con \(0\log0=0\),
\begin{align}
 \Tr\rho^2&=\prod_j\nu_j^{-1}=(\det W)^{-1/2},
 \label{app:eq:pureza-general}\\
 -\Tr(\rho\log\rho)&=\sum_j
 \left[\frac{\nu_j+1}{2}\log\frac{\nu_j+1}{2}
 -\frac{\nu_j-1}{2}\log\frac{\nu_j-1}{2}\right].
 \label{app:eq:entropia-general}
\end{align}
\end{theorem}
\begin{proof}
Williamson da \(W=SDS^{\mathsf T}\). El producto de las densidades
isotrópicas tiene función característica \(e^{-\xi^{\mathsf T}D\xi/4}\).
La implementación unitaria y la unicidad de Weyl prueban la
equivalencia y la fórmula espectral. En un modo,
\[
 \sum_{n\geq0}(1-t)^2t^{2n}=\frac{1-t}{1+t}=\nu^{-1},
\]
y
\[
 -\sum_{n\geq0}(1-t)t^n\log((1-t)t^n)
 =-\log(1-t)-\frac{t}{1-t}\log t.
\]
La sustitución \(t=(\nu-1)/(\nu+1)\) prueba las fórmulas; el
límite en \(t=0\) es cero. En el producto finito, Tonelli justifica
la suma de los términos de entropía no negativos, y
\(\det W=\prod_j\nu_j^2\) da la expresión determinantal.
\end{proof}

Si el estado es la reducción de un vector puro bipartito, sus
autovalores son los cuadrados de los coeficientes de Schmidt.
La descomposición singular del operador de Hilbert--Schmidt definido
por ese vector demuestra la igualdad de espectros no nulos de ambas
reducciones. La entropía calculada es entonces la de entrelazamiento.

\subsection{Preparación y transporte de las dos covarianzas}
La aplicación posterior utiliza dos entradas gaussianas puras idénticas.
En general, para \(M\in\operatorname{Sp}(2d,\RR)\), su covarianza es
\(V_0=(\hbar_a/2)MM^{\mathsf T}\), y la entrada conjunta tiene
covarianza \(V_0\oplus V_0\). El transporte físico se obtiene
conjugando \(U_H\) por \(M\oplus M\). En la carta equilibrada,
la covarianza normalizada de entrada es \(I_{4d}\); con
\[
 T=\frac{8I+H}{9},\qquad D=\frac{\sqrt8(H-I)}9,
\]
la primera reducción tiene
\[
 W=TT^{\mathsf T}+DD^{\mathsf T}=\frac{8I+HH^{\mathsf T}}9.
\]
Su función característica se obtiene anulando los argumentos del
segundo grupo de modos. Por ello permanece gaussiana y los teoremas
anteriores calculan todo su espectro, sin truncar el registro de ocupación.

\section{Acciones físicas y espectro de la memoria reducida}
\label{nuclear:12}
\subsection{Genealogía de las secciones dimensionales}

Las coordenadas \(\pi,\varphi,e,\alpha\) son las realizaciones de las salidas correlacionadas construidas en las secciones~\ref{sec:alpha} y~\ref{sec:electron}. En esta sección se utiliza \(\alpha=\alpha_A\), con la orientación y la periodicización de la definición~\ref{def:alpha-operativa}:



\begin{equation}
\alpha=\pi+e-\varphi-4-\kappa_{\rm per},\qquad
\kappa_{\rm per}=\frac{N_K}{1000^{12}-1},
\quad N_K=234543140729659824621058914794146601.
\label{nuc12:eq:2.1}
\end{equation}



La identidad~\eqref{eq:alpha-identidad-completa} se deduce de la recurrencia con acarreo compatible y de su telescopía, demostradas en la proposición~\ref{prop:alpha-telescopia}. El período doce corresponde al registro: una ventana de doce bloques con condición de frontera nula y una sección infinita compatible son objetos diferentes. La igualdad \eqref{nuc12:eq:2.1} utiliza esta última. El registro ordenado \(K\) conserva sus incidencias además de admitir la lectura escalar \(\kappa_{\rm per}\).

El lector de acción, construido en la sección~\ref{iv:action:section}, determina



\begin{equation}
\begin{aligned}
H_5&=\tfrac12\sqrt{1000\alpha/\varphi}-\alpha+
\tfrac9{16}\alpha^2-\tfrac59\alpha^3+
\tfrac7{48}\alpha^4-\tfrac1{54}\alpha^5,\\
D_A&=\exp(-100\pi\alpha/9),\\
\mathcal C_\pi&=169\alpha^6+\frac{D_A\alpha^7}{1-D_A\alpha},\\
\eta_{\rm ret}&=H_5-\frac{90}{\pi}\mathcal C_\pi,\\
\hbar_{\rm pre}&=H_5\mathcal A_0,\qquad
\hbar_{\rm ret}=\eta_{\rm ret}\mathcal A_0,\qquad
R_{\rm act}=H_5/\eta_{\rm ret},
\end{aligned}
\label{nuc12:eq:2.2}
\end{equation}



Aquí \(\mathcal A_0=10^{-34}\mathcal U_S\). La década procede del lector determinantal de \(\varphi\alpha^{16}\), cuya construcción y acotación figuran en el teorema~\ref{iv:action:decada}; \(H_5\) y \(\eta_{\rm ret}\) son las mantisas de las dos secciones positivas de acción reducida. En cada sección \(a\in\{\mathrm{pre},\mathrm{ret}\}\) se cumple \(h_a=2\pi\hbar_a\).

La construcción del electrón central, desarrollada en la sección~\ref{sec:electron}, determina la coordenada



\begin{equation}
\begin{aligned}
\Delta_4&=\frac{\pi-e\log\pi}{270}\left(1+\frac{\pi}{729}\right),\\
B_e&=\frac{\sqrt3}{4}
\left[\exp(\varphi/\pi^2)-22\alpha^3\right](1+15\Delta_4),\\
\kappa_e&=80-54\alpha+6\Delta_4,\qquad
\varepsilon_e^\beta=B_eR_{\rm act}^{\kappa_e}.
\end{aligned}
\label{nuc12:eq:2.3}
\end{equation}



El prefactor \(\sqrt3/4\) procede de la norma del conmutador de los proyectores de suma y producto sobre las fibras especificadas en \eqref{eq:el-prefactor}. La ruta central y sus registros determinan los coeficientes de retorno de \eqref{eq:el-kappa} y~\eqref{eq:el-beta}. La realización dimensional es \(E_e^\beta=\varepsilon_e^\beta\mathcal U_E\), \(m_e^\beta=E_e^\beta/c^2\), con la carta energética de \eqref{eq:el-dimensional}.

El reloj dodecafásico determina la base energética de \eqref{eq:el-cuanto-reloj}:



\[
\omega_{108}=\frac{2\pi}{108t_0},\quad
E_{\rm clk}=\hbar_{\rm ret}\omega_{108},\quad
b_e=B_e\frac{\mathcal U_E}{E_{\rm clk}}.
\]



La sustitución de \eqref{nuc12:eq:2.3}, conservando la misma energía física, demuestra



\begin{equation}
\boxed{\frac{m_e^\beta c^2}{\hbar_{\rm ret}}
=b_eR_{\rm act}^{\kappa_e}\omega_{108},\qquad
\frac{m_e^\beta c^2}{\hbar_{\rm pre}}
=b_eR_{\rm act}^{\kappa_e-1}\omega_{108}.}
\label{nuc12:eq:2.4}
\end{equation}



La división por \(h_a\), en lugar de \(\hbar_a\), convierte la frecuencia angular en frecuencia cíclica. La conversión \(\mathcal U_E/E_{\rm clk}\) permanece explícita: una coordenada expresada en \(\mathrm{MeV}\) cambia al sustituir su base energética por el cuanto del reloj.

\subsection{Normalización de la acción Maxwell--Dirac}

La respuesta constitutiva de la sección~\ref{iii:sec:constitutiva} produce la velocidad \(c\), la impedancia \(Z\), la permitividad \(\varepsilon\) y la permeabilidad \(\mu\), todas positivas, con \(\varepsilon c=Z^{-1}\) y \(\mu c=Z\). La sección positiva de carga de \eqref{iii:eq:charge} satisface



\begin{equation}
e_a^2=\frac{2\alpha h_a}{Z}
=4\pi\alpha\varepsilon\hbar_a c .
\label{nuc12:eq:3.1}
\end{equation}



Aquí \(e_a>0\) designa la carga elemental de la sección \(a\); \(e\) en \eqref{nuc12:eq:2.1}--\eqref{nuc12:eq:2.3} designa la coordenada de Euler. Fijamos una carta homogénea con \(c,Z,\varepsilon,\mu,\hbar_a,e_a\) constantes positivos, la masa \(m=m_e^\beta>0\) y la carga \(q\in\{-e_a,e_a\}\); entonces \(\sigma=q/e_a\in\{-1,1\}\). En la realización de Minkowski de signatura \((+,-,-,-)\), con \(x^0=ct\), la acción Maxwell--Dirac es



\begin{equation}
\mathcal I=\int d^4x\left[
-\frac{F_{\mu\nu}F^{\mu\nu}}{4\mu c}
+\bar\psi\left(i\hbar_a\gamma^\mu
(\partial_\mu+iqA_\mu/\hbar_a)-mc\right)\psi
\right].
\label{nuc12:eq:3.2}
\end{equation}



La componente \(A_0\) es el potencial escalar dividido por \(c\), y \(F=dA\). La representación espinorial se fija mediante las matrices \(g^I\) de \eqref{vii:eq:clifford-material-explicito}, que satisfacen \(\{g^I,g^J\}=2\eta^{IJ}I_4\) para \(\eta=\operatorname{diag}(-1,1,1,1)\). Definimos \(\Gamma^I=-ig^I\), de modo que \(\{\Gamma^I,\Gamma^J\}=-2\eta^{IJ}I_4\), y, en esta carta plana, \(\gamma^\mu=\delta^\mu_I\Gamma^I\), \(\bar\psi=\psi^\dagger\Gamma^0\). Así se fijan conjuntamente la firma, el adjunto y el factor imaginario; las matrices \(\gamma^\mu\) se distinguen del parámetro escalar de Holst utilizado después. Se consideran campos suaves y variaciones de soporte compacto, o condiciones de frontera que anulen los términos de borde. La acción realiza los coeficientes ya generados y la representación material de la sección~\ref{vii:subsec:accion-espinorial-afin}.

Definimos



\[
\ell_a=\frac{\hbar_a}{mc},\quad X=x/\ell_a,\quad
\Psi=\ell_a^{3/2}\psi,\quad
\mathsf a_\mu=\frac{e_a\ell_a}{\hbar_a}A_\mu,\quad
\mathsf f=d\mathsf a.
\]



Entonces



\begin{equation}
\boxed{\frac{\mathcal I}{\hbar_a}=
\int d^4X\left[
-\frac{\mathsf f_{\mu\nu}\mathsf f^{\mu\nu}}{16\pi\alpha}
+\bar\Psi\bigl(i\gamma^\mu(\partial_\mu+i\sigma\mathsf a_\mu)-1\bigr)\Psi
\right].}
\label{nuc12:eq:3.3}
\end{equation}



\begin{proof}
El cambio de volumen aporta \(\ell_a^4\), cada derivada aporta \(\ell_a^{-1}\) y el producto de los campos fermiónicos aporta \(\ell_a^{-3}\). El coeficiente del término de masa es \(mc\ell_a/\hbar_a=1\). El término electromagnético aporta \(\hbar_a/(4\mu c e_a^2)\). Como \(\mu c=Z\) y \(Ze_a^2=4\pi\alpha\hbar_a\), ese coeficiente es \(1/(16\pi\alpha)\). Quedan determinados ambos términos y sus normalizaciones.
\end{proof}

La variación respecto de \(\mathsf a_\nu\), seguida de la sustitución de \eqref{nuc12:eq:2.1}, da



\begin{equation}
\partial_\mu\mathsf f^{\mu\nu}
=4\pi\left(\pi+e-\varphi-4-\kappa_{\rm per}\right)
\sigma\bar\Psi\gamma^\nu\Psi .
\label{nuc12:eq:3.4}
\end{equation}



La cancelación compone las secciones de acción, la construcción electrónica y la respuesta constitutiva; el lector regional y dodecafásico determina el acoplamiento restante. La masa persiste en la longitud \(\ell_a\), y la acción en la conversión dimensional. La normalización de Compton es una realización posterior de esas salidas. En la realización coulombiana del mismo electrón, la longitud de Bohr es \(\ell_a/\alpha\): sustituir \(e_a^2=4\pi\alpha\varepsilon\hbar_a c\) en \(4\pi\varepsilon\hbar_a^2/(m e_a^2)\) da exactamente esa expresión.

La deducción utiliza coeficientes homogéneos fijos. Para coeficientes dependientes del espacio-tiempo, la regla del producto introduce las derivadas de las escalas y del acoplamiento en las ecuaciones transformadas. La variación de una acción con parámetros variables debe conservar esos términos.

\subsection{Escala areal de la acción gravitatoria y de la entropía}

La realización circular del ciclo de \(108\) pasos, construida en la sección~\ref{v:ant:sec:gravity}, determina



\begin{equation}
t_P=\frac{54}{\pi}t_0,\qquad
G_a=\frac{c^5t_P^2}{\hbar_a},\qquad
\ell_P^2=\frac{\hbar_aG_a}{c^3}=c^2t_P^2.
\label{nuc12:eq:4.1}
\end{equation}



La condición de coincidencia de radios de esa realización es \(G_a m_{108,a}/c^2=R_{108}\), con \(m_{108,a}=\hbar_a\omega_{108}/c^2\). Su solución única fija el factor \((54/\pi)^2\); el radio empleado es \(G_a m/c^2\). La reciprocidad \(G_{\rm pre}/G_{\rm ret}=R_{\rm act}^{-1}\) conserva \(\ell_P^2\) entre ambas secciones.

La sección~\ref{vii:sec:realizacion-cartan} construye la cotetrada y la conexión a partir del transporte ya realizado. En la acción de la sección~\ref{vii:sec:corriente-respuesta-cartan}, fijamos una cotetrada no degenerada \(e^I\), la orientación, la firma \(\eta=\operatorname{diag}(-1,1,1,1)\) y una conexión métrica \(\omega\), con curvatura \(F^{IJ}=d\omega^{IJ}+\omega^I{}_{K}\wedge\omega^{KJ}\). Tomamos \(\gamma\in\mathbb R\setminus\{0\}\) y mantenemos constantes \(\gamma,\Lambda,c,G_a,\hbar_a\) en la variación. Ponemos \(\Sigma^{IJ}=e^I\wedge e^J\) y \(\mathsf P_\gamma=\star+\gamma^{-1}I\); la dualidad interna lorentziana satisface \(\star^2=-I\) sobre bivectores. Las variaciones tienen soporte compacto o condiciones de frontera que anulan el término superficial. La sustitución de \eqref{nuc12:eq:4.1} en \eqref{vii:eq:accion-cartan-holst} da



\begin{equation}
\boxed{\frac{\mathcal I_{\rm grav}}{\hbar_a}
=\frac1{16\pi\ell_P^2}
\left[\int\Sigma_{IJ}\wedge(\mathsf P_\gamma F)^{IJ}
-2\Lambda\int\operatorname{vol}_e\right].}
\label{nuc12:eq:4.2}
\end{equation}



Con la cotetrada y los campos materiales fijos durante la variación de la conexión, sea \(\sigma_{IJ}\) la corriente material definida por
\(\delta_\omega\mathcal I_m=-\tfrac12\int\delta\omega^{IJ}\wedge\sigma_{IJ}\), según \eqref{vii:eq:corriente-variacional}. La misma sustitución en el teorema~\ref{vii:thm:ecuacion-cartan-holst} produce



\begin{equation}
D_\omega(\mathsf P_\gamma\Sigma)
=8\pi\ell_P^2\,\frac{\sigma}{\hbar_a}.
\label{nuc12:eq:4.3}
\end{equation}



\begin{proof}
En la sección \(a\), el coeficiente original es \(1/(2\kappa_a c)\), con \(\kappa_a=8\pi G_a/c^4\). Dividir por \(\hbar_a\) lo convierte en \(1/(16\pi\ell_P^2)\); el término cosmológico conserva su factor dos. En la ecuación de conexión se cumple \(\kappa_a c\hbar_a=8\pi\ell_P^2\). Ambas identidades utilizan la misma sección dimensional.
\end{proof}

El funcional de Barbero de la sección~\ref{v:ant:sec:barbero} conserva su expresión íntegra. Sus canales son \(q_\pm=e^{-x\mp y}\), con \(x=\pi A/180\), \(y=\pi C^*/180\), donde \(A\) y \(C^*\) son las coordenadas angulares expresadas en grados. El dominio \(x>|y|\) asegura \(0<q_\pm<1\):



\begin{equation}
\gamma=\frac{180}{\pi}xy+
12[S_{90}(q_-)-S_{90}(q_+)]
-[S_{120}(q_-)-S_{120}(q_+)],\qquad
S_m(q)=\frac{q^m}{1-q^{3m}}.
\label{nuc12:eq:4.4}
\end{equation}



La incidencia hexádica y octádica fija las multiplicidades \(90\) y \(120\), y la cadena orientada fija los coeficientes \(12\) y \(-1\), como demuestra \eqref{v:ant:eq:gamma}. La construcción areal de la sección~\ref{v:ant:sec:area} da además



\begin{equation}
a_0=\frac{1+2\cos(\pi/18)}3\,\frac y6,\qquad
\widehat{\mathcal A}=\gamma a_0\ell_P^2(N_{90}+N_{120}).
\label{nuc12:eq:4.5}
\end{equation}



La inversión \(y\mapsto-y\) cambia los signos de \(\gamma\) y \(a_0\); su producto es par. El operador \(\mathsf P_\gamma\) conserva, en cambio, su dependencia de la orientación. La invariancia del producto areal se distingue así de la invariancia de cada coeficiente del operador. El espectro areal se realiza en la sección positiva \(\gamma a_0>0\), con los operadores de ocupación \(N_{90},N_{120}\geq0\) construidos en la sección~\ref{v:ant:sec:area}.

Para el horizonte cosmológico esférico de \eqref{vii:eq:horizonte-cosmologico}, de radio \(R>0\) y área \(\mathcal A_H=4\pi R^2\), fijamos una sección \(a\) y abreviamos \(\hbar=\hbar_a\), \(G=G_a\) durante este cálculo y el del rebote que sigue. Entonces



\begin{equation}
\frac{S_H}{k_B}=\frac{\mathcal A_H}{4\ell_P^2},\qquad
T_{\rm cos}=\frac{\hbar c}{2\pi k_BR},\qquad
E_\Lambda=\frac{c^4R}{2G},\qquad T_{\rm cos}S_H=E_\Lambda .
\label{nuc12:eq:4.6}
\end{equation}



La escala areal de \eqref{nuc12:eq:4.2}, \eqref{nuc12:eq:4.3} y \eqref{nuc12:eq:4.6} es la misma. La temperatura corresponde al horizonte cosmológico esférico; la normalización del horizonte de Schwarzschild constituye otra realización. Al variar \(R\) manteniendo \(c,G,\hbar,k_B\) fijos, la diferenciación de \eqref{nuc12:eq:4.6} conserva ambos términos: \(T_{\rm cos}\,dS_H=2\,dE_\Lambda\) y \(S_H\,dT_{\rm cos}=-\,dE_\Lambda\), como en \eqref{vii:eq:diferencial-horizonte}. Se obtiene el diferencial de una identidad entre estados de esa familia, no una sustitución del balance dinámico con otros parámetros variables.

\subsubsection{Una aplicación dinámica explícita}

El teorema~\ref{vii:thm:amplitud-material-explicita} determina, para un estado espinorial con segundo momento conservado \(q_\varrho>0\), la siguiente amplitud. La conservación se expresa mediante \eqref{vii:eq:conservacion-momento-cuartico}; el corolario~\ref{vii:cor:dominio-positivo-conservado} proporciona una realización concreta en el sector de tres ocupaciones:



\begin{equation}
s_0=\frac{3\pi G\hbar^2}{2c^2V_c^2}
\frac{\gamma^2}{1+\gamma^2}q_\varrho .
\label{nuc12:eq:4.7}
\end{equation}



En la realización espacialmente plana con materia de presión nula y \(\Lambda=0\), la constante \(\rho_0>0\) es una densidad de energía, \(V_c>0\) es el volumen comóvil de referencia y \(E_0=\rho_0V_c>0\). La solución del teorema~\ref{vii:thm:rebote-polvo} es



\[
a_b^3=s_0/\rho_0,\qquad
\tau^2=\frac{s_0c^2}{6\pi G\rho_0^2},\qquad
a(t)=a_b\left(1+\frac{(t-t_b)^2}{\tau^2}\right)^{1/3}.
\]



La restitución de \(c\) en \eqref{vii:eq:dinamica-homogenea} utiliza el coeficiente \(8\pi G/(3c^2)\) cuando la densidad es energética. En particular, \((\dot a/a)^2=(8\pi G/(3c^2))(\rho_0a^{-3}-s_0a^{-6})\). Sustituir \eqref{nuc12:eq:4.7} en los parámetros de la solución demuestra



\begin{equation}
\boxed{\left(\frac{\tau E_0}{\hbar}\right)^2
=\frac{q_\varrho}{4}\frac{\gamma^2}{1+\gamma^2},\qquad
\mathcal V_b=\frac{3\pi}{2}q_\varrho
\frac{\gamma^2}{1+\gamma^2}\ell_P^2\frac{\hbar c}{E_0}.}
\label{nuc12:eq:4.8}
\end{equation}



Aquí \(\mathcal V_b=V_c a_b^3\) es el volumen físico en el rebote. Se cancelan \(G\) y \(V_c\) en la primera igualdad; \(E_0\) permanece como energía material. En el sector conservado \(\Lambda^3\mathbb C^4\), la proposición~\ref{vii:prop:operador-espin-car} demuestra \(\widehat{\mathscr Q}_{\rm sp}=4I\), por lo que \(q_\varrho=4\) para todo estado normalizado soportado en ese sector, y
\(\tau E_0/\hbar=|\gamma|/\sqrt{1+\gamma^2}\).
Se obtiene así una escala dinámica para el estado espinorial y la realización cosmológica especificados.

\subsection{Prolongación del transporte de memoria: realización cuántica}

La realización cuántica utiliza el acoplamiento reversible de la sección~\ref{nuclear:10}. Consideramos primero un plano canónico real con matriz alternante \(\Omega=\left(\begin{smallmatrix}0&1\\-1&0\end{smallmatrix}\right)\). La partición nonádica determina



\begin{equation}
Q=\frac13\begin{pmatrix}\sqrt8 I&-I\\I&\sqrt8 I\end{pmatrix},\qquad
U_H=Q^{\mathsf T}\operatorname{diag}(I,H)Q
=\frac19\begin{pmatrix}8I+H&\sqrt8(H-I)\\
\sqrt8(H-I)&I+8H\end{pmatrix}.
\label{nuc12:eq:5.1}
\end{equation}



El operador \(H\in\operatorname{Sp}(2,\mathbb R)\) transporta ese plano. La matriz \(Q\) conserva el producto euclídeo y la forma \(\Omega\oplus\Omega\); por ello \(U_H\) es simpléctica. La primera componente define el sistema observado y la segunda el registro de memoria; la transformación completa actúa sobre ambos.

El lazo ordenado construido en la sección~\ref{nuclear:11} es



\begin{equation}
C=\begin{pmatrix}0&-1\\1&-1\end{pmatrix},\quad
P=\begin{pmatrix}1&1\\0&1\end{pmatrix},\quad
H_n=C^{-1}P^{-n}CP^n
=\begin{pmatrix}1+n&n^2\\n&n^2-n+1\end{pmatrix}.
\label{nuc12:eq:5.2}
\end{equation}



Para cada \(n\in\mathbb Z\) se cumple \(\det H_n=1\). En particular, para \(n=1\),


\begin{equation}
H_1=\begin{pmatrix}2&1\\1&1\end{pmatrix}
=F_1^2,\quad F_1=\begin{pmatrix}1&1\\1&0\end{pmatrix}
=P F_{\rm av}P^{-1}.
\label{nuc12:eq:5.3}
\end{equation}



Su espectro es \(\{\varphi^2,\varphi^{-2}\}\), por el reconocimiento posterior del modo áureo ya generado. El entero \(n\) especifica el número orientado de iteraciones en el lazo elegido. La familia de transportes está determinada algebraicamente; su realización como protocolo físico conserva además los datos de preparación y de evolución correspondientes.

\subsubsection{Estado inicial, transporte y lectura}

La proposición~\ref{prop:ii-covarianza-gaussiana} construye la realización gaussiana de la elipse de acción. En la sección positiva \(\hbar_a>0\), con \(A>0\) y \(|C^*/A|<1\), su covarianza es



\begin{equation}
V_\eta=\frac{\hbar_a}{2}M_\eta M_\eta^{\mathsf T},\qquad
M_\eta=\operatorname{diag}(e^{\eta/2},e^{-\eta/2}),\qquad
\eta=\operatorname{artanh}(C^*/A).
\label{nuc12:eq:5.4}
\end{equation}



En la representación de Schrödinger sobre \(L^2(\mathbb R,dQ)\), con \(\widehat Q\) multiplicación por \(Q\) y \(\widehat P=-i\hbar_a\,d/dQ\), la relación \([\widehat Q,\widehat P]=i\hbar_a I\) vale en el espacio de Schwartz. El estado puro normalizado es



\[
\psi_\eta(Q)=(\pi\hbar_a e^\eta)^{-1/4}
\exp[-Q^2/(2\hbar_a e^\eta)].
\]



La preparación es el producto \(\psi_\eta\otimes\psi_\eta\). El transporte de \eqref{nuc12:eq:5.1} se expresa en la misma carta de covarianza por



\begin{equation}
\widetilde U=(M_\eta\oplus M_\eta)U_H
(M_\eta^{-1}\oplus M_\eta^{-1}).
\label{nuc12:eq:5.5}
\end{equation}



El lema~\ref{app:lem:implementacion-simplectica} proporciona una implementación unitaria de esta matriz simpléctica en dos modos. Las implementaciones que difieren por una fase global producen la misma conjugación del estado. La covarianza se transforma por \(V\mapsto\widetilde U V\widetilde U^{\mathsf T}\); el estado conjunto permanece puro y su traza parcial sobre el segundo modo determina la lectura visible, conforme a la construcción de la sección~\ref{app:gaussianas-espectro}.

\begin{theorem}[Mezcla reducida por deformación relativa]

La covarianza de la lectura visible es



\begin{equation}
\boxed{V_{\rm vis}=
\frac{\hbar_a}{2}M_\eta
\frac{8I+HH^{\mathsf T}}9
M_\eta^{\mathsf T}.}
\label{nuc12:eq:5.6}
\end{equation}



Su valor simpléctico normalizado es



\begin{equation}
\boxed{\nu(H):=\frac{2\sqrt{\det V_{\rm vis}}}{\hbar_a}
=\sqrt{1+\frac8{81}\left[\operatorname{tr}(HH^{\mathsf T})-2\right]}.}
\label{nuc12:eq:5.7}
\end{equation}
\end{theorem}



\begin{proof}
Sean \(T=(8I+H)/9\) y \(D=\sqrt8(H-I)/9\). Los dos modos iniciales son independientes y tienen la misma covarianza. En la carta equilibrada del estado inicial, la covarianza visible normalizada es \(TT^{\mathsf T}+DD^{\mathsf T}\). La expansión da



\[
\begin{aligned}
TT^{\mathsf T}+DD^{\mathsf T}
&=\frac{64I+8(H+H^{\mathsf T})+HH^{\mathsf T}
+8[HH^{\mathsf T}-H-H^{\mathsf T}+I]}{81}\\
&=\frac{8I+HH^{\mathsf T}}9.
\end{aligned}
\]



Esto demuestra \eqref{nuc12:eq:5.6}. Puesto que \(\det M_\eta=1\) y \(\det(HH^{\mathsf T})=1\),
\(\det(8I+HH^{\mathsf T})=65+8\operatorname{tr}(HH^{\mathsf T})\).
La división por \(81\) demuestra \eqref{nuc12:eq:5.7}. Los dos autovalores positivos de \(HH^{\mathsf T}\) son recíprocos; su suma es al menos dos. En consecuencia, \(\nu\geq1\), con igualdad exactamente cuando \(HH^{\mathsf T}=I\).
\end{proof}

La escala \(\hbar_a\) y la deformación \(\eta\) se conservan hasta efectuar el determinante. Su cancelación es covariante porque se transportan conjuntamente estado y operador. La expresión \(HH^{\mathsf T}\) utiliza el producto euclídeo de la carta equilibrada del estado inicial; un cambio de carta transporta también ese producto métrico.

Un transporte \(H\) ortogonal y simpléctico puede tener \(D\neq0\) y, a la vez, conservar el producto gaussiano isotrópico, con \(\nu=1\). El complemento operatorio y el entrelazamiento del estado tienen, por tanto, criterios distintos. En esta preparación, la deformación relativa no isométrica caracteriza exactamente el caso \(\nu>1\).

\subsubsection{Evaluación exacta del lazo áureo}

Para \eqref{nuc12:eq:5.2},



\begin{equation}
\operatorname{tr}(H_nH_n^{\mathsf T})-2
=n^2(2n^2-2n+5),\qquad
\nu_n^2=1+\frac8{81}n^2(2n^2-2n+5).
\label{nuc12:eq:5.8}
\end{equation}



La identidad resulta de sumar los cuadrados de las cuatro entradas. Para \(n=1\), la suma es siete y



\begin{equation}
\boxed{\nu_1=\frac{11}{9},\quad
\operatorname{Tr}\rho_{\rm vis}^2=\frac9{11},\quad
\bar n_{\rm W}=\frac19,\quad
\frac{S_{\rm vis}}{k_B}=\frac{10}{9}\log10-\log9.}
\label{nuc12:eq:5.9}
\end{equation}



El teorema~\ref{app:thm:espectro-gaussiano} demuestra que el estado gaussiano centrado de un modo con valor simpléctico \(\nu\) es unitariamente equivalente a la densidad geométrica de \eqref{app:eq:estado-geometrico}, cuya ocupación media es \(\bar n_{\rm W}=(\nu-1)/2\). Para \(\nu=11/9\), sus autovalores son



\begin{equation}
\lambda_j=\frac9{10}\,10^{-j},\qquad j=0,1,2,\ldots.
\label{nuc12:eq:5.10}
\end{equation}



La serie suma uno. Su cuadrado suma
\((9/10)^2/(1-10^{-2})=9/11\). Finalmente,


\[
-\sum_j\lambda_j\log\lambda_j
=-\log(9/10)+\log10\sum_jj\lambda_j
=\frac{10}{9}\log10-\log9.
\]



La reducción gaussiana y su espectro están demostrados en la sección~\ref{app:gaussianas-espectro}; las dos series anteriores evalúan exactamente la pureza y la entropía de esta realización. Como el estado conjunto es puro, \(S_{\rm vis}/k_B\) es también su entropía adimensional de entrelazamiento. La transformación global es unitaria, mientras la restricción al primer modo tiene entropía positiva.

El resultado corresponde a la preparación y al transporte especificados. Su identificación con una medición del vacío requiere la realización experimental de ese protocolo. La entropía de una reducción y el calor transferido se mantienen como magnitudes distintas.

\subsubsection{Composición con el transductor térmico}

En la representación normal de Williamson se elige una frecuencia \(\omega>0\) y se define el lector energético
\(\widehat H_{\rm W}=\hbar_a\omega(\widehat N+\tfrac12)\). La frecuencia forma parte de esta realización energética; la covarianza determina las poblaciones. Con esta elección, \eqref{nuc12:eq:5.10} es un estado de Gibbs cuyo cociente de poblaciones satisface



\[
e^{-\hbar_a\omega/(k_BT_{\rm W})}=\frac1{10},\qquad
T_{\rm W}=\frac{\hbar_a\omega}{k_B\log10}.
\]



Si se toma la frecuencia del reloj de la misma sección, \(\omega=1/t_P=2\pi/T_{108}\), la identidad del transductor térmico \eqref{v:ant:eq:ii-kb-aut-par-termico},
\(\hbar_a/t_P=k_B\Theta_{{\rm clk},a}\log3\), da



\begin{equation}
\boxed{\frac{T_{\rm W}}{\Theta_{{\rm clk},a}}=\frac{\log3}{\log10}.}
\label{nuc12:eq:5.11}
\end{equation}



La sección de acción y el transductor térmico se cancelan en el cociente. El índice \(a\) de \(\Theta_{{\rm clk},a}\) conserva la compatibilidad con la sección \(\hbar_a\), como en \eqref{v:ant:eq:ii-kb-aut-orientaciones}. La temperatura \(T_{\rm W}\) es la equivalente de Gibbs para el Hamiltoniano declarado; la evolución unitaria aislada conserva su carácter reversible. Esta identificación espectral determina un cociente térmico y no postula una dinámica de termalización.

\subsection{Extensión a dimensión canónica finita arbitraria}

Sea \(H\in\operatorname{Sp}(2d,\mathbb R)\), con \(d\geq1\) finito. Se preparan dos estados gaussianos puros, centrados e idénticos, de \(d\) modos cada uno, con covarianza \(V_0=(\hbar_a/2)MM^{\mathsf T}\), donde \(M\in\operatorname{Sp}(2d,\mathbb R)\). La sección~\ref{app:gaussianas-espectro} especifica esta preparación y el transporte conjunto de la carta por \(M\oplus M\). En la carta equilibrada, el mismo acoplamiento \(U_H\) produce \(W=(8I+HH^{\mathsf T})/9\). El producto \(A=HH^{\mathsf T}\) es positivo, simétrico y simpléctico.

Existe una base ortonormal simpléctica en la que


\[
A=\operatorname{diag}(e^{2r_1},e^{-2r_1},\ldots,
e^{2r_d},e^{-2r_d}),\qquad r_j\ge0.
\]



La identidad \(A\Omega A=\Omega\), con \(\Omega=\bigoplus_{j=1}^d\left(\begin{smallmatrix}0&1\\-1&0\end{smallmatrix}\right)\), implica que \(\Omega\) lleva el autoespacio de \(\lambda\) al de \(\lambda^{-1}\). Se eligen bases ortonormales emparejadas en esos espacios. El autoespacio de \(\lambda=1\) es \(\Omega\)-invariante y admite parejas canónicas ortonormales. Su reunión construye la base requerida.

La matriz \(W\) es diagonal en esa base. Los productos de cada pareja determinan sus valores simplécticos:



\begin{equation}
\boxed{\nu_j^2=
\frac{(8+e^{2r_j})(8+e^{-2r_j})}{81}
=1+\frac{32}{81}\sinh^2r_j.}
\label{nuc12:eq:6.1}
\end{equation}



El teorema~\ref{app:thm:espectro-gaussiano} compone la reducción por traza parcial con esos valores simplécticos:



\begin{equation}
\boxed{\operatorname{Tr}\rho_{\rm vis}^2=\prod_{j=1}^d\nu_j^{-1},
\qquad
\frac{S_{\rm vis}}{k_B}=\sum_{j=1}^d
\left[\frac{\nu_j+1}{2}\log\frac{\nu_j+1}{2}
-\frac{\nu_j-1}{2}\log\frac{\nu_j-1}{2}\right].}
\label{nuc12:eq:6.2}
\end{equation}



Se adopta \(0\log0=0\) por continuidad. La prueba del teorema~\ref{app:thm:espectro-gaussiano} transforma unitariamente el estado en un producto de \(d\) densidades geométricas: la pureza del producto es el producto de las purezas y su entropía es la suma de las entropías. El resultado vale para cada \(d\) finito y cada transporte \(H\) del dominio, manteniendo todas las ocupaciones del espacio de Fock.

En el lazo \eqref{nuc12:eq:5.3}, \(r=2\log\varphi\) y \(\sinh^2r=5/4\), lo que reproduce \eqref{nuc12:eq:5.9}. Así, el espectro del transporte y la partición nonádica \(8+1\) determinan el espectro de la lectura reducida; las secciones de acción fijan su realización dimensional.

La dimensión de una realización excepcional y su transporte proceden del constructor correspondiente. La ecuación \eqref{nuc12:eq:6.1} determina la ley espectral después de realizar sus modos canónicos y la preparación gaussiana especificada.

\subsubsection{Identidad operatoria entre memoria y estructura compleja}

En la carta canónica equilibrada, sea \(J=\Omega\), con \(J^2=-I\) y \(J^{\mathsf T}=-J\). Esta matriz realiza la estructura compleja canónica del espacio de fases, y \(D=\sqrt8(H-I)/9\) es el bloque de memoria. Para la misma preparación y el mismo transporte de carta de la sección~\ref{app:gaussianas-espectro}, se cumple



\begin{equation}
\boxed{-(JW)^2=I+[D,J][D,J]^{\mathsf T},\qquad
W=\frac{2}{\hbar_a}M^{-1}V_{\rm vis}M^{-\mathsf T}.}
\label{nuc12:eq:6.3}
\end{equation}



\begin{proof}
Las identidades simplécticas \(HJH^{\mathsf T}=J\) y \(H^{\mathsf T}JH=J\), junto con \(A=HH^{\mathsf T}\), dan \(JAJ^{\mathsf T}=A^{-1}\). Si \(C=HJ-JH\), la multiplicación de los cuatro sumandos produce



\[
CC^{\mathsf T}
=A+JAJ^{\mathsf T}
-HJH^{\mathsf T}J^{\mathsf T}
-JHJ^{\mathsf T}H^{\mathsf T}
=A+A^{-1}-2I.
\]



Por otra parte,



\[
-(JW)^2
=\frac{(8I+A^{-1})(8I+A)}{81}
=I+\frac8{81}(A+A^{-1}-2I).
\]



Como \([D,J]=\sqrt8\,C/9\), se obtiene \eqref{nuc12:eq:6.3}. La identidad vale en cada dimensión canónica finita y para cada \(H\) del dominio, incluidos los transportes que mezclan los modos.
\end{proof}

Los autovalores de ese operador son los \(\nu_j^2\), cada uno con multiplicidad dos. La fórmula~\eqref{app:eq:pureza-general} expresa entonces la pureza directamente en términos del bloque de memoria:



\begin{equation}
\boxed{\operatorname{Tr}\rho_{\rm vis}^2
=\det\!\left(I+[D,J][D,J]^{\mathsf T}\right)^{-1/4}.}
\label{nuc12:eq:6.4}
\end{equation}



En un modo,



\begin{equation}
\nu^2=1+\frac12\|[D,J]\|_F^2.
\label{nuc12:eq:6.5}
\end{equation}



En este protocolo, \(D\) conmuta con \(J\) exactamente cuando la lectura reducida conserva la pureza. Para el lazo \(H_1\), \(\|[D,J]\|_F^2=80/81\), y \eqref{nuc12:eq:6.5} da \(\nu^2=121/81\), de donde resulta la pureza \(9/11\).

El conmutador utiliza la estructura compleja compatible con la covarianza inicial y su producto métrico. Un cambio canónico transporta conjuntamente los tres datos; la transpuesta de \eqref{nuc12:eq:6.3} se refiere a la carta equilibrada. La construcción previa del elemento imaginario en el núcleo formal precede a esta representación, que especifica su actuación canónica.

La evolución conjunta conserva la pureza mediante su implementación unitaria; el conmutador determina la pérdida de pureza de la lectura parcial. La información del estado conjunto permanece en los dos factores y en sus correlaciones. La fórmula determina la pureza de la reducción, que es un funcional no lineal de su operador densidad.

La descomposición real respecto de \(J\) distingue sus componentes compleja lineal y compleja antilineal:



\[
D_{\rm lin}=\frac{D-JDJ}{2},\qquad
D_{\rm anti}=\frac{D+JDJ}{2}.
\]



La primera componente conmuta con \(J\) y la segunda anticonmuta con \(J\). Las identidades
\([D,J]=2D_{\rm anti}J\) y \(JJ^{\mathsf T}=I\) implican



\begin{equation}
\boxed{-(JW)^2=I+4D_{\rm anti}D_{\rm anti}^{\mathsf T},\qquad
\operatorname{Tr}\rho_{\rm vis}^2
=\det(I+4D_{\rm anti}D_{\rm anti}^{\mathsf T})^{-1/4}.}
\label{nuc12:eq:6.6}
\end{equation}



La componente compleja antilineal del bloque de memoria, definida respecto de la estructura compleja del estado inicial, determina exactamente la pureza y el espectro de mezcla de este protocolo. La componente compleja lineal puede transportar memoria conservando la pureza. La distinción se refiere a la linealidad respecto de \(J\); la complejidad combinatoria y la clasificación aritmética de constantes son propiedades de otros objetos.

\begin{HMTCompactSection}
\fontsize{10}{12}\selectfont
\titlespacing*{\subsection}{0pt}{8pt}{4pt}
\section{Síntesis de los resultados y consecuencias}
\label{nuclear:sintesis}
La estructura generadora conserva operaciones e historia antes de su realización numérica. En el transporte realizado, esa conservación se formula como una identidad de formas. Una lectura parcial distribuye la carga entre componente visible y complemento; la inclusión de ambos define una transformación recuperable. El refinamiento prolonga la misma identidad sin identificar profundidades diferentes ni sustituir el registro por su intensidad.

\subsection{Conservación, reconstrucción y dimensión}
La igualdad \(\mathcal U^*(b'\oplus b')\mathcal U=b\oplus b\) reúne lectura y memoria en una operación reversible. Para compresiones con registro separado, la telescopía conserva el terminal positivo y la serie completa de complementos. Para una holonomía unitaria fija, el terminal es la proyección sobre su espacio de vectores fijos; para una sucesión variable, lo determina el límite fuerte de los operadores de Gram.

La memoria dodecafásica permite recuperar el registro centrado mediante dos observaciones especificadas, con norma de inversa \(9/4\). La incidencia transporta esa recuperación al proyector dimensional. Su prolongación exterior conserva cada grado mediante los determinantes de Gram y exige incluir componentes mixtas. Así, las relaciones entre longitud, área y volumen se expresan mediante mapas definidos sobre el mismo registro, con sus invariantes respectivos.

\subsection{Acción y consecuencias físicas}
La acción, el reloj y la base energética mantienen sus tipos durante la composición. El cociente electrónico \(m_ec^2/\hbar\) se expresa mediante la frecuencia del reloj y el carácter electrónico, conservando el factor de conversión de base. La normalización de Maxwell--Dirac elimina las escalas dimensionales escogidas y deja el acoplamiento adimensional. El transporte gravitatorio conserva la escala areal \(\hbar G/c^3\); la solución homogénea de rebote mantiene sus condiciones de materia, torsión y constante cosmológica.

La relación de Noether se manifiesta cuando el observable transportado es una carga de la acción declarada. El cambio de forma añade un término de conexión al generador hamiltoniano y conserva exactamente la acción orientada. La holonomía de un retorno conserva, además, la información que distingue retorno de fase y retorno del estado completo.

\subsection{Memoria y espectro cuántico}
La identidad \eqref{nuclear:eq:espectro} determina el espectro simpléctico de la covarianza reducida por la componente antilineal del operador de memoria. Para el lazo
\[
H_1=\begin{pmatrix}2&1\\1&1\end{pmatrix},
\qquad
W_1=\frac19\begin{pmatrix}13&3\\3&10\end{pmatrix},
\]
se obtiene \(\det W_1=121/81\), valor simpléctico \(11/9\) y pureza \(9/11\). El espectro de la densidad reducida es
\[
\lambda_j=\frac9{10}\,10^{-j},\qquad j=0,1,2,\ldots,
\]
y su entropía satisface
\[
\frac{S}{k_B}=\frac{10}{9}\log10-\log9.
\]
Las cifras proceden del transporte y del estado fijados en el teorema. La equivalencia espectral con una densidad de Gibbs asigna una temperatura al Hamiltoniano elegido; esa equivalencia conserva su distinción respecto de un proceso dinámico de termalización.

\subsection{Contenido fundacional}
La articulación entre número, incidencia y dinámica reside en la conservación de un antecedente común y en las operaciones que relacionan sus realizaciones. El artículo hace explícita esa articulación mediante reconstrucciones, cotas, espectros y acciones. La generación de valores y el transporte de simetrías conservan una misma procedencia; sus consecuencias observables se calculan después con el lector y el estado que intervienen efectivamente. Este orden proporciona una formulación matemática precisa de la memoria dinámica y de sus realizaciones físicas.

\end{HMTCompactSection}
\HMTDivision{\(\kappa_{\mathrm{ag}}\)}{Acoplamiento aritmético-geométrico: conservación bilateral y refinamiento reversible}{Registro orientado, norma ternaria, refinamiento reversible y ley bilateral de conservación y recuperación.}
\HMTNarrativeHeading{Proporción, refinamiento y restitución}

La extrema y media razón holográfica reúne aquí la completación
aritmética y la recuperación operatoria. El refinamiento con acarreo
redistribuye una unidad entre componentes y conserva la suma
normalizada. La división euclídea, aplicada con la profundidad y
la convención de acarreo declaradas, recupera las etiquetas
anteriores. La conservación y la inversa reciben así operaciones
concretas.

La norma ternaria aporta el entero setenta y tres y su interpretación
algebraica y reticular. La relación con la completación decimal
se expresa mediante identidades exactas. El balance bilateral
extiende el análisis a dos transportes isométricos entre espacios
de dimensión arbitraria. Los términos cruzados se cancelan en
una combinación ponderada y restituyen la norma cuadrada inicial.

La especialización nonádica determina el reparto inicial de pesos
setenta y dos sobre setenta y tres y uno sobre setenta y tres.
Una mezcla ortogonal produce después los dos canales normalizados.
Cuando los transportes admiten el relativo unitario declarado,
el lector efectivo combina el observable con su conjugado; su
invariancia equivale entonces a la conmutación con ese relativo. El transporte
del estado completo conserva simultáneamente las relaciones
que permite evaluar cada observable.

La política de archivo prolonga este balance. Las memorias
sucesivas y el terminal mantienen la totalidad a profundidad
finita. La estimación uniforme del terminal permite obtener
un archivo infinito recuperable mediante el adjunto del operador
isométrico. Las cotas describen la estabilidad de esa recuperación
cuando aumenta la profundidad.

La constante de acoplamiento participa en la misma construcción.
Su órbita y su forma de Gram permiten identificar operadores;
la isometría de archivo conserva esas respuestas y sus cotas.
El refinamiento de etiquetas se levanta a un transporte unitario
y se compone con el balance bilateral. La unidad conserva así
su contenido aritmético, su realización lineal y la posibilidad
de restitución. Ésta es la culminación operativa del argumento
rector: una estructura generada admite transformaciones que
redistribuyen su lectura y conservan los datos de su recuperación.
\HMTMathematicalPaper
\section{La constante holográfica de acoplamiento aritmético-geométrico}
\label{lib01:acoplamiento}

El registro dodecafásico enlaza dos realizaciones del mismo estado: la
lectura posicional que interviene en la clausura aritmética y la lectura
que conserva incidencia, orientación y dirección. El desarrollo siguiente
determina las operaciones que permite realizar el registro ya producido,
la información que retienen sus memorias y los errores bajo los cuales
pueden recuperarse sus relaciones.
Su función de referencia admite una formulación precisa: la órbita de fase
del registro es una base completa de su portador, y la memoria isométrica
conserva exactamente el Gram de esa base.

El punto de partida sigue siendo la composición
\[
 \APP\longrightarrow\mathrm{TRIT}\longrightarrow\TPK
 \longrightarrow\text{estado enriquecido}
 \longrightarrow\text{estructura discreta conjunta del continuo}.
\]
APP aporta las hojas y sus residuos y cocientes; TRIT conserva régimen y
orientación; TPK selecciona, transporta y actualiza conservando acarreo,
ruta, frontera, supervivencia y memoria. Las cinco construcciones del
continuo pertenecen a esa producción conjunta, desarrollada en el cuerpo
precedente. Aquí se trabaja después de ella, sobre su publicación entera
dodecafásica y sus lectores tipados. Ni el portador lineal posterior ni el
conjunto de candidatos de un decodificador reemplazan el dominio de
historias generadas.

\subsection{Objeto, carta y producción del registro}
\label{lib01:registro}

La definición~\ref{lib01:definicion} fija el objeto generado y sus lectores. Se desarrolla ahora su codificación reversible, su capacidad de identificación operatoria y la estabilidad de sus realizaciones incidenciales.

La denominación expresa una función constitutiva del objeto, no una
identificación automática con un coeficiente de interacción física. La
canonicidad se refiere al perfil, origen y orientación declarados. Una
transformación de carta y una actualización que incorpora nuevos eventos
no son la misma operación. Tampoco se afirma que el mismo vector numérico
deba ser un invariante necesario de todo sistema.

El calendario de ciento ocho posiciones consta de doce ventanas nonádicas.
La extensión
\[
 0\longrightarrow C_{12}\longrightarrow C_{108}
 \longrightarrow C_9\longrightarrow0
\]
conserva el acarreo de fase; su no escisión y la distinción entre retorno de
fase y retorno del estado están demostradas en
\cref{prop:k-extension}. El registro de una historia retiene, en cada
ventana, la subruta, su orientación, el acarreo y el registro de incidencias.
El orden causal de la extracción es
\begin{equation}
 x_{\rm term}\longmapsto\mathcal R_{12}(x_{\rm term})
 \longmapsto(b^{90},b^{120},Q_{\rm TPK})
 \longmapsto U_{\rm sgn}\longmapsto K.
 \label{lib01:cadena-registro}
\end{equation}
Las coordenadas regionales y el registro firmado son dos lecturas del
mismo terminal; no se postula que la matriz ternaria visible determine por
sí sola la coordenada firmada.

Para fijar materialmente la última composición, los canales producidos son
\begin{align*}
 b^{90}&=(-495,-116,-684,108,601,-90,
                -173,-88,313,560,-397,461),\\
 b^{120}&=(-425,-281,-481,671,-255,30,
                 475,-543,680,251,6,-128),\qquad Q_{\rm TPK}=6263.
\end{align*}
Sea \(B_r=\sum_{j=0}^3b^{120}_{r+3j}\) y
\(d_{r,j}=b^{90}_{r+3j}\), para \(r=1,2,3\). El lector armónico da
\[
 s_1=\frac{Q_{\rm TPK}+2B_1+B_2}{3},\qquad
 s_2=s_1-B_1,\qquad s_3=s_2-B_2,
\]
\[
 U^{(r)}=(s_r,d_{r,1}-d_{r,3},d_{r,0}+d_{r,2},d_{r,0}-d_{r,2}).
\]
En esta salida compatible, \((B_1,B_2,B_3)=(972,-1073,101)\),
\((s_1,s_2,s_3)=(2378,1406,2479)\), y
\begin{align*}
 U^{(1)}&=(2378,-452,-668,-322),\\
 U^{(2)}&=(1406,998,-204,-28),\\
 U^{(3)}&=(2479,-551,-371,-997).
\end{align*}
La división por tres pertenece a las tres órbitas del lector y la división
por cuatro que sigue a la inversa de Hadamard; no son coeficientes ajustados
a una salida física. Sobre las órbitas
\((1,4,7,10)\), \((2,5,8,11)\), \((3,6,9,12)\), se aplica
\[
 H_4=\begin{pmatrix}
 1&1&1&1\\1&1&-1&-1\\1&-1&1&-1\\1&-1&-1&1
 \end{pmatrix},\qquad H_4^2=4I.
\]
Los productos \(H_4U^{(r)}\) son, respectivamente,
\[
 (936,2916,2484,3176),\quad(2172,2636,232,584),
 \quad(560,3296,3656,2404).
\]
Su divisibilidad por cuatro verifica la pertenencia a la subred de
Hadamard. Tras dividir y restituir el orden natural se obtiene
\begin{equation}
 K=(234,543,140,729,659,824,621,58,914,794,146,601),
 \qquad\mathbf1^*K=6263.
 \label{lib01:K}
\end{equation}
La inversión y sus congruencias, incluido el índice \(4096\) de la imagen
firmada, están probadas en \cref{lem:k-hadamard,prop:k-sello}. La
reconstrucción es única porque \(H_4^{-1}=H_4/4\), no porque se haya
buscado un registro próximo a un valor conocido de alfa.

Para evitar una colisión de símbolos, escribiremos
\(\Delta_j=I-S^j\) para los extractores enteros antecedentes, con
\((Sx)_i=x_{i+1}\) e índices módulo doce. Entonces
\[
 \Delta_3K=b^{90},\qquad \Delta_4K=b^{120},\qquad Q(K)=6263.
\]
El extractor \((\Delta_3,\Delta_4,Q)\) es inyectivo, con factores
invariantes \(1,\ldots,1,12\), según \cref{prop:k-transversal}. Sus
condiciones de imagen integral y su ley de acumulación permanecen las del
registro antecedente. Los complementos normalizados de memoria que se
denotarán por \(D_j\) son otros operadores, derivados de esos
desplazamientos y con un codominio métrico declarado.

\subsection{Codificación periódica, clausura e inversión exacta}
\label{lib01:kappa}

Fijemos \(B=1000\), doce posiciones y
\[
 N(k)=\sum_{j=1}^{12}k_jB^{12-j},\qquad d_B=B^{12}-1,
 \qquad \kappa(k)=\frac{N(k)}{d_B}.
\]
En la carta canónica,
\begin{equation}
 \begin{aligned}
 \kappa_{\rm ag}
 &=\frac{234543140729659824621058914794146601}{10^{36}-1},\\
 \alpha_A&=A_{reg}-\kappa_{\rm ag},\qquad
 A_{reg}=\pi_{\rm HMT}+e_{\rm HMT}-\varphi_{\rm HMT}-4.
 \end{aligned}
 \label{lib01:clausura}
\end{equation}
Las tres coordenadas regionales se publican antes de actuar en este lector;
su origen común con el cierre dodecafásico no obliga a confundir el orden de
evaluación de las coordenadas correlacionadas. La igualdad es un lector
posterior del estado generado. Despejarla una vez conocidas sus salidas es
una operación inversa legítima, distinta de utilizar un valor objetivo para
producir el registro.

\begin{proposition}[Reversibilidad y radio de la codificación]
\label{lib01:inversion-kappa}
En el dominio \(\{0,\ldots,999\}^{12}\), la fracción \(\kappa(k)\),
con base y longitud declaradas, determina exactamente \(k\). Lo mismo
ocurre con el par exacto \((A_{reg},\alpha_A)\) cuando satisface
\eqref{lib01:clausura}. Si
\[
 |\widetilde\kappa-\kappa(k)|<\frac1{2d_B},
\]
el redondeo de \(d_B\widetilde\kappa\) al entero más próximo recupera
\(N(k)\) sin error. Si se aproximan por separado los dos términos de la
clausura, basta que la suma de sus errores absolutos sea menor que
\(1/(2d_B)\).
\end{proposition}
\begin{proof}
La multiplicación por \(d_B\) recupera el entero \(N(k)\).
Doce divisiones euclídeas por \(B\), conservando los restos y el número
de posiciones, recuperan sus bloques, incluidos los iniciales nulos. La
hipótesis de error sitúa el entero verdadero a distancia estrictamente
menor que media unidad. Para la clausura se aplica la desigualdad
triangular al error de la diferencia. Dos enteros consecutivos muestran
que el radio no puede ampliarse uniformemente en todo el dominio del
código: su punto medio es ambiguo.
\end{proof}

Si se recibe una fracción reducida \(p/q\), se calcula
\(N=d_Bp/q\) y se comprueban integralidad y \(0\leq N\leq d_B\).
El denominador reducido no sustituye los metadatos de base, longitud y
carta. El código contiene \(10^{36}\) palabras y necesita
\(\lceil\log_2(10^{36})\rceil=120\) bits en una representación binaria
de longitud fija. Doce bloques de diez bits ocupan también 120 bits: no se
ha demostrado una compresión adicional, autenticación ni integridad
criptográfica. Una fracción exacta o el entero \(N\) evita perder la
precisión necesaria para esta inversión.

La lectura periódica no se confunde con el truncamiento \(N/B^{12}\).
Si \(\kappa_j=\kappa(S^jK)\), una división posicional da
\begin{equation}
 \kappa_{j+1}=1000\kappa_j-K_{j+1},\qquad
 K_{j+1}=1000\kappa_j-\kappa_{j+1},\qquad
 \sum_{j=0}^{11}\kappa_j=\frac{6263}{999}.
 \label{lib01:orbita-kappa}
\end{equation}
Así, la órbita escalar conserva todas las diferencias de fase y permite
leer \(B_K=\{j+1:1000\kappa_j-\kappa_{j+1}\geq729\}\).
La reversibilidad de estas publicaciones no implica que un escalar
reconstruya toda la historia enriquecida que produjo el registro.

\subsection{La órbita completa como referencia de su portador}
\label{lib01:orbita}

En \(V=\mathbb R^{12}\), con su producto euclídeo canónico, consideremos
\[
 O_K=(K,SK,\ldots,S^{11}K):\mathbb R^{12}\longrightarrow V.
\]
El desplazamiento \(S\) actúa sobre posiciones del registro. No se
identifica por la sola periodicidad con la evolución completa TPK o con
una simetría excepcional posterior.

\begin{theorem}[Referencia cíclica completa y Gram exacto]
\label{lib01:marco-ciclico}
La matriz \(O_K\) es invertible. Su Gram tiene los valores propios
\[
\begin{array}{c|c}
 \text{valor}&\text{multiplicidad}\\\hline
 39225169=6263^2&1\\
 1248025-345420\sqrt3&2\\
 589609&2\\
 1407529&2\\
 1053157&2\\
 1248025+345420\sqrt3&2\\
 697225&1
\end{array}
\]
y satisface
\begin{equation}
 589609\|c\|^2\leq\|O_Kc\|^2\leq39225169\|c\|^2,
 \quad
 \|O_K^{-1}\|=\frac1{\sqrt{589609}},\quad
 \operatorname{cond}(O_K)=\frac{6263}{\sqrt{589609}}.
 \label{lib01:cotas-orbita}
\end{equation}
La órbita del registro centrado tiene rango once.
\end{theorem}
\begin{proof}
Las entradas del Gram son las autocorrelaciones
\(\langle K,S^jK\rangle\). En orden \(j=0,\ldots,11\) valen
\begin{align*}
 (&4251257,2999323,3163385,3287920,3216553,3344743,\\
  &2950064,3344743,3216553,3287920,3163385,2999323).
\end{align*}
Como \(S\) es ortogonal, el Gram es circulante. Sustituir los doce
caracteres \(\exp(2\pi i m/12)\) en esa fila da la tabla real
anterior. Todos los valores son positivos. El menor de los dos valores
con \(\sqrt3\) es mayor que \(589609\), pues
\(658416>345420\sqrt3\), desigualdad que se comprueba elevando al
cuadrado sus dos miembros positivos. Los demás valores de la tabla son
mayores o iguales a \(589609\). El modo uniforme tiene valor
\((\sum K_i)^2=6263^2\); la desigualdad triangular de la transformada
de Fourier, usando \(K_i\geq0\), acota todos los demás por ese mismo
valor. El teorema espectral da \eqref{lib01:cotas-orbita}. Centrar
el registro elimina sólo el modo uniforme; los once restantes conservan
los valores no nulos de la tabla.
\end{proof}

El determinante, con el orden de columnas y la orientación de \(S\)
fijados aquí, es
\[
 \det O_K=5483119259214020183850397930008283625.
\]
La ausencia de un período propio del vector, por sí sola, no probaría la
invertibilidad: lo decisivo es que no se anula ninguno de sus modos de
Fourier. La condición numérica es aproximadamente \(8.15643\). Si se
normaliza \(K\) a norma uno, los extremos del Gram se dividen por
\(\|K\|^2=4251257\); la amplitud de sus componentes no debe
confundirse con una mejora intrínseca frente a toda referencia de igual
norma.

\begin{corollary}[Separación polar de forma y orientación]
\label{lib01:polar}
Sean
\[
 G_K=O_K^*O_K,\qquad P_K=G_K^{1/2},\qquad
 F_K=O_KG_K^{-1/2}.
\]
Entonces \(P_K\) es positivo definido, \(F_K\) es ortogonal y
\(O_K=F_KP_K\). Una realización isométrica \(J:V\to\mathcal Y\)
conserva \(G_K\) y \(P_K\), y transporta la parte isométrica a
\(JF_K\).
\end{corollary}
\begin{proof}
La positividad del Gram permite definir su raíz e inversa mediante el
teorema espectral. La identidad
\(F_K^*F_K=G_K^{-1/2}G_KG_K^{-1/2}=I\) prueba la afirmación.
Además, \((JO_K)^*(JO_K)=O_K^*O_K\).
\end{proof}

Esta polar se define a partir del Gram completo: no se obtiene sólo de la
lista de sus valores propios ni de la norma de \(K\). El símbolo
\(P_K\) de este corolario no es el proyector dimensional de rango diez
que se construirá más abajo. Esta separación es una consecuencia de la
referencia cíclica, no una identificación de todos sus lectores.

\subsection{Identificación de operadores y holonomía relativa}
\label{lib01:identificacion}

\begin{theorem}[Doce respuestas generales o una respuesta cíclica]
\label{lib01:operadores}
Para cualquier operador lineal \(H:V\to V\), sus doce respuestas
\(Y=(HK,HSK,\ldots,HS^{11}K)\) determinan
\[
 H=YO_K^{-1}.
\]
Si \(HS=SH\), una única respuesta vectorial completa \(y=HK\)
determina
\begin{equation}
 H=O_yO_K^{-1},\qquad
 H=\sum_{j=0}^{11}h_jS^j,\qquad h=O_K^{-1}y.
 \label{lib01:filtro}
\end{equation}
Con errores \(E\) en la matriz de respuestas o \(\delta y\) en el
vector, respectivamente,
\begin{align}
 \|\widehat H-H\|&\leq\frac{\|E\|}{\sqrt{589609}},
 \label{lib01:error-matriz}\\
 \|\widehat h-h\|&\leq\frac{\|\delta y\|}{\sqrt{589609}},\qquad
 \|\widehat H-H\|\leq\sqrt{\frac{12}{589609}}\,\|\delta y\|.
 \label{lib01:error-filtro}
\end{align}
\end{theorem}
\begin{proof}
La primera igualdad es \(Y=HO_K\). Si \(H\) conmuta con \(S\),
entonces \(HS^jK=S^jy\), de modo que \(Y=O_y\). La conmutación con
el ciclo determina todas las columnas de \(H\) desde una de ellas;
las doce potencias \(I,S,\ldots,S^{11}\) son independientes y generan
ese espacio de operadores. Por tanto \(y=O_Kh\), lo que da la
segunda fórmula. Las dos primeras cotas usan
\(\|O_K^{-1}\|=1/\sqrt{589609}\). La última usa
\(\|O_{\delta y}\|\leq\|O_{\delta y}\|_{\rm F}
=\sqrt{12}\|\delta y\|\).
\end{proof}

Una respuesta significa aquí un vector de doce componentes, no una
medición escalar. Como ejemplo, si \(H=2I+3S-S^5\), la respuesta
\(y=HK\) devuelve exactamente los coeficientes \(2,3,-1\) en sus
posiciones mediante \eqref{lib01:filtro}; el inversor necesita \(y,K,S\),
no los coeficientes del filtro. En cambio, el operador no nulo cuya primera
fila es \((543,-234,0,\ldots,0)\) y cuyas restantes filas son nulas
anula \(K\). Una sola respuesta no identifica un operador arbitrario.

La función de referencia admite una generalización exacta. En un ciclo de
\(n\) posiciones, \(H\mapsto Hg\) identifica todos los operadores
que conmutan con el ciclo si y sólo si
\(O_g=(g,Sg,\ldots,S^{n-1}g)\) es invertible. La suficiencia es la
prueba anterior. Si \(O_g\) es singular, un vector no nulo de su núcleo
da un polinomio de grado menor que \(n\) que anula \(g\); el operador
es no nulo porque las potencias del ciclo son independientes. El impulso
unitario también tiene órbita ortonormal y sirve de referencia. La
consecuencia específica para \(K\) es que el mismo objeto ya usado en
clausura e incidencia satisface esa condición, sin haber sido elegido de
nuevo para optimizar el ensayo. No se deduce una identificación de
dinámicas no lineales ni de estados ocultos arbitrarios; una aplicación a
una linealización necesita que ésta se haya definido y justificado.

\begin{corollary}[Recuperación de holonomía relativa desde el complemento]
\label{lib01:holonomia}
Sean \(H_B\) conocido e invertible y \(H_C\) otro transporte del
mismo tipo. El complemento
\[
 D_\gamma=\frac{\sqrt8}{9}(H_C-H_B)
\]
determina
\[
 H_B^{-1}H_C=I+\frac9{\sqrt8}H_B^{-1}D_\gamma.
\]
Si \(D_\gamma\) conmuta con \(S\), en particular si lo hacen ambos
transportes, basta \(y=D_\gamma K\) para recuperar
\begin{equation}
 H_B^{-1}H_C
 =I+\frac9{\sqrt8}H_B^{-1}O_yO_K^{-1}.
 \label{lib01:holonomia-una-respuesta}
\end{equation}
El error en esta reconstrucción está acotado por
\[
 \frac9{\sqrt8}\|H_B^{-1}\|
 \sqrt{\frac{12}{589609}}\,\|\delta y\|.
\]
\end{corollary}
\begin{proof}
Se despeja \(H_C\) en la definición del complemento y se aplica
\cref{lib01:operadores}. La cota resulta de la submultiplicatividad.
\end{proof}

Para \(H_B=S^3\) y \(H_C=S^4\), el dato racional
\(y/\sqrt8=(S^4-S^3)K/9\) permite reconstruir la holonomía relativa
\(S\) sin proporcionar \(H_C\) al inversor. Si \(H_B\) es
unitario, el factor de error es aproximadamente \(0.01435510\). Sin
conmutación se mantienen las doce respuestas generales, no una hipótesis
ficticia de simetría. La selección de un lazo físico y su representación
siguen perteneciendo a su dinámica generadora; este protocolo no las
sustituye.

\subsection{La memoria completa conserva la referencia y los observables}
\label{lib01:memoria-completa}

Sean \(C_n\) operadores unitarios de un espacio de Hilbert \(H\),
\[
 T_n=\frac{8I+C_n}{9},\qquad
 D_n=\frac{\sqrt8(C_n-I)}9,\qquad
 R_0=I,\quad R_{n+1}=T_nR_n.
\]
La escritura incluye secuencias no estacionarias. No exige conmutación
entre \(C_n\) de etapas distintas.

\begin{theorem}[Registro isométrico, terminal y transporte]
\label{lib01:isometria}
Para todo \(N\),
\[
 R_N^*R_N+\sum_{n=0}^{N-1}R_n^*D_n^*D_nR_n=I.
\]
Existe el límite fuerte positivo
\(A_\infty=\lim_N R_N^*R_N\), y
\begin{equation}
 \mathcal Vx=
 \bigl(A_\infty^{1/2}x,(D_nR_nx)_{n\geq0}\bigr)
 \quad\text{satisface}\quad\mathcal V^*\mathcal V=I.
 \label{lib01:V-infinita}
\end{equation}
Su imagen \(\mathcal M\) es cerrada y la inversa sobre ella es
\(\mathcal V^*\), de norma uno. En el caso finito se sustituye el
primer componente por \(R_Nx\) y se conservan los \(N\)
complementos. Para cualquier operador acotado \(A\),
\[
 A^{\mathcal V}=\mathcal VA\mathcal V^*\big|_{\mathcal M}
\]
preserva productos, adjuntos, norma y espectro en esa imagen, y
\[
 \langle\mathcal Vx,A^{\mathcal V}\mathcal Vy\rangle
 =\langle x,Ay\rangle.
\]
\end{theorem}
\begin{proof}
Expandir los dos términos de una etapa da
\(T_n^*T_n+D_n^*D_n=I\). Por tanto
\(R_n^*R_n-R_{n+1}^*R_{n+1}=R_n^*D_n^*D_nR_n\), y la suma
telescópica demuestra la identidad finita. La sucesión de operadores
positivos \(R_N^*R_N\) es decreciente y acotada inferiormente; su
límite fuerte existe. Tomando productos con \(x\), la suma de las
normas cuadradas de los complementos es
\(\|x\|^2-\langle x,A_\infty x\rangle\). Esto prueba que la
serie pertenece a la suma hilbertiana y que \(\mathcal V\) es una
isometría. Una isometría tiene imagen cerrada e inversa adjunta en ella.
Insertar \(\mathcal V^*\mathcal V=I\) en las composiciones demuestra
las propiedades del transporte. No se ha supuesto convergencia de
\(R_Nx\), ni que \(A_\infty\) sea un proyector.
\end{proof}

Al aplicar este resultado al portador del registro,
\begin{equation}
 (\mathcal VO_K)^*(\mathcal VO_K)=O_K^*O_K.
 \label{lib01:gram-memoria}
\end{equation}
Se conservan las doce direcciones y todas sus correlaciones, con las
mismas cotas de \cref{lib01:marco-ciclico}, a cualquier profundidad de
memoria completa. No se afirma que el continuo completo tenga dimensión
doce: se conserva esta referencia de doce dimensiones en su imagen.
Los operadores correctos allí son
\(S^{\mathcal V}=\mathcal VS\mathcal V^*\) y
\(H^{\mathcal V}=\mathcal VH\mathcal V^*\). No se presume que un
desplazamiento ambiente de las casillas de almacenamiento sea
\(S^{\mathcal V}\).

Si \(z=\mathcal Vy+\delta z\), la estimación
\[
 \widehat h=O_K^{-1}\mathcal V^*z
\]
mantiene la primera cota de \eqref{lib01:error-filtro} con
\(\|\delta z\|\) en lugar de \(\|\delta y\|\). La profundidad
no empeora esa constante porque se conserva la isometría y su terminal.
Una selección parcial de canales tiene otro Gram y requiere sus propias
cotas. Las potencias exteriores \(\bigwedge^p\mathcal V\) conservan
los determinantes de Gram; también se conservan trazas de rango finito y
medidas espectrales de los observables transportados. Las normas
cuadradas no se interpretan por ello automáticamente como probabilidades
de Born: tal lectura exige su realización posterior.

\subsection{Incidencia excepcional y dirección seleccionada}
\label{lib01:direccion}

Sea \(e=\mathbf1\), \(\Pi=I-ee^*/12=P_{11}\). En la carta
icosaédrica \(A_5/C_5\) ya construida, la representación ortogonal
\(\rho\) y el carácter tridimensional fijan
\begin{equation}
 P_3=\frac3{60}\sum_{g\in A_5}\chi_3(g^{-1})\rho(g),
 \quad P_3^*=P_3=P_3^2,\quad
 \operatorname{rank}P_3=3,\quad P_3e=0.
 \label{lib01:P3}
\end{equation}
Este proyector se recibe con su representación y sus marcas, no se elige
para adaptar una dirección al registro. Pongamos
\[
 k=\Pi K,\qquad u=P_3k,\qquad
 E_u=\frac{uu^*}{\|u\|^2},\qquad Q_K=\Pi-E_u=P_{10}(K).
\]
Las evaluaciones exactas son
\begin{align}
 \|k\|^2&=\frac{11789915}{12},&
 \|u\|^2&=\frac{6638585+2275584\sqrt5}{20},
 \label{lib01:normas-direccion}\\
 \eta_K:=\frac{\|u\|^2}{\|k\|^2}
 &=\frac{3(6638585+2275584\sqrt5)}{5\cdot11789915}
 \simeq0.5967954228259091.
 \label{lib01:eta}
\end{align}
En particular, \(u\ne0\). Como \(u^*k=\|u\|^2\),
\begin{equation}
 E_uk=u,\qquad Q_Kk=k-u,\qquad
 k=u+Q_Kk,\qquad \|k\|^2=\|u\|^2+\|Q_Kk\|^2.
 \label{lib01:particion-direccion}
\end{equation}
La reducción \(12\to11\to10\) elimina primero la dirección uniforme
y después una recta seleccionada. No elimina todo \(\operatorname{im}P_3\):
en \(\operatorname{im}Q_K\) permanecen las dos direcciones de ese
subespacio ortogonales a \(u\). Estas dimensiones son las de los
espacios algebraicos escritos; su identificación con dimensiones físicas
requiere sus mapas de realización.

\begin{proposition}[Transporte exacto de dirección y descriptor]
\label{lib01:transporte-direccion}
Para una isometría completa \(\mathcal V\), sea
\(\Pi^{\mathcal V}=\mathcal V\Pi\mathcal V^*\),
\(P_3^{\mathcal V}=\mathcal VP_3\mathcal V^*\) y
\(u^{\mathcal V}=\mathcal Vu\). En \(\mathcal M\),
\[
 Q_K^{\mathcal V}
 =\mathcal VQ_K\mathcal V^*
 =\Pi^{\mathcal V}
 -\frac{u^{\mathcal V}(u^{\mathcal V})^*}{\|u^{\mathcal V}\|^2},
 \quad \operatorname{rank}Q_K^{\mathcal V}=10,
 \quad \eta_K^{\mathcal V}=\eta_K.
\]
Además, \(\mathcal V^*u^{\mathcal V}=u\) y
\(\mathcal V^*Q_K^{\mathcal V}\mathcal V=Q_K\).
\end{proposition}
\begin{proof}
Se sustituyen las definiciones y \(\mathcal V^*\mathcal V=I\).
Las normas y el rango se conservan por la isometría. La identidad de la
imagen es \(\mathcal V\mathcal V^*\), no la identidad de un
almacenamiento ambiente mayor. Extender por cero fuera de la imagen no
añade direcciones físicas.
\end{proof}

La pantalla incidencial completa ofrece otra realización exacta. Si
\(\mathcal I\) es el operador de incidencia ya construido y
\(E_{\rm inc}=\mathcal I(e^\perp)\), entonces
\begin{equation}
 F:V\longrightarrow Y=\mathbb R\oplus E_{\rm inc},\quad
 Fx=\left(\mu(x),\frac{\mathcal I\Pi x}{6}\right),\quad
 \mu(x)=\frac{\mathbf1^*x}{12},\qquad
 \langle(\mu,y),(\nu,z)\rangle_Y=12\mu\nu+\langle y,z\rangle
 \label{lib01:incidencia-F}
\end{equation}
es una isometría sobreyectiva, pues
\(\Pi\mathcal I^*\mathcal I\Pi=36\Pi\). Para la incidencia
cruda de las 132 hexadas, la identidad completa es
\(\mathcal I^*\mathcal I=36I+30\mathbf1\mathbf1^*\); el
centrado elimina exactamente su componente uniforme. La imagen incidencial tiene
dimensión once; \(Y\) no es todo \(\mathbb R\oplus\mathbb R^{132}\).
Aplicar \(F\) término a término a la memoria conserva tanto el Gram
como los operadores anteriores y sus dominios. No se introduce por esta
composición una acción del Monstruo sobre \(K\).

\subsection{Estabilidad continua y necesidad de conservar correlaciones}
\label{lib01:ruido}

\begin{theorem}[Error de la dirección desde memoria completa]
\label{lib01:error-proyectores}
Sea \(z=\mathcal VK+\delta\), \(\|\delta\|\leq\varepsilon\),
y reconstruyamos
\(\widehat K=\mathcal V^*z\), \(\widehat k=\Pi\widehat K\),
\(\widehat u=P_3\widehat k\). Los tres errores son menores o iguales
a \(\varepsilon\). Si \(\varepsilon<\|u\|\),
\[
 \|Q_{\widehat K}-Q_K\|\leq\frac\varepsilon{\|u\|}.
\]
Si \(\varepsilon<\|k\|\),
\[
 |\eta_{\widehat K}-\eta_K|\leq\frac\varepsilon{\|k\|}.
\]
\end{theorem}
\begin{proof}
El adjunto de la isometría y los proyectores tienen norma uno. El ruido
ortogonal a \(\mathcal M\) desaparece al aplicar \(\mathcal V^*\).
Para dos rectas no nulas, la norma de la diferencia de sus proyectores es
el seno del ángulo entre ellas. La distancia de \(u\) a la recta de
\(\widehat u\) no supera \(\|u-\widehat u\|\), por lo que ese
seno es menor o igual que \(\varepsilon/\|u\|\). La hipótesis
impide \(\widehat u=0\).

Para la segunda afirmación se normalizan \(k\) y \(\widehat k\).
La diferencia de sus proyectores de rango uno tiene, en su plano, valores
propios \(\pm\sin\theta\). Al emparejarla con \(0\leq P_3\leq I\)
el valor absoluto no supera \(\sin\theta\): en una base de vectores
propios se obtiene \(\sin\theta(a-b)\), con \(a,b\in[0,1]\).
El mismo argumento de distancia da
\(\sin\theta\leq\varepsilon/\|k\|\). Si las rectas coinciden,
la diferencia de los cocientes es cero.
\end{proof}

Aquí \(\|u\|\simeq765.733\), por lo que el factor para el proyector
es aproximadamente \(0.00130594\). En el lector de sólo dos memorias,
la inversa tiene norma \(9/4\) y el factor correspondiente es
\(9/(4\|u\|)\simeq0.00293836\). Son datos observados distintos;
la diferencia no es una mejora gratuita del mismo canal. La corrección
entera de \cref{lib02:radios-memoria} permitirá recuperar primero el
registro sin error dentro de su radio, en lugar de propagar un error
continuo hasta el proyector.

\begin{proposition}[Los bloques cruzados forman parte del observable]
\label{lib01:correlaciones}
Escribamos \(\mathcal Vx=(W_ax)_a\), incluyendo el terminal. Entonces
el bloque \((a,b)\) de \(A^{\mathcal V}\) es
\[
 (A^{\mathcal V})_{ab}=W_aAW_b^*.
\]
Su forma cuadrática se calcula como límite de truncamientos finitos
\[
 \langle y,A^{\mathcal V}y\rangle
 =\lim_F\sum_{a,b\in F}\langle y_a,W_aAW_b^*y_b\rangle.
\]
La conservación general no se obtiene descartando todos los términos
\(a\ne b\).
\end{proposition}
\begin{proof}
La fórmula de bloques resulta de multiplicar \(\mathcal VA\mathcal V^*\).
Los proyectores de truncamiento convergen fuertemente a la identidad y
\(A^{\mathcal V}\) es acotado, lo que prueba el límite sin afirmar
convergencia absoluta de una serie doble arbitraria. Para un
contraejemplo, tome \(H=\mathbb R\), \(C=-I\),
\(T=7/9\), \(D=-2\sqrt8/9\), \(\mathcal Vx=(Tx,Dx)\) y
\(A=I\). La parte diagonal, evaluada en \(\mathcal Vx\), vale
\[
 (T^4+D^4)x^2=\frac{3425}{6561}x^2.
\]
Los términos cruzados aportan \(3136x^2/6561\), y la suma es
\(x^2\), como exige la isometría.
\end{proof}

\subsection{Corrección discreta antes de la selección incidencial}
\label{lib01:correccion-incidencia}

Los lectores canónicos de dos memorias son
\[
 T_j=\frac{8I+S^j}{9},\qquad
 D_j=\frac{\sqrt8(S^j-I)}9,\qquad
 Mx=(D_3x,D_4T_3x).
\]
La segunda memoria actúa después de la primera compresión; no se reinyecta
en este protocolo el primer complemento. El teorema general de
\cref{lib02:radios-memoria} demostrará, sin utilizar el registro como
objetivo, los radios estrictos
\[
 r_0=\frac{2\sqrt{146}}{81},\qquad
 r_q=\frac{4\sqrt{73}}{81}.
\]
El primer radio recupera \(\Pi K\) a partir de \(MK\), módulo la
traslación entera uniforme. El segundo recupera \(K\) íntegro cuando
se conserva además \(q=\mathbf1^*K\). Si se observa una carga real
con error menor que \(1/2\), su redondeo recupera antes el canal entero
de carga. Ninguno de estos procedimientos necesita el terminal
\(T_4T_3K\).

La selección de bloque alto en la carta marcada es
\begin{equation}
 B_K=\{i:K_i\geq729\}=\{4,6,9,10\}.
 \label{lib01:umbral}
\end{equation}
La cuarta componente está exactamente en el umbral. Esta circunstancia
separa la estabilidad continua de la discreta.

\begin{proposition}[Margen continuo nulo y restitución entera]
\label{lib01:umbral-estable}
El selector directo de \eqref{lib01:umbral}, aplicado a la reconstrucción
real de las memorias, no tiene radio euclídeo positivo de estabilidad en
\(MK\), incluso con carga exacta. En cambio, la decodificación entera
previa con error menor que \(r_q\) restituye exactamente \(B_K\)
sin modificar el umbral.
\end{proposition}
\begin{proof}
Sea \(L\) la inversa de mínimos cuadrados, que satisface
\(LM=\Pi\). Para \(t>0\), ponga
\[
 e_t=-tM\Pi e_4,
 \qquad \widetilde K_t=\frac q{12}\mathbf1+L(MK+e_t)
 =K-t\Pi e_4.
\]
Se tiene \(\|e_t\|\to0\), pero la cuarta componente es
\(729-11t/12<729\). Para \(t\) pequeño sólo cambia esa pertenencia.
La segunda afirmación sigue de recuperar primero \(K\) exactamente
mediante \cref{lib02:radios-memoria} y aplicar después el mismo selector.
\end{proof}

Si los restantes datos marcados \(P_\pi,H_e,H_\varphi,H_{\rm APP}\)
se conservan exactos o se transportan por su ley, la unicidad de la
construcción excepcional devuelve la misma bandera
\(B_K\subset H^-_\alpha\) y el mismo origen marcado. Corregir el
registro no equivale a corregir errores independientes de esas marcas.
El testigo
\[
 K'=K-e_4+e_6,
 \qquad Q(K')=6263,\qquad B_{K'}=\{6,9,10\}
\]
pertenece también a \(\{0,\ldots,999\}^{12}\) y cumple
\(\|MK'-MK\|^2=4672/6561\). El punto medio de ambas memorias está
a distancia \(r_q\) de cada una. Prueba que el radio estricto es
óptimo uniformemente en el dominio algebraico de carga fija, no que
\(K'\) tenga una historia terminal HMT construida.

Sin carga, \(M(K-\mathbf1)=MK\) y las selecciones altas son distintas.
Las dos memorias pueden recuperar exactamente \(Q_K\) y \(\eta_K\),
que ignoran la media, sin recuperar ese selector absoluto. Para el
proyector sigue siendo necesaria la condición \(P_3\Pi K\ne0\);
la corrección entera no suprime esa condición de definición. Un ensayo
finito concreto usa \(z=MK/\sqrt8\) y añade \(1/10\) a su primera
componente. El error en la norma de \(M\) es \(\sqrt8/10<r_q\).
El algoritmo de suelo y techo de \cref{lib02:decodificador} deja 924
candidatos de carga 6263 y recupera \(K\) y \(B_K\) exactamente.

Aplicando la isometría \eqref{lib01:incidencia-F} se obtiene
\(M_{\rm inc}=(F\oplus F)M\). Todas las distancias entre palabras
y las normas de ruido se conservan, de modo que los mismos radios y
decodificadores valen en \(Y\oplus Y\). Ésta es la conexión
cuantitativa entre memoria e incidencia; no depende de atribuir al
desplazamiento \(S\) una simetría excepcional adicional.

\subsection{Covarianza, refinamiento y carga variacional}
\label{lib01:covarianza}

Si \(G\) es ortogonal y se transportan conjuntamente
\(K'=GK\), \(\Pi'=G\Pi G^*\) y \(P_3'=GP_3G^*\), entonces
\[
 u'=Gu,\qquad Q_{K'}'=GQ_KG^*,\qquad \eta_{K'}'=\eta_K.
\]
En una carta fija, \(K\mapsto aK+b\mathbf1\), \(a\ne0\), deja
invariantes \(Q_K\) y \(\eta_K\); no deja invariante todo el objeto.
Por ejemplo, \(\kappa(K+\mathbf1)-\kappa(K)=1/999\), y la lectura
\(\alpha_A\) cambiaría en \(-1/999\) si se mantuviesen fijas las
regiones. Es un contraejemplo algebraico a identificar el descriptor con
todo el registro, no una afirmación de admisibilidad de esa actualización
como historia HMT.

Tampoco se puede mover sólo el registro y mantener arbitrariamente el
proyector. La fórmula finita \eqref{lib01:P3} da exactamente
\[
 \|P_3S^3K\|^2=\frac{369035}{4}-\frac{281861\sqrt5}{10},
 \qquad
 \|P_3S^4K\|^2=\frac{1327717}{4}+\frac{694618\sqrt5}{5}.
\]
Son distintas de \eqref{lib01:normas-direccion}; en particular,
\([P_3,S^3]\) y \([P_3,S^4]\) no se anulan. Al transportar
\(P_3\) por la misma conjugación sí se conserva \(\eta_K\).

La decodificación también es covariante. Para una permutación \(R\),
\[
 S'=RSR^{-1},\qquad M'=(R\oplus R)MR^{-1},\qquad
 M'R=(R\oplus R)M.
\]
La red, la carga y las distancias se conservan. En el régimen de unicidad,
la salida del decodificador se transforma por \(R\); fuera de él se
transporta el conjunto de minimizadores, sin elegir un empate de forma
arbitraria. Un marco ortogonal general exige transportar también
\(\mathbb Z^{12}\) a \(R\mathbb Z^{12}\) y la dirección uniforme
a \(R\mathbf1\). Redondear en la red antigua después de cambiar la
carta no representa el mismo problema.

\begin{proposition}[Naturalidad bajo refinamiento con marcos locales]
\label{lib01:refinamiento}
Sea \(z\) un cilindro con hijos \(z'\), pesos positivos compatibles
\(\sum_{\tau z'=z}\mu(z')=\mu(z)\), y marcos ortogonales \(G_z\).
El refinamiento
\[
 \mathcal R_n(e_z\otimes x)=
 \sum_{\tau z'=z}\sqrt{\frac{\mu(z')}{\mu(z)}}
 e_{z'}\otimes G_{z'}G_z^*x
\]
es isométrico. Para el campo de proyectores
\(Q_z=G_zQ_KG_z^*\) y su operador por bloques \(Q_n\),
\[
 Q_{n+1}\mathcal R_n=\mathcal R_nQ_n,
 \qquad \mathcal R_n^*Q_{n+1}\mathcal R_n=Q_n.
\]
Las identidades se componen entre niveles y pasan al límite inductivo.
\end{proposition}
\begin{proof}
Los hijos tienen soportes ortogonales y los pesos suman uno; los cambios
de marco son ortogonales. En cada hijo,
\(Q_{z'}G_{z'}G_z^*=G_{z'}G_z^*Q_z\), que demuestra el primer
cuadrado. Multiplicar por el adjunto del refinamiento da el segundo.
La compatibilidad por composición define el operador en el límite de
las inclusiones isométricas. El mismo argumento vale para \(E_u\).
\end{proof}

El enunciado conserva la publicación de un prefijo bajo refinamientos
compatibles. Nuevos eventos pueden producir otro registro; no se deduce
una constancia temporal de \(K\). Los rangos de campos sobre fibras
tensoriales son los de esas fibras, no un recuento de dimensiones físicas.

\begin{proposition}[Una carga de transporte con acción explícita]
\label{lib01:carga-noether}
Sean \(U_n:E_n\to E_{n+1}\) transportes invertibles. La acción
auxiliar discreta con variables cotangentes es
\[
 \mathcal A(q,p)=\sum_n p_{n+1}^*(q_{n+1}-U_nq_n).
\]
Las variaciones de soporte finito dan
\(q_{n+1}=U_nq_n\) y \(p_n=U_n^*p_{n+1}\). Si
\(G_0=I\), \(G_{n+1}=U_nG_n\), \(A_0=Q_K\) y
\(A_n=G_nA_0G_n^{-1}\), la cantidad
\[
 J_n=p_n^*A_nq_n
\]
es independiente de \(n\). La acción es invariante bajo el grupo
continuo
\(q_n\mapsto\exp(tA_n)q_n\),
\(p_n\mapsto\exp(-tA_n^*)p_n\).
\end{proposition}
\begin{proof}
Las ecuaciones indicadas resultan al variar \(p_{n+1}\) y \(q_n\).
La definición de \(A_n\) implica
\(A_{n+1}U_n=U_nA_n\). Por tanto
\[
 p_{n+1}^*A_{n+1}q_{n+1}
 =p_{n+1}^*U_nA_nq_n=p_n^*A_nq_n.
\]
El mismo entrelazamiento vale para las exponenciales y hace que cada
término de la acción sea invariante. No se requiere que \(A_n\)
sea antisimétrico: la transformación del covector es la contragrediente.
\end{proof}

Para transportes ortogonales puede elegirse \(p_n=q_n\); entonces
\(J_n=\|A_nq_n\|^2\), y el complemento proporciona la carga
restante. En una realización simpléctica sin forma positiva, se conserva
el emparejamiento escrito y no se lo llama energía positiva. Si se
archivan las fibras mediante isometrías \(\mathcal V_n\), las
definiciones
\[
 \widehat U_n=\mathcal V_{n+1}U_n\mathcal V_n^*,\qquad
 \widehat A_n=\mathcal V_nA_n\mathcal V_n^*,\qquad
 \widehat q_n=\mathcal V_nq_n,
\]
con \(\widehat p_n(y)=p_n(\mathcal V_n^*y)\), conservan el
entrelazamiento y la misma carga en las imágenes. Esta composición
proporciona una acción y una simetría continuas concretas para una ley de
conservación; no atribuye corrientes de Noether automáticas a un grupo
finito, ni identifica la acción auxiliar con toda acción física.

\subsection{Codificación, identificación operatoria y estabilidad de la incidencia}
\label{lib01:alcance}

La aportación reunida es una cadena operativa: el registro generado tiene
una codificación reversible; su órbita es una referencia completa;
terminal y complementos conservan esa referencia y los observables
transportados; la separación entera de dos memorias permite restituir
incidencia y dirección aun cuando el selector real tiene margen cero.
Las pruebas de linealidad, de proyección y de redes se utilizan como
lenguaje de realización posterior, sin cambiar la producción del objeto.

Cada hipótesis tiene un testigo que impide ampliarla inadvertidamente.
Una respuesta aislada no identifica un operador que no conmuta con el
ciclo; sólo la memoria completa conserva automáticamente el Gram original;
los términos cruzados no pueden omitirse; la selección absoluta exige
carga; el proyector requiere dirección no nula; las cotas de ruido
dependen de la normalización; una transformación de carta exige transportar
forma, red y lectores. Una perturbación de lectura no se convierte por
decreto en una historia admisible. Estas delimitaciones preservan las
funciones de \(K\), sin reducirlo a doce números desconectados ni
promover una propiedad canónica a invariancia universal.

\section{Norma algebraica ternaria y radios óptimos de recuperación}
\label{lib02:norma-ternaria}

El radio de recuperación de las dos memorias contiene el factor
\(\sqrt{73}\). Su origen puede seguirse sin introducir una magnitud
objetivo: procede del Gram del lector nonádico y, en otra realización del
ciclo ternario ya construido, es la escala de una semejanza. Su cuadrado,
\(73\), es el índice de la inclusión reticular de rango dos. Estas
apariciones se reúnen mediante operaciones explícitas;
no se identifican por una semejanza decimal con alfa. La comparación
posterior con \(10^4\alpha_A\) se tratará como un problema inverso,
con sus datos, sus transformaciones permitidas y sus conclusiones exactas.

La materia prima conserva su producción APP--TRIT--TPK y el corte
posterior descrito en \cref{lib01:registro}. El registro entero, la
orientación de los desplazamientos, la carga y los lectores se reciben
antes de calcular distancias. El caso canónico utiliza \(8+1\) y el
denominador nueve. El parámetro \(a>0\) que se introduce a continuación
es una familia algebraica posterior, no una nueva semilla ni un parámetro
físico seleccionado por alfa. Los errores pertenecen al canal de lectura;
no se los identifica con eventos o historias admisibles.

\subsection{Familia de dos memorias y su Gram}
\label{lib02:familia-memoria}

En \(V=\mathbb R^{12}\), sea \((Sx)_i=x_{i+1}\), con índices módulo
doce. Para \(a>0\) y \(j=3,4\), definimos
\begin{equation}
 T_{j,a}=\frac{aI+S^j}{a+1},\qquad
 D_{j,a}=\frac{\sqrt a(S^j-I)}{a+1},\qquad
 M_ax=(D_{3,a}x,D_{4,a}T_{3,a}x).
 \label{lib02:lectores}
\end{equation}
La segunda memoria se produce después de la primera compresión. Archivar
ambas no equivale a reinyectar el primer complemento en una evolución
reversible de dos canales. Sean
\[
 p(a)=a^2+a+1,\quad A(a)=a(a^2+1),\quad B(a)=a^2,
 \qquad\Pi=I-\frac{\mathbf1\mathbf1^*}{12},\quad P_u=I-\Pi.
\]

\begin{proposition}[Balance y forma de diferencias]
\label{lib02:gram}
Para todo \(a>0\),
\[
 T_{j,a}^*T_{j,a}+D_{j,a}^*D_{j,a}=I,\qquad
 G_a:=M_a^*M_a=I-(T_{4,a}T_{3,a})^*(T_{4,a}T_{3,a}),
\]
y
\begin{align}
 (a+1)^4G_a={}&4ap(a)I
 -A(a)(S^3+S^4+S^8+S^9)\notag\\
 &-B(a)(S+S^5+S^7+S^{11}).
 \label{lib02:gram-explicito}
\end{align}
Equivalentemente,
\begin{align}
 (a+1)^4\|M_ah\|^2={}&A(a)\sum_i
 \bigl((h_i-h_{i+3})^2+(h_i-h_{i+4})^2\bigr)\notag\\
 &+B(a)\sum_i
 \bigl((h_i-h_{i+1})^2+(h_i-h_{i+5})^2\bigr).
 \label{lib02:energia-aristas}
\end{align}
\end{proposition}
\begin{proof}
La ortogonalidad de \(S\) da
\[
 T_{j,a}^*T_{j,a}
 =\frac{(a^2+1)I+a(S^j+S^{-j})}{(a+1)^2},\qquad
 D_{j,a}^*D_{j,a}
 =\frac{a(2I-S^j-S^{-j})}{(a+1)^2}.
\]
Su suma es \(I\); aplicar la identidad a las dos etapas cancela el
término intermedio \(T_{3,a}^*T_{3,a}\). Como todos los operadores
son polinomios de \(S\), el producto de los dos Grams individuales
tiene diagonal \((a^2+1)^2\), términos simples de peso
\(a(a^2+1)\) y cruzados de peso \(a^2\). La identidad
\((a+1)^4-(a^2+1)^2=4a(a^2+a+1)\) da
\eqref{lib02:gram-explicito}. Expandir las diferencias cuadradas da
\eqref{lib02:energia-aristas}; cada arista no orientada de cada paso
aparece una sola vez y el grado ponderado es \(4A+4B=4ap(a)\).
\end{proof}

En \(a=8\) se obtiene
\begin{align}
 6561G_8={}&2336I-520(S^3+S^4+S^8+S^9)\notag\\
 &-64(S+S^5+S^7+S^{11}),
 \label{lib02:gram-nonadico}
\end{align}
cuya primera fila es
\((2336,-64,0,-520,-520,-64,0,-64,-520,-520,0,-64)/6561\).
También
\[
 T_{4,8}T_{3,8}=\frac{64I+8S^3+8S^4+S^7}{81}.
\]
Los cuatro desplazamientos distintos producen la contribución diagonal
\[
 1-\frac{64^2+8^2+8^2+1}{81^2}
 =\frac{2336}{6561}
 =\frac{32(8^2+8+1)}{9^4}.
\]

\begin{theorem}[Núcleo, espectro e inversa completa]
\label{lib02:inversa-memorias}
Para todo \(a>0\), \(\ker M_a=\mathbb R\mathbf1\). Los valores
propios de \(G_a\), indexados por los modos duodécimos \(m\), son
\[
\begin{array}{c|c|c}
 m&\lambda_m(G_a)&\text{multiplicidad}\\\hline
 0&0&1\\
 3,9&2a/(a+1)^2&2\\
 4,8&3a/(a+1)^2&2\\
 6&4a/(a+1)^2&1\\
 1,5,7,11&a(5a^2+4a+5)/(a+1)^4&4\\
 2,10&a(7a^2+2a+7)/(a+1)^4&2
\end{array}
\]
En particular,
\begin{equation}
 L_a^{\rm rec}=(G_a+P_u)^{-1}M_a^*,\qquad
 L_a^{\rm rec}M_a=\Pi,\qquad
 \|L_a^{\rm rec}\|=\frac{a+1}{\sqrt{2a}}.
 \label{lib02:inversa}
\end{equation}
\end{theorem}
\begin{proof}
Si \(M_ah=0\), la primera memoria implica \(S^3h=h\), luego
\(T_{3,a}h=h\); la segunda implica \(S^4h=h\). Como
\(S=S^4(S^3)^{-1}\), el vector es uniforme. La recíproca es
inmediata. En el modo \(m\), el balance da
\[
 \lambda_m(G_a)=1-
 \frac{[a^2+1+2a\cos(\pi m/2)]
 [a^2+1+2a\cos(2\pi m/3)]}{(a+1)^4}.
\]
Los doce pares de cosenos dan la tabla. Respecto de
\(2a/(a+1)^2\), las diferencias de los dos últimos valores son
\[
 \frac{3a(a^2+1)}{(a+1)^4}>0,\qquad
 \frac{a(5a^2-2a+5)}{(a+1)^4}>0.
\]
Los otros valores positivos son múltiplos mayores. Por tanto el mínimo
en \(\mathbf1^\perp\) es \(2a/(a+1)^2\). El operador
\(G_a+P_u\) vale uno sobre la uniforme y es positivo invertible en
su complemento. La fórmula escrita es la inversa de mínimos cuadrados
y su norma es el inverso de la menor singularidad no nula de \(M_a\).
\end{proof}

En \(a=1\), \(T_{3,1}=(I+S^3)/2\) es singular, pero
\eqref{lib02:inversa} sigue siendo válida. La singularidad de una
compresión individual no elimina la recuperación mediante sus dos
memorias. En \(a=8\), la norma de la inversa es \(9/4\).

\subsection{Distancias mínimas: todos los registros enteros}
\label{lib02:redes}

Los dominios algebraicos de decodificación son
\[
 \mathcal C_0(a)=M_a\mathbb Z^{12},\qquad
 \Lambda_q=\{k\in\mathbb Z^{12}:\mathbf1^*k=q\},\qquad
 \mathcal C_q(a)=M_a\Lambda_q.
\]
La primera imagen se identifica con la red \(\Pi\mathbb Z^{12}\)
o con clases módulo \(\mathbb Z\mathbf1\); la segunda conserva una
carga entera fijada. Son dominios que contienen las publicaciones enteras
del tipo indicado. Su uso no afirma que cada candidato corresponda a una
historia HMT producida.

\begin{lemma}[Corte mínimo del grafo de memoria]
\label{lib02:cortes}
Todo corte propio no vacío de las doce posiciones tiene al menos cuatro
aristas de pasos \(1,5\) y al menos cuatro de pasos \(3,4\). Su peso
en \eqref{lib02:energia-aristas} es, por tanto, al menos
\(4A(a)+4B(a)=4ap(a)\); un corte de una posición alcanza esa cota.
\end{lemma}
\begin{proof}
Cada paso \(1\) y \(5\) forma un ciclo conexo de doce posiciones;
cada ciclo cruza un corte no trivial al menos dos veces. Para los pasos
\(3,4\), si el conjunto no es invariante por ninguno, cada familia
aporta al menos dos aristas. Si es invariante por \(S^3\), es unión
de clases módulo tres; cada uno de los cuatro triángulos de \(S^4\)
lo cruza dos veces, dando ocho aristas. Si es invariante por \(S^4\),
es unión de clases módulo cuatro; cada uno de los tres ciclos de longitud
cuatro de \(S^3\) lo cruza al menos dos veces, dando al menos seis.
La invariancia simultánea obligaría a invariancia por \(S\), imposible
para un conjunto propio no vacío. Todos los pesos son positivos; el
corte de una sola posición tiene exactamente cuatro aristas de cada
peso.
\end{proof}

\begin{theorem}[Distancias exactas y minimizadores de carga fija]
\label{lib02:distancias}
Para todo \(a>0\) y toda carga entera \(q\),
\begin{equation}
 d_0(a)^2=\min_{h\notin\mathbb Z\mathbf1}\|M_ah\|^2
 =\frac{4ap(a)}{(a+1)^4},\qquad
 d_q(a)^2=\frac{8ap(a)}{(a+1)^4}.
 \label{lib02:distancias-exactas}
\end{equation}
En el primer mínimo \(h\) recorre \(\mathbb Z^{12}\); el segundo
es la distancia mínima entre palabras distintas de \(\mathcal C_q(a)\).
Las diferencias que alcanzan el segundo mínimo son exactamente
\(e_i-e_j\) con \(j-i\equiv2,6,10\pmod{12}\).
\end{theorem}
\begin{proof}
Para un entero \(h\) no constante,
\((h_i-h_j)^2\geq|h_i-h_j|\). Cada diferencia absoluta se
descompone como suma de los cortes de los conjuntos de nivel
\(\{i:h_i>t\}\). Hay al menos un nivel no trivial; por
\cref{lib02:cortes}, la energía ponderada es al menos \(4ap(a)\).
El vector \(e_i\) alcanza esa energía. Dividir por \((a+1)^4\)
prueba el primer mínimo sobre todos los enteros, no sólo sobre
indicadores de subconjuntos.

Para carga fija, la diferencia \(h\ne0\) es entera y de suma cero.
Su norma cuadrada es par porque \(h_i^2\equiv h_i\pmod2\). Si
\(\|h\|^2\geq4\), el mínimo espectral da
\[
 \|M_ah\|^2\geq\frac{8a}{(a+1)^2}
 >\frac{8ap(a)}{(a+1)^4},
\]
ya que \((a+1)^2-p(a)=a>0\). Sólo queda
\(\|h\|^2=2\), es decir, \(h=e_i-e_j\). Su energía es
\(2((G_a)_{ii}-(G_a)_{ij})\). Los elementos no diagonales son cero
en separaciones \(2,6,10\), negativos de peso \(B(a)/(a+1)^4\)
en \(1,5,7,11\), y negativos de peso \(A(a)/(a+1)^4\) en
\(3,4,8,9\). La diagonal es \(4ap(a)/(a+1)^4\); esto prueba el
mínimo y la clasificación exhaustiva para toda la semirrecta \(a>0\).
\end{proof}

\begin{theorem}[Radios estrictos de corrección de las memorias]
\label{lib02:radios-memoria}
Para un dato \(y=M_ak+e\), con error en la norma conjunta de las dos
memorias, los radios uniformes óptimos son
\begin{equation}
 r_0(a)=\frac{\sqrt{ap(a)}}{(a+1)^2},\qquad
 r_q(a)=\frac{\sqrt{2ap(a)}}{(a+1)^2}.
 \label{lib02:radios}
\end{equation}
Si \(\|e\|<r_0(a)\), el vecino más próximo de
\(\mathcal C_0(a)\) recupera exactamente \(\Pi k\). Con carga
\(q\) conocida y \(\|e\|<r_q(a)\), el vecino más próximo de
\(\mathcal C_q(a)\) recupera exactamente \(k\).
\end{theorem}
\begin{proof}
La red \(\Pi\mathbb Z^{12}\subset(1/12)\mathbb Z^{12}\) es
discreta y \(M_a\) está acotado inferiormente en
\(\mathbf1^\perp\). Sus imágenes son discretas y admiten vecino
más próximo. Si \(c'\) es otra palabra de código y \(d\) es la
distancia mínima correspondiente,
\(\|y-c'\|\geq d-\|e\|>\|e\|\) cuando \(\|e\|<d/2\).
El mínimo es único y la preimagen es única en cada dominio indicado.
El punto medio de dos palabras a distancia mínima está a \(d/2\)
de ambas: ningún decodificador puede distinguir uniformemente ambas
entradas desde ese mismo dato. La optimalidad se refiere a los dominios
algebraicos completos, no a un subconjunto adicional de historias que
pudiera tener una separación mayor.
\end{proof}

En el caso canónico,
\begin{equation}
 r_0(8)=\frac{2\sqrt{146}}{81}\simeq0.298346814162829,
 \qquad
 r_q(8)=\frac{4\sqrt{73}}{81}\simeq0.421926110879878.
 \label{lib02:radio73}
\end{equation}
El \(4\) procede de \(\sqrt{2\cdot8}\) al tomar la mitad de la
distancia entre registros de carga fija; \(73=8^2+8+1\) procede del
Gram; \(81=9^2\) procede de las dos normalizaciones nonádicas. El mismo
producto \(9\times9\) cuenta las celdas APP, pero esos dos papeles
no se identifican sin su aplicación tipada. Tampoco el factor cuatro
prueba por sí solo una afirmación sobre direcciones espaciales.

Si las memorias se almacenan como \(z=M_ak/\sqrt a\), la norma de
ruido cambia. En \(a=8\), los radios de ese dato racional son
\(r_0/\sqrt8=\sqrt{73}/81\) y
\(r_q/\sqrt8=\sqrt{146}/81\). Usar los radios de
\eqref{lib02:radio73} sin esta conversión es cambiar el problema.
La aplicación a \(Q_K\), \(\eta_K\), \(B_K\) y la pantalla
incidencial se ha probado en \cref{lib01:correccion-incidencia}; en
particular, el umbral exacto \(K_4=729\) no impide su restitución
discreta.

\subsection{Decodificador finito y dualidad del parámetro}
\label{lib02:computacion}

\begin{proposition}[Suelo y techo sin tabla objetivo]
\label{lib02:decodificador}
Con carga conocida, bajo el radio \(r_q(a)\), basta examinar a lo
sumo \(2^{12}=4096\) candidatos enteros. Sin carga, bajo
\(r_0(a)\), basta examinar a lo sumo \(12\cdot2^{12}=49152\)
candidatos repartidos entre doce residuos de carga. El mínimo residual
devuelve respectivamente \(k\) o su clase uniforme.
\end{proposition}
\begin{proof}
Con carga \(q\), forme
\(\widetilde k=L_a^{\rm rec}y+(q/12)\mathbf1\). Entonces
\[
 \|\widetilde k-k\|
 <\|L_a^{\rm rec}\|r_q(a)
 =\frac{\sqrt{p(a)}}{a+1}<1.
\]
Cada coordenada entera verdadera pertenece a
\(\{\lfloor\widetilde k_i\rfloor,
\lceil\widetilde k_i\rceil\}\). Se enumeran esos productos,
se eliminan los de suma distinta de \(q\) y se minimiza
\(\|M_ak'-y\|\). El candidato correcto está presente y
\cref{lib02:radios-memoria} asegura su unicidad. Coordenadas ya enteras
reducen el número de candidatos.

Sin carga, escriba \(\rho\in\{0,\ldots,11\}\) para el residuo
de la carga. Cada clase módulo \(\mathbb Z\mathbf1\) tiene un
único representante cuya suma es \(\rho\), obtenido restando
\(\lfloor q/12\rfloor\mathbf1\). Para cada \(\rho\), forme
\(\widetilde k_\rho=L_a^{\rm rec}y+(\rho/12)\mathbf1\).
En el residuo correcto,
\[
 \|\widetilde k_\rho-k_\rho\|
 <\|L_a^{\rm rec}\|r_0(a)
 =\frac{\sqrt{p(a)}}{\sqrt2(a+1)}<1.
\]
Se aplica suelo y techo, se conserva la suma \(\rho\) y se minimiza
el residual entre los doce grupos. La unicidad de la clase de código
da una única salida centrada. Estos procedimientos usan el lector y la
integralidad, no reciben \(K\) como tabla de comparación.
\end{proof}

Para \(a\) racional se puede conservar aritmética racional escribiendo
\(R_a=M_a/\sqrt a\), \(z=y/\sqrt a\). La reconstrucción centrada
es
\[
 (R_a^*R_a+P_u)^{-1}R_a^*z=L_a^{\rm rec}y.
\]
El residual puede compararse mediante sus cuadrados racionales. Estas
garantías se refieren a doce coordenadas: no afirman complejidad
polinómica para una dimensión arbitraria, autenticación criptográfica ni
ganancia cuántica. El ejemplo perturbado de
\cref{lib01:correccion-incidencia} ilustra el algoritmo canónico.

\begin{proposition}[Dualidad recíproca y máximo del radio]
\label{lib02:dualidad-a}
La sustitución \(a\mapsto1/a\) conserva el Gram, las distancias,
los radios y la norma de la inversa. Se tienen
\[
 D_{j,1/a}=D_{j,a},\qquad
 T_{j,1/a}=S^jT_{j,a}^*.
\]
Los radios máximos de esta familia se alcanzan únicamente en \(a=1\):
\[
 r_q(a)\leq\frac{\sqrt6}{4},\qquad
 r_0(a)\leq\frac{\sqrt3}{4}.
\]
\end{proposition}
\begin{proof}
Las dos identidades operatorias resultan al multiplicar numerador y
denominador por \(a\). Los Grams individuales y la fórmula
\eqref{lib02:gram} son invariantes bajo la reciprocidad. Ponga
\(t=a/(a+1)^2\in(0,1/4]\). Entonces
\[
 r_q(a)^2=2t(1-t),\qquad r_0(a)^2=t(1-t).
\]
Ambas expresiones crecen para \(0<t\leq1/4\), y \(t=1/4\)
equivale a \(a=1\). La reciprocidad deja \(t\) fijo. Al tender
\(a\) a cero o a infinito, \(t\) y los radios tienden a cero.
\end{proof}

La igualdad de Grams no significa que \(M_a=M_{1/a}\): el segundo
complemento incluye una compresión diferente. Los extremos \(a=0\)
y \(a=\infty\) no pertenecen al dominio del teorema. Optimizar la
familia en \(a=1\), o utilizar el representante recíproco \(1/8\),
no cambia retroactivamente el origen canónico \(a=8\).

\subsection{La norma en el módulo de aumento ternario}
\label{lib02:modulo-ternario}

El ciclo \(r\mapsto4r\pmod9\), con órbitas \((1,4,7)\) y
\((2,8,5)\), proporciona el módulo de aumento cuya carta y forma son
\[
 C=\begin{pmatrix}0&-1\\1&-1\end{pmatrix},\qquad
 G_{A_2}=\begin{pmatrix}2&-1\\-1&2\end{pmatrix},\qquad
 C^2+C+I=0,\qquad C^TG_{A_2}C=G_{A_2}.
\]
La acción y su orientación contragrediente pertenecen a las fibras
ternarias antecedentes. Sobre ese módulo ya construido se componen
\[
 Q_-(a)=aI-C,\qquad Q_+(a)=(a+1)I+C.
\]

\begin{proposition}[Semejanza ternaria e índice entero]
\label{lib02:norma-A2}
Para \(a>0\),
\[
 Q_-Q_+=p(a)I,\qquad Q_-^{\dagger_{G_{A_2}}}=Q_+,
 \qquad Q_-^TG_{A_2}Q_-=p(a)G_{A_2},\qquad \det Q_-=p(a).
\]
En \(a=8\), ambas cartas conjugadas son
\[
 Q_-=
 \begin{pmatrix}8&1\\-1&9\end{pmatrix},\qquad
 Q_+=\begin{pmatrix}9&-1\\1&8\end{pmatrix},\qquad Q_-Q_+=73I.
\]
Cada cociente entero es cíclico de orden 73. En el cociente de \(Q_-\),
\(C\) actúa por \(8\); en el de \(Q_+\), por
\(-9\equiv64\equiv8^{-1}\pmod{73}\).
\end{proposition}
\begin{proof}
Expandir el producto y sustituir \(C^2=-C-I\) da \(p(a)I\).
La invariancia de la forma implica
\(C^{\dagger_{G_{A_2}}}=C^{-1}=C^2=-I-C\), de donde se obtiene
el adjunto. Componer producto y adjunto da la semejanza. El determinante
de \(\bigl(\begin{smallmatrix}a&1\\-1&a+1\end{smallmatrix}\bigr)\)
es \(p(a)\). En \(a=8\), el máximo común divisor de las entradas
de cada matriz es uno y su determinante es 73; la forma de Smith es
\(\operatorname{diag}(1,73)\). La relación \(8I-C=0\) en el
primer cociente y \(9I+C=0\) en el segundo da las acciones indicadas.
Como \(8^3\equiv1\pmod{73}\) y \(8\not\equiv1\), ambas
acciones tienen orden tres y orientaciones inversas.
\end{proof}

Así, \(\sqrt{73}\) es una escala de semejanza en la métrica ternaria
heredada, no sólo una cifra dentro de un radio. Su reconocimiento complejo
posterior es \(|8-\omega|^2=73\), con
\(\omega^2+\omega+1=0\). La unidad compleja describe la realización
del operador; no fue utilizada para seleccionar el registro. El operador
\(Q_-\) no se identifica por estas igualdades con \(K\), con toda
la dinámica TPK ni con el conúcleo del lector de escala de acción.

\subsection{El ciclo completo y el origen del término fijo}
\label{lib02:ciclo-completo}

En las doce posiciones, sea \(U=S^4\). Ahora hay cuatro ciclos de
longitud tres. Sus proyectores son
\[
 P_0=\frac{I+U+U^2}{3},\qquad P=I-P_0.
\]
El rango de \(P_0\) es cuatro y el de \(P\) ocho. En particular,
\(P_0\ne P_u\), pues este último sólo conserva la uniforme global;
tampoco se confunde \(P\) con el proyector excepcional \(P_3\).
Utilizaremos ahora \(Q_{-,a}=aI-U\) y \(Q_{+,a}=(a+1)I+U\)
para los operadores del ciclo completo.

\begin{theorem}[Gram ternario con sector fijo]
\label{lib02:gram-ternario}
Se tienen
\begin{equation}
 Q_{-,a}^*Q_{-,a}
 =(a-1)^2P_0+p(a)P=p(a)I-3aP_0.
 \label{lib02:gram-Qmenos}
\end{equation}
Para \(v\ne0\), sean
\[
 w(v)=\frac{\|P_0v\|^2}{\|v\|^2},\qquad
 s_{-,a}(v)=\frac{\|Q_{-,a}v\|^2}{\|v\|^2}.
\]
Entonces \(s_{-,a}(v)=p(a)-3aw(v)\). En el caso canónico,
\begin{equation}
 Q_-^*Q_-=49P_0+73P=73I-24P_0,
 \qquad s(v)=73-24w(v)\in[49,73].
 \label{lib02:73-24w}
\end{equation}
\end{theorem}
\begin{proof}
De \(U^3=I\) y \(U^*=U^2\) resulta que \(P_0\) es el
proyector ortogonal sobre \(\ker(U-I)\). Al expandir el Gram,
\[
 Q_{-,a}^*Q_{-,a}=(a^2+1)I-a(U+U^2).
\]
La relación \(U+U^2=3P_0-I\) da
\eqref{lib02:gram-Qmenos}. Evaluar sobre \(v\) y dividir por su
norma cuadrada demuestra las demás identidades.
\end{proof}

El término \(24\) es exactamente \(3\cdot8\): es la diferencia
\(p(8)-(8-1)^2\) entre la lectura de aumento y la fija. Ya aparece
en un solo ciclo de tres posiciones. Para \(m\) ciclos, el rango fijo
es \(m\), la dimensión es \(3m\) y
\[
 \operatorname{tr}(Q_{-,a}^*Q_{-,a})=3m(a^2+1).
\]
En \(m=4,a=8\), la traza es 780 y la corrección total de la traza
es \(24\cdot4=96\). La igualdad numérica \(24=2\cdot12\) no es
el origen operatorio del coeficiente.

Las dos cartas completas distinguen también sus determinantes. En un
solo ciclo,
\[
 \det(8I-U)=8^3-1=511=7\cdot73,
 \qquad
 \det(9I+U)=9^3+1=730=10\cdot73.
\]
El factor fijo es, respectivamente, 7 y 10; el módulo de aumento aporta
73. Sobre las doce posiciones,
\(\det(8I-U)=511^4=68184176641\). Se conserva la identidad
\begin{equation}
 73=8^2+8+1=9^2-9+1=\frac{9^3+1}{9+1}
 =\frac{729+1}{10}.
 \label{lib02:completacion73}
\end{equation}
Su interpretación como índice y factor de determinante precede a
cualquier lectura de completación o refinamiento posterior.

\subsection{Seis estados recibidos y conservación de la lectura completa}
\label{lib02:seis-estados}

Fijamos de nuevo \(a=8\) y omitimos ese subíndice en los lectores.
Los seis objetos que se comparan son
\[
 K,\quad k=\Pi K,\quad m_0=D_3K,\quad m_1=D_4T_3K,
 \quad u=P_3k,\quad q_K=Q_KK=k-u.
\]
Ninguno se selecciona a partir del objetivo del ensayo. Para dar las
evaluaciones sin sustituirlas por decimales, escribamos
\(a_v=\|v\|^2\), \(b_v=\|P_0v\|^2\). Se obtiene
\begin{equation}
\begin{array}{c|c|c}
 v&a_v&b_v\\\hline
 K&4251257&10684363/3\\
 k&11789915/12&1170761/4\\
 m_0&15413392/81&16838032/243\\
 m_1&1105759232/6561&0\\
 u&1327717/4+(568896/5)\sqrt5
   &3029657/20+(145120/3)\sqrt5\\
 q_K&1951691/3-(568896/5)\sqrt5
   &2292404/15-(145120/3)\sqrt5
\end{array}
 \label{lib02:tabla-normas}
\end{equation}
La tabla se obtiene aplicando \(P_0=(I+S^4+S^8)/3\), las fórmulas
de \(D_3,D_4T_3\) y el proyector finito de carácter
\eqref{lib01:P3} a \eqref{lib01:K}. Todas sus entradas pertenecen a
\(\mathbb Q(\sqrt5)\); las primeras cuatro son racionales. En
particular, los tres primeros pesos y lecturas son
\begin{align*}
 w(K)&=\frac{10684363}{12753771},&
 s(K)&=\frac{224866857}{4251257},\\
 w(k)&=\frac{3512283}{11789915},&
 s(k)&=\frac{776369003}{11789915},\\
 w(m_0)&=\frac{1052377}{2890011},&
 s(m_0)&=\frac{61904585}{963337}.
\end{align*}
Para \(u\) y \(q_K\), las expresiones exactas son
\(w=b_v/a_v\), \(s=73-24b_v/a_v\), con los numeradores y
denominadores de \eqref{lib02:tabla-normas}. Las evaluaciones posteriores
son
\[
\begin{array}{c|c|c}
 v&w(v)\text{, aproximadamente}&s(v)\text{, aproximadamente}\\\hline
 K&0.8377414805&52.89420446705527\\
 k&0.2979057101&65.85026295779062\\
 m_0&0.3641429046&64.26057028848679\\
 m_1&0&73\text{ exacto}\\
 u&0.4428244530591254&62.37221312658099\\
 q_K&0.1127385160275117&70.29427561533972
\end{array}
\]
La igualdad de la segunda memoria es forzada:
\(P_0D_4T_3=\sqrt8P_0(U-I)T_3/9=0\), mientras
\(\|m_1\|^2>0\). No es una coincidencia de evaluación.

\begin{proposition}[La memoria completa no modifica el peso fijo]
\label{lib02:memoria-ternaria}
La aplicación
\[
 V_2v=(t(v),m_0(v),m_1(v))
 =(T_4T_3v,D_3v,D_4T_3v)
\]
es isométrica y entrelaza \(P_0\) y \(Q_-=8I-U\) con sus
copias diagonales. Conserva por separado \(\|v\|^2\),
\(\|P_0v\|^2\) y \(\|Q_-v\|^2\). Por tanto las lecturas
conjuntas \(w\) y \(s\) no cambian.
\end{proposition}
\begin{proof}
El balance de dos etapas da \(V_2^*V_2=I\). Cada canal es un
polinomio en \(S\), por lo que conmuta con \(U,P_0,Q_-\).
Aplicar la isometría a \(v\), \(P_0v\) y \(Q_-v\) da las tres
igualdades de normas. Al dividir, la lectura total es la media de las
lecturas de los canales ponderada por sus normas cuadradas, no su media
aritmética sin pesos. Los canales nulos no aportan al cociente total.
\end{proof}

El terminal confirma materialmente esa partición. Para \(K\),
\[
 \|t(K)\|^2=\frac{25538253193}{6561},\qquad
 \|P_0t(K)\|^2=\frac{848595371}{243},
\]
\[
 w(t(K))=\frac{22912075017}{25538253193},\qquad
 s(t(K))=\frac{1314402682681}{25538253193}
 \simeq51.46799480558349.
\]
La suma de las tres normas cuadradas es \(4251257\), y la de las
normas fijas es
\((848595371+16838032)/243=10684363/3\). Para \(k=\Pi K\),
\[
 \|t(k)\|^2=\frac{16367568169}{26244},\qquad
 \|P_0t(k)\|^2=\frac{217142795}{972}.
\]
Los complementos no cambian al centrar y sus sumas recuperan las normas
de la segunda fila de \eqref{lib02:tabla-normas}. Descartar el terminal
o un complemento altera la lectura parcial; conservarlos no introduce
una corrección oculta.

En una memoria infinita completa se mantiene la misma conclusión cuando
los canales conmutan con \(U\): también el límite positivo y su raíz
conmutan con él. Sin esa conmutación se usa el transporte
\(\mathcal VP_0\mathcal V^*\) en la imagen, conforme a
\cref{lib01:isometria}, no una copia diagonal supuesta.

\subsection{Inversión, cartas y balance bilateral}
\label{lib02:bidireccional}

\begin{proposition}[Qué conserva una inversión o un cambio de marco]
\label{lib02:inversion-carta}
Sustituir \(U\) por \(U^{-1}\) intercambia cada \(Q_{\pm,a}\)
con su adjunto y conserva \(P_0\), sus Grams, valores singulares,
determinantes y distancias inducidas. Para todo \(v\ne0\), conserva
\(w(v)\) y \(s_{-,a}(v)\). No es lo mismo que invertir
\(Q_{-,a}\), cuya inversa, para \(a\ne1\), es
\begin{equation}
 Q_{-,a}^{-1}
 =\frac1{a-1}P_0+\frac1{p(a)}Q_{+,a}P.
 \label{lib02:inversa-Q}
\end{equation}
Para una carta invertible \(R\), transportando
\(U'=RUR^{-1}\), \(v'=Rv\) y
\(g'=(R^{-1})^*R^{-1}\), se conservan esas razones de normas en
la métrica \(g'\).
\end{proposition}
\begin{proof}
Como \(U^*=U^{-1}=U^2\), la inversión produce los adjuntos.
Los operadores son polinomios normales en \(U\); sus Grams dependen
de \(U+U^2\), que no cambia. La inversa escrita se verifica por
sectores: \(Q_{-,a}\) actúa por \(a-1\) en el fijo y el producto
con \(Q_{+,a}\) vale \(p(a)I\) en el aumento. Para la carta,
\(Q'v'=RQv\) y \(R^*g'R=I\); numerador y denominador recuperan
exactamente sus valores originales. El adjunto respecto de \(g'\)
es también el transporte del adjunto original.
\end{proof}

La reversión del orden de las doce posiciones satisface
\(RSR^{-1}=S^{-1}\) y conserva \(P_0\). Las traslaciones \(S^j\),
incluida la antipodal \(S^6\), conmutan con \(U\); una negación
global tampoco altera las normas. Ninguna de esas operaciones convierte
por sí sola 73 en una lectura distinta.

Con la métrica euclídea fija, un marco ortogonal aplicado sólo al estado
conserva \(s_{-,a}(v)\) para todos los \(v\) si y sólo si
\(R^*P_0R=P_0\), por \eqref{lib02:gram-Qmenos} y polarización.
El grupo de estas transformaciones es
\(O(\operatorname{im}P_0)\times O(\operatorname{im}P)
\cong O(4)\times O(8)\), mayor que el conmutante de \(U\).
En cambio, un filtro de estado
\(c_0P_0+c_1P\), cuando su salida es no nula, produce
\begin{equation}
 w'=
 \frac{|c_0|^2w}{|c_0|^2w+|c_1|^2(1-w)}.
 \label{lib02:filtro-pesos}
\end{equation}
Puede cambiar la lectura. Elegir esos coeficientes desde un objetivo es
un ajuste declarado del filtro, no un cambio de carta de la misma
lectura.

\begin{theorem}[Balance bilateral y recuperación explícita]
\label{lib02:balance-bilateral}
Para todo \(a>0\), en el ciclo completo se cumplen
\begin{align}
 Q_{+,a}-Q_{-,a}^*&=3P_0,&
 Q_{+,a}Q_{-,a}&=p(a)I-3P_0,\notag\\
 Q_{+,a}^*Q_{+,a}&=p(a)I+3(a+1)P_0,
 \label{lib02:dos-cartas}
\end{align}
y
\begin{equation}
 (a+1)Q_{-,a}^*Q_{-,a}
 +aQ_{+,a}^*Q_{+,a}=(2a+1)p(a)I.
 \label{lib02:balance-pesado}
\end{equation}
La aplicación
\begin{equation}
 W_av=\frac{(\sqrt{a+1}\,Q_{-,a}v,
                    \sqrt a\,Q_{+,a}v)}{\sqrt{(2a+1)p(a)}}
 \label{lib02:W}
\end{equation}
es isométrica y se invierte en su imagen mediante \(W_a^*\), de
norma uno. Para los canales sin normalizar \(x=Q_{-,a}v\),
\(y=Q_{+,a}v\),
\begin{equation}
 v=\frac{x+y}{2a+1},\qquad
 Uv=\frac{ay-(a+1)x}{2a+1}.
 \label{lib02:reconstruccion-bilateral}
\end{equation}
Un par \((x,y)\) pertenece a esa imagen sin normalizar si y sólo si
\(Q_{+,a}x=Q_{-,a}y\).
\end{theorem}
\begin{proof}
En \(Q_+-Q_-^*\) aparece \(I+U+U^2=3P_0\). Expandir
\(Q_+Q_-\), \(Q_+^*Q_+\) y usar
\(U^2=3P_0-I-U\) da \eqref{lib02:dos-cartas}. La combinación
ponderada cancela los términos de \(P_0\) en
\eqref{lib02:gram-Qmenos} y \eqref{lib02:dos-cartas}; esto prueba
\(W_a^*W_a=I\). Como \(Q_-+Q_+=(2a+1)I\), la suma de los
canales da \(v\); su combinación \(ay-(a+1)x\) da
\((2a+1)Uv\).

La condición de imagen es necesaria porque ambos operadores conmutan.
Si se cumple, defina \(v=(x+y)/(2a+1)\). Entonces
\[
 Q_-v=\frac{Q_-x+Q_-y}{2a+1}
 =\frac{(Q_-+Q_+)x}{2a+1}=x,
\]
y análogamente \(Q_+v=y\). La fórmula de la inversa adjunta es la
de toda isometría; para el par crudo, la inversa lineal
\((x,y)\mapsto(x+y)/(2a+1)\) tiene norma
\(\sqrt2/(2a+1)\) en el espacio producto.
\end{proof}

En \(a=8\), los sectores fijo y de aumento tienen lecturas
\((49,100)\) y \((73,73)\), respectivamente, y
\[
 9\|Q_-v\|^2+8\|Q_+v\|^2=17\cdot73\|v\|^2=1241\|v\|^2,
 \qquad v=\frac{Q_-v+Q_+v}{17}.
\]
Para \(s_{+,a}(v)=\|Q_{+,a}v\|^2/\|v\|^2\),
\[
 s_{+,a}=p(a)+3(a+1)w,
 \qquad (a+1)s_{-,a}+a s_{+,a}=(2a+1)p(a).
\]
La cancelación vale aun con componente fija. La composición bilateral es
una formalización sobre operadores antecedentes, no la afirmación de que
todo retorno TPK sea automáticamente este par.

Si \(\mathcal V:H\to\mathcal M\) es una memoria isométrica
completa, \(U^{\mathcal V}=\mathcal VU\mathcal V^*\) satisface
\((U^{\mathcal V})^3=\Pi_{\mathcal M}\), donde
\(\Pi_{\mathcal M}=\mathcal V\mathcal V^*\) es la identidad en
la imagen. Allí el Gram transportado es
\[
 (Q_{-,a}^{\mathcal V})^*Q_{-,a}^{\mathcal V}
 =p(a)\Pi_{\mathcal M}-3a\mathcal VP_0\mathcal V^*.
\]
Conserva las lecturas y la recuperación bilateral. Usar copias diagonales
de \(U\) sólo es legítimo cuando los canales entrelazan esa acción;
en caso contrario el transporte por \(\mathcal V\) es la definición
correcta sobre la imagen.

\subsection{El dato de alfa y las familias de iteraciones enteras}
\label{lib02:ensayo-inverso}

El lector completo de alfa, con acarreo, proporciona el prefijo por
bloques
\[
 (007,297,352,569,283,800,997,285,105,472,380,662,999,\ldots).
\]
La ventana de frontera cero que termina en 663 es otro objeto; no se la
intercambia con la salida completa. Para el ensayo basta el intervalo
cerrado heredado
\begin{equation}
 \frac{7297352569283800997285105472380662}{10^{36}}
 \leq\alpha_A\leq
 \frac{7297352569283800997285105472380663}{10^{36}}.
 \label{lib02:intervalo-alpha}
\end{equation}
No se identifica la constante con ninguno de sus extremos racionales.
Pongamos \(s_\alpha=10^4\alpha_A\), de modo que
\begin{equation}
 \begin{gathered}
 \ell\leq s_\alpha\leq h,\\
 \ell:=72.97352569283800997285105472380662,\\
 h:=72.97352569283800997285105472380663.
 \end{gathered}
 \label{lib02:intervalo-s}
\end{equation}
La condición inversa del lector fijo es exactamente
\begin{equation}
 s(v)=s_\alpha
 \quad\Longleftrightarrow\quad
 w(v)=\frac{73-s_\alpha}{24}
 \simeq0.00110309613175.
 \label{lib02:criterio-alpha}
\end{equation}

Ninguno de los seis estados de \cref{lib02:seis-estados} satisface esa
condición. La segunda memoria tiene \(s=73\), mientras que las otras
cinco tienen \(s<72\). Esta última desigualdad se comprueba exactamente
como \(24b_v>a_v\). En las tres filas racionales basta multiplicar
enteros; para \(u\) y \(q_K\) las diferencias son
\[
 24b_u-a_u=\frac{66073183}{20}+\frac{5235904}{5}\sqrt5>0,
\]
\[
 24b_{q_K}-a_{q_K}
 =\frac{45259241}{15}-\frac{5235904}{5}\sqrt5>0.
\]
Para la segunda, \(\sqrt5<5/2\) da ya una cota inferior positiva.
La exclusión concierne a este lector y estos estados; no declara una
ausencia global de relaciones con alfa.

Como \(0<s_\alpha<73\), el radio
\[
 r_\alpha=\frac{4\sqrt{s_\alpha}}{81}<\frac{4\sqrt{73}}{81}
\]
es una garantía conservadora válida, pero no es el radio óptimo del
lector canónico. Cualquier número entre cero y 73 daría de igual modo
una garantía menor. Si se impusiera \(a^2+a+1=s_\alpha\) cambiando
\(a\), cambiarían también \(\sqrt{2a}\) y \((a+1)^2\); no se
puede sustituir sólo el 73 y conservar 4 y 81 como si el operador no
hubiera cambiado.

\begin{theorem}[Iteración exacta del peso y recurrencia de la lectura]
\label{lib02:iteraciones}
Sea \(v_n=Q_{-,a}^nv_0\), \(a>0\), y
\(b=(a-1)^2\). Cuando \(v_n\ne0\),
\begin{equation}
 w_n=\frac{b^nw_0}{b^nw_0+p(a)^n(1-w_0)}.
 \label{lib02:pesos-iterados}
\end{equation}
Para \(a\ne1\) y \(0<w_0<1\),
\[
 \frac{w_n}{1-w_n}
 =\left(\frac b{p(a)}\right)^n\frac{w_0}{1-w_0},
 \quad w_n\downarrow0,\quad s_n\uparrow p(a),
\]
y
\begin{equation}
 s_{n+1}=p(a)+b-\frac{bp(a)}{s_n},\qquad
 s_{n+1}-s_n=\frac{(s_n-b)(p(a)-s_n)}{s_n}>0.
 \label{lib02:riccati}
\end{equation}
\end{theorem}
\begin{proof}
El operador y sus adjuntos preservan los dos sectores ortogonales. Por
\eqref{lib02:gram-Qmenos}, la norma cuadrada fija se multiplica por
\(b\) en cada paso y la complementaria por \(p(a)\). Esto da
\eqref{lib02:pesos-iterados} y la razón de pesos. Como
\(p(a)-b=3a>0\), la monotonía y el límite son estrictos en el caso
mixto. Finalmente,
\(s_n=\|v_{n+1}\|^2/\|v_n\|^2\); la sucesión de normas es
una combinación de \(b^n\) y \(p(a)^n\), que satisface
\(a_{n+2}=(p(a)+b)a_{n+1}-bp(a)a_n\). Dividir por
\(a_{n+1}\) da la recurrencia y factorizar su diferencia da
\eqref{lib02:riccati}.
\end{proof}

Si \(a=1\), el sector fijo se anula en el primer paso. Cuando el
complemento no es nulo, la lectura posterior vale exactamente tres;
un estado puramente fijo se transforma en cero y su cociente deja de
estar definido. Para \(a\ne1\), la iteración inversa satisface
\(s_{n-1}=bp(a)/(p(a)+b-s_n)\) y disminuye la lectura mixta.
Invertir la orientación \(U\mapsto U^{-1}\), en cambio, deja
invariada la sucesión de normas de las iteraciones hacia delante.
Las potencias de \(Q_{+,a}\) multiplican la razón de pesos por
\(((a+2)^2/p(a))^n>1\) y disminuyen \(s_{-,a}\) en estados
mixtos. Son filtros reales del estado, no cambios de carta ni una
trayectoria temporal TPK identificada por definición.

\begin{proposition}[Ensayo completo de todas las iteraciones enteras]
\label{lib02:cruces-enteros}
Para \(a=8\), ninguna iteración entera \(Q_-^nv\), \(n\geq0\),
con \(v=K,k,m_0\), satisface \(s=s_\alpha\). Los únicos cruces
del intervalo objetivo se producen entre los índices \(21,22\),
\(14,15\) y \(15,16\), respectivamente. Para \(v=m_1\), la
lectura permanece exactamente 73 en todas las iteraciones.
\end{proposition}
\begin{proof}
Definamos los extremos racionales del peso objetivo
\[
 w_L=\frac{73-h}{24}
 =\frac{2647430716199002714894527619337}
 {2400000000000000000000000000000000},
\]
\[
 w_H=\frac{73-\ell}{24}
 =\frac{441238452699833785815754603223}
 {400000000000000000000000000000000}.
\]
Para \(w_0=A/(A+B)\), la fórmula exacta es
\(w_n=A49^n/(A49^n+B73^n)\). Los datos y los cruces son
\[
\begin{array}{c|r|r|c}
 v&A&B&N\\\hline
 K&10684363&2069408&21\\
 k&3512283&8277632&14\\
 m_0&1052377&1837634&15
\end{array}
\]
y satisfacen, en cada fila,
\[
 \frac{A49^N}{A49^N+B73^N}>w_H>w_L>
 \frac{A49^{N+1}}{A49^{N+1}+B73^{N+1}}.
\]
Son desigualdades de enteros: por ejemplo, si
\(w_H=p_H/q_H\), la primera equivale a
\((q_H-p_H)A49^N>p_HB73^N\), y la última se verifica
análogamente con \(w_L\). La sustitución de las tres filas y de
los enteros escritos prueba los cruces estrictos sin un redondeo de
alfa. La monotonía estricta de \cref{lib02:iteraciones} excluye
todos los índices anteriores y posteriores; entre \(N\) y
\(N+1\) no hay otro entero. Para \(m_1\), \(w_0=0\) permanece
cero por \eqref{lib02:pesos-iterados}.
\end{proof}

Como referencia numérica, las lecturas a ambos lados de los cruces son
\[
\begin{array}{c|c|c}
 v&s_N&s_{N+1}\\\hline
 K&72.97136264842803&72.98077012438719\\
 k&72.96167997982852&72.97426483341773\\
 m_0&72.96527892867391&72.97668298511348
\end{array}
\]
La prueba cubre todos los \(n\geq0\), no sólo una muestra finita.
La orientación inversa da los mismos cruces; las potencias positivas de
\(Q_+\) y las potencias inversas de \(Q_-\) disminuyen las
lecturas iniciales, ya menores que 72, y tampoco alcanzan el objetivo
en estas tres familias. Los productos mixtos \(Q_-^nQ_+^m\)
constituyen otra clase de dos índices, con razón de pesos
\((49/73)^n(100/73)^m w_0/(1-w_0)\); no quedan excluidos por
la prueba uniparamétrica anterior ni se eligen aquí índices para ajustar
alfa.

\subsection{Reciprocidad angular y corrección regional existente}
\label{lib02:correccion-regional}

Los cambios de carta de \cref{lib02:inversion-carta} incluyen la
reciprocidad y la deformación elíptica ya construidas. En un plano
realizado por \(M_\eta\), la forma transportada es
\(M_\eta^{-T}M_\eta^{-1}\); la semejanza no cambia los
autovalores del ciclo. Mantener el producto euclídeo anterior después de
cambiar la carta definiría otro observable. La involución
\(q_+\leftrightarrow q_-\) conserva \(A\), revierte \(C^*\)
y envía \(\eta_{\rm el}\) a \(-\eta_{\rm el}\). La relación
polar de la elipse
\(\tan\phi=e^{-\eta_{\rm el}}\tan\theta\) cambia la
parametrización de la curva, no la fase espectral de un transporte
semejante. No añade una fase pequeña a \(S^4\).

La conjugación incidencial tiene otro tipo. Su ley de bloques conserva
\[
 z_m^\dagger=z_{\rho(m)}
 +100\mathbf1_{\{c_{\rho(m)}\ne0\}}-20c_{\rho(m)},
 \qquad c_m=[C_m]_{10},
\]
con los términos de frontera antes de la reducción decimal. No es en
general una isometría del vector de bloques. Aplicarla al registro
requiere su empalme con el registro de incidencias y sus extremos; no se la
reemplaza por una corrección angular inventada.

\begin{proposition}[Corrección regional del contraángulo]
\label{lib02:retorno-regional}
Las coordenadas antecedentes de acción y ángulo cumplen
\[
 A_{\rm deg}=1000\alpha_A,\qquad
 C^*_{\rm ret,deg}=2(\eta_{\rm ret}+\alpha_A),\qquad
 C^*_{\rm pre,deg}=2(H_5+\alpha_A),
\]
\[
 \eta_{\rm ret}=H_5-\frac{90}{\pi}\mathcal C_\pi(\alpha_A),
 \quad
 \mathcal C_\pi(\alpha_A)=169\alpha_A^6+
 \frac{D_A\alpha_A^7}{1-D_A\alpha_A},\qquad
 D_A=\exp(-100\pi\alpha_A/9).
\]
La diferencia en radianes es exactamente
\[
 \theta_{C,\rm pre}-\theta_{C,\rm ret}=\mathcal C_\pi(\alpha_A).
\]
\end{proposition}
\begin{proof}
La diferencia en grados es
\(2(H_5-\eta_{\rm ret})=180\mathcal C_\pi/\pi\).
Multiplicar una sola vez por \(\pi/180\) da el resultado.
\end{proof}

Aquí \(\pi\) y \(\alpha_A\) son las coordenadas ya generadas;
\(H_5\) y los restantes lectores conservan sus definiciones
antecedentes, no una recalibración. Las evaluaciones posteriores son
\[
\begin{array}{c|c}
 \text{coordenada}&\text{radianes, aproximadamente}\\\hline
 \theta_A&0.127362829013\\
 \theta_C&0.0370662264658\\
 \theta_C/6\text{ en cada pareja antipodal}&0.00617770441097\\
 \mathcal C_\pi&2.55207421805\cdot10^{-11}
\end{array}
\]
Los canales \(q_\pm\) son atenuaciones reales positivas; su exponente
no se convierte automáticamente en una fase unitaria mediante un factor
\(i\). No se introducen divisores adicionales para aproximar las
escalas del ensayo que sigue.

\subsection{Ensayo inverso angular en las doce posiciones}
\label{lib02:ensayo-angular}

Para un objetivo general \(s\in[49,81]\), la orientación principal
del sector puro se obtiene de
\[
 |8-e^{i\theta}|^2=65-16\cos\theta,
 \qquad \theta(s)=\arccos\frac{65-s}{16}\in[0,\pi].
\]
En particular, el dato posterior \(s_\alpha\) determina
\[
 \theta_\alpha=\theta(s_\alpha),\qquad
 \delta_\alpha=\frac{2\pi}{3}-\theta_\alpha.
\]
Esta fórmula resuelve un problema inverso con una salida ya producida;
no pretende producir alfa de nuevo.

\begin{lemma}[Estructura angular del complemento ternario]
\label{lib02:J}
El operador \(J=(U-U^2)/\sqrt3\) satisface
\[
 J^*=-J,\quad J^2=-P,\quad JP_0=P_0J=0,\quad JP=PJ=J,
 \qquad U=P_0-\frac12P+\frac{\sqrt3}{2}J.
\]
\end{lemma}
\begin{proof}
El adjunto intercambia \(U,U^2\). Además,
\((U-U^2)^2=U+U^2-2I=-3P\), y
\((I+U+U^2)(U-U^2)=0\). Las demás relaciones se deducen de
\(P=I-P_0\); sumar \((U+U^2)/2\) y \((U-U^2)/2\)
da la última expresión.
\end{proof}

\begin{theorem}[Familia angular y propiedades conservadas]
\label{lib02:Utheta}
Para \(\theta\in\mathbb R\), defina
\[
 U_\theta=P_0+\cos\theta\,P+\sin\theta\,J,\qquad
 Q_{-,\theta}=8I-U_\theta.
\]
Se tienen
\[
 U_\theta^*U_\theta=I,\qquad
 U_\theta U_\varphi=U_{\theta+\varphi},\qquad U_{2\pi/3}=U,
\]
\begin{equation}
 Q_{-,\theta}^*Q_{-,\theta}
 =49P_0+(65-16\cos\theta)P.
 \label{lib02:gram-angular}
\end{equation}
Los lectores
\(T_\theta=(8I+U_\theta)/9\),
\(D_\theta=\sqrt8(U_\theta-I)/9\) conservan
\(T_\theta^*T_\theta+D_\theta^*D_\theta=I\). En cambio,
\[
 U_\theta^3=I\ \Longleftrightarrow\ 3\theta\in2\pi\mathbb Z,
 \qquad
 \det Q_{-,\theta}=7^4(65-16\cos\theta)^4.
\]
\end{theorem}
\begin{proof}
Las relaciones de \cref{lib02:J} anulan los términos entre sectores.
En el complemento \(J^2=-I\), de modo que
\(U_\theta^*U_\theta=P_0+(\cos^2\theta+\sin^2\theta)P=I\).
Multiplicar dos expresiones y utilizar las fórmulas de adición da la
ley de composición. La expresión de \(U\) en el lema da el valor
ternario. Puesto que
\(U_\theta+U_\theta^*=2P_0+2\cos\theta P\), expandir el
Gram da \eqref{lib02:gram-angular}. Las expansiones
\[
 81T_\theta^*T_\theta=65I+8(U_\theta+U_\theta^*),\qquad
 81D_\theta^*D_\theta=16I-8(U_\theta+U_\theta^*)
\]
prueban el balance. La ley de composición da
\(U_\theta^3=P_0+\cos(3\theta)P+\sin(3\theta)J\), que es
la identidad exactamente bajo la condición indicada.

Para el determinante, \(Q_{-,\theta}\) actúa por 7 en las cuatro
dimensiones fijas. En el complemento, un vector unitario \(e\) y
\(Je\) forman un par ortonormal; su complemento ortogonal es estable
bajo \(J\). Repitiendo se obtienen cuatro planos, en cada uno de los
cuales el determinante es
\((8-\cos\theta)^2+\sin^2\theta=65-16\cos\theta\).
Multiplicar las contribuciones demuestra la fórmula.
\end{proof}

\begin{corollary}[Solución inversa y separación de correcciones]
\label{lib02:fase-alpha}
En \(\operatorname{im}P\), el lector
\(Q_{-,\theta_\alpha}\) tiene razón cuadrática exactamente
\(s_\alpha\). Las orientaciones \(\pm\theta_\alpha\) dan la
misma razón y
\[
 73-s_\alpha
 =8(1-\cos\delta_\alpha+\sqrt3\sin\delta_\alpha).
\]
Se cumple \(U_{\theta_\alpha}^3\ne I\), y
\(\delta_\alpha\ne\mathcal C_\pi(\alpha_A)\).
\end{corollary}
\begin{proof}
La definición de la fase resuelve el coeficiente complementario de
\eqref{lib02:gram-angular}; el coseno es par. Sustituir
\(\theta_\alpha=2\pi/3-\delta_\alpha\) da la identidad estable
para la diferencia. Como \(72<s_\alpha<73\),
\(\pi/2<\theta_\alpha<2\pi/3\) y
\(0<3\delta_\alpha<\pi/2\); por tanto \(3\theta_\alpha\)
no es múltiplo de \(2\pi\).

La separación de correcciones puede probarse sin confiar en cifras
aproximadas. Del intervalo resulta \(73-s_\alpha>0.026\).
Si \(\delta_\alpha\leq10^{-3}\), las desigualdades
\(1-\cos t\leq t^2/2\), \(\sin t\leq t\) y
\(\sqrt3<2\) darían
\[
 73-s_\alpha\leq8(10^{-6}/2+2\cdot10^{-3})=0.016004,
\]
contradicción. Así, \(\delta_\alpha>10^{-3}\). Por otro lado,
\(0<\alpha_A<10^{-2}\), \(0<D_A<1\), implican
\[
 0<\mathcal C_\pi(\alpha_A)
 <169\cdot10^{-12}+\frac{10^{-14}}{1-10^{-2}}
 <170\cdot10^{-12}<10^{-7}.
\]
Las dos cantidades son distintas.
\end{proof}

La evaluación diagnóstica da
\[
 \delta_\alpha\simeq0.00190956706826\ {\rm rad}
 \simeq0.109410133709^\circ,\qquad
 \theta_\alpha\simeq119.890589866291^\circ.
\]
La corrección regional es aproximadamente \(7.48\cdot10^7\) veces
menor. La familia inversa conserva ortogonalidad y balance, pero pierde
el orden tres en esa fase y cambia el determinante original. Su carta
real no aporta por sí sola una matriz entera ni el cociente reticular de
índice 73. El hecho de resolver el dato prescrito no demuestra
\(73=10^4\alpha_A\).

\begin{proposition}[Alcanzabilidad al conservar el estado]
\label{lib02:alcanzabilidad}
Para \(v\ne0\), con \(w=w(v)\) fijo,
\[
 s_\theta(v)=49w+(65-16\cos\theta)(1-w).
\]
Si \(w<1\), su imagen es exactamente
\([49,81-32w]\). Por tanto existe una fase que alcanza
\(s_\alpha\) si y sólo si
\[
 w\leq\frac{81-s_\alpha}{32}\simeq0.250827322099.
\]
Cuando existe, su orientación principal es
\[
 \theta_v=\arccos\left[
 \frac1{16}\left(65-\frac{s_\alpha-49w}{1-w}\right)\right].
\]
Si \(w=1\), la lectura es siempre 49 y no alcanza el objetivo.
\end{proposition}
\begin{proof}
Evaluar \eqref{lib02:gram-angular} y dividir por \(\|v\|^2\)
da la fórmula. El coseno recorre \([-1,1]\), por lo que el
coeficiente complementario recorre \([49,81]\). Su combinación con
el peso fijo da el intervalo completo, sin huecos. Resolver la igualdad
y comprobar el dominio del arcocoseno da la condición y la fase.
El caso \(w=1\) se evalúa sin dividir por \(1-w\).
\end{proof}

\begin{corollary}[Cuatro imposibilidades y dos fases posibles]
\label{lib02:seis-angulares}
Manteniendo el estado y el sector fijo, la familia angular no alcanza
\(s_\alpha\) sobre \(K,k,m_0,u\). Sí lo alcanza sobre
\(m_1\) y \(q_K\).
\end{corollary}
\begin{proof}
De la tabla exacta \eqref{lib02:tabla-normas} se obtiene
\[
\begin{array}{c|c}
 v&32b_v-9a_v\\\hline
 K&227115677/3\\
 k&2094607/4\\
 m_0&122655440/243\\
 u&37201759/20+(7859008/15)\sqrt5
\end{array}
\]
Todas las entradas son positivas. Así \(w>9/32\), y el máximo
\(81-32w\) es menor que \(72<s_\alpha\). Para \(m_1\),
\(b_{m_1}=0\) y \(a_{m_1}>0\), luego el intervalo es
\([49,81]\). Para \(q_K\),
\[
 a_{q_K}-4b_{q_K}=\frac{588839+1195712\sqrt5}{15}>0.
\]
Por tanto \(w(q_K)<1/4\) y su máximo es mayor que
\(73>s_\alpha\). En ambos casos el objetivo pertenece al
intervalo alcanzable.
\end{proof}

Las evaluaciones de los máximos y, cuando existen, las fases son
\[
\begin{array}{c|c|c}
 v&\max_\theta s_\theta(v)&\theta_v\\\hline
 K&54.1922726227&\text{no existe}\\
 k&71.4670172771&\text{no existe}\\
 m_0&69.3474270513&\text{no existe}\\
 u&66.8296175021&\text{no existe}\\
 m_1&81&119.8905898663^\circ\\
 q_K&77.3923674871&133.5296853010^\circ
\end{array}
\]
La fase del sector puro no sirve indiscriminadamente para un estado
mixto: su lectura sería
\(49w+s_\alpha(1-w)=s_\alpha-(s_\alpha-49)w\).
Los lectores originales de \(V_2\) son polinomios en \(S\), como
\(P_0,P,J,U_\theta\). Por ello entrelazan también la familia
angular, y \(s_\theta(V_2v)=s_\theta(v)\). Archivar los
complementos conserva la lectura conjunta, pero no elimina las
restricciones de alcanzabilidad.

\subsection{Resultado de la composición y alcance de los ensayos}
\label{lib02:alcance}

El factor \(73\) ha quedado determinado por tres operaciones
compatibles y diferenciadas: el Gram de dos lectores nonádicos, la
semejanza del módulo de aumento ternario y los determinantes enteros de
las dos cartas conjugadas. El balance bilateral conserva el estado
completo y admite inversión explícita; la memoria completa conserva las
lecturas transportadas, mientras la selección de un canal puede cambiar
sus pesos. Los radios de corrección son uniformes en todos los registros
enteros del dominio algebraico y en toda la familia \(a>0\), con
pruebas de cortes y paridad que no dependen de un muestreo.

El ensayo inverso añade condiciones comprobables, no una calibración
oculta. La inversión de orientación y la covarianza métrica no cambian
la lectura. Las iteraciones enteras especificadas quedan resueltas para
todos sus índices por un cruce exacto y monotonía. La variación angular
produce una familia que realiza el objetivo en los estados permitidos,
pero el dato \(s_\alpha\) selecciona su fase y se pierde el índice
ternario original. Su identificación con una realización previa de
\(\Gamma_9\) exigiría una composición tipada sobre el estado
enriquecido; este ensayo no aporta tal identificación ni declara su
ausencia en todo el corpus.

Las nuevas formalizaciones reúnen operadores y preguntas antecedentes;
no se atribuye una prioridad mundial a la codificación, al uso de
vectores cíclicos o a la descomposición polar. Las cifras decimales
expuestas son evaluaciones posteriores de fórmulas exactas. Las
condiciones de no nulidad, carga, métrica, red y normalización son parte
del resultado, y sus contraejemplos permanecen junto a las pruebas.
Así se conserva la función de la constante holográfica de acoplamiento
aritmético-geométrico: un objeto producido que articula lectores, memoria
e incidencia, sin convertir cada coincidencia numérica en una nueva ley
física.

\section{Refinamiento conservativo de la unidad y memoria de acarreos}
\label{lib03:refinamiento}

La completación dimensional y el registro de las operaciones admiten una
composición aritmética explícita. La arquitectura APP--TRIT--TPK proporciona
previamente el estado enriquecido, las lecturas orientadas y las
completaciones de la unidad. Sobre esas salidas se construye aquí un
transporte afín que conserva la suma normalizada y permite recuperar, por
divisiones euclídeas, toda palabra finita de acarreos canónicos. Su límite
posee cilindros de diámetro controlado; las marcas de frontera distinguen
las historias que una evaluación escalar identifica.

El bloque central, sus concatenaciones y la partición
\(1000=729+243+27+1\) son antecedentes de esta construcción. La composición
de esos objetos con el transporte afín, la inversa por residuos y su
prolongación constituye la formalización desarrollada en esta sección.
La aplicación actúa sobre un dominio aritmético declarado. Cada realización
sobre una historia enriquecida particular conserva, además, la regla
prospectiva que produce sus acarreos.

\subsection{De la norma cuadrática a la completación cúbica}
\label{lib03:completacion}

Sean \(b\geq1\) entero y \(B=b+1\geq2\). Definimos
\begin{equation}
 p_b=b^2-b+1,\qquad q_b=B^2-p_b=3b.
 \label{lib03:pb}
\end{equation}

\begin{proposition}[Transferencia de la completación]
\label{lib03:prop-completacion}
Las identidades
\begin{equation}
 Bp_b=b^3+1,\qquad Bq_b+1=B^3-b^3
 \label{lib03:cuadratica-cubica}
\end{equation}
determinan la transformación
\begin{equation}
 (p_b,q_b)\longmapsto(Bp_b-1,Bq_b+1)
       =(b^3,B^3-b^3).
 \label{lib03:transferencia-cubica}
\end{equation}
La suma pasa de \(B^2\) a \(B^3\), mientras que la suma de las
coordenadas divididas por sus respectivas escalas permanece igual a uno.
\end{proposition}

\begin{proof}
La factorización \(b^3+1=(b+1)(b^2-b+1)\) proporciona la primera
identidad. La segunda resulta de
\((b+1)^3-b^3=3b^2+3b+1\). Sumarlas prueba el cambio de escala y,
después de dividir por \(B^3\), la conservación de la unidad normalizada.
\end{proof}

Para \(b=9\),
\begin{equation}
 (73,27)\longmapsto(10\cdot73-1,10\cdot27+1)=(729,271).
 \label{lib03:caso-729}
\end{equation}
El complemento conserva la descomposición \(271=243+27+1\). En la
completación precedente, la última contribución corresponde al ápice
añadido. La corrección \((-1,+1)\) tiene suma cero, y la multiplicación
por diez transforma la escala de cien en la de mil. Para la familia
completa, el transporte de las proporciones es
\begin{equation}
 \left(\frac{p_b}{B^2},\frac{q_b}{B^2}\right)
 \longmapsto
 \left(\frac{p_b}{B^2}-\frac1{B^3},
       \frac{q_b}{B^2}+\frac1{B^3}\right).
 \label{lib03:fracciones-completacion}
\end{equation}
La cantidad conservada en esta identidad es una suma lineal. Las formas
cuadráticas y las magnitudes físicas poseen sus lectores propios, que se
compondrán posteriormente con este transporte.

\subsection{Concatenaciones del bloque central y restitución unitaria}
\label{lib03:centro}

Para \(b\geq2\), las cifras \(b-2\) y \(2\) pertenecen al alfabeto de
base \(B=b+1\). Consideremos la familia algebraica
\begin{equation}
 B_c^{(b)}=\begin{pmatrix}b-2&2\\2&b-2\end{pmatrix},\qquad
 C_B(u,v)=Bu+v.
 \label{lib03:bloque-central}
\end{equation}
El caso \(b=9\) recupera el bloque central documentado
\(\bigl(\begin{smallmatrix}7&2\\2&7\end{smallmatrix}\bigr)\), con
autovalores \(9\) y \(5\). Sus concatenaciones satisfacen
\(72+27=77+22=99\). Concatenación, espectro y posición central son
lecturas distintas del mismo antecedente y mantienen sus dominios.

La aplicación del lector a las filas de \eqref{lib03:bloque-central} da
\begin{align}
 C_B(b-2,2)&=B(b-2)+2=b(b-1)=p_b-1,\\
 C_B(2,b-2)&=2B+(b-2)=3b=q_b.
 \label{lib03:concatenaciones}
\end{align}
La suma es \(B^2-1\). La restitución de una unidad en la primera
coordenada y la transferencia siguiente producen
\begin{equation}
 (p_b-1,3b)\xrightarrow{\ +(1,0)\ }(p_b,3b)
 \xrightarrow{\ L_1\ }(b^3,B^3-b^3).
 \label{lib03:cadena-restauracion}
\end{equation}
En particular,
\begin{equation}
 (72,27)\longmapsto(73,27)\longmapsto(729,271).
 \label{lib03:cadena-central}
\end{equation}
La primera flecha incrementa la suma en uno. La segunda cambia la escala
y redistribuye una unidad con corrección de suma cero. Para representar
explícitamente la contribución previa a su incorporación, el triple
\((p_b-1,3b,1)\) ya tiene suma \(B^2\), y la aplicación
\((x,y,r)\mapsto(x+r,y,0)\) conserva esa suma. Su inversión requiere
conservar el dato \(r=1\): sobre triples arbitrarios la aplicación no
es inyectiva.

Así queda registrada la restitución sin atribuirle una conservación
cerrada de la pareja antes de incorporar la unidad. La familia en \(b\)
extiende algebraicamente el caso central recibido; la selección de los
estados del atlas continúa siendo la establecida por APP--TRIT--TPK.
La composición \eqref{lib03:cadena-restauracion} determina un acarreo
igual a uno en esa transición, conservando abierta la especificación
prospectiva de las operaciones siguientes.

\subsection{Transporte afín y dominio de positividad}
\label{lib03:afin}

\begin{definition}[Transporte con acarreo]
\label{lib03:def-transporte}
Para una base entera \(B\geq2\) y \(c\in\mathbb Z\), se define
\begin{equation}
 L_c(x,y)=(Bx-c,By+c).
 \label{lib03:lc}
\end{equation}
El entero \(c\) especifica la transferencia orientada entre las dos
coordenadas. Su signo positivo reduce la primera respecto de \(Bx\)
y aumenta la segunda en la misma cantidad.
\end{definition}

\begin{proposition}[Inversa y positividad]
\label{lib03:prop-inversa}
En la realización afín posterior \(\mathbb R^2\), con \(B,c\) fijados,
\begin{equation}
 L_c^{-1}(X,Y)=\left(\frac{X+c}{B},\frac{Y-c}{B}\right),\qquad
 X+Y=B(x+y).
 \label{lib03:inversa-afin}
\end{equation}
La preimagen es entera exactamente cuando \(X+c\) e \(Y-c\) son
divisibles por \(B\). Para \(x,y\geq0\), la imagen es no negativa
exactamente cuando
\begin{equation}
 -By\leq c\leq Bx.
 \label{lib03:positividad}
\end{equation}
\end{proposition}

\begin{proof}
La sustitución directa demuestra las dos identidades de
\eqref{lib03:inversa-afin} y su criterio de integralidad. Las
desigualdades \(Bx-c\geq0\), \(By+c\geq0\) equivalen a
\eqref{lib03:positividad}.
\end{proof}

Para el alfabeto canónico \(c\in\{0,\ldots,B-1\}\), todas las
prolongaciones son admisibles si \(x\geq1\), \(y\geq0\). En la
frontera \(x=0\), la no negatividad admite únicamente \(c=0\).
El dominio positivo y el dominio entero general quedan así
diferenciados por una condición explícita.

\Needspace{12\baselineskip}
\subsection{Recuperación canónica mediante el residuo}
\label{lib03:residuo}

\begin{theorem}[Inversa canónica por residuo]
\label{lib03:teo-residuo}
Sean \(M\geq1\) entero y \(X,Y\geq0\) enteros tales que
\(X+Y=BM\). Existe una única terna admisible \((x,y,c)\) que satisface
\[
 x+y=M,\qquad 0\leq c\leq B-1,\qquad L_c(x,y)=(X,Y).
\]
Se obtiene mediante
\begin{equation}
 c=Y\bmod B=(-X)\bmod B,\qquad
 y=\frac{Y-c}{B},\qquad x=\frac{X+c}{B}.
 \label{lib03:residuo-inversor}
\end{equation}
\end{theorem}

\begin{proof}
Toda preimagen satisface \(Y\equiv c\pmod B\). El representante en
\(\{0,\ldots,B-1\}\) es único. La congruencia \(X+Y\equiv0\pmod B\)
implica la divisibilidad de \(X+c\). La división euclídea de \(Y\)
produce \(y\geq0\), y \(X+c\geq0\) produce \(x\geq0\). Si
\(c>0\), el caso \(x=0\) exigiría \(X=-c<0\), de modo que la
terna obtenida pertenece al dominio admisible. Su suma es \(M\), y
la sustitución en \eqref{lib03:lc} recupera \((X,Y)\). La unicidad
del residuo y de los dos cocientes prueba la unicidad de la terna.
\end{proof}

La pareja entera y su escala contienen el último acarreo canónico. La
recuperación deriva de su residuo y no requiere conservar otra copia
independiente de ese acarreo. Hay \(M+1\) ternas con \(c=0\) y
\(M\) ternas en cada una de las \(B-1\) clases restantes: el total
es \(BM+1\), igual al número de parejas no negativas de suma \(BM\).

Para un acarreo entero general, escribamos
\(c=r+B\ell\), con \(0\leq r\leq B-1\). El residuo determina
\(r\), mientras que la reconstrucción completa utiliza también la
elevación \(\ell\):
\begin{equation}
 y=\frac{Y-r}{B}-\ell,\qquad
 x=\frac{X+r}{B}+\ell.
 \label{lib03:elevacion}
\end{equation}
La conservación de la elevación distingue el acarreo firmado de su
reducción residual.

\subsection{Iteración variable y memoria finita completa}
\label{lib03:memoria}

Fijemos \((x_0,y_0)\in\mathbb Z^2\), de suma \(M_0>0\), y una
secuencia de operaciones \(c_1,c_2,\ldots\). Definimos
\begin{equation}
 (x_n,y_n)=L_{c_n}\cdots L_{c_1}(x_0,y_0),\qquad
 C_n=\sum_{j=1}^n c_jB^{n-j},\qquad C_0=0.
 \label{lib03:registro-acumulado}
\end{equation}
Los acarreos se reciben de la regla de operación correspondiente; el
valor límite no interviene en su selección.

\begin{theorem}[Registro acumulado exacto]
\label{lib03:teo-acumulacion}
Para todo \(n\geq0\),
\begin{equation}
 \begin{aligned}
 x_n&=B^nx_0-C_n,& y_n&=B^ny_0+C_n,\\
 x_n+y_n&=B^nM_0,& C_{n+1}&=BC_n+c_{n+1}.
 \end{aligned}
 \label{lib03:iteracion-exacta}
\end{equation}
\end{theorem}

\begin{proof}
Para \(n=0\) las identidades son inmediatas. Si son válidas en
\(n\), la aplicación de \(L_{c_{n+1}}\) multiplica por \(B\) todas
las contribuciones anteriores y añade \(c_{n+1}\) al acumulado.
Se obtiene la recurrencia y, por sustitución, ambas coordenadas de la
etapa siguiente. La suma elimina \(C_n\).
\end{proof}

Cuando los acarreos son canónicos,
\begin{equation}
 0\leq C_n\leq B^n-1.
 \label{lib03:intervalo-acumulado}
\end{equation}
En particular, \(x_0\geq1\), \(y_0\geq0\) garantizan
\(x_n\geq B^n(x_0-1)+1\) e \(y_n\geq0\) en todas las etapas.

\begin{corollary}[Recuperación de la palabra completa]
\label{lib03:cor-palabra}
Conocidos \(B,n\) y la pareja final, una historia canónica finita
recupera su entrada y todos sus acarreos mediante
\begin{align}
 C_n&=y_n\bmod B^n,&
 y_0&=\left\lfloor\frac{y_n}{B^n}\right\rfloor,&
 x_0&=\frac{x_n+C_n}{B^n},
 \label{lib03:recuperacion-inicial}\\
 c_j&=\left\lfloor\frac{C_n}{B^{n-j}}\right\rfloor\bmod B,
 &&1\leq j\leq n.
 \label{lib03:recuperacion-digitos}
\end{align}
\end{corollary}

\begin{proof}
La cota \eqref{lib03:intervalo-acumulado} hace de
\(y_n=B^ny_0+C_n\) una división euclídea. Su cociente y residuo
proporcionan \(y_0,C_n\); la primera coordenada da \(x_0\).
Las divisiones sucesivas de \(C_n\) por potencias de \(B\) recuperan
los dígitos. Equivalentemente se aplica
\eqref{lib03:residuo-inversor} desde el último paso hasta el primero.
\end{proof}

La profundidad conserva los ceros iniciales de la palabra, aunque estos
no alteren el entero \(C_n\). Si se guarda \(M_0\) en lugar de
\(n\), la igualdad exacta \((x_n+y_n)/M_0=B^n\) recupera la
profundidad. Normalizar la pareja para que sume uno elimina ese contador
de escala cuando no se conserva por otra lectura.

\subsection{Composición de bloques y compatibilidad con la partición}
\label{lib03:composicion}

Un bloque de \(n\) operaciones, con acumulado \(C\), actúa por
\begin{equation}
 L_C^{[n]}(x,y)=(B^nx-C,B^ny+C).
 \label{lib03:bloque}
\end{equation}
La composición de dos bloques satisface
\begin{equation}
 L_D^{[m]}\circ L_C^{[n]}=L_{B^mC+D}^{[n+m]}.
 \label{lib03:ley-bloques}
\end{equation}
En efecto, sustituir la primera imagen en la segunda multiplica su
acumulado por \(B^m\) y añade \(D\). La ley
\[
 (n,C),(m,D)\longmapsto(n+m,B^mC+D)
\]
es asociativa por composición de aplicaciones. Expresa la contribución
del primer bloque en la escala del segundo y mantiene la salida al
reagrupar los mismos eventos ordenados. Agregar eventos nuevos modifica
el acumulado según esta misma ley.

\subsection{Convergencia y control cuantitativo de la lectura}
\label{lib03:limite}

Escribamos \(M_n=B^nM_0\) y
\begin{equation}
 f_n=\frac{x_n}{M_n},\quad g_n=\frac{y_n}{M_n},\quad
 \theta_n=\frac{C_n}{B^n}=\sum_{j=1}^n\frac{c_j}{B^j}.
 \label{lib03:lecturas-normalizadas}
\end{equation}
Las identidades exactas son
\begin{equation}
 f_n=\frac{x_0}{M_0}-\frac{\theta_n}{M_0},\qquad
 g_n=\frac{y_0}{M_0}+\frac{\theta_n}{M_0},\qquad f_n+g_n=1.
 \label{lib03:unidad-normalizada}
\end{equation}

\begin{theorem}[Convergencia de acarreos acotados]
\label{lib03:teo-limite}
Si \(|c_j|\leq C\) para todo \(j\), la serie
\(\theta=\sum_{j\geq1}c_jB^{-j}\) converge absolutamente. Las
lecturas \(f_n,g_n\) convergen a \(f,g\), con \(f+g=1\), y
\begin{equation}
 |f-f_n|=|g-g_n|\leq\frac{C}{M_0(B-1)B^n}.
 \label{lib03:cota-cola}
\end{equation}
\end{theorem}

\begin{proof}
La cola absoluta satisface
\[
 \sum_{j>n}\frac{|c_j|}{B^j}
 \leq C\sum_{j>n}B^{-j}=\frac{C}{(B-1)B^n}.
\]
La convergencia de la serie geométrica y
\eqref{lib03:unidad-normalizada} prueban todas las afirmaciones.
\end{proof}

Para acarreos de ambos signos, una condición suficiente de positividad
en todas las etapas es \(x_0,y_0\geq C/(B-1)\). Para el alfabeto
canónico basta \(x_0\geq1,y_0\geq0\). En este segundo caso,
\(0\leq\theta\leq1\), \(f_n\) decrece y \(g_n\) crece, y
\begin{equation}
 0\leq f_n-f=g-g_n\leq\frac1{M_0B^n}.
 \label{lib03:cota-canonica}
\end{equation}
La suma permanece fija y las proporciones individuales evolucionan.
La historia finita sigue siendo recuperable cuando se conservan las
coordenadas enteras y la escala.

\subsection{Cilindros compatibles y multiplicidad de frontera}
\label{lib03:cilindros}

Consideremos el alfabeto canónico completo y una entrada
\(x_0\geq1,y_0\geq0\). Un prefijo de profundidad \(n\), con
acumulado \(C_n\), admite exactamente los valores límite
\begin{equation}
 I_n(C_n)=\left[\frac{C_n}{B^n},\frac{C_n+1}{B^n}\right]
 \label{lib03:cilindro}
\end{equation}
de la lectura \(\theta\).

\begin{theorem}[Cilindros anidados y distinción de historias]
\label{lib03:teo-cilindros}
Los \(B\) cilindros hijos de \(I_n(C_n)\) cubren el padre y se
solapan únicamente en extremos. Los cilindros de una historia infinita
tienen diámetro \(B^{-n}\) y una intersección de un solo punto.
Dos historias canónicas diferentes producen el mismo valor si y sólo
si, en su primera posición distinta, sus dígitos difieren en uno y las
colas posteriores son, respectivamente, todo \(B-1\) para el dígito
menor y todo cero para el mayor.
\end{theorem}

\begin{proof}
El hijo asociado a \(c\) tiene extremos
\((BC_n+c)/B^{n+1}\) y \((BC_n+c+1)/B^{n+1}\).
Estas expresiones demuestran inclusión, cobertura y solapamiento sólo
en extremos. La cola de una expansión recorre el intervalo
\([0,B^{-n}]\); la elección sucesiva de uno de los hijos produce una
expansión de cada punto de ese intervalo. Los intervalos cerrados
anidados tienen diámetros que tienden a cero, lo que prueba existencia
y unicidad de la intersección para cada historia.

Sean \(k\) la primera posición distinta y \(d_k>c_k\). Si las
lecturas coinciden,
\[
 \frac{d_k-c_k}{B^k}
 =\sum_{j>k}\frac{c_j-d_j}{B^j}
 \leq\sum_{j>k}\frac{B-1}{B^j}=\frac1{B^k}.
\]
La diferencia positiva entera \(d_k-c_k\) debe ser uno. La igualdad
en la cota de la cola exige \(c_j=B-1\), \(d_j=0\) para todo
\(j>k\). A la inversa, esas colas tienen diferencia exactamente
\(B^{-k}\), que compensa el primer dígito. Queda probada la
caracterización completa.
\end{proof}

La imagen del cilindro para \(f\) tiene diámetro
\(1/(M_0B^n)\), coincidente con \eqref{lib03:cota-canonica}.
Las historias \((1,0,0,\ldots)\) y
\((0,B-1,B-1,\ldots)\) producen ambas \(\theta=1/B\), conservando
prefijos y estados enteros finitos distintos. Una marca de frontera
mantiene ambas procedencias. Una convención que excluya colas finalmente
constantes \(B-1\) selecciona una expansión para \(\theta<1\);
el extremo \(\theta=1\) se conserva por separado, pues en el
alfabeto fraccionario corresponde a la cola constante \(B-1\).

Historia de operaciones, sistema de cilindros y valor límite son
objetos relacionados mediante estas aplicaciones. La unicidad de la
lectura de una historia coexiste con la multiplicidad de historias en
los valores de frontera.

\subsection{Acarreos firmados y especificación prospectiva}
\label{lib03:firmados}

Permitir acarreos fuera del alfabeto canónico exige conservar sus
elevaciones. Ya en dos pasos, las palabras
\[
 (c_1,c_2)=(1,-1),\qquad(d_1,d_2)=(0,B-1)
\]
producen \(C_2=B-1\) y la misma pareja final. Ambas son acotadas por
\(B-1\) y pueden preservar la positividad, por ejemplo desde
\((73,27)\) en base diez. La palabra firmada, o las elevaciones de
\eqref{lib03:elevacion}, distingue esa información. El teorema de
convergencia permanece válido para ambas palabras; la inversa canónica
tiene, en cambio, el dominio indicado en su enunciado.

La transición \eqref{lib03:caso-729} fija \(c_1=1\). Una regla
adicional que imponga \(c_j=1\) para todo \(j\) produce
\begin{equation}
 \theta_n=\frac{1-B^{-n}}{B-1},\qquad
 \theta=\frac1{B-1}.
 \label{lib03:acarreo-constante}
\end{equation}
En base diez, desde \((73,27)\), la secuencia es
\[
 (73,27)\longmapsto(729,271)\longmapsto(7289,2711)
 \longmapsto(72889,27111),
\]
y la primera fracción converge a
\begin{equation}
 \frac{73}{100}-\frac1{900}=\frac{164}{225}
 =0{,}728888\ldots.
 \label{lib03:limite-constante}
\end{equation}
La conclusión pertenece a esa regla adicional de acarreos constantes.
Con sólo \(c_1=1\) y prolongaciones canónicas no negativas, el cilindro
del primer paso proporciona exactamente
\begin{equation}
 \frac{b^3-1}{B^3}\leq f\leq\frac{b^3}{B^3}.
 \label{lib03:primer-cilindro}
\end{equation}
En base diez es \([0{,}728,0{,}729]\). Cambiar el signo o el alfabeto
cambia esta familia y requiere conservar la memoria adicional
correspondiente. La presencia de \(729\) en la completación determina
la transición aritmética escrita; una identificación con otro lector de
constantes utiliza además su regla efectiva de generación y su
entrelazamiento.

\subsection{Contenido operativo y transición a las amplitudes}
\label{lib03:operativo}

El procedimiento resultante comprende actualización mediante
\eqref{lib03:lc}, comprobación local por residuo, reconstrucción de
toda palabra canónica mediante \eqref{lib03:recuperacion-digitos},
control de la lectura límite mediante \eqref{lib03:cota-cola} y
reagrupación de eventos mediante \eqref{lib03:ley-bloques}. Cada
operación mantiene el orden y el contador de escala que su inversa
utiliza.

La conservación de suma y la conservación cuadrática poseen formas
diferentes. Para una transferencia normalizada
\(\delta=c/(BM)\),
\begin{equation}
 (f-\delta)^2+(g+\delta)^2-(f^2+g^2)
 =2\delta(g-f)+2\delta^2.
 \label{lib03:variacion-cuadratica}
\end{equation}
Esta expresión determina la variación cuadrática y generalmente es
distinta de cero. El enlace con una isometría se obtendrá sobre las
amplitudes de las etiquetas admisibles, mediante la biyección por
residuos de la \cref{lib05:levantamiento}. Así, el transporte aritmético
y su realización unitaria se componen conservando sus respectivas
cantidades, en lugar de identificarlas por una coincidencia escalar.

\section{Conservación bilateral y transducción aritmético-geométrica}
\label{lib04:bilateral}

La constante holográfica de acoplamiento aritmético-geométrico designa
el registro generado \(K\), con sus ventanas ordenadas, origen de fase,
orientación e incidencias. Su producción por el registro firmado y la
inversión de Hadamard está desarrollada en la \cref{sec:k-registro}.
La lectura racional \(\kappa\), el marco orbital \(O_K\) y los
operadores que transportan sus realizaciones son objetos diferentes.
Esta distinción permite componer sus funciones: generación del registro,
codificación aritmética, transporte de información y recuperación de las
lecturas geométricas.

Las coordenadas regionales y \(\alpha\) se reciben de la genealogía
APP--TRIT--TPK, del mismo estado enriquecido y de su estructura discreta
conjunta del continuo. Las construcciones siguientes actúan sobre esas
realizaciones posteriores. Su contenido adicional es una conservación
bilateral independiente de la dimensión, una extensión reversible que
admite memoria entrante y una política de archivo con recuperación
uniformemente controlada.

\subsection{Ley bilateral para dos transportes isométricos}
\label{lib04:ley}

Sean \(H,H'\) espacios de Hilbert y
\(\mathcal B,\mathcal C:H\to H'\) isometrías lineales, de modo que
\(\mathcal B^*\mathcal B=\mathcal C^*\mathcal C=I_H\).
La tipografía caligráfica distingue estos transportes de la base
aritmética del refinamiento. Para \(a>0\), definimos
\begin{equation}
 \begin{aligned}
 p&=a^2+a+1=a(a+1)+1,\\
 N&=(2a+1)p=(a+1)^3+a^3,\\
 Q_-&=a\mathcal B-\mathcal C,\qquad
 Q_+=(a+1)\mathcal B+\mathcal C.
 \end{aligned}
 \label{lib04:def-q}
\end{equation}

\begin{theorem}[Balance de dos realizaciones isométricas]
\label{lib04:teo-balance}
Para las aplicaciones \eqref{lib04:def-q},
\begin{equation}
 (a+1)Q_-^*Q_-+aQ_+^*Q_+=NI_H.
 \label{lib04:balance}
\end{equation}
La identidad vale en toda dimensión, incluidos espacios de Hilbert
infinitos, y admite transportes rectangulares con dominio y codominio
comunes.
\end{theorem}

\begin{proof}
Escribamos \(X=\mathcal B^*\mathcal C+\mathcal C^*\mathcal B\).
La isometría de ambos transportes da
\begin{align*}
 Q_-^*Q_-&=(a^2+1)I-aX,\\
 Q_+^*Q_+&=((a+1)^2+1)I+(a+1)X.
\end{align*}
Las ponderaciones \(a+1\) y \(a\) cancelan exactamente los términos
de \(X\). El coeficiente restante es
\[
 (a+1)(a^2+1)+a((a+1)^2+1)
 =2a^3+3a^2+3a+1=N.
\]
Las composiciones son de operadores acotados y el cálculo utiliza
únicamente sus productos y adjuntos. Por ello la prueba se aplica
también en dimensión infinita.
\end{proof}

Para \(a=8\), recibido de la partición nonádica precedente,
\begin{equation}
 9\|Q_-v\|^2+8\|Q_+v\|^2
 =(9^3+8^3)\|v\|^2=17\cdot73\,\|v\|^2.
 \label{lib04:balance-canonico}
\end{equation}
El factor \(73=8\cdot9+1\) normaliza esta identidad general. En el
módulo ternario se realiza además como determinante y factor de
semejanza; cada una de esas funciones conserva sus hipótesis propias.
El balance \eqref{lib04:balance} utiliza las isometrías y admite tanto
transportes con orden finito como transportes sin esa propiedad.

\subsection{Factorización, recuperación y compatibilidad de los canales}
\label{lib04:factorizacion}

Definamos
\begin{equation}
 E_av=\frac1{\sqrt p}
 \begin{pmatrix}\sqrt{a(a+1)}\,\mathcal Bv\\\mathcal Cv\end{pmatrix},
 \qquad
 R_a=\frac1{\sqrt{2a+1}}
 \begin{pmatrix}\sqrt a&-\sqrt{a+1}\\
                 \sqrt{a+1}&\sqrt a\end{pmatrix}.
 \label{lib04:factorizacion-er}
\end{equation}
La matriz \(R_a\) actúa en los dos canales y se entiende tensorizada
con \(I_{H'}\). Se comprueba
\(E_a^*E_a=I_H\), \(R_a^*R_a=I_{H'\oplus H'}\), y
\begin{equation}
 W_av=R_aE_av
 =\frac1{\sqrt N}
 \begin{pmatrix}\sqrt{a+1}\,Q_-v\\\sqrt a\,Q_+v\end{pmatrix},
 \qquad W_a^*W_a=I_H.
 \label{lib04:w}
\end{equation}
En el caso \(a=8\), el primer reparto tiene pesos \(72/73\) y
\(1/73\); la segunda operación es una rotación entre los canales.
Esta factorización expresa el papel de \(72+1\) dentro de la norma
total, distinguiéndolo de la restitución lineal de
\eqref{lib03:cadena-restauracion}.

\begin{proposition}[Recuperación bilateral]
\label{lib04:prop-recuperacion}
Los canales sin normalizar \(x=Q_-v\), \(y=Q_+v\) determinan
\begin{equation}
 \mathcal Bv=\frac{x+y}{2a+1},\qquad
 \mathcal Cv=\frac{ay-(a+1)x}{2a+1},\qquad
 v=\mathcal B^*\frac{x+y}{2a+1}.
 \label{lib04:inversa-canales}
\end{equation}
Desde los canales normalizados se recupera
\(v=W_a^*W_av\), con \(\|W_a^*\|=1\) cuando \(H\neq\{0\}\).
\end{proposition}

\begin{proof}
Sumar las dos definiciones de \(Q_\pm\) elimina \(\mathcal Cv\).
La combinación \(aQ_+-(a+1)Q_-\) elimina \(\mathcal Bv\).
La primera isometría tiene inversa izquierda \(\mathcal B^*\), lo
que demuestra las fórmulas. La última afirmación se sigue de
\eqref{lib04:w}: el adjunto de una isometría tiene norma uno.
\end{proof}

La imagen \(\operatorname{im}W_a\) es cerrada y su proyector
ortogonal es \(W_aW_a^*\). Un par normalizado es compatible con una
entrada única exactamente cuando pertenece a esa imagen; su componente
ortogonal cuantifica el defecto de compatibilidad.

Si \(\mathcal B=I\) y \(\mathcal C=U\) es unitario, los dos
\(Q_\pm\) son polinomios del mismo operador. Para los canales sin
normalizar, el criterio se simplifica a
\begin{equation}
 Q_+x=Q_-y.
 \label{lib04:compatibilidad-unitaria}
\end{equation}
En efecto, la igualdad equivale a
\(U(x+y)=ay-(a+1)x\). Definiendo
\(v=(x+y)/(2a+1)\), se obtiene por sustitución
\(Q_-v=x\), \(Q_+v=y\). La necesidad resulta de la conmutación
de los dos polinomios. Esta forma simplificada pertenece al dominio
unitario indicado; el criterio de imagen cubre los transportes
generales entre espacios distintos.

\subsection{Transformación de lectores y simetría relativa}
\label{lib04:lectores}

\begin{theorem}[Lector común en ambos canales]
\label{lib04:teo-lector}
Para todo operador acotado autoadjunto \(F\) sobre \(H'\),
\begin{equation}
 W_a^*\operatorname{diag}(F,F)W_a
 =\frac{a(a+1)\mathcal B^*F\mathcal B+
                  \mathcal C^*F\mathcal C}{a(a+1)+1}.
 \label{lib04:lector-general}
\end{equation}
Si \(\mathcal B\) es unitario y \(\mathcal C=\mathcal BR\), con
\(R\) unitario, la expresión, para \(G=\mathcal B^*F\mathcal B\),
es
\begin{equation}
 \mathcal T_a(G)=\frac{a(a+1)G+R^*GR}{a(a+1)+1}.
 \label{lib04:lector-relativo}
\end{equation}
Se cumple \(\mathcal T_a(G)=G\) si y sólo si \(GR=RG\).
\end{theorem}

\begin{proof}
La expansión de \((a+1)Q_-^*FQ_-+aQ_+^*FQ_+\) cancela los
términos \(\mathcal B^*F\mathcal C+\mathcal C^*F\mathcal B\).
Los coeficientes de los términos restantes son
\(a(a+1)(2a+1)\) para \(\mathcal B^*F\mathcal B\) y \(2a+1\)
para \(\mathcal C^*F\mathcal C\). Dividir por \(N\) demuestra
\eqref{lib04:lector-general}. Sustituir \(\mathcal C=\mathcal BR\)
produce \eqref{lib04:lector-relativo}. La igualdad con \(G\) equivale
a \(R^*GR=G\), y la unitariedad de \(R\) la hace equivalente a
\(GR=RG\).
\end{proof}

El lector \(F=I\) recupera la norma. Para un lector general, la
identidad determina su transformación, y la coincidencia
\(\mathcal B^*F\mathcal B=\mathcal C^*F\mathcal C=G\) garantiza
su conservación en todos los estados. En el caso canónico,
\begin{equation}
 \mathcal T_8(G)=\frac{72G+R^*GR}{73}.
 \label{lib04:lector-72}
\end{equation}
Esta fórmula relaciona conservación y simetría mediante un conmutador
concreto. El observable que transporta exactamente un operador
\(A\) del dominio a la imagen completa es \(W_aAW_a^*\); su
identidad ambiente es \(W_aW_a^*\). El observable diagonal de
\eqref{lib04:lector-general} es otra elección, cuya diferencia queda
calculada por el mismo teorema.

\subsection{Actualización reversible con memoria entrante}
\label{lib04:reentrada}

Sean ahora \(\mathcal B=I\), \(\mathcal C=U\), con \(U\)
unitario en un mismo espacio, y escribamos
\begin{equation}
 A=\frac{\sqrt{a+1}}{\sqrt N}Q_-,\qquad
 D=\frac{\sqrt a}{\sqrt N}Q_+.
 \label{lib04:ad}
\end{equation}
Los operadores \(A,D\) son normales y conmutan entre sí y con sus
adjuntos, por ser polinomios de \(U\) y \(U^*\).

\begin{theorem}[Extensión unitaria]
\label{lib04:teo-unitaria}
La actualización
\begin{equation}
 \mathscr R_a=
 \begin{pmatrix}A&-D^*\\D&A^*\end{pmatrix}:H\oplus H\to H\oplus H
 \label{lib04:unitaria}
\end{equation}
es unitaria. Su primera columna es \(W_a\), y una entrada
\((v,m)\) conserva \(\|v\|^2+\|m\|^2\).
\end{theorem}

\begin{proof}
El balance da \(A^*A+D^*D=I\). La normalidad proporciona
\(AA^*+DD^*=I\). Los productos cruzados de
\(\mathscr R_a^*\mathscr R_a\) y
\(\mathscr R_a\mathscr R_a^*\) se anulan por las conmutaciones
anteriores; sus dos diagonales son la identidad. La primera columna
reproduce \eqref{lib04:w}, y la conservación de la norma conjunta
es la identidad unitaria aplicada a \((v,m)\).
\end{proof}

Para unos datos \(z=(z_1,z_2)\), la inversión explícita es
\begin{equation}
 v_{\mathrm{rec}}=A^*z_1+D^*z_2,\qquad
 m_{\mathrm{rec}}=-Dz_1+Az_2.
 \label{lib04:inversa-unitaria}
\end{equation}
La hipótesis de memoria entrante nula se comprueba mediante
\(m_{\mathrm{rec}}=0\). Si \(e\) es una perturbación de los datos,
la unitariedad da
\begin{equation}
 \|e\|^2=\|e_{\mathrm{rec}}\|^2+\|e_{\mathrm{mem}}\|^2.
 \label{lib04:descomposicion-ruido}
\end{equation}
La segunda componente detecta la parte incompatible con esa entrada.
Un error perteneciente a \(\operatorname{im}W_a\) tiene residuo de
memoria nulo y altera la entrada recuperada. La corrección del registro
entero utiliza, además, sus distancias discretas de separación.

\subsection{Reducción causal de la actualización bilateral}
\label{lib04:reduccion-causal}

La actualización unitaria anterior realiza el transporte del estado y de su
memoria en un mismo espacio. La observación de la primera componente admite
una descripción equivalente que conserva explícitamente la contribución del
registro. Esta descripción se obtiene del transporte ya construido, con los
mismos coeficientes y la misma orientación; la memoria eliminada de la lista
de variables reaparece como dependencia de los estados anteriores.

Fijemos los operadores \(A,D\) de \eqref{lib04:ad} y una realización
\(H\oplus H\). Para \(n\geq0\), sean
\begin{equation}
 v_{n+1}=Av_n-D^*m_n,\qquad
 m_{n+1}=Dv_n+A^*m_n.
 \label{lib04:recurrencia-memoria}
\end{equation}
La actualización conjunta es \(\mathscr R_a\) y conserva
\(\|v_n\|^2+\|m_n\|^2\).

\begin{proposition}[Eliminación exacta del registro auxiliar]
\label{lib04:prop-memoria-reducida}
Para todo \(n\geq0\), con la suma vacía interpretada como cero,
\begin{align}
 m_n&=(A^*)^n m_0+
       \sum_{k=0}^{n-1}(A^*)^{n-1-k}Dv_k,
       \label{lib04:memoria-explicita}\\
 v_{n+1}&=Av_n-D^*(A^*)^n m_0
       +\sum_{k=0}^{n-1}\mathscr K_{n-1-k}v_k,
       \label{lib04:evolucion-reducida}\\
 \mathscr K_j&=-D^*(A^*)^jD\quad(j\geq0).
       \label{lib04:nucleo-retardo}
\end{align}
La ecuación reducida, junto con \eqref{lib04:memoria-explicita}, es
equivalente a la recurrencia conjunta para los mismos datos \((v_0,m_0)\).
\end{proposition}
\begin{proof}
La primera identidad vale para \(n=0\). Si vale para \(n\), la segunda
ecuación de \eqref{lib04:recurrencia-memoria} proporciona
\[
 m_{n+1}=(A^*)^{n+1}m_0+
 \sum_{k=0}^{n-1}(A^*)^{n-k}Dv_k+Dv_n,
\]
que es la identidad en \(n+1\). Sustituirla, en profundidad \(n\), en
la primera ecuación demuestra \eqref{lib04:evolucion-reducida}.
Recíprocamente, si \(v_n\) satisface esta última y se define \(m_n\)
por \eqref{lib04:memoria-explicita}, la misma diferencia de sumas prueba
\(m_{n+1}=Dv_n+A^*m_n\); la otra ecuación se recupera por sustitución.
\end{proof}

El núcleo \(\mathscr K_j\) describe una transferencia hacia el registro,
\(j\) actualizaciones dentro de él y una reentrada en la componente leída.
El término \(-D^*(A^*)^nm_0\) conserva la memoria inicial. Su supresión
corresponde al caso particular \(m_0=0\), no a una identidad del transporte
general. En \eqref{lib04:evolucion-reducida} intervienen únicamente el
estado actual y los estados anteriores. La causalidad de esta reducción
procede del orden de la recurrencia conjunta.

La normalización del balance proporciona además una estimación explícita.
Como \(Q_-=aI-U\), se tiene
\(\|A\|^2\leq r_a=(a+1)^3/((a+1)^3+a^3)<1\), mientras que
\(D^*D\leq I\). Por tanto,
\begin{equation}
 \|\mathscr K_j\|\leq r_a^{j/2},\qquad
 \sum_{j\geq0}\|\mathscr K_j\|
 \leq\frac{1}{1-\sqrt{r_a}}.
 \label{lib04:cota-nucleo-retardo}
\end{equation}
Esta cota controla el peso de la memoria en la ecuación reducida; la norma
del estado conjunto sigue conservándose por la actualización unitaria.

La proposición es la realización temporal de la reducción de Schur de
\cref{mem:prop-schur}. En las series formales
\(V(z)=\sum_{n\geq0}v_nz^n\) y
\(M(z)=\sum_{n\geq0}m_nz^n\), las recurrencias dan
\(V=v_0+zAV-zD^*M\) y \(M=m_0+zDV+zA^*M\). Eliminando \(M\),
\begin{equation}
 V(z)=\bigl[I-zA+z^2D^*(I-zA^*)^{-1}D\bigr]^{-1}
       \bigl[v_0-zD^*(I-zA^*)^{-1}m_0\bigr].
 \label{lib04:schur-bilateral}
\end{equation}
Las inversas formales existen porque sus términos constantes son identidades;
las series de los estados y el resolvente conjunto también convergen para
\(|z|<1\). La fuente inicial y el operador efectivo se transforman juntos.
Un lector que actúe también sobre \(m_n\) conserva su contribución mediante
\eqref{lib04:memoria-explicita}, conforme a \eqref{mem:schur-lector}.

La reentrada aquí descrita actualiza un registro entrante. La política de
archivo que sigue conserva, en cambio, cada complemento en una coordenada
distinta. Ambas operaciones poseen su propia ley de conservación. El núcleo
\(\mathscr K_j\), los propagadores de Green y una fuerza dimensional son
objetos de tipos diferentes: el primero elimina una componente de estado,
los segundos resuelven ecuaciones de fuente con condiciones de soporte, y
la fuerza añade una realización dinámica y sus unidades.


\subsection{Archivo no estacionario y agotamiento uniforme del terminal}
\label{lib04:archivo}

En cada etapa sean \(H_n,H_{n+1}\) espacios de Hilbert y
\(\mathcal B_n,\mathcal C_n:H_n\to H_{n+1}\) isometrías. El
parámetro \(a>0\) permanece fijo. Sean \(A_n,D_n\) los bloques
normalizados de \eqref{lib04:w}; se continúa por el primero y se
registra el segundo:
\begin{equation}
 R_0=I,\qquad R_{n+1}=A_nR_n,\qquad m_n=D_nR_nv.
 \label{lib04:politica-archivo}
\end{equation}
Esta política posee un dominio distinto de la reintroducción de
memoria de \eqref{lib04:unitaria}; cada una conserva su actualización.

\begin{lemma}[Transferencia uniforme]
\label{lib04:lem-tasa}
En toda etapa,
\begin{equation}
 \|A_n\|^2\leq r_a:={\frac{(a+1)^3}{(a+1)^3+a^3}}<1,
 \quad D_n^*D_n\geq(1-r_a)I,
 \quad\|R_n\|^2\leq r_a^n.
 \label{lib04:tasa}
\end{equation}
\end{lemma}

\begin{proof}
La desigualdad triangular da
\(\|a\mathcal B_n-\mathcal C_n\|\leq a+1\). Su normalización
proporciona la primera cota. La diferencia
\(N-(a+1)^3=a^3>0\) prueba \(r_a<1\). El balance local implica
\(D_n^*D_n=I-A_n^*A_n\), y multiplicar las cotas de los canales
continuados produce la estimación de \(R_n\).
\end{proof}

\begin{theorem}[Memoria completa a profundidad arbitraria]
\label{lib04:teo-archivo}
La aplicación finita
\begin{equation}
 Z_Nv=(R_Nv,D_0v,D_1R_1v,\ldots,D_{N-1}R_{N-1}v)
 \label{lib04:zn}
\end{equation}
es isométrica. El archivo infinito
\begin{equation}
 Z_\infty v=(D_nR_nv)_{n\geq0}
 \label{lib04:zinfty}
\end{equation}
pertenece a la suma hilbertiana de los espacios de salida y satisface
\begin{equation}
 \|v\|^2=\sum_{n\geq0}\|m_n\|^2,\qquad
 v=\sum_{n\geq0}R_n^*D_n^*m_n,
 \qquad Z_\infty^*Z_\infty=I.
 \label{lib04:recuperacion-infinita}
\end{equation}
La serie inversora converge en norma y el adjunto tiene norma uno.
\end{theorem}

\begin{proof}
La identidad \(A_n^*A_n+D_n^*D_n=I\), conjugada por \(R_n\),
produce
\[
 R_n^*R_n-R_{n+1}^*R_{n+1}=R_n^*D_n^*D_nR_n.
\]
Sumar desde cero hasta \(N-1\) da
\begin{equation}
 I=R_N^*R_N+\sum_{n<N}R_n^*D_n^*D_nR_n.
 \label{lib04:telescopia}
\end{equation}
Esta es la identidad de Gram de \(Z_N\). Por
\eqref{lib04:tasa}, el primer término tiende a cero en norma
operatoria. La identidad límite prueba la convergencia cuadrática del
archivo y su isometría. El adjunto de esa isometría es acotado; las
truncaciones de todo archivo cuadrado-sumable convergen en la suma
hilbertiana, y su imagen por el adjunto prueba la convergencia de la
serie de \eqref{lib04:recuperacion-infinita}. Aplicarla a los datos
exactos recupera \(v\).
\end{proof}

Con sólo los primeros \(N\) complementos, el adjunto truncado tiene
residuo exacto
\begin{equation}
 v-\sum_{n<N}R_n^*D_n^*m_n=R_N^*R_Nv,
 \qquad
 \left\|v-\sum_{n<N}R_n^*D_n^*m_n\right\|
 \leq r_a^N\|v\|.
 \label{lib04:sesgo-truncado}
\end{equation}
Si el error conjunto de los datos es a lo sumo \(\varepsilon\), la
cota total es \(\varepsilon+r_a^N\|v\|\). Con \(Z_N\), incluido
su terminal, o con \(Z_\infty\), el adjunto recupera con error a lo
sumo \(\varepsilon\).

La política anterior agota el terminal por la cota uniforme probada.
Las evoluciones precedentes que poseen un terminal positivo
\(A_\infty\neq0\) conservan en su archivo la componente
\(A_\infty^{1/2}v\). La profundidad arbitraria y la eliminación del
terminal son, por tanto, conclusiones diferentes y mantienen sus
hipótesis respectivas.

\begin{proposition}[Inversa de las memorias finitas sin terminal]
\label{lib04:prop-inversa-finita}
Para \(N\geq1\), sea \(M_Nv=(D_nR_nv)_{n<N}\). Entonces
\begin{equation}
 M_N^*M_N=I-R_N^*R_N\geq(1-r_a^N)I,
 \label{lib04:gram-finito}
\end{equation}
y la inversa de mínimos cuadrados
\begin{equation}
 L_N=(I-R_N^*R_N)^{-1}M_N^*,\qquad
 L_NM_N=I,\qquad
 \|L_N\|\leq(1-r_a^N)^{-1/2}
 \label{lib04:inversa-finita}
\end{equation}
recupera exactamente la entrada con datos exactos.
\end{proposition}

\begin{proof}
La identidad de Gram es \eqref{lib04:telescopia}; la cota procede de
\eqref{lib04:tasa}. El operador positivo es invertible por su cota
inferior estricta. La composición demuestra \(L_NM_N=I\), y
\(L_NL_N^*=(M_N^*M_N)^{-1}\) prueba la norma declarada.
\end{proof}

Esta inversa finita elimina el sesgo de
\eqref{lib04:sesgo-truncado}; su amplificación de ruido queda
especificada por \eqref{lib04:inversa-finita}. El adjunto truncado,
la inversa finita y el adjunto del archivo completo son así tres
operadores con comportamientos diferenciados.

Para \(a=8\), la cota general es \(r_a=729/1241\). En el caso
\(\mathcal B=I\), \(\mathcal C=S^4\), con \(S\) el ciclo de
doce posiciones, el proyector fijo
\(P_0=(I+S^4+S^8)/3\) da
\begin{equation}
 A^*A=\frac{a+1}{(2a+1)p}
       \bigl((a-1)^2P_0+p(I-P_0)\bigr).
 \label{lib04:gram-tasa-ternaria}
\end{equation}
Como \(p-(a-1)^2=3a>0\), su norma es
\((a+1)/(2a+1)\). En \(a=8\) resulta \(9/17\); el valor propio
del sector fijo es \(441/1241\). Estas tasas pertenecen a los
transportes y a la política declarados.

Si se refinan ambos canales en un árbol, cada nivel completo es también
isométrico por composición. Sumar las normas de todas las capas
completas contaría repetidamente la misma norma inicial. La conservación
se formula mediante las aplicaciones de refinamiento entre niveles o
mediante el archivo de complementos de \eqref{lib04:zinfty}.

\subsection{Formas geométricas y conservación de su transporte}
\label{lib04:formas}

Sean \(b,b'\) formas no degeneradas, simétricas o alternantes, sobre
los espacios correspondientes, y supóngase
\(\mathcal B^*b'\mathcal B=\mathcal C^*b'\mathcal C=b\).
El mismo desarrollo bilineal demuestra
\begin{equation}
 (a+1)Q_-^*b'Q_-+aQ_+^*b'Q_+=Nb.
 \label{lib04:balance-formas}
\end{equation}
En efecto, los términos mixtos
\(\mathcal B^*b'\mathcal C+\mathcal C^*b'\mathcal B\) se cancelan
con las mismas ponderaciones; los términos puros suman \(Nb\).
Así se conserva la forma transportada sin exigir que sea positiva.
Las cotas contractivas y de ruido de la subsección anterior utilizan,
en cambio, normas de Hilbert positivas. Una forma lorentziana o
alternante conserva \eqref{lib04:balance-formas}, con sus propias
interpretaciones geométricas, y no sustituye la positividad en esas
desigualdades.

\subsection{Correlación entre canales y ensayo inverso de una lectura}
\label{lib04:correlacion}

Para \(v\neq0\), definamos
\begin{equation}
 \chi(v)=\frac{\operatorname{Re}\langle\mathcal Bv,\mathcal Cv\rangle}
                      {\|v\|^2},\qquad -1\leq\chi(v)\leq1.
 \label{lib04:chi}
\end{equation}
La cota es Cauchy--Schwarz aplicada a las dos imágenes isométricas.
Las lecturas normalizadas son
\begin{equation}
 s_-=a^2+1-2a\chi,\qquad
 s_+=(a+1)^2+1+2(a+1)\chi.
 \label{lib04:s-pm}
\end{equation}
Para \(a=8\),
\begin{equation}
 s_-=65-16\chi,\qquad s_+=82+18\chi,\qquad
 9s_-+8s_+=1241.
 \label{lib04:s-canonico}
\end{equation}
Las dos normas individuales contienen una correlación real común. El
balance es una relación exacta entre ellas, no dos determinaciones
aritméticas independientes de una constante.

Usar la salida previa \(\alpha_A\) en el ensayo inverso
\(s_-=10^4\alpha_A\) determina la lectura complementaria
\begin{equation}
 s_+=73+\frac98(73-10^4\alpha_A)
       \simeq73{,}029783595557.
 \label{lib04:ensayo-alpha}
\end{equation}
La igualdad se obtiene sustituyendo en
\eqref{lib04:s-canonico}. El promedio ponderado
\((9s_-+8s_+)/17\) permanece exactamente igual a \(73\).
La comparación conserva la genealogía previa de \(\alpha_A\): es
una condición del ensayo inverso, no otro productor independiente de
esa salida. Un cambio unitario puede alterar \(\chi\) y las dos
lecturas manteniendo el balance, aunque cambie una relación particular
de orden ternario.

\subsection{Marco cíclico y transporte de las respuestas}
\label{lib04:marco}

Sobre el portador del registro generado, sean
\(H=\mathbb R^{12}\), \((Sx)_i=x_{i+1}\), con índices cíclicos,
y
\begin{equation}
 O_K=(K,SK,\ldots,S^{11}K),\qquad
 589609I\leq O_K^*O_K\leq39225169I.
 \label{lib04:gram-k}
\end{equation}
Las cotas son las del marco cíclico del registro canónico,
demostradas en la \cref{lib01:marco-ciclico}.
Sea \(Z\) una de las isometrías completas
\(Z_N,Z_\infty\), o una composición con los archivos precedentes
que conserve sus terminales y complementos.

\begin{theorem}[Conservación del condicionamiento del marco]
\label{lib04:teo-marco}
Si \(x=O_K\beta\), el dato \(z=Zx\) determina
\begin{equation}
 \beta=O_K^{-1}Z^*z,\qquad
 \|\delta\beta\|\leq\frac{\|\delta z\|}{\sqrt{589609}}.
 \label{lib04:inversa-marco}
\end{equation}
Las doce columnas de \(ZO_K\) forman una base de su imagen, incluso
cuando el archivo tenga infinitas coordenadas.
\end{theorem}

\begin{proof}
La igualdad \(Z^*Z=I\) da
\((ZO_K)^*(ZO_K)=O_K^*O_K\). Las cotas de Gram se conservan y
\(O_K\) es invertible por la cota inferior estricta. Componer sus
inversas demuestra la recuperación, y la norma
\(\|O_K^{-1}\|\leq589609^{-1/2}\), junto con
\(\|Z^*\|=1\), da la estabilidad. El número de columnas y su
independencia determinan la dimensión de la imagen, sin identificarla
con todo el espacio de almacenamiento.
\end{proof}

Para un operador \(H_0\) que conmute con \(S\), la respuesta
vectorial completa \(y=H_0K\) determina
\begin{equation}
 H_0=O_yO_K^{-1}.
 \label{lib04:identificacion-filtro}
\end{equation}
En efecto, la conmutación implica
\(H_0O_K=(y,Sy,\ldots,S^{11}y)=O_y\). Si se recibe únicamente
\(z=Zy+e\), se define \(\widehat y=Z^*z\) y
\(\widehat H_0=O_{\widehat y}O_K^{-1}\). Entonces
\begin{equation}
 \|\widehat H_0-H_0\|
 \leq\sqrt{\frac{12}{589609}}\,\|e\|.
 \label{lib04:error-filtro}
\end{equation}
La matriz orbital de \(\delta y\) tiene doce columnas de norma
\(\|\delta y\|\), por lo que su norma operatoria no supera
\(\sqrt{12}\|\delta y\|\). Esto y
\eqref{lib04:gram-k} prueban \eqref{lib04:error-filtro}.
Para un operador general el protocolo requiere las doce respuestas
\(H_0S^jK\); una respuesta en el caso cíclico significa sus doce
componentes, no una medición escalar.

En la imagen transportada, el desplazamiento es
\(S_Z=ZSZ^*\), y satisface \(S_Z^jZ=ZS^j\). Esta igualdad
determina el avance de fase en el archivo; una traslación de casillas
del espacio ambiente sólo representa ese avance cuando coincide con
el operador transportado.

\subsection{Lectura posicional y precisión escalar}
\label{lib04:kappa}

Fijemos \(B_0=1000\), \(d=B_0^{12}-1\), y la fila
\begin{equation}
 \ell=\frac{(B_0^{11},B_0^{10},\ldots,1)}{d}.
 \label{lib04:fila-posicional}
\end{equation}
Su extensión lineal actúa sobre \(H\), y en el dominio de palabras
\(\kappa(K)=\ell K\). El representante del lector transportado es
\(Z\ell^*\), por lo que
\begin{equation}
 \kappa(K)=\langle Z\ell^*,ZK\rangle,\qquad
 |\delta\kappa|\leq\|\ell\|\,\|e\|.
 \label{lib04:lectura-kappa}
\end{equation}
La suma geométrica proporciona
\begin{equation}
 \|\ell\|^2
 =\frac{B_0^{12}+1}{(B_0^2-1)(B_0^{12}-1)}.
 \label{lib04:norma-fila}
\end{equation}
En efecto, el numerador de la suma de cuadrados es
\((B_0^{24}-1)/(B_0^2-1)\), que al dividirse por
\((B_0^{12}-1)^2\) da \eqref{lib04:norma-fila}.

Con sólo \(N\) complementos y el adjunto truncado, aparece además
un error a lo sumo \(\|\ell\|r_a^N\|K\|\). La inversa finita
\eqref{lib04:inversa-finita} elimina ese sesgo y conserva su cota
explícita de amplificación de ruido.

En el dominio \(\{0,\ldots,999\}^{12}\), la lectura exacta
recupera \(N_K=d\kappa\) y sus doce bloques mediante divisiones
euclídeas sucesivas. Si
\begin{equation}
 |\widetilde\kappa-\kappa|<\frac1{2d},
 \label{lib04:radio-escalar}
\end{equation}
el redondeo de \(d\widetilde\kappa\) al entero más próximo
recupera \(N_K\): su distancia al entero verdadero es menor que
media unidad. La base, longitud, origen y orientación acompañan a
esa representación. El umbral escalar de
\eqref{lib04:radio-escalar} se distingue de las cotas vectoriales
de los registros de memoria.

\subsection{Incidencia, carga y restitución exacta del registro}
\label{lib04:incidencia}

Sea \(\mathbf1=(1,\ldots,1)^T\) y
\(\Pi=I-\mathbf1\mathbf1^*/12\). La isometría incidencial de
la \cref{nuc05:isometria} es
\begin{equation}
 Fv=(\mu(v),\mathcal I\Pi v/6),\qquad
 \mu(v)=\frac{\mathbf1^*v}{12},\qquad
 \|(\mu,y)\|_Y^2=12\mu^2+\|y\|^2.
 \label{lib04:f-incidencia}
\end{equation}
Su codominio es su imagen declarada. El Gram
\(\mathcal I^*\mathcal I=36I+30\mathbf1\mathbf1^*\) da
\(\|\mathcal I\Pi v\|^2=36\|\Pi v\|^2\), y
\(12\mu(v)^2+\|\Pi v\|^2=\|v\|^2\), de donde \(F^*F=I\).
Así, a partir de \(z=Zv+e\),
\begin{equation}
 \widehat{Fv}=FZ^*z,\qquad \|FZ^*e\|_Y\leq\|e\|.
 \label{lib04:error-incidencia}
\end{equation}
La incidencia sin normalizar \(\mathcal I\Pi v\) tiene error a
lo sumo \(6\|e\|\); el registro firmado \(H_{12}v\), cuya
normalización satisface \(H_{12}^*H_{12}=4I\), tiene error a lo sumo
\(2\|e\|\).

La dirección \(u=P_3\Pi K\) y el proyector
\(P_{10}=\Pi-uu^*/\|u\|^2\), con \(u\neq0\), se recuperan
con sus lectores y carta previos. La selección absoluta
\(\{j:K_j\geq729\}\) tiene una discontinuidad en \(K_4=729\).
La restitución entera se efectúa antes de aplicar ese umbral.

\begin{proposition}[Decodificación después del archivo completo]
\label{lib04:prop-redondeo}
Sea \(K\) entero y \(z=ZK+e\), con \(Z^*Z=I\). Si
\(\|e\|<1/2\), redondear cada coordenada de \(Z^*z\) recupera
\(K\) exactamente. Por composición se restituyen \(\kappa\), las
firmas y los lectores geométricos definidos sobre ese registro.
\end{proposition}

\begin{proof}
La contractividad de \(Z^*\) da
\(\|Z^*z-K\|\leq\|e\|<1/2\). Cada error coordenado queda
acotado por la norma conjunta, por lo que la única coordenada entera
a distancia menor que media unidad es la original. Una vez
recuperado \(K\), evaluar cualquiera de sus lectores exactos
restituye su valor.
\end{proof}

En particular, se recupera el soporte \(\{4,6,9,10\}\), incluida
su componente de frontera. Las otras marcas que intervienen en una
bandera excepcional se conservan por sus canales propios: la
restitución de \(K\) corrige exactamente la información que tiene
en él sus argumentos suficientes.

\subsection{Conservación del condicionamiento de dos memorias}
\label{lib04:dos-memorias}

Mantengamos los lectores antecedentes
\begin{equation}
 T_j=\frac{8I+S^j}{9},\qquad
 D_j=\frac{\sqrt8}{9}(S^j-I),\qquad
 Mv=(D_3v,D_4T_3v).
 \label{lib04:m-antecedente}
\end{equation}
Los \(D_j\) de esta ecuación son los complementos precedentes;
el segundo bloque bilateral \(D\) de \eqref{lib04:ad} conserva
su propia definición. Se tiene \(\ker M=\mathbb R\mathbf1\)
y una inversa \(L_M\) de norma \(9/4\) sobre el sector centrado,
como se demuestra en la \cref{nuclear:11}.

Sea \(E=W_a\oplus W_a\) en el espacio de las dos memorias.
Su isometría implica
\begin{equation}
 (EM)^*(EM)=M^*M,\qquad
 \ker(EM)=\mathbb R\mathbf1.
 \label{lib04:gram-m-preservado}
\end{equation}
Para \(z=EMv+e\) y carga exacta
\(q=\mathbf1^*v\), la reconstrucción es
\begin{equation}
 \widehat v=\frac q{12}\mathbf1+L_ME^*z,\qquad
 \|\widehat v-v\|\leq\frac94\|e\|.
 \label{lib04:inversa-m}
\end{equation}
La prueba consiste en aplicar \(E^*E=I\), usar
\(L_MM=\Pi\) y \(\|E^*\|=1\). Si la carga tiene error
\(\delta q\), los sectores uniforme y centrado son ortogonales y
\begin{equation}
 \|\widehat v-v\|^2
 \leq\frac{|\delta q|^2}{12}+\frac{81}{16}\|e\|^2.
 \label{lib04:error-carga}
\end{equation}

Las distancias de todas las palabras de la imagen entera se conservan
exactamente por \(E\). Por tanto, los radios de restitución de
la \cref{lib02:radios-memoria} permanecen
\begin{equation}
 r_0=\frac{2\sqrt{146}}{81}\quad\hbox{sin carga},\qquad
 r_q=\frac{4\sqrt{73}}{81}\quad\hbox{con carga exacta}.
 \label{lib04:radios}
\end{equation}
Dentro del primer radio se recupera \(\Pi K\), y con ello la
dirección y su proyector cuando están definidos. Dentro del segundo,
conservando \(q\), se recuperan \(K\) y el selector absoluto de
umbral. Componer después otra isometría completa, finita o infinita,
mantiene las mismas garantías: se aplica su adjunto antes del
decodificador. El radio siempre se expresa en la norma conjunta
de los datos con la normalización fijada.

Esta composición redistribuye la información de \(M\) sin deteriorar
su condicionamiento y sin crear el canal uniforme que \(M\) elimina.
La conservación de Gram distingue esta propiedad de una mejora
numérica obtenida cambiando el lector o la escala del ruido.

\subsection{Norma orbital y normalización polar del marco}
\label{lib04:normalizacion}

La inversión precedente cumple
\(K=H_{12}U_{\mathrm{sgn}}/4\),
\(H_{12}^*H_{12}=4I\). Por tanto
\begin{equation}
 \|K\|^2=\frac14\|U_{\mathrm{sgn}}\|^2=4251257,
 \qquad \|U_{\mathrm{sgn}}\|^2=17005028.
 \label{lib04:norma-k}
\end{equation}
Cada traslación \(S^j\) es ortogonal y conserva
\(\nu=\sqrt{4251257}\). De ahí
\begin{equation}
 \operatorname{tr}(O_K^*O_K)=12\nu^2=51015084.
 \label{lib04:traza-marco}
\end{equation}
La división por \(\nu\) normaliza cada columna orbital a norma
uno y conserva su orden. Es uniforme en la órbita y bajo sus
transportes isométricos. Los registros distintos mantienen sus normas
propias, y las columnas normalizadas continúan teniendo un Gram
no escalar: sus valores propios son los anteriores divididos por
\(4251257\).

La normalización operatoria utiliza el Gram completo:
\begin{equation}
 G_K=O_K^*O_K,\qquad
 F_K=O_KG_K^{-1/2},\qquad
 F_K^*F_K=I.
 \label{lib04:polar}
\end{equation}
La raíz inversa existe por \eqref{lib04:gram-k}. La igualdad se
comprueba multiplicando por \(G_K^{-1/2}\) a ambos lados del Gram.
Así \(ZF_K\) es isométrico a cualquier profundidad.
La recuperación íntegra del marco es
\begin{equation}
 O_K=F_KG_K^{1/2}.
 \label{lib04:polar-inversa}
\end{equation}
El par \((F_K,G_K)\) conserva la información completa.
El factor unitario conserva un marco ortonormal orientado en las
cartas fijadas, mientras que las amplitudes residen también en
\(G_K\). Multiplicar \(K\) por un escalar positivo mantiene
\(F_K\) y cambia \(G_K\); \(K\) y \(-K\) comparten
Gram y poseen factores polares opuestos, con distinta lectura lineal
\(\kappa\). Cada ejemplo muestra por qué un factor aislado no
sustituye al par.

La traslación del origen conserva \(\nu\), pero en la carta
posicional fija
\begin{equation}
 \kappa(SK)=1000\kappa(K)-K_1.
 \label{lib04:cambio-origen}
\end{equation}
Se obtiene desplazando el numerador posicional y usando
\(d=1000^{12}-1\). Para conservar la misma lectura bajo un cambio
de carta se transporta también la fila \(\ell\) por la inversa
de ese cambio. Una carta unitaria general puede dejar la red de
bloques enteros; la decodificación por divisiones euclídeas se realiza
en la carta de bloques recuperada.

\subsection{Ejemplo exacto de transducción y corrección de errores}
\label{lib04:ejemplo}

Tomemos \(a=8\), \(U=S^4\) y el registro canónico
\begin{equation}
 \begin{split}
 K=(&234,543,140,729,659,824,\\
    &621,58,914,794,146,601).
 \end{split}
 \label{lib04:k-ejemplo}
\end{equation}
Indicando por \((v)_{1:6}\) y \((v)_{7:12}\) sus dos bloques de
seis coordenadas, las lecturas sin normalizar son
\begin{align}
 (Q_-K)_{1:6}&=(1213,3520,499,5774,4358,5798),\nonumber\\
 (Q_-K)_{7:12}&=(4822,-137,7078,5809,1028,4079),\nonumber\\
 (Q_+K)_{1:6}&=(2765,5711,1881,6619,6845,8210),\nonumber\\
 (Q_+K)_{7:12}&=(5735,1123,8460,7689,1454,6138).
 \label{lib04:canales-ejemplo}
\end{align}
Su suma, dividida por \(17\), recupera todos los bloques de \(K\).
La identidad matricial
\(9Q_-^*Q_-+8Q_+^*Q_+=1241I\) antecede a su evaluación en el
vector concreto.

Para introducir un error controlado, se racionalizan las dos memorias
como \(z_0=MK/\sqrt8\), y se añade \(1/10\) a la primera
coordenada. Cada bloque de doce componentes se transporta después
por los dos canales bilaterales. La inversa por suma recupera los
datos de memoria perturbados, y el decodificador de carga recibe
\(q=6263\) junto con esos datos, sin recibir el registro buscado.
La norma del error en las memorias originales es
\(\sqrt8/10<r_q\). La normalización isométrica conserva esa norma;
los canales sin normalizar sólo facilitan el cálculo racional exacto.

En esta instancia, los candidatos de la reconstrucción por suelo y
techo que satisfacen la carga son \(924\). El mínimo de distancia
cuadrática racionalizada es \(1/100\), y el decodificador recupera
\eqref{lib04:k-ejemplo}. La unicidad se sigue del radio estricto
\eqref{lib04:radios}, no del mero tamaño de la enumeración.
La lectura aritmética recuperada es
\begin{equation}
 \kappa=
 \frac{234543140729659824621058914794146601}
      {999999999999999999999999999999999999},
 \label{lib04:kappa-ejemplo}
\end{equation}
y la selección incidencial vuelve a ser \(\{4,6,9,10\}\), con
el mismo umbral \(729\). El ejemplo introduce el ruido antes del
transporte. Las pruebas de contractividad del adjunto cubren también
errores posteriores arbitrarios cuya norma conjunta pertenezca al
radio declarado.

La capacidad operativa reunida es transportar y recuperar un registro
generado, mantener el condicionamiento de sus respuestas y restituir
sus lecturas aritméticas e incidenciales antes de aplicar selecciones
discontinuas. Sus premisas efectivas son la producción previa de
\(K\), el marco cíclico invertible, el balance bilateral, la
separación discreta de las memorias y el teorema de archivo. El
registro de doce componentes conserva el alcance de sus lectores;
la historia enriquecida completa mantiene sus otros registros y
dependencias.

\section{Levantamiento unitario del refinamiento y transporte de lectores}
\label{lib05:levantamiento}

La biyección por residuos de la \cref{lib03:teo-residuo} y el
balance bilateral de la \cref{lib04:teo-balance} admiten una
composición explícita. Una terna aritmética admisible, que contiene
una pareja y su acarreo, determina una pareja refinada. La
linealización de esta biyección sobre bases ortonormales proporciona
un transporte unitario; dos transportes de ese tipo pueden actuar
como las dos realizaciones del balance bilateral.

La realización se establece después de los objetos generados por
APP--TRIT--TPK. Su dominio combinatorio comprende las parejas y los
acarreos que satisfacen las condiciones escritas abajo. La norma
conservada por su linealización es la suma de cuadrados de amplitudes
sobre esas etiquetas; la conservación aritmética de
\(L_c(x,y)=(Bx-c,By+c)\) continúa siendo la suma de las dos
coordenadas tras reescalar. La composición relaciona ambas
conservaciones mediante una aplicación concreta.

\subsection{Biyección de los estados aritméticos admisibles}
\label{lib05:biyeccion}

Sean \(B\geq2\) y \(M\geq1\) enteros. Definimos
\begin{align}
 \mathcal D_{B,M}
 &=\bigl\{(x,y,c)\in\mathbb Z_{\geq0}^2
                 \times\{0,\ldots,B-1\}:x+y=M,\ c\leq Bx\bigr\},
 \label{lib05:dominio}\\
 \mathcal E_{BM}
 &=\bigl\{(X,Y)\in\mathbb Z_{\geq0}^2:X+Y=BM\bigr\}.
 \label{lib05:codominio}
\end{align}
La condición \(c\leq Bx\) mantiene la no negatividad de la
primera coordenada de la salida. Para \(x=0\) sólo admite
\(c=0\); para \(x\geq1\) admite todo el alfabeto canónico.

\begin{theorem}[Correspondencia exacta de etiquetas]
\label{lib05:teo-biyeccion}
La aplicación
\begin{equation}
 \ell_{B,M}:\mathcal D_{B,M}\longrightarrow\mathcal E_{BM},
 \qquad (x,y,c)\longmapsto(Bx-c,By+c)
 \label{lib05:mapa-etiquetas}
\end{equation}
es biyectiva. Ambos conjuntos tienen cardinal \(BM+1\). Su inversa
está dada por
\begin{equation}
 c=Y\bmod B,\qquad y=\frac{Y-c}{B},\qquad x=\frac{X+c}{B}.
 \label{lib05:inversa-etiquetas}
\end{equation}
\end{theorem}

\begin{proof}
La condición de dominio y la identidad
\(Bx-c+By+c=BM\) aseguran que la imagen pertenece a
\(\mathcal E_{BM}\). Dada una pareja de ese conjunto, el residuo
\(c\) es el único elemento del alfabeto que satisface
\(Y\equiv c\pmod B\). Como \(X+Y\) es divisible por \(B\),
\(X+c\) también lo es. La división euclídea de \(Y\) da
\(y\geq0\), y \(X+c\geq0\) da \(x\geq0\). Si \(c>0\),
\(x=0\) exigiría \(X=-c<0\), por lo que se conserva la
admisibilidad. La suma de las preimágenes es \(M\), y la
sustitución verifica la imagen original. La unicidad del residuo y
de los cocientes demuestra la biyección.

Hay \(M+1\) etiquetas con \(c=0\) y \(M\) con cada acarreo
positivo; por tanto el cardinal es
\((M+1)+(B-1)M=BM+1\), coincidente con el de la salida.
\end{proof}

La prueba explicita el extremo de frontera del dominio, condición
necesaria para la igualdad de cardinales. La nueva etiqueta de acarreo
es una parte tipada de la entrada al refinamiento.

\subsection{Unitariedad sobre las amplitudes de las etiquetas}
\label{lib05:unitariedad}

Tomemos los espacios de Hilbert reales o complejos con medida de conteo
\begin{equation}
 \mathscr H_D=\ell^2(\mathcal D_{B,M}),\qquad
 \mathscr H_E=\ell^2(\mathcal E_{BM}),
 \label{lib05:hilbert}
\end{equation}
y sus bases ortonormales etiquetadas. Definimos linealmente
\begin{equation}
 J_{B,M}|x,y,c\rangle=|Bx-c,By+c\rangle.
 \label{lib05:j}
\end{equation}
Se escribirá \(J\) cuando los parámetros estén fijados.

\begin{proposition}[Levantamiento unitario]
\label{lib05:prop-unitaria}
El operador \(J:\mathscr H_D\to\mathscr H_E\) es unitario:
\(J^*J=I_D\), \(JJ^*=I_E\). El adjunto recupera cada etiqueta
mediante \eqref{lib05:inversa-etiquetas}.
\end{proposition}

\begin{proof}
La \cref{lib05:teo-biyeccion} asegura que \(J\) transforma
biyectivamente dos bases ortonormales completas. Para
\(\psi=\sum_{d\in\mathcal D}\psi_d|d\rangle\),
\[
 \|J\psi\|^2
 =\sum_{e\in\mathcal E}|\psi_{\ell_{B,M}^{-1}(e)}|^2
 =\sum_{d\in\mathcal D}|\psi_d|^2=\|\psi\|^2.
\]
La sobreyectividad prueba las dos identidades unitarias y la
biyección inversa describe el adjunto sobre cada vector de base.
\end{proof}

La etiqueta \(|73,27,1\rangle\) tiene norma uno; los enteros
\(73,27,1\) describen sus valores aritméticos, no sus amplitudes.
La construcción representa un procesamiento reversible de etiquetas
o amplitudes. Una realización física concreta conserva, además,
el protocolo que prepare y transporte tales amplitudes.

\subsection{Lectores diagonales de proporción y acarreo}
\label{lib05:diagonales}

Para una función real \(h\) sobre un conjunto finito, denotemos
por \(M_h\) su operador diagonal: cada vector de base se multiplica
por el valor de \(h\) en su etiqueta. La biyección satisface
\begin{equation}
 J^*M_hJ=M_{h\circ\ell_{B,M}}.
 \label{lib05:lector-diagonal}
\end{equation}
Se demuestra actuando sobre todos los vectores de la base.

Definamos en la salida
\begin{equation}
 f_E(X,Y)=\frac{X}{BM},\quad
 g_E(X,Y)=\frac{Y}{BM},\quad
 \gamma_E(X,Y)=Y\bmod B,
 \label{lib05:lectores-salida}
\end{equation}
y en la entrada \(f_D=x/M\), \(g_D=y/M\), \(\gamma_D=c\).
La identidad \eqref{lib05:lector-diagonal} proporciona
\begin{align}
 J^*M_{f_E}J&=M_{\,x/M-c/(BM)},\nonumber\\
 J^*M_{g_E}J&=M_{\,y/M+c/(BM)},\nonumber\\
 J^*M_{\gamma_E}J&=M_c.
 \label{lib05:transporte-fracciones}
\end{align}
Los dos primeros lectores transportan la transferencia de unidad.
Como \(f_E+g_E=1\), su suma es \(I_E\) y su compresión es
\(I_D\). El tercer lector conserva exactamente el acarreo como
observable entero, con su representación residual en la salida.

\subsection{Rutas completas a profundidad arbitraria}
\label{lib05:rutas}

Sea \(\mathcal P_n(B,M)\) el conjunto de rutas de longitud \(n\)
que parten de todas las parejas no negativas de suma \(M\) y, en
cada etapa, satisfacen las condiciones de
\eqref{lib05:dominio} para la suma correspondiente. Una ruta conserva
su pareja inicial y sus \(n\) acarreos admisibles. Sea
\(\Phi_n\) la aplicación que asigna a cada ruta su pareja final.

\begin{theorem}[Recuperación de rutas a toda profundidad finita]
\label{lib05:teo-rutas}
Para todo \(n\geq0\),
\begin{equation}
 \Phi_n:\mathcal P_n(B,M)\longrightarrow\mathcal E_{B^nM}
 \quad\hbox{es biyectiva},\qquad
 |\mathcal P_n(B,M)|=B^nM+1.
 \label{lib05:biyeccion-rutas}
\end{equation}
Su linealización
\begin{equation}
 J_n:\ell^2(\mathcal P_n(B,M))\longrightarrow
          \ell^2(\mathcal E_{B^nM}),\qquad
 J_n|\pi\rangle=|\Phi_n(\pi)\rangle
 \label{lib05:jn}
\end{equation}
es unitaria.
\end{theorem}

\begin{proof}
Desde una pareja de suma \(B^nM\), el residuo de
\eqref{lib05:inversa-etiquetas} recupera el último acarreo y una
única pareja previa de suma \(B^{n-1}M\). Repetir \(n\) veces
produce una única ruta cuya entrada suma \(M\). La composición
directa recupera el extremo original, lo que prueba biyección e
inversa. La cuenta también se obtiene desde las entradas: las
\(M\) parejas con \(x_0\geq1\) admiten \(B^n\) palabras cada
una, mientras que \((0,M)\) sólo admite la palabra de todos
ceros. El total es \(MB^n+1\). La transformación biyectiva de
bases ortonormales demuestra la unitariedad de \(J_n\).
\end{proof}

Conocida la profundidad, los lectores sobre la pareja final son
\begin{equation}
 c_j=\left\lfloor\frac{Y}{B^{n-j}}\right\rfloor\bmod B,
 \quad y_0=\left\lfloor\frac{Y}{B^n}\right\rfloor,
 \quad x_0=M-y_0,
 \qquad 1\leq j\leq n.
 \label{lib05:lectores-ruta}
\end{equation}
Estas funciones son diagonales y recuperan la memoria completa de
las operaciones finitas. La compatibilidad con prolongación es
\begin{equation}
 \Phi_{n+1}(\pi,c)=L_c(\Phi_n(\pi)).
 \label{lib05:compatibilidad-prolongacion}
\end{equation}
Al pasar de \(n\) a \(n+1\), se incorpora la nueva etiqueta
admisible de acarreo. Los espacios aumentan su dimensión y la
etiqueta de frontera sólo admite \(c=0\); por ello el dominio del
paso siguiente conserva la admisibilidad, en lugar de reemplazarla
por un producto irrestricto con \(B\) dígitos.

\subsection{Entrada fijada e imagen exacta}
\label{lib05:entrada-fija}

Si se fija una pareja inicial \((x_0,y_0)\), con \(x_0\geq1\),
hay \(B^n\) rutas y su imagen es exactamente
\begin{equation}
 \bigl\{(B^nx_0-C,B^ny_0+C):0\leq C\leq B^n-1\bigr\}.
 \label{lib05:imagen-fija}
\end{equation}
La \cref{lib03:cor-palabra} demuestra que los \(B^n\)
acumulados son distintos y recuperan las respectivas palabras. La
linealización es unitaria hacia el subespacio generado por esta
imagen. Si \(x_0=0\), la imagen consta de una sola etiqueta.

Los cilindros de la \cref{lib03:teo-cilindros} corresponden a una
entrada fijada. El dominio completo de
\eqref{lib05:biyeccion-rutas} reúne todas las parejas de suma \(M\).
Ambas construcciones poseen inversas explícitas, pero sus imágenes
y cardinales son diferentes. La profundidad arbitraria expresa la
validez de la prueba para todo \(n\); el límite de la lectura y
sus fronteras sigue siendo el establecido en la
\cref{lib03:limite}.

\subsection{Acoplamiento bilateral del refinamiento}
\label{lib05:bilateral}

Sea \(R\) una permutación previamente declarada de
\(\mathcal D_{B,M}\), linealizada unitariamente en
\(\mathscr H_D\). Definimos los dos transportes
\begin{equation}
 \mathcal B=J,\qquad \mathcal C=JR:
            \mathscr H_D\longrightarrow\mathscr H_E.
 \label{lib05:dos-transportes}
\end{equation}
La permutación forma parte de la especificación del procedimiento
y se fija antes de evaluar una lectura de salida.
Para \(a>0\), \(p=a^2+a+1\), \(N=(2a+1)p\),
\begin{equation}
 Q_-=J(aI-R),\qquad Q_+=J((a+1)I+R).
 \label{lib05:q-refinamiento}
\end{equation}

\begin{proposition}[Distribución y recuperación de amplitudes refinadas]
\label{lib05:prop-bilateral}
Los operadores \eqref{lib05:q-refinamiento} cumplen
\begin{equation}
 (a+1)Q_-^*Q_-+aQ_+^*Q_+=NI_D,
 \label{lib05:balance}
\end{equation}
y la aplicación
\begin{equation}
 W_a\psi=\frac1{\sqrt N}
    \bigl(\sqrt{a+1}\,Q_-\psi,\sqrt a\,Q_+\psi\bigr)
 \label{lib05:w}
\end{equation}
es isométrica. Para los canales sin normalizar
\(u=Q_-\psi\), \(v=Q_+\psi\),
\begin{equation}
 J\psi=\frac{u+v}{2a+1},\qquad
 \psi=J^*\frac{u+v}{2a+1}.
 \label{lib05:inversa-bilateral}
\end{equation}
\end{proposition}

\begin{proof}
Las dos aplicaciones \(\mathcal B,\mathcal C\) son unitarias.
Los términos cruzados de los Grams de \(Q_-\) y \(Q_+\) tienen
coeficientes \(-a\) y \(a+1\), respectivamente. Las ponderaciones
del balance los cancelan, mientras los términos diagonales suman
\(N\), como en la \cref{lib04:teo-balance}. La normalización
demuestra la isometría. La suma de los canales es
\((2a+1)J\psi\), y la unitariedad de \(J\) da la inversa.
\end{proof}

Desde los canales normalizados se recupera directamente la entrada
por \(W_a^*\), con norma inversora uno; este adjunto ya incorpora
\(J^*\). La fórmula alternativa
\eqref{lib05:inversa-bilateral} separa la recuperación de \(J\psi\)
por suma de la aplicación de \(J^*\), cuya acción sobre etiquetas
se calcula mediante residuos. Quedan compuestos el refinamiento
aritmético reversible y el transporte bilateral de amplitudes.

\begin{proposition}[Compatibilidad en el espacio refinado]
\label{lib05:prop-compatibilidad}
Un par sin normalizar \((u,v)\) pertenece a la imagen de
\(\psi\mapsto(Q_-\psi,Q_+\psi)\) si y sólo si
\begin{equation}
 \bigl[(a+1)I_E+JRJ^*\bigr]u
 =\bigl[aI_E-JRJ^*\bigr]v.
 \label{lib05:compatibilidad}
\end{equation}
\end{proposition}

\begin{proof}
La igualdad equivale a
\(JRJ^*(u+v)=av-(a+1)u\). Si se cumple, definamos
\(\psi=J^*(u+v)/(2a+1)\). Al sustituir en
\eqref{lib05:q-refinamiento}, esa identidad da
\(Q_-\psi=u\) y \(Q_+\psi=v\). La necesidad resulta de
insertar los canales verdaderos. Todas las composiciones actúan
en \(\mathscr H_E\), por lo que sus tipos están definidos.
\end{proof}

\subsection{Ley de los observables sobre los dos canales}
\label{lib05:ley-lectores}

\begin{theorem}[Compresión bilateral de un lector refinado]
\label{lib05:teo-lector}
Para \(F=F^*\) sobre \(\mathscr H_E\), escribamos
\(F_J=J^*FJ\). Entonces
\begin{equation}
 W_a^*\operatorname{diag}(F,F)W_a
 =\frac{a(a+1)F_J+R^*F_JR}{p}.
 \label{lib05:lector-bilateral}
\end{equation}
El lector \(F_J\) se conserva exactamente si y sólo si conmuta
con \(R\).
\end{theorem}

\begin{proof}
Expandir \((a+1)Q_-^*FQ_-+aQ_+^*FQ_+\) cancela los
términos mixtos. El coeficiente de \(\mathcal B^*F\mathcal B\)
es \(a(a+1)(2a+1)\); el de
\(\mathcal C^*F\mathcal C\) es \(2a+1\). Dividir por
\(N\), sustituir \(\mathcal C=JR\) y usar
\(F_J=J^*FJ\) prueba la fórmula. La igualdad con \(F_J\)
equivale a \(R^*F_JR=F_J\), o a \([F_J,R]=0\) por
unitariedad.
\end{proof}

En el caso \(a=8\), los pesos son \(72/73\) y \(1/73\).
El lector identidad recupera la norma. Un lector no constante
puede cambiar aun cuando la norma se conserve; su variación está
calculada por \eqref{lib05:lector-bilateral}.
El operador que conserva exactamente un lector original \(A\)
en la imagen completa es \(W_aAW_a^*\), cuya identidad de imagen
es \(W_aW_a^*\). Leer \(\operatorname{diag}(F,F)\) corresponde
a la compresión específica del teorema, y no sustituye
automáticamente ese observable transportado.

\subsection{Ejemplo aritmético de dimensión mil uno}
\label{lib05:ejemplo}

Tomemos \(B=10\), \(M=100\). Los espacios de etiquetas tienen
dimensión \(1001\), y
\begin{equation}
 d_1=(73,27,1)\longmapsto e_1=(729,271).
 \label{lib05:ejemplo-etiqueta}
\end{equation}
Los lectores diagonales dan
\begin{equation}
 f_E(e_1)=\frac{729}{1000}
 =\frac{73}{100}-\frac1{1000},\qquad
 g_E(e_1)=\frac{271}{1000},\qquad \gamma_E(e_1)=1.
 \label{lib05:ejemplo-lectores}
\end{equation}
Declaremos \(R\) como la transposición de las etiquetas
admisibles \(d_1=(73,27,1)\) y \(d_2=(73,27,2)\), dejando
fijas las demás. La segunda imagen es \(e_2=(728,272)\).
La regla de transposición se ha fijado sobre las etiquetas, antes
de cualquier comparación con una constante.

Para \(a=8\) y \(\psi=|d_1\rangle\),
\begin{equation}
 Q_-\psi=8|e_1\rangle-|e_2\rangle,\qquad
 Q_+\psi=9|e_1\rangle+|e_2\rangle.
 \label{lib05:ejemplo-canales}
\end{equation}
Las normas cuadradas son \(65\) y \(82\), y
\begin{equation}
 9\cdot65+8\cdot82=585+656=1241=17\cdot73.
 \label{lib05:ejemplo-balance}
\end{equation}
La suma de los canales es \(17|e_1\rangle\). La inversa
\eqref{lib05:inversa-bilateral} y el residuo recuperan
\(|d_1\rangle\), incluido el acarreo uno.

Si se evalúa \(F=M_{f_E}\) diagonalmente en ambos canales
normalizados, la \cref{lib05:teo-lector} da
\begin{equation}
 \begin{split}
 \langle W_8\psi,\operatorname{diag}(F,F)W_8\psi\rangle
 &=\frac{72\cdot729+728}{73\cdot1000}\\
 &=\frac{729}{1000}-\frac1{73000}.
 \end{split}
 \label{lib05:ejemplo-fraccion}
\end{equation}
La lectura complementaria es uno menos esta cantidad. Para el
lector diagonal de acarreo se obtiene
\begin{equation}
 \frac{72\cdot1+2}{73}=\frac{74}{73}.
 \label{lib05:ejemplo-acarreo}
\end{equation}
El acarreo original sigue siendo recuperable exactamente: el
observable \(W_8M_cW_8^*\) tiene sobre \(W_8|d_1\rangle\)
el valor propio uno. El valor \(74/73\) es la respuesta del
lector diagonal aplicado a ambos canales, cuyo transporte relativo
está especificado por \eqref{lib05:lector-bilateral}. El ejemplo
distingue conservación de la información y variación de una lectura
particular.

\subsection{Articulación con el registro generado y dominios de aplicación}
\label{lib05:dominios}

La unitaria \(J\) actúa sobre \(BM+1\) etiquetas, mientras que
la realización de \(K\) tiene doce coordenadas. En el ejemplo,
\(1001\) cuenta parejas refinadas de suma mil, y doce cuenta
ventanas del registro producido. Ambas dimensiones conservan su
origen combinatorio específico.

Una composición particular entre esos dos portadores puede formularse
mediante un selector material de doce rutas distintas, con su
procedencia y sus operadores. Si ese selector proporciona una
isometría
\(J_{12}:\mathbb R^{12}\to\ell^2(\mathcal D_{B,M})\), entonces
\(JJ_{12}\) es isométrica, y el entrelazamiento de los lectores se
comprueba sobre ella. Esta es la hipótesis de esa composición
adicional, no una identificación implícita por el número de
coordenadas. Las aplicaciones bilaterales ya definidas directamente
sobre \(K\) en la \cref{lib04:bilateral} conservan sus dominios
completos sin utilizar ese selector adicional.

El resultado de esta sección reúne una biyección aritmética, su
unitaria, lectores diagonales entrelazados, recuperación de rutas a
cualquier profundidad finita y una composición bilateral explícita.
El dominio de rutas es el definido en
\eqref{lib05:biyeccion-rutas}. Las aplicaciones a historias
filtradas por TPK conservan sus reglas de supervivencia y de
prolongación; la correspondencia de etiquetas prueba aquí el
transporte declarado, sin reemplazar esos operadores por un dominio
aritmético irrestricto.

\section{Unidad dimensional, simetría relativa y recuperación de la estructura}
\label{lib06:sintesis}

La extrema y media razón holográfica articula la conservación evolutiva
de la unidad con sus realizaciones aritméticas y dimensionales. La
composición desarrollada permite precisar ese principio mediante
cantidades conservadas, dominios e inversas. El estado enriquecido
generado por APP--TRIT--TPK precede a las lecturas numéricas y
geométricas; sus operaciones conservan la orientación, el acarreo,
la incidencia y la memoria que cada realización utiliza.

\subsection{Conservación aritmética, unitariedad y descomposición isométrica}
\label{lib06:tres-conservaciones}

El refinamiento aritmético conserva la unidad normalizada:
\begin{equation}
 (x,y,c)\longmapsto(Bx-c,By+c),\qquad
 \frac{Bx-c}{B(x+y)}+\frac{By+c}{B(x+y)}=1.
 \label{lib06:unidad}
\end{equation}
La inversa por residuos recupera el acarreo canónico y la pareja
anterior. La composición de bloques conserva las contribuciones de
los eventos con su escala, y las divisiones sucesivas recuperan toda
historia finita. La lectura límite se controla mediante cilindros
decrecientes; las marcas de frontera mantienen la distinción entre
historias que comparten una evaluación escalar.

El levantamiento sobre etiquetas transforma esa biyección en una
unitaria \(J\). Conserva la norma de las amplitudes y transporta
los lectores de fracción y acarreo mediante
\(J^*M_hJ=M_{h\circ\ell}\). La cantidad \(x^2+y^2\) y la norma
de \(|x,y,c\rangle\) pertenecen a realizaciones distintas. La
prueba de unitariedad establece el enlace preciso entre la
reversibilidad aritmética y la conservación de amplitudes.

La tercera conservación corresponde a dos realizaciones isométricas:
\begin{equation}
 (a+1)Q_-^*Q_-+aQ_+^*Q_+
 =(2a+1)(a^2+a+1)I.
 \label{lib06:balance}
\end{equation}
La cancelación de los términos mixtos permite la distribución
bilateral y su reconstrucción de norma uno. La ley vale para los
transportes del dominio especificado en la
\cref{lib04:teo-balance}. Su especialización al ciclo ternario
aporta propiedades adicionales de norma e índice; la ley general
conserva su validez para isometrías que carezcan de ese orden finito.

La cadena formada por las tres aplicaciones es efectiva: la biyección
produce \(J\); una permutación previamente declarada produce el
segundo transporte \(JR\); los dos transportes satisfacen el
balance, y sus canales recuperan la entrada por
\eqref{lib05:inversa-bilateral}. El ejemplo de la
\cref{lib05:ejemplo} muestra el procedimiento completo con
\((73,27,1)\), \((729,271)\) y \(1001\) etiquetas admisibles.
Los números de esa realización no se introducen como sustitutos de
la estructura que genera el registro \(K\).

\subsection{Memoria completa y contribución lógica al límite de Landauer}
\label{landauer:puente-archivo}

La contribución lógica nula de Landauer, establecida en
\eqref{v:ant:eq:ii-kb-aut-reversible}, se aplica a los transportes
recuperables construidos en este artículo. La biyección de
\cref{lib05:teo-biyeccion} conserva el acarreo por
\eqref{lib05:inversa-etiquetas}; su prolongación conserva la ruta
mediante \cref{lib05:teo-rutas}. Los archivos de
\cref{lib04:teo-archivo} conservan las amplitudes y correlaciones
necesarias para invertir el transporte.

\begin{corollary}[Recuperación y conservación espectral del archivo]
\label{landauer:cor-archivo}
Sea \(Z:H\to\mathcal K\) uno de los archivos completos
\(Z_N,Z_\infty\) de \cref{lib04:teo-archivo}, o el levantamiento
unitario \(J_n\) de \cref{lib05:teo-rutas}, en su respectivo dominio.
Para un operador de densidad de rango finito \(\rho\) sobre \(H\),
la salida \(\rho_Z=Z\rho Z^*\) satisface
\begin{align}
 \operatorname{Tr}\rho_Z&=1,& Z^*\rho_ZZ&=\rho,\\
 s(\rho_Z)&=s(\rho),&
 s(\rho)&=-\operatorname{Tr}(\rho\log\rho),
 \label{landauer:espectro}
\end{align}
con \(0\log0=0\). Si \(X\) es una variable aleatoria sobre las rutas
finitas de \cref{lib05:teo-rutas}, entonces
\begin{equation}
 \mathsf H(X\mid\Phi_n(X))=0.
 \label{landauer:recuperacion-ruta}
\end{equation}
\end{corollary}
\begin{proof}
Escribamos la descomposición espectral finita
\(\rho=\sum_jp_j|v_j\rangle\langle v_j|\), con los \(v_j\)
ortonormales, \(p_j>0\) y \(\sum_jp_j=1\). La igualdad
\(Z^*Z=I\) implica que los \(Zv_j\) son ortonormales. Por tanto,
\(\rho_Z=\sum_jp_j|Zv_j\rangle\langle Zv_j|\) tiene los mismos
autovalores no nulos, lo que prueba la traza y la entropía.
Multiplicar por \(Z^*\) y \(Z\) recupera \(\rho\). Finalmente,
\(\Phi_n^{-1}\) recupera cada ruta de su pareja final; la
distribución condicional está concentrada en una sola ruta y su
entropía es cero.
\end{proof}

La conservación corresponde al archivo completo. Una lectura visible
\(q\), acompañada de memoria \(M\), recupera el estado bajo la
hipótesis de \eqref{v:ant:eq:ii-kb-aut-fibra}; la información no
publicada por \(q\) permanece en
\(\mathsf H(M(X)\mid q(X))\). Descartar un complemento o un
terminal necesario define otra operación, cuya recuperabilidad se
determina sobre la información efectivamente retenida. La igualdad
espectral anterior conserva la distinción entre entropía del estado
completo, entropía de una reducción parcial y calor transferido.

En el modelo térmico de la sección
\ref{v:ant:sec:ii-kb-aut-desarrollo}, con estados lógicos degenerados,
entorno isotermo y contabilidad de registros auxiliares, el término
asociado al borrado es
\(k_B\Theta\,\mathsf H(X\mid F(X))\), para un mapa lógico
determinista \(F\). Su contribución es nula para los transportes
recuperables anteriores; el reinicio de un TRIT uniforme satisface
\(Q_{\rm reset}\geq k_B\Theta\log3\). La preparación, el control,
la velocidad finita y el acoplamiento al entorno conservan sus
balances físicos propios. El calor total de una implementación
requiere componer esas contribuciones con la operación lógica.

\subsection{La conservación de un observable y su simetría relativa}
\label{lib06:simetria}

El balance de la norma es el caso identidad de una ley operatoria
más informativa. Cuando los transportes son unitarios, con relativo
\(R\), el lector efectivo es
\begin{equation}
 \mathcal T_a(G)=\frac{a(a+1)G+R^*GR}{a(a+1)+1}.
 \label{lib06:lector}
\end{equation}
El observable permanece fijo exactamente cuando \([G,R]=0\).
En el caso canónico,
\begin{equation}
 \mathcal T_8(G)-G=\frac{R^*GR-G}{73}.
 \label{lib06:defecto-lector}
\end{equation}
Esta identidad determina tanto la conservación como la variación:
el defecto mide la diferencia entre dos realizaciones del mismo
lector, ponderada por el reparto \(72+1\).

La memoria completa transporta además el observable mediante
\(A\mapsto ZAZ^*\) en \(\operatorname{im}Z\). Los productos,
adjuntos y emparejamientos se conservan en esa imagen. La lectura
diagonal de canales y el observable transportado son operaciones
diferentes; el ejemplo del acarreo muestra la diferencia entre
\(74/73\) para la primera y el valor propio uno para el segundo.
La información permanece recuperable aunque una lectura diagonal
particular cambie de valor.

Las cargas variacionales de la \cref{nuc11:noether} conservan su
fundamento en una acción y en un entrelazamiento explícitos. Para
\(A_{n+1}U_n=U_nA_n\), la acción de transporte da la carga
\(p_n^*A_nq_n\). Su conservación y la identidad de norma son
resultados compatibles, con hipótesis diferentes. De igual modo,
la conservación de una forma alternante o lorentziana en
\eqref{lib04:balance-formas} mantiene su significado geométrico;
las estimaciones de ruido y contracción utilizan la positividad
de Hilbert declarada en sus pruebas.

\subsection{Función del registro en la transducción}
\label{lib06:registro}

La constante holográfica de acoplamiento aritmético-geométrico tiene
una función definida por sus composiciones. El registro \(K\)
se obtiene previamente de la extracción orientada y de Hadamard.
Su lectura \(\kappa\) es reversible con base, longitud y origen
fijados. Su marco \(O_K\) identifica respuestas cíclicas porque
su Gram tiene una cota inferior estrictamente positiva. La
incidencia y la dirección geométrica utilizan sus lectores
declarados sobre el mismo registro.

La isometría de archivo conserva el Gram de ese marco, por lo que
mantiene las cotas de recuperación a cualquier profundidad. El
registro completo admite inversión de norma uno; las dos memorias
parciales conservan su factor \(9/4\) y su núcleo uniforme. La
carga adicional distingue la restitución del registro absoluto de
la de su componente centrada. La decodificación entera restituye
después la selección incidencial incluso en un umbral situado
exactamente sobre una coordenada.

Esta capacidad es más específica que conservar una magnitud total:
permite recuperar operaciones y relaciones que actúan sobre un
objeto generado. Sus garantías descansan en el Gram, en las
distancias de la imagen entera y en las inversas demostradas, no
en conservar por definición otra copia de la entrada. A la vez,
el número finito de ventanas de \(K\) permanece distinto de la
historia enriquecida completa; los otros registros y marcas
acompañantes mantienen sus propios transportes.

La órbita y sus transportes isométricos conservan la norma
\(\nu=\|K\|\). Normalizar sus columnas a norma uno exige
conservar \(\nu\) para invertir esa operación. La normalización
polar utiliza el par
\((F_K,G_K)\). La fórmula
\(O_K=F_KG_K^{1/2}\) mantiene orientación y amplitudes.
El lector posicional debe viajar con el cambio de carta, pues la
norma orbital común no vuelve invariante \(\kappa\) bajo una
renumeración aislada de sus componentes.

\subsection{Profundidad, dimensiones y realizaciones físicas}
\label{lib06:profundidad}

La política bilateral que continúa un canal y archiva el otro
tiene agotamiento uniforme del terminal. Su límite es una
isometría hacia la suma de todas las memorias, y su inversa es
la serie adjunta convergente de
\eqref{lib04:recuperacion-infinita}. Un número finito de
complementos posee dos reconstrucciones diferenciadas: el
adjunto truncado conserva un sesgo explícito y la inversa de
mínimos cuadrados lo elimina con la amplificación de ruido
calculada. Las evoluciones precedentes con terminal positivo
conservan esa componente en el archivo completo.

La dimensión del espacio de etiquetas, el rango de un proyector y
la dimensión de una realización física se relacionan mediante
sus mapas concretos. La unitaria de \(1001\) etiquetas y el
registro de doce ventanas son dos realizaciones explícitas de
dominios distintos. La conservación exterior del desarrollo
precedente prolonga los productos de Gram a cada grado finito,
incluyendo las contribuciones mixtas entre terminal y memorias.
Las realizaciones excepcionales y físicas permanecen vinculadas
a sus representaciones y operadores ya construidos; los rangos
numéricos por sí solos no sustituyen esa vinculación.

\Needspace{18\baselineskip}
\subsection{Síntesis de la conservación evolutiva}
Las acciones físicas y el protocolo gaussiano del manuscrito
conservan los acoplamientos, las formas, las unidades y las
preparaciones que intervienen en sus pruebas. El balance
bilateral aquí añadido amplía la familia de transportes
conservativos y sus lecturas; su aplicación a una realización
física se establece componiendo los operadores de esa
realización, sin reemplazarlos por una coincidencia de
normalización.

La unidad del desarrollo reside, por tanto, en una secuencia
demostrada: generación y completación, refinamiento con memoria,
realización unitaria, simetría relativa del lector y recuperación
de la estructura. La extrema y media razón holográfica recibe
así una exposición en la que cambio de escala, incidencia,
conservación y transducción poseen operaciones explícitas y
mantienen conjuntamente la información necesaria para invertir
sus lecturas.

\section*{Formalización aritmética del balance bilateral}
\addcontentsline{toc}{section}{Formalización aritmética del balance bilateral}
\label{libro:formalizacion-lean}

La especialización nonádica de la ley bilateral admite una formalización
aritmética independiente. El punto de aplicación se encuentra después de la
construcción de sus coeficientes y transportes: la partición de base diez fija
\(a=8\), y los dos registros son \(8x-y\) y \(9x+y\). Sobre ese dominio,
Lean~4.21.0 verifica siete teoremas con su biblioteca \texttt{Std}. Las fuentes
formales y el procedimiento de ejecución se proporcionan como suplemento.
El presente complemento enuncia su contenido matemático y lo relaciona con
las demostraciones operatorias del artículo.

\subsection*{Identidad polinómica y conservación de la norma}

Para cualesquiera \(x,y\in\mathbb Z\), se cumple
\[
 9(8x-y)^2+8(9x+y)^2=17(72x^2+y^2).
\]
La expansión de los dos cuadrados cancela los productos cruzados. Cuando
\(y^2=x^2\), la misma igualdad da
\[
 9(8x-y)^2+8(9x+y)^2=17\cdot73\,x^2.
\]
Éstos son los dos primeros teoremas formalizados. El primero expresa un
balance para entradas enteras arbitrarias; el segundo aplica la hipótesis
explícita de conservación de la norma. El entero \(73=72+1\) ocupa
exactamente el lugar que tiene en la \cref{lib04:teo-balance}.

\subsection*{Dimensión finita arbitraria}

Sean \(n\in\mathbb N\) y \(x,y:\mathbb N\to\mathbb Z\). Definamos
\[
 S_n(x)=\sum_{j<n}x_j^2,
 \qquad
 E_n(x,y)=\sum_{j<n}\bigl(9(8x_j-y_j)^2+8(9x_j+y_j)^2\bigr).
\]
Los dos teoremas siguientes establecen
\[
 E_n(x,y)=17\bigl(72S_n(x)+S_n(y)\bigr)
\]
y, bajo la hipótesis \(S_n(y)=S_n(x)\),
\[
 E_n(x,y)=17\cdot73\,S_n(x).
\]
La prueba formal procede por inducción sobre \(n\). El caso inicial es la
suma vacía; el paso de inducción añade una coordenada y aplica la identidad
polinómica ya demostrada. El cuantificador universal sobre \(n\) incluye
todas las dimensiones finitas. La hipótesis utiliza la suma total de
cuadrados, por lo que admite redistribuciones entre coordenadas.

\subsection*{Recuperación de las entradas}

Escribamos \(q_-=8x-y\) y \(q_+=9x+y\). Los teoremas quinto y sexto
son las igualdades exactas
\[
 q_-+q_+=17x,\qquad 8q_+-9q_-=17y.
\]
En la imagen de la aplicación \((x,y)\mapsto(q_-,q_+)\), ambos miembros
izquierdos son múltiplos de \(17\). Las divisiones recuperan las dos
entradas. El séptimo teorema formula su consecuencia inyectiva:
\[
 \left.
 \begin{aligned}
 8x-y&=8x'-y',\\
 9x+y&=9x'+y'
 \end{aligned}
 \right\}
 \quad\Longrightarrow\quad x=x'\ \text{y}\ y=y'.
\]
La cancelación en \(\mathbb Z\) justifica esta conclusión. El contenido
formal coincide con la recuperación del par de registros, antes de
normalizarlos como canales de una realización métrica.

\subsection*{Relación entre las pruebas}

La formalización utiliza variables enteras arbitrarias e inducción
matemática. Por ello constituye una prueba de los enunciados anteriores,
distinta de la comprobación de un conjunto finito de ejemplos. El sistema
acepta los términos de prueba sin declaraciones axiomáticas añadidas ni
demostraciones diferidas. El suplemento registra asimismo las dependencias
lógicas comunicadas por la herramienta para cada declaración.

La ley operatoria del cuerpo del artículo posee un dominio más amplio:
isometrías entre espacios de Hilbert, archivo de profundidad arbitraria y
lectores conjugados. Sus demostraciones se encuentran en las
\cref{lib04:teo-balance,lib06:simetria,lib06:profundidad}. La prueba Lean
aquí incorporada certifica la especialización aritmética enunciada, con
sus cuantificadores y su estructura de partida. La procedencia de los
coeficientes, la generación del registro y las realizaciones excepcionales
conservan las demostraciones anteriores de las que dependen.

\HMTDivision{\(\mathcal C\)}{Recuperación compatible y familias del continuo}{Memoria suficiente, naturalidad, medida y transporte de pertenencia en el continuo generado.}
\HMTNarrativeHeading{De recuperar un estado a transportar una familia}

La recuperación compatible permite extender la exposición desde
un estado hasta familias de estados. Una lectura enriquecida
individualiza sus antecedentes cuando su memoria distingue los
elementos que comparten el mismo valor visible. La inversa
restituye entonces el estado sobre la imagen de esa lectura.
La compatibilidad con el refinamiento reúne las inversas finitas
en una reconstrucción de historias completas.

El transporte de familias exige conservar la pertenencia. Una
biyección de historias lleva cada familia a su imagen y cada
elemento a su correspondiente. Las uniones, intersecciones y
complementos relativos se transportan mediante esas operaciones.
La restricción de la biyección a una familia conserva su cardinal.
La teoría de memoria adquiere así un enlace preciso con las
operaciones sobre conjuntos de historias.

La medida requiere además el transporte de los pesos. Una
simetría compatible con el árbol de refinamiento conserva la
medida de cilindros cuando respeta las emisiones y sus
multiplicidades. La naturalidad expresa la concordancia entre
transportar y truncar. Las demostraciones siguientes explicitan
estas relaciones sobre los objetos definidos en el artículo.

El estudio específico de la hipótesis del continuo desarrolla
la clausura genealógica y su jerarquía de conjuntos. Su resultado
cardinal se formula sobre esa clausura y su potencia interna.
La relación que esta parte incorpora es la recuperación de
estados y el transporte de familias que mantienen la estructura
del continuo generado. El enunciado cardinal conserva su
desarrollo especializado y su dominio expresamente identificado.

Esta articulación sitúa el continuo dentro del argumento de
conservación. La historia completa, la familia a la que pertenece
y la medida de su cilindro requieren operaciones relacionadas,
con contenidos diferentes. El artículo reúne sus condiciones de
compatibilidad para mostrar cómo la doble proyección se mantiene
al pasar de una lectura individual a la organización de familias.
La unidad de la exposición reside en esos mapas y en las pruebas
que preservan sus relaciones.

\HMTMathematicalPaper
\section{Registros de acarreo e incidencia y reconstrucción proyectiva de historias}
\label{lib07:armonizacion}
La construcción APP--TRIT--TPK fija un estado enriquecido antes de
sus realizaciones numéricas, incidenciales y operatorias. Los registros
que intervienen en esa construcción se distinguen por la operación
que permiten invertir y por las relaciones que conservan. Esta sección
reúne esas distinciones y compone la recuperación con el refinamiento.

\subsection{Acarreo, historia, incidencia y archivo}
El acarreo de las secciones~\ref{lib03:residuo} y
\ref{lib03:memoria} es un dato aritmético entero. Su conservación,
junto con la profundidad y el estado final, restituye el antecedente
por la inversa explícita~\eqref{lib03:inversa-afin}.
La historia de transporte contiene la sucesión ordenada de operaciones,
hojas, orientaciones y retornos. En ella, dos etapas con igual fase
pueden tener registros acumulados distintos. La holonomía actúa sobre
esa historia y sobre el marco en que se efectúan sus lecturas.

El registro de incidencia conserva las relaciones de frontera que
utilizan las construcciones excepcionales. El objeto ordenado \(K\)
y su lectura \(\kappa_{\rm ag}\) intervienen con las normalizaciones
y los mapas de la sección~\ref{lib01:identificacion}.
La suficiencia para recuperar un operador se expresa allí mediante
las respuestas del marco y la inversión correspondiente.
La memoria de archivo de la sección~\ref{lib04:archivo}, por su parte,
es una sucesión de componentes de un espacio de Hilbert. La identidad
telescópica~\eqref{lib04:telescopia} conserva la norma entre el estado
remanente y las componentes archivadas; la inversa
\eqref{lib04:recuperacion-infinita} reconstruye la amplitud inicial.

Estas cuatro funciones se componen conservando sus tipos.
Un entero de acarreo, una palabra de transporte, una incidencia
y una amplitud archivada constituyen datos distintos del sistema.
Las pruebas de recuperación especifican qué combinación de ellos
es suficiente para cada lectura. La figura~\ref{lib07:fig-registros}
representa esa jerarquía funcional.

\begin{figure}[htbp]
\centering
\begin{tikzpicture}[
 node distance=9mm,
 box/.style={draw=ink,rounded corners=2pt,fill=white,align=center,
 text width=58mm,minimum height=15mm,font=\small},
 arr/.style={-{Latex[length=2mm]},draw=ink,thick}]
\node[box,text width=125mm] (s) {Estado enriquecido APP--TRIT--TPK\\
 operación, hoja, orientación, frontera y memoria};
\node[box] (a) at (-33mm,-30mm)
 {Acarreo aritmético\\Residuo, profundidad e inversa};
\node[box] (h) at (33mm,-30mm)
 {Historia de transporte\\Composición ordenada y holonomía};
\node[box] (q) at (-33mm,-62mm)
 {Registro de incidencia\\Lecturas de frontera y marco};
\node[box] (m) at (33mm,-62mm)
 {Archivo de recuperación\\Componentes y reconstrucción};
\node[box,text width=125mm] (r) at (0,-94mm)
 {Recuperación compatible con el refinamiento\\
 Inversas tipadas y relaciones conservadas};
\draw[arr] (s.south -| a.north)--(a);
\draw[arr] (s.south -| h.north)--(h);
\draw[arr] (a)--(q);
\draw[arr] (h)--(m);
\draw[arr] (q.south)--(r.north -| q.south);
\draw[arr] (m.south)--(r.north -| m.south);
\end{tikzpicture}
\caption{Organización de los registros según su función.
Las flechas inferiores indican composición de condiciones de
recuperación, y no identificación de los cuatro tipos de dato.}
\label{lib07:fig-registros}
\end{figure}

\subsection{Árbol de refinamiento y pesos locales}
Sea \(H_k\) el conjunto finito de prefijos admisibles de profundidad
\(k\), con truncamientos \(p_{k+1,k}\), del sistema de historias ya
construido. Una emisión se cuenta con su multiplicidad y
\(d(\zeta)>0\) denota el número de emisiones desde \(\zeta\).
La distribución uniforme de las semillas y el peso local
\(1/d(\zeta)\) definen \(\mu_k\). La suma de los pesos de las
prolongaciones de un prefijo coincide con su peso: las medidas son
compatibles y determinan la medida cilíndrica \(\mu\).
Se demuestran a continuación las proposiciones de memoria de fibra y
naturalidad de la medida, con estas definiciones explícitas.

\subsection{Lectores enriquecidos y memoria de fibra}
\label{sec:ch-memoria-fibra}

La diferencia entre historia y valor admite un criterio operativo. En un
nivel finito \(H_k\), sea \(q_k:H_k\to Z_k\) un lector y sea
\(M_k:H_k\to\mathcal M_k\) el registro conservado. El lector enriquecido
es \(\widehat q_k=(q_k,M_k)\). La memoria distingue estados de una misma
fibra del lector; su suficiencia se comprueba en el dominio considerado.
Esta formulación incorpora al propio artículo el desarrollo completo de la
información de fibra que se utiliza a continuación.

\begin{proposition}[Recuperabilidad en un nivel finito]
\label{prop:ch-recuperabilidad-fibra}
Existe una inversa de \(\widehat q_k\) sobre su imagen si y sólo si
\(\widehat q_k\) es inyectivo. Para una variable aleatoria \(Z\) con
valores en \(H_k\) y distribución \(\mu_k\), se cumple
\[
 \mathsf H(Z)=\mathsf H(q_k(Z))+\mathsf H(Z\mid q_k(Z)).
\]
Si \(\widehat q_k\) es inyectivo sobre el soporte de \(Z\), entonces
\[
 \mathsf H(Z\mid q_k(Z),M_k(Z))=0.
\]
Aquí \(\mathsf H(Z)=-\sum_z\Pr(Z=z)\log\Pr(Z=z)\), con
\(0\log0=0\), es la entropía finita en nats; la entropía condicional es
el promedio de la entropía de las distribuciones condicionadas.
\end{proposition}

\begin{proof}
Una inversa recupera un único antecedente, lo que exige inyectividad;
recíprocamente, ésta permite asignar a cada par publicado su único
antecedente. Puesto que \(q_k(Z)\) es función determinista de \(Z\),
agrupar la suma que define \(\mathsf H(Z)\) por fibras de \(q_k\) da la
regla de la cadena. Si el par \((q_k(Z),M_k(Z))\) distingue el soporte,
cada distribución condicionada correspondiente está concentrada en un solo
estado y su entropía es cero.
\end{proof}

No publicar una coordenada y borrar el registro que permite recuperarla son
operaciones distintas. Este resultado no identifica entropía con cardinalidad
del continuo ni introduce una interpretación térmica: su dominio, lector y
distribución son los que se acaban de fijar.

\subsection{Equivarianza de la medida bajo simetrías del árbol}
\label{sec:ch-naturalidad-medida}

\begin{proposition}[Transporte conjunto del árbol y de su medida]
\label{prop:ch-naturalidad-medida}
Sea \(g=(g_k)_{k\geq0}\) una familia de biyecciones \(g_k:H_k\to H_k\)
que conmuta con los truncamientos, fija la raíz y transporta biyectivamente
las emisiones admisibles, conservando sus multiplicidades. Entonces
\[
 d(g_k\zeta)=d(\zeta),\qquad
 (g_k)_*\mu_k=\mu_k,\qquad (g_\infty)_*\mu=\mu.
\]
Más generalmente, si \(\mu_\zeta\) es la ley de continuaciones de una
historia \(\zeta\in H_k\), se tiene
\[
 (g_\infty)_*\mu_\zeta=\mu_{g_k\zeta}.
\]
\end{proposition}

\begin{proof}
La biyección de emisiones conserva su número con multiplicidades. La
distribución uniforme de semillas también es invariante. Cada trayectoria
y su imagen reciben, por tanto, el mismo producto de pesos locales
\(1/d(\zeta)\). Se obtiene la igualdad de medidas en cada nivel y en
cada cilindro. La coherencia con los truncamientos define \(g_\infty\),
y la unicidad de la medida determinada por cilindros prueba la igualdad
en el límite. El mismo argumento, comenzando en \(\zeta\), prueba la
afirmación sobre continuaciones.
\end{proof}

Se obtiene así la naturalidad del árbol y de la medida: un cambio coherente
de representación transporta los pesos y no obliga a escogerlos de nuevo.
La conservación de emisiones y multiplicidades es una hipótesis esencial de
la proposición, no un dato añadido después de la construcción.

\subsection{Recuperación de historias y transporte de pertenencia}
\begin{proposition}[Compatibilidad de inversas y familias]
\label{lib07:prop-familias}
Sean \(H_k\) los niveles anteriores y \(Y_k=\widehat q_k(H_k)\).
Supóngase que \(\widehat q_k:H_k\to Y_k\) es biyectivo y que existen
truncamientos \(s_{k+1,k}\) tales que
\[
 s_{k+1,k}\widehat q_{k+1}=\widehat q_k p_{k+1,k}.
\]
Entonces las inversas finitas inducen una biyección
\[
 \widehat q_\infty:\varprojlim H_k\longrightarrow\varprojlim Y_k.
\]
Si \(H=\varprojlim H_k\), \(Y=\varprojlim Y_k\) y
\(\mathcal T(A)=\widehat q_\infty[A]\), la aplicación
\(\mathcal T:\mathcal P(H)\to\mathcal P(Y)\) es biyectiva y
preserva uniones, intersecciones, complementos relativos y pertenencia.
Cada restricción \(A\to\mathcal T(A)\) es una biyección.
\end{proposition}
\begin{proof}
Al componer la identidad de compatibilidad con las inversas se obtiene
\[
 p_{k+1,k}\widehat q_{k+1}^{-1}
   =\widehat q_k^{-1}s_{k+1,k}.
\]
Por tanto, una familia compatible \(y=(y_k)\) determina la familia
compatible \(x=(\widehat q_k^{-1}(y_k))\). Las dos aplicaciones
coordenadas son inversas. La imagen directa por una biyección tiene
como inversa la imagen directa por la biyección inversa. Para toda
familia \((A_i)_{i\in I}\) de subconjuntos de \(H\),
\[
 \begin{aligned}
 \mathcal T\Bigl(\bigcup_i A_i\Bigr)&=\bigcup_i\mathcal T(A_i),\\
 \mathcal T\Bigl(\bigcap_i A_i\Bigr)&=\bigcap_i\mathcal T(A_i),\\
 \mathcal T(H\setminus A)&=Y\setminus\mathcal T(A).
 \end{aligned}
\]
La primera igualdad usa la definición de imagen. Para la segunda y
la tercera, la unicidad del antecedente permite trasladar las
condiciones de pertenencia. Finalmente,
\(x\in A\) equivale a
\(\widehat q_\infty(x)\in\mathcal T(A)\), lo que prueba también
la afirmación sobre restricciones.
\end{proof}

El teorema de refinamiento de rutas~\ref{lib05:teo-rutas}
proporciona una realización de esta compatibilidad sobre las etiquetas
aritméticas y sus acarreos. El resultado distingue la recuperación
de cada estado y la conservación de las operaciones sobre familias.
En el nivel de amplitudes, el levantamiento unitario
de~\ref{lib05:prop-unitaria} transporta además el producto interno.
En el nivel probabilístico, la proposición
\ref{prop:ch-naturalidad-medida} conserva los pesos cuando se
transportan las emisiones con sus multiplicidades.

\subsection{Articulación con la clausura genealógica del continuo}
La sección~\ref{sec:cierre-cardinal-nativo} construye una clausura
genealógica \(\mathfrak G\) y establece la identidad cardinal
\(2^{\aleph_0}=\aleph_1\) sobre su potencia interna. La clausura,
el código generador y la historia realizada forman parte de su
definición; su representación conjuntista se establece después de la
construcción genealógica, mediante las operaciones y demostraciones
incorporadas en dicha sección.

La presente articulación incorpora las pruebas de memoria,
naturalidad y transporte de familias que utiliza esa relación con
la conservación. La cadena demostrativa queda discriminada:
las inversas recuperan estados; su compatibilidad recupera
historias; la biyección transporta las operaciones de pertenencia;
la clausura genealógica de la sección~\ref{sec:cierre-cardinal-nativo} fija el dominio
del enunciado cardinal. Así, la recuperación y la organización
de familias se integran en la doble proyección con sus respectivos
objetos, operaciones y conclusiones.

\section{Clausura genealógica y cardinalidad del continuo generado}
\label{sec:cierre-cardinal-nativo}

Esta sección reúne en una demostración contigua la prolongación, la doble
proyección y la profundidad arbitraria; construye después la clausura
semántica, el levantamiento uniforme, la maquinaria genealógica y el cierre
cardinal. La jerarquía conjuntista se identifica únicamente al término de la
construcción y no puede ser leída como origen causal. La dependencia causal es:

\[
 \mathrm{APP}\longrightarrow\mathrm{TRIT}\longrightarrow\mathrm{TPK}
 \longrightarrow X_\infty^{\rm enr}
 \longrightarrow\mathcal C_{\rm cont}^{\rm disc}
 \longrightarrow m_\infty
 \longrightarrow\mathfrak G_{\rm HMT}(h,m_\infty).
\tag{C.32}\label{eq:ch-cadena-causal}
\]

Primero se construyen el estado, sus prolongaciones, el continuo y el registro genealógico.
Después se prolonga \textbf{esa genealogía producida} mediante una regla semántica
explícita a todos sus objetos hereditarios. La identificación posterior con
una jerarquía conjuntista conocida no selecciona las semillas, los operadores,
el registro genealógico, la recta ni sus proyecciones. Sí permite contrastar, con total
transparencia, la fuerza lógico-cardinal exacta de la regla semántica adoptada.

\subsection{Prolongación de estados y extensión semántica transfinita}

Sea

\[
 \Sigma_{\rm HMT}
 =\bigl(\dot R_{\rm APP},\dot R_{\rm TRIT},\dot R_{\rm TPK},
 \dot R_{U},\dot R_{\Gamma_9},\dot R_{\rm Ext},\dot R_\rho,
 \dot R_{\pi},\dot R_{\varphi},\dot R_e,\dot R_{\alpha},
 \dot R_{\rm ar},\dot R_{\rm exc}\bigr)
\tag{C.33}\label{eq:ch-signatura}
\]

la signatura \textbf{relacional} numerable cuyo código efectivo es \(h\). Cada
\(\dot R_S\) es el grafo del operador \(S\) sobre los códigos donde está
definido; no se presupone que una función externa sea total dentro de cada
nivel. Los grafos \(\dot R_{\pi},\dot R_{\varphi},\dot R_e,\dot R_{\alpha}\)
son los lectores de salida construidos por la generación coinductiva y el
cierre dodecafásico; no reciben como argumentos los valores convencionales de
las constantes. Su presencia en \(\Sigma_{\rm HMT}\) asegura que la clausura
conserve el estado completo: integra simultáneamente las rutas que publican
valor y la incidencia que se prolonga hacia la realización excepcional.

En cada horizonte, el código del estado registra las coordenadas APP--TRIT--TPK,
la ruta, la hoja, el acarreo, la frontera y la memoria; en sus etapas de
publicación registra asimismo \(P_N,\Phi_N,E_N,A_N\), donde
\(A_N\) es el prefijo de \(\alpha_{\rm HMT}\) producido por el mismo estado.
Estas coordenadas no seleccionan la rama desde fuera: son parte de la historia
que se codifica. Fijada una numeración primitivo-recursiva de las coordenadas finitas de
un estado, toda historia inversa
\(x_\infty=(x_N)_{N\geq1}\in X_\infty^{\rm enr}\) determina el código
\[
 \operatorname{code}(x_\infty)
 :=\{\langle N,j,c\rangle:(x_N)_j\text{ tiene código }c\}\subseteq\omega.
\]
Las ecuaciones de coherencia
\(\rho_{N+1,N}(x_{N+1})=x_N\) permiten decodificar de manera única la
historia a partir de ese real. Escribimos \(m_\infty\) para el código de la
historia realizada y \(m_\infty(n)\) para su función característica bajo la
numeración fijada. El código operatorio \(h\) y este registro completo se
reúnen, mediante una función primitiva recursiva de emparejamiento, en el real

\[
 u_{h,m}:=
 \{\langle0,\langle n,h(n)\rangle\rangle:n<\omega\}
 \cup
 \{\langle1,\langle n,m_\infty(n)\rangle\rangle:n<\omega\}.
\tag{C.33a}\label{eq:ch-codigo-conjunto}
\]

Las dos proyecciones primitivas recuperan \(h\) y \(m_\infty\) desde
\(u_{h,m}\).

La primera extensión es la prolongación métrica y operatoria de los estados
finitos. Para \(N\ge1\), sea
\(\widehat{\mathcal X}^{\rm enr,raw}_{6N}\) el conjunto de estados
enriquecidos brutos de longitud \(6N\), y defínase conjuntistamente el
subsistema superviviente por
\[
\begin{aligned}
 X_N
 &:=\widehat{\mathcal X}^{\rm enr,surv}_{6N}\\
 &:=\Bigl\{x\in\widehat{\mathcal X}^{\rm enr,raw}_{6N}:\\[-.2ex]
 &\qquad\forall M>N\ \exists z\in
   \widehat{\mathcal X}^{\rm enr,raw}_{6M}\;\text{ tal que}\\[-.2ex]
 &\qquad\operatorname{tr}_{6M,6N}(z)=x\ \wedge\
   z\text{ satisface toda transición TPK intermedia}\Bigr\}.
\end{aligned}
\]
y póngase
\[
 \rho_{N+1,N}:=\operatorname{tr}_{6N+6,6N}\!\upharpoonright X_{N+1},
\]
y
\[
 \operatorname{Ext}_N(x):=
 \{y\in X_{N+1}:(x,y)\in
 \operatorname{Ext}_{6N+6,6N}\}.
\]
El nivel raíz se incluye sin excepción implícita:
\[
 X_0:=\{\ast\},\qquad
 \rho_{1,0}:X_1\to X_0\text{ es constante},\qquad
 \operatorname{Ext}_0(\ast):=X_1.
\]
El superíndice \(\mathrm{surv}\) denota el
árbol podado de los estados que poseen descendientes TPK compatibles a todo
horizonte finito. La ley de prolongación TPK establecida en las secciones
precedentes prueba que las semillas
emitidas pertenecen a ese árbol; la ramificación finita y el lema de Kőnig
convierten la compatibilidad a todo horizonte en una rama infinita. En
particular,

\[
 \forall N\ge0\;\forall x\in X_N\quad
 \varnothing\ne\operatorname{Ext}_N(x)
 \subseteq\rho_{N+1,N}^{-1}(x).
\tag{C.34}\label{eq:ch-extension-no-vacia}
\]

Por el mismo lema de árbol finitamente ramificado, toda semilla superviviente
posee una prolongación compatible

\[
 x_\infty=(x_N)_{N\ge1}\in
 X_\infty^{\rm enr}:=\varprojlim_N(X_N,\rho_{N+1,N}).
\tag{C.35}\label{eq:ch-limite-inverso}
\]

La segunda extensión actúa \textbf{después} sobre el código completo de esa
historia realizada.
No se infiere de la mera existencia del operador finito
\(\operatorname{Ext}_N\): formaliza a escala conjuntista el principio HMT de
que sólo son objetos del dominio íntegro los que poseen una derivación
genealógica desde el estado y sus reglas. La extensión semántica transfinita
es:

\[
 \begin{aligned}
  \mathfrak G_0(h,m_\infty)
  &:=\operatorname{TC}\bigl(\{\omega,u_{h,m},h,m_\infty\}\bigr),\\
  \mathfrak G_{\beta+1}(h,m_\infty)
  &:=\operatorname{Def}_{\Sigma_{\rm HMT}}
       \bigl(\mathfrak G_\beta(h,m_\infty)\bigr),\\
  \mathfrak G_\lambda(h,m_\infty)
  &:=\bigcup_{\beta<\lambda}\mathfrak G_\beta(h,m_\infty)
       \quad(\lambda\text{ límite}),\\
  \mathfrak G_{\rm HMT}(h,m_\infty)
  &:=\bigcup_{\beta\in\operatorname{Ord}}
       \mathfrak G_\beta(h,m_\infty).
 \end{aligned}
\tag{C.36}\label{eq:ch-clausura-transfinita}
\]

La hipótesis del continuo y una biyección objetivo no se introducen entre las
condiciones que definen el operador sucesor; el operador está dado
prospectivamente por

\[
 \operatorname{Def}_{\Sigma_{\rm HMT}}(M)
 :=M\cup
 \Bigl\{
   \{x\in M:\langle M,\in,
       (\dot R_S\cap M^{k_S})_{S\in\Sigma_{\rm HMT}}\rangle
       \models\varphi(x,\bar p)\}:
   \ulcorner\varphi\urcorner\in\omega,
   \ \bar p\in M^{<\omega}
 \Bigr\}.
\tag{C.37}\label{eq:ch-operador-def}
\]

Es decir: publica simultáneamente todos los subconjuntos definibles por una
fórmula finita del lenguaje relacional HMT y parámetros que la genealogía ya
produjo. No
consulta un valor real objetivo, una constante convencional, un cardinal
objetivo ni un modelo exterior. La fórmula~\eqref{eq:ch-operador-def} define el cierre transfinito
sobre las semillas, los operadores, el estado enriquecido, el registro y la
doble proyección construidos previamente.

\begin{lemma}[Codificación conjunta del programa HMT y del registro realizado]
\label{thm:ch-6}
El par formado por
el código operatorio \(h\) y el registro \(m_\infty\) se codifica por un único
real \(u_{h,m}\subseteq\omega\), y toda fórmula de \(\Sigma_{\rm HMT}\) se
traduce uniformemente a una fórmula del lenguaje \(\{\in,u_{h,m}\}\).

\end{lemma}
\begin{proof}
Tanto \(h\) como \(m_\infty\) son sucesiones de naturales: \(h\)
codifica la signatura, las tablas finitas y los programas de los operadores;
\(m_\infty\) codifica una sola historia realizada, no el espacio completo de
historias. La fórmula~\eqref{eq:ch-codigo-conjunto} y sus dos decodificadores prueban la primera
afirmación. Cada símbolo de~\eqref{eq:ch-signatura} designa su \textbf{relación de grafo codificada en \(h\)}, nunca el grafo externo de todas sus posibles evaluaciones. APP,
TRIT, \(U_t\), \(\Gamma_9\), las restricciones \(\rho\) y las extensiones
finitas tienen grafos primitivo-recursivos sobre códigos enteros y, por tanto,
fórmulas \(\Delta_0\) absolutas entre niveles transitivos. Las dos
publicaciones se expresan como relaciones de grafo mediante las fórmulas

\[
 \begin{aligned}
 \dot R_{\rm ar}(x,y)
 &\;\Longleftrightarrow\;
 \forall k\in\omega\;\exists N\;\forall n\ge N\quad
 |\mu_n(x)-y|<2^{-k},\\
 \dot R_{\rm exc}(x,c,z)
 &\;\Longleftrightarrow\;
 \forall N\in\omega\quad z\!\upharpoonright N
   =E_N(x\!\upharpoonright N,c\!\upharpoonright N),
 \end{aligned}
\tag{C.37c$_0$}\label{eq:ch-grafos-publicacion}
\]

donde \(\mu_n\) y \(E_N\) son las funciones racionales/finitas cuyos códigos
figuran en \(h\). Se sustituye inductivamente cada fórmula atómica que contiene
un símbolo HMT por su fórmula de grafo con parámetro \(u_{h,m}\). Los
conectivos y cuantificadores se conservan. Esta inducción sobre la complejidad
sintáctica produce una traducción \(\varphi\mapsto\varphi^*\) tal que, para
todo nivel transitivo \(M\) de~\eqref{eq:ch-clausura-transfinita} que contiene los parámetros,

\[
 \langle M,\in,(\dot R_S\cap M^{k_S})_S\rangle
   \models\varphi(\bar a)
 \quad\Longleftrightarrow\quad
 M\models\varphi^*(\bar a;u_{h,m}).
\tag{C.37c}\label{eq:ch-eliminacion-uniforme}
\]

Así, \eqref{eq:ch-operador-def} no puede consultar un oráculo, una tabla decimal objetivo ni una
relación no codificada que introduzca incontables objetos en un solo paso.
\end{proof}

La semántica HMT se define ahora \textbf{independientemente de identificarla con \(\mathfrak G\)}. Sea \(\mathcal T_0^{\rm HMT}\) el conjunto de nombres
canónicos de los elementos de
\(\mathfrak G_0(h,m_\infty)=
\operatorname{TC}(\{\omega,u_{h,m},h,m_\infty\})\). En particular, el nivel
inicial contiene el registro realizado \(m_\infty\), no todas las historias de
\(X_\infty^{\rm enr}\) como parámetros primitivos.
Los términos admisibles son los árboles bien fundados obtenidos por

\[
 \begin{aligned}
 \mathcal T^{\rm HMT}_{\beta+1}
 &:={\mathcal T}^{\rm HMT}_\beta\cup
 \{\langle\beta,\ulcorner\varphi\urcorner,
       t_0,\ldots,t_{k-1}\rangle:
       t_i\in\mathcal T^{\rm HMT}_\beta\},\\
 \mathcal T^{\rm HMT}_\lambda&:=
       \bigcup_{\beta<\lambda}\mathcal T^{\rm HMT}_\beta,
 \qquad
 \mathcal T^{\rm HMT}:=
       \bigcup_{\beta\in\operatorname{Ord}}\mathcal T^{\rm HMT}_\beta.
 \end{aligned}
\tag{C.37a}\label{eq:ch-terminos}
\]

El valor de un término sucesor es el subconjunto del nivel anterior definido
por \(\varphi\) en la estructura relacional inducida, con los valores de
\(t_0,\ldots,t_{k-1}\) como parámetros. Se define, antes de cualquier
comparación conjuntista,

\[
 \operatorname{Ob}_{\rm sem}^{\rm HMT}(h,m_\infty)
 :=\{\llbracket t\rrbracket:t\in\mathcal T^{\rm HMT}\}.
\tag{C.37b}\label{eq:ch-objetos-semanticos}
\]

Ésta es la regla de \textbf{exhaustividad generativa} adoptada por HMT: todo
objeto del dominio íntegro debe poseer uno de esos árboles de derivación y todo
árbol admisible publica un objeto. No es consecuencia de la ramificación
finita. La hipótesis del continuo no se introduce en esta regla; se obtiene
mediante la demostración desarrollada en el artículo. Tampoco se introducen
como datos una enumeración de los reales ni una biyección objetivo.

\begin{theorem}[Equivalencia entre la semántica de términos y la jerarquía genealógica]
\label{thm:ch-6a}
La evaluación de
los términos anteriores satisface

\[
 \operatorname{Ob}_{\rm sem}^{\rm HMT}(h,m_\infty)
 =\mathfrak G_{\rm HMT}(h,m_\infty).
\tag{C.37d}\label{eq:ch-equivalencia-semantica}
\]

\end{theorem}
\begin{proof}
Por inducción transfinita. En el nivel inicial las dos clases
coinciden. Si todo valor de un término de
\(\mathcal T^{\rm HMT}_\beta\) pertenece a
\(\mathfrak G_\beta\), la semántica del término sucesor es uno de los
subconjuntos de~\eqref{eq:ch-operador-def}. Recíprocamente, cada subconjunto de~\eqref{eq:ch-operador-def} viene dado por
un código de fórmula y una tupla finita de parámetros; por la hipótesis
inductiva esos parámetros poseen términos y~\eqref{eq:ch-terminos} forma el término que lo
denota. En un límite ambas construcciones toman la misma unión.
\end{proof}

\Needspace{9\baselineskip}
El codificador de estados es compatible con los truncamientos y envía cada
prolongación TPK a un término sucesor. Sea
\(\mathcal T_N^{\rm st}\subset\mathcal T^{\rm HMT}\) la subclase de términos
que codifica estados completos de profundidad \(N\), y sea
\(\tau_{N+1,N}:\mathcal T_{N+1}^{\rm st}\to\mathcal T_N^{\rm st}\) el borrado
del último bloque junto con el truncamiento correspondiente del registro genealógico,
acarreo y memoria. Si \(J_N:X_N\to\mathcal T_N^{\rm st}\) es el codificador
completo, la compatibilidad de \(\operatorname{Ext}_N\) prueba la identidad
correctamente tipada

\[
 \forall x\in X_N\ \forall y\in\operatorname{Ext}_N(x)\qquad
 \tau_{N+1,N}\bigl(J_{N+1}(y)\bigr)=J_N(x).
\tag{C.37e$_0$}\label{eq:ch-naturalidad-codificador}
\]

Defínase dentro de la semántica
\[
 \operatorname{Succ}^{\rm TPK}_{N+1}(J_N(x)):=
 \{t\in\mathcal T_{N+1}^{\rm st}:\exists y\in X_{N+1}\,
   (y\in\operatorname{Ext}_N(x)\ \wedge\ t=J_{N+1}(y))\}.
\]
Elíjase \(\beta\) tal que \(\mathfrak G_\beta\) contenga \(x\), \(X_N\),
\(X_{N+1}\), los grafos finitos de
\(\operatorname{Ext}_N,J_N,J_{N+1}\) y sus códigos en \(h\). Entonces la
fórmula primitivo-recursiva del grafo de \(\operatorname{Ext}_N\) da

\[
 J_{N+1}[\operatorname{Ext}_N(x)]
 =\operatorname{Succ}^{\rm TPK}_{N+1}(J_N(x)),
 \qquad
 J_{N+1}[\operatorname{Ext}_N(x)]
 \in\mathfrak G_{\beta+1},
\tag{C.37e}\label{eq:ch-sucesores-semanticos}
\]

y ambos miembros son formados por la fórmula de grafo de la misma extensión
TPK. Las ecuaciones~\eqref{eq:ch-naturalidad-codificador}--\eqref{eq:ch-sucesores-semanticos} son el enlace semántico explícito: la
subclase de términos sucesores es exactamente la imagen de las prolongaciones
admisibles y su truncamiento conmuta. No se identifica la fibra completa de
\(\rho\) con \(\operatorname{Ext}_N\), ni se identifica
\(\operatorname{Def}\) con TPK por nomenclatura.

\begin{lemma}[Transitividad, axiomas y buen orden del modelo genealógico]
\label{thm:ch-6b}
\(\mathfrak G_{\rm HMT}(h,m_\infty)\) es una clase transitiva, contiene todos
los ordinales, satisface la teoría de conjuntos de Zermelo--Fraenkel (ZF) y
posee un buen orden global definible; satisface,
por tanto, elección y todos los axiomas usados en la comparación cardinal posterior.

\end{lemma}
\begin{proof}
Abreviemos \(B=\mathfrak G_0\) y \(u=u_{h,m}\). La inducción
sobre~\eqref{eq:ch-clausura-transfinita} prueba simultáneamente que cada nivel es transitivo y que
\(\mathfrak G_\alpha\in\mathfrak G_{\alpha+1}\); si todos los ordinales
menores que \(\alpha\) están presentes, su conjunto es definible en el nivel
correspondiente y publica \(\alpha\). Por ello la unión contiene todos los
ordinales. Extensionalidad y Fundación se heredan del universo ambiente;
Vacío e Infinito están en \(B\). Elegidos un nivel que contenga los parámetros,
las fórmulas que definen \(\{a,b\}\), \(\bigcup a\) y los subconjuntos por
Separación los publican en un sucesor.

Para Colección--Reemplazo, sea \(F\) una relación funcional definible en
\(\mathfrak G\) sobre el conjunto \(a\). El esquema de reflexión relativo a
la jerarquía definible \((\mathfrak G_\xi,\in,u_{h,m})_{\xi\in\mathrm{Ord}}\),
aplicado en la metateoría ambiente a la lista finita de fórmulas que define
\(F\), produce un nivel
\(\mathfrak G_\theta\) que contiene \(a\), los parámetros y es correcto para
esas fórmulas. Los rangos de los testigos de \(F[a]\) están entonces acotados
por \(\theta\); Separación sobre \(\mathfrak G_\theta\) forma la imagen en
\(\mathfrak G_{\theta+1}\). Para Potencia no se presupone el conjunto interno
que se quiere construir. Fijado \(a\in\mathfrak G_\gamma\), se considera en
la metateoría ambiente el conjunto ya existente \(\mathcal P^V(a)\). Para
cada \(b\in\mathcal P^V(a)\cap\mathfrak G\), sea \(r(b)\) su primera etapa
genealógica. Como \(\mathfrak G\) es una clase definible, Separación en la
metateoría ambiente forma primero el conjunto
\(\mathcal P^V(a)\cap\mathfrak G\). Reemplazo \textbf{en la metateoría ambiente}
aplicado a ese conjunto acota los ordinales \(r(b)\) por algún \(\theta\).
En consecuencia,
\[
 \mathcal P^{\mathfrak G}(a)
 =\{b\in\mathfrak G_\theta:b\subseteq a\},
\]
y el miembro derecho se publica por Separación en
\(\mathfrak G_{\theta+1}\). No introduce en \(\mathfrak G\) los subconjuntos
exteriores a la semántica~\eqref{eq:ch-objetos-semanticos}. Éste es el argumento no circular de
acotación por rangos.

Finalmente, a cada objeto se le asigna su primer nivel de nacimiento y su
menor \textbf{nombre de derivación}: un árbol de derivación bien fundado y
finitamente ramificado (y, por~\eqref{eq:ch-terminos}, finito para cada nombre individual), etiquetado por un
ordinal de nivel, un código de fórmula y la tupla de nombres de sus parámetros.
Por recursión transfinita se ordenan esos nombres: primero la etiqueta ordinal,
después el código de fórmula y después, lexicográficamente, los nombres de los
parámetros ya ordenados. El menor nombre de cada valor define una relación de
clase \(<_{\mathfrak G}\) que compara todos los objetos. Es un buen orden
global definible y, tomando el menor elemento de cada conjunto no vacío,
prueba Elección.
\end{proof}

\subsection{Elevación TPK de la fibra ternaria y partición cilíndrica completa}
\label{subsec:elevacion-catalogal-ch}

La fibra de \(729=3^6\) subcilindros de refinamiento no procede del cociente de transferencia
\(\mathcal K_{\rm ph}\). Ese operador es una publicación posterior de memoria
reducida y no puede generar historias. La cobertura se construye directamente
en el sistema \((H_k,\rho_k,\operatorname{Ext}_k)\) mediante las tres
prolongaciones de una vuelta de la holonomía nonádica.

\subsubsection{Inversión de Hensel en la torre aritmética y realización por historias}
\label{sec:ch-capacidad-hensel}

Antes de levantar las ramas a historias enriquecidas, conviene recuperar el
argumento finito que determina la capacidad de la torre aritmética. Se
incorpora la prueba completa del transductor no filtrado de la matriz; su
dominio se distingue del árbol superviviente.

\begin{proposition}[Conjugación de la torre de Hensel no filtrada]
\label{prop:ch-hensel-bruto}
Sean \(p\) primo, \(d\geq1\) y \(A\in\operatorname{GL}_d(\mathbb Z)\).
Para vectores fila, tomando representantes de residuos en
\(\{0,\ldots,p-1\}^d\), se define
\[
 x_tA=u_t+px_{t+1},\qquad u_t=x_tA\bmod p.
\]
La aplicación
\[
 \Psi_T:(\mathbb Z/p^T\mathbb Z)^d\longrightarrow
             (\mathbb F_p^d)^T,\qquad
 x_0\longmapsto(u_0,\ldots,u_{T-1})
\]
es biyectiva, con inversa
\[
 x_0\equiv\sum_{t=0}^{T-1}p^tu_tA^{-(t+1)}\pmod{p^T}.
\]
Las biyecciones conmutan con las reducciones y dan un homeomorfismo
\(\Psi:\mathbb Z_p^d\to(\mathbb F_p^d)^{\mathbb N}\), donde
\(\mathbb Z_p:=\varprojlim_T\mathbb Z/p^T\mathbb Z\). La actualización
\(x_t\mapsto x_{t+1}\) corresponde a suprimir el primer bloque.
\end{proposition}

\begin{proof}
La inversa entera de \(A\) existe porque su determinante es \(\pm1\).
Iterando \(x_t=(u_t+px_{t+1})A^{-1}\), se obtiene
\[
 x_0=\sum_{t=0}^{T-1}p^tu_tA^{-(t+1)}+p^Tx_TA^{-T}.
\]
Módulo \(p^T\), una palabra determina una única clase inicial. Dada una
palabra arbitraria, fijar \(x_T=0\) y reconstruir hacia atrás produce una
clase que la emite, lo que prueba la sobreyectividad. La misma fórmula
prueba compatibilidad al truncar. El paso al límite inverso da el
homeomorfismo y la relación con el desplazamiento de la cinta.
\end{proof}

Para \(p=3\) y \(d=6\), el alfabeto de bloques tiene \(3^6=729\)
elementos. Esta capacidad aritmética no se identifica por definición con
la realizabilidad de cada bloque tras cada prefijo superviviente. El
enlace con las aristas \(\operatorname{Ext}\) es precisamente el cometido
del lema~\ref{lem:elevacion-catalogal-nonadica}. Se mantienen así las dos
operaciones que distingue la matriz: representación en la torre bruta y
realización bajo las compatibilidades APP--TRIT--TPK.



\subsubsection{Realización arista por arista de las tres prolongaciones nonádicas}
\label{subsubsec:tres-vueltas-ext-ch}

Esta subsección especifica el modelo autónomo generado por la relación de
extensión TPK que interviene en la prueba cardinal. No se infiere la existencia de
una arista por coincidencia de tipos: se construyen sus estados inicial y final,
su relación APP, su transporte de fibra y la actualización unaria del registro.
La construcción utiliza exclusivamente las dos hojas de APP, el selector TRIT
y la composición \(\mathsf{Sel}\)--\(\mathsf{Tra}\)--\(\operatorname{Upd}\) del TPK.

\paragraph{Células completas y tres reducciones balanceadas.}
En la carta \(D_9=\{1,\ldots,9\}\), mantendremos siempre la célula APP
completa
\[
 a_{p,q}:=(p,q;R^\Sigma_{p,q},q^\Sigma_{p,q},
                 R^\Pi_{p,q},q^\Pi_{p,q}),
\]
con
\[
 p+q=R^\Sigma_{p,q}+9q^\Sigma_{p,q},\qquad
 pq=R^\Pi_{p,q}+9q^\Pi_{p,q}.
\]
La hoja activa se registra en una coordenada separada; no se escinde
\(a_{p,q}\) en una supuesta célula aditiva y otra multiplicativa. Para
\(r\in D_9\), sean
\begin{equation}
 a(r):=
 \begin{cases}
  a_{1,9},&r=1,\\
  a_{1,r-1},&2\le r\le9,
 \end{cases}
 \qquad
 \tau(r):=M_{\rm ph}(r),
 \qquad
 m(r):=\tau(r)\in\{-1,0,+1\}.
 \label{eq:celula-completa-fase-ch}
\end{equation}
El residuo aditivo de \(a(r)\) proporciona aquí un representante APP de la
fase; esto no convierte \(\Sigma\) en hoja activa universal. El selector
dinámico activa \(\Sigma\) en \(r\in\{1,4,7\}\), activa \(\Pi\) en
\(r\in\{2,5,8\}\) y conserva ambas hojas, sin desplazar cursor, en
\(r\in\{3,6,9\}\). En los tres casos se transportan conjuntamente
\(R^\Sigma,q^\Sigma,R^\Pi,q^\Pi\). Se tiene
\[
 (m(1),\ldots,m(9))=(1,-1,0,1,-1,0,1,-1,0).
\]
Las tres parejas ordenadas por la base nonádica son
\begin{equation}
 \mathfrak p_{+1}:=(1,4),\qquad
 \mathfrak p_{-1}:=(2,5),\qquad
 \mathfrak p_{0}:=(3,6).
 \label{eq:tres-pares-internos-ch}
\end{equation}
Estas parejas son la sección normalizada que toma las dos primeras
ocurrencias de cada tipo TRIT bajo el desplazamiento \(r\mapsto r+3\); el
tercer bloque de fases \(7,8,9\) conserva el control de cierre de la vuelta.
Sus dos miembros pertenecen a una misma fibra de \(M_{\rm ph}\). La
reducción balanceada del capítulo TRIT da, sin introducir un coeficiente
objetivo,
\begin{equation}
 \begin{array}{c|c|c|c}
 \varepsilon&m(\mathfrak p_\varepsilon^{(1)})
                  +m(\mathfrak p_\varepsilon^{(2)})
   &r_3\bigl(2\varepsilon\bigr)&
   \chi(\varepsilon,\varepsilon)\\ \hline
 +1&1+1=(-1)+3\cdot1&-1&+1\\
 0&0+0=0+3\cdot0&0&0\\
 -1&(-1)+(-1)=1+3\cdot(-1)&+1&-1
 \end{array}
 \label{eq:tres-carry-derivados-ch}
\end{equation}
Por tanto,
\begin{equation}
 \chi(\varepsilon,\varepsilon)
 =\frac{2\varepsilon-r_3(2\varepsilon)}3=\varepsilon.
 \label{eq:carry-no-parametrico-ch}
\end{equation}
El signo de la rama es así la salida del cociclo balanceado aplicado a dos
emisiones TRIT producidas por APP. Cada pareja proporciona una de las tres
salidas de acarreo. Por claridad, indexaremos después la pareja mediante
\(\varepsilon(\mathfrak p):=
\chi(m(\mathfrak p^{(1)}),m(\mathfrak p^{(2)}))\); el símbolo
\(\varepsilon\) no suministra a \(\operatorname{Upd}^{\rm cal}\) un acarreo libre.

La arquitectura de fase, el levantamiento balanceado y el transporte de
memoria se realizan aquí mediante veintisiete incidencias locales y sus
grafos de extensión. La sección normalizada fija representantes de las
fibras de \(M_{\rm ph}\); la cobertura que se demostrará utiliza la existencia
de esos representantes y permanece invariante al sustituirlos por otra
sección compatible.

\paragraph{Datos tipados de las nueve aristas.}
Sea
\(K_{\rm Carr}:=\langle\kappa_{\rm Carr}\rangle_{\rm Ab}\) el módulo de
acarreo. Las acciones están tipadas por
\[
 H_{\rm Carr}:\mathscr P_{\rm TPK}^{\rm op}
   \longrightarrow\operatorname{End}(K_{\rm Carr}),\qquad
 Y_{\rm Carr}:\mathscr P_{\rm TPK}
   \longrightarrow\operatorname{End}(K_{\rm Carr}).
\]
En este modelo generado se toman las acciones triviales
\(H_{\rm Carr}(e)=Y_{\rm Carr}(e)=\operatorname{id}_{K_{\rm Carr}}\): la
hoja activa y la ruta se conservan como coordenadas propias del registro y no
actúan sobre el sumando libre de acarreo. Sea además
\((\mathsf M_d,\jmath_{d',d})\) el sistema directo graduado de memoria. Para
no confundir dos representantes identificados por el colímite, el estado
conserva además el grado como coordenada y se usa el portador etiquetado
\[
 \widetilde{\mathsf M}:=
 \bigsqcup_{d\ge1}\bigl(\{d\}\times\mathsf M_d\bigr),
 \qquad
 \operatorname{gr}_{\mathsf M}(d,\mu)=d.
\]
La publicación en la memoria límite viene dada por
\[
\begin{aligned}
 \mathsf M&:=\varinjlim_d\mathsf M_d,&
 \jmath_d&:\mathsf M_d\longrightarrow\mathsf M,\\
 u&:=\jmath_d(u_d),&
 \iota_{\mathsf M}:\mathbb Z&\longrightarrow\mathsf M,
 \qquad n\longmapsto nu .
\end{aligned}
\]
Aquí \(\jmath_d\) es el mapa canónico del colímite y los \(u_d\) forman una
familia compatible. Su componente libre permite
el lector \(\deg_u(nu)=n\).
Las proyecciones de configuración, fase, memoria, acarreo, ruta y
supervivencia se tiparán sobre el portador calibrado construido
recursivamente más abajo; no se importará para ellas ningún espacio ambiental.
Para cada arista construida \(e\), la
componente calibrada
\(\operatorname{Upd}^{\rm cal}_e:=\mathsf{Mem}^{\rm cal}_e\) denotará la
tercera aplicación de la factorización; su ley se fija campo por campo más
abajo y no se presupone como restricción de una relación ambiental.

Fijemos la profundidad inicial \(k_0=1\), una configuración completa
\(q_\star\in\mathcal Q_{\rm obs}\) de fase uno y
\[
 q_{n,r}:=F_{\rm obs}^{9n+r-1}(q_\star),\qquad
 q_{n,10}:=q_{n+1,1}.
\]
Como el calendario observable es periódico, su valor \(q_{n,r}\) no
distingue por sí solo apariciones separadas por doce vueltas. El modelo
calibrado emplea, por ello, la elevación etiquetada
\[
 \widetilde q_{n,r}:=(q_{n,r},n,r),\qquad
 \widetilde q_{n,10}:=\widetilde q_{n+1,1},\qquad
 \pi_{\rm obs}(\widetilde q_{n,r})=q_{n,r}.
\]
La dinámica elevada
\(\widetilde F_{\rm obs}(\widetilde q_{n,r})
=\widetilde q_{n,r+1}\) proyecta a \(F_{\rm obs}\); la etiqueta no modifica
la fase física, sino que hace distintos sus niveles de aparición.
Se etiqueta la copia de la célula en cada nivel mediante
\[
 a_{n,r}:=(n,r;a(r)),\qquad
 (n,r)^+:=
 \begin{cases}
  (n,r+1),&1\le r<9,\\
  (n+1,1),&r=9.
 \end{cases}
\]
La etiqueta sólo evita identificar dos apariciones de una misma célula; las
dos divisiones euclídeas y sus cuatro coordenadas son las de \(a(r)\).
Para la vuelta \(n\ge0\) y la fase \(r\in D_9\), el símbolo
\[
 e_{\varepsilon,n,r}:\widetilde q_{n,r}
 \longrightarrow\widetilde q_{n,r+1},
 \qquad r^+:=1+(r\bmod9),
\]
ocupa el nivel
\(d_{n,r}:=k_n+r-1\), donde \(k_n:=k_0+9n\). Su registro local contiene
la célula completa \(a(r)\), la hoja activa determinada por \(M_{\rm ph}(r)\),
el tipo \(\tau(r)\),
la marca de pertenencia a \(\mathfrak p_\varepsilon\) y, al alcanzar el segundo miembro
de esa pareja, la identidad certificada
\eqref{eq:tres-carry-derivados-ch}. Definimos sus operadores únicamente a
partir de ese registro:
\begin{align}
 \mathsf C_{\mathsf M}(e_{\varepsilon,n,r})
   &=\jmath_{d_{n,r}+1,d_{n,r}},&
 \kappa_{\mathsf M}(e_{\varepsilon,n,r})
   &=\mathbf1_{\{r=9\}}u_{d_{n,r}+1},
 \label{eq:memoria-arista-ch}\\
       \operatorname{Carr}(e_{\varepsilon,n,r})
   &=d_{\varepsilon,r}\kappa_{\rm Carr},&
 d_{\varepsilon,r}
   &:=\mathbf1_{\{r=\mathfrak p_\varepsilon^{(2)}\}}
       \chi\!\left(m(\mathfrak p_\varepsilon^{(1)}),
                    m(\mathfrak p_\varepsilon^{(2)})\right).
 \label{eq:carry-arista-derivado-ch}
\end{align}
Para las identidades y para toda composición admisible se define, además,
\begin{align}
 \mathsf C_{\mathsf M}(\operatorname{id}_{\widetilde q};d)
   &:=\operatorname{id}_{\mathsf M_d},&
 \kappa_{\mathsf M}(\operatorname{id}_{\widetilde q};d)&:=0,
 \label{eq:memoria-identidad-ch}\\
 \mathsf C_{\mathsf M}(e_t\cdots e_1)
   &:=\mathsf C_{\mathsf M}(e_t)\circ\cdots\circ
      \mathsf C_{\mathsf M}(e_1),&
 \kappa_{\mathsf M}(e_2e_1)
   &:=\mathsf C_{\mathsf M}(e_2)
        \bigl(\kappa_{\mathsf M}(e_1)\bigr)
      +\kappa_{\mathsf M}(e_2).
 \label{eq:cociclo-memoria-calibrado-ch}
\end{align}
La identidad lleva el grado \(d\) porque una misma configuración observable
puede reaparecer en profundidades distintas. La última igualdad, iterada, define
\(\kappa_{\mathsf M}(e_t\cdots e_1)\) para cualquier longitud y es la ley de
cociclo del sistema directo; las inclusiones satisfacen
\(\jmath_{d+2,d+1}\jmath_{d+1,d}=\jmath_{d+2,d}\).
Además se fija
\(s(e_{\varepsilon,n,r})=\top\). Las siguientes veintisiete entradas
certifican las coordenadas que distinguen las tres prolongaciones; los campos
universales del transporte se especifican inmediatamente después. Los símbolos \(I\) y \(II\)
designan, respectivamente, la primera emisión conservada y la segunda emisión
que efectúa la reducción; un guion indica incremento de acarreo nulo, pero no
ausencia de arista.
\begin{center}
\small
\setlength{\tabcolsep}{3.2pt}
\begin{tabular}{c c cc cc cc c}
\toprule
\(r\)&\(M_{\rm ph}(r)\)&\multicolumn{2}{c}{\(\mathfrak p_{+1}\)}&
\multicolumn{2}{c}{\(\mathfrak p_0\)}&
\multicolumn{2}{c}{\(\mathfrak p_{-1}\)}&
\(\Delta\operatorname{mem}\)\\
& &marca&\(d_{+1,r}\)&marca&\(d_{0,r}\)&marca&\(d_{-1,r}\)&\\
\midrule
1&\(+1\)&I&0&--&0&--&0&0\\
2&\(-1\)&--&0&--&0&I&0&0\\
3&0&--&0&I&0&--&0&0\\
4&\(+1\)&II&\(+1\)&--&0&--&0&0\\
5&\(-1\)&--&0&--&0&II&\(-1\)&0\\
6&0&--&0&II&0&--&0&0\\
7&\(+1\)&--&0&--&0&--&0&0\\
8&\(-1\)&--&0&--&0&--&0&0\\
9&0&--&0&--&0&--&0&\(u\)\\
\bottomrule
\end{tabular}
\end{center}
Así, \(d_{\varepsilon,r}\) se calcula a partir de dos salidas del selector y
es distinto de los cocientes \(q^\Sigma,q^\Pi\) de la célula. El registro
conserva también el par marcado y el residuo balanceado; en particular, la
rama nula no se confunde con una ausencia de evento.

La ley de composición del acarreo empleada por la categoría TPK es
\begin{equation}
 \operatorname{Carr}(\operatorname{id}_{\widetilde q})=0,\qquad
 \operatorname{Carr}(e_2e_1)
 =H_{\rm Carr}(e_1)\bigl(\operatorname{Carr}(e_2)\bigr)
  +Y_{\rm Carr}(e_2)\bigl(\operatorname{Carr}(e_1)\bigr).
 \label{eq:ley-carry-correcta-ch}
\end{equation}
En identidades y caminos se prolongan functorialmente:
\[
\begin{aligned}
 H_{\rm Carr}(\operatorname{id})
   &=Y_{\rm Carr}(\operatorname{id})=\operatorname{id},\\
 H_{\rm Carr}(e_t\cdots e_1)
   &=H_{\rm Carr}(e_1)\circ\cdots\circ H_{\rm Carr}(e_t),\\
 Y_{\rm Carr}(e_t\cdots e_1)
   &=Y_{\rm Carr}(e_t)\circ\cdots\circ Y_{\rm Carr}(e_1).
\end{aligned}
\]
En este modelo todas estas acciones siguen siendo la identidad, pero
su tipado explicita qué endomorfismo actúa sobre cada sumando de
\eqref{eq:ley-carry-correcta-ch}.
En las aristas calibradas, \(H_{\rm Carr}=Y_{\rm Carr}=\operatorname{id}\),
por lo que la composición suma los incrementos ya producidos. La ecuación
\eqref{eq:ley-carry-correcta-ch}, y no una identificación entre orientación y
cociente, gobierna el registro acumulado.

\paragraph{Construcción independiente de las tuplas de transición.}
Un estado de la fibra se escribirá, en el orden de sus campos pertinentes,
\[
 x=(a;q;\tau,o,h;b_{\rm ref};
       R^\Sigma,q^\Sigma,R^\Pi,q^\Pi;c;
       \operatorname{Sig};C,\partial C;p_{\APP},p_{\TRIT};
       \lambda;\operatorname{gr}_{\mathsf M};\operatorname{mem};
       s;\operatorname{Led}).
\]
La notación de actualización \(x[f\leftarrow v]\) cambia sólo el campo
indicado. Para \(h\in\{(\Sigma,\Pi),(\Pi,\Sigma)\}\), defínase
\[
 h_r(h)=
 \begin{cases}
  (\Sigma,\Pi),&r\in\{1,4,7\},\\
  (\Pi,\Sigma),&r\in\{2,5,8\},\\
 h,&r\in\{3,6,9\}.
 \end{cases}
\]
La coordenada completa de hoja y aritmética se abrevia por
\[
 \widehat h(x):=
 \bigl(h(x),R^\Sigma(x),q^\Sigma(x),R^\Pi(x),q^\Pi(x)\bigr).
\]
Sea \(\widehat{\mathcal H}^{\rm cal}_{\widetilde q_{n,r}}\) el conjunto de
las quíntuplas
\[
 \bigl(h,R^\Sigma_{a_{n,r}},q^\Sigma_{a_{n,r}},
          R^\Pi_{a_{n,r}},q^\Pi_{a_{n,r}}\bigr),
 \qquad
 h\in\{(\Sigma,\Pi),(\Pi,\Sigma)\},
\]
cuyos cuatro datos aritméticos son las dos divisiones euclídeas de la célula
APP etiquetada \(a_{n,r}\). Esta definición directa no presupone un estado
en un portador todavía no construido. Sean
\(\mathsf{Hist}^{\rm cal}_{\widetilde q}\) el conjunto de caminos finitos
admisibles del grafo de aristas etiquetadas que terminan en
\(\widetilde q\), incluida la identidad, y
\(\mathsf B^{\rm cal}_{\widetilde q}\) el conjunto de sus sucesiones
ordenadas de extremos. Definimos
\[
 \mathsf S^{\rm cal}_{\widetilde q}:=
 \widehat{\mathcal H}^{\rm cal}_{\widetilde q}
 \times\{o_{\rightarrow},o_{\circ},o_{\leftarrow}\}
 \times K_{\rm Carr}\times
 \mathsf{Hist}^{\rm cal}_{\widetilde q}\times
 \mathsf B^{\rm cal}_{\widetilde q}\times\mathsf M .
\]
La firma calibrada no se postula: es la aplicación de reunión tipada
\[
 \operatorname{Sig}^{\rm cal}_{\widetilde q}:
 \widehat{\mathcal H}^{\rm cal}_{\widetilde q}
 \times\{o_{\rightarrow},o_{\circ},o_{\leftarrow}\}
 \times K_{\rm Carr}\times
 \mathsf{Hist}^{\rm cal}_{\widetilde q}\times
 \mathsf B^{\rm cal}_{\widetilde q}\times\mathsf M
 \longrightarrow\mathsf S^{\rm cal}_{\widetilde q},
\quad
 (\widehat h,o,c,\lambda,\partial C,\mu)
 \longmapsto(\widehat h,o,c,\lambda,\partial C,\mu).
\]
Asimismo, \(\pi_{\mathcal F}\) denota la proyección que elimina únicamente
las dos primeras coordenadas \((a;q)\) de la tupla de estado; conserva
\(\tau,o,h,b_{\rm ref}\), los cuatro datos euclídeos, acarreo, firma,
cilindro, frontera, prefijos, ruta, grado de memoria, memoria, supervivencia y
registro.
La orientación correspondiente queda fijada por
\(o(+1)=o_{\rightarrow}\), \(o(0)=o_{\circ}\) y
\(o(-1)=o_{\leftarrow}\).
La fibra APP calibrada sobre \(\widetilde q_{n,r}\) contendrá la copia etiquetada
\(a_{n,r}\) ya definida. El estado inicial queda fijado sin coordenadas libres:
\[
 \widehat h_\star:=
 \bigl((\Sigma,\Pi),R^\Sigma_{1,9},q^\Sigma_{1,9},
       R^\Pi_{1,9},q^\Pi_{1,9}\bigr).
\]
\begin{equation}
\begin{aligned}
 x_\star:=\bigl(&a_{0,1};\widetilde q_{0,1};
 +1,o_{\rightarrow},(\Sigma,\Pi);\bot;
 R^\Sigma_{1,9},q^\Sigma_{1,9},R^\Pi_{1,9},q^\Pi_{1,9};0;\\
 &\operatorname{Sig}^{\rm cal}_{\widetilde q_{0,1}}
       (\widehat h_\star,o_{\rightarrow},0,
        \operatorname{id}_{\widetilde q_{0,1}},\varnothing,
        \jmath_{k_0}(0_{\mathsf M_{k_0}}));
 \varnothing,\varnothing;
 (a_{0,1}),(+1);\operatorname{id}_{\widetilde q_{0,1}};
 k_0;0_{\mathsf M_{k_0}};\top;\varnothing\bigr).
\end{aligned}
\label{eq:estado-inicial-calibrado-ch}
\end{equation}
La tupla no contiene coordenadas indeterminadas y satisface
\(s(x_\star)=\top\). Sea \(G_0:=\{x_\star\}\).
Supuesto construido \(G_n\), fijemos \(x\in G_n\),
\(\varepsilon\in\{-1,0,+1\}\) y
\(x_{\varepsilon,0}(x):=x\). Para \(r=1,\ldots,9\), se construye primero
el registro elemental
\begin{equation}
\begin{aligned}
 \ell_{\varepsilon,n,r}:=\bigl(&e_{\varepsilon,n,r},a_{n,r},a_{(n,r)^+},
 \widetilde q_{n,r},\widetilde q_{n,r+1},M_{\rm ph}(r),\\
 &h_r(h(x_{\varepsilon,r-1}(x))),o(M_{\rm ph}(r)),\mathfrak p_\varepsilon,
 d_{\varepsilon,r},\\
 &\mathsf C_{\mathsf M}(e_{\varepsilon,n,r}),
 \kappa_{\mathsf M}(e_{\varepsilon,n,r}),
 \operatorname{Carr}(e_{\varepsilon,n,r})\bigr),
\end{aligned}
 \label{eq:registro-elemental-ext-ch}
\end{equation}
y se pone
\(\operatorname{Led}^{\rm new}_{\varepsilon,n,r}:=
\operatorname{Led}(x_{\varepsilon,r-1}(x))^\frown
\ell_{\varepsilon,n,r}\). Se construye entonces la tupla
\begin{equation}
 y_{\varepsilon,r}(x):=
 x_{\varepsilon,r-1}(x)
 [\tau\leftarrow M_{\rm ph}(r),\
  h\leftarrow h_r(h(x_{\varepsilon,r-1}(x))),\
  o\leftarrow o(M_{\rm ph}(r)),\
  b_{\rm ref}\leftarrow\mathfrak p_\varepsilon],
 \label{eq:tupla-sel-cal-ch}
\end{equation}
donde la selección deja intactos célula, cocientes, acarreo, ruta, memoria,
cilindro, frontera, supervivencia y registro; la coordenada
\(b_{\rm ref}\) hace explícita la reducción balanceada escogida. No se presupone aquí
pertenencia a una relación ambiental.

A continuación se define \(z_{\varepsilon,r}(x)\) campo por campo: su
configuración es \(\widetilde q_{n,r+1}\); su célula completa es
\(a_{(n,r)^+}\), con los
cuatro residuos y cocientes recalculados por las dos divisiones euclídeas;
conserva la hoja seleccionada y la orientación; y efectúa las seis
actualizaciones siguientes, en las que sólo se omite el argumento común \(x\):
\begin{equation}
\begin{aligned}
 C(z_{\varepsilon,r})&:=C(y_{\varepsilon,r})^\frown
       e_{\varepsilon,n,r},\\
 \partial C(z_{\varepsilon,r})&:=\partial C(y_{\varepsilon,r})^\frown
       (\widetilde q_{n,r},\widetilde q_{n,r+1}),\\
 p_{\APP}(z_{\varepsilon,r})&:=p_{\APP}(y_{\varepsilon,r})^\frown
       a_{(n,r)^+},\\
 p_{\TRIT}(z_{\varepsilon,r})&:=p_{\TRIT}(y_{\varepsilon,r})^\frown
       \tau(r),\\
 \operatorname{Led}(z_{\varepsilon,r})&:=
       \operatorname{Led}^{\rm new}_{\varepsilon,n,r},\\
 \lambda(z_{\varepsilon,r})&:=e_{\varepsilon,n,r}
       \lambda(y_{\varepsilon,r}),
\end{aligned}
 \label{eq:tupla-tra-cal-ch}
\end{equation}
y mantiene todavía los campos
\(c,\operatorname{gr}_{\mathsf M},\operatorname{mem},s\) de
\(y_{\varepsilon,r}(x)\). Su firma, en cambio, se vuelve a tipar en la
configuración de destino mediante la prefirma transportada
\begin{equation}
\begin{aligned}
 \operatorname{Sig}(z_{\varepsilon,r}(x))
  :=\operatorname{Sig}^{\rm cal}_{\widetilde q_{n,r+1}}
       \Bigl(&\widehat h(z_{\varepsilon,r}(x)),
             o(z_{\varepsilon,r}(x)),c(y_{\varepsilon,r}(x)),
             \lambda(z_{\varepsilon,r}(x)),\\[-2pt]
             &\partial C(z_{\varepsilon,r}(x)),
             \jmath_{d_{n,r}}
                (\operatorname{mem}_{d_{n,r}}
                   (y_{\varepsilon,r}(x)))\Bigr).
\end{aligned}
\label{eq:prefirma-transporte-cal-ch}
\end{equation}
Así, y sólo en
\(\mathsf{Tra}^{\rm cal}\), se añaden la arista, sus extremos, la marca
\((\mathfrak p_\varepsilon,r,d_{\varepsilon,r})\) y sus cociclos al registro de
incidencias.

El estado inicial satisface
\[
 c(x_\star)=0=\operatorname{Carr}
   (\operatorname{id}_{\widetilde q_{0,1}})
   =\operatorname{Carr}(\lambda(x_\star)).
\]
En los estados completos \(x_{\varepsilon,r-1}(x)\), y también después de la
selección \(y_{\varepsilon,r}(x)\), se conserva el invariante tipado
\(c(v)=\operatorname{Carr}(\lambda(v))\), pues
\(\mathsf{Sel}^{\rm cal}\) no modifica ninguno de esos dos campos. El
transporte avanza la ruta pero deja pendiente el acarreo; por ello la tupla de
preactualización satisface exactamente
\[
 c(z_{\varepsilon,r}(x))=c(y_{\varepsilon,r}(x))
 =\operatorname{Carr}(\lambda(y_{\varepsilon,r}(x))).
\]
La aplicación \(\operatorname{Upd}^{\rm cal}\) restaura entonces
\(c(x_{\varepsilon,r}(x))=
\operatorname{Carr}(\lambda(x_{\varepsilon,r}(x)))\) mediante la ley de
composición, sin atribuir ese invariante al estado intermedio \(z\).

Finalmente se aplica la función unaria
\(\operatorname{Upd}^{\rm cal}_{e_{\varepsilon,n,r}}
=\mathsf{Mem}^{\rm cal}_{e_{\varepsilon,n,r}}\) a esa tupla. Su
resultado \(x_{\varepsilon,r}(x)\) viene dado por
\begin{equation}
\begin{aligned}
 c\bigl(x_{\varepsilon,r}(x)\bigr)
   &=\operatorname{Carr}\bigl(\lambda(z_{\varepsilon,r}(x))\bigr)\\
   &=H_{\rm Carr}(\lambda(y_{\varepsilon,r}(x)))
       \bigl(\operatorname{Carr}(e_{\varepsilon,n,r})\bigr)
     +Y_{\rm Carr}(e_{\varepsilon,n,r})
       \bigl(c(y_{\varepsilon,r}(x))\bigr),\\
 \lambda\bigl(x_{\varepsilon,r}(x)\bigr)
   &=\lambda(z_{\varepsilon,r}(x)),\\
 \operatorname{gr}_{\mathsf M}
      \bigl(x_{\varepsilon,r}(x)\bigr)
   &=d_{n,r}+1,\\
 \operatorname{mem}_{d_{n,r}+1}
       \bigl(x_{\varepsilon,r}(x)\bigr)
   &=\mathsf C_{\mathsf M}(e_{\varepsilon,n,r})
       \operatorname{mem}_{d_{n,r}}(z_{\varepsilon,r}(x))
     +\kappa_{\mathsf M}(e_{\varepsilon,n,r}),\\
 s\bigl(x_{\varepsilon,r}(x)\bigr)
   &=s(e_{\varepsilon,n,r})\wedge s(z_{\varepsilon,r}(x)),\\
 \operatorname{Sig}\bigl(x_{\varepsilon,r}(x)\bigr)
   &=\operatorname{Sig}^{\rm cal}_{\widetilde q_{n,r+1}}
       \Bigl(\widehat h(z_{\varepsilon,r}(x)),
             o(z_{\varepsilon,r}(x)),
             c(x_{\varepsilon,r}(x)),
             \lambda(z_{\varepsilon,r}(x)),
             \partial C(z_{\varepsilon,r}(x)),\\[-2pt]
   &\hspace{112pt}
             \jmath_{d_{n,r}+1}
              \bigl(\operatorname{mem}_{d_{n,r}+1}
                 (x_{\varepsilon,r}(x))\bigr)\Bigr).
\end{aligned}
\label{eq:upd-unaria-tres-vueltas-ch}
\end{equation}
Todos los campos no enumerados en
\eqref{eq:upd-unaria-tres-vueltas-ch} coinciden con los de
\(z_{\varepsilon,r}(x)\). En particular, la firma se calcula con la hoja,
orientación, frontera y ruta ya fijadas por \(\mathsf{Sel}^{\rm cal}\) y
\(\mathsf{Tra}^{\rm cal}\), y con el acarreo y la memoria calculados en las
líneas anteriores; no se define circularmente a partir del estado final.
Ni \(\varepsilon\), ni \(d_{\varepsilon,r}\), ni \(u\) son argumentos
adicionales de \(\operatorname{Upd}^{\rm cal}\): la función lee los tres datos de la
arista ya construida.

\paragraph{Inducción simultánea de los niveles.}
Completadas las nueve etapas anteriores para un \(G_n\) ya construido,
se define inmediatamente
\begin{equation}
 \gamma_{\varepsilon,k_n}(x):=x_{\varepsilon,9}(x),\qquad
 G_{n+1}:=\{\gamma_{\varepsilon,k_n}(x):
       x\in G_n,\ \varepsilon\in\TRIT\}.
 \label{eq:injerto-total-tres-vueltas-ch}
\end{equation}
\label{eq:tres-vueltas-tpk-ch}
La base \(G_0=\{x_\star\}\), las tuplas intermedias y
\eqref{eq:injerto-total-tres-vueltas-ch} constituyen una única definición por
recursión sobre \(n\). Por tanto, ningún portador posterior se utiliza para
suponer la existencia de \(G_n\): el nivel siguiente sólo se introduce
después de construir todas sus aristas desde el nivel precedente.

\paragraph{Relación de extensión TPK y pertenencia efectiva.}
Sea \(\mathcal E^{\rm cal}_{\widetilde q}\) la menor familia etiquetada que contiene
\(x_\star\) y, en cada paso de la recurrencia anterior, las cuatro tuplas
\(x_{\varepsilon,r-1}(x),y_{\varepsilon,r}(x),
z_{\varepsilon,r}(x),x_{\varepsilon,r}(x)\) en sus configuraciones de origen
y destino. Sea asimismo
\(\mathcal A^{\rm cal}_{\widetilde q_{n,r}}:=\{a_{n,r}\}\), y póngase
\[
 \widetilde{\mathcal Q}_{\rm obs}:=
   \{\widetilde q_{n,r}:n\ge0,\ 1\le r\le9\},
 \qquad
 \mathcal E^{\rm cal}:=\bigsqcup_{\widetilde q}
     \mathcal E^{\rm cal}_{\widetilde q},
 \qquad
 \mathcal F^{\rm cal}_{\widetilde q}:=
     \pi_{\mathcal F}(\mathcal E^{\rm cal}_{\widetilde q}).
\]
Para cada profundidad \(d\), sea
\(\mathcal E^{\rm cal}_d:=
\operatorname{gr}_{\mathsf M}^{-1}(d)\subset\mathcal E^{\rm cal}\). La
coordenada explícita de grado, y no la mera pertenencia de un representante a
alguna imagen del sistema directo, hace disjuntos estos portadores. Sobre ellos
quedan tipadas las proyecciones
\[
\begin{aligned}
 \operatorname{cfg}&:\mathcal E^{\rm cal}\longrightarrow
     \widetilde{\mathcal Q}_{\rm obs},\\
 g(x)&:=\operatorname{fase}\!
       \bigl(\pi_{\rm obs}(\operatorname{cfg}(x))\bigr)\in C_9,\\
 c&:\mathcal E^{\rm cal}\longrightarrow K_{\rm Carr},\\
 s&:\mathcal E^{\rm cal}\longrightarrow\{\bot,\top\},\\
 \lambda&:\mathcal E^{\rm cal}_{\widetilde q}
       \longrightarrow\mathsf{Hist}^{\rm cal}_{\widetilde q},\\
 \operatorname{gr}_{\mathsf M}&:\mathcal E^{\rm cal}
       \longrightarrow\mathbb N_{\ge1},\\
 \operatorname{mem}_d&:\mathcal E^{\rm cal}_d
       \longrightarrow\mathsf M_d,\\
 \operatorname{mem}(x)&:=
       \jmath_d\bigl(\operatorname{mem}_d(x)\bigr)\in\mathsf M
       \quad(x\text{ de profundidad }d).
\end{aligned}
\]
Aquí \(\operatorname{cfg}(x)\) es la segunda coordenada de la tupla y
\(\operatorname{fase}\) lee la fase del calendario observable. Para las tres
clases de portador se fijan las diagonales
\[
\begin{aligned}
 \Delta_{\APP,\widetilde q}
   &:=\{(a,a):a\in\mathcal A^{\rm cal}_{\widetilde q}\},\\
 \Delta_{{\rm fib},\widetilde q}
   &:=\{(f,f):f\in\mathcal F^{\rm cal}_{\widetilde q}\},\\
 \Delta_{{\rm enr},\widetilde q}
   &:=\{(x,x):x\in\mathcal E^{\rm cal}_{\widetilde q}\}.
\end{aligned}
\]
En particular,
\[
 a_{0,1}\in\mathcal A^{\rm cal}_{\widetilde q_{0,1}},
 \qquad
 x_\star\in\mathcal E^{\rm cal}_{\widetilde q_{0,1}}.
\]
Sobre estos portadores se
definen por sus grafos completos, y no mediante implicaciones necesarias, las
relaciones
\begin{align}
 \mathsf{Sel}^{\rm cal}_{e_{\varepsilon,n,r}}
   &:=\{(x_{\varepsilon,r-1}(x),y_{\varepsilon,r}(x)):
          x\in G_n\},\\
 \mathsf{Tra}^{\rm cal}_{e_{\varepsilon,n,r}}
   &:=\{(y_{\varepsilon,r}(x),z_{\varepsilon,r}(x)):
          x\in G_n\},\\
 \operatorname{Ext}^{\APP,\rm cal}_{e_{\varepsilon,n,r}}
   &:=\{(a_{n,r},a_{(n,r)^+})\},
 \label{eq:ext-app-calibrada-ch}\\
 \operatorname{Ext}^{\rm fib,cal}_{e_{\varepsilon,n,r}}
   &:=\{(\pi_{\mathcal F}x_{\varepsilon,r-1}(x),
          \pi_{\mathcal F}x_{\varepsilon,r}(x)):x\in G_n\}.
 \label{eq:ext-fibra-calibrada-ch}
\end{align}
La extensión enriquecida sincronizada es
\begin{equation}
 \operatorname{Ext}^{\rm enr,cal}_{e_{\varepsilon,n,r}}
 :=\{(x_{\varepsilon,r-1}(x),x_{\varepsilon,r}(x)):
       x\in G_n\}.
 \label{eq:ext-enriquecida-calibrada-ch}
\end{equation}
La tercera aplicación queda definida por
\begin{equation}
 \operatorname{graph}
 \bigl(\operatorname{Upd}^{\rm cal}_{e_{\varepsilon,n,r}}\bigr)
 :=\{(z_{\varepsilon,r}(x),x_{\varepsilon,r}(x)):
       x\in G_n\},
 \label{eq:grafo-upd-calibrado-ch}
\end{equation}
donde el segundo componente es exactamente el de
\eqref{eq:upd-unaria-tres-vueltas-ch}. Por consiguiente,
\begin{equation}
 \mathcal U^{(1),\rm cal}_{e_{\varepsilon,n,r}}
 :=\operatorname{Upd}^{\rm cal}_{e_{\varepsilon,n,r}}\circ
   \mathsf{Tra}^{\rm cal}_{e_{\varepsilon,n,r}}\circ
   \mathsf{Sel}^{\rm cal}_{e_{\varepsilon,n,r}}
 =\{(x_{\varepsilon,r-1}(x),x_{\varepsilon,r}(x)):
      x\in G_n\}.
 \label{eq:transicion-cal-en-tpk-ch}
\end{equation}
Cada operador está indexado por la arista
\(e_{\varepsilon,n,r}\) producida por la reducción
\(\mathfrak p_\varepsilon\), y no por un trit suministrado desde fuera. La
coordenada \(b_{\rm ref}\) distingue además sus codominios. Por ello
\(\mathsf{Sel}^{\rm cal}_{e}\),
\(\mathsf{Tra}^{\rm cal}\) y \(\operatorname{Upd}^{\rm cal}\) son funciones
sobre sus respectivos dominios, mientras
la unión
\(\bigcup_{\varepsilon\in\TRIT}
\mathcal U^{(1),\rm cal}_{e_{\varepsilon,n,r}}\) conserva deliberadamente
las tres salidas de la trisección.

Para \(\bullet\in\{\APP,{\rm fib},{\rm enr}\}\), las identidades y las
composiciones no quedan implícitas. Se definen
\begin{align}
 \operatorname{Ext}^{\bullet,\rm cal}_{\operatorname{id}_{\widetilde q}}
   &:=\Delta_{\bullet,\widetilde q},\\
 \operatorname{Ext}^{\bullet,\rm cal}_{e_t\cdots e_1}
   &:=\operatorname{Ext}^{\bullet,\rm cal}_{e_t}\circ\cdots\circ
      \operatorname{Ext}^{\bullet,\rm cal}_{e_1}
 \label{eq:ext-palabra-calibrada-ch}
\end{align}
para toda palabra admisible de aristas. Las dos clausuras de caminos APP y
fibra se forman por separado,
\begin{equation}
 \operatorname{Path}(\operatorname{Ext}^{\APP,\rm cal}),\qquad
 \operatorname{Path}(\operatorname{Ext}^{\rm fib,cal}),
 \label{eq:dos-clausuras-ext-ch}
\end{equation}
y \(\mathsf{Tra}^{\rm cal}\) las sincroniza arista por arista.

Finalmente, los niveles de estados completos y sus truncamientos se fijan
por la propia recurrencia:
\begin{align}
 H^{\rm cal}_{d_{n,r}}
   &:=\{x_{\varepsilon,r-1}(x):
          x\in G_n,\ \varepsilon\in\TRIT\},\\
 H^{\rm cal}_{d_{n,r}+1}
   &:=\{x_{\varepsilon,r}(x):
          x\in G_n,\ \varepsilon\in\TRIT\},\\
 \rho^{\rm cal}_{d_{n,r}+1,d_{n,r}}
   \bigl(x_{\varepsilon,r}(x)\bigr)
   &:=x_{\varepsilon,r-1}(x).
 \label{eq:truncamiento-calibrado-definido-ch}
\end{align}
El último registro conserva \((n,r,\mathfrak p_\varepsilon)\), de modo que
el predecesor del miembro izquierdo es único y \(\rho^{\rm cal}\) está bien
definida. Para \(b>a\), se pone
\(\rho^{\rm cal}_{b,a}:=\rho^{\rm cal}_{a+1,a}\circ\cdots\circ
\rho^{\rm cal}_{b,b-1}\). Esta aplicación recupera el estado anterior
completo, incluida su firma; no pretende invertir una coordenada sobrescrita
sin consultar el registro.

\begin{proposition}[Realización autónoma del TPK mediante transiciones explícitas]
\label{prop:modelo-tpk-calibrado-ch}
La familia
\[
\begin{aligned}
 \mathfrak T_{\rm TPK}^{\rm cal}:=\bigl(
 &(\widetilde q_{n,r})_{n,r},\pi_{\rm obs},
   \widetilde F_{\rm obs},M_{\rm ph},\\
 &(\mathcal A_{\widetilde q}^{\rm cal})_{\widetilde q},
   (\mathcal E_{\widetilde q}^{\rm cal})_{\widetilde q},
   (\mathcal F_{\widetilde q}^{\rm cal})_{\widetilde q},
   \operatorname{Ext}^{\APP,\rm cal},
   \operatorname{Ext}^{\rm fib,cal},
   \operatorname{Ext}^{\rm enr,cal},\rho^{\rm cal},\\
 &\mathsf{Sel}^{\rm cal},\mathsf{Tra}^{\rm cal},
   \operatorname{Upd}^{\rm cal},\mathsf C_{\mathsf M},
   \kappa_{\mathsf M},\operatorname{gr}_{\mathsf M},
   H_{\rm Carr},Y_{\rm Carr},\\
 &\operatorname{Sig}^{\rm cal},\pi_{\mathcal F}\bigr).
\end{aligned}
\]
es un modelo enriquecido del TPK. En particular, cada par del miembro derecho
de \eqref{eq:transicion-cal-en-tpk-ch} es una arista TPK efectiva, y no la
consecuencia inferida de una condición solamente necesaria.
\end{proposition}

\begin{proof}
Las etiquetas \((n,r)\) hacen disjuntas las copias de los portadores y fijan
sin ambigüedad cada origen y destino. La relación
\(\operatorname{Ext}^{\APP,\rm cal}\) transporta la célula completa y
preserva las dos identidades euclídeas; la relación
\(\operatorname{Ext}^{\rm fib,cal}\) relaciona las fibras de los estados
completos anterior y posterior: transporta hoja, orientación, prefijo,
cilindro, frontera y registro e incorpora el acarreo, la memoria y la firma
actualizados. La selección no modifica el registro y el transporte concatena
exactamente una entrada y produce la prefirma de destino
\eqref{eq:prefirma-transporte-cal-ch}. La actualización recalcula la firma mediante la
aplicación tipada \(\operatorname{Sig}^{\rm cal}_{\widetilde q}\) desde
hoja, orientación, acarreo, ruta, frontera y memoria ya determinados. La aplicación
\(\operatorname{Upd}^{\rm cal}\) es unaria porque la arista contenida en
\(z_{\varepsilon,r}(x)\) determina sus cociclos.

Las identidades son las diagonales y las extensiones por palabras son las
composiciones de \eqref{eq:ext-palabra-calibrada-ch}. La concatenación de
prefijos y registros es asociativa; la memoria satisface la ley de cociclo y
el acarreo satisface
\eqref{eq:ley-carry-correcta-ch}. Por tanto, las dos parentizaciones de tres
aristas producen el mismo estado. La fórmula
\eqref{eq:truncamiento-calibrado-definido-ch} define el truncamiento sobre
los estados completos y la unicidad del último registro prueba que recupera
el origen en todos los campos. Éstas son las leyes de identidad, fuente,
destino, composición, memoria, acarreo y truncamiento del TPK enriquecido.
En particular, las relaciones de extensión quedan determinadas como grafos
de aplicaciones efectivas y no como consecuencias de meras condiciones
necesarias.
\end{proof}

Cada \(\gamma_{\varepsilon,k_n}\) comprende exactamente nueve aplicaciones
de \(\operatorname{Upd}^{\rm cal}_{e_{\varepsilon,n,r}}\circ
\mathsf{Tra}^{\rm cal}_{e_{\varepsilon,n,r}}\circ
\mathsf{Sel}^{\rm cal}_{e_{\varepsilon,n,r}}\). La
supervivencia se probará mediante una cola construida y no se utilizará para
definir \(G_n\).

\begin{lemma}[Trisección interna efectiva de una vuelta]
\label{lem:elevacion-catalogal-nonadica}
Para todo \(n\ge0\), todo \(x\in G_n\) y todo
\(\varepsilon\in\{-1,0,+1\}\), la expresión
\(\gamma_{\varepsilon,k_n}(x)\) está definida por nueve transiciones TPK,
con sus relaciones sincronizadas
\(\operatorname{Ext}^{\APP,\rm cal}\) y
\(\operatorname{Ext}^{\rm fib,cal}\), y satisface
\begin{align}
 \rho^{\rm cal}_{k_n+9,k_n}\gamma_{\varepsilon,k_n}(x)&=x,
 \label{eq:truncamiento-injerto-ch}\\
 g\bigl(\gamma_{\varepsilon,k_n}(x)\bigr)
   &=g(x),&
 \jmath_{k_n+9}\!\left(\operatorname{mem}_{k_n+9}
        (\gamma_{\varepsilon,k_n}(x))\right)
   &=\jmath_{k_n}\!\left(\operatorname{mem}_{k_n}(x)\right)+u,
 \label{eq:fase-memoria-injerto-ch}\\
 c\bigl(\gamma_{\varepsilon,k_n}(x)\bigr)
   &=c(x)+\varepsilon\kappa_{\rm Carr}.&&
 \label{eq:carry-injerto-ch}
\end{align}
\label{eq:tres-vueltas-propiedades-ch}
\label{eq:tres-vueltas-acarreo-ch}
El registro ordenado de las nueve aristas determina \(\varepsilon\) de manera única.
Las tres imágenes son, por tanto, disjuntas. Cada extremo pertenece a
\(G_{n+1}\) y vuelve a admitir las tres aplicaciones
\(\gamma_{-1,k_{n+1}},\gamma_{0,k_{n+1}},
\gamma_{+1,k_{n+1}}\).
\end{lemma}

\begin{proof}
Las fórmulas \eqref{eq:celula-completa-fase-ch} recorren una vez las nueve
fases y devuelven la copia \(a_{n+1,1}\) de \(a(1)\). En cada paso, la relación APP transporta el
séxtuplo completo; \(\mathsf{Sel}^{\rm cal}\) fija el tipo TRIT desde la marca de fase
\(r\) conservada en el estado;
y \(\mathsf{Tra}^{\rm cal}\) añade exactamente un símbolo al prefijo, un tramo al
cilindro y su par ordenado de extremos a la frontera. Por inducción sobre
\(r\), las tuplas explícitas \(y_{\varepsilon,r}(x)\) y
\(z_{\varepsilon,r}(x)\) satisfacen los predicados de las relaciones
indicadas, y la actualización unaria produce el estado del nivel siguiente.
La identidad
\(\rho^{\rm cal}_{d_{n,r}+1,d_{n,r}}
(x_{\varepsilon,r}(x))=x_{\varepsilon,r-1}(x)\) es la definición
\eqref{eq:truncamiento-calibrado-definido-ch}; su buena definición procede de
la unicidad del último registro. La composición de las nueve aplicaciones de
truncamiento recupera \(x\), lo que prueba
\eqref{eq:truncamiento-injerto-ch} sin invertir por separado los campos
sobrescritos.

La fase \(9\to1\) aparece una sola vez. Por
\eqref{eq:memoria-arista-ch} y la ley de cociclo, la parte lineal de memoria
transporta la memoria previa y su parte traslacional suma exactamente \(u\)
en el límite directo. La compatibilidad
\(\jmath_{d+1}\jmath_{d+1,d}=\jmath_d\) prueba
\eqref{eq:fase-memoria-injerto-ch}; el retorno concierne a la fase, no al
estado enriquecido.

Por \eqref{eq:carry-arista-derivado-ch}, ocho aristas poseen incremento nulo y
la arista que completa \(\mathfrak p_\varepsilon\) posee el incremento calculado
\(\chi(\varepsilon,\varepsilon)\kappa_{\rm Carr}\). La ley
\eqref{eq:ley-carry-correcta-ch} y
\eqref{eq:carry-no-parametrico-ch} dan
\eqref{eq:carry-injerto-ch}. El registro conserva
\((\mathfrak p_\varepsilon,2\varepsilon,r_3(2\varepsilon),
\chi(\varepsilon,\varepsilon))\); los tres valores de este registro son
distintos incluso cuando el incremento escalar es cero. Por ello las imágenes
no se identifican y el decodificador de bloque recupera una única
\(\varepsilon\).

Las definiciones de las aristas dependen sólo de la nueva fase y concatenan,
en vez de sustituir, prefijo, cilindro, frontera y registro. También son
invariantes bajo
\(\jmath_d\operatorname{mem}_d\mapsto
  \jmath_d\operatorname{mem}_d+nu\)
y bajo la traslación del acarreo de
entrada. En consecuencia, la misma construcción se aplica al extremo de
cualquier vuelta. Para cada estado finito, la sucesión explícita
\[
 x,\quad\gamma_{0,k_n}(x),\quad
 \gamma_{0,k_{n+1}}\gamma_{0,k_n}(x),\quad\ldots
\]
es una cola infinita compatible. Esto demuestra la supervivencia de todos los
estados construidos y la disponibilidad reiterada de las tres ramas, sin
postular ninguna de ambas propiedades. Por consiguiente podemos, finalmente,
identificar, de acuerdo con las definiciones de nivel anteriores,
\begin{equation}
 H_{k_n}^{\rm cal}:=G_n
 \label{eq:subarbol-calibrado-superviviente-ch}
\end{equation}
para todo \(n\ge0\).
\end{proof}

El cierre anterior distingue los tres relojes implicados. Nueve emisiones
devuelven la fase en \(C_9\) y avanzan la memoria. Doce vueltas, es decir,
\(108\) emisiones, devuelven además la posición del calendario observable;
tampoco entonces se borran el registro ni la memoria. El operador
\(\mathcal K_{\rm ph}\) actúa únicamente después, como reconocimiento de
memoria reducida de las vueltas ya construidas; nunca genera las aristas de
\(\operatorname{Ext}^{\rm enr,cal}\).



Sea
\[
 \upsilon:\mathbb F_3\longrightarrow\{-1,0,+1\},\qquad
 \upsilon(0)=0,\quad\upsilon(1)=+1,\quad\upsilon(2)=-1
\]
el representante balanceado. Para
\(b=(b_1,\ldots,b_6)\in\mathbb F_3^6\), la prolongación de bloque es la
composición ordenada de seis vueltas
\[
 \Gamma_{b,k}^{[54]}:=
 \gamma_{\upsilon(b_6),k+45}\circ\gamma_{\upsilon(b_5),k+36}\circ
 \gamma_{\upsilon(b_4),k+27}\circ\gamma_{\upsilon(b_3),k+18}\circ
 \gamma_{\upsilon(b_2),k+9}\circ\gamma_{\upsilon(b_1),k}.
\tag{C.39d}
\label{eq:gamma-bloque-54-ch}
\]
En el nivel de bloque se emplea siempre
\(k=k_{6N}=1+54N\); por ello cada factor tiene el dominio construido por el
factor anterior.
Fijemos una semilla APP canónica \(x_\star\in H^{\rm cal}_1\) y definamos
recursivamente
\[
 \begin{aligned}
  Y_0&:=\{x_\star\},\qquad Y_N\subseteq H_{1+54N}^{\rm cal},\\
  Y_{N+1}&:=\{\Gamma_{b,1+54N}^{[54]}(x):
                 x\in Y_N,\ b\in\mathbb F_3^6\},\\
  \rho^{\rm blk}_{N+1,N}
    &:=\rho^{\rm cal}_{1+54(N+1),1+54N}
       \!\upharpoonright Y_{N+1},\\
  \operatorname{Ext}^{\rm cal}_N(x)
    &:=\{\Gamma_{b,1+54N}^{[54]}(x):b\in\mathbb F_3^6\}.
 \end{aligned}
\tag{C.39e}
\label{eq:sincronizacion-54-ch}
\]
La notación \(\rho^{\rm cal}_{b,a}\) designa la composición de todas las supresiones
entre los niveles \(b\) y \(a\). Por construcción, cada prolongación de
\(\operatorname{Ext}^{\rm cal}_N(x)\) es una cadena de cincuenta y cuatro
aristas TPK, no un salto del cociente de Markov.

El lector se define directamente sobre el registro, antes de apelar a una
representación mediante \(\Gamma_b\). El lema precedente proporciona un
decodificador
\[
 \operatorname{dec}_9:\{\text{segmentos generados de nueve aristas}\}
 \longrightarrow\{-1,0,+1\},
\]
que lee la pareja TRIT y su reducción en el segmento. Si
\(\operatorname{Led}^{(9)}_{j,i}(x)\) denota el \(i\)-ésimo segmento nonádico
del bloque \(j\) del registro de \(x\), se fija
\begin{equation}
 \beta_{{\rm blk},N}(x):=
 \Bigl(
   \bigl(\upsilon^{-1}\operatorname{dec}_9
          (\operatorname{Led}^{(9)}_{j,i}(x))\bigr)_{i=1}^{6}
 \Bigr)_{j=0}^{N-1}.
 \label{eq:lector-bloques-ch}
\end{equation}
Para \(N=0\) el valor es la palabra vacía. Esta definición no presupone que
la descomposición de \(x\) como composición de aplicaciones \(\Gamma_b\) sea
única.

Especializamos ahora la conjugación aritmética precedente a
\(p=3\), \(d=6\) y \(A=I_6\). Si
\(\pi_{N+1,N}:(\mathbb Z/3^{N+1}\mathbb Z)^6\to
(\mathbb Z/3^N\mathbb Z)^6\) es la reducción canónica, definimos
\begin{equation}
 \Phi_N:=\Psi_N^{-1}\circ\beta_{{\rm blk},N}:
 Y_N\longrightarrow(\mathbb Z/3^N\mathbb Z)^6.
 \label{eq:phi-hensel-bloques-ch}
\end{equation}
Esta definición no asigna una coordenada después de observar una rama: el
lector \(\beta_{{\rm blk},N}\) se obtiene del registro TPK ya construido y
\(\Psi_N^{-1}\) es la inversa explícita, fijada con anterioridad, de la torre
aritmética no filtrada.

\begin{proposition}[Conjugación finita y cobertura completa de la fibra]
\label{prop:cobertura-729-seccion}
Para todo \(N\geq0\),
\[
 \beta_{{\rm blk},N}:Y_N\xrightarrow{\ \cong\ }
 (\mathbb F_3^6)^N
\]
es una biyección compatible con los truncamientos, y para cada \(x\in Y_N\)
\[
 \begin{gathered}
 \beta_{{\rm blk},N+1}
 \bigl[\operatorname{Ext}^{\rm cal}_N(x)\bigr]
 =\{\beta_{{\rm blk},N}(x)^\frown b:
       b\in\mathbb F_3^6\},\\
 \operatorname{pref}_N\circ\beta_{{\rm blk},N+1}
 =\beta_{{\rm blk},N}\circ\rho^{\rm blk}_{N+1,N}.
 \end{gathered}
\tag{C.40$_0$}
\label{eq:cobertura-729-ch}
\]
Las aplicaciones \(\Phi_N\) son biyecciones y satisfacen la identidad de
entrelazamiento
\begin{equation}
 \pi_{N+1,N}\circ\Phi_{N+1}
 =\Phi_N\circ\rho^{\rm blk}_{N+1,N}.
 \label{eq:entrelazamiento-hensel-ext-ch}
\end{equation}
Por consiguiente, la torre de historias generadas y la torre aritmética de
Hensel son isomorfas como sistemas inversos finitos.
\end{proposition}

\begin{proof}
La afirmación es inmediata para \(N=0\). Supongamos que vale en el nivel
\(N\). El lema anterior construye, para cada palabra
\(b\in\mathbb F_3^6\), una prolongación
    \(\Gamma_{b,1+54N}^{[54]}(x)\). Si \(b\ne b'\), la primera coordenada en la
que difieren determina una pareja TRIT y un cociclo distintos en el primer
segmento nonádico divergente del registro genealógico. Ese registro conserva
las aristas ordenadas, no sólo la suma final de sus acarreos; por tanto, las dos
historias son distintas incluso si acarreos de vueltas posteriores se
compensan. La lectura directa \eqref{eq:lector-bloques-ch} da entonces
\begin{equation}
 \beta_{{\rm blk},N+1}
 \bigl(\Gamma_{b,1+54N}^{[54]}(x)\bigr)
 =\beta_{{\rm blk},N}(x)^\frown b.
 \label{eq:beta-concatenacion-demostrada-ch}
\end{equation}
Las seis trisecciones sucesivas producen exactamente \(3^6\) prolongaciones
distintas. La supresión de los últimos cincuenta y cuatro pasos recupera
\(x\), y la supresión del último bloque del registro recupera
\(\beta_{{\rm blk},N}(x)\). Esto prueba simultáneamente la biyección y la
compatibilidad de \(\beta_{{\rm blk},N}\) con los truncamientos.

La proposición~\ref{prop:ch-hensel-bruto} prueba, independientemente, que
\(\Psi_N\) es biyectiva y que
\[
 \operatorname{pref}_N\circ\Psi_{N+1}
 =\Psi_N\circ\pi_{N+1,N}.
\]
Al componer esta igualdad con la compatibilidad de \(\beta_{\rm blk}\) y
aplicar \(\Psi_N^{-1}\), se obtiene
\[
 \pi_{N+1,N}\Psi_{N+1}^{-1}\beta_{{\rm blk},N+1}
 =\Psi_N^{-1}\beta_{{\rm blk},N}\rho^{\rm blk}_{N+1,N},
\]
que es exactamente \eqref{eq:entrelazamiento-hensel-ext-ch}. La igualdad no
procede de definir una coordenada ad hoc sobre cada rama: resulta de componer
dos sistemas de biyecciones construidos y compatibles a todas las
profundidades.
\end{proof}

El límite inverso
\[
 \Omega_\Gamma^{\rm cal}:=
 \varprojlim_N(Y_N,\rho^{\rm blk}_{N+1,N})
\tag{C.40a$_0$}
\label{eq:omega-gamma-catalogal-ch}
\]
es el límite de historias del modelo TPK construido arista por
arista en el lema precedente. Su proyección por \(\pi_{\rm obs}\) conserva
las configuraciones del calendario observable, pero la afirmación no depende
de identificar este modelo autónomo con una relación ambiental definida sólo
por condiciones necesarias. Las biyecciones finitas pasan
al límite y producen el homeomorfismo cilíndrico
\[
 \Lambda:(\mathbb F_3^6)^{\mathbb N}
 \xrightarrow{\ \cong\ }\Omega_\Gamma^{\rm cal},
 \qquad
 \beta_{\rm blk}:=\Lambda^{-1}.
\tag{C.40b$_0$}
\label{eq:homeomorfismo-calibrado-ch}
\]
La proyección canónica se denota
\(\rho^{\rm blk}_{\infty,N}:\Omega_\Gamma^{\rm cal}\to Y_N\).
Por tanto, \(\Lambda\) establece una equivalencia compatible con los
truncamientos: cada cinta ternaria codifica una historia del TPK y cada
prefijo se prolonga mediante aristas del mismo generador.

Los subcilindros se definen ahora como subconjuntos del límite, no como pares
de sufijos. Para \(x\in Y_N\) y \(b\in\mathbb F_3^6\), sean
\begin{align}
 \mathscr Z_N(x)
   &:=\{\xi\in\Omega_\Gamma^{\rm cal}:
       \rho^{\rm blk}_{\infty,N}(\xi)=x\},
 \label{eq:cilindro-real-limite-ch}\\
 \mathscr Z_{N+1}(x;b)
   &:=\{\xi\in\Omega_\Gamma^{\rm cal}:
       \rho^{\rm blk}_{\infty,N+1}(\xi)
       =\Gamma_{b,1+54N}^{[54]}(x)\}.
 \label{eq:subcilindro-real-limite-ch}
\end{align}

\begin{theorem}[Partición efectiva en \(729\) subcilindros de refinamiento]
\label{thm:particion-729-subcilindros-ch}
Para todo \(N\ge0\) y todo \(x\in Y_N\), se tiene la unión disjunta
\begin{equation}
 \mathscr Z_N(x)=
 \bigsqcup_{b\in\mathbb F_3^6}\mathscr Z_{N+1}(x;b).
 \label{eq:particion-real-729-ch}
\end{equation}
Los \(3^6=729\) términos son subcilindros de refinamiento no vacíos y
clopen; cada uno está realizado por cincuenta y cuatro aristas del TPK
construido.
\end{theorem}

\begin{proof}
Sea \(\xi\in\mathscr Z_N(x)\). Su coordenada de nivel \(N+1\) pertenece a
\(Y_{N+1}\), trunca a \(x\) y, por la biyección
\(\beta_{{\rm blk},N+1}\), tiene un último bloque único
\(b\in\mathbb F_3^6\). La identidad
\eqref{eq:beta-concatenacion-demostrada-ch} implica entonces
\(\rho^{\rm blk}_{\infty,N+1}(\xi)
=\Gamma_{b,1+54N}^{[54]}(x)\); por tanto, \(\xi\) pertenece a uno de los
términos del miembro derecho. La inclusión inversa es inmediata al truncar.
Si dos términos se intersecaran, sus coordenadas de nivel \(N+1\) serían
iguales y el decodificador del registro daría \(b=b'\), de modo que la unión
es disjunta.

Cada extremo \(\Gamma_b^{[54]}(x)\) admite una cola infinita obtenida
iterando la prolongación de bloque nulo; por ello todos los términos son no
vacíos. Como \(Y_{N+1}\) es finito y discreto, son preimágenes de puntos por
una proyección continua del límite inverso y, por tanto, clopen. La propia
definición de \(\Gamma_b^{[54]}\) acredita sus cincuenta y cuatro aristas.
\end{proof}

La familia \((\Phi_N)_N\) induce finalmente el homeomorfismo de límites
\begin{equation}
 \Phi_\infty:\Omega_\Gamma^{\rm cal}
 \xrightarrow{\ \cong\ }\mathbb Z_3^6,\qquad
 \Phi_\infty(\xi):=
 \bigl(\Phi_N(\rho^{\rm blk}_{\infty,N}(\xi))\bigr)_{N\ge0},
 \label{eq:phi-infinito-hensel-ext-ch}
\end{equation}
donde
\(\mathbb Z_3^6=\varprojlim_N(\mathbb Z/3^N\mathbb Z)^6\).
La ecuación \eqref{eq:entrelazamiento-hensel-ext-ch} garantiza que el miembro
derecho es una clase compatible, y las biyecciones finitas prueban
inyectividad, sobreyectividad y continuidad en ambos sentidos. Bajo
\(\Phi_\infty\), \(\mathscr Z_N(x)\) es exactamente la clase de congruencia
\[
 \{z\in\mathbb Z_3^6:z\equiv\Phi_N(x)\pmod{3^N}\},
\]
y los conjuntos \(\mathscr Z_{N+1}(x;b)\) son sus \(729\) clases de refinamiento
módulo \(3^{N+1}\). Éste es el empalme efectivo entre prolongación TPK,
cilindros del límite y torre de Hensel.

Para \(b=(b_1,\ldots,b_6)\in\mathbb F_3^6\), póngase
\[
 \nu(b):=\sum_{i=1}^{6}[b_i]_3\,3^{6-i}
 \in\{0,\ldots,728\},
\tag{C.40c$_0$}
\label{eq:valor-bloque-ch}
\]
donde \([b_i]_3\in\{0,1,2\}\) es el representante posicional. Las
\(243\) palabras visibles del censo inicial son la imagen de la primera
emisión; no se confunden con la capacidad \(3^6\) de cada bloque de
prolongación. Finalmente,
\[
 X_{\infty,\mathbb Z}^{\rm cal}
 :=\mathbb Z_{\rm HMT}\times\Omega_\Gamma^{\rm cal}.
\]



Para un prefijo \(s=(m;b_0,\ldots,b_{N-1})\), defínanse en la aritmética
generada de \(\mathbb Q_{\rm HMT}\)
\[
 a_N(s):=m+\sum_{q<N}\frac{\nu(b_q)}{729^{q+1}},
 \qquad
 \operatorname{Succ}(s):=
 \{(m;b_0,\ldots,b_{N-1},b):b\in\mathbb F_3^6\}.
\]
Si \((m;\beta_{{\rm blk},N}(x))\) denota el prefijo obtenido anteponiendo
\(m\) a la tupla \(\beta_{{\rm blk},N}(x)\),
\eqref{eq:cobertura-729-ch} da la igualdad tipada
\[
 \operatorname{Succ}(m;\beta_{{\rm blk},N}(x))
 =\{(m;\beta_{{\rm blk},N+1}(y)):
       y\in\operatorname{Ext}^{\rm cal}_N(x)\}.
\]
Los cilindros semiabiertos son, desde su primera definición, subconjuntos de
la completación HMT:
\[
 I_N^{\rm HMT}(s):=
 \{r\in\mathbb R_{\rm HMT}:
      a_N(s)\le_{\rm HMT}r<_{\rm HMT}a_N(s)+729^{-N}\}.
\]
Satisfacen materialmente
\[
 I_N^{\rm HMT}(s)=
 \bigsqcup_{y\in\operatorname{Succ}(s)}I_{N+1}^{\rm HMT}(y),
 \qquad |\operatorname{Succ}(s)|=729,
 \qquad \operatorname{diam}_{\rm HMT}I_N^{\rm HMT}(s)
      =729^{-N}\longrightarrow0.
\tag{C.40a}\label{eq:ch-particion-cilindros}
\]
La convención de frontera identifica las dos escrituras sólo después de
construirlas. La partición~\eqref{eq:ch-particion-cilindros}, el anidamiento y la contracción producen

\[
\begin{aligned}
 P_{\rm ar}&:X_{\infty,\mathbb Z}^{\rm cal}
     \twoheadrightarrow\mathbb R_{\rm HMT},\\
 P_{\rm ar}(m,\xi)
   &:=\text{el único elemento de }
     \bigcap_N C_N^{\rm cl}
       \bigl(m;\beta_{{\rm blk},N}(\rho^{\rm blk}_{\infty,N}(\xi))\bigr),\\
 C_N^{\rm cl}(s)&:=\overline{I_N^{\rm HMT}(s)}.
\end{aligned}
\tag{C.40}\label{eq:ch-proyeccion-arquimediana}
\]

Por tanto,

\[
 X_{\infty,\mathbb Z}^{\rm cal}/\!\sim_{P_{\rm ar}}
 \;\cong\;\mathbb R_{\rm HMT}.
\tag{C.41}\label{eq:ch-cociente-recta}
\]

La otra publicación \(P_{\rm exc}\) conserva incidencia, frontera y memoria
hacia Paley--Witt--Golay--Leech y, desde allí, hacia
\(V^\natural\)--Monster. Ambas proceden del
mismo antecedente de~\eqref{eq:ch-limite-inverso}; ni la coordenada arquimediana ni el reconocimiento
excepcional seleccionan las semillas. Las cinco construcciones
\(sp,sol,gau,coh,det\) permanecen consustanciales dentro de
\(\mathcal C_{\rm cont}^{\rm disc}\) y no se disgregan para efectuar el paso
cardinal.

\subsubsection{Descenso de operaciones al cociente de evaluación}
\label{sec:ch-descenso-fibras}

El cociente anterior identifica historias que publican el mismo valor, no
sus registros completos. El criterio siguiente caracteriza exactamente
cuándo una operación de historias induce una operación en ese cociente sin
escoger una historia privilegiada dentro de cada fibra; su formulación
continúa la construcción genealógica establecida en las secciones anteriores.

\begin{proposition}[Criterio de descenso por fibras de evaluación]
\label{prop:ch-descenso-fibras}
Sea \(q:\Omega\to J\) la evaluación sobre su imagen de una carta
genealógica. Sea \(\star:D\to\Omega\) una operación con
\(D\subseteq\Omega^2\) saturado por las fibras de \(q\times q\).
Existe una única operación
\(\overline\star:(q\times q)(D)\to J\) tal que
\[
 q(x\star y)=q(x)\,\overline\star\,q(y)
\]
si y sólo si, para pares del dominio,
\[
 q(x)=q(x'),\quad q(y)=q(y')
 \quad\Longrightarrow\quad q(x\star y)=q(x'\star y').
\]
\end{proposition}

\begin{proof}
Si existe la operación visible, su resultado depende sólo de los dos
valores, lo que prueba la necesidad. Para la suficiencia, se define
\(s\,\overline\star\,t=q(x\star y)\), donde \(q(x)=s\) y
\(q(y)=t\). La saturación conserva el dominio al cambiar representantes;
la condición enunciada conserva el resultado. La sobreyectividad de
\(q\times q\) sobre la imagen del dominio prueba la unicidad.
\end{proof}

Una coordenada no reemplaza automáticamente una historia como argumento de
todas sus operaciones: el criterio determina cuándo esa sustitución es
legítima. En particular, transportar una operación a través de una biyección
no basta para identificarla con una operación nativa anterior a la lectura.
Si una carta \(\Omega\) es compacta y \(q\) es continua y sobreyectiva
sobre un espacio Hausdorff \(J\), la evaluación induce además un
homeomorfismo \(\Omega/{\sim_q}\to J\): es una biyección continua de un
compacto a un Hausdorff. Es el \emph{cociente}, y no necesariamente el
espacio de historias, el que se identifica topológicamente con \(J\).



\subsubsection{Reconstrucción interna de la carta y levantamiento uniforme}

Dentro de \(\mathfrak G_{\rm HMT}(h,m_\infty)\), la reconstrucción de
\eqref{eq:ch-limite-inverso}--\eqref{eq:ch-cociente-recta} se realiza con los mismos códigos y operadores. Denotemos por
\(\mathbb R_{\rm HMT}^{\mathfrak G}\) la completación HMT interna. La
unicidad del cuerpo ordenado completo arquimediano proporciona, \textbf{después} de
construirlo, el isomorfismo interno
\[
 \eta:\mathbb R_{\rm HMT}^{\mathfrak G}
 \xrightarrow{\ \cong\ }\mathbb R^{\mathfrak G}.
\tag{C.41$_0$}\label{eq:ch-isomorfismo-interno}
\]
Para el intervalo interno correspondiente
\(I_N^{\mathfrak G}(s):=
 I_N^{\rm HMT}(s)\cap\mathbb R_{\rm HMT}^{\mathfrak G}\), el isomorfismo
preserva extremos y orden:
\[
 \eta[I_N^{\mathfrak G}(s)]
 =\{r\in\mathbb R^{\mathfrak G}:
      a_N(s)\le r<a_N(s)+729^{-N}\}.
\tag{C.41$_0$a}\label{eq:ch-intervalos-internos}
\]
La sobreyectividad de~\eqref{eq:ch-proyeccion-arquimediana} cubre todos los elementos de esa completación, no
sólo los obtenidos en una ejecución finita.

Para que «todos» no quede oculto en la definición del codominio, se construyen
las dos inyecciones que relacionan la recta con su potencia interna. Sea

\[
 E(A)=\sum_{n\ge0}\frac{2\mathbf1_A(n)}{3^{2n+2}}
 \qquad
 \bigl(A\in\mathcal P^{\mathfrak G}(\omega)\bigr).
\tag{C.41a}\label{eq:ch-incrustacion-cantor}
\]

La escritura coloca \(0\) o \(2\) únicamente en las posiciones ternarias
pares; su imagen está contenida en \([0,1/4]\). Por tanto nunca alcanza el
extremo derecho del cilindro raíz y dos subconjuntos distintos tienen una
primera posición par distinta: \(E\) es inyectiva y no sufre la ambigüedad de
frontera. Escribamos \(E_{\rm HMT}:=\eta^{-1}\circ E\).

Dado \(A\), se parte del cilindro raíz que contiene \(E_{\rm HMT}(A)\). Si
\(s_N(A)\) es el prefijo ya elegido, escríbase
\(s_N(A)=(0;b_0(A),\ldots,b_{N-1}(A))\). La partición~\eqref{eq:ch-particion-cilindros}, transportada
por \(\eta^{-1}\), entrega un único subcilindro semiabierto
\(s_{N+1}(A)\) que contiene \(E_{\rm HMT}(A)\). Por~\eqref{eq:cobertura-729-ch} esa prolongación es
también una prolongación TPK admisible. Por inducción,
\[
 s_{N+1}(A)\!\upharpoonright N=s_N(A),\qquad
 E_{\rm HMT}(A)\in I_N^{\rm HMT}(s_N(A))\quad(N<\omega).
\tag{C.41b$_0$}\label{eq:ch-prefijos-uniformes}
\]
La recursión es una fórmula uniforme con parámetros \(A,h\). Como ambos son
internos, Reemplazo en el modelo del lema~\ref{thm:ch-6b} forma la sucesión completa de
bloques \((b_q(A))_{q<\omega}\). Defínase, sin confundir la cinta con la
sucesión de prefijos,
\[
 \xi_A:=\Lambda\bigl((b_q(A))_{q<\omega}\bigr)
 \in\Omega_\Gamma^{\rm cal,\mathfrak G}.
\]
Se obtiene así el \textbf{levantamiento uniforme de la incrustación de Cantor}

\[
 Q:\mathcal P^{\mathfrak G}(\omega)
 \longrightarrow X_{\infty,\mathbb Z}^{\rm cal,\mathfrak G},
 \qquad
 Q(A):=(0,\xi_A),\qquad
 P_{\rm ar}(Q(A))=E_{\rm HMT}(A).
\tag{C.41b}\label{eq:ch-levantamiento-uniforme}
\]

No se llama a \(Q\) inverso de \(P_{\rm ar}\): su composición publica la
incrustación \(E_{\rm HMT}\), no la identidad en toda la recta. En la otra
dirección, sea \((q_n)_{n<\omega}\) la enumeración generada de
\(\mathbb Q_{\rm HMT}\) y defínase
\[
 C(r):=\{n<\omega:\eta(q_n)<r\}
 \qquad(r\in\mathbb R^{\mathfrak G}).
\]
El corte racional pertenece a \(\mathcal P^{\mathfrak G}(\omega)\) y
\(C\) es inyectivo por densidad. Las inyecciones \(E\) y \(C\), seguidas de
Cantor--Bernstein interno, prueban

\[
 P_{\rm ar}
 \bigl[X_{\infty,\mathbb Z}^{\rm cal,\mathfrak G}\bigr]
 =\mathbb R_{\rm HMT}^{\mathfrak G},
 \qquad
 \mathbb R_{\rm HMT}^{\mathfrak G}
 \cong\mathbb R^{\mathfrak G}
 \approx\mathcal P^{\mathfrak G}(\omega).
\tag{C.41c}\label{eq:ch-recta-potencia}
\]

Aquí un real arbitrario se usa únicamente como variable universal para probar
sobreyectividad de un árbol ya construido. No selecciona APP, TRIT, TPK, las
constantes ni las transiciones de ese árbol.

\subsection{Enumeración por códigos genealógicos y condensación HMT}

La regla~\eqref{eq:ch-operador-def} entrega nombres de derivación a partir de

\[
 (\ulcorner\varphi\urcorner,
  \sigma(p_0),\ldots,\sigma(p_{k-1})).
\tag{C.42}\label{eq:ch-codigo-nombres}
\]

Para cualquier objeto \(a\) de la jerarquía, sean \(\beta(a)\) su primera
etapa y \(\sigma(a)\) su menor árbol de derivación en el orden recursivo del
lema~\ref{thm:ch-6b}. No se exige que \(\sigma(a)\) sea un natural cuando la etapa es no
numerable. La relación

\[
 a\prec_{\mathfrak G}b
 \quad\Longleftrightarrow\quad
 (\beta(a),\sigma(a))<_{\rm lex}
 (\beta(b),\sigma(b))
\tag{C.42a}\label{eq:ch-orden-genealogico}
\]

es el buen orden genealógico global \(<_{\mathfrak G}\) inducido por el
constructor. Los nombres repetidos se eliminan tomando el menor; ningún real
ni ninguna enumeración de la recta se inserta como dato.

\begin{lemma}[Numerabilidad interna de los niveles anteriores a \(\omega_1^{\mathfrak G}\)]
\label{thm:ch-7}
Para todo
\(\beta<\omega_1^{\mathfrak G}\), el nivel
\(\mathfrak G_\beta(h,m_\infty)\) es numerable dentro de
\(\mathfrak G_{\rm HMT}(h,m_\infty)\), y posee una sobreyección canónica

\[
 e_\beta:\omega\twoheadrightarrow
 \mathfrak G_\beta(h,m_\infty).
\tag{C.43}\label{eq:ch-enumeracion-niveles}
\]

\end{lemma}
\begin{proof}
El nivel inicial tiene un código numerable. Si \(e_\beta\)
enumera el nivel \(\beta\), los códigos de fórmula y las tuplas finitas de
parámetros de~\eqref{eq:ch-codigo-nombres} forman un conjunto numerable y enumeran el sucesor. Si
\(\lambda<\omega_1^{\mathfrak G}\) es límite, existe internamente alguna
sobreyeción \(s:\omega\twoheadrightarrow\lambda\); se toma la
\(<_{\mathfrak G}\)-menor, orden que ya quedó construido en el lema~\ref{thm:ch-6b}, y se
diagonalizan \(e_{s(n)}\). Finalmente, entre todas las sobreyecciones así
obtenidas se define \(e_\beta\) como la \(<_{\mathfrak G}\)-menor. La
recursión \(\beta\mapsto e_\beta\) es una clase uniformemente definible de
\(\mathfrak G\); cada \(e_\beta\), y cada restricción de la recursión a un
ordinal, pertenece a \(\mathfrak G\). Sólo usa \(h,u_{h,m}\) y el buen orden
global.
\end{proof}

\begin{lemma}[Condensación de los reales internos por debajo de \(\omega_1^{\mathfrak G}\)]
\label{thm:ch-8}
Todo real de
\(\mathfrak G_{\rm HMT}(h,m_\infty)\) aparece en un nivel de índice menor que
\(\omega_1^{\mathfrak G}\).

\end{lemma}
\begin{proof}
Sea \(r\subseteq\omega\) un real de \(\mathfrak G\) y sea
\(B:=\mathfrak G_0(h,m_\infty)\). Esta base es transitiva y numerable
internamente. Con los códigos de naturales y pares finitos ya fijados,
\(\operatorname{Ord}\cap B=\omega+1\). Por inducción,
\[
 \operatorname{Ord}\cap\mathfrak G_\alpha=(\omega+1)+\alpha.
\]
El sucesor añade el ordinal formado por los ordinales del nivel anterior;
no puede añadir uno mayor, pues no sería un subconjunto de ese nivel. Los
límites son uniones. Elegimos \(\Theta\) límite, \(\Theta>\omega^2\),
con \(r\in\mathfrak G_\Theta\). Así
\(\operatorname{Ord}\cap\mathfrak G_\Theta=\Theta\), y todo nivel
\(\mathfrak G_\xi\), con \(\xi<\Theta\), pertenece a
\(\mathfrak G_\Theta\). Formamos internamente la estructura expandida
\[
 \mathcal A_\Theta=
 \langle\mathfrak G_\Theta,\in,u_{h,m},
         H_\Theta,<_{\mathfrak G}\!\upharpoonright\mathfrak G_\Theta\rangle
\]
con la relación de niveles como predicado de su signatura,
\(H_\Theta=\{\langle\xi,x\rangle:\xi<\Theta\ \&\
x\in\mathfrak G_\xi\}\), y sus funciones de Skolem definibles. Sea
\[
 Y=\operatorname{Hull}^{\mathcal A_\Theta}
   (B\cup\{B,r\})\prec\mathcal A_\Theta .
\]
La clausura bajo la familia numerable de funciones de Skolem muestra dentro
de \(\mathfrak G\) que \(Y\) es numerable. Sea
\(c:Y\rightarrow M\) su colapso de Mostowski y
\[
 \bar\Theta:=\operatorname{otp}(Y\cap\Theta)
 <\omega_1^{\mathfrak G}.
\]
Como \(B\subseteq Y\) es transitivo, la inducción de pertenencia muestra
que \(c\) fija cada elemento de \(B\) y también \(B\). En particular,
fija \(\omega,u_{h,m},h,m_\infty\). Puesto que
\(r\subseteq\omega\subseteq Y\), también \(c(r)=r\).

Probamos por inducción sobre \(\xi\in Y\cap\Theta\) que
\[
 c``\bigl(Y\cap\mathfrak G_\xi(h,m_\infty)\bigr)
 =\mathfrak G_{c(\xi)}(h,m_\infty).
\tag{C.43a}\label{eq:ch-condensacion}
\]
Para \(\xi\in Y\cap\Theta\), el predicado \(H_\Theta\) define el
conjunto \(\mathfrak G_\xi\) desde \(\xi\), de modo que
\(\mathfrak G_\xi\in Y\). Relativizar cada fórmula a ese conjunto da
\(Y\cap\mathfrak G_\xi\prec\mathfrak G_\xi\), en el lenguaje al que
el lema~\ref{thm:ch-6} traduce las definiciones HMT.

La base de~\eqref{eq:ch-condensacion} ya está probada porque el colapso fija \(B\). En un
sucesor, si \(\xi+1\in Y\), también \(\xi\in Y\) y
\(c(\xi+1)=c(\xi)+1\). Para
\(a\in Y\cap\mathfrak G_{\xi+1}\), fijamos una fórmula finita que lo
define sobre \(\mathfrak G_\xi\). La afirmación
\[
 a=\{z\in\mathfrak G_\xi:
                  \varphi^{\mathfrak G_\xi}(z,\bar p)\}
\]
es una fórmula ordinaria sobre \(\mathfrak G_\Theta\), con parámetros
\(a,\mathfrak G_\xi\). La elementalidad proporciona los parámetros
\(\bar p\) dentro de \(Y\cap\mathfrak G_\xi\). La hipótesis inductiva
y la elementalidad de ese nivel muestran que \(c(a)\) es el subconjunto
definido por la misma fórmula y los parámetros colapsados sobre
\(\mathfrak G_{c(\xi)}\). Recíprocamente, dados una fórmula y parámetros
en el nivel colapsado, se levantan los parámetros al nivel original.
El subconjunto allí definido existe, es único y pertenece a \(Y\) por
elementalidad. Su colapso es el subconjunto solicitado. Esto prueba ambas
inclusiones en el sucesor, sin introducir un predicado uniforme de verdad.

Para \(\lambda\in Y\cap\Theta\) límite, la elementalidad proporciona,
para cada \(x\in Y\cap\mathfrak G_\lambda\), un testigo
\(\xi\in Y\cap\lambda\) con \(H_\Theta(\xi,x)\). Por tanto,
\[
 Y\cap\mathfrak G_\lambda
 =\bigcup_{\xi\in Y\cap\lambda}(Y\cap\mathfrak G_\xi),
 \qquad c``(Y\cap\lambda)=c(\lambda).
\]
La continuidad de la jerarquía y la hipótesis inductiva dan~\eqref{eq:ch-condensacion} en el
límite. No se requiere cofinalidad de \(Y\cap\lambda\) en el ordinal
ambiental \(\lambda\): sus índices colapsados recorren \(c(\lambda)\).

Finalmente, cada \(x\in Y\) posee un testigo \(\xi\in Y\cap\Theta\)
con \(H_\Theta(\xi,x)\), y \(\bar\Theta\) es límite. Tomando uniones,
\[
 M=\bigcup_{\eta<\bar\Theta}\mathfrak G_\eta(h,m_\infty)
  =\mathfrak G_{\bar\Theta}(h,m_\infty).
\]
Como \(c(r)=r\in M\), se concluye
\(r\in\mathfrak G_{\bar\Theta}\) con
\(\bar\Theta<\omega_1^{\mathfrak G}\).
\end{proof}

El lema no dice que un prefijo de longitud finita enumere la recta. Dice que
cada objeto hereditario generado posee un primer nivel transfinito y un primer
código genealógico dentro de ese nivel. Ésta es exactamente la información que
el reducto extensional pierde cuando conserva sólo el valor publicado.

\subsection{Inyección de la recta generada en el primer ordinal no numerable de la clausura genealógica}

Sea \(\iota(r)\subseteq\omega\) el corte racional de
\(r\in\mathbb R_{\rm HMT}^{\mathfrak G}\) respecto de la enumeración de
\(\mathbb Q_{\rm HMT}\), construida previamente desde los enteros y cocientes
APP. El mapa \(r\mapsto\iota(r)\) es definible desde \(r\) y esa enumeración,
de modo que \(\iota(r)\in\mathfrak G\); el lema~\ref{thm:ch-8} lo sitúa en una etapa de
índice numerable. Definimos

\[
 \begin{aligned}
  \beta(r)&:=\min\{\beta:
       \iota(r)\in\mathfrak G_{\beta+1}(h,m_\infty)\},\\
  n(r)&:=\min\{n:e_{\beta(r)+1}(n)=\iota(r)\},\\
  j_{\rm HMT}(r)&:=\omega\cdot\beta(r)+n(r).
 \end{aligned}
\tag{C.44}\label{eq:ch-indices-real}
\]

La definición~\eqref{eq:ch-indices-real} es uniforme en los parámetros \(h,m_\infty\). Como
\(\mathbb R_{\rm HMT}^{\mathfrak G}\) es un conjunto de \(\mathfrak G\),
Separación fija el dominio de la definición y Reemplazo en \(\mathfrak G\)
forma su grafo:
\[
 \{\langle r,\omega\cdot\beta(r)+n(r)\rangle:
 r\in\mathbb R_{\rm HMT}^{\mathfrak G}\}\in\mathfrak G.
\]
Por tanto, \(j_{\rm HMT}\) es una función interna de \(\mathfrak G\).

Por los lemas~\ref{thm:ch-7} y~\ref{thm:ch-8},
\(j_{\rm HMT}(r)<\omega_1^{\mathfrak G}\) para todo
real generado. Si \(j_{\rm HMT}(r)=j_{\rm HMT}(s)\), la división ordinal por
\(\omega\) recupera el mismo par \((\beta,n)\); \eqref{eq:ch-enumeracion-niveles} recupera el mismo corte
racional y, por completitud, \(r=s\). Por ello

\[
 j_{\rm HMT}:\mathbb R_{\rm HMT}^{\mathfrak G}
 \hookrightarrow\omega_1^{\mathfrak G}
\tag{C.45}\label{eq:ch-inyeccion-j}
\]

es una inyección construida por el rango y el nombre de la genealogía HMT. Su
existencia utiliza exactamente la clausura semántica~\eqref{eq:ch-operador-def}, el buen orden de
nombres del lema~\ref{thm:ch-6b} y la condensación~\eqref{eq:ch-condensacion}; no recibe la hipótesis del continuo para elegir la
etapa. La identificación conjuntista posterior describe esta misma
maquinaria en otra notación, sin invertir su procedencia causal.

\subsection{Inyección del primer ordinal no numerable en la recta y equivalencia cardinal}

Para cada \(\beta<\omega_1^{\mathfrak G}\), el lema~\ref{thm:ch-7} permite codificar en un
subconjunto \(w_\beta\subseteq\omega\) un buen orden numerable de tipo
\(\beta\). Para no confundir el orden vacío con el unipuntual, se codifican
tanto el dominio \(D\subseteq\omega\) como su relación de orden estricto
\(R\subseteq D^2\). Fijado un emparejamiento biyectivo
\(\langle\cdot,\cdot\rangle:\omega^2\to\omega\), el código es
\[
 w(D,R)=\{2n:n\in D\}\cup
        \{2\langle a,b\rangle+1:(a,b)\in R\}.
\]
Los bits pares recuperan el dominio y los impares la relación. Se toma el
\(<_{\mathfrak G}\)-menor real que codifica tal buen orden. Como la propiedad
«codifica un buen orden de tipo \(\beta\)» y el buen orden global son internos,
Reemplazo forma la función \(\beta\mapsto w_\beta\) dentro de
\(\mathfrak G\). Los tipos de orden distintos producen códigos distintos. La
aplicación ternaria, ya usada en~\eqref{eq:ch-incrustacion-cantor}, se define por

\[
 E(w)=\sum_{n\ge0}\frac{2\mathbf1_w(n)}{3^{2n+2}}
\tag{C.46}\label{eq:ch-incrustacion-ternaria}
\]

y es inyectiva y evita el extremo del cilindro raíz. El levantamiento~\eqref{eq:ch-levantamiento-uniforme}
asegura que \(E_{\rm HMT}(w_\beta)\) pertenece a la recta publicada por una
historia HMT. Así

\[
 k_{\rm HMT}(\beta):=E_{\rm HMT}(w_\beta),
 \qquad
 k_{\rm HMT}:\omega_1^{\mathfrak G}
 \hookrightarrow\mathbb R_{\rm HMT}^{\mathfrak G}.
\tag{C.47}\label{eq:ch-inyeccion-k}
\]

Las inyecciones~\eqref{eq:ch-inyeccion-j} y~\eqref{eq:ch-inyeccion-k}, construidas sin utilizar la igualdad buscada,
dan por Cantor--Bernstein

\[
 \bigl|\mathbb R_{\rm HMT}^{\mathfrak G}\bigr|
 =\aleph_1^{\mathfrak G}.
\tag{C.48}\label{eq:ch-igualdad-cardinal}
\]

Como~\eqref{eq:ch-recta-potencia} identifica la recta HMT completa con
\(\mathcal P^{\mathfrak G}(\omega)\), \eqref{eq:ch-igualdad-cardinal} se reescribe exactamente como

\[
 \boxed{
  \mathfrak G_{\rm HMT}(h,m_\infty)
  \models 2^{\aleph_0}=\aleph_1.}
\tag{C.49}\label{eq:ch-decision}
\]

La igualdad no procede de \(729\), de un producto informal entre cardinalidad
y ordinalidad, de las cifras de una constante ni de una horquilla numérica.
Procede de la cobertura de la recta por el límite inverso, la condensación de
la clausura genealógica y las dos inyecciones explícitas.

\subsection{Cuantificador completo sobre todos los subconjuntos internos}

Sea
\(A\in\mathcal P^{\mathfrak G}
(\mathbb R_{\rm HMT}^{\mathfrak G})\). Si \(A\) es numerable, se satisface la
primera alternativa. Si no lo es, \(j_{\rm HMT}[A]\) es un subconjunto no
numerable de \(\omega_1^{\mathfrak G}\). No puede estar acotado por un ordinal
numerable, pues cada segmento inicial de \(\omega_1\) es numerable. Su
enumeración creciente tiene tipo \(\omega_1^{\mathfrak G}\), y la inversa de
\(j_{\rm HMT}\) sobre su imagen produce una inyección
\(\omega_1^{\mathfrak G}\hookrightarrow A\). La inclusión
\(A\subseteq\mathbb R_{\rm HMT}^{\mathfrak G}\), \eqref{eq:ch-igualdad-cardinal} y
Cantor--Bernstein dan

\[
 \boxed{
 \forall A\in\mathcal P^{\mathfrak G}
   (\mathbb R_{\rm HMT}^{\mathfrak G}),\qquad
 |A|\le\aleph_0
 \quad\text{o}\quad
 |A|=\bigl|\mathbb R_{\rm HMT}^{\mathfrak G}\bigr|.}
\tag{C.50}\label{eq:ch-dicotomia}
\]

Aquí se cuantifica la potencia completa del continuo integral generado por
HMT. El transporte biyectivo del plegado con memoria lleva~\eqref{eq:ch-dicotomia} a cualquiera
de sus representaciones sin perder partes; no sustituye las inyecciones
\eqref{eq:ch-inyeccion-j}--\eqref{eq:ch-inyeccion-k}.

\subsection{Representación de la clausura genealógica como \texorpdfstring{\(L[h,m_\infty]\)}{L[h,m-infinito]}}

Sólo después de~\eqref{eq:ch-decision}--\eqref{eq:ch-dicotomia} se asigna
a~\eqref{eq:ch-clausura-transfinita} su nombre en la notación
conjuntista usual. Sea \(L[u_{h,m}]\) la jerarquía constructible relativa al
real ya producido en~\eqref{eq:ch-codigo-conjunto}. Como la signatura HMT es numerable y todos sus
operadores son relaciones de grafo codificadas por \(h\), el lema~\ref{thm:ch-6} elimina
sus símbolos uniformemente. Se obtiene la igualdad de clases

% Se separa el contador hipervinculado de la etiqueta editorial (51) para
% evitar dos destinos PDF homónimos en la última ecuación numerada del capítulo.
\stepcounter{equation}
\[
 \mathfrak G_{\rm HMT}(h,m_\infty)
 =L[u_{h,m}]=L[h,m_\infty].
\tag{C.51}\label{eq:ch-representacion-L}
\]

Para probar la primera inclusión se fortalece la inducción: para cada
ordinal \(\beta\), tanto \(\mathfrak G_\beta\) como la secuencia
\(\langle\mathfrak G_\xi:\xi\leq\beta\rangle\) pertenecen a
\(L[u_{h,m}]\). La base es la clausura transitiva de un conjunto construido
desde \(\omega\) y \(u_{h,m}\); los parámetros \(h,m_\infty\) se recuperan
de \(u_{h,m}\). En el sucesor, el nivel precedente es un \emph{conjunto
interno}, no sólo una subclase de \(L[u_{h,m}]\). Su relación de satisfacción
y la eliminación uniforme~\eqref{eq:ch-eliminacion-uniforme} permiten formar internamente el conjunto de
sus subconjuntos definibles. En un límite, la recursión transfinita,
Reemplazo y Unión de \(L[u_{h,m}]\) forman la secuencia anterior y su unión.
La transitividad da entonces \(\mathfrak G\subseteq L[u_{h,m}]\).

Recíprocamente, el lema~\ref{thm:ch-6b} proporciona un modelo interno
transitivo de ZF, con todos los ordinales y con \(u_{h,m}\) como elemento.
La recursión que construye \(L[u_{h,m}]\) se efectúa en ese modelo: en cada
etapa, satisfacción sobre el nivel anterior y definibilidad relativa son
absolutas entre los dos dominios transitivos. Por inducción, sus niveles
coinciden con los de la construcción ambiental y pertenecen a
\(\mathfrak G\). Así \(L[u_{h,m}]\subseteq\mathfrak G\). Finalmente,
\eqref{eq:ch-codigo-conjunto} hace a \(u_{h,m}\) y al par \((h,m_\infty)\)
primitivo-recursivamente interdefinibles, de lo cual se sigue la segunda
igualdad.

La ecuación~\eqref{eq:ch-representacion-L} identifica explícitamente la
clausura con la jerarquía constructible relativa al parámetro generado. Esta
representación no modifica la procedencia causal: \(L[u]\) ni selecciona el
registro genealógico ni define las prolongaciones o la recta; la decisión de
la hipótesis del continuo resulta de la regla de exhaustividad y de los mapas probados entre
\eqref{eq:ch-clausura-transfinita} y~\eqref{eq:ch-dicotomia}. La composición
\eqref{eq:ch-codigo-conjunto}, \eqref{eq:ch-clausura-transfinita},
\eqref{eq:ch-equivalencia-semantica} y
\eqref{eq:ch-decision}--\eqref{eq:ch-representacion-L}, junto con el
teorema~\ref{thm:decision-continuo-hmt}, ofrece una expresión completa y
uniforme del universo realizado, la coincidencia semántica y el cierre
cardinal; ninguna de estas representaciones sustituye la generación
APP--TRIT--TPK.

\subsection{Teorema de resolución afirmativa de la hipótesis del continuo}

\begin{theorem}[Decisión del continuo en HMT]
\label{thm:decision-continuo-hmt}
Para toda historia realizada
\(m_\infty\) producida por la acción conjunta
\(\mathrm{APP}\to\mathrm{TRIT}\to\mathrm{TPK}\to
X_\infty^{\rm enr}\to\mathcal C_{\rm cont}^{\rm disc}\), la extensión
genealógica HMT~\eqref{eq:ch-clausura-transfinita} genera la reconstrucción interna de la recta integral
\eqref{eq:ch-proyeccion-arquimediana}, \(\mathbb R_{\rm HMT}^{\mathfrak G}\), construye las inyecciones
\eqref{eq:ch-inyeccion-j} y~\eqref{eq:ch-inyeccion-k}, y satisface
\(2^{\aleph_0}=\aleph_1\) sobre toda su potencia interna. Por tanto, HMT
demuestra afirmativamente la hipótesis del continuo para el continuo
íntegro que define y genera. La notación \(L[h,m_\infty]\) es una
representación posterior de esa jerarquía, no la premisa que decide el
resultado.
\end{theorem}


\subsection{Criterios de falsación del cierre cardinal}

La lista siguiente identifica los fallos que invalidarían un paso del cierre
anterior. Se distinguen los controles de construcción, semántica y
cardinalidad:

\begin{enumerate}
\item existe un nivel \(N\) o un estado superviviente sin sucesor compatible en
\eqref{eq:ch-extension-no-vacia};
\item \(u_{h,m}\) no codifica hereditariamente \(h,m_\infty\), o alguna relación
de grafo de \(\Sigma_{\rm HMT}\) no admite la eliminación uniforme~\eqref{eq:ch-eliminacion-uniforme};
\item un objeto admitido por las reglas semánticas HMT carece de término en
\eqref{eq:ch-terminos}, o un término de~\eqref{eq:ch-terminos} publica un objeto ajeno a esas reglas, rompiendo
\eqref{eq:ch-equivalencia-semantica};
\item \(\mathfrak G\) no satisface alguno de los axiomas usados, falla su buen
orden definible o una supuesta enumeración \(e_\beta\) no es interna y
uniforme;
\item los subcilindros de~\eqref{eq:ch-particion-cilindros} no forman una partición disjunta exhaustiva, sus
cilindros dejan de contraerse o la recursión~\eqref{eq:ch-prefijos-uniformes} no produce una
historia compatible interna;
\item el colapso del lema~\ref{thm:ch-8} no conmuta con una etapa sucesora o límite, o existe
un real interno que no aparece en ningún nivel de índice menor que
\(\omega_1^{\mathfrak G}\);
\item la selección \(\beta\mapsto w_\beta\) no es interna, uniforme o conserva
el tipo de orden;
\item dos reales diferentes tienen el mismo par genealógico \((\beta,n)\) en
\eqref{eq:ch-indices-real}, dos ordinales diferentes producen el mismo código en~\eqref{eq:ch-inyeccion-k}, o falla
alguna de las inyecciones \(E,C,j_{\rm HMT},k_{\rm HMT}\);
\item existe una parte interna no numerable de la recta cuyo rango bajo
\(j_{\rm HMT}\) sea acotado en \(\omega_1^{\mathfrak G}\);
\item alguna etapa generativa anterior a la cobertura~\eqref{eq:ch-particion-cilindros} utiliza la hipótesis del continuo,
\(\mathbb R\), una constante objetivo, el Comité de Datos del Consejo
Internacional de Ciencia (CODATA), el Particle Data Group (PDG) o una identificación física
posterior como antecedente causal para definir una semilla, una relación
\(\mathscr L_\tau\), \(\operatorname{Ext}\), un coeficiente, la admisibilidad
o una prolongación del árbol. No constituye tal inversión el levantamiento
posterior~\eqref{eq:ch-prefijos-uniformes}: allí un real arbitrario ya construido es una variable
universal y selecciona, dentro de la partición exhaustiva previamente
demostrada, el único subcilindro que lo contiene.
\end{enumerate}

Ninguna ejecución limitada a un número prefijado de niveles sustituye el
cuantificador universal del sistema inverso; lo falsificaría una ruptura
efectiva de compatibilidad en alguna profundidad. Ninguna notación conjuntista
posterior prueba por sí sola~\eqref{eq:ch-decision}. El cierre utiliza conjuntamente el
lema~\ref{thm:ch-6b}, la identificación recta--potencia de~\eqref{eq:ch-recta-potencia}, los
lemas~\ref{thm:ch-7}--\ref{thm:ch-8} y las dos inyecciones construidas desde
la genealogía.


\section{Censo de firmas del emisor residual}
\label{app:firmas-emision}

Este censo especifica las fibras finitas utilizadas en la sección~\ref{sec:extension}.
Cada condición inicial se identifica por
\[
 \nu=4(9i+j)+d,\qquad 0\leq i,j\leq8,\quad
 d=0,1,2,3\ \longleftrightarrow\ N,E,S,O.
\]
La división euclidiana recupera \(d=[\nu]_4\), \(j=[\lfloor\nu/4\rfloor]_9\)
e \(i=\lfloor\nu/36\rfloor\). Por tanto, \(0\leq\nu<324\) enumera
cada condición exactamente una vez, y las parejas de tales identificadores
enumeran las \(104\,976\) condiciones de \(\Omega_{\rm seed}\).

Sea \(f_+(\nu)\), respectivamente \(f_\times(\nu)\), la firma de seis
componentes obtenida por la recurrencia de emisión. En cada paso activo se lee
la casilla antes de avanzar; los vectores de dirección son
\[
 v_N=(-1,0),\quad v_E=(0,1),\quad v_S=(1,0),\quad v_O=(0,-1).
\]
Las coordenadas se reducen módulo nueve y el vector se invierte tras los
pasos \(27\) y \(54\). Los pasos de fase cero no desplazan ninguno de los
cursores. Estas reglas, junto con \eqref{eq:emisor-seis}, determinan
\(f_+\) y \(f_\times\) sin un catálogo de valores objetivo.

La tabla da todas sus imágenes. La multiplicidad \(m_\sigma\) es el número
de identificadores en la fibra \(f_\sigma^{-1}(a)\); \(\nu_{\min}\) es su
menor elemento. Cada palabra conserva los ceros iniciales. La fibra completa
queda determinada por la condición
\[
 f_\sigma^{-1}(a)=\{\nu\in\{0,\ldots,323\}:f_\sigma(\nu)=a\}.
\]
\begin{center}\small
\begin{tabular}{@{}lrr@{\qquad}lrr@{}}
\toprule
\(f_+\) & \(m_+\) & \(\nu_{\min}\) & \(f_\times\) & \(m_\times\) & \(\nu_{\min}\)\\
\midrule
\texttt{122564} & 18 & 17 & \texttt{000000} & 108 & 8 \\
\texttt{122898} & 18 & 24 & \texttt{177393} & 8 & 1 \\
\texttt{212645} & 18 & 5 & \texttt{177555} & 8 & 54 \\
\texttt{212988} & 18 & 0 & \texttt{177771} & 8 & 11 \\
\texttt{221456} & 18 & 29 & \texttt{339555} & 8 & 6 \\
\texttt{221889} & 18 & 12 & \texttt{339771} & 8 & 15 \\
\texttt{456221} & 18 & 28 & \texttt{339933} & 8 & 59 \\
\texttt{465898} & 18 & 21 & \texttt{393177} & 8 & 0 \\
\texttt{546988} & 18 & 9 & \texttt{393393} & 8 & 45 \\
\texttt{564122} & 18 & 16 & \texttt{393555} & 8 & 40 \\
\texttt{645212} & 18 & 4 & \texttt{555177} & 8 & 52 \\
\texttt{654889} & 18 & 33 & \texttt{555339} & 8 & 4 \\
\texttt{889221} & 18 & 13 & \texttt{555393} & 8 & 41 \\
\texttt{889654} & 18 & 32 & \texttt{555555} & 24 & 84 \\
\texttt{898122} & 18 & 25 & \texttt{555717} & 8 & 14 \\
\texttt{898465} & 18 & 20 & \texttt{555771} & 8 & 18 \\
\texttt{988212} & 18 & 1 & \texttt{555933} & 8 & 24 \\
\texttt{988546} & 18 & 8 & \texttt{717555} & 8 & 12 \\
 & &  & \texttt{717717} & 8 & 21 \\
 & &  & \texttt{717933} & 8 & 19 \\
 & &  & \texttt{771177} & 8 & 9 \\
 & &  & \texttt{771339} & 8 & 13 \\
 & &  & \texttt{771555} & 8 & 16 \\
 & &  & \texttt{933339} & 8 & 57 \\
 & &  & \texttt{933555} & 8 & 26 \\
 & &  & \texttt{933717} & 8 & 17 \\
\bottomrule
\end{tabular}
\end{center}

La evaluación de la recurrencia para los \(324\) identificadores produce
las dos particiones de la tabla: \(18\cdot18=324\) y
\(24\cdot8+24+108=324\). El producto de sus fibras da
\(432\) fibras de tamaño \(144\), \(18\) de tamaño \(432\) y \(18\)
de tamaño \(1944\), por la inyectividad de \eqref{eq:acoplamiento}.
Finalmente, aplicar \eqref{eq:psi-catalogo} a ese producto de imágenes
da \(243\) palabras distintas. El censo es finito y pertenece al emisor
residual; el refinamiento coinductivo posterior conserva su propia ley.


\HMTNarrativePaper
\HMTNarrativeHeading{La unidad conservada en sus realizaciones}

El recorrido del artículo vuelve a su tesis inicial con operaciones
ya desarrolladas. APP conserva la evaluación y su cociente;
TRIT distingue orientación y régimen; el TPK compone transporte,
selección y prolongación. La conexión nonádica actualiza la
historia durante el retorno de fase. El refinamiento produce
las lecturas numéricas y conserva las incidencias que sostienen
su realización excepcional.

La constante holográfica de acoplamiento aritmético-geométrico
reúne tres funciones precisas. Su registro admite una codificación
posicional reversible; interviene en la clausura de alfa y en
la incidencia marcada; su órbita permite identificar operadores
con cotas de recuperación. La ley bilateral transporta esa
capacidad mediante la conservación del estado y de sus formas.

Las vacancias explican cómo la historia avanza cuando una
capacidad permanece fija. La sección electrónica muestra una
persistencia con trayectoria y registro. Las secciones de acción,
los canales orientados, los lectores de área y la transducción
térmica convierten los resultados anteriores en relaciones
físicas calculables. Las realizaciones excepcionales conservan
sus portadores, incidencias y mapas dimensionales.

La conservación evolutiva de la unidad dimensional queda
expresada en varios niveles enlazados: partición normalizada,
refinamiento reversible, transporte de formas, archivo con
terminal, recuperación estable y transformación de observables.
Cada nivel tiene su objeto y su prueba; las aplicaciones entre
ellos construyen la continuidad del argumento. La doble
proyección mantiene el centro de esa organización: valor e
incidencia proceden del mismo estado, cuya transformación
conserva relaciones recuperables.


\section{Conclusiones}
\label{paper:conclusiones}
La extrema y media razón holográfica queda formulada mediante una estructura generadora, una evaluación aritmética reversible y una ley de transporte con recuperación de información. El registro dodecafásico \(K\), producido por APP--TRIT--TPK, conserva su orden en la carta de fase y enlaza la clausura aritmética con la representación incidencial. Su lectura
\[
\kappa=
\frac{234543140729659824621058914794146601}{10^{36}-1}
\]
recupera exactamente las doce componentes: la multiplicación por el denominador restituye el entero y la base mil determina sus bloques. Este resultado fija el sentido matemático de la constante holográfica de acoplamiento aritmético--geométrico. Su valor es una representación exacta del registro ordenado y participa, junto con las coordenadas regionales de \(\pi,\varphi,e\), en la determinación de \(\alpha_A\).

La misma estructura permite identificar operaciones. Los doce desplazamientos cíclicos de \(K\) forman una base cuya matriz de Gram tiene extremos espectrales \(589609\) y \(6263^2\). Las respuestas a esa base determinan un operador lineal con factor de amplificación del error \(1/\sqrt{589609}\); una respuesta vectorial basta en la clase que conmuta con el desplazamiento. La representación incidencial transporta el registro a su imagen especificada con inversa explícita. Dos observaciones de memoria recuperan su parte centrada con cota óptima \(9/4\), y la carga uniforme completa el registro. Evaluación, incidencia e identificación se componen así sobre un mismo objeto, conservando en cada paso los datos requeridos por la inversa.

La coincidencia de las publicaciones entera y analítica de \(\alpha\) se establece a toda profundidad. La carta analítica prolonga el truncamiento de orden nueve mediante una recurrencia de razón \(1/729\), cuyo coeficiente inicial se obtiene de la precoordenada ya generada. La convergencia uniforme, la cota positiva de la derivada y los intervalos racionales encajados identifican la misma raíz simple y los mismos prefijos canónicos en ambas publicaciones. El teorema determina su compatibilidad completa dentro de esa construcción correlacionada.

El resultado general que sostiene la conservación es la identidad bilateral
\[
 (a+1)(aB-C)^*(aB-C)
 +a\bigl((a+1)B+C\bigr)^*\bigl((a+1)B+C\bigr)
 =\bigl((a+1)^3+a^3\bigr)I,
\]
para \(a>0\) y dos isometrías \(B,C\) con dominio y codominio comunes. La cancelación de los términos cruzados demuestra la igualdad en dimensión arbitraria. En la especialización \(a=8\), el factor ternario \(73\) conserva su interpretación como norma algebraica e índice reticular, y satisface \(10\cdot73=729+1\). El refinamiento de etiquetas enteras admisibles tiene una inversa por división euclídea, con profundidad y convención de acarreo especificadas; su levantamiento a las bases de estados es unitario y se compone con el balance.

La ley permite determinar también qué cambia al observar. Para transporte relativo unitario \(R\), el lector diagonal transforma el observable de la carta común según
\[
 \mathcal T(G)=\frac{72}{73}G+\frac1{73}R^*GR.
\]
Sus puntos fijos son exactamente los observables que conmutan con \(R\). El observable transportado sobre el estado completo conserva, en cambio, la lectura inicial por su fórmula de restitución. Las identidades exteriores prolongan este resultado a cada grado admisible, conservando los sectores puramente observados, los registrados y los mixtos. Las contribuciones de área y volumen quedan incluidas en una ley de formas, con las métricas y dimensiones correspondientes.

La recuperación se mantiene a profundidad arbitraria mediante una política explícita de archivo. En el caso nonádico, el operador terminal satisface \(\|R_N\|^2\leq(729/1241)^N\). El archivo infinito es, por tanto, una isometría y su adjunto reconstruye el estado con norma operatoria uno. Para cada profundidad positiva, el archivo finito admite una inversa exacta corregida por su Gram, con cota
\[
 \|L_N\|\leq\bigl(1-(729/1241)^N\bigr)^{-1/2}.
\]
Estas fórmulas cuantifican la estabilidad y separan la inversión exacta de la aproximación obtenida al truncar una serie de recuperación. En las dinámicas generales, la memoria y el terminal superviviente conservan conjuntamente el estado; la política indicada demuestra el régimen en que toda su norma se transfiere al archivo.

Las realizaciones físicas desarrolladas aplican estas operaciones con sus datos explícitos de acción, orientación, escala y representación. En particular, para el protocolo simpléctico declarado con dos entradas gaussianas puras, centradas e idénticas y un número finito de modos, la componente antilineal de la memoria fija el espectro simpléctico de la reducción, manteniendo todas las ocupaciones de Fock. El lazo elíptico--parabólico produce pureza \(9/11\), espectro \((9/10)10^{-j}\) y entropía adimensional \((10/9)\log10-\log9\). La conservación global y la mezcla de la observación reducida se describen dentro del mismo protocolo. Las realizaciones electromagnéticas, gravitatorias, térmicas y variacionales mantienen los operadores y normalizaciones de sus respectivas pruebas.

La construcción común del continuo conserva las cinco operaciones correlativas y la compatibilidad de sus refinamientos. La clausura de familias de historias de la sección~\ref{sec:cierre-cardinal-nativo} establece, en el universo genealógico allí definido, la igualdad cardinal interna de su recta con \(\aleph_1\) de ese universo. La afirmación conserva el dominio y las operaciones de pertenencia que intervienen en sus dos inyecciones y en su identificación final.

El censo del emisor residual hace explícita la base finita de una parte de esa cadena: reúne \(44\) firmas, \(18\) aditivas y \(26\) multiplicativas, con sus multiplicidades y la regla de reconstrucción de las fibras. Las \(648\) pertenencias de condiciones iniciales, \(324\) en cada hoja, se conservan en los datos del suplemento. Los siete teoremas comprobados en Lean certifican las identidades aritméticas y sus consecuencias finitas declaradas; las demostraciones de operadores, límites e incidencias se desarrollan en el texto. Censo, cálculo y formalización fijan así sus respectivos objetos de comprobación.

El resultado central es una caracterización calculable de la conservación evolutiva de la unidad: la estructura determina una evaluación, sus transformaciones distribuyen la información y las inversas precisan cómo recuperarla. La constante de acoplamiento, las leyes de observables y las cotas de reconstrucción expresan esa continuidad mediante ecuaciones y operaciones, desde el registro aritmético hasta las realizaciones especificadas.

\clearpage\HMTSectionPagefalse
\begin{thebibliography}{99}
\phantomsection
\addcontentsline{toc}{section}{Referencias}
\bibitem{feynman1948} R. P. Feynman, ``Space-Time Approach to Non-Relativistic Quantum Mechanics'', \emph{Reviews of Modern Physics} \textbf{20}, 367--387 (1948). \href{https://doi.org/10.1103/RevModPhys.20.367}{doi:10.1103/RevModPhys.20.367}.
\bibitem{feynman1949} R. P. Feynman, ``The Theory of Positrons'', \emph{Physical Review} \textbf{76}, 749--759 (1949). \href{https://doi.org/10.1103/PhysRev.76.749}{doi:10.1103/PhysRev.76.749}.
\bibitem{feynman1965} R. P. Feynman, ``The Development of the Space-Time View of Quantum Electrodynamics'', conferencia Nobel, 11 de diciembre de 1965. \href{https://www.nobelprize.org/prizes/physics/1965/feynman/lecture/}{Texto de la conferencia}.
\bibitem{mustovicary} B. Musto y J. Vicary, ``Quantum Latin Squares and Unitary Error Bases'', \emph{Quantum Information and Computation} \textbf{16}, 1318--1332 (2016). \href{https://doi.org/10.26421/QIC16.15-16-4}{doi:10.26421/QIC16.15-16-4}.
\bibitem{delascuevas2020} G. De las Cuevas, T. Drescher y T. Netzer, ``Quantum Magic Squares: Dilations and Their Limitations'', \emph{Journal of Mathematical Physics} \textbf{61}, 111704 (2020). \href{https://doi.org/10.1063/5.0022344}{doi:10.1063/5.0022344}.
\bibitem{delascuevas2023} G. De las Cuevas, T. Netzer e I. Valentiner-Branth, ``Magic Squares: Latin, Semiclassical, and Quantum'', \emph{Journal of Mathematical Physics} \textbf{64}, 022201 (2023). \href{https://doi.org/10.1063/5.0127393}{doi:10.1063/5.0127393}.
\bibitem{dirac1933} P. A. M. Dirac, ``The Lagrangian in Quantum Mechanics'', \emph{Physikalische Zeitschrift der Sowjetunion} \textbf{3}, 64--72 (1933). \href{https://www.informationphilosopher.com/solutions/scientists/dirac/Lagrangian_1933.pdf}{Reproducción del artículo original}.
\bibitem{pauli1946} W. Pauli, ``Exclusion Principle and Quantum Mechanics'', conferencia Nobel, 13 de diciembre de 1946, en \emph{Nobel Lectures, Physics 1942--1962}, Elsevier, 1964, pp. 27--43, especialmente p. 42. \href{https://www.nobelprize.org/uploads/2018/06/pauli-lecture.pdf}{Texto de la conferencia}.
\bibitem{codata2018} E. Tiesinga, P. J. Mohr, D. B. Newell y B. N. Taylor, ``CODATA Recommended Values of the Fundamental Physical Constants: 2018'', \emph{Reviews of Modern Physics} \textbf{93}, 025010 (2021). \href{https://doi.org/10.1103/RevModPhys.93.025010}{doi:10.1103/RevModPhys.93.025010}. \href{https://physics.nist.gov/cuu/Constants/ArchiveASCII/allascii_2018.txt}{Tablas archivadas de 2018}.
\bibitem{codata2022} P. J. Mohr, D. B. Newell, B. N. Taylor y E. Tiesinga, ``CODATA Recommended Values of the Fundamental Physical Constants: 2022'', \emph{Reviews of Modern Physics} \textbf{97}, 025002 (2025). \href{https://doi.org/10.1103/RevModPhys.97.025002}{doi:10.1103/RevModPhys.97.025002}. \href{https://physics.nist.gov/cuu/Constants/Table/allascii.txt}{Tabla de constantes}.
\bibitem{flm} I. B. Frenkel, J. Lepowsky y A. Meurman, \emph{Vertex Operator Algebras and the Monster}, Pure and Applied Mathematics, vol. 134, Academic Press, 1988.
\bibitem{borcherds} R. E. Borcherds, ``Monstrous Moonshine and Monstrous Lie Superalgebras'', \emph{Inventiones Mathematicae} \textbf{109}, 405--444 (1992). \href{https://doi.org/10.1007/BF01232032}{doi:10.1007/BF01232032}.
\bibitem{conwaysloane} J. H. Conway y N. J. A. Sloane, \emph{Sphere Packings, Lattices and Groups}, 3.\textsuperscript{a} ed., Grundlehren der mathematischen Wissenschaften 290, Springer, 1999. \href{https://doi.org/10.1007/978-1-4757-6568-7}{doi:10.1007/978-1-4757-6568-7}.
\bibitem{bipm2018} Conférence générale des poids et mesures, \emph{Resolution 1 of the 26th CGPM: On the revision of the International System of Units}, 2018. \href{https://www.bipm.org/en/committees/cg/cgpm/26-2018/resolution-1}{Resolución oficial}. \href{https://doi.org/10.59161/CGPM2018RES1E}{doi:10.59161/CGPM2018RES1E}.
\bibitem{dirac1931} P. A. M. Dirac, ``Quantised singularities in the electromagnetic field'', \emph{Proceedings of the Royal Society of London. Series A} \textbf{133}, 60--72 (1931). \href{https://doi.org/10.1098/rspa.1931.0130}{doi:10.1098/rspa.1931.0130}.
\bibitem{lep2006} ALEPH, DELPHI, L3, OPAL y SLD Collaborations; LEP Electroweak Working Group; SLD Electroweak Group y SLD Heavy Flavour Group, ``Precision electroweak measurements on the Z resonance'', \emph{Physics Reports} \textbf{427}, 257--454 (2006). \href{https://doi.org/10.1016/j.physrep.2005.12.006}{doi:10.1016/j.physrep.2005.12.006}.
\bibitem{chenlamshimakura2016} H.-Y. Chen, C. H. Lam y H. Shimakura, ``\(\mathbb Z_3\)-orbifold construction of the Moonshine vertex operator algebra and some maximal \(3\)-local subgroups of the Monster'', prepublicación, arXiv:1606.05961 (2016), versión 2 (2017). \href{https://arxiv.org/abs/1606.05961}{arXiv:1606.05961}.
\bibitem{abelamyamada2017} T. Abe, C. H. Lam y H. Yamada, ``A remark on \(\mathbb Z_p\)-orbifold constructions of the Moonshine vertex operator algebra'', prepublicación, arXiv:1705.09022 (2017), versión 4. \href{https://arxiv.org/abs/1705.09022v4}{arXiv:1705.09022}.
\bibitem{kmconwaynorton1979} J. H. Conway y S. P. Norton, ``Monstrous Moonshine'', \emph{Bulletin of the London Mathematical Society} \textbf{11}, 308--339 (1979). \href{https://doi.org/10.1112/blms/11.3.308}{doi:10.1112/blms/11.3.308}.
\bibitem{bennett2003} C. H. Bennett, ``Notes on Landauer's principle, reversible computation, and Maxwell's Demon'', \emph{Studies in History and Philosophy of Modern Physics} \textbf{34}, 501--510 (2003). \href{https://doi.org/10.1016/S1355-2198(03)00039-X}{doi:10.1016/S1355-2198(03)00039-X}.
\end{thebibliography}

\section*{Referencias complementarias de conservación y representación}
\addcontentsline{toc}{section}{Referencias complementarias de conservación y representación}
\begin{enumerate}
\item E. Noether, \emph{Invariante Variationsprobleme}, Nachrichten von der Gesellschaft der Wissenschaften zu Göttingen, Mathematisch-Physikalische Klasse (1918), 235--257. Traducción de M. A. Tavel: \href{https://arxiv.org/abs/physics/0503066}{Invariant Variation Problems}.
\item C. Weedbrook, S. Pirandola, R. García-Patrón, N. J. Cerf, T. C. Ralph, J. H. Shapiro y S. Lloyd, \emph{Gaussian Quantum Information}, Reviews of Modern Physics 84 (2012), 621--669. \href{https://arxiv.org/abs/1110.3234}{arXiv:1110.3234}.
\end{enumerate}
\begin{enumerate}[resume]
\item J. Williamson, \emph{On the Algebraic Problem Concerning the Normal Forms of Linear Dynamical Systems}, American Journal of Mathematics 58 (1936), 141--163. \href{https://doi.org/10.2307/2371062}{doi:10.2307/2371062}.
\item A. S. Holevo, M. Sohma y O. Hirota, \emph{Capacity of Quantum Gaussian Channels}, Physical Review A 59 (1999), 1820--1828. \href{https://doi.org/10.1103/PhysRevA.59.1820}{doi:10.1103/PhysRevA.59.1820}.
\end{enumerate}

% La procedencia inédita se conserva en el inventario externo de fuentes.
% Las demostraciones utilizadas figuran en el cuerpo del artículo.

\begingroup
\renewcommand{\refname}{Fuentes de la integración estructural}
\begin{thebibliography}{9}
\phantomsection\addcontentsline{toc}{section}{Fuentes de la integración estructural}
\bibitem{Manivel2005} L. Manivel, \emph{Configurations of lines and models of Lie algebras}, arXiv:math/0507118, 2005.
\bibitem{Bochner1955} S. Bochner, \emph{Harmonic Analysis and the Theory of Probability}, University of California Press, Berkeley, 1955.
\end{thebibliography}
\endgroup

\end{document}
